%!TEX TS-program = xelatex 
%!TEX encoding = UTF-8 Unicode






% Options for packages loaded elsewhere
% \PassOptionsToPackage{unicode}{hyperref}
% \PassOptionsToPackage{hyphens}{url}


%

% \documentclass[
% % % % % % ]{report}

% PEARSON
\documentclass[twoside,10pt,letterpaper,usenames]{newstyle-PearsonGeneric-7-38-ROMEO}
\usepackage[twoside]{geometry} % to set the dimensions of the page

% PEARSON
\geometry{                                  % setting all the necessary dimensions...
paperwidth=7.375in,                  % Do not change these settings.
paperheight=9.125in,                 % These set up the trim size for the book.
lmargin=1in,
rmargin=.875in,
bmargin=.650in,
tmargin=.95in,
width=5.5in,
height=7.4in,
marginparwidth=0in,
marginparsep=0in,
headheight=0.2in,
headsep=.25in,
footskip=.025in}


% PEARSON
%%%%%%%% Skip over this part.  These lines should remain unaltered.
%
%%    These lines input packages that are used for this template.
\usepackage[letter,cam,center]{crop-AlteredCamMarks-Romeo}  % to add trim marks; updated Nov 2015 to allow more space between marks and trim edge
\usepackage{graphicx} % for graphics layout

%%%%%%%%%%%%%%%%% ROMEO- Running Heads
% Set up RHs for frontmatter and Ch0; they will have to be altered again for Chs 1-3 and then revert to this again for Ch4
\usepackage{fancyhdr}
\usepackage{extramarks}
\newcommand{\cppxx}{C++xx}
\newcounter{cppxx}[chapter]
\pagestyle{fancy}
% for fancy pages, i.e., most pages
\rfoot[\fancyplain{}{\hspace{\fill}\\ \sfnine\thepage\hfill}]{\fancyplain{\hspace{\fill}\\ \sfnine\thepage}{\hspace{\fill}\\ \sfnine\thepage}}
\lfoot[\fancyplain{\bf\thepage}{}]{\fancyplain{}{}}
\cfoot[\fancyplain{}{}]{\fancyplain{}{}}
%\renewcommand{\sectionmark}[1]{\markright{{\sfnine\thesection\ #1}}} 
\renewcommand{\sectionmark}[1]{\markright{{\sfnine\ #1}}} % to remove section number   \sfnine
\renewcommand{\chaptermark}[1]{\markboth{{\sfb Chapter~\thechapter\quad {\sfb #1}}}{}}
\renewcommand{\headrulewidth}{0pt}
\renewcommand{\chaptersmark}[1]{\markboth{}{\sfb #1}}
% for plain pages, like chapter openers
\fancypagestyle{plain}{%
\fancyhf{} % clear all header and footer fields
\fancyfoot[RO]{\sfnine{\thepage}} 
\renewcommand{\headrulewidth}{0pt}
\renewcommand{\footrulewidth}{0pt}
}




%%%% CHANGED FOR ROMEO
% Changing code font to Cousine
\usepackage{amsmath}
\usepackage{mathspec}
\setmonofont[Ligatures=NoCommon,Scale=.9,SlantedFont={Cousine Italic},BoldFont={Cousine Bold},BoldSlantedFont={Cousine Bold Italic}]{Cousine Regular} % to set our code font
\newfontface{\SecCode}[Ligatures=NoCommon,Scale=1.3]{Cousine Bold}
%\newfontface{\SecCode}[Ligatures=NoCommon,Scale=1.105]{Cousine Bold}
\newfontface{\SubsecCode}[Ligatures=NoCommon,Scale=1.2]{Cousine Bold}
%\newfontface{\SubsecCode}[Ligatures=NoCommon,Scale=1.02]{Cousine Bold}
\newfontface{\SubsubsecCode}[Ligatures=NoCommon,Scale=1]{Cousine Bold}
%\newfontface{\SubsubsecCode}[Ligatures=NoCommon,Scale=.935]{Cousine Bold}
\newfontface{\ParaCode}[Ligatures=NoCommon,Scale=.9]{Cousine Bold}
\newfontface{\SubparaCode}[Ligatures=NoCommon,Scale=.9]{Cousine Regular}
%\newfontface{\RHCode}[Ligatures=NoCommon,Scale=.81]{Cousine Regular}
%\newfontface{\RHCode}[Ligatures=NoCommon,Scale=.85]{Cousine Bold}
\newfontface{\RHCode}[Ligatures=NoCommon,Scale=1.17]{Cousine Bold} % set at scaled .9 of 13pts
\newfontface{\RHCodeEleven}[Ligatures=NoCommon,Scale=.935]{Cousine Bold}
%\newfontface{\RHCodeTwelve}[Ligatures=NoCommon,Scale=1.02]{Cousine Bold}
\newfontface{\RHCodeTwelve}[Ligatures=NoCommon,Scale=1.1]{Cousine Bold}
\newfontface{\TOCCode}[Ligatures=NoCommon,Scale=.85]{Cousine Bold}
%\newfontface{\TOCCode}[Ligatures=NoCommon,Scale=.9]{Cousine Bold}
\newfontface{\TOCCodeR}[Ligatures=NoCommon,Scale=.85]{Cousine Regular}
\newfontface{\TOCCodeEleven}[Ligatures=NoCommon,Scale=.935]{Cousine Regular}
\newfontface{\TOCCodeElevenB}[Ligatures=NoCommon,Scale=.935]{Cousine Bold}
%\newfontface{\TOCCodeElevenB}[Ligatures=NoCommon,Scale=1.1]{Cousine Bold}
\newfontface{\TOCCodeTwelve}[Ligatures=NoCommon,Scale=1.02]{Cousine Regular}
\newfontface{\TOCCodeThirteen}[Ligatures=NoCommon,Scale=1.105]{Cousine Regular}

% Changing sans serif font to Gill Sans Nova
\setsansfont[Ligatures=Common,Scale=1,SlantedFont={Gill Sans Nova Medium Italic},BoldFont={Gill Sans Nova Bold},BoldSlantedFont={Gill Sans Nova Bold Italic}]{Gill Sans Nova Medium} % to set our sans font
%%
\newfontface{\sfbHugeRomeo}[Ligatures=Common,Scale=2.5]{Gill Sans Nova Bold}
\newfontface{\cmssbxparttocRomeo}[Ligatures=Common,Scale=1.2]{Gill Sans Nova Bold}
\newfontface{\cmssbxsectionRomeo}[Ligatures=Common,Scale=1.4]{Gill Sans Nova Bold}
\newfontface{\cmssbxelevenRomeo}[Ligatures=Common,Scale=1.1]{Gill Sans Nova Bold}
\newfontface{\cmssbxchaptocRomeo}[Ligatures=Common,Scale=1]{Gill Sans Nova Bold}
\newfontface{\cmssbxchaptitleRomeo}[Ligatures=Common,Scale=3.0]{Gill Sans Nova Bold}
\newfontface{\cmsschapnameRomeo}[Ligatures=Common,Scale=2.6]{Gill Sans Nova Medium}
\newfontface{\cmssbxpartRomeo}[Ligatures=Common,Scale=4.4]{Gill Sans Nova Bold}
\newfontface{\cmssparttitleRomeo}[Ligatures=Common,Scale=3.4]{Gill Sans Nova Medium}
\newfontface{\sfiHuge}[Ligatures=Common,Scale=2.5]{Gill Sans Nova Medium Italic}
\newfontface{\sfititleRomeo}[Ligatures=Common,Scale=2.4]{Gill Sans Nova Medium Italic}
\newfontface{\sfihalftitleRomeo}[Ligatures=Common,Scale=2.0]{Gill Sans Nova Medium Italic}
\newfontface{\sfiauthorRomeo}[Ligatures=Common,Scale=1.6]{Gill Sans Nova Medium Italic}
\newfontface{\sfblargeRomeo}[Ligatures=Common,Scale=1.2]{Gill Sans Nova Bold}
\newfontface{\sfbelevenRomeo}[Ligatures=Common,Scale=1.1]{Gill Sans Nova Bold}
\newfontface{\sfelevenRomeo}[Ligatures=Common,Scale=1.1]{Gill Sans Nova Medium}
\newfontface{\sftwelveRomeo}[Ligatures=Common,Scale=1.2]{Gill Sans Nova Medium}
\newfontface{\sfthirteenRomeo}[Ligatures=Common,Scale=1.3]{Gill Sans Nova Medium}
\newfontface{\sfbRomeo}[Ligatures=Common,Scale=1]{Gill Sans Nova Bold}
\newfontface{\sfbiRomeo}[Ligatures=Common,Scale=1]{Gill Sans Nova Bold Italic}
\newfontface{\sfeightRomeo}[Ligatures=Common,Scale=.8]{Gill Sans Nova Medium}
\newfontface{\sfnineRomeo}[Ligatures=Common,Scale=.9]{Gill Sans Nova Medium}
\newfontface{\sfbnineRomeo}[Ligatures=Common,Scale=.9]{Gill Sans Nova Bold}
\newfontface{\defnemRomeo}[Ligatures=Common,Scale=1]{Gill Sans Nova Bold}
\newfontface{\sfbitwelveRomeo}[Ligatures=Common,Scale=1.2]{Gill Sans Nova Bold Italic}
%%%%%
\newfontface{\sfbsectionRomeo}[Ligatures=Common,Scale=1.3]{Gill Sans Nova Bold}
\newfontface{\sfbsectionitalRomeo}[Ligatures=Common,Scale=1.4]{Gill Sans Nova Bold Italic}
\newfontface{\sfbsubsecRomeo}[Ligatures=Common,Scale=1.2]{Gill Sans Nova Bold}
\newfontface{\sfbsubsubRomeo}[Ligatures=Common,Scale=1.1]{Gill Sans Nova Bold}
\newfontface{\sfbparaRomeo}[Ligatures=Common,Scale=1]{Gill Sans Nova Bold}
\newfontface{\sfbsubsecitalRomeo}[Ligatures=Common,Scale=1.2]{Gill Sans Nova Bold Italic}
\newfontface{\sfbsubsubsecitalRomeo}[Ligatures=Common,Scale=1.1]{Gill Sans Nova Bold Italic}
\newfontface{\sfbiRHRomeo}[Ligatures=Common,Scale=1.3]{Gill Sans Nova Bold Italic}
\newfontface{\sfiRHRomeo}[Ligatures=Common,Scale=1.3]{Gill Sans Nova Medium Italic}
\newfontface{\sfRHRomeo}[Ligatures=Common,Scale=1.3]{Gill Sans Nova Medium}

%% redefining existing sans serif font commands
\renewcommand{\sfbHuge}{\sfbHugeRomeo}       
\renewcommand{\cmssbxparttoc}{\cmssbxparttocRomeo}
\renewcommand{\cmssbxsection}{\cmssbxsectionRomeo}
\renewcommand{\cmssbxeleven}{\cmssbxelevenRomeo}
\renewcommand{\cmssbxchaptoc}{\cmssbxchaptocRomeo}    
\renewcommand{\cmssbxchaptitle}{\cmssbxchaptitleRomeo}
\renewcommand{\cmsschapname}{\cmsschapnameRomeo}
\renewcommand{\cmssbxpart}{\cmssbxpartRomeo}
\renewcommand{\cmssparttitle}{\cmssparttitleRomeo} 
\renewcommand{\sfiHuge}{\sfiHugeRomeo}
\renewcommand{\sfititle}{\sfititleRomeo}
\renewcommand{\sfihalftitle}{\sfihalftitleRomeo} 
\renewcommand{\sfiauthor}{\sfiauthorRomeo} 
\renewcommand{\sfblarge}{\sfblargeRomeo} 
\renewcommand{\sfbeleven}{\sfbelevenRomeo}  
\renewcommand{\sfb}{\sfbRomeo} 
\renewcommand{\sfeight}{\sfeightRomeo}
\renewcommand{\sfnine}{\sfnineRomeo}
\renewcommand{\defnem}{\defnemRomeo}                  
%%%%
\renewcommand{\sfbsection}{\sfbsectionRomeo}
\renewcommand{\sfbsubsec}{\sfbsubsecRomeo}
\renewcommand{\sfbsubsub}{\sfbsubsubRomeo}
\renewcommand{\sfbpara}{\sfbparaRomeo}

%%%%%%% emoji fonts
\newcommand{\martini}{\includegraphics[scale=.0125]{1F378_color}}
% this is an open source emoji image.  Make sure to note for Right & Reprints. 

%%%%%%% strikeout text function
\usepackage[normalem]{ulem}  % used for a meredith citation
%%%%%%%%%%%%%%

%%%%%%%%%% getting picky about hyphenation
\hyphenation{name-space}
\hyphenation{an-aly-sis}
%%%%%%%%%%%%%

\newcommand\subsubemdash{\sfbsubsubRomeo\char"2014} %% for whatever reason, emdashes are not displaying properly in subsubsection heads. This creates a direct command to create them. 

\newcommand\romeovalue[1]{\textit{#1}}  % for normal text usage
\newcommand\romeovalueinside{\itshape} % for inside other environments, like table cells
\newcommand\romeogloss[1]{\textbf{#1}}
\newcommand\intraref[2]{\textit{\titleref{#1}~\intrarefdelim~\titleref{#2}} on page~\pageref{#2}}
\newcommand\intrarefsimple[1]{\textit{\titleref{#1}} on page~\pageref{#1}}
\newcommand\intrarefnopage[1]{\textit{\titleref{#1}}}
%\newcommand\seealsoref[2]{``\titleref{#1}"~({#2}, p.~\pageref{#1})~\seealsodelim~}
\newcommand\seealsoref[2]{``\titleref{#1}"~({#2}, p.~\pageref{#1})~}
\newcommand\featureref[2]{{#1}.``\titleref{#2}" on page~\pageref{#2}}
\newcommand{\intrarefdelim}{---}
\newcommand{\seealsodelim}{~$\blacklozenge$~}
\usepackage{ifthen}

\usepackage{cite}
%%%%%%%%%%%%%%%%% to set citations and biblabels bold 
\makeatletter
\renewcommand\@cite[2]{%
{\bfseries #1}}\ifthenelse{\boolean{@tempswa}}{,\nolinebreak[3] #2}{}
%\makeatletter
\renewcommand\@biblabel[1]{\MakeLowercase{\bfseries #1}}  % sets bib label citations are bold
\makeatother
%%%%%%%%%%%%%%%%

% Because the cppxx in the RH takes the moniker of "section" per the Au's requirement even though the Au requires that it not exist on the page and because each "feature" is a section and is unnumbered and because the Au decided when the design was almost complete that a section xref was needed for each feature, these manual commands will take the place of a working \ref for each cppxx fake section. Readdress this for the second edition and come up with a better solution; that will probably involve redefining the section command, so cppxx can be an actual section, or creating a new section-like environment for cppxx.   
\newcommand{\locationa}{Section~1.1}
\newcommand{\locationb}{Section~1.2}
\newcommand{\locationc}{Section~2.1}
\newcommand{\locationd}{Section~2.2}
\newcommand{\locatione}{Section~3.1}
\newcommand{\locationf}{Section~3.2}
\newcommand{\seealsolocationa}{{\S}1.1}
\newcommand{\seealsolocationb}{{\S}1.2}
\newcommand{\seealsolocationc}{{\S}2.1}
\newcommand{\seealsolocationd}{{\S}2.2}
\newcommand{\seealsolocatione}{{\S}3.1}
\newcommand{\seealsolocationf}{{\S}3.2}




%%% LORI
\usepackage{cprotect}

\usepackage{amsfonts} % to ensure attractive mathematics
\usepackage{amssymb} % to ensure attractive mathematics
\usepackage{float}      % to allow things to float even if they normally don't
\usepackage{wrapfig} % to allow text to wrap around small figures
\usepackage{outline} % to allow outline creation
\usepackage{framed-PearsonGeneric} % to create the leftside vertical bar for the example environment

\usepackage{listings}  % for setting attractive computer code
\makeatletter
\@addtoreset{lstlisting}{section}  % to renumber listings per section
\makeatother

% All code samples that should be tested are going to be written within the
% following environment
%
% Source for initial implementation for dumping to files:
% https://tex.stackexchange.com/questions/284189/how-to-dump-listings-to-a-file
%
% Change the following in existing features:
%    \begin{lstlisting}[language=C++]   \begin{emcppslisting}
%    \end{lstlisting}                   \end{emcppslisting}
%
%\usepackage{currfile}
\makeatletter

% This counter is used to generate unique ids for the files that will be output.
% Ideally it would also be reset per feature, but that is not strictly needed.
\newcounter{emcppslistings}
\setcounter{emcppslistings}{0}

% A batch id for code snippets that should be concatenated together to compile.
\lst@Key{emcppsbatch}{}{\def\emcppsBatch{#1}}

% A reason why this code snippet should not be compiled, such as code that is
% explained outside of the example as an error, or code that is more
% descriptive than actually valid code.  Useless and irrelevant on
% emcppshiddenslistings.
\lst@Key{emcppsignore}{}{\def\emcppsIgnore{#1}}

% Which standards the code example should be compiled for, which will be passed
% to the -std= argument of gcc or clang  This must be on the first snippet in a
% batch, otherwise a default oc "c++11,c++14" will be used.
% Example value [emcppsstandards={c++11,c++17}]
\lst@Key{emcppsstandards}{}{\def\emcppsStandards{#1}}

% Provides a list of line numbers within this snippet that are errors.
% Useless and irrelevant on emcppshiddenslistings.
\lst@Key{emcppserrorlines}{}{\def\emcppsErrorLines{#1}}

% Provides an alternate output file name for this snippet.  Multiple snippets
% in the same batch with the same output file name will be concatenated.
% This is equivalent to a comment "// <filename.(h|cpp)>" in the snippet, but
% is not visible in the text.
\lst@Key{emcppsoutputfile}{}{\def\emcppsOutputFile{#1}}

% This environment should be used for all C++ listings/code examples that
% readers should see.  Note that in other builds this will write snippets to
% files in order to validate compilation.
\lstnewenvironment{emcppslisting}[1][]
{%
%  \lstset{language=C++,#1}%
  \lstset{style=Romeo,#1}%  %% change from josh 5 march 2021
  \stepcounter{emcppslistings}%
} 
{}

% This environment should be used for adding "hidden" parts of a listing batch 
% are needed to make it compile (use [emcppsbatch=...] on all related snippets}
% This should take up no space and nothing should be visible in the text.  Note
% that in other builds this will write snippets to files in order to validate
% compilation.
\lstnewenvironment{emcppshiddenlisting}[1][]
{%
%  \lstset{language=C++,aboveskip=0pt,belowskip=0pt,gobble=1000,#1}%  
    \lstset{style=Romeo,aboveskip=0pt,belowskip=0pt,gobble=1000,#1}%  %% change from josh 5 march 2021
  \stepcounter{emcppslistings}%
}
{}

\makeatother



\definecolor{maroon}{cmyk}{.26,1,.63,.17}
\definecolor{skyblue}{cmyk}{.62,.07,0,0}
\definecolor{forestgreen}{cmyk}{.9,0,.9,0}
\definecolor{mustard}{cmyk}{.22,.14,1,0}
\definecolor{purple}{cmyk}{.83,.81,0,0}
\definecolor{codegray}{cmyk}{0,0,0,.70}

%%%%% CHANGED FOR ROMEO
% changed font to Cousine, changed to 2ex indent

%%% PRINT BOOK SETTINGS
%\lstset{language=C++,basicstyle=\ttfamily\small,framerule=0.5pt,
%emph=[1]{const},
%  emphstyle=[1]{\color{forestgreen}\ttfamily},
%  moreemph=[2]{int,void,char},
%  emphstyle=[2]{\color{forestgreen}\ttfamily},
%    moreemph=[3]{decltype},
%      emphstyle=[3]{\color{forestgreen}\ttfamily},
%    keywordstyle={\color{forestgreen}\ttfamily\bfseries},
%  stringstyle={\color{forestgreen}},
%commentstyle={\color{codegray}\itshape},
%numberstyle=\color{forestgreen},
%xleftmargin=2ex,
%    escapeinside={(ù}{ù)},
%    showstringspaces=false, 
%} 
% setting the code in small sans serif font and
% setting the margin for the code.  Nov 2015: setting the rules to be 0.5pt
%%%%%%%%%%%%%
%\newcommand{\codeincomments}{\color{codegray}\ttfamily}
%%%%%%%%%%%%%%%%%%%%%%%%%%%%%


%%% EBOOK SETTINGS
\lstset{language=C++,basicstyle=\ttfamily\small,framerule=0.5pt,
emph=[1]{const},
  emphstyle=[1]{\color{mustard}\ttfamily},
  moreemph=[2]{int,void,char},
  emphstyle=[2]{\color{maroon}\ttfamily},
    moreemph=[3]{decltype},
      emphstyle=[3]{\color{forestgreen}\ttfamily},
    keywordstyle={\color{forestgreen}\ttfamily\bfseries},
  stringstyle={\color{purple}},
commentstyle={\color{skyblue}\itshape},
numberstyle=\color{purple},
xleftmargin=2ex,
    escapeinside={(ù}{ù)},
    showstringspaces=false, 
} 

\lstdefinestyle{plain}{basicstyle=\ttfamily\small,framerule=0.5pt,xleftmargin=2ex,
emph=[1]{const},
  emphstyle=[1]{\ttfamily},
  moreemph=[2]{int,void,char},
  emphstyle=[2]{\ttfamily},
    moreemph=[3]{decltype},
      emphstyle=[3]{\ttfamily},
    keywordstyle={\ttfamily},
  stringstyle={},
commentstyle={\ttfamily},
numberstyle={},
    escapeinside={(ù}{ù)},
    showstringspaces=false, 
}

% setting the code in small sans serif font and
%% setting the margin for the code.  Nov 2015: setting the rules to be 0.5pt
%%%%%%%%%%%%%%
\newcommand{\codeincomments}{\color{skyblue}\ttfamily}
\newcommand{\emphincomments}{\color{skyblue}\ttfamily}

\lstdefinestyle{footcode}{language=C++,
basicstyle=\ttfamily\footnotesize,framerule=0.5pt,
keywordstyle={\ttfamily\bfseries\footnotesize},
commentstyle={\color{grey}\ttfamily},
xleftmargin=2ex,
escapeinside={(ù}{ù)},
showstringspaces=false, 
}  
%%%%%%%%%%%%%%%%%%%%%%%%%%%%%%%%%%%%%%%%
\newcommand{\fncodeincomments}{\color{grey}\ttfamily}
%%%%%%%%%%%%%%%%%


%%%%%%%%%%%%%%%%%%%%%%%%%%%%%%%

\newcommand{\authorsnote}{Authors' Note: }
%\newcommand{\intraref}[1]{{\textit{#1}}}
\newcommand{\intrarefsub}[1]{\textit{#1}}  % needed for V's markdown conversion


\usepackage[stable,bottom,multiple]{footmisc} % to stabilize the footnote environment and to ensure footnotes
							  % appear at the bottom of the page, even if a figure or table
							  % appears at the bottom
\usepackage[flushleft]{threeparttable} % to set attractive footnotes within table environments
\usepackage{titleref} % to allow cross-referencing by title
\usepackage{xr} % to allow cross-referencing between files
\usepackage{colortbl} % to allow the use of color/shading in tables
\usepackage{caption} % to set the appearance of captions
\captionsetup[table]{singlelinecheck=off,labelfont={sf,bf,small},font={sf,bf,small}}
\captionsetup[figure]{labelfont={sf,bf,small},font={sf,bf,small}}
\captionsetup[lstlisting]{labelfont={sf,bf,small},font={sf,bf,small}}
\usepackage{tocloft-hacked-PearsonGeneric-ROMEO} % to allow manipulation of the TOC
\frenchspacing



%%%%%%%% Determine your indexing plan.
%
% Most of the time, the publisher will hire an independent indexing expert to create an
%    index for your book.  Your production editor can let you know the plan for your book.
% Use the line below only if your production plan includes the creation of a LaTeX index.
% Note that this situation is uncommon.
%\makeindex           % if your book will include a LaTeX-generated index
%\usepackage{See-makeidx}

%%%%%%%%%% These lines set up the Preamble of the document. You can skip over this part.
%% Preamble:
%%%%%%%%%%%%%%%%%%%%%%%%%%%%%
%%% Altered for Lakos & Romeo book, 2020
% Changes marked with ROMEO-


% New commands and/or command redefinitions
%
% You can also place such commands in a macro file (e.g. mydefs.tex)
% and load them in the preamble with "input mydefs"
% Here are some examples:

\newcommand{\be}{\begin{equation}}           %---- begin numbered equation
\newcommand{\ee}{\end{equation}}             %---- end numbered equation
\newcommand{\dst}{\everymath{\displaystyle}} %---- use displaystle in eqs.
\newcommand{\Nuclide}[2]{${}^{#2}$#1}        %---- Nuclide macro
\renewcommand{\SS}[1]{${}^{#1}$}             %---- superscript in text
\newcommand{\STRUT}{\rule{0in}{3ex}}         %--- small strut
\newcommand{\dotprod}{{\scriptscriptstyle \stackrel{\bullet}{{}}}}
\renewcommand{\tilde}{\widetilde}
\renewcommand{\hat}{\widehat}
\renewcommand{\theequation}{\arabic{chapter}.\arabic{equation}}
%\renewcommand{\thefigure}{\arabic{chapter}.\arabic{section}--\arabic{figure}}  %%% Romeo - reset per section


\renewcommand{\thetable}{\arabic{table}}   %%% Romeo - reset per section
\newcommand{\ul}{\underline}
\newcommand{\dul}{\underline{\underline}}
\newcommand{\ol}{\overline}
\newcommand{\beq}{\begin{equation}}
\newcommand{\eeq}{\end{equation}}
\newcommand{\bea}{\begin{eqnarray}}
\newcommand{\eea}{\end{eqnarray}}
\newcommand{\beas}{\begin{eqnarray*}}
\newcommand{\eeas}{\end{eqnarray*}}
\newcommand{\ba}{\begin{array}}
\newcommand{\ea}{\end{array}}

%%%%%%%% set chapter author commands
\newcommand{\chapauthor}[1]{\hspace*{\fill}{\textsf{\large 
 #1}}\\[1.5ex]}
%%%%%%%%%%%%%%

%%%%%%% set epigraph commands
\newcommand{\epigraph}[1]{\vspace*{2ex}{\raggedleft{\textit{#1}}\\[6ex]}}
%%%%%%%%%%%%

%%%%%% set definition commands
\newcounter{definitionctr}[chapter]
\def\thedefinitionctr{\arabic{chapter}--\arabic{definitionctr}}
\newcommand{\definition}[1]{\refstepcounter{definitionctr}\par\vskip2ex \indent\parbox[]{.96\textwidth}{\indent{\textbf{Definition~\thedefinitionctr.}~~#1}\par\vskip2ex}}
%%%%%%%%%%%%%%%%%

%%%%%%%%%%%%%%% create a thick line for tables
\newcommand{\thickhline}{\noalign{\hrule height 1.5pt}} 
%%%%%%%%%%%%%

%%%%%%%%%%%%%% create a medium line for tables
\newcommand{\medhline}{\noalign{\hrule height 1.25pt}}
%%%%%%%%%%%%%

%%%%%%%%%%%%%% define the unnumbered list
\newenvironment{unnumlist}{\begin{list}{}% empty for no label
{}} % empty since we don't need a counter and the standard lengths/indentations are fine
{\end{list}} 
%%%%%%%%%%%%%%%%%%%

%%%%%%%%%%% define a simple command for the quote environment
\newenvironment{italquote}{\begin{quote}{}% empty for no label
{}\itshape} % empty since we don't need a counter and the standard lengths/indentations are fine
{\end{quote}} 
%%%%%%%%%%%

%%%%%%%%%%% define a simple command for the quote environment
\newenvironment{dialogue}[2]{\begin{quote}{\bfseries #1:~}% empty for no label
{} #2} % empty since we don't need a counter & the standard lengths/indentations and text style are fine
{\end{quote}} 
%%%%%%%%%%%

%%%%%%%%%%%%% define the multicolumn list environment
\newcommand{\mclhead}{\bfseries} % define a multi-column list (mcl) heading
\newlength{\mcltopsep}% define a spacing command
\setlength{\mcltopsep}{\topsep} % 
\newenvironment{mcl}{\par\vspace*{\mcltopsep}
\begin{tabular}{ll}}{\end{tabular}\linebreak[4] \par\vspace*{.5\mcltopsep}}
%%%%%%%%%%%%%%%%

%%%%%%%%%%%% define some sidebar commands
% the sidebar is created via the framed.sty package
\definecolor{shadecolor}{gray}{0.9} % define the color to be used for the sidebar
% define the sidebar head  ---  corrected bug 4/19/2011
\newcommand{\sidebarhead}[1]{\noindent{\cmssbxsection #1}\hspace*{\fill}\linebreak\nopagebreak} 
%%%%%%%%%%%%%%

%%%%%%%%%%% define note, tip, and warning environments
\newcommand{\note}[2]{\pagebreak[3]\par\vspace*{2ex} % define a new note command 
\noindent{\sffamily\bfseries\small NOTE}\linebreak\vspace*{-1.5ex}\nopagebreak
\rule[1.5ex]{\textwidth}{.75pt}\linebreak\nopagebreak
{\sffamily\small  {\bfseries{\raggedright #1}\\}\nopagebreak
\noindent{#2}}\par
\noindent\rule[1ex]{\textwidth}{.75pt}\par\vspace*{1ex}\pagebreak[3] 
}

\newcommand{\tip}[2]{\pagebreak[3]\par\vspace*{2ex} % define a new tip command 
\noindent{\sffamily\bfseries\small TIP}\linebreak\vspace*{-1.5ex}\nopagebreak
\rule[1.5ex]{\textwidth}{.75pt}\linebreak\nopagebreak
{\sffamily\small  {\bfseries{\raggedright #1}\\}\nopagebreak
\nopagebreak
\noindent{#2}}\par
\noindent\rule[1ex]{\textwidth}{.75pt}\par\vspace*{1ex}\pagebreak[3]
}

\newcommand{\warning}[2]{\pagebreak[3]\par\vspace*{2ex} % define a new warning command 
\noindent{\sffamily\bfseries\small WARNING}\linebreak\vspace*{-1.5ex}\nopagebreak
\rule[1.5ex]{\textwidth}{.75pt}\linebreak\nopagebreak
{\sffamily\small  {\bfseries{\raggedright #1}\\}\nopagebreak
\noindent{#2}}\par
\noindent\rule[1ex]{\textwidth}{.75pt}\par\vspace*{1ex}\pagebreak[3]
}
%%%%%%%%%%%%%%

%%%%%%%%%%% define glossary command
\newcommand{\gloss}[2]{\pagebreak[3]\par
\noindent{\sffamily\bfseries #1}\\
\noindent{#2}\hspace*{\fill}\\ 
\par\pagebreak[3]}
%%%%%%%%%%%%%%%%

%%%%%%%%%%%%%% ROMEO - reduce indent of nested itemize
%%%% NOTE than any nested enumerate will need to be set within the list itself
\setlength{\leftmarginii}{4mm}




  % Loading the definitions needed for this book:

%%%%%%%%%%CHANGED FOR ROMEO
\setcounter{secnumdepth}{0}

% Required for PANDOC conversion
\newcommand{\passthrough}[1]{#1}
%
%% Hack to allow PANDOC to parse Markdown in inline LaTeX environments
\let\Begin\begin
\let\End\end

% PANDOC stuff
%\IfFileExists{upquote.sty}{\usepackage{upquote}}{}
%\IfFileExists{microtype.sty}{% use microtype if available
%  \usepackage[]{microtype}
%  \UseMicrotypeSet[protrusion]{basicmath} % disable protrusion for tt fonts
%}{}
%\makeatletter
%\@ifundefined{KOMAClassName}{% if non-KOMA class
%  \IfFileExists{parskip.sty}{%
%    \usepackage{parskip}
%  }{% else
%    \setlength{\parindent}{0pt}
%    \setlength{\parskip}{6pt plus 2pt minus 1pt}}
%}{% if KOMA class
%  \KOMAoptions{parskip=half}}
%\makeatother
%\usepackage{xcolor}
% \IfFileExists{xurl.sty}{\usepackage{xurl}}{} % add URL line breaks if available

% \IfFileExists{bookmark.sty}{\usepackage{bookmark}}{\usepackage{hyperref}}
% \hypersetup{
% % % % % % %   hidelinks,
% %   pdfcreator={LaTeX via pandoc}}


% \urlstyle{same} % disable monospaced font for URLs


% 
% 


%\setlength{\emergencystretch}{3em} % prevent overfull lines
%\providecommand{\tightlist}{%
%\setlength{\itemsep}{0pt}\setlength{\parskip}{0pt}}

% 




%%%%%%%%% Now the document begins

\begin{document}
%%%%%% commands to fine-tune to table of contents
%\cftnodots
%\renewcommand{\@dotsep}{10000}


\renewcommand{\cfttoctitlefont}{\cmsschapname}
\renewcommand{\cftpartfont}{\cmssbxparttoc}
\renewcommand{\cftchappagefont}{\cmssbxparttoc}
\renewcommand{\cftchapfont}{\cmssbxchaptoc}
\renewcommand{\cftchapnumwidth}{5.5em}
\renewcommand{\cftchappagefont}{\cmssbxchaptoc}
\renewcommand{\cftsecfont}{\sf}
\renewcommand{\cftsecindent}{1.25em}
%\renewcommand{\cftsecnumwidth}{0em}
\renewcommand{\cftsecpagefont}{\sf}
\renewcommand{\cftsubsecfont}{\sf}
\renewcommand{\cftsubsecindent}{4em}
%\renewcommand{\cftsubsecnumwidth}{0em}
\renewcommand{\cftsubsecpagefont}{\sf}
\renewcommand{\contentsname}{{Contents}} 
\setcounter{tocdepth}{1}



\setcounter{secnumdepth}{0} 
%\renewcommand{\numberline}[2]{} 
%%%%% end toc commands % Loading a few commands related to the tocloft package

\renewcommand{\thelstlisting}{\thesection--\arabic{lstlisting}}  %%% Romeo - reset listings per section


%%%%% The front matter begins:
\pagenumbering{roman}             % Roman numbering
\setcounter{page}{1}                    % Starting page number (specified by the publisher)
\cleardoublepage                         % This command is defined in our documentclass file. It tells 
                                                      % our compiler that we want the next item 
                                                       % to start on a right-hand page, and it inserts a blank left-hand 
                                                       % page if needed.    
\thispagestyle{empty}
\begin{center}
{\sfititle Embracing Modern C++ Safely}
\end{center}                        % The half-title page
\thispagestyle{empty}
This is simply a placeholder. Your production team will replace this page with the real series page.                              % A placeholder for the series page. The publisher normally 
%                                                       % creates this page and just adds it into the final PDF. 
\thispagestyle{empty}
\begin{center}
{\sfititle Embracing Modern C++ Safely}\\
\vspace*{54pt}
%{\sfihalftitle Subtitle of the Book}\\
\vspace*{\fill}    % This space will automatically adjust
\noindent{\sfiauthor John Lakos}\\[3ex]
\noindent{\sfiauthor Vittorio Romeo}\\        % Add or delete lines as necessary for author names

\vspace*{\fill}    % This space will automatically adjust
% Authors:  For the next line, use either the AW line or the PH line per your 
%   production editor's instructions. You will find both PC and Mac versions of the logos 
%   in the template folder. Use the one appropriate for your system.
%\includegraphics{PH_title_page_Mac} \\   %PH Mac
%\includegraphics{PH_stack_BLACK-PC} \\  %PH PC
\includegraphics{AWlogoTitlePage-Mac}\\   % AW Mac
%\includegraphics{AWlogoTitlePage_PC}  \\  %AW PC
{\footnotesize 
Boston~$\bullet$~Columbus~$\bullet$~Indianapolis~$\bullet$~New York~$\bullet$~San Francisco~$\bullet$~Amsterdam~$\bullet$~Cape Town\\[1ex]
Dubai~$\bullet$~London~$\bullet$~Madrid~$\bullet$~Milan~$\bullet$~Munich~$\bullet$~Paris~$\bullet$~Montreal~$\bullet$~Toronto~$\bullet$~Delhi~$\bullet$~Mexico City\\[1ex]
Sao Paulo~$\bullet$~Sidney~$\bullet$~Hong Kong~$\bullet$~Seoul~$\bullet$~Singapore~$\bullet$~Taipei~$\bullet$~Tokyo\\[1ex]
}
\end{center}                               % The title page
\thispagestyle{empty}
\vspace*{\fill}
{\footnotesize\raggedright 
\noindent Many of the designations used by manufacturers and sellers to distinguish their products are claimed as trademarks. Where those designations appear in this book, and the publisher was aware of a trademark claim, the designations have been printed with initial capital letters or in all capitals.\\
\hspace{\fill}\\
\noindent The authors and publisher have taken care in the preparation of this book, but make no expressed or implied warranty of any kind and assume no responsibility for errors or omissions. No liability is assumed for incidental or consequential damages in connection with or arising out of the use of the information or programs contained herein.\\   % change to authors (plural) if this book has multiple authors
\hspace{\fill}\\
\noindent For information about buying this title in bulk quantities, or for special sales opportunities (which may include electronic versions; custom cover designs; and content particular to your business, training goals, marketing focus, or branding interests), please contact our corporate sales department\linebreak[4] at \mbox{corpsales}@\mbox{pearsoned.com} or (800) 382-3419.\\
\hspace{\fill}\\
\noindent For government sales inquiries, please contact governmentsales@pearsoned.com.\\
\hspace{\fill}\\
\noindent For questions about sales outside the U.S., please contact intlcs@pearson.com.\\
\hspace{\fill}\\
\noindent Visit us on the Web: informit.com/aw\\     % use this line if you are publishing with AW
%\noindent Visit us on the Web: informit.com/ph\\   % use this line if you are publishing with PH
\hspace{\fill}\\
\noindent Library of Congress Control Number: TO COME FROM ITP\\  % update this with info provided by your production editor
\hspace{\fill}\\
\noindent Copyright~\copyright~2021 Pearson Education, Inc.\\
\hspace{\fill}\\
\noindent Cover image: [TO COME FROM ITP]\\
\hspace*{\fill}\\
\noindent All rights reserved. This publication is protected by copyright, and permission must be obtained from the publisher prior to any prohibited reproduction, storage in a retrieval system, or transmission in any form or by any means, electronic, mechanical, photocopying, recording, or likewise. For information regarding permissions, request forms and the appropriate contacts within the Pearson Education Global Rights \& Permissions Department, please visit www.pearson.com/permissions.\\
\hspace*{\fill}\\
\noindent ISBN-13: 978-0-13-738035-0\\[-.5ex] % update this with info provided by your production editor 
\noindent ISBN-10: 0-13-738035-6 \\ % update this with info provided by your production editor
\hspace*{\fill}\\
\noindent{\sf ScoutAutomatedPrintCode}
%\noindent Text printed in the United States on recycled paper at PRINTER INFO HERE.\\[-.5ex] % update this with info provided by your production editor
%\noindent First printing, MONTH YEAR % update this with info provided by your production editor
}                                % The cataloging-in-publication page
\thispagestyle{empty}
\begin{center}The lines of text of the dedication will be broken\\ 
with the help of the production team\\ 
to ensure an attractive\\
layout.\\[1ex] 
Author Initials\\
\vspace*{3pc}
Should there be a second author, that author's lines of\\  % These lines are 
dedication text will go here and will also be broken\\        % commented out 
in whatever way is attractive and logical.\\[1ex]                % if unneeded.
Author Initials\\
\end{center}                      % The dedication page
\cleardoublepage
\parskip .5ex                       % Add space between Contents items
\tableofcontents                  % Make the Table of Contents 
\cleardoublepage                
\parskip 0in                         % Reset to zero interpar spacing
\cleardoublepage
\chapter*{Foreword}
\markboth{Foreword}{Foreword}

The text of the foreword will go here.                 % Add foreword
\cleardoublepage
\chapter*{Preface}

The text of the preface will go here.  All the elements and commands offered within the text will work here as well. 

Sed ut perspiciatis unde omnis iste natus error sit voluptatem accusantium doloremque laudantium, totam rem aperiam, eaque ipsa quae ab illo inventore veritatis et quasi architecto beatae vitae dicta sunt explicabo. Nemo enim ipsam voluptatem quia voluptas sit aspernatur aut odit aut fugit, sed quia consequuntur magni dolores eos qui ratione voluptatem sequi nesciunt. Neque porro quisquam est, qui dolorem ipsum quia dolor sit amet, consectetur, adipisci velit, sed quia non numquam eius modi tempora incidunt ut labore et dolore magnam aliquam quaerat voluptatem. Ut enim ad minima veniam, quis nostrum exercitationem ullam corporis suscipit laboriosam, nisi ut aliquid ex ea commodi consequatur? Quis autem vel eum iure reprehenderit qui in ea voluptate velit esse quam nihil molestiae consequatur, vel illum qui dolorem eum fugiat quo voluptas nulla pariatur? 

At vero eos et accusamus et iusto odio dignissimos ducimus qui blanditiis praesentium voluptatum deleniti atque corrupti quos dolores et quas molestias excepturi sint occaecati cupiditate non provident, similique sunt in culpa qui officia deserunt mollitia animi, id est laborum et dolorum fuga. Et harum quidem rerum facilis est et expedita distinctio. Nam libero tempore, cum soluta nobis est eligendi optio cumque nihil impedit quo minus id quod maxime placeat facere possimus, omnis voluptas assumenda est, omnis dolor repellendus. Temporibus autem quibusdam et aut officiis debitis aut rerum necessitatibus saepe eveniet ut et voluptates repudiandae sint et molestiae non recusandae. Itaque earum rerum hic tenetur a sapiente delectus, ut aut reiciendis voluptatibus maiores alias consequatur aut perferendis doloribus asperiores repellat.

Sed ut perspiciatis unde omnis iste natus error sit voluptatem accusantium doloremque laudantium, totam rem aperiam, eaque ipsa quae ab illo inventore veritatis et quasi architecto beatae vitae dicta sunt explicabo. Nemo enim ipsam voluptatem quia voluptas sit aspernatur aut odit aut fugit, sed quia consequuntur magni dolores eos qui ratione voluptatem sequi nesciunt.Neque porro quisquam est, qui dolorem ipsum quia dolor sit amet, consectetur, adipisci velit, sed quia non numquam eius modi tempora incidunt ut labore et dolore magnam aliquam quaerat voluptatem. Ut enim ad minima veniam, quis nostrum exercitationem ullam corporis suscipit laboriosam, nisi ut aliquid ex ea commodi consequatur? Quis autem vel eum iure reprehenderit qui in ea voluptate velit esse quam nihil molestiae consequatur, vel illum qui dolorem eum fugiat quo voluptas nulla pariatur? 

At vero eos et accusamus et iusto odio dignissimos ducimus qui blanditiis praesentium voluptatum deleniti atque corrupti quos dolores et quas molestias excepturi sint occaecati cupiditate non provident, similique sunt in culpa qui officia deserunt mollitia animi, id est laborum et dolorum fuga. Et harum quidem rerum facilis est et expedita distinctio. Nam libero tempore, cum soluta nobis est eligendi optio cumque nihil impedit quo minus id quod maxime placeat facere possimus, omnis voluptas assumenda est, omnis dolor repellendus. Temporibus autem quibusdam et aut officiis debitis aut rerum necessitatibus saepe eveniet ut et voluptates repudiandae sint et molestiae non recusandae. Itaque earum rerum hic tenetur a sapiente delectus, ut aut reiciendis voluptatibus maiores alias consequatur aut perferendis doloribus asperiores repellat.

Sed ut perspiciatis unde omnis iste natus error sit voluptatem accusantium doloremque laudantium, totam rem aperiam, eaque ipsa quae ab illo inventore veritatis et quasi architecto beatae vitae dicta sunt explicabo. Nemo enim ipsam voluptatem quia voluptas sit aspernatur aut odit aut fugit, sed quia consequuntur magni dolores eos qui ratione voluptatem sequi nesciunt.Neque porro quisquam est, qui dolorem ipsum quia dolor sit amet, consectetur, adipisci velit, sed quia non numquam eius modi tempora incidunt ut labore et dolore magnam aliquam quaerat voluptatem. Ut enim ad minima veniam, quis nostrum exercitationem ullam corporis suscipit laboriosam, nisi ut aliquid ex ea commodi consequatur? Quis autem vel eum iure reprehenderit qui in ea voluptate velit esse quam nihil molestiae consequatur, vel illum qui dolorem eum fugiat quo voluptas nulla pariatur? 

At vero eos et accusamus et iusto odio dignissimos ducimus qui blanditiis praesentium voluptatum deleniti atque corrupti quos dolores et quas molestias excepturi sint occaecati cupiditate non provident, similique sunt in culpa qui officia deserunt mollitia animi, id est laborum et dolorum fuga. Et harum quidem rerum facilis est et expedita distinctio. Nam libero tempore, cum soluta nobis est eligendi optio cumque nihil impedit quo minus id quod maxime placeat facere possimus, omnis voluptas assumenda est, omnis dolor repellendus. Temporibus autem quibusdam et aut officiis debitis aut rerum necessitatibus saepe eveniet ut et voluptates repudiandae sint et molestiae non recusandae. Itaque earum rerum hic tenetur a sapiente delectus, ut aut reiciendis voluptatibus maiores alias consequatur aut perferendis doloribus asperiores repellat.

Sed ut perspiciatis unde omnis iste natus error sit voluptatem accusantium doloremque laudantium, totam rem aperiam, eaque ipsa quae ab illo inventore veritatis et quasi architecto beatae vitae dicta sunt explicabo. Nemo enim ipsam voluptatem quia voluptas sit aspernatur aut odit aut fugit, sed quia consequuntur magni dolores eos qui ratione voluptatem sequi nesciunt. Neque porro quisquam est, qui dolorem ipsum quia dolor sit amet, consectetur, adipisci velit, sed quia non numquam eius modi tempora incidunt ut labore et dolore magnam aliquam quaerat voluptatem. Ut enim ad minima veniam, quis nostrum exercitationem ullam corporis suscipit laboriosam, nisi ut aliquid ex ea commodi consequatur? Quis autem vel eum iure reprehenderit qui in ea voluptate velit esse quam nihil molestiae consequatur, vel illum qui dolorem eum fugiat quo voluptas nulla pariatur? 

At vero eos et accusamus et iusto odio dignissimos ducimus qui blanditiis praesentium voluptatum deleniti atque corrupti quos dolores et quas molestias excepturi sint occaecati cupiditate non provident, similique sunt in culpa qui officia deserunt mollitia animi, id est laborum et dolorum fuga. Et harum quidem rerum facilis est et expedita distinctio. Nam libero tempore, cum soluta nobis est eligendi optio cumque nihil impedit quo minus id quod maxime placeat facere possimus, omnis voluptas assumenda est, omnis dolor repellendus. Temporibus autem quibusdam et aut officiis debitis aut rerum necessitatibus saepe eveniet ut et voluptates repudiandae sint et molestiae non recusandae. Itaque earum rerum hic tenetur a sapiente delectus, ut aut reiciendis voluptatibus maiores alias consequatur aut perferendis doloribus asperiores repellat.


Sed ut perspiciatis unde omnis iste natus error sit voluptatem accusantium doloremque laudantium, totam rem aperiam, eaque ipsa quae ab illo inventore veritatis et quasi architecto beatae vitae dicta sunt explicabo. Nemo enim ipsam voluptatem quia voluptas sit aspernatur aut odit aut fugit, sed quia consequuntur magni dolores eos qui ratione voluptatem sequi nesciunt. Neque porro quisquam est, qui dolorem ipsum quia dolor sit amet, consectetur, adipisci velit, sed quia non numquam eius modi tempora incidunt ut labore et dolore magnam aliquam quaerat voluptatem. Ut enim ad minima veniam, quis nostrum exercitationem ullam corporis suscipit laboriosam, nisi ut aliquid ex ea commodi consequatur? Quis autem vel eum iure reprehenderit qui in ea voluptate velit esse quam nihil molestiae consequatur, vel illum qui dolorem eum fugiat quo voluptas nulla pariatur? 

At vero eos et accusamus et iusto odio dignissimos ducimus qui blanditiis praesentium voluptatum deleniti atque corrupti quos dolores et quas molestias excepturi sint occaecati cupiditate non provident, similique sunt in culpa qui officia deserunt mollitia animi, id est laborum et dolorum fuga. Et harum quidem rerum facilis est et expedita distinctio. Nam libero tempore, cum soluta nobis est eligendi optio cumque nihil impedit quo minus id quod maxime placeat facere possimus, omnis voluptas assumenda est, omnis dolor repellendus. Temporibus autem quibusdam et aut officiis debitis aut rerum necessitatibus saepe eveniet ut et voluptates repudiandae sint et molestiae non recusandae. Itaque earum rerum hic tenetur a sapiente delectus, ut aut reiciendis voluptatibus maiores alias consequatur aut perferendis doloribus asperiores repellat.


Sed ut perspiciatis unde omnis iste natus error sit voluptatem accusantium doloremque laudantium, totam rem aperiam, eaque ipsa quae ab illo inventore veritatis et quasi architecto beatae vitae dicta sunt explicabo. Nemo enim ipsam voluptatem quia voluptas sit aspernatur aut odit aut fugit, sed quia consequuntur magni dolores eos qui ratione voluptatem sequi nesciunt. Neque porro quisquam est, qui dolorem ipsum quia dolor sit amet, consectetur, adipisci velit, sed quia non numquam eius modi tempora incidunt ut labore et dolore magnam aliquam quaerat voluptatem. Ut enim ad minima veniam, quis nostrum exercitationem ullam corporis suscipit laboriosam, nisi ut aliquid ex ea commodi consequatur? Quis autem vel eum iure reprehenderit qui in ea voluptate velit esse quam nihil molestiae consequatur, vel illum qui dolorem eum fugiat quo voluptas nulla pariatur? 

At vero eos et accusamus et iusto odio dignissimos ducimus qui blanditiis praesentium voluptatum deleniti atque corrupti quos dolores et quas molestias excepturi sint occaecati cupiditate non provident, similique sunt in culpa qui officia deserunt mollitia animi, id est laborum et dolorum fuga. Et harum quidem rerum facilis est et expedita distinctio. Nam libero tempore, cum soluta nobis est eligendi optio cumque nihil impedit quo minus id quod maxime placeat facere possimus, omnis voluptas assumenda est, omnis dolor repellendus. Temporibus autem quibusdam et aut officiis debitis aut rerum necessitatibus saepe eveniet ut et voluptates repudiandae sint et molestiae non recusandae. Itaque earum rerum hic tenetur a sapiente delectus, ut aut reiciendis voluptatibus maiores alias consequatur aut perferendis doloribus asperiores repellat.
                 % Add preface 
\cleardoublepage
\chapter*{Acknowledgements}

\noindent The text of the author's acknowledgements will go here.                         % Add acknowledgements 
\cleardoublepage
\chapter*{About the Authors}

 % wrapfig doesn't like to be immediately after a chapter/section command 
 % so we are tricking it a bit with the \hspace and the \par commands.
\hspace*{\fill}\\  
\par\begin{wrapfigure}[12]{L}[0pt]{0pt}\includegraphics{AuthorPhoto}\end{wrapfigure}\noindent AUTHOR BIO WILL APPEAR HERE.  Lorem ipsum dolor sit amet, consectetur adipisicing elit, sed do eiusmod tempor incididunt ut labore et dolore magna aliqua. Ut enim ad minim veniam, quis nostrud exercitation ullamco laboris nisi ut aliquip ex ea commodo consequat. Duis aute irure dolor in reprehenderit in voluptate velit esse cillum dolore eu fugiat nulla pariatur. Excepteur sint occaecat cupidatat non proident, sunt in culpa qui officia deserunt mollit anim id est laborum. Lorem ipsum dolor sit amet, consectetur adipisicing elit, sed do eiusmod tempor incididunt ut labore et dolore magna aliqua. Ut enim ad minim veniam, quis nostrud exercitation ullamco laboris nisi ut aliquip ex ea commodo consequat. Lorem ipsum dolor sit amet, consectetur adipisicing elit, sed do eiusmod tempor incididunt ut labore et dolore magna aliqua. Ut enim ad minim veniam, quis nostrud exercitation ullamco laboris nisi ut aliquip ex ea commodo consequat. 

\vspace*{18pt} % we add some space for a nice layout
\begin{wrapfigure}[12]{L}[0pt]{0pt}\includegraphics{AuthorPhoto}\end{wrapfigure} \noindent AUTHOR BIO WILL APPEAR HERE.  Lorem ipsum dolor sit amet, consectetur adipisicing elit, sed do eiusmod tempor incididunt ut labore et dolore magna aliqua. Ut enim ad minim veniam, quis nostrud exercitation ullamco laboris nisi ut aliquip ex ea commodo consequat. Duis aute irure dolor in reprehenderit in voluptate velit esse cillum dolore eu fugiat nulla pariatur. Excepteur sint occaecat cupidatat non proident, sunt in culpa qui officia deserunt mollit anim id est laborum. Lorem ipsum dolor sit amet, consectetur adipisicing elit, sed do eiusmod tempor incididunt ut labore et dolore magna aliqua. Ut enim ad minim veniam, quis nostrud exercitation ullamco laboris nisi ut aliquip ex ea commodo consequat. Lorem ipsum dolor sit amet, consectetur adipisicing elit, sed do eiusmod tempor incididunt ut labore et dolore magna aliqua. Ut enim ad minim veniam, quis nostrud exercitation ullamco laboris nisi ut aliquip ex ea commodo consequat. 
                   % Add about the author page
\cleardoublepage                  
%%%%%%%The front matter ends

\pagenumbering{arabic}		  % Arabic numbering
%%\part{This is the Part Title}
\vspace*{4pt}
\sffamily % this sets this page in sans serif font

\noindent Chapter \ref{ExampleChap}, "\titleref{ExampleChap},"  is followed by a one-sentence description of each chapter. 

Chapter \ref{samplechap2}, "\titleref{samplechap2},"  is followed by a one-sentence description of each chapter and if the sentence is longer what happens when it runs to two lines.

Chapter \ref{samplechap3}, "\titleref{samplechap3}," is followed by a one-sentence description of each chapter and if the sentence is longer what happens when it runs to two lines. 

\normalfont % this resets the fonts to our normal settings




            % Part I
\cleardoublepage
\setcounter{chapter}{-1}
\chapter[Introduction]{Introduction}\label{ch-intro}
% Andrei's version

Welcome! \textit{Embracing Modern C++ Safely} is a \textit{reference book} dedicated to professionals who want to leverage modern C++ features in the development and maintenance of large-scale, complex C++ software systems.

This book deliberately concentrates on the productive value afforded by each new language feature added by C++ starting with C++11, particularly when the systems and organizations involved are considered at scale. We left aside ideas and idioms, however clever and intellectually intriguing, that could hurt the bottom line when applied at large. Instead, we focus on what is objectively true and relevant to making wise economic and design decisions, with an understanding of the inevitable tradeoffs that arise in any engineering discipline. In doing so, we do our best to steer clear of subjective opinions and recommendations.

Richard Feynman famously said: ``If it disagrees with experiment, it’s wrong. In that simple statement is the key to science."\footnote{Richard Feynman, lecture at Cornell University, 1964. Video and commentary available at https://\linebreak[3]fs.\linebreak[3]blog/\linebreak[3]2009/12/mental-model-scientific-method.} There is no better way to experiment with a language feature than letting time do its work. We took that to heart by dedicating \textit{Embracing Modern C++ Safely} to only the features of Modern C++ that have been part of the Standard for at least five years, which grants enough perspective to properly evaluate its practical impact. Thus, we are able to provide you with a thorough and comprehensive treatment based on practical experience and worthy of your limited professional development time. If you're out there looking for tried and true ways to better use modern C++ features for improving your productivity, we hope this book will be the one you'll reach for.

What's missing from a book is as important as what's present. \textit{Embracing Modern C++ Safely} is not a tutorial on programming, on C++, or even on new features of C++. We assume you are an experienced developer, team lead, or manager, that you already have a good command of ``classic" C++98/03, and that you are looking for clear, goal-driven ways to integrate modern C++ features within your and your team's toolbox.

\section{What Makes This Book Different}

The book you're now reading aims very strongly at being objective, empirical, and practical. We simply present features, their applicability, and their potential pitfalls as reflected by the analysis of millions of human-hours of using C++11 and C++14 in the development of varied large-scale software systems; personal preference matters have been neutralized to our, and our reviewers', best ability. We wrote down the distilled truth that remains, which should shape your understanding of what modern C++ has to offer to you without being skewed by our subjective opinions or domain-specific inclinations.

The final analysis and interpretation of what is appropriate for your context is left to you, the reader. Hence, this book is, by design, not a C++ style or coding-standards guide; it would, however, provide valuable input to any development organization seeking to author or enhance one.

Practicality is a topic very important to us, too, and in a very real-world, economic sense. We examine modern C++ features through the lens of a large company developing and using software in a competitive environment. In addition to showing you how to best utilize a given C++ language feature in practice, our analysis takes into account the costs associated with having that feature employed routinely in the ecosystem of a software development organization. (We believe that costs of using language features are sadly neglected by most texts.) In other words, we weigh the benefits of successfully using a feature against the risk of its widespread ineffective use (or misuse) and/or the costs associated with training and code review required to reasonably ensure that such ill-conceived use does not occur. We are acutely aware that what applies to one person or small crew of like-minded individuals is quite different from what works with a large, distributed team. The outcome of this analysis is our signature categorization of features in terms of safety of adoption --- namely \textit{safe}, \textit{conditionally safe}, or \textit{unsafe} features.

We are not aware of any similar text amid the rich offering of C++ textbooks; in a very real sense, we wrote it because we needed it.

\section{Scope for the First Edition}

Given the vastness of C++'s already voluminous and rapidly growing standardized libraries, we have chosen to limit this book’s scope to just the language features themselves. A companion book, \textit{Embracing Modern C++ Standard Libraries Safely}, is a separate project that we hope to tackle in the future. However, to be effective, this book must remain small, concise, and focused on what expert C++ developers need to know well to be successful right now.

In this first of an anticipated series of periodically extended volumes, we characterize, dissect, and elucidate most of the modern language features introduced into the C++ Standard starting with C++11. We chose to limit the scope of this first edition to only those features that have been in the language Standard and widely available in practice for at least five years. This limited focus enables us to more fully evaluate the real-world impact of these features and to highlight any caveats that might not have been anticipated prior to standardization and sustained, active, and widespread use in industry.

We assume you are quite familiar with essentially all of the basic and important special-purpose features of classic C++98/03, so in this book we confined our attention to just the subset of C++ language features introduced in C++11 and C++14. This book is best for you if you need to know how to safely incorporate C++11/14 language features into a predominately C++98/03 code base, today.

Over time, we expect, and hope, that practicing senior developers will emerge entirely from the postmodern C++ era. By then, a book that focuses on all of the important features of modern C++ would naturally include many of those that were around before C++11. With that horizon in mind, we are actively planning to cover pre-C++11 material in future editions. For the time being, however, we highly recommend \textit{Effective C++} by Scott Meyers\footnote{\textbf{meyers05}} as a concise, practical treatment of many important and useful C++98/03 features.

\section[The {\sffamily\itshape EMC++S} White Paper]{The {\sfbsectionitalRomeo EMC++S} White Paper}

Throughout the writing of \textit{Embracing Modern C++ Safely}, we have followed a set of guiding principles, which collectively drive the style and content of this book.

\subsection{Facts (Not Opinions)}

This book describes only beneficial uses and potential pitfalls of modern C++ features. The content presented is based on objectively verifiable facts, either derived from standards documents or from extensive practical experience; we explicitly avoid subjective opinion such as our evaluation of the relative merits of design tradeoffs (restraint that admittedly is a good exercise in humility). Although such opinions are often valuable, they are inherently biased toward the author's area of expertise.

Note that \textit{safety} --- the rating we use to segregate features by chapter --- is the one exception to this objectivity guideline. Although the analysis of each feature aims at being entirely objective, its chapter classification --- indicating the relative safety of its quotidian use in a large software-development environment --- reflects our combined accumulated experience totaling decades of real-world, hands-on experience with developing a variety of large-scale C++ software systems.

\subsection{Elucidation (Not Prescription)}

We deliberately avoid prescribing any cut-and-dried solutions to address specific feature pitfalls. Instead, we merely describe and characterize such concerns in sufficient detail to equip you to devise a solution suitable for your own development environment. In some cases, we might reference techniques or publicly available libraries that others have used to work around such speed bumps, but we do not pass judgment about which workaround should be considered a best practice.

\subsection{Brevity (Not Verbosity)}

\textit{Embracing Modern C++ Safely} is neither designed nor intended to be an introduction to modern C++. It is a handy reference for experienced C++ programmers who may have a passing knowledge of the recently added C++ features and a desire to perfect their understanding. Our writing style is intentionally tight, with the goal of providing you with facts, concise objective analysis, and cogent, real-world examples. By doing so we spare you the task of wading through introductory material. If you are entirely unfamiliar with a feature, we suggest you start with a more elementary and language-centric text such as \textit{The C++ Programming Language} by Bjarne Stroustrup.\footnote{\textbf{stroustrup13}}

\subsection{Real-World (Not Contrived) Examples}

We hope you will find the examples in this book useful in multiple ways. The primary purpose of examples is to illustrate productive use of each feature as it might occur in practice. We stay away from contrived examples that give equal importance to seldom-used aspects of the feature, as to the intended, idiomatic uses. Hence, many of our examples are based on simplified code fragments extracted from real-world codebases. Though we typically change identifier names to be more appropriate to the shortened example (rather than the context and the process that led to the example), we keep the code structure of each example as close as possible to its original real-world counterpart.

\subsection{At Scale (Not Overly Simplistic) Programs}

By scale, we attempt to simultaneously capture two distinct aspects of size: (1) the sheer product size (e.g., in bytes, source lines, separate units of release) of the programs, systems, and libraries developed and maintained by a software organization; and (2) the size of an organization itself as measured by the number of software developers, quality assurance engineers, site reliability engineers, operators, and so on that the organization employs. As with many aspects of software development, what works for small programs simply doesn’t scale to larger development efforts.

What’s more, powerful new language features that are handled perfectly well by a few expert programmers working together in the archetypal garage on a prototype for their new start-up don't always fare as well when they are wantonly exercised by numerous members of a large software development organization. Hence, when we consider the relative safety of a feature, as defined in the next section, we do so with mindfulness that any given feature might be used, and occasionally misused, in very large programs and by a very large number of programmers having a wide range of knowledge, skill, and ability.

\section[{What Do We Mean by {\sffamily\itshape Safely}?}]{What Do We Mean by {\sfbsectionitalRomeo Safely}?}

The ISO C++ Standards Committee, of which we are members, would be remiss --- and downright negligent --- if it allowed any feature of the C++ language to be standardized if that feature were not reliably safe when used as intended. Still, we have chosen the word ``safely" as the moniker for the signature aspect of our book, by which we indicate a comparatively favorable risk-to-reward ratio for using a given feature in a large-scale development environment. By contextualizing the meaning of the term ``safe," we get to apply it to a real-world economy in which everything has a cost in multiple dimensions: risk of misuse, added maintenance burden borne by using a new feature in an older code base, and training needs for developers who might not be familiar with that feature.

Several aspects conspire to offset the value added by the adoption and widespread use of any new language feature, thereby reducing its intrinsic safety. By categorizing features in terms of safety, we strive to capture an appropriately weighted combination of the following factors:
\begin{enumerate}
\item{Number and severity of known deficiencies}
\item{Difficulty in teaching consistent proper use}
\item{Experience level required for consistent proper use}
\item{Risks associated with widespread misuse}
\end{enumerate}
Bottom line: In this book, the degree of safety of a given feature is the relative likelihood that widespread use of that feature will have positive impact and no adverse effect on a large software company’s codebase.

\section[A {\sffamily\itshape Safe} Feature]{A {\sfbsectionitalRomeo Safe} Feature}

Some of the new features of modern C++ add considerable value, are easy to use, and are decidedly hard to misuse unintentionally; hence, ubiquitous adoption of such features is productive, relatively unlikely to become a problem in the context of a large-scale development organization, and to be generally encouraged --- even without training. We identify such staunchly helpful, unflappable C++ features as \textit{safe}.

For example, we categorize the \texttt{override} contextual keyword as a safe feature because it prevents bugs, serves as documentation, cannot easily be misused, and has no serious deficiencies. If someone has heard of this feature and tried to use it and the software compiles, the code base is likely better for it. Using \texttt{override} wherever applicable is always a sound engineering decision.

\section[A {\sffamily\itshape Conditionally Safe} Feature]{A {\sfbsectionitalRomeo Conditionally Safe} Feature}

The preponderance of new features available in modern C++ has important, frequently occurring, and valuable uses, yet how these features are used appropriately, let alone optimally, might not be obvious. What’s more, some of these features are fraught with inherent dangers and deficiencies, requiring explicit training and extra care to circumnavigate their pitfalls.

For example, we deem default member initializers a \textit{conditionally safe} feature because, although they are easy to use per se, the perhaps less-than-obvious unintended consequences of doing so (e.g., tight compile-time coupling) might be prohibitively costly in certain circumstances (e.g., might prevent relink-only patching in production).

\section[An {\sffamily\itshape Unsafe} Feature]{An {\sfbsectionitalRomeo Unsafe} Feature}

When an expert programmer uses any C++ feature appropriately, the feature typically does no direct harm. Yet other developers --- seeing the feature’s use in the code base but failing to appreciate the highly specialized or nuanced reasoning justifying it --- might attempt to use it in what they perceive to be a similar way, yet with profoundly less desirable results. Similarly, maintainers may change the use of a fragile feature altering its semantics in subtle but damaging ways.

Features that are classified as unsafe are those that might have valid, and even very important, use cases, yet our experience indicates that routine or widespread use thereof would be counterproductive in a typical large-scale software-development enterprise.

For example, we deem the final contextual keyword an unsafe feature because the situations in which it would be misused overwhelmingly outnumber those vanishingly few isolated cases in which it is appropriate, let alone valuable. Furthermore, its widespread use would inhibit fine-grained (e.g., hierarchical) reuse, which is critically important to the success of a large organization.

\section{Modern C++ Feature Catalog}

As an essential aspect of its design, this first edition of \textit{Embracing Modern C++ Safely} aims to serve as a comprehensive catalog of C++11 and C++14 language features, presenting vital information for each of them in a clear, concise, consistent, and predictable format to which experienced engineers can readily refer during development or technical discourse.

\subsection{Organization}

This book is divided into five chapters, the middle three of which form the catalog characterizing modern C++ language features grouped by their respective safety classifications:
\begin{itemize}
\item{ Chapter 0: Introduction}
\item{Chapter 1: Safe Features}
\item{Chapter 2: Conditionally Safe Features}
\item{Chapter 3: Unsafe Features}
\item{Chapter 4: Parting Thoughts}
\end{itemize}
For this first edition, the language-feature chapters (1, 2, and 3) each consist of two sections containing, respectively, C++11 and C++14 features having the safety level (\textit{safe}, \textit{conditionally safe}, or \textit{unsafe}) corresponding to that chapter. Recall, however, that Standard Library features are outside the scope of this book.

Each feature resides in its own subsection, rendered in a canonical format:
\begin{itemize}
\item{Description}
\item{Use Cases}
\item{Potential Pitfalls}
\item{Annoyances}
\item{See Also}
\item{Further Reading}
\end{itemize}
By constraining our treatment of each individual feature to this canonized format, we avoid gratuitous variations in rendering, thereby facilitating rapid discovery of whatever particular aspects of a given language feature you are searching for.

\section{How To Use This Book}

Depending on your needs, \textit{Embracing Modern C++ Safely} can be handy in a variety of ways.
\begin{enumerate}
\item{\textbf{Read the entire book from front to back.} If you are conversant with classic C++, consuming this book in its entirety all at once will provide a complete and nuanced practical understanding of each of the language features introduced by C++11 and C++14.}
\item{\textbf{Read the chapters in order but slowly over time.} An incremental, priority-driven approach is also possible and recommended, especially if you're feeling less sure-footed. Understanding and applying first the safe features of Chapter 1 gets you the low-hanging fruit. In time, the conditionally safe features of Chapter 2 will allow you to ease into the breadth of useful modern C++ language features, prioritizing those that are least likely to prove problematic.}
\item{\textbf{Read the first sections of each of the three catalog chapters first.} If you are a developer whose organization uses C++11 but not yet C++14, you can focus on learning everything that can be applied now and then circle back and learn the rest later when it becomes relevant to your evolving organization.}
\item{\textbf{Use the book as a quick-reference guide if and as needed.} Random access is great, too, especially now that you've made it through Chapter 0. If you prefer not to read the book in its entirety (or simply want to refer to it periodically as a refresher), reading any arbitrary individual feature subsection in any order will provide timely access to all relevant details of whichever feature is of immediate interest.}
\end{enumerate}

We wish you would derive value in several ways from the knowledge imbued into \textit{Embracing Modern C++ Safely}, irrespective of how you read it. In addition to helping you become a more knowledgeable and therefore safer developer, this book aims to clarify (whether you are a developer, a lead, or a manager) which features demand more training, attention to detail, experience, peer review, and such. The factual, objective presentation style also makes for excellent input into the preparation of coding standards and style guides that suit the particular needs of a company, project, team, or even just a single discriminating developer (which, of course, we all aim at being). Finally, any C++ software development organization that adopts this book will be taking the first steps toward leveraging modern C++ in a way that maximizes reward while minimizing risks, i.e., by embracing modern C++ \textit{safely}. We are very much looking forward to getting feedback and suggestions for future editions of \textit{Embracing Modern C++ Safely} at www.TODOTODOTODO.com. Happy coding!

  % Chapter 0
\cleardoublepage

%%%%%% RHs redefined for Chs 1-3
%\fancyhead[RO]{\sectionmark}
\fancyhead[LO]{\sfnine\cppxx}
%\fancyhead[RE]{\sectionmark}
%\fancyhead[LE]{\chaptermark}
%%%%%%%%%%

\chapter[Safe Features]{Safe Features}\label{ch-safe}
\emcppschapterstart{}

Intro text should be here. 

%%%%%%%%%%%% C++11
\stepcounter{cppxx}
\addtocontents{toc}{\protect\renewcommand*\protect\cftsecindent{1.25em}}
\cftaddnumtitleline{toc}{section}{\thechapter.\thecppxx}{C++11}{\thepage}
% to add to the TOC at the section level
\renewcommand{\cppxx}{C++11}

\newpage
\addtocontents{toc}{\protect\renewcommand*\protect\cftsecindent{4.5em}}
%\section[Attribute Syntax]{Generalized Attribute Support}\label{attributes}
% 6 Feb 2021, copyedits in and proofed
% 8 Feb 2021 JMB, compilation fixes
% 9 Feb 2021, FPPs set
% 17 Feb 2021, glossary commands updated 
% 3 March 2021, proofreading correx in and proofed 

\emcppsFeature{
    short={Attribute Syntax},
    long={Generalized Attribute Support},
}{attributes}
\setcounter{table}{0}
\setcounter{footnote}{0}
\setcounter{lstlisting}{0}
%\section[Attribute Syntax]{Generalized Attribute Support}\label{attributes}

An \emph{attribute} is an annotation (e.g., of a statement or named
\emcppsgloss{entity}) used to provide supplementary information.

\subsection[Description]{Description}\label{description}

Developers are often aware of information that cannot be easily deduced directly
from the source code within a given translation unit. Some of this
information might be useful to certain compilers, say, to inform
diagnostics or optimizations; typical attributes, however, are designed
to avoid affecting the semantics of a
well-written program. By \emph{semantics}, here
we typically mean any observable behavior apart from runtime
performance. Generally, ignoring an attribute is a valid (and safe)
choice for a compiler to make. Sometimes, however, an
attribute will not affect the behavior of a \emph{correct} program but might affect the behavior of a well-formed yet incorrect one (see \intraref{attribute-use-cases}{delineating-explicit-assumptions-in-code-to-achieve-better-optimizations}).
%  {\it\titleref{attribute-use-cases}:} {\it\titleref{delineating-explicit-assumptions-in-code-to-achieve-better-optimizations}} on page~\pageref{delineating-explicit-assumptions-in-code-to-achieve-better-optimizations}).
  Customized annotations targeted at external tools
%{\cprotect\footnote{Such tools include static analyzers such as Clang sanitizers.}}
might be beneficial as well.

\subsubsection[C++ attribute syntax]{C++ attribute syntax}\label{c++-attribute-syntax}

C++ supports a standard syntax for attributes, introduced via a matching
pair of \lstinline![[! and \lstinline!]]!, the simplest of which is a single
attribute represented using a simple identifier, e.g.,
\lstinline!attribute_name!:

\begin{emcppslisting}[emcppsbatch=e1]
[[attribute_name]]
\end{emcppslisting}

\noindent A single annotation can consist of zero or more attributes:

\begin{emcppslisting}[emcppsbatch=e2]
[[]]           // permitted in every position where any attribute is allowed
[[foo, bar]]   // equivalent to (ù{\codeincomments{[[foo]] [[bar]]}}ù)
\end{emcppslisting}

\noindent An attribute may have an argument list consisting of an
arbitrary sequence of tokens:

\begin{emcppslisting}[emcppsbatch=e3]
[[attribute_name()]]            // same as (ù{\codeincomments{attribute\_name}}ù)
[[deprecated("bad API")]]       // single-argument attribute
[[theoretical(1, "two", 3.0)]]  // multiple-argument attribute
[[complicated({1, 2, 3} + 5)]]  // arbitrary tokens(ù{\footnotemark}ù)
\end{emcppslisting}
\footnotetext{GCC offered no support for certain tokens in the attributes until GCC~v9.3 (c. 2020).}

\noindent Note that having an incorrect number of arguments or an incompatible
argument type is a compile-time error for all  attributes defined by the Standard; %\pagebreak%%%%%%%
the behavior for all other attributes,
 however, is
\emcppsgloss[implementation defined]{implementation-defined} (see \intraref{attribute-potential-pitfalls}{unrecognized-attributes-have-implementation-defined-behavior}).
%{\it\titleref{attribute-potential-pitfalls}:} {\it\titleref{unrecognized-attributes-have-implementation-defined-behavior}} on page~\pageref{unrecognized-attributes-have-implementation-defined-behavior}).

Any attribute may be qualified with an attribute namespace{\cprotect\footnote{Attributes
having a namespace-qualified name (e.g., \lstinline![[gnu::const]]!) were
only \emcppsgloss{conditionally supported} in C++11 and C++14, but
historically they were supported by all major compilers, including both
Clang and GCC; all C++17-conforming compilers \emph{must} support attribute namespaces.}} (a single arbitrary\linebreak[4] \mbox{identifier}):

\begin{emcppslisting}[emcppsbatch=e4]
[[gnu::const]]  // (GCC-specific) namespace-(ù{\codeincomments{gnu}}ù)-qualified (ù{\codeincomments{const}}ù) attribute
[[my::own]]     // (user-specified) namespace-(ù{\codeincomments{my}}ù)-qualified (ù{\codeincomments{own}}ù) attribute
\end{emcppslisting}

\begin{emcppshiddenlisting}[emcppsbatch={e1,e2,e3,e4}]
// the above listings with just atttributes need to be on something, like a function
void foo();
\end{emcppshiddenlisting}


\subsubsection[C++ attribute placement]{C++ attribute placement}\label{c++-attribute-placement}

%%%%% old section
%Attributes can, in principle, be introduced almost anywhere within the
%C++ syntax to annotate almost anything, including an entity,
%statement, code block, and even entire translation
%unit; however, most contemporary compilers do not support arbitrary
%placement of attributes (see {\it\titleref{attribute-use-cases}:} {\it\titleref{probing-where-attributes-are-permitted-in-the-compiler’s-c++-grammar}} on page~\pageref{probing-where-attributes-are-permitted-in-the-compiler’s-c++-grammar}) outside of a
%declaration statement. Furthermore, in some cases, the entity to
%which an unrecognized attribute pertains might not be clear from its
%syntactic placement alone.
%
%In the case of a declaration statement, however, the intended entity is
%well specified; an attribute placed in front of the statement applies to
%every entity being declared, whereas an attribute placed immediately
%after the named entity applies to just that one entity:
%
%\begin{emcppslisting}
%[[noreturn]] void f(), g();  // Both (ù{\codeincomments{f()}}ù) and (ù{\codeincomments{g()}}ù) are (ù{\codeincomments{noreturn}}ù).
%void u(), v() [[noreturn]];  // Only (ù{\codeincomments{v()}}ù) is (ù{\codeincomments{noreturn}}ù).
%\end{emcppslisting}
%
%\noindent Attributes placed in front of a declaration statement and immediately
%behind the name{\cprotect\footnote{There are rare edge cases in which an
%entity (e.g., an anonymous \texttt{union} or \texttt{enum}) is
%declared without a name:
%
%%\begin{emcppslisting}[basicstyle={\ttfamily\footnotesize}]
%\begin{lstlisting}[style=footcode]
%struct S { union [[attribute_name]] { int a; float b }; };
%enum [[attribute_name]] { SUCCESS, FAIL } result;
%\end{emcppslisting} \vspace*{-1ex}
%      }} of an individual entity in the same statement are additive (for
%that entity). The behavior of attributes associated with an entity
%across multiple declaration statements, however, depends on the
%attributes themselves. As an example, \texttt{[[noreturn]]} is required
%to be present on the \emph{first} declaration of a function. Other
%attributes might be additive, such as the hypothetical \texttt{foo} and
%\texttt{bar} shown here:
%
%\begin{emcppslisting}
%[[foo]] void f(), g();  // declares both (ù{\codeincomments{f()}}ù) and (ù{\codeincomments{g()}}ù) to be (ù{\codeincomments{foo}}ù)
%void f [[bar]](), g();  // Now (ù{\codeincomments{f()}}ù) is both (ù{\codeincomments{foo}}ù) and (ù{\codeincomments{bar}}ù) while
%                        //     (ù{\codeincomments{g()}}ù) is still just (ù{\codeincomments{foo}}ù).
%\end{emcppslisting}
%\pagebreak%%%%%%%

%\noindent Redundant attributes are not themselves necessarily considered an error;
%however, most standard attributes do consider redundancy an
%error{\cprotect\footnote{Redundancy of standard
%attributes might no longer be an error in future revisions of the
%  C++ Standard; see \textbf{{iso20a}}.}}:
%
%\begin{emcppslisting}
%[[attr1]] void f [[attr2]](), f [[attr3]](int);
%                                      // (ù{\codeincomments{f()}}ù)    is (ù{\codeincomments{attr1}}ù) and (ù{\codeincomments{attr2}}ù).
%                                      // (ù{\codeincomments{f(int)}}ù) is (ù{\codeincomments{attr1}}ù) and (ù{\codeincomments{attr3}}ù).
%
%[[a1]][[a1]] int [[a1]][[a1]] & x;    // (ù{\codeincomments{x}}ù) (the reference itself) is (ù{\codeincomments{a1}}ù).
%
%void g [[noreturn]] [[noreturn]]();   // (ù{\codeincomments{g()}}ù) is (ù{\codeincomments{noreturn}}ù).
%
%void h [[noreturn, noreturn]]();      // error: repeated (standard) attribute
%\end{emcppslisting}
%
%\noindent In most other cases, an attribute will typically apply to the statement
%(including a block statement) that immediately (apart from other
%attributes) follows it:
%
%\begin{emcppslisting}
%[[attr1]];                                // null statement
%[[attr2]] return 0;                       // return statement
%[[attr3]] for (int i = 0; i < 10; ++i);   // for statement
%[[attr4]] [[attr5]] { /* ... */ }         // block statement
%\end{emcppslisting}
%
%\noindent The valid positions of any particular attribute, however, will be
%constrained by whatever entities to which it applies. That is, an
%attribute such as \texttt{noreturn}, which pertains only to functions,
%would be valid syntactically but not semantically were it placed so as
%to annotate any other kind of entity or syntactic element. Misplacement of
%standard attributes results in an ill-formed
%program{\cprotect\footnote{As of this writing, GCC is lax and merely
%warns when it sees the standard \texttt{noreturn} attribute in an
%unauthorized syntactic position, whereas Clang (correctly) fails to
%compile. Hence creative use of even a standard attribute might
%  lead to different behavior on different compilers.}}:
%
%\begin{emcppslisting}
%void [[noreturn]] g() { throw; }  // error: appertains to type specifier
%void i() [[noreturn]] { throw; }  // error: appertains to type specifier
%\end{emcppslisting}


%%%%%% new section
Attributes can be placed in a variety of locations within the C++ grammar.  For each such location, the Standard defines the entity or statement to which the attribute is said to \emph{appertain}. For example, an attribute in front of a simple declaration statement appertains to each of the entities declared by the statement, whereas an attribute placed immediately after the declared name appertains only to that entity:

\begin{emcppslisting}[language=C]
[[foo]] void f(), g();   // foo appertains to both f() and g().
void u(), v [[foo]] ();  // foo appertains only to v().
\end{emcppslisting}

\noindent Attributes can apply to an entity without a name (e.g., anonymous \lstinline!union! or \lstinline!enum!):

\begin{emcppslisting}
struct S { union [[attribute_name]] { int a; float b; }; };
enum [[attribute_name]] { SUCCESS, FAIL } result;
\end{emcppslisting}

\noindent The valid positions for any particular attribute are constrained to only those locations
where the attribute appertains to the entity to which it applies. That is, an attribute such as \lstinline!noreturn!,
which applies only to functions, would be valid syntactically but not semantically valid if it were used to annotate
any other kind of entity or syntactic element.  Misplacement of a standard attribute results
in an ill-formed program\footnote{As of this writing, GCC is lax and merely
warns when it sees the standard \lstinline!noreturn! attribute in an
unauthorized syntactic position, whereas Clang (correctly) fails to
compile. Hence, using even a standard attribute might
  lead to a different behavior on different compilers.}:

\begin{emcppslisting}[emcppserrorlines={4}]
void [[noreturn]] x() {       // Error, cannot be applied to a type
     [[noreturn]] int i;      // Error, cannot be applied to a variable
     [[noreturn]] { throw; }  // Error, cannot be applied to a statement
}
\end{emcppslisting}

\noindent The empty attribute specifier sequence \lstinline![[]]! is allowed to appear anywhere the C++ grammar allows attributes.

%%%%%%%%%%%%%%%%%%%%%%


\subsubsection[Common compiler-dependent attributes]{Common compiler-dependent attributes}\label{common-compiler-dependent-attributes}

Prior to C++11, no standardized syntax for attributes was available and nonportable compiler intrinsics
(such as \mbox{\lstinline!__attribute__((fallthrough))!}, which is
GCC-specific syntax) had to be used instead. Given the new standard
syntax, vendors are now able to express these extensions in a syntactically consistent manner. If an unknown attribute is
encountered during compilation, it is ignored, emitting a likely{\cprotect\footnote{Prior to C++17, a conforming implementation was
permitted to treat an unknown attribute as ill-formed and terminate
  translation; to the authors' knowledge, however, none of them did.}} nonfatal
diagnostic.

Table \ref{attribute-table1} provides several examples of popular compiler-specific attributes
that have been standardized or have migrated to the standard syntax. (For
additional compiler-specific attributes, see \intrarefsimple{attribute-further-reading}).
%{\it\titleref{attribute-further-reading}} on page~\pageref{attribute-further-reading}.)
%\pagebreak%%%%%%%%

\begin{table}[h!]
\begin{center}
%\vspace{-3ex}
\begin{threeparttable}
\caption{Some standardized compiler-specific attributes}\label{attribute-table1}\vspace{1.5ex}
{\small \begin{tabular}{c|c|c}\thickhline
\rowcolor[gray]{.9}   {\sffamily\bfseries Compiler} & {\sffamily\bfseries Compiler-Specific} &
{\sffamily\bfseries Standard-Conforming} \\ \hline
GCC &\lstinline!__attribute__((pure))! & \lstinline![[gnu::pure]]! \\ \hline
Clang & \lstinline!__attribute__((no_sanitize))! &\lstinline![[clang::no_sanitize]]! \\ \hline
MSVC & \lstinline!declspec(deprecated)! & \lstinline![[deprecated]]! \\ \thickhline
\end{tabular}
}
\end{threeparttable}
\end{center}
\end{table}

Portability is the biggest advantage of preferring standard syntax when it is available for compiler- and external-tool-specific attributes. Because most compilers will simply ignore unknown attributes that use standard attribute syntax (and, as of C++17, they are required to do so), conditional compilation is no longer required.

\subsection[Use Cases]{Use Cases}\label{attribute-use-cases}

\subsubsection[Prompting useful compiler diagnostics]{Prompting useful compiler diagnostics}\label{eliciting-useful-compiler-diagnostics}

Decorating entities with certain attributes can give compilers enough
additional con-\linebreak %%%%%%%% to reduce the wacky linespacing caused by long code term next line
text to provide more detailed diagnostics. For example,
the GCC-specific\linebreak[4] \lstinline![[gnu::warn_unused_result]]!
attribute{\cprotect\footnote{For compatibility with GCC,
Clang supports \lstinline![[gnu::warn_unused_result]]! as
  well.}} can be used to inform the compiler (and developers) that a
function's return value should not be ignored{\cprotect\footnote{The
C++17 Standard \lstinline![[nodiscard]]! attribute serves the same
  purpose and is portable.}}:

\begin{emcppslisting}[emcppsbatch=e5]
struct UDPListener
{
    [[gnu::warn_unused_result]] int start();
        // Start the UDP listener's background thread (which can fail for a
        // variety of reasons). Return 0 on success and a nonzero value
        // otherwise.

    void bind(int port);
        // The behavior is undefined unless (ù{\codeincomments{start}}ù) was called successfully.
};
\end{emcppslisting}
\pagebreak%%%%%%

\noindent Such annotation of the client-facing declaration can prevent defects
caused by a client forgetting to inspect the result of a
function{\cprotect\footnote{Because the
\lstinline![[gnu::warn_unused_result]]! attribute does not affect code
generation, it is explicitly \emph{not} ill formed for a client to
make use of an unannotated declaration and yet compile its
corresponding definition in the context of an annotated one (or vice
versa); such is not always the case for other attributes, however, and
  best practice might argue in favor of consistency regardless.}}:

\begin{emcppslisting}[emcppsbatch=e5]
void init()
{
    UDPListener listener;
    listener.start();      // Might fail; return value must be checked!
    listener.bind(27015);  // Possible undefined behavior (BAD IDEA)
}
\end{emcppslisting}

\noindent For the code above, GCC produces a useful warning:

\begin{lstlisting}[style=plain]
warning: ignoring return value of 'int UDPListener::start()' declared
         with attribute 'warn_unused_result' [-Wunused-result]
\end{lstlisting}


\subsubsection[Hinting at additional optimization opportunities]{Hinting at additional optimization opportunities}\label{hinting-at-additional-optimization-opportunities}

Some annotations can affect compiler optimizations leading to more
efficient or smaller binaries. For example, decorating the function
\lstinline!reportError! below with the GCC-specific
\lstinline![[gnu::cold]]! attribute (also available on Clang) tells the
compiler that the developer believes the function is unlikely to be
called often:

\begin{emcppslisting}[emcppsbatch=e6]
[[gnu::cold]] void reportError(const char* message) { /* ... */ }
\end{emcppslisting}

\noindent Not only might the definition of \lstinline!reportError! itself be
optimized differently (e.g., for space over speed), any use of this
function will likely be given lower priority during branch \mbox{prediction}:

\begin{emcppslisting}[emcppsbatch=e6]
void checkBalance(int balance)
{
    if (balance >= 0)  // likely branch
    {
        // ...
    }
    else  // unlikely branch
    {
        reportError("Negative balance.");
    }
}
\end{emcppslisting}
% \pagebreak%%%%

\noindent Because the (annotated) \lstinline!reportError(const!~\lstinline!char*)!
appears on the \lstinline!else! branch of the \lstinline!if! statement (above), the compiler
knows to expect that \lstinline!balance! is likely \pagebreak%%%%%%%  
\emph{not} to be negative  
and therefore optimizes its predictive branching accordingly.
Note that even if our hint to the compiler turns out to be misleading at
run time, the semantics of every well-formed program remain the same.

\subsubsection[Stating explicit assumptions in code to achieve better optimizations]{Stating explicit assumptions in code to achieve better optimizations}\label{delineating-explicit-assumptions-in-code-to-achieve-better-optimizations}

Although the presence of an attribute usually has no effect
on the behavior of any well-formed program besides its runtime performance,
an attribute sometimes imparts knowledge to the compiler,
which, if incorrect, could alter the intended behavior of the program. As an
example of this more forceful form of attribute, consider the
GCC-specific \lstinline![[gnu::const]]! attribute (also available in
Clang). When applied to a function, this attribute instructs the compiler to \emph{assume}
that the function is a \emcppsgloss{pure function}, which has no
\emcppsgloss{side effects}. In other words, the function always returns the same value for a given set of arguments, and the globally reachable state of the program is not altered by the function. For example, a function performing a linear interpolation between two values may be annotated with \lstinline![[gnu::const]]!:

\begin{emcppslisting}[emcppsbatch={e7,e8}]
[[gnu::const]]
double linearInterpolation(double start, double end, double factor)
{
    return (start * (1.0 - factor)) + (end * factor);
}
\end{emcppslisting}

\noindent More generally, the return value of a function annotated with \lstinline![[gnu::const]]! is not permitted to depend on any state that might change between its successive invocations. For example, it is not allowed to examine contents of memory supplied to it by address. In contrast, functions annotated with a similar but more lenient \lstinline![[gnu::pure]]! attribute are allowed to return values that depend on any nonvolatile state.  Therefore, functions such as \lstinline!strlen! or \lstinline!memcmp!, which read but do not modify the observable state, may be annotated with \lstinline![[gnu::pure]]! but not \lstinline![[gnu::const]]!.

The \lstinline!vectorLerp! function below performs linear
interpolation (referred to as LERP) between two bidimen\-sional vectors. The body
of this function comprises two invocations to the
\mbox{\lstinline!linearInterpolation!} function (above) --- one per vector
component:

\begin{emcppshiddenlisting}[emcppsbatch={e7,e8}]
struct Vector2D {
    double x, y;
    Vector2D(double, double);
};
\end{emcppshiddenlisting}
\begin{emcppslisting}[emcppsbatch=e7]
Vector2D vectorLerp(const Vector2D& start, const Vector2D& end, double factor)
{
    return Vector2D(linearInterpolation(start.x, end.x, factor),
                    linearInterpolation(start.y, end.y, factor));
}
\end{emcppslisting}
%\pagebreak%%%%%%

\noindent In the case where the values of the two components
are the same, the compiler is allowed to invoke
\lstinline!linearInterpolation! only once --- even if its body is not
visible in \lstinline!vectorLerp!'s translation unit:\newpage%%%%%%

\begin{emcppslisting}[emcppsbatch=e8]
// pseudocode (hypothetical compiler transformation)
Vector2D vectorLerp(const Vector2D& start, const Vector2D& end, double factor)
{
    if (start.x == start.y && end.x == end.y)
    {
        const double cache = linearInterpolation(start.x, end.x, factor);
        return Vector2D(cache, cache);
    }

    return Vector2D(linearInterpolation(start.x, end.x, factor),
                    linearInterpolation(start.y, end.y, factor));
}
\end{emcppslisting}

\noindent If the implementation of a function tagged with the
\lstinline![[gnu::const]]! attribute does not satisfy limitations imposed by
the attribute, however, the compiler will not be able to detect this, and a runtime
defect will be the likely result;
%{\cprotect\footnote{The briefly adopted
%--- and then \emph{unadopted} --- contract-checking facility proposed
%for C++20 contemplated incorporating a feature similar in spirit to
%\texttt{[[gnu::const]]} in which preconditions (in addition to being
%runtime checked or ignored) could be \emph{assumed} to be true by the
%compiler for the purposes of optimization; this unique use of
%attribute-like syntax also required that a conforming implementation
%could not unilaterally ignore these precondition-checking attributes
%since that would make attempting to test them result in hard
%  (\emph{language}) \textbf{undefined behavior}.}};
  see \intraref{attribute-potential-pitfalls}{some-attributes,-if-misused,-can-affect-program-correctness}.
%{\it\titleref{attribute-potential-pitfalls}:} {\it\titleref{some-attributes,-if-misused,-can-affect-program-correctness}} on page~\pageref{some-attributes,-if-misused,-can-affect-program-correctness}.

\subsubsection[Using attributes to control external static analysis]{Using attributes to control external static analysis}\label{using-attributes-to-control-external-static-analysis}

Since unknown attributes are ignored by the compiler, external
static-analysis tools can define their own custom attributes that can be
used to embed detailed information to influence or control those tools
without affecting program semantics. For example, the Microsoft-specific
\lstinline![[gsl::suppress(/*!~\lstinline!rules!~\lstinline!*/)]]! attribute can
be used to suppress unwanted warnings from static-analysis tools that
verify \emph{Guidelines Support Library}{\cprotect\footnote{\emph{Guidelines
Support Library} (see \cite{microsoft}) is an open-source library, developed by Microsoft,
that implements functions and types suggested for use by the ``C++
  Core Guidelines'' (see \cite{stroustrup20}).}}
rules. In particular, consider GSL C26481 (Bounds rule \#1),\footnote{\cite{microsoftC26481}} which forbids any pointer arithmetic, instead
suggesting that users rely on the \lstinline!gsl::span!
type{\cprotect\footnote{\lstinline!gsl::span! is a lightweight reference
type that observes a contiguous sequence (or subsequence) of objects
of homogeneous type. \lstinline!gsl::span! can be used in interfaces as
an alternative to both pointer/size or iterator-pair arguments and in
implementations as an alternative to (raw) pointer arithmetic. Since
  C++20, the standard \lstinline!std::span! template can be used instead.}}:

\begin{emcppshiddenlisting}[emcppsbatch={e9,e10}]
void printElements(int*,int*);
\end{emcppshiddenlisting}
\begin{emcppslisting}[emcppsbatch=e9]
void hereticalFunction()
{
    int array[] = {0, 1, 2, 3, 4, 5};

    printElements(array, array + 6);  // elicits warning C26481
}
\end{emcppslisting}
%\pagebreak%%%%%%%%%

\noindent Any block of code for which validating rule C26481 is considered
undesirable can be decorated with the
\lstinline![[gsl::suppress(bounds.1)]]! attribute:\pagebreak%%%%%%%%%

% gsl::suppress is ignored by gcc, requires a string on clang, and an identifier
% on MSVC.  sigh.
\begin{emcppslisting}[emcppsbatch=e10,emcppserrorlines=5]
void hereticalFunction()
{
    int array[] = {0, 1, 2, 3, 4, 5};

    [[gsl::suppress(bounds.1)]]           // Suppress GSL C26481.
    {
        printElements(array, array + 6);  // Silence!
    }
}
\end{emcppslisting}


\subsubsection[Creating new attributes to express semantic properties]{Creating new attributes to express semantic properties}\label{creating-new-attributes-to-express-semantic-properties}

Other uses of attributes for static analysis include statements of
properties that cannot otherwise be deduced. Consider a function, \lstinline!f!, that takes two pointers,
\lstinline!p1! and \lstinline!p2!, and has a \emcppsgloss{precondition} that both
pointers must refer to the same contiguous block of memory. Using the standard attribute to inform the analyzer of such a precondition has a distinct advantage of requiring nothing other than the agreement between the developer and the static analyzer regarding the namespace and the name of the attribute. For example, we could choose to designate \lstinline!home_grown::in_same_block(p1, p2)! for this purpose:

\begin{emcppslisting}[emcppsbatch=e11]
// lib.h:

[[home_grown::in_same_block(p1, p2)]]
int f(double* p1, double* p2);
\end{emcppslisting}
%\pagebreak%%%%%%%%%

\noindent The compiler will simply ignore this unknown attribute. However, because our static-analysis tool knows the meaning of the \lstinline!home_grown::in_same_block! attribute, it will report, at analysis time, defects that might otherwise have resulted in \emcppsgloss{undefined behavior} at run time:

\begin{emcppslisting}[emcppsbatch=e11]
// client.cpp:
#include <lib.h>

void client()
{
    double a[10], b[10];
    f(a, b);  // Unrelated pointers --- Our static analyzer reports an error.
}
\end{emcppslisting}

%\noindent Because our static-analysis tool knows from the
%\texttt{home\_grown::in\_same\_block} attribute that \texttt{a} and
%\texttt{b} must point into the same contiguous block, however, it has
%enough information to report, at compile time, what might otherwise have
%resulted in \textbf{undefined behavior} at run~time.

%\subsubsection[Probing where attributes are permitted in the compiler’s C++ grammar]{Probing where attributes are permitted in the compiler’s C++ grammar}\label{probing-where-attributes-are-permitted-in-the-compiler’s-c++-grammar}
%
%An attribute can generally appear syntactically at the beginning of any
%statement --- e.g.,
%\texttt{[[attr]]}~\texttt{x}~\texttt{=}~\texttt{5;} --- or in almost any
%position relative to a type or expression (e.g.,
%\texttt{const}~\texttt{int\&}) but typically cannot be associated
%within named objects outside of a declaration statement:
%
%\begin{emcppslisting}
%[[]] static [[]] int [[]] a [[]], /*[[]]*/ b [[]];  // declaration statement
%\end{emcppslisting}
%
%\noindent Notice how we have used the empty attribute syntax \texttt{[[]]} above
%to probe for positions allowed for arbitrary attributes by the compiler
%(in this case, GCC) --- the only invalid one being immediately following
%the comma, shown above as \texttt{/*[[]]*/}. Outside of a declaration
%statement, however, viable attribute locations are typically far more
%limited:
%
%\begin{emcppslisting}[label=attribute-gcc-example]
%[[]] void [[]] f [[]] ( [[]] int [[]] n [[]] )
%[[]] {
%    [[]] n /**/ *= /**/ sizeof /**/ ( [[]] const [[]] int [[]] & [[]] ) /**/;
%    [[]] for ([[]] int [[]] i [[]] = /**/ 0 /**/ ;
%                       /**/ i  /**/ < /**/ n /**/ ;
%               /**/ ++ /**/ i /**/ )
%    [[]] {
%        [[]] ;             // (ù{\codeincomments{[[]]}}ù) denotes viable attribute location (on GCC)
%    /**/ }
%/**/ }                     // (ù{\codeincomments{/**/}}ù) denotes no attribute allowed (on GCC)
%\end{emcppslisting}
%
%\noindent Type expressions --- e.g., the argument to \texttt{sizeof} (above) ---
%are a notable exception; see {\it\titleref{attribute-potential-pitfalls}:} {\it\titleref{not-every-syntactic-location-is-viable-for-an-attribute}} on page~\pageref{not-every-syntactic-location-is-viable-for-an-attribute}.
%%\newpage%%%%%%%%%%%%

\subsection[Potential Pitfalls]{Potential Pitfalls}\label{attribute-potential-pitfalls}

\subsubsection[Unrecognized attributes have implementation-defined behavior]{Unrecognized attributes have implementation-defined behavior}\label{unrecognized-attributes-have-implementation-defined-behavior}

Although standard attributes work well and are portable across all
platforms, the behavior of compiler-specific and user-specified
attributes is entirely implementation defined, with unrecognized
attributes typically resulting in compiler warnings. Such warnings can typically be disabled (e.g., on GCC using
\lstinline!-Wno-attributes!), but, if they are, misspellings in even standard
attributes will go unreported.{\cprotect\footnote{Ideally, every relevant platform would offer a way to silently ignore a specific
  attribute on a case-by-case basis.}}

\subsubsection[Some attributes, if misused, can affect program correctness]{Some attributes, if misused, can affect program correctness}\label{some-attributes,-if-misused,-can-affect-program-correctness}

Many attributes are benign in that they might improve diagnostics or
performance but cannot themselves cause a program to behave incorrectly.
However, misuse of some attributes can lead to incorrect
results and/or \emcppsgloss{undefined behavior}.

For example, consider the \lstinline!myRandom! function that is intended to
return a new random number between $[0.0$ and $0.1]$ on each successive
call:
\begin{emcppshiddenlisting}[emcppsbatch=e12]
#include <random>  // (ù{\codeincomments{std::random\_device}}ù), (ù{\codeincomments{std::mt19937}}ù), (ù{\codeincomments{std::uniform\_real\_distribution}}ù)
\end{emcppshiddenlisting}
\begin{emcppslisting}[emcppsbatch=e12]
double myRandom()
{
    static std::random_device randomDevice;
    static std::mt19937 generator(randomDevice());

    std::uniform_real_distribution<double> distribution(0, 1);
    return distribution(generator);
}
\end{emcppslisting}

\noindent Suppose that we somehow observed that decorating \lstinline!myRandom! with
the \lstinline![[gnu::const]]! attribute occasionally improved runtime
performance and innocently but naively decided to use it in production. This
is clearly a misuse of the \lstinline![[gnu::const]]! attribute because the
function doesn't inherently satisfy the requirement of producing the
same result when invoked with the same arguments (in this case, none).
Adding this attribute tells the compiler that it need not call this
function repeatedly and is free to treat the first value returned as a
constant for all time.

%\subsubsection[Not every syntactic location is viable for an attribute]{Not every syntactic location is viable for an attribute}\label{not-every-syntactic-location-is-viable-for-an-attribute}
%
%For a fairly limited subset of syntactic locations, most
%conforming implementations are likely to tolerate the double-bracketed
%attribute-list syntax. The ubiquitously available locations include the
%beginning of any statement, immediately following a named entity in a
%declaration statement, and (typically) arbitrary positions relative to a
%\textbf{type expression} but, beyond that, caveat emptor. For example, GCC
%allowed all of the positions indicated in the example shown in
%{\it\titleref{attribute-use-cases}:} {\it\titleref{probing-where-attributes-are-permitted-in-the-compiler’s-c++-grammar}} on page~\pageref{attribute-gcc-example}, yet Clang had
%issues with the third line in two places:
%
%\begin{emcppslisting}
%<source>:3:39: error: expected variable name or 'this' in lambda capture list
%    [[]] n /**/ *= /**/ sizeof /**/ ([[]] const [[]] int [[]] & [[]] ) /**/;
%                                      ^
%
%<source>:3:48: error: an attribute list cannot appear here
%    [[]] n /**/ *= /**/ sizeof /**/ (/**/ const [[]] int [[]] & [[]] ) /**/;
%                                                ^~~~
%\end{emcppslisting}
%
%\noindent Hence, just because an arbitrary syntactic location is valid for an
%attribute on one compiler doesn't mean that it is necessarily valid on
%another.a

\subsection[Annoyances]{Annoyances}\label{annoyances}

\hspace{\fill}

\subsection[See Also]{See Also}\label{see-also}

\begin{itemize}
\item{\seealsoref{the-noreturn-attribute}{\seealsolocationa}presents a standard attribute used to indicate that a particular function never returns control flow to its caller.}
\item{\seealsoref{deprecated}{\seealsolocationb}presents a standard attribute that discourages the use of an entity via compiler diagnostics.}
\item{\seealsoref{carriesdependency}{\seealsolocatione}presents a standard attribute used to communicate release-consume dependency-chain information to the compiler to avoid unnecessary memory-fence instructions.}
%\item{\locationa, ``\titleref{alignas}" on page~\pageref{alignas} --- Safe C++11 attribute (with a keyword-like syntax) used to widen the alignment of a type or an object}
\end{itemize}

\subsection[Further Reading]{Further Reading}\label{attribute-further-reading}
\begin{itemize}
\item{For more information on commonly supported function attributes, see
section~6.33.1, ``Common Function Attributes," \cite{freesoftwarefdn20}.}
\end{itemize}



%\newpage
%

Lorem ipsum dolor sit amet, consectetur
adipiscing elit. Suspendisse non nisi ex. Ut non sollicitudin est. Donec
laoreet, risus pretium egestas condimentum, ex eros tempor diam, quis
congue nisi metus in tellus. Proin tempor ac lectus nec elementum.
Maecenas augue turpis, pellentesque sed eros sit amet, tincidunt pretium
nulla. Nunc lacus ligula, ullamcorper a porta eget, fringilla sed orci.
Maecenas eget ultricies risus. Donec varius vehicula diam.

\subsection[Description]{Description}\label{description}

Lorem ipsum dolor sit amet, consectetur
adipiscing elit. Suspendisse non nisi ex. Ut non sollicitudin est. Donec
laoreet, risus pretium egestas condimentum, ex eros tempor diam, quis
congue nisi metus in tellus. Proin tempor ac lectus nec elementum.
Maecenas augue turpis, pellentesque sed eros sit amet, tincidunt pretium
nulla. Nunc lacus ligula, ullamcorper a porta eget, fringilla sed orci.
Maecenas eget ultricies risus. Donec varius vehicula diam.

\subsubsection[C++ attribute syntax]{C++ attribute syntax}\label{c++-attribute-syntax}

Lorem ipsum dolor sit amet, consectetur
adipiscing elit.{\cprotect\footnote{\authorsnote We will have
some footnotes that are authors' notes.}} Suspendisse non nisi ex. Ut non sollicitudin est. Donec
laoreet, risus pretium egestas condimentum, ex eros tempor diam, quis
congue nisi metus in tellus. Proin tempor ac lectus nec elementum.
Maecenas augue turpis, pellentesque sed eros sit amet, tincidunt pretium
nulla. Nunc lacus ligula, ullamcorper a porta eget, fringilla sed orci.
Maecenas eget ultricies risus. Donec varius vehicula diam.

Nested unordered list:

\begin{itemize}
\item{One\begin{itemize}
\item{Sub}
\item{Sub}
\end{itemize}
}
\item{Two}
\item{Three}
\end{itemize}

Donec
laoreet, risus pretium egestas condimentum, ex eros tempor diam, quis
congue nisi metus in tellus. Proin tempor ac lectus nec elementum.
Maecenas augue turpis, pellentesque sed eros sit amet, tincidunt pretium
nulla:

Nested ordered list:

\begin{enumerate}\setlength{\leftmarginii}{6mm}
\item{One\begin{enumerate}
\item{Sub}
\item{Sub}
\end{enumerate}
}
\item{Two}
\item{Three}
\end{enumerate}

This feature, when used in conjunction with \emph{explicit instantiation
definitions}, can significantly improve compilation times for a set of
translation units that often instantiate common templates:

\noindent\begin{minipage}{.45\textwidth}
\begin{lstlisting}[language=C++, caption=code 1,frame=tb]{Name}
void code()
{

}
\end{lstlisting}
\end{minipage}\hfill
\begin{minipage}{.45\textwidth}
\begin{lstlisting}[language=C++, caption=code 2,frame=tb]{Name}
void code()
{

}
\end{lstlisting}
\end{minipage}

Nullam nibh tortor, finibus ut lectus eu, convallis vehicula libero. Cras maximus ligula nisl, a eleifend mauris venenatis nec. Suspendisse potenti. Maecenas viverra laoreet mauris ut laoreet. Donec non felis risus. Ut faucibus, justo id luctus lobortis, nisi lacus pharetra est, id pretium arcu ligula ac magna. Morbi a nisl sit amet lacus vehicula tempor nec sit amet nibh. 

\noindent\begin{minipage}{.45\textwidth}
\begin{lstlisting}[language=C++, caption=code 1 with a long and wrapping title,frame=tb]{Name}
void code()
{

}
\end{lstlisting}
\end{minipage}\hfill
\begin{minipage}{.45\textwidth}
\begin{lstlisting}[language=C++, caption=code 2 with a long and wrapping title,frame=tb]{Name}
void code()
{

}
\end{lstlisting}
\end{minipage}

Pellentesque id lorem ac sem ullamcorper semper. Donec imperdiet sapien in nisi faucibus, et tempus lacus hendrerit. Nulla laoreet risus eu tortor ultrices, sed bibendum neque sodales. Ut faucibus ipsum id convallis convallis. Vivamus vehicula rutrum metus in semper. Morbi fringilla ex vel vulputate vehicula. Maecenas ut porttitor massa.

\subsection[Appendix to the Feature]{Appendix to the Feature}
Suspendisse lorem libero, egestas non semper eleifend, finibus ut purus. Quisque leo nulla, lacinia vel dui et, vestibulum facilisis erat. Vestibulum scelerisque auctor diam, ut fringilla diam elementum non. Vivamus sed mauris lobortis, blandit urna quis, iaculis neque. Phasellus sit amet venenatis sapien, vel vestibulum augue. In pellentesque sit amet enim a rutrum. Fusce nec nulla et nisl faucibus efficitur vel non diam. Nulla lobortis feugiat augue, a lobortis nunc imperdiet eget. Sed in neque ultricies, efficitur mi eu, sollicitudin massa. Vivamus sed dui convallis, mattis sem a, scelerisque augue. Sed non ultricies neque. Praesent tincidunt feugiat lorem, a sagittis dui eleifend ac. Nullam molestie nulla quis risus molestie dignissim. Aenean egestas enim ut tellus vulputate finibus. Integer iaculis sodales gravida.

\begin{table}[h!] 
\begin{center}
\begin{threeparttable}
\caption{Testing Table Numbering}\label{test-table1}\vspace{1.5ex} 
{\small \begin{tabular}{c|c|c}\thickhline 
\rowcolor[gray]{.9}   {\sffamily\bfseries Compiler} & {\sffamily\bfseries Compiler-Specific} &
{\sffamily\bfseries Standard-Conforming} \\ \hline 
GCC &\texttt{\_\_attribute\_\_((pure))} & \texttt{[[gnu::pure]]} \\ \hline
Clang & \texttt{\_\_attribute\_\_((no\_sanitize))} &\texttt{[[clang::no\_sanitize]]} \\ \hline 
MSVC & \texttt{declspec(deprecated)} & \texttt{[[deprecated]]} \\ \thickhline
\end{tabular}
}
\end{threeparttable} 
\end{center}
\end{table}

Donec mattis, ex a ornare auctor, tortor felis pulvinar odio, eu mollis tortor orci id ligula. Vestibulum dignissim magna vitae lectus cursus, ac sollicitudin mi pulvinar. Praesent vulputate lorem lectus, eu pulvinar mi finibus eu. Donec sit amet massa massa. Sed eget molestie quam. Praesent iaculis diam ut nunc condimentum auctor. Morbi porttitor varius ante vel venenatis. Integer ut porttitor elit, id lacinia turpis. Integer id iaculis ante, placerat ornare tortor. Sed et consectetur est, sit amet rhoncus turpis.

\begin{table}[h!] 
\begin{center}
\begin{threeparttable}
\caption{Testing Table Numbering}\label{test-table1}\vspace{1.5ex} 
{\small \begin{tabular}{c|c|c}\thickhline 
\rowcolor[gray]{.9}   {\sffamily\bfseries Compiler} & {\sffamily\bfseries Compiler-Specific} &
{\sffamily\bfseries Standard-Conforming} \\ \hline 
GCC &\texttt{\_\_attribute\_\_((pure))} & \texttt{[[gnu::pure]]} \\ \hline
Clang & \texttt{\_\_attribute\_\_((no\_sanitize))} &\texttt{[[clang::no\_sanitize]]} \\ \hline 
MSVC & \texttt{declspec(deprecated)} & \texttt{[[deprecated]]} \\ \thickhline
\end{tabular}
}
\end{threeparttable} 
\end{center}
\end{table}

Phasellus ac tellus ipsum. Donec ornare id urna a blandit. Pellentesque finibus nulla augue, vitae pretium justo imperdiet ac. Aenean condimentum nibh eget turpis viverra gravida. Nullam diam lectus, egestas id ex quis, blandit facilisis libero. Proin laoreet ante dictum tristique finibus. Vestibulum pretium varius ipsum, sed blandit ex venenatis a. Nam pharetra pretium accumsan. Integer accumsan purus elementum tortor aliquam, commodo semper leo commodo. Nunc finibus varius erat, non hendrerit leo dapibus at. Donec mattis porta ex, eu ornare nibh condimentum id. Sed sit amet erat sit amet urna volutpat ullamcorper. Aenean feugiat eget orci ac feugiat. Nulla facilisi. Cras augue lacus, placerat sit amet vestibulum at, mollis sit amet nisl.

Nam id sem sagittis, placerat justo eget, eleifend justo. Sed ultrices rhoncus quam ut auctor. Nam pellentesque risus orci, a rhoncus arcu lobortis in. Vestibulum tristique nisi sed sollicitudin aliquam. Praesent ut odio sapien. Ut pretium sollicitudin nisi sit amet blandit. Cras ut eleifend elit. Aliquam vel tincidunt lacus, non sodales justo. Sed euismod sapien non diam euismod, sit amet vulputate nisi rhoncus. Vestibulum leo diam, dapibus suscipit blandit molestie, suscipit eget ante. Donec tristique rhoncus purus, nec pellentesque lorem volutpat non. Quisque purus dui, egestas ut ultricies sed, pharetra sit amet turpis. Praesent ut ligula porttitor, convallis velit sit amet, rhoncus urna. 
  %%% to create a sample file for the Proofreader

\newpage
%\section[Consecutive {\tt >}s]{Consecutive Right-Angle Brackets}\label{consecutive-right-angle-brackets}
% 13 Jan 2021, copyedits in and proofed
% 27 Jan, code checked and replaced as needed by Josh
% 01 Feb, 2021 JMB - added a missing include to hidden listing (line 88)
% 09 Feb 2021, FPPs set
% 17 Feb 2021, glossary tags would have been updated, but there are zero glossary entries in this feature. 

\emcppsFeature{
toclong={Consecutive Right-Angle Brackets},
tocshort={Consecutive {\TOCCode >}s},
short={Consecutive \lstinline!>!s},
long={Consecutive Right-Angle Brackets},
rhshort={Consecutive {\RHCode >}s},
}{consecutive-right-angle-brackets}
%\section[Consecutive {\tt >}s]{Consecutive Right-Angle Brackets\sectionmark{Consecutive {\RHCode >}s}}\label{consecutive-right-angle-brackets}\sectionmark{Consecutive {\RHCode >}s}
\setcounter{table}{0}
\setcounter{footnote}{0}
\setcounter{lstlisting}{0}


In the context of template argument lists, \lstinline!>>! is parsed as two
separate closing angle brackets.

\subsection[Description]{Description}\label{description}

Prior to C++11, a pair of consecutive right-pointing angle brackets anywhere in
the source code was always interpreted as a bitwise right-shift
operator, making an intervening space mandatory for them to be treated
as separate closing-angle-bracket tokens:

\begin{emcppshiddenlisting}[emcppsbatch=e1]
#include <vector>  // (ù{\codeincomments{std::vector}}ù)
\end{emcppshiddenlisting}
\begin{emcppslisting}[emcppsbatch=e1]
// C++03
std::vector<std::vector<int>> v0;   // annoying compile-time error in C++03
std::vector<std::vector<int> > v1;  // OK
\end{emcppslisting}

\noindent To facilitate the common use case above, a special rule was added
whereby, when parsing a template-argument expression, \emph{non-nested}
(i.e., within parentheses) appearances of \lstinline!>!, \lstinline!>>!,
\lstinline!>>>!, and so on are to be treated as separate closing angle brackets:

\begin{emcppshiddenlisting}[emcppsbatch=e2]
#include <vector>  // (ù{\codeincomments{std::vector}}ù)
\end{emcppshiddenlisting}
\begin{emcppslisting}[emcppsbatch=e2]
// C++11
std::vector<std::vector<int>> v0;               // OK
std::vector<std::vector<std::vector<int>>> v1;  // OK
\end{emcppslisting}

\subsubsection[Using the greater-than or right-shift operators within template-argument expressions]{Using the greater-than or right-shift operators within\\[0.5ex] template-argument expressions}\label{using-the-greater-than-or-right-shift-operators-within-template-argument-expressions}

For templates that take only type parameters, there's no issue. When the
template parameter is a non-type, however, the greater-than or right-shift operators might be
useful. In the unlikely event that we need either the greater-than
operator (\lstinline!>!) or the right-shift operator (\lstinline!>>!) within a
non-type template-argument expression, we can achieve our goal by
nesting that expression within parentheses:

\begin{emcppslisting}
const int i = 1, j = 2;  // arbitrary integer values (used below)

template <int I> class C { /*...*/ };
    // class (ù{\codeincomments{C}}ù) taking non-type template parameter (ù{\codeincomments{I}}ù) of type (ù{\codeincomments{int}}ù)

C<i > j>    a1;  // Error, always has been
C<i >> j>   b1;  // Error, in C++11, OK in C++03
C<(i > j)>  a2;  // OK
C<(i >> j)> b2;  // OK
\end{emcppslisting}


\noindent In the definition of \lstinline!a1! above, the first \lstinline!>! is
interpreted as a closing angle bracket, and the subsequent \lstinline!j! is
(and always has been) a syntax error. In the case of \lstinline!b1!, the
\lstinline!>>! is, as of C++11, parsed as a pair of separate tokens in this
context, so the second \lstinline!>! is now considered an error. For
both \lstinline!a2! and \lstinline!b2!, however, the would-be operators appear
nested within parentheses and thus are blocked from matching any
active open angle bracket to the left of the parenthesized expression.

\subsection[Use Cases]{Use Cases}\label{use-cases}

\subsubsection[Avoiding annoying whitespace when composing template types]{Avoiding annoying whitespace when composing template types}\label{avoiding-annoying-whitespace-when-composing-template-types}

When using nested templated types (e.g.,~nested containers) in C++03,
having to remember to insert an intervening space between trailing angle
brackets added no value. What made it even more galling was that every
popular compiler was able to tell you confidently that you had forgotten
to leave the space. With this new
feature (rather, this repaired defect), we can now render closing angle
brackets contiguously, just like parentheses and square brackets:

\begin{emcppshiddenlisting}[emcppsbatch=e3]
#include <list>    // (ù{\codeincomments{std::list}}ù)
#include <vector>  // (ù{\codeincomments{std::vector}}ù)
#include <string>  // (ù{\codeincomments{std::string}}ù)
#include <map>     // (ù{\codeincomments{std::map}}ù)
\end{emcppshiddenlisting}
\begin{emcppslisting}[emcppsbatch=e3]
// OK in both C++03 and C++11
std::list<std::map<int, std::vector<std::string> > > idToNameMappingList1;

// OK in C++11, compile-time error in C++03
std::list<std::map<int, std::vector<std::string>>>   idToNameMappingList2;
\end{emcppslisting}


\subsection[Potential Pitfalls]{Potential Pitfalls}\label{potential-pitfalls}

\subsubsection[Some C++03 programs may stop compiling in C++11]{Some C++03 programs may stop compiling in C++11}\label{some-c++03-programs-may-stop-working-in-c++11}

If a right-shift operator is used in a template expression, the newer
parsing rules may result in a compile-time error where before there was
none:

\begin{emcppshiddenlisting}[emcppsbatch=ped1,emcppsstandards=c++03]
template <int i> class T {};
\end{emcppshiddenlisting}
\begin{emcppslisting}[emcppsbatch=ped1]
T<1 >> 5> t;  // worked in C++03, compile-time error in C++11
\end{emcppslisting}

\noindent The easy fix is simply to parenthesize the expression:

\begin{emcppshiddenlisting}[emcppsbatch=ped2,emcppsstandards={c++11,c++14}]
template <int i> class T {};
\end{emcppshiddenlisting}
\begin{emcppslisting}[emcppsbatch=ped2]
T<(1 >> 5)> t;  // OK
\end{emcppslisting}

\noindent This rare syntax error is invariably caught at compile time, avoiding undetected surprises at run time.

\subsubsection[The meaning of a C++03 program can, in theory, silently change in C++11]{The meaning of a C++03 program can, in theory, silently change in C++11}\label{the-meaning-of-a-c++03-program-can-(in-theory)-silently-change-in-c++11}

Though pathologically rare, the same valid expression can, in theory, have a different interpretation in C++11 than it had when compiled for C++03.
Consider the case{\cprotect\footnote{Example adapted from \cite{gustedt13}}} where the
\lstinline!>>! token is embedded as part of an expression involving
templates:

\begin{emcppslisting}[emcppsignore={Pathological}]
S<G< 0 >>::c>::b>::a
//   ^~~~~~~
\end{emcppslisting}

\noindent In the expression above, \lstinline!0!~\lstinline!>>::c! will be interpreted
as a \emph{bitwise right-shift operator} in C++03 but not in C++11. Writing a program that (1) compiles under both C++03 and
C++11 and (2) exposes the difference in parsing rules is possible:

\begin{emcppshiddenlisting}[emcppsbatch=e4]
#include <iostream>    // (ù{\codeincomments{std::cout}}ù)
\end{emcppshiddenlisting}
\begin{emcppslisting}[emcppsbatch=e4]
enum Outer { a = 1, b = 2, c = 3 };

template <typename> struct S
{
    enum Inner { a = 100, c = 102 };
};

template <int> struct G
{
    typedef int b;
};

int main()
{
    std::cout << (S<G< 0 >>::c>::b>::a) << '\n';
}
\end{emcppslisting}

\noindent The program above will print \lstinline!100! when compiled for C++03 and
\lstinline!0! for C++11:

\begin{emcppslisting}[emcppsignore={Pathological}]
// C++03

//     (2) instantiation of (ù{\codeincomments{G<0>}}ù)
//    (ù{\codeincomments{$\|$}}ù)~~~~~~~~~~~~
//    (ù{\codeincomments{$\|\:\,|\:\,\|$}}ù)   (4) instantiation of (ù{\codeincomments{S<int>}}ù)
//  ~~(ù{\codeincomments{$\|\downarrow\|$}}ù)~~~~~~~~~~~~~~(ù{\codeincomments{$\downarrow$}}ù)
    S< G< 0 >>::c > ::b >::a
//    ~~(ù{\codeincomments{$\|\,\,\uparrow\,\,\|$}}ù)~~~~~~~~~(ù{\codeincomments{$\uparrow$}}ù)
//      (ù{\codeincomments{$\|\:\,\,\,|\:\,\,\,\|$}}ù) (3) type alias for (ù{\codeincomments{int}}ù)
//      (ù{\codeincomments{$\|$}}ù)~~~~~~~
// (1) bitwise right-shift ((ù{\codeincomments{0 >> 3}}ù))
\end{emcppslisting}


\begin{emcppslisting}[emcppsignore={Pathological}]
// C++11

//
//
//  (2) compare ((ù{\codeincomments{>}}ù)) (ù{\codeincomments{Inner::c}}ù) and (ù{\codeincomments{Outer::b}}ù)
//  (ù{\codeincomments{$\downarrow$}}ù) ~~~~~~~~~~~~~~~~~~
    S< G< 0 >>::c > ::b >::a
//  (ù{\codeincomments{$\uparrow$}}ù) ~~~~~~~~~
//  (1) instantiation of (ù{\codeincomments{S<G<0>>}}ù)
//
//
\end{emcppslisting}

\pagebreak[4]%%%%%%%%
\noindent Though theoretically possible, programs that (1) are syntactically valid
in both C++03 and C++11 and (2) have distinct semantics have not
emerged in practice anywhere that we are aware of.

\subsection[Annoyances]{Annoyances}\label{annoyances}

\hspace{\fill}

\subsection[See Also]{See Also}\label{see-also}

\hspace{\fill}

\subsection[Further Reading]{Further Reading}\label{further-reading}

\begin{itemize}
\item{Daveed Vandevoorde, \textit{Right Angle Brackets,} \cite{vandevoorde05}}
\end{itemize}




\newpage
%\section[\tt{decltype}]{Operator for Extracting Expression Types}\label{decltype}
% copyedits in and proofed
% 27 Jan, code checked and replaced as needed by Josh
% 9 Feb 2021, FPPs set
% 17 Feb 2021, glossary tags updated

\emcppsFeature{
    toclong={Operator for Extracting Expression Types},
    tocshort={\TOCCode decltype},
    short={\lstinline!decltype!},
    long={Operator for Extracting Expression Types},
    rhshort={\RHCode decltype},
    }{decltype}
%\section[{\tt decltype}]{Operator for Extracting Expression Types\sectionmark{\RHCode decltype}}\label{decltype}\sectionmark{\RHCode decltype}
\setcounter{table}{0}
\setcounter{footnote}{0}
\setcounter{lstlisting}{0}


The keyword \lstinline!decltype! enables the compile-time inspection of the \emcppsgloss{declared type} of an \emcppsgloss{entity} \emph{or} the type and
\emcppsgloss{value category} of an expression. Note that the special construct \lstinline!decltype(auto)! has a separate meaning; see \featureref{\locationf}{decltypeauto}.
%``\titleref{decltypeauto}" on page~\pageref{decltypeauto}.}

\subsection[Description]{Description}\label{description}

What results from the use of \lstinline!decltype! depends on the nature of
its operand.

\subsubsection[Use with entities]{Use with entities}\label{use-with-(typically-named)-entities}

%If an unparenthesized operand is either an \textbf{id-expression} that names an entity (or a non-type template parameter) or a \textbf{class member access expression}\glossary{class member access expression} (that identifies a class member), \texttt{decltype} yields the \emph{declared type} (the type of the \emph{entity} indicated by the operand):
If the operand is an unparenthesized \emcppsgloss[id expression]{id-expression} or unparenthesized member access, \lstinline!decltype! yields the \emph{declared type}, meaning the type of the \emph{entity} indicated by the operand:


\begin{emcppshiddenlisting}[emcppsbatch=e1]
#include <string>  // (ù{\codeincomments{std::string}}ù)
\end{emcppshiddenlisting}
\begin{emcppslisting}[emcppsbatch=e1]
int i;                // (ù{\codeincomments{decltype(i)}}ù)   -> (ù{\codeincomments{int}}ù)
std::string s;        // (ù{\codeincomments{decltype(s)}}ù)   -> (ù{\codeincomments{std::string}}ù)
int* p;               // (ù{\codeincomments{decltype(p)}}ù)   -> (ù{\codeincomments{int}}ù)*
const int& r = *p;    // (ù{\codeincomments{decltype(r)}}ù)   -> (ù{\codeincomments{const int\&}}ù)
struct { char c; } x; // (ù{\codeincomments{decltype(x.c)}}ù) -> (ù{\codeincomments{char}}ù)
double f();           // (ù{\codeincomments{decltype(f)}}ù)   -> (ù{\codeincomments{double()}}ù)
double g(int);        // (ù{\codeincomments{decltype(g)}}ù)   -> (ù{\codeincomments{double(int)}}ù)
\end{emcppslisting}


\subsubsection[Use with expressions]{Use with expressions}\label{use-with-(unnamed)-expressions}

When \lstinline!decltype! is used with any other expression \lstinline!E! of type \lstinline!T!, including parenthesized \emcppsgloss[id expression]{id-expression} or parenthesized member access, the result incorporates both the expression's type and
its \emcppsgloss{value category} (see \featureref{\locationc}{Rvalue-References}):
%\mclhead{Value category of \texttt{E}} & \mclhead{Result of
%\texttt{decltype(E)}} \\
%\emph{prvalue} & \texttt{T} \\
%\emph{lvalue} & \texttt{T\&} \\
%\emph{xvalue} & \texttt{T\&\&} \\
%%%%%%%%%%% changing this per Greg & design discussion
\begin{center}
{\small \begin{tabular}{c|c}\thickhline
\rowcolor[gray]{.9}   {\sffamily\bfseries Value category of \lstinline!E!}
& {\sffamily\bfseries Result of {\lstinline!decltype(E)!}} \\ \hline
{\romeovalueinside prvalue} &\lstinline!T! \\ \hline
{\romeovalueinside lvalue} & \lstinline!T&! \\ \hline
{\romeovalueinside xvalue} & \lstinline!T&&! \\ \thickhline
\end{tabular}
}
\end{center}
%\noindent The three integer expressions below illustrate the various value
%categories:
%
%begin{emcppslisting}
%decltype(0)   // -> (ù{\codeincomments{int}}ù)   (*prvalue* category)
%
%int i;
%decltype((i)) // -> (ù{\codeincomments{int\&}}ù)  (*lvalue* category)
%
%int&& g();
%decltype(g()) // -> (ù{\codeincomments{int\&\&}}ù) (*xvalue* category)
%\end{emcppslisting}
%begin{emcppslisting}
%int g();
%decltype (g) g1; // id-expression,  (ù{\codeincomments{int ()();}}ù)
%decltype (g()) g2; // function call, (ù{\codeincomments{int}}ù)
%decltype ((g)) g3 = g; // id expression in parentheses, treat as lvalue,  (ù{\codeincomments{int (&)()}}ù)
%decltype ((g())) g4; // function call, parentheses have no effect, prvalue, (ù{\codeincomments{int}}ù)
%\end{emcppslisting}
%%%% new change from Josh
\noindent In general, \romeovalue{prvalues} can be passed to \lstinline!decltype! in a number of ways, including numeric literals, function calls that return by value, and explicitly created temporaries:

\begin{emcppslisting}
decltype(0)   i; // -> (ù{\codeincomments{int}}ù)
int f();
decltype(f()) j; // -> (ù{\codeincomments{int}}ù)
struct S{};
decltype(S()) k; // -> (ù{\codeincomments{S}}ù)
\end{emcppslisting}

\noindent An entity name passed to \lstinline!decltype!, as mentioned above, produces the type of the entity. If an entity name is enclosed in an additional set of parentheses, however, \lstinline!decltype! interprets its argument as an expression and its result incorporates the value category:

\begin{emcppslisting}[emcppsbatch=e2]
int i;
decltype(i)   l = i; // -> (ù{\codeincomments{int}}ù)
decltype((i)) m = i; // -> (ù{\codeincomments{int\&}}ù)
\end{emcppslisting}

\noindent Similarly, for all other \romeovalue{lvalue} expressions, the result of \lstinline!decltype! will be an \romeovalue{lvalue} reference:

\begin{emcppslisting}[emcppsbatch=e2]
int* pi = &i;
decltype(*pi) j = *pi; // -> (ù{\codeincomments{int\&}}ù)
decltype(++i) k = ++i; // -> (ù{\codeincomments{int\&}}ù)
\end{emcppslisting}

\noindent Finally, the value category of the expression will be an \romeovalue{xvalue} if it is a cast to or a function returning an \romeovalue{rvalue} reference:

\begin{emcppslisting}
int i;
decltype(static_cast<int&&>(i)) j = static_cast<int&&>(i); // -> (ù{\codeincomments{int\&\&}}ù)
int&& g();
decltype(g()) k = g();                                     // -> (ù{\codeincomments{int\&\&}}ù)
\end{emcppslisting}


\noindent Much like the \lstinline!sizeof!
operator (which is also resolved at compile time), the expression operand of \lstinline!decltype! is not evaluated:

%begin{emcppslisting}
%int i = 0;
%decltype(i++) j;  // equivalent to (ù{\codeincomments{int j;}}ù)
%assert(i == 0);   // The function (ù{\codeincomments{next}}ù) is never invoked.
%\end{emcppslisting}
\begin{emcppshiddenlisting}[emcppsbatch=e3]
#include <cassert>  // standard C (ù{\codeincomments{assert}}ù) macro
\end{emcppshiddenlisting}
\begin{emcppslisting}[emcppsbatch=e3]
void test1()
{
    int i = 0;
    decltype(i++) j;  // equivalent to (ù{\codeincomments{int j;}}ù)
    assert(i == 0);   // The expression (ù{\codeincomments{i++}}ù) was not evaluated.
}
\end{emcppslisting}
Note that the choice of using the postfix increment is significant; the prefix increment yields a different type:
\begin{emcppslisting}[emcppsbatch=e3]
void test2()
{
    int i = 0;
    int m = 1;
    decltype(++i) k = m; // equivalent to (ù{\codeincomments{int\& k = m;}}ù)
    assert(i == 0); // The expression (ù{\codeincomments{++1}}ù) is not evaluated.
}
\end{emcppslisting}

\subsection[Use Cases]{Use Cases}\label{use-cases-decltype}

\subsubsection[Avoiding unnecessary use of explicit typenames]{Avoiding unnecessary use of explicit typenames}\label{avoiding-unnecessary-use-of-explicit-typenames}

Consider two logically equivalent ways of declaring a vector of
iterators into a list of \lstinline!Widgets!:

\begin{emcppshiddenlisting}[emcppsbatch=e4]
#include <list>    // (ù{\codeincomments{std::list}}ù)   
#include <vector>  // (ù{\codeincomments{std::vector}}ù)
struct Widget {};
\end{emcppshiddenlisting}
\begin{emcppslisting}[emcppsbatch=e4]
std::list<Widget> widgets;
std::vector<std::list<Widget>::iterator> widgetIterators;
    // (1) The full type of (ù{\codeincomments{widgets}}ù) needs to be restated, and (ù{\codeincomments{iterator}}ù)
    // needs to be explicitly named.
\end{emcppslisting}
\begin{emcppshiddenlisting}[emcppsbatch=e4-1]
#include <list>    // (ù{\codeincomments{std::list}}ù)   
#include <vector>  // (ù{\codeincomments{std::vector}}ù)
struct Widget {};
\end{emcppshiddenlisting}
\begin{emcppslisting}[emcppsbatch=e4-1]
std::list<Widget> widgets;
std::vector<decltype(widgets.begin())> widgetIterators;
    // (2) Neither (ù{\codeincomments{std::list}}ù) nor (ù{\codeincomments{Widget}}ù) nor (ù{\codeincomments{iterator}}ù) need be named
    // explicitly.
\end{emcppslisting}

\noindent Notice that, when using \lstinline!decltype!, if the C++ type representing
the widget changes (e.g., from \lstinline!Widget! to, say,
\lstinline!ManagedWidget!) or the container used changes (e.g., from
\lstinline!std::list! to \lstinline!std::vector!), the declaration of
\lstinline!widgetIterators! does not necessarily need to change.

\subsubsection[Expressing type-consistency explicitly]{Expressing type-consistency explicitly}\label{expressing-type-consistency-explicitly}

In some situations, repetition of explicit type names might
inadvertently result in latent defects caused by mismatched types during
maintenance. For example, consider a \lstinline!Packet! class exposing a
\lstinline!const! member function that returns a value of type \lstinline!std::uint8_t!
representing the length of the packet's checksum:

\begin{emcppshiddenlisting}[emcppsbatch=e6]
#include <cstdint>  // (ù{\codeincomments{std::uint8\_t}}ù)
\end{emcppshiddenlisting}
\begin{emcppslisting}[emcppsbatch=e6]
class Packet
{
    // ...
    
public:
    std::uint8_t checksumLength() const;
};
\end{emcppslisting}

%%%% lots of manual linebreaking to balance the whole paragraph due to the code term at the end
\noindent This unsigned 8-bit type was selected to minimize bandwidth usage
as the checksum\linebreak length is sent over the network. Next, picture a loop
that computes the checksum of a \lstinline!Packet!, using the same type for its iteration variable to match the
type returned by\linebreak%%%%%% 
\lstinline!Packet::checksumLength!:

\begin{emcppshiddenlisting}[emcppsbatch=e7]
#include <cstdint>  // (ù{\codeincomments{std::uint8\_t}}ù)
struct Packet
{
public:
    std::uint8_t checksumLength() const;
    std::uint8_t nthByte(std::uint8_t i) const;
};
struct Checksum {
    void appendByte(std::uint8_t b);
};
\end{emcppshiddenlisting}
\begin{emcppslisting}[emcppsbatch=e7]
void f()
{
    Checksum sum;
    Packet data;

    for (std::uint8_t i = 0; i < data.checksumLength(); ++i)  // brittle
    {
        sum.appendByte(data.nthByte(i));
    }
}
\end{emcppslisting}


Now suppose that, over time, the data transmitted by the \lstinline!Packet!
type grows to the point where the range of an \lstinline!std::uint8_t!
value might not be enough to ensure a sufficiently reliable checksum. If
the type returned by \lstinline!checksumLength()! is changed to, say,
\lstinline!std::uint16_t! without updating the type of the iteration
variable \lstinline!i! in lockstep, the loop might
silently{\cprotect\footnote{As of this writing, neither
GCC~9.3 nor Clang~10.0.0 provide
a warning (using \lstinline!-Wall!, \lstinline!-Wextra!, and
\lstinline!-Wpedantic!) for the comparison between \lstinline!std::uint8_t!
and \lstinline!std::uint16_t! --- even if (1) the value returned by
\lstinline!checksumLength! does not fit in a 8-bit integer, and (2) the
body of the function is visible to the compiler. Decorating
\lstinline!checksumLength! with \lstinline!constexpr! causes
\lstinline!clang++! to issue a warning, but this is clearly not a general
solution.}} become infinite.{\cprotect\footnote{The loop
variable is promoted to an \lstinline!unsigned! \lstinline!int! for
comparison purposes but wraps to 0 whenever its value prior to
  being incremented is 255.}}

Had \lstinline!decltype(packet.checksumLength())! been used to express the
type of \lstinline!i!, the types would have remained consistent, and the
ensuing defect would naturally have been avoided:

\begin{emcppshiddenlisting}[emcppsbatch=e7]
void g(const Packet &data)
{
\end{emcppshiddenlisting}
\begin{emcppslisting}[emcppsbatch=e7]
// ...
for (decltype(data.checksumLength()) i = 0; i < data.checksumLength(); ++i)
// ...
\end{emcppslisting}
\begin{emcppshiddenlisting}[emcppsbatch=e7]
{}
} // end g
\end{emcppshiddenlisting}    

\subsubsection[Creating an auxiliary variable of generic type]{Creating an auxiliary variable of generic type}\label{creating-an-auxiliary-variable-of-generic-type}

Consider the task of implementing a generic \lstinline!loggedSum! function
template that returns the sum of two arbitrary objects \lstinline!a! and
\lstinline!b! after logging both the operands and the result value (e.g.,
for debugging or monitoring purposes). To avoid computing the possibly
expensive sum twice, we decide to create an auxiliary function-scope
variable, \lstinline!result!. Since the type of the sum depends on both
\lstinline!a! and \lstinline!b!, we can use
\lstinline!decltype(a!~\lstinline!+!~\lstinline!b)! to infer the type for both
the trailing return type of the
function (see %Section~\ref{trailing-function-return-types},
%``\titleref{trailing-function-return-types}" on page~\pageref{trailing-function-return-types})
\featureref{\locationa}{trailing-function-return-types}) and the auxiliary variable:

\begin{emcppshiddenlisting}[emcppsbatch=e8]
struct LogTrace {};  
constexpr LogTrace LOG_TRACE; // Quick and dirty syntactic simulation of logger
template <typename T>
const LogTrace& operator<<(const LogTrace& lhs, const T&) { return lhs; }
\end{emcppshiddenlisting}
\begin{emcppslisting}[emcppsbatch=e8]
template <typename A, typename B>
auto loggedSum(const A& a, const B& b)
    -> decltype(a + b)                 // (1) exploiting trailing return types
{
    decltype(a + b) result = a + b;    // (2) auxiliary generic variable
    LOG_TRACE << a << " + " << b << " = " << result;
    return result;
}
\end{emcppslisting}
Using \lstinline!decltype(a!~\lstinline!+!~\lstinline!b)! as a return type is significantly different from relying on automatic \emcppsgloss[return type deduction]{\lstinline!return!-type
deduction}; see \featureref{\locationc}{auto-feature}. Note that this particular use involves significant repetition of the expression \lstinline!a+b!.  See
%\textit{\titleref{annoyances-decltype}: \titleref{decltype-mechanical}} on page~\pageref{decltype-mechanical}
\intraref{annoyances-decltype}{decltype-mechanical}
for a discussion of ways in which this might be avoided.

\subsubsection[Determining the validity of a generic expression]{Determining the validity of a generic expression}\label{determining-the-validity-of-a-generic-expression}

In the context of generic-library development, \lstinline!decltype! can be
used in conjunction with \emcppsgloss{SFINAE} (``Substitution Failure Is Not An Error") to validate an expression involving a
template parameter.

For example, consider the task of writing a generic \lstinline!sortRange!
function template that, given a \emcppsgloss{range}, either invokes the
\lstinline!sort! member function of the argument (the one specifically
optimized for that type) if available or falls back to the more
general \lstinline!std::sort!:

\begin{emcppslisting}
template <typename Range>
void sortRange(Range& range)
{
    sortRangeImpl(range, 0);
}
\end{emcppslisting}

\noindent The client-facing \lstinline!sortRange! function (above) delegates its
behavior to an overloaded\linebreak%%%%
 \lstinline!sortRangeImpl! function (below),
invoking the latter with the \lstinline!range! and a \emcppsgloss{disambiguator}
of type \lstinline!int!. The type of this additional parameter, whose value
is arbitrary, is used to give priority to the \lstinline!sort! member
function at compile time by exploiting overload resolution rules in
the presence of an implicit (\emph{standard}) conversion from
\lstinline!int! to \lstinline!long!:

\begin{emcppshiddenlisting}[emcppsbatch=e9]
#include <algorithm>  // (ù{\codeincomments{std::sort}}ù)
\end{emcppshiddenlisting}
\begin{emcppslisting}[emcppsbatch=e9]
template <typename Range>
void sortRangeImpl(Range& range,
                   long)                  // low priority: standard conversion
{
    // fallback implementation
    std::sort(std::begin(range), std::end(range));
}
\end{emcppslisting}

\noindent The fallback overload of \lstinline!sortRangeImpl! (above) will accept a
\lstinline!long! \emcppsgloss{disambiguator}, requiring a standard conversion from
\lstinline!int!, and will simply invoke \lstinline!std::sort!. The more
specialized overload of \lstinline!sortRangeImpl! (below) will accept an
\lstinline!int! \emcppsgloss{disambiguator} requiring no conversions and thus
will be a better match, provided a range-specific sort is available:

\begin{emcppslisting}
template <typename Range>
void sortRangeImpl(Range& range,
                   int,                          // high priority: exact match
                   decltype(range.sort())* = 0)  // check expression validity
{
    // optimized implementation
    range.sort();
}
\end{emcppslisting}

\noindent Note that, by exposing \lstinline!decltype(range.sort())! as part of
\lstinline!sortRangeImpl!'s declaration, the more specialized overload will
be discarded during template substitution if \lstinline!range.sort()! is
not a valid expression for the deduced \lstinline!Range!
type.{\cprotect\footnote{The technique of exposing a possibly unused
unevaluated expression --- e.g., using \lstinline!decltype! --- in a function's
declaration for the purpose of expression-validity detection prior to
template instantiation is commonly known as \emcppsgloss{expression SFINAE},
which is a restricted form of the more general (classical) SFINAE, and acts exclusively on
expressions visible in a function's signature rather than on frequently obscure
  template-based type computations.}}
\pagebreak%%%%%%%  

The relative position of
\lstinline!decltype(range.sort())! in the signature of
\lstinline!sortRangeImpl! is not significant, as long as it is visible to
the compiler
during template substitution. The example shown in the main text uses a function parameter that is defaulted to
\lstinline!nullptr!. Alternatives involving a trailing return type or a
default template argument are also viable:

%\begin{emcppslisting}[basicstyle={\ttfamily\footnotesize}]
%template <typename Range>
%auto sortRangeImpl(Range& range, int) -> decltype(range.sort(), void());
%    // The comma operator is used to force the return type to (ù{\codeincomments{void}}ù),
%    // regardless of the return type of (ù{\codeincomments{range.sort()}}ù).
%
%template <typename Range, typename = decltype(std::declval<Range&>().sort()>
%auto sortRangeImpl(Range& range, int);
%    // (ù{\codeincomments{std::declval}}ù) is used to generate a reference to (ù{\codeincomments{Range}}ù) that can
%    // be used in an unevaluated expression.
%\end{emcppslisting}
\begin{emcppslisting}
#include <utility>  // (ù{\codeincomments{declval}}ù)
template <typename Range>
auto sortRangeImpl(Range& range, int) -> decltype(range.sort(), void());
    // The comma operator is used to force the return type to (ù{\codeincomments{void}}ù),
    // regardless of the return type of (ù{\codeincomments{range.sort()}}ù).

template <typename Range, typename = decltype(std::declval<Range&>().sort())>
auto sortRangeImpl(Range& range, int) -> void;
    // (ù{\codeincomments{std::declval}}ù) is used to generate a reference to (ù{\codeincomments{Range}}ù) that can
    // be used in an unevaluated expression.
\end{emcppslisting}

\noindent Putting it all together, we see that exactly two possible
outcomes exist for the original client-facing \lstinline!sortRange! function
invoked with a range argument of type \lstinline!R!:

\begin{itemize}
\item{If \lstinline!R! does have a \lstinline!sort! member function, the more specialized overload of\linebreak[4] \lstinline!sortRangeImpl! will be viable as \lstinline!range.sort()! is a well-formed expression and preferred because the \emcppsgloss{disambiguator} \lstinline!0! (of type \lstinline!int!) requires no conversion.}
\item{Otherwise, the more specialized overload will be discarded during template substitution as \lstinline!range.sort()! is not a well-formed expression, and the only remaining more general \lstinline!sortRangeImpl! overload will be chosen instead.}
\end{itemize}

\subsection[Potential Pitfalls]{Potential Pitfalls}\label{potential-pitfalls}

Perhaps surprisingly, \lstinline!decltype(x)! and \lstinline!decltype((x))!
will sometimes yield different results for the same expression
\lstinline!x!:

\begin{emcppslisting}[emcppsbatch=e10]
int i = 0; // (ù{\codeincomments{decltype(i)}}ù) yields (ù{\codeincomments{int}}ù).
           // (ù{\codeincomments{decltype((i))}}ù) yields (ù{\codeincomments{int\&}}ù).
\end{emcppslisting}

\noindent In the case where the unparenthesized operand is an entity having a
declared type \lstinline!T! and the parenthesized operand is an expression
whose value category is represented (by \lstinline!decltype!) as the same
type \lstinline!T!, the results will coincidentally be the same:

\begin{emcppslisting}[emcppsbatch=e10]
int& ref = i;  // (ù{\codeincomments{decltype(ref)}}ù) yields (ù{\codeincomments{int\&}}ù).
               // (ù{\codeincomments{decltype((ref))}}ù) yields (ù{\codeincomments{int\&}}ù).
\end{emcppslisting}

\noindent Wrapping its operand with parentheses ensures \lstinline!decltype! yields
the \emcppsgloss{value category} of a given expression. This technique can be
useful in the context of metaprogramming --- particularly in the case of
\emcppsgloss{value category} propagation.

\subsection[Annoyances]{Annoyances}\label{annoyances-decltype}

\subsubsection{Mechanical repetition of expressions might be required}\label{decltype-mechanical}

As mentioned in
%\textit{\titleref{use-cases-decltype}: \titleref{creating-an-auxiliary-variable-of-generic-type}} on page~\pageref{creating-an-auxiliary-variable-of-generic-type},
\intraref{use-cases-decltype}{creating-an-auxiliary-variable-of-generic-type},
using \lstinline!decltype! to capture a value of an expression that is about to be used, or for the return value of an expression, can often lead to repeating the same expression in multiple places (three distinct ones in the earlier example).

An alternate solution to this problem is to capture the result of the \lstinline!decltype! expression in a \lstinline!typedef!, \lstinline!using! type alias, or as a defaulted template parameter --- but that runs into the problem that it can be used only once the expression is valid.  A defaulted \lstinline!template! parameter cannot reference parameter names as it is written before them, and a type alias cannot be introduced prior to the return type being needed.  A solution to this problem lies in using standard library function \lstinline!std::declval! to create expressions of the appropriate type without needing to reference the actual function parameters by name:

\begin{emcppshiddenlisting}[emcppsbatch=e11]
#include <utility>  // (ù{\codeincomments{std::declval}}ù)
struct LogTrace {};  
constexpr LogTrace LOG_TRACE; // Quick and dirty syntactic simulation of logger
template <typename T>
const LogTrace& operator<<(const LogTrace& lhs, const T&) { return lhs; }
\end{emcppshiddenlisting}
\begin{emcppslisting}[emcppsbatch=e11]
template <typename A, typename B,
          typename Result = decltype(std::declval<const A&>() +
                                     std::declval<const B&>())>
Result loggedSum(const A& a, const B& b)
{
    Result result = a + b;  // no duplication of the (ù{\codeincomments{decltype}}ù) expression
    LOG_TRACE << a << " + " << b << "=" << result;
    return result;
}
\end{emcppslisting}

Here, \lstinline!std::declval!, a function that cannot be executed at runtime and is only appropriate for use in \emcppsgloss{unevaluated contexts}, produces an expression of the specified type.  When mixed with \lstinline!decltype!, this lets us determine the result types for expressions without needing to construct (or even being able to construct) objects of the needed types.

\subsection[See Also]{See Also}\label{see-also}

\begin{itemize}
\item{%``\titleref{alias-declarations-and-alias-templates}" on page~\pageref{alias-declarations-and-alias-templates} ---
\seealsoref{alias-declarations-and-alias-templates}{\locationa}Oftentimes, it is useful to give a name to the type yielded by \lstinline!decltype!, which is done with a \lstinline!using! alias.}
\item{%``\titleref{auto-feature}" on page~\pageref{auto-feature} ---
\seealsoref{auto-feature}{\locationc}The type computed by \lstinline!decltype! is similar to, but distinct from, the type deduction used by \lstinline!auto!.}
\item{%``\titleref{Rvalue-References}" on page~\pageref{Rvalue-References} ---
\seealsoref{Rvalue-References}{\locationc}The \lstinline!decltype! operator yields precise information on whether an expression is an \romeovalue{lvalue} or \romeovalue{rvalue}.}
\newpage%%%%%%%
\item{%``\titleref{decltypeauto}" on page~\pageref{decltypeauto} ---
\seealsoref{decltypeauto}{\locationf}\lstinline!decltype! type computation rules can be useful in conjunction with an \lstinline!auto! variable.}
\end{itemize}

\subsection[Further Reading]{Further Reading}\label{further-reading}






\newpage
%\section[Defaulted Functions]{Using {\SecCode =default} for Special Member Functions}\label{Defaulted-Special-Member-Functions}\label{defaulted-special-member-functions}
% copyedits in and proofed
% 8 Feb 2021 JMB, compilation fixes
% 9 Feb 2021, FPPs set
% 17 Feb 2021, glossary tags updated

\emcppsFeature{
    short={Defaulted Functions},
    long={Using {\SecCode =}\,{\SecCode default} for Special Member Functions},
    toclong={Using \lstinline!=!\,\lstinline!default! for Special Member Functions},
}{defaulted-special-member-functions}
\label{Defaulted-Special-Member-Functions}
\setcounter{table}{0}
\setcounter{footnote}{0}
\setcounter{lstlisting}{0}
%\section[Defaulted Functions]{Using {\SecCode =default} for Special Member Functions}\label{Defaulted-Special-Member-Functions}\label{defaulted-special-member-functions}


Use of \lstinline!=!\,\lstinline!default! in a \emcppsgloss[special member function]{special member function}'s
declaration instructs the compiler to attempt to generate the function
automatically.

\subsection[Description]{Description}\label{description}

An important aspect of C++ class design is the understanding that the
compiler will attempt to generate certain member functions to
\emph{create}, \emph{copy}, \emph{destroy}, and now
\emph{move} (see \featureref{\locationc}{Rvalue-References})
%Section~\ref{Rvalue-References}, ``\titleref{Rvalue-References}")
an object unless
developers implement some or all of these functions
themselves. Determining which of the \emcppsgloss[special member function]{special member functions}
will continue to be generated and which will be suppressed in the
presence of \emcppsgloss[user provided special member function]{user-provided special member functions}
requires remembering the numerous rules the compiler uses.

\subsubsection[Declaring a special member function explicitly]{Declaring a special member function explicitly}\label{declaring-a-special-member-function-explicitly}

The rules specifying what happens in the presence of one or more
user-provided special member functions are inherently complex and not
necessarily intuitive; in fact, some have been
deprecated. Specifically, even in the presence of a user-provided
destructor, both the copy constructor and the copy-assignment operator
have historically been generated implicitly. Relying on such generated
behavior is problematic because it is unlikely that a class
requiring a user-provided destructor will function correctly without
corresponding user-provided copy operations. As of C++11, reliance on
  such dubious implicitly generated behavior is deprecated.

Here, we
will briefly illustrate a few common cases and then refer you to Howard
Hinnant's now famous table (see page~\pageref{default-table1} of \intrarefnopage{appendix:-implicit-generation-of-special-member-functions})
%{\it\titleref{appendix:-implicit-generation-of-special-member-functions}})
to
demystify what's going on under the hood.

\paragraph[Example 1: Providing just the default constructor]{Example 1: Providing just the default constructor}\label{example-1:-providing-just-the-default-constructor}

Consider a \lstinline!struct! with a user-provided default constructor:

\begin{emcppslisting}
struct S1
{
    S1();  // user-provided default constructor
};
\end{emcppslisting}

\noindent A user-provided default constructor has no effect on other special
member functions. Providing any other constructor, however, will
suppress automatic generation of the default constructor. We can, however, use
\lstinline!=!\,\lstinline!default! to restore the constructor as a \emcppsgloss{trivial
operation}; see \intraref{default-use-cases}{restoring-the-generation-of-a-special-member-function-suppressed-by-another}.
%{\it\titleref{default-use-cases}: \titleref{restoring-the-generation-of-a-special-member-function-suppressed-by-another}} on page~\pageref{restoring-the-generation-of-a-special-member-function-suppressed-by-another}.
Note that a nondeclared
function is nonexistent, which means that it will \emph{not}
participate in overload resolution at all. In contrast, a
\emcppsgloss{deleted function} participates in overload resolution and, if
selected, results in a compilation failure; see \featureref{\locationa}{deleted-functions}.
%Section~\ref{deleted-functions}, ``\titleref{deleted-functions}."

\paragraph[Example 2: Providing just a copy constructor]{Example 2: Providing just a copy constructor}\label{example-2:-providing-just-a-copy-constructor}

Now, consider a \lstinline!struct! with a user-provided copy constructor:

\begin{emcppslisting}
struct S2
{
    S2(const S2&);  // user-provided copy constructor
};
\end{emcppslisting}

\noindent A user-provided copy constructor (1) suppresses the generation of the
default constructor and both move operations and (2) allows implicit
generation of both the copy-assignment operator and the destructor.
Similarly, providing just the copy-assignment operator would allow the
compiler to implicitly generate both the copy constructor and the
destructor, but, in this case, it would also generate the default
constructor. Note that --- in either of these
cases --- relying on the compiler's implicitly generated copy operation
is deprecated.

\paragraph[Example 3: Providing just the destructor]{Example 3: Providing just the destructor}\label{example-3:-providing-just-the-destructor}

Finally, consider a \lstinline!struct! with a user-provided destructor:


\begin{emcppslisting}
struct S3
{
    ~S3();  // user-provided destructor
};
\end{emcppslisting}

\noindent A user-provided destructor suppresses the generation of move
operations but still allows copy operations to be generated. Again,
relying on either of these implicitly compiler-generated copy
operations is deprecated.

\paragraph[Example 4: Providing more than one special member function]{Example 4: Providing more than one special member function}\label{example-4:-providing-more-than-one-special-member}

When more than one special member function is declared explicitly, the
\emph{union} of their respective declaration suppressions and the
\emph{intersection} of their respective implicit generations pertain ---
e.g., if just the default constructor and destructor are provided
(\lstinline!S1! + \lstinline!S3! in Examples 1 and 3), then the declarations of both move
operations are suppressed, and both copy operations are generated
implicitly.

\subsubsection[Defaulting the first declaration of a special member function explicitly]{Defaulting the first declaration of a special member function explicitly}\label{defaulting-the-first-declaration-of-a-special-member-function-explicitly}

Using the \lstinline!=!\,\lstinline!default! syntax with the first declaration
of a special member function instructs the compiler to synthesize such a
function automatically without treating it as being
user provided. The compiler-generated version for a
special member function is required to call the corresponding special
member functions on every base class in base-class-declaration
order and then every data member of the encapsulating type in
declaration order (regardless of any access specifiers). Note that the
destructor calls will be in exactly the opposite order of the
  other special-member-function calls.

For example, consider struct \lstinline!S4! (in the code snippet below) in which we have chosen
to make explicit that the copy operations are to be autogenerated by
the compiler; note, in particular, that implicit declaration and generation
of each of the other special member functions is left unaffected.

\begin{emcppslisting}
struct S4
{
    S4(const S4&) = default;             // copy constructor
    S4& operator=(const S4&) = default;  // copy-assignment operator

    // has no effect on other other four special member functions, i.e.,
    // implicitly generates the default constructor, the destructor,
    // the move constructor, and the move-assignment operator
};
\end{emcppslisting}

\noindent A defaulted declaration may appear with any \emcppsgloss{access specifier} ---
i.e., \lstinline!private!,\linebreak[4] \lstinline!protected!, or \lstinline!public! --- and access to that generated function will be regulated accordingly:

\begin{emcppslisting}
struct S5
{
private:
    S5(const S5&) = default;             // private copy constructor
    S5& operator=(const S5&) = default;  // private copy-assignment operator

protected:
    ~S5() = default;                     // protected destructor

public:
     S5() = default;                     // public default constructor
};
\end{emcppslisting}

\noindent In the example above, copy operations exist for use by \emph{member} and
\emph{friend} functions only. Declaring the destructor
\lstinline!protected! or \lstinline!private! limits which functions can create
automatic variables of the specified type to those functions with the
appropriately privileged access to the class. Declaring the default
constructor \lstinline!public! is necessary to avoid its declaration's
being suppressed by another constructor (e.g., the private copy
constructor in the code snippet above) or \emph{any} move operation.

In short, using \lstinline!=!\,\lstinline!default! on the first declaration
denotes that a special member function is intended to be generated by
the compiler --- irrespective of any user-provided declarations; in
conjunction with {\lstinline!=!\,\lstinline!delete!} (see \featureref{\locationa}{deleted-functions}),
%Section~\ref{deleted-functions}, ``\titleref{deleted-functions}"),
using
\lstinline!=!\,\lstinline!default! affords the fine-grained control over which
special member functions are to be generated and/or made publicly
available.

\subsubsection[Defaulting the implementation of a user-provided special member function]{Defaulting the implementation of a user-provided special member function}\label{defaulting-the-implementation-of-a-user-provided-special-member-function}

The \lstinline!=!\,\lstinline!default! syntax can also be used after the first
declaration, but with a distinctly different meaning: The compiler will
treat the first declaration as a \emcppsgloss[user provided special member
function]{user-provided special member
function} and thus will suppress the generation of other
\emcppsgloss[special member function]{special member functions} accordingly.

\begin{emcppslisting}[language=C++,label=default-exampleh-code]
// example.h:

struct S6
{
    S6& operator=(const S6&);  // user-provided copy-assignment operator

    // suppresses the declaration of both move operations
    // implicitly generates the default and copy constructors, and destructor
};

inline S6& S6::operator=(const S6&) = default;
    // Explicitly request the compiler to generate the default implementation
    // for this copy-assignment operator. This request might fail (e.g., if (ù{\codeincomments{S6}}ù)
    // were to contain a non-copy-assignable data member).
\end{emcppslisting}

\noindent Alternatively, an explicitly defaulted noninline implementation of this
copy-assignment operator may appear in a separate (\lstinline!.cpp!) file;
see \intraref{default-use-cases}{physically-decoupling-the-interface-from-the-implementation}.
%{\it\titleref{default-use-cases}: \titleref{physically-decoupling-the-interface-from-the-implementation}} on page~\pageref{physically-decoupling-the-interface-from-the-implementation}.

\subsection[Use Cases]{Use Cases}\label{default-use-cases}

\subsubsection[Restoring the generation of a special member function suppressed by another]{Restoring the generation of a special member function suppressed by another}\label{restoring-the-generation-of-a-special-member-function-suppressed-by-another}

Incorporating \lstinline!=!\,\lstinline!default! in the declaration of a
special member function instructs the compiler to generate its
definition regardless of any other user-provided special member
functions. As an example, consider a \emcppsgloss[value semantic]{value-semantic}
\lstinline!SecureToken! class that wraps a standard string
(\lstinline!std::string!) and an arbitrary-precision-integer
(\lstinline!BigInt!) token code that together satisfy certain invariants:

\begin{emcppshiddenlisting}[emcppsbatch={e1,e2,e3}]
#include <string>   // (ù{\codeincomments{std::string}}ù)
#include <cassert>  // standard C (ù{\codeincomments{assert}}ù) macro
struct BigInt {};
bool operator==(const BigInt&,const BigInt&);
\end{emcppshiddenlisting}
\begin{emcppslisting}[emcppsbatch=e1]
class SecureToken
{
    std::string d_value;  // The default-constructed value is the empty string.
    BigInt      d_code;   // The default-constructed value is the integer zero.

public:
    // All six special member functions are (implicitly) defaulted.

    void setValue(const char* value);
    const char* value() const;
    BigInt code() const;
};
\end{emcppslisting}

\noindent By default, a secure token's \lstinline!value! will be the empty-string
value, and the token's \lstinline!code! will be the numerical value of zero
because those are, respectively, the \emcppsgloss[default initialized]{default-initialized} values
of the two data members, \lstinline!d_value! and \lstinline!d_code!:

\begin{emcppslisting}[emcppsbatch=e1,label=default-voidf-code]
void f()
{
    SecureToken token;                       // default constructed         (1)
    assert(token.value() == std::string());  // default value: empty string (2)
    assert(token.code() == BigInt());        // default value: zero         (3)
}
\end{emcppslisting}

\noindent Now suppose that we get a request to add a \emcppsgloss{value constructor}
that creates and initializes a \lstinline!SecureToken! from a specified
token string:

\begin{emcppslisting}[emcppsbatch=e2]
class SecureToken
{
    std::string d_value;  // The default-constructed value is the empty string.
    BigInt      d_code;   // The default-constructed value is the integer zero.

public:
    SecureToken(const char* value);  // newly added value constructor

    // suppresses the declaration of just the default constructor --- i.e.,
    // implicitly generates all of the other five special member functions

    void setValue(const char* value);
    const char* value() const;
    const BigInt& code() const;
};
\end{emcppslisting}

\noindent Attempting to compile function \lstinline!f! (from page~\pageref{default-voidf-code}) would now fail on the
first line, where it attempts to default-construct the token. Using
the \lstinline!=!\,\lstinline!default! feature, however, we can reinstate the
default constructor to work trivially, just as
it did before:

\begin{emcppslisting}[emcppsbatch=e3]
class SecureToken
{
    std::string d_value;  // The default-constructed value is the empty string.
    BigInt d_code;        // The default-constructed value is the integer zero.

public:
    SecureToken() = default;         // newly defaulted default constructor
    SecureToken(const char *value);  // newly added value constructor

    // implicitly generates all of the other five special member functions

    void setValue(const char *value);
    const char *value() const;
    const BigInt& code() const;
};
\end{emcppslisting}


\subsubsection[Making class APIs explicit at no runtime cost]{Making class APIs explicit at no runtime cost}\label{making-class-apis-explicit-at-no-runtime-cost}

In the early days of C++, coding standards sometimes
required that each special member function be declared explicitly so
that it could be documented or even just to know that it hadn't been
forgotten:

\begin{emcppslisting}
class C1
{
     // ...

public:
   C1();
       // Create an empty object.

   C1(const C1& rhs);
       // Create an object having the same value as the specified (ù{\codeincomments{rhs}}ù) object.

   ~C1();
       // Destroy this object.

   C1& operator=(const C1& rhs);
       // Assign to this object the value of the specified (ù{\codeincomments{rhs}}ù) object.
};
\end{emcppslisting}

\noindent Over time, explicitly writing out what the compiler
could do more reliably itself became more clearly an inefficient use of developer time and a maintenance burden.
What's more, even if the function definition was empty, implementing it
explicitly often degraded performance compared to a \emcppsgloss{trivial} default. Hence, such
standards tended to evolve toward conventionally commenting out (e.g.,
using \lstinline|//!|)
the declarations of functions having an empty
body rather than providing it explicitly:

\begin{emcppslisting}
class C2
{
     // ...

public:
   //! C2();
       // Create an empty object.

   //! C2(const C2& rhs);
       // Create an object having the same value as the specified (ù{\codeincomments{rhs}}ù) object.

   //! ~C2();
       // Destroy this object.

   //! C2& operator=(const C2& rhs);
       // Assign to this object the value of the specified (ù{\codeincomments{rhs}}ù) object.
};
\end{emcppslisting}

\noindent Note, however, that the compiler does not check the commented code, which is easily susceptible to copy-paste and other errors. By
uncommenting the code and defaulting it explicitly in class scope, we
regain the compiler's syntactic checking of the function signatures
without incurring the cost of turning what would have been
\emcppsgloss{trivial} functions into equivalent
non-\emcppsgloss{trivial} ones:

\begin{emcppslisting}
class C3
{
     // ...

public:
   C3() = default;
       // Create an empty object.

   C3(const C3& rhs) = default;
       // Create an object having the same value as the specified (ù{\codeincomments{rhs}}ù) object.

   ~C3() = default;
       // Destroy this object.

   C3& operator=(const C3& rhs) = default;
       // Assign to this object the value of the specified (ù{\codeincomments{rhs}}ù) object.
};
\end{emcppslisting}


%\subsubsection[Preserving trivial copyability]{Preserving trivial copyability}\label{preserving-trivial-copyability}
%
%In some situations, a particular type \textit{must} be
%usable with \texttt{std::memcpy} (e.g., runtime performance,
%serialization to binary, or interoperability with C code). Only
%\textbf{trivially copyable} types are safe to use with
%\texttt{std::memcpy}; use with any other types results in
%\textbf{undefined behavior}. A type \texttt{T} is \textbf{trivially
%copyable} if it exposes a \textbf{trivial copy constructor}:
%\begin{enumerate}
%\item{the copy constructor for \texttt{T} is not user provided}
%\item{the type \texttt{T} itself has no \texttt{virtual} member functions or
%\texttt{virtual} base classes}
%\item{any member or base class of
%\texttt{T} is itself \textbf{trivially copyable} (recursively).}
%\end{enumerate}
%
%As an example, the \texttt{EntityHandle} class (in the code snippet below) represents an
%integer handle (to an entity of opaque type) that must be usable with
%\texttt{std::memcpy} for the purpose of efficient serialization (the
%capacity of the encapsulated fundamental integral type is subject to
%change){\cprotect\footnote{Objects of this type are sometimes said to
%  hold ``dumb data''; see \textbf{{lakos20}}, section~3.5.5,
%  pp.~629--633.}}:
%
%\begin{emcppslisting}
%class EntityHandle
%{
%    short int d_id;  // Note: Implementation size may increase over time.
%
%public:
%    EntityHandle(int id);  // value constructor
%
%    // suppresses the declaration of just the default constructor --- i.e.,
%    // implicitly generates all of the other five special member functions
%
%    // ...
%}
%\end{emcppslisting}
%
%\noindent The presence of any other constructor, except a \emph{move constructor},
%never affects implicit generation of a copy constructor, and
%\texttt{short}~\texttt{int} (like all \emph{enumerated}, \emph{pointer},
%and other \emph{fundamental} types) is a \textbf{trivial type}, thus
%establishing the \emph{triviality} of copying an \texttt{EntityHandle}.
%Now imagine that, to monitor the places in the codebase where
%\emph{temporary} entity handles are exchanged (with the goal of
%ultimately optimizing those), a user-provided \emph{move constructor} is
%added{\cprotect\footnote{Note that a move constructor will be preferred
%over a copy constructor when the type category of the argument is an
%\emph{xvalue} (i.e., expiring value) or \emph{prvalue} (i.e., pure rvalue), which
%are the value categories to which a temporary can pertain. See Section~\ref{Rvalue-References}, ``\titleref{Rvalue-References}," for more information.}}:
%
%\begin{emcppslisting}
%class EntityHandle
%{
%    short int d_id;  // Note: Implementation size may increase over time.
%
%public:
%    EntityHandle(int id);              // value constructor
%    EntityHandle(EntityHandle&& rhs);  // user-provided move constructor
%
%    // implicitly generates only the destructor
%    // suppresses synthesis of all four other special member functions
%
%    // ...
%};
%\end{emcppslisting}
%
%\noindent As illustrated by Table~\ref{default-table1} on page~\pageref{default-table1}, the
%presence of a user-provided \emph{move constructor} suppressed the
%automatic generation of a copy constructor along with the destructor
%and both the copy- and move-assignment operators, thereby rendering the
%\texttt{EntityHandle} unusable. Replacing these four previously
%generated functions with seemingly equivalent user-provided ones might
%appear to work as intended:
%
%\begin{emcppslisting}
%class EntityHandle
%{
%    short int d_id;  // Note: Implementation size may increase over time.
%
%public:
%    EntityHandle(int id);  // value constructor
%
%    EntityHandle(const EntityHandle& rhs);  // user-provided copy constructor
%    EntityHandle(EntityHandle&& rhs);       // user-provided move constructor
%
%    EntityHandle& operator=(const EntityHandle& rhs);
%        // user-provided copy-assignment operator
%
%    EntityHandle& operator=(EntityHandle&& rhs);
%        // user-provided move-assignment operator
%
%    // implicitly generates only the destructor
%    // suppresses synthesis of the default constructor
%
%    // ...
%};
%\end{emcppslisting}
%
%\noindent The user-provided nature of the copy constructor, however, renders the
%\texttt{EntityHandle} class ineligible for copy triviality --- even if
%the definitions are identical! Hence, any direct use of
%\texttt{std::memcpy} with an \texttt{EntityHandle} object will result in
%\textbf{undefined behavior}. We could have instead explicitly requested
%that these four special member functions be generated using
%\texttt{=}~\texttt{default}:
%
%\begin{emcppslisting}
%class EntityHandle
%{
%    short int d_id;  // Note: Implementation size may increase over time.
%
%public:
%    EntityHandle(int id);  // value constructor
%
%    EntityHandle(const EntityHandle& rhs) = default;
%        // defaulted (trivial) copy constructor
%
%    EntityHandle(EntityHandle&& rhs);
%        // user-provided move constructor
%
%    EntityHandle& operator=(const EntityHandle& rhs) = default;
%        // default (trivial) copy-assignment operator
%
%    EntityHandle& operator=(EntityHandle&& rhs);
%        // user-provided move-assignment operator
%
%    // Implicitly generates only the destructor.
%    // suppresses synthesis of the default constructor
%
%    // ...
%};
%\end{emcppslisting}
%
%\noindent By explicitly defaulting these three special member functions {in class
%scope}, we (1) re-enable their generation and (2) preserve the
%\emph{copy triviality} of the class.


%%%% new material from Slava
\subsubsection{Preserving type triviality}

It can be beneficial if a particular type is \emcppsgloss{trivial}.  The type is considered \emcppsgloss{trivial} if its default constructor is \emcppsgloss{trivial} and it is \emcppsgloss{trivially copyable} --- i.e., it has no non-trivial copy or move constructors, no non-trivial copy or move assignment operators, at least one of those nondeleted, and a trivial destructor.  As an example, consider a simple \emcppsgloss{trivial}  \lstinline!Metrics! type in the code snippet below containing certain collected metrics for our application:

\begin{emcppslisting}
struct Metrics
{
    int d_numRequests;  // number of requests to the service
    int d_numErrors;    // number of error responses

    // All special member functions are generated implicitly.
};
\end{emcppslisting}

\noindent Now imagine that we would like to add a constructor to this struct to make its use more convenient:\newpage%%%%%%

\begin{emcppslisting}
struct Metrics
{
    int d_numRequests;  // number of requests to the service
    int d_numErrors;    // number of error responses

    Metrics(int, int);  // user-provided value constructor

    // Generation of default constructor is suppressed.
};
\end{emcppslisting}

\noindent As illustrated in \intrarefsimple{appendix:-implicit-generation-of-special-member-functions}, the presence of a user-provided constructor suppressed the implicit generation of the default constructor.  Replacing the default constructor with a seemingly equivalent user-provided one might appear to work as intended:

\begin{emcppslisting}
struct Metrics
{
    int d_numRequests;  // number of requests to the service
    int d_numErrors;    // number of error responses

    Metrics(int, int);  // user-provided value constructor
    Metrics() {}        // user-provided default constructor

    // Default constructor is user-provided: (ù{\codeincomments{Metrics}ù) is (ù{\emphincomments{not}ù) trivial.
};
\end{emcppslisting}

\noindent The user-provided nature of the default constructor, however, renders the \lstinline!Metrics! type non-trivial --- even if the definitions are identical!  In contrast, explicitly requesting the default constructor be generated using \lstinline!=!\,\lstinline!default! restores the triviality of the type:

\begin{emcppslisting}
struct Metrics
{
    int d_numRequests;  // number of requests to the service
    int d_numErrors;    // number of error responses

    Metrics(int, int);    // user-provided value constructor
    Metrics() = default;  // defaulted (trivial) default constructor

    // Default constructor is defaulted: (ù{\codeincomments{Metrics}ù) (ù{\emphincomments{is}ù) trivial.
};
\end{emcppslisting}
%%%%%%%%%%%%% new material ends

\subsubsection[Physically decoupling the interface from the implementation]{Physically decoupling the interface from the implementation}\label{physically-decoupling-the-interface-from-the-implementation}

Sometimes, especially during large-scale development, avoiding compile-time coupling clients to the
implementations of individual methods offers distinct maintenance advantages.\linebreak\newpage%%%%%%
Specifying that a special member
function is defaulted on its first declaration (i.e., in class scope)
implies that making any change to this implementation will force all
clients to recompile:

\begin{emcppslisting}
// smallscale.h:

struct SmallScale
{
    SmallScale() = default;  // explicitly defaulted default constructor
};
\end{emcppslisting}

\noindent The important issue regarding recompilation here is not merely compile time per se but compile-time coupling{\cprotect\footnote{See
  \cite{lakos20}, section 3.10.5, pp.~783--789.}}

Alternatively, we can choose to declare the function but deliberately
\emph{not} default it in class scope (or anywhere in the \lstinline!.h!
file):

\begin{emcppslisting}[emcppsbatch=e4]
// largescale.h:

struct LargeScale
{
    LargeScale();  // user-provided default constructor
};
\end{emcppslisting}

\noindent We can then default just the non-inline implementation in a
corresponding{\cprotect\footnote{In practice, every \lstinline!.cpp! file
(other than the one containing \lstinline!main!) typically has a unique
associated header (\lstinline!.h!) file and often vice versa  (with the \lstinline!.cpp! and \lstinline!.h! pair of files constituting a component); see \cite{lakos20}, sections 1.6 and 1.11,
  pages 209--216 and 256--259, respectively.}}~\lstinline!.cpp! file:

\begin{emcppslisting}[emcppsbatch=e4]
// largescale.cpp:
#include <largescale.h>

LargeScale::LargeScale() = default;
    // Generate the default implementation for this default destructor.
\end{emcppslisting}

\noindent Using this \emph{insulation} technique, we are free to change our minds
and implement the default constructor ourselves in any way we see fit
without necessarily forcing our clients to recompile.

\subsection[Potential Pitfalls]{Potential Pitfalls}\label{potential-pitfalls}

%%%% REMEMBER TO NOTE THIS CHANGE ON PDF SO PROOFREADER KNOWS
%%%%%%%%%%%%%%%% Moved from Annoyances to Pitfalls
\subsubsection[Defaulted special member functions cannot restore trivial copyability]{Defaulted special member functions cannot restore trivial copyability}

Library classes often rely on whether the type on which they are operating is eligible for being copied with \lstinline!memcpy! for optimization purposes.  Such could be the case for implementing, \pagebreak%%%%% 
say, \lstinline!vector!, which would make a single call to \lstinline!memcpy! when growing its buffer.  
For the \lstinline!memcpy! or \lstinline!memmove! to be well-defined, however, the type of the object that is stored in the buffer must be \emcppsgloss{trivially copyable}.  One might assume that this trait means that, as long as the copy constructor of the type is trivial, this optimization will apply.  Defaulting the copy operations would then allow us to achieve this goal, while allowing the type to have a non-trivial destructor or move operation.  Such, however, is not the case.

The requirements for a type to be considered \emcppsgloss{trivially copyable} --- and thus eligible for use with \lstinline!memcpy! --- include triviality of all of its nondeleted copy and move operations as well as of its destructor.  Furthermore, library authors cannot perform fine-grained dispatch based on which operations on the type are in fact trivial.  Even if we detect that the type is trivially copy-constructible with the \lstinline!std::is_trivially_copy_constructible! trait and know that our code would use only copy constructors (and not copy assignment nor any move operations), we still would not be able to use \lstinline!memcpy! unless the more restrictive \lstinline!std::is_trivially_copyable! trait is also \lstinline!true!.
%%%%%%%%%%%%%%%%

\subsection{Annoyances}

%%%%%%%%%%%%%%%%%%%% Moved from Pitfalls to Annoyances
\subsubsection[Generation of defaulted functions is not guaranteed]{Generation of defaulted functions is not guaranteed}\label{generation-of-defaulted-functions-is-not-guaranteed}

Using \lstinline!=!\,\lstinline!default! does not guarantee that the special
member function of a type, \lstinline!T!, will be generated. For example, a
noncopyable member variable (or base class) of \lstinline!T! will inhibit
generation of \lstinline!T!'s copy constructor even when
\lstinline!=!\,\lstinline!default! is used. Such behavior can be observed in
the presence of a \lstinline!std::unique_ptr!{\cprotect\footnote{\lstinline!std::unique_ptr<T>!
is a move-only (movable but noncopyable) class template introduced in
C++11. It models unique ownership over a dynamically allocated
  \lstinline!T! instance, leveraging rvalue references (see \featureref{\locationc}{Rvalue-References})
%  Section~\ref{Rvalue-References}, ``\titleref{Rvalue-References}")
  to represent ownership transfer between instances:
  %%% question

  \begin{emcppshiddenlisting}[emcppsbatch=e5,style=footcode]
  #include <memory>  // (ù{\codeincomments{std::unique\_ptr}}ù)
  \end{emcppshiddenlisting}
  \begin{emcppslisting}[emcppsbatch=e5,basicstyle={\ttfamily\footnotesize}]
  int* p = new int(42);
  std::unique_ptr<int> up(p);                   // OK, take ownership of (ù{\codeincomments{p}}ù).
  std::unique_ptr<int> upCopy = up;             // Error, copy is deleted
  std::unique_ptr<int> upMove = std::move(up);  // OK, transfer ownership.
  \end{emcppslisting}\vspace*{-1ex}
      }} data member:

\begin{emcppslisting}
#include <memory>  // (ù{\codeincomments{std::unique\_ptr}}ù)
class Connection
{
private:
    class Impl;                    // nested implementation class
    std::unique_ptr<Impl> d_impl;  // noncopyable data member

public:
    Connection() = default;
    Connection(const Connection&) = default;
};
\end{emcppslisting}
\pagebreak%%%%%%%

\noindent Despite the defaulted copy constructor, \lstinline!Connection! will not be
copy-constructible as\linebreak[4] \lstinline!std::unique_ptr! is a noncopyable type.
Some compilers \emph{may} produce a warning on the declaration of
{\lstinline!Connection(const! \lstinline!Connection&)!, but they are not
required to do so since the example code above is well formed and would produce a
compilation failure only if an attempt were made to default-construct or
copy a \lstinline!Connection!.{\cprotect\footnote{Clang 8.x
and later produces a diagnostic with no warning flags specified.
MSVC produces a diagnostic if \lstinline!/Wall! is specified. As of this writing, GCC produces no warning, even with both
  \lstinline!-Wall! and \lstinline!-Wextra! enabled.}}

If desired, a possible way to ensure that a defaulted special member
function has indeed been generated is to use
{\lstinline!static_assert!} (see \featureref{\locationa}{compile-time-assertions-(static_assert)})
%Section~\ref{compile-time-assertions-(static_assert)}, ``\titleref{compile-time-assertions-(static_assert)}")
in conjunction with an
appropriate trait from the \lstinline!<type_traits>! header:

\begin{emcppshiddenlisting}[emcppsbatch=e6]
#include <vector>  // (ù{\codeincomments{std::vector}}ù)
\end{emcppshiddenlisting}
\begin{emcppslisting}[emcppsbatch=e6]
class IdCollection
{
    std::vector<int> d_ids;

public:
    IdCollection() = default;
    IdCollection(const IdCollection&) = default;
    // ...
};

static_assert(std::is_default_constructible<IdCollection>::value,
              "IdCollection must be default constructible.");

static_assert(std::is_copy_constructible<IdCollection>::value,
              "IdCollection must be copy constructible.");

// ...
\end{emcppslisting}

\noindent Routinely using such compile-time testing techniques can help to ensure
that a type will continue to behave as expected (at no additional
runtime cost) even when member and base types evolve as a result
of ongoing software maintenance.
%%%%%%%%%%%%%%%%%%%%%%%

\subsection[See Also]{See Also}\label{see-also}

\begin{itemize}
\item{\seealsoref{deleted-functions}{\seealsolocationa}describes a companion feature, \lstinline!=!\,\lstinline!delete!, that can be used to suppress access to implicitly generated \emcppsgloss[special member function]{special member functions}.}
\item{\seealsoref{compile-time-assertions-(static_assert)}{\seealsolocationa}describes a facility that can be used to verify at compile time that undesirable copy and move operations are declared to be accessibly.}
\item{\seealsoref{Rvalue-References}{\seealsolocationc}provides the bases for \emcppsgloss{move operations}, namely, the move-constructor and move-assignment \emcppsgloss[special member function]{special member functions}, which too can be defaulted.}
\end{itemize}

\subsection[Further Reading]{Further Reading}\label{further-reading}

\begin{itemize}
\item{Howard Hinnant, ``Everything You Ever Wanted to Know About Move Semantics (and Then Some)," \cite{hinnant14}}
\item{Howard Hinnant, ``Everything You Ever Wanted to Know About Move Semantics," \cite{hinnant16}}
\end{itemize}

\subsection[Appendix: Implicit Generation of Special Member Functions]{Appendix: Implicit Generation of Special Member Functions}\label{appendix:-implicit-generation-of-special-member-functions}

The rules a compiler uses to decide if a special member function
should be generated implicitly are not entirely intuitive. Howard
Hinnant, lead designer and author of the C++11 proposal for move
semantics\footnote{\cite{hinnant02}} (among other proposals), produced a tabular representation\footnote{\cite{hinnant16}} of
such rules in the situation where the user provides a single special
member function and leaves the rest to the compiler. To understand Table~\ref{default-table1}, after picking a special member function in the first
column, the corresponding row will show what is implicitly generated by
the compiler.

\begin{table}[h!]  %%% PRODUCTION: REQUIRES REPRINT PERMISSION
\begin{center}
\begin{threeparttable}
\caption{Implicit generation of special member functions}\label{default-table1}
{\footnotesize %%% setting smaller to fit
\begin{tabular}{p{.127\textwidth}|p{.1\textwidth}|p{.11\textwidth}|p{.1\textwidth}|p{.127\textwidth}|p{.1\textwidth}|p{.127\textwidth}}
\thickhline
% header
\rowcolor[gray]{.9}    & {\sffamily\bfseries Default\linebreak[4] Ctor} & {\sffamily\bfseries Destructor} & {\sffamily\bfseries Copy\linebreak[4] Ctor} & {\sffamily\bfseries Copy\linebreak[4] Assignment} & {\sffamily\bfseries Move\linebreak[4] Ctor} & {\sffamily\bfseries Move\linebreak[4] Assignment} \tabularnewline \hline
 % row 1
\cellcolor[gray]{.9}{\sffamily\bfseries Nothing} & Defaulted & Defaulted & Defaulted & Defaulted
& Defaulted & Defaulted \\ \hline
 % row 2
 \cellcolor[gray]{.9}{\sffamily\bfseries Any\linebreak[4] Ctor} & Not\linebreak[4] Declared &
Defaulted & Defaulted & Defaulted & Defaulted & Defaulted \\ \hline
 % row 3
 \cellcolor[gray]{.9}{\sffamily\bfseries Default\linebreak[4] Ctor} & User\linebreak[4] Declared & Defaulted & Defaulted &
Defaulted & Defaulted & Defaulted \\ \hline
 % row 4
 \cellcolor[gray]{.9}{\sffamily\bfseries Destructor} &
Defaulted & User\linebreak[4] Declared & Defaulted\tnote{a} & Defaulted\tnote{a} & Not\linebreak[4] Declared &
Not\linebreak[4] Declared \tabularnewline \hline
 % row 5
 \cellcolor[gray]{.9}{\sffamily\bfseries Copy\linebreak[4] Ctor} & Not\linebreak[4] Declared & Defaulted &
User\linebreak[4] Declared & Defaulted\tnote{a} & Not\linebreak[4] Declared & Not\linebreak[4] Declared \tabularnewline \hline
 % row 6
 \cellcolor[gray]{.9}{\sffamily\bfseries Copy\linebreak[4] Assignment} & Defaulted & Defaulted & Defaulted\tnote{a} & User\linebreak[4]
Declared & Not\linebreak[4] Declared & Not\linebreak[4] Declared \tabularnewline \hline
 % row 7
 \cellcolor[gray]{.9}{\sffamily\bfseries Move\linebreak[4] Ctor} &
Not\linebreak[4] Declared & Defaulted & Deleted & Deleted & User\linebreak[4] Declared & Not\linebreak[4]
Declared \tabularnewline \hline
 % row 8
 \cellcolor[gray]{.9}{\sffamily\bfseries Move\linebreak[4] Assignment} & Defaulted & Defaulted & Deleted
& Deleted & Not\linebreak[4] Declared & User\linebreak[4] Declared \tabularnewline
\thickhline
\end{tabular}
} % end small
\begin{tablenotes}{\footnotesize
\item[a]{Deprecated behavior: compilers might warn upon reliance of this
implicitly generated member \mbox{function}.}
} % end footnotesize
\end{tablenotes} % we end the table footnotes
\end{threeparttable}
\end{center}
\end{table}

As an example, explicitly declaring a copy-assignment operator
would result in the default constructor, destructor, and
copy constructor being defaulted and in the move
operations not being declared. If more than one \emcppsgloss[special member function]{special member
function} is user declared (regardless of whether or how it is
implemented), the remaining generated member functions are those in the
intersection of the corresponding rows. For example, explicitly
declaring both the~destructor and the default constructor would still
result in the copy constructor and the copy-assignment operator being
defaulted and both move operations not being declared. Relying on the
compiler-generated copy operations when the destructor is anything but
defaulted is dubious; if correct, defaulting them explicitly makes both
their existence and intended definition clear.





\newpage
%\section[Delegating Ctors]{Constructors Calling Other Constructors}\label{delegating-constructors}
% 5 Feb 2021, copyedits in and proofed
% 8 Feb 2021 JMB, compilation fixes
% 9 Feb 2021, FPPs set
% 17 Feb 2021, glossary tags updated

\emcppsFeature{
    short={Delegating Ctors},
    long={Constructors Calling Other Constructors},
}{delegating-constructors}
\setcounter{table}{0}
\setcounter{footnote}{0}
\setcounter{lstlisting}{0}
%\section[Delegating Ctors]{Constructors Calling Other Constructors}\label{delegating-constructors}

Delegating constructors are constructors of a class that delegate initialization to another
constructor of the same class.

\subsection[Description]{Description}\label{description}

A \emcppsgloss{delegating constructor} is a constructor of a
\emcppsgloss[user defined type (UDT)]{user-defined type (UDT)} --- i.e., \lstinline!class!, \lstinline!struct!, or
\lstinline!union! --- that invokes another constructor defined for the same
\emcppsgloss[user defined type (UDT)]{UDT} as part of its initialization of an object of that type.
The syntax for invoking another constructor is to specify
the name of the type as the only element in the \emcppsgloss{member initializer list}:

%\begin{emcppslisting}
%struct S0
%{
%    int d_i;
%
%    S0(int i) : d_i(i)        { }  // non-delegating constructor
%    S0()      : S0(0)         { }  // OK, delegates to (ù{\codeincomments{S0(int)}}ù)
%    S0(bool)  : S0(0), d_i(0) { }  // error: delegation must be on its own
%};
%\end{emcppslisting}
\begin{emcppslisting}
#include <string>  // (ù{\codeincomments{std::string}}ù)

struct S0
{
  int         d_i;
  std::string d_s;

  S0(int i)         : d_i(i)        {} // nondelegating constructor
  S0()              : S0(0)         {} // OK, delegates to (ù{\codeincomments{S0(int)}}ù)
  S0(const char *s) : S0(0), d_s(s) {} // Error, delegation must be on its own
};
\end{emcppslisting}

\noindent Multiple delegating constructors can be chained together (one calling
exactly one other) so long as cycles are avoided (see \intraref{ctordelegrating-potential-pitfalls}{delegation-cycles}).
%{\it\titleref{ctordelegrating-potential-pitfalls}:} {\it\titleref{delegation-cycles}} on page~\pageref{delegation-cycles}).
Once a \emph{target} (i.e., invoked via delegation) constructor returns,
the body of the delegator is invoked:

\begin{emcppslisting}
#include <iostream>  // (ù{\codeincomments{std::cout}}ù)

struct S1
{
    S1(int, int)            { std::cout << 'a'; }
    S1(int)      : S1(0, 0) { std::cout << 'b'; }
    S1()         : S1(0)    { std::cout << 'c'; }
};

void f()
{
    S1 s;  // OK, prints (ù{\codeincomments{"abc"}}ù) to (ù{\codeincomments{stdout}}ù)
}
\end{emcppslisting}
\newpage%%%%%

\noindent If an exception is thrown while executing a nondelegating constructor,
the object being initialized is considered only \emcppsgloss{partially
constructed} (i.e., the object is not yet known to be in a valid state),
and hence its destructor will \emph{not} be
invoked:

\begin{emcppslisting}
#include <iostream>  // (ù{\codeincomments{std::cout}}ù)

struct S2
{
    S2()  { std::cout << "S2() ";  throw 0; }
    ~S2() { std::cout << "~S2() ";          }
};

void f() try { S2 s; } catch(int) { }
    // prints only "(ù{\codeincomments{S2() }}ù)" to (ù{\codeincomments{stdout}}ù) (the destructor of (ù{\codeincomments{S2}}ù) is never invoked)
\end{emcppslisting}

\noindent Although the destructor of a \emcppsgloss{partially
constructed} object will not be invoked, the destructors of
each successfully constructed base and of data members will still be
invoked:

\begin{emcppslisting}
#include <iostream>  // (ù{\codeincomments{std::string}}ù)

using std::cout;
struct A { A() { cout << "A() "; } ~A() { cout << "~A() "; } };
struct B { B() { cout << "B() "; } ~B() { cout << "~B() "; } };

struct C : B
{
    A d_a;

    C()  { cout << "C() "; throw 0; }  // nondelegating constructor that throws
    ~C() { cout << "~C() ";         }  // destructor that never gets called
};

void f() try { C c; } catch(int) { }
    // prints (ù{\codeincomments{"B() A() C() \textasciitilde A() \textasciitilde B()"}}ù) to (ù{\codeincomments{stdout}}ù)
\end{emcppslisting}

\noindent Notice that base-class \lstinline!B! and member \lstinline!d_a! of type
\lstinline!A! were fully constructed, and so their respective destructors
are called, even though the destructor for class \lstinline!C! itself is never \mbox{executed}. 

However, if an exception is thrown in the body of a delegating
constructor, the object being initialized is considered \emcppsgloss{fully
constructed}, as the target constructor must have returned control to
the delegator; hence, the object's destructor \emph{is} invoked:
\newpage%%%%%

\begin{emcppslisting}
#include <iostream>  // (ù{\codeincomments{std::cout}}ù)

struct S3
{
    S3()           { std::cout << "S3() ";             }
    S3(int) : S3() { std::cout << "S3(int) "; throw 0; }
    ~S3()          { std::cout << "~S3() ";            }
};

void f() try { S3 s(0); } catch(int) { }
    // prints "(ù{\codeincomments{S3() S3(int) \texttt{\~{}}S3() }}ù)" to (ù{\codeincomments{stdout}}ù)
\end{emcppslisting}


\subsection[Use Cases]{Use Cases}\label{ctordelegating-use-cases}

\subsubsection[Avoiding code duplication among constructors]{Avoiding code duplication among constructors}\label{avoiding-code-duplication-among-constructors}

%Avoiding gratuitous code duplication is considered by many to be a best practice. 
Many consider avoiding gratuitous code duplication a best
practice. Having one ordinary member function call another has always
been an option, but having one constructor invoke another constructor
directly has not. Classic workarounds included repeating the code or
else factoring the code into a private member function that would be
called from multiple constructors. The drawback with this workaround is
that the private member function, not being a constructor, would be unable to
make use of \emcppsgloss[member initializer list]{member initializer lists} to initialize base classes and data members efficiently. As of C++11, \emph{delegating
constructors} can be used to minimize code duplication when some of
the same operations are performed across multiple constructors without
having to forgo efficient initialization. 
%As an example, 
Consider an \lstinline!IPV4Host! class representing a
network endpoint that can be constructed either (1) by a 32-bit address
and a 16-bit port or (2) by an IPV4 string with
\lstinline!XXX.XXX.XXX.XXX:XXXXX! format{\cprotect\footnote{Note that
this initial design might itself be suboptimal in that the
representation of the IPV4 address and port value might profitably be
factored out into a separate \emcppsgloss[value semantic]{value-semantic} class, say,
\lstinline!IPV4Address!, that itself might be constructed in multiple ways;
  see \intraref{ctordelegrating-potential-pitfalls}{suboptimal-factoring}.}}:
%  {\it\titleref{ctordelegrating-potential-pitfalls}:} {\it\titleref{suboptimal-factoring}} on page~\pageref{suboptimal-factoring}.}}:

%\enlargethispage*{2ex}%%%%% to fit code
\begin{emcppshiddenlisting}[emcppsbatch={e1,e2,e3,e4}]
#include <cstdint> // (ù{\codeincomments{std::uint16\_t}}ù), (ù{\codeincomments{std::uint32\_t}}ù)
#include <string>   // (ù{\codeincomments{std::string}}ù)
struct ConnectionException {
    ConnectionException(std::uint32_t address, std::uint16_t port);
};
std::uint32_t extractAddress(const std::string &);
std::uint16_t extractPort(const std::string &);
\end{emcppshiddenlisting}
\begin{emcppslisting}[emcppsbatch=e1]
#include <cstdint>  // (ù{\codeincomments{std::uint16\_t}}ù), (ù{\codeincomments{std::uint32\_t}}ù)
#include <string>   // (ù{\codeincomments{std::string}}ù)

class IPV4Host
{
     // ...
private:
    int connect(std::uint32_t address, std::uint16_t port);

public:
    IPV4Host(std::uint32_t address, std::uint16_t port)
    {
        if (!connect(address, port))  // code duplication: BAD IDEA
        {
            throw ConnectionException{address, port};
        }
    }

    IPV4Host(const std::string& ip)
    {
        std::uint32_t address = extractAddress(ip);
        std::uint16_t port = extractPort(ip);

        if (!connect(address, port))  // code duplication: BAD IDEA
        {
            throw ConnectionException{address, port};
        }
    }
};
\end{emcppslisting}

\noindent Prior to C++11, working around such code duplication would require the
introduction of a separate, private helper function that
would be called by each of the constructors:

\enlargethispage*{2ex}%%%% to fit the code
\begin{emcppslisting}[emcppsbatch=e2]
// C++03 (obsolete)
#include <cstdint>  // (ù{\codeincomments{std::uint16\_t}}ù), (ù{\codeincomments{std::uint32\_t}}ù)

class IPV4Host
{
    // ...

private:
    int connect(std::uint32_t address, std::uint16_t port);
    void init(std::uint32_t address, std::uint16_t port)  // helper function
    {
        if (!connect(address, port))  // factored implementation of needed logic
        {
            throw ConnectionException{address, port};
        }
    }

public:
    IPV4Host(std::uint32_t address, std::uint16_t port)
    {
        init(address, port);  // Invoke factored private helper function.
    }

    IPV4Host(const std::string& ip)
    {
        std::uint32_t address = extractAddress(ip);
        std::uint16_t port = extractPort(ip);

        init(address, port);  // Invoke factored private helper function.
    }
};
\end{emcppslisting}

\noindent With C++11 delegating constructors, the constructor accepting a string can be rewritten to
delegate to the one accepting \lstinline!address! and \lstinline!port!,
avoiding repetition without having to use a private function:

\begin{emcppslisting}[emcppsbatch=e3]
#include <cstdint>  // (ù{\codeincomments{std::uint16\_t}}ù), (ù{\codeincomments{std::uint32\_t}}ù)
#include <string>   // (ù{\codeincomments{std::string}}ù)

class IPV4Host
{
     // ...
private:
    int connect(std::uint32_t address, std::uint16_t port);

public:
    IPV4Host(std::uint32_t address, std::uint16_t port)
    {
        if(!connect(address, port))
        {
            throw ConnectionException{address, port};
        }
    }

    IPV4Host(const std::string& ip)
        : IPV4Host{extractAddress(ip), extractPort(ip)}
    {
    }
};
\end{emcppslisting}

\noindent Using delegating constructors results in less boilerplate and fewer runtime
operations, as data members and base classes can be initialized
directly through the \emcppsgloss{member initializer list}.

\subsection[Potential Pitfalls]{Potential Pitfalls}\label{ctordelegrating-potential-pitfalls}

\subsubsection[Delegation cycles]{Delegation cycles}\label{delegation-cycles}

If a constructor delegates to itself either directly or indirectly, the
program is \emcppsgloss[ill formed, no diagnostic required (IFNDR)]{ill formed, no diagnostic required (IFNDR)}. While some compilers can, under certain conditions, detect delegation cycles at compile time, they are neither
required nor necessarily able to do so. For example, even the simplest delegation cycles might not result in a diagnostic from a compiler{\cprotect\footnote{GCC 10.x does not detect this delegation
cycle at compile time and produces a binary that, if run, will
necessarily exhibit \emcppsgloss{undefined behavior}. Clang 10.x, on the
other hand, halts compilation with a helpful error message:

\begin{lstlisting}[language=bash, style=footcodeplain]
error: constructor for S creates a delegation cycle
\end{lstlisting}\vspace*{-1ex}
      }}:

\begin{emcppslisting}[emcppsignore={Intentionally Fails Everywhere}]
struct S  // Object
{
    S(int)  : S(true) { }  // delegating constructor
    S(bool) : S(0)    { }  // delegating constructor
};
\end{emcppslisting}

\subsubsection[Suboptimal factoring]{Suboptimal factoring}\label{suboptimal-factoring}

The need for delegating constructors might result from initially
suboptimal factoring --- e.g., in the case where the same \emcppsgloss{value}
is being presented in different forms to a variety of different
\emcppsgloss{mechanisms}. For example, consider the \lstinline!IPV4Host! class
in \intrarefsimple{ctordelegating-use-cases}. %{\titleref{ctordelegating-use-cases}}
%(which starts on page~\pageref{ctordelegating-use-cases}).
While having two constructors to
initialize the host might be appropriate, if either (1) the number of
ways of expressing the same value increases or (2) the number of
consumers of that value increases, we might be well advised to create a
separate \emcppsgloss[value semantic]{value-semantic} type, e.g., \lstinline!IPV4Address!, to
represent that value{\cprotect\footnote{The notion that each component
in a subsystem ideally performs one focused function well is sometimes
referred to as separation of (logical) concerns or
  fine-grained (physical) factoring; see \cite{dijkstra82} and see \cite{lakos20},
  sections 0.4, 3.2.7, and 3.5.9, pp.~20--28, 529--530, and 674--676,
  respectively.}}:

\begin{emcppslisting}[emcppsbatch=e4]
#include <cstdint>  // (ù{\codeincomments{std::uint16\_t}}ù), (ù{\codeincomments{std::uint32\_t}}ù)
#include <string>   // (ù{\codeincomments{std::string}}ù)


class IPV4Address
{
    std::uint32_t d_address;
    std::uint16_t d_port;

public:
    IPV4Address(std::uint32_t address, std::uint16_t port)
        : d_address{address}, d_port{port}
    {
    }

    IPV4Address(const std::string& ip)
        : IPV4Address{extractAddress(ip), extractPort(ip)}
    {
    }
};
\end{emcppslisting}

\noindent Note that \lstinline!IPV4Address! itself makes use of delegating
constructors but as a purely private, encapsulated implementation
detail. With the introduction of \lstinline!IPV4Address! into the codebase,
\lstinline!IPV4Host! (and similar components requiring an
\lstinline!IPV4Address! value) can be redefined to have a single
constructor (or other previously overloaded member function) taking
an \lstinline!IPV4Address! object as an argument.

\subsection[Annoyances]{Annoyances}\label{annoyances}

\hspace*{\fill}

\subsection[See Also]{See Also}\label{see-also}

\begin{itemize}
\item{\seealsoref{forwardingref}{\seealsolocationc}provides perfect forwarding of arguments from one ctor to another.}%\pagebreak%%%%%
\item{\seealsoref{variadictemplate}{\seealsolocationc}describes how to implement constructors that forward an arbitrary list of arguments to other constructors.}
\end{itemize}

\subsection[Further Reading]{Further Reading}\label{further-reading}

\hspace*{\fill}





\newpage
%\section[Deleted Functions]{Using {\SecCode =}\,{\SecCode delete} for Arbitrary Functions}\label{deleted-functions}
% 13 Jan, copyedits in and proofed
% 27 Jan, code compiled by Josh
% 27 Jan, new section added
% 27 Jan, new section copyedited and copyedits entered
% 8 Feb 2021, Josh: Cleanup in comments
% 9 Feb 2021, FPPs set
% 17 Feb 2021, glossary tags updated

\emcppsFeature{
    toclong={Using \lstinline!=!\,\lstinline!delete! for Arbitrary Functions},
    short={Deleted Functions},
    long={Using {\SecCode =}\,{\SecCode delete} for Arbitrary Functions},
}{deleted-functions}
%\section[Deleted Functions]{Using {\SecCode =delete} for Arbitrary Functions}\label{deleted-functions}
\setcounter{table}{0}
\setcounter{footnote}{0}
\setcounter{lstlisting}{0}

Using \lstinline!=!\,\lstinline!delete! in a function's first declaration
forces a compilation error upon any attempt to use or access it.

\subsection[Description]{Description}\label{description}

Declaring a particular function or function overload to result in a
fatal diagnostic upon invocation can be useful --- e.g., to suppress the
generation of a \emcppsgloss[special member function]{special member function} or to limit the types of
arguments a particular overload set is able to accept. In such cases,
\lstinline!=!\,\lstinline!delete! followed by a semicolon (\lstinline!;!) can be used in place of the body of any
function on first declaration only to force a compile-time error if
any attempt is made to invoke it or take its address.

\begin{emcppslisting}
void g(double) { }
void g(int) = delete;

void f()
{
    g(3.14);  // OK, (ù{\codeincomments{f(double)}}ù) is invoked.
    g(0);     // Error, (ù{\codeincomments{f(int)}}ù) is deleted.
}
\end{emcppslisting}

\noindent Notice that deleted functions participate in \emcppsgloss{overload resolution}
and produce a compile-time error when selected as the best candidate.

\subsection[Use Cases]{Use Cases}\label{use-cases}

\subsubsection[Suppressing special member function generation]{Suppressing special member function generation}\label{suppressing-special-member-function-generation}

When instantiating an object of user-defined type, \emcppsgloss[special member function]{special member functions} that have not been declared explicitly are
often generated automatically by the compiler.
The generation of individual special member
functions can be affected by the existence of other user-defined special
member functions or by limitations imposed by the specific types of
any data members or base types; see
%Section~\ref{Defaulted-Special-Member-Functions}, ``\titleref{Defaulted-Special-Member-Functions} on page~\pageref{Defaulted-Special-Member-Functions}."
\featureref{\locationa}{Defaulted-Special-Member-Functions}.
%For certain kinds of types, the notion of \romeogloss{copy semantics}
%including \romeogloss{move semantics} is not meaningful, and
%hence permitting the generation of copy operations is contraindicated. The two special member functions controlling \romeogloss{\emph{move} operations}, introduced in C++11, are sometimes implemented as effective optimizations of \romeogloss{\emph{copy} operations} and much less frequently with \romeogloss{\emph{copy} operations} explicitly deleted; see \featureref{\locationc}{Rvalue-References}.
For certain kinds of types, the notion of \emph{copying} is not meaningful, and hence permitting the compiler to generate \emph{copy} operations would be inappropriate. The two special member functions controlling \emcppsgloss{move operations}, introduced in C++11, are typically implemented as effective optimizations of \emcppsgloss[copy operations]{\emph{copy} operations} and thus would be similarly contraindicated. Much less frequently, a useful notion of moving exists where copying does not, and so we may choose to have move operations generated, while \emcppsgloss[copy operations]{\emph{copy} operations} are explicitly deleted; see \featureref{\locationc}{Rvalue-References}.
\newpage%%%%%%%

Consider a class, \lstinline!FileHandle!, that uses the \emcppsgloss{RAII} idiom
to safely acquire and release an I/O stream. As \emcppsgloss{copy semantics}
are typically not meaningful for such resources, we will want to
suppress generation of both the \emcppsgloss{copy constructor} and \emcppsgloss{copy
assignment operator}. Prior to C++11, there was no direct way to express
suppression of \emcppsgloss{special functions} in C++. The commonly
recommended workaround was to declare the two methods \lstinline!private!
and leave them unimplemented, typically resulting in a compile-time or
link-time error when accessed:

\begin{emcppslisting}
#include <cstdio>  // (ù{\codeincomments{FILE}}ù)
class FileHandle
{
private:
    // ...

    FileHandle(const FileHandle&);             // not implemented
    FileHandle& operator=(const FileHandle&);  // not implemented

public:
    explicit FileHandle(FILE* filePtr);
    ~FileHandle();

    // ...
};
\end{emcppslisting}

\noindent Not implementing a special member function that is declared to be private ensures that there will be at least a link-time error in case that function is inadvertently accessed from within the implementation of the class itself. With the \lstinline!=!\,\lstinline!delete! syntax, we are able to (1)
explicitly express our intention to make these special member
functions unavailable, (2) do so directly in the \lstinline!public! region
of the class, and (3) enable clearer compiler diagnostics:

%begin{emcppslisting}
%class FileHandle
%{
%private:
%    // ...
%
%public:
%    explicit FileHandle(FILE* filePtr);
%    ~FileHandle();
%
%    FileHandle(const FileHandle&) = delete;             // make unavailable
%    FileHandle& operator=(const FileHandle&) = delete;  // make unavailable
%
%    // ...
%};
%\end{emcppslisting}
\begin{emcppshiddenlisting}[emcppsbatch=e1]
#include <cstdio>  // (ù{\codeincomments{FILE}}ù)
\end{emcppshiddenlisting}
\begin{emcppslisting}[emcppsbatch=e1]
class FileHandle
{
private:
    // ...
    // Declarations of copy constructor and copy assignment are now public.

public:
    explicit FileHandle(FILE* filePtr);
    ~FileHandle();

    FileHandle(const FileHandle&) = delete;             // make unavailable
    FileHandle& operator=(const FileHandle&) = delete;  // make unavailable

    // ...
};
\end{emcppslisting}
\newpage%%%%%

\noindent Using the \lstinline!=!\,\lstinline!delete! syntax on declarations that are private results in error messages concerning privacy, not the use of deleted functions. Care must be exercised to make \emph{both} changes when converting code from the old style to the new syntax.

\subsubsection[Preventing a particular implicit conversion]{Preventing a particular implicit conversion}\label{preventing-a-particular-implicit-conversion}

Certain functions --- especially those that take a \lstinline!char! as an
argument --- are prone to inadvertent misuse. As a truly classic
example, consider the C library function \lstinline!memset!, which may be used
to write the character \lstinline!*! five times in a row, starting at a
specified memory address, \lstinline!buf!:

%begin{emcppslisting}
%#include <cstring>  // (ù{\codeincomments{puts}}ù)
%#include <cstdio>   // (ù{\codeincomments{memset}}ù)
%
%void f()
%{
%    char buf[] = "Hello World!";
%    memset(buf, 5, '*');  // undefined behavior
%    puts(buf);            // expected output: "***** World!"
%}
%\end{emcppslisting}
\begin{emcppslisting}
#include <cstdio>   // (ù{\codeincomments{puts}}ù)
#include <cstring>  // (ù{\codeincomments{memset}}ù)

void f()
{
    char buf[] = "Hello World!";
    memset(buf, 5, '*');  // undefined behavior: buffer overflow
    puts(buf);            // expected output: "***** World!"
}
\end{emcppslisting}


\noindent Sadly, inadvertently reversing the order of the last two arguments is a commonly
recurring error, and the C language provides no help. As shown above, \lstinline!memset! writes the nonprinting character \lstinline!5! (e.g., the integer value of ASCII \lstinline!'*'!) 42 times --- way past the end of \lstinline!buf!. In C++, we
can target such observed misuse using an extra deleted overload:

%begin{emcppslisting}
%#include <cstring>  // (ù{\codeincomments{memset}}ù)
%void* memset(void* str, int ch, size_t n);      // standard library function
%void* memset(void* str, int n, char) = delete;  // defensive against misuse
%\end{emcppslisting}
\begin{emcppslisting}
#include <cstring>  // (ù{\codeincomments{memset}}ù)
void* memset(void* str, int ch, size_t n);      // standard library function
void* memset(void* str, int n, char) = delete;  // defense against misuse
\end{emcppslisting}

\noindent Pernicious user errors can now be reported during compilation:

\begin{emcppslisting}
// ...
memset(buf, 5, '*');  // Error, (ù{\codeincomments{memset(void\*, int, char)}}ù) is deleted.
// ...
\end{emcppslisting}


\subsubsection[Preventing all implicit conversions]{Preventing all implicit conversions}\label{preventing-all-implicit-conversions}

The \lstinline!ByteStream::send! member function below is designed to work
with 8-bit unsigned integers only. Providing a deleted overload
accepting an \lstinline!int! forces a caller to ensure that the argument is
always of the appropriate type:
\newpage%%%%%%

\begin{emcppslisting}
class ByteStream
{
public:
    void send(unsigned char byte) { /* ... */ }
    void send(int) = delete;

    // ...
};

void f()
{
    ByteStream stream;
    stream.send(0);   // Error, (ù{\codeincomments{send(int)}}ù) is deleted.     (1)
    stream.send('a'); // Error, (ù{\codeincomments{send(int)}}ù) is deleted.     (2)
    stream.send(0L);  // Error, ambiguous                 (3)
    stream.send(0U);  // Error, ambiguous                 (4)
    stream.send(0.0); // Error, ambiguous                 (5)
    stream.send(
        static_cast<unsigned char>(100));  // OK          (6)
}
\end{emcppslisting}

\noindent Invoking \lstinline!send! with an \lstinline!int! (noted with (1) in the code above) or any integral type
(other than \lstinline!unsigned!~\lstinline!char!) that promotes to \lstinline!int! (2)
will map exclusively to the deleted \lstinline!send(int)! overload; all
other integral (3 and 4) and floating-point types (5) are convertible to
both via a \emcppsgloss{standard conversion} and hence will be ambiguous. Note that
implicitly converting from \lstinline!unsigned!~\lstinline!char! to either a
\lstinline!long! or \lstinline!unsigned! integer involves a \emcppsgloss{standard
conversion} (not just an \emcppsgloss{integral promotion}), the same as
  converting to a \lstinline!double!.
An explicit cast to \lstinline!unsigned!~\lstinline!char! (6) can always be
pressed into service if needed.

\subsubsection[Hiding a structural (nonpolymorphic) base class's member function]{Hiding a structural (nonpolymorphic) base class’s member function}\label{hiding-a-structural-(nonpolymorphic)-base-class's-member-function}

It is commonly advised to avoid deriving publicly from concrete classes because by doing so, we do not hide the underlying capabilities, which can easily be accessed (potentially breaking any invariants the derived class may want to keep) via assignment to a pointer or reference to a base class, with no casting required.  Worse, inadvertently passing such a class to a function taking the base class by value will result in slicing, which can be~especially problematic when the derived class holds data. A more robust approach would be to use layering or at least private inheritance.\footnote{For more on improving compositional designs at scale, see
  \cite{lakos20}, sections 3.5.10.5 and 3.7.3, pp.~687--703 and
  726--727, respectively.} Best practices notwithstanding,\footnote{See ``Inheritance and Object-Oriented Programming," Item 38, \cite{meyers92}, pp. 132--135.} it can be cost-effective in the short
term to provide an elided ``view'' on a concrete class for trusted
clients. Imagine a class \lstinline!AudioStream! designed to play sounds
and music that --- in addition to providing basic ``play'' and
``rewind'' operations --- sports a large, robust interface:
\newpage%%%%%%

\begin{emcppshiddenlisting}[emcppsbatch=e2]
struct ForwardAudioStream;
struct FMRadio {
    static ForwardAudioStream getStream();
};
\end{emcppshiddenlisting}
\begin{emcppslisting}[emcppsbatch=e2]
struct AudioStream
{
    void play();
    void rewind();
    // ...
    // ... (large, robust interface)
    // ...
};
\end{emcppslisting}


Now suppose that, on short notice, we need to whip up a similar
class,\linebreak[4] \lstinline!ForwardAudioStream!, to use with audio samples that
cannot be rewound (e.g., coming directly from a live feed). Realizing
that we can readily reuse most of \lstinline!AudioStream!'s interface, we
pragmatically decide to prototype the new class simply by exploiting
public \emcppsgloss{structural inheritance} and then deleting just the lone
unwanted \lstinline!rewind! member function:

\begin{emcppslisting}[emcppsbatch=e2]
struct ForwardAudioStream : AudioStream
{
    void rewind() = delete; // Make just this one function unavailable.
};

void f()
{
    ForwardAudioStream stream = FMRadio::getStream();
    stream.play();   // fine
    stream.rewind(); // Error, (ù{\codeincomments{rewind()}}ù) is deleted.
}
\end{emcppslisting}

\noindent If the need for a \lstinline!ForwardAudioStream! type persists, we can always
consider reimplementing it more carefully later.{\cprotect\footnote{\cite{lakos20}, sections~3.5.10.5 and 3.7.3, pp.~687--703 and 726--727}} As discussed at the beginning of this section, the protection provided by this example is easily circumvented:
\begin{emcppslisting}[emcppsbatch=e2]
void g(const ForwardAudioStream &stream)
{
    AudioStream fullStream = stream;
    fullStream.play();   // OK
    fullStream.rewind(); // compiles OK, but what happens at run time?
}
\end{emcppslisting}

\noindent Hiding non\lstinline!virtual! functions is something one undertakes only after attaining a complete understanding of what makes such an unorthodox endeavor \emph{safe}; see, in particular, the appendix of \featureref{\locatione}{final}.

\subsection[Potential Pitfalls]{Potential Pitfalls}\label{potential-pitfalls}

\hspace*{\fill}
\newpage%%%%%%%

\subsection[Annoyances]{Annoyances}\label{annoyances}

%%%%%%%%%%%%%%% new material added 27 Jan 2021
\subsubsection[Deleting a function declares it]{Deleting a function declares it}

It should come as no surprise that when we ``declare" a \emcppsgloss{free function}
followed by \lstinline!=!\,\lstinline!delete!, we \emph{are} in fact \emph{declaring} it.  For example, consider
the pair of overloads of functions \lstinline!f! declared taking a \lstinline!char! and \lstinline!int!,
respectively:

\begin{emcppslisting}
int f(char);          // (1) accessible declaration of (ù{\codeincomments{f}ù) taking a (ù{\codeincomments{char}ù)
int f(int) = delete;  // (2) inaccessible declaration of (ù{\codeincomments{f}ù) taking an (ù{\codeincomments{int}ù)

int x = f('a');       // OK, exact match for (1) (ù{\codeincomments{f(char)}ù), which is accessible
int y = f(123);       // Error, exact match for (2) (ù{\codeincomments{f(int)}ù), which is deleted
\end{emcppslisting}

\noindent It is necessary that both functions above are \emph{declared} so that both of them
can participate in overload resolution; it is only after the inaccessible
overload is selected that it will be reported as a compile-time error.

When it comes to deleting certain \emcppsgloss[special member function]{special member functions} of a class (or
class  template), however, what might seem like a tiny bit of extra,
self-documenting code can have subtle, unintended consequences as evidenced
below.

Let's begin by considering an empty \lstinline!struct!, \lstinline!S0!:

\begin{emcppslisting}[emcppsbatch=e3]
struct S0 { };  // The default constructor is declared implicitly.

S0 x0;  // OK, invokes the implicitly generated default constructor
\end{emcppslisting}

\noindent As \lstinline!S0! defines not constructors, destructors, or assignment operators,
the compiler will generate (\emcppsgloss{declare} and \emcppsgloss{define}), for \lstinline!S0!, all
\emph{six} of  the \emcppsgloss[special member function]{special member functions} available as of C++11; see
\featureref{\locationa}{defaulted-special-member-functions}.

Next, suppose we create a second \lstinline!struct!, \lstinline!S1!, that differs from \lstinline!S0!
only in that \lstinline!S1! declares a \emph{value} constructor taking an \lstinline!int!:

\begin{emcppslisting}[emcppsbatch=e3]
struct S1  // Implicit declaration of the default constructor is suppressed.
{
    S1(int);  // explicit declaration of (ù{\emphincomments{value}ù) constructor
};

S1 y1(5);  // OK, invokes the explicitly declared (ù{\emphincomments{value}ù) constructor
S1 x1;     // Error, no declaration for default constructor (ù{\codeincomments{S1::S1()}ù)
\end{emcppslisting}

\noindent By explicitly declaring a \emph{value} constructor (or any other constructor
for that matter), we automatically suppress the implicit declaration of
the default constructor for \lstinline!S1!. If suppressing the default destructor is
\emph{not} our intention, we can always reinstate it via an explicit declaration
followed by \lstinline!=!\,\lstinline!default;! --- see \featureref{\locationa}{defaulted-special-member-functions}.
\newpage%%%%%%%

Let's now suppose it \emph{is} our intention to suppress generation of the default
constructor and, to make our intention clear, we elect to explicitly \emcppsgloss{declare}
and \emcppsgloss{delete} it:

\begin{emcppslisting}[emcppsbatch=e3]
struct S2  // Implicit declaration of the default constructor is suppressed.
{
    S2() = delete;  // explicit declaration of inaccessible default constructor
    S2(int);        // explicit declaration of (ù{\emphincomments{value}ù) constructor
};

S2 y2(5);  // OK, invokes the explicitly declared (ù{\emphincomments{value}ù) constructor
S2 x2;     // Error, use of deleted function, (ù{\codeincomments{S2::S2()}ù)
\end{emcppslisting}

\noindent By declaring and then deleting the default constructor we have, it would
appear that we (1) made our intentions clear and (2) improved diagnostics for our
clients at the cost of a single extra line of self-documenting code.  Ah, if
C++ were only that straightforward.

Deleting certain \emcppsgloss[special member function]{special member functions} --- i.e., \emph{default} constructor,
\emph{move} constructor, or \emph{move}-assignment operator --- that are not necessarily
implicitly declared can have nonobvious consequence that adversely affect
subtle compile-time properties of a class. One such subtle property is whether
the compiler considers it to be a \emcppsgloss{literal type} --- i.e., a type whose
\emph{value} is eligible for use as part of a \emcppsgloss{constant expression}.  This same
property of being a \emcppsgloss{literal type} is what determines whether an arbitrary
type may be passed by value in the interface of a \lstinline!constexpr! function; see
\featureref{\locationc}{constexprfunc}.

As a simple illustration of a subtle compile-time difference between \lstinline!S1! and
\lstinline!S2!, consider this practically useful \emph{pattern} for a developer's ``test"
function that will compile if and only if its by-value parameter, \lstinline!x!, is of a
literal type:

\begin{emcppslisting}[emcppsbatch=e3]
constexpr int test(S0 x) { return 0; }  // OK,    (ù{\codeincomments{S0}ù) (ù{\emphincomments{is}ù)     a literal type.
constexpr int test(S1 x) { return 0; }  // Error, (ù{\codeincomments{S1}ù) is (ù{\emphincomments{not}ù) a literal type.
constexpr int test(S2 x) { return 0; }  // OK,    (ù{\codeincomments{S2}ù) (ù{\emphincomments{is}ù)     a literal type.
\end{emcppslisting}

\noindent For the compiler to treat a given class type as a \emcppsgloss{literal type}, it must,
among other things, have at least one constructor (other than the \emph{copy} or
\emph{move} constructor) declared as \lstinline!constexpr!.

In the case of the empty \lstinline!S0! class, the implicitly generated default constructor
is  \emcppsgloss{trivial} and so it is implicitly \emph{declared} \lstinline!constexpr! too.  Class \lstinline!S1!'s
explicitly declared \emph{non}-\lstinline!constexpr! value constructor suppresses the declaration
of its only \lstinline!constexpr! constructor, the  default constructor; hence, \lstinline!S1! does not
qualify as a \emph{literal type}.

Finally, by conspicuously declaring and deleting \lstinline!S2!'s default constructor, we
\emph{declare} it nonetheless.  What's more, the declaration brought about by deleting it
is the same as if it had been generated implicitly (or declared explicitly and then
defaulted); hence, \lstinline!S2!, unlike \lstinline!S1!, \emph{is} a \emcppsgloss{literal type}.  Go figure!
%%%%%%%%%%%%%%%%%%%%%% new material ends

\newpage%%%%%%%

\subsection[See Also]{See Also}\label{see-also}

\begin{itemize}
\item{%Section~\ref{Defaulted-Special-Member-Functions}, ``\titleref{Defaulted-Special-Member-Functions}" on page~\pageref{Defaulted-Special-Member-Functions} —
\seealsoref{Defaulted-Special-Member-Functions}{\seealsolocationa}Companion feature that enables defaulting, as opposed to deleting, special member functions.}
\item{%Section~\ref{Rvalue-References}, ``\titleref{Rvalue-References}" on page~\pageref{Rvalue-Referencess} —
\seealsoref{Rvalue-References}{\seealsolocationc}The two \emph{move} variants of special member functions, which use \romeovalue{rvalue} references in their signatures, may also be subject to deletion.}
\end{itemize}

\subsection[Further Reading]{Further Reading}\label{further-reading}

\begin{itemize}
\item{``Item 27" of \cite{meyers15b}}
\end{itemize}



\newpage
%\section[{\tt explicit} Operators]{Explicit Conversion Operators}\label{explicit-conversion-operators}
% 12 Feb 2021, revisions in; ready for Josh's code fixes
% 16 Feb 2021 JMB; code compiles

\emcppsFeature{
    short={\lstinline!explicit! Operators},
    tocshort={{\TOCCode explicit} Operators},
    long={Explicit Conversion Operators},
    toclong={Explicit Conversion Operators},
    rhshort={{\RHCode explicit} Operators},
}{explicit-conversion-operators}
\setcounter{table}{0}
\setcounter{footnote}{0}
\setcounter{lstlisting}{0}
%\section[{\tt explicit} Operators]{Explicit Conversion Operators}\label{explicit-conversion-operators}

Ensure that a user-defined type is convertible to another type only in
contexts where the conversion is made obvious in the code.

\subsection[Description]{Description}\label{description-explicitconv}

Though sometimes desirable, implicit conversions achieved via user-defined \emph{conversion
functions} --- either \romeogloss{converting constructors} accepting a
single argument or \romeogloss{conversion operators} --- can also be problematic, especially when the
conversion involves a commonly used type (e.g., \lstinline!int! or
\lstinline!double!){\cprotect\footnote{Use of a conversion operator to
calculate distance from the origin in this unrealistically simple \lstinline!Point!
example is for didactic purposes only. In practice, we would typically
use a named function for this purpose; see \intraref{potential-pitfalls-explicitconv}{sometimes-a-named-function-is-better}.}}:
%{\it\titleref{potential-pitfalls-explicitconv}: \titleref{sometimes-a-named-function-is-better}} on page~\pageref{sometimes-a-named-function-is-better}.}}:

\begin{emcppslisting}[emcppsbatch=e1]
class Point  // implicitly convertible from an (ù{\codeincomments{int}}ù) or to a (ù{\codeincomments{double}}ù)
{
    int d_x, d_y;

public:
    Point(int x = 0, int y = 0);  // default, conversion, & value constructor
    // ...
    operator double() const;  // Return distance from origin as a (ù{\codeincomments{double}}ù).
};
\end{emcppslisting}
    
\noindent As ever, calling a function that takes a \lstinline!Point! but
accidentally passing an \lstinline!int! can lead to surprises:

\begin{emcppslisting}[emcppsbatch=e1]
void g0(Point p);         // arbitrary function taking a (ù{\codeincomments{Point}}ù) object by value
void g1(const Point& p);  // arbitrary function taking a (ù{\codeincomments{Point}}ù) by (ù{\codeincomments{const}}ù) reference

void f1(int i)
{
    g0(i);  // oops, called (ù{\codeincomments{g0}}ù) with (ù{\codeincomments{Point(i, 0)}}ù) by mistake
    g1(i);  // oops, called (ù{\codeincomments{g1}}ù) with (ù{\codeincomments{Point(i, 0)}}ù) by mistake
}
\end{emcppslisting}
    
\noindent This problem could have been solved even in C++98 by declaring the
constructor to be \lstinline!explicit!:

\begin{emcppshiddenlisting}[emcppsbatch=e2]
class Point;
void g0(Point p);       
void g1(const Point& p);
class Point {
public:
\end{emcppshiddenlisting}
\begin{emcppslisting}[emcppsbatch=e2]
explicit Point(int x = 0, int y = 0);  // explicit converting constructor
\end{emcppslisting}
\begin{emcppshiddenlisting}[emcppsbatch=e2]
};
\end{emcppshiddenlisting}
\noindent If the conversion is desired, it must now be specified explicitly:

\begin{emcppslisting}[emcppsbatch=e2]
void f2(int i)
{
    g0(i);         // Error, could not convert (ù{\codeincomments{i}}ù) from (ù{\codeincomments{int}}ù) to (ù{\codeincomments{Point}}ù)
    g1(i);         // Error, invalid initialization of reference type
    g0(Point(i));  // OK
    g1(Point(i));  // OK
}
\end{emcppslisting}
    
\noindent The companion problem stemming from an \emph{implicit conversion
operator}, albeit less severe, remained:

\begin{emcppslisting}[emcppsbatch=e1]
void h(double d);

double f3(const Point& p)
{
    h(p);      // OK? Or maybe called (ù{\codeincomments{h}}ù) with a "hypotenuse" by mistake
    return p;  // OK? Or maybe this is a mistake too.
}
\end{emcppslisting}

\noindent As of C++11, we can now use the \romeogloss{\lstinline!explicit! specifier}
when declaring \romeogloss{conversion operators} (as well as
\romeogloss{converting constructors}), thereby forcing the client to request
conversion explicitly --- e.g., using \romeogloss{direct initialization} or
\lstinline!static_cast!:

\begin{emcppslisting}
struct S0 { explicit operator int(); };

void g()
{
    S0 s0;
    int i = s0;                    // Error, copy initialization
    int k(s0);                     // OK, direct initialization
    double d = s0;                 // Error, copy initialization
    int j = static_cast<int>(s0);  // OK, static cast
    if (s0) { }                    // Error, contextual conversion to (ù{\codeincomments{bool}}ù)
    double e(s0);                  // Error, direct initialization
}
\end{emcppslisting}
    
\noindent In contrast, had the conversion operator above not been declared to be
\lstinline!explicit!, all conversions shown above would compile:

\begin{emcppslisting}
struct S1 { /* implicit */ operator int(); };

void f()
{
    S1 s1;
    int i = s1;                    // OK (copy initialization)
    double d = s1;                 // OK (copy initialization)
    int j = static_cast<int>(s1);  // OK (static cast)
    if (s1) { }                    // OK (contextual conversion to (ù{\codeincomments{bool}}ù))
    int k(s1);                     // OK (direct initialization)
    double e(s1);                  // OK (direct initialization)
}
\end{emcppslisting}
    
\noindent Additionally, the notion of \romeogloss{contextual convertibility to
\lstinline!bool!} 
%{\cprotect\footnote{Since the early days of C++, a common
%  idiom to test for validity of an object has been to use it in a
%  context where it can (implicitly) convert itself to a type whose value
%  can be interpreted (contextually) as a boolean, with \lstinline!true!
%  implying validity (and \lstinline!false! otherwise). Implicit conversion
%  to \lstinline!bool! (an integral type) was considered too dangerous, 
%  so the cumbersome \romeogloss{safe-\lstinline!bool! idiom} was used instead,
%  converting to a type that --- while contextually convertible to
%  \lstinline!bool! --- could not (by design) participate in any other
%  operations. While making the conversion to \lstinline!bool! (or
%  \lstinline!const!~\lstinline!bool!) \lstinline!explicit! solves the safety
%  issue, the benefit of the idiom would be entirely lost if an explicit
%  cast would have to be performed to test for validity. To address this,
%  C++11 extends contextual conversion to \lstinline!bool! for a given
%  expression \lstinline!E! to include an application of
%  \lstinline!static_cast<const!~\lstinline!volatile!~\lstinline!bool>! to
%  \lstinline!E!, thus enabling explicit conversion to \lstinline!bool! to be
%  used in lieu of the (now deprecated) \romeogloss{safe-\lstinline!bool! idiom}; see
%  \cite{sharpe13}.}} 
applicable to arguments of logical operations
(e.g., \lstinline!&&!, \lstinline!||!, and \lstinline|!|) and conditions of
most control-flow constructs (e.g., \lstinline!if!, \lstinline!while!) was
extended in C++11 to admit \emph{explicit} (user-defined) \lstinline!bool!
conversion operators (see \intraref{use-cases-explicitconv}{enabling-contextual-conversions-to-bool-as-a-test-for-validity}):
%{\it\titleref{use-cases-explicitconv}: \titleref{enabling-contextual-conversions-to-bool-as-a-test-for-validity}} on page~\pageref{enabling-contextual-conversions-to-bool-as-a-test-for-validity}):

\begin{emcppslisting}
struct S2 { explicit operator bool(); };

void h()
{
    S2 s2;
    int i = s2;                    // Error, copy initialization
    double d = s2;                 // Error, copy initialization
    int j = static_cast<int>(s2);  // Error, static cast
    if (s2) { }                    // OK, contextual conversion to (ù{\codeincomments{bool}}ù)
    int k(s2);                     // Error, direct initialization
    double fd(s2);                 // Error, direct initialization
    bool b0 = s2;                  // Error, copy initialization
    bool b1(s2);                   // OK, direct initialization)                                     
    !s2;                           // OK, contextual conversion to (ù{\codeincomments{bool}}ù)                                               -
    s2 && s2;                      // OK, contextual conversion to (ù{\codeincomments{bool}}ù)
}
\end{emcppslisting}
    
%\noindent Prior to C++11, essentially the same effect as having an \emph{explicit}
%\lstinline!operator!~\lstinline!bool()! member was achieved (albeit far less
%conveniently) via the \romeogloss{safe-\lstinline!bool! idiom}.

\subsection[Use Cases]{Use Cases}\label{use-cases-explicitconv}

\subsubsection[Enabling contextual conversions to {\tt bool} as a test for validity]{Enabling contextual conversions to {\SubsubsecCode bool} as a test for validity}\label{enabling-contextual-conversions-to-bool-as-a-test-for-validity}

Having a conventional test for validity that involves testing whether the object itself evaluates to \lstinline!true! or \lstinline!false! is an idiom that goes back to the
origins of C++. The Standard input/output library, for example, uses this
idiom to determine if a given stream is valid:

\begin{emcppslisting}
// C++03
#include <ostream>  // (ù{\codeincomments{std::ostream}}ù)

std::ostream& printTypeValue(std::ostream& stream, double value)
{
    if (stream)  // relies on an implicit conversion to (ù{\codeincomments{bool}}ù)
    {
        stream << "double(" << value << ')';
    }
    else
    {
        // ... (handle stream failure)
    }

    return stream;
}
\end{emcppslisting}
    
\noindent Implementing the implicit conversion to \lstinline!bool! was, however,
problematic as the straight\-forward approach of using a
\romeogloss{conversion operator} could easily allow accidental misuse to go
undetected:

\begin{emcppslisting}
class ostream
{
    // ...

    public:

    /* implicit */ operator bool();  // hypothetical (bad) idea
};

int client(ostream& out)
{
    // ...
    return out + 1;  // likely a latent runtime bug: always returns 1 or 2
}
\end{emcppslisting}
    
\noindent The classic workaround, the \romeogloss{safe-\lstinline!bool! idiom},\footnote{https://www.artima.com/cppsource/safebool.html} 
was to return some obscure pointer\linebreak[4] %%%%%
 type (e.g., \romeogloss{pointer to
member}) that could not possibly be useful in any context other\linebreak[4] %%%%%
 than one in which \lstinline!false! and a null pointer-to-member value are
treated equivalently. With explicit conversion operators, such workarounds are no longer required. As discussed in \intrarefsimple{description-explicitconv}, 
%\textit{\titleref{description-explicitconv}} on page~\pageref{description-explicitconv}, 
a conversion operator to type
\lstinline!bool! that is declared \lstinline!explicit! continues to act as if
it were \emph{implicit} only in those places where we might want it to
do so and nowhere else --- i.e., exactly those places that enable
\romeogloss{contextual conversion to \lstinline!bool!}.{\cprotect\footnote{Note that two
consecutive \lstinline|!| operators can be used to synthesize a
\romeogloss{contextual conversion to \lstinline!bool!} --- i.e., if \lstinline!X! is an
expression that is explicitly convertible to \lstinline!bool!, then
\lstinline|(!!(X))| will be \lstinline!(true)! or \lstinline!(false)!
  accordingly.}}

As a concrete example, consider a \lstinline!ConnectionHandle! class that
can be in either a \emph{valid} or \emph{invalid} state. For the user's
convenience and consistency with other proxy types (e.g., raw pointers)
that have a similar \emph{invalid} state, representing the invalid (or null) state via an explicit conversion to
\lstinline!bool! might be desirable:


\begin{emcppslisting}[emcppsbatch=e2]
#include <cstddef>  // (ù{\codeincomments{std::size\_t}ù)                                            
#include <iostream> // (ù{\codeincomments{std::cerr}ù)
struct ConnectionHandle
{
    std::size_t maxThroughput() const;
        // Return the maximum throughput (in bytes) of the connection.

    explicit operator bool() const;
        // Return (ù{\codeincomments{true}}ù) if the handle is valid and (ù{\codeincomments{false}}ù) otherwise.
};
\end{emcppslisting}
    
\noindent Instances of \lstinline!ConnectionHandle! will convert to \lstinline!bool!
only where one might reasonably want them to do so, say, as the
predicate of an \lstinline!if! statement:

\begin{emcppslisting}[emcppsbatch=e2]
int ping(const ConnectionHandle& handle)
{
    if (handle)  // OK (contextual conversion to bool)
    {
        // ...
        return 0;  // success
    }

    std::cerr << "Invalid connection handle.\n";
    return -1;  // failure
}
\end{emcppslisting}
    
\noindent Having an \lstinline!explicit! conversion operator prevents unwanted
conversions to \lstinline!bool! that might otherwise happen inadvertently:

\begin{emcppslisting}[emcppsbatch=e2]
bool hasEnoughThroughput(const ConnectionHandle& ingress,
                         const ConnectionHandle& egress)
{
    return ingress.throughput() <= egress;  // Error, thankfully
//                                    ^~~~~~
}
\end{emcppslisting}
    
\noindent After the relational operator (\lstinline!<=!) in the example above, the
programmer mistakenly wrote \lstinline!egress! instead of
\mbox{\lstinline!egress.maxThroughput()!}. Fortunately the conversion operator of\linebreak[4] %%%%%
\mbox{\lstinline!ConnectionHandle!} was declared to be \lstinline!explicit! and a
compile-time error ensued; if the conversion had
been \emph{implicit}, the example code above would have compiled, and,
if executed, the above (faulty) implementation of the \mbox{\lstinline!hasEnoughThroughput!}
function would have silently exhibited well-defined but incorrect
behavior.

\subsection[Potential Pitfalls]{Potential Pitfalls}\label{potential-pitfalls-explicitconv}

\subsubsection[Sometimes implicit conversion \emph{is} indicated]{Sometimes implicit conversion {\sfbsubsubsecitalRomeo is} indicated}\label{sometimes-implicit-conversion-is-indicated}

Implicit conversions to and from common arithmetic types, especially
\lstinline!int!, are generally ill advised given the likelihood of
accidental misuse. 
%However, sometimes \emph{implicit}
%conversion is exactly what is needed. Such cases occur frequently with
%wrapper and proxy types that might need to interoperate with a large
%legacy codebase. Consider, for example, an initial implementation of
%memory allocators in which each constructor takes, as an optional
%trailing argument, a pointer to an abstract memory resource that itself
%provides pure virtual \lstinline!allocate! and \lstinline!deallocate! member
%functions. Later, we decide to move in the direction of the
%\lstinline!std::pmr! (C++17) standard and wrap those pointers in classes
%that support additional operations. Making such constructors on the
%wrapper explicit would force every client supplying an allocator to a
%container to rework their code (e.g., by using \lstinline!static_cast!).
%An implicit conversion in this case is further justified because the
%likelihood of accidental spontaneous conversion to an \lstinline!Allocator!
%is all but nonexistent.
%
%The same sort of stability argument favors implicit conversion for proxy
%types intended to be dropped in and used in existing codebases. If, for
%example, we wanted to provide a proxy for a writeable
%\lstinline!std::string! that, say, also logged, we might want an implicit
%conversion to a \lstinline!std::string&&! (perhaps using reference qualifiers; 
%see \featureref{\locatione}{refqualifiers}). In such cases, making the
%conversion explicit would entirely defeat the purpose of the proxy,
%which is to achieve new functionality with minimal effect on existing
%client code.
However, for proxy types that are intended to be drop-in replacements for the types they represent, implicit conversions are precisely what we want. Consider, for example, a \lstinline!NibbleConstReference!
 proxy type that represents the 4-bit integer elements of a \lstinline!PackedNibbleVector!:


\begin{emcppslisting}[emcppsbatch=e3]
class NibbleConstReference
{
    // ...
public:
    operator int() const; // IMPLICIT

    // ...
};

class PackedNibbleVector 
{
    // ...
public:
    bool empty() const;
    NibbleConstReference operator[](int index) const;
    
    // ...
};
\end{emcppslisting}

\noindent The \lstinline!NibbleConstReference! proxy is intended to interoperate well with other integral types in various expressions and making its conversion operator \lstinline!explicit! hinders its intended use as a drop-in replacement by requiring an explicit conversion (a.k.a. cast):

\begin{emcppslisting}[emcppsbatch=e3]
int firstOrZero(const PackedNibbleVector& values)
{
    return values.empty() 
        ? 0 
        : values[0]; // Compiles only if conversion operator is implicit.
}
\end{emcppslisting}


\subsection[Sometimes a named function is better]{Sometimes a named function is better}\label{sometimes-a-named-function-is-better}

Other kinds of overuses of even \emph{explicit} conversion
operators exist. Like any user-defined operator, when the operation being
implemented is not somehow either canonical or ubiquitously idiomatic
for that operator, expressing that operation by a
named (i.e., non-operator) function is often better. Recall from \intrarefsimple{description-explicitconv} 
% {\it\titleref{description-explicitconv}} on page~\pageref{description-explicitconv} 
that we used a conversion operator of
class \lstinline!Point! to represent the distance from the origin. This
example serves both to illustrate how conversion operators \emph{can} be
used and also how they probably should \emph{not} be. Consider that (1)
many mathematical operations on a 2-D integral point might return a \lstinline!double! (e.g., \lstinline!magnitude!,
\lstinline!angle!) and (2) we might want to represent the same
information but in different units (e.g., \lstinline!angleInDegrees!,
\lstinline!angleInRadians!).{\cprotect\footnote{Another valid design
decision is returning an object of type \lstinline!Angle! that captures
the amplitude and provides named accessory to the different units
  (e.g., \lstinline!asDegrees!, \lstinline!asRadians!).}}

Rather than employing any conversion \emph{operator} (\lstinline!explicit!
or otherwise), consider instead providing a named function, which (1) is
automatically \lstinline!explicit! and (2) affords both flexibility (in
writing) and clarity (in reading) for a variety of domain-specific
functions --- now and in the future --- that might well have had
overlapping return types:

\begin{emcppslisting}
class Point  // only explicitly convertible (and from only an (ù{\codeincomments{int}}ù))
{
    int d_x, d_y;

public:
    explicit Point(int x = 0, int y = 0);  // explicit converting constructor
    // ...
    double magnitude() const;  // Return distance from origin as a (ù{\codeincomments{double}}ù).
};
\end{emcppslisting}
    
\noindent Note that defining \romeogloss{nonprimitive functionality}, like
\lstinline!magnitude!, in a separate \emph{utility} at a higher level in the
physical hierarchy (e.g., \lstinline!PointUtil::magnitude(const!~\lstinline!Point&!~\lstinline!p)!) might be better still.{\cprotect\footnote{For more on
separating out \romeogloss{nonprimitive functionality}, see
  \cite{lakos20}, sections 3.2.7--3.2.8, pp 529--552.}}

\subsection[Annoyances]{Annoyances}\label{annoyances}

None so far

\subsection[See Also]{See Also}\label{see-also}

None so far

\subsection[Further Reading]{Further Reading}\label{further-reading}

None so far



\newpage
%\section[Function {\tt static} '11]{Threadsafe Function-Scope {\SecCode static} Variables}\label{thread-safe-function-static-variables}

% 01 Feb, 2021 JMB - packet 4 compilation fixes
% 18 Feb 2021, handing off to Josh for revisions
% 6 March 2021, to copyediting
% 19 March 2021, copyedits in and proofed
% 21 March 2021, in FPPs



\emcppsFeature{
    short={Function \lstinline!static! '11},
    tocshort={Function {\TOCCode static} '11},
    long={Thread-Safe Function-Scope {\SecCode static} Variables},
    toclong={Thread-Safe Function-Scope \lstinline!static! Variables},
    rhshort={Function {\RHCode static} '11},
}{thread-safe-function-static-variables}
\setcounter{table}{0}
\setcounter{footnote}{0}
\setcounter{lstlisting}{0}
%\section[Function {\tt static} '11]{Threadsafe Function-Scope {\SecCode static} Variables}\label{thread-safe-function-static-variables}


Initialization of function-scope \lstinline!static! objects is now guaranteed to be
free of data races in the presence
of multiple concurrent threads.

\subsection[Description]{Description}\label{description-functionstatic}

A variable declared at function (a.k.a. local) scope has \emcppsgloss{automatic storage duration}, 
except when it is marked \lstinline!static!, in which case it has \emcppsgloss{static storage duration}.
Variables having \emcppsgloss{automatic storage duration} are allocated on the stack
each time the function is invoked, and initialized when that invocation's \emcppsgloss{flow of control}
passes through the \emcppsgloss{definition} of that object.
In contrast, variables with \emcppsgloss{static storage duration} (e.g., \lstinline!iLocal!)
defined at function scope (e.g., \lstinline!f!) are instead allocated once per program and are initialized
only the first time the \emcppsgloss{flow of control} passes
through the \emcppsgloss{definition} of that object:

\begin{emcppslisting}
#include <cassert>  // standard C (ù{\codeincomments{assert}}ù) macro

int f(int i) // function returning the first argument with which it is called
{
    static int iLocal = i;  // object initialized once only, on the first call
    return iLocal;          // the same (ù{\codeincomments{iLocal}}ù) value is returned on every call
}

int main()
{
    int a = f(10);  assert(a == 10);  // Initialize and return (ù{\codeincomments{iLocal}}ù).
    int b = f(20);  assert(b == 10);  // Return (ù{\codeincomments{iLocal}}ù).
    int c = f(30);  assert(c == 10);  // Return (ù{\codeincomments{iLocal}}ù).

    return 0;
}
\end{emcppslisting}

\noindent In the simple example above, the function, \lstinline!f!, initializes its
\lstinline!static! object, \lstinline!iLocal!, with its argument, \lstinline!i!, only
the first time it is called and then always returns the same value
(e.g., 10). Hence, when that function is called repeatedly with distinct
arguments to initialize~the \lstinline!a!, \lstinline!b!, and
\lstinline!c! variables, all three of them are initialized to the same
value, 10, supplied to the~first invocation of \lstinline!f!. Although the function-scope \lstinline!static!
object, \lstinline!iLocal!, was created after \lstinline!main! was entered, it
will not be destroyed until after \lstinline!main! exits.

\subsubsection[Concurrent initialization]{Concurrent initialization}\label{concurrent-initialization}

Historically, initialization of \emcppsgloss[function scope]{function-scope}
\emcppsgloss{static storage duration} objects was not guaranteed to be safe in a
\emcppsgloss{multithreading context} because it was subject to \emcppsgloss{data races}   
if the function was called concurrently from multiple threads.   These \emcppsgloss{data races}
around initialization can lead to the initializer being invoked multiple times, object
construction running concurrently on the same object, and control flow continuing past
the variable definition before initialization had completed at all.  All of these
variations would result in critical software flaws.
One common but nonportable pre-C++11 workaround was the 
\emph{double-checked lock pattern}; see \intraref{appendix-functionstatic}{appendix:-c++03-double-checked-lock-pattern}.

As of C++11, a conforming compiler is now required to ensure that
initialization of \emcppsgloss[function scope]{function-scope} \emcppsgloss{static storage duration}
objects is performed safely, and exactly once, before execution continues past the initializer, even when the function is called
concurrently from multiple threads. 

\subsubsection[Destruction]{Destruction}\label{destruction}

\emcppsgloss[automatic]{Automatic} objects within a local scope are destroyed when control
leaves the scope in which they are declared.
In contrast, \lstinline!static! local objects that have been initialized are
not destroyed until normal program termination,  either after the \lstinline!main! function returns
normally or when the \lstinline!std::exit! function is
called. 
The order of destruction
of these objects will be the reverse of the order in which they completed construction.
Note that programs can terminate in several other ways,
such as a call to \lstinline!std::quick_exit!, \lstinline!_Exit!, or
\lstinline!std::abort!, that explicitly do \emph{not} destroy
  \emcppsgloss{static storage duration} objects. This behavior is as if each static object is scheduled for destruction by using the C Standard Library function \lstinline!std::atexit! right after construction.


\subsubsection[Logger example]{Logger example}\label{logger-example}

Let's now consider a 
real-world example in which a single object --- e.g.,
\lstinline!localLogger! in the example below --- is used widely throughout a program (see
also \intraref{use-cases-functionstatic}{meyers-singleton}){\cprotect\footnote{An
eminently useful, full-featured logger, known as the \texttt{ball} logger, can be found in the
\texttt{ball} package of the \texttt{bal} package group of Bloomberg's open-source
  BDE libraries (\cite{bde14}, subdirectory \texttt{/groups/bal/ball}).}}:

\begin{emcppshiddenlisting}[emcppsbatch=e3]
struct Logger {
  Logger(const char *f) {}
};
template <typename T>
void operator<<(Logger&,const T&);
\end{emcppshiddenlisting}
\begin{emcppslisting}[emcppsbatch=e3]
Logger& getLogger()  // ubiquitous pattern commonly known as "Meyers Singleton"
{
    static Logger localLogger("log.txt");  // function-local (ù{\codeincomments{static}}ù) definition
    return localLogger;
}

int main()
{
    getLogger() << "hello";
        // OK, invokes (ù{\codeincomments{Logger}}ù)'s constructor for the first (and only) time

    getLogger() << "world";
        // OK, uses the previously constructed (ù{\codeincomments{Logger}}ù) instance
}
\end{emcppslisting}

Here we have an example of the ``Singleton pattern''\footnote{\cite{gamma95}, Chapter~3, section~``Singleton,'' pp.~127--134} being used to create the shared \lstinline!Logger! instance
and provide access to it through the \lstinline!getLogger()! function.  The \lstinline!static! local
instance of \lstinline!Logger!, \lstinline!localLogger!, will be initialized exactly once and then
destroyed after normal program termination.  In C++03, it would not be safe to call
this function concurrently from multiple threads. 
Conversely, C++11 \emph{guarantees} that the initialization  
of \lstinline!localLogger! will happen exactly once even when multiple threads call
\lstinline!getLogger! \mbox{concurrently}.

\subsubsection[Multithreaded contexts]{Multithreaded contexts}\label{multithreaded-contexts}

The C++11 Standard Library provides several utilities and abstractions related
to multi\-threading. The \lstinline!std::thread! class is a portable
wrapper for a platform-specific thread handle provided by the
operating system. When constructing an \lstinline!std::thread! object
with a \emcppsgloss{callable object}, a new thread
invoking that \emcppsgloss{callable object} will be spawned.   Prior to destroying such
\lstinline!std::thread! objects it is necessary to invoke the
\lstinline!join! member function on the thread object, which will block until
the background thread of execution completes invoking its \emcppsgloss{callable object}.

This threading facility from the standard library can be used
with our earlier \lstinline!Logger! example from \intrarefsimple{logger-example},
to concurrently attempt to access the \lstinline!getLogger! function:

\begin{emcppshiddenlisting}[emcppsbatch={e6,e15}]
struct Logger {
  Logger(const char *f) {}
};
template <typename T>
void operator<<(Logger&,const T&);
Logger& getLogger();
\end{emcppshiddenlisting}
\begin{emcppslisting}[emcppsbatch=e6]
#include <thread>  // (ù{\codeincomments{std::thread}}ù)

void useLogger() { getLogger() << "example"; }  // concurrently called function

int main()
{
    std::thread t0(&useLogger);
    std::thread t1(&useLogger);
        // Spawn two new threads, each of which invokes (ù{\codeincomments{useLogger}}ù).

    // ...

    t0.join();  // Wait for (ù{\codeincomments{t0}}ù) to complete execution.
    t1.join();  // Wait for (ù{\codeincomments{t1}}ù) to complete execution.

    return 0;
}
\end{emcppslisting}

Such use prior to the C++11 thread-safety guarantees (with pre-C++11 threading libraries) could
have led to a \emcppsgloss[data races]{data race} during the initialization of
\lstinline!localLogger!, which was defined as a local \lstinline!static!
object in \lstinline!getLogger!.   This \emcppsgloss{undefined behavior} might have resulted in
invoking the constructor of \lstinline!localLogger! multiple times, returning from
\lstinline!localLogger! \pagebreak%%%% 
 before that constructor had actually been completed, or
any other form of misbehavior over which the developer has no control.

\noindent As of C++11, the example above has no \emcppsgloss{data races} provided that \linebreak%%%%%%
\lstinline!Logger::operator<<(const!~\lstinline!char*)! is designed properly
for multithreaded use, even though the
\lstinline!Logger::Logger(const!~\lstinline!char*!~\lstinline!logFilePath)!
constructor (i.e., the one used to configure the singleton instance of
the logger) is not. That is to say, the implicit \emcppsgloss{critical section}
that is guarded by the compiler includes evaluation of the
initializer, which is why a recursive call to initialize a
function-scope \lstinline!static! variable is \emcppsgloss{undefined behavior} and is
likely to result in deadlock; see \intraref{potential-pitfalls-functionstatic}{dangerous-recursive-initialization}.
Such use of function-scope \lstinline!static!s, however, is not foolproof;
see \intraref{potential-pitfalls-functionstatic}{depending-on-order-of-destruction-of-local-objects-after-main-returns}.

The destruction of \emcppsgloss[function scope]{function-scope} \lstinline!static! objects is
and always has been guaranteed to be safe \emph{provided} (1) no threads
are running after returning from \lstinline!main! and (2)
\emcppsgloss[function scope]{function-scope} \lstinline!static! objects do not depend on each
other during destruction; see \intraref{potential-pitfalls-functionstatic}{depending-on-order-of-destruction-of-local-objects-after-main-returns}.

\subsection[Use Cases]{Use Cases}\label{use-cases-functionstatic}

\subsubsection[Meyers Singleton]{Meyers Singleton}\label{meyers-singleton}

The guarantees surrounding access across \emcppsgloss[translation unit (TU)]{translation units} to runtime
initialized objects at file or namespace scope are few and weak ---
especially when that access might occur prior to entering \lstinline!main!.
Consider a library component, \lstinline!libcomp!, that defines a
file-scope \lstinline!static! singleton, \lstinline!globalS!, that is
initialized at run time:

\begin{emcppshiddenlisting}[emcppsbatch=e7]
// --- Replace
struct S { /*... */ };
struct S {
  bool isInitialized();
};
// --- End
\end{emcppshiddenlisting}
\begin{emcppslisting}[emcppsbatch=e7]
// libcomp.h:
#ifndef INCLUDED_LIBCOMP
#define INCLUDED_LIBCOMP

struct S { /*... */ };
S& getGlobalS();  // access to global singleton object of type (ù{\codeincomments{S}}ù)

#endif
\end{emcppslisting}
\vspace*{2ex} % for some extra space

%%% ask if these are supposed to be one.  JMB: They should be separate

\begin{emcppslisting}[emcppsbatch=e7]
// libcomp.cpp:
#include <libcomp.h>

static S globalS;
S& getGlobalS() { return globalS; }  // access into this translation unit
\end{emcppslisting}
\newpage%%%%%

\noindent The interface in the \lstinline!libcomp.h! file comprises the definition of
\lstinline!S! along with the declaration of an accessor function,
\lstinline!getGlobalS!. Code outside the \lstinline!libcomp.cpp! file can access the singleton object \lstinline!globalS! only by calling the free function \lstinline!getGlobalS()!. Now consider the 
 \lstinline!main.cpp! 
  file in the example below, which  implements
\lstinline!main! and also makes use of \lstinline!globalS! prior to entering
\lstinline!main!:


\begin{emcppslisting}[emcppsbatch=e7]
// main.cpp:
#include <cassert>    // standard C (ù{\codeincomments{assert}}ù) macro
#include <libcomp.h>  // (ù{\codeincomments{getGlobalS()}}ù)

bool globalInitFlag = getGlobalS().isInitialized();

int main()
{
    assert(globalInitFlag);   // Bug, or at least potentially so
    return 0;
}
\end{emcppslisting}

\noindent Depending on the compiler or the link line, the call
initializing \lstinline!globalInitFlag! may occur and return \emph{prior}
to the initialization of  \lstinline!globalS!.  C++ does not guarantee that objects at
file or namespace scope in separate \emcppsgloss[translation unit (TU)]{translation units} will be
initialized just because a function located within that \emcppsgloss[translation unit (TU)]{translation unit}
happens to be called.

An effective pattern for helping to ensure that a nonlocal object
\emph{is} initialized before it is used from a separate \emcppsgloss[translation unit (TU)]{translation unit}
--- especially when that use might occur prior to entering \lstinline!main!
--- is simply to move the \lstinline!static! object from file or namespace
scope to the scope of the function accessing it, making it a
function-scope \lstinline!static! instead:

\begin{emcppshiddenlisting}[emcppsbatch=e8]
struct S {
};
\end{emcppshiddenlisting}
\begin{emcppslisting}[emcppsbatch=e8]
S& getGlobalS()  // access into this translation unit
{
    static S globalS;  // singleton is now function-scope (ù{\codeincomments{static}}ù)
    return globalS;
}
\end{emcppslisting}

\noindent Commonly known as the \emcppsgloss[meyers singleton]{Meyers Singleton}, for the
author Scott Meyers who popularized it, this pattern ensures that the
singleton object will \emph{necessarily} be initialized on the first
call to the accessor function that envelopes it, irrespective of when
and where that call is made. Moreover, that singleton object will also
live past the end of \lstinline!main!. The \emcppsgloss[meyers singleton]{Meyers Singleton} pattern also gives us a chance to catch and respond to
exceptions thrown when constructing the \lstinline!static! object, rather
than immediately terminating the program, as would be the case if declared as a
\lstinline!static! global variable. Much more importantly, however, since C++11, the
\emcppsgloss[meyers singleton]{Meyers Singleton} pattern automatically inherits the benefits of
effortless race-free initialization of \emph{reusable} program-wide
singleton objects.  The \emcppsgloss[meyers singleton]{Meyers Singleton} can be safely used 
both in the programs where the singleton initialization might happen before \lstinline!main!
and those where it might happen after additional threads have already been started.
\newpage%%%%

As discussed in \intrarefsimple{description-functionstatic}, the augmentation of a
thread-safety guarantee for the runtime initialization of
\emcppsgloss[function scope]{function-scope} \lstinline!static! objects in C++11 minimizes the
effort required to create a thread-safe singleton.  Note that, prior to C++11,
the simple function-scope \lstinline!static!
implementation would not be safe if concurrent threads were trying to
initialize the logger; see \intraref{appendix-functionstatic}{appendix:-c++03-double-checked-lock-pattern}.

The \emcppsgloss[meyers singleton]{Meyers Singleton} is also seen in a slightly different form
where the singleton type's constructor is made \lstinline!private! to
prevent more than just the one singleton object from being created:

%%%% note ~ in code
\begin{emcppslisting}[emcppsbatch=e10]
class Logger
{
private:
    Logger(const char* logFilePath);  // configures the singleton
    ~Logger();                        // suppresses copy construction too

public:
    static Logger& getInstance()
    {
        static Logger localLogger("log.txt");
        return localLogger;
    }
};
\end{emcppslisting}
\begin{emcppshiddenlisting}[emcppsbatch=e10]
template <typename T>
void operator<<(Logger&,const T&);
\end{emcppshiddenlisting}

\noindent This variant of the function-scope-\lstinline!static! singleton pattern
prevents users from manually creating rogue \lstinline!Logger! objects; the
only way to get one is to invoke the logger's \lstinline!static!
\lstinline!Logger::getInstance()! member function:

\begin{emcppslisting}[emcppsbatch=e10]
void client()
{
    Logger::getInstance() << "Hi";  // OK
    Logger myLogger("myLog.txt");   // Error, (ù{\codeincomments{Logger}}ù) constructor is (ù{\codeincomments{private}}ù).
}
\end{emcppslisting}

\noindent This formulation of the singleton pattern, however, conflates the type
of the singleton object with its use and purpose as a singleton. Once we
find a use of a singleton object, finding another and perhaps even a
third is not uncommon.  

Consider, for example, an application on an early
model of mobile phone where we want to refer to the phone's camera.
Let's presume that a \lstinline!Camera! class is a fairly involved and
sophisticated mechanism. Initially we use the variant of the Meyers
Singleton pattern where at most one \lstinline!Camera! object can be
present in the entire program. The next generation of the phone, however, turns out to have more than
one camera, say, a front \lstinline!Camera! and a back \lstinline!Camera!. Our
brittle design doesn't
admit the dual-singleton use of the same fundamental \lstinline!Camera!
type. A more finely factored solution would be to implement the
\lstinline!Camera! type\linebreak[4] 
 \mbox{separately} and then to provide a thin wrapper,
e.g., perhaps using the \emcppsgloss[strong typedef idiom]{strong-\lstinline!typedef!\linebreak[4]%%%%% 
\newpage
 idiom} (see \featureref{\locationa}{inheriting-constructors}), corresponding to each singleton\linebreak[4]%%%%% 
 use:


\begin{emcppshiddenlisting}[emcppsbatch=e11]
struct Camera {
  Camera();
};    
\end{emcppshiddenlisting}
\begin{emcppslisting}[emcppsbatch=e11]
class PrimaryCamera
{
private:
    Camera& d_camera_r;
    PrimaryCamera(Camera& camera)  // implicit constructor
      : d_camera_r(camera) { }  

public:
    static PrimaryCamera getInstance()
    {
        static Camera localCamera{/*...*/};
        return localCamera;
    }
};
\end{emcppslisting}

\noindent With this design, adding a second and even a third singleton that is
able to reuse the underlying \lstinline!Camera! mechanism is facilitated.

Although this function-scope-\lstinline!static! approach provides stronger guarantees than the file-scope-\lstinline!static! one, it does have its limitations. In
particular, when one global facility object, such as a logger, is used
in the destructor of another function-scope static object, the logger
object may possibly have already been destroyed when it is
used.\footnote{An amusing workaround, the so-called \emph{Phoenix
Singleton}, is proposed in \cite{alexandrescu01}, section~6.6, pp.~137--139.} One approach is to construct the logger object by
explicitly allocating it and never deleting it:

\begin{emcppshiddenlisting}[emcppsbatch={e12,e13,e16,e17}]
struct Logger {
  Logger(const char *f) {}
};
\end{emcppshiddenlisting}
\begin{emcppslisting}[emcppsbatch=e12]
Logger& getLogger()
{
    static Logger& l = *new Logger("log.txt");  // dynamically allocated
    return l;  // Return a reference to the logger (on the heap).
}
\end{emcppslisting}

\noindent A distinct advantage of this approach is that once an object is created, it
\emph{never} goes away before the process ends. The disadvantage is
that, for many classic and current profiling tools (e.g., \emph{Purify},
\emph{Coverity}), this intentionally never-freed dynamic allocation is
indistinguishable from a \emcppsgloss{memory leak}. The ultimate workaround
is to create the object itself in \lstinline!static! memory, in an
appropriately sized and aligned region of memory:

\begin{emcppslisting}[emcppsbatch=e13]
#include <new>  // placement (ù{\codeincomments{new}}ù)

Logger& getLogger()
{
    static std::aligned_storage<sizeof(Logger), alignof(Logger)>::type buf;
    static Logger& logger = *new(&buf) Logger("log.txt");  // allocate in place
    return logger;
}
\end{emcppslisting}

\noindent Note that any memory that the \lstinline!Logger! itself manages would still come from the global heap and be recognized as memory leaks.{\cprotect\footnote{If the global heap is to be entirely avoided,
we could leverage a polymorphic-allocator implementation such as
\lstinline!std::pmr! in C++17. We would first create a fixed-size array
of memory having \emcppsgloss{static storage duration}. Then we would
create a \lstinline!static! memory-allocation mechanism (e.g.,
\lstinline!std::pmr::monotonic_buffer_resource!). Next we would use
placement \lstinline!new! to construct the logger within the static
memory pool using our static allocation mechanism and supply that same
mechanism to the \lstinline!Logger! object so that it could get all its
  internal memory from that static pool as well; see \cite{lakos22}.}} 
  
In this final incarnation of a decidedly non-Meyers-Singleton pattern,
we first reserve a block of memory of sufficient size and the correct
alignment for \lstinline!Logger! using \linebreak[4]%%%%%
\lstinline!std::aligned_storage!. Next
we use that storage in conjunction with placement \lstinline!new! to create
the logger directly in that static memory. Notice that this allocation
is not from the dynamic store, so typical profiling tools will not track
and will not provide a false warning when we fail to destroy this object
at program termination time. Now we can return a reference to the logger
object embedded safely in static memory knowing that it will be there
until application exit.

\subsection[Potential Pitfalls]{Potential Pitfalls}\label{potential-pitfalls-functionstatic}

\subsubsection[\lstinline!static! storage duration objects are not guaranteed to be initialized]{{\SubsubsecCode static} storage duration objects are not guaranteed to be initialized}\label{static-storage duration-objects-are-not-guaranteed-to-be-initialized}

Despite C++11's guarantee that each individual function-scope
\lstinline!static! initialization will occur at most once and before control can reach
a point where the variable can be referenced, no analogous guarantees are made of nonlocal objects of \emcppsgloss{static storage duration}.
This makes any interdependency in the initialization of such objects, especially across
\emcppsgloss[translation unit (TU)]{translation units (TUs)}, an abundant source of insidious
errors.

Objects that undergo \emcppsgloss{constant initialization} have no such issue:
such objects will never be accessible at run time before having their
initial values. Objects that are not constant
initialized{\cprotect\footnote{C++20 added a new keyword,
\lstinline!constinit!, that can be placed on a variable declaration to
\emph{require} that the variable in question undergo constant
initialization and thus can never be accessed at run time prior to the
  start of its lifetime.}} will instead be \emcppsgloss{zero initialized}
until their constructors run, which itself might lead to \emcppsgloss{undefined behavior} that is not necessarily conspicuous.

\enlargethispage*{1ex}%%%%%
As a demonstration of what can happen when we depend on the relative
order of initialization of variables at file or
namespace scope used before \lstinline!main!, consider the
\emcppsgloss{cyclically dependent} pair of source files, \lstinline!a.cpp! and
\lstinline!b.cpp!:

\begin{emcppslisting}
// a.cpp:
extern int setB(int);  // declaration (only) of setter in other TU
int *a = new int;      // runtime initialization of file-scope variable
int setA(int i)        // Initialize (ù{\codeincomments{a}}ù); then (ù{\codeincomments{b}}ù).
{
    *a = i;            // Populate the allocated heap memory.
    setB(i);           // Invoke setter to populate the other one.
    return 0;          // Return successful status.
}
\end{emcppslisting}
\pagebreak%%%%%%%
\begin{emcppslisting}
// b.cpp:
int *b = new int;      // runtime initialization of file-scope variable
int setB(int i)        // Initialize (ù{\codeincomments{b}}ù).
{
    *b = i;            // Populate the allocated heap memory.
    return 0;          // Return successful status.
}

extern int setA(int);  // declaration (only) of setter in other TU
int x = setA(5);       // Initialize (ù{\codeincomments{a}}ù) and (ù{\codeincomments{b}}ù).
int main()             // main program entry point
{
    return 0;          // Return successful status.
}
\end{emcppslisting}

\noindent These two \emcppsgloss[translation unit (TU)]{translation units} will be initialized before \lstinline!main! is
entered in some order, but --- regardless of that order --- the program
in the example above will wind up dereferencing a null pointer
before entering \lstinline!main!:

\begin{lstlisting}[language=bash]
$ g++ a.cpp b.cpp main.cpp
$ ./a.out
  Segmentation fault (core dumped)
\end{lstlisting}

\noindent Suppose we were to instead move the file-scope \lstinline!static! pointers,
corresponding to both \lstinline!setA! and \lstinline!setB!, inside their
respective function bodies:

\begin{emcppslisting}
// a.cpp:
extern int setB(int);  // declaration (only) of setter in other TU
int setA(int i)        // Initialize this (ù{\codeincomments{static}}ù) variable; then that one.
{
    static int *p = new int;  // runtime init. of function-scope (ù{\codeincomments{static}}ù)
    *p = i;                   // Populate this (ù{\codeincomments{static}}ù)-owned heap memory.
    setB(i);                  // Invoke setter to populate the other one.
    return 0;                 // Return successful status.
}

// b.cpp: (make analogous changes)
\end{emcppslisting}

\noindent Now the program reliably executes without incident:

\begin{lstlisting}[language=bash]
$ g++ a.cpp b.cpp main.cpp
$ ./a.out
$
\end{lstlisting}

\noindent In other words, even though no order exists in which the \emcppsgloss[translation unit (TU)]{translation units}
as a whole could have been initialized prior to entering
\lstinline!main! such that the \emph{file}-scope variables
would be valid before they were used, by instead making them
\emph{function}-scope \lstinline!static!, we are able to guarantee that
each variable is itself initialized before it is used, regardless of
translation-unit-initialization order.

While on the surface it may seem as though local and nonlocal objects of \emcppsgloss{static storage duration}
are effectively interchangeable, this is clearly not the case.  Even when
clients cannot directly access the nonlocal object due to giving it
\emcppsgloss{internal linkage} by marking it \lstinline!static! or putting it in an \emcppsgloss{unnamed namespace},
the initialization behaviors make such objects behave very differently.

\subsubsection[Dangerous recursive initialization]{Dangerous recursive initialization}\label{dangerous-recursive-initialization}

As with all other initialization, control flow does not continue
\emph{past} the \emcppsgloss{definition} of a \lstinline!static! local object until
after the initialization is complete, making recursive \lstinline!static!
initialization --- or any initializer that might eventually call back to the same function --- dangerous:

\begin{emcppslisting}
int fz(int i)  // The behavior of the first call is undefined unless (ù{\codeincomments{i}}ù) is 0.
{
    static int dz = i ? fz(i - 1) : 0;  // Initialize recursively. (BAD IDEA)
    return dz;
}

int main()  // The program is ill-formed.
{
    int x = fz(5);  // Bug, e.g., due to possible deadlock
}
\end{emcppslisting}

\noindent In the example above, the second recursive call of \lstinline!fz!
to initialize \lstinline!dz! has \emcppsgloss{undefined behavior} because the control flow reached the
same definition again before the initialization of the
\lstinline!static! object was completed; hence, control flow cannot continue to the
\lstinline!return! statement in \lstinline!fz!. Given a likely implementation
with a nonrecursive mutex or similar lock, the program can potentially deadlock,
though many implementations provide better diagnostics with an exception or
assertion violation when this form of error is encountered.{\cprotect\footnote{Prior to standardization (see \cite{ellis90}, section~6.7, p.~92),
C++ allowed control to flow past a
\lstinline!static! function-scope variable even during a recursive call
made as part of the initialization of that variable. This would result
in the rest of such a function executing with a zero-initialized and
possibly partially constructed local object. Even modern compilers,
such as GCC with \lstinline!-fno-threadsafe-statics!, allow turning off
the locking and protection from concurrent initialization and
retaining some of the pre-C++98 behavior. This optional behavior is,
however, fraught with peril and unsupported in any standard version of
  C++.}}

\subsubsection[Subtleties with recursion]{Subtleties with recursion}\label{subtleties-with-recursion}

Even when not recursing within the initializer itself, the rule for the
initialization of \lstinline!static! objects at function
scope becomes more subtle for self-recursive functions.
Notably, the initialization happens based
on when flow of control first passes the variable definition and \emph{not}
based on the first invocation of the containing function.  Due to this, when a recursive call
happens in relation to the definition of a \lstinline!static! local variable
impacts which values might be used for the initialization:\pagebreak%%%%%


\begin{emcppshiddenlisting}[emcppsbatch=e2]
#include <cassert>  // standard C (ù{\codeincomments{assert}}ù) macro
\end{emcppshiddenlisting}
\begin{emcppslisting}[emcppsbatch=e2]
int fx(int i)  // self-recursive after creating function-(ù{\codeincomments{static}}ù) variable, (ù{\codeincomments{dx}}ù)
{
    static int dx = i;     // Create (ù{\codeincomments{dx}}ù) first.
    if (i) { fx(i - 1); }  // Recurse second.
    return dx;             // Return (ù{\codeincomments{dx}}ù) third.
}

int fy(int i)  // self-recursive before creating function-(ù{\codeincomments{static}}ù) variable, (ù{\codeincomments{dy}}ù)
{
    if (i) { fy(i - 1); }  // Recurse first.
    static int dy = i;     // Create (ù{\codeincomments{dy}}ù) second.
    return dy;             // Return (ù{\codeincomments{dy}}ù) third.
}

int main()
{
    int x = fx(5);  assert(x == 5);  // (ù{\codeincomments{dx}}ù) is initialized before recursion.
    int y = fy(5);  assert(y == 0);  // (ù{\codeincomments{dy}}ù) is initialized after recursion.
    return 0;
}
\end{emcppslisting}

\noindent If the self-recursion takes place \emph{after} the \lstinline!static!
variable is initialized (e.g., \lstinline!fx! in the example above), then the
\lstinline!static! object (e.g., \lstinline!dx!) is initialized on the
\emph{first} recursive call; if the recursion occurs \emph{before}
(e.g., \lstinline!fy! in the example above), the initialization (e.g., of \lstinline!dy!)
occurs on the \emph{last} recursive call.


\subsubsection[Depending on order-of-destruction of local objects after \lstinline!main! returns]{Depending on order-of-destruction of local objects after {\SubsubsecCode main} returns}\label{depending-on-order-of-destruction-of-local-objects-after-main-returns}

Within any given \emcppsgloss[translation unit (TU)]{translation unit}, the relative order of initialization
of objects at file or namespace scope having \emcppsgloss{static storage duration} is
well defined and predictable. As soon as we have a way to reference an
object outside of the current \emcppsgloss[translation unit (TU)]{translation unit}, before \lstinline!main! is
entered, we are at risk of using the object before it has been
initialized. Provided the initialization itself is not cyclic in nature,
we can make use of function-scope \lstinline!static! objects (see \intraref{use-cases-functionstatic}{meyers-singleton}) to
ensure that no such uninitialized use occurs, even across \emcppsgloss[translation unit (TU)]{translation units}
before \lstinline!main! is entered. The relative order of destruction
of such function-scope \lstinline!static! variables --- even when they
reside within the same \emcppsgloss[translation unit (TU)]{translation unit} --- is not clearly known at compile time,
as it will be the reverse of the order in which they are initialized, and
reliance on such order can easily lead to \emcppsgloss{undefined behavior} in
practice.

This specific problem occurs when a \lstinline!static! object at file,
namespace, or function scope uses (or might use) in its destructor
another \lstinline!static! object that is either (1) at file or namespace
scope and resides in a separate \emcppsgloss[translation unit (TU)]{translation unit} or (2) any other
function-scope \lstinline!static! object (i.e., including one in the same
\emcppsgloss[translation unit (TU)]{translation unit}). For example, suppose we have implemented a low-level
logging facility as a Meyers Singleton:

\newpage%%%%%%
\begin{emcppslisting}[emcppsbatch=e15]
Logger& getLogger()
{
    static Logger local("log.txt");
    return local;
}
\end{emcppslisting}

\noindent Now suppose we implement a higher-level file-manager type that depends
on the function-scope \lstinline!static! logger object:

%%% note ~ in code
\begin{emcppslisting}[emcppsbatch=e15]
struct FileManager
{
    FileManager()
    {
         getLogger() << "Starting up file manager...";
         // ...
    }

    ~FileManager()
    {
        getLogger() << "Shutting down file manager...";
        // ...
    }
};
\end{emcppslisting}

\noindent Now, consider a Meyers Singleton implementation for
\lstinline!FileManager!:

\begin{emcppslisting}[emcppsbatch=e15]
FileManager& getFileManager()
{
    static FileManager fileManager;
    return fileManager;
}
\end{emcppslisting}

\noindent Whether \lstinline!getLogger! or \lstinline!getFileManager! is called first
doesn't really matter; if\linebreak[4]%%%%%
 \lstinline!getFileManager! is called first, the
logger will be initialized as part of \lstinline!FileManager!'s
constructor. However, whether the \lstinline!Logger! or
\lstinline!FileManager! object is destroyed first \emph{is} important:

\begin{itemize}
\item{If the \lstinline!FileManager! object is destroyed prior to the \lstinline!Logger! object, the program will have well-defined behavior.}
\item{Otherwise, the program will have \emcppsgloss{undefined behavior} because the destructor of\linebreak[4]%%%%%
 \lstinline!FileManager! will invoke \lstinline!getLogger!, which will now return a reference to a previously destroyed object.}
\end{itemize}

%\enlargethispage*{-1ex}%%%%%%%
Logging in the constructor of the \lstinline!FileManager! 
makes it certain that the logger's function-local
\lstinline!static! will be initialized before that of the file manager;
hence, since destruction occurs in reverse relative order of creation,
the logger's function-local \lstinline!static! will be destroyed after\linebreak[4]
\newpage%%%%% 
that
of the file manager. But suppose that \lstinline!FileManager! didn't always
log at construction and was created before anything else logged. In that
case, we have no reason to think that the logger would be around for the
\lstinline!FileManager! to log during its destruction after \lstinline!main!.

In the case of low-level, widely used facilities, such as a logger, a
conventional Meyers Singleton is contraindicated. The two most common
alternatives discussed at the end of \intraref{use-cases-functionstatic}{meyers-singleton} involve never ending the lifetime of
the mechanism at all. It is worth noting that truly global objects ---
such as \lstinline!cout!, \lstinline!cerr!, and \lstinline!clog! --- from the
Standard \lstinline!iostream! Library are typically not implemented using
conventional methods and are in fact treated specially by the runtime system.

\subsection[Annoyances]{Annoyances}\label{annoyances}

\subsubsection[Overhead in single-threaded applications]{Overhead in single-threaded applications}\label{overhead-in-single-threaded-applications}

A single-threaded application invoking a function containing a
\emcppsgloss[function scope]{function-scope} \emcppsgloss{static storage duration} variable might
have unnecessary synchronization overhead, such as an \emcppsgloss{atomic}
load operation. For example, consider a program that accesses
a simple  \emcppsgloss[meyers singleton]{Meyers Singleton} for a user-defined type with a
\emcppsgloss[user provided]{user-provided} default constructor:

\begin{emcppslisting}
struct S  // user-defined type
{
    S() { }  // inline default constructor
};

S& getS()  // free function returning (ù{\codeincomments{local}}ù) object
{
    static S local;  // function-scope local object
    return local;
}

int main()
{
    getS();    // Initialize the file-scope (ù{\codeincomments{static}}ù) singleton.
    return 0;  // successful status
}
\end{emcppslisting}

%\enlargethispage*{3ex}
\noindent Although it is clearly visible to the compiler that \lstinline!getS()! is
invoked by only one thread, the generated assembly instructions might
still contain \emcppsgloss{atomic} operations or other forms of
synchronization and the call to \lstinline!getS()! might not be 
inlined.{\cprotect\footnote{Both GCC 10.x and Clang 10.x, using the
\lstinline!-Ofast! optimization level, generate assembly instructions for
an \emcppsgloss{acquire/release memory barrier} and fail to inline the call
to \lstinline!getS!. Using \lstinline!-fno-threadsafe-statics! reduces the
number of operations performed considerably but still does not lead to
the compilers' inlining of the function call. Both popular compilers
will, however, reduce the program to just two x86 assembly
instructions if the \emcppsgloss[user provided]{user-provided} constructor of \lstinline!S!
  is either removed or defaulted
   (see \featureref{\locationa}{defaulted-special-member-functions});
  doing so will turn \lstinline!S! into a \emcppsgloss[trivially constructible]{trivially-constructible}
  type, implying that no code needs to be executed during
  initialization:\pagebreak[4]

  \begin{lstlisting}[language=bash,style=footcodeplain]
  xor eax, eax  ; zero out 'eax' register
  ret           ; return from 'main'
  \end{lstlisting}

\noindent A sufficiently smart compiler might, however, not generate
  synchronization code in a single-threaded context or else provide a
  flag to control this behavior.
  }}

\subsection[See Also]{See Also}\label{see-also}

\hspace*{\fill}

\subsection[Further Reading]{Further Reading}\label{further-reading}

\begin{itemize}
\item{\cite{meyers04a}}
\item{For an in-depth discussion of the difficulties of implementing double-checked locking in C++03, see \cite{meyers04b}.}
\item{\cite{stroustrup20}}
\item{For a discussion of  the Singleton pattern and a variety of implementations in C++03, see Chapter~6 of \cite{alexandrescu01}.}
\end{itemize}

\subsection[Appendix]{Appendix}\label{appendix-functionstatic}

\subsubsection[C++03 double-checked-lock pattern]{C++03 double-checked-lock pattern}\label{appendix:-c++03-double-checked-lock-pattern}

Prior to the introduction of the \emcppsgloss[function scope]{function-scope} \lstinline!static!
object initialization guarantees discussed in \intrarefsimple{description-functionstatic}, preventing multiple initializations of \lstinline!static! objects
and use before initialization of those same objects was still needed.
Guarding access using a \lstinline!mutex! was often a significant performance cost, so
using the unreliable, double-checked lock pattern was often attempted to
avoid the overhead:

\enlargethispage*{2ex}
\begin{emcppshiddenlisting}[emcppsbatch=e16]
#include <mutex>  // (ù{\codeincomments{std::mutex}}ù), (ù{\codeincomments{std::lock\_guard}}ù)
\end{emcppshiddenlisting}
\begin{emcppslisting}[emcppsbatch=e16]
Logger& getInstance()
{
    static Logger* volatile loggerPtr = 0;  // hack, used to simulate *atomics*

    if (!loggerPtr)  // Does the logger need to be initialized?
    {
        static std::mutex m;
        std::lock_guard<std::mutex> guard(m);  // Lock the mutex.

        if (!loggerPtr)  // We are first, as the logger is still uninitialized.
        {
            static Logger logger("log.txt");
            loggerPtr = &logger;
        }
    }                    // Either way, the lock (ù{\codeincomments{guard}}ù) unlocks the mutex here.

    return *loggerPtr;
}
\end{emcppslisting}

\noindent In this example, we are using a \lstinline!volatile! pointer
as a partial
substitute for an atomic variable, a non-portable solution that is not correct
in standard C++ but has historically been moderately effective.  The C++11 standard
library does, however, provide the \lstinline!<atomic>! header, which is a far superior alternative,
and many implementations have historically
provided extensions to support atomic types even prior to C++11.  Where available, compiler extensions are typically preferable over home-grown solutions.

In addition to
being difficult to write, this decidedly complex workaround would often
prove unreliable. The problem is that, even though the logic appears
sound, architectural changes in widely used CPUs allowed for the CPU
itself to optimize and reorder the sequence of instructions. Without
additional support, the hardware would not see the dependency that the
second test of \lstinline!loggerPtr! has on the locking behavior of the
mutex and would do the read of \lstinline!loggedPtr! prior to acquiring the
lock. This reordering of instructions
would then allow multiple threads to acquire the lock while each thinking the
\lstinline!static! variable still needs to be initialized.

To solve this subtle issue, concurrency library authors are expected to
issue ordering hints such as \emcppsgloss{fences} and \emcppsgloss{barriers}. A
well-implemented threading library would provide atomics equivalent to
the modern \lstinline!std::atomic! that would issue the correct
instructions when accessed and modified. The C++11 Standard makes the
compiler aware of these concerns and provides portable \emph{atomics}
and support for threading that enables users to handle such issues
correctly. The above \lstinline!getInstance! function could be corrected by
changing the type of \lstinline!loggerPtr! to
\lstinline!std::atomic<Logger*>!. Prior to C++11, despite being
complicated, the same function would reliably implement the Meyers
Singleton in C++03 on contemporary hardware.

%\enlargethispage*{2ex} %%%%%% just to get that final } on this page so external reviewers don't freak out
So the final recommended solution for portable thread-safe
initialization in modern C++ is to simply let the compiler do the work
and to use the simplest implementation that gets the job done, e.g., a
Meyers Singleton (see \intraref{use-cases-functionstatic}{meyers-singleton}):

\begin{emcppslisting}[emcppsbatch=e17]
Logger& getInstance()
{
    static Logger logger("log.txt");
    return logger;
}
\end{emcppslisting}


 

\newpage
%\section[Local Types '11]{Local/Unnamed Types as Template Arguments}\label{local-types-as-template-arguments}
% 5 Feb 2021, copyedits in and proofed
% 8 Feb 2021 JMB, compilation fixes
% 22 March 2021, in FPPs

\emcppsFeature{
    short={Local Types '11},
    long={Local/Unnamed Types as Template Arguments},
}{local-types-as-template-arguments}
\setcounter{table}{0}
\setcounter{footnote}{0}
\setcounter{lstlisting}{0}
%\section[Local Types '11]{Local/Unnamed Types as Template Arguments}\label{local-types-as-template-arguments}



C++11 allows function-scope and unnamed types to be used as template arguments.

\subsection[Description]{Description}\label{description}

Historically, types without \emcppsgloss{linkage} (i.e., local and unnamed
types) were forbidden as template arguments due to implementability
concerns using the compiler technology available at that
time.{\cprotect\footnote{\cite{narkive}}} Modern C++ lifts this
restriction, making use of local or unnamed types consistent with
nonlocal, named ones, thereby obviating the need to gratuitously name
or enlarge the scope of a type.

\begin{emcppslisting}
template <typename T>
void f(T) { };            // function template

template <typename T>
class C { };              // class template

struct { } obj;           // object (ù{\codeincomments{obj}}ù) of unnamed C++ type

void g()
{
    struct S { };         // local type

    f(S());               // OK in C++11; was error in C++03
    f(obj);               // OK in C++11; was error in C++03

    C<S>             cs;  // OK in C++11; was error in C++03
    C<decltype(obj)> co;  // OK in C++11; was error in C++03
}
\end{emcppslisting}

\noindent Notice that we have used the \emph{\lstinline!decltype!
keyword} (see \featureref{\locationa}{decltype}) to extract the unnamed type of the object \lstinline!obj!.

These new relaxed rules for template arguments are essential to the
ergonomics of \emcppsgloss{lambda expressions} (see \featureref{\locationc}{lambda}), as such types are both
unnamed and local in typical usage: 

\newpage%%%%
\begin{emcppslisting}
#include <algorithm>  // (ù{\codeincomments{std::sort}}ù)
#include <string>     // (ù{\codeincomments{std::string}}ù)
#include <vector>     // (ù{\codeincomments{std::vector}}ù)

struct Person { std::string d_name; };

void sortByName(std::vector<Person>& people)
{
    std::sort(people.begin(), people.end(),
              [](const Person& lhs, const Person& rhs)
              {
                  return lhs.d_name < rhs.d_name;
              });
}
\end{emcppslisting}

\noindent In the example above, the lambda expression passed to the
\lstinline!std::sort! algorithm is a local unnamed type, and the algorithm
itself is a function template.

\subsection[Use Cases]{Use Cases}\label{use-cases}

\subsubsection[Encapsulating a type within a function]{Encapsulating a type within a function}\label{encapsulating-a-type-within-a-function}

Limiting the scope and visibility of an \emcppsgloss{entity} to the body of a
function actively prevents its direct use, even when the function body
is exposed widely --- say, as an \lstinline!inline! function or function
template defined within a header file.

Consider, for instance, an implementation of Dijkstra's algorithm that
uses a local type to keep track of metadata for each vertex in the input
graph:

\begin{emcppshiddenlisting}[emcppsbatch=e1]
// dijkstra.h:
#include <vector>  // (ù{\codeincomments{std::vector}}ù)
struct Vertex {};
struct VertexMetadata {};
struct Graph { std::size_t numNodes() const; };
\end{emcppshiddenlisting}
\begin{emcppslisting}[emcppsbatch=e1]
// dijkstra.h:

#include <vector>  // (ù{\codeincomments{std::vector}}ù)

inline int dijkstra(std::vector<Vertex>* path, const Graph& graph)
{
    struct VertexMetadata         // implementation-specific helper class
    {
        int  d_distanceFromSource;
        bool d_inShortestPath;
    };

    std::vector<VertexMetadata> vertexMetadata(graph.numNodes());
        // standard vector of local (ù{\codeincomments{VertexMetadata}}ù) objects -- one per vertex

    // ... (body of algorithm)
}
\end{emcppslisting}

\newpage%%%%
\noindent Defining \lstinline!VertexMetadata! outside of the body of
\lstinline!dijkstra! --- e.g., to comply with C++03 restrictions --- would
make that implementation-specific helper class directly accessible to
anyone including the \lstinline!dijkstra.h! header file. As Hyrum's
law{\cprotect\footnote{``With a sufficient number
of users of an API, it does not matter what you promise in the
contract: all observable behaviors of your system will be depended on
  by somebody'': see \cite{wight}.}} suggests, if the
imple\-men\-tation-specific \lstinline!VertexMetadata! detail is defined
outside the function body, it is to be expected that some user somewhere
will depend on it in its current form, making it problematic, if not
impossible, to change.\footnote{The C++20 \emph{modules} facility enables the
encapsulation of helper types (such as \lstinline!metadata! in the \mbox{\lstinline!dijkstra.h!} example on this page) used in
the implementation of other locally defined types or functions, even
when the helper types appear at namespace scope within the module.} Conversely, encapsulating
the type within the function body avoids unintended use by clients,
while improving human cognition by colocating the definition of the
type with its sole purpose.{\cprotect\footnote{For a detailed
  discussion of malleable  versus stable software, see \cite{{lakos20}},
  section 0.5, pp.~29--43.}}

\subsubsection[Instantiating templates with local function objects as type arguments]{Instantiating templates with local function objects as type arguments}\label{instantiating-templates-with-local-function-objects-as-type-arguments}

Suppose that we have a program that makes wide use of an aggregate data
type, \lstinline!City!:

\begin{emcppslisting}[emcppsbatch={e2,e3}]
#include <algorithm>  // (ù{\codeincomments{std::copy}}ù)
#include <iostream>   // (ù{\codeincomments{std::ostream}}ù)
#include <iterator>   // (ù{\codeincomments{std::ostream\_iterator}}ù)
#include <set>        // (ù{\codeincomments{std::set}}ù)       
#include <string>     // (ù{\codeincomments{std::string}}ù)
#include <vector>     // (ù{\codeincomments{std::vector}}ù)

struct City
{
    int         d_uniqueId;
    std::string d_name;
};
std::ostream& operator<<(std::ostream& stream,
                         const City&   object);
\end{emcppslisting}

\noindent Consider now the task of writing a function to print unique elements of
an\linebreak[4] \lstinline!std::vector<City>!, ordered by name:

\newpage%%%%%
\begin{emcppslisting}[emcppsbatch=e2]
void printUniqueCitiesOrderedByName(const std::vector<City>& cities)
{
    struct OrderByName
    {
        bool operator()(const City& lhs, const City& rhs) const
        {
            return lhs.d_name < rhs.d_name;
                // increasing order (subject to change)
        }
    };

    const std::set<City, OrderByName> tmp(cities.begin(), cities.end());

    std::copy(tmp.begin(), tmp.end(),
              std::ostream_iterator<City>(std::cout, "\n"));
}
\end{emcppslisting}

\noindent Absent reasons to make the \lstinline!OrderByName!
function object more generally available, rendering its definition
alongside the one place where it is used --- i.e., directly within
function scope --- again enforces and readily communicates its tightly
encapsulated (and therefore \emph{malleable}) status.

As an aside, note that using a lambda (see \featureref{\locationc}{lambda}) in such scenario requires using \lstinline!decltype! and passing the closure to the set's constructor:
\begin{emcppslisting}[emcppsbatch=e3]
void printUniqueCitiesOrderedByName(const std::vector<City>& cities)
{
    auto compare = [](const City& lhs, const City& rhs) {
        return lhs.d_name < rhs.d_name;
    };
    const std::set<City, decltype(compare)>
        tmp(cities.begin(), cities.end(), compare);
}
\end{emcppslisting}

\noindent We discuss the topic of lambda expressions further in the very next section; see \intrarefnopage{configuring-algorithms-via-lambda-expressions}.

\subsubsection[Configuring algorithms via lambda expressions]{Configuring algorithms via lambda expressions}\label{configuring-algorithms-via-lambda-expressions}

Suppose we are representing a 3D environment using a \emph{scene
graph} and managing
the graph's nodes via an \lstinline!std::vector! of \lstinline!SceneNode!
objects (a \emph{scene graph} data structure, commonly
used in computer games and 3D-modeling software, represents the
  logical and spatial hierarchy of objects in a scene). Our \lstinline!SceneNode! class supports a variety of
\lstinline!const! member functions used to query its status (e.g.,
\lstinline!isDirty! and \lstinline!isNew!). Our task is to implement a
\emcppsgloss{predicate function}, \lstinline!mustRecalculateGeometry!, that
returns \lstinline!true! if and only if at least one of the nodes is either
``dirty'' or ``new.''

These days, we might reasonably elect to implement this functionality
using the C++11 standard algorithm
\lstinline!std::any_of!{\cprotect\footnote{\cite{cpprefa}}}:

\begin{emcppslisting}[emcppsbatch=e4]
template <typename InputIterator, typename UnaryPredicate>
bool any_of(InputIterator first, InputIterator last, UnaryPredicate pred);
    // Return (ù{\codeincomments{true}}ù) if any of the elements in the range satisfies (ù{\codeincomments{pred}}ù).
\end{emcppslisting}

\noindent Prior to C++11, however, using a function template, such as
\lstinline!any_of!, would have required a separate function or
function object (defined \emph{outside} of the scope of the function):

\begin{emcppshiddenlisting}[emcppsbatch={e4,e5}]
#include <vector>     // (ù{\codeincomments{std::vector}}ù)
struct SceneNode {
    bool isDirty() const;
    bool isNew() const;
};
template <typename InputIterator, typename UnaryPredicate>
bool any_of(InputIterator first, InputIterator last, UnaryPredicate pred)
{
   // this needs an implementation
   for (; first != last; ++first) {
      if (pred(*first)) {
          return true;
      }
  }
  return false;
}
\end{emcppshiddenlisting}
\begin{emcppslisting}[emcppsbatch=e4]
// C++03 (obsolete)
namespace {

struct IsNodeDirtyOrNew
{
    bool operator()(const SceneNode& node) const
    {
        return node.isDirty() || node.isNew();
    }
};

}  // close unnamed namespace

bool mustRecalculateGeometry(const std::vector<SceneNode>& nodes)
{
    return any_of(nodes.begin(), nodes.end(), IsNodeDirtyOrNew());
}
\end{emcppslisting}

\noindent Because unnamed types can serve as arguments to this function template, we can also employ a lambda expression instead of a function object that would be required in C++03:

\begin{emcppslisting}[emcppsbatch=e5]
#include <algorithm> // (ù{\codeincomments{'std::any\_of'}}ù)
bool mustRecalculateGeometry(const std::vector<SceneNode>& nodes)
{
    return std::any_of(nodes.begin(),             // start of range
                       nodes.end(),               // end of range
                       [](const SceneNode& node)  // lambda expression
                       {
                           return node.isDirty() || node.isNew();
                       }
                      );
}
\end{emcppslisting}

\noindent By creating a \emcppsgloss{closure} of unnamed type via a lambda
expression, unnecessary boilerplate, excessive scope, and even local
symbol visibility are avoided.

\subsection[Potential Pitfalls]{Potential Pitfalls}\label{potential-pitfalls}

\hspace{\fill}

\subsection[Annoyances]{Annoyances}\label{annoyances}

\hspace{\fill}

\subsection[See Also]{See Also}\label{see-also}

\begin{itemize}
\item{\seealsoref{decltype}{\seealsolocationa}describes how developers may query the type of any expression or entity, including objects with unnamed types.}
\item{\seealsoref{lambda}{\seealsolocationc}provides strong practical motivation for the relaxations discussed here.}
\end{itemize}

\subsection[Further Reading]{Further Reading}\label{further-reading}

\hspace{\fill}





\newpage
%\section[{\tt long} {\tt long}]{The {\SecCode long} {\SecCode long} ({\large $\boldsymbol{\geq}$}64 bits) Integral Type}\label{long-long}
% 5 Feb 2021, copyedits in and proofed
% 8 Feb 2021 JMB, compilation fixes
% 22 March 2021, in FPPs

\emcppsFeature{
    short={\lstinline!long!~\lstinline!long!},
    tocshort={{\TOCCode long}~{\TOCCode long}},
    long={The {\SecCode long}~{\SecCode long} ({\large $\boldsymbol{\geq}$}64 bits) Integral Type},
    toclong={The \lstinline!long!~\lstinline!long! ({\small$\boldsymbol{\geq}$}64 bits) Integral Type},
    rhshort={{\RHCode long}~{\RHCode long}},
}{long-long}
\setcounter{table}{0}
\setcounter{footnote}{0}
\setcounter{lstlisting}{0}
%\section[{\tt long} {\tt long}]{The {\SecCode long} {\SecCode long} ({\large $\boldsymbol{\geq}$}64 bits) Integral Type}\label{long-long}


\lstinline!long!~\lstinline!long! is a \emcppsgloss{fundamental integral type} guaranteed to have at least 64 bits on all
platforms.

\subsection[Description]{Description}\label{description}

The \emcppsgloss{integral type} \lstinline!long!~\lstinline!long! and its companion
type \lstinline!unsigned!~\lstinline!long!~\lstinline!long! are the only two
\emcppsgloss{fundamental integral types} in C++ that are guaranteed to have
at least 64~bits on all conforming platforms{\cprotect\footnote{\lstinline!long!~\lstinline!long!
has been available in C since the C99 standard, and many C++ compilers
  supported it as an extension prior to C++11.}}:

\begin{emcppslisting}
#include <climits>  // (ù{\codeincomments{CHAR\_BIT}}ù) (a.k.a.~8, see below)

long long          a;  // (ù{\codeincomments{sizeof(a) * CHAR\_BIT >= 64}}ù)
unsigned long long b;  // (ù{\codeincomments{sizeof(b) * CHAR\_BIT >= 64}}ù)

static_assert(sizeof(a) == sizeof(b), "");
    // I.e., (ù{\codeincomments{a}}ù) and (ù{\codeincomments{b}}ù) necessarily have the same size in every program.
\end{emcppslisting}

\noindent On all conforming platforms, \lstinline!CHAR_BIT! --- the number of bits
in a byte --- is at least~8 and,~on virtually all commonly available
commercial platforms today, is exactly~8.

The corresponding integer-literal suffixes indicating type
\lstinline!long!~\lstinline!long! are \lstinline!ll! and \lstinline!LL!; for
\lstinline!unsigned!~\lstinline!long!~\lstinline!long!, any of eight alternatives
are accepted: \lstinline!ull!, \lstinline!ULL!, \lstinline!uLL!, \lstinline!Ull!,
\lstinline!llu!, \lstinline!LLU!, \lstinline!LLu!,
\lstinline!llU!:

\begin{emcppslisting}
auto i = 0LL;   // (ù{\codeincomments{long long}}ù), (ù{\codeincomments{sizeof(i) * CHAR\_BIT >= 64}}ù)
auto u = 0uLL;  // (ù{\codeincomments{unsigned long long}}ù), (ù{\codeincomments{sizeof(u) * CHAR\_BIT >= 64}}ù)
\end{emcppslisting}

\noindent Note that \lstinline!long!~\lstinline!long!
and \lstinline!unsigned!~\lstinline!long!~\lstinline!long! are also candidates
for the type of an integer literal having a large enough value. As an
example, the type of the literal \lstinline!2147483648! (one more than
the upper bound of a 32-bit integer) is likely to be
  \lstinline!long!~\lstinline!long! on a 32-bit platform. For a historical perspective on how integral types have evolved (and
continue to evolve) over time, see \intraref{longlong-appendix}{historical-perspective-on-the-evolution-of-use-of-fundamental-integral-types}.
%{\it\titleref{historical-perspective-on-the-evolution-of-use-of-fundamental-integral-types}} on page~\pageref{historical-perspective-on-the-evolution-of-use-of-fundamental-integral-types}.

\subsection[Use Cases]{Use Cases}\label{use-cases}

\subsubsection[Storing values that won't safely fit in 32 bits]{Storing values that won{\subsubapos}t safely fit in 32 bits}\label{when-your-pedestrian-four-byte-int-might-not-cut-it}\label{Storing-values-that-won't-safely-fit-in-32-bits}

For many quantities that need to be represented as an integral value in a program, plain \lstinline!int! is a natural choice. For example, this could be the case for years of a person's age, score in a ten-pin bowling game, or number of stories in a building. For efficient storage in a \lstinline!class!
or \lstinline!struct!, however, we may well decide to represent such
quantities more compactly using a \lstinline!short! or \lstinline!char!; see
  also the aliases found in C++11's \lstinline!<cstdint>!.

Sometimes the size of the virtual address space for the
underlying architecture itself dictates how large an integer you will
need. For example, on a 64-bit platform, specifying the \emph{distance} between two pointers
into a contiguous array or the size of the array itself could well exceed the size of an \lstinline!int! or
\lstinline!unsigned!~\lstinline!int!, respectively. Using either
\lstinline!long!~\lstinline!long! or
\lstinline!unsigned!~\lstinline!long!~\lstinline!long! here would, however, not
be indicated as the respective platform-dependent integer types
(\lstinline!typedef!s) \lstinline!std::ptrdiff_t! and \lstinline!std::size_t!
are provided expressly for such use, and avoid wasting space where it
cannot be used by the underlying hardware.

Occasionally, however, the decision of whether to use an \lstinline!int! is
neither platform dependent nor clear cut, in which case using an
\lstinline!int! is almost certainly a bad idea. As part of a financial
library, suppose we were asked to provide a function that, given a date,
returns the number of shares of some particular stock, identified by its
security id (\lstinline!SecId!) traded on the New York Stock Exchange
(NYSE).\footnote{There are more than 3,200 listed symbols on the NYSE. Composite daily volume of NYSE-listed securities across all exchanges  ranges from 3.5 to 6~billion shares, with a high reached in March 2020 of more than 9~billion shares.} Since the average daily volume of even the most heavily traded stocks (roughly 70 million shares) appears to
be well under the maximum value a signed \lstinline!int! supports (more
than 2 billion on our production platforms), we might at first think to write the function
to return \lstinline!int!:

\begin{emcppshiddenlisting}[emcppsbatch={e1,e2,e3,e4,e5,e6}]
struct SecId {};
\end{emcppshiddenlisting}
\begin{emcppslisting}[emcppsbatch=e1]
int volYMD(SecId equity, int year, int month, int day);  // (1) BAD IDEA
\end{emcppslisting}

\noindent One obvious problem with this interface is that the daily fluctuations
in turbulent times might exceed the maximum value representable by a
32-bit \lstinline!int!, which, unless detected internally, would result in
\emcppsgloss{signed integer overflow}, which is both \emcppsgloss{undefined
behavior} and potentially a pervasive defect enabling avenues of deliberate attack from outside sources.{\cprotect\footnote{For an overview of integer overflow in C++, see \cite{ballman}. For a more focused discussion of secure coding in CPP using CERT standards, see \cite{seacord13}, Chapter~5, ``Integer Security," pp. 225--307.}} What's more, the growth rate of some companies,
especially technology startups, has
been at times seemingly exponential. To gain an extra
insurance factor of two, we might opt to replace the return type
\lstinline!int! with an \lstinline!unsigned!~\lstinline!int!:

\begin{emcppslisting}[emcppsbatch=e2]
unsigned volYMD(SecId stock, int year, int month, int day);  // (2) BAD IDEA!
\end{emcppslisting}

\noindent Use of an \lstinline!unsigned!~\lstinline!int!, however, simply delays the
inevitable as the number of shares being traded is almost certainly
going to grow over time.

Furthermore, the algebra for unsigned quantities is entirely different
from what one would normally expect from an \lstinline!int!. For example, if we
were to try to express the day-over-day change in volume by
subtracting two calls to this function and if the number of shares
traded were to have decreased, then the \lstinline!unsigned!~\lstinline!int!
difference would wrap, and the result would be a typically large,
erroneous value. Because integer
literals are themselves of type \lstinline!int! and not \lstinline!unsigned!,
comparing an unsigned value with a negative signed one does not
typically go well; hence, many compilers will warn when the two types
  are mixed, which itself is problematic.

If we happen to be on a 64-bit platform, we might choose to return a
\lstinline!long!:

\begin{emcppslisting}[emcppsbatch=e3]
long volYMD(SecId stock, int year, int month, int day);  // (3) NOT A GOOD IDEA
\end{emcppslisting}

\noindent The problems using \lstinline!long! as the return type are that it (1) is
not yet generally considered a \emcppsgloss{vocabulary type} (see
\intraref{longlong-appendix}{historical-perspective-on-the-evolution-of-use-of-fundamental-integral-types}),
%{\it\titleref{longlong-appendix}} on page~\pageref{longlong-appendix}),
and (2) would reduce portability (see
\intraref{longlong-potential-pitfalls}{relying-on-the-relative-sizes-of-int,-long,-and-long-long}).
%{\it\titleref{longlong-potential-pitfalls}: \titleref{relying-on-the-relative-sizes-of-int,-long,-and-long-long}} on page~\pageref{relying-on-the-relative-sizes-of-int,-long,-and-long-long}).

Prior to C++11, we might have considered returning a \lstinline!double!:

\begin{emcppslisting}[emcppsbatch=e4]
double volYMD(SecId stock, int year, int month, int day);  // (4) OK
\end{emcppslisting}

\noindent At least with \lstinline!double! we know that we will have sufficient precision (53 bits) to express
integers accurately into the quadrillions, which will certainly cover us
for any foreseeable future. The main drawback is that \lstinline!double!
doesn't properly describe the nature of the type that we are returning
--- i.e., a whole integer number of shares --- and so its algebra,
although not as dubious as \lstinline!unsigned!~\lstinline!int!, isn't ideal
either.

With the advent of C++11, we might consider using one of the type aliases
in \lstinline!<cstdint>!:

\begin{emcppshiddenlisting}[emcppsbatch={e5}]
#include <cstdint>  // (ù{\codeincomments{std::int64\_t}}ù)
\end{emcppshiddenlisting}
\begin{emcppslisting}[emcppsbatch=e5]
std::int64_t volYMD(SecId stock, int year, int month, int day);  // (4) OK
\end{emcppslisting}

\noindent This choice addresses most of the issues discussed above except that,
instead of being a specific C++ type, it is a platform-dependent alias
that is likely to be a \lstinline!long! on a 64-bit platform and almost
certainly a \lstinline!long!~\lstinline!long! on a 32-bit one. Such exact size
requirements are often necessary for packing data in structures and
arrays but are not as useful when reasoning about them in the
interfaces of functions where having a common set of fundamental
\emcppsgloss{vocabulary types} becomes much more important (e.g., for
interoperability).

All of this leads us to our final alternative, \lstinline!long!~\lstinline!long!:

\begin{emcppslisting}[emcppsbatch=e6]
long long volYMD(SecId stock, int year, int month, int day);  // (5) GOOD IDEA
\end{emcppslisting}

\noindent In addition to being a signed fundamental integral type of sufficient
capacity on all platforms, \lstinline!long!~\lstinline!long! is the same C++
type \emph{relative} to other C++ types on all platforms.

\subsection[Potential Pitfalls]{Potential Pitfalls}\label{longlong-potential-pitfalls}

\subsubsection[Relying on the relative sizes of \lstinline!int!, \lstinline!long!, and \lstinline!long!~\lstinline!long!]{Relying on the relative sizes of {\SubsubsecCode int}, {\SubsubsecCode long}, and {\SubsubsecCode long}~{\SubsubsecCode long}}\label{relying-on-the-relative-sizes-of-int,-long,-and-long-long}

As discussed at some length in \intraref{longlong-appendix}{historical-perspective-on-the-evolution-of-use-of-fundamental-integral-types},
%{\it\titleref{longlong-appendix}} on page~\pageref{longlong-appendix},
the
fundamental integral types have historically been a moving target. On
older, 32-bit platforms, a \lstinline!long! was often 32~bits and, \lstinline!long!~\lstinline!long!, which was nonstandard prior to C++11, or its platform-dependent equivalent was needed to ensure that 64~bits were
available. When the correctness of code depends on either
\lstinline!sizeof(int)!~\lstinline!<!~\lstinline!sizeof(long)! or
\lstinline!sizeof(long)!~\lstinline!<!~\lstinline!sizeof(long!~\lstinline!long)!,
portability is needlessly restricted. Relying instead on only the
guaranteed{\cprotect\footnote{Due to the unfathomable amount of software
that would stop working if an \lstinline!int! were ever anything but
exactly \emph{four} bytes, we --- along with the late Richard Stevens
  of Unix fame (see \cite{stevens93}, section~2.5.1., pp.~31--32, specifically row~6, column~4, Figure~2.2, p.~32) --- are prepared
  to \emph{guarantee} that it will never become as large as a
  \lstinline!long!~\lstinline!long! for any general-purpose computer.}}
property that\linebreak[4]
\lstinline!sizeof(int)!~\lstinline!<!~\lstinline!sizeof(long!~\lstinline!long)!}
avoids such portability issues since the relative sizes of the \lstinline!long! and \lstinline!long!~\lstinline!long! integral
types continue to evolve.

When precise control of size \emph{in the implementation} (as opposed to
in the interface) matters, consider using one of the standard signed
(\texttt{int{\itshape n}\_t}) or unsigned (\texttt{uint{\itshape n}\_t}) integer
%%%%% cannot escape within lstinline, so those two terms are set \texttt{}
aliases provided, since C++11, in \lstinline!<cstdint>! and
summarized here in Table~\ref{longlong-table1}.%\newpage

\begin{table}[h!]
\begin{center}
\begin{threeparttable}
\caption{Useful {\bfseries\lstinline!typedef!}s found in {\bfseries \lstinline!<cstdint>!} (since
C++11)}\label{longlong-table1} \vspace{1.5ex}
{\small \begin{tabular}{c|c|c}\thickhline
\rowcolor[gray]{.9} {\sffamily\bfseries Exact Size (optional)\tnote{a}} & {\sffamily\bfseries Fastest integral type } & {\sffamily\bfseries Smallest integer type }\\
\rowcolor[gray]{.9}  & {\sffamily\bfseries having at least N bits} & {\sffamily\bfseries  having at least N bits}\\\hline
\multicolumn{1}{r|}{\lstinline!int8_t!} & \multicolumn{1}{r|}{\lstinline!int_fast8_t!} & \multicolumn{1}{r}{\lstinline!int_least8_t!}\\ \hline
\multicolumn{1}{r|}{\lstinline!int16_t!} & \multicolumn{1}{r|}{\lstinline!int_fast16_t!} & \multicolumn{1}{r}{\lstinline!int_least16_t!}\\ \hline
\multicolumn{1}{r|}{\lstinline!int32_t!} & \multicolumn{1}{r|}{\lstinline!int_fast32_t!} & \multicolumn{1}{r}{\lstinline!int_least32_t!}\\ \hline
\multicolumn{1}{r|}{\lstinline!int64_t!} & \multicolumn{1}{r|}{\lstinline!int_fast64_t!} & \multicolumn{1}{r}{\lstinline!int_least64_t!}\\ \medhline
\multicolumn{1}{r|}{\lstinline!uint8_t!} & \multicolumn{1}{r|}{\lstinline!uint_fast8_t!} & \multicolumn{1}{r}{\lstinline!uint_least8_t!}\\ \hline
\multicolumn{1}{r|}{\lstinline!uint16_t!\tnote{a}} & \multicolumn{1}{r|}{\lstinline!uint_fast16_t!} & \multicolumn{1}{r}{\lstinline!uint_least16_t!}\\ \hline
\multicolumn{1}{r|}{\lstinline!uint32_t!} & \multicolumn{1}{r|}{\lstinline!uint_fast32_t!} & \multicolumn{1}{r}{\lstinline!uint_least32_t!}\\ \hline
\multicolumn{1}{r|}{\lstinline!uint64_t!} & \multicolumn{1}{r|}{\lstinline!uint_fast64_t!} & \multicolumn{1}{r}{\lstinline!uint_least64_t!}\\ \thickhline
\end{tabular} }
\begin{tablenotes}{\footnotesize
\item[a]{The compiler doesn't need to fabricate the exact-width type if the target platform doesn't support it.\\
Note: Also see \lstinline!intmax_t!, the maximum width integer type, which might be different from all of the above.}
}
\end{tablenotes}
\end{threeparttable}
    \end{center}
\end{table}

\subsection[See Also]{See Also}\label{see-also}

\begin{itemize}
\item{\seealsoref{binary-literals}{\seealsolocationb}explains how programmers can specify binary constants directly in the source code; large binary values might only fit in a \lstinline!long!~\lstinline!long! or even \lstinline!unsigned!~\lstinline!long!~\lstinline!long!.}
\item{\seealsoref{digit-separators}{\seealsolocationb}describes visually separating digits of large \lstinline!long!~\lstinline!long! literals.}
\end{itemize}

\subsection[Further Reading]{Further Reading}\label{further-reading}

\begin{itemize}
\item{\cite{adamczyk05}}
\end{itemize}

\subsection[Appendix]{Appendix}\label{longlong-appendix}

\subsubsection[Historical perspective on the evolution of use of fundamental integral types]{Historical perspective on the evolution of use of fundamental integral types}
\label{historical-perspective-on-the-evolution-of-use-of-fundamental-integral-types}

The designers of C got it right back in 1972 when
they created a portable \lstinline!int! type that could act as a bridge
from a single-word (16-bit) integer, \lstinline!short!, to a double-word
(32-bit) integer, \lstinline!long!. Just by using \lstinline!int!, one would
get the optimal space versus speed trade-off as the 32-bit computer \emph{word}
was on its way to becoming the norm. As an example, the Motorola
68000 series (c.~1979) was a hybrid \emph{CISC} architecture employing
a 32-bit instruction set with 32-bit registers and a 32-bit external
data bus; internally, however, it used only 16-bit ALUs and a 16-bit
  data bus.

During the late 1980s and into the 1990s, the word size of the machine and
the size of an \lstinline!int! were synonymous. Some of
the earlier mainframe computers, such as IBM 701 (c.~1954), had a word
size of 36~characters (1) to allow accurate representation of a signed
10-digit decimal number or (2) to hold up to six 6-bit characters.
Smaller computers, such as Digital Equipment Corporation's
PDP-1/PDP-9/PDP-15 used 18-bit words (so a double word held 36~bits);
memory addressing, however, was limited to just 12--18~bits (i.e., a
maximum 4K--256K 18-bit words of \emph{DRAM}). With the standardization
of 7-bit ASCII (c.~1967), its adoption throughout the 1970s, and its most
recent update (c.~1986), the common typical notion of character size
moved from 6 to 7~bits. Some early conforming implementations (of C)
would choose to set \lstinline!CHAR_BIT! to 9 to allow two characters
per half word. (On some early vector-processing computers,
\lstinline!CHAR_BIT! is 32, making every type, including a
\lstinline!char!, at least a 32-bit quantity.) As double-precision floating-point calculations --- enabled by type \lstinline!double! and supported by floating-point coprocessors --- became typical in the scientific community, machine architectures
naturally evolved from 9-, 18-, and 36-bit words to the familiar 8-,
16-, 32-, and now 64-bit addressable integer words we have today.
Apart from embedded systems and \emph{DSPs}, a \lstinline!char! is now
almost universally considered to be exactly 8~bits. Instead of
scrupulously and actively using \lstinline!CHAR_BIT! for the number of bits
in a \lstinline!char!, consider statically asserting it instead:

\begin{emcppshiddenlisting}[emcppsbatch=e7]
#include <climits>  // (ù{\codeincomments{CHAR\_BIT}}ù)
\end{emcppshiddenlisting}
\begin{emcppslisting}[emcppsbatch=e7]
static_assert(CHAR_BIT == 8, "A (ù{\codeincomments{char}}ù) is not 8-bits on this CrAzY platform!");
\end{emcppslisting}

\noindent As cost of \emph{main memory} was decreasing exponentially
throughout the final two decades of the 20th
century,{\cprotect\footnote{Moore's law (c.~1965) --- the observation
that the number of transistors in densely packed integrated circuits
(e.g., DRAM) grows exponentially over time, doubling every 1--2~years
or so --- held for nearly a half century, until finally saturating in
  the 2010s.}} the need for a much larger \emph{virtual address space}
quickly followed. Intel began its work on 64-bit architectures in the
early 1990s and realized one a decade later. As we progressed into
the 2000s, the common notion of \emph{word size} --- i.e., the width (in
bits) of typical registers within the CPU itself --- began to shift from
``the size of an \lstinline!int!'' to ``the size of a simple (nonmember)
pointer type,'' e.g., \lstinline!8!~\lstinline!*!~\lstinline!sizeof(void*)!, on the host platform. By this time, 16-bit \lstinline!int! types --- like 16-bit architectures for \emcppsgloss[general purpose machines]{general-purpose machines} (i.e., excluding \emcppsgloss{embedded systems}) --- were long gone but a \lstinline!long!~\lstinline!int! was still expected to be 32~bits on a 32-bit platform. Embedded systems are designed specifically to work with high-performance hardware, such as digital-signal processors (DSPs). Sadly,
\lstinline!long! was often used (improperly) to hold an address; hence,
the size of \lstinline!long! is associated with a de facto need (due to immeasurable amounts of legacy code) to remain in lockstep with pointer size.

Something new was needed to mean at least 64 bits on all platforms.
Enter \lstinline!long!~\lstinline!long!. We have now come full circle. On
64-bit platforms, an \lstinline!int! is still 4~bytes, but a \lstinline!long!
is now --- for practical reasons --- typically 8~bytes unless requested
explicitly{\cprotect\footnote{On 64-bit systems, \lstinline!sizeof(long)!
is typically 8~bytes. Compiling with the \lstinline!-m32! flag on either
GCC or Clang emulates compiling on a 32-bit platform:
\lstinline!sizeof(long)! is likely to be~4, while
  \lstinline!sizeof(long!~\lstinline!long)! remains~8.}} to be otherwise. To ensure portability until 32-bit machines go the way of 16-bit ones,
we have \lstinline!long!~\lstinline!long! to (1) provide a common \emph{vocabulary type}, (2)
make our intent clear, and (3) avoid the portability issue for
at least the next decade or two; still, see \intraref{longlong-potential-pitfalls}{relying-on-the-relative-sizes-of-int,-long,-and-long-long}
%{\it\titleref{longlong-potential-pitfalls}: \titleref{relying-on-the-relative-sizes-of-int,-long,-and-long-long}} on page~\pageref{relying-on-the-relative-sizes-of-int,-long,-and-long-long}
for some alternative ideas.





\newpage
%\section[{\tt noreturn}]{The {\SecCode [[noreturn]]} Attribute}\label{the-noreturn-attribute}
%%% 15 Jan: copyedits in and proofed
%% 27 Jan, code checked and replaced as needed by Josh
% 01 Feb, 2021 JMB - listing that does compile marked as non-compiling

\emcppsFeature{
    toclong={The \lstinline![[noreturn]]! Attribute},
    tocshort={\TOCCode noreturn},
    short={\lstinline!noreturn!},
    long={The {\SecCode [[noreturn]]} Attribute},
    rhshort={{\RHCode noreturn}},
}{the-noreturn-attribute}
%\section[{\tt noreturn}]{The {\SecCode [[noreturn]]} Attribute\sectionmark{\RHCode noreturn}}\label{the-noreturn-attribute}\sectionmark{\RHCode noreturn}
\setcounter{table}{0}
\setcounter{footnote}{0}
\setcounter{lstlisting}{0}


The \lstinline![[noreturn]]! attribute promises that the function to which
it pertains never \nobreak{returns}.

\subsection[Description]{Description}\label{description}

The presence of the standard \lstinline![[noreturn]]! attribute as part of
a function declaration informs both the compiler and human readers that
such a function never returns control flow to the caller:

%begin{emcppslisting}
%[[noreturn]] void f()
%{
%    throw;
%}
%\end{emcppslisting}
\begin{emcppslisting}[emcppsbatch=e1]
[[noreturn]] void f()
{
    throw 1;
}
\end{emcppslisting}


\noindent The \lstinline![[noreturn]]! attribute is not part of a function's type
and is also, therefore, not part of the type of a function pointer. Applying \lstinline![[noreturn]]! to a function
pointer is not an error, though doing so has no actual effect in standard C++; see
%{\it\titleref{noreturn-potential-pitfalls}: \titleref{misuse-of-[[noreturn]]-on-function-pointers}} on page~\pageref{misuse-of-[[noreturn]]-on-function-pointers}
\intraref{noreturn-potential-pitfalls}{misuse-of-[[noreturn]]-on-function-pointers}. Using it on a pointer might have
benefits for external tooling, code expressiveness, and future language
evolution:

\begin{emcppslisting}[emcppsbatch=e1,emcppsignore={Only works on gcc -fpermissive}]
void (*fp [[noreturn]])() = f;
\end{emcppslisting}


\subsection[Use Cases]{Use Cases}\label{use-cases}

\subsubsection[Better compiler diagnostics]{Better compiler diagnostics}\label{better-compiler-diagnostics}

Consider the task of creating an assertion handler that, when invoked,
always aborts execution of the program after printing some useful
information about the source of the assertion. Since this specific
handler will never return because it unconditionally invokes a \linebreak[4]%%%%%
\mbox{\lstinline![[noreturn]]std::abort!} function, it is a viable candidate for
\lstinline![[noreturn]]!:

\begin{emcppshiddenlisting}[emcppsbatch=e2]
#include <cstdlib>  // (ù{\codeincomments{std::abort}}ù)
struct LogError {};
constexpr LogError LOG_ERROR; // Quick and dirty syntactic simulation of logger
template <typename T>
const LogError& operator<<(const LogError& lhs, const T&) { return lhs; }
\end{emcppshiddenlisting}
\begin{emcppslisting}[emcppsbatch=e2]
[[noreturn]] void abortingAssertionHandler(const char* filename, int line)
{
    LOG_ERROR << "Assertion fired at " << filename << ':' << line;
    std::abort();
}
\end{emcppslisting}


\noindent The additional information provided by the attribute will allow a
compiler to warn if it determines that a code path in the
function would allow it to return normally:

%[[noreturn]] void abortingAssertionHandler(const char* filename, int line)
%{
%    if (filename)  // just being safe, but see "Further Reading," below
%    {
%        LOG_ERROR << "Assertion fired at " << filename << ':' << line;
%        std::abort();
%    }
%}  // compile-time warning made possible
%\end{emcppslisting}
\begin{emcppshiddenlisting}[emcppsbatch=e3]
#include <cstdlib>   // (ù{\codeincomments{std::abort}}ù)
#include <iostream>  // (ù{\codeincomments{std::cout}}ù)
struct LogError {};
constexpr LogError LOG_ERROR;
template <typename T>
const LogError& operator<<(const LogError& lhs, const T&) { return lhs; }
\end{emcppshiddenlisting}
\begin{emcppslisting}[emcppsbatch=e3]
[[noreturn]] void abortingAssertionHandler(const char* filename, int line)
{
    if (filename)
    {
        LOG_ERROR << "Assertion fired at " << filename << ':' << line;
        std::abort();
    }
}  // compile-time warning made possible
\end{emcppslisting}



\noindent This information can also be used to warn in case unreachable code is
present after\linebreak[4] \lstinline!abortingAssertionHandler! is invoked:

\begin{emcppslisting}[emcppsbatch=e3]
int main()
{
    // ...
    abortingAssertionHandler("main.cpp", __LINE__);
    std::cout << "We got here.\n";  // compile-time warning made possible
    // ...
}
\end{emcppslisting}


\noindent Note that this warning is made possible by decorating just the
declaration of the handler function --- i.e., even if the definition of
the function is not visible in the current translation unit.

\subsubsection[Improved runtime performance]{Improved runtime performance}\label{improved-runtime-performance}

If the compiler knows that it is going to invoke a function
that is guaranteed not to return, the compiler is within its rights to optimize
that function by removing what it can now determine to be dead code. As
an example, consider a utility component, \lstinline!util!, that defines a
function, \lstinline!throwBadAlloc!, that is used to \romeogloss{insulate} the
throwing of an \lstinline!std::bad\_alloc! exception in what would
otherwise be template code fully exposed to clients:

%begin{emcppslisting}
%// util.h:
%[[noreturn]] void throwBadAlloc();
%
%// util.cpp:
%#include <util.h>  // (ù{\codeincomments{[[noreturn]] void throwBadAlloc()}}ù)
%
%void throwBadAlloc()  // This redeclaration is also (ù{\codeincomments{[[noreturn]]}}ù).
%{
%    throw std::bad_alloc();
%}
%\end{emcppslisting}
\begin{emcppslisting}[emcppsbatch=e4]
// util.h:
[[noreturn]] void throwBadAlloc();

// util.cpp:
#include <util.h>  // (ù{\codeincomments{[[noreturn]] void throwBadAlloc()}}ù)

#include <new>    // (ù{\codeincomments{std::bad\_alloc}}ù)

void throwBadAlloc()  // This redeclaration is also (ù{\codeincomments{[[noreturn]]}}ù).
{
    throw std::bad_alloc();
}
\end{emcppslisting}


\noindent The compiler is within its
rights to elide code that is rendered unreachable by the call to the
\lstinline!throwBadAlloc! function due to the function being decorated with the
\lstinline![[noreturn]]! attribute on its declaration:

\begin{emcppslisting}[emcppsbatch=e4]
// client.cpp:
#include <util.h>  // [[noreturn]] void throwBadAlloc()

void client()
{
    // ...
    throwBadAlloc();
    // ... (Everything below this line can be optimized away.)
}
\end{emcppslisting}


\noindent Notice that even though \lstinline![[noreturn]]! appeared only on the first
declaration --- that in the \lstinline!util.h! header --- the \lstinline![[noreturn]]!
attribute carries over to the redeclaration used in the\linebreak[4] \lstinline!throwBadAlloc!
function's definition because the header was included in the
corresponding \lstinline!.cpp! file.

\subsection[Potential Pitfalls]{Potential Pitfalls}\label{noreturn-potential-pitfalls}

\subsubsection[\tt{[[noreturn]]} can inadvertently break an otherwise working program]{{\SubsubsecCode [[noreturn]]} can inadvertently break an otherwise working program}\label{[[noreturn]]-can-inadvertently-break-an-otherwise-working-program}

Unlike many attributes, using \lstinline![[noreturn]]! \emph{can} alter
the semantics of a well-formed program, potentially introducing a
runtime defect and/or making the program ill-formed. If a function that
can potentially return is decorated with \lstinline![[noreturn]]! and then,
in the course of executing a program, it ever does return, that behavior
is \romeogloss{undefined}.

Consider a \lstinline!printAndExit! function whose role is to print a fatal
error message before aborting the program:

\begin{emcppshiddenlisting}[emcppsbatch=e5]
#include <iostream>  // (ù{\codeincomments{std::cout}}ù)
#include <cassert>   // standard C (ù{\codeincomments{assert}}ù) macro
\end{emcppshiddenlisting}
\begin{emcppslisting}[emcppsbatch=e5]
[[noreturn]] void printAndExit()
{
    std::cout << "Fatal error. Exiting the program.\n";
    assert(false);
}
\end{emcppslisting}

\noindent The programmer chose to (sloppily) implement termination by using an
assertion, which would not be incorporated into a program compiled with
the preprocessor definition \lstinline!NDEBUG! active, and thus
\lstinline!printAndExit! would return normally in such a build mode. If the
compiler of the client is informed that function will not return, the compiler is
free to optimize accordingly. If the function then does return, any
number of hard-to-diagnose defects (e.g., due to incorrectly elided
code) might materialize as a consequence of the ensuing
\romeogloss{undefined behavior}. Furthermore, if a function is declared \lstinline![[noreturn]]! in some
translation units within a program but not in others, that program is
\romeogloss{ill-formed, no diagnostic required (IFNDR)}.

\subsubsection[Misuse of {\tt[[noreturn]]} on function pointers]{Misuse of {\SubsubsecCode [[noreturn]]} on function pointers}\label{misuse-of-[[noreturn]]-on-function-pointers}

Although the \lstinline![[noreturn]]! attribute is permitted to syntactically appertain
to a function pointer for the benefit of external tools,
it has no effect in standard C++; fortunately, most compilers will issue a warning:

%begin{emcppslisting}
%void (*fp [[noreturn]])();  // not supported by standard C++ (will likely warn)
%\end{emcppslisting}
\begin{emcppslisting}[emcppsbatch=e6,emcppsignore={Only works on gcc -fpermissive}]
void (*fp [[noreturn]])();  // no effect in standard C++ (will likely warn)
\end{emcppslisting}

\noindent What's more, assigning the address of a function
that is not decorated with \lstinline![[noreturn]]! to an otherwise
suitable function pointer that is so decorated is perfectly fine:

\begin{emcppslisting}[emcppsbatch=e6,emcppsignore={Only works on gcc -fpermissive}]
void f() { return; };  // function that always returns

void g()
{
    fp = f;  // (ù{\codeincomments{[[noreturn]]}}ù) on (ù{\codeincomments{fp}}ù) is silently ignored.
}
\end{emcppslisting}
% code line changed from
% fp = f;  // "OK" -- that (ù{\codeincomments{fp}}ù) is (ù{\codeincomments{[[noreturn]]}}ù) is silently ignored
% to
% fp = f;  // (ù{\codeincomments{[[noreturn]]}}ù) on (ù{\codeincomments{fp}}ù) is silently ignored.
% per JD Garcia and Slava

\noindent Any reliance on \lstinline![[noreturn]]! to have any effect in standard C++
when applied to other than a function's declaration is misguided.

\subsection[Annoyances]{Annoyances}\label{annoyances}

\hspace*{\fill}

% per email, Nina would omit.
%The contradictory declaration:
%begin{emcppslisting}
%[[noreturn]] int f();
%\end{emcppslisting}
%compiles without error or even a warning.

\subsection[See Also]{See Also}\label{see-also}

\begin{itemize}
\item{%``\titleref{attributes}" on page~\pageref{attributes} —
\seealsoref{attributes}{\locationa}\lstinline![[noreturn]]! is a built-in attribute that follows the general syntax and placement rules of C++ attributes.}
\end{itemize}

\subsection[Further Reading]{Further Reading}\label{further-reading}

\begin{itemize}
\item{\cite{svoboda10}}
\item{\cite{sutter12}}
\end{itemize}






\newpage
%\section[{\tt nullptr}]{The Null-Pointer-Literal Keyword}\label{null-pointer-literal-(nullptr)}
% 13 Jan: copyedits in and proofed
%% 27 Jan, code checked and replaced as needed by Josh

The keyword \lstinline!nullptr! unambiguously denotes the
null-pointer-value literal.

\subsection[Description]{Description}\label{description}

The \lstinline!nullptr! keyword is a \romeovalue{prvalue} (pure rvalue) of type
\lstinline!std::nullptr_t! representing the implementation-defined
bit pattern corresponding to a \romeogloss{null address} on the host platform;
\lstinline!nullptr! and other values of type \lstinline!std::nullptr_t!, along with the integer literal \lstinline!0! and the macro \lstinline!NULL!, can be converted implicitly to any pointer or pointer-to-member type:

\begin{emcppslisting}
#include <cstddef> // (ù{\codeincomments{NULL}}ù)
int data;  // nonmember data

int *pi0 = &data;    // Initialize with non-null address.
int *pi1 = nullptr;  // Initialize with null address.
int *pi2 = NULL;     //  "          "    "    "
int *pi3 = 0;        //  "          "    "    "

double f(int x);  // nonmember function

double (*pf0)(int) = &f;       // Initialize with non-null address.
double (*pf1)(int) = nullptr;  // Initialize with null address.

struct S
{
    short d_data;    // member data
    float g(int y);  // member function
};

short S::*pmd0 = &S::d_data;  // Initialize with non-null address.
short S::*pmd1 = nullptr;     // Initialize with null address.

float (S::*pmf0)(int) = &S::g;    // Initialize with non-null address.
float (S::*pmf1)(int) = nullptr;  // Initialize with null address.
\end{emcppslisting}


\noindent Because \lstinline!std::nullptr_t! is its own distinct type, overloading on
it is possible:

\begin{emcppslisting}
#include <cstddef>  // (ù{\codeincomments{std::nullptr\_t}}ù)

void g(void*);           // (1)
void g(int);             // (2)
void g(std::nullptr_t);  // (3)

void f()
{
    char buf[] = "hello";
    g(buf);      // OK, (1) (ù{\codeincomments{void g(void*)}}ù)
    g(0);        // OK, (2) (ù{\codeincomments{void g(int)}}ù)
    g(nullptr);  // OK, (3) (ù{\codeincomments{void g(std::nullptr\_t)}}ù)
    g(NULL);     // Error, ambiguous --- (1), (2), or (3)
}
\end{emcppslisting}


\subsection[Use Cases]{Use Cases}\label{use-cases}

\subsubsection[Improvement of type safety]{Improvement of type safety}\label{improve-type-safety}

In pre-C++11 codebases, using the \lstinline!NULL!
macro was a common way of indicating, mostly to the human
reader, that the literal value the macro conveys is intended specifically to
represent a \emph{null address} rather than the literal \lstinline!int!
value \lstinline!0!. In the C Standard, the macro \lstinline!NULL! is
defined as an \romeogloss{implementation-defined} integral or \lstinline!void*!
constant. Unlike C, C++ forbids conversions from \lstinline!void*! to
arbitrary pointer types and instead, prior to C++11, defined
\lstinline!NULL! as an ``integral constant expression rvalue of
integer type that evaluates to zero''; any integer literal (e.g.,
\lstinline!0!, \lstinline!0L!, \lstinline!0U!, \lstinline!0LLU!) satisfies this
  criterion. From a type-safety perspective, its
implementation-defined definition, however,
makes using \lstinline!NULL! only marginally better suited than a raw
literal \lstinline!0! to represent a null pointer. It is worth noting that as of C++11, the definition of \lstinline!NULL! has been expanded to --- in theory --- permit \lstinline!nullptr! as a conforming definition; as of this writing, however, no major compiler vendors do so.{\cprotect\footnote{Both GCC and
Clang default to \lstinline!0L! (\lstinline!long!~\lstinline!int!), while MSVC
defaults to \lstinline!0! (\lstinline!int!). Such definitions are unlikely
to change since existing code could cease to compile or (possibly
  silently) present altered runtime behavior.}}

As just one specific illustration of the added type safety provided by
\lstinline!nullptr!, imagine that the coding standards of a large software company historically required that
values returned via output parameters (as opposed to a \lstinline!return!
statement) are always returned via pointer to a modifiable
object. Functions that return via argument typically do so to reserve the function’s return value to communicate status.{\cprotect\footnote{See \cite{lakos96}, section~9.1.11, pp.~621--628, specifically the \emph{Guideline} at the bottom of p.~621: ``Be consistent about returning values through arguments (e.g., avoid declaring non\lstinline!const! reference parameters)."}} A function in this codebase might ``zero" the output parameter's local pointer variable to indicate and ensure that nothing more is to be written. The function below illustrates three different ways of doing this:

\begin{emcppshiddenlisting}[emcppsbatch=e1]
#include <cstdlib>  // (ù{\codeincomments{NULL}}ù)
// --- Replace
    if (/*...*/)
    if (true)
// --- End
\end{emcppshiddenlisting}
\begin{emcppslisting}[emcppsbatch=e1]
int illustrativeFunction(int* x)   // pointer to modifiable integer
{
    // ...
    if (/*...*/)
    {
        x = 0;       // OK, Set pointer (ù{\codeincomments{x}}ù) to null address.
        x = NULL;    // OK, Set pointer (ù{\codeincomments{x}}ù) to null address.
        x = nullptr; // Bug, Set pointer (ù{\codeincomments{x}}ù) to null address.
    }
    // ...
    return 0;    // success
}
\end{emcppslisting}


Now suppose that the function signature is changed (e.g., due to a
change in coding standards in the organization) to accept a reference
instead of a pointer:

\begin{emcppshiddenlisting}[emcppsbatch=e2]
#include <cstdlib>  // (ù{\codeincomments{NULL}}ù)
// --- Replace
    if (/*...*/)
    if (true)
// --- End
\end{emcppshiddenlisting}
\begin{emcppslisting}[emcppsbatch=e2]
int illustrativeFunction(int& x)  // reference to modifiable integer
{
    // ...
    if (/*...*/)
    {
        x = 0;       // OK, always compiles; makes what (ù{\codeincomments{x}}ù) refers to 0
        x = NULL;    // OK, implementation-defined (might warn)
        x = nullptr; // Error, always a compile-time error
    }
    // ...
    return 0;    // SUCCESS
}
\end{emcppslisting}


As the example above demonstrates, how we represent the notion of a
null address matters:
\begin{enumerate}
\item{\lstinline!0! — Portable across all implementations but minimal type safety}
\item{\lstinline!NULL! — Implemented as a macro; added type safety (if any) is platform specific}
\item{\lstinline!nullptr! — Portable across all implementations and fully type-safe}
\end{enumerate}
Using \lstinline!nullptr! instead of \lstinline!0! or \lstinline!NULL! to denote
a null address maximizes type safety and readability, while avoiding
both macros and implementation-defined behavior.

\subsubsection[Disambiguation of \lstinline!(int)0! vs. \lstinline!(T*)0! during overload resolution]{Disambiguation of {\SubsubsecCode (int)0} vs. {\SubsubsecCode (T*)0} during overload resolution}\label{disambiguation-of-(int)-0-versus-(t*)-0-during-overload-resolution}

The platform-dependent nature of \lstinline!NULL! presents additional
challenges when used to call a function whose overloads differ only in
accepting a pointer or an integral type as the same positional argument, which might be the case, e.g., in a poorly designed third-party
library:

\begin{emcppshiddenlisting}[emcppsbatch=e3]
#include <cstdlib>  // (ù{\codeincomments{NULL}}ù)
\end{emcppshiddenlisting}
\begin{emcppslisting}[emcppsbatch=e3]
void uglyLibraryFunction(int* p);  // (1)
void uglyLibraryFunction(int  i);  // (2)
\end{emcppslisting}


\noindent Calling this function with the literal \lstinline!0! will always invoke
overload (2), but that might not always be what casual clients expect:

\begin{emcppslisting}[emcppsbatch=e3]
void f()
{
    uglyLibraryFunction(0);         // unambiguously invokes (2)
    uglyLibraryFunction((int*) 0);  // unambiguously invokes (1)
    uglyLibraryFunction(nullptr);   // unambiguously invokes (1)
    uglyLibraryFunction(NULL);      // Error, anything! (platform-defined)
    uglyLibraryFunction(0L);        // Error, ambiguous call (on all platforms)
    uglyLibraryFunction(0U);        // Error, ambiguous call (on all platforms)
}
\end{emcppslisting}


\noindent\lstinline!nullptr! is especially useful when such problematic overloads
are unavoidable because it obviates explicit
casts. (Note that explicitly casting \lstinline!0! to an
appropriately typed pointer --- other than \lstinline!void*! --- was at one
  time considered by some to be a best practice, especially in C.)

\subsubsection[Overloading for a literal null pointer]{Overloading for a literal null pointer}\label{overloading-for-a-literal-null-pointer}

Being a distinct type, \lstinline!std::nullptr_t! can itself participate
in an overload set:

\begin{emcppslisting}
#include  <cstddef> // (ù{\codeincomments{std::nullptr\_t}}ù)
void f(int* v);          // (1)
void f(std::nullptr_t);  // (2)

void g()
{
    int* ptr = nullptr;
    f(ptr);      // unaemmbiguously invokes (1)
    f(nullptr);  // unambiguously invokes (2)
}
\end{emcppslisting}


\noindent Given the relative ease with which a \lstinline!nullptr! can be
converted to a typed pointer having the same null-address value, such
overloads are dubious when used to control essential behavior.
Nonetheless, we can envision such use to, say, aid in compile-time
diagnostics when passing a \romeogloss{null address} would otherwise result in
a runtime error (see
%``\titleref{deleted-functions}" on page~\pageref{deleted-functions}
\featureref{\locationb}{deleted-functions}):

\begin{emcppshiddenlisting}[emcppsbatch=e4]
#include <cstdlib>  // (ù{\codeincomments{std::size\_t}}ù)
\end{emcppshiddenlisting}
\begin{emcppslisting}[emcppsbatch=e4]
std::size_t strlen(const char* s);
    // The behavior is undefined unless s is null-terminated.

std::size_t strlen(std::nullptr_t) = delete;
    // Function is not defined but still participates in overload resolution.
\end{emcppslisting}


\noindent Another arguably safe use of such an overload for a \lstinline!nullptr! is
to avoid a null-pointer check. However, for cases where the client knows the address is null at compile time, better ways typically exist for avoiding the (often
insignificant) overhead of testing for a null pointer at run time.

\subsection[Potential Pitfalls]{Potential Pitfalls}\label{potential-pitfalls}

\hspace*{\fill}

\subsection[Annoyances]{Annoyances}\label{annoyances}

\hspace*{\fill}

\subsection[See Also]{See Also}\label{see-also}

\hspace*{\fill}

\subsection[Further Reading]{Further Reading}\label{further-reading}

\begin{itemize}
\item{\cite{meyers15a}}
\end{itemize}



\newpage
%\section[{\tt override}]{The {\SecCode override} Member-Function Specifier}\label{override}
% 13 Jan: copyedits in and proofed
%% 27 Jan, code checked and replaced as needed by Josh

The \lstinline!override! keyword ensures that a member function overrides a corresponding \lstinline!virtual! member function in a base class.

\subsection[Description]{Description}\label{description}

The \romeogloss{contextual keyword} \lstinline!override! can be provided at the
end of a member-function declaration to ensure that the decorated
function is indeed \emph{overriding} a corresponding \lstinline!virtual!
member function in a base class, as opposed to \emph{hiding} it or
otherwise inadvertently introducing a distinct function declaration:

% original
%begin{emcppslisting}
%struct Base
%{
%    virtual void f(int);
%    void g(int);
%};
%
%struct Derived : Base
%{
%    void f();              // hides (ù{\codeincomments{Base::f(int)}}ù) (likely mistake)
%    void f() override;     // error: (ù{\codeincomments{Base::f()}}ù) not found
%
%    void f(int);           // implicitly overrides (ù{\codeincomments{Base::f(int)}}ù)
%    void f(int) override;  // explicitly overrides (ù{\codeincomments{Base::f(int)}}ù)
%
%    void g();              // hides (ù{\codeincomments{Base::g(int)}}ù) (likely mistake)
%    void g() override;     // error: (ù{\codeincomments{Base::g()}}ù) not found
%
%    void g(int);           // hides (ù{\codeincomments{Base::g(int)}}ù) (likely mistake)
%    void g(int) override;  // Error: (ù{\codeincomments{Base::g()}}ù) is not (ù{\codeincomments{virtual}}ù).
%};
%\end{emcppslisting}
% Slava's replacement
\begin{emcppslisting}
struct Base
{
    virtual void f(int);
            void g(int);
    virtual void h(int) const;
    virtual void i(int) = 0;
};

struct DerivedWithoutOverride : Base
{
    void f();             // hides (ù{\codeincomments{Base::f(int)}}ù) (likely mistake)
    void f(int);          // OK, implicitly overrides (ù{\codeincomments{Base::f(int) }}ù)

    void g();             // hides (ù{\codeincomments{Base::g(int)}}ù) (likely mistake)
    void g(int);          // hides (ù{\codeincomments{Base::g(int)}}ù) (likely mistake)

    void h(int);          // hides (ù{\codeincomments{Base::h(int)}}ù) const (likely mistake)
    void h(int) const;    // OK, implicitly overrides (ù{\codeincomments{Base::h(int)}}ù) const

    void i(int);          // OK, implicitly overrides (ù{\codeincomments{Base::i(int)}}ù)
};

struct DerivedWithOverride : Base
{
    void f()          override;    // Error, (ù{\codeincomments{Base::f()}}ù) not found
    void f(int)       override;    // OK, explicitly overrides (ù{\codeincomments{Base::f(int)}}ù)

    void g()          override;    // Error, (ù{\codeincomments{Base::g()}}ù) not found
    void g(int)       override;    // Error, (ù{\codeincomments{Base::g()}}ù) is not virtual.

    void h(int)       override;    // Error, (ù{\codeincomments{Base::h(int)}}ù) not found
    void h(int) const override;    // OK, explicitly overrides (ù{\codeincomments{Base::h(int)}}ù)

    void i(int)       override;    // OK, explicitly overrides (ù{\codeincomments{Base::i(int)}}ù)
};
\end{emcppslisting}

\noindent Using this feature expresses design intent so that (1) human readers
are aware of it and (2) compilers can validate it.

As noted, \lstinline!override! is a contextual keyword. C++11 introduces keywords that have special meaning only in certain contexts.   In this case, \lstinline!override! is a keyword in the context of a declaration, but not otherwise using it as the identifier for a variable name, for example, is perfectly fine:

%%%%%%%% EXCEPTION! override here should NOT be bold
\begin{emcppslisting}
int override = 1;  // OK
\end{emcppslisting}

\subsection[Use Cases]{Use Cases}\label{use-cases}

\subsubsection[Ensuring that a member function of a base class is being overridden]{Ensuring that a member function of a base class is being overridden}\label{ensuring-that-a-member-function-of-a-base-class-is-being-overridden}

Consider the following polymorphic hierarchy of error-category classes, as we might have defined them using C++03:

\begin{emcppshiddenlisting}[emcppsbatch=e1]
struct ErrorCode;
struct ErrorCondition;
\end{emcppshiddenlisting}
\begin{emcppslisting}[emcppsbatch=e1]
struct ErrorCategory
{
    virtual bool equivalent(const ErrorCode& code, int condition);
    virtual bool equivalent(int code, const ErrorCondition& condition);
};

struct AutomotiveErrorCategory : ErrorCategory
{
    virtual bool equivalent(const ErrorCode& code, int condition);
    virtual bool equivolent(int code, const ErrorCondition& condition);
};
\end{emcppslisting}

\noindent Notice that there is a defect in the last line of the example above:
\lstinline!equivalent! has been misspelled. Moreover, the compiler did not
catch that error. Clients calling \lstinline!equivalent! on
\lstinline!AutomotiveErrorCategory! will incorrectly invoke the base-class
function. If the function in the base class happens to be defined, the
code might compile and behave unexpectedly at run time. Now, suppose
that over time the interface is changed by marking the
equivalence-checking function \lstinline!const! to bring the interface
closer to that of \lstinline!std::error_category!:

\begin{emcppshiddenlisting}[emcppsbatch=e2]
struct ErrorCode;
struct ErrorCondition;
\end{emcppshiddenlisting}
\begin{emcppslisting}[emcppsbatch=e2]
struct ErrorCategory
{
    virtual bool equivalent(const ErrorCode& code, int condition) const;
    virtual bool equivalent(int code, const ErrorCondition& condition) const;
};
\end{emcppslisting}

Without applying the corresponding modification to all classes deriving
from\linebreak[4] \lstinline!ErrorCategory!, the semantics of the program change due to
the derived classes now hiding
the base class's
\lstinline!virtual! member function instead of overriding it. Both errors discussed above
would be detected automatically if the \lstinline!virtual!
functions in all derived classes were decorated with \lstinline!override!:

\begin{emcppslisting}[emcppsbatch=e2]
struct AutomotiveErrorCategory : ErrorCategory
{
    bool equivalent(const ErrorCode& code, int condition) override;
        // Error, failed when base class changed

    bool equivolent(int code, const ErrorCondition& code) override;
        // Error, failed when first written
};
\end{emcppslisting}

\noindent What's more, \lstinline!override! serves as a clear indication of the derived-class author's intent to customize the
behavior of \lstinline!ErrorCategory!. For any given member function, using \lstinline!override! necessarily renders any use of \lstinline!virtual! for
that function syntactically and semantically redundant. The only
cosmetic reason for retaining \lstinline!virtual! in the presence of
\lstinline!override! would be that \lstinline!virtual! appears to the left of
the function declaration, as it always has, instead of all the way to
the right, as \lstinline!override! does now.

\subsection[Potential Pitfalls]{Potential Pitfalls}\label{potential-pitfalls}

\subsubsection[Lack of consistency across a codebase]{Lack of consistency across a codebase}\label{lack-of-consistency-across-a-codebase}

Relying on \lstinline!override! as a means of ensuring that changes to
base-class interfaces are propagated across a codebase can prove
unreliable if this feature is used inconsistently --- i.e., not applied in every circumstance where its use would be appropriate. In
particular, altering the signature of a \lstinline!virtual! member function
in a base class and then compiling ``the world'' will always flag as an
error any nonmatching derived-class function where \lstinline!override!
was used but might fail even to warn where it is not.

\subsection[Annoyances]{Annoyances}\label{annoyances}

\hspace*{\fill}

\subsection[See Also]{See Also}\label{see-also}

\hspace*{\fill}

\subsection[Further Reading]{Further Reading}\label{further-reading}

\begin{itemize}
\item{\cite{meyers15b}}
\item{\cite{fluentcpp}}
\end{itemize}




\addtocontents{toc}{\protect\newpage}
\newpage
%\section[Raw String Literals]{Syntax for Unprocessed String Contents}\label{raw-string-literals}
% 12 Feb 2021, revisions in; ready for Josh's code fixes
% 16 Feb 2021 JMB; code compiles

\emcppsFeature{
    short={Raw String Literals},
    long={Syntax for Unprocessed String Contents},
}{raw-string-literals}
\setcounter{table}{0}
\setcounter{footnote}{0}
\setcounter{lstlisting}{0}
%\section[Raw String Literals]{Syntax for Unprocessed String Contents}

Raw string literals obviate the need to escape each contained special
character individually.

\subsection[Description]{Description}\label{description}

A \emph{raw} string literal is a new form of syntax for string literals
that allows developers to embed arbitrary character sequences in a
program's source code, without having to modify them by escaping
individual special characters. As an introductory example, suppose that
we want to write a small program that outputs the following text into the standard output stream:

% this is intentionally left plain - this is not code, this is what we are saying
% we  want to output
\begin{lstlisting}[style=plain]
printf("Hello, %s%c\n", "World", '!');
\end{lstlisting}
    
\noindent In C++03, capturing the line of C code above in a string literal would
require five escape (\lstinline!\!) characters distributed
throughout the string:

\begin{emcppslisting}
#include <iostream>  // (ù{\codeincomments{std::cout}}ù), (ù{\codeincomments{std::endl}}ù)

int main()
{
    std::cout << "printf(\"Hello, %s%c\\n\", \"World\", '!');" << std::endl;
    return 0;  //        ^            ^  ^   ^      ^
               //             escape characters
}
\end{emcppslisting}
    
\noindent If we use C++11's \emph{raw} string-literal syntax, no escaping is
required:

\begin{emcppslisting}
#include <iostream>  // (ù{\codeincomments{std::cout}}ù), (ù{\codeincomments{std::endl}}ù)

int main()
{
    std::cout << R"(printf("Hello, %s%c\n", "World", '!');)" << std::endl;
    return 0;  //^ ^                                      ^
               //   additional raw string-literal syntax (C++11)
}
\end{emcppslisting}
    
\noindent To represent the original character
data as a raw string literal, we typically need only to add a capital \lstinline!R! immediately
(adjacently) before the starting quote (\lstinline!"!) and nest the
character data within parentheses, \lstinline!(!~\lstinline!)! (with some exceptions; see \intrarefsimple{collisions}).  
%\textit{\titleref{collisions}} on page~\pageref{collisions}). 
Sequences of
characters that would be escaped in a regular string literal are instead
interpreted verbatim:

\begin{emcppslisting}
const char s0[] = R"({ "key": "value" })";
    // OK, equivalent to (ù{\codeincomments{"\{ $\backslash$"key$\backslash$": $\backslash$"value$\backslash$" \}"}}ù)
\end{emcppslisting}
    
\noindent In contrast to conventional string literals, \emph{raw} string literals
(1) treat unescaped embedded double quotes (\lstinline!"!) as literal data,
(2) do not interpret special-character escape sequences (e.g.,
\lstinline!\n!, \lstinline!\t!), and (3) interpret both vertical{\cprotect\footnote{To incorporate a newline character into a conventional string literal one must represent that newline using the escape sequence \lstinline!\n!.  Attempting to do so by actually entering a newline into the source (i.e., making the string literal span lines of source code) is an error.}} and horizontal 
whitespace characters present in
the source file as part of the string contents{\cprotect\footnote{In
this example, we assume that all trailing whitespace has been stripped
since even trailing whitespace in a raw literal would be captured.}}:


\begin{emcppslisting}
const char s1[] = R"(line one
line two
    line three)";
    // OK
\end{emcppslisting}
    
\noindent Note that any literal tab characters are treated the same as a
\lstinline!\t! and hence can be problematic, especially when
developers have inconsistent tab settings; see \intraref{potential-pitfalls-rawstringliteral}{unexpected-indentation}.  
%\textit{\titleref{potential-pitfalls-rawstringliteral}: \titleref{unexpected-indentation}} on page~\pageref{unexpected-indentation}. 
Finally, all string literals are concatenated with
adjacent ones in the same way the conventional ones are in C++03:

\begin{emcppslisting}
const char s2[] = R"(line one)"         "\n"
                    "line two"          "\n"
                  R"(    line three)";
    // OK, equivalent to (ù{\codeincomments{"line one$\backslash$nline two$\backslash$n~~~~line three"}}ù)    
\end{emcppslisting}

\noindent These same rules apply to both raw \emph{wide} string literals and raw \emph{Unicode} ones (see \featureref{\locationa}{unicode-string-and-character-literals}) that are introduced by placing their corresponding prefix before the \lstinline!R! character:

\begin{emcppslisting}
const wchar_t  ws [] =  LR"(Raw\tWide\tLiteral)";  
    // Represents "Raw\tWide\tLiteral", not "Raw    Wide    Literal".

const char     u8s[] = u8R"(\U0001F378)";  // Represents "\U0001F378", *not* "(ù{\martini}ù)"
const char16_t us [] =  uR"(\U0001F378)";  //       "         "          "     " 
const char32_t Us [] =  UR"(\U0001F378)";  //       "         "          "     " 
\end{emcppslisting}

\subsubsection[Collisions]{Collisions}\label{collisions}

Although unlikely, the data to be expressed within a
string literal might itself contain the character sequence \lstinline!)"!
embedded within it:

\begin{emcppslisting}
#include <cstdio>   // (ù{\codeincomments{printf}}ù)

void emitHelloWorld()
{
    printf("printf(\"Hello, World!\")");
  //                                ^^
  // The (ù{\codeincomments{)"}}ù) character sequence terminates a typical raw string literal.
}
\end{emcppslisting}
    
\noindent If we use the basic syntax for a \emph{raw} string literal we will get a
syntax error:


\begin{emcppslisting}[emcppsignore={syntax error}]
const char s3[] = R"(printf("printf(\"Hello, World!\")"))";  (ù{\codeincomments{// collision}}ù)
(ù{\codeincomments{//}}ù)                                                   (ù{\codeincomments{\textasciicircum\textasciicircum}}ù)
(ù{\codeincomments{//}}ù)                      (ù{\codeincomments{Syntax error after literal ends}}ù)
\end{emcppslisting}
    
\noindent To circumvent this problem, we could escape every special character in
the string separately, as in C++03, but the result is difficult to read
and error prone:

\begin{emcppslisting}
const char s4[] = "printf(\"printf(\\\"Hello, World!\\\")\")";  // error prone
\end{emcppslisting}
    
\noindent Instead, we can use the extended disambiguation syntax of \emph{raw}
string literals to resolve the issue:

%\begin{emcppslisting}
%const char s5[] = R"###(printf("printf(\"Hello, World!\")"))###";  (ù{\codeincomments{\color{skyblue}\ttfamily\itshape // cleaner}}ù)
%\end{emcppslisting}
\begin{emcppslisting}
const char s5[] = R"###(printf("printf(\"Hello, World!\")"))###";  (ù{\codeincomments{// cleaner}}ù)
\end{emcppslisting}

    
\noindent This disambiguation syntax allows us to insert an essentially
arbitrary{\cprotect\footnote{The delimiter of a raw string literal can
comprise any member of the \romeogloss{basic source character set} except
space, backslash, parentheses, and the control characters representing
  horizontal tab, vertical tab, form feed, and new line.}} sequence of
characters between the outermost quote/parenthesis pairs that avoids the
collision with the literal data when taken as a combined sequence (e.g.,
\lstinline!)###"!):

\begin{emcppslisting}
//                       delimiter and parenthesis
//                  v~~~                           ~~~v
const char s6[] = R"xyz(<-- Literal String Data -->)xyz";
//                ^     ^~~~~~~~~~~~~~~~~~~~~~~~~~^
//                |          string contents
//                |
//                | uppercase (ù{\codeincomments{R}}ù)
\end{emcppslisting}
    
\noindent The value of \lstinline!s6! above is equivalent to
\lstinline!"<--!~\lstinline!Literal!~\lstinline!String!~\lstinline!Data!~\lstinline!-->"!.
Every raw string literal comprises these syntactical elements, in order:
\begin{itemize}
\item{an uppercase \lstinline!R!}
\item{the opening double quotes, \lstinline!"!} 
\item{an optional arbitrary sequence of characters called the \emph{delimiter} (e.g., \lstinline!xyz!)}
\item{an opening parenthesis, \lstinline!(!}
\item{the contents of the string}
\item{a closing parenthesis, \lstinline!)!}
\item{the same delimiter (if any) specified previously (i.e., \lstinline!xyz!, not reversed)} 
\item{the closing double quotes, \lstinline!"!}
\end{itemize}

The delimiter can be --- and, in practice, very often is --- an empty character
sequence:

\begin{emcppslisting}
const char s7[] = R"("Hello, World!")";
    // OK, equivalent to (ù{\codeincomments{$\backslash$"Hello, World!$\backslash$"}}ù)
\end{emcppslisting}
    
\noindent A nonempty delimiter (e.g.,~\lstinline|!|) can be used to disambiguate any
appearance of the \lstinline!)"! character sequence within the literal data

\begin{emcppslisting}
const char s8[] = R"!("--- R"(Raw literals are not recursive!)" ---")!";
    // OK, equivalent to (ù{\codeincomments{$\backslash$"--- R$\backslash$"(Raw literals are not recursive!)$\backslash$" ---$\backslash$"}}ù)
\end{emcppslisting}
    
\noindent Had an empty delimiter been used to initialize \lstinline!s8! (above), the
compiler would have produced a (perhaps obscure) compile-time error:

%\begin{emcppslisting}
% // error: decrement of read-only location
%
% const char s8[] = R"("--- R"(Raw literals are not recursive!)" ---")";
% //                                                              ^~
%\end{emcppslisting}
\begin{emcppslisting}[emcppserrorlines={1}]
const char s8a[] = R"("---R"( Raw literals are not recursive!)" ---")";         
    //                                                            ^~            
    // Error, decrement of read-only location
\end{emcppslisting}
    
\noindent In fact, it could turn out that a program with an unexpectedly
terminated \emph{raw} string literal could still be valid and compile
quietly:

\begin{emcppshiddenlisting}[emcppsbatch=e1]
#include <cstdio>   // (ù{\codeincomments{printf}}ù)
\end{emcppshiddenlisting}
\begin{emcppslisting}[emcppsbatch=e1]
void emitPith()
{
    printf(R"("Live-Free, don't (ever)","Die!");
        (ù{\codeincomments{// Prints: "Live-Free, don't (ever}}ù)

    printf((R"("Live-Free, don't (ever)","Die!"));
        (ù{\codeincomments{// Prints: Die!}}ù)
}
\end{emcppslisting}
    
\noindent Fortunately, examples like the one above are invariably contrived, not
accidental.

\subsection[Use Cases]{Use Cases}\label{use-cases}

\subsubsection[Embedding code in a C++ program]{Embedding code in a C++ program}\label{embedding-code-in-a-c++-program}

When a source code snippet needs to be embedded as part of the source
code of a C++ program, use of a \emph{raw} string literal can
significantly reduce the syntactic noise that would otherwise be caused
by repeated escape sequences. As an example, consider a regular
expression for an online shopping product ID represented as a
conventional string literal:

\begin{emcppslisting}
const char* productIdRegex = "[0-9]{5}\\(\".*\"\\)";
    // This regular expression matches strings like (ù{\codeincomments{12345("Product")}}ù).
\end{emcppslisting}
    
\noindent Not only do the backslashes obscure the meaning to human readers, a
mechanical translation is often needed{\cprotect\footnote{Such as when
you want to copy the contents of the string literal into an online
  regular-expression validation tool.}} when transforming between source
and data, introducing significant opportunities for human error. Using a
raw string literal solves these problems:

\begin{emcppslisting}
const char* productIdRegex = R"([0-9]{5}\(".*"\))";
\end{emcppslisting}
    
\noindent Another format that benefits from raw string literals is JSON, due to
its frequent use of double quotes:

\begin{emcppslisting}
const char* testProductResponse = R"!(
{
    "productId": "58215(\"Camera\")",
    "availableUnits": 5,
    "relatedProducts": ["59214(\"CameraBag\")", "42931(\"SdStorageCard\")"]
})!";
\end{emcppslisting}
    
\noindent With a conventional string literal, the JSON string above would require
every occurrence of \lstinline!"! and \lstinline!\! to be escaped
and every new line to be represented as \lstinline!\n!, resulting
in visual noise, less interoperability with other tools accepting or
producing JSON, and heightened risk during manual maintenance.

Finally, raw string literals can also be helpful for
whitespace-sensitive languages, such as Python (but see \intraref{potential-pitfalls-rawstringliteral}{encoding-of-newlines-and-whitespace}): 
%\textit{\titleref{potential-pitfalls-rawstringliteral}: \titleref{encoding-of-newlines-and-whitespace}} on page~\pageref{encoding-of-newlines-and-whitespace}):

\begin{emcppslisting}
const char* testPythonInterpreterPrint = R"(def test():
    print("test printing from Python")
)";
\end{emcppslisting}
    

\subsection[Potential Pitfalls]{Potential Pitfalls}\label{potential-pitfalls-rawstringliteral}

\subsubsection[Unexpected indentation]{Unexpected indentation}\label{unexpected-indentation}

Consistent indentation and formatting of source code facilitates human
comprehension of program structure. Space and tabulation
(\lstinline!\t!) characters{\cprotect\footnote{Always
representing indentation as the precise number of spaces (instead of
tab characters) --- especially when committed to source-code control
  systems --- goes a long way to avoiding this issue.}} used for the
purpose of source code formatting are, however, always interpreted as
part of a raw string literal's contents:

\begin{emcppshiddenlisting}[emcppsbatch={e2,e3,e4}]
#include <iostream>  // (ù{\codeincomments{std::cout}}ù)
\end{emcppshiddenlisting}
\begin{emcppslisting}[emcppsbatch=e2]
void emitPythonEvaluator0(const char* expression)
{
    std::cout << R"(
        def evaluate():
            print("Evaluating...")
            return )" << expression << '\n';
}
\end{emcppslisting}
    
\noindent Despite the intention of the programmer to aid readability by indenting
the above raw string literal consistently with the rest of the code, the
streamed data will contain a large number of spaces (or tabulation
characters), resulting in an invalid Python program:

\begin{lstlisting}[language=python]
        def evaluate():
            print("Evaluating...")
            return (ù{\emph{someExpression}}ù)
# ^~~~~~
# Error: excessive indentation
\end{lstlisting}
    
\noindent Correct Python code would start unindented and then be indented the same number
of spaces (e.g.,~exactly four):

\begin{lstlisting}[language=python]
def evaluate():
    print("Evaluating...")
    return (ù{\emph{someExpression}}ù)
\end{lstlisting}
    
\noindent Correct --- albeit visually jarring --- Python code can be expressed with a
single \emph{raw} string literal, but visualizing the final output requires some effort:

\begin{emcppslisting}[emcppsbatch=e3]
void emitPythonEvaluator1(const char* expression)
{
    std::cout << R"(def evaluate():
    print("Evaluating...")
    return )" << expression << '\n';
}
\end{emcppslisting}
    
\noindent When more explicit control is desired, we can use a mixture of
\romeogloss{raw string literals} and explicit new lines represented as
\romeogloss{conventional string literals}:

\begin{emcppslisting}[emcppsbatch=e4]
void emitPythonEvaluator2(const char *expression)
{
    std::cout <<
        R"(def evaluate():)"                 "\n"
        R"(    print("Evaluating..."))"      "\n"
        R"(    return )" << expression << '\n';
}
\end{emcppslisting}
    

\subsubsection[Encoding of new lines and whitespace]{Encoding of new lines and whitespace}\label{encoding-of-newlines-and-whitespace}

The intent of the feature is that new lines should map to a single
\lstinline!\n! character regardless of how new lines are encoded
in the platform-specific encoding of the source file (e.g., \lstinline!\r\n!). The wording of the C++ Standard, however, is not
entirely clear.\footnote{\cite{miller13}} While all major compiler
implementations act in accordance with the original intent of the
feature, relying on a specific new line encoding may lead to nonportable
code until clarity is achieved.

In a similar fashion, the type of whitespace characters (e.g., tabs
versus spaces) used as part of a raw string literal can be significant.
As an example, consider a unit test verifying that a string representing
the status of the system is as expected:

\begin{emcppshiddenlisting}[emcppsbatch=e5]
#include <string>   // (ù{\codeincomments{std::string}}ù)
#include <cassert>  // standard C (ù{\codeincomments{assert}}ù) macro
struct System
{
    static std::string outputStatus();
};
\end{emcppshiddenlisting}
\begin{emcppslisting}[emcppsbatch=e5]
void verifyDefaultOutput()
{
    const std::string output = System::outputStatus();
    const std::string expected = R"(Current status:
    - No violations detected.)";

    assert(output == expected);
}
\end{emcppslisting}
    
\noindent The unit test might pass for years, until, for instance, the
company's indentation style changes from tabulation characters to
spaces, leading the \lstinline!expected! string to contain spaces instead of tabs, and thus test failures.{\cprotect\footnote{A
well-designed unit test will typically be imbued with \emph{expected
values}, rather than values that were produced by the previous run.
The latter is sometimes referred to as a \romeogloss{benchmark test}, and such tests are often implemented as \emph{diffs} against
a file containing output from a previous run. This file has
presumably been reviewed and is known (believed) to be correct and is sometimes called the \romeogloss{golden file}. Though ill advised, when trying to get a new version of the software to pass the benchmark test and when the precise format of the output of a system
changes subtly, the \romeogloss{golden file} may be summarily jettisoned --- and the new
output installed in its stead --- with little if any detailed review.
Hence, well-designed unit tests will often hard code exactly what is
to be expected (nothing more or less) directly in the
  \romeogloss{test-driver} source code.}}

\subsection[Annoyances]{Annoyances}\label{annoyances}

None so far

\subsection[See Also]{See Also}\label{see-also}

None so far

\subsection[Further Reading]{Further Reading}\label{further-reading}

None so far




\newpage
%\section[{\tt static\_assert}]{Compile-Time Assertions}\label{compile-time-assertions-(static_assert)}
%%%% 15 Jan: copyedits in and proofed
%% 27 Jan, code checked and replaced as needed by Josh
% 22 March 2021, in FPPs

\emcppsFeature{
    short={\lstinline!static_assert!},
    tocshort={\TOCCode static\_assert},
    long={Compile-Time Assertions},
    rhshort={\RHCode static\_assert},
}{compile-time-assertions-(static_assert)}
%\section[{\tt static\_assert}]{Compile-Time Assertions\sectionmark{\RHCode static\_assert}}\label{compile-time-assertions-(static_assert)}\sectionmark{\RHCode static\_assert}
\setcounter{table}{0}
\setcounter{footnote}{0}
\setcounter{lstlisting}{0}


The \lstinline!static_assert! keyword allows programmers to intentionally terminate compilation whenever a given compile-time predicate
 evaluates to \lstinline!false!.

\subsection[Description]{Description}\label{description}

Assumptions are
inherent in every program, whether we explicitly document them or not. A common way of validating certain
assumptions at run time is to use the classic \lstinline!assert! macro found
in \lstinline!<cassert>!. Such runtime assertions are not always ideal
because (1) the program must already be built and running for them to
even have a chance of being triggered and (2) executing a
\emcppsgloss{redundant check} at run time typically{\cprotect\footnote{It is
not unheard of for a program having runtime assertions to run faster with them
enabled than disabled. For example, asserting that a pointer is not
null enables the optimizer to elide all code branches that
  can be reached only if that pointer were null.}} results in a slower
program. Being able to validate an assertion at compile time avoids
several drawbacks:

\begin{enumerate}
\item{Validation occurs at compile time within a single translation unit and therefore doesn’t need to wait until a complete program is linked and executed.}
\item{Compile-time assertions can exist in many more places than runtime assertions and are unrelated to program control flow.}
\item{No runtime code will be generated due to a \lstinline!static_assert!, so program performance will not be impacted.}
\end{enumerate}

\subsubsection[Syntax and semantics]{Syntax and semantics}\label{syntax-and-semantics}

We can use \emcppsgloss{static assertion declarations} to conditionally trigger
controlled compilation failures depending on the truthfulness of a
\emcppsgloss{constant expression}. Such declarations are introduced by the
\lstinline!static_assert! keyword, followed by a parenthesized list
consisting of (1) a constant Boolean expression and (2) a mandatory (see
%{\it\titleref{static-annoyances}: \titleref{mandatory-string-literal}} on page~\pageref{mandatory-string-literal})
\intraref{static-annoyances}{mandatory-string-literal})
\emcppsgloss{string literal}, which will be
part of the compiler diagnostics if the compiler determines that the
assertion fails to hold:

\begin{emcppslisting}[emcppserrorlines=2]
static_assert(true, "Never fires.");
static_assert(false, "Always fires.");  
\end{emcppslisting}

\newpage%%%%
\noindent Static assertions can be placed anywhere in the scope of a namespace,
block, or class:

\begin{emcppshiddenlisting}[emcppsbatch=e1]
template <typename T> struct Predicate
{
   constexpr static int value = 1;
};
\end{emcppshiddenlisting}
\begin{emcppslisting}[emcppsbatch=e1,emcppserrorlines=8]
static_assert(1 + 1 == 2, "Never fires.");  // (global) namespace scope

template <typename T>
struct S
{
    void f0()
    {
        static_assert(1 + 1 == 3, "Always fires.");  // block scope
    }

    static_assert(!Predicate<T>::value, "Might fire.");  // class scope
};
\end{emcppslisting}

\noindent Providing a nonconstant expression to a \lstinline!static_assert! is
itself a compile-time error:

\begin{emcppslisting}
extern bool x;
static_assert(x, "Nice try.");  // Error, (ù{\codeincomments{x}}ù) is not a compile-time constant.
\end{emcppslisting}


\subsubsection[Evaluation of static assertions in templates]{Evaluation of static assertions in templates}\label{evaluation-of-static-assertions-in-templates}

The C++ Standard does not explicitly specify at precisely what point during
the compilation process the expressions tested by static assertion declarations are
evaluated. In particular, when used within the body of a template, the expression tested by a
\lstinline!static_assert! declaration might not be evaluated until
\emcppsgloss{template instantiation time}. In practice, however, a
\lstinline!static_assert! that does not depend on any template parameters
is essentially always{\cprotect\footnote{E.g.,
GCC~10.1, Clang~10.0, and MSVC~19.24}} evaluated immediately --- i.e., as
soon as it is parsed and irrespective of whether any subsequent template
instantiations occur:

\begin{emcppslisting}[emcppserrorlines={3,9}]
void f1()
{
    static_assert(false, "Impossible!");  // always evaluated immediately...
}                                         // even if (ù{\codeincomments{f1()}}ù) is never invoked

template <typename T>
void f2()
{
    static_assert(false, "Impossible!");  // always evaluated immediately...
}                                         // even if (ù{\codeincomments{f2()}}ù) is never instantiated
\end{emcppslisting}

\noindent The evaluation of a static assertion that is located within the
body of a class or function template and depends on at least one
template parameter is almost always bypassed during its initial parse since
the assertion predicate might evaluate to true or false depending on the template argument:

\begin{emcppshiddenlisting}[emcppsbatch=e2]
#include <complex>  // (ù{\codeincomments{std::complex}}ù)
\end{emcppshiddenlisting}
\begin{emcppslisting}[emcppsbatch=e2]
template <typename T>
void f3()
{
    static_assert(sizeof(T) >= 8, "Size < 8.");  // depends on (ù{\codeincomments{T}}ù)
}
\end{emcppslisting}

\noindent However, see
%{\it\titleref{static-potential-pitfalls}: \titleref{static-assertions-in-templates-can-trigger-unintended-compilation-failures}} on page~\pageref{static-potential-pitfalls}.
\intraref{static-potential-pitfalls}{static-assertions-in-templates-can-trigger-unintended-compilation-failures}.
In the example above, the compiler has no choice but to wait until each
time \lstinline!f3! is instantiated because the truth of the predicate will
vary depending on the type provided:

\begin{emcppslisting}[emcppsbatch=e2]
void g()
{
    f3<double>();               // OK
    f3<long double>();          // OK
    f3<std::complex<float>>();  // OK
    f3<char>();                 // Error, static assertion failed: Size < 8.
}
\end{emcppslisting}

\noindent The standard does, however, specify that a program containing any
template definition for which no valid specialization exists is
\emcppsgloss{ill formed, no diagnostic required (IFNDR)}, which was the case for
\lstinline!f2! but not \lstinline!f3!, above. Contrast each of the
\lstinline!h*!\lstinline[basicstyle=\ttfamily\itshape]!n!\lstinline!*! definitions below with its correspondingly numbered
\lstinline!f*!\lstinline[basicstyle=\ttfamily\itshape]!n!\lstinline!*!  definition above\footnote{The formula used --- \lstinline!int!~\lstinline!a[-1];! --- leads to $-1$, not $0$, to avoid a nonconforming extension to GCC that allows \lstinline!a[0]!.}:

\begin{emcppslisting}
void h1()
{
    int a[!sizeof(int) - 1];  // Error, same as (ù{\codeincomments{int a[-1];}}ù)
}

template <typename T>
void h2()
{
    int a[!sizeof(int) - 1];  // Error, always reported
}

template <typename T>
void h3()
{
    int a[!sizeof(T) - 1];    // typically reported only if instantiated
}
\end{emcppslisting}

\noindent Both \lstinline!f1! and \lstinline!h1! are ill-formed, nontemplate functions,
and both will always be reported at compile time, albeit typically with
decidedly different error messages as demonstrated by GCC 10.x's output:

\enlargethispage*{1ex}
\begin{lstlisting}[language=bash,style=plain]
f1: error: static assertion failed: Impossible!
h1: error: size (ù{\codeincomments{-1}}ù) of array (ù{\codeincomments{a}}ù) is negative
\end{lstlisting}

\noindent Both \lstinline!f2! and \lstinline!h2! are ill-formed template functions; the cause of their being ill-formed has nothing to do with the
template type and hence will always be reported as a compile-time error
in practice. Finally, \lstinline!f3! can be only contextually ill-formed,
whereas \lstinline!h3! is always necessarily ill-formed, and yet neither is
reported by typical compilers as such unless and until it has been
instantiated. Reliance on a compiler not to notice that a program is
ill-formed is dubious; see
%{\it\titleref{static-potential-pitfalls}: \titleref{static-assertions-in-templates-can-trigger-unintended-compilation-failures}} on page~\pageref{static-assertions-in-templates-can-trigger-unintended-compilation-failures}.
\intraref{static-potential-pitfalls}{static-assertions-in-templates-can-trigger-unintended-compilation-failures}.

\subsection[Use Cases]{Use Cases}\label{use-cases}

\subsubsection[Verifying assumptions about the target platform]{Verifying assumptions about the target platform}\label{verifying-assumptions-about-the-target-platform}

Some programs rely on specific properties of the native types provided
by their target platform. Static assertions can help ensure portability
and prevent such programs from being compiled into a malfunctioning
binary on an unsupported platform. As an example, consider a
program that relies on the size of an \lstinline!int! to be exactly 32
bits (e.g., due to the use of inline \lstinline!asm! blocks). Placing a
\lstinline!static_assert! in namespace scope in any of the program's
translation units will ensure that the assumption regarding the size
of \lstinline!int! is valid, and also serve as documentation for readers:

\begin{emcppslisting}
#include <climits>  // (ù{\codeincomments{CHAR\_BIT}}ù)

static_assert(sizeof(int) * CHAR_BIT == 32,
    "An (ù{\codeincomments{int}}ù) must have exactly 32 bits for this program to work correctly.");
\end{emcppslisting}

\noindent More typically, statically asserting the \emph{size} of an \lstinline!int!
avoids having to write code to handle an \lstinline!int! type's having
greater or fewer bytes when no such platforms are likely ever to
materialize:

\begin{emcppslisting}
static_assert(sizeof(int) == 4, "An (ù{\codeincomments{int}}ù) must have exactly 4 bytes.");
\end{emcppslisting}


\subsubsection[Preventing misuse of class and function templates]{Preventing misuse of class and function templates}\label{preventing-misuse-of-class-and-function-templates}

Static assertions are often used in practice to constrain class or function templates to prevent their being instantiated with unsupported types. If a type is not syntactically compatible with the template, static assertions provide clear customized error messages that replace compiler-issued diagnostics, which are often absurdly long and notoriously hard-to-read. More critically, static assertions actively avoid erroneous runtime behavior.
%%%  VR WANTS TO COME BACK TO THIS AND ADD A SEE ALSO.

As an example, consider the \lstinline!SmallObjectBuffer<N>! class
templates, which provide storage, aligned properly using \lstinline!alignas! (see
%``\titleref{alignas}" on page~\pageref{alignas}),
\featureref{\locationc}{alignas}),
for arbitrary objects whose size does
not exceed \lstinline!N!{\cprotect\footnote{A \lstinline!SmallObjectBuffer! is
  similar to C++17's \lstinline!std::any! (\cite{cpprefstdany}) in
  that it can store any object of any type. Instead of performing
  dynamic allocation to support arbitrarily sized objects, however,
  \lstinline!SmallObjectBuffer! uses an internal fixed-size buffer, which
  can lead to better performance and cache locality provided the
  maximum size of all of the types involved is known.}}:

%begin{emcppslisting}
%template <std::size_t N>
%class SmallObjectBuffer
%{
%private:
%    char d_buffer[N];
%
%public:
%    template <typename T>
%    void set(const T& object);
%
%    // ...
%};
%\end{emcppslisting}
\begin{emcppslisting}[emcppsbatch=e3]
#include <cstddef> // (ù{\codeincomments{std::size\_t}}ù), (ù{\codeincomments{std::max\_align\_t}}ù)
#include <new>     // placement (ù{\codeincomments{new}}ù)

template <std::size_t N>
class SmallObjectBuffer
{
private:
    alignas(std::max_align_t) char d_buffer[N];

public:
    template <typename T>
    void set(const T& object);

    // ...
};
\end{emcppslisting}


\noindent To prevent buffer overruns, it is important that \lstinline!set! accepts
only those objects that will fit in \lstinline!d_buffer!. The use of a
static assertion in the \lstinline!set! member function template catches
--- at compile time --- any such misuse:

%begin{emcppslisting}
%template <std::size_t N>
%template <typename T>
%void SmallObjectBuffer<N>::set(const T& object)
%{
%    static_assert(sizeof(T) <= N, "(ù{\codeincomments{object}}ù) does not fit in the small buffer.");
%    new (&d_buffer) T(object);
%}
%\end{emcppslisting}
\begin{emcppslisting}[emcppsbatch=e3]
template <std::size_t N>
template <typename T>
void SmallObjectBuffer<N>::set(const T& object)
{
    static_assert(sizeof(T) <= N, "(ù{\codeincomments{object}}ù) does not fit in the small buffer.");
    // Destroy existing object, if any; store how to destroy this new object of
    // type (ù{\codeincomments{T}}ù) later; then...
    new (&d_buffer) T(object);
}
\end{emcppslisting}



The principle of constraining inputs can be applied to most class and
function templates. \lstinline!static_assert! is particularly useful in
conjunction with standard \emcppsgloss[type trait]{type traits} provided in
\lstinline!<type_traits>!. In the \lstinline!rotateLeft! function template
(below), we have used two static assertions to ensure that only unsigned
integral types will be accepted:

\begin{emcppslisting}
#include <climits>      // (ù{\codeincomments{CHAR\_BIT}}ù)
#include <type_traits>  // (ù{\codeincomments{std::is\_integral}}ù), (ù{\codeincomments{std::is\_unsigned}}ù)

template <typename T>
T rotateLeft(T x)
{
    static_assert(std::is_integral<T>::value, "(ù{\codeincomments{T}}ù) must be an integral type.");
    static_assert(std::is_unsigned<T>::value, "(ù{\codeincomments{T}}ù) must be an unsigned type.");

    return (x << 1) | (x >> (sizeof(T) * CHAR_BIT - 1));
}
\end{emcppslisting}


\subsection[Potential Pitfalls]{Potential Pitfalls}\label{static-potential-pitfalls}

\subsubsection[Static assertions in templates can trigger unintended compilation failures]{Static assertions in templates can trigger unintended compilation failures}\label{static-assertions-in-templates-can-trigger-unintended-compilation-failures}

As mentioned in the description, any program containing a template for
which no valid specialization can be generated is
\emcppsgloss[ill formed, no diagnostic required (IFNDR)]{IFNDR}. Attempting to prevent the
use of, say, a particular function template overload by using a static
assertion that never holds produces such a program:

\begin{emcppslisting}[emcppserrorlines=10]
template <bool>
struct SerializableTag { };

template <typename T>
void serialize(char* buffer, const T& object, SerializableTag<true>);  // (1)

template <typename T>
void serialize(char* buffer, const T& object, SerializableTag<false>)  // (2a)
{
    static_assert(false, "T must be serializable.");  // independent of (ù{\codeincomments{T}}ù)
        // too obviously ill-formed: always a compile-time error
}
\end{emcppslisting}

\noindent In the example above, the second overload (2a) of \lstinline!serialize! is
provided with the intent of eliciting a meaningful compile-time error
message in the event that an attempt is made to serialize a
nonserializable type. The program, however, is technically
\emcppsgloss{ill formed} and, in this simple case, will likely result in a
compilation failure --- irrespective of whether either overload of
\lstinline!serialize! is ever instantiated.

A commonly attempted workaround
is to make the predicate of the assertion somehow dependent on a
template parameter, ostensibly forcing the compiler to withhold
evaluation of the \lstinline!static_assert! unless and until the template
is actually instantiated (a.k.a. \emcppsgloss{instantiation time}):

\begin{emcppshiddenlisting}[emcppsbatch=e4]
template <bool>
struct SerializableTag { };
\end{emcppshiddenlisting}
\begin{emcppslisting}[emcppsbatch=e4]
template <typename>  // N.B., we make no use of the (nameless) type parameter:
struct AlwaysFalse   // This class exists only to "outwit" the compiler.
{
    enum { value = false };
};

template <typename T>
void serialize(char* buffer, const T& object, SerializableTag<false>)  // (2b)
{
    static_assert(AlwaysFalse<T>::value, "T must be serializable.");  // OK
        // less obviously ill-formed: compile-time error when instantiated
}
\end{emcppslisting}

\noindent To implement this version of the second overload, we have provided an
intermediary class template \lstinline!AlwaysFalse! that, when instantiated
on any type, contains an enumerator named \lstinline!value!, whose value is
\lstinline!false!. Although this second implementation is more likely to
produce the desired result (i.e., a controlled compilation failure
only when \lstinline!serialize! is invoked with unsuitable arguments),
sufficiently ``smart'' compilers looking at just the current translation
unit would still be able to know that no valid instantiation of
\lstinline!serialize! exists and would therefore be well within their
rights to refuse to compile this still technically \emcppsgloss{ill formed}
program.

Equivalent workarounds achieving the same result without a
helper class are possible.

\begin{emcppshiddenlisting}[emcppsbatch=e5]
template <bool>
struct SerializableTag { };
\end{emcppshiddenlisting}
\begin{emcppslisting}[emcppsbatch=e5]
template <typename T>
void serialize(char* buffer, const T& object, SerializableTag<false>)  // (2c)
{
    static_assert(0 == sizeof(T), "T must be serializable.");  // OK
         // not too obviously ill-formed: compile-time error when instantiated
}
\end{emcppslisting}

\noindent Using this sort of obfuscation is not guaranteed to be either
portable or future-proof.

\subsubsection[Misuse of static assertions to restrict overload sets]{Misuse of static assertions to restrict overload sets}\label{misuse-of-static-assertions-to-restrict-overload-sets}

Even if we are careful to \emph{fool} the compiler into thinking that a
specialization is wrong \emph{only} if instantiated, we still cannot use
this approach to remove a candidate from an overload set because translation
will terminate if the static assertion is triggered. Consider this
flawed attempt at writing a \lstinline!process! function that will behave
differently depending on the size of the given argument:

\begin{emcppslisting}[emcppserrorlines={8,9,10,11,12,13}]
template <typename T>
void process(const T& x)  // (1) first definition of (ù{\codeincomments{process}}ù) function
{
    static_assert(sizeof(T) <= 32, "Overload for small types");  // BAD IDEA
    // ... (process small types)
}

template <typename T>
void process(const T& x)  // (2) compile-time error: redefinition of function
{
    static_assert(sizeof(T) > 32, "Overload for big types");     // BAD IDEA
    // ... (process big types)
}
\end{emcppslisting}

\noindent While the intention of the developer might have been to statically
dispatch to one of the two mutually exclusive overloads, the ill-fated
implementation above will not compile because the signatures of the two
overloads are identical, leading to a redefinition error. The semantics
of \lstinline!static_assert! are not suitable for the purposes of
\emcppsgloss[compile time dispatch]{compile-time dispatch}, and \emcppsgloss{SFINAE}-based approaches should be used instead.

\enlargethispage*{1ex}
To achieve the goal of removing up
front a specialization from consideration, we will need~to employ
\emcppsgloss{SFINAE}. To do that, we must instead find a way to get the
failing compile-time expression to be part of the function's
\emcppsgloss{declaration}:

\begin{emcppslisting}[emcppsbatch=e6]
template <bool> struct Check { };
    // helper class template having a (non-type) boolean template parameter
    // representing a compile-time predicate

template <> struct Check<true> { typedef int Ok; };
    // specialization of (ù{\codeincomments{Check}}ù) that makes the type (ù{\codeincomments{Ok}}ù) manifest *only* if
    // the supplied predicate (boolean template argument) evaluates to (ù{\codeincomments{true}}ù)

template <typename T,
          typename Check<(sizeof(T) <= 32)>::Ok = 0>  // SFINAE
void process(const T& x)  // (1)
{
    // ... (process small types)
}

template <typename T,
          typename Check<(sizeof(T) > 32)>::Ok = 0>  // SFINAE
void process(const T& x)  // (2)
{
    // ... (process big types)
}
\end{emcppslisting}

\noindent The empty \lstinline!Check! helper class template above in conjunction with
just one of its two possible specializations conditionally
exposes the \lstinline!Ok! type alias \emph{only} if the provided boolean
template parameter evaluates to \lstinline!true!. (Otherwise, by default, it
does not.) C++11 provides a library function, \lstinline!std::enable_if!, that more directly addresses this use case.\footnote{\emcppsgloss[concepts]{Concepts} --- a
language feature introduced in C++20 --- provides a far less baroque
alternative to \emcppsgloss{SFINAE} that allows for overload sets to be governed by
  the syntactic properties of their (compile-time) template arguments.}

During the substitution phase of template instantiation,
exactly one of the two overloads of the \lstinline!process! function will
attempt to access a nonexisting \lstinline!Ok! type alias via the
\lstinline!Check<false>! instantiation, which again, by default, is
nonexistent. Although such an error would typically result in a
compilation failure, in the context of template argument substitution it
will instead result in only the offending overload's being discarded,
giving other valid overloads a chance to be selected:

\begin{emcppshiddenlisting}[emcppsbatch=e6]
struct SmallType { char d_data[16]; };
struct BigType   { char d_data[32]; };
\end{emcppshiddenlisting}
\begin{emcppslisting}[emcppsbatch=e6]
void client()
{
    process(SmallType());  // discards (2), selects (1)
    process(BigType());    // discards (1), selects (2)
}
\end{emcppslisting}

\noindent This general technique of pairing template specializations is used widely
in modern C++ programming. For another, often more convenient way of
constraining overloads using \emcppsgloss{expression SFINAE}, see
%Section~\ref{trailing-function-return-types}, ``\titleref{trailing-function-return-types} on page~\pageref{trailing-function-return-types}."
\featureref{\locationa}{trailing-function-return-types}.

\subsection[Annoyances]{Annoyances}\label{static-annoyances}

\subsubsection[Mandatory string literal]{Mandatory string literal}\label{mandatory-string-literal}

Many compilation failures caused by static assertions are
self-explanatory since the offending line (which necessarily contains the
predicate code) is displayed as part of the compiler diagnostic. In
those situations, the message required{\cprotect\footnote{As of C++17,
  the message argument of a static assertion is optional.}} as part of
\lstinline!static_assert!'s grammar is redundant:

\begin{emcppshiddenlisting}[emcppsbatch=e7]
#include <type_traits>  // (ù{\codeincomments{std::is\_integral}}ù)
using T = int;
\end{emcppshiddenlisting}
\begin{emcppslisting}[emcppsbatch=e7]
static_assert(std::is_integral<T>::value, "(ù{\codeincomments{T}}ù) must be an integral type.");
\end{emcppslisting}

\noindent Developers commonly provide an empty string literal in these
cases:

\begin{emcppslisting}[emcppsbatch=e7]
static_assert(std::is_integral<T>::value, "");
\end{emcppslisting}

There is no universal consensus as to the ``parity" of the user-supplied error message.  Should it restate the asserted condition, or should it state what went amiss?
\begin{emcppshiddenlisting}[emcppsbatch=e8]
constexpr int x = 1;
\end{emcppshiddenlisting}
\begin{emcppslisting}[emcppsbatch=e8]
static_assert(0 < x, "x is negative");
   // misleading when (ù{\codeincomments{0 == x}}ù)
\end{emcppslisting}

\subsection[See Also]{See Also}\label{see-also}

\begin{itemize}
\item{%Section~\ref{trailing-function-return-types}, ``\titleref{trailing-function-return-types}" on page~\pageref{trailing-function-return-types} —
\seealsoref{trailing-function-return-types}{\locationa}Enabling expression \emcppsgloss{SFINAE} directly as part of a function's \emcppsgloss{declaration}  allows simple and fine-grained control over overload resolution.}
\end{itemize}

\subsection[Further Reading]{Further Reading}\label{further-reading}

\begin{itemize}
\item{\cite{klarer04}}
\end{itemize}




\newpage
%\section[Trailing Return]{Trailing Function Return Types}\label{trailing-function-return-types}
%%%% 15 Jan: copyedits in and proofed
%% 27 Jan, code checked and replaced as needed by Josh

\emcppsFeature{
    short={Trailing Return},
    long={Trailing Function Return Types},
}{trailing-function-return-types}
%\section[Trailing Return]{Trailing Function Return Types}\label{trailing-function-return-types}
\setcounter{table}{0}
\setcounter{footnote}{0}
\setcounter{lstlisting}{0}

Trailing return types provide a new alternate syntax in which the return type of a function is specified at
the end of a function declaration as opposed to at the beginning,
thereby allowing it to reference function parameters by name and to reference class
or namespace members without explicit qualification.

\subsection[Description]{Description}\label{description}

C++11 offers an alternative function-declaration syntax in which the
return type of a function is located to the right of its
\romeogloss{signature} (name, parameters, and qualifiers), offset by the
arrow token (\lstinline!->!); the function itself is introduced by the
keyword \lstinline!auto!, which acts as a type placeholder:
%\cprotect\footnote{Note that, when using the alternative trailing return syntax for a function (e.g., one returning a \texttt{double}), the override keyword would be inserted after call qualifiers and before the arrow:
%
%\begin{emcppslisting}[basicstyle=\ttfamily\footnotesize]
%virtual f(int value) const override -> double;
%\end{emcppslisting}\vspace*{-1ex}}:

\begin{emcppslisting}
auto f() -> void;  // equivalent to (ù{\codeincomments{void f();}}ù)
\end{emcppslisting}

When using the alternative, trailing-\lstinline!return!-type syntax, any \lstinline!const!,
\lstinline!volatile!, and reference qualifiers (see
%``\titleref{refqualifiers}" on page~\pageref{refqualifiers}
\featureref{\locatione}{refqualifiers}) are placed to the left of
the \lstinline!->!~\lstinline!*<return-type>*!, and any contextual keywords, such as
\lstinline!override! and \lstinline!final! (see
%``\titleref{override}" on page~\pageref{override}
\featureref{\locationa}{override}
and
%``\titleref{final}" on page~\pageref{final}
\featureref{\locatione}{final}), are placed to its right:

\begin{emcppslisting}
struct Base
{
    virtual int e() const;     // (ù{\codeincomments{const}}ù) qualifier
    virtual int f() volatile;  // (ù{\codeincomments{volitile}}ù) qualfier
    virtual int g() &;         // *lvalue*-reference qualifier
    virtual int h() &&;        // *rvalue*-reference qualifier
};

struct Derived : Base
{
    auto e() const    -> int override;  //  (ù{\codeincomments{override}}ù) contextual keyword
    auto f() volatile -> int final;     //  (ù{\codeincomments{final}}ù)         "        "
    auto g() &        -> int override;  //  (ù{\codeincomments{override}}ù)      "        "
    auto h() &&       -> int final;     //  (ù{\codeincomments{final}}ù)         "        "
};
\end{emcppslisting}

\noindent Using a trailing return type allows the parameters of a function to be
named as part of the specification of the return type, which can be
useful in conjunction with \lstinline!decltype!:

\begin{emcppslisting}
auto g(int x) -> decltype(x);  // equivalent to (ù{\codeincomments{int g(int x);}}ù)
\end{emcppslisting}

\noindent When using the trailing-return-type syntax in a member function
definition outside the class definition, names appearing in the return
type, unlike with the classic notation, will be looked up in class scope
by default:

\begin{emcppslisting}
struct S
{
    typedef int T;
    auto h1() -> T;  // trailing syntax for member function
    T h2();          // classical syntax for member function
};

auto S::h1() -> T { /*...*/ }  // equivalent to (ù{\codeincomments{S::T S::h1() \{ /\*...\*/ \}}}ù)
T    S::h2()      { /*...*/ }  // Error, (ù{\codeincomments{T}}ù) is unknown in this context.
\end{emcppslisting}

\noindent The same advantage would apply to a nonmember
function{\cprotect\footnote{A \lstinline!static! member function of a
\lstinline!struct! can be a viable alternative implementation to a free
function declared within a namespace; see \cite{lakos20}, section~1.4,
pp.~190--201, especially Figure~1-37c (p.~199), and section~2.4.9, pp.~312--321, especially Figure~2-23 (p.~316).}} defined outside of the namespace in
which it is declared:

\begin{emcppslisting}
namespace N
{
    typedef int T;
    auto h3() -> T;  // trailing syntax for free function
    T h4();          // classical syntax for free function
};

auto N::h3() -> T { /*...*/ }  // equivalent to (ù{\codeincomments{N::T N::h3() \{ /\*...\*/ \}}}ù)
T    N::h4()      { /*...*/ }  // Error, (ù{\codeincomments{T}}ù) is unknown in this context.
\end{emcppslisting}


Finally, since the syntactic element to be provided after the arrow
token is a separate type unto itself, return types involving pointers to
functions are somewhat simplified. Suppose, for example, we want to
describe a \romeogloss{higher-order function}, \lstinline!f!, that takes as its
argument a \lstinline!long!~\lstinline!long! and returns a pointer to a
function that takes an \lstinline!int! and returns a
\lstinline!double!{\cprotect\footnote{Co-author John Lakos first used the shown verbose declaration notation
while teaching Advanced Design and Programming using C++ at Columbia
  University (1991--1997).}}:

\begin{emcppslisting}
// [function(long long) returning]
//     [pointer to] [function(int x) returning] double   f;
//     [pointer to] [function(int x) returning] double   f(long long);
//                  [function(int x) returning] double  *f(long long);
//                                              double (*f(long long))(int x);
\end{emcppslisting}

\noindent Using the alternate trailing syntax, we can conveniently break the
declaration of \lstinline!f! into two parts: (1) the declaration of the
function's signature, \lstinline!auto!~\lstinline!f(long!~\lstinline!long)!, and (2) that of the return type, say, \lstinline!R! for now:

\begin{emcppslisting}
// [pointer to] [function (int) returning] double   R;
//              [function (int) returning] double  *R;
//                                         double (*R)(int);
\end{emcppslisting}

\noindent The two equivalent forms of the same declaration are shown below:

\begin{emcppslisting}
double (*f(long long))(int x);         // classic return-type syntax
auto f(long long) -> double (*)(int);  // trailing return-type syntax
\end{emcppslisting}

\noindent Note that both syntactic forms of the same declaration may appear
together within the same scope. Note also that not all functions that
can be represented in terms of the trailing syntax have a convenient
equivalent representation in the classic one:

\begin{emcppslisting}
#include <utility>  // (ù{\codeincomments{declval}}ù)

template <typename A, typename B>
auto foo(A a, B b) -> decltype(a.foo(b));
    // trailing return-type syntax

template <typename A, typename B>
decltype(std::declval<A&>().foo(std::declval<B&>())) foo(A a, B b);
    // classic return-type syntax (using C++11's (ù{\codeincomments{std::declval}}ù))
\end{emcppslisting}

\noindent In the example above, we were essentially forced to use the C++11
standard library template \lstinline!std::declval!\footnote{\cite{cpprefdeclval}} to express our intent with the classic
return-type syntax.

\subsection[Use Cases]{Use Cases}\label{use-cases}

\subsubsection[Function template whose return type depends on a parameter type]{Function template whose return type depends on a parameter type}\label{function-template-whose-return-type-depends-on-a-parameter-type}

Declaring a function template whose return type depends on the types of
one or more of its parameters is not uncommon in generic programming.
For example, consider a mathematical function that linearly interpolates
between two values of possibly different types:

\begin{emcppslisting}
template <typename A, typename B, typename F>
auto linearInterpolation(const A& a, const B& b, const F& factor)
    -> decltype(a + factor * (b - a))
{
    return a + factor * (b - a);
}
\end{emcppslisting}

\noindent The return type of \lstinline!linearInterpolation! is the type of
expression inside the \emph{\lstinline!decltype! specifier}, which is
identical to the expression returned in the body of the function. Hence,
this interface necessarily supports any set of input types for which
\lstinline!a!~\lstinline!+!~\lstinline!factor!~\lstinline!*!~\lstinline!(b!~\lstinline!-!~\lstinline!a)!
is valid, including types such as mathematical vectors, matrices, or
expression templates. As an added benefit, the presence of the
expression in the function's declaration enables \romeogloss{expression
SFINAE}, which is typically desirable for generic template functions
(see %Section~\ref{decltype}, ``\titleref{decltype}" on page~\pageref{decltype}
\featureref{\locationa}{decltype}).

\subsubsection[Avoiding having to qualify names redundantly in return types]{Avoiding having to qualify names redundantly in return types}\label{avoiding-having-to-qualify-names-redundantly-in-return-types}

When defining a function outside the \lstinline!class!, \lstinline!struct!, or
\lstinline!namespace! in which it is first declared, any unqualified names
present in the return type might be looked up differently depending on
the particular choice of function-declaration syntax used. When the
return type precedes the qualified name of the function definition as
is the case with classic syntax, all references to types declared in
the same scope where the function itself is declared must also be
qualified. By contrast, when the return type follows the
qualified name of the function,
the return type is looked up in the same scope in which the
function was first declared, just like its parameter types would. Avoiding redundant qualification of the return type  can be beneficial, especially when the qualifying name is long.

As an illustration, consider a class representing an abstract syntax
tree node that exposes a type alias:

\begin{emcppslisting}[emcppsbatch=e1]
struct NumericalASTNode
{
    using ElementType = double;
    auto getElement() -> ElementType;
};
\end{emcppslisting}

\noindent Defining the \lstinline!getElement! member function using traditional
function-declaration syntax would require repetition of the
\lstinline!NumericalASTNode! name:

\begin{emcppslisting}[emcppsbatch=e1]
NumericalASTNode::ElementType NumericalASTNode::getElement() { /*...*/ }
\end{emcppslisting}

\noindent Using the trailing-return-type syntax handily avoids the repetition:

\begin{emcppshiddenlisting}[emcppsbatch=e2]
// duplicated as hidden to show alternate implementation below
struct NumericalASTNode
{
    using ElementType = double;
    auto getElement() -> ElementType;
};
\end{emcppshiddenlisting}
\begin{emcppslisting}[emcppsbatch=e2]
auto NumericalASTNode::getElement() -> ElementType { /*...*/ }
\end{emcppslisting}

\noindent By ensuring that name lookup within the return type is the same as for
the parameter types, we avoid needlessly having to qualify names that
should be found correctly by default.

\subsubsection[Improving readability of declarations involving function pointers]{Improving readability of declarations involving function pointers}\label{improving-readability-of-declarations-involving-function-pointers}

Declarations of functions returning a pointer to either a function, a member function, or a data member are notoriously hard to parse --- even for seasoned programmers. As an example, consider a function called
\lstinline!getOperation! that takes a \lstinline!kind! of
enumerated \lstinline!Operation! as its argument and returns a pointer to a member
function of \lstinline!Calculator! that takes a \lstinline!double! and
returns a \lstinline!double!:

\begin{emcppshiddenlisting}[emcppsbatch=e3]
struct Calculator;
struct Operation;
\end{emcppshiddenlisting}
\begin{emcppslisting}[emcppsbatch=e3]
double (Calculator::*getOperation(Operation kind))(double);
\end{emcppslisting}

\noindent As we saw in the description, such declarations can be constructed
systematically but do not exactly roll off the fingers. On the other
hand, by partitioning the problem into (1) the declaration of the
function itself and (2) the type it returns, each individual problem
becomes far simpler than the original:

\begin{emcppslisting}[emcppsbatch=e3]
auto getOperation(Operation kind)  // (1) function taking a kind of (ù{\codeincomments{Operation}}ù)
    -> double (Calculator::*)(double);
        // (2) returning a pointer to a (ù{\codeincomments{Calculator}}ù) member function taking a
        //     (ù{\codeincomments{double}}ù) and returning a (ù{\codeincomments{double}}ù)
\end{emcppslisting}

\noindent Using this divide-and-conquer approach, writing such functions becomes fairly straightforward. Declaring a \romeogloss{higher-order function} that
takes a function pointer as an argument might be even easier to read
if a type alias is used via \lstinline!typedef! or, as of C++11,
  \lstinline!using!.

\subsection[Potential Pitfalls]{Potential Pitfalls}\label{potential-pitfalls}

\hspace{\fill}

\subsection[Annoyances]{Annoyances}\label{annoyances}

\hspace{\fill}

\subsection[See Also]{See Also}\label{see-also}

\begin{itemize}
\item{%Section~\ref{decltype}, ``\titleref{decltype}" on page~\pageref{decltype} —
\seealsoref{decltype}{\locationa}Function declarations may use \lstinline!decltype! either in conjunction with, or as an alternative to, trailing return types.}
\item{%Section~\ref{Function-Return-Type-Deduction}, ``\titleref{Function-Return-Type-Deduction}" on page~\pageref{Function-Return-Type-Deduction} —
\seealsoref{Function-Return-Type-Deduction}{\locationf}Leaving the return type to deduction shares syntactical similarities with trailing return types but brings with it significant pitfalls when migrating from C++11 to C++14.}
\end{itemize}

\subsection[Further Reading]{Further Reading}\label{further-reading}

\begin{itemize}
\item{\cite{mertz18}}
\end{itemize}




\newpage
%\section[Unicode Literals]{Unicode String Literals}\label{unicode-string-and-character-literals}
% 12 Feb 2021, revisions in; ready for Josh's code fixes
% 16 Feb 2021 JMB; code compiles
% 24 Feb 2021, copyedits in and proofed
% 22 March 2021, in FPPs

\emcppsFeature{
    short={Unicode Literals},
    long={Unicode String Literals},
}{unicode-string-and-character-literals}
\setcounter{table}{0}
\setcounter{footnote}{0}
\setcounter{lstlisting}{0}
%\section[Unicode Literals]{Unicode String Literals}\label{unicode-string-and-character-literals}

C++11 introduces a portable mechanism for ensuring that a literal is
encoded as UTF-8, UTF-16, or UTF-32.

\subsection[Description]{Description}\label{description-unicodestring}

According to the C++ Standard, the character encoding of string literals
is unspecified and can vary with the target platform or the configuration of
the compiler. In essence, the C++ Standard does not guarantee that the
string literal \lstinline!"Hello"! will be encoded as the
ASCII{\cprotect\footnote{In fact, C++ still fully supports platforms
  using EBCDIC, a rarely used alternative encoding to ASCII, as their primary text encoding.}} sequence
\lstinline!0x48,! \lstinline!0x65,! \lstinline!0x6C,! \lstinline!0x6C,! \lstinline!0x6F!
or that the character literal \lstinline!'X'! has the value
\lstinline!0x58!.

Table~\ref{unicodestring-table1} illustrates three new kinds of \emph{Unicode}-compliant
\emph{string literals}, each delineating the precise encoding of each
character.
 \begin{table}[h!]
\begin{center}
\begin{threeparttable}
\caption{Three new Unicode-compliant\\ literal strings}\label{unicodestring-table1}\vspace{1.5ex}
{\small \begin{tabular}{c|c|c}\thickhline
\rowcolor[gray]{.9}{\sffamily\bfseries Encoding} & {\sffamily\bfseries Syntax} & {\sffamily\bfseries Underlying Type}\\ \hline
UTF-8   &  \lstinline!u8"Hello"! & \lstinline!char!\tnote{a}\\ \hline
UTF-16 & \lstinline!u"Hello"!  & \lstinline!char16_t! \\ \hline
UTF-32  & \lstinline!U"Hello"! & \lstinline!char32_t!\\ \thickhline
\end{tabular}
} % end small
\begin{tablenotes}{\footnotesize
\item[a]{\lstinline!char8_t! in C++20}
}
\end{tablenotes}
\end{threeparttable}
\end{center} 
\end{table} 

\noindent A Unicode literal value is guaranteed to be encoded in UTF-8, UTF-16, or
UTF-32, for \lstinline!u8!, \lstinline!u!, and \lstinline!U! literals,
respectively:

\begin{emcppslisting}
char s0[] = "Hello";
    // unspecified encoding (albeit very likely ASCII)

char s1[] = u8"Hello";
    // guaranteed to be encoded as (ù{\codeincomments{\{0x48, 0x65, 0x6C, 0x6C, 0x6F, 0x0\}}}ù)
\end{emcppslisting}
    
\noindent C++11 also introduces \emph{universal character names} that provide a
reliably portable way of embedding Unicode code points in a C++ program.
They can be introduced by the \lstinline!\u! character
sequence followed by four hexadecimal digits or by the
\lstinline!\U! character sequence followed by eight
hexadecimal digits:

\newpage%%%%%%
\begin{emcppslisting}
#include <cstdio>  // (ù{\codeincomments{std::puts}ù)                                               
void f()                                                                        
{                                                                               
    std::puts(u8"\U0001F378"); // Unicode code point in a UTF8-encoded literal  
}
\end{emcppslisting}
    
\noindent This output statement is guaranteed to emit the cocktail emoji
(\martini) to \lstinline!stdout!, assuming that the receiving end is configured to
interpret output bytes as UTF-8.

\subsection[Use Cases]{Use Cases}\label{use-cases}

\subsubsection[Guaranteed-portable encodings of literals]{Guaranteed-portable encodings of literals}\label{guaranteed-portable-encodings-of-literals}

The encoding guarantees provided by the Unicode literals can be useful,
such as in communication with other programs or network/IPC protocols that
expect character strings having a particular encoding.

As an example, consider an instant-messaging program in which both the
client and the server expect messages to be encoded in UTF-8. As part of
broadcasting a message to all clients, the server code uses UTF-8
Unicode literals to guarantee that every client will receive a sequence
of bytes they are able to interpret and display as human-readable text:

\begin{emcppshiddenlisting}[emcppsbatch=e1]
#include <string>   // (ù{\codeincomments{std::string}}ù)
#include <ostream>  // (ù{\codeincomments{std::ostream}}ù)

class Packet                                                                    
{                                                                               
    // ...                                                                      
public:                                                                         
    // ...                                                                      
    std::ostream&  operator<<(const char *message);                             
};                                                                              
                                                                                
                                                                                
class Server                                                                    
{                                                                               
    void broadcastPacket(const Packet& data);                                   
    // ...                                                                      
public:                                                                         
    void broadcastServerMessage(const std::string& message);                    
    // ...                                                                      
                                                                                
};                                                                              
                                                                                
\end{emcppshiddenlisting}
\begin{emcppslisting}[emcppsbatch=e1]
void Server::broadcastServerMessage(const std::string& utf8Message)
{
    Packet data;
    data << u8"Message from the server: '" << utf8Message << u8"'\n";

    broadcastPacket(data);
}
\end{emcppslisting}
    
\noindent Not using \lstinline!u8! literals in the code snippet above could potentially result in
nonportable behavior and might require compiler-specific flags to
ensure that the source is UTF-8 encoded.

\subsection[Potential Pitfalls]{Potential Pitfalls}\label{potential-pitfalls}

\subsubsection[Embedding Unicode graphemes]{Embedding Unicode graphemes}\label{embedding-unicode-graphemes}

The addition of Unicode string literals to the language did not bring
along an extension of~the \emcppsgloss{basic source character set}: Even in C++11, the
default \emcppsgloss{basic source character set} is a subset of
ASCII.{\cprotect\footnote{Implementations are free to map characters
outside the basic source character set to sequences of its members,
resulting in the possibility of embedding other characters, such as emojis, in a C++ source file.}}

Developers might incorrectly assume that \lstinline!u8"!\martini\lstinline!"! is a
portable way of embedding a string literal representing the cocktail
emoji in a C++ program. The representation of the string literal, however, depends on what
encoding the compiler assumes for the source file, which can generally
be controlled through compiler flags. The only portable way of embedding
the cocktail emoji is to use its corresponding Unicode code point escape
sequence (\lstinline!u8"\U0001F378"!).

\subsubsection[Lack of library support for Unicode]{Lack of library support for Unicode}\label{lack-of-library-support-for-unicode}

Essential \emcppsgloss{vocabulary types}, such as \lstinline!std::string!, are
completely unaware of encoding. They treat any stored string as a
sequence of bytes. Even when correctly using Unicode string literals,
programmers unfamiliar with Unicode might be surprised by seemingly
innocent operations, such as asking for the size of a string
representing the cocktail emoji:

\begin{emcppslisting}
#include <cassert>  // standard C (ù{\codeincomments{assert}}ù) macro
#include <string>   // (ù{\codeincomments{std::string}}ù)
                                                                                
void f()                                                                        
{                                                                               
    std::string cocktail(u8"\U0001F378"); // big character (!)                  
    assert(cocktail.size() == 1);         // assertion failure (!)              
}
\end{emcppslisting}
    
\noindent Even though the cocktail emoji is a \emph{single} code point,
\lstinline!std::string::size! returns the number of code units (bytes) required to
encode it. The lack of Unicode-aware vocabulary types and utilities in
the Standard Library can be a source of defects and misunderstandings,
especially in the context of international program localization.

\subsubsection[Problematic treatment of UTF-8 in the type system]{Problematic treatment of UTF-8 in the type system}\label{utf-8-quirks}

UTF-8 string literals use \lstinline!char! as their \emcppsgloss{underlying type}. Such a
choice is inconsistent with~UTF-16 and UTF-32 literals, which provide
their own distinct character types, \lstinline!char16_t! and
\lstinline!char32_t!, respectively. This precludes providing distinct behavior for UTF-8 encoded strings using function overloading or template specialization because they are indistinguishable from strings having the encoding of the execution character set. Furthermore, whether the underlying
type of \lstinline!char! is a \lstinline!signed! or \lstinline!unsigned! type is
itself implementation defined. Note that
\lstinline!char! is distinct from both \lstinline!signed!~\lstinline!char! and
\lstinline!unsigned!~\lstinline!char!, but its behavior is guaranteed to be
  the same as one of those.
  
C++20
fundamentally changes how UTF-8 string literals work, by introducing a
new non\-ali\-as\-ing \lstinline!char8_t! character type whose representation
is guaranteed to match \lstinline!unsigned!~\lstinline!char!. The new
character type provides several benefits:
\begin{itemize}
\item{Ensures an \lstinline!unsigned! and distinct type for UTF-8 character data}
\item{Enables overloading for regular string literals versus UTF-8 string literals}
\item{Potentially achieves better performance due to the lack of special aliasing rules}
\end{itemize}
Unfortunately, the changes brought by C++20 are not
backward-compatible and might cause code targeting previous versions
of the language using \lstinline!u8! literals either to fail to compile
or to silently change its behavior when targeting C++20:

\begin{emcppslisting}
template <typename T> void print(const T*);  // (0)
void print(const char*);                     // (1)

void f()
{
    print(u8"text");  // invokes (1) prior to C++20, (0) afterwards
}
\end{emcppslisting}
        

\subsection[Annoyances]{Annoyances}\label{annoyances}

\hspace{\fill}

\subsection[See Also]{See Also}\label{see-also}

\hspace{\fill}

\subsection[Further Reading]{Further Reading}\label{further-reading}

\hspace{\fill}



\newpage
%\section[{\tt using} Aliases]{Type/Template Aliases (Extended {\SecCode typedef})}\label{alias-declarations-and-alias-templates}
% 2 Feb 2021, copyedits in and proofed
% 8 Feb 2021 JMB, compilation fixes
% 22 March 2021, in FPPs

\emcppsFeature{
    short={\lstinline!using! Aliases},
    tocshort={{\TOCCode using} Aliases},
    long={Type/Template Aliases (Extended {\SecCode typedef})},
    toclong={Type/Template Aliases (Extended \lstinline!typedef!)},
    rhshort={{\RHCode using} Aliases},
}{alias-declarations-and-alias-templates}
\setcounter{table}{0}
\setcounter{footnote}{0}
\setcounter{lstlisting}{0}
%\section[{\tt using} Aliases]{Type/Template Aliases (Extended {\SecCode typedef})}\label{alias-declarations-and-alias-templates}



%Alias declarations and alias templates provide an expanded use of the
%\lstinline!using! keyword, offering an alternative syntax (to
%\lstinline!typedef!) for creating a \romeogloss{type alias} that can itself be
%a template.

Alias declarations and alias templates provide an expanded use of the \lstinline!using! keyword~to introduce aliases for types and templates, thus providing a more general alternative to~\lstinline!typedef!.

\subsection[Description]{Description}\label{description}

The keyword \lstinline!using! has historically supported the introduction
of an alias for a named entity (e.g., type, function, or data) from some
named scope into the current one.
%see {\it\titleref{appendix:-brief-review-of-(c++03)-using-declarations}} on page~\pageref{appendix:-brief-review-of-(c++03)-using-declarations}.
%%%%% This appendix is removed and saved for a future edition. See file cpp03appendix.tex.
As of C++11, we can
employ the \lstinline!using! keyword to achieve everything that could
previously be accomplished with a \lstinline!typedef! declaration but in a
syntactic form that many people find more natural and intuitive (but
that offers nothing profoundly new):

\begin{emcppslisting}
using Type1 = int;     // equivalent to (ù{\codeincomments{typedef int Type1;}}ù)
using Type2 = double;  // equivalent to (ù{\codeincomments{typedef double Type2;}}ù)
\end{emcppslisting}

\noindent In contrast to \lstinline!typedef!, the name of the synonym created via the
\lstinline!using! syntax always appears on the left side of the \lstinline!=!
token and separate from the type declaration itself --- the advantage of
which becomes apparent with more involved types, such as
\emph{pointer-to-function}, \emph{pointer-to-member-function}, or
\emph{pointer-to-data-member}:

\begin{emcppslisting}
struct S { int i; void f(); };  // user-defined type (ù{\codeincomments{S}}ù) defined at file scope

using Type3 = void(*)();        // equivalent to (ù{\codeincomments{typedef void(*Type3)();}}ù)
using Type4 = void(S::*)();     // equivalent to (ù{\codeincomments{typedef void(S::*Type4)();}}ù)
using Type5 = int S::*;         // equivalent to (ù{\codeincomments{typedef int S::*Type5;}}ù)
\end{emcppslisting}

\noindent Just as with a \lstinline!typedef!, the name representing the type can be
qualified, but the symbol representing the synonym cannot:

\begin{emcppslisting}
namespace N { struct S { }; }  // original type (ù{\codeincomments{S}}ù) defined with namespace (ù{\codeincomments{N}}ù)

using Type6 = N::S;            // equivalent to (ù{\codeincomments{typedef N::S Type6;}}ù)
using ::Type7 = int;           // Error, the alias's name must be unqualified.
\end{emcppslisting}

\noindent Unlike a \lstinline!typedef!, however, a type alias introduced via
\lstinline!using! can itself be a template, known as an \emph{alias
template}:

\begin{emcppslisting}
template <typename T>
using Type8 = T;  // "identity" alias template

Type8<int>    i;  // equivalent to (ù{\codeincomments{int i;}}ù)
Type8<double> d;  // equivalent to (ù{\codeincomments{double d;}}ù)
\end{emcppslisting}

\newpage%%%%%%
\noindent Note, however, that neither partial nor full specialization of alias
templates is supported:

\begin{emcppslisting}[emcppserrorlines={4,7}]
template <typename, typename>   // general alias template
using Type9 = char;             // OK

template <typename T>           // attempted partial specialization of above
using Type9<T, int> = char;     // Error, expected (ù{\codeincomments{=}}ù) before (ù{\codeincomments{<}}ù) token

template <>                     // attempted full specialization of above
using Type10<int, int> = char;  // Error, expected unqualified-id before (ù{\codeincomments{using}}ù)
\end{emcppslisting}

\noindent Used in conjunction with existing class templates, alias templates allow
programmers to \emph{bind} one or more template parameters to a fixed
type, while leaving others open:

\begin{emcppslisting}
#include <utility>  // (ù{\codeincomments{std::pair}}ù)

template <typename T>
using PairOfCharAnd = std::pair<char, T>;
    // alias template that binds (ù{\codeincomments{char}}ù) to the first type parameter of (ù{\codeincomments{std::pair}}ù)

PairOfCharAnd<int>    pci;  // equivalent to (ù{\codeincomments{std::pair<char, int> pci;}}ù)
PairOfCharAnd<double> pcd;  // equivalent to (ù{\codeincomments{std::pair<char, double> pcd;}}ù)
\end{emcppslisting}

\noindent Finally, note that similar functionality can be
achieved in C++03, it suppresses type deduction and requires additional boilerplate code at both the
point of definition and the call~site:

\begin{emcppshiddenlisting}[emcppsbatch=e1]
#include <utility>  // (ù{\codeincomments{std::pair}}ù)
\end{emcppshiddenlisting}
\begin{emcppslisting}[emcppsbatch=e1]
// C++03 (obsolete)
template <typename T>
struct PairOfCharAnd
    // template class holding an alias, (ù{\codeincomments{Type}}ù), to (ù{\codeincomments{std::pair<char, T>}}ù)
{
    typedef std::pair<char, T> Type;
        // type alias binding (ù{\codeincomments{char}}ù) to the first type parameter of (ù{\codeincomments{std::pair}}ù)
};

PairOfCharAnd<int>::Type    pci;  // equivalent to (ù{\codeincomments{std::pair<char, int> pci;}}ù)
PairOfCharAnd<double>::Type pcd;  // equivalent to (ù{\codeincomments{std::pair<char, double> pcd;}}ù)
\end{emcppslisting}


\subsection[Use Cases]{Use Cases}\label{use-cases}

\subsubsection[Simplifying convoluted \lstinline!typedef! declarations]{Simplifying convoluted {\SubsubsecCode typedef} declarations}\label{simplifying-convoluted-typedef-declarations}

Complex \lstinline!typedef! declarations involving pointers to functions,
member functions, or data members require looking in the middle of the
declaration to find the alias name. As an example, consider a
\emph{callback} type alias intended to be used with asynchronous functions:

\newpage%%%%
\begin{emcppslisting}
typedef void(*CompletionCallback)(void* userData);
\end{emcppslisting}

\noindent Developers coming from a background other than C or C++03 might find the
above declaration hard to parse since the name of the alias
(\lstinline!CompletionCallback!) is embedded in the function pointer type.
Replacing \lstinline!typedef! with \lstinline!using! results in a simpler,
more consistent formulation of the same alias:

\begin{emcppslisting}
using CompletionCallback = void(*)(void* userData);
\end{emcppslisting}

\noindent The \lstinline!CompletionCallback! alias declaration (above) reads almost
completely left-to-right, and the name of the alias is clearly
specified after the \lstinline!using! keyword. To make the\linebreak[4]
\lstinline!CompletionCallback! alias read left-to-right, a
  trailing return (see \featureref{\locationa}{trailing-function-return-types})
%  Section~\ref{trailing-function-return-types}, ``\titleref{trailing-function-return-types}")
can be used:

  \begin{emcppslisting}
  using CompletionCallback = auto(*)(void* userData) -> void;
  \end{emcppslisting}

\noindent The alias declaration above can be read as,
  ``\lstinline!CompletionCallback! is an alias for a pointer to a
  function taking a \lstinline!void*! parameter named \lstinline!userData! and
  returning \lstinline!void!.''

\subsubsection[Binding arguments to template parameters]{Binding arguments to template parameters}\label{binding-template-arguments}

An alias template can be used to \emph{bind} one or more template
parameters of, say, a commonly used class template, while leaving the
other parameters open to variation. Suppose, for example, we have a class, \lstinline!UserData!, that contains
several distinct instances of \lstinline!std::map! --- each
having the same key type, \lstinline!UserId!, but with different payloads:

\begin{emcppshiddenlisting}[emcppsbatch={e2,e3}]
#include <map>  // (ù{\codeincomments{std::map}}ù)
#include <set>  // (ù{\codeincomments{std::set}}ù)
struct UserId {};
struct Message {};
struct Photos {};
struct Article {};
\end{emcppshiddenlisting}
\begin{emcppslisting}[emcppsbatch=e2]
class UserData  // class having excessive code repetition (BAD IDEA)
{
private:
    std::map<UserId, Message>          d_messages;
    std::map<UserId, Photos>           d_photos;
    std::map<UserId, Article>          d_articles;
    std::map<UserId, std::set<UserId>> d_friends;
};
\end{emcppslisting}

\noindent The example above, though clear and regular, involves significant
repetition, making it more difficult to maintain should we later opt to
change data structures. If we were to instead use an \emcppsgloss{alias
template} to bind the \lstinline!UserId! type to the first type
parameter of \lstinline!std::map!, we could both reduce code repetition
and enable the programmer to consistently replace \lstinline!std::map!
to another container
(e.g., \lstinline!std::unordered_map!{\cprotect\footnote{An
\lstinline!std::unordered_map! is an STL container type that became
available on all conforming platforms along with C++11. The
functionality is similar except that since it is not required to
support ordered traversal or (worst case) \lstinline!O[log(n)]! lookups
and \lstinline!O[n*log(n)]! insertions, \lstinline!std::unordered_map! can
be implemented as a hash table instead of a balanced tree, yielding
significantly faster average access times. See
  \cite{cpprefb}.}}) by performing the change in only
one place:

\newpage%%%%%
\begin{emcppslisting}[emcppsbatch=e3]
class UserData  // class with well-factored implementation (GOOD IDEA)
{
private:
    template <typename V>                  // using a template alias to bind
    using Mapping = std::map<UserId, V>;   // (ù{\codeincomments{UserId}}ù) as the key type

    Mapping<Message>          d_messages;
    Mapping<Photos>           d_photos;
    Mapping<Article>          d_articles;
    Mapping<std::set<UserId>> d_friends;
};
\end{emcppslisting}


\subsubsection[Providing a shorthand notation for type traits]{Providing a shorthand notation for type traits}\label{providing-a-shorthand-notation-for-type-traits}

Alias templates can provide a shorthand notation for \emcppsgloss[type trait]{type
traits}, avoiding \emcppsgloss{boilerplate code} in the usage site. As an
example, consider a simple type trait that adds a pointer to a given
type (akin to \lstinline!std::add_pointer!):

\begin{emcppslisting}[emcppsbatch=e4]
template <typename T>
struct AddPointer
{
    typedef T* Type;
};
\end{emcppslisting}

\noindent To use the trait above, the \lstinline!AddPointer! class template must be instantiated, and its nested
\lstinline!Type! alias must be accessed. Furthermore, in the generic context, it has to be prepended with the \lstinline!typename! keyword::

\begin{emcppslisting}[emcppsbatch=e4]
template <typename T>void f()
{
    T t;
    typename AddPointer<T>::Type p = t;
}
\end{emcppslisting}

\noindent The syntactical overhead of \lstinline!AddPointer! can be removed by
creating an alias template for its nested type alias, such as
\lstinline!AddPointer_t!:

\begin{emcppslisting}[emcppsbatch=e4]
template <typename T>
using AddPointer_t = typename AddPointer<T>::Type;
\end{emcppslisting}

\noindent Using \lstinline!AddPointer_t! instead of \lstinline!AddPointer! results in
shorter code devoid of boilerplate:

\begin{emcppslisting}[emcppsbatch=e4]
void g()
{
    int i;
    AddPointer_t<int> p = &i;
}
\end{emcppslisting}

\noindent Note that, since C++14, all
the standard type traits defined in the \lstinline!<type_traits>! header
provide a corresponding alias template with the goal of reducing
boilerplate code. For instance, C++14 introduces the
\lstinline!std::remove_reference_t! alias template for the C++11
\lstinline!std::remove_reference! type trait:

\begin{emcppshiddenlisting}[emcppsbatch=e5]
#include <type_traits>  // (ù{\codeincomments{std::remove\_reference}}ù)
\end{emcppshiddenlisting}
\begin{emcppslisting}[emcppsbatch=e5,emcppsstandards={c++14}]
typename std::remove_reference<int&>::type i0 = 5; // OK in both C++11 and C++14
std::remove_reference_t<int&> i1 = 5;              // OK in C++14
\end{emcppslisting}


\subsection[Potential Pitfalls]{Potential Pitfalls}\label{potential-pitfalls}

\hspace*{\fill}

\subsection[Annoyances]{Annoyances}\label{annoyances}

\hspace*{\fill}

\subsection[See Also]{See Also}\label{see-also}

\begin{itemize}
\item{\seealsoref{ctorinheriting}{\seealsolocationc}provides another meaning for the \lstinline!using! keyword to allow base-class constructors to be invoked as part of the derived class.}
\item{\seealsoref{trailing-function-return-types}{\seealsolocationa}provides an alternative syntax for function declaration, which can help improve readability in type aliases and alias templates involving function types.}
\end{itemize}

\subsection[Further Reading]{Further Reading}\label{further-reading}








%%%%%%%%%%%%%%% C++14
\stepcounter{cppxx}
\addtocontents{toc}{\protect\renewcommand*\protect\cftsecindent{1.25em}}
\cftaddnumtitleline{toc}{section}{\thechapter.\thecppxx}{C++14}{\thepage}
% to add to the TOC at the section level
\renewcommand{\cppxx}{C++14}

\newpage
\addtocontents{toc}{\protect\renewcommand*\protect\cftsecindent{4.5em}}
%\section[Aggregate Init '14]{Aggregates Having Default Member Initializers}\label{aggregate-member-initialization-relaxation}
\renewcommand{\cppxx}{C++14}%%%%%%
% 2 Feb 2021, copyedits entered and proofed
% 8 Feb 2021 JMB, compilation fixes
% 22 March 2021, in FPPs

\emcppsFeature{
    short={Aggregate Init '14},
    long={Aggregates Having Default Member Initializers},
}{aggregate-member-initialization-relaxation}
\setcounter{table}{0}
\setcounter{footnote}{0}
\setcounter{lstlisting}{0}
%\section[Aggregate Init '14]{Aggregates Having Default Member Initializers}\label{aggregate-member-initialization-relaxation}


C++14 enables the use of \emcppsgloss{aggregate initialization} with classes
employing default member initializers (see \featureref{\locationa}{Default-Member-Initializers}).

\subsection[Description]{Description}\label{description}

Prior to C++14, classes that used default member initializers (see \featureref{\locationa}{Default-Member-Initializers}) --- i.e., initializers that appear directly within the
scope of the class --- were not considered \emcppsgloss{aggregate} types:

\begin{emcppslisting}[emcppsbatch={e1,e2,e3}]
struct S                // aggregate type in C++14 but not C++11
{
    int i;
    bool b = false;     // uses default member initializer
};

struct A                // aggregate type in C++11 and C++14
{
    int  i;
    bool b;             // does not use default member initializer
};
\end{emcppslisting}

\noindent Because \lstinline!A! (but not \lstinline!S!) is considered an \emcppsgloss{aggregate} in
C++11, instances of \lstinline!A! can be created via \emcppsgloss{aggregate
initialization} (whereas instances of \lstinline!S! cannot):

\begin{emcppslisting}[emcppsbatch=e1]
A a={100, true};  // OK, in both C++11 and C++14
S s={100, true};  // Error, in C++11; OK, in C++14
\end{emcppslisting}

\noindent Note that since C++11, direct-list-initialization may be used to perform aggregate initialization (see \featureref{\locationc}{bracedinit}):

\begin{emcppslisting}[emcppsbatch=e2]
A a{100, true};  // OK in both C++11 and C++14 but not in C++03
\end{emcppslisting}

\enlargethispage*{2ex}
\noindent As of C++14, the requirements for a type to be categorized as an
\emcppsgloss{aggregate} are relaxed, allowing classes employing default
member initializers to be considered as such; hence, both \lstinline!A! and
\lstinline!S! are considered \emcppsgloss[aggregate]{aggregates} in C++14 and eligible for
\emcppsgloss{aggregate initialization}:

\begin{emcppshiddenlisting}[emcppsbatch=e3]
#include <cassert>  // standard C (ù{\codeincomments{assert}}ù) macro
\end{emcppshiddenlisting}
\begin{emcppslisting}[emcppsbatch=e3,emcppsstandards={c++14}]
void f()
{
    S s0{100, true};        // OK in C++14 but not in C++11
    assert(s0.i == 100);    // set via explicit aggregate initialization (above)
    assert(s0.b == true);   // set via explicit aggregate initialization (above)

    S s1{456};              // OK in C++14 but not in C++11
    assert(s1.i == 456);    // set via explicit aggregate initialization (above)
    assert(s1.b == false);  // set via default member initializer
}
\end{emcppslisting}

\noindent In the code snippet above, the C++14 aggregate \lstinline!S! is initialized
in two ways: \lstinline!s0! is created using aggregate initialization for
both data members, and \lstinline!s1! is created using aggregate
initialization for only the first data member (and the second is set via its
default member \nobreak{initializer}).

\subsection[Use Cases]{Use Cases}\label{use-cases}

\subsubsection[Configuration \lstinline!struct!s]{Configuration {\SubsubsecCode struct}s}\label{configuration-structs}

\emcppsgloss[aggregate]{Aggregates} in conjunction with default member initializers (see \featureref{\locationa}{Default-Member-Initializers}) can be used to provide concise customizable
configuration \lstinline!struct!s, packaged with typical default values. As
an example, consider a configuration \lstinline!struct! for an HTTP request
handler:

\begin{emcppslisting}[emcppsbatch=e4,emcppsstandards={c++14}]
struct HTTPRequestHandlerConfig
{
    int maxQueuedRequests = 1024;
    int timeout           = 60;
    int minThreads        = 4;
    int maxThreads        = 8;
};
\end{emcppslisting}

\enlargethispage*{2ex}
\noindent \emcppsgloss[aggregate initialization]{Aggregate initialization} can be used when creating objects of
type\linebreak[4]
\lstinline!HTTPRequestHandlerConfig! (above) to override one or more
of the defaults in definition order{\cprotect\footnote{In C++20, the
designated initializers feature adds flexibility (e.g., for
configuration \lstinline!struct!s, such as
\lstinline!HTTPRequestHandlerConfig!) by enabling explicit specification
of the names of the data members:

\begin{emcppslisting}[emcppsbatch=e4,style=footcode]
HTTPRequestHandlerConfig lowTimeout{.timeout = 15};
    // (ù{\fncodeincomments{maxQueuedRequests}}ù), (ù{\fncodeincomments{minThreads}}ù), and (ù{\fncodeincomments{maxThreads}}ù) have their default value.

HTTPRequestHandlerConfig highPerformance{.timeout = 120, .maxThreads = 16};
    // (ù{\fncodeincomments{maxQueuedRequests}}ù) and (ù{\fncodeincomments{minThreads}}ù) have their default value.
\end{emcppslisting} 
\vspace*{-3ex} }}:

\begin{emcppslisting}[emcppsbatch=e4]
HTTPRequestHandlerConfig getRequestHandlerConfig(bool inLowMemoryEnvironment)
{
    if (inLowMemoryEnvironment)
    {
        return HTTPRequestHandlerConfig{128};
            // (ù{\codeincomments{timeout}}ù), (ù{\codeincomments{minThreads}}ù), and (ù{\codeincomments{maxThreads}}ù) have their default value.
    }
    else
    {
        return HTTPRequestHandlerConfig{2048, 120};
            // (ù{\codeincomments{minThreads}}ù), and (ù{\codeincomments{maxThreads}}ù) have their default value.
    }
}

// ...
\end{emcppslisting}


\subsection[Potential Pitfalls]{Potential Pitfalls}\label{potential-pitfalls}

%%%%%%%% New section added 19 March 2021
\subsubsection{Empty list initializing members does not use the default initializer}

When we add default member initializers to members of an aggregate,
those initializers are in effect \emph{only} if that member has no corresponding
initializer in the braced initializer list. We cannot explicitly \emph{request}
the default value by placing empty braces into the list because that will
value-initialize the member with an empty initializer, not
use its default member initializer.  Any case where we want to explicitly initialize a later member variables value during aggregate initialization means we must manually determine the proper default for all prior members:

\begin{emcppslisting}[emcppsstandards={c++14}]
struct A 
{
     int i{1};
     int j{2};
     int k{3};
};

A a1{};            // OK, result is i=1, j=2,         k=3
A a2{ 4 };         // OK, result is i=4, j=2,         k=3
A a3{ 4, {}, 8 };  // Bug, result is i=4, j=0 (not 2), k=3
\end{emcppslisting}
%%%%%%%%%%%%%%%%%%%%%%

\subsection[Annoyances]{Annoyances}\label{annoyances}

\subsubsection[Syntactical ambiguity in the presence of brace elision]{Syntactical ambiguity in the presence of brace elision}\label{syntactical-ambiguity-in-the-presence-of-brace-elision}

During the initialization of multilevel \emcppsgloss[aggregate]{aggregates}, braces
around the initialization of a nested aggregate can be omitted
(\emcppsgloss{brace elision}):

\begin{emcppslisting}
struct S
{
    int arr[3];
};

S s0{{0, 1, 2}};  // OK, nested (ù{\codeincomments{arr}}ù) initialized explicitly
S s1{0, 1, 2};    // OK, brace elision for nested (ù{\codeincomments{arr}}ù)
\end{emcppslisting}

\noindent The possibility of \emcppsgloss{brace elision} creates an interesting
syntactical ambiguity when used alongside \emcppsgloss[aggregate]{aggregates} with
default member initializers (see \featureref{\locationa}{Default-Member-Initializers}). Consider a
\lstinline!struct!~\lstinline!X! containing three data members, one of which
has a default value:

\newpage%%%%%
\begin{emcppslisting}[emcppsbatch=e5]
struct X
{
    int a;
    int b;
    int c = 0;
};
\end{emcppslisting}

\noindent Now, consider various ways in which an array of elements of type
\lstinline!X! can be initialized:

\begin{emcppslisting}[emcppsbatch=e5,emcppsstandards={c++14},emcppserrorlines={9}]
X xs0[] = {{0, 1}, {2, 3}, {4, 5}};
    // OK, clearly 3 elements having the respective values:
    // (ù{\codeincomments{\{0, 1, 0\}}}ù), (ù{\codeincomments{\{2, 3, 0\}}}ù), (ù{\codeincomments{\{4, 5, 0\}}}ù)

X xs1[] = {{0, 1, 2}, {3, 4, 5}};
    // OK, clearly 2 elements with values:
    // (ù{\codeincomments{\{0, 1, 2\}}}ù), (ù{\codeincomments{\{3, 4, 5\}}}ù)

X xs2[] = {0, 1, 2, 3, 4, 5};
    // ...?
\end{emcppslisting}

\noindent Upon seeing the definition of \lstinline!xs2!, a programmer not versed in
the details of the C++ Language Standard might be unsure as to whether the
initializer of \lstinline!xs2! is three elements (like \lstinline!xs0!) or two
elements (like \lstinline!xs1!). The Standard is, however, clear that the
compiler will interpret \lstinline!xs2! the same as \lstinline!xs1!, and,
thus, the default values of \lstinline!X::c! for the two array elements
will be replaced with \lstinline!2! and \lstinline!5!, respectively.

\subsection[See Also]{See Also}\label{see-also}

\begin{itemize}
\item{\seealsoref{Default-Member-Initializers}{\seealsolocationa}allows developers to provide a default initializer for a data member directly in the definition of a class.}
\item{\seealsoref{bracedinit}{\seealsolocationc}introduces a syntactically similar feature for initializing objects in a uniform manner.}
\end{itemize}

\subsection[Further Reading]{Further Reading}\label{further-reading}





\newpage
%\section[Binary Literals]{Binary Literals: The {\SecCode 0b} Prefix}\label{binary-literals}
%% copyedits in and proofed
%% 27 Jan, code checked and replaced as needed by Josh

\emcppsFeature{
    toclong={Binary Literals: The \lstinline!0b! Prefix},
    short={Binary Literals},
    long={Binary Literals: The {\SecCode 0b} Prefix},
}{binary-literals}
%\section[Binary Literals]{Binary Literals}\label{binary-literals}\sectionmark{Binary Literals}
%\section[Binary Literals]{Binary Literals: The {\SecCode 0b} Prefix}\label{binary-literals}
\setcounter{table}{0}
\setcounter{footnote}{0}
\setcounter{lstlisting}{0}


\textit{Binary literals} are \textbf{integer literals}\glossary{integer literals} representing their values in base 2.

\subsection[Description]{Description}\label{description-binary}

A binary literal is an integral value represented in code in a binary numeral system. A binary literal consists of a \lstinline!0b! or \lstinline!0B! prefix followed by a nonempty sequence of binary digits, namely, \lstinline!0! and \lstinline!1!{\cprotect\footnote{Prior to being introduced in C++14, GCC supported binary literals (with the same syntax as the standard feature) as a nonconforming extension since version~4.3.0, released in March 2008;
%4.3 (released between March 2008 and May 2010);
for more details, see https://gcc.gnu.org/gcc-4.3/.}}:

%A \emph{binary literal} (e.g., \texttt{0b1110}) --- much like a
%hexadecimal literal (e.g., \texttt{0xE}) or an octal literal (e.g.,
%\texttt{016}) --- is a kind of \emph{integer literal} (in this case,
%having the \emph{decimal} value \texttt{14}). A binary literal consists
%of a \texttt{0b} (or \texttt{0B}) prefix followed by a nonempty
%sequence of binary digits (\texttt{0} or \texttt{1})

\begin{emcppslisting}[emcppsbatch=e1]
int i = 0b11110000;  // equivalent to (ù{\codeincomments{240}}ù), (ù{\codeincomments{0360}}ù), or (ù{\codeincomments{0xF0}}ù)
int j = 0B11110000;  // same value as above
\end{emcppslisting}

\noindent The first digit after the \lstinline!0b! prefix is the most significant
one:

\begin{emcppslisting}[emcppsbatch=e1]
static_assert(0b0     ==  0, "");  // 0*2^0
static_assert(0b1     ==  1, "");  // 1*2^0
static_assert(0b10    ==  2, "");  // 1*2^1 + 0*2^0
static_assert(0b11    ==  3, "");  // 1*2^1 + 1*2^0
static_assert(0b100   ==  4, "");  // 1*2^2 + 0*2^1 + 0*2^0
static_assert(0b101   ==  5, "");  // 1*2^2 + 0*2^1 + 1*2^0
// ...
static_assert(0b11010 == 26, "");  // 1*2^4 + 1*2^3 + 0*2^2 + 1*2^1 + 0*2^0
\end{emcppslisting}


\noindent Leading zeros --- as with octal and hexadecimal (but not decimal)
literals --- are ignored but can be added for readability:

\begin{emcppslisting}
static_assert(0b00000000 ==   0, "");
static_assert(0b00000001 ==   1, "");
static_assert(0b00000010 ==   2, "");
static_assert(0b00000100 ==   4, "");
static_assert(0b00001000 ==   8, "");
static_assert(0b10000000 == 128, "");
\end{emcppslisting}


\noindent The type of a binary literal
%{\cprotect\footnote{Its \emph{value category} is \emph{prvalue} like every other integer literal.}}
is by
default an \lstinline!int! unless that value cannot fit in an
\lstinline!int!. In that case, its type is the first type in the sequence
\{\lstinline!unsigned!~\lstinline!int!, %\pagebreak%%%%%
 \lstinline!long!,
 \lstinline!unsigned!~\lstinline!long!, \lstinline!long!~\lstinline!long!,
\lstinline!unsigned!~\lstinline!long!~\lstinline!long!\} in which it will fit. This
same type list applies for both \lstinline!octal! and \lstinline!hex!
literals but not for decimal literals, which, if initially
\lstinline!signed!, skip over any \lstinline!unsigned! types, and vice versa;
see \intrarefsimple{description-binary}.
  If neither of those is applicable, the compiler may use implementation-defined extended integer types such as \lstinline!__int128! to represent the literal if it fits --- otherwise the program is ill-formed:

% If neither of those is applicable, then the program is \emph{ill-formed}

\begin{emcppslisting}[emcppsignore={Invalid Descriptive Literals}]
// example platform 1:
// ((ù{\codeincomments{sizeof(int)}}ù): 4; (ù{\codeincomments{sizeof(long)}}ù): 4; (ù{\codeincomments{sizeof(long long)}}ù): 8)
auto i32  = 0b0111...[ 24 1-bits]...1111;  // (ù{\codeincomments{i32}}ù) is (ù{\codeincomments{int}}ù).
auto u32  = 0b1000...[ 24 0-bits]...0000;  // (ù{\codeincomments{u32}}ù) is (ù{\codeincomments{unsigned int}}ù).
auto i64  = 0b0111...[ 56 1-bits]...1111;  // (ù{\codeincomments{i64}}ù) is (ù{\codeincomments{long long}}ù).
auto u64  = 0b1000...[ 56 0-bits]...0000;  // (ù{\codeincomments{u64}}ù) is (ù{\codeincomments{unsigned long long}}ù).
auto i128 = 0b0111...[120 1-bits]...1111;  // Error, integer literal too large
auto u128 = 0b1000...[120 0-bits]...0000;  // Error, integer literal too large

// example platform 2:
// ((ù{\codeincomments{sizeof(int)}}ù): 4; (ù{\codeincomments{sizeof(long)}}ù): 8; (ù{\codeincomments{sizeof(long long)}}ù): 16)
auto i32  = 0b0111...[ 24 1-bits]...1111;  // (ù{\codeincomments{i32}}ù)  is (ù{\codeincomments{int}}ù).
auto u32  = 0b1000...[ 24 0-bits]...0000;  // (ù{\codeincomments{u32}}ù)  is (ù{\codeincomments{unsigned int}}ù).
auto i64  = 0b0111...[ 56 1-bits]...1111;  // (ù{\codeincomments{i64}}ù)  is (ù{\codeincomments{long}}ù).
auto u64  = 0b1000...[ 56 0-bits]...0000;  // (ù{\codeincomments{u64}}ù)  is (ù{\codeincomments{unsigned long}}ù).
auto i128 = 0b0111...[120 1-bits]...1111;  // (ù{\codeincomments{i128}}ù) is (ù{\codeincomments{long long}}ù).
auto u128 = 0b1000...[120 0-bits]...0000;  // (ù{\codeincomments{u128}}ù) is (ù{\codeincomments{unsigned long long}}ù).
\end{emcppslisting}


\noindent (Purely for convenience of
exposition, we have employed the C++11 \lstinline!auto! feature to
conveniently capture the type implied by the literal itself; see \featureref{\locationc}{auto-feature}.) Separately, the precise initial type of a binary literal, like any other
literal, can be controlled explicitly using the common integer-literal
suffixes \{\lstinline!u!, \lstinline!l!, \lstinline!ul!, \lstinline!ll!,
\lstinline!ull!\} in either lower- or uppercase:

\begin{emcppslisting}
auto i   = 0b101;        // type: (ù{\codeincomments{int}}ù);                 value: 5
auto u   = 0b1010U;      // type: (ù{\codeincomments{unsigned int}}ù);        value: 10
auto l   = 0b1111L;      // type: (ù{\codeincomments{long}}ù);                value: 15
auto ul  = 0b10100UL;    // type: (ù{\codeincomments{unsigned long}}ù);       value: 20
auto ll  = 0b11000LL;    // type: (ù{\codeincomments{long long}}ù);           value: 24
auto ull = 0b110101ULL;  // type: (ù{\codeincomments{unsigned long long}}ù);  value: 53
\end{emcppslisting}


\noindent Finally, note that affixing a minus sign to a binary
literal (e.g., \lstinline!-b1010!) --- just like any other integer literal
(e.g., \lstinline!-10!, \lstinline!-012!, or \lstinline!-0xa!) --- is parsed as a
non-negative value first, after which a unary minus is applied:

%\begin{emcppslisting}
%static_assert(sizeof(int) == 4, "");  // true on virtually all machines today
%static_assert(-0b1010 == -10, "");    // as if: (ù{\codeincomments{0 - 0b1010 == 0 - 10}}ù)
%static_assert(0x7fffffff != -0x7fffffff, "");  // Both values are (ù{\codeincomments{signed int}}ù).
%static_assert(0x80000000 == -0x80000000, "");  // Both values are (ù{\codeincomments{unsigned int}}ù).
%\end{emcppslisting}
\begin{emcppslisting}[emcppsignore={Invalid Descriptive Literals}]
static_assert(sizeof(int) == 4, "");  // true on virtually all machines today
static_assert(-0b1010 == -10, "");    // as if: (ù{\codeincomments{0 - 0b1010 == 0 - 10}}ù)
static_assert( 0b0111...[ 24 1-bits]...1111       //   signed
           != -0b0111...[ 24 1-bits]...1111, ""); //   signed

static_assert( 0b1000...[ 24 0-bits]...0000       // unsigned
           != -0b1000...[ 24 0-bits]...0000, ""); // unsigned
\end{emcppslisting}



%\newpage%%%%%

\subsection[Use Cases]{Use Cases}\label{use-cases}

\subsubsection[Bit masking and bitwise operations]{Bit masking and bitwise operations}\label{bit-masking-and-bitwise-operations}

Prior to the introduction of binary literals, hexadecimal and octal literals were commonly used to represent bit masks or
specific bit constants in source code.
As an example, consider a
function that returns the least significant four bits of a given
\lstinline!unsigned!~\lstinline!int! value:

\begin{emcppslisting}
unsigned int lastFourBits(unsigned int value)
{
    return value & 0xFu;
}
\end{emcppslisting}


\noindent The correctness of the \emph{bitwise and} operation above might not be
immediately obvious to a developer inexperienced with
hexadecimal literals. In contrast, using a binary literal more directly
states our intent to mask all but the four least-significant bits of the
input:

%\begin{emcppslisting}
%unsigned int lastFourBits(unsigned int value)
%{
%    return value & 0b1111u;  // The (ù{\codeincomments{u}}ù) literal suffix here is entirely optional.
%}
%\end{emcppslisting}
\begin{emcppslisting}
unsigned int lastFourBits(unsigned int value)
{
    return value & 0b1111u;
}
\end{emcppslisting}



Similarly, other bitwise operations, such as setting or getting
individual bits, might benefit from the use of binary literals. For
instance, consider a set of flags used to represent the state of an
avatar in a game:
%\enlargethispage{2ex}

%\begin{emcppslisting}
%struct AvatarStateFlags
%{
%    enum Enum
%    {
%        e_ON_GROUND    = 0b0001,
%        e_INVULNERABLE = 0b0010,
%        e_INVISIBLE    = 0b0100,
%        e_SWIMMING     = 0b1000,
%    };
%};
%
%class Avatar
%{
%    unsigned char d_state;  // power set of possible state flags
%
%public:
%    bool isOnGround() const
%    {
%        return d_flags & AvatarStateFlags::e_ON_GROUND;
%    }
%
%    // ...
%};
%\end{emcppslisting}
\begin{emcppslisting}
struct AvatarStateFlags
{
    enum Enum
    {
        e_ON_GROUND    = 0b0001,
        e_INVULNERABLE = 0b0010,
        e_INVISIBLE    = 0b0100,
        e_SWIMMING     = 0b1000,
    };
};

class Avatar
{
    unsigned char d_state;

public:
    bool isOnGround() const
    {
        return d_state & AvatarStateFlags::e_ON_GROUND;
    }

    // ...
};
\end{emcppslisting}

\noindent Note that the choice of using a nested classic \lstinline!enum! was deliberate;
see\linebreak[4] \featureref{\locationc}{enumclass}.


\subsubsection[Replicating constant binary data]{Replicating constant binary data}\label{replicating-constant-binary-data}

Especially in the context of \romeogloss{embedded development} or emulation,
a programmer will commonly write code that needs to deal
with specific ``magic'' constants (e.g., provided as part of the
specification of a CPU or virtual machine) that must be incorporated in
the program's source code. Depending on the original format of such
constants, a binary representation can be the most convenient or most
easily understandable one.

As an example, consider a function decoding instructions of a virtual
machine whose opcodes are specified in binary format:

%\begin{emcppslisting}
%#include <cstdint>  // (ù{\codeincomments{std::uint8\_t}}ù)
%
%void VirtualMachine::decodeInstruction(std::uint8_t instruction)
%{
%    switch(instruction)
%    {
%        case 0b00000000u:  // no-op
%            break;
%
%        case 0b00000001u:  // (ù{\codeincomments{add(register0, register1)}}ù)
%            d_register0 += d_register1;
%            break;
%
%        case 0b00000010u:  // (ù{\codeincomments{jmp(register0)}}ù)
%            jumpTo(d_register0);
%            break;
%
%        // ...
%    }
%}
%\end{emcppslisting}
\begin{emcppshiddenlisting}[emcppsbatch=e2]
// needed here for virtual machine, also below for visibility
#include <cstdint>  // // (ù{\codeincomments{std::uint8\_t}}ù)
struct VirtualMachine {
    std::uint8_t d_register0;
    std::uint8_t d_register1;
    void decodeInstruction(std::uint8_t instruction);
    void jumpTo(std::uint8_t address);
};
\end{emcppshiddenlisting}
\begin{emcppslisting}[emcppsbatch=e2]
#include <cstdint>  // (ù{\codeincomments{std::uint8\_t}}ù)

void VirtualMachine::decodeInstruction(std::uint8_t instruction)
{
    switch (instruction)
    {
        case 0b00000000u:  // no-op
            break;

        case 0b00000001u:  // (ù{\codeincomments{add(register0, register1)}}ù)
            d_register0 += d_register1;
            break;

        case 0b00000010u:  // (ù{\codeincomments{jmp(register0)}}ù)
            jumpTo(d_register0);
            break;

        // ...
    }
}
\end{emcppslisting}

\noindent Replicating the same binary constant specified as part of the CPU's or
virtual machine's manual or documentation directly in the source avoids the need to
mentally convert such constant data to and from, say, a hexadecimal
number.

Binary literals are also suitable for capturing bitmaps. For instance,
consider a bitmap representing the uppercase letter \textit{C}\/:

\begin{emcppslisting}
const unsigned char letterBitmap_C[] =
{
    0b00011111,
    0b01100000,
    0b10000000,
    0b10000000,
    0b10000000,
    0b01100000,
    0b00011111
};
\end{emcppslisting}


\noindent Using \emph{binary} literals makes the shape of the image that the
bitmap represents apparent directly in the source code.

\subsection[Potential Pitfalls]{Potential Pitfalls}\label{potential-pitfalls}

\hspace*{\fill}

\subsection[Annoyances]{Annoyances}\label{annoyances}

\hspace*{\fill}

\subsection[See Also]{See Also}\label{see-also}

\begin{itemize}
\item{%``\titleref{digitseparator}" on page~\pageref{digitseparator}{\seealsodelim}
%\seealsoref{digitseparator}{\locationb}Long binary literals are made much more readable by grouping digits visually.}
\seealsoref{digitseparator}{\seealsolocationb}explains grouping digits visually to make long binary literals much more readable.}
\end{itemize}

\subsection[Further Reading]{Further Reading}\label{further-reading}

\begin{itemize}
\item{\cite{informit}}
\end{itemize}




\newpage
%\section[{\tt deprecated}]{The {\SecCode [[deprecated]]} Attribute}\label{deprecated}\label{the-standard-[[deprecated]]-attribute}
% 16 Feb 2021 JMB; code compiles

\emcppsFeature{
    short={\lstinline!deprecated!},
    tocshort={\TOCCode deprecated},
    long={The {\SecCode [[deprecated]]} Attribute},
    toclong={The \lstinline![[deprecated]]! Attribute},
    rhshort={\RHCode deprecated},
}{the-standard-[[deprecated]]-attribute}
\label{deprecated}
\setcounter{table}{0}
\setcounter{footnote}{0}
\setcounter{lstlisting}{0}
%\section[{\tt deprecated}]{The {\SecCode [[deprecated]]} Attribute}\label{deprecated}\label{the-standard-[[deprecated]]-attribute}



The \lstinline![[deprecated]]! attribute discourages the use of a decorated
\romeogloss{entity}, typically via the emission of a compiler warning.

\subsection[Description]{Description}\label{description}

The standard \lstinline![[deprecated]]! attribute is used to portably
indicate that a particular \romeogloss{entity} is no longer recommended and
to actively discourage its use. Such deprecation typically follows the
introduction of alternative constructs that are
superior to the original one, providing time for clients to migrate to
them (\emph{asynchronously}{\cprotect\footnote{A process for ongoing
improvement of legacy codebases, sometimes referred to as
\romeogloss{continuous refactoring}, often allows time for clients to
migrate --- on their own respective schedules and time frames --- from
existing \emph{deprecated} constructs to newer ones, rather than
having every client change in lock step. Allowing clients time to move
\emph{asynchronously} to newer alternatives is often the only viable
approach unless (1) the codebase is a closed system, (2) all of the
relevant code is governed by a single authority, and (3) there is a mechanical way to make the change.}}) before the deprecated
one is removed in some subsequent release. Although not strictly
required, the Standard explicitly encourages{\cprotect\footnote{The C++
Standard characterizes what constitutes a well-formed program, but
compiler vendors require a great deal of leeway to facilitate the
needs of their users. In case any feature induces warnings,
command-line options are typically available to disable those warnings
(\lstinline!-Wno-deprecated! in GCC) or methods are in place to suppress those warnings
locally (e.g.,
  \lstinline!#pragma!~\lstinline!GCC!~\lstinline!diagnostic!~\lstinline!ignored!~\lstinline!"-Wdeprecated"!).}}
conforming compilers to produce a diagnostic message in case a program
refers to any \romeogloss{entity} to which the \lstinline![[deprecated]]!
attribute appertains. For instance, most popular compilers emit a
warning whenever a \lstinline![[deprecated]]! function or
object{\cprotect\footnote{The \lstinline![[deprecated]]! attribute can be
used portably to decorate other entities: \lstinline!class!,
\lstinline!struct!, \lstinline!union!, type alias, variable, data member,
function, enumeration, template specialization. Applying
\lstinline![[deprecated]]! to a specific enumerator or namespace,
however, is guaranteed to be supported only since C++17; see
  \cite{smith15a} for more information.}} is used:

\begin{emcppslisting}
               void f();
[[deprecated]] void g();

               int a;
[[deprecated]] int b;

void h()
{
    f();
    g();  // Warning: (ù{\codeincomments{g}}ù) is deprecated.
    a;
    b;    // Warning: (ù{\codeincomments{b}}ù) is deprecated.
}
\end{emcppslisting}
    
\noindent A programmer can supply a \romeogloss{string literal} as an
argument to the \mbox{\lstinline![[deprecated]]!} attribute (e.g.,
\mbox{\lstinline![[deprecated("message")]]!}) to inform human users regarding the
reason for the deprecation:

\begin{emcppslisting}
[[deprecated("too slow, use (ù{\codeincomments{algo1}}ù) instead")]] void algo0();
                                              void algo1();

void f()
{
    algo0();  // Warning: (ù{\codeincomments{algo0}}ù) is deprecated; too slow, use (ù{\codeincomments{algo1}}ù) instead.
    algo1();
}
\end{emcppslisting}
    
\noindent An \romeogloss{entity} that is initially \emph{declared} without
\lstinline![[deprecated]]! can later be redeclared with the attribute and
vice versa:

\begin{emcppslisting}
void f();
void g0() { f(); }  // OK, likely no warnings

[[deprecated]] void f();
void g1() { f(); }  // Warning: (ù{\codeincomments{f}}ù) is deprecated.

void f();
void g2() { f(); }  // Warning: (ù{\codeincomments{f}}ù) is deprecated (still).
\end{emcppslisting}
    
\noindent As shown in \lstinline!g2! (above), redeclaring an \romeogloss{entity} that was
previously decorated with\linebreak[4] \mbox{\lstinline![[deprecated]]!} without the attribute leaves the entity still deprecated.

\subsection[Use Cases]{Use Cases}\label{use-cases}

\subsubsection[Discouraging use of an obsolete or unsafe \romeogloss{entity}]{Discouraging use of an obsolete or unsafe \romeogloss{entity}}\label{discouraging-use-of-an-obsolete-or-unsafe-entity}

Decorating any \romeogloss{entity} with the \lstinline![[deprecated]]! attribute serves both
to indicate a particular feature should not be used in the future and
to actively encourage migration of existing uses to a better
alternative. Obsolescence, lack of safety, and poor performance are
common motivators for deprecation.

As an example of productive deprecation, consider the
\lstinline!RandomGenerator! class having a static \lstinline!nextRandom!
member function to generate random numbers:

\begin{emcppslisting}
struct RandomGenerator
{
    static int nextRandom();
        // Generate a random value between 0 and 32767 (inclusive).
};
\end{emcppslisting}
    
\noindent Although such a simple random number generator can be very useful, it
might become unsuitable for heavy use because good pseudorandom number
generation requires more state (and the overhead of synchronizing such
state for a single \lstinline!static! function can be a significant
performance bottleneck) while good random number generation requires
potentially very high overhead access to external sources of
entropy.{\cprotect\footnote{The C Standard Library provides
\lstinline!rand!, available in C++ through the \lstinline!<cstdlib>! header.
It has similar issues to our \lstinline!RandomGenerator::nextRandom!
function, and similarly developers are guided to use the facilities
  provided in the \lstinline!<random>! header since C++11.}} One
solution is to provide an alternative random number generator that
maintains more state, allows users to decide where to store that state
(the random number generator objects), and overall offers more
flexibility for clients. The downside of such a change is that it comes
with a functionally distinct API, requiring that users update their code
to move away from the inferior solution:

\begin{emcppslisting}
class StatefulRandomGenerator
{
    // ... (internal state of a quality pseudorandom number generator) ...

public:
    int nextRandom();
       // Generate a quality random value between 0 and 32767 (inclusive).
};
\end{emcppslisting}
    
\noindent Any user of the original random number generator can migrate to the new
facility with little effort, but that is not a completely trivial
operation, and migration will take some time before the original feature is no
longer in use. The empathic maintainers of \lstinline!RandomGenerator! can
decide to use the
\lstinline![[deprecated]]! attribute to discourage continued use
of\linebreak%%%%%%%
 \lstinline!RandomGenerator::nextRandom()! instead of removing it completely:

\begin{emcppslisting}
struct RandomGenerator
{
    [[deprecated("Use (ù{\codeincomments{StatefulRandomGenerator class}}ù) instead.")]]
    static int nextRandom();
        // ...
};
\end{emcppslisting}
    
\noindent By using \lstinline![[deprecated]]! as shown above, existing clients of
\lstinline!RandomGenerator! are informed that a superior
alternative, \lstinline!BetterRandomGenerator!, is available, yet they are granted time to
migrate their code to the new solution rather than having their code broken by the
removal of the old solution. When clients are notified of the
deprecation (thanks to a compiler diagnostic), they can schedule time to rewrite their applications to consume the new interface.{\cprotect\footnote{\romeogloss{Continuous refactoring} is
an essential responsibility of a development organization, and
deciding when to go back and fix what's suboptimal instead of writing
new code that will please users and contribute more immediately to the
bottom line will forever be a source of tension. Allowing disparate
development teams to address such improvements in their own respective
time frames (perhaps subject to some reasonable overall deadline
date) is a proven real-world practical way of ameliorating this
  tension.}}

\subsection[Potential Pitfalls]{Potential Pitfalls}\label{potential-pitfalls}

\subsubsection[Interaction with treating warnings as errors]{Interaction with treating warnings as errors}\label{interaction-with--werror-(e.g,-gcc,-clang)-or-/wx-(msvc)}

In some code bases, compiler warnings are promoted to errors using compiler flags, such as \lstinline!-Werror! for GCC and Clang or \lstinline!/WX! for MSVC, to ensure that their builds are warning-clean. For such code bases, use of the \lstinline![[deprecated]]! attribute by their dependencies as part of the API might introduce unexpected compilation failures. 

Having the compilation process completely stopped due to use of a
deprecated \romeogloss{entity} defeats the purpose of the attribute because users of such an \romeogloss{entity} are given no time to adapt their code
to use a newer alternative. On GCC and Clang, users can selectively demote deprecation errors back to warnings by using the
\lstinline!-Wno-error=deprecated-declarations! compiler flag. On MSVC,
however, such demotion of warnings is not possible, and the available workarounds, such as entirely disabling the effects of the \lstinline!/WX! flag or the deprecation diagnostics using the  \lstinline!-wd4996! flag, are often unsuitable.

Furthermore, this interaction between \lstinline![[deprecated]]! and
treating warnings as errors makes it impossible for owners of a low-level library
to deprecate a function when releasing their code requires that they do
not break the ability for \emph{any} of their higher-level clients to compile;
a single client using the to-be-deprecated function in a code base that treats warnings as errors prevents the release of the code that uses the
\lstinline![[deprecated]]! attribute. With the frequent advice given in practice to  aggressively treat warnings as errors, the use of
\lstinline![[deprecated]]! might be completely unfeasible.

\subsection[Annoyances]{Annoyances}\label{annoyances}

None so far

\subsection[See Also]{See Also}\label{see-also}

None so far

\subsection[Further Reading]{Further Reading}\label{further-reading}

None so far





\newpage
%\section[Digit Separators]{The Digit Separator: {\SecCode '}}
%\label{digitseparator} % used with placeholder file
%\label{digit-separators} % label had changed when actual file arrived
% 5 Feb 2021, copyedits in and proofed
% 8 Feb 2021 JMB, compilation fixes
% 22 March 2021, in FPPs

\emcppsFeature{
    toclong={The Digit Separator: \lstinline!'!},
    tocshort={Digit Separators},
    short={Digit Separators},
    long={The Digit Separator: {\SecCode '}},
}{digitseparator}
%\section[Digit Separators]{The Digit Separator: {\SecCode '}}
%\label{digitseparator} % used with placeholder file
\label{digit-separators} % label had changed when actual file arrived
\setcounter{table}{0}
\setcounter{footnote}{0}
\setcounter{lstlisting}{0}



A digit separator is a single-character token (\lstinline!'!) that can appear as part of a
numeric literal without altering its value.

\subsection[Description]{Description}\label{description}

A digit separator --- i.e., an instance of the single-quote
character (\lstinline!'!) --- may be placed anywhere within a numeric
literal to visually separate its digits without affecting its value:

%%%%%%%%%%%%%%%% ALL code snippets in this subsection
%%%%%%%%  setting this code coloring manually
% Everything after // (including //) itself is a comment, and should be light blue (like -12345).
% `int`, `unsigned int`, `long`, `long long`, `float`, `double`, `long double`, `int`, `int`, `int` should be in the reddish color like the first occurrence of `int`.
% The variable names (e.g., i,u,j,k,f,d,e,hex,oct,bin) should be black.
% All equal signs should be black.
% All numeric literals (e.g., -12'345, 1'000'000u) should be black as well.
%%%%%%%%%%%%%%%%
% 8 Feb 2021: Our final design does not use color in this way. Removing these color and shape commands and asking authors to review.

%\begin{emcppslisting}
%int          (ù{\color{black} i = -12'345;}ù)                    (ù{\itshape\color{skyblue}// same as}ù) (ù{\codeincomments{-12345}}ù)
%(ù{\color{maroon}unsigned int}ù) (ù{\color{black}u = 1'000'000u;}ù)                 (ù{\itshape\color{skyblue}// same as}ù) (ù{\codeincomments{1000000u}}ù)
%(ù{\color{maroon}long}ù)         (ù{\color{black}j = 500'000L;}ù)                   // same as (ù{\codeincomments{500000L}}ù)
%(ù{\color{maroon}long long}ù)    (ù{\color{black}k = 9'223'372'036'854'775'807;}ù)  // same as (ù{\codeincomments{9223372036854775807}}ù)
%(ù{\color{maroon}float}ù)        (ù{\color{black}f = 3.14159'26535f}ù)              // same as (ù{\codeincomments{3.1415926535f}}ù)
%(ù{\color{maroon}double}ù)       (ù{\color{black}d = 3.14159'26535'89793;}ù)        (ù{\itshape\color{skyblue}// same as}ù) (ù{\codeincomments{3.141592653589793}}ù)
%(ù{\color{maroon}long double}ù)  (ù{\color{black}e = 20'812.80745'23204}ù)          (ù{\itshape\color{skyblue}// same as}ù) (ù{\codeincomments{20812.8074523204}}ù)
%(ù{\color{maroon}int}ù)        (ù{\color{black}hex = 0x8C'25'00'F9;}ù)              // same as (ù{\codeincomments{0x8C2500F9}}ù)
%int        (ù{\color{black}oct = 044'73'26;}ù)                  // same as (ù{\codeincomments{0447326}}ù)
%int        (ù{\color{black}bin = 0b1001'0110'1010'0111;}ù)      (ù{\itshape\color{skyblue}// same as}ù) (ù{\codeincomments{0b1001011010100111}}ù)
%\end{emcppslisting}
\begin{emcppslisting}[emcppsstandards={c++14}]
int          i = -12'345;                    (ù{\commentcolor{// same as}ù) (ù{\codeincomments{-12345}}ù)
unsigned int u = 1'000'000u;                 (ù{\commentcolor{// same as}ù) (ù{\codeincomments{1000000u}}ù)
long         j = 500'000L;                   // same as (ù{\codeincomments{500000L}}ù)
long long    k = 9'223'372'036'854'775'807;  // same as (ù{\codeincomments{9223372036854775807}}ù)
float        f = 3.14159'26535f;             (ù{\commentcolor{// same as}ù) (ù{\codeincomments{3.1415926535f}}ù)
double       d = 3.14159'26535'89793;        (ù{\commentcolor{// same as}ù) (ù{\codeincomments{3.141592653589793}}ù)
long double  e = 20'812.80745'23204;         (ù{\commentcolor{// same as}ù) (ù{\codeincomments{20812.8074523204}}ù)
int        hex = 0x8C'25'00'F9;              // same as (ù{\codeincomments{0x8C2500F9}}ù)
int        oct = 044'73'26;                  // same as (ù{\codeincomments{0447326}}ù)
int        bin = 0b1001'0110'1010'0111;      (ù{\commentcolor{// same as}ù) (ù{\codeincomments{0b1001011010100111}}ù)
\end{emcppslisting}

\noindent Multiple digit separators within a single literal are allowed, but they
cannot be contiguous, nor can they appear either before or after the
\emph{numeric} part (i.e., digit sequence) of the
literal:

%\begin{emcppslisting}
%int (ù{\color{black}e0 = 10''00;}ù)  // Error, consecutive digit separators
%int (ù{\color{black}e1 = -'1000;}ù)  // Error, before numeric part
%(ù{\color{maroon}int}ù) (ù{\color{black}e2 = 1000'u;}ù)  // Error, after numeric part
%int (ù{\color{black}e3 = 0x'abc;}ù)  // Error, before numeric part
%(ù{\color{maroon}int}ù) (ù{\color{black}e4 = 0'xdef;}ù)  // Error, way before numeric part
%int (ù{\color{black}e5 = 0'89;}ù)    // Error, non-octal digits
%(ù{\color{maroon}int}ù) (ù{\color{black}e6 = 0'67;}ù)    // OK, valid octal literal
%\end{emcppslisting}
\begin{emcppslisting}[emcppsstandards={c++14},emcppserrorlines={2,4,6}]
int e0 = 10''00;  // Error, consecutive digit separators
int e1 = -'1000;  (ù{\commentcolor{// Error, before numeric part}ù)
int e2 = 1000'u;  // Error, after numeric part
int e3 = 0x'abc;  (ù{\commentcolor{// Error, before numeric part}ù)
int e4 = 0'xdef;  // Error, way before numeric part
int e5 = 0'89;    (ù{\commentcolor{// Error, non-octal digits}ù)
int e6 = 0'67;    // OK, valid octal literal
\end{emcppslisting}

\noindent Although the leading \lstinline!0x! and
\lstinline!0b! prefixes for hexadecimal and binary literals,
respectively, are not considered part of the \emph{numeric} part of
  the literal, a leading \lstinline!0! in an octal literal is. As a side note, remember that on some platforms an integer literal
that is too large to fit in a \lstinline!long!~\lstinline!long!~\lstinline!int! but that does fit in an
\lstinline!unsigned!~\lstinline!long!~\lstinline!long!~\lstinline!int! might
generate a warning{\cprotect\footnote{Tested on GCC 7.4.0.}}:

%\begin{emcppslisting}
%(ù{\color{maroon}unsigned long long}ù) (ù{\color{black}big1 = 9'223'372'036'854'775'808;}ù)  // (ù{\codeincomments{2\^{}63}}ù)
%    // warning: integer constant is so large that it is an
%    // unsigned long long big1 = 9'223'372'036'854'775'808;
%    //                           ^~~~~~~~~~~~~~~~~~~~~~~~~~
%\end{emcppslisting}
\begin{emcppslisting}[emcppsstandards={c++14}]
unsigned long long big1 = 9'223'372'036'854'775'808;  // (ù{\codeincomments{2\^{}63}}ù)
    // warning: integer constant is so large that it is an
    // unsigned long long big1 = 9'223'372'036'854'775'808;
    //                           ^~~~~~~~~~~~~~~~~~~~~~~~~~
\end{emcppslisting}

\newpage%%%%%%
\noindent Such warnings can typically be suppressed by adding a \lstinline!ull!
suffix to the literal:

%\begin{emcppslisting}
%(ù{\color{maroon}unsigned long long}ù) (ù{\color{black}big2 = 9'223'372'036'854'775'808ull;}ù)  // OK
%\end{emcppslisting}
\begin{emcppslisting}[emcppsstandards={c++14}]
unsigned long long big2 = 9'223'372'036'854'775'808ull;  // OK
\end{emcppslisting}

\noindent Warnings like the one above, however, are not typical when the implied
precision of a floating-point literal exceeds what can be represented:

%\begin{emcppslisting}
%(ù{\color{maroon}float}ù) (ù{\color{black}reallyPrecise = 3.141'592'653'589'793'238'462'643'383'279'502'884;}ù)  // OK
%    // Everything after (ù{\codeincomments{3.141'592'6}}ù) is typically ignored silently.
%\end{emcppslisting}
\begin{emcppslisting}[emcppsstandards={c++14}]
float reallyPrecise = 3.141'592'653'589'793'238'462'643'383'279'502'884;  (ù{\commentcolor{// OK}ù)
    (ù{\commentcolor{// Everything after}ù) (ù{\codeincomments{3.141'592'6}}ù) (ù{\commentcolor{is typically ignored silently.}ù)
\end{emcppslisting}

\noindent For more information, see \intraref{appendix-digitsep}{appendix:-silent-loss-of-precision-in-floating-point-literals}.
%{\it\titleref{appendix:-silent-loss-of-precision-in-floating-point-literals}} on page~\pageref{appendix:-silent-loss-of-precision-in-floating-point-literals}.

\subsection[Use Cases]{Use Cases}\label{use-cases}

\subsubsection[Grouping digits together in large constants]{Grouping digits together in large constants}\label{grouping-digits-together-in-large-constants}

When embedding large constants in source code, consistently placing
digit separators (e.g., every thousand) might improve readability, as
illustrated in Table~\ref{digitseparator-table1}.

\begin{table}[h!]
\begin{center}
\begin{threeparttable}
\caption{Use of digit separators to improve readability}\label{digitseparator-table1}\vspace{1.5ex}
{\small \begin{tabular}{c|c}\thickhline
\rowcolor[gray]{.9}{\sffamily\bfseries Without Digit Separator} & {\sffamily\bfseries With Digit Separators} \\ \hline
\multicolumn{1}{r|}{\lstinline!10000!} & \multicolumn{1}{r}{\lstinline!10'000!} \\ \hline
\multicolumn{1}{r|}{\lstinline!100000!} & \multicolumn{1}{r}{\lstinline!100'000!} \\ \hline
\multicolumn{1}{r|}{\lstinline!1000000!} & \multicolumn{1}{r}{\lstinline!1'000'000!} \\ \hline
\multicolumn{1}{r|}{\lstinline!1000000000!} & \multicolumn{1}{r}{\lstinline!1'000'000'000!} \\ \hline
\multicolumn{1}{r|}{\lstinline!18446744073709551615ULL!} & \multicolumn{1}{r}{\lstinline!18'446'744'073'709'551'615ULL!} \\ \hline
\multicolumn{1}{r|}{\lstinline!1000000.123456!} & \multicolumn{1}{r}{\lstinline!1'000'000.123'456!} \\ \hline
\multicolumn{1}{r|}{\lstinline!3.141592653589793238462L!} & \multicolumn{1}{r}{\lstinline!3.141'592'653'589'793'238'462L!} \\ \thickhline
\end{tabular}
} % end small
\end{threeparttable}
\end{center}
\end{table}
\noindent Use of digit separators is especially useful with binary
literals to group bits in octets (\emcppsgloss{bytes}) or quartets
(\emcppsgloss{nibbles}), as shown in Table~\ref{digitseparator-table2}. In addition, using a binary literal with digits grouped in triplets instead of an octal literal to represent UNIX file permissions might improve code readability --- e.g., \lstinline!0b111'101'101! instead of \lstinline!0755!.\begin{table}[h!]
\begin{center}
\begin{threeparttable}
\caption{Use of digit separators in binary data}\label{digitseparator-table2}\vspace{1.5ex}
{\small \begin{tabular}{c|c}
\thickhline
\rowcolor[gray]{.9} {\sffamily\bfseries Without Digit Separator} & {\sffamily\bfseries With Digit Separators} \\\hline
\multicolumn{1}{r|}{\lstinline!0b1100110011001100!} & \multicolumn{1}{r}{\lstinline!0b1100'1100'1100'1100!} \\ \hline
\multicolumn{1}{r|}{\lstinline!0b0110011101011011!} & \multicolumn{1}{r}{\lstinline!0b0110'0111'0101'1011!} \\ \hline
\multicolumn{1}{r|}{\lstinline!0b1100110010101010!} & \multicolumn{1}{r}{\lstinline!0b11001100'10101010!} \\
\thickhline
\end{tabular}
} % end small
\end{threeparttable}
\end{center}
\end{table}

\newpage%%%%%%
\subsection[Potential Pitfalls]{Potential Pitfalls}\label{potential-pitfalls}

\hspace{\fill}

\subsection[See Also]{See Also}\label{see-also}

\begin{itemize}
\item{\seealsoref{binary-literals}{\seealsolocationb}represents a binary constant for which digit separators are commonly used to group bits in octets (\emcppsgloss{bytes}) or quartets (\emcppsgloss{nibbles})}
\end{itemize}

\subsection[Further Reading]{Further Reading}\label{further-reading}

\begin{itemize}
\item{William Kahan. ``Lecture Notes on the Status of
IEEE Standard 754 for Binary Floating-Point Arithmetic,'' \cite{kahan97}}
\item{{\textit{IEEE Standard for Floating-Point Arithmetic}}, \cite{ieee19}}
\end{itemize}

\subsection[Appendix]{Appendix}\label{appendix-digitsep}

\subsubsection[Silent loss of precision in floating-point literals]{Silent loss of precision in floating-point literals}\label{appendix:-silent-loss-of-precision-in-floating-point-literals}

Just because we can keep track of precision in floating-point literals
doesn't mean that the compiler can. As an aside, it is worth
pointing out that the binary representation of floating-point types is
not mandated by the Standard, nor are the precise minimums on the ranges
and precisions they must support. Although the C++ Standard says little
that is normative, the macros in \lstinline!<cfloat>! are defined by
reference to the C Standard.\footnote{\cite{iso20b}, section~6.8.2, ``Fundamental types [basic.fundamental]," pp. 73--75; section~17.3.5.2, ``\texttt{numeric\_limits} members [numeric.limits.members]," pp. 513--516; and section~17.3.7, ``Header \texttt{<cfloat>} synopsis [cfloat.syn]," p. 519}\footnote{\cite{iso18b}, section~7.7,
``Characteristics of floating types \texttt{<float.h>}," p. 157}
% 5.2.4.2.2 Characteristics of floating types \texttt{<float.h>}

There are, however, normal and customary minimums that one can typically
rely upon in practice. On conforming compilers that employ the IEEE 754
floating-point standard representation\footnote{\cite{ieee19}} (as most do),
a \lstinline!float! can typically represent up to 7 significant decimal
digits accurately, while a \lstinline!double! typically has nearly 15 decimal
digits of precision. For any given program,
\lstinline!long!~\lstinline!double! is required to hold whatever a
\lstinline!double! can hold, but is typically larger (e.g., 10, 12, or 16
bytes) and typically adds at least 5 decimal digits of precision (i.e.,
supports a total of at last 20 decimal digits). A notable exception is Microsoft Visual C++ where \lstinline!long! \lstinline!double! is a distinct type whose representation is identical to \lstinline!double!.\footnote{\cite{microsoft-types}} A table summarizing
typical precisions for various IEEE-conforming floating-point types is
presented for convenient reference in Table~\ref{digitseparator-table3}. The actual bounds on a
given platform can be found using the standard
\lstinline!std::numeric_limits! class template found in \lstinline!<limits>!.
\pagebreak%%%%%



\begin{table}[t]
\begin{center}
\vspace*{-3ex}
\begin{threeparttable}
\caption{Available precisions for various IEEE-754 floating-point types}\label{digitseparator-table3}\vspace{1.5ex}
{\small \begin{tabular}{c|c|c|c|c|c}\thickhline
\rowcolor[gray]{.9}   {\sffamily\bfseries Name} & {\sffamily\bfseries Common} &
{\sffamily\bfseries Significant } & {\sffamily\bfseries Decimal }& {\sffamily\bfseries Exponent } & {\sffamily\bfseries Dynamic}\\
\rowcolor[gray]{.9}    & {\sffamily\bfseries Name} &
{\sffamily\bfseries Bits\tnote{a}} & {\sffamily\bfseries Digits}& {\sffamily\bfseries Bits} & {\sffamily\bfseries Range}\\ \hline
binary16 & Half precision & \multicolumn{1}{r|}{11} & \multicolumn{1}{r|}{3.31} & \multicolumn{1}{r|}{5} & \multicolumn{1}{r}{$\sim6.50\text{e}5$}\\ \hline
binary32 & Single precision & \multicolumn{1}{r|}{24} & \multicolumn{1}{r|}{7.22} & \multicolumn{1}{r|}{8} & \multicolumn{1}{r}{$\sim3.4\text{e}38$}\\ \hline
binary64 & Double precision & \multicolumn{1}{r|}{53} & \multicolumn{1}{r|}{15.95} & \multicolumn{1}{r|}{11} & \multicolumn{1}{r}{$\sim1.\text{e}308$}\\ \hline
\rule{0pt}{3.5mm}binary80 & Extended precision & \multicolumn{1}{r|}{69} & \multicolumn{1}{r|}{20.77} & \multicolumn{1}{r|}{11} & \multicolumn{1}{r}{$\sim10^{308}$}\\ \hline
\rule{0pt}{3.5mm}binary128 & Quadruple precision & \multicolumn{1}{r|}{113} & \multicolumn{1}{r|}{34.02} & \multicolumn{1}{r|}{15} & \multicolumn{1}{r}{$\sim10^{4932}$}\\ 
\thickhline
\end{tabular}
} % end small
\begin{tablenotes}{\footnotesize
\item[a]{Note that the most significant bit of the \emcppsgloss{mantissa}
is always a \lstinline!1! for normalized numbers, and 0 for denormalized ones and, hence, is not stored explicitly. This leaves 1
additional bit to represent the sign of the overall floating-point value
(the sign of the exponent is encoded using \emcppsgloss[excess n]{excess-$\mathbf{n}$}
notation).}
} % end footnotesize
\end{tablenotes} % we end the table footnotes
\end{threeparttable}
\end{center}
\end{table}

Determining the minimum number of decimal digits needed to accurately
approximate a transcendental value, such as $\pi$, for a given type on
a given platform can be tricky (requiring some binary-search-like
detective work), which is likely why overshooting the precision without
warning is the default on most platforms. One way to establish that
\emph{all} of the decimal digits in a given floating-point literal are
relevant for a given floating-point type is to compare that literal and
a similar one with its least significant decimal digit
removed{\cprotect\footnote{Note that affixing the \lstinline!f!
(\emph{literal suffix}) to a floating-point literal is
equivalent to applying a \lstinline!static_cast<float>! to the
(unsuffixed) literal:

\begin{emcppslisting}[emcppsstandards={c++14}, basicstyle={\ttfamily\footnotesize}]
static_assert(3.14'159'265'358f == static_cast<float>(3.14'159'265'358));
\end{emcppslisting}\vspace*{-1ex}
      }}:

\begin{emcppslisting}[emcppsignore={Bold numbers prevent assertions from not compiling properly}]
static_assert(3.141592653(ù{\bfseries 5}ù)f != 3.141592653f, "too precise for float");
    // This assert will fire on a typical platform.

static_assert(3.14159265(ù{\bfseries 3}ù)f != 3.14159265f, "too precise for float");
    // This assert too will fire on a typical platform.

static_assert(3.1415926(ù{\bfseries 5}ù)f != 3.1415926f, "too precise for float");
    // This assert will NOT fire on a typical platform.

static_assert(3.141592(ù{\bfseries 6}ù)f != 3.141592f, "too precise for float");
    // This assert too will NOT fire on a typical platform.
\end{emcppslisting}

\noindent If the values are \emph{not} the same, then that floating-point type
can make use of the precision suggested by the original literal; if
they \emph{are} the same, however, then it is likely that the available
precision has been exceeded. Iterative use of this technique by
developers can help them to empirically narrow down the maximal number
of decimal digits a particular platform\linebreak[4]  
\newpage will support for a particular
floating-point type and value. Note, however, that because the compiler is not required to use the floating-point arithmetic of the target platform \emph{during compilation}, this approach might not be applicable for a cross-compilation scenario.

One final useful tidbit pertains to the safe (lossless) conversion
between binary and decimal floating-point representations; note that
``Single'' (below) corresponds to a single-precision IEEE-754-conforming
(32-bit) \lstinline!float!{\cprotect\footnote{\cite{kahan97}, section ``Representable Numbers," p.~4}}:

\begin{quote}
If a decimal string with at most 6 sig. dec. is converted to Single and then converted back to the
 same number of sig. dec., then the final string should match the original. Also, ...

 If a Single Precision floating-point number is converted to a decimal string with at least 9 sig.
 dec. and then converted back to Single, then the final number must match the original.
\end{quote}

\noindent The ranges corresponding to 6--9 for a single-precision (32-bit) \lstinline!float! (described above), when applied to a double-precision (64-bit) \lstinline!double! and a quad-precision (128-bit) \linebreak[4]%%%%%
\lstinline!long!~\lstinline!double!, are 15--17 and 33--36, respectively.





\newpage
%\section[Variable Templates]{Templated Variable Declarations/Definitions}
%\label{variable-templates}
% 12 Feb 2021, revisions in; ready for Josh's code fixes
% 16 Feb 2021 JMB; code compiles
% 22 feb 2021, copyedits in and proofed
% 22 March 2021, in FPPs


\emcppsFeature{
    short={Variable Templates},
    long={Templated Variable Declarations/Definitions},
}{variable-templates}
\setcounter{table}{0}
\setcounter{footnote}{0}
\setcounter{lstlisting}{0}
%\section[Variable Templates]{Templated Variable Declarations/Definitions}\label{variable-templates}


Variable templates extend traditional template syntax to define, in namespace or
class (but not function) scope, a family of like-named variables that
can subsequently be instantiated explicitly.

\subsection[Description]{Description}\label{variabletemplate-description}

%By beginning a variable declaration with the familiar
%\textbf{template-head} syntax --- e.g.,
%\texttt{template}~\texttt{<typename}~\texttt{T>} --- we can create a
%\emph{variable template}, which defines a family of variables having the
%same name (e.g., \texttt{typeid}):
%
%\begin{emcppslisting}
%template <typename> int typeId;  // template variable defined at file scope
%\end{emcppslisting}
%    
%\noindent Like any other kind of template, a variable template can be instantiated
%(explicitly) by providing an appropriate number (one or more) of type or
%non-type arguments:
%
%\begin{emcppslisting}
%void f1()
%{
%    typeId<bool> = -1;    // (ù{\codeincomments{typeId<bool>}}ù) is an (ù{\codeincomments{int}}ù)
%    typeId<char> = 1000;  // (ù{\codeincomments{typeId<char>}}ù) is an (ù{\codeincomments{int}}ù)
%    typeId<void> = -666;  // (ù{\codeincomments{typeId<void>}}ù) is an (ù{\codeincomments{int}}ù)
%
%    assert(typeId<bool> ==   -1);
%    assert(typeId<char> == 1000);
%    assert(typeId<void> == -666);
%}
%\end{emcppslisting}
%    
%\noindent In the example above, the type of each instantiated variable --- i.e.,
%\texttt{typeId<bool>} and \texttt{typeId<char>} --- is \texttt{int}.
%Such need not be the case{\cprotect\footnote{Use of
%  \texttt{constexpr} variables would allow the instantiated
%  variables to be usable as a constant in a compile-time context (see
%  {\it\titleref{variabletemplate-use-cases}:} {\it\titleref{parametrized-constants}} on page~\pageref{parametrized-constants}).}}:
%
%\begin{emcppslisting}
%template <typename T> const T pi(3.1415926535897932385);  // distinct types
%\end{emcppslisting}
%    
%\noindent In the example above, the type of the instantiated non-\texttt{const}
%variable is that of its (type) argument, and its (mutable) value is
%initialized to the best approximation of \(\pi\) offered by that type:
%
%\begin{emcppslisting}
%void f2()
%{
%    bool        pi_as_bool        = 1;                      // ( 1 bit)
%    int         pi_as_int         = 3;                      // (32 bits)
%    float       pi_as_float       = 3.1415927;              // (32 bits)
%    double      pi_as_double      = 3.141592653589793;      // (64 bits)
%    long double pi_as_long_double = 3.1415926535897932385;  // (80 bits)
%
%    assert(pi<bool>        == pi_as_bool);
%    assert(pi<int>         == pi_as_int);
%    assert(pi<float>       == pi_as_float);
%    assert(pi<double>      == pi_as_double);
%    assert(pi<long double> == pi_as_long_double);
%}
%\end{emcppslisting}
%    
%\noindent For examples involving immutable variable templates, see {\it\titleref{variabletemplate-use-cases}:} {\it\titleref{parametrized-constants}} on page~\pageref{parametrized-constants}.

%%%%%%%%%%%%%%%% new material from Slava
By beginning a variable declaration with the familiar \emcppsgloss[template head]{template-head} syntax --- e.g.,\linebreak[4] \lstinline!template! \lstinline!<typename! \lstinline!T>! --- we can create a \emph{variable template}, which defines a family of variables
having the same name (e.g., \lstinline!exampleOf!):

\begin{emcppslisting}[emcppsbatch=e1]
template <typename T> T exampleOf;  // variable template defined at file scope
\end{emcppslisting}

\noindent Like any other kind of template, a variable template can be instantiated explicitly by
providing an appropriate number of type or non-type arguments:

\begin{emcppslisting}[emcppsbatch=e1]
#include <iostream>  // (ù{\codeincomments{std::cout}}ù)

void initializeExampleValues()
{
    exampleOf<int>   = -1;
    exampleOf<char>  = 'a';
    exampleOf<float> = 12.3f;
}

void printExampleValues()
{
    initializeExampleValues();
    std::cout << "int = "   << exampleOf<int>   << "; "
              << "char = "  << exampleOf<char>  << "; "
              << "float = " << exampleOf<float> << ';';

    // outputs "int = -1; char = a; float = 12.3;"
}
\end{emcppslisting}

\noindent In the example above, the type of each instantiated variable is the same as its template parameter, but this matching is not required. For example, the type might be the same for all instantiated variables or derived from its parameters, such as by adding \lstinline!const! qualification:

\newpage%%%%%%
\begin{emcppshiddenlisting}[emcppsbatch={e2,e3}]
#include <iostream>  // (ù{\codeincomments{std::cout}}ù)
\end{emcppshiddenlisting}
\begin{emcppslisting}[emcppsbatch=e2]
#include <type_traits>  // (ù{\codeincomments{std::is\_floating\_point}}ù)
#include <cassert>      // standard C (ù{\codeincomments{assert}}ù) macro

template <typename T> 
const bool sane_for_pi = std::is_floating_point<T>::value;  // same type

template <typename T> const T pi(3.1415926535897932385);    // distinct types

void testPi()
{
    assert(!sane_for_pi<bool>);
    assert(!sane_for_pi<int>);

    assert( sane_for_pi<float>);
    assert( sane_for_pi<double>);
    assert( sane_for_pi<long double>);

    const float       pi_as_float       = 3.1415927; 
    const double      pi_as_double      = 3.141592653589793; 
    const long double pi_as_long_double = 3.1415926535897932385; 

    assert(pi<float>       == pi_as_float);
    assert(pi<double>      == pi_as_double); 
    assert(pi<long double> == pi_as_long_double); 
}
\end{emcppslisting}
%%%%%%%%%%%%% new material ends

\noindent Variable templates, like \emcppsgloss[C style function]{C-style functions}, may be declared at
namespace-scope or as \lstinline!static! members of a
\lstinline!class!, \lstinline!struct!, or \lstinline!union! but are not
permitted as non\lstinline!static! members nor at all in function scope:

\enlargethispage*{-10ex}
\begin{emcppslisting}
template <typename T> T vt1;             // OK, external linkage
template <typename T> static T vt2;      // OK, internal linkage

namespace N
{
    template <typename T> T vt3;         // OK, external linkage
    template <typename T> static T vt4;  // OK, internal linkage
}

struct S
{
    template <typename T> T vt5;         // Error, not (ù{\codeincomments{static}}ù)
    template <typename T> static T vt6;  // OK, external linkage
};

void f3()  // Variable templates cannot be defined in functions.
{
    template <typename T> T vt7;         // Error
    template <typename T> static T vt8;  // Error

    vt1<bool> = true;                    // OK, to use them
    N::vt3<bool> = true;                                                        
    N::vt4<bool> = true;                                                        
    S::vt6<bool> = true;  
}
\end{emcppslisting}
    
\noindent Like other templates, variable templates may be defined with multiple
parameters consisting of arbitrary combinations of type and non-type
parameters, including a \emcppsgloss{parameter pack} in the last position:

\begin{emcppslisting}
namespace N
{
    template <typename V, int I, int J> V factor;  // namespace scope
}
\end{emcppslisting}
    
\noindent Variable templates can even be defined recursively (but see \intraref{variabletemplate-potential-pitfalls}{recursive-variable-template-initializations-require-const-or-constexpr}): 
%{\it\titleref{variabletemplate-potential-pitfalls}:} {\it\titleref{recursive-variable-template-initializations-require-const-or-constexpr}} on page~\pageref{recursive-variable-template-initializations-require-const-or-constexpr}):

\begin{emcppslisting}[emcppsbatch=e3]
namespace {
template <int N>
const int sum = N + sum<N - 1>;    // recursive general template

template <> const int sum<0> = 0;  // base case specialization
}  // close unnamed namespace

void f()
{
    std::cout << sum<4> << '\n';  // prints 10
    std::cout << sum<5> << '\n';  // prints 15
    std::cout << sum<6> << '\n';  // prints 21
}
\end{emcppslisting}
    
\noindent Note that while variable templates do not add new functionality, they significantly reduce the boilerplate associated with achieving the same goals without them.  For example, compare the definition of \lstinline!pi! above with the pre-C++14 code:

\newpage%%%%%%
\begin{emcppslisting}
// C++03 (obsolete)
#include <cassert>  // standard C (ù{\codeincomments{assert}}ù) macro

template <typename T>
struct Pi {
    static const T value;
};

template <typename T>
const T Pi<T>::value(3.1415926535897932385);  // separate definition

void testCpp03Pi()
{
    const float       piAsFloat      = 3.1415927;
    const double      piAsDouble     = 3.141592653589793;     
    const long double piAsLongDouble = 3.1415926535897932385; 

    // additional boilerplate on use ((ù{\codeincomments{::value}ù))
    assert(Pi<float>::value       == piAsFloat);
    assert(Pi<double>::value      == piAsDouble);
    assert(Pi<long double>::value == piAsLongDouble);
}
\end{emcppslisting}    

\subsection[Use Cases]{Use Cases}\label{variabletemplate-use-cases}

\subsubsection[Parameterized constants]{Parameterized constants}\label{parametrized-constants}

A common effective use of variable templates is in the definition of
type-parameterized constants. As discussed in \intrarefsimple{variabletemplate-description},
%{\it\titleref{variabletemplate-description}} on page~\pageref{variabletemplate-description}, 
the mathematical
constant $\pi$ serves as our example. Here we want to
initialize the constant as part of the variable template (the literal
chosen is the shortest decimal string to do so accurately for an 80-bit
\lstinline!long!~\lstinline!double!):
%  Section~\ref{digit-separators}, ``\titleref{digit-separators}."}}:

\begin{emcppslisting}[emcppsbatch=e4]
template <typename T>
constexpr T pi(3.1415926535897932385);
    // smallest digit sequence to accurately represent pi as a (ù{\codeincomments{long double}}ù)
\end{emcppslisting}
    
\noindent For
portability, a floating-point literal value of \(\pi\) that provides
sufficient precision for the longest \lstinline!long!~\lstinline!double! on
any relevant platform (e.g., 34 decimal digits for 128~bits of precision:
\lstinline!3.141'592'653'589'793'238'462'643'383'279'503!) should be
  used; see \featureref{\locationb}{digit-separators}.

Notice that we have elected to use \lstinline!constexpr! variables in place of  \lstinline!const! to guarantee that the floating-point \lstinline!pi! is a compile-time constant that will be usable as part of a constant expression.

\newpage%%%%%%%%
With the definition above, we can provide a
\lstinline!toRadians! function template that performs at maximum runtime
efficiency by avoiding needless type conversions during the computation:

\begin{emcppslisting}[emcppsbatch=e4]
template <typename T>
constexpr T toRadians(T degrees)
{
    return degrees * (pi<T> / T(180));
}
\end{emcppslisting}
    

\subsubsection[Reducing verbosity of type traits]{Reducing verbosity of type traits}\label{reducing-verbosity-of-type-traits}

A \emcppsgloss{type trait} is an empty type carrying compile-time information
about one or more aspects of another type. The way in which type traits
have been specified historically has been to define a class template
having the trait name and a public \lstinline!static! 
data member conventionally called \lstinline!value!, which is
initialized in the primary template to \lstinline!false!. Then, for each
type that wants to advertise that it has this trait, the header defining
the trait is included, and the trait is specialized for that type,
initializing \lstinline!value! to \lstinline!true!. We can achieve precisely
this same usage pattern replacing a trait \lstinline!struct! with a
variable template whose name represents the type trait and whose type of
variable itself is always \lstinline!bool!. Preferring variable templates
in this use case decreases the amount of \emcppsgloss{boilerplate code} ---
both at the point of definition and at the call
site.{\cprotect\footnote{As of C++17, the Standard Library provides a
more convenient way of inspecting the result of a type trait, by
introducing variable templates named the same way as the corresponding
traits but with an additional \lstinline!_v! suffix:

\begin{emcppshiddenlisting}[emcppsbatch=e5,emcppsstandards={c++17},style=footcode]
#include <type_traits>  // (ù{\codeincomments{std::is\_default\_constructible}}ù)
using T=int;
\end{emcppshiddenlisting}
\begin{emcppslisting}[emcppsbatch=e5,emcppsstandards={c++14},style=footcode]
// C++11/14
bool dc1 = std::is_default_constructible<T>::value;

// C++17
bool dc2 = std::is_default_constructible_v<T>;
\end{emcppslisting}
This delay is a consequence of the train release model of the Standard: Thoughtful application of the new feature throughout the vast Standard Library required significant effort that could not be completed before the next release date for the Standard and thus was delayed until C++17.
      }}

Consider, for example, a boolean trait designating whether a particular
type \lstinline!T! can be serialized to JSON:

\begin{emcppshiddenlisting}[emcppsbatch=e6]
// isSerializableToJson.h:
#pragma once
\end{emcppshiddenlisting}
\begin{emcppslisting}[emcppsbatch=e6]
// isSerializableToJson.h:

template <typename T>
constexpr bool isSerializableToJson = false;
\end{emcppslisting}
    
\noindent The header above contains the general variable template trait that, by
default, concludes that a given type is not serializable to JSON. Next
we consider the streaming utility itself:

\newpage%%%%%%
\begin{emcppshiddenlisting}[emcppsbatch=e6]
// serializeToJson.h:
struct JsonObject {};
\end{emcppshiddenlisting}
\begin{emcppslisting}[emcppsbatch=e6]
// serializeToJson.h:
#include <isSerializableToJson.h>  // general trait variable template

template <typename T>
JsonObject serializeToJson(const T& object)  // serialization function template
{
    static_assert(isSerializableToJson<T>,
                  "(ù{\codeincomments{T}}ù) must support serialization to JSON.");

    // ...

    return { /*..*/ };
}
\end{emcppslisting}
    
\noindent Notice that we have used the C++11
\lstinline!static_assert! feature to ensure that any type
used to instantiate this function will have specialized (see the next code snippet) the
general variable template associated with the specific type to be
\lstinline!true!.

Now imagine that we have a type, \lstinline!CompanyData!, that we would
like to advertise at compile time as being serializable to JSON. Like
other templates, variable templates can be specialized explicitly:

\begin{emcppslisting}[emcppsbatch=e6]
// companyData.h:
#include <isSerializableToJson.h>  // general trait variable template

struct CompanyData { /* ... */ };  // type to be JSON serialized

template <>
constexpr bool isSerializableToJson<CompanyData> = true;
    // Let anyone who needs to know that this type is JSON serializable.
\end{emcppslisting}
    
\noindent Finally, our \lstinline!client! function incorporates all of the above and
attempts to serialize both a \lstinline!CompanyData! object and an
\lstinline!std::map<int,!~\lstinline!char>!:

\enlargethispage*{-22ex}
\begin{emcppslisting}[emcppsbatch=e6]
// client.h:
#include <isSerializableToJson.h>  // general trait template
#include <companyData.h>           // JSON serializable type
#include <serializeToJson.h>       // serialization function
#include <map>                     // (ù{\codeincomments{std::map}}ù) (not JSON serializable)

void client()
{
    JsonObject jsonObj0 = serializeToJson(CompanyData());          // OK
    JsonObject jsonObj1 = serializeToJson(std::map<int, char>());  // Error
}
\end{emcppslisting}

     
\noindent In the \lstinline!client()! function above, \lstinline!CompanyData! works
fine, but because the variable template \lstinline!isSerializableToJson!
was never specialized to be \lstinline!true! for type
\mbox{\lstinline!std::map<int,! \lstinline!char>!}, the client header will --- as
desired --- fail to compile.

\subsection[Potential Pitfalls]{Potential Pitfalls}\label{variabletemplate-potential-pitfalls}

\subsubsection[Recursive variable template initializations require \lstinline!const! or \lstinline!constexpr!]{Recursive variable template initializations require {\SubsubsecCode const} or {\SubsubsecCode constexpr}}\label{recursive-variable-template-initializations-require-const-or-constexpr}

Instantiating variable templates that are defined recursively might have a subtle issue that could produce different results{\cprotect\footnote{For
example, GCC version 4.7.0 (2017) produces the expected results, whereas
  Clang version 10.x (2020) produces 1, 3, and 4, respectively.}} despite having no undefined behavior:

\begin{emcppslisting}
#include <iostream>  // (ù{\codeincomments{std::cout}}ù)

template <int N>
int fib = fib<N - 1> + fib<N - 2>;

template <> int fib<2> = 1;
template <> int fib<1> = 1;

int main()
{
    std::cout << fib<4> << '\n';  // 3 expected
    std::cout << fib<5> << '\n';  // 5 expected
    std::cout << fib<6> << '\n';  // 8 expected

    return 0;
}
\end{emcppslisting}
    
\noindent The root cause of this instability is
that the relative order of the initialization of the recursively generated variable template instantiations 
is not guaranteed because they are not defined explicitly \emph{within the same translation
unit}.  Therefore, a similar issue might have occurred in C++03 using \lstinline!static! members of a \lstinline!struct!: 
%The didactic value in answering this question dwarfs any potential
%practical value that recursive template variable instantiation can
%offer. First, consider that this same issue could, in theory, have
%occurred in C++03 using nested \texttt{static} members of a
%\texttt{struct}:

\newpage%%%%%%
\begin{emcppslisting}
#include <iostream>  // (ù{\codeincomments{std::cout}}ù)

template <int N> struct Fib
{
    static int value;                             // BAD IDEA: not (ù{\codeincomments{const}}ù)
};

template <> struct Fib<2> { static int value; };  // BAD IDEA: not (ù{\codeincomments{const}}ù)
template <> struct Fib<1> { static int value; };  // BAD IDEA: not (ù{\codeincomments{const}}ù)

template <int N> int Fib<N>::value = Fib<N - 1>::value + Fib<N - 2>::value;
int Fib<2>::value = 1;
int Fib<1>::value = 1;

int main()
{
    std::cout << Fib<4>::value << '\n';  // 3 expected
    std::cout << Fib<5>::value << '\n';  // 5 expected
    std::cout << Fib<6>::value << '\n';  // 8 expected

    return 0;
};
\end{emcppslisting}
    
\noindent However, this is not an issue when using \lstinline!enum!s due to enumerators always being compile-time constants:
%The problem did not manifest, however, because the simpler solution of
%using \texttt{enum}s (below) obviated separate initialization of the
%local \texttt{static} and didn't admit the possibility of failing to
%make the initializer a compile-time constant:

\begin{emcppslisting}
#include <iostream>  // (ù{\codeincomments{std::cout}}ù)

template <int N> struct Fib
{
    enum { value = Fib<N - 1>::value + Fib<N - 2>::value };  // OK, (ù{\codeincomments{const}}ù)
};

template <> struct Fib<2> { enum { value = 1 }; };           // OK, (ù{\codeincomments{const}}ù)
template <> struct Fib<1> { enum { value = 1 }; };           // OK, (ù{\codeincomments{const}}ù)

int main()
{
    std::cout << Fib<4>::value << '\n';  // 3 guaranteed
    std::cout << Fib<5>::value << '\n';  // 5 guaranteed
    std::cout << Fib<6>::value << '\n';  // 8 guaranteed

    return 0;
};
\end{emcppslisting}
    
\enlargethispage*{2ex}    
\noindent For integral variable templates, this issue can be resolved simply by adding a \lstinline!const! qualifier because the C++ Standard requires that any integral variable declared as \lstinline!const! and initialized with a compile-time constant is itself to be treated as a compile-time constant within the translation unit. 
%It was not until C++14 that the variable templates feature
%readily exposed this latent pitfall involving recursive initialization
%of non-\texttt{const} variables. The root cause of the instability is
%that the relative order of the initialization of the (recursively
%generated) variable instantiations is not guaranteed because they are
%not defined explicitly \emph{within the same translation unit}. The
%magic sauce that makes everything work is the C++ language requirement
%that any variable that is declared \texttt{const} and initialized with a
%compile-time constant is itself to be treated as a compile-time constant
%within the translation unit. This compile-time-constant propagation
%requirement imposes the needed ordering to ensure that the expected
%results are portable to all conforming compilers:

\begin{emcppslisting}
#include <iostream>  // (ù{\codeincomments{std::cout}}ù)

template <int N>
const int fib = fib<N - 1> + fib<N - 2>;  // OK, compile-time (ù{\codeincomments{const}}ù)

template <> const int fib<2> = 1;         // OK, compile-time (ù{\codeincomments{const}}ù)
template <> const int fib<1> = 1;         // OK, compile-time (ù{\codeincomments{const}}ù)

int main()
{
    std::cout << fib<4> << '\n';  // guaranteed to print out (ù{\codeincomments{3}}ù)
    std::cout << fib<5> << '\n';  // guaranteed to print out (ù{\codeincomments{5}}ù)
    std::cout << fib<6> << '\n';  // guaranteed to print out (ù{\codeincomments{8}}ù)

    return 0;
}
\end{emcppslisting}
    
\noindent Note that replacing each of the three \lstinline!const! keywords with
\lstinline!constexpr! in the example above also achieves the
desired goal, does not consume memory in the \emcppsgloss{static data
space}, and would also be applicable to nonintegral constants.

\subsection[Annoyances]{Annoyances}\label{annoyances}

\subsubsection[Variable templates do not support template template parameters]{Variable templates do not support template template parameters}\label{variable-templates-do-not-support-template-template-parameters}

Although a class or function template can accept a
\emcppsgloss{template template class parameter}, no equivalent
construct is available for variable templates{\cprotect\footnote{Mateusz Pusz has proposed for C++23 a way to increase consistency between
variable templates and class templates when used as \emcppsgloss[template template parameter]{template template parameters}; see \cite{pusz20}.}}:

\begin{emcppslisting}[emcppsbatch=e7]
template <typename T> T vt(5);

template <template <typename> class>
struct S { };

S<vt> s1;  // Error
\end{emcppslisting}
    
\noindent Providing a wrapper \lstinline!struct!
around a variable template might therefore be necessary in case the variable template needs to be passed to an interface
accepting a \emcppsgloss{template template parameter}:

\begin{emcppslisting}[emcppsbatch=e7]
template <typename T>
struct Vt { static constexpr T value = vt<T>; };

S<Vt> s2;  // OK
\end{emcppslisting}
    

\subsection[See Also]{See Also}\label{see-also}

\begin{itemize}
\item{%Section~\ref{constexprvar}, ``\titleref{constexprvar}" — 
\seealsoref{constexprvar}{\seealsolocationc}Conditionally safe C++11 feature providing an alternative to \lstinline!const! template variables that can reduce unnecessary consumption of the \emcppsgloss{static data space}.}
\end{itemize}

\subsection[Further Reading]{Further Reading}\label{further-reading}

\hspace*{\fill}







  % Chapter 1
\cleardoublepage

\chapter[Conditionally Safe Features]{Conditionally Safe Features}\label{ch-conditional}
\emcppschapterstart{}

PRODUCTION: Material copied from Ch 0; AUs will edit/replace on FPPs. 

AUTHORS: Remember to add material about whether these features are easy to teach. 

The preponderance of new features available in modern C++ has important, frequently occurring, and valuable uses, yet how these features are used appropriately, let alone optimally, might not be obvious. What’s more, some of these features are fraught with inherent dangers and deficiencies, requiring explicit training and extra care to circumnavigate their pitfalls.

For example, we deem default member initializers a \emph{conditionally safe} feature because, although they are easy to use per se, the perhaps less-than-obvious unintended consequences of doing so (e.g., tight compile-time coupling) might be prohibitively costly in certain circumstances (e.g., might prevent relink-only patching in production).

%%%%%%%%%%%%%%% C++11
\stepcounter{cppxx}
\addtocontents{toc}{\protect\renewcommand*\protect\cftsecindent{1.25em}}
\cftaddnumtitleline{toc}{section}{\thechapter.\thecppxx}{C++11}{\thepage}
% to add to the TOC at the section level
\renewcommand{\cppxx}{C++11}
%%%%%  

\newpage
\addtocontents{toc}{\protect\renewcommand*\protect\cftsecindent{4.5em}}
%\section[{\tt alignas}~]{The {\SecCode alignas} Decorator}\label{alignas}
% 18 March 2021, revisions in and proofed; commands updated
% 18 March 2021, sent to Slava for a second review and sent to Josh for code review
% 19 March 2021 JMB - code compiles take 2
% 27 March 2021, revisions from Slava in and proofed. 
% 28 March 2021, sent to CE 

\emcppsFeature{
    short={\lstinline!alignas!},
    tocshort={\TOCCode alignas},
    long={The {\SecCode alignas} Decorator},
    toclong={The \lstinline!alignas! Decorator},
    rhshort={\RHCode alignas},
}{alignas}
%\section[{\tt alignas}]{The {\SecCode alignas} Decorator\sectionmark{\RHCode alignas}}\label{alignas}\sectionmark{\RHCode alignas}
\setcounter{table}{0}
\setcounter{footnote}{0}
\setcounter{lstlisting}{0}

The \lstinline!alignas! storage specifier is used to strengthen
the \emcppsgloss{alignment} of a \lstinline!struct!, \lstinline!class!, \lstinline!union!, \lstinline!enum!, \emcppsgloss{variable}, or \emcppsgloss{data member}.

\subsection[Description]{Description}\label{description-alignas}

Each object type in C++ has an \emcppsgloss{alignment requirement} that restricts the addresses at which an object of that type is permitted to reside within the virtual-memory-address space. The \emcppsgloss{alignment requirement} is imposed by the object type on all objects of that type. The \lstinline!alignas! specifier provides a means of specifying a stricter alignment requirement than dictated by the type itself for a particular variable of the type or an individual data member of a \emcppsgloss[user defined type (UDT)]{user-defined type (UDT)}. The \lstinline!alignas! specifier can also be applied to a \emcppsgloss[user defined type (UDT)]{UDT} itself, but see \intraref{potential-pitfalls-alignas}{applying-alignas-to-a-type-might-be-misleading}.

\subsubsection[Supported alignments]{Supported alignments}\label{description-supported-alignments}

An alignment value is an integer of type \lstinline!std::size_t! that represents the number of bytes between the addresses at which a given object may be allocated.  All alignment values in C++ are always non-negative powers of two and are divided into two categories depending on whether they are larger than the alignment requirement of the \lstinline!std::max_align_t! type.  The \lstinline!std::max_align_t! type's alignment requirement is at least as strict as that of every \emcppsgloss{scalar type}.  An alignment value of less than or equal to the alignment requirement of \lstinline!std::max_align_t! is a \emcppsgloss{fundamental alignment}; otherwise, it is an \emcppsgloss{extended alignment}.  The \lstinline!std::max_align_t! type is typically an alias to the largest scalar type, which is \lstinline!long double! on most platforms, and its alignment requirement is usually 8 or 16.

\emcppsgloss[fundamental alignment]{Fundamental alignments} are required to be supported in \emph{all} contexts --- i.e., for variables with automatic, static, and dynamic storage durations as well as for nonstatic data members of a class and for function arguments.  While all fundamental and pointer types have \emcppsgloss[fundamental alignment]{fundamental alignments}, their specific values are \emcppsgloss[implementation defined]{implementation defined} and may differ between platforms.  For example, the alignment requirement of type \lstinline!long! might be \lstinline!4! on MSVC and \lstinline!8! on GCC.

%%%%%%%%% Revisit xref in FN during FPP creation
In contrast, whether \emph{any} \emcppsgloss{extended alignment} is supported at all and, if so, in which contexts, is \emcppsgloss{implementation defined}.\footnote{Implementations may warn when the alignment of a global object exceeds some maximal hardware threshold (such as the size of a physical memory page, e.g., 4096 or 8192).  For automatic variables (defined on the program stack), making alignment more restrictive than what would naturally be employed is seldom desired because at most one thread is able to access proximately located variables there unless explicitly passed in via address to separate threads; see \intraref{alignas-use-cases}{avoiding-false-sharing-among-distinct-objects-in-a-multi-threaded-program}. Note that, in the case of \lstinline!i0! in the \lstinline!alignas(32)! line in the last code snippet on page~\pageref{fnref}, a conforming platform that did not support an extended alignment of 32 would be required to report an error at compile time. [PRODUCTION: update the page xref as needed at FPPs. We may need to say "on this page" or otherwise clarify once pagination is set.]} For example, the strictest supported \emcppsgloss{extended alignment} for a variable with static storage duration might be as large as $2^{28}$ or $2^{29}$ or as small as $2^{13}$.

Since many aspects pertaining to the alignment requirements are \emcppsgloss{implementation defined}, we will use a specific platform to illustrate the behavior of \lstinline!alignas! throughout this feature's section.  Accordingly, the examples below show the behavior observed for the Clang compiler targetting desktop x86-64 Linux.

%The \texttt{alignas} specifier provides a means of further restricting
%the granularity at which (1) a particular object of arbitrary type, (2)
%a user-defined type (\texttt{class}, \texttt{struct}, \texttt{union}, or
%\texttt{enum}), or (3) an individual data member is permitted to reside
%within the virtual-memory-address space.

\subsubsection[Strengthening the alignment of a particular object]{Strengthening the alignment of a particular object}\label{restricting-the-alignment-of-a-particular-object}\label{strengthening-the-alignment-of-a-particular-object}

In its most basic form, the \lstinline!alignas! specifier strengthens the alignment requirement of a particular object. The desired alignment requirement is an \emcppsgloss{integral constant expression} provided as an argument to \lstinline!alignas!:

\begin{emcppslisting}
alignas(8) int i;     // OK, i is aligned on an 8 byte address boundary.
int j alignas(8), k;  // OK, j is 8 byte aligned; alignment of k is unchanged.
\end{emcppslisting}

\noindent If more than one alignment pertains to a given object, the strictest alignment value is applied:

\begin{emcppslisting}
alignas(4) alignas(8) alignas(2) char m;  // OK, (ù{\codeincomments{m}}ù) is 8-byte aligned.
alignas(8) int n alignas(16);             // OK, (ù{\codeincomments{n}}ù) is 16-byte aligned.
\end{emcppslisting}

\noindent For a program to be \emcppsgloss{well formed}, a specified alignment value
must satisfy several\linebreak[4] \mbox{requirements}:

\begin{enumerate}
\item{Be either zero or a non-negative integral power of two of type \lstinline!std::size_t! (0, 1, 2, 4, 8, 16\dots).}
%\item{Be at least the minimum alignment\cprotect\footnote{The minimum alignment of an entity is the least restrictive memory-address boundary at which the entity can be placed and have the program continue to work properly. This value is platform dependent and often subject to compiler controls but, by default, is often well approximated by \textbf{natural alignment}; see \textit{\titleref{alignas-appendix}: \titleref{natural-alignment}} on page \pageref{natural-alignment}.} required by the decorated entity.}
\item{Be larger or equal to what the alignment requirement would be without the \lstinline!alignas! specifier.}
\item{Be supported on the platform in the context in which the entity appears.}
%\cprotect\footnote{The notion of the largest supported alignment is characterized by both \textbf{maximal alignment} and the maximum \textbf{extended alignment}. \textbf{Maximal alignment} is defined as that most restrictive alignment that is valid in \emph{all} contexts on the current platform. All fundamental and pointer types necessarily have a minimal alignment that is less than or equal to \texttt{alignof(std::max\_align\_t)} — typically 8 or 16. Any alignment value greater than \textbf{maximal alignment} is an \textbf{extended alignment} value. Whether any extended alignment is supported (and in which contexts) is implementation defined. On typical platforms, extended alignment will often be as large as $2^{18}$ or $2^{19}$, however implementations may warn when the alignment of a global object exceeds some maximal hardware threshold (such as the size of a physical memory page, e.g., 4096 or 8192). For \textbf{automatic variables} (defined on the program stack), making alignment more restrictive than what would naturally be employed  is seldom desired because at most one thread is able to access proximately located variables there unless explicitly passed in via address to separate threads; see \textit{\titleref{alignas-use-cases}: \titleref{avoiding-false-sharing-among-distinct-objects-in-a-multi-threaded-program}} on page~\pageref{avoiding-false-sharing-among-distinct-objects-in-a-multi-threaded-program}. Note that, in the case of \texttt{i} in the first code snippet on page~\pageref{restricting-the-alignment-of-a-particular-object}, a conforming platform that did not support an extended alignment of 64 would be required to report an error at compile time.}
\end{enumerate}

\noindent Additionally, if the specified alignment value is zero, the
\lstinline!alignas! specifier is ignored:

%\begin{lstlisting}[language=C++]
%// Static variables declared at namespace scope
%alignas(32) int i0;  // OK, aligned on a 32-byte boundary (extended alignment)
%alignas(16) int i1;  // OK, aligned on a 16-byte boundary (maximum alignment)
%alignas(8)  int i2;  // OK, aligned on an 8-byte boundary
%alignas(7)  int i3;  // Error, not a power of two
%alignas(4)  int i4;  // OK, no change to alignment boundary
%alignas(2)  int i5;  // Error, less than minimum alignment on this platform
%alignas(0)  int i6;  // OK, (ù{\codeincomments{alignas}}ù) specifier ignored
%
%alignas(1024 * 16) int i7;
%    // OK, might warn: e.g., exceeds (physical) page size on current platform
%
%alignas(1024 * 1024 * 512) int i8;
%    // (likely) compile-time error: e.g., exceeds maximum size of object file
%
%alignas(8) char buf[128];  // create 8-byte-aligned, 128-byte character buffer
%
%void f()
%{
%    // automatic variables declared at function scope
%    alignas(4)  double e0;  // Error, less than minimum alignment on this platform
%    alignas(8)  double e1;  // OK, no-change to (8-byte) alignment boundary
%    alignas(16) double e2;  // OK, aligned to maximum (fundamental) alignment value
%    alignas(32) double e3;  // OK, maximum alignment value exceeded; might warn
%}
%\end{lstlisting}
\begin{emcppslisting}[label=fnref]
// Static variables declared at namespace scope
alignas(32) int i0; // OK, 32-byte aligned (extended alignment)
alignas(16) int i1; // OK, 16-byte aligned (strictest fundamental alignment)
alignas(8)  int i2; // OK,  8-byte aligned (fundamental alignment)
alignas(7)  int i3; // Error, not a power of two
alignas(4)  int i4; // OK, no change to alignment requirement
alignas(2)  int i5; // Error, less than alignment would be without (ù{\codeincomments{alignas}}ù)
alignas(0)  int i6; // OK, (ù{\codeincomments{alignas}}ù) specifier ignored

alignas(1024 * 16) int i7;
    // OK, might warn on other platforms, e.g., exceeds physical page size

alignas(1 << 30) int i8;
    // Error, exceeds maximum supported extended alignment

alignas(8) char buf[128]; // OK, 8-byte-aligned, 128-byte character buffer

void f()
{
  // automatic variables declared at function scope
  alignas(4)  double e0; // Error, less than alignment would be without (ù{\codeincomments{alignas}}ù)
  alignas(8)  double e1; // OK, no change to (8-byte) alignment requirement
  alignas(16) double e2; // OK, 16-byte aligned (fundamental alignment)
  alignas(32) double e3; // OK, 32-byte aligned (extended alignment)
}
\end{emcppslisting}



\subsubsection[Strengthening the alignment of individual data members]{Strengthening the alignment of individual data members}\label{restricting-the-alignment-of-individual-data-members}\label{strengthening-the-alignment-of-individual-data-members}

Within a user-defined type (\lstinline!class!, \lstinline!struct!, or
\lstinline!union!), using the
\lstinline!alignas! keyword to specify the alignments of individual data
members is possible:

\begin{emcppslisting}
struct T2
{
    alignas(8)  char   x;  // size 1; alignment 8
    alignas(16) int    y;  // size 4; alignment 16
    alignas(64) double z;  // size 8; alignment 64
};  // size 128; alignment 64
\end{emcppslisting}

\noindent The effect here is the same as if we had added the padding explicitly
and then set the alignment of the structure overall:

\begin{emcppslisting}
struct alignas(64) T3
{
    char   x;      // size 1; alignment 8
    char   a[15];  // padding
    int    y;      // size 4; alignment 16
    char   b[44];  // padding
    double z;      // size 8; alignment 64
    char   c[56];  // padding (optional)
};  // size 128; alignment 64
\end{emcppslisting}

\noindent Again, if more than one alignment specifier pertains to a given data member,
the strictest applicable alignment value is applied:

\begin{emcppslisting}
struct T4
{
    alignas(2) char
        c1 alignas(1),  // size 1; alignment 2
        c2 alignas(2),  // size 1; alignment 2
        c4 alignas(4);  // size 1; alignment 4
};                      // size 8; alignment 4
\end{emcppslisting}

\subsubsection[Strengthening the alignment of a user-defined type]{Strengthening the alignment of a user-defined type}\label{restricting-the-alignment-of-a-user-defined-type}
\label{strengthening-the-alignment-of-a-user-defined-type}

The \lstinline!alignas! specifier can also be used to specify alignment for
user-defined types (UDTs), such as a \lstinline!class!, \lstinline!struct!,
\lstinline!union!, or \lstinline!enum!. When specifying the alignment of a UDT,
 the \lstinline!alignas! keyword is placed \emph{after} the
type specifier (e.g., \lstinline!class!) and just before the name of the
type (e.g., \lstinline!C!):

%\begin{lstlisting}[language=C++]
%class alignas(2)  C { };  // OK, aligned on a 2-byte boundary; size = 2
%struct alignas(4) S { };  // OK, aligned on a 4-byte boundary; size = 4
%union alignas(8)  U { };  // OK, aligned on an 8-byte boundary; size = 8
%enum alignas(16)  E { };  // OK, aligned on a 16-byte boundary; size = 4
%\end{lstlisting}
\begin{emcppslisting}
class  alignas( 2) C { };  // OK, aligned on a  2-byte boundary; size = 2
struct alignas( 4) S { };  // OK, aligned on a  4-byte boundary; size = 4
union  alignas( 8) U { };  // OK, aligned on an 8-byte boundary; size = 8
enum   alignas(16) E { };  // OK, aligned on a 16-byte boundary; size = 4
\end{emcppslisting}

\noindent Notice that, for each of \lstinline!class!, \lstinline!struct!, and
\lstinline!union! above, the \lstinline!sizeof! objects of that type increased
to match the alignment; in the case of the \lstinline!enum!, however, the
size remains that of the default \emcppsgloss[underlying type (UT)]{underlying type} (e.g., 4
bytes) on the current platform.{\cprotect\footnote{When \lstinline!alignas!
is applied to an enumeration \lstinline!E!, the Standard does not
indicate whether padding bytes are added to \lstinline!E!'s object
representation or not, affecting the result of \lstinline!sizeof(E)!. The
implementation variance resulting from this lack of clarity in the
  Standard was captured in \cite{miller17}. The outcome of the core
  issue was to completely remove permission for \lstinline!alignas! to be
  applied to enumerations (see \cite{iso18a}). Therefore, conforming implementations will
  eventually stop accepting the \lstinline!alignas! specifier on
  enumerations in the future.}}

Again, specifying an alignment that is less than what would be without the \lstinline!alignas! specifier is ill formed:

\begin{emcppslisting}
struct alignas(2) T0 { int i; };
    // Error, Alignment of (ù{\codeincomments{T0}}ù) (2) is less than that of (ù{\codeincomments{int}}ù) (4).
struct alignas(1) T1 { C c; };
    // Error, Alignment of (ù{\codeincomments{T1}}ù) (1) is less than that of (ù{\codeincomments{C}}ù) (2).
\end{emcppslisting}



\subsubsection[Matching the alignment of another type]{Matching the alignment of another type}\label{matching-the-alignment-of-another-type}

The \lstinline!alignas! specifier also accepts (as an argument) a type
identifier. In its alternate form, \lstinline!alignas(T)! is strictly
equivalent to \lstinline!alignas(alignof(T))! (see \featureref{\locationc}{alignof}):

\begin{emcppslisting}
alignas(int) char c;  // equivalent to (ù{\codeincomments{alignas(alignof(int)) char c;}}ù)
\end{emcppslisting}


\subsection[Use Cases]{Use Cases}\label{alignas-use-cases}

\subsubsection[Creating a sufficiently aligned object buffer]{Creating a sufficiently aligned object buffer}\label{creating-a-sufficiently-aligned-object-buffer}

%When writing low-level, system-infrastructure code, constructing an object within a raw buffer is sometimes useful. As a minimal
%example, consider a function that uses a local character buffer to
%create an object of type \lstinline!std::complex<long!~\lstinline!double>! on
%the program stack using placement \lstinline!new!:
%
%\begin{emcppslisting}
%#include <complex>  // (ù{\codeincomments{std::complex}}ù)
%
%void f()
%{
%    // ...
%    char objectBuffer[sizeof(std::complex<long double>)];    // BAD IDEA
%    // ...
%    new (objectBuffer) std::complex<long double>(1.0, 0.0);  // Might dump core!
%    // ...
%}
%\end{emcppslisting}
%
%\noindent The essential problem with the code above is that \lstinline!objectBuffer!,
%being an array of characters (each having an alignment of 1), is itself
%byte aligned. The compiler is therefore free to place it on any address
%boundary. On the other hand, \lstinline!std::complex<long!~\lstinline!double>! is an aggregate consisting of two \lstinline!long!~\lstinline!double! objects
%and therefore necessarily requires (at least) the same strict alignment
%(typically 16) as the two \lstinline!long!~\lstinline!double! objects it comprises. Previous
%solutions to this problem involved creating a \lstinline!union! of the
%object buffer and some maximally aligned type (e.g.,
%\lstinline!std::max_align_t!):
%
%\begin{emcppslisting}
%#include <complex>  // (ù{\codeincomments{std::complex}}ù)
%#include <cstddef>  // (ù{\codeincomments{std::max\_align\_t}}ù)
%
%void f()
%{
%    // ...
%
%    union {  // awkward workaround
%        std::max_align_t dummy;  // (ù{\codeincomments{typedef}}ù) to maximally aligned type
%        char objectBuffer[sizeof(std::complex<long double>)];
%    } objectBuffer;
%
%    // ...
%
%    new (&objectBuffer) std::complex<long double>(1.0, 0.0);  // OK
%
%    // ...
%}
%\end{emcppslisting}
%
%\noindent Using the alternate syntax for \lstinline!alignas!, we can avoid gratuitous
%complexity and just state our intentions explicitly:
%
%\begin{emcppslisting}
%#include <complex>  // (ù{\codeincomments{std::complex}}ù)
%void f()
%{
%    // ...
%
%    alignas(std::complex<long double>) char objectBuffer[
%                              sizeof(std::complex<long double>)];  // GOOD IDEA
%
%    // ...
%
%    new (objectBuffer) std::complex<long double>(1.0, 0.0);  // OK
%
%    // ...
%}
%\end{emcppslisting}

One of motivating use cases for introducing the \lstinline!alignas! feature was creating static capacity, dynamic size containers that do not use dynamic allocations:

\begin{emcppslisting}
#include <cassert>  // standard C (ù{\codeincomments{assert}}ù) macro
#include <new>      // placement new

template <typename TYPE, std::size_t CAPACITY>
class FixedVector {
\end{emcppslisting}

To avoid the overhead of initializing the unused elements, such a generic container needs to have a character buffer data member and employ placement \lstinline!new! to construct the elements as needed. This buffer needs to be of sufficient size to store the elements, which one can easily compute using \lstinline!CAPACITY!~\lstinline!*!~\lstinline!sizeof(T)!.  In addition, however, ensuring that the buffer is sufficiently aligned to store elements of the supplied \lstinline!TYPE! is important.  With \lstinline!alignas!, ensuring this requirement is straightforward:

\begin{emcppslisting}
    alignas(TYPE) char d_buffer[CAPACITY * sizeof(TYPE)];
        // raw memory buffer of proper size and alignment for (ù{\codeincomments{TYPE}}ù) elements

    std::size_t        d_size; 
        // current size of the vector

    TYPE *rawElementPtr(std::size_t index)
        // Return the pointer to the element with the specified (ù{\codeincomments{index}}ù).
    { 
        return reinterpret_cast<TYPE *>(d_buffer) + index;
    }

public:
    // ...

    void resize(std::size_t size)
    {
        assert(size <= CAPACITY);
        while (d_size < size) new (rawBufferPtr(d_size++)) TYPE;
        while (d_size > size) rawElementPtr(--d_size)->~TYPE();
    }

    // ...
};
\end{emcppslisting}

Without the use of \lstinline!alignas!, \lstinline!d_buffer! (in the code snippet above), which is an array of characters, would itself have an alignment requirement of 1.  The compiler would therefore be free to place it on any address boundary, which is problematic for any \lstinline!TYPE! argument with an alignment requirement stricter than 1.  

\subsubsection[Ensuring proper alignment for architecture-specific instructions]{Ensuring proper alignment for architecture-specific instructions}\label{ensuring-proper-alignment-for-architecture-specific-instructions}

Architecture-specific instructions or compiler intrinsics might require
the data they act on to have a specific alignment. One example of such
intrinsics is the \emph{Streaming SIMD Extensions (SSE)}\footnote{\cite{inteliig}, ``Technologies: SSE"} instruction set available on the x86
architecture. SSE instructions operate on groups of four 32-bit
single-precision floating-point numbers at a time, which are required to
be 16-byte aligned.{\cprotect\footnote{``Data must be
16-byte aligned when loading to and storing from the 128-bit XMM
  registers used by SSE/SSE2/SSE3/SSSE3'': see \cite{{intel16}},
  section 4.4.4, ``Data Alignment for 128-Bit Data," pp. 4-19--4-20.}} The
\lstinline!alignas! specifier can be used to create a type satisfying this requirement:

\begin{emcppslisting}[emcppsbatch=e1]
struct SSEVector
{
    alignas(16) float d_data[4];
};
\end{emcppslisting}

\noindent Each object of the \lstinline!SSEVector! type above is guaranteed always to
be aligned to a 16-byte boundary and can therefore be safely (and
conveniently) used with SSE intrinsics:

\begin{emcppslisting}[emcppsbatch=e1]
#include <cassert>     // standard C (ù{\codeincomments{assert}}ù) macro
#include <xmmintrin.h> // (ù{\codeincomments{\_\_m128}}ù) and (ù{\codeincomments{\_mm\_XXX}}ù) functions

void f()
{
    const SSEVector v0 = {0.0f, 1.0f, 2.0f, 3.0f};
    const SSEVector v1 = {10.0f, 10.0f, 10.0f, 10.0f};

    __m128 sseV0 = _mm_load_ps(v0.d_data);
    __m128 sseV1 = _mm_load_ps(v1.d_data);
        // (ù{\codeincomments{\_mm\_load\_ps}}ù) requires the given (ù{\codeincomments{float}}ù) array to be 16-byte aligned.
        // The data is loaded into a dedicated 128-bit CPU register.

    __m128 sseResult = _mm_add_ps(sseV0, sseV1);
        // sum two 128-bit registers; typically generates an (ù{\codeincomments{addps}}ù) instruction

    SSEVector vResult;
    _mm_store_ps(vResult.d_data, sseResult);
        // Store the result of the sum back into a (ù{\codeincomments{float}}ù) array.

    assert(vResult.d_data[0] == 10.0f);
    assert(vResult.d_data[1] == 11.0f);
    assert(vResult.d_data[2] == 12.0f);
    assert(vResult.d_data[3] == 13.0f);
}
\end{emcppslisting}


\subsubsection[Avoiding false sharing among distinct objects in a multi-threaded program]{Avoiding false sharing among distinct objects in a multi-threaded program}\label{avoiding-false-sharing-among-distinct-objects-in-a-multi-threaded-program}

In the context of an application where multithreading has been employed
to improve performance, seeing
a previously single-threaded workflow become even less performant after
a parallelization attempt can be surprising (and disheartening). One possible insidious cause of such
disappointing results comes from \emcppsgloss{false sharing} --- a situation
in which multiple threads unwittingly harm each other's performance
while writing to logically independent variables that happen to reside
on the same \emcppsgloss{cache line}; see \intraref{alignas-appendix}{cache-lines,-l1,-l2,-and-l3-cache,-pages,-and-virtual-memory}.

As a simple illustration of the potential
performance degradation resulting from \emcppsgloss{false sharing}, consider
a function that spawns separate threads to repeatedly increment
(concurrently) logically distinct variables that happen to reside in
close proximity on the program stack:

\begin{emcppslisting}
#include <thread>  // (ù{\codeincomments{std::thread}}ù)

void incrementJob(int* p);
    // Repeatedly increment (ù{\codeincomments{*p}}ù) a large, fixed number of times.

void f()
{
    int i0 = 0;  // Here, (ù{\codeincomments{i0}}ù) and (ù{\codeincomments{i1}}ù) likely share the same cache line,
    int i1 = 0;  // i.e., byte-aligned memory block on the program stack.

    std::thread t0(&incrementJob, &i0);
    std::thread t1(&incrementJob, &i1);
        // Spawn two parallel jobs incrementing the respective variables.

    t0.join();
    t1.join();
        // Wait for both jobs to be completed.
}
\end{emcppslisting}

\noindent In the simplistic example above, the proximity in memory between
\lstinline!i0! and \lstinline!i1! can result in their belonging to the same
\emcppsgloss{cache line}, thus leading to \emcppsgloss{false sharing}. By
using \lstinline!alignas! to strengthen the alignment requirement of both integers to the cache line size, we
ensure that the two variables reside on distinct cache lines:

%\begin{lstlisting}[language=C++]
%// ...
%
%void f()
%{
%    alignas(64) int i0 = 0;   // Assuming a cache line on this platform is 64
%    alignas(64) int i1 = 0;   // bytes, (ù{\codeincomments{i0}}ù) and (ù{\codeincomments{i1}}ù) will be on separate ones.
%
%    // ...
%\end{lstlisting}
\begin{emcppslisting}[emcppsbatch=e2]
// ...

enum { k_CACHE_LINE_SIZE = 64 };  // A cache line on this platform is 64 bytes.

void f()
{
    alignas(k_CACHE_LINE_SIZE) int i0 = 0; // (ù{\codeincomments{i1}}ù) and (ù{\codeincomments{i2}}ù) are on separate
    alignas(k_CACHE_LINE_SIZE) int i1 = 0; // cache lines.

    // ...
}
\end{emcppslisting}

\noindent As an empirical demonstration of the effects of \emcppsgloss{false sharing},
a benchmark program repeatedly calling \lstinline!f! completed its
execution seven times faster on average when compared to the same
program without use of \lstinline!alignas!.{\cprotect\footnote{The benchmark
program was compiled using Clang 11.0.0 using \lstinline!-Ofast!,
\lstinline!-march=native!, and \mbox{\lstinline!-std=c++11!}. The program was then
executed on a machine running Windows 10 x64, equipped with an Intel
Core i7-9700k CPU (8 cores, 64-byte cache line size). Over the
course of multiple runs, the version of the benchmark without
\lstinline!alignas! took 18.5967ms to complete (on average), while the
version with \lstinline!alignas! took 2.45333ms to complete (on average).
  See \textbf{{[PRODUCTION: CODE PROVIDED WITH BOOK] alignasbenchmark}} for the source code of the program.}} Note that because supported extended alignments are implementation defined, using \lstinline!alignas! is not a strictly portable solution. Opting for less elegant and more wasteful padding approach instead of \lstinline!alignas! might be preferrable for portability.

\subsubsection[Avoiding false sharing within a single thread-aware object]{Avoiding false sharing within a single thread-aware object}\label{avoiding-false-sharing-within-a-single-thread-aware-object}

A real-world scenario where the need for preventing \emcppsgloss{false
sharing} is fundamental occurs in the implementation of high-performance
concurrent data structures. As an example, a thread-safe ring buffer
might make use of \lstinline!alignas! to ensure that the indices of the
head and tail of the buffer are aligned at the start of a cache line
(typically 64, 128, or 256 bytes),\footnote{In C++17, one can portably retrieve the minimum offset between two objects to avoid false sharing through the \lstinline!std::hardware_destructive_interference_size! constant defined in the \lstinline!<new>! header.}  thereby preventing them from
occupying the same one.

\begin{emcppslisting}[emcppsbatch=e2]
#include <atomic>  // (ù{\codeincomments{std::atomic}}ù)
class ThreadSafeRingBuffer
{
    alignas(k_CACHE_LINE_SIZE) std::atomic<std::size_t> d_head;
    alignas(k_CACHE_LINE_SIZE) std::atomic<std::size_t> d_tail;

    // ...
};
\end{emcppslisting}

\noindent Not aligning \lstinline!d_head! and \lstinline!d_tail! (above) to the CPU
cache size might result in poor performance of the
\lstinline!ThreadSafeRingBuffer! because CPU cores that need to access only one
of the variables will inadvertently load the other one as well,
triggering expensive hardware-level coherency mechanisms between the
cores' caches. On the other hand, specifying such substantially stricter
alignment on consecutive data members necessarily increases the size of
the object; see \intraref{potential-pitfalls-alignas}{stricter-alignment-might-reduce-cache-utilization}.
%{\it\titleref{potential-pitfalls}:} {\it\titleref{stricter-alignment-might-reduce-cache-utilization}} on page~\pageref{stricter-alignment-might-reduce-cache-utilization}.

\subsection[Potential Pitfalls]{Potential Pitfalls}\label{potential-pitfalls-alignas}

\subsubsection[Underspecifying alignment is not universally reported]{Underspecifying alignment is not universally reported}\label{underspecifying-alignment-is-not-universally-reported}

The Standard is clear when it comes to underspecifying
alignment{\cprotect\footnote{\cite{{cpp11}}, section 7.6.2, ``Alignment Specifier," paragraph~5, pp. 179}}:
\begin{quote}
The combined effect of all \emph{alignment-specifiers} in a declaration
shall not specify an alignment that is less strict than the alignment
that would be required for the entity being declared if all
\emph{alignment-specifiers} were omitted (including those in other
declarations).
\end{quote}
The compiler is required to honor the
specified value if it is a \emcppsgloss{fundamental
alignment},{\cprotect\footnote{``If the constant expression evaluates to
a fundamental alignment, the alignment requirement of the declared
entity shall be the specified fundamental alignment'': \cite{cpp11}, section~7.6.2, ``Alignment Specifier," paragraph~2, item~2, p. 178.}} so
imagining how this would lead to anything other than an ill-formed
program is difficult:

\begin{emcppslisting}
alignas(4) void* p;              // Error, (ù{\codeincomments{alignas(4)}}ù) is below minimum, 8.

struct alignas(2) S { int x; };  // Error, (ù{\codeincomments{alignas(2)}}ù) is below minimum, 4.

struct alignas(2) T { };
struct alignas(1) U { T e; };    // Error, (ù{\codeincomments{alignas(1)}}ù) is below minimum, 2.
\end{emcppslisting}

\noindent Each of the three errors above are reported by Clang. MSVC and ICC issue a warning, whereas
 GCC
provides no diagnostic at all, even in the most pedantic warning mode. Thus, one could write a program,
involving statements like those above, that happens to work on one
platform (e.g., GCC) but fails to compile on another (e.g.,
Clang).{\cprotect\footnote{Underspecifying alignment is not reported at
all by GCC 10.2, using the
\lstinline!-std=c++11! \lstinline!-Wall! \lstinline!-Wextra! \mbox{\lstinline!-Wpedantic!}
flags. This behavior is reported as a compiler defect; see \cite{wakely15}. With the same set of options, Clang 10.1 produces a compilation
failure. ICC 2021.1.2 and MSVC v19.28 will produce a warning and ignore any alignment
  less than the minimum one.}}

\subsubsection[Incompatibly specifying alignment is IFNDR]{Incompatibly specifying alignment is IFNDR}\label{incompatibly-specifying-alignment-is-ifndr}

It is permissible to forward declare a user-defined type (UDT)
without an \lstinline!alignas! specifier so long as all defining
declarations of the type have either no \lstinline!alignas! specifier or
have the same one. Similarly, if any forward declaration of a
user-defined type has an \lstinline!alignas! specifier, then all defining
declarations of the type must have the same specifier and that specifier
must be \emph{equivalent to} (not necessarily \emph{the same as}) that
in the forward declaration:

\begin{emcppslisting}
struct Foo;                  // OK, does not specify an alignment
struct alignas(double) Foo;  // OK, must be equivalent to every definition
struct alignas(8) Foo;       // OK, all definitions must be identical.

struct alignas(8) Foo { };   // OK, def. equiv. to each decl. specifying alignment

struct Foo;                  // OK, has no effect
struct alignas(8) Foo;       // OK, has no effect; might warn after definition
\end{emcppslisting}

\noindent Specifying an alignment in a forward declaration without specifying an
equivalent one in the defining declaration is \emcppsgloss{ill formed, no diagnostic
required (IFNDR)} if the two declarations appear in distinct translation
units:

\begin{emcppslisting}
struct alignas(4) Bar;      // OK, forward declaration
struct Bar { };             // Error, missing (ù{\codeincomments{alignas}}ù) specifier

struct alignas(4) Baz;      // OK, forward declaration
struct alignas(8) Baz { };  // Error, nonequivalent (ù{\codeincomments{alignas}}ù) specifier
\end{emcppslisting}

\noindent Both of the errors above are flagged by Clang. MSVC and ICC warn on the first one and produce an error on the second one, whereas neither of them is
reported by GCC. Note that when the inconsistency \emph{occurs across
translation units}, no mainstream compiler is likely to diagnose the
problem:

\begin{emcppslisting}
// file1.cpp:
struct Bam { char ch; } bam, *p = &bam;

// file2.cpp:
struct alignas(int) Bam;  // Error, definition of (ù{\codeincomments{Bam}}ù) lacks alignment specifier.
extern Bam* p;            //        (no diagnostic required)
\end{emcppslisting}

\noindent Any program incorporating both translation units above is
\emcppsgloss[ill formed, no diagnostic required (IFNDR)]{IFNDR}.

\subsubsection[Applying \lstinline!alignas! to a \emph{type} might be misleading]{Applying {\SubsubsecCode alignas} to a {\sfbsubsubsecitalRomeo type} might be misleading}\label{applying-alignas-to-a-type-might-be-misleading}

When applying the \lstinline!alignas! specifier to a user-defined type having no base classes, one might be convinced that it is equivalent to applying \lstinline!alignas! to its first declared data member:

\begin{emcppslisting}
struct S0 {
    alignas(16) char d_buffer[128];  // guaranteed to be 16-byte aligned
                int  d_index;
};

struct alignas(16) S1 {
    char d_buffer[128];              // guaranteed to be 16-byte aligned
    int  d_index;
};
\end{emcppslisting}

Indeed, for all objects of the \lstinline!S0! and \lstinline!S1! types (in the example above), their respective \lstinline!d_buffer! data members will be aligned on a 16-byte boundary. Such equivalency, however, holds only for \emcppsgloss{standard layout types}.  Adding a virtual function or even simply changing the access control for some of the data members\footnote{According to the C++20 Standard, compilers are allowed to reorder data members having different access control.  However, no compilers take advantage of this ability in practice, and C++23 might mandate that the data members are always laid out in declaration order; see \cite{balog20}.} might break this equivalency:

\begin{emcppslisting}
struct S2 {
    alignas(16) char d_buffer[128];  // guaranteed to be 16-byte aligned
                int  d_index;

    virtual ~S2();
};

struct alignas(16) S3 {
    char d_buffer[128];              // (ù{\emphincomments{not}}ù) guaranteed to be 16-byte aligned
    int  d_index;

    virtual ~S3();
};

struct S4 {
    alignas(16) char d_buffer[128];  // guaranteed to be 16-byte aligned
private:
                int  d_index;
};

struct alignas(16) S5 {
    char d_buffer[128];              // (ù{\emphincomments{not}}ù) guaranteed to be 16-byte aligned
private:
    int  d_index;
};
\end{emcppslisting}

Any code that relies on the \lstinline!d_buffer! member of instances of the \lstinline!S3! and \lstinline!S5! types (in the code example above) being 16-byte aligned is defective.

\subsubsection[Overlooking alternative approaches to avoid false sharing]{Overlooking alternative approaches to avoid false sharing}\label{stricter-alignment-might-reduce-cache-utilization}\label{overlooking-alternative-approaches-to-avoid-false-sharing}

User-defined types having artificially stricter alignments than would
naturally occur on the host platform means that fewer of them can fit
within any given level of physical cache within the hardware. Types
having data members whose alignment is artificially strengthened tend to be
larger and thus suffer the same lost cache utilization. As an
alternative to enforcing stricter alignment to avoid \emcppsgloss{false
sharing}, consider organizing a multithreaded program such that tight
clusters of repeatedly accessed objects are always acted upon by only a
single thread at a time, e.g., using local (arena) memory allocators;
see \intraref{alignas-appendix}{cache-lines,-l1,-l2,-and-l3-cache,-pages,-and-virtual-memory}.
%{\it\titleref{alignas-appendix}:} {\it\titleref{cache-lines,-l1,-l2,-and-l3-cache,-pages,-and-virtual-memory}} on page~\pageref{cache-lines,-l1,-l2,-and-l3-cache,-pages,-and-virtual-memory}.

\subsection[See Also]{See Also}\label{see-also}

\begin{itemize}
\item{%Section~\ref{alignof}, ``\titleref{alignof}" — Safe C++11 feature that
\seealsoref{alignof}{\seealsolocationc}inspects the alignment of a given type.}
%\item{Section~\ref{attributes}, ``\titleref{attributes}" — Safe C++11 feature that shows how other attributes (following the conventional attribute notation) are used to annotate source code, improve error diagnostics, and implicitly code generation} %%% WORDS MISSING!
\end{itemize}

\subsection[Further Reading]{Further Reading}\label{further-reading}

None so far

\subsection[Appendix]{Appendix}\label{alignas-appendix}

\subsubsection[Natural Alignment]{Natural Alignment}\label{natural-alignment}

%By default, fundamental, pointer, and enumerated types typically reside
%on an address boundary that divides the size of the object; we refer to
%such alignment as \textbf{natural alignment}{\cprotect\footnote{Sizes
%and alignment shown here are typical but not specifically required by
%the standard. On some platforms, one can request that all
%types be \textbf{byte aligned}. While such a representation is more
%compact, entities that span memory boundaries can require multiple
%fetch operations leading to run times that are typically
%significantly (sometimes as much as an order of magnitude) slower when
%  run in this ``packed'' mode.}}:
%
%\begin{lstlisting}[language=C++]
%char   c;  // size 1; alignment 1; boundaries: 0x00, 0x01, 0x02, 0x03, ...
%short  s;  // size 2; alignment 2: boundaries: 0x00, 0x02, 0x04, 0x06, ...
%int    i;  // size 4; alignment 4; boundaries: 0x00, 0x04, 0x08, 0x0c, ...
%float  f;  // size 4; alignment 4; boundaries: 0x00, 0x04, 0x08, 0x0c, ...
%double d;  // size 8; alignment 8; boundaries: 0x00, 0x08, 0x10, 0x18, ...
%\end{lstlisting}

Many micro-architectures are optimized for working with data that has \emcppsgloss{natural alignment}, i.e., objects reside on an address boundary that divides their size rounded up to the nearest power of two.  With the additional restriction that no padding is allowed between C++ array elements, the alignment requirements of fundamental types are often equal to their respective size on most platforms:

\begin{emcppslisting}
char        c;  // size 1;  alignment 1;  boundaries: 0x00, 0x01, 0x02, ...
short       s;  // size 2;  alignment 2:  boundaries: 0x00, 0x02, 0x04, ...
int         i;  // size 4;  alignment 4;  boundaries: 0x00, 0x04, 0x08, ...
float       f;  // size 4;  alignment 4;  boundaries: 0x00, 0x04, 0x08, ...
double      d;  // size 8;  alignment 8;  boundaries: 0x00, 0x08, 0x10, ...
long double l;  // size 16; alignment 16; boundaries: 0x00, 0x10, 0x20, ...
\end{emcppslisting}

\noindent The alignment requirement of an array of objects is the same as that of its elements:

\begin{emcppslisting}
char  arrC[4];  // size 4; alignment 1
short arrS[4];  // size 8; alignment 2
\end{emcppslisting}

For user-defined types, compilers compute the alignment and add appropriate padding between the data members and after the last one such that all alignment requirements of the data members are satisfied and no padding would be required should an array of the type be created. Typically, this results in the alignment requirement of a \emcppsgloss[user defined type (UDT)]{UDT} being the same as that of the most strictly aligned nonstatic data member:

%For aggregates (including arrays) or user-defined types, the alignment
%is typically that of the most strictly aligned subelement:

\begin{emcppslisting}[emcppsbatch=e3]
struct S0
{
    char a;     // size 1; alignment 1
    char b;     // size 1; alignment 1
    int  c;     // size 4; alignment 4
};              // size 8; alignment 4

struct S1
{
    char a;     // size  1; alignment 1
    int  b;     // size  4; alignment 4
    char c;     // size  1; alignment 1
};              // size 12; alignment 4

struct S2
{
    int  a;     // size 4; alignment 4
    char b;     // size 1; alignment 1
    char c;     // size 1; alignment 1
};              // size 8; alignment 4

struct S3
{
    char a;     // size 1; alignment 1
    char b;     // size 1; alignment 1
};              // size 2; alignment 1

struct S4
{
    char a[2];  // size 2; alignment 1
};              // size 2; alignment 1
\end{emcppslisting}

\noindent Size and alignment behave similarly with respect to \emcppsgloss{structural
inheritance}:

\begin{emcppslisting}[emcppsbatch=e3]
struct D0 : S0
{
    double d;  // size  8; alignment 8
};             // size 16; alignment 8

struct D1 : S1
{
    double d;  // size  8; alignment 8
};             // size 24; alignment 8

struct D2 : S2
{
    int d;     // size  4; alignment 4
};             // size 12; alignment 4

struct D3 : S3
{
    int d;     // size 4; alignment 4
};             // size 8; alignment 4

struct D4 : S4
{
    char d;    // size 1; alignment 1
};             // size 3; alignment 1
\end{emcppslisting}

\noindent Finally, virtual functions and virtual base classes invariably introduce an implicit
virtual-table-pointer member having a size and alignment corresponding
to that of a memory address (e.g., 4 or 8) on the target platform:

\begin{emcppslisting}
struct S5
{
    virtual ~S5();
};             // size 8; alignment 8

struct D5 : S5
{
    char d;    // size  1; alignment 1
};             // size 16; alignment 8
\end{emcppslisting}


\subsubsection[Cache lines; L1, L2, and L3 cache; pages; and virtual memory]{Cache lines; L1, L2, and L3 cache; pages; and virtual memory}\label{cache-lines,-l1,-l2,-and-l3-cache,-pages,-and-virtual-memory}

Modern computers are highly complex systems, and a detailed
understanding of their intricacies is unnecessary to achieve most of
the performance benefits. Still, certain general themes and
rough thresholds aid in understanding how to squeeze just a bit
more out of the underlying hardware. In this section, we sketch
fundamental concepts that are common to all modern computer hardware;
although the precise details will vary, the general ideas remain
essentially the same.

In its most basic form, a computer consists of central processing unit
(CPU) having internal registers that access main memory (MM).
Registers in the CPU (on the order of hundreds of bytes) are among the
fastest forms of memory, while main memory (typically many gigabytes) is
orders of magnitude slower. An almost universally observed phenomenon is
that of \emcppsgloss{locality of reference}, which suggests that data that
resides in close proximity (in the virtual address space) is more likely
to be accessed together in rapid succession than more distant data.

To exploit the phenomenon of \emcppsgloss{locality of reference}, computers
introduce the notion of a cache that, while much faster than main
memory, is also much smaller. Programs that attempt to amplify
\emcppsgloss{locality of reference} will, in turn, often be rewarded with
faster run times. The organization of a cache and, in fact, the number of
levels of cache (e.g., L1, L2, L3,$\ldots$) will vary, but the basic
design parameters are, again, more or less the same. A given level of
cache will have a certain total size in bytes (invariably an integral
power of two). The cache will be segmented into what are called
\emcppsgloss[cache line]{cache lines} whose size (a smaller power of two) divides that of
the cache itself. When the CPU accesses main memory, it first looks to
see if that memory is in the cache; if it is, the value is returned
quickly (known as a \emcppsgloss{cache hit}). Otherwise, the cache line(s)
containing that data is (are) fetched (from the next higher level of cache
or from main memory) and placed into the cache (known as a \emcppsgloss{cache
miss}), possibly ejecting other less recently used
ones.{\cprotect\footnote{Conceptually, the cache is often thought of as
being able to hold any arbitrary subset of the most recently accessed
cache lines. This kind of cache is known as \emcppsgloss{fully
associative}. Although it provides the best hit rate, a \emcppsgloss{fully
associative} cache requires the most power along with significant
additional chip area to perform the fully parallel lookup. \emcppsgloss[direct mapped]{Direct-mapped} cache associativity is at the
other extreme. In direct mapped, each memory location has exactly one location
available to it in the cache. If another memory location mapping to
that location is needed, the current cache line must be flushed from
the cache. Although this approach has the lowest hit rate, lookup
times, chip area, and power consumption are all minimized (optimally).
Between these two extremes is a continuum that is referred to as
\emcppsgloss{set associative}. A \emcppsgloss[set associative]{set associate} cache has more than
  one (typically 2, 4, or 8; see \cite{solihin15}, section~5.2.1, ``Placement Policy," pp. 136--141, and \cite{hruska20})
  location in which each memory location in main memory can reside.
  Note that, even with a relatively small $N$, as $N$ increases, an $N$-way
  \emcppsgloss{set associative} cache quickly approaches the hit rate of a fully
  associative cache at greatly reduced collateral cost; for most
  software-design purposes, any loss in hit rate due to set
  associativity of a cache can be safely ignored.}}

Data residing in distinct cache lines is physically independent and can
be written concurrently by multiple threads, possibly running on separate cores or even processors. Logically unrelated data
residing in the same cache line, however, is nonetheless physically
coupled; two threads that write to such logically unrelated data will find
themselves synchronized by the hardware. Such unexpected and typically
undesirable sharing of a cache line by unrelated data acted upon by two
concurrent threads is known as \emcppsgloss{false sharing}. One way of
avoiding \emcppsgloss{false sharing} is to align such data on a cache-line
boundary, thus rendering accidental colocation of such data on the same
cache line impossible. Another (more broad-based) design approach that
avoids lowering cache utilization is to ensure that data acted upon by a
given thread is kept physically separate --- e.g., through the use of
local (arena) memory allocators.{\cprotect\footnote{\cite{{lakos17}},
  \cite{{lakos19}}, \cite{{lakos22}}}}

Finally, even data that is not currently in cache but resides nearby in
MM can benefit from locality. The virtual address space, synonymous with
the size of a \lstinline!void*! (typically 64-bits on modern general-purpose hardware), has historically well exceeded the physical memory
available to the CPU. The operating system must therefore maintain a
mapping (in main memory) from what is resident in physical memory and
what resides in secondary storage (e.g., on disc). In addition,
essentially all modern hardware provides a
\emcppsgloss{TLB}{\cprotect\footnote{A translation-lookaside buffer (TLB) is
a kind of address-translation cache that is typically part of a chip's
memory management unit (MMU). A TLB holds a recently accessed subset
of the complete mapping (itself maintained in MM) from virtual memory
address to physical ones. A TLB is used to reduce access time when the
requisite pages are already resident in memory; its size (e.g., 4K) is
capped at the number of bytes of physical memory (e.g., 32Gb) divided
by the number of bytes in each physical page (e.g., 8Kb), but could be
smaller. Because it resides on chip, is typically an order of
magnitude faster (SRAM versus DRAM), and requires only a single lookup
(as opposed to two or more when going out to MM), there is an enormous
  premium on minimizing TLB misses.}} that caches the addresses of the
most recently accessed physical pages, providing yet another advantage
to having the \emcppsgloss{working set} (i.e., the current set of frequently
accessed pages) remain small and densely packed with relevant
data.{\cprotect\footnote{Note that memory for handle-body
types (e.g., \lstinline!std::vector! or \lstinline!std::deque!) and
especially node-based containers (e.g., \lstinline!std::map! and
\lstinline!std::unordered_map!), originally allocated within a single
page, can --- through deallocation and reallocation (or even move
operations) --- become scattered across multiple (perhaps many)
 pages, thus causing what was originally a relatively small \emcppsgloss{working set}
to no longer fit within physical memory. This phenomenon, known as
\emcppsgloss{diffusion} (which is a distinct concept from
\emcppsgloss{fragmentation}), is what typically leads to a substantial
runtime performance degradation (due to cache line \emcppsgloss{thrashing}) in large,
long-running programs. Such \emcppsgloss{diffusion} can be mitigated by
judicious use of local arena memory allocators (and deliberate
avoidance of \emcppsgloss{move operations} across disparate localities of
  frequent memory usage).}} What's more, dense working sets, in addition
to facilitating hits for repeat access, increase the likelihood that
data that is coresident on a page (or cache line) will be needed soon
(i.e., in effect acting as a prefetch).{\cprotect\footnote{Beneficial prefetching of unrelated data that
is accidentally needed subsequently (e.g., within a single thread) due
to high locality within a cache line (or a physical page) is sometimes referred to as
  \emcppsgloss{true sharing}.}} Table~\ref{table-alignas-appendix} provides a summary of typical physical parameters found in modern computers today.

%%%%%%TODO: TBD: @Lori: Format this as a figure with a caption. @jslakos: Double-check accuracy of values.

\begin{table}[h!]
\begin{center}
\begin{threeparttable}
\caption{Various sizes and access speeds of typical memory for modern \mbox{computers}}\label{table-alignas-appendix} \vspace{1.5ex}
{\small \begin{tabular}{c|c|c}\thickhline
\rowcolor[gray]{.9}    {\sffamily\bfseries Memory Type} & {\sffamily\bfseries Typical Memory Size (Bytes)} & {\sffamily\bfseries Typical Access Times}\\\hline
CPU Registers & 512 \ldots~2048 & $\sim$250ps\\ \hline
Cache Line  & 64 \ldots~256 & NA\\ \hline
L1 Cache & 16Kb \ldots~64Kb & $\sim$1ns\\ \hline
L2 Cache & 1Mb \ldots~2Mb & $\sim$10ns\\ \hline
L3 Cache & 8Mb \ldots~32Mb & $\sim$80ns--120ns\\ \hline
L4 Cache & 32Mb \ldots~128Mb & $\sim$100ns--200ns\\ \hline
Set Associativity  &  2 \ldots~64 & NA\\ \hline
TL & 4 words \ldots~65536 & 10ns \ldots~50ns\\ \hline
Physical Memory Page & 512 \ldots~8192 & 100ns \ldots~500ns\\ \hline
\rule{0pt}{3.5mm}Virtual Memory & $2^{32}$ bytes \ldots~$2^{64}$ bytes & $\sim$10$\mu$s--50$\mu$s\\ \hline
Solid-State Disc (SSD) & 256Gb \ldots~16Tb & $\sim$25$\mu$s--100$\mu$s\\ \hline
Mechanical Disc & Huge & $\sim$5ms--10ms\\ \hline
Clock Speed & NA & $\sim$4GHz \\ \thickhline
\end{tabular}
}
\end{threeparttable}
    \end{center}
\end{table}
\clearpage



\newpage
%\section[{\tt alignof}]{The (Compile-Time) {\SecCode alignof} Operator}\label{alignof}
% 26 March 2021, updating commands
% 27 March 2021, Slava's revisions in and proofed. 
% 28 March 2021, sent to CE, external review, and Josh for code check
% 28 March 2021 JMB, code check done

\emcppsFeature{
    short={\lstinline!alignof!},
    tocshort={\TOCCode alignof},
    long={The (Compile-Time) {\SecCode alignof} Operator},
    toclong={The (Compile-Time) \lstinline!alignof! Operator},
    rhshort={\RHCode alignof},
}{alignof}
\setcounter{table}{0}
\setcounter{footnote}{0}
\setcounter{lstlisting}{0}
%\section[{\tt alignof}]{The (Compile-Time) {\SecCode alignof} Operator}\label{alignof}

The keyword \lstinline!alignof! serves as a compile-time operator used to
query the \emcppsgloss[alignment requirement]{alignment requirements} of a type on the current
platform.

\subsection[Description]{Description}\label{description}

The \lstinline!alignof! operator, when applied to a type, evaluates to an
\emcppsgloss{integral constant expression} that represents the
\emcppsgloss{alignment requirement} of its argument type. Similar to
\lstinline!sizeof!, the (compile-time) value of \lstinline!alignof! is of type
\lstinline!std::size_t!; unlike \lstinline!sizeof! that can accept an
arbitrary expression, \lstinline!alignof! is defined for only type identifiers but often works on expressions
anyway (see \intrarefsimple{annoyances-alignof}). The argument type,
\lstinline!T!, supplied to \lstinline!alignof! must be a
\emcppsgloss{complete type}, a \emcppsgloss{reference type}, or an \emcppsgloss{array
type}. If \lstinline!T! is a \emcppsgloss{complete type}, the result is the
alignment requirement for \lstinline!T!. If \lstinline!T! is a
\emcppsgloss{reference type}, the result is the alignment requirement for the
referenced type. If \lstinline!T! is an \emcppsgloss{array type}, the
result is the alignment requirement for every element in the array. For example, on a platform where \lstinline!sizeof(short)!~\lstinline!==!~\lstinline!2! and \lstinline!alignof(short)!~\lstinline!==!~\lstinline!2!, the following assertions pass:{\cprotect\footnote{According to the C++11 Standard, ``An object of \emcppsgloss{array type} contains a contiguously allocated
  non-empty set of \lstinline!N! subobjects of type \lstinline!T!''  (\cite{cpp11}, section 8.3.4, ``Arrays," paragraph 1, p. 188). Note
  that, for every type \lstinline!T!, \lstinline!sizeof(T)! is always a
  multiple of \lstinline!alignof(T)!; otherwise, storing multiple \lstinline!T! instances in an array would be impossible without padding,
  and the Standard explicitly prohibits padding between array elements.}}:

\begin{emcppslisting}
static_assert(alignof(short)    == 2, "");  // complete type   ((ù{\codeincomments{sizeof}}ù) is 2)
static_assert(alignof(short&)   == 2, "");  // reference type  ((ù{\codeincomments{sizeof}}ù) is 2)
static_assert(alignof(short[5]) == 2, "");  // array type      ((ù{\codeincomments{sizeof}}ù) is 10)
static_assert(alignof(short[])  == 2, "");  // array type      ((ù{\codeincomments{sizeof}}ù) fails)
\end{emcppslisting}
    

\subsubsection[\lstinline!alignof! Fundamental Types]{{\SubsubsecCode alignof} Fundamental Types}\label{alignof-fundamental-types}

Like their size, the alignment requirements of a \lstinline!char!,
\lstinline!signed!~\lstinline!char!, and \lstinline!unsigned!~\lstinline!char! are
guaranteed to be 1 on every conforming
platform. For any other fundamental or pointer type \lstinline!FPT!,
\lstinline!alignof(FPT)! is platform-dependent but is typically approximated well by the type's \emcppsgloss{natural
alignment} --- i.e.,
\lstinline!sizeof(FPT)!~\lstinline!==!~\lstinline!alignof(FPT)!:

\begin{emcppslisting}
static_assert(alignof(char)   == 1, "");  // guaranteed to be 1
static_assert(alignof(short)  == 2, "");  // platform-dependent
static_assert(alignof(int)    == 4, "");  //    "         "
static_assert(alignof(double) == 8, "");  //    "         "
static_assert(alignof(void*)  >= 4, "");  //    "         "
\end{emcppslisting}
    

\subsubsection[\lstinline!alignof! User-Defined Types]{{\SubsubsecCode alignof} User-Defined Types}\label{alignof-user-defined-types}

When applied to user-defined types, alignment is always at least that of
the strictest alignment of any of its arguments' base or member objects. 
%Empty types are defined to have a size (and alignment) of 1 to
%ensure that every object has a unique address.{\cprotect\footnote{An
%exception is made for an object of a type derived from an empty (base)
%class in that neither the size nor the alignment of the derived object
%is affected by the derivation:
%
%\begin{emcppslisting}[style=footcode]
%struct S { int i; }      // size = 4; alignment = 4
%struct E { };            // size = 1; alignment = 1
%struct D : E { int i };  // size = 4; alignment = 4
%\end{emcppslisting}
%      }} 
Compilers will (by default) avoid nonessential padding because any
extra padding would be wasteful of (e.g., cache)
memory:

%\begin{emcppslisting}
%struct S0 { };                          // (ù{\codeincomments{sizeof(S0)}}ù) is  (ù{\codeincomments{1}}ù); (ù{\codeincomments{alignof(S0)}}ù) is  (ù{\codeincomments{1}}ù)
%struct S1 { char c; };                  // (ù{\codeincomments{sizeof(S1)}}ù) is  (ù{\codeincomments{1}}ù); (ù{\codeincomments{alignof(S1)}}ù) is  (ù{\codeincomments{1}}ù)
%struct S2 { short s; };                 // (ù{\codeincomments{sizeof(S2)}}ù) is  (ù{\codeincomments{2}}ù); (ù{\codeincomments{alignof(S2)}}ù) is  (ù{\codeincomments{2}}ù)
%struct S3 { char c; short s; };         // (ù{\codeincomments{sizeof(S3)}}ù) is  (ù{\codeincomments{4}}ù); (ù{\codeincomments{alignof(S3)}}ù) is  (ù{\codeincomments{2}}ù)
%struct S4 { short s1; short s2; };      // (ù{\codeincomments{sizeof(S4)}}ù) is  (ù{\codeincomments{4}}ù); (ù{\codeincomments{alignof(S4)}}ù) is  (ù{\codeincomments{2}}ù)
%struct S5 { int i; char c; };           // (ù{\codeincomments{sizeof(S5)}}ù) is  (ù{\codeincomments{8}}ù); (ù{\codeincomments{alignof(S5)}}ù) is  (ù{\codeincomments{4}}ù)
%struct S6 { char c1; int i; char c2};   // (ù{\codeincomments{sizeof(S6)}}ù) is (ù{\codeincomments{12}}ù); (ù{\codeincomments{alignof(S6)}}ù) is  (ù{\codeincomments{4}}ù)
%struct S7 { char c; short s; int i; };  // (ù{\codeincomments{sizeof(S7)}}ù) is  (ù{\codeincomments{8}}ù); (ù{\codeincomments{alignof(S7)}}ù) is  (ù{\codeincomments{4}}ù)
%struct S8 { double d; };                // (ù{\codeincomments{sizeof(S8)}}ù) is  (ù{\codeincomments{8}}ù); (ù{\codeincomments{alignof(S8)}}ù) is  (ù{\codeincomments{8}}ù)
%struct S9 { double d; char c};          // (ù{\codeincomments{sizeof(S9)}}ù) is (ù{\codeincomments{16}}ù); (ù{\codeincomments{alignof(S9)}}ù) is  (ù{\codeincomments{8}}ù)
%struct SA { long double; };             // (ù{\codeincomments{sizeof(SA)}}ù) is (ù{\codeincomments{16}}ù); (ù{\codeincomments{alignof(SA)}}ù) is (ù{\codeincomments{16}}ù)
%struct SB { long double; char c};       // (ù{\codeincomments{sizeof(SB)}}ù) is (ù{\codeincomments{32}}ù); (ù{\codeincomments{alignof(SB)}}ù) is (ù{\codeincomments{16}}ù)
%\end{emcppslisting}
\begin{emcppslisting}[emcppsbatch=e1]
struct S0 { };                           // sizeof(S0) is  1; alignof(S0) is  1   
struct S1 { char c; };                   // sizeof(S1) is  1; alignof(S1) is  1   
struct S2 { short s; };                  // sizeof(S2) is  2; alignof(S2) is  2   
struct S3 { char c; short s; };          // sizeof(S3) is  4; alignof(S3) is  2   
struct S4 { short s1; short s2; };       // sizeof(S4) is  4; alignof(S4) is  2   
struct S5 { int i; char c; };            // sizeof(S5) is  8; alignof(S5) is  4   
struct S6 { char c1; int i; char c2; };  // sizeof(S6) is 12; alignof(S6) is  4   
struct S7 { char c; short s; int i; };   // sizeof(S7) is  8; alignof(S7) is  4   
struct S8 { double d; };                 // sizeof(S8) is  8; alignof(S8) is  8   
struct S9 { double d; char c; };         // sizeof(S9) is 16; alignof(S9) is  8   
struct SA { long double ld; };           // sizeof(SA) is 16; alignof(SA) is 16   
struct SB { long double ld; char c; };   // sizeof(SB) is 32; alignof(SB) is 16
\end{emcppslisting}

\noindent The size of empty types, such as \lstinline!S0! above, are defined to have the size and alignment of 1 to ensure that each object and member subobject of type \lstinline!S0! has a unique address. However, if an empty type is used as a base, the size of the derived type will not be affected (with some exceptions) due to the  \emcppsgloss{empty base optimization}:

\begin{emcppslisting}[emcppsbatch=e1]
struct D0 : S0 { int i; };  // sizeof(D0) is 4; alignof(D0) is 4
\end{emcppslisting}

The alignment of the base type always affects the derived type's alignment. However, this is observable only for an empty base if it is \emcppsgloss{over-aligned} (see \featureref{\locationc}{alignas}):

\begin{emcppslisting}
struct alignas(8) E { };   // sizeof(E)  is 8; alignof(E)  is 8
struct D1 : E { int i; };  // sizeof(D1) is 8; alignof(D1) is 8
\end{emcppslisting}

Compilers are permitted to increase alignment
(e.g., in the presence of virtual functions, which typically implies a virtual function table pointer) but have certain restrictions
on padding. For example, they must ensure that each comprised type is
itself sufficiently aligned.  Furthermore, sufficient padding must be added so that the alignment of the parent type divides its size, ensuring that storing multiple instances in an
array does not require any padding between array elements (which is explicitly prohibited by the Standard). In other words, the following identities hold for all types, \lstinline!T!, and positive integers,
\lstinline!N!:

\begin{emcppslisting}
#include <cstddef>  // (ù{\codeincomments{std::size\_t}}ù)

template <typename T, std::size_t N>                                            
void f()                                                                        
{                                                                               
    static_assert(0 == sizeof(T) % alignof(T), "guaranteed");
                                 
    T a[N];                                                                     
    static_assert(N == sizeof(a) / sizeof(*a), "guaranteed");                                 
}

\end{emcppslisting}
    
\noindent The alignment of user-defined types can be made artificially stricter
  (but not weaker) using the \lstinline!alignas! specifier (see \featureref{\locationc}{alignas}).
  Also note that, for \emcppsgloss[standard layout types]{standard-layout types}, the address of the
  first member object is guaranteed to be the same as that of the parent
  object (see \featureref{\locationc}{gpods}):

  \begin{emcppslisting}
struct S { int i; };                                                            
class T { public: S s; };                                                       
T t;                                                                            
static_assert(static_cast<void*>(&t.s) == &t, "guaranteed");                                                                             
static_assert(static_cast<void*>(&t.s) == &t.s.i, "guaranteed");
\end{emcppslisting}
      
\noindent This property also holds for unions:

\begin{emcppslisting}
struct { union { char c; float f; double d; }; } u;                             
static_assert(static_cast<void*>(&u) == &u.c, "guaranteed");          
static_assert(static_cast<void*>(&u) == &u.f, "guaranteed");          
static_assert(static_cast<void*>(&u) == &u.d, "guaranteed");
\end{emcppslisting}

    

\subsection[Use Cases]{Use Cases}\label{use-cases}

\subsubsection[Probing the alignment of a type during development]{Probing the alignment of a type during development}\label{probing-the-alignment-of-a-type-during-development}

Both \lstinline!sizeof! and \lstinline!alignof! are often used informally
during development and debugging to confirm the values
of those attributes for a given type on the current platform. For
example:

\begin{emcppslisting}
#include <iostream>  // (ù{\codeincomments{std::cout}}ù)

void f()
{
    std::cout << " sizeof(double): " <<  sizeof(double) << '\n';  //  always 8
    std::cout << "alignof(double): " << alignof(double) << '\n';  // usually 8
}
\end{emcppslisting}
    
\noindent Printing the size and alignment of a \lstinline!struct! along with those of
each of its individual data members can lead to the discovery of
suboptimal ordering of data members, resulting in wasteful extra
padding. As an example, consider two \lstinline!struct!s,
\lstinline!Wasteful! and \lstinline!Optimal!, having the same three data members but
in different order:

\begin{emcppslisting}
struct Wasteful
{
    char   d_c;  // size =  1;  alignment = 1
    double d_d;  // size =  8;  alignment = 8
    int    d_i;  // size =  4;  alignment = 4
};               // size = 24;  alignment = 8

struct Optimal
{
    double d_d;  // size =  8;  alignment = 8
    int    d_i;  // size =  4;  alignment = 4
    char   d_c;  // size =  1;  alignment = 1
};               // size = 16;  alignment = 8
\end{emcppslisting}
    
\noindent Both \lstinline!alignof(Wasteful)! and \lstinline!alignof(Optimal)! are
\lstinline!8! on our platform but \lstinline!sizeof(Wasteful)! is \lstinline!24!, whereas
\lstinline!sizeof(Optimal)! is only \lstinline!16!. Even though these two
\lstinline!struct!s contain the very same data members, the individual
alignment requirements of these members forces the compiler to insert
more total padding between the data members in \lstinline!Wasteful! than is
necessary in \lstinline!Optimal!:

\begin{emcppslisting}
struct Wasteful
{
    char   d_c;           // size =  1;  alignment = 1
    char   padding_0[7];  // size =  7
    double d_d;           // size =  8;  alignment = 8
    int    d_i;           // size =  4;  alignment = 4
    char   padding_1[4];  // size =  4
};                        // size = 24;  alignment = 8

struct Optimal
{
    double d_d;           // size =  8;  alignment = 8
    int    d_i;           // size =  4;  alignment = 4
    char   d_c;           // size =  1;  alignment = 1
    char   padding_0[3];  // size =  3
};                        // size = 16;  alignment = 8
\end{emcppslisting}
    

\subsubsection[Determining if a given buffer is sufficiently aligned]{Determining if a given buffer is sufficiently aligned}\label{determining-if-a-given-buffer-is-sufficiently-aligned}

The \lstinline!alignof! operator can be used to determine if a given (e.g.,
\lstinline!char!) buffer is suitably aligned for storing an object of
arbitrary type. As an example, consider the task of creating a
\emcppsgloss[value semantic]{value-semantic} class, \lstinline!MyAny!, that represents an object
of arbitrary type{\cprotect\footnote{The C++17 Standard Library provides
the (nontemplate) class \lstinline!std::any!, which is a type-safe
container for single values of \emph{any} \emcppsgloss{regular type}. The
implementation strategies surrounding alignment for \lstinline!std::any!
in both \lstinline!libstdc++! and \lstinline!libc++! closely mirror those
used to implement the simplified \lstinline!MyAny! class presented here.
Note that \lstinline!std::any! also records the current \lstinline!typeid!
(on construction or assignment) which can be queried with the \lstinline!type! member
function to determine, at runtime, 
whether a specified type is currently the active one:

\begin{emcppslisting}[emcppsstandards={c++17},style=footcode]
#include <any>  // (ù{\codeincomments{std::any}}ù)
void f(const std::any& object)
{
    if (object.type() == typeid(int)) { /* ... */ }
}
\end{emcppslisting}
      }}:

\begin{emcppslisting}[emcppsbatch=e2]
#include <cassert>   // standard C (ù{\codeincomments{assert}}ù) macro
#include <string>    // (ù{\codeincomments{std::string}}ù)
#include <my_any.h>  // (ù{\codeincomments{MyAny}}ù)
void f()
{
    MyAny obj = 10;               // can be initialized with values of any type
    assert(obj.as<int>() == 10);  // inner data can be retrieved at runtime

    obj = std::string{"hello"};   // can be reassigned from a value of any type
    assert(obj.as<std::string>() == "hello");
}
\end{emcppslisting}
    
\noindent A straightforward implementation of \lstinline!MyAny! would be to allocate
an appropriately sized block of dynamic memory each time a value of a
new type is assigned. Such a naive implementation would force memory
allocations even though the vast majority of values assigned in practice
are small (e.g., fundamental types), most of which would fit within the
space that would otherwise be occupied by just the pointer needed to
refer to dynamic memory. As a practical optimization, we might instead
consider reserving a small buffer  
%(say, roughly{\cprotect\footnote{We
%would likely choose a slightly larger value, e.g., 35 or 39, if that
%space would otherwise be filled with essential padding due to overall
%  alignment requirements.}} 32 bytes) 
  within the footprint of the
\lstinline!MyAny! object to hold the value provided (1) it will fit and (2)
the buffer is sufficiently aligned. The natural implementation of this
type --- typically having a \lstinline!union! of a \lstinline!char! array and a \lstinline!char! pointer as a data member ) --- will naturally result in the alignment requirement of at least that of the \lstinline!char*! (e.g., 4 on a 32-bit
platform and 8 on a 64-bit one){\cprotect\footnote{We could, in
  addition, use the \lstinline!alignas! attribute to ensure that
  the minimal alignment of \lstinline!d_buffer! was at least 8 (or even
  16):

  \begin{emcppslisting}[style=footcode]
  // ...
  alignas(8) char d_buffer[39];  // small buffer aligned to (at least) 8
  // ...
  \end{emcppslisting}
      }}:

%\begin{emcppslisting}
%#include <cstddef>  // (ù{\codeincomments{std::size\_t}}ù)
%
%class MyAny  // nontemplate class
%{
%    union
%    {
%        struct
%        {
%            char*       d_buf_p;  // pointer to dynamic memory if needed
%            std::size_t d_size;   // for (ù{\codeincomments{d\_buf\_p}}ù); same alignment as (ù{\codeincomments{(char*)}}ù)
%        } d_imp;  // Size/alignment of (ù{\codeincomments{d\_imp}}ù) is (ù{\codeincomments{sizeof(d\_buf\_p)}}ù) (e.g., 4 or 8).
%
%        char d_buffer[39];         // small buffer aligned as a (ù{\codeincomments{(char*)}}ù)
%    };  // Size of (ù{\codeincomments{union}}ù) is (ù{\codeincomments{39}}ù); alignment of (ù{\codeincomments{union}}ù) is (ù{\codeincomments{alignof(char*)}}ù).
%
%    bool d_onHeapFlag;             // boolean (discriminator) for union (above)
%
%public:
%    template <typename T>
%    MyAny(const T& x);               // (member template) constructor
%
%    template <typename T>
%    MyAny& operator=(const T& rhs);  // (member template) assignment operator
%
%    template <typename T>
%    const T& as() const;             // (member template) accessor
%
%    // ...
%
%};  // Size of (ù{\codeincomments{MyAny}}ù) is (ù{\codeincomments{40}}ù); alignment of (ù{\codeincomments{MyAny}}ù) is (ù{\codeincomments{alignof(char*)}}ù) (e.g., 8).
%\end{emcppslisting}
\begin{emcppslisting}[emcppsbatch=e2]
// my_any.h:
#include <cstddef>  // (ù{\codeincomments{std::size\_t}}ù)

class MyAny  // nontemplate class
{
    union {
        char *d_buf_p;       // pointer to dynamic memory if needed
        char  d_buffer[39];  // small buffer

    };  // Size of (ù{\codeincomments{union}}ù) is (ù{\codeincomments{39}}ù); alignment of (ù{\codeincomments{union}}ù) is (ù{\codeincomments{alignof(char*)}}ù).

    char d_onHeapFlag;             // boolean (discriminator) for union (above)

public:
    template <typename T>
    MyAny(const T& x);               // (member template) constructor

    template <typename T>
    MyAny& operator=(const T& rhs);  // (member template) assignment operator

    template <typename T>
    const T& as() const;             // (member template) accessor

    // ...

};  // Size of (ù{\codeincomments{MyAny}}ù) is (ù{\codeincomments{40}}ù); alignment of (ù{\codeincomments{MyAny}}ù) is (ù{\codeincomments{alignof(char*)}}ù) (e.g., 8).
\end{emcppslisting}

    
\noindent We chose the size of \lstinline!d_buffer! in the example above to be 39 for two reasons. First, we decided that we want 32-byte types to fit into the buffer, meaning that the size of \lstinline!d_buffer! should be at least 32. Combined with the use of \lstinline!char! for the \lstinline!d_onHeapFlag!, which is guaranteed to have the size of 1, this means that \lstinline!sizeof(MyAny) >= 33!. Second, we want to ensure that no space is wasted on padding. On platforms where \lstinline!alignof(MyAny)! is 8, which will be the case for many 64-bit platforms, \lstinline!sizeof(MyAny)! would be 40, which we choose to achieve by increasing the useful capacity to 39 instead of having the compiler add unused padding.

The (templated) constructor{\cprotect\footnote{In a real-world
implementation, among other improvements, a \emph{forwarding reference} would be used as the
parameter type of \lstinline!MyAny!'s constructor to
\emph{perfectly forward} the argument object into the appropriate
  storage; see \featureref{\locationc}{forwardingref}.}} of \lstinline!MyAny!
can then decide (potentially at compile time) whether to store the given
object \lstinline!x! in the internal small buffer storage or on the heap,
depending on \lstinline!x!'s size and alignment:

\begin{emcppslisting}[emcppsbatch=e2]
template <typename T>
MyAny::MyAny(const T& x)
{
    if (sizeof(x) <= 39 && alignof(T) <= alignof(char*))
    {
        // Store (ù{\codeincomments{x}}ù) in place in the small buffer.
        new(d_buffer) T(x);
        d_onHeapFlag = false;
    }
    else
    {
        // Store (ù{\codeincomments{x}}ù) on the heap and its address in the buffer.
        d_buf_p = reinterpret_cast<char*>(new T(x));
        d_onHeapFlag = true;
    }
}
\end{emcppslisting}
    
\noindent Using the \lstinline!alignof! operator in the constructor
above to check whether the alignment of \lstinline!T! is compatible with
the alignment of the small buffer is necessary to avoid attempting to
store overly aligned objects in place --- even if they would fit in the
39-byte buffer. As an example, consider \lstinline!long!~\lstinline!double!,
which on typical platforms has both a size and alignment of
16. Even though \lstinline!sizeof(long!~\lstinline!double)! (16) is not greater
than 39, \lstinline!alignof(long!~\lstinline!double)! (16) is greater than
that of \lstinline!d_buffer! (8); hence, attempting to store an instance
of \lstinline!long!~\lstinline!double! in the small buffer, \lstinline!d_buffer!, might
--- depending on where the \lstinline!MyAny! object resides in memory --- result
in \emcppsgloss{undefined behavior}. User-defined types that either contain
a \lstinline!long!~\lstinline!double! or have had their alignments
artificially extended beyond 8 bytes are also unsuitable candidates for the
internal buffer even if they might otherwise fit:

\begin{emcppslisting}
struct Unsuitable1 { long double d_value; };
    // Size is 16 (<= 39), but alignment is 16 (> 8).

struct alignas(32) Unsuitable2 { };
    // Size is  1 (<= 39), but alignment is 32 (> 8).
\end{emcppslisting}

\subsubsection[Monotonic memory allocation]{Monotonic memory allocation}\label{monotonic-memory-allocation}

A common pattern in software --- e.g., request/response in client/server
architectures --- is to quickly build up a complex data structure,
use it, and then quickly destroy it. A \emcppsgloss{monotonic allocator} is
a special-purpose memory allocator that returns a monotonically
increasing sequence of addresses into an arbitrary buffer, subject to
specific size and alignment requirements.{\cprotect\footnote{C++17
introduces an alternate interface to supply memory allocators via an
abstract base class. The C++17 Standard Library provides a complete
version of standard containers using this more interoperable design in
a sub-namespace, \lstinline!std::pmr!, where \lstinline!pmr! stands for
\emcppsgloss{polymorphic memory resource}. Also adopted as part of C++17
are two concrete memory resources,
\lstinline!std::pmr::monotonic_buffer_resource! and
  \lstinline!std::pmr::unsynchronized_pool_resource!.}} Especially when
the memory is allocated by a single thread, there are prodigious\footnote{see
\cite{lakos16}} performance benefits to having unsynchronized raw
memory be taken directly off the (always hot) program stack. In what
follows, we will provide the building blocks of a monotonic memory
allocator wherein the \lstinline!alignof! operator plays an essential role.

As a practically useful example, suppose that we want to create a
lightweight\linebreak[4] \mbox{\lstinline!MonotonicBuffer!} class template that will allow us
to allocate raw memory directly from the footprint of the object. Just
by creating an object of an (appropriately sized) instance of this type
on the program stack, memory will naturally come from the stack. For
didactic reasons, we will start with a first pass at this class ---
ignoring alignment --- and then go back and fix it using
\lstinline!alignof! so that it returns properly aligned memory:

\begin{emcppslisting}
#include <cstddef>  // (ù{\codeincomments{std::size\_t}}ù) 

template <std::size_t N>
struct MonotonicBuffer  // first pass at a monotonic memory buffer
{
    char  d_buffer[N];  // fixed-size buffer
    char* d_top_p;      //  next available address

    MonotonicBuffer() : d_top_p(d_buffer) { }
       // Initialize the next available address to be the start of the buffer.

    template <typename T>
    void* allocate()             // BAD IDEA --- doesn't address alignment
                                 // doesn't check buffer limit
    {
        void* result = d_top_p;  // Remember the current next-available address.
        d_top_p += sizeof(T);    // Reserve just enough space for this type.
        return result;           // Return the address of the reserved space.
    }
};
\end{emcppslisting}
    
\noindent \lstinline!MonotonicBuffer! is a class template with one integral template
parameter that controls the size of the \lstinline!d_buffer! member from
which it will dispense memory. Note that, while \lstinline!d_buffer! has
an alignment of 1, the \lstinline!d_top_p! member, used to keep track of
the next available address, has an alignment that is typically 4 or 8
(corresponding to 32-bit and 64-bit architectures, respectively). The
constructor merely initializes the next-address pointer,
\lstinline!d_top_p!, to the start of the local memory pool,
\lstinline!d_buffer!. The interesting part is how the \lstinline!allocate!
function manages to return a monotonically increasing sequence of addresses corresponding to
objects allocated sequentially from the local pool:

\begin{emcppshiddenlisting}[emcppsbatch=e3]
#include <cstddef>  // (ù{\codeincomments{std::size\_t}}ù) 

template <std::size_t N>
struct MonotonicBuffer  // first pass at a monotonic memory buffer
{
    char  d_buffer[N];  // fixed-size buffer
    char* d_top_p;      //  next available address

    MonotonicBuffer() : d_top_p(d_buffer) { }
       // Initialize the next available address to be the start of the buffer.

    template <typename T>
    void* allocate();
};
\end{emcppshiddenlisting}
\begin{emcppslisting}[emcppsbatch=e3]
void test1()
{
    MonotonicBuffer<20> mb;  // On a 64-bit platform, the alignment will be 8.
    char*   cp = static_cast<char*  >(mb.allocate<char  >());  // &d_buffer[ 0]
    double* dp = static_cast<double*>(mb.allocate<double>());  // &d_buffer[ 1]
    short*  sp = static_cast<short* >(mb.allocate<short >());  // &d_buffer[ 9]
    int*    ip = static_cast<int*   >(mb.allocate<int   >());  // &d_buffer[11]
    float*  fp = static_cast<float* >(mb.allocate<float >());  // &d_buffer[15]
}
\end{emcppslisting}
    
\noindent The predominant problem with this first attempt at an implementation of
\lstinline!allocate! is that the addresses returned do not necessarily
satisfy the alignment requirements of the supplied type. A
secondary concern is that there is no internal check to see if
sufficient room remains. To patch this faulty implementation,
we will need a function that, given an initial address and an alignment
requirement, returns the amount by which the address must be rounded up
(i.e., necessary padding) for an object having that alignment
requirement to be properly aligned:

%\begin{emcppslisting}
%std::size_t calculatePadding(const char* address, std::size_t alignment)
%    // Requires: alignment is a (non-negative, integral) power of 2.
%{
%    // rounding up (ù{\codeincomments{X}}ù) to (ù{\codeincomments{N}}ù) (where (ù{\codeincomments{N}}ù) is a power of 2): (ù{\codeincomments{(x + N - 1) \& \textasciitilde(N - 1)}}ù)
%    const std::size_t maxA = alignof(std::max_align_t);
%    const std::size_t a = reinterpret_cast<std::size_t>(address) & (maxA - 1);
%    const std::size_t am1 = alignment - 1;
%    const std::size_t alignedAddress = (a + am1) & ~am1;  // round up
%    return alignedAddress - a;                            // return padding
%}
%\end{emcppslisting}
\begin{emcppslisting}[emcppsbatch=e3]
#include <cstdint>  // (ù{\codeincomments{std::uintptr\_t}}ù)
std::size_t calculatePadding(const char* address, std::size_t alignment)
    // Requires: alignment is a (non-negative, integral) power of 2.
{
    return (alignment - reinterpret_cast<std::uintptr_t>(address)) & 
           (alignment - 1);
}
\end{emcppslisting}
    
\noindent Armed with the
\lstinline!calculatePadding! helper function (above), we are all set to write the final (correct) version of the \lstinline!allocate! method of the \lstinline!MonotonicBuffer! class template:

\begin{emcppslisting}[emcppsbatch=e3]
template <std::size_t N>                                                        
template <typename T>
void* MonotonicBuffer<N>::allocate()
{
    // Calculate just the padding space needed for alignment.
    const std::size_t padding = calculatePadding(d_top_p, alignof(T));

    // Calculate the total amount of space needed.
    const std::size_t delta = padding + sizeof(T);

    // Check to make sure the properly aligned object will fit.
    if (delta > d_buffer + N - d_top_p)  // if (Needed > Total - Used)
    {
        return 0;  // not enough properly aligned unused space remaining
    }

    // Reserve needed space; return the address for a properly aligned object.
    void* alignedAddress = d_top_p + padding;  // Align properly for (ù{\codeincomments{T}}ù) object.
    d_top_p += delta;                          // Reserve memory for (ù{\codeincomments{T}}ù) object.
    return alignedAddress;                     // Return memory for (ù{\codeincomments{T}}ù) object.
}
\end{emcppslisting}
    
\noindent Using this corrected implementation that uses \lstinline!alignof! to pass
the alignment of the supplied type \lstinline!T! to the
\lstinline!calculatePadding! function, the addresses returned from the
benchmark example (above) would be different{\cprotect\footnote{Note
that on a 32-bit architecture, the \lstinline!d_top_p! character
pointer would be only four-byte aligned, which means that the entire
buffer might be only four-byte aligned. In that case, the respective
offsets for \lstinline!cp!, \lstinline!dp!, \lstinline!sp!, \lstinline!ip!, and
\lstinline!bp! in the example for the aligned use case might sometimes
instead be 0, 4, 12, 16, and \lstinline!nullptr!, respectively. If
  desired, we can use the \lstinline!alignas!
  attribute/keyword to artificially constrain the \lstinline!d_buffer!
  data member always to reside on a maximally aligned address boundary,
  thereby improving consistency of behavior, especially on 32-bit
  platforms.}}:

\begin{emcppslisting}[emcppsbatch=e3]
void test2()
{
    MonotonicBuffer<20> mb;  // Assume 64-bit platform (8-byte aligned).
    char*   cp = static_cast<char*  >(mb.allocate<char  >());  // &d_buffer[ 0]
    double* dp = static_cast<double*>(mb.allocate<double>());  // &d_buffer[ 8]
    short*  sp = static_cast<short* >(mb.allocate<short >());  // &d_buffer[16]
    int*    ip = static_cast<int*   >(mb.allocate<int   >());  // 0 (no space)
    bool*   bp = static_cast<bool*  >(mb.allocate<bool  >());  // &d_buffer[18]
}
\end{emcppslisting}
    
\noindent In practice, an object that allocates memory, such as a \lstinline!vector!
or a \lstinline!list!, will be constructed to use an allocator that provides memory that is guaranteed to have either \emcppsgloss{maximal fundamental alignment}, \emcppsgloss{natural alignment}, or alignment that satisfies.

Finally, instead of returning a null pointer when the buffer was
exhausted, we would typically have the allocator fall back to a
geometrically growing sequence of dynamically allocated blocks; the
\lstinline!allocate! method would then fail (i.e., a
\lstinline!std::bad_alloc! exception would somehow be thrown) only if all
available memory were exhausted and the \emcppsgloss[new handler]{\lstinline!new! handler}
were unable to acquire more memory yet still opted to return control
to its caller.

\subsection[Annoyances]{Annoyances}\label{annoyances-alignof}

\subsubsection[\lstinline!alignof! (unlike \lstinline!sizeof!) is defined only on types]{{\SubsubsecCode alignof} (unlike {\SubsubsecCode sizeof}) is defined only on types}\label{alignof-(unlike-sizeof)-is-defined-only-on-types}

Unlike the \lstinline!sizeof! operator, the \lstinline!alignof! operator can accept only a \emph{type}, not an \emph{expression}, as its argument{\cprotect\footnote{Although
the Standard does not require \lstinline!alignof! to work on arbitrary
expressions, \lstinline!alignof! is a common GNU extension and most compilers support
  it. Both Clang and GCC will warn only if \lstinline!-Wpedantic! is set.}}:

\begin{emcppslisting}
static_assert(sizeof(int)  == 4, "");     // OK, (ù{\codeincomments{int}}ù) is a type.
static_assert(alignof(int) == 4, "");     // OK, (ù{\codeincomments{int}}ù) is a type.
static_assert(sizeof(3 + 2) == 4, "");    // OK, (ù{\codeincomments{3 + 2}}ù) is an expression.           
static_assert(alignof(3 + 2) == 4, "");   // Error, (ù{\codeincomments{3 + 2}}ù) is not a type.
\end{emcppslisting}
    
\noindent This asymmetry can result in a need to leverage
\lstinline!decltype! (see \featureref{\locationa}{decltype}) when inspecting an expression instead
of a type:

\begin{emcppshiddenlisting}[emcppsbatch={e4,e5,e6}]
#include <iostream>  // (ù{\codeincomments{std::cout}}ù)
\end{emcppshiddenlisting}
\begin{emcppslisting}[emcppsbatch=e4]
void f()
{
    enum { e_SUCCESS, e_FAILURE } result;
    std::cout << "size: " << sizeof(result) << '\n';
    std::cout << "alignment:" << alignof(decltype(result)) << '\n';
}
\end{emcppslisting}
    
\noindent The same sort of issue occurs in conjunction with modern \emcppsgloss{type
inference} features such as \lstinline!auto! (see \featureref{\locationc}{auto-feature}) and
generic lambdas (see \featureref{\locationd}{genericlambda}). As a real-world example, consider the
generic lambda (C++14) being used to introduce a small
\emph{local function} that prints out information regarding the size and
alignment of a given \lstinline!object!, likely for debugging purposes:

\begin{emcppslisting}[emcppsbatch=e5,emcppsstandards={c++14}]
auto printTypeInformation = [](auto object)
{
    std::cout << "     size: " << sizeof(object) << '\n'
              << "alignment: " << alignof(decltype(object)) << '\n';
};
\end{emcppslisting}
    
\noindent Because there is no explicit type available within the body of the
\lstinline!printTypeInformation! lambda,{\cprotect\footnote{In C++20, referring to the type of a generic lambda parameter explicitly is possible (due to the addition to lambdas of some familiar template syntax):

\begin{emcppslisting}[emcppsbatch=e6,style=footcode,emcppsstandards={c++20}]
auto printTypeInformation = []<typename T>(T object)
{
    std::cout << "     size: " << sizeof(T) << '\n'
              << "alignment: " << alignof(T) << '\n';
};
\end{emcppslisting}
      }} a programmer wishing to remain entirely within the C++
Standard{\cprotect\footnote{Note that \lstinline!alignof(object)! will work
on every major compiler (GCC~10.x, Clang~10.x, and MSVC~19.x) as a
  nonstandard extension.}} is forced to use the
\lstinline!decltype! construct explicitly to first obtain
the type of \lstinline!object! before passing it on to \lstinline!alignof!.

\subsection[See Also]{See Also}\label{see-also}

\begin{itemize}
\item{\seealsoref{decltype}{\seealsolocationa}helps work around \lstinline!alignof!'s limitation of accepting only a type, not an expression (see \intrarefsimple{annoyances-alignof}).}
\item{\seealsoref{alignas}{\seealsolocationc}can be used to provide an artificially stricter alignment (e.g., more than \emcppsgloss{natural alignment}).}

\end{itemize}

\subsection[Further Reading]{Further Reading}\label{further-reading}

None so far





\newpage
%\cftaddtitleline{toc}{section}{\protect\hspace*{-5em}\framebox{\sffamily\itshape\cppxx}\vspace*{-3ex}}{}
% \section[{\tt auto}]{{\SecCode auto}\checkoddpage\Cbox}\label{auto-feature}
%

\emcppsFeature{
    short={\lstinline!auto! Variables},
    tocshort={{\TOCCode auto} Variables},
    long={Variables of Automatically Deduced Type},
    rhshort={{\RHCode auto} Variables},
}{auto}
\label{auto-feature}
\setcounter{table}{0}
\setcounter{footnote}{0}
\setcounter{lstlisting}{0}
%\section[{\tt auto} Variables]{Variables of Automatically Deduced Type}\label{auto}\label{auto-feature}

placeholder text.........
 
 
%\cftaddtitleline{toc}{section}{\protect\hspace*{-5em}\framebox{\sffamily\itshape\cppxx}\vspace*{-3ex}}{}
%\section[{\tt auto} Variables]{Variables of Automatically Deduced Type}\label{auto}\label{auto-feature}
% 3 March 2021, See Also lines need to be updated in current style; waiting to hear back from Vittorio
% 4 March 2021, ready for Josh's code check
% 6 March 2021, back from Josh's code check
% 23 March 2021, copyedits in and proofed



\emcppsFeature{
    short={\lstinline!auto! Variables},
    tocshort={{\TOCCode auto} Variables},
    long={Variables of Automatically Deduced Type},
    rhshort={{\RHCode auto} Variables},
}{auto}
\label{auto-variables}
\label{auto-feature}
\setcounter{table}{0}
\setcounter{footnote}{0}
\setcounter{lstlisting}{0}
%\section[{\tt auto} Variables]{Variables of Automatically Deduced Type}\label{auto}\label{auto-feature}
%\subsection[\tt{auto} Variables]{{\SubsecCode auto} Variables}\label{auto-variables}

The keyword \lstinline!auto! was repurposed in C++11 to act as a \emcppsgloss{placeholder type} such that, when used in place of a type as part of a variable declaration, the compiler will deduce the variable
type from its initializer.

\subsection[Description]{Description}\label{description}

Prior to
C++11, the \lstinline!auto! keyword could be used as a \emcppsgloss{storage class specifier} for objects declared at block scope and in function
parameter lists to indicate automatic storage duration, which is the
default for these kinds of declarations. The \lstinline!auto! keyword was
repurposed in C++11 alongside the deprecation of the \lstinline!register! keyword as a storage class
  specifier.\footnote{The deprecated keyword \lstinline!register! was removed as of C++17, but the name remains reserved for future use.}
  
  Starting with C++11, the \lstinline!auto! keyword may be used instead of an explicit type's name
when declaring variables. In such cases, the variable's type is deduced
from its initializer by the compiler applying the \emcppsgloss{placeholder type} deduction rules, which, apart from a single exception for list
initializers (see \intraref{potential-pitfalls-auto}{surprising-deduction-for-list-initialization}),
are the same as the rules for \emcppsgloss{function template argument type deduction}:

\begin{emcppslisting}
auto two = 2;     // Type of (ù{\codeincomments{two}}ù) is deduced to be (ù{\codeincomments{int}}ù)
auto pi = 3.14f;  // Type of (ù{\codeincomments{pi}}ù) is deduced to be (ù{\codeincomments{float}}ù).
\end{emcppslisting}
    
\noindent The types of the \lstinline!two! and \lstinline!pi! variables above are
deduced in the same manner as they would be if the same initializer were
passed to a function template taking a single argument of the template
type by value:

\begin{emcppslisting}
template <typename T> void deducer(T);

void testDeduction()
{
    deducer(2);      // (ù{\codeincomments{T}}ù) is deduced to be (ù{\codeincomments{int}}ù).
    deducer(3.14f);  // (ù{\codeincomments{T}}ù) is deduced to be (ù{\codeincomments{float}}ù).
}
\end{emcppslisting}
    
\noindent If the variable declared with \lstinline!auto! does not have an
initializer, if its name appears in the expression used to initialize
it, or if the initializers of multiple variables in the same declaration
don't deduce the same type, the program is ill formed:

\begin{emcppslisting}
auto x;                // Error, declaration of (ù{\codeincomments{auto x}}ù) has no initializer.
auto n = sizeof(n);    // Error, use of (ù{\codeincomments{n}}ù) before deduction of (ù{\codeincomments{auto}}ù)
auto i = 3, f = 0.3f;  // Error, inconsistent deduction for (ù{\codeincomments{auto}}ù)
\end{emcppslisting}
    
\noindent Just like the function template argument deduction never deduces a
reference type for its by-value argument, a variable declared with
unqualified \lstinline!auto! is never deduced to have a reference type:

\begin{emcppslisting}
int  val = 3;
int& ref = val;
auto tmp = ref;  // Type of (ù{\codeincomments{tmp}}ù) is deduced to be (ù{\codeincomments{int}}ù), not (ù{\codeincomments{int\&}}ù).
\end{emcppslisting}
    
\noindent Augmenting \lstinline!auto! with reference and cv-qualifiers, however,
allows controlling whether the deduced type is a reference and whether
it is \lstinline!const! and/or \lstinline!volatile!:

\begin{emcppslisting}
auto val = 3;
    // Type of (ù{\codeincomments{val}}ù) is deduced to be (ù{\codeincomments{int}}ù),
    // the same as the argument for (ù{\codeincomments{template <typename T> void deducer(T)}}ù).

const auto cval = val;
    // Type of (ù{\codeincomments{cval}}ù) is deduced to be (ù{\codeincomments{const int}}ù),
    // the same as the argument for (ù{\codeincomments{template <typename T> void deducer(const T)}}ù).

auto& ref = val;
    // Type of (ù{\codeincomments{ref}}ù) is deduced to be (ù{\codeincomments{int\&}}ù),
    // the same as the argument for (ù{\codeincomments{template <typename T> void deducer(T\&)}}ù).

auto& cref1 = cval;
    // Type of (ù{\codeincomments{cref1}}ù) is deduced to be (ù{\codeincomments{const int\&}}ù),
    // the same as the argument for (ù{\codeincomments{template <typename T> void deducer(T\&)}}ù).

const auto& cref2 = val;
    // Type of (ù{\codeincomments{cref2}}ù) is deduced to be (ù{\codeincomments{const int\&}}ù),
    // the same as the argument for (ù{\codeincomments{template <typename T> void deducer(const T\&)}}ù).
\end{emcppslisting}
    
\noindent Note that qualifying \lstinline!auto! with \lstinline!&&!
does \emph{not} result in deduction of an \romeovalue{rvalue} reference (see
\featureref{\locationc}{Rvalue-References}), but, in line with function template
argument deduction rules, would be treated as a \emph{forwarding
reference} (see \featureref{\locationc}{forwardingref}). This means that a
variable declared with \lstinline!auto&&! will result in an \romeovalue{lvalue} or an
\romeovalue{rvalue} reference depending on the value category of its initializer:

\begin{emcppslisting}[emcppsbatch={e1,e2}]
double doStuff();

      int val  = 3;
const int cval = 7;

// Deduction rules are the same as for (ù{\codeincomments{template <typename T> void deducer(T\&\&)}}ù):

auto&& lref1 = val;
    // Type of (ù{\codeincomments{lref1}}ù) is deduced to be (ù{\codeincomments{int\&}}ù).

auto&& lref2 = cval;
    // Type of (ù{\codeincomments{lref2}}ù) is deduced to be (ù{\codeincomments{const int\&}}ù).

auto&& rref = doStuff();
    // Type of (ù{\codeincomments{rref}}ù) is deduced to be (ù{\codeincomments{double\&\&}}ù).
\end{emcppslisting}
    
\noindent Similarly to references, explicitly specifying that a
pointer type is to be deduced is possible. If the supplied initializer is not of a
pointer type, the compiler will issue an error:

\begin{emcppslisting}[emcppsbatch=e1]
const auto* cptr = &cval;
    // Type of (ù{\codeincomments{cptr}}ù) is deduced to be (ù{\codeincomments{const int*}}ù).
    // same as argument for (ù{\codeincomments{template <typename T> void deducer(const T*)}}ù)

auto* cptr2 = cval;  // Error, cannot deduce (ù{\codeincomments{auto*}}ù) from (ù{\codeincomments{cval}}ù)
\end{emcppslisting}
    
\noindent The compiler can also be instructed to deduce pointers to functions,
data members, and member functions (but see \intraref{annoyances-auto}{not-all-template-argument-deduction-constructs-are-allowed-for-auto}):

\begin{emcppslisting}[emcppsbatch=e2]
float freeF(float);

struct S
{
    double d_data;
    int memberF(long);
};

auto (*fptr)(float) = &freeF;
  // Type of (ù{\codeincomments{fptr}}ù) is deduced to be (ù{\codeincomments{float (*)(float)}}ù).
  // same as argument for (ù{\codeincomments{template <typename T> void deducer(T (*)(float))}}ù)

const auto S::* mptr = &S::d_data;
  // Type of (ù{\codeincomments{mptr}}ù) is deduced to be (ù{\codeincomments{const double S::*}}ù).
  // same as argument for (ù{\codeincomments{template <typename T> void deducer(T S::*)}}ù)

auto (S::* mfptr)(long) = &S::memberF;
  // Type of (ù{\codeincomments{mfptr}}ù) is deduced to be (ù{\codeincomments{int (S::*)(long)}}ù).
  // same as argument for (ù{\codeincomments{template <typename T> void deducer(T (S::*)(long))}}ù)
\end{emcppslisting}
    
\noindent Unlike references, pointer types may be deduced by \lstinline!auto! alone.
Therefore, different forms of \lstinline!auto! can be used to declare a
variable of a pointer type:

\begin{emcppslisting}[emcppsbatch=e2]
auto  cptr1 = &cval;  // (ù{\codeincomments{const int*}}ù)
auto *cptr2 = &cval;  //   "     "

auto   fptr1         = &freeF;  // (ù{\codeincomments{float (*)(float)}}ù)
auto  *fptr2         = &freeF;  //   "         "
auto (*fptr3)(float) = &freeF;  //   "         "

auto      mptr1 = &S::d_data;  // (ù{\codeincomments{double S::*}}ù)
auto S::* mptr2 = &S::d_data;  //    "     "

auto       mfptr1        = &S::memberF;  // (ù{\codeincomments{int (S::*)(long)}}ù)
auto (S::* mfptr2)(long) = &S::memberF;  //  "     "     "
\end{emcppslisting}
    
\noindent Note, however, that because regular and member pointers are
incompatible, \lstinline!auto*! cannot be used to deduce pointers to data
members and member functions:

\begin{emcppslisting}[emcppsbatch=e2]
auto *mptr3  = &S::d_data;   // Error, cannot deduce (ù{\codeincomments{auto*}}ù) from (ù{\codeincomments{\&S::d\_data}}ù)
auto *mfptr3 = &S::memberF;  // Error, cannot deduce (ù{\codeincomments{auto*}}ù) from (ù{\codeincomments{\&S::memberF}}ù)
\end{emcppslisting}
    
\noindent In addition, storage class specifiers as well as the \lstinline!constexpr!
(see \featureref{\locationc}{constexprvar}) specifier can be
applied to variables that use \lstinline!auto! in their declaration:

\begin{emcppslisting}
thread_local     const auto* logPrefix = "mylib";
static constexpr       auto  pi        = 3.1415926535f;
\end{emcppslisting}
    
\noindent Finally, \lstinline!auto! variables may be declared in any location that
allows declaring a variable supplied with an initializer with a single
exception of nonstatic data members (see \intraref{annoyances-auto}{auto-not-allowed-for-non-static-data-members}):


\begin{emcppshiddenlisting}[emcppsbatch=e3]
#include <vector>     // (ù{\codeincomments{std::vector}}ù)
int sendRequest();
int responseStatus();
bool haveMoreWork();
std::vector<int> vec;
\end{emcppshiddenlisting}
\begin{emcppslisting}[emcppsbatch=e3]
// namespace scope
auto globalNamespaceVar = 3.;

namespace ns
{
    static auto nsNamespaceVar = "...";
}

void f()
{
    // block scope
    constexpr auto blockVar = 'a';

    // condition of (ù{\codeincomments{if}}ù), (ù{\codeincomments{switch}}ù), and (ù{\codeincomments{while}}ù) statements
    if     (auto rc        = sendRequest())    { /*...*/ }
    switch (auto status    = responseStatus()) { /*...*/ }
    while  (auto keepGoing = haveMoreWork())   { /*...*/ }

    // init-statement of (ù{\codeincomments{for}}ù) loops
    for (auto it = vec.begin(); it != vec.end(); ++it) { /*...*/ }

    // range declaration of range-based (ù{\codeincomments{for}}ù) loops
    for (const auto& constVal : vec) { /*...*/ }
}

struct S
{
    // static data members
    static const auto k_CONSTANT = 11u;
};
\end{emcppslisting}
    

\subsection[Use Cases]{Use Cases}\label{use-cases-auto}

\subsubsection[Ensuring variable initialization]{Ensuring variable initialization}\label{ensuring-variable-initialization}

Consider a defect introduced due to mistakenly leaving a variable
uninitialized:

\begin{emcppshiddenlisting}[emcppsbatch=e4]
struct Cursor {
    bool next();
};
Cursor cursor;
\end{emcppshiddenlisting}
\begin{emcppslisting}[emcppsbatch=e4]
void testUninitializedInt()
{
    int recordCount;
    while (cursor.next()) { ++recordCount; }  // Bug, undefined behavior
}
\end{emcppslisting}
    
\noindent Variables declared with \lstinline!auto! must be initialized. Using
\lstinline!auto! might, therefore, prevent such defects:

\begin{emcppslisting}[emcppsbatch=e4,emcppserrorlines={4}]
void testUnitializedAuto()
{
    auto recordCount; // Error, declaration of (ù{\codeincomments{recordCount}}ù) has no initializer
    while (cursor.next()) { ++recordCount; }
}
\end{emcppslisting}
    
\noindent In addition, the initialization requirement might encourage a good
practice of reducing the scope of local variables.

\subsubsection[Avoiding redundant type name repetition]{Avoiding redundant type name repetition}\label{avoiding-redundant-type-name-repetition}

Certain function templates require that the caller explicitly specify
the type that the function uses as its return type. For example, the
\lstinline!std::make_shared<TYPE>! function returns a
\lstinline!std::shared_ptr<TYPE>!.{\cprotect\footnote{Often, such
functions are associated with \lstinline!optional! and \lstinline!variant!
types, which were standardized in C++17:

\begin{emcppshiddenlisting}[emcppsbatch=e5]
#include <string>    // (ù{\codeincomments{std::string}}ù)
#include <variant>   // (ù{\codeincomments{std::variant}}ù)
#include <optional>  // (ù{\codeincomments{std::optional}}ù)
\end{emcppshiddenlisting}
\begin{emcppslisting}[emcppsbatch=e5,emcppsstandards={c++17},style=footcode]
std::variant<std::string, int, double> valueVariant;

// Without (ù{\fncodeincomments{auto}}ù):
std::optional<std::string> greeting1 = std::make_optional<std::string>("Hello");
const std::string& valueString1 = std::get<std::string>(valueVariant);

// With (ù{\fncodeincomments{auto}}ù):
auto greeting2 = std::make_optional<std::string>("Hello");
const auto& valueString2 = std::get<std::string>(valueVariant);
\end{emcppslisting}
      }} If a variable's type is specified explicitly, such declarations
redundantly repeat the type. The use of \lstinline!auto! obviates this repetition:

\begin{emcppshiddenlisting}[emcppsbatch=e5]
#include <memory>  // (ù{\codeincomments{std::shared\_ptr}}ù), (ù{\codeincomments{std::unique\_ptr}}ù)
struct RequestContext{};
struct Socket{};
\end{emcppshiddenlisting}
\begin{emcppslisting}[emcppsbatch=e5]
// Without (ù{\codeincomments{auto}}ù):
std::shared_ptr<RequestContext> context1 = std::make_shared<RequestContext>();
std::unique_ptr<Socket>         socket1  = std::make_unique<Socket>();

// With (ù{\codeincomments{auto}}ù):
auto context2 = std::make_shared<RequestContext>();
auto socket2  = std::make_unique<Socket>();
\end{emcppslisting}
    

\subsubsection[Preventing unexpected implicit conversions]{Preventing unexpected implicit conversions}\label{preventing-unexpected-implicit-conversions}

Using \lstinline!auto! might prevent defects arising from explicitly
specifying a variable's type that is distinct --- yet implicitly
convertible --- from its initializer. As an example, the code below has
a subtle defect that can lead to performance degradation or incorrect
semantics:

\begin{emcppshiddenlisting}[emcppsbatch={e6,e7}]
#include <map>      // (ù{\codeincomments{std::map}}ù)
#include <utility>  // (ù{\codeincomments{std::pair}}ù)
struct User {};
\end{emcppshiddenlisting}
\begin{emcppslisting}[emcppsbatch=e6]
void testManualForLoop()
{
    std::map<int, User> users{/*...*/};
    for (const std::pair<int, User>& idUserPair : users)
    {
        // ...
    }
}
\end{emcppslisting}
    
\noindent On every iteration, the \lstinline!idUserPair! will be bound to a
\emph{copy} of the corresponding pair in the \lstinline!users! map. This
happens because the type returned by dereferencing the \lstinline!map!'s
iterator is \lstinline!std::pair<const!~\lstinline!int,!~\lstinline!User>!, which
is implicitly convertible to \lstinline!std::pair<int,!~\lstinline!User>!.
Using \lstinline!auto! would allow the compiler to deduce the correct type
and avoid this unnecessary and potentially expensive copy:

\begin{emcppslisting}[emcppsbatch=e7]
void testAutoForLoop()
{
    std::map<int, User> users{/*...*/};
    for (const auto& idUserPair : users)
    {
        // (ù{\codeincomments{auto}}ù) is deduced as (ù{\codeincomments{std::pair<const int, User>}}ù).
    }
}
\end{emcppslisting}
    

\subsubsection[Declaring variables of implementation-defined or compiler-synthesized types]{Declaring variables of implementation-defined or compiler-synthesized types}\label{declaring-variables-of-implementation-defined-or-compiler-synthesized-types}

Using \lstinline!auto! is the only way to declare variables of
implementation-defined or compiler-synthesized types, such as
lambda expressions (see \featureref{\locationc}{lambda}). While in some cases 
using type erasure to avoid the need to spell out the type is possible, doing so typically
comes with additional overhead. For example, storing a lambda closure in
a \lstinline!std::function! might entail an allocation on construction and
virtual dispatch upon every call:

\begin{emcppshiddenlisting}[emcppsbatch=e8]
#include <functional>      // (ù{\codeincomments{std::function}}ù)
void saveCurrentWork();
\end{emcppshiddenlisting}
\begin{emcppslisting}[emcppsbatch=e8]
void testCallbacks()
{
    std::function<void()> errorCallback0 = [&]{ saveCurrentWork(); };
        // OK, implicit conversion from anonymous closure type to
        // (ù{\codeincomments{std::function<void()>}}ù), which incurs additional overhead

    auto errorCallback1 = [&]{ saveCurrentWork(); };
        // Better, deduces the compiler-synthesized type
}
\end{emcppslisting}
    

\subsubsection[Declaring variables of complex and deeply nested types]{Declaring variables of complex and deeply nested types}\label{declaring-variables-of-complex-and-deeply-nested-types}

\lstinline!auto! can be used to declare variables of types that are impractical to spell and/or do not convey useful information to the reader.  A typical example is avoiding the
need for spelling out the iterator type of a container:

\begin{emcppshiddenlisting}[emcppsbatch=e9]
#include <vector>      // (ù{\codeincomments{std::vector}}ù)
\end{emcppshiddenlisting}
\begin{emcppslisting}[emcppsbatch=e9]
void doWork(const std::vector<int>& data)
{
    // without (ù{\codeincomments{auto}}ù):
    for (std::vector<int>::const_iterator it = data.begin();
         it != data.end();
         ++it) {
        // ...
    }

    // with (ù{\codeincomments{auto}}ù):
    for (auto it = data.begin(); it != data.end(); ++it) { /*...*/ }
}
\end{emcppslisting}
    
\noindent Furthermore, the need for such types can arise, for example, when
storing intermediate results of \emcppsgloss[expression template]{expression templates} whose types can
be deeply nested and unreadable and might even differ between versions of
the same library:

\begin{emcppshiddenlisting}[emcppsbatch=e10]
struct Employee {
  int name;
};

struct JoinedInYear {
JoinedInYear(int i) {}
};

template <typename T, typename U>
struct FilterRange {};

template <typename T, int (Employee::*)>
struct TransformRange {};

struct Es {};;
Es employees;

int filter(const JoinedInYear&);
struct T {};
T transform(int (Employee::*));

struct A {};
A operator|(const Es&,int);

TransformRange<
    FilterRange<decltype(employees), JoinedInYear>,
    &Employee::name>
operator|(const A&, const T&);

\end{emcppshiddenlisting}
\begin{emcppslisting}[emcppsbatch=e10]
// without (ù{\codeincomments{auto}}ù):
TransformRange<
    FilterRange<decltype(employees), JoinedInYear>,
    &Employee::name
> newEmployeeNames1 =
    employees | filter(JoinedInYear(2019))
              | transform(&Employee::name);

// with (ù{\codeincomments{auto}}ù):
auto newEmployeeNames2 =
    employees | filter(JoinedInYear(2019))
              | transform(&Employee::name);
\end{emcppslisting}
    

\subsubsection[Improving resilience to library code changes]{Improving resilience to library code changes}\label{improving-resilience-to-library-code-changes}

\lstinline!auto! may be used to indicate that code using the variable
doesn't rely on a specific type but rather on certain requirements that
the type must satisfy. Such an approach might give library implementers
more freedom to change return types without affecting the semantics of
their clients' code in projects where automated large-scale refactoring
tools are not available (but see \intraref{potential-pitfalls-auto}{lack-of-interface-restrictions}). As an example,
consider the following library function:

\begin{emcppshiddenlisting}[emcppsbatch=e11]
#include <vector>     // (ù{\codeincomments{std::vector}}ù)
struct Node {};
\end{emcppshiddenlisting}
\begin{emcppslisting}[emcppsbatch=e11]
std::vector<Node> getNetworkNodes();
    // Return a sequence of nodes in the current network.
\end{emcppslisting}
    
\noindent As long as the return value of the \lstinline!getNetworkNodes! function is
only used for iteration, it is not pertinent that a \lstinline!std::vector! is
returned. If clients use \lstinline!auto! to initialize
variables storing the return value of this function, the implementers of
\lstinline!getNetworkNodes! can migrate from \lstinline!std::vector! to,
for example, \lstinline!std::deque! without breaking their clients'
code.

\begin{emcppslisting}[emcppsbatch=e11]
// without (ù{\codeincomments{auto}}ù):
void testConcreteContainer()
{
    const std::vector<Node> nodes = getNetworkNodes();
    for (const Node& node : nodes) { /*...*/ }
        // prevents migration
}

void testDeducedContainer()
{
    // with (ù{\codeincomments{auto}}ù):
    const auto nodes = getNetworkNodes();
    for (const Node& node : nodes) { /*...*/ }
        // The return type of (ù{\codeincomments{getNetworkNodes}}ù) can be silently
        // changed while retaining correctness of the user code.
}
\end{emcppslisting}
    

\subsection[Potential Pitfalls]{Potential Pitfalls}\label{potential-pitfalls-auto}

\subsubsection[Compromised readability]{Compromised readability}\label{compromised-readability}

Using \lstinline!auto! might sometimes hide important semantic information contained in a
variable's type, which could increase the cognitive load for readers. In conjunction
with unclear variable naming, disproportionate use of \lstinline!auto! can
make code difficult to read and maintain.

\begin{emcppshiddenlisting}[emcppsbatch=e12]
#include <vector>  // (ù{\codeincomments{std::vector}}ù)
#include <string>  // (ù{\codeincomments{std::string}}ù)
std::vector<std::string> parseArgs(int argc, char** argv);
\end{emcppshiddenlisting}
\begin{emcppslisting}[emcppsbatch=e12]
int main(int argc, char** argv)
{
    const auto args0 = parseArgs(argc, argv);
        // The behavior of (ù{\codeincomments{parseArgs}}ù) and operations available on (ù{\codeincomments{args0}}ù) is unclear.

    const std::vector<std::string> args1 = parseArgs(argc, argv);
        // It is clear what (ù{\codeincomments{parseArgs}}ù) does and what can be done with (ù{\codeincomments{args1}}ù).
}
\end{emcppslisting}
    
\noindent Although reading the contract of the \lstinline!parseArgs!
function at least once may be necessary to fully understand its behavior, using an
explicit type's name at the call site helps readers understand its
purpose.

\subsubsection[Unintentional copies]{Unintentional copies}\label{unintentional-copies}

Since the rules for function template type deduction apply to
\lstinline!auto!, appropriate cv-qualifiers and declarator modifiers
(\lstinline!&!, \lstinline!&&!, \lstinline!*!, etc.) must be applied to avoid
unnecessary copies that might negatively affect both code performance
and correctness. For example, consider a function that capitalizes a
user's name:

\begin{emcppshiddenlisting}[emcppsbatch={e13,e14}]
#include <string>  // (ù{\codeincomments{std::string}}ù)
#include <cctype>  // (ù{\codeincomments{std::toupper}}ù)
struct User {
   std::string &name();
};
\end{emcppshiddenlisting}
\begin{emcppslisting}[emcppsbatch=e13]
void capitalizeName(User& user)
{
    if (user.name().empty())
    {
        return;
    }

    user.name()[0] = std::toupper(user.name()[0]);
}
\end{emcppslisting}
    
\noindent This function was then incorrectly refactored to avoid repetition of
the \lstinline!user.name()! invocation. However, a missing reference
qualification leads not only to an unnecessary copy of the string, but
also to the function failing to perform its job:

\begin{emcppslisting}[emcppsbatch=e14]
void capitalizeName(User& user)
{
    auto name = user.name();  // Bug, unintended copy

    if (name.empty())
    {
         return;
    }

    name[0] = std::toupper(name[0]);  // Bug, changes the copy
}
\end{emcppslisting}
    
\noindent Furthermore, even a fully cv-ref-qualified \lstinline!auto! might still
prove inadequate in cases as simple as introducing a variable for a
returned-temporary value. As an example, consider refactoring the
contents of this simple function:

\begin{emcppshiddenlisting}[emcppsbatch={e15,e15a,e16}]
struct V {};
V getValue();
void useValue(const V&v);
\end{emcppshiddenlisting}
\begin{emcppslisting}[emcppsbatch=e15]
void testExpression()
{
    useValue(getValue());
}
\end{emcppslisting}

For debugging or readablity, it can help to use an intermediate variable to
store the results of \lstinline!getValue()!:

\begin{emcppslisting}[emcppsbatch=e15a]
void testRefactoredExpression()
{
    auto&& tempValue = getValue();
    useValue(tempValue);
}
\end{emcppslisting}
    
\noindent The above invocation of \lstinline!useValue! is not equivalent to the
original expression; the semantics of the program might have changed because \lstinline!tempValue! is an \romeovalue{lvalue} expression. To get
close to the original semantics, \lstinline!std::forward! and
\lstinline!decltype! must be used to propagate the original value category
of \lstinline!getValue()! to the invocation of \lstinline!useValue! (see
\featureref{\locationc}{forwardingref}):

\begin{emcppslisting}[emcppsbatch=e16]
#include <utility>  // (ù{\codeincomments{std::forward}}ù)
void testBetterRefactoredExpression()
{
    auto&& tempValue = getValue();
    useValue(std::forward<decltype(tempValue)>(tempValue));
}
\end{emcppslisting}
    
\noindent Note that, even with the latest changes, the code above achieves the same result but in a somewhat different way because
 \lstinline!std::forward<decltype(tempValue)>(tempValue)! is
an \romeovalue{xvalue} expression whereas \lstinline!getValue()! is a
\romeovalue{prvalue} expression; see \featureref{\locationc}{Rvalue-References}.

\subsubsection[Unexpected conversions (and lack of expected ones)]{Unexpected conversions (and lack of expected ones)}\label{unexpected-conversions-(and-lack-of-expected-ones)}

Compulsively declaring variables using \lstinline!auto!, even in cases
where the desired type has to be spelled out in the initializer, allows
explicit conversions to be used where they would not be applicable
otherwise. For an example of unintentional consequences of allowing such
conversions, consider a function template that is intended to combine
two \lstinline!chrono! duration values:

\begin{emcppslisting}
#include <chrono>  // (ù{\codeincomments{std::chrono::seconds}}ù)
template <typename Duration1, typename Duration2>
std::chrono::seconds combine_durations(Duration1 d1, Duration2 d2)
{
    auto d = std::chrono::seconds{d1 + d2};
    // ...
}
\end{emcppslisting}
    
\noindent This function template will be successfully instantiated and compiled
even when its arguments are two \lstinline!int!s. Were \lstinline!auto! not
used in this situation --- i.e., were \lstinline!d! declared as
\lstinline!std::chrono::seconds!~\lstinline!d!~\lstinline!=!~\lstinline!d1!~\lstinline!+!~\lstinline!d2;!
--- the code would fail to compile because the conversion from
\lstinline!int! to \lstinline!seconds! is explicit, which would indicate a
likely defect in the caller's code.

In fairness, a better solution for this particular problem would be to declare \lstinline!combined_durations! in a more restricted manner:

\begin{emcppslisting}
template <typename R1, typename P1, typename R2, typename P2>
std::chrono::seconds combine_durations(std::chrono::duration<R1, P1> d1,
                                       std::chrono::duration<R2, P2> d2)
{
    auto d = std::chrono::seconds{d1 + d2};
    // ...
}
\end{emcppslisting}

%\begin{emcppslisting}
%template <typename R1, typename P1, typename R2, typename P2>
%std::chrono::seconds combine_durations(std::chrono::duration<R1, P1> d1,
%                                       std::chrono::duration<R2, P2> d2)
%{
%    auto d = std::chrono::seconds{d1 + d2};
%    // ...
%}
%\end{emcppslisting}

\noindent The definition above obviates any possible risks related to the use of \lstinline!auto! because calls with undesirable argument types are not accepted in the first place. 

Conversely, some conversions that would be expected to happen might be
missed when using \lstinline!auto! instead of an explicitly specified type.
For example, \lstinline!auto! might deduce a proxy type that might lead to
difficult-to-diagnose defects:

\begin{emcppshiddenlisting}[emcppsbatch=e17]
#include <vector>     // (ù{\codeincomments{std::vector}}ù)
std::vector<bool> loadFlags();
\end{emcppshiddenlisting}
\begin{emcppslisting}[emcppsbatch=e17]
void testProxyDeduction()
{
    std::vector<bool> flags = loadFlags();

    auto firstFlag = flags[0];  // deduces a proxy type, not (ù{\codeincomments{bool}}ù)
    flags.clear();

    if (firstFlag)  // Bug, use-after-free: (ù{\codeincomments{flags}}ù) vector released its memory.
    {
        // ...
    }
}
\end{emcppslisting}
    

\subsubsection[Lack of interface restrictions]{Lack of interface restrictions}\label{lack-of-interface-restrictions}

Lack of any restrictions placed by \lstinline!auto! on the type that is
deduced might result in defects that could otherwise be detected at
compile time. Consider refactoring the \lstinline!getNetworkNodes! function
illustrated in \intraref{use-cases-auto}{improving-resilience-to-library-code-changes} to return \lstinline!std::deque<Node>! instead of
\lstinline!std::vector<Node>!:

\begin{emcppshiddenlisting}[emcppsbatch=e18]
#include <deque>     // (ù{\codeincomments{std::deque}}ù)
struct Node {};
void CLibraryProcessNodes(Node *, std::size_t);
\end{emcppshiddenlisting}
\begin{emcppslisting}[emcppsbatch=e18]
std::deque<Node> getNetworkNodes();  // Return type changed from (ù{\codeincomments{std::vector<Node>}}ù).
    // Return a sequence of nodes in the current network.
\end{emcppslisting}
    
\noindent While code that uses \lstinline!auto! to store the result returned by
\lstinline!getNetworkNodes! only to subsequently iterate over it with a
range-based \lstinline!for! wouldn't be affected, the behavior of code that
relied on the contiguous layout of elements in \lstinline!std::vector!
objects \emph{silently} becomes undefined:

\begin{emcppslisting}[emcppsbatch=e18]
void testUseContiguousMemory()
{
    auto nodes = getNetworkNodes();
    CLibraryProcessNodes(&nodes[0], nodes.size());  
        // exhibits UB after (ù{\codeincomments{std::vector}}ù)-to-(ù{\codeincomments{std::deque}}ù) change
}
\end{emcppslisting}
    
\noindent While specifying constraints on types deduced by
\lstinline!auto! with \lstinline!static_assert! is possible, doing so is often
cumbersome{\cprotect\footnote{C++20 introduced \emcppsgloss{concepts} ---
named type requirements --- as well as means to constrain
\lstinline!auto! with a specific concept, which can be used instead of
  \lstinline!static_assert! in such circumstances.}}:

\begin{emcppshiddenlisting}[emcppsbatch=e19]
#include <type_traits>
#include <vector>
struct Packet {};
struct PacketCache {
    std::vector<Packet> d_packets;
    const Packet* findFirstCorruptPacket() const;
};
\end{emcppshiddenlisting}
\begin{emcppslisting}[emcppsbatch=e19]
const Packet* PacketCache::findFirstCorruptPacket() const
{
    auto it = std::begin(this->d_packets);

    static_assert(
        std::is_base_of<
            std::random_access_iterator_tag,
            std::iterator_traits<decltype(it)>::iterator_category>::value,
        "'it' must satisfy the requirements of a random access iterator.");

    // ...

    return it == std::end(this->d_packets) ? nullptr : &*it;
}
\end{emcppslisting}
    

\subsubsection[Obscuration of important properties of fundamental types]{Obscuration of important properties of fundamental types}\label{obscuration-of-important-properties-of-fundamental-types}

Using \lstinline!auto! for variables of fundamental types might hide
important, context-sensitive considerations, such as overflow behavior
or a mix of signed and unsigned arithmetic. In the example below, the
\lstinline!lowercaseEncode! function will either work correctly or enter an
infinite loop depending on whether the type returned by
\lstinline!Encoder::encodedLengthFor! is signed.

\begin{emcppshiddenlisting}[emcppsbatch=e20]
#include <string>  // (ù{\codeincomments{std::string}}ù)
#include <cctype>  // (ù{\codeincomments{std::tolower}}ù)
struct Encoder {
    static int encodedLengthFor(const std::string &);
    static void encode(std::string::iterator, const std::string &);
};
\end{emcppshiddenlisting}
\begin{emcppslisting}[emcppsbatch=e20]
void lowercaseEncode(std::string* result, const std::string& input)
{
    auto encodedLength = Encoder::encodedLengthFor(input);

    result->resize(encodedLength);
    Encoder::encode(result->begin(), input);

    while (--encodedLength >= 0)  // infinite loop if (ù{\codeincomments{encodedLength}}ù) is unsigned
    {
        (*result)[encodedLength] = std::tolower((*result)[encodedLength]);
    }
}
\end{emcppslisting}
    

\subsubsection[Surprising deduction for list initialization]{Surprising deduction for list initialization}\label{surprising-deduction-for-list-initialization}

\lstinline!auto! type-deduction rules differ from those of function
templates if brace-enclosed initializer list are used. Function template
argument deduction will always fail, whereas, according to C++11 rules,
\lstinline!std::initializer_list! will be deduced for \lstinline!auto!.

\begin{emcppslisting}
auto example0 = 0; // copy initialization, deduced as (ù{\codeincomments{int}}ù)
auto example1(0);  // direct initialization, deduced as (ù{\codeincomments{int}}ù)
auto example2{0};  // list initialization, deduced as (ù{\codeincomments{std::initializer\_list<int>}}ù)

template <typename T> void func(T);
void testFunctionDeduction()
{
    func(0);    // (ù{\codeincomments{T}}ù) deduced as (ù{\codeincomments{int}}ù)
    func({0});  // Error
}
\end{emcppslisting}
    
\noindent This surprising behavior was, however, widely regarded as a mistake.\footnote{This erroneous behavior was formally rectified in C++17 with, e.g., \lstinline!auto!~\lstinline!i{0}!
deducing \lstinline!int!. Furthermore, mainstream compilers had applied
this deduction-rule change retroactively as early as GCC~5.1 and Clang~3.8, with the revised rule being applied even if \lstinline!std=c++11!
flag is explicitly supplied.}

Nonetheless, even with this retroactive fix, the effects of the
deduction rules when applied to braced initializer lists might be
puzzling. In particular, \lstinline!std::initializer_list! is deduced when
\emcppsgloss{copy initialization} is used instead of
\emcppsgloss{direct initialization}, and this requires including \lstinline!<initializer_list>!:

\begin{emcppslisting}
auto x1 = 1;                 // (ù{\codeincomments{int}}ù)
auto x2(1);                  //  "
auto x3{1};                  //  "

#include <initializer_list>  // (ù{\codeincomments{std::initializer\_list}}ù)
auto x4 = {1};               // OK, deduced as (ù{\codeincomments{std::initializer\_list<int>}}ù)

auto x5{1, 2};               // Error, direct-list-init requires exactly 1 element.
auto x6 = {1, 2};            // OK, deduced as (ù{\codeincomments{std::initializer\_list<int>}}ù)
\end{emcppslisting}
    

\subsubsection[Deducing built-in arrays is problematic]{Deducing built-in arrays is problematic}\label{deducing-built-in-arrays-is-problematic}

Deducing built-in array types using \lstinline!auto! presents multiple
challenges. First, declaring an array of \lstinline!auto! is ill formed:

\begin{emcppslisting}
auto arr1[]  = {1, 2};  // Error, array of (ù{\codeincomments{auto}}ù) is not allowed.
auto arr2[2] = {1, 2};  // Error, array of (ù{\codeincomments{auto}}ù) is not allowed.
\end{emcppslisting}
    
\noindent Second, if the array bound is not specified, either the program does not
compile or\linebreak[4]%%%%%%
 \lstinline!std::initializer_list! is deduced instead of a
built-in array:

\begin{emcppslisting}
#include <initializer_list>  // (ù{\codeincomments{std::initializer\_list}}ù)
auto arr3 = {1, 2};  // OK, deduced as (ù{\codeincomments{std::initializer\_list<int>}}ù)
auto arr4{1, 2};     // Error, direct-list-init requires exactly 1 element.
\end{emcppslisting}
    
\noindent Finally, attempting to circumvent this deficiency by using an alias
template (see \featureref{\locationa}{alias-declarations-and-alias-templates}) will result in code that
compiles but has undefined behavior:

\begin{emcppshiddenlisting}[emcppsbatch=e21]
#include <cstddef>  // (ù{\codeincomments{std::size\_t}}ù)
\end{emcppshiddenlisting}
\begin{emcppslisting}[emcppsbatch=e21]
template <typename TYPE, std::size_t SIZE>
using BuiltInArray = TYPE[SIZE];

auto arr5 = BuiltInArray<int, 2>{1, 2};
    // Error, taking the address of a temporary array
\end{emcppslisting}
    
\noindent Note that in this case such code also almost entirely defeats the
purpose of \lstinline!auto! since neither the array element's type nor the
array's bound is deduced.

With that said, using \lstinline!auto! to deduce
references to built-in arrays is straightforward:

\begin{emcppslisting}[emcppsbatch=e21]
int data[] = {1, 2};

      auto&  arr6 = data;                        //       (ù{\codeincomments{int (\&) [2]}}ù)
const auto&  arr7 = BuiltInArray<int, 2>{1, 2};  // (ù{\codeincomments{const int (\&) [2]}}ù)
      auto&& arr8 = BuiltInArray<int, 2>{1, 2};  //       (ù{\codeincomments{int (\&\&)[2]}}ù)
\end{emcppslisting}
    
\noindent Note that the \lstinline!arr7! and \lstinline!arr8! references in the code snippet immediately above extend
the lifetime of the temporary arrays that they bind to, so subscripting
them does not have the undefined behavior that subscripting
\lstinline!arr5! (in the previous code snippet) has.

\subsection[Annoyances]{Annoyances}\label{annoyances-auto}

\subsubsection[\lstinline!auto! not allowed for nonstatic data members]{{\SubsubsecCode auto} not allowed for nonstatic data members}\label{auto-not-allowed-for-non-static-data-members}

Despite C++11 allowing nonstatic data members to be initialized within
class definitions, \lstinline!auto! cannot be used to declare them:

\begin{emcppslisting}
class C
{
    auto d_i = 1;  // Error, nonstatic data member is declared with (ù{\codeincomments{auto}}ù).
};
\end{emcppslisting}
    

\subsubsection[Not all template argument deduction constructs are allowed for \lstinline!auto!]{Not all template argument deduction constructs are allowed for {\SubsubsecCode auto}}\label{not-all-template-argument-deduction-constructs-are-allowed-for-auto}

Despite \lstinline!auto! type deduction largely following the template
argument deduction rules, certain constructs that are allowed for
templates are not allowed for \lstinline!auto!. For example, when deducing
a pointer-to-data-member type, templates allow for deducing both the
data member type and the class type, whereas \lstinline!auto! can 
deduce only the former:

\begin{emcppslisting}
struct Node
{
    int   d_data;
    Node* d_next;
};

template <typename TYPE>
void deduceMemberTypeFn(TYPE Node::*);

void testDeduceMemberType()
{
                 deduceMemberTypeFn   (&Node::d_data);  // OK, (ù{\codeincomments{int Node::*}}ù)
    auto Node::* deduceMemberTypeVar = &Node::d_data;   // OK,  "     "
}

template <typename TYPE>
void deduceClassTypeFn(int TYPE::*);

void test1DeduceClassType()
{
                deduceClassTypeFn   (&Node::d_data);  // OK, (ù{\codeincomments{int Node::*}}ù)
    int auto::* deduceClassTypeVar = &Node::d_data;   // Error, not allowed
}

template <typename TYPE>
void deduceBothTypesFn(TYPE* TYPE::*);

void testDeduceBothTypes()
{
                  deduceBothTypesFn   (&Node::d_next);  // OK, (ù{\codeincomments{Node* Node::*}}ù)
    auto* auto::* deduceBothTypesVar = &Node::d_next;   // Error, not allowed
}
\end{emcppslisting}
    
\noindent Furthermore, deducing the parameter of a class template is also not
allowed:

\begin{emcppshiddenlisting}[emcppsbatch=e22]
#include <vector>     // (ù{\codeincomments{std::vector}}ù)
\end{emcppshiddenlisting}
\begin{emcppslisting}[emcppsbatch=e22]
std::vector<int> vectorOfInt;

template <typename TYPE>
void deduceVectorArgFn(const std::vector<TYPE>&);

void testDeduceVectorArg()
{
                      deduceVectorArgFn   (vectorOfInt); // OK, (ù{\codeincomments{TYPE}}ù) is (ù{\codeincomments{int}}ù)
    std::vector<auto> deduceVectorArgVar = vectorOfInt;  // Error, not allowed
}
\end{emcppslisting}
    
\noindent Instead, if \lstinline!auto! type deduction is desired in such cases,
\lstinline!auto! alone is suitable to deduce the type from the initializer:

\begin{emcppshiddenlisting}[emcppsbatch=e22]
struct Node {
int d_data;
Node *d_next;
};
\end{emcppshiddenlisting}
\begin{emcppslisting}[emcppsbatch=e22]
auto deduceClassTypeVar = &Node::d_data;  // OK, (ù{\codeincomments{int   Node::*}}ù)
auto deduceBothTypesVar = &Node::d_next;  // OK, (ù{\codeincomments{Node* Node::*}}ù)

auto deduceVectorArgVar = vectorOfInt;    // OK, (ù{\codeincomments{std::vector<int>}}ù)
\end{emcppslisting}
    

\subsection[See Also]{See Also}\label{see-also}

\begin{itemize}
\item{\seealsoref{trailing-function-return-types}{\seealsolocationa}explains how the \lstinline!auto! placeholder can be used to specify a function’s return type at the end of its signature.}
\item{\seealsoref{genericlambda}{\seealsolocationd} illustrates how the \lstinline!auto! placeholder can be used in the argument list of a lambda to make its function call operator a template.}
\item{\seealsoref{Function-Return-Type-Deduction}{\seealsolocationf}describes how the \lstinline!auto! placeholder can be used to deduce a function’s return type.}
\end{itemize}

\subsection[Further Reading]{Further Reading}

\begin{itemize}
\item{``Item 1: Understand template type deduction," pp. xx--yy; ``Item 2: Understand auto type deduction," pp. xx--yy; ``Item 5: Prefer auto to explicit type declarations," pp. xx--yy; and ``Item 6: Use the explicitly typed initializer idiom when auto deduces undesired types," pp. xx-yy, \cite{meyers15b}}
\end{itemize}
 
 


\newpage
 %\section[Braced Init]{Brace-Initialization Syntax: {\SecCode \{\}}}\label{bracedinit}
% 22 March 2021, cleaned up LaTeX, updated commands. 
% 23 March 2021, out to CE, Steven for tech review, external reviewers, and Josh for code check. 



\emcppsFeature{
    short={Braced Init},
    long={Braced-Initialization Syntax: {\SecCode \{\}}},
    toclong={Braced-Initialization Syntax: \lstinline!\{\}!},
}{bracedinit}
\setcounter{table}{0}
\setcounter{footnote}{0}
\setcounter{lstlisting}{0}
%\section[Braced Init]{Brace-Initialization Syntax: {\SecCode \{\}}}\label{bracedinit}
%\subsection[Braced-Initialization Syntax: \tt{\{ \}}]{Braced-Initialization Syntax: {\SubsecCode \{ \}}}\label{braced-initialization-syntax:-{-}}

Braced initialization, a generalization of C++03 initialization syntax,
was designed with the intention designed to be used safely and uniformly
in any initialization context.

\subsection[Description]{Description}\label{description}

\emcppsgloss[list initialization]{List initialization}, originally dubbed \emcppsgloss{uniform
initialization}, was conceived to enable a uniform syntax (having the
same meaning) to be used generically to initialize objects irrespective
of (1) the context in which the syntax is used or (2) the type of the
object being initialized. Braced-initialization syntax is the language
mechanism that --- in close collaboration with the C++ Standard
Library's \lstinline!std::initializer_list! template (see \featureref{\locationc}{initlist}) --- is used to implement
\emcppsgloss{list initialization} generally. As we will see, this design goal
was largely achieved albeit with some idiosyncrasies and rough edges.

\subsubsection[C++03 initialization syntax review]{C++03 initialization syntax review}\label{c++03-initialization-syntax-review}

Classic C++ affords several forms of initialization, each sporting its
own custom syntax, some of which is syntactically interchangeable yet
belying subtle differences. At the highest level, there are two dual
categories of initialization: (1) \emcppsgloss[copy direct]{copy/direct} (when you have
something from which to initialize) and (2) \emcppsgloss[default value]{default/value} (when
you don't).

The first dual category of syntactic/semantic initialization comprises
\emcppsgloss[copy initialization]{\emph{copy} initialization} and \emcppsgloss[direct initialization]{\emph{direct}
initialization}. \emph{Direct} initialization is produced when
initializing an object with one or more arguments within parentheses,
such as initializing a data member or base class in a constructor's
initializer list, or in a \lstinline!new! expression. \emph{Copy}
initialization happens when initializing from a value without using
parentheses, such as passing an argument to a function, or returning a
value from a function. Both forms may be used to initialize a variable:

\begin{emcppslisting}
#include <cassert>  // (ù{\codeincomments{assert}}ù) macro

int i = 23;  // copy initialization
int j(23);   // direct initialization

assert(i == j);
\end{emcppslisting}
    
\noindent In both cases above, we are initializing the variable (\lstinline!i! or
\lstinline!j!) with the literal value \lstinline!23!.

For scalar types, there is no observable difference between these two
dual forms of initialization in C++03, but, for user-defined types,
there are. First \emcppsgloss[direct initialization]{\emph{direct} initialization} \emph{considers} (as
part of the overload set) all valid user-defined conversion sequences,
whereas \emcppsgloss[copy initialization]{\emph{copy} initialization} excludes explicit conversion:

\begin{emcppslisting}
struct S
{
    explicit S(int);       //     explicit value constructor (from (ù{\codeincomments{int}}ù))
             S(double);    // non-explicit value constructor (from (ù{\codeincomments{double}}ù))
             S(const S&);  // non-explicit copy  constructor
};

S s1(1);     // direct init of (ù{\codeincomments{s1}}ù): calls (ù{\codeincomments{S(int)}}ù)   ; copy constructor is not called
S s2(1.0);   // direct init of (ù{\codeincomments{s2}}ù): calls (ù{\codeincomments{S(double)}}ù);   "      "        "   "     "

S s3 = 1;    // copy init of (ù{\codeincomments{s3}}ù):   calls (ù{\codeincomments{S(double)}}ù); copy constructor may be called
S s4 = 1.0;  // copy init of (ù{\codeincomments{s4}}ù):   calls (ù{\codeincomments{S(double)}}ù);   "      "         "    "   "
\end{emcppslisting}
    
\noindent What's more, \emcppsgloss[copy initialization]{\emph{copy} initialization} is defined as if a
temporary object is constructed; the compiler is permitted to elide this
temporary and, in practice, typically does. Note, however, that
\emcppsgloss[copy initialization]{\emph{copy} initialization} is not permitted unless there is an
accessible \emph{copy} (or \emph{move}{\cprotect\footnote{If the
\emph{move} constructor for a user-defined type is declared and not
accessible, copy initialization is ill formed; see \featureref{\locationc}{Rvalue-References} and \featureref{\locationa}{deleted-functions}.}}) constructor,
even if the temporary would have been elided.{\cprotect\footnote{In
C++17, direct materialization replaces \emcppsgloss[copy initialization]{\emph{copy}
initialization} in some contexts, thereby obviating the temporary
object construction and also the need for an accessible \emph{copy} or
  \emph{move} constructor.}} Note that function arguments and return
values are initialized using \emcppsgloss[copy initialization]{\emph{copy} initialization}.

Reference types are also initialized by copy and direct initialization,
binding the declared reference to an object (or function). For a
reference to a non\lstinline!const! qualified type, the referenced type
must match exactly. However, if binding a reference to a \lstinline!const!
qualified type, the compiler may copy initialize a temporary object of
the target type of the reference and bind the reference to that
temporary; in such cases, the lifetime of the temporary object is
extended to the end of the lifetime of the reference.

\begin{emcppslisting}
void ref_inits()
{
    int i = 0;          // OK, copy initialization of (ù{\codeincomments{int}}ù)
    int& x(i);          // OK, direct initialization of reference
    const long& y = x;  // OK, (ù{\codeincomments{y}}ù) binds to a temporary and extends lifetime
    long& z = x;        // Error, incompatible types
}
\end{emcppslisting}
    

\noindent The second dual category of syntactic/semantic initialization comprises
\emcppsgloss[default initialization]{\emph{default} initialization} and \emcppsgloss[value initialization]{\emph{value}
initialization}. Both default and value initialization pertain to
situations in which \emph{no} argument is supplied, and are
distinguished by the presence or absence of parentheses, where the
absence of parentheses indicates default initialization and the presence
indicates value initialization. Note that in simple contexts such as
declaring a variable, empty parentheses may also indicate a function
declaration instead (see \intraref{use-cases-bracedinit}{avoiding-the-most-vexing-parse}):

\begin{emcppslisting}
int i;          // default initialization
int j();        // Oops, function declaration
int k = int();  // value initialization

int *pd = new int;    // default initialization of dynamic (ù{\codeincomments{int}}ù) object
int *pv = new int();  // value        "         "     "      "    "
\end{emcppslisting}
    

\noindent For \emcppsgloss[scalar type]{scalar types}, \emph{default} initialization does not
actually initialize an object, and \emph{value} initialization will
initialize that object as if by the literal \lstinline!0!. Note that the
representation of this value is not necessarily all zero bits, as some
platforms use distinct trap values for the null pointer value for
pointers and for pointer-to-member objects.

For class types with an accessible \emcppsgloss[user provided]{user-provided} default
constructor, default initialization and value initialization behave
identically, calling the default constructor. If there is no accessible
default constructor, both forms produce a compilation error. For objects
of class types with an implicitly defined default constructor, each base
and member subobject will be default initialized or value initialized
according to the form of initialization indicated for the complete
object; if any of those initializations produces an error, then the
program is ill formed. Note that for a union with an implicitly defined
default constructor, the first member of the union will be value
initialized as the active member of that union when a union object is
value initialized.

\begin{emcppslisting}
struct B
{
    int i;
    B() : i() { }  // user-provided default constructor
};

struct C
{
    int i;
    C() { }  // user-provided default constructor
};

struct D : B { int j; };  // derived class with no user-provided constructors

int *pdi = new int;    // default initialization of dynamic (ù{\codeincomments{int}}ù) object, (ù{\codeincomments{*pdi}}ù) is uninitialized
int *pvi = new int();  // value   initialization of dynamic (ù{\codeincomments{int}}ù) object, (ù{\codeincomments{*pvi}}ù) is 0
B *pdb = new B;        // default initialization of dynamic (ù{\codeincomments{B}}ù) object, (ù{\codeincomments{b::i}}ù) is 0
B *pvb = new B();      // value   initialization of dynamic (ù{\codeincomments{B}}ù) object, (ù{\codeincomments{b::i}}ù) is 0
C *pda = new C;        // default initialization of dynamic (ù{\codeincomments{C}}ù) object, (ù{\codeincomments{a::i}}ù) is uninitialized
C *pva = new C();      // value   initialization of dynamic (ù{\codeincomments{C}}ù) object, (ù{\codeincomments{a::i}}ù) is uninitialized
D *pdc = new D;        // default initialization of dynamic (ù{\codeincomments{D}}ù) object, (ù{\codeincomments{c::j}}ù) is uninitialized
D *pvc = new D();      // value   initialization of dynamic (ù{\codeincomments{D}}ù) object, (ù{\codeincomments{c::j}}ù) is 0
\end{emcppslisting}
    

\noindent In the case of an object of type \lstinline!B!, both default and value
initialization will invoke the user-provided default constructor, which
initializes the subobject \lstinline!i! to \lstinline!0!. In the case of an
object of type \lstinline!C!, both default and value initialization will
invoke the user-provided default constructor, which does not initialize
the subobject \lstinline!i!. In the case of an object of type \lstinline!D!,
which has an implicitly defined default constructor, the initialization
of the subobject \lstinline!j! depends on whether the \lstinline!D! object is
initialized by default initialization or by value initialization.

Any attempt to read the value of an \emph{uninitialized} object will
result in \emcppsgloss{undefined behavior}. It is a diagnosable error to
default initialize a constant object that does not execute a
user-provided constructor to initialize each base and member. Note that
the top level object, like \lstinline!D! above, need not have a
user-provided constructor as long as all of its bases and members can
recursively apply this rule.

\begin{emcppslisting}
struct D2 : B { B j; };  // derived class with no user-provided constructors

const D  w;         // Error, (ù{\codeincomments{w.j}}ù) is not initialized.
const D  x = D();   // OK, (ù{\codeincomments{x}}ù) is value initialized.
const D2 y;         // OK, (ù{\codeincomments{y}}ù) is default initialized; sub-objects invoke default ctor.
const D2 z = D2();  // OK, (ù{\codeincomments{z}}ù) is value initialized.
\end{emcppslisting}
    

\noindent Objects of static storage duration at file, namespace, or function scope
(see \featureref{\locationa}{thread-safe-function-static-variables}) are
\emcppsgloss[zero initialized]{\emph{zero} initialized} before any other initialization takes
place. Note that \emcppsgloss[default initialized]{\emph{default} initialized} static storage
duration pointer objects are \emcppsgloss[zero initialized]{\emph{zero} initialized} to have a
\emph{null} address value (see \featureref{\locationa}{null-pointer-literal-(nullptr)}), even if that
representation on the host platform is not numerically zero:

\begin{emcppslisting}
struct A
{
    int i;
};

struct B
{
    int i;
    B() i(42) { }
};

A globalA;
    // Zero initialization also zero initializes (ù{\codeincomments{globalA::i}}ù).
    // Default initialization provides no further initializations.

B globalB;
    // zero initialization initializes globalB.i.
    // After that, default constructor is invoked.

int globalI;
    // Zero initialization initializes globalI.
\end{emcppslisting}
    

\noindent Note the implication that default initialization for a static storage
duration object will always initialize an object ready for use, either
calling the default constructor of a type with a user-provided default
constructor or zero-initializing scalars. 

\begin{table}[h!]
\begin{center}
\begin{threeparttable}
\caption{Helpful summary of C++03 rules}\label{table-bracedinit-cpp3rules} \vspace{1.5ex}
{\small \begin{tabular}{c|c|c}\thickhline
\rowcolor[gray]{.9}    {\sffamily\bfseries Initialization Type} & {\sffamily\bfseries No Arguments} & {\sffamily\bfseries >= 1 Arguments}\\\hline
\emcppsgloss[with parentheses]{With Parentheses} & \emph{value} & \emph{direct} \\ 
 & \lstinline!int!~\lstinline!i!~\lstinline!=!~\lstinline!int();! & \lstinline!int!~\lstinline!i(23);! \\ \hline
\emcppsgloss[without parentheses]{Without Parentheses} & \emph{default} & \emph{copy} \\  
&  \lstinline!int!~\lstinline!i;! & \lstinline!int!~\lstinline!i!~\lstinline!=!~\lstinline!23;! \\ \thickhline
\end{tabular}
}
\end{threeparttable}
    \end{center}
\end{table}


\subsubsection[C++03 aggregate initialization]{C++03 aggregate initialization}\label{c++03-aggregate-initialization}

Aggregates are a special kind of object in C++03 that generally do not
use constructors but follow a different set of rules for initialization,
typically denoted by braces. There are two varieties of aggregates: (1)
arrays and (2) user-defined class types that have no nonpublic data
members that are not static, no base classes, no user-declared
constructors, and no virtual functions. Aggregates are very similar to a
classic C \lstinline!struct!, potentially with additional,
non-\lstinline!virtual! member functions. Note that members of an aggregate
are not themselves required to be an aggregate.

\begin{emcppslisting}
int a[5];           // Arrays are aggregates.

struct A
{
    int         i;  // public data member
    std::string s;  // (ù{\codeincomments{A}}ù) is an aggregate even though (ù{\codeincomments{std::string}}ù) is not.

private:
    static int j;   // Private data member is static.
    void f();       // Member functions are OK, even if private.
};
\end{emcppslisting}
    

\noindent A quick note on terminology: Strictly speaking, arrays comprise
\emph{elements} and classes comprise \emph{members}, but for ease of
exposition in this text, we refer to both as \emph{members}.

When an aggregate is copied by either direct or copy initialization,
rather than calling the copy constructors, the corresponding members
(elements for an array) of each aggregate are copied using direct
initialization, which corresponds to the behavior of an implicitly
defined copy constructor for a class. Note that this process may be
applied recursively, if members are aggregates themselves. Further note
that in most cases, arrays do not copy because the argument supplied to
the copy operation will undergo \emcppsgloss[array to pointer decay]{array-to-pointer decay} and so
will no longer be an appropriate type to initialize from. However,
arrays as data members of classes follow the rules for aggregate
initialization and so will copy array data members. This array-copy
behavior is one of the motivations for the addition of the
\lstinline!std::array! template in C++11.

When an aggregate is \emph{default} initialized, each of its
members/elements is \emph{default} initialized. When an aggregate is
\emph{value} initialized, each of its members/elements is \emph{value}
initialized. This follows the usual rules for an implicitly defined
constructor for a class type and defines the corresponding behavior for
array initialization.

\begin{emcppslisting}
int *pid = new int[N];   // default initialization of dynamic array object and its elements
int *piv = new int[N](); // value        "         "     "      "     "     "   "     "

A *pd = new A;           // default initialization of dynamic (ù{\codeincomments{A}}ù) object and its members
A *pv = new A();         // value        "         "     "     "    "     "   "     "
\end{emcppslisting}
    

\noindent Otherwise, an aggregate must be \emph{aggregate} initialized by a braced
list in the form\linebreak[4] 
\lstinline!=!~\lstinline!{!~\lstinline!list-of-values!~\lstinline!};!, where members
of the aggregate will be initialized by \emph{copy} initialization from
the corresponding value in the list of values; if the aggregate has more
members than are provided by the list, the remaining members are
\emph{value} initialized; it is an error to provide more values in the
list than there are members in the aggregate. Note that, because a union
has only one active member, a union will be initialized by no more than
a single value from the list; this becomes relevant for unions as data
members of an aggregate initialized by \emph{brace elision}:

\begin{emcppslisting}
union U
{
    int i;
    const char* s;
};

U x = {    };  // OK, value initializes (ù{\codeincomments{x.i}}ù) = 0
U y = { 1  };  // OK, initializes (ù{\codeincomments{x.i}}ù) = 1
U z = { "" };  // Error, cannot aggregate initialize (ù{\codeincomments{z.s}}ù)
\end{emcppslisting}
    

\noindent Let's review the various ways in which we might attempt to initialize an
object of aggregate type \lstinline!A! in the body of a function,
\lstinline!test!, i.e., defined at function scope:

\begin{emcppslisting}
struct A2 { int i };  // aggregate with a single data member

void test()
{
           A2  a1;             // default init: (ù{\codeincomments{i}}ù) is not initialized!
    const  A2& a2 = A();       // value init followed by copy init: (ù{\codeincomments{i}}ù) is 0.
           A2  a3 = A();       // value init followed by copy init: (ù{\codeincomments{i}}ù) is 0.
           A2  a4();           // Oops, function declaration!
           A2  a5 = { 5 };     // aggregate initialization employing copy init
           A2  a6 = { };       //    "            "            "     value  "
           A2  a7 = { 5, 6 };  // Error, too many initializers for aggregate (ù{\codeincomments{A}}ù)
    static A2  a8;             // default init after (ù{\codeincomments{i}}ù) is zero initialized.
}
\end{emcppslisting}
    

\noindent In the sample code above:
\begin{itemize}
\item{\lstinline!a1!: \lstinline!a1! is \emcppsgloss[default initialized]{\emph{default} initialized}, which means that each data member within the aggregate is itself independently \emcppsgloss[default initialized]{\emph{default} initialized}. For scalar types, such as an \lstinline!int!, the effect of default initialization at function scope is a no-op — i.e., \lstinline!a1.i! is not initialized. Any attempt to access the contents of \lstinline!a1.i! will result in \emcppsgloss{undefined behavior}.}
\item{\lstinline!a2! and \lstinline!a3!: In the cases of both \lstinline!a2! and \lstinline!a3!, a temporary of type \lstinline!A! is first \emcppsgloss[value initialized]{\emph{value} initialized} and then that temporary is used to \emcppsgloss[copy initialized]{\emph{copy} initialize} the named variable: Both \lstinline!a2.i! and \lstinline!a3.i! are initialized to the value \lstinline!0!.}
\item{\lstinline!a4!: Notice that we are not able to create a \emcppsgloss[value initialized]{\emph{value} initialized} local variable \lstinline!a4! by applying parentheses since that would interpreted as declaring a function taking no arguments and returning an object of type \lstinline!A! by value; see \intraref{use-cases-bracedinit}{} [\textbf{AUs: there is no section with this name ``Avoiding accidentally declaring a function taking no arguments"}] 
and \intraref{use-cases-bracedinit}{avoiding-the-most-vexing-parse}.}
\item{\lstinline!a5!, \lstinline!a6!, and \lstinline!a7!: C++03 supports \emcppsgloss{aggregate initialization} using braced syntax as illustrated by \lstinline!a5!, \lstinline!a6!, and \lstinline!a7! in the code snippet above. The local variable \lstinline!a5! is \emcppsgloss[copy initialization]{\emph{copy} initialized} such that \lstinline!a5.i! has the \emph{user supplied} value \lstinline!5! whereas \lstinline!a6! is \emcppsgloss[value initialized]{\emph{value} initialized} since there are no supplied initializers; hence, \lstinline!a6.i! is initialized to 0. Attempting to pass \lstinline!a7! two values to initialize a single data member results in a compile-time error. Note that had class \lstinline!A! held a second data member, the line initializing \lstinline!a5! would have resulted in \emcppsgloss[copy initialization]{\emph{copy} initialization} of the first and \emcppsgloss[value initialized]{\emph{value} initialization} of the second.}
\item{\lstinline!a8! has static storage duration therefore it is first \emph{zero} initialized (\lstinline!a8.i! is 0) then it is \emph{default} initialized, which is a no-op for the same reasons that \lstinline!a1! is not initialized at all.}
\end{itemize}

Finally, note that a scalar can be thought of as though it were an array
of a single element (though note that scalars never suffer
\emph{array-to-pointer decay}); in fact, if we were to take the address
of any scalar and add \lstinline!1! to it, the new pointer value would
represent the one-past-the-end iterator for that scalar's implied array
(of length \lstinline!1!). Similarly, scalars can be initialized using
aggregate initialization, just as if they were single-element arrays,
where the braced list for a scalar may contain zero or one elements.

\begin{emcppslisting}
int    i = { };       // OK, (ù{\codeincomments{i}}ù) is 0.
int    j = { 1 };     // OK, (ù{\codeincomments{i}}ù) is 1.
double k = { 3.14 };  // OK, (ù{\codeincomments{k}}ù) is 3.14.
\end{emcppslisting}
    

\subsubsection[Braced initialization in C++11]{Braced initialization in C++11}\label{braced-initialization-in-c++11}

Everything we've discussed so far, including braced initialization of
aggregates, is well defined in C++03. This same braced initialization
syntax --- modified slightly so as to preclude narrowing conversions
(see the next section) --- is extended in C++11 to work consistently and
uniformly in many new situations. This enhanced braced initialization
syntax is designed to better support the two dual initialization
categories discussed in \textbf{C++03 initialization syntax review}%%%%%%
 above
as well as entirely new capabilities including language-level support
for lists of initial values implemented using the C++ Standard Library's
\lstinline!std::initializer_list! class template. As the opportunity arose
by touching the rules for initialization, a few more potential errors,
such as narrowing conversions, become diagnosable by the compiler when
using braced initialization syntax.

\subsubsection[C++11 restrictions on narrowing conversions]{C++11 restrictions on narrowing conversions}\label{c++11-restrictions-on-narrowing-conversions}

\emcppsgloss[narrowing conversion]{Narrowing conversions} (a.k.a. \emcppsgloss[lossy conversion]{lossy conversions}) are a
notorious source of runtime errors. One of the important properties of
list initializations implemented using the C++11 braced-initialization
syntax is that error-prone narrowing conversions are no longer
permitted. Consider, for example, an \lstinline!int! array, \lstinline!a!,
initialized with various built-in (compile-time constant) \emph{literal}
values:

\begin{emcppslisting}
int ai[] =  //      C++03    C++11
{
    5,      // (0)   OK       OK
    5.0,    // (1)   OK      Error, (ù{\codeincomments{double}}ù) to (ù{\codeincomments{int}}ù) conversion is not permitted.
    5.5,    // (2)   OK      Error, (ù{\codeincomments{double}}ù) to (ù{\codeincomments{int}}ù) conversion is not permitted.
    "5",    // (3)  Error    Error, no (ù{\codeincomments{const char*}}ù) to (ù{\codeincomments{int}}ù) conversion exists.
};
\end{emcppslisting}
    

\noindent In C++03, floating-point literals would be coerced to fit within an
integer even if the conversion was known to be lossy --- e.g., element
(2) in the code snippet above. By contrast, C++11 disallows \emph{any}
such implicit conversions in braced initializations even when the
conversion is known \emph{not} to be lossy --- e.g., element
\lstinline!ai[1]! above.

\emcppsgloss[narrowing conversion]{Narrowing conversions} within the integral and floating-point
type families, respectively, are generally disallowed except where it
can be verified at compile-time that overflow does not occur and, in the
case of integers and (classic) \lstinline!enum!s , the initializer value
can be represented exactly{\cprotect\footnote{As of C++20, implicit
conversion from either a pointer or pointer-to-member type to
  \lstinline!bool! is generally supported in braced initializations.}}:

\begin{emcppslisting}
const unsigned long ulc = 1;  // compile-time integral constant: (ù{\codeincomments{1UL}}ù)

short as[] =  //       C++03     C++11                        Stored Value
{
     32767,   // (0)     OK        OK                        (ù{\codeincomments{as[0] == 32767}}ù)
     32768,   // (1)     OK      Error, overflow
    -32768,   // (2)     OK        OK                        (ù{\codeincomments{as[0] == -32768}}ù)
    -32769,   // (3)   Warning?  Error, underflow
    1UL,      // (4)     OK        OK                        (ù{\codeincomments{as[0] == 1}}ù)
    ulc,      // (5)     OK        OK                        (ù{\codeincomments{as[0] == 1}}ù)
    1.0       // (6)     OK      Error, narrowing
};
\end{emcppslisting}
    

\noindent Notice that both \emph{overflow} (1) and \emph{underflow} (3) are
rejected for integral values in C++11, whereas neither is ill formed in
C++03. An integral literal (4) (or an \emcppsgloss{integral constant} (5)) of
a wider type (e.g., \lstinline!unsigned!~\lstinline!long!) can be used to
initialize a smaller one (e.g., \lstinline!signed!~\lstinline!short!) provided
that the value can be represented exactly; however, even a floating
point literal that can be \emph{represented exactly} (6) is nonetheless
rejected in C++11 when used to initialize any integral scalar.

Floating-point initializers, on the other hand, need not be represented
precisely so long as overflow does not occur; if, however, what is being
initialized is a floating-point scalar and the initializer is integral,
then the value must be represented exactly:

\begin{emcppslisting}
float af[] =     //       C++03   C++11                  Stored Value
{
     3L,        // (0)     OK     OK                   (ù{\codeincomments{af[0] == 3}}ù)
     16777216,  // (1)     OK     OK                   (ù{\codeincomments{af[1] == 1<<24}}ù)
     16777217,  // (2)     OK     Error, lossy
     0.75,      // (3)     OK     OK                   (ù{\codeincomments{af[1] == 0.75}}ù)
     2.4,       // (4)     OK     OK, but lossy        (ù{\codeincomments{af[2] != 2.4}}ù)
     0.4,       // (5)     OK     OK, but lossy        (ù{\codeincomments{af[3] != 0.4}}ù)
     1e-39,     // (6)     OK     OK, but underflow    (ù{\codeincomments{af[4] != 1e-39}}ù)
     1e+39,     // (7)     OK     Error, overflow
};
\end{emcppslisting}
    

\noindent In the example above elements (0) -- (2) represent initialization from
an integral type (\lstinline!int!), which requires that the initialized
value be represented exactly. Elements (3) -- (7) are instead
initialized from a floating-point type (\lstinline!double!) and therefore
are restricted only from overflow.

When an initializer is \emph{not} a \emcppsgloss{constant expression}, braced
initialization precludes any possibility of such \emph{narrowing}
initializations at run time --- e.g., initializing a \lstinline!float! with
a double or a \lstinline!long!~\lstinline!double!, a \lstinline!double! with a
\lstinline!long!~\lstinline!double!, or \emph{any} floating-point type with an
integral one. By the same token, an integral type (e.g., \lstinline!int!)
is not permitted to be initialized by non-\emcppsgloss{constant expression}
integer value of any other potentially larger integral type (e.g.,
\lstinline!long!) --- even if the number of bits in the representation for
the two types on the current platform is the same. Finally,
non-\emcppsgloss{constant expression} of integral type (e.g., \lstinline!short!)
cannot be used to initialize an unsigned version of the same type (e.g.,
\lstinline!unsigned!~\lstinline!short!) and vice versa.

To illustrate the constraints imposed on
non-\emcppsgloss[constant expression]{constant-expressions} described above, consider a simple
aggregate class, \lstinline!S!, comprising an \lstinline!int!, \lstinline!i!, and
a \lstinline!double!, \lstinline!d!:

\begin{emcppslisting}
struct S  // aggregated class
{
    int    i;  // *integral* scalar type
    double j;  // *floating-point* scalar type
};
\end{emcppslisting}
    

\noindent A function, \lstinline!test!, declaring a variety of arithmetic parameter
types illustrates restrictions imposed by C++11 braced initialization on
narrowing initializations that were well formed in C++03:

\begin{emcppslisting}
void test(short s, int i, long j, unsigned u, float f, double d, long double e)
{                     //     C++03  C++11
    S s0 = { i, d };  // (0)  OK     OK
    S s1 = { s, f };  // (1)  OK     OK
    S s2 = { u, d };  // (2)  OK    Error, (ù{\codeincomments{u}}ù) causes narrowing.
    S s3 = { i, e };  // (3)  OK    Error, (ù{\codeincomments{e}}ù) causes narrowing.
    S s4 = { f, d };  // (4)  OK    Error, (ù{\codeincomments{f}}ù) causes narrowing.
    S s5 = { i, s };  // (5)  OK    Error, (ù{\codeincomments{s}}ù) causes narrowing (theoretically).
};
\end{emcppslisting}
    

\noindent In the \lstinline!test! function above, lines (0) and (1) are OK because
there is no possibility of narrowing on any conforming platform unlike
lines (2) through (5) --- despite the fact that, in practice, it is more
than likely that a \lstinline!double! will be able to represent exactly
every value representable by a \lstinline!short!~\lstinline!int!. Note that,
just as with the array example above, when the initializing value is a
\emcppsgloss{constant expression}, it is sufficient that that value be
representable exactly in the target type and produce the original value
when converted back.

\subsubsection[C++11 Aggregate initialization]{C++11 Aggregate initialization}\label{c++11-aggregate-initialization}

Aggregate initialization in C++11, including initialization of arrays,
is subject to the rules prohibiting narrowing conversions.

\begin{emcppslisting}
int  i   = { 1 };        // OK
long j   = { 2 };        // OK

int  a[] = { 0, 1, 2 };  // OK
int  b[] = { 0, i, j };  // Error, cannot narrow (ù{\codeincomments{j}}ù) from (ù{\codeincomments{long}}ù) to (ù{\codeincomments{int}}ù)

struct S { int a; };
S s1 = { 0  };  // OK
S s2 = { i  };  // OK
S s3 = { 0L };  // OK, (ù{\codeincomments{0L}}ù) is an integer constant expression.
S s4 = { j  };  // Error, narrowing
\end{emcppslisting}
    

\noindent In addition, the rules for \emcppsgloss{value initialization} now state that
members without a specific initializer value in the braced list are
``as-if'' \emcppsgloss{copy initialized} from \lstinline!{}!
{\cprotect\footnote{From C++ 14 onwards, if the member doesn't have an
initializer value, but has a default member initializer, it is
initialized from the default member initializer (see \featureref{\locationc}{aggregate-member-initialization-relaxation}).}}. This will result in an error when
initializing a member that has an \lstinline!explicit! default constructor
according to the new \emcppsgloss[copy list initialization]{copy \emph{list} initialization} rules in
the next section, which give a meaning for explicit constructors. Note
that if the member is of reference type and no initializer is provided,
the initialization is ill formed.

Regardless of whether the aggregate itself is initialized using a copy
initialization or direct initialization, the members of the aggregate
will be copy initialized from the corresponding initializer.

\begin{emcppslisting}
struct E { };                         // empty type
struct AE { int x; E y; E z; };       // aggregate comprising several empty objects
struct S { explicit S(int = 0) {} };  // class with (ù{\codeincomments{explicit}}ù) default constructor
struct AS{ int x; S y; S z; };        // aggregate comprising several (ù{\codeincomments{S}}ù) objects

AE aed;
AE ae0 = {};               // OK
AE ae1 = { 0 };            // OK
AE ae2 = { 0, {} };        // OK
AE ae3 = { 0, {}, {} };    //

AS asd;                    //
AS as0 = {};               // OK in 03; error in 11 calling explicit ctor for (ù{\codeincomments{S}}ù)
AS as1 = { 0 };            // OK in 03; error in 11 calling explicit ctor for (ù{\codeincomments{S}}ù)
AS as2 = { 0, S() };       // OK in 03; error in 11 calling explicit ctor for (ù{\codeincomments{S}}ù)
AS as3 = { 0, S(), S() };  // OK, all aggregate members have an initializer.
\end{emcppslisting}
    

\noindent To better support generalizing the syntax of brace initialization in a
style similar to aggregate initialization, an aggregate can make a copy
of itself through \emph{aggregate} initialization in C++11 as well as
through \emph{direct} initialization per C++03:

\begin{emcppslisting}
S x{};      // OK, value initialization
S y = {x};  // OK in C++11; copy initialization via aggregate initialization syntax
\end{emcppslisting}
    

\noindent Otherwise, initialization of aggregates in C++11 is exactly the same
where it would have a meaning in C++03 and is correspondingly extended
into new places where braced initialization is permitted, as documented
in the following subsections.

\subsubsection[Copy list initialization]{Copy list initialization}\label{copy-list-initialization}

For C++03, only aggregates and scalars could be initialized via
braced-initialization syntax:

\begin{emcppslisting}
Type var = { /*...*/ };  // C++03-style aggregate (only) initialization
\end{emcppslisting}
    

\noindent The first part of generalizing braced initialization syntax for C++11 is
to allow the same syntactic form used to initialize aggregates to be
used for \emph{all} user-defined types. This extended form of braced
initialization --- known as \emcppsgloss[copy list initialization]{copy \emph{list} initialization} ---
follows the rules of \emcppsgloss{copy initialization}:

\begin{emcppslisting}
Class var1 = val;      // (C++03) copy initialization
Class var2 = { val };  // (C++11) copy list initialization
\end{emcppslisting}
    

\noindent For a non-aggregate class type, C++11 allows the use of a braced list
provided that its sequence of values serves as a suitable argument to a
non-\lstinline!explicit! constructor of the class being initialized.
Importantly, this use of \emcppsgloss[copy list initialization]{\emph{copy} list initialization}
provides meaning to \lstinline!explicit! constructors when taking other
than a single argument. For example, consider a
\lstinline!struct!~\lstinline!S! having constructors with 0-3 parameters, only
that last of which is \lstinline!explicit!:

\begin{emcppslisting}
struct S
{
             S();                          // default cor
             S(int};                       // 1-value ctor
             S(int, const char*);          // 2-value ctor
    explicit S(int, const char*, double);  // 3-value ctor
};
\end{emcppslisting}
    

\noindent We can use \emcppsgloss[copy list initialization]{\emph{copy} list initialization} only if the selected
constructor is \emph{not} declared to be \lstinline!explicit!, e.g.,
\lstinline!s0!, \lstinline!s1!, and \lstinline!s2! but not \lstinline!s3!:

\begin{emcppslisting}
S s0 = { };                 // OK, copy list initialization
S s1 = { 1 };               // OK, copy list initialization
S s2 = { 1, "two" };        // OK, copy list initialization
S s3 = { 1, "two", 3.14 };  // Error, constructor is explicit
\end{emcppslisting}
    

\noindent Had we instead declared our default constructor or any of the others to
be \lstinline!explicit!, the corresponding \emph{copy} (or
\emph{copy}-\emph{list}) initialization above would have failed too.

Another important difference between C++11 \emcppsgloss[copy list initialization]{copy \emph{list}
initialization} and C++03 \emcppsgloss{copy initialization} is that the
braced-list syntax considers all constructors, including those that are
declared to be \lstinline!explicit!. Consider a \lstinline!struct!~\lstinline!X!
having two overloaded single-argument constructors, i.e., (1) one taking
an \lstinline!int! and (2) the other a member template taking a single
(deduced) type, \lstinline!T!, by value:

\begin{emcppslisting}
struct Q  // class containing both (ù{\codeincomments{explicit}}ù) and implicit constructor overloads
{
    explicit Q(int);          // (1) value constructor taking a (ù{\codeincomments{int}}ù)
    template <class T> Q(T);  // (2) value constructor taking a (ù{\codeincomments{T}}ù)
};
\end{emcppslisting}
    

\noindent Employing \emcppsgloss[direct initialization]{\emph{direct} initialization} (e.g., \lstinline!x0! in
the code snippet below) selects the most appropriate constructor,
regardless of whether it is declared to be \lstinline!explicit!, and
successfully uses that one; employing \emcppsgloss[copy initialization]{\emph{copy}
initialization} (e.g., \lstinline!x1!) drops explicit constructors from the
overload set before determining a best match; and employing \emcppsgloss[copy list initialization]{copy
\emph{list} initialization} (e.g., \lstinline!x2!) again includes all
constructors in the overloads set but is \emcppsgloss{ill formed} if the
selected constructor is \lstinline!explicit!:

\begin{emcppslisting}
Q x0(0);     // OK, direct initialization calls (ù{\codeincomments{Q(int)}}ù).
Q x1 = 1;    // OK, copy initialization calls (ù{\codeincomments{Q(T)}}ù).
Q x2 = {2};  // Error, copy list initialization selects but cannot call (ù{\codeincomments{Q(int)}}ù).
Q x3{3};     // Same idea as (ù{\codeincomments{x0}}ù); direct list initialization calls (ù{\codeincomments{Q(int)}}ù).
\end{emcppslisting}
    

\noindent In other words, the presence of the \lstinline!=! coupled with the braced
notation (e.g., \lstinline!x2! in the code example above) forces the
compiler to choose the constructor \emph{as if} it were direct
initialization (e.g., \lstinline!x0!) but then forces a compilation failure
if the selected constructor turns out to be \lstinline!explicit!. This
``consider-but-fail-if-selected'' behavior of \emcppsgloss{copy list
initialization} is analogous to that of functions declared using
\lstinline!=!~\lstinline!delete!; see \featureref{\locationa}{deleted-functions}. Using
braces but omitting the \lstinline!=! (e.g., \lstinline!x3!) puts us back in
the realm of \emph{direct} rather than \emph{copy} initialization; see
\intrarefsimple{direct-list-initialization}.

When initializing references, \emcppsgloss{copy list initialization} (braced
syntax) behaves similarly to \emcppsgloss{copy initialization} (no braces)
with respect to the generation of temporaries. For example, when using a
braced list to initialize an \romeovalue{lvalue} reference --- e.g.,
\lstinline!int&!~\lstinline!ri! or \lstinline!const!~\lstinline!int&!~\lstinline!cri!
in the code example below --- to a scalar of a type that exactly matches
it (e.g., \lstinline!int!~\lstinline!i!), no temporary is created (just as it
would not have been without the braces); otherwise, a temporary will be
created, provided that a viable conversion exists and is not narrowing:

\begin{emcppslisting}
void test()
{
    int  i = 2;               assert(i   == 2);
    int& ri = { i };          assert(ri  == 2);  // OK, no temporary created
    ri = 3;                   assert(i   == 3);  // original is affected

    const int& cri = { i };   assert(cri == 3);  // OK, no temporary created
    ri = 4;                   assert(cri == 4);  // other reference is affected

    short s = 5;              assert(s   == 5);
    const int& crs = { s };   assert(crs == 5);  // OK, temporary is created
    s = 6;                    assert(crs == 5);  // note temporary is unchanged

    long j = 7;               assert(j   == 7);
    const int& crj = { j };   // Error, narrowing conversion from (ù{\codeincomments{long}}ù) to (ù{\codeincomments{int}}ù)
}
\end{emcppslisting}
    

\noindent As evidenced by the C-style asserts above, no temporary is created when
initializing either \lstinline!ri! or \lstinline!cri! since modifying the
reference affects the underlying variable, and vice versa. The C++ type
of \lstinline!crs!, on the other hand, does \emph{not} match exactly that
of the type to which it is bound, a temporary \emph{is} created and
hence changing the underlying object does \emph{not} affect its
referenced value. Lastly, unlike \lstinline!s! (of type \lstinline!short!),
attempting to initialize a \lstinline!const! \romeovalue{lvalue} reference of
type \lstinline!int!, \lstinline!crj!, with \lstinline!j! (of type \lstinline!long!)
is a narrowing conversion and thus ill formed.

Another consideration involves the standard
\lstinline!typedef!~\lstinline!std::size_t! found in the standard header
\lstinline!<cstddef>!, which must have sufficient bits to represent the
unsigned difference between any two pointers (into contiguous memory)
and is typically, but not necessarily, an alias for an
\lstinline!unsigned!~\lstinline!long!:

\begin{emcppslisting}
std::size_t k = 8;  // alias to implementation-defined type, say, (ù{\codeincomments{unsigned long}}ù)

unsigned int&       ra1 = { k };  // Note: Only one of these three lines will
unsigned long&      ra2 = { k };  // compile on any given platform; the other two
unsigned long long& ra3 = { k };  // will necessarily be ill formed and not compile.
\end{emcppslisting}
    

\noindent Historically, a \lstinline!long! has always been of sufficient size to
\emcppsgloss{pun} a pointer value (BAD IDEA) yet, back in the day when
\lstinline!int! and \lstinline!long! were the same number of bytes on a given
platform, \lstinline!size_t! was often an alias to an
\lstinline!unsigned!~\lstinline!int! rather than an
\lstinline!unsigned!~\lstinline!long!. Moving forward, one might expect
\lstinline!size_t! to be an alias for an \lstinline!unsigned!~\lstinline!long!
on a 64-bit platform, but there is no such assurance in the C++ Standard
and \lstinline!unsigned!~\lstinline!long!~\lstinline!long! (see
\featureref{\locationa}{long-long}) is another viable option on
any standards-conforming platform.

Finally, \lstinline!const! \romeovalue{lvalue} references to scalars and
aggregates initialized via braced lists of literal values follow the
rules of aggregates (see \intrarefsimple{c++11-aggregate-initialization}); a temporary is materialized having the indicated value and bound
permanently to the reference with its lifetime extended coterminously:

\begin{emcppslisting}
const int& i0 = { };    // OK, (ù{\emphincomments{materialized}}ù) temporary is (ù{\emphincomments{value}}ù) initialized.
const int& i1 = { 5 };  // OK,       "           "       " (ù{\emphincomments{copy}}ù)      "
\end{emcppslisting}
    

\noindent In the example above, a temporary is \emph{value} initialized (to 0) and
bound to \lstinline!i0!; another temporary is then \emph{copy} initialized
(to 5) and bound to \lstinline!i1!.

Non-modifiable references to \emph{aggregate UDT}s exploit the
generalization of \emph{copy} and \emph{direct} list initialization.
Consider an aggregate \lstinline!A! that comprises three \lstinline!int! data
members, \lstinline!i!, \lstinline!j!, and \lstinline!k!:

\begin{emcppslisting}
struct A          // (ù{\codeincomments{struct A}}ù) is an aggregate data type.
{
    int i, j, k;  // This (ù{\codeincomments{struct}}ù) contains three data members of type (ù{\codeincomments{int}}ù).
};
\end{emcppslisting}
    

\noindent We can now use braced initialization to materialize a temporary object
of aggregate type \lstinline!A! using aggregate initialization.

\begin{emcppslisting}
const A& s0 = { };           // (ù{\codeincomments{i}}ù), (ù{\codeincomments{j}}ù), and (ù{\codeincomments{k}}ù) are  (ù{\emphincomments{value}}ù) initialized.
const A& s1 = { 1 };         // (ù{\codeincomments{i}}ù) (ù{\emphincomments{copy}}ù) and (ù{\codeincomments{j}}ù) and (ù{\codeincomments{k}}ù) are (ù{\emphincomments{value}}ù) initialized.
const A& s2 = { 1, 2 };      // (ù{\codeincomments{i}}ù) and (ù{\codeincomments{j}}ù) (ù{\emphincomments{copy}}ù) and (ù{\codeincomments{k}}ù) are (ù{\emphincomments{value}}ù) initialized.
const A& s3 = { 1, 2, 3 };   // (ù{\codeincomments{i}}ù),(ù{\codeincomments{j}}ù), and (ù{\codeincomments{k}}ù) are (ù{\emphincomments{copy}}ù) initialized.
\end{emcppslisting}
    

In the example above, each of the references, \lstinline!s0! thru
\lstinline!s3!, is initialized to a temporary \lstinline!struct! of type
\lstinline!A! holding the respective aggregate value \lstinline!{0,0,0}!,
\lstinline!{1,0,0}!, \lstinline!{1,2,0}!, and \lstinline!{1,2,3}!.

\subsubsection[Direct list initialization]{Direct list initialization}\label{direct-list-initialization}

Of the two dual forms of initialization, \emph{direct} versus
\emph{copy}, \emph{direct} initialization is the stronger since it
enables use of all accessible constructors, i.e., including those
declared to be \lstinline!explicit!. The next step in generalizing the use
of braced initialization is to allow use of a braced list without the
intervening \lstinline!=! character between the variable and the opening
brace to denote \emcppsgloss{direct initialization} too:

\begin{emcppslisting}
Class var(/*...*/);  // C++03-style direct initialization
Class var{/*...*/};  // C++11-style direct list initialization
\end{emcppslisting}
    

\noindent Note that C++ does \emph{not} similarly relax the rules to allow for
initialization of aggregates with parentheses.\footnote{C++20 finally allows for aggregates to be
initialized with parentheses.}

The syntax suggested in the previous example is known as \emcppsgloss{direct
list initialization} and follows the rules of \emcppsgloss{direct
initialization} rather than \emcppsgloss{copy initialization} in that all
constructors of the named \lstinline!type! are both considered and
accessible in the initial overload set:

\begin{emcppslisting}
struct Q  // class containing explicit constructor
{
    explicit Q(int);  // value constructor taking a (ù{\codeincomments{int}}ù)
    // ...
};

Q x(5);    // OK  direct initialization can call explicit constructors.
Q y{5};    // OK, direct list initialization can call explicit constructors.
Q z = {5}; // Error, copy list initialization can't call explicit constructors.
\end{emcppslisting}
    

\noindent Either form of direct initialization (shown for \lstinline!x! and
\lstinline!y! in the code example above) may invoke an \lstinline!explicit!
constructor of class \lstinline!Q!, whereas \emph{copy list} initialization
will necessarily result in a compile-time error.

However, following the rules of C++11 braced initialization, narrowing
conversions are rejected by direct \emph{list} initialization:

\begin{emcppslisting}
long a = 3L;

Q b(a);  // OK, direct initialization
Q c{a};  // Error, direct list initialization cannot use a narrowing conversion.
\end{emcppslisting}
    

\noindent Similarly, \lstinline!explicit! conversion operators (see
\featureref{\locationa}{explicit-conversion-operators}) can be considered when
\emcppsgloss[direct list initialization]{\emph{direct} list initialization} (or \emcppsgloss{direct
initialization}) is employed on a scalar, but not so with
\emcppsgloss[copy list initialization]{\emph{copy} list initialization} (or \emcppsgloss{copy
initialization}). Consider, for example, a class, \lstinline!W! that can
covert to either an \lstinline!int! or a \lstinline!long! where conversion to
\lstinline!int! is explicit and therefore must be \emph{direct}:

\begin{emcppslisting}
struct W
{
    explicit operator int()  const;  // used via direct initialization only
             operator long() const;  // used via direct or copy initialization
};
\end{emcppslisting}
    

\noindent Initializing a \lstinline!long! variable with an expression of type
\lstinline!W! can be accomplished via either \emph{direct} or \emph{copy}
initialization (e.g., \lstinline!jDirect! and \lstinline!jCopy!, respectively,
in the code snippet below), but initializing an \lstinline!int! variable
with such an expression can be accomplished only via \emph{direct}
initialization (e.g., \lstinline!iDirect!):

\begin{emcppslisting}
long jDirect {W()};  // OK, considers both operators, calls (ù{\codeincomments{operator long}}ù)
long jCopy = {W()};  // OK, considers implicit op only, calls (ù{\codeincomments{operator long}}ù)
int  iDirect {W()};  // OK, considers both operators, calls (ù{\codeincomments{operator int}}ù)
int  iCopy = {W()};  // Error, considers implicit op only, narrowing conversion
\end{emcppslisting}
    

\noindent In the example above, attempting to use \emcppsgloss[copy list initialization]{\emph{copy} list
initialization} (e.g., \lstinline!iCopy!) forces the conversion to
\lstinline!long! as the only option, which results in a narrowing
conversion and an ill-formed program.

Note that, for \emcppsgloss{aggregate} types, even \emcppsgloss{direct list
initialization} will not allow the explicit constructors of the
individual member types to be considered since such \emph{member-wise}
initialization is invariably \emcppsgloss[copy initialization]{\emph{copy} initialization}; see
\intrarefsimple{c++11-aggregate-initialization}.

We may use \emcppsgloss{direct list initialization} as part of \emcppsgloss{member
initialization lists} for base classes and member data of a class (note
that there is no equivalent to allow \emcppsgloss{copy list initialization}
in such a context). Consider, for example, an aggregate class,
\lstinline!B!, a non-aggregate class, \lstinline!C!, and a derived class,
\lstinline!D!, that inherits from \lstinline!B! and has an object of type
\lstinline!C! as a data member, \lstinline!m!:

\begin{emcppslisting}
struct B { int i; };        // aggregate base type
struct C { C(); C(int); };  // non-aggregate member type

struct D : B  // class publicly derived from (ù{\codeincomments{B}}ù) containing (ù{\codeincomments{C}}ù)
{
    C m;  // non-aggregate data member

    D()      : B{},  m{}  { }  // direct initialized base/member objects
    D(int x) : B{x}, m{x} { }  //    "       "        "     "       "
};
\end{emcppslisting}
    

\noindent In the definition of class \lstinline!D! above, both constructors employ
\emcppsgloss{direct list initialization}; the first is also an example of
\emph{value} initialization for both (aggregate) class \lstinline!B! and
(non-aggregate) class \lstinline!C!. Note that in C++03, aggregate bases
and members could only be default initialized, value initialized, or
direct initialized, and not initialized to another value.

New expressions are another context in which \emcppsgloss{direct list
initialization} (braces) or \emcppsgloss{direct initialization} (parentheses)
can occur and applies similarly to both \emcppsgloss{aggregate} and
non\emcppsgloss{aggregate} types. If no initializer is provided, the
allocated object is \emcppsgloss{default initialized}; if empty braces or
parentheses are supplied, the object is \emcppsgloss{value initialized};
otherwise, the object is initialized from the contents of the braced or
(where permitted) parenthesized list.

As an illustrative example, let's consider the scalar type \lstinline!int!
which itself can be \emcppsgloss{default initialized} (not initialized),
\emcppsgloss{value initialized} (to 0) via empty braces or parentheses, or
\emcppsgloss{direct initialized} via a single element within either
parentheses or braces:

\begin{emcppslisting}
int* s0  = new int;               // default initialized (no initializer)

int* t0 = new int();              // direct (value) initialized from (ù{\codeincomments{()}}ù)
int* t1 = new int{};              // direct (value) list initialized from (ù{\codeincomments{\{\}}}ù)

int* u0 = new int(7);             // direct initialized from (ù{\codeincomments{7}}ù)
int* u1 = new int{7};             // direct list initialized from (ù{\codeincomments{\{7\}}}ù)

int* v0 = new int[5];             // All 5 elements are default initialized.

int* w0 = new int[5]();           // All 5 elements are value initialized.
int* w1 = new int[5]{};           // Array is direct list initialized from (ù{\codeincomments{\{\}}}ù).

int* x0 = new int[5](9);          // Error, invalid initializer for an array
int* x1 = new int[5]{9};          // Array is direct list initialized from (ù{\codeincomments{\{9\}}}ù).

int* y1 = new int[5]{1,2,3};      // array direct list initialized from (ù{\codeincomments{\{1,2,3\}}}ù)

int* z1 = new int[5]{1,2,3,4,5};  // direct list initialized from (ù{\codeincomments{\{1,2,3,4,5\}}}ù)
\end{emcppslisting}
    

\noindent All the comments above apply to the object being created in the
\lstinline!new! expression; the pointer set to the address of the
dynamically allocated object is copy initialized in all cases. Note
that, in C++03, we could default initialize (e.g., \lstinline!v0!) or value
initialize (e.g., \lstinline!w0!) the elements of an array in a
\lstinline!new! expression but there was no way to initialize the elements
of such an array to anything other than their default value (e.g.,
\lstinline!x0!); as of C++11, direct list initialization with braces (e.g.,
\lstinline!x1!, \lstinline!y1!, \lstinline!z1!) makes this more flexible,
heterogeneous initialization of array elements in \lstinline!new!
expressions possible.

\subsubsection[Contrasting copy and direct list initialization]{Contrasting copy and direct list initialization}\label{contrasting-copy-and-direct-list-initialization}

The difference between \emcppsgloss[copy list initialization]{\emph{copy} list initialization} and
\emcppsgloss[direct list initialization]{\emph{direct} list initialization} can be seen in this example:

\begin{emcppslisting}
struct C
{
    explicit C() { }
    explicit C(int) { }
};

struct A  // aggregate of (ù{\codeincomments{C}}ù)
{
    C x;
    C y;
};

int main()
{
    C c1;           // OK, default initialization
    C c2{};         // OK, value initialization
    C c3{1};        // OK, direct list initialization
    C c4 = {};      // Error, copy list initialization cannot use explicit default ctor.
    C c5 = {1};     // Error, copy list initialization cannot use explicit ctor.

    C c6[5];        // OK, default initialization
    C c7[5]{};      // Error, aggregate initialization requires a non-(ù{\codeincomments{explicit}}ù) default ctor.
    C c8[5]{1};     // Error, aggregate initialization requires non-(ù{\codeincomments{explicit}}ù) ctors.
    C c9[5] = {};   // Error, aggregate initialization requires a non-(ù{\codeincomments{explicit}}ù) default ctor.
    C ca[5] = {1};  // Error, aggregate initialization requires non-(ù{\codeincomments{explicit}}ù) ctors.

    A a1;           // OK, default initialization
    A a2{};         // Error, aggregate initialization requires a non-(ù{\codeincomments{explicit}}ù) default ctor.
    A a3{1};        // Error, aggregate initialization requires non-(ù{\codeincomments{explicit}}ù) ctors.
    A a4 = {};      // Error, aggregate initialization requires a non-(ù{\codeincomments{explicit}}ù) default ctor.
    A a5 = {1};     // Error, aggregate initialization requires non-(ù{\codeincomments{explicit}}ù) ctors.
}
\end{emcppslisting}
    

\noindent Note that if the constructors for \lstinline!C! were not marked explicit,
then all of the variables in the example above would be safely
initialized. If only the \lstinline!int! constructor of \lstinline!C! were
explicit, then the initializations that did not depend on the
\lstinline!int! constructor would be valid:

\begin{emcppslisting}
struct C
{
    C() { }
    explicit C(int) { }
};

struct A  // aggregate of (ù{\codeincomments{C}}ù)
{
    C x;
    C y;
};

int main()
{
    C c1;           // OK, default initialization
    C c2{};         // OK, value initialization
    C c3{1};        // OK, direct list initialization
    C c4 = {};      // OK, copy list initialization
    C c5 = {1};     // Error, copy list initialization cannot use explicit ctor.

    C c6[5];        // OK, default initialization
    C c7[5]{};      // OK, value initialization
    C c8[5]{1};     // Error, aggregate initialization requires non-(ù{\codeincomments{explicit}}ù) ctors.
    C c9[5] = {};   // OK, copy list initialization
    C ca[5] = {1};  // Error, aggregate initialization requires non-(ù{\codeincomments{explicit}}ù) ctors.

    A a1;           // OK, default initialization
    A a2{};         // OK, value initialization
    A a3{1};        // Error, aggregate initialization requires non-(ù{\codeincomments{explicit}}ù) ctors.
    A a4 = {};      // OK, copy list initialization
    A a5 = {1};     // Error, aggregate initialization requires non-(ù{\codeincomments{explicit}}ù) ctors.
}
\end{emcppslisting}
    

\subsubsection[Integrating default member initialization with braced initialization]{Integrating default member initialization with braced initialization}\label{integrating-default-member-initialization-with-braced-initialization}

Another new feature for C++11 is \emph{default member initializers} for
data members in a class. This new syntax supports both copy list
initialization and value list initialization. However, initialization
with parentheses is not permitted in this context.

\begin{emcppslisting}
struct S
{
    int i = { 13 };

    S() { }                       // OK, i == 13.
    explicit S(int x) : i(x) { }  // OK, i == x.
};


struct W
{
    S a{};        // OK, by default (ù{\codeincomments{j.i == 13}}ù).
    S b{42};      // OK, by default (ù{\codeincomments{j.i == 42}}ù).
    S c = {42};   // Error, constructor for (ù{\codeincomments{S}}ù) is explicit.
    S d = S{42};  // OK, direct initialization of temporary for initializer
    S e(42);      // Error, fails to parse as a function declaration
    S f();        // OK, declares member function (ù{\codeincomments{f}}ù)
};
\end{emcppslisting}
    

\subsubsection[List initialization where the list itself is a single argument to a constructor]{List initialization where the list itself is a single argument to a constructor}\label{list-initialization-where-the-list-itself-is-a-single-argument-to-a-constructor}

Another new form of initialization for C++11 is \emcppsgloss{list
initialization} with a braced list of arguments to populate a container.
See \featureref{\locationc}{initlist} for details. If a braced list of arguments are
all of the same type, then the compiler will look for a constructor
taking an \lstinline!std::initializer_list<T>! argument, where \lstinline!T!
is that common type. Similarly, if a braced list of values can be
implicitly converted to a common type, then a constructor for an
\lstinline!std::initializer_list! of that common type will be preferred.
When initializing from a nonempty braced initializer list, a matching
initializer list constructor always wins overload resolution. However,
\emph{value} initializing from a pair of empty braces will prefer a
default constructor.

\begin{emcppslisting}
struct S
{
  S() {}
  S(std::initializer_list<int>) {}
  S(int, int);
};

S a;             // default initialization with default constructor
S b();           // function declaration!
S c{};           // value initialization with default constructor
S d = {};        // copy list initialization with (ù{\codeincomments{std::initializer\_list}}ù)
S e{1,2,3,4,5};  // direct list initialization with (ù{\codeincomments{std::initializer\_list}}ù)
S f{1,2};        // direct list initialization with (ù{\codeincomments{std::initializer\_list}}ù)
S g = {1,2};     // copy list initialization with (ù{\codeincomments{std::initializer\_list}}ù)
S h(1,2);        // direct initialization with two (ù{\codeincomments{int}}ù)s
\end{emcppslisting}
    

\subsubsection[Omitting the type name when braced initializing a temporary]{Omitting the type name when braced initializing a temporary}\label{omitting-the-type-name-when-braced-initializing-a-temporary}

In addition to supporting new forms of initialization, C++11 allows for
braced lists to implicitly construct an object where the type is known
by context, such as for function arguments and return values. This use
of a braced list without explicitly specifying a type is just like
constructing a temporary object by copy initialization, hence using copy
list initialization, for the implicit type. As such construction is copy
list initialization, it will reject explicit constructors:

\begin{emcppslisting}
struct S
{
    S(int, int) {}
    explicit S(const char*, const char*) {}
};

S foo(bool b)
{
    if (b)
    {
        return S{ "hello", "world" };  // OK, direct list initialization of temporary
        return { "hello", "world" };   // Error, copy list initialization cannot call explicit ctor.
    }
    else
    {
        return {0, 0};  // OK, (ù{\codeincomments{int}}ù) constructor is not explicit.
    }
}

void bar(S s) { }

int main()
{
   bar( S{0,0} );  // OK, direct list initialization, then copy initialization
   bar(  {0,0} );  // OK, copy list initialization
   bar( S{"Hello", "world"} );  // OK, direct list initialization
   bar(  {"Hello", "world"} );  // Error, copy list initialization cannot use explicit ctor.
}
\end{emcppslisting}
    

\subsubsection[Initializing variables in conditional expressions]{Initializing variables in conditional expressions}\label{initializing-variables-in-conditional-expressions}

As a final tweak to make initialization consistent across the language,
initializing a variable in the condition of a \lstinline!while! or
\lstinline!if! statement in C++03 supported only copy initialization and
required use of the \lstinline!=! token. For C++11, those rules are relaxed
to allow any valid form of braced initialization. Conversely, the
declaration of a control variable in a \lstinline!for! loop has supported
all forms of initialization permitted for a variable declaration since
the original C++ standard{\cprotect\footnote{Note that GCC would
traditionally accept the C++11-only syntaxes, even when using
  C++98/03.}}:

\begin{emcppslisting}
void f()
{
   for (int i = 0; ; ) {}    // OK in all versions of C++
   for (int i = {0}; ; ) {}  // OK (ù{\emphincomments{for aggregates}}ù) in C++03 and all types from C++11
   for (int i{0}; ; ) {}     // OK from C++11, direct list initialization
   for (int i{}; ; ) {}      // OK from C++11, value initialization
   for (int i(0); ; ) {}     // OK in all versions of C++
   for (int i(); ; ) {}      // OK in all versions of C++
   for (int i; ; ) {}        // OK in all versions of C++

   if (int i = 0) {}    // OK in all versions of C++
   if (int i = {0}) {}  // OK from C++11, copy list initialization
   if (int i{0}) {}     // OK from C++11, direct list initialization
   if (int i{}) {}      // OK from C++11, value initialization
   if (int i(0)) {}     // Error in all versions of C++
   if (int i()) {}      // Error in all versions of C++
   if (int i) {}        // Error in all versions of C++

   while (int i = 0) {}    // OK in all versions of C++
   while (int i = {0}) {}  // OK from C++11, copy list initialization
   while (int i{0}) {}     // OK from C++11, direct list initialization
   while (int i{}) {}      // OK from C++11, value initialization
   while (int i(0)) {}     // Error in all versions of C++
   while (int i()) {}      // Error in all versions of C++
   while (int i) {}        // Error in all versions of C++
}
\end{emcppslisting}
    

\subsubsection[Default initialization and \lstinline!=!\,\lstinline!default!]{Default initialization and {\SubsubsecCode =}\,{\SubsubsecCode default}}\label{default-initialization-and-=-default}

Another new feature for C++11 is the notion of defaulted constructors,
defined by \lstinline!=!~\lstinline!default! (see \featureref{\locationa}{Defaulted-Special-Member-Functions}) to have
the same definition as the implicitly defined constructor. This is
especially useful when you want to provide constructors for your class
without losing the \emcppsgloss{triviality} otherwise associated with the
implicit constructors, notably for the default constructor that would no
longer be declared as soon as another constructor is declared.

One important property of \lstinline!=!~\lstinline!default! constructors is
that, while they are user \emph{declared}, they are not user
\emph{provided} as long as this definition occurs within the class
definition itself:

\begin{emcppslisting}
struct Trivial
{
    int i;
    Trivial() = default;            // user declared but not user provided
};

struct NonTrivial
{
    int j;
    NonTrivial();                    // user declared
};

NonTrivial::NonTrivial() = default;  // user provided
\end{emcppslisting}
    

\noindent Following the rules derived from C++03, a class with no user-provided
default constructor may behave differently under \emph{default}
initialization compared to \emph{value} initialization:

\begin{emcppslisting}
void demo()
{
    const Trivial    a;    // Error, (ù{\codeincomments{a.i}}ù) not initialized
    const Trivial    b{};  // OK, (ù{\codeincomments{b.i}}ù) = 0.
    const NonTrivial c;    // OK, but (ù{\codeincomments{c.j}}ù) never initialized
    const NonTrivial d{};  // OK, but (ù{\codeincomments{d.j}}ù) never initialized
}
\end{emcppslisting}
    

\noindent Note that we use function local variables for this example to avoid
confusion with \emph{zero} initialization for global variables.

\subsubsection[Copy initialization and scalars]{Copy initialization and scalars}\label{copy-initialization-and-scalars}

With the addition of explicit conversion operators to C++11 (see \featureref{\locationa}{explicit-conversion-operators}), it becomes possible for \emph{copy}
initialization and \emph{copy list} initialization to fail for scalars
and similarly for \emph{direct list} initialization:

\begin{emcppslisting}
struct S
{
    explicit operator int() const { return 1; }
};

S one{};

int a(one);     // OK, (ù{\codeincomments{a}}ù) = 1.
int b{one};     // OK, (ù{\codeincomments{b}}ù) = 1.
int c = {one};  // Error, copy list initialization used with
                // with explicit conversion operator.
int d = one;    // Error, copy initialization used with
                // with explicit conversion operator.

class C {
    int x;
    int y;

public:
    C(const S& value) : x(value)  // OK, (ù{\codeincomments{x}}ù) = 1
                      , y{value}  // OK, (ù{\codeincomments{y}}ù) = 1
    {
    }
};
\end{emcppslisting}
    

\subsection[Use Cases]{Use Cases}\label{use-cases-bracedinit}

\subsubsection[Defining a value-initialized variable]{Defining a value-initialized variable}\label{defining-a-value-initialized-variable}

The C++ parser has a pitfall where an attempt to value initialize a
variable turns out to be a function declaration:

\begin{emcppslisting}
struct S{};

void foo()
{
    S s1();      // Oops! function declaration
    S s2 = S();  // variable declaration using value initialization and then copy initialization
}
\end{emcppslisting}
    

\noindent The declaration of \lstinline!s1! looks like an attempt to value initialize
a local variable of type \lstinline!S!, but, in fact, it is a forward
declaration for a function \lstinline!s1! that takes no arguments and
returns an \lstinline!S! object by value. This is particularly surprising
for folks who did not realize we could declare (but not define) a
function within the body of another function, a feature retained from
the original C Standard. Clearly, there would be an ambiguity in the
grammar at this point unless the language provided a rule to resolve the
ambiguity, and the grammar opts in favor of the function declaration in
all circumstances, including at function local scope. Whilst this rule
would be essential at namespace/class scope, otherwise functions taking
no arguments could not be so easily declared, the same rule also applies
at function local scope, first for consistency and second for
compatibility with pre-existing C code that we might want to compile
more strictly\footnote{While C and C++ both enforce type safety, C++
enforces strict type safety, where all declared types are distinct; C
enforces structural conformance, where two distinct structs with the
same sequence of members are treated as the same type.} with a C++ compiler.

By switching from parentheses to braces, there is no more risk of
confusion between a vexing parse and a variable declaration:

\begin{emcppslisting}
S s{};  // object of type (ù{\codeincomments{S}}ù)
\end{emcppslisting}
    
\noindent As a digression, it is worth noting that \lstinline!S! is an empty type.
Just as it is ill-formed to rely on default initialization for a
\lstinline!const! object of a trivial type such as an \lstinline!int! or an
aggregate of just an \lstinline!int!, it was also, prior to Defect Report
CWG 253 for C++17 and the introduction of
\emcppsgloss[const default constructible]{const-default-constructible} types, ill-formed to rely on
default initialization for a \lstinline!const! object of an empty type.
Hence, it is often desirable to value initialize such objects, running
into the vexing parse. As this Defect Report was resolved at the end of
2016, and applies retroactively to earlier dialects, most current
compilers no longer enforce this restriction, and some compilers
(notably GCC) had already stopped enforcing this rule several years
previously:

\begin{emcppslisting}
const S cs1;        // Error on some compilers (as described above)
const S cs2{};      // OK, value initialization
const S cs3 = {};   // OK, aggregate initialization
const S cs4 = S();  // OK, copy initialization
\end{emcppslisting}
    

\subsubsection[Avoiding the most vexing parse]{Avoiding the most vexing parse}\label{avoiding-the-most-vexing-parse}

Value initializing function arguments can lead to another pitfall, often
called \emph{the most vexing parse}. C++ will parse the intended value
initialization of a function argument as the declaration of an unnamed
parameter of a function type instead, which would otherwise not be legal
but for the language rule that such a parameter implicitly decays to a
pointer to a function of that type:

\begin{emcppslisting}
struct V { V(const S&) { } };

void foo()
{
    V v1(S());    // most vexing parse, declares function (ù{\codeincomments{v1}}ù) taking a function pointer
    V v2((S()));  // workaround, object of type (ù{\codeincomments{V}}ù) due to non-redundant parentheses on argument

    S x = S();    // declare a variable of type (ù{\codeincomments{S}}ù)
    V v3(x);      // workaround, object of type (ù{\codeincomments{V}}ù) but argument is now named, with longer lifetime
}
\end{emcppslisting}
    

\noindent In the example above, \lstinline!v1! is the forward declaration of a
function in the surrounding namespace that returns an object of type
\lstinline!V! and has an (unnamed) parameter of type
pointer-to-function-returning-\lstinline!S!-and-taking-no-arguments,
\lstinline!S(*)()!. That is, the declaration is equivalent to:

\begin{emcppslisting}
V v1(S(*)());
\end{emcppslisting}
    

\noindent This most vexing of parses can be disambiguated by having the arguments
clearly form an expression, not a type. One simple way to force the
argument to be parsed as an expression is to add an otherwise redundant
pair of parentheses. Note that declaring the constructor of \lstinline!V!
as \lstinline!explicit! in the hopes of forcing a compile error is no help
here since the declaration of \lstinline!v1! is not interpreted as a
declaration of an object of type \lstinline!V!, so the \lstinline!explicit!
constructor is never considered.

With the addition of generalized braced initialization in C++11, a
coding convention to prefer empty braces, rather than parentheses, for
all value initializations avoids the question of the most vexing parse
arising:

\begin{emcppslisting}
V v4(S{});  // direct initialize object of type (ù{\codeincomments{V}}ù) with value initialized temporary
\end{emcppslisting}
    

\noindent Note that the most vexing parse can also apply to constructors taking
multiple arguments, but the issue arises less often since any one
supplied argument clearly being an expression, rather than a function
type, resolves the whole parse:

\begin{emcppslisting}
struct W { W(const S&, const S&) { } };

void foo() {
  W w1( S(),  S());  // most vexing parse, declares function (ù{\codeincomments{w1}}ù) taking two function pointers
  W w2((S()), S());  // workaround, object of type (ù{\codeincomments{V}}ù) due to non-redundant parentheses on argument
  W w3( S{},  S());  // workaround, even a single use of (ù{\codeincomments{S\{\}}}ù) disambiguates further use of (ù{\codeincomments{S()}}ù)
}
\end{emcppslisting}
    

\subsubsection[Uniform initialization in generic code]{Uniform initialization in generic code}\label{uniform-initialization-in-generic-code}

One of the design concerns facing an author of generic code is which
form of syntax to choose to initialize objects of a type dependent on
template parameters. Different C++ types behave differently and accept
different syntaxes, so providing a single consistent syntax for all
cases is not possible. Consider the following example of a simple test
harness for a unit testing framework:

\begin{emcppslisting}
#include <initializer_list>

template <class T, class U>
bool run_test(bool (*test)(const T&), std::initializer_list<U> iL)
{
    for (const auto& val : iL)
    {
        T obj = val;      // initialize the test value
        if (!test(obj))
        {
            return false;
        }
    }

    return true;
}
\end{emcppslisting}
    

\noindent In this example, a test function is provided for an object of parameter
type \lstinline!T! along with an \lstinline!initializer_list! of test values.
The \lstinline!for! loop will construct a test object with each test value,
in turn, and call the test function, returning early if any value fails.
The question is which syntax to use to create the test object
\lstinline!obj!.

\begin{itemize}
\item{As written, the example uses \emph{copy} initialization --- \lstinline!T!~\lstinline!obj!~\lstinline!=!~\lstinline!val;! --- and so will fail to compile if a non-\lstinline!explicit! constructor cannot be found or if \lstinline!T! is an aggregate that is not \lstinline!U!.}
\item{If we switched to \emph{direct} initialization --- \lstinline!T!~\lstinline!obj(val);! --- then explicit constructors would also be considered.}
\item{If we switched to \emph{direct list} initialization --- \lstinline!T!~\lstinline!obj{val};! --- then aggregates would be supported as well as explicit constructors but not narrowing conversions;\linebreak[4] \lstinline!initializer_list! constructors are also considered and preferred.}
\item{If we switched to \emph{copy list} initialization --- \lstinline!T!~\lstinline!obj!~\lstinline!=!~\lstinline!{val};! --- then aggregates would be supported, but it would be an error if an explicit constructor is the best match, rather than considering the non-\lstinline!explicit! constructors for the best viable match, and it would be an error to rely on a narrowing conversion; \lstinline!initializer_list! constructors are also considered and preferred.}
\end{itemize}

Table \ref{table-bracedinit-inittypes} summarizes the different initialization types and highlights the
options and trade-offs.
\begin{table}[h!]
\begin{center}
\begin{threeparttable}
\caption{Summary of the different initialization types}\label{table-bracedinit-inittypes} \vspace{1.5ex}
{\small \begin{tabular}{c|c|c|c|c|c}\thickhline
\rowcolor[gray]{.9}    {\sffamily\bfseries Initialization} & {\sffamily\bfseries Syntax } & {\sffamily\bfseries Aggregate} & {\sffamily\bfseries Explicit} & {\sffamily\bfseries Narrowing} & {\sffamily\bfseries \lstinline!initializer_list!} \\ 
 \rowcolor[gray]{.9}    {\sffamily\bfseries Type} & & {\sffamily\bfseries Support} & {\sffamily\bfseries Constructor} & & {\sffamily\bfseries Constructor} \\
\rowcolor[gray]{.9}   & & &  {\sffamily\bfseries Used} & & {\sffamily\bfseries Used}\\ \hline
\emcppsgloss[copy]{Copy} & \lstinline!T!~\lstinline!obj!~\lstinline!=!~\lstinline!val;! & only if \lstinline!U!  & no & allow & no \\ \hline 
\emcppsgloss[direct]{Direct} & \lstinline!T!~\lstinline!obj(val);! & only if \lstinline!U!  & yes & allow & no \\ \hline 
\emcppsgloss[direct list]{Direct List} & \lstinline!T!~\lstinline!obj{val};! & yes  & yes & error & yes \\ \hline 
\emcppsgloss[copy list]{Copy List} & \lstinline!T!~\lstinline!obj!~\lstinline!=!~\lstinline!{val}! & yes  & error if  & error & yes \\ 
& & & best match & & \\ \thickhline 
\end{tabular}
}
\end{threeparttable}
    \end{center}
\end{table}
In general, there is no one true, universal syntax for initialization in
generic (template) code. The library author should make an intentional
choice among the trade-offs described in this section and document that
as part of their contract.

\subsubsection[Uniform initialization in factory functions]{Uniform initialization in factory functions}\label{uniform-initialization-in-factory-functions}

One of the design concerns facing an author of generic code is which
form of syntax to choose to initialize objects of a type dependent on
template parameters. Different C++ types behave differently and accept
different syntaxes, so providing a single consistent syntax for all
cases is not possible. Here we present the different trade-offs to
consider when writing a factory function that takes an arbitrary set of
parameters to create an object of a user-specified type:

\begin{emcppslisting}
template <class T, class... ARGS>
T factory1(ARGS&&... args)
{
    return T(std::forward<ARGS>(args)...);
}

template <class T, class... ARGS>
T factory2(ARGS&&... args)
{
    return T{std::forward<ARGS>(args)...};
}

template <class T, class... ARGS>
T factory3(ARGS&&... args)
{
    return {std::forward<ARGS>(args)...};
}
\end{emcppslisting}
    

\noindent All three factory functions are defined using \emcppsgloss{perfect
forwarding} (see \featureref{\locationc}{forwardingref}) but support
different subsets of C++ types and may interpret their arguments
differently.

\lstinline!function1! returns a value created by direct initialization but,
because it uses parentheses, cannot return an aggregate unless (as a
special case) the \lstinline!args! list is empty or contains exactly one
argument of the same type \lstinline!T!; otherwise, the attempt to
construct the return value will parse as an error.{\cprotect\footnote{Note
that C++20 will allow aggregates to be initialized with parentheses as
well as with braces, which will result in this form being accepted for
  aggregates as well.}}

\lstinline!function2! returns an object created by \emcppsgloss{direct list
initialization}. Hence, \lstinline!function2! supports the same types as
\lstinline!function1!, plus aggregates. However, due to the use of braced
initialization, \lstinline!function2! will reject any types in
\lstinline!ARGS! that require narrowing conversion when passed to the
constructor (or to initialize the aggregate member) of the return value.
Also, if the supplied arguments can be converted into a homogeneous
\lstinline!std::initializer_list! that matches a constructor for
\lstinline!T!, then that constructor will be selected, rather than the
constructor best matching that list of arguments, despite
\lstinline!function2! being called using parentheses (as for any function
call).

\lstinline!function3! behaves the same as \lstinline!function2!, except that
it uses \emph{copy list initialization} so will also produce a compile
error if the selected constructor for the return value is declared as
\lstinline!explicit!.

There is no one true form of initialization that works best in all
circumstances for such a factory function, and it is for library
developers to choose (and document in their contract) the form that best
suits their needs. Note that the Standard Library runs into this same
problem when implementing factory functions like
\lstinline!std::make_shared! or the \lstinline!emplace! function of any
container. The Standard Library consistently chooses parentheses
initialization like \lstinline!function1! in the code example above, and so
these functions do not work for aggregates prior to C++20.

\subsubsection[Uniform member initialization in generic code]{Uniform member initialization in generic code}\label{uniform-member-initialization-in-generic-code}

With the addition of general braced initialization to C++11, class
authors should consider whether constructors should use \emph{direct}
initialization or \emph{direct list} initialization to initialize their
bases and members. Note that as copy initialization and copy list
initialization are not options, whether or not the constructor for a
given base or member is \lstinline!explicit! will never be a concern.

Prior to C++11, writing code that initialized aggregate subobjects
(including arrays) with a set of data in the constructor's member
initializer list was not really possible. We could only \emph{default}
initialize, \emph{value} initialize, or \emph{direct} initialize from
another aggregate of the same type.

Starting with C++11, we are able to initialize aggregate members with a
list of values, using aggregate initialization in place of direct list
initialization for members that are aggregates:

\begin{emcppslisting}
struct S
{
    int         i;
    std::string str;
};

class C
{
    int j;
    int a[3];
    S   s;

public:
    C(int x, int y, int z, int n, const std::string t)
    : j(0)
    , a{ x, y, z }  // Ill-formed in C++03, OK in C++11
    , s{ n, t }     // Ill-formed in C++03, OK in C++11
    {
    }
};
\end{emcppslisting}
    

\noindent Note that as the initializer for \lstinline!C.j! shows in the code example
above, there is no requirement to consistently use either braces or
parentheses for all member initializers.

As with the case of factory functions, the class author must make a
choice for constructors between adding support for initializing
aggregates vs.~reporting errors for narrowing conversion. Since member
initialization supports only \emph{direct} list initialization, there is
never a concern regarding \lstinline!explicit! conversions in this context:

\begin{emcppslisting}
template <class T>
class Wrap
{
    T data;

    template <class... ARGS>
    Wrap(ARGS&&... args)
    : data(std::forward<ARGS>(args))  // must be empty list or copy for aggregate (ù{\codeincomments{T}}ù)
    {
    }
};

template <class T>
class WrapAggregate
{
    T data;

    template <class... ARGS>
    WrapAggregate(ARGS&&... args)
    : data{std::forward<ARGS>(args)}  // no narrowing conversions
    {
    }
};
\end{emcppslisting}
    

\noindent Again, there is no universal best answer, and an explicit choice should
be made and documented so that consumers of the class know what to
expect.

\subsection[Potential Pitfalls]{Potential Pitfalls}\label{potential-pitfalls-bracedinit}

\subsubsection[Inadvertently calling initializer-list constructors]{Inadvertently calling initializer-list constructors}\label{inadvertently-calling-initializer-list-constructors}

Classes with an \lstinline!std::initializer_list! constructors (see
\featureref{\locationc}{initlist}) follow some special rules to disambiguate
overload resolution, which contain subtle pitfalls for the unwary. This
pitfall describes how overload resolution might (or might not) select
those constructors in surprising ways.

When an object is initialized by braced initialization, the compiler
will first look to find an \lstinline!std::initializer_list! constructor
that could be called, with the exception that if the braced list is
empty, a default constructor (if available) would have priority:

\begin{emcppslisting}
#include <initializer_list>

struct S {
  explicit S() {}
  explicit S(int) {}
  S(std::initializer_list<int> iL) { if (0 == iL.size()) {throw 13;} }
};

S a{};            // OK, value initialization
S b = {};         // Error, default constructor is (ù{\codeincomments{explicit}}ù)
S c{1};           // OK, (ù{\codeincomments{std::initializer\_list}}ù)
S d = {1};        // OK, (ù{\codeincomments{std::initializer\_list}}ù)
S e{1, 2, 3};     // OK, (ù{\codeincomments{std::initializer\_list}}ù)
S f = {1, 2, 3};  // OK, (ù{\codeincomments{std::initializer\_list}}ù)
\end{emcppslisting}
    

\noindent In the presence of initializer list constructors, the overload
resolution to select which constructor to call will be a two-step
process. First, all initializer-list constructors are considered, and
only if no matching \lstinline!std::initializer_list! constructor has been
found, non--initializer-list constructors will be considered. This
process has some possibly surprising consequences since implicit
conversions are allowed when performing the overload matching. It is
possible that an \lstinline!std::initializer_list! constructor requiring
an implicit conversion will be selected over a non--initializer-list
constructor that does not require a conversion:

\begin{emcppslisting}
#include <initializer_list>

struct S
{
    S(std::initializer_list<int>); // #1
    S(int i, char c);              // #2
};

S s1{1, 'a'};  // calls #1, even though #2 would be a better match
\end{emcppslisting}
    

\noindent In the example above, due to braced initialization preferring
initializer-list constructors and because \lstinline!S! has an
\lstinline!initializer_list! constructor that can match the initializer of
\lstinline!s1!, the constructor that would have been a better match
otherwise is not considered.

The other possibly surprising consequence is related to narrowing
conversion being checked for only \emph{after} the constructor has been
selected. This means that an \lstinline!initializer_list! constructor that
matches but requires a narrowing conversion will cause an error even in
the presence of a non\lstinline!initializer_list! constructor that would
be a match without requiring a narrowing conversion:

\begin{emcppslisting}
#include <initializer_list>

struct S
{
    S(std::initializer_list<int>); // #1
    S(int i, double d);            // #2
};

S s2{1, 3.2}; // narrowing conversion when attempting to call #1, 
              //     even though invoking #2 would be well formed
\end{emcppslisting}
    

\noindent In the example above, due to braced initialization first selecting a
constructor and then checking for narrowing conversion, the
non-initializer-list constructor, which would not require a narrowing
conversion, is not considered.

Both of these situations can be resolved by using parentheses or other
forms of initialization than brace lists, which do not prefer
initializer-list constructors:

\begin{emcppslisting}
#include <initializer_list>

struct S
{
    S(std::initializer_list<int>); // #1
    S(int i, char c);              // #2
    S(int i, double d);            // #3
};

S s3(1, 'c'); // calls #2
S s4(1, 3.2); // calls #3
\end{emcppslisting}
    
\noindent This problem often comes up when talking about \lstinline!std::vector!:

\begin{emcppslisting}
vector<size_t> v1{5u, 13u};  // Possible bug here.
  // If trying to construct a vector of 5 (ù{\codeincomments{size\_t}}ù) with value 13,
  // The (ù{\codeincomments{std::initializer\_list}}ù) constructor is preferred over exact match,
  // so we actually construct a vector with 2 values, 5 and 13.

vector<size_t> v2(5u, 13u);  // OK, calls the normal constructor
\end{emcppslisting}
    

\subsubsection[Classes with default member initializers lose aggregate status]{Classes with default member initializers lose aggregate status}\label{classes-with-default-member-initializers-lose-aggregate-status}

An aggregate class is a class with no user-provided constructors, no
base classes, no virtual functions, and all public data members. Braced
initialization of an aggregate matches each member of the brace list to
the corresponding member of the aggregate class in the order of
declaration. However, if any of the members has a \featureref{\locationc}{Default-Member-Initializers}, then the class ceases to be an aggregate in
C++11, and braced initialization will fail to compile because no
matching constructor will be found. C++14 fixes this oversight, so this
pitfall should occur only when supporting code across multiple versions
of the language, which is typically more of a concern for library
maintainers than application developers:

\begin{emcppslisting}
struct S
{
    int a;
    int b;
    int c;
};

struct A  // not an aggregate in C++11 (aggregate in C++14, see next section)
{
    int a{1};
    int b{2};
    int c{3};
};

S s1 = {};         // OK, aggregate initialization, value initializes each member
S s2 = {1, 2, 3};  // OK, aggregate initialization

A a1 = {};         // OK, value initialization with implicit default ctor in C++11
A a2 = {4, 5, 6};  // Error, no matching constructor in C++11
\end{emcppslisting}
    

\subsubsection[Implicit move and named return value optimization may be disabled in \lstinline!return! statements]{Implicit move and named return value optimization may be disabled in {\SubsubsecCode return} statements}\label{implicit-move-and-named-return-value-optimization-may-be-disabled-in-return-statements}

Using extra braces in a return statement around a value may disable the
named return value optimization or an implicit move into the returned
object.

Named return value optimization (NRVO) is an optimization that compilers
are allowed to perform when the operand of a return statement is just
the name (id-expression) of a nonvolatile local variable (an object of
automatic storage duration that is not a parameter of the function or a
catch clause) and the type of that variable, ignoring cv-qualification,
is the same as the function return type. In such cases, the compiler is
allowed to elide the copy implied by the return expression and
initialize the return value directly where the local variable is
defined. Naturally this applies only to functions returning objects, not
pointers or references. Note that this optimization is allowed to change
the meaning of programs that may rely on observable side effects on the
elided copy constructor. Most modern compilers are capable of performing
this optimization in at least simple circumstances, such as where there
is only one return expression for the whole function.

In the example below, we see that the \lstinline!no_brace()! function
returns using the name of a local variable from within that function. As
we call \lstinline!no_brace()! we can observe (with a compiler that
performs the optimization) that only one object of the \lstinline!S! class
is created, using its default constructor. There is no copy, no move,
and no other object created. Essentially the local variable, \lstinline!a!,
inside the \lstinline!no_brace()! function is created directly in the
memory region of the variable \lstinline!m1! of the \lstinline!main()!
function.

In the \lstinline!braced()! function, we use the exact same local variable,
but in the return statement we put braces around its name; therefore,
the operand of the return is no longer a name (id-expression), and so
the rules that allow NRVO do not apply. By calling \lstinline!braced()!, we
see that now two copies, and so two objects, are created, the first
being \lstinline!a!, the local variable, using the default constructor, and
the second being \lstinline!m2!, which is created as a copy of \lstinline!a!,
demonstrating that NRVO is not in effect:

\begin{emcppslisting}
#include <iostream>

struct S
{
    S()          { std::cout << "S()\n"; }
    S(const S &) { std::cout << "S(copy)\n"; }
    S(S &&)      { std::cout << "S(move)\n"; }
};

S no_brace()
{
    S a;
    return a;
}

S braced()
{
    S a;
    return { a };  // disables NRVO
}

int main()
{
    S m1 = no_brace();  // S()
    S m2 = braced();    // S(), S(copy)
}
\end{emcppslisting}
    

\noindent Implicit move (see \featureref{\locationc}{Rvalue-References}) in a return statement
is a more subtle operation, so much so that it required a defect
report{\cprotect\footnote{\cite{yasskin12}}} to actually
make it work as the original intention. We demonstrate implicit moves in
a return statement from a local variable by using two types. The class
type \lstinline!L! will be used for the local variable, whereas the class
type \lstinline!R!, which can be move- or copy-constructed from \lstinline!L!,
is used as the return type. Essentially, we are forcing a type
conversion in the return statement, one that may be a copy or a move.

The \lstinline!no_brace()! function just creates a local variable and
returns it. By calling the function, we observe that an \lstinline!L!
object is created, which is then moved into an \lstinline!R! object. Note
that the wording of the ISO standard allows this implicit move only if
the return statement's operand is a name (an id-expression).

The \lstinline!braced()! function is identical to the previous one, except
for adding curly braces around the operand of the return statement.
Calling the function shows that the \emph{move-from-\lstinline!L!} return
expression turned into a \emph{copy-from-\lstinline!L!} return expression
because a braced initializer is not a name of an object:

\begin{emcppslisting}
#include <iostream>

struct L
{
    L()        { std::cout << "L()\n"; }
};

struct R
{
    R(const L &) { std::cout << "R(L-copy)\n"; }
    R(L &&)      { std::cout << "R(L-move)\n"; }
};

R no_brace()
{
    L a;
    return a;
}

R braced()
{
    L a;
    return { a };  // disables implicit move from l
}

int main()
{
    R r1 = no_brace();  // L(), R(L-move)
    R r2 = braced();    // L(), R(L-copy)
}
\end{emcppslisting}
    

\subsubsection[Surprising behavior of aggregates having deleted constructors]{Surprising behavior of aggregates having deleted constructors}\label{surprising-behavior-of-aggregates-having-deleted-constructors}

Value initialization of aggregates is allowed with a braced initializer
list, even if the default constructor is deleted{\cprotect\footnote{Note
that C++20 finally addresses the issue so the presence of deleted
  constructors cause a class to no longer qualify as an aggregate.}}:

\begin{emcppslisting}
struct S
{
    int data;
    S() = delete; // don't want "empty"
};

S s{}; // surprisingly works (until C++20), and (ù{\codeincomments{0 == s.data}}ù)
\end{emcppslisting}
    
\noindent This surprising pitfall occurs for two reasons:
\begin{enumerate}
\item{A deleted constructor is \emph{user declared} but not \emph{user provided}, so it does not feature in the list of things that stop a class being an aggregate.}
\item{The rules state that aggregate initialization is not defined in terms of constructors but directly in terms of the initialization of its members.}
\end{enumerate}

\subsection[Annoyances]{Annoyances}\label{annoyances}

\subsubsection[Narrowing aggregate initialization may break C++03 code]{Narrowing aggregate initialization may break C++03 code}\label{narrowing-aggregate-initialization-may-break-c++03-code}

When compiling existing C++03 code with a C++11 compiler, previously
valid code may report errors for narrowing conversion in aggregate (and,
therefore, also array) initialization.

\begin{emcppslisting}
unsigned u;
u = 128;                    // (ù{\codeincomments{u}}ù) is computed to an (ù{\codeincomments{int}}ù)-friendly value
int ia[] = { 1, 2, u, 9 };  // OK in C++03, narrowing is allowed.
                            // Error in C++11, narrowing conversion.
\end{emcppslisting}
    

\noindent Suppose that the computation in the above code ensures that the value
\lstinline!u! holds at the point of initialization is in the range of
values an \lstinline!int! is able to represent. Yet, the code will not
compile in C++11 or later modes. Unfortunately each and every case has
to be fixed by applying the appropriate type cast or changing the types
involved to be ``compatible''.

\subsubsection[No easy way to allow narrowing conversions]{No easy way to allow narrowing conversions}\label{no-easy-way-to-allow-narrowing-conversions}

In generic code, curly braces have to be used if support for aggregates
is required, but, if our interface definition requires supporting
narrowing conversions (for example \lstinline!std::tuple!), there is no
direct way to enable them:

\begin{emcppslisting}
struct S
{
    short m;
};

class X
{
    S m;

public:
    template <class U>
    X(const U& a) : m{a}  // no narrowing allowed
    {
    }
};

int i;
X x(i);  // Error, would narrow in initializing S.m
\end{emcppslisting}
    

\noindent The workaround is to \lstinline!static_cast! to the target type if it is
known or to use parentheses and give up aggregate support in the generic
code.{\cprotect\footnote{C++20 enables the use of parentheses to
  initialize aggregates.}}

\subsubsection[Breaks macro-invocation syntax]{Breaks macro-invocation syntax}\label{breaks-macro-invocation-syntax}

The macro-invocation syntax of the C++ preprocessor (inherited from C)
``understands'' parentheses and thus ignores commas within parentheses
but does not understand any other list markers, such as braces for
braced initialization, square brackets, or the angle bracket notation of
templates. If we attempt to use commas in other contexts, the macro
parsing will interpret such commas as separators for multiple macro
arguments and will likely complain that the macro does not support that
many arguments:

\begin{emcppslisting}
#define MACRO(oneArg) /*...*/

struct C
{
    C(int, int, int);
};

struct S
{
    int i1, i2, i3;
};

MACRO(C x(a, b, c))         // OK, commas inside parentheses ignored
MACRO(S y{a, b, c})         // Error, 3 arguments but MACRO needs 1
MACRO(std::map<int, int> z) // Error, 2 arguments but MACRO needs 1
\end{emcppslisting}
    

\noindent As the example above demonstrates, on the first macro invocation, the
commas within the parentheses are ignored, and the macro is invoked with
one argument: \lstinline!Demo!~\lstinline!x(a,!~\lstinline!b,!~\lstinline!c)!.

In the second macro invocation, we attempt to use braced initialization,
but, because the syntax of the preprocessor does not recognize curly
braces as special delimiters, the commas are interpreted as separating
macro arguments, so we end up with three unusual arguments: first
\lstinline!Demo!~\lstinline!y{a!, second \lstinline!b!, and finally the third
\lstinline!c}!. This problematic interaction between braced initialization
and macros has existed forever, even back in C code when initializing
arrays or \lstinline!struct!s. However, with braced initialization becoming
used more widely in C++, it is much more likely that a programmer will
encounter this annoyance.

The third invocation of \lstinline!MACRO! in the example is just a reminder
that the same issue exists in C++ with the angle brackets of templates.

The workaround, as is so often the case with the C preprocessor, is more
use of the C preprocessor! We need to define macros to help us hide the
commas. Such macros will use the \emcppsgloss{variadic macros} C99
preprocessor feature that was adopted for C++11 to turn a
comma-separated list into a braced-initializer list (and similarly for a
template instantiation):

\begin{emcppslisting}
#define BRACED(...) { __VA_ARGS__ }
#define TEMPLATE(name, ...) name<__VA_ARGS__>

MACRO(X y BRACED(a, b, c));             // OK, (ù{\codeincomments{X y \{ a, b, c \}}}ù)
MACRO(TEMPLATE(std::map, int, int) z);  // OK, (ù{\codeincomments{std::map<int, int> z}}ù)
\end{emcppslisting}
    

\noindent A common way this annoyance might show up is using the Standard Library
\lstinline!assert! macro:

\begin{emcppslisting}
bool operator==(const C&, const C&);

void f(const C& x, int i, int j, int k)
{
   assert(C(i, j, k) == x);  // OK
   assert(C{i, j, k} == x);  // Error, too many arguments to (ù{\codeincomments{assert}}ù) macro
}
\end{emcppslisting}
    

\subsubsection[Default member initializer does not deduce array size]{Default member initializer does not deduce array size}\label{default-member-initializer-does-not-deduce-array-size}

Although the syntax looks the same, default member initializers using
braced-initializer lists do not deduce the size of an array member:

\begin{emcppslisting}
struct S
{
    char s[]{"Idle"};  // Error, must specify array size
};
\end{emcppslisting}
    

\noindent The rationale is that there is no guarantee that the default member
initializer will be used to initialize the member; hence, it cannot be a
definitive source of information about the size of such a member in the
object layout.

\subsubsection[No copy list initialization in member initializer lists]{No copy list initialization in member initializer lists}\label{no-copy-list-initialization-in-member-initializer-lists}

The syntax for base and member initializers allows for both direct
initialization with parentheses (since C++03) and direct list
initialization with braces (since C++11). However, there is no syntax
corresponding to copy list initialization, which would allow member
initializers to report errors for using an unintended explicit
constructor or conversion operator. It would seem relatively intuitive
to extend the syntax to support
\lstinline!=!~\lstinline!{!~\lstinline!...!~\lstinline!}! for member initializers to
support such use, but so far there have been no proposals to add this
feature to the language. That may be a sign that there is simply no
demand, and the authors of this book are the only ones annoyed since
this is the only part of the language that supports \emph{direct}
initializations without a corresponding syntax for \emph{copy}
initializations.

\begin{emcppslisting}
class C
{
    explicit C(int);
    C(int, int);
};

class X
{
    C a;
    C b;
    C c;

public:
    X(int i)
    : a(i)       // OK, direct initialization
    : b(i)       // OK, direct list initialization
    : c = (i,i)  // Error, copy list initialization is not allowed.
    {
    }
};
\end{emcppslisting}
    

\subsubsection[Accidental meaning for explicit constructors passed multiple arguments]{Accidental meaning for explicit constructors passed multiple arguments}\label{accidental-meaning-for-explicit-constructors-passed-multiple-arguments}

In C++03, marking a default or multi-argument constructor explicit,
typically as a result of supplying default arguments, had no useful
meaning, and compilers did not warn about them because they were
harmless. However, C++11 takes notice of the \lstinline!explicit! keyword
for such constructors when invoked by copy list initialization. This
design point is generally not considered when migrating code from C++03
to C++11 and may require programmers to invest more thought, and
potentially split constructors with multiple default arguments into
multiple constructors, applying \lstinline!explicit! to only the intended
overloads:

\begin{emcppslisting}
class C
{
public:
    explicit C(int = 0, int = 0, int = 0);
};

C c0 = {};       // Error, constructor is (ù{\codeincomments{explicit}}ù).
C c1 = {1};      // Error, constructor is (ù{\codeincomments{explicit}}ù).
C c2 = {1,2};    // Error, constructor is (ù{\codeincomments{explicit}}ù).
C c3 = {1,2,3};  // Error, constructor is (ù{\codeincomments{explicit}}ù).

class D
{
public:
    D();
    explicit D(int i) : D(i, 0) { }  // delegating constructor
    D(int, int, int = 0);
};

D d0 = {};       // OK
D d1 = {1};      // Error, constructor is (ù{\codeincomments{explicit}}ù).
D d2 = {1,2};    // OK
D d3 = {1,2,3};  // OK

C f(int i, C arg)
{
    switch (i)
    {
        case 0: return {};         // Error, constructor is (ù{\codeincomments{explicit}}ù).
        case 1: return {1};        // Error, constructor is (ù{\codeincomments{explicit}}ù).
        case 2: return {1, 2};     // Error, constructor is (ù{\codeincomments{explicit}}ù).
        case 3: return {1, 2, 3};  // Error, constructor is (ù{\codeincomments{explicit}}ù).
    }
}

D g(int i, D arg)
{
    switch (i)
    {
      case 0: return {};         // OK
      case 1: return {1};        // Error, constructor is (ù{\codeincomments{explicit}}ù).
      case 2: return {1, 2};     // OK
      case 3: return {1, 2, 3};  // OK
    }
}

void test()
{
    f(0, {});         // Error, constructor is (ù{\codeincomments{explicit}}ù).
    f(0, {1});        // Error, constructor is (ù{\codeincomments{explicit}}ù).
    f(0, {1, 2});     // Error, constructor is (ù{\codeincomments{explicit}}ù).
    f(0, {1, 2, 3});  // Error, constructor is (ù{\codeincomments{explicit}}ù).

    g(0, {});         // OK
    g(0, {1});        // Error, constructor is (ù{\codeincomments{explicit}}ù).
    g(0, {1, 2});     // OK
    g(0, {1, 2, 3});  // OK
}
\end{emcppslisting}
    

\noindent Note that this topic is deemed an annoyance, rather than a pitfall,
because it affects only newly written C++11 (or later) code using the
new forms of initialization syntax, so it does not break existing C++03
code recompiled with a more modern language dialect. However, also note
that many containers and other types in the C++ Standard Library
inherited such a design and have not been refactored into multiple
constructors (although some such refactoring occurs in later standards).

\subsubsection[Obfuscation due to opaque use of braced-list]{Obfuscation due to opaque use of braced-list}\label{obfuscation-due-to-opaque-use-of-braced-list}

Use of braced initializers for function arguments, omitting any hint of
the expected object type at the call site, requires deep familiarity
with functions being called in order to understand the actual types of
arguments being initialized, especially when overload resolution must
disambiguate several viable candidates. Such usage may produce more
fragile code as further overloads are added, silently changing the type
initialized by the brace list as a different function wins overload
resolution. Such code is also much harder for a subsequent maintainer,
or casual code reader, to understand:

\begin{emcppslisting}
#include <initializer_list>

struct C
{
    C(int, int) { }
};

int test(C, long) { return 0; }

int main()
{
    int a = test({1, 2}, 3);
    return a;
}
\end{emcppslisting}
    

\noindent This program compiles and runs, returning the intended result. However,
consider how the behavior changes if we add a second overload during
subsequent maintenance:

\begin{emcppslisting}
#include <initializer_list>

struct C
{
    C(int, int) { }
};

int test(C, long) { return 0; }

struct A  // additional aggregate class
{
    int x;
    int y;
};

int test(A, int) { return -1; }  // overload for the aggregate class

int main()
{
    int a = test({1, 2}, 3);     // overload resolution prefers the aggregate
    return a;
}
\end{emcppslisting}
    

\noindent Because the overload for \lstinline!A! must now be considered, overload
resolution may pick a different result. If we are lucky, then the choice
of the \lstinline!A! and \lstinline!C! overloads becomes ambiguous, and an
error is diagnosed. However, in this case, there was an integer
promotion on the second argument, and the new \lstinline!A! overload is now
the stronger match, producing a different program result. If this
overload is added through maintenance of an included header file, this
code will have silently changed meaning without touching the file. If
the above flexibility is not the desired intent, the simple way to avoid
this risk is to always name the type of any temporary variables:

\begin{emcppslisting}
int main()
{
    int a = test(C{1,2}, 3);  // Overload resolution prefers (ù{\codeincomments{struct C}}ù).
    return a;
}
\end{emcppslisting}
    

\subsubsection[\lstinline!auto! deduction and braced initialization]{{\SubsubsecCode auto} deduction and braced initialization}\label{auto-deduction-and-braced-initialization}

C++11 introduces type inference, where an object's type is deduced from
its initialization, using the \lstinline!auto! keyword (see \featureref{\locationc}{auto-feature}). When presented with a homogeneous, nonempty list using
\emph{copy} list initialization, \lstinline!auto! will deduce the type of
the supplied argument list as an \lstinline!std::initializer_list! of the
same type as the list values. When presented with a braced list of a
single value using \emph{direct} list initialization, \lstinline!auto! will
deduce the variable type as the same type as the list value:

\begin{emcppslisting}
#include <initializer_list>

auto g{1};         // OK, deduces (ù{\codeincomments{g}}ù) is (ù{\codeincomments{int}}ù)
auto h{1, 2, 3};     // Error, (ù{\codeincomments{auto}}ù) requires exactly one element in brace list
auto i = {1};      // OK, deduces (ù{\codeincomments{i}}ù) is (ù{\codeincomments{initializer\_list<int>}}ù)
auto j = {1, 2, 3};  // OK, deduces (ù{\codeincomments{j}}ù) is (ù{\codeincomments{initializer\_list<int>}}ù)
\end{emcppslisting}
    

\noindent Note that the declarations of \lstinline!i! and \lstinline!j! in the code
example above would also be errors if the \lstinline!<initializer_list>!
header had not been included to supply the
\lstinline!std::initializer_list! class template.

Finally, observe that for \lstinline!auto! deduction from \emph{direct}
list initialization, an \lstinline!initializer_list! constructor may still
be called in preference to copy constructors, even though the syntax
seems restricted to making copies:

\begin{emcppslisting}
#include <iostream>
#include <initializer_list>

struct S
{
    S() { }
    S(std::initializer_list<S>) { std::cout << "init list\n";}
    S(const S&) { std::cout << "copy\n";}
};

int main()
{
    S s;
    auto s2{s};  // (ù{\codeincomments{std::initializer\_list<S>}}ù) constructor is called after
                 // deduction. (Note: (ù{\codeincomments{s2}}ù) is deduced to be of type (ù{\codeincomments{S}}ù).)
}
\end{emcppslisting}
    

\subsubsection[Compound assignment but not arithmetic operators accept braced lists]{Compound assignment but not arithmetic operators accept braced lists}\label{compound-assignment-but-not-arithmetic-operators-accept-braced-lists}

Braced initializers can be used to provide arguments to the assignment
operator and additionally to compound assignment operators such as
\lstinline!+=!, where they are treated as calls to the overloaded operator
function for class types, or as \lstinline!+=!~\lstinline!T{value}! for a
scalar type \lstinline!T!.{\cprotect\footnote{Although valid, the two
  \lstinline!x!~\lstinline!+=!~\lstinline!{3}! and
  \lstinline!x!~\lstinline!*=!~\lstinline!{3}! lines in the example compile
  successfully on Clang but not on GCC or MSVC (applies to all versions
  at the time of writing). The C++11 standard currently states: 
  \begin{quote}
  A braced-init-list may appear on the right-hand side of
  \begin{itemize}
  \item[---~]{an assignment to a scalar, in which case the initializer list
  shall have at most a single element. The meaning of
  \lstinline!x={v}!, where \lstinline!T! is the scalar
  type of the expression \lstinline!x!, is that of
  \lstinline!x=T{v}!. The meaning of
  \lstinline!x={}! is
  \lstinline!x=T{}!}
  \end{itemize}
  \end{quote} 
  (\cite{cpp11}, paragraph~9, section~5.17, ``Assignment and Compound Assignment Operators," p.~126). There is currently a defect
  report due to clarify the Standard and explicitly state that this rule
  also applies to compound assignments (see \cite{miller12} and
  \cite{miller21}).}} Note
that assigning to scalars supports brace lists of no more than a single
element and does not support compound assignment for pointer types,
since the brace lists are converted to a pointer type, which cannot
appear on the right-hand side of a compound assignment operator.

While the intent of compounded assignment is to be semantically
equivalent to the expression
\lstinline!a!~\lstinline!=!~\lstinline!a!~\lstinline!+!~\lstinline!b! (or
\lstinline!*!~\lstinline!b!, or \lstinline!-!~\lstinline!b!, and so on), brace lists
cannot be used in regular arithmetic expressions since the grammar does
not support brace lists as arbitrary expressions:

\begin{emcppslisting}
#include <initializer_list>

struct S
{
    S(std::initializer_list<int>) { }
    S& operator+=(const S&) { return *this; }
};

S operator+(const S&, const S&) { return S{}; }

void demo()
{
    S s1{};               // OK, calls (ù{\codeincomments{initializer\_list}}ù) constructor
    s1 += {1,2,3};        // OK, equivalent to (ù{\codeincomments{s1.operator+=(\{1,2,3\})}}ù)
    s1 = s1 + {1, 2, 3};  // Error, expecting an expression, not an (ù{\codeincomments{std::initializer\_list}}ù)
    s1 = operator+(s1, {1,2,3}); // OK, braces are allowed as function argument.

    int x = 0;
    x += {3};    // OK, equivalent to (ù{\codeincomments{x += int\{3\};}}ù).[^cwg_1542].
    x *= {5};    // OK, equivalent to (ù{\codeincomments{x *= int\{5\};}}ù).[^cwg_1542].

    char y[4] = {1, 2, 3, 4};
    char*p = +y;
    p += {3};    // Error,  equivalent to (ù{\codeincomments{p += (char*)\{3\};}}ù)
}
\end{emcppslisting}
    

\subsection[See Also]{See Also}\label{see-also}

TODO: Add see also items. Pulled from the feature:
\begin{itemize}
\item{\seealsoref{deleted-functions}{\seealsolocationa}}
\item{\seealsoref{explicit-conversion-operators}{\seealsolocationa}}
\item{\seealsoref{thread-safe-function-static-variables}{\seealsolocationa}}
\item{\seealsoref{long-long}{\seealsolocationa}}
\item{\seealsoref{null-pointer-literal-(nullptr)}{\seealsolocationa}}
\item{\seealsoref{initlist}{\seealsolocationc}}
\item{\seealsoref{Rvalue-References}{\seealsolocationc}}
\item{Implicit move (unknown what this was intended to be)}
\item{\seealsoref{Default-Member-Initializers}{\seealsolocationc}}
\item{perfect forwarding (unknown what this was intended to be) }
\item{\seealsoref{variadictemplate}{\seealsolocationc}}
\item{\seealsoref{auto-feature}{\seealsolocationc}}
\end{itemize}

\subsection[Further Reading]{Further Reading}\label{further-reading}

TBD




\newpage
%\section[{\tt constexpr} Functions]{Compile-Time Evaluatable Functions}\label{constexprfunc}


\emcppsFeature{
    short={\lstinline!constexpr! Functions},
    tocshort={{\TOCCode constexpr} Functions},
    long={Compile-Time Evaluatable Functions},
    rhshort={{\RHCode constexpr} Functions},
}{constexprfunc}
\setcounter{table}{0}
\setcounter{footnote}{0}
\setcounter{lstlisting}{0}
%\section[{\tt constexpr} Functions]{Compile-Time Evaluatable Functions}\label{constexprfunc}

placeholder text.........
 
 

\newpage
%\section[{\tt constexpr} Variables]{Compile-Time Accessible Variables}\label{constexprvar}
% 3 March 2021, to Josh for code check and addition to Annoyances 
% 6 March 2021, back from Josh
% 23 March 2021, copyedits in... 



\emcppsFeature{
    short={\lstinline!constexpr! Variables},
    tocshort={{\TOCCode constexpr} Variables},
    long={Compile-Time Accessible Variables},
    rhshort={{\RHCode constexpr} Variables},
}{constexprvar}
\setcounter{table}{0}
\setcounter{footnote}{0}
\setcounter{lstlisting}{0}
%\section[{\tt constexpr} Variables]{Compile-Time Accessible Variables}\label{constexprvar}

A variable or \emcppsgloss{variable template} of \emcppsgloss{literal type} may be
declared to be \lstinline!constexpr!, ensuring its successful construction
and enabling its use \emph{at compile-time}.

\subsection[Description]{Description}\label{description}

Variables of \emph{all} built-in types and certain user-defined types,
collectively known as \emcppsgloss[literal type]{literal types}, may be declared
\lstinline!constexpr!, allowing them to be initialized \emph{at
compile-time} and subsequently used in \emcppsgloss[constant expression]{constant expressions}:

\begin{emcppshiddenlisting}[emcppsbatch=e1]
#include <array>  // (ù{\codeincomments{std::array}}ù)
\end{emcppshiddenlisting}
\begin{emcppslisting}[emcppsbatch=e1]
          int i0 = 5;
    const int i1 = 5;               // (ù{\codeincomments{i1}}ù) is a compile-time constant.
constexpr int i2 = 5;               // (ù{\codeincomments{i2}}ù)  " "    "      "     "

    const double d1 = 5.0;
constexpr double d2 = 5.0;          // (ù{\codeincomments{d2}}ù) is a compile-time constant.

          const char *s1 = "help";
constexpr const char *s2 = "help";  // (ù{\codeincomments{s2}}ù) is a compile-time constant.
\end{emcppslisting}

Although \lstinline!const! integers initialized in the view of the compiler
can be used within constant expressions --- e.g., as the first argument
to \lstinline!static_assert!, as the size of an array, or as a non-type
template parameter --- such is not the case for any other types:

\begin{emcppslisting}[emcppsbatch=e1]
static_assert(i0 == 5, "");         // Error, (ù{\codeincomments{i0}}ù) is not a compile-time constant.
static_assert(i1 == 5, "");         // OK, (ù{\codeincomments{const}}ù) is "magical" for integers (only).
static_assert(i2 == 5, "");         // OK

static_assert(d1 == 5, "");         // Error, (ù{\codeincomments{d1}}ù) is not a compile-time constant.
static_assert(d2 == 5, "");         // OK

static_assert(s1[1] == 'e', "");    // Error, (ù{\codeincomments{s1}}ù) is not a compile-time constant.
static_assert(s2[1] == 'e', "");    // OK, the ASCII code for (ù{\codeincomments{e}}ù) is decimal 101.

int a[s2[1]];                       // OK, defines an array, (ù{\codeincomments{a}}ù), of 101 integers
static_assert(sizeof a == 404, ""); // OK, we are assuming the size of (ù{\codeincomments{int}}ù) is 4.

std::array<int, s1[1]> f;           // Error, (ù{\codeincomments{s1}}ù) is not (ù{\codeincomments{constexpr}}ù).
std::array<int, s2[1]> e;           // OK, defines a (ù{\codeincomments{std::array}}ù), (ù{\codeincomments{e}}ù), of 101 integers
\end{emcppslisting}
    
\noindent Prior to C++11, the types of variables usable in a \emcppsgloss{constant expression} were quite limited:

\begin{emcppslisting}
const int b;         // Error, (ù{\codeincomments{const}}ù) scalar variable must be initialized.
extern const int c;  // OK, declaration
const int d = c;     // OK, not constant initialized ((ù{\codeincomments{c}}ù) initializer not seen)

int ca1[c];          // Error, (ù{\codeincomments{c}}ù) initializer not visible
int da1[d];          // Error, (ù{\codeincomments{d}}ù) initializer not constant

const int c = 7;
int ca2[c];          // OK, initializer is visible
int da2[d];          // Error, (ù{\codeincomments{d}}ù) initializer (still) not constant

const int e = 17;    // OK
int ea[e];           // OK
\end{emcppslisting}
    
\noindent For an \emcppsgloss{integral constant} to be usable at compile time (e.g., as
part of \emcppsgloss{constant expression}), certain requirements must be
satisfied:

\begin{enumerate}
\item{The variable must be marked \lstinline!const!.}
\item{The initializer for a variable must have been seen by the time it is used, and it must be a \emcppsgloss{constant expression} — this is the information needed for a compiler to be able to make use of the variable in other \emcppsgloss[constant expression]{constant expressions}.}
\item{The variable must be of \emph{integral} type, e.g., \lstinline!bool!, \lstinline!char!, \lstinline!short!, \lstinline!int!, \lstinline!long!, \lstinline!long!~\lstinline!long!, as well as the \lstinline!unsigned! variations on these and any additional \lstinline!char! types; see also \featureref{\locationa}{long-long}.}
\end{enumerate}

This restriction to integral types provides support for those values
where compile-time constants are most frequently needed while limiting
the complexity of what compilers were required to support at compile
time.

Use of \lstinline!constexpr! when declaring a variable (or \emcppsgloss{variable template}; see \featureref{\locationb}{variable-templates}) enables a much richer
category of types to participate in \emcppsgloss[constant expression]{constant expressions}. This
generalization, however, was not made for mere \lstinline!const! variables
because they are not \emph{required} to be initialized by compile-time
constants:

\begin{emcppslisting}
int f() { return 0; }  // (ù{\codeincomments{f()}}ù) is not a compile-time (ù{\emphincomments{constant expression}ù).

                int x0 = f();  // OK
          const int x1 = f();  // OK, but (ù{\codeincomments{x1}}ù) is not a compile-time constant.
constexpr       int x2 = f();  // Error, (ù{\codeincomments{f()}}ù) "    "   "     "     "     "
constexpr const int x3 = f();  // Error, (ù{\codeincomments{f()}}ù) "    "   "     "     "     "
\end{emcppslisting}
    
\noindent As the example code above demonstrates, variables marked
\lstinline!constexpr! \emph{must} satisfy the same requirements needed for
integral constants to be usable as \emcppsgloss{constant expression}. Unlike
other integral constants, their initializers \emph{must} be constant
expressions or else the program is \emcppsgloss{ill formed}. 

For a variable of other than \lstinline!const! integral type to be usable
in a \emcppsgloss{constant expression}, certain criteria must hold:

\begin{enumerate}
\item{The variable must be annotated with \lstinline!constexpr!, which implicitly also declares the variable to be \lstinline!const!\footnote{C++20 added the \lstinline!constinit! keyword to identify a variable that is initialized at compile time (with a constant expression) but may then be modified at runtime.}:

\begin{emcppslisting}
struct S  // simple (aggregate) literal type
{
    int i;  // built-in integer data member
};

void test1()
{
    constexpr S s{1};       // OK, literal type (ù{\emphincomments{constant expression}}ù) initialized
    s = S();                // Error, (ù{\codeincomments{constexpr}}ù) implies (ù{\codeincomments{const}}ù).
    constexpr int j = s.i;  // OK, subobjects are usable in (ù{\emphincomments{constant expressions}}ù).
}
\end{emcppslisting}
    
\noindent In the example above, we have, for expedience of exposition, used brace
initialization to initialize the aggregate; see \featureref{\locationc}{bracedinit}. Note that subobjects of \lstinline!constexpr! objects
are also effectively \lstinline!constexpr! and can be used freely in
\emcppsgloss[constant expression]{constant expressions} even though they themselves might not be
explicitly declared \lstinline!constexpr!.}

\item{All \lstinline!constexpr! variables must be initialized with a \emcppsgloss{constant expression} when they are defined. Hence, every \lstinline!constexpr! variable has an initializer, and that initializer must be a valid \emcppsgloss{constant expression} (see \featureref{\locationc}{constexprfunc}):

\begin{emcppslisting}
          int g() { return 17; }  // a non-(ù{\codeincomments{constexpr}}ù) function
constexpr int h() { return 34; }  // a (ù{\codeincomments{constexpr}}ù) function

constexpr int v1;        // Error, (ù{\codeincomments{v1}}ù) is not initialized.
constexpr int v2 = 17;   // OK
constexpr int v3 = g();  // Error, (ù{\codeincomments{g()}}ù) is not (ù{\codeincomments{constexpr}}ù).
constexpr int v4 = h();  // OK

void test2(int c)
{
    constexpr int v5 = c;          // Error, (ù{\codeincomments{c}}ù) not a compile-time constant.
    constexpr int v6 = sizeof(c);  // OK, (ù{\codeincomments{c}}ù) is not evaluated.
}
\end{emcppslisting}
}
    
\item{Any variable declared \lstinline!constexpr! must be of \emcppsgloss{literal type}; all literal types are, among other things, \emcppsgloss{trivially destructible}:

\begin{emcppslisting}
struct Lt  // literal type
{
    constexpr  Lt()  { }  // (ù{\codeincomments{constexpr}}ù) constructor
    ~Lt() = default;      // (ù{\codeincomments{default}}ù) trivial destructor
};

constexpr Lt lt;  // OK, (ù{\codeincomments{Lt}}ù) is a (ù{\emphincomments{literal type}ù).

struct Nlt  // (ù{\emphincomments{non}ù)literal type.
{
    Nlt()  { }  // cannot initialize at compile-time
    ~Nlt() { }  // cannot skip non-trivial destruction
};

constexpr Nlt nlt;  // Error, (ù{\codeincomments{Nlt}}ù) is not a (ù{\emphincomments{literal type}ù).
\end{emcppslisting}
    
\noindent Since all \emcppsgloss[literal type]{literal types} are trivially destructible, the
compiler does not need to emit any special code to manage the end of the
lifetime of a \lstinline!constexpr! variable, which can essentially live
``forever'' --- i.e., until the program exits.{\cprotect\footnote{In
C++20, literal types can have non-trivial destructors, and the
destructors for \lstinline!constexpr! variables will be invoked under the
same conditions that a destructor would be invoked for a
  non\lstinline!constexpr! global or \lstinline!static! variable.}}}

\item{Unlike integral constants, non\lstinline!static! data members cannot be \lstinline!constexpr!. Only variables at global or \lstinline!namespace! scope, automatic variables, or \lstinline!static! data members of a \lstinline!class! or \lstinline!struct! may be declared \lstinline!constexpr!. Consequently, any given \lstinline!constexpr! variable is a top-level object, never a subobject of another, possibly non\lstinline!constexpr!, object:

\begin{emcppslisting}
               constexpr int i = 17;    // OK, file scope
namespace ns { constexpr int j = 34; }  // OK, namespace scope

struct C
{
    static constexpr int k = 51;  // OK, (ù{\codeincomments{static}}ù) data member
           constexpr int l = 68;  // Error, (ù{\codeincomments{constexpr}}ù) non(ù{\codeincomments{static}}ù) data member
};

void g()
{
    static constexpr int m = 85;  // OK
           constexpr int n = 92;  // OK
};
\end{emcppslisting}
    
\noindent Recall, however, that \emcppsgloss[nonstatic data member]{nonstatic data members} of
\lstinline!constexpr! objects are implicitly \lstinline!constexpr! and
therefore can be used directly in any \emcppsgloss[constant expression]{constant expressions}:

\begin{emcppslisting}
constexpr struct D { int i; } x{1};  // brace-initialized aggregate (ù{\codeincomments{x}}ù)
constexpr int k = x.i;  // Subobjects of (ù{\codeincomments{constexpr}}ù) objects are (ù{\codeincomments{constexpr}}ù).
\end{emcppslisting}
}
\end{enumerate}

\subsubsection[Initializer Undefined Behavior]{Initializer Undefined Behavior}

It is important to note the significance of one of the
differences between a \lstinline!constexpr! integral variable and a
\lstinline!const! integral variable. Because the initializer of a
\lstinline!constexpr! variable is \emph{required} to be a \emcppsgloss{constant expression}, it is not subject to the possibility of undefined behavior
(e.g., integer overflow or out-of-bounds array access) at run time and will
instead result in a compile-time error:

\begin{emcppslisting}
          const int iA = 1 << 15;  // 2^15 = 32,768 fits in 2 bytes.
          const int jA = iA * iA;  // OK

          const int iB = 1 << 16;  // 2^16 = 65,536 doesn't fit in 2 bytes.
          const int jB = iB * iB;  // Bug, overflow (might warn)

constexpr const int iC = 1 << 16;
constexpr const int jC = iC * iC;  // Error, overflow in constant expression

constexpr       int iD = 1 << 16;  // Example (ù{\codeincomments{D}}ù) is the same as (ù{\codeincomments{C}}ù), above.
constexpr       int jD = iD * iD;  // Error, overflow in constant expression
\end{emcppslisting}
    
\noindent The code example above shows that an integer constant-expression
overflow, absent \lstinline!constexpr!, is not required by the C++ Standard
to be treated as ill formed, whereas the initializer for a
\lstinline!constexpr! expression is required to report overflow as an error
(not just a warning).

There is a strong association between \lstinline!constexpr! variable and
functions; see\linebreak[4]%%%%%  
\featureref{\locationc}{constexprfunc}. Using a
\lstinline!constexpr! variable rather than just a \lstinline!const! one forces the
compiler to detect overflow within the body of \lstinline!constexpr!
function and report it --- as a compile-time error --- in a way that it
would not otherwise be authorized to do.

For example, suppose we have two similar functions, \lstinline!squareA! and
\lstinline!square!, defined for the built-in type \lstinline!int! that each return the integral product of multiplying the
argument with itself:

\begin{emcppslisting}[emcppsbatch=e2]
          int squareA(int i) { return i * i; }  // non-(ù{\codeincomments{constexpr}}ù) function
constexpr int squareB(int i) { return i * i; }  // (ù{\codeincomments{constexpr}}ù) function
\end{emcppslisting}
    
\noindent Declaring a variable to be just \lstinline!const! does nothing to force the
compiler to check the evaluation of either function for overflow:

\begin{emcppslisting}[emcppsbatch=e2]
                int xA0 = squareA(1 << 15);  // OK
          const int xA1 = squareA(1 << 15);  // OK
constexpr       int xA2 = squareA(1 << 15);  // Error, (ù{\codeincomments{squareA}}ù), not (ù{\codeincomments{constexpr}}ù)

                int yB0 = squareB(1 << 15);  // OK
          const int yB1 = squareB(1 << 15);  // OK
constexpr       int yB2 = squareB(1 << 15);  // OK

                int zC0 = squareB(1 << 16);  // Bug, (ù{\codeincomments{zC0}}ù) is likely 0.
          const int zC1 = squareB(1 << 16);  // Bug, (ù{\codeincomments{zC1}}ù)  "   "    "
constexpr       int zC2 = squareB(1 << 16);  // Error, (ù{\codeincomments{int}}ù) overflow detected!
\end{emcppslisting}
    
\noindent By declaring a variable to be not just \lstinline!const! but
\lstinline!constexpr!, we make the compiler evaluate that (necessarily
\lstinline!constexpr!) function for those specific arguments at compile
time and, if there is overflow, immediately report the defect as being
\emcppsgloss{ill formed}.

\subsubsection[Internal Linkage]{Internal Linkage}

When a variable at file or namespace scope is either \lstinline!const! or \lstinline!constexpr! and nothing explicitly gives it external linkage (e.g., being marked \lstinline!extern!), it
will have \emcppsgloss{internal linkage}.  This means each translation unit will have its own copy of the
variable.{\cprotect\footnote{In C++17, all \lstinline!constexpr! variables are instead automatically \lstinline!inline! as well,
guaranteeing that there is only one instance in a program.}}

In general, only the values of such variables are relevant: their initializers are seen, they are used at compile time, and there is no effect if different translation units use different objects having the same
name and a different value.  Often, after compile-time evaluation is completed, the variables themselves will no longer be
needed, and no actual address will be allocated for them at run time.   Only in cases where the address of
the variable is used will this difference be observable.

Notably, \lstinline!static! member variables do not get \emcppsgloss{internal linkage}.  Because of this, if they
are going to be used in a way that requires they have an address allocated at run time, they need to have a
definition outside of their class; see \intraref{annoyances-constexprvar}{static-member-variables-require-external-definitions}.

\subsection[Use Cases]{Use Cases}\label{use-cases}

\subsubsection[Alternative to enumerated compile-time integral constants]{Alternative to enumerated compile-time integral constants}\label{alternative-to-enumerated-compile-time-integral-constants}

It is not uncommon to want to express specific integral constants at
compile time --- e.g., for precomputed operands to be used in
algorithms, mathematical constants, configuration variables, or any
number of other reasons. A naive, brute-force approach might be to
hard-code the constants where they are used:

\begin{emcppslisting}
int hoursToSeconds0(int hours)
    // Return the number of seconds in the specified (ù{\codeincomments{hours}}ù).  The behavior is
    // undefined unless the result can be represented as an (ù{\codeincomments{int}}ù).
{
    return hours * 3600;
}
\end{emcppslisting}
    
\noindent This use of \emph{magic constants} has, however, long been
known\footnote{\cite{kernighan99}, section~1.5, pp.~19--22} to make finding uses of the constants
and the relationships between related ones needlessly difficult. For
integral values only, we could always represent such compile-time
constants symbolically by using a classic \lstinline!enum! type, in
deliberate preference to the modern, type-safe enumerator; see \featureref{\locationc}{enumclass}:

\begin{emcppslisting}
struct TimeRatios1  // explicit scope for single classic anonymous (ù{\codeincomments{enum}}ù) type
{
    enum  // anonymous enumeration comprising related symbolic constants
    {
        k_SECONDS_PER_MINUTE = 60,     // Underlying type (UT) (ù{\emphincomments{might}ù) be (ù{\codeincomments{int}}ù).
        k_MINUTES_PER_HOUR   = 60,
        k_SECONDS_PER_HOUR   = 60*60,  // these enumerators have the same UT
    };
};

int hoursToSeconds1(int hours)
    // ...
{
    return hours * TimeRatios1::k_SECONDS_PER_HOUR;
}
\end{emcppslisting}
    
\noindent This traditional solution, while often effective, gives little control to
the underlying integral type of the enumerator used to represent the
symbolic compile-time constant, leaving it at the mercy of the totality
of values used to initialize its members. Such inflexibility might lead
to compiler warnings and nonintuitive behavior resulting from
\lstinline!enum!-specific ``integral-promotion'' rules, especially when the
\emcppsgloss[underlying type]{underlying type (UT)} used to represent the time ratios differs
from the integral type with which they are ultimately used; see \featureref{\locationc}{explicit-enumeration-underlying-type}.

In this particular example, extending the \lstinline!enum! to cover ratios
up to a week and conversions down to nanoseconds would manifestly change
its \emcppsgloss[underlying type (UT)]{UT} (because there are far more than $2^{32}$ nanoseconds in a week),
altering how all of the enumerators behave when used in expressions
with, say, values of type \lstinline!int! (e.g., to \lstinline!long!); see \featureref{\locationa}{long-long}:

\begin{emcppslisting}
struct TimeRatios2  // explicit scope for single classic anonymous (ù{\codeincomments{enum}}ù) type
{
    enum  // Anonymous enumeration --- UT is governed by all of the enumerators.
    {
        k_SECONDS_PER_MINUTE = 60,    // UT (ù{\emphincomments{might}ù) be (ù{\codeincomments{long}}ù) (or (ù{\codeincomments{long long}}ù)).
        k_MINUTES_PER_HOUR   = 60,
        k_SECONDS_PER_HOUR   = 3600,
        // ...
        k_USEC_PER_WEEK = 1000L*1000*60*60*24*7,  // same UT as all of the above
    };
};
\end{emcppslisting}
    
\noindent The original \emph{values} will remain unchanged, but the burden of all
of the warnings resulting from the change in \emcppsgloss[underlying type (UT)]{UT} and rippling throughout a large codebase could be expensive to
repair.

We would like the original values to remain unchanged (e.g., remain as
\lstinline!int! if that's what they were), and we want only those values
that do \emph{not} fit in an \lstinline!int! to morph into a larger
integral type. We might achieve this effect by placing each enumerator
in its own separate anonymous enumeration:

\begin{emcppslisting}
struct TimeRatios3  // explicit scope for multiple classic anonymous (ù{\codeincomments{enum}}ù) types
{
    enum { k_SECONDS_PER_MINUTE = 60               };  // UT: (ù{\codeincomments{int}}ù) (likely)
    enum { k_MINUTES_PER_HOUR   = 60               };  //  "    "      "
    enum { k_SECONDS_PER_HOUR   = 60*60            };  //  "    "      "
    // ...
    enum { k_USEC_PER_SEC  = 1000*1000             };  // UT: (ù{\codeincomments{int}}ù) (v. likely)
    enum { k_USEC_PER_MIN  = 1000*1000*60          };  // UT: (ù{\codeincomments{int}}ù) (v. likely)
    enum { k_USEC_PER_HOUR = 1000U*1000*60*60      };  // UT: (ù{\codeincomments{unsigned int}}ù)
    enum { k_USEC_PER_DAY  = 1000L*1000*60*60*24   };  // UT: (ù{\codeincomments{long}}ù)
    enum { k_USEC_PER_WEEK = 1000L*1000*60*60*24*7 };  // UT: (ù{\codeincomments{long}}ù)
};
\end{emcppslisting}
    
\noindent In this case, the original values as well as their respective \emcppsgloss[underlying type (UT)]{UTs} will
remain unchanged, and each new enumerated value will independently take
on its own independent \emcppsgloss[underlying type (UT)]{UT}, which is either implementation defined or else
dictated by the number of bits required to represent the value, which
is, in this case, non-negative.

A modern alternative to having separate anonymous \lstinline!enum!s for
each distinct value (or class of values) is to instead encode each ratio
as an explicitly typed \lstinline!constexpr! variable:

\begin{emcppslisting}
struct TimeRatios4
{
    static constexpr int k_SECONDS_PER_MINUTE = 60;
    static constexpr int k_MINUTES_PER_HOUR   = 60;
    static constexpr int k_SECONDS_PER_HOUR   = k_MINUTES_PER_HOUR *
                                                k_SECONDS_PER_MINUTE;
    // ...
    static constexpr long k_NANOS_PER_SECOND  = 1000*1000*1000;
    static constexpr long k_NANOS_PER_HOUR    = k_NANOS_PER_SECOND *
                                                k_SECONDS_PER_HOUR;
};

int hoursToSeconds(int hours)
    // ...
{
    return hours * TimeRatios4::k_SECONDS_PER_HOUR;
}

long hoursToNanos(int hours)
    // Return the number of nanoseconds in the specified (ù{\codeincomments{hours}}ù).  The behavior
    // is undefined unless the result can be represented as a (ù{\codeincomments{long}}ù).
{
    return hours * TimeRatios4::k_NANOS_PER_HOUR;
}
\end{emcppslisting}
    
\noindent In the example above, we've rendered the \lstinline!constexpr! variables as
\lstinline!static! members of a \lstinline!struct! rather than placing them at
namespace scope primarily to show that, from a user perspective, the two
are syntactically indistinguishable --- the substantive difference here
being that a client would be prevented from unilaterally adding logical
content to the ``namespace'' of a \lstinline!TimeRatio!~\lstinline!struct!.
When it comes to C-style \emcppsgloss[free function]{free functions}, however, the
advantages for \lstinline!static! members of a struct over namespace scope
are many and \mbox{unequivocal}.{\cprotect\footnote{\cite{lakos20}, section~2.4.9, pp.~312--321, specifically Figure 2-23}}

\subsubsection[Nonintegral symbolic numeric constants]{Nonintegral symbolic numeric constants}\label{non-integral-symbolic-numeric-constants}

Not all symbolic numeric constants that are needed at compile time are
necessarily integral. Consider, for example, the mathematical constants
\lstinline!pi! and \lstinline!e!, which are typically represented as a
floating-point type, such as \lstinline!double! (or
\lstinline!long!~\lstinline!double!).

The classical solution to avoid encoding this type of constant values as
a \emph{magic number} is to instead use a macro, such as is done in the
\lstinline!math.h! header on most operating systems:

\begin{emcppslisting}
#define M_E   2.7182818284590452354  /* e */
#define M_PI  3.14159265358979323846 /* pi */

double areaOfCircle(double radius)
{
    return 2 * M_PI * radius;
}
\end{emcppslisting}
    
\noindent While this approach can be effective, it comes with all the well-known
downsides of using the C preprocessor. This approach is also fraught
with risk such that the headers standard on many systems make these
macros available in C or C++ only upon request by \lstinline!#define!ing
specific macros before including those headers.\footnote{See your library
documentation for details.}

Another safer and far less error-prone solution to name collisions is to
instead use a \lstinline!constexpr! variable for this form of nonintegral
constant. Note that, while the macro is defined with more precision to
be able to initialize variables of possibly higher-precision
floating-point types, here we need only enough digits to uniquely
identify the appropriate \lstinline!double! constant:

\begin{emcppslisting}[emcppsstandards={c++14}]
struct NumericConstants
{
    static constexpr double k_E  = 2.718'281'828'459'045;  /* e */
    static constexpr double k_PI = 3.141'592'653'589'793;  /* pi */
};

double areaOfCircle(double radius)
{
    return 2 * NumericConstants::k_PI * radius;
}
\end{emcppslisting}
    
\noindent In the example above, we have made valuable use of a safe C++14
feature to help identify the needed precision of the numeric literal;
see \featureref{\locationb}{digitseparator}. Beyond the potential name
collisions and global name pollution, preferring a constexpr variable
over a C preprocessor macro has the added benefit of making explicit the
C++ type of constant being defined. Note that supplying digits beyond
what are significant will nonetheless be silently ignored.

\subsubsection[Storing \lstinline!constexpr! data structures]{Storing {\SubsubsecCode constexpr} data structures}\label{storing-constexpr-data-structures}

Precomputing values at compile time for subsequent use at run time
is one impactful use of \lstinline!constexpr! \emph{functions}, but see \intraref{potential-pitfalls-constexprvar}{constexpr-function-pitfalls}. Storing these values in explicitly
\lstinline!constexpr! variables ensures that the values are (1) guaranteed
to be computed at compile time and not, for example, at startup as the
result of a dangerous runtime initialization of a file- or
namespace-scoped variable and (2) usable as part of the evaluation of
any \emcppsgloss{constant expression}; see \featureref{\locationc}{constexprfunc}.

Rather than attempting to circumvent the draconian limitations of the
C++11 version of \lstinline!constexpr! functions, we will make use of the
relaxed restrictions of C++14. For this, we will define a class template
that initializes an array member with the results of a
\lstinline!constexpr! functor applied to each array
index{\cprotect\footnote{Note that, in C++17, most of the manipulators
of \lstinline!std::array! have been changed to be \lstinline!constexpr! and,
when combined with the relaxation of the rules for \lstinline!constexpr!
evaluation in C++14 (see \featureref{\locationd}{relaxed-constexpr-restrictions}), this compile-time-friendly container provides a simple way to
  define functions that populate tables of values.}}:

\begin{emcppshiddenlisting}[emcppsbatch=e3,emcppsstandards={c++14}]
#include <cstddef>  // (ù{\codeincomments{std::size\_t}}ù)
\end{emcppshiddenlisting}
\begin{emcppslisting}[emcppsbatch=e3]
template <typename T, std::size_t N>
struct ConstexprArray
{
private:
    T d_data[N];  // data initialized at construction

public:
    template <typename F>
    constexpr ConstexprArray(const F &func)
    : d_data{}
    {
        for (int i = 0; i < N; ++i)
        {
            d_data[i] = func(i);
        }
    }

    constexpr const T& operator[](std::size_t ndx) const
    {
        return d_data[ndx];
    }
};
\end{emcppslisting}
    
\noindent The numerous alternative approaches to writing such data structures
vary in their complexity, trade-offs, and understandability. In this
case, we default initialize our elements before populating them but do
not need to rely on any other significant new language infrastructure.
Other approaches could be taken; see \featureref{\locationc}{constexprfunc}.

Given this utility class, we can then precompute at compile time any
function that we can express as a \lstinline!constexpr! function, such as a
simple table of the first \lstinline!N! squares:

\begin{emcppslisting}[emcppsbatch=e3]
constexpr int square(int x) { return x * x; }

constexpr ConstexprArray<int, 500> squares(square);

static_assert(squares[1]   == 1,    "");
static_assert(squares[289] == 83521,"");
\end{emcppslisting}
    
\noindent Note that, as with many applications of \lstinline!constexpr! functions,
attempting to initialize a large array of \lstinline!constexpr! variables
will quickly bump up against compiler-imposed template instantiation
limits. What's more, attempting to perform more complex arithmetic at
compile time would be likely to exceed computation limits as well.

\subsubsection[Diagnosing undefined behavior at compile time]{Diagnosing undefined behavior at compile time}\label{diagnosing-undefined-behavior-at-compile-time}

Avoiding overflow during intermediate calculations is an important
consideration, especially from a security perspective, and yet is a
generally difficult-to-solve problem. Forcing computations to occur at
compile time brings the full power of the compiler to bear in addressing
such undefined behavior.

As an academically interesting example of this eminently practical
security problem, suppose we want to write a compile-time function in
C++ to compute the \emcppsgloss[Collatz length]{Collatz length} of an arbitrary positive
integer and generate a compilation error if any intermediate calculation
would result in signed integer overflow.

First let's take a step back now to understand what we are talking about
here with respect to \emcppsgloss[Collatz length]{Collatz length}. Suppose we have a function
\lstinline!cf! that takes a positive \lstinline!int!, \lstinline!n!, and for even
\lstinline!n! returns \lstinline!n/2! and for odd \lstinline!n! returns
\lstinline!3n+1!.

\begin{emcppslisting}
int cf(int n) { return n % 2 ? 3 * n + 1 : n / 2; }  // Collatz function
\end{emcppslisting}
    
\noindent Given a positive integer, \lstinline!n!, the \emcppsgloss[Collatz sequence]{Collatz sequence},
\lstinline!cs(n)!, is defined as the sequence of integers generated by
repeated application of the \emcppsgloss[Collatz function]{Collatz function} --- e.g.,\linebreak[4]%%%%
 \lstinline!cs(1)!~\lstinline!=!~\lstinline!{!~\lstinline!4,!~\lstinline!2,!~\lstinline!1,!~\lstinline!4,!~\lstinline!2,!~\lstinline!1,!~\lstinline!4,!~\lstinline!...!~\lstinline!};!,
\lstinline!cs(3)!~\lstinline!=!~\lstinline!{!~\lstinline!10,!~\lstinline!5,!~\lstinline!16,!~\lstinline!8,!~\lstinline!4,!~\lstinline!2,!~\lstinline!1,!~\lstinline!4,!~\lstinline!...!~\lstinline!}!,
and so on. A classic (but as yet unproven) conjecture in mathematics
states that, for every positive integer, $n$, the \emcppsgloss[Collatz sequence]{Collatz sequence} for $n$ will eventually reach 1. The \emcppsgloss[Collatz length]{Collatz length} of
the positive integer $n$ is the number of iterations of the
\emcppsgloss[Collatz function]{Collatz function} needed to reach 1, starting from $n$.~Note that
the \emcppsgloss[Collatz sequence]{Collatz sequence} for \lstinline!n!~\lstinline!=!~\lstinline!1! is
\lstinline!{!~\lstinline!1,!~\lstinline!1,!~\lstinline!1,!~\lstinline!...!~\lstinline!}!
and its \emcppsgloss[Collatz length]{Collatz length} is 0.

This example showcases the need for a \lstinline!constexpr! variable in
that we need to create a \emcppsgloss{constexpr context} -- i.e., one
requiring a \emcppsgloss{constant expression} to require that the
evaluation of a \lstinline!constexpr! function occur at compile time.
Again, to avoid distractions related to implementing more complex
functionality within the limitations of C++11 \lstinline!constexpr!
functions, we will make use of the relaxed restrictions of C++14; see
\featureref{\locationc}{constexprfunc}:

\begin{emcppslisting}[emcppsstandards={c++14}]
constexpr int collatzLength(long long number)
    // Return the length of the Collatz sequence of the specified (ù{\codeincomments{number}}ù). The
    // behavior is undefined unless each intermediate sequence member can be
    // expressed as an (ù{\codeincomments{unsigned long long}}ù).
{
    int length = 0;         // collatLength(1) is 0.

    while (number > 1)      // The current value of number is not 1.
    {
        ++length;           // Keep track of the length of the sequence so far.

        if (number % 2)     // if the current number is odd
        {
            number = 3 * number + 1;    // advance from odd sequence value
        }
        else
        {
            number /= 2;                // advance from even sequence value
        }
    }

    return length;
}

const     int c1 = collatzLength(942488749153153);  // OK, 1862
constexpr int x1 = collatzLength(942488749153153);  // OK, 1862

const int c2 = collatzLength(104899295810901231);
    // Bug, program aborts at runtime.

constexpr int x2 = collatzLength(104899295810901231);
    // Error, overflow in constant evaluation
\end{emcppslisting}
    
\noindent In the example above, the variables \lstinline!c1! and \lstinline!x1! can be
initialized correctly at compile time, but \lstinline!c2! and \lstinline!x2!
cannot. The non\lstinline!constexpr! nature of \lstinline!c2! allows the
overflow to occur and exhibit undefined behavior --- integer overflow
--- at run time. On the other hand, the variable \lstinline!x2!, due to
its being declared \lstinline!constexpr!, forces the computation to
occur at compile time, thereby discovering the overflow and dutifully
invoking the \emph{non}\lstinline!constexpr!~\lstinline!std::abort! function,
which in turn generates the desired error at compile time.

\subsection[Potential Pitfalls]{Potential Pitfalls}\label{potential-pitfalls-constexprvar}

\subsubsection[\lstinline!constexpr! function pitfalls]{{\SubsubsecCode constexpr} function pitfalls}\label{constexpr-function-pitfalls}

Many of the uses of \lstinline!constexpr! variables involve a corresponding
use of \lstinline!constexpr! functions (see \featureref{\locationc}{constexprfunc}). The pitfalls related to \lstinline!constexpr! functions are
similarly applicable to the variables in which their results might be
stored. In particular, it can be profoundly advantageous to forgo use of
the \lstinline!constexpr! function feature altogether, do the
precomputation externally to program, and embed the calculated result
--- along with a comment containing source text of the (e.g., Perl or
Python) script that performed the calculation --- into the C++ program
source itself.

\subsection[Annoyances]{Annoyances}\label{annoyances-constexprvar}

\subsubsection[\lstinline!static! member variables require external definitions]{{\SubsubsecCode static} member variables require external definitions}\label{static-member-variables-require-external-definitions}

In most situations there is little behavioral difference between a variable at file or namespace scope and a
\lstinline!static! member variable; primarily they differ only in name lookup and access control restrictions.
When they are \lstinline!constexpr!, however, they behave very differently when they need to exist
at runtime.  A file or namespace scope variable will have \emcppsgloss{internal linkage}, allowing free use
of its address with the understanding that that address will be different in different translation units:

\begin{emcppslisting}
// common.h:

constexpr int c = 17;
const int *f1();
const int *f2();

// file1.cpp:
#include <common.h>

const int *f1() { return &c; }

// file2.cpp:
#include <common.h>

const int *f2() { return &c; }

// main.cpp:
#include <common.h>
#include <cassert>  // standard C (ù{\codeincomments{assert}}ù) macro

int main()
{
    assert( f1() != f2() );    // Different addresses in memory per TU.
    assert( *f1() == *f2() );  // Same value.
    return 0;
}
\end{emcppslisting}

For \lstinline!static! data members, however, things become more difficult.  While the \emcppsgloss{declaration} in the
class \emcppsgloss{definition} needs to have an initializer, that is not itself a \emcppsgloss{definition} and will
not result in static storage being allocated at runtime for the object, ending in a linker error when you try to
build an application that tries to reference it{\cprotect\footnote{The C++17 change to make \lstinline!constexpr! variables \lstinline!inline! also applies
to \lstinline!static! member variables, removing the need to provide external definitions when they are used.}:


\begin{emcppslisting}[emcppsbatch=e4]
struct S {
    static constexpr int d_i = 17;
};
void useByReference(const int& i) { /* ... */ }

int main()
{
    const int local = S::d_i;  // OK, value is only used at compile time.
    useByReference(S::d_i);    // Link-Time Error, (ù{\codeincomments{S::d\_i}}ù) not defined
    return 0;
}
\end{emcppslisting}

This link-time error would be removed by adding a \emcppsgloss{definition} of \lstinline!S::d_i!.  Note that the
initializer needs to be skipped, as it has already been specified in the definition of \lstinline!S!:

\begin{emcppslisting}[emcppsbatch=e4]
constexpr int S::d_i;  // Define (ù{\codeincomments{constexpr}}ù) data member.
\end{emcppslisting}

\subsubsection[No static \lstinline!constexpr! member variables defined in their own class]{No static {\SubsubsecCode constexpr} member variables defined in their own class}\label{no-static-constexpr-member-variables-defined-in-their-own-class}

When implementing a class using the singleton pattern, it may seem
desirable to have the single object of that type be a
\lstinline!constexpr!~\lstinline!private!~\lstinline!static! member of the class
itself, with guaranteed compile-time (data-race-free) initialization and
no direct accessibility outside the class. This does not work as easily
as planned because \lstinline!constexpr!~\lstinline!static! data members must
have a complete type and the class being defined is not complete until
its closing brace:

\begin{emcppslisting}[emcppsbatch=e5]
class S
{
    static const     S constVal;      // OK, initialized outside class below
    static constexpr S constexprVal;  // Error, (ù{\codeincomments{constexpr}}ù) must be initialized.
    static constexpr S constInit{};   // Error, (ù{\codeincomments{S}}ù) is not complete.
};

const S S::constVal{};  // OK, initialize static (ù{\codeincomments{const}}ù) member.
\end{emcppslisting}
    
\noindent The ``obvious'' workaround of applying a more traditional singleton
pattern, where the singleton object is a static local variable in a
function call, also fails (see \featureref{\locationa}{thread-safe-function-static-variables}) because \lstinline!constexpr! functions are not allowed to have
static variables (see \featureref{\locationc}{constexprfunc}):

\begin{emcppslisting}[emcppsbatch=e5,emcppsignore={Invalid function}]
constexpr const S& singleton()
{
   static constexpr S object{};  // Error, even in C++14, (ù{\codeincomments{static}}ù) is not allowed.
   return object;
}
\end{emcppslisting}
    
\noindent The only fully featured solution available for \lstinline!constexpr!
objects of static storage duration is to put them outside of their type,
either at global scope, at \lstinline!namespace! scope, or nested within a
befriended helper class{\cprotect\footnote{C++20 provides an alternate
partial solution with the \lstinline!constinit! keyword, allowing for
compile-time initialization of static data members, but that still
does not make such objects usable in a \emcppsgloss{constant expression}.}}:

\begin{emcppslisting}
class S
{
    friend struct T;
    S() = default;  // (ù{\codeincomments{private}}ù)
    // ...
};

struct T
{
    static constexpr S constexprS{};
};
\end{emcppslisting}
    

\subsection[See Also]{See Also}\label{see-also}

\begin{itemize}
\item{\seealsoref{constexprfunc}{\seealsolocationc}can be used to initialize a \lstinline!constexpr! variable.}
\item{\seealsoref{userdeflit}{\seealsolocationc}provides a convenient way of initializing a \lstinline!constexpr! variable of a UDT with a compile-time value.}
\end{itemize}

\subsection[Further Reading]{Further Reading}\label{further-reading}

\hspace*{\fill}



\newpage
%\section[Default Member Init]{Default {\SecCode class}/{\SecCode union} Member Initializers}\label{Default-Member-Initializers}


\emcppsFeature{
    short={Default Member Init},
    long={Default {\SecCode class}/{\SecCode union} Member Initializers},
    toclong={Default \lstinline!class!/\lstinline!union! Member Initializers},
}{Default-Member-Initializers}
\setcounter{table}{0}
\setcounter{footnote}{0}
\setcounter{lstlisting}{0}
%\section[Default Member Init]{Default {\SecCode class}/{\SecCode union} Member Initializers}\label{Default-Member-Initializers}


placeholder text.........




\newpage
%\section[{\tt enum} {\tt class}]{Strongly Typed Scoped Enumerations}\label{enumclass}
% 01 Feb, 2021 JMB - packet 4 compilation fixes
% 10 Feb 2021, cleaned up a bit before sending for revisions



\emcppsFeature{
    short={\lstinline!enum!~\lstinline!class!},
    tocshort={{\TOCCode enum}~{\TOCCode class}},
    long={Strongly Typed Scoped Enumerations},
    rhshort={{\RHCode enum}~{\RHCode class}},
}{enumclass}
\setcounter{table}{0}
\setcounter{footnote}{0}
\setcounter{lstlisting}{0}
%\section[{\tt enum} {\tt class}]{Strongly Typed Scoped Enumerations}\label{enumclass}

\lstinline!enum!~\lstinline!class! is an alternative to the classic
\lstinline!enum! construct that simultaneously provides both stronger
typing and an enclosing scope for its enumerated values.

\subsection[Description]{Description}\label{description-enumclass}

Classic, C-style enumerations are useful and continue to fulfill
important engineering needs:

\begin{emcppslisting}
enum EnumName { e_Enumerator0 /*= value0 */, e_EnumeratorN /* = valueN */ };
    // classic, C-style (ù{\codeincomments{enum}}ù): enumerators are neither type-safe nor scoped
\end{emcppslisting}

\noindent For more examples where the classic \lstinline!enum! shines, see \intraref{potential-pitfalls-enumclass}{strong-typing-of-an-enum-class-can-be-counterproductive} and \intraref{annoyances-enumclass}{scoped-enumerations-do-not-necessarily-add-value}. 
%\textit{\titleref{potential-pitfalls-enumclass}: \titleref{strong-typing-of-an-enum-class-can-be-counterproductive}} on page~\pageref{strong-typing-of-an-enum-class-can-be-counterproductive} and \textit{\titleref{annoyances-enumclass}: \titleref{scoped-enumerations-do-not-necessarily-add-value}} on page~\pageref{scoped-enumerations-do-not-necessarily-add-value}. 
Still,
innumerable practical situations occur in which enumerators that are
both scoped and more type-safe would be preferred; see \intrarefsimple{introducing-the-c++11-enum-class}.
%\textit{\titleref{introducing-the-c++11-enum-class}} on page~\pageref{introducing-the-c++11-enum-class}.

\subsubsection[Drawbacks and workarounds relating to unscoped C++03 enumerations]{Drawbacks and workarounds relating to unscoped C++03 enumerations}\label{drawbacks-and-workarounds-relating-to-unscoped-c++03-enumerations}

Since the enumerators of a classic \lstinline!enum! leak out into the
enclosing scope, if two unrelated enumerations that happen to use the
same enumerator name appear in the same scope, an ambiguity could ensue:

\begin{emcppslisting}
enum Color { e_RED, e_ORANGE, e_YELLOW };    // OK
enum Fruit { e_APPLE, e_ORANGE, e_BANANA };  // Error, (ù{\codeincomments{e\_ORANGE}}ù) is redefined.
\end{emcppslisting}

\noindent Note that
we use a lowercase, single-letter prefix, such as \lstinline!e_!, to
ensure that the uppercase enumerator name is less likely to collide
  with a legacy macro, which is especially useful in header files. The problems associated with the use of unscoped enumerations is
exacerbated when those enumerations are placed in their own respective
header files in the global or some other large namespace scope, such as
\lstinline!std!, for general reuse. In such cases, latent defects will
typically not manifest unless and until the two enumerations are
included in the same translation unit.

If the only issue were the leakage of the enumerators into the enclosing
scope, then the long-established workaround of enclosing the enumeration
within a \lstinline!struct! would suffice:

\begin{emcppslisting}[emcppsbatch=e1]
struct Color { enum Enum { e_RED, e_ORANGE, e_YELLOW };  };   // OK
struct Fruit { enum Enum { e_APPLE, e_ORANGE, e_BANANA }; };  // OK (scoped)
\end{emcppslisting}

\noindent Employing the C++03 workaround in the above code snippet implies that,
when passing such an explicitly scoped, classical \lstinline!enum! into a
function, the distinguishing name of the \lstinline!enum! is subsumed by
its enclosing \lstinline!struct! and the \lstinline!enum! name itself, such as
\lstinline!Enum!, becomes \romeogloss{boilerplate code}:

\begin{emcppslisting}[emcppsbatch=e1]
int enumeratorValue1 = Color::e_ORANGE;  // OK
int enumeratorValue2 = Fruit::e_ORANGE;  // OK

void colorFunc(Color::Enum color);  // enumerated (scoped) (ù{\codeincomments{Color}}ù) parameter
void fruitFunc(Fruit::Enum fruit);  // enumerated (scoped) (ù{\codeincomments{Fruit}}ù) parameter
\end{emcppslisting}

\noindent Hence, adding \emph{just} scope to a classic, C++03 \lstinline!enum! is
easily doable and might be exactly what is indicated; see \intraref{potential-pitfalls-enumclass}{strong-typing-of-an-enum-class-can-be-counterproductive}. 
%\textit{\titleref{potential-pitfalls-enumclass}: \titleref{strong-typing-of-an-enum-class-can-be-counterproductive}} on page~\pageref{strong-typing-of-an-enum-class-can-be-counterproductive}.

\subsubsection[Drawbacks relating to weakly typed, C++03 enumerators]{Drawbacks relating to weakly typed, C++03 enumerators}\label{drawbacks-relating-to-weakly-typed,-c++03-enumerators}

Historically, C++03 enumerations have been employed to represent at
least two distinct concepts:

\begin{enumerate}
\item{A collection of related, but not necessarily unique, named integral values}
\item{A pure, perhaps ordered, set of named entities in which cardinal value has no relevance}
\end{enumerate}

\noindent It will turn out that the modern \lstinline!enum!~\lstinline!class! feature,
which we will discuss in \intraref{description-enumclass}{introducing-the-c++11-enum-class}, 
%\textit{\titleref{description-enumclass}: \titleref{introducing-the-c++11-enum-class}}, 
is more closely aligned with this second
concept.

A classic enumeration, by default, has an implementation-defined
\romeogloss{underlying type} (see \featureref{\locationc}{explicit-enumeration-underlying-type},
%``\titleref{explicit-enumeration-underlying-type}" on page~\pageref{explicit-enumeration-underlying-type}), 
which it uses to represent
variables of that enumerated type as well as the values of its
enumerators. While implicit conversion \emph{to} an enumerated type is
never permitted, when implicitly converting \emph{from} a classical
\lstinline!enum! type to some arithmetic type, the \lstinline!enum! promotes
to integral types in a way similar to how its underlying type would
promote using the rules of \romeogloss{integral promotion} and
\romeogloss{standard conversion}:

\begin{emcppslisting}
void f()
{
    enum A { e_A0, e_A1, e_A2 };  // classic, C-style C++03 (ù{\codeincomments{enum}}ù)
    enum B { e_B0, e_B1, e_B2 };  //    "       "       "     "

    A a;  // Declare object (ù{\codeincomments{a}}ù) to be of type (ù{\codeincomments{A}}ù).
    B b;  //   "       "    (ù{\codeincomments{b}}ù)  " "  "       (ù{\codeincomments{B}}ù).

    a = e_B2;  // Error, cannot convert (ù{\codeincomments{e\_B2}}ù) to (ù{\codeincomments{enum}}ù) type (ù{\codeincomments{A}}ù)
    b = e_B2;  // OK, assign the value (ù{\codeincomments{e\_B2}}ù) (numerically 2) to (ù{\codeincomments{b}}ù).
    a = b;     // Error, cannot convert enumerator (ù{\codeincomments{b}}ù) to (ù{\codeincomments{enum}}ù) type (ù{\codeincomments{A}}ù)
    b = b;     // OK, self-assignment
    a = 1;     // Error, invalid conversion from (ù{\codeincomments{int 1}}ù) to (ù{\codeincomments{enum}}ù) type (ù{\codeincomments{A}}ù)
    a = 0;     // Error, invalid conversion from (ù{\codeincomments{int 0}}ù) to (ù{\codeincomments{enum}}ù) type (ù{\codeincomments{A}}ù)

    bool     v = a;     // OK
    char     w = e_A0;  // OK
    unsigned y = e_B1;  // OK
    float    x = b;     // OK
    double   z = e_A2;  // OK
    char*    p = e_B0;  // Error, unable to convert (ù{\codeincomments{e\_B0}}ù) to (ù{\codeincomments{char*}}ù)
    char*    q = +e_B0; // Error, invalid conversion of (ù{\codeincomments{int}}ù) to (ù{\codeincomments{char*}}ù)
}
\end{emcppslisting}

\noindent Notice that, in this example, the final two diagnostics for the
attempted initializations of \lstinline!p! and \lstinline!q!, respectively,
differ slightly. In the first, we are trying to initialize a pointer,
\lstinline!p!, with an enumerated type, \lstinline!B!. In the second, we have
creatively used the built-in unary-plus operator to explicitly promote
the enumerator to an integral type before attempting to assign it to a
pointer, \lstinline!q!. Even though the numerical value of the enumerator
is \lstinline!0! and such is known at compile time, implicit
conversion to a pointer type from anything but the literal integer
constant \lstinline!0! is not
permitted. Excluding esoteric user-defined
types, only a literal \lstinline!0! or, as of C++11, a value of type
\lstinline!std::nullptr_t! is implicitly convertible to an arbitrary
  pointer type; see \featureref{\locationa}{null-pointer-literal-(nullptr)}.
%  ``\titleref{null-pointer-literal-(nullptr)}" on page~\pageref{null-pointer-literal-(nullptr)}.

C++ fully supports comparing values of \emph{classic} \lstinline!enum!
types with values of arbitrary \romeogloss{arithmetic type} as well as those
of the same enumerated type; the operands of a comparator will be
promoted to a sufficiently large integer type and the comparison will be
done with those values. Comparing values having distinct enumerated
types, however, is deprecated and will typically elicit a
warning.{\cprotect\footnote{As of C++20, attempting to compare two values of distinct classically enumerated
types is a compile-time error. Note that explicitly converting at least one of them to an
integral type --- for example, using built-in unary plus --- both
makes our intentions clear and avoids warnings.

\begin{emcppshiddenlisting}[emcppsbatch=e2]
    enum A { e_A0, e_A1, e_A2 };
    enum B { e_B0, e_B1, e_B2 };
\end{emcppshiddenlisting}
\begin{emcppslisting}[emcppsbatch=e2,style=footcode]
void test()
{
  if (e_A0 < 0)       { /* ... */ }  // OK, comparison with integral type
  if (1.0 != e_B1)    { /* ... */ }  // OK, comparison with arithmetic type
  if (A() <= e_A2)    { /* ... */ }  // OK, comparison with same enumerated type
  if (e_A0 == e_B0)   { /* ... */ }  // warning, deprecated (error as of C++20)
  if ( e_A0 == +e_B0) { /* ... */ }  // OK, unary (ù{\codeincomments{+}}ù) converts to integral type
  if (+e_A0 ==  e_B0) { /* ... */ }  // OK,   "        "     "     "      "
  if (+e_A0 == +e_B0) { /* ... */ }  // OK,   "        "     "     "      "
}
\end{emcppslisting}
      }}

\subsubsection[Introducing the C++11 \lstinline!enum! \lstinline!class!]{Introducing the C++11 {\SubsubsecCode enum class}}\label{introducing-the-c++11-enum-class}

With the advent of modern C++, we now have a new, alternative
enumeration construct, \lstinline!enum!~\lstinline!class!, that simultaneously
addresses strong type safety and lexical scoping, two distinct and often
desirable properties:

\begin{emcppslisting}
enum class Name { e_Enumerator0 /* = value0 */, e_EnumeratorN /* = valueN */ };
    // (ù{\codeincomments{enum class}}ù) enumerators are both type-safe and scoped
\end{emcppslisting}

\noindent Another major distinction is that the default
\romeogloss{underlying type} for a C-style \lstinline!enum! is
\romeogloss{implementation defined}, whereas, for an
\lstinline!enum!~\lstinline!class!, it is always an \lstinline!int!. See \intraref{description-enumclass}{enum-class-and-underlying-type} and \intraref{potential-pitfalls-enumclass}{external-use-of-opaque-enumerators-enumclass}. 
%\textit{\titleref{description-enumclass}: \titleref{enum-class-and-underlying-type}} on page~\pageref{enum-class-and-underlying-type} and
%\textit{\titleref{potential-pitfalls-enumclass}: \titleref{external-use-of-opaque-enumerators-enumclass}} on page~\pageref{external-use-of-opaque-enumerators-enumclass}.

The enumerators within an \lstinline!enum!~\lstinline!class! are all scoped by
its name, while classic enumerations leak the enumerators into the
enclosing scope:

\begin{emcppslisting}
enum Vehicle  { e_CAR,   e_TRAIN, e_PLANE };
enum Geometry { e_POINT, e_LINE,  e_PLANE };  // Error, (ù{\codeincomments{e\_PLANE}}ù) is redefined.
\end{emcppslisting}

\noindent Unlike unscoped enumerations, \lstinline!enum!~\lstinline!class! does not leak
its enumerators into the enclosing scope and can therefore help avoid
collisions with other enumerations having like-named enumerators defined
in the same scope:

\begin{emcppslisting}[emcppsbatch=e3]
enum       VehicleUnscoped  { e_CAR, e_TRAIN, e_PLANE };
struct     VehicleScopedExplicitly { enum Enum { e_CAR, e_TRAIN, e_PLANE }; };
enum class VehicleScopedImplicitly { e_CAR, e_BOAT,  e_PLANE };
\end{emcppslisting}

\noindent Just like an unscoped \lstinline!enum! type, an object of type
\lstinline!enum!~\lstinline!class! is passed as a parameter to a function
using the enumerator name itself:

\begin{emcppslisting}[emcppsbatch=e3]
void f1(VehicleUnscoped value);          // classic enumeration passed by value
void f2(VehicleScopedImplicitly value);  // modern enumeration passed by value
\end{emcppslisting}

If we use the approach for
adding scope to enumerators that is described in \intraref{description-enumclass}{drawbacks-relating-to-weakly-typed,-c++03-enumerators}, 
%\textit{\titleref{description-enumclass}: \titleref{drawbacks-relating-to-weakly-typed,-c++03-enumerators}} on page~\pageref{drawbacks-relating-to-weakly-typed,-c++03-enumerators}, 
the name of the enclosing \lstinline!struct! together with a
consistent name for the enumeration, such as \lstinline!Enum!, has to be used to indicate an
enumerated type:

\begin{emcppslisting}[emcppsbatch=e3,]
void f3(VehicleScopedExplicitly::Enum value);
    // classically scoped (ù{\codeincomments{enum}}ù) passed by value
\end{emcppslisting}

\noindent Qualifying the enumerators of a scoped enumeration is the same,
irrespective of whether the scoping is explicit or implicit:

\begin{emcppslisting}[emcppsbatch=e3]
void g()
{
    f1(VehicleUnscoped::e_PLANE);
        // call (ù{\codeincomments{f1}}ù) with an explicitly scoped enumerator

    f2(VehicleScopedImplicitly::e_PLANE);
        // call (ù{\codeincomments{f2}}ù) with an implicitly scoped enumerator
}
\end{emcppslisting}

\noindent Apart from implicit scoping, the modern, C++11
\lstinline!enum!~\lstinline!class! deliberately does \emph{not} support
implicit conversion, in any context, to its \romeogloss{underlying type}:

\begin{emcppslisting}[emcppsbatch=e3]
void h()
{
    int i1 = VehicleScopedExplicitly::e_PLANE;
       // OK, scoped C++03 (ù{\codeincomments{enum}}ù) (implicit conversion)
 
    int i2 = VehicleScopedImplicitly::e_PLANE;
        // Error, no implicit conversion to underlying type

    if (VehicleScopedExplicitly::e_PLANE > 3) {} // OK
    if (VehicleScopedImplicitly::e_PLANE > 3) {} // Error, implicit conversion
}
\end{emcppslisting}

\noindent Enumerators of an \lstinline!enum!~\lstinline!class! do, however, admit
equality and ordinal comparisons within their own type:

\begin{emcppslisting}[emcppsbatch=e4]
enum class E { e_A, e_B, e_C };  // By default, enumerators increase from 0.

static_assert(E::e_A < E::e_C, "");  // OK, comparison between same-type values
static_assert(0 == E::e_A, "");      // Error, no implicit conversion from (ù{\codeincomments{E}}ù)
static_assert(0 == static_cast<int>(E::e_A), "");  // OK, explicit conversion

void f(E v)
{
    if (v > E::e_A) { /* ... */ }  // OK, comparing values of same type, (ù{\codeincomments{E}}ù)
}
\end{emcppslisting}

\noindent Note that incrementing an enumerator variable from one strongly typed
enumerator's value to the next requires an explicit cast; see \intraref{potential-pitfalls-enumclass}{strong-typing-of-an-enum-class-can-be-counterproductive}. 
%\textit{\titleref{potential-pitfalls-enumclass}: \titleref{strong-typing-of-an-enum-class-can-be-counterproductive}} on page~\pageref{strong-typing-of-an-enum-class-can-be-counterproductive}.

\subsubsection[\lstinline!enum! \lstinline!class! and underlying type]{{\SubsubsecCode enum} {\SubsubsecCode class} and underlying type}\label{enum-class-and-underlying-type}

Since C++11, both scoped and unscoped enumerations permit explicit
specification of their integral \romeogloss{underlying type}:

\begin{emcppslisting}
enum Ec : char { e_X, e_Y, e_Z };
    // underlying type is (ù{\codeincomments{char}}ù)

static_assert(1 == sizeof(Ec),     "");
static_assert(1 == sizeof Ec::e_X, "");

enum class Es : short { e_X, e_Y, e_Z };
    // underlying type is (ù{\codeincomments{short int}}ù)

static_assert(sizeof(short) == sizeof(Es),     "");
static_assert(sizeof(short) == sizeof Es::e_X, "");
\end{emcppslisting}

\noindent Unlike a classic \lstinline!enum!, which has an
\romeogloss{implementation-defined} default \romeogloss{underlying type}, the
default \romeogloss{underlying type} for an \lstinline!enum!~\lstinline!class! is
always \lstinline!int!:

\begin{emcppslisting}
enum class Ei { e_X, e_Y, e_Z };
    // When not specified the underlying type of an (ù{\codeincomments{enum class}}ù) is (ù{\codeincomments{int}}ù).

static_assert(sizeof(int) == sizeof(Ei),     "");
static_assert(sizeof(int) == sizeof Ei::e_X, "");
\end{emcppslisting}

\noindent Note that, because the default \romeogloss{underlying type} of an
\lstinline!enum!~\lstinline!class! is specified by the Standard, eliding the
enumerators of an \lstinline!enum!~\lstinline!class! in
a local redeclaration is \emph{always} possible; see \intraref{potential-pitfalls-enumclass}{external-use-of-opaque-enumerators-enumclass} and \featureref{\locationc}{enumopaque}. 
% \textit{\titleref{potential-pitfalls-enumclass}: \titleref{external-use-of-opaque-enumerators-enumclass}} on page~\pageref{external-use-of-opaque-enumerators-enumclass} and ``\titleref{enumopaque}" on page~\pageref{enumopaque}.

\subsection[Use Cases]{Use Cases}\label{use-cases-enumclass}

\subsubsection[Avoiding unintended implicit conversions to arithmetic types]{Avoiding unintended implicit conversions to arithmetic types}\label{avoiding-unintended-implicit-conversions-to-arithmetic-types}

Suppose that we want to represent the result of selecting one of a fixed
number of alternatives from a drop-down menu as a simple unordered set
of uniquely valued named integers. For example, this might be the case
when configuring a product, such as a vehicle, for purchase:

\begin{emcppslisting}
struct Trans
{
    enum Enum { e_MANUAL, e_AUTOMATIC };  // classic, C++03 scoped (ù{\codeincomments{enum}}ù)
};
\end{emcppslisting}

\noindent Although automatic promotion of a classic enumerator to \lstinline!int!
works well when typical use of the enumerator involves knowing its
cardinal value, such promotions are less than ideal when cardinal values
have no role in intended usage:

\begin{emcppshiddenlisting}[emcppsbatch=e5]
void attachDoor(int i);
\end{emcppshiddenlisting}
\begin{emcppslisting}[emcppsbatch=e5]
class Car { /* ... */ };

struct Trans
{                                         // explicitly scoped
    enum Enum { e_MANUAL, e_AUTOMATIC };  // classic (ù{\codeincomments{enum}}ù)
};                                        // (BAD IDEA)

int buildCar(Car* result, int numDoors, Trans::Enum trans)
{
    int status = Trans::e_MANUAL;    // Bug, accidental misuse!

    for (int i = 0; i < trans; ++i)  // Bug, accidental misuse!
    {
         attachDoor(i);
    }

    return status;
}
\end{emcppslisting}

\noindent As shown in the example above, it is never correct for a value of type
\lstinline!Trans::Enum! to be assigned to, compared with, or otherwise
modified like an integer; hence, \emph{any} such use would necessarily
be considered a mistake and, ideally, flagged by the compiler as an
error. The stronger typing provided by \lstinline!enum!~\lstinline!class!
achieves this goal:

\begin{emcppslisting}[emcppsignore={Fails to build with all Declarations marked Error}]
class Car { /* ... */ };

enum class Trans { e_MANUAL, e_AUTOMATIC };  // modern (ù{\codeincomments{enum class}}ù) (GOOD IDEA)

int buildCar(Car* result, int numDoors, Trans trans)
{
    int status = Trans::e_MANUAL;    // Error, incompatible types

    for (int i = 0; i < trans; ++i)  // Error, incompatible types
    {
         attachDoor(i);
    }

    return status;
}
\end{emcppslisting}

\noindent By deliberately choosing the \lstinline!enum!~\lstinline!class! over the
\emph{classic} \lstinline!enum! above, we automate the detection of many common
kinds of accidental misuse. Secondarily, we slightly simplify the interface of the function signature by removing the extra \lstinline!::Enum! boilerplate qualifications required of an explicitly
scoped, less-type-safe, classic \lstinline!enum!, but see \intraref{potential-pitfalls-enumclass}{strong-typing-of-an-enum-class-can-be-counterproductive}.
%\textit{\titleref{potential-pitfalls-enumclass}: \titleref{strong-typing-of-an-enum-class-can-be-counterproductive}} on page~\pageref{strong-typing-of-an-enum-class-can-be-counterproductive}.

In an unlikely event that the numeric value of a strongly typed
enumerator is needed (e.g., for serialization), it can be extracted
explicitly via a \lstinline!static_cast!:

\begin{emcppshiddenlisting}[emcppsbatch=e6]
class Car { /* ... */ };

enum class Trans { e_MANUAL, e_AUTOMATIC };  // modern (ù{\codeincomments{enum class}}ù) (GOOD IDEA)
\end{emcppshiddenlisting}
\begin{emcppslisting}[emcppsbatch=e6]
const int manualIntegralValue    = static_cast<int>(Trans::e_MANUAL);
const int automaticIntegralValue = static_cast<int>(Trans::e_AUTOMATIC);
static_assert(0 == manualIntegralValue,    "");
static_assert(1 == automaticIntegralValue, "");
\end{emcppslisting}

\subsubsection[Avoiding namespace pollution]{Avoiding namespace pollution}\label{avoiding-namespace-pollution}

Classic, C-style enumerations do not provide scope for their
enumerators, leading to unintended latent name collisions:

\begin{emcppslisting}
// vehicle.h:
// ...
enum Vehicle  { e_CAR, e_TRAIN, e_PLANE };  // classic, C-style (ù{\codeincomments{enum}}ù)
// ...

// geometry.h:
// ...
enum Geometry { e_POINT, e_LINE,  e_PLANE };  // classic, C-style (ù{\codeincomments{enum}}ù)
// ...

// client.cpp:
#include <vehicle.h>  // OK
#include <geometry.h> // Error, (ù{\codeincomments{e\_PLANE}}ù) redefined
// ...
\end{emcppslisting}

\noindent The common workaround is to wrap the \lstinline!enum! in a \lstinline!struct!
or \lstinline!namespace!:

\begin{emcppslisting}
// vehicle.h:
// ...
struct Vehicle {                            // explicitly scoped
    enum Enum { e_CAR, e_TRAIN, e_PLANE };  // classic, C-style (ù{\codeincomments{enum}}ù)
};
// ...

// geometry.h:
// ...
struct Geometry {                            // explicitly scoped
    enum Enum { e_POINT, e_LINE, e_PLANE };  // classic, C-style (ù{\codeincomments{enum}}ù)
};
// ...

// client.cpp:
#include <vehicle.h>    // OK
#include <geometry.h>   // OK, enumerators are scoped explicitly.
// ...
\end{emcppslisting}

\noindent If implicit conversions of enumerators to integral types are not
required, we can achieve the same scoping effect with much more type
safety and slightly less boilerplate --- i.e., without the
\lstinline!::Enum! when declaring a variable --- by employing
\lstinline!enum!~\lstinline!class! instead:

\begin{emcppslisting}
// vehicle.h:
// ...
enum class Vehicle { e_CAR, e_TRAIN, e_PLANE };
// ...

// geometry.h:
// ...
enum class Geometry { e_POINT, e_LINE, e_PLANE };
// ...

// client.cpp:
#include <vehicle.h>  // OK
#include <geometry.h>   // OK, enumerators are scoped implicitly.
// ...
\end{emcppslisting}

\subsubsection[Improving overloading disambiguation]{Improving overloading disambiguation}\label{improving-overloading-disambiguation}

Overloaded functions are notorious for providing opportunities for
misuse. Maintenance difficulties are exacerbated when arguments for
these overloads are convertible to more than a single parameter in the
function. As an illustration of the compounding of such maintenance
difficulties, suppose that we have a widely used, named type,
\lstinline!Color!, and the numeric values of its enumerators are small,
unique, and irrelevant. Imagine we have chosen to represent
\lstinline!Color! as a \emph{classic} \lstinline!enum!:

\begin{emcppslisting}[emcppsbatch=e7]
struct Color
{                                           // explicitly scoped
    enum Enum { e_RED, e_BLUE /*, ...*/ };  // classic, C-style (ù{\codeincomments{enum}}ù)
};                                          // (BAD IDEA)
\end{emcppslisting}

\noindent Suppose further that we have provided two overloaded functions, each
having two parameters, with one signature's parameters including the
enumeration \lstinline!Color!:

\begin{emcppslisting}[emcppsbatch=e7]
void clearScreen(int pattern, int orientation);          // (0)
void clearScreen(Color::Enum background, double alpha);  // (1)
\end{emcppslisting}

\noindent Depending on the types of the arguments supplied, one or the other
functions will be selected or else the call will be ambiguous and the
program will fail to compile{\cprotect\footnote{GCC version 7.4.0
incorrectly diagnoses both ambiguity errors as warnings, although it
states in the warning that it is an error:

\begin{emcppslisting}[language=bash,style=footcodeplain]
warning: ISO C++ says that these are ambiguous, even though the worst conversion for the
         first is better than the worst conversion for the second:

note: candidate 1: void clearScreen(int, int)
void clearScreen(int pattern, int orientation);
     ^~~~~~~~~~~
note: candidate 2: void clearScreen(Color::Enum, double)
void clearScreen(Color::Enum background, double alpha;
     ^~~~~~~~~~~
\end{emcppslisting}
      }}:

\begin{emcppslisting}[emcppsbatch=e7]
void f0()
{
    clearScreen(1           , 1           );  // calls (0) above
    clearScreen(1           , 1.0         );  // calls (0) above
    clearScreen(1           , Color::e_RED);  // calls (0) above

    clearScreen(1.0         , 1           );  // calls (0) above
    clearScreen(1.0         , 1.0         );  // calls (0) above
    clearScreen(1.0         , Color::e_RED);  // calls (0) above

    clearScreen(Color::e_RED, 1           );  // Error, ambiguous call
    clearScreen(Color::e_RED, 1.0         );  // calls (1) above
    clearScreen(Color::e_RED, Color::e_RED);  // Error, ambiguous call
}
\end{emcppslisting}

\noindent Now suppose that we had instead defined our \lstinline!Color! enumeration
as a modern \lstinline!enum!~\lstinline!class!:

\begin{emcppslisting}[emcppsbatch=e8]
enum class Color { e_RED, e_BLUE /*, ...*/ };

void clearScreen(int pattern, int orientation);    // (2)
void clearScreen(Color background, double alpha);  // (3)
\end{emcppslisting}

\noindent The function that will be called from a given set of arguments becomes
clear:

\begin{emcppslisting}[emcppsbatch=e8]
void f1()
{
    clearScreen(1           , 1           );  // calls (2) above
    clearScreen(1           , 1.0         );  // calls (2) above
    clearScreen(1           , Color::e_RED);  // Error, no matching function

    clearScreen(1.0         , 1           );  // calls (2) above
    clearScreen(1.0         , 1.0         );  // calls (2) above
    clearScreen(1.0         , Color::e_RED);  // Error, no matching function

    clearScreen(Color::e_RED, 1           );  // calls (3) above
    clearScreen(Color::e_RED, 1.0         );  // calls (3) above
    clearScreen(Color::e_RED, Color::e_RED);  // Error, no matching function
}
\end{emcppslisting}

\noindent Returning to our original, classic-\lstinline!enum! design, suppose that we
find we need to add a third parameter, \lstinline!bool!~\lstinline!z!, to the
second overload:

\begin{emcppshiddenlisting}[emcppsbatch=e9]
struct Color
{
    enum Enum { e_RED, e_BLUE /*, ...*/ };
};
\end{emcppshiddenlisting}
\begin{emcppslisting}[emcppsbatch=e9]
void clearScreen(int pattern, int orientation);                  // (0)
void clearScreen(Color::Enum background, double alpha, bool z);  // (4) classic
\end{emcppslisting}

\noindent If our plan is that any existing client calls involving
\lstinline!Color::Enum! will now be flagged as errors, we are going to be
very disappointed:

\begin{emcppslisting}[emcppsbatch=e9]
void f2()
{
    clearScreen(Color::e_RED, 1.0);  // calls (0) above
}
\end{emcppslisting}

\noindent In fact, every combination of arguments above --- all nine of them ---
will call function (0) above with no warnings at all:

\begin{emcppslisting}[emcppsbatch=e9]
void f3()
{
    clearScreen(1           , 1           );  // calls (0) above
    clearScreen(1           , 1.0         );  // calls (0) above
    clearScreen(1           , Color::e_RED);  // calls (0) above

    clearScreen(1.0         , 1           );  // calls (0) above
    clearScreen(1.0         , 1.0         );  // calls (0) above
    clearScreen(1.0         , Color::e_RED);  // calls (0) above

    clearScreen(Color::e_RED, 1           );  // calls (0) above
    clearScreen(Color::e_RED, 1.0         );  // calls (0) above
    clearScreen(Color::e_RED, Color::e_RED);  // calls (0) above
}
\end{emcppslisting}

\noindent Finally, let's suppose again that we have used
\lstinline!enum!~\lstinline!class! to implement our \lstinline!Color!\linebreak[4]
\mbox{enumeration}:

\begin{emcppshiddenlisting}[emcppsbatch=e10]
enum class Color { e_RED, e_BLUE /*, ...*/ };
\end{emcppshiddenlisting}
\begin{emcppslisting}[emcppsbatch=e10]
void clearScreen(int pattern, int orientation);            // (2)
void clearScreen(Color background, double alpha, bool z);  // (5) modern

void f4()
{
    clearScreen(Color::e_RED, 1.0);  // Error, no matching function
}
\end{emcppslisting}

\noindent And in fact, the \emph{only} calls that succeed unmodified are precisely
those that do not involve the enumeration \lstinline!Color!, as desired:

\begin{emcppslisting}[emcppsbatch=e10]
void f5()
{
    clearScreen(1           , 1           );  // calls (2) above
    clearScreen(1           , 1.0         );  // calls (2) above
    clearScreen(1           , Color::e_RED);  // Error, no matching function

    clearScreen(1.0         , 1           );  // calls (2) above
    clearScreen(1.0         , 1.0         );  // calls (2) above
    clearScreen(1.0         , Color::e_RED);  // Error, no matching function

    clearScreen(Color::e_RED, 1           );  // Error, no matching function
    clearScreen(Color::e_RED, 1.0         );  // Error, no matching function
    clearScreen(Color::e_RED, Color::e_RED);  // Error, no matching function
}
\end{emcppslisting}

\noindent Bottom line: Having a \emph{pure} enumeration --- such as
\lstinline!Color!, used widely in function signatures --- be strongly typed
can only help to expose accidental misuse but, again, see \intraref{potential-pitfalls-enumclass}{strong-typing-of-an-enum-class-can-be-counterproductive}.
%\textit{\titleref{potential-pitfalls-enumclass}: \titleref{strong-typing-of-an-enum-class-can-be-counterproductive}} on page~\pageref{strong-typing-of-an-enum-class-can-be-counterproductive}.

Note that strongly typed enumerations help to avoid accidental misuse by
requiring an explicit \emph{cast} should conversion to an arithmetic
type be desired:

\begin{emcppslisting}[emcppsbatch=e10]
void f6()
{
    clearScreen(Color::e_RED, 1.0);                   // Error, no match
    clearScreen(static_cast<int>(Color::e_RED), 1.0); // OK, calls (2) above
    clearScreen(Color::e_RED, 1.0, false);            // OK, calls (5) above
}
\end{emcppslisting}


\subsubsection[Encapsulating implementation details within the enumerators themselves]{Encapsulating implementation details within the enumerators themselves}\label{encapsulating-implementation-details-within-the-enumerators-themselves}

In rare cases, providing a pure, ordered enumeration having unique (but
not necessarily contiguous) numerical values that exploit lower-order
bits to categorize and make readily available important
individual properties might offer an advantage, such as in performance. Note that in order to preserve the ordinality of the enumerators
overall, the higher-level bits must encode their relative order. The
lower-level bits are then available for arbitrary use in the
  implementation..

For example, suppose that we have a \lstinline!MonthOfYear! enumeration
that encodes the months that have 31 days in their least-significant bit
and an accompanying \lstinline!inline! function to quickly determine
whether a given enumerator represents such a month:

\begin{emcppslisting}
#include <type_traits>  // (ù{\codeincomments{std::underlying\_type}}ù)

enum class MonthOfYear : unsigned char  // optimized to flag long months
{
    e_JAN = ( 1 << 4) + 0x1,
    e_FEB = ( 2 << 4) + 0x0,
    e_MAR = ( 3 << 4) + 0x1,
    e_APR = ( 4 << 4) + 0x0,
    e_MAY = ( 5 << 4) + 0x1,
    e_JUN = ( 6 << 4) + 0x0,
    e_JUL = ( 7 << 4) + 0x1,
    e_AUG = ( 8 << 4) + 0x1,
    e_SEP = ( 9 << 4) + 0x0,
    e_OCT = (10 << 4) + 0x1,
    e_NOV = (11 << 4) + 0x0,
    e_DEC = (12 << 4) + 0x1
};

bool hasThirtyOneDays(MonthOfYear month)
{
    return static_cast<std::underlying_type<MonthOfYear>::type>(month) & 0x1;
}
\end{emcppslisting}

\noindent In the example above, we are using a new cross-cutting
feature of all enumerated types that allows the client defining the
type to specify its underlying type precisely. In this case, we have
chosen an \lstinline!unsigned!~\lstinline!char! to maximize the number of
flag bits while keeping the overall size to a single byte. Three bits
remain available. Had we needed more flag bits, we could have just as
easily used a larger underlying type, such as
  \lstinline!unsigned!~\lstinline!short!; see \featureref{\locationc}{explicit-enumeration-underlying-type}.
  %``\titleref{explicit-enumeration-underlying-type}" on page~\pageref{explicit-enumeration-underlying-type}.

\noindent In case \lstinline!enum!s are used for encoding purposes, the public clients are not intended to make use of the
cardinal values; hence clients are well advised to treat them as
implementation details, potentially subject to change without notice.
Representing this enumeration using the modern
\lstinline!enum!~\lstinline!class!, instead of an explicitly scoped classic
\lstinline!enum!, deters clients from making any use (apart from same-type
comparisons) of the cardinal values assigned to the enumerators. Notice
that implementors of the \lstinline!hasThirtyOneDays! function will require
a verbose but efficient \lstinline!static_cast! to resolve the
cardinal value of the enumerator and thus make the requested
determination as efficiently as possible.

\subsection[Potential Pitfalls]{Potential Pitfalls}\label{potential-pitfalls-enumclass}

\subsubsection[Strong typing of an \lstinline!enum! \lstinline!class! can be counterproductive]{Strong typing of an {\SubsubsecCode enum} {\SubsubsecCode class} can be counterproductive}\label{strong-typing-of-an-enum-class-can-be-counterproductive}

The additive value in using a modern \lstinline!enum!~\lstinline!class! is
governed \emph{solely} by whether its stronger typing, \emph{not} its
implicit scoping, of its enumerators would be beneficial in its
anticipated typical usage. If the expectation is that the client will
never need to know the specific values of the enumerators, then use
of the modern \lstinline!enum!~\lstinline!class! is often just what's needed.
But if the cardinal values themselves are ever needed during typical
use, extracting them will require the client to perform an explicit
cast. Beyond mere inconvenience, encouraging clients to use casts
invites defects.

Suppose, for example, we have a function, \lstinline!setPort!, from an
external library that takes an integer port number:

\begin{emcppslisting}[emcppsbatch=e11]
int setPort(int portNumber);
    // Set the current port; return (ù{\codeincomments{0}}ù) on success and a nonzero value otherwise.
\end{emcppslisting}

\noindent Suppose further that we have used the modern
\lstinline!enum!~\lstinline!class! feature to implement an enumeration,
\lstinline!SysPort!, that identifies well-known ports on our system:

\begin{emcppslisting}[emcppsbatch=e11]
enum class SysPort { e_INPUT = 27, e_OUTPUT = 29, e_ERROR = 32, e_CTRL = 6 };
    // enumerated port values used to configure our systems
\end{emcppslisting}

\noindent Now suppose we want to call the function \lstinline!f! using one of these
enumerated values:

\begin{emcppslisting}[emcppsbatch=e11]
void setCurrentPortToCtrl()
{
    setPort(SysPort::e_CTRL);  // Error, cannot convert (ù{\codeincomments{SetPort}}ù) to (ù{\codeincomments{int}}ù)
}
\end{emcppslisting}

\noindent Unlike the situation for a \emph{classic} \lstinline!enum!, no implicit
conversion occurs from an \lstinline!enum!~\lstinline!class! to its underlying
integral type, so anyone using this enumeration will be forced to
somehow explicitly \romeogloss{cast} the enumerator to some arithmetic type.
There are, however, multiple choices for performing this cast:

\begin{emcppslisting}[emcppsbatch=e11]
#include <type_traits>  // (ù{\codeincomments{std::underlying\_type}}ù)

void test()
{
    setPort(int(SysPort::e_CTRL));                                       // (1)
    setPort((int)SysPort::e_CTRL);                                       // (2)
    setPort(static_cast<int>(SysPort::e_CTRL));                          // (3)
    setPort(static_cast<std::underlying_type<SysPort>::type>(            // (4)
                                                         SysPort::e_CTRL));
    setPort(static_cast<int>(                                            // (5)
             static_cast<std::underlying_type<SysPort>::type>(SysPort::e_CTRL)));
}
\end{emcppslisting}

\noindent Any of the above casts would work in this case, but consider a future
where a platform changed \lstinline!setPort! to take a \lstinline!long! and
the control port was changed to a value that cannot be represented as an
\lstinline!int!:

\begin{emcppslisting}
int setPort(long portNumber);
enum class SysPort : unsigned { e_INPUT = 27, e_OUTPUT = 29, e_ERROR = 32,
                                e_CTRL = 0x80000000 };
    // enumerated port values used to configure our systems
\end{emcppslisting}

\noindent Only casting method (4) above will pass the correct value for
\lstinline!e_CTRL! to this new \lstinline!setPort! implementation. The other
variations will all pass a negative number for the port, which would
certainly not be the intention of the user writing this code. A classic
C-style \lstinline!enum! would have avoided any manually written cast
entirely and the proper value would propagate into \lstinline!setPort! even
as the range of values used for ports changes:

\begin{emcppshiddenlisting}[emcppsbatch=e12]
#include <type_traits>  // (ù{\codeincomments{std::is\_same}}ù)
int setPort(long portNumber);
\end{emcppshiddenlisting}
\begin{emcppslisting}[emcppsbatch=e12]
struct SysPort  // explicit scoping for a classic, C-style (ù{\codeincomments{enum}}ù)
{
    enum Enum { e_INPUT = 27, e_OUTPUT = 29, e_ERROR = 32,
                e_CTRL = 0x80000000 };

    // Note that the underlying type of (ù{\codeincomments{Enum}}ù) is implicit and will be
    // large enough to represent all of these values.
    static_assert(
        std::is_same<std::underlying_type<Enum>::type,unsigned>::value, "");
};

void setCurrentPortToCtrl()
{
    setPort(SysPort::e_CTRL);  // OK, (ù{\codeincomments{SysPort::Enum}}ù) promotes to (ù{\codeincomments{long}}ù).
}
\end{emcppslisting}

\noindent When the intended client will depend on the cardinal values of the
enumerators during routine use, we can avoid tedious, error-prone, and
repetitive casting by instead employing a classic, C-style
\lstinline!enum!, possibly nested within a \lstinline!struct! to achieve
explicit scoping of its enumerators. The subsections that follow
highlight specific cases in which classic, C-style, C++03
\lstinline!enum!s are appropriate.

\subsubsection[Misuse of \lstinline!enum! \lstinline!class! for collections of named constants]{Misuse of {\SubsubsecCode enum} {\SubsubsecCode class} for collections of named constants}\label{misuse-of-enum-class-for-collections-of-named-constants}

When constants are truly independent, we are often encouraged to avoid
enumerations altogether, preferring instead individual constants; see \featureref{\locationc}{Default-Member-Initializers}. 
%``\titleref{Default-Member-Initializers}" on page~\pageref{Default-Member-Initializers}. 
On the other hand, when the constants all participate within a coherent
theme, the expressiveness achieved using a \emph{classic} \lstinline!enum!
to aggregate those values is compelling. Another
advantage of an enumerator over an individual constant is that the
enumerator is guaranteed to be a \romeogloss{compile-time constant} (see \featureref{\locationc}{constexprvar})
%  ``\titleref{constexprvar}" on page~\pageref{constexprvar}) 
  and a \romeovalue{prvalue} (see \featureref{\locationc}{Rvalue-References}), 
%  ``\titleref{Rvalue-References}" on page~\pageref{Rvalue-References}), 
  which never needs static storage
  and cannot have its address taken.

For example, suppose we want to collect the coefficients for various
numerical suffixes representing \emph{thousands}, \emph{millions}, and
\emph{billions} using an enumeration:

\begin{emcppslisting}[emcppsbatch=e13]
enum class S0 { e_K = 1000, e_M = e_K * e_K, e_G = e_M * e_K };  // (BAD IDEA)
\end{emcppslisting}

\noindent A client trying to access one of these enumerated values would need to
cast it explicitly:

\begin{emcppslisting}[emcppsbatch=e13]
void client0()
{
    int distance = 5 * static_cast<int>(S0::e_K);  // casting is error-prone
    // ...
}
\end{emcppslisting}

\noindent By instead making the enumeration an explicitly scoped, \emph{classic}
\lstinline!enum! nested within a \lstinline!struct!, no casting is needed
during typical use:

\begin{emcppslisting}
struct S1  // scoped
{
    enum Enum { e_K = 1000, e_M = e_K * e_K, e_G = e_M * e_K };
        // *classic* (ù{\codeincomments{enum}}ù) (GOOD IDEA)
};

void client1()
{
    int distance = 5 * S1::e_K;  // no casting required during typical use
    // ...
}
\end{emcppslisting}

\noindent If the intent is that these constants will be specified and used in a
purely local context, we might choose to drop the enclosing scope, along
with the name of the enumeration itself; see \featureref{\locationa}{local-types-as-template-arguments}:
%``\titleref{local-types-as-template-arguments}" on page~\pageref{local-types-as-template-arguments}:

\begin{emcppslisting}
void client2()
{
    enum { e_K = 1000, e_M = e_K * e_K, e_G = e_M * e_K };  // function scoped

    double salary = 95 * e_K;
    double netWorth = 0.62 * e_M;
    double companyRevenue = 47.2 * e_G;
    // ...
}
\end{emcppslisting}

\noindent We sometimes use the
lowercase prefix \lstinline!k_! instead of \lstinline!e_! to indicate
salient \romeogloss{compile-time constants} that are not considered part
of an enumerated set, irrespective of whether they are implemented as
enumerators:

\begin{emcppslisting}[language=C++]
enum { k_NUM_PORTS = 500, k_PAGE_SIZE = 512 };     // compile-time constants
static const double k_PRICING_THRESHOLD = 0.03125; // compile-time constant
\end{emcppslisting}

\subsubsection[Misuse of \lstinline!enum! \lstinline!class! in association with bit flags]{Misuse of {\SubsubsecCode enum} {\SubsubsecCode class} in association with bit flags}\label{misuse-of-enum-class-in-association-with-bit-flags}

Using \lstinline!enum!~\lstinline!class! to implement enumerators that are
intended to interact closely with arithmetic types will typically
require the definition of arithmetic and bitwise operator overloads
between values of the same enumeration and between the enumeration and
arithmetic types, leading to yet more code that needs to be written,
tested, and maintained. This is often the case for bit flags. Consider, for example, an
enumeration used to control a file system:

\begin{emcppslisting}[emcppsbatch=e14]
enum class Ctrl { e_READ = 0x1, e_WRITE = 0x2, e_EXEC = 0x4 };  // (BAD IDEA)
    // low-level bit flags used to control file system

void chmodFile(int fd, int access);
    // low-level function used to change privileges on a file
\end{emcppslisting}

\noindent We could conceivably write a series of functions to combine the
individual flags in a type-safe manner:

\begin{emcppslisting}[emcppsbatch=e14]
#include <type_traits>  // (ù{\codeincomments{std::underlying\_type}}ù)

int flags() { return 0; }
int flags(Ctrl a) { return static_cast<std::underlying_type<Ctrl>::type>(a); }
int flags(Ctrl a, Ctrl b) { return flags(a) | flags(b); }
int flags(Ctrl a, Ctrl b, Ctrl c) { return flags(a, b) | flags(c); }

void setRW(int fd)
{
    chmodFile(fd, flags(Ctrl::e_READ, Ctrl::e_WRITE));  // (BAD IDEA)
}
\end{emcppslisting}

\noindent Alternatively, a \emph{classic}, C-style \lstinline!enum! nested within a
\lstinline!struct! achieves what's needed:

\begin{emcppslisting}
struct Ctrl // scoped
{
    enum Enum { e_READ = 0x1, e_WRITE = 0x2, e_EXEC = 0x4 };  // classic (ù{\codeincomments{enum}}ù)
        // low-level bit flags used to control file system (GOOD IDEA)
};

void chmodFile(int fd, int access);
    // low-level function used to change privileges on a file

void setRW(int fd)
{
    chmodFile(fd, Ctrl::e_READ | Ctrl::e_WRITE);  // (GOOD IDEA)
}
\end{emcppslisting}


\subsubsection[Misuse of \lstinline!enum! \lstinline!class! in association with iteration]{Misuse of {\SubsubsecCode enum} {\SubsubsecCode class} in association with iteration}\label{misuse-of-enum-class-in-association-with-iteration}

Sometimes the relative values of enumerators are considered important as
well. For example, let's again consider enumerating the months of the
year:

\begin{emcppslisting}[emcppsbatch={e15,e16}]
enum class MonthOfYear  // modern, strongly typed enumeration
{
    e_JAN, e_FEB, e_MAR,  // winter
    e_APR, e_MAY, e_JUN,  // spring
    e_JUL, e_AUG, e_SEP,  // summer
    e_OCT, e_NOV, e_DEC,  // autumn
};
\end{emcppslisting}

\noindent If all we need to do is compare the ordinal values of the enumerators,
there's no problem:

\begin{emcppslisting}[emcppsbatch=e15]
bool isSummer(MonthOfYear month)
{
    return MonthOfYear::e_JUL <= month && month <= MonthOfYear::e_AUG;
}
\end{emcppslisting}

\noindent Although the \lstinline!enum!~\lstinline!class! features allow for relational
and equality operations between like-typed enumerators, no arithmetic
operations are supported directly, which becomes problematic when we
need to iterate over the enumerated values:

\begin{emcppslisting}[emcppsbatch=e15,emcppsignore={Entire for block is marked Error}]
void doSomethingWithEachMonth()
{
    for (MonthOfYear i =  MonthOfYear::e_JAN;
                     i <= MonthOfYear::e_DEC;
                   ++i)  // Error, no match for (ù{\codeincomments{++}}ù)
    {
        // ...
    }
}
\end{emcppslisting}

\noindent To make this code compile, an explicit cast from and to the enumerated
type will be required:

\begin{emcppslisting}[emcppsbatch=e15]
void doSomethingWithEachMonth()
{
    for (MonthOfYear i =  MonthOfYear::e_JAN;
                     i <= MonthOfYear::e_DEC;
                     i = static_cast<MonthOfYear>(static_cast<int>(i) + 1))
    {
        // ...
    }
}
\end{emcppslisting}

\noindent Alternatively, an auxiliary, helper function could be supplied to allow
clients to bump the enumerator:

\begin{emcppslisting}[emcppsbatch=e16]
MonthOfYear nextMonth(MonthOfYear value)
{
    return static_cast<MonthOfYear>(static_cast<int>(value) + 1);
}

void doSomethingWithEachMonth()
{
    for (MonthOfYear i =  MonthOfYear::e_JAN;
                     i <= MonthOfYear::e_DEC;
                     i = nextMonth(i))
    {
        // ...
    }
}
\end{emcppslisting}

\noindent If, however, the cardinal value of the \lstinline!MonthOfYear! enumerators
is likely to be relevant to clients, an explicitly scoped \emph{classic}
\lstinline!enum! should be considered as a viable alternative:

\begin{emcppslisting}
struct MonthOfYear  // explicit scoping for (ù{\codeincomments{enum}}ù)
{
    enum Enum
    {
        e_JAN, e_FEB, e_MAR,  // winter
        e_APR, e_MAY, e_JUN,  // spring
        e_JUL, e_AUG, e_SEP,  // summer
        e_OCT, e_NOV, e_DEC,  // autumn
    };
};

bool isSummer(MonthOfYear::Enum month)  // must now pass nested (ù{\codeincomments{Enum}}ù) type
{
    return MonthOfYear::e_JUL <= month && month <= MonthOfYear::e_AUG;
}

void doSomethingWithEachMonth()
{
    for (int i =  MonthOfYear::e_JAN;  // iteration variable is now an (ù{\codeincomments{int}}ù)
             i <= MonthOfYear::e_DEC;
           ++i)  // OK, convert to underlying type
    {
        // ... (might require cast back to enumerated type)
    }
}
\end{emcppslisting}

\noindent Note that such code presumes that the enumerated values will (1) remain
in the same order and (2) have contiguous numerical values irrespective
of the implementation choice.

\subsubsection[External use of opaque enumerators]{External use of opaque enumerators}\label{external-use-of-opaque-enumerators-enumclass}

Since \lstinline!enum!~\lstinline!class! types have an \romeogloss{underlying
type} of \lstinline!int! by default, clients are always able to (re)declare
it, as a \romeogloss{complete type}, without its enumerators. Unless the
opaque form of an \lstinline!enum!~\lstinline!class!'s definition is exported
in a header file separate from the one implementing the publicly
accessible full definition, external clients wishing to exploit the
opaque version will experience an \emph{attractive nuisance} in that
they can provide it locally, along with its \romeogloss{underlying type}, if
any.

If the underlying type of the full definition were to subsequently
change, any program incorporating the original elided definition locally
and also the new, full one from the header would become silently \romeogloss{ill formed, no diagnostic required (IFNDR)}; see \featureref{\locationc}{enumopaque}.
%``\titleref{enumopaque}" on page~\pageref{enumopaque}.

\subsection[Annoyances]{Annoyances}\label{annoyances-enumclass}

\subsubsection[Scoped enumerations do not necessarily add value]{Scoped enumerations do not necessarily add value}\label{scoped-enumerations-do-not-necessarily-add-value}

When the enumeration is local, say, within the scope of a given
function, forcing an additional scope on the enumerators is superfluous.
For example, consider a function that returns an integer status
\lstinline!0! on success and a nonzero value otherwise:

\begin{emcppslisting}[emcppserrorlines={5,7}]
int f()
{
    enum { e_ERROR = -1, e_OK = 0 } result = e_OK;
    // ...
    if (/* error 1 */) { result = e_ERROR; }
    // ...
    if (/* error 2 */) { result = e_ERROR; }
    // ...
    return result;
}
\end{emcppslisting}

\noindent Use of \lstinline!enum!~\lstinline!class! in this context would require
potentially needless qualification --- and perhaps even casting ---
where it might not be warranted:

\begin{emcppslisting}[emcppserrorlines={5,7}]
int f()
{
    enum class RC { e_ERROR = -1, e_OK = 0 } result = RC::e_OK;
    // ...
    if (/* error 1 */) { result = RC::e_ERROR; } // undesirable qualification
    // ...
    if (/* error 2 */) { result = RC::e_ERROR; } // undesirable qualification
    // ...
    return static_cast<int>(result);  // undesirable explicit cast
}
\end{emcppslisting}

\subsection[See Also]{See Also}\label{see-also}

\begin{itemize}
\item{\seealsoref{enumopaque}{\seealsolocationc}Sometimes it is useful to entirely insulate individual enumerators from clients.}
\item{\seealsoref{explicit-enumeration-underlying-type}{\seealsolocationc}Absent implicit conversion to integrals, \lstinline!enum! class values may use \lstinline!static_cast! in conjunction with their underlying type.}
\end{itemize}

\subsection[Further Reading]{Further Reading}\label{further-reading}

\begin{itemize}
\item{\cite{miller19}}
\item{\cite{meyers15}}
\end{itemize}




\newpage
%\section[{\tt extern} {\tt template}]{Explicit Instantiation Declarations}\label{explicit-template-instantiations}

% 01 Feb, 2021 JMB - packet 4 compilation fixes

\emcppsFeature{
    short={\lstinline!extern!~\lstinline!template!},
    tocshort={{\TOCCode extern}~{\TOCCode template}},
    long={Explicit Instantiation Declarations},
    rhshort={{\RHCode extern}~{\RHCode template}},
}{explicit-template-instantiations}
\setcounter{table}{0}
\setcounter{footnote}{0}
\setcounter{lstlisting}{0}
%\section[{\tt extern} {\tt template}]{Explicit Instantiation Declarations}\label{explicit-template-instantiations}

The \lstinline!extern!~\lstinline!template! prefix can be used to suppress
\emph{implicit} generation of local object code for the definitions of
particular specializations of class, function, or
variable templates used within a translation unit, with the
expectation that any suppressed object-code-level definitions will be
provided elsewhere within the program by template definitions that are
instantiated \emph{explicitly}.

\subsection[Description]{Description}\label{description}

Inherent in the current ecosystem for supporting template programming in
C++ is the need to generate redundant definitions of fully specified
function and variable templates within \lstinline!.o! files. For common
instantiations of popular templates, such as \lstinline!std::vector!, the
increased object-file size --- a.k.a. \romeogloss{code bloat} --- and
potentially extended link times might become significant:

\begin{emcppslisting}
#include <vector>    // (ù{\codeincomments{std::vector}}ù) is a popular template.
std::vector<int> v;  // (ù{\codeincomments{std::vector<int>}}ù) is a common instantiation.

#include <string>    // (ù{\codeincomments{std::basic\_string}}ù) is a popular template.
std::string s;       // (ù{\codeincomments{std::string}}ù), an alias for (ù{\codeincomments{std::basic\_string<char>}}ù), is
                     // a common instantiation.
\end{emcppslisting}

\noindent The intent of the \lstinline!extern!~\lstinline!template! feature is to
\emph{suppress} the implicit generation of duplicative object code
within each and every translation unit in which a fully specialized
class template, such as \lstinline!std::vector<int>! in the code snippet
above, is used. Instead, \lstinline!extern!~\lstinline!template! allows
developers to choose a single translation unit in which to explicitly
\emph{generate} object code for all the definitions pertaining to that
specific template specialization as explained in the next
subsection, \intrarefnopage{explicit-instantiation-definition}.

\subsubsection[Explicit-instantiation definition]{Explicit-instantiation definition}\label{explicit-instantiation-definition}

The ability to create an \romeogloss{explicit-instantiation
definition} has been available since C++98.{\cprotect\footnote{The
  C++03 Standard term for the syntax used to create an
  \romeogloss{explicit-instantiation \emph{definition}}, though rarely used,
  was \romeogloss{explicit-instantiation \emph{directive}}. The term
  \romeogloss{explicit-instantiation directive} was clarified in
  C++11 and can now also refer to syntax that is used to create a
  \emph{declaration} --- i.e., \romeogloss{explicit-instantiation
  \emph{declaration}}.}} The requisite syntax is to place the keyword
\lstinline!template! in front of the name of the fully specialized class
template, function template, or --- in C++14 --- variable
template (see \featureref{\locationb}{variable-templates}):

\begin{emcppslisting}
#include <vector>  // (ù{\codeincomments{std::vector}}ù) (general template)

template class std::vector<int>;
    // Deposit all definitions for this specialization into the (ù{\codeincomments{.o}}ù) for this
    // translation unit.
\end{emcppslisting}

\noindent This \romeogloss{explicit-instantiation directive} compels the compiler to
instantiate \emph{all} functions defined by the named
\lstinline!std::vector! class template having the specified \lstinline!int!
template argument; any collateral object code resulting from these
instantiations will be deposited in the resulting \lstinline!.o! file for
the current translation unit. Importantly, even functions that are never
used are still specialized, so this might not be the correct solution
for many classes; see \intraref{potential-pitfalls-externtemplate}{accidentally-making-matters-worse}.

\subsubsection[Explicit-instantiation declaration]{Explicit-instantiation declaration}\label{explicit-instantiation-declaration}

C++11 introduced the \romeogloss{explicit-instantiation \emph{declaration}},
complement to the \romeogloss{explicit-instantiation \emph{definition}}. The
newly provided syntax allows us to place
\lstinline!extern!~\lstinline!template! in front of the declaration of the
explicit specialization of a class template, a function template, or a
variable template: %(see \featureref{\locationb}{variable-templates}):

\begin{emcppslisting}
#include <vector>  // (ù{\codeincomments{std::vector}}ù) (general template)

extern template class std::vector<int>;
    // Suppress depositing of any object code for (ù{\codeincomments{std::vector<int>}}ù) into the
    // (ù{\codeincomments{.o}}ù) file for this translation unit.
\end{emcppslisting}

\noindent The use of the modern \lstinline!extern!~\lstinline!template! syntax above
instructs the compiler to specifically \emph{not} deposit any object
code for the named specialization in the current translation unit and
instead to rely on some other translation unit to provide any missing
object-level definitions that might be needed at link time; see
\intraref{annoyances-externtemplate}{no-good-place-to-put-definitions-for-unrelated-classes}.

Note, however, that declaring an explicit instantiation to be an
\lstinline!extern!~\lstinline!template! \emph{in no way} affects the ability
of the compiler to instantiate and to inline visible function-definition
bodies for that template specialization in the translation unit:

\begin{emcppslisting}
// client.cpp:
#include <vector>  // (ù{\codeincomments{std::vector}}ù) (general template)

extern template class std::vector<int>;  // specialization for (ù{\codeincomments{int}}ù) elements

void client(std::vector<int>& inOut)  // fully specialized instance of a (ù{\codeincomments{vector}}ù)
{
    if (inOut.size())          // This invocation of (ù{\codeincomments{size}}ù) can inline.
    {
        int value = inOut[0];  // This invocation of (ù{\codeincomments{operator[]}}ù) can inline.
    }
}
\end{emcppslisting}

\noindent In the example above, the two tiny member functions of \lstinline!vector!,
namely \lstinline!size! and \lstinline!operator[]!, will typically be
substituted inline --- in precisely the same way they would have been
had the \lstinline!extern!~\lstinline!template! declaration been omitted. The
\emph{only} purpose of an \lstinline!extern!~\lstinline!template! declaration
is to suppress object-code generation for this particular template
instantiation for the current translation unit.

Finally, note that the use of \romeogloss{explicit-instantiation
\emph{directives}} have absolutely no affect on the logical meaning of a
well-formed program; in particular, when applied to specializations of
function templates, they have no affect whatsoever on overload
resolution:

\begin{emcppslisting}
template <typename T> bool f(T v) {/*...*/}  // general template definition

extern template bool f(char c);  // specialization of (ù{\codeincomments{f}}ù) for (ù{\codeincomments{char}}ù)
extern template bool f(int v);   // specialization of (ù{\codeincomments{f}}ù) for (ù{\codeincomments{int}}ù)

char     c;
short    s;
int      i;
unsigned u;

bool bc = f(c);  // exact match: Object code is suppressed locally.
bool bs = f(s);  // not exact match: Object code is generated locally.
bool bi = f(i);  // exact match: Object code is suppressed locally.
bool bu = f(u);  // not exact match: Object code is generated locally.
\end{emcppslisting}

\noindent As the example above illustrates, overload resolution and template
parameter deduction occur independently of any
\romeogloss{explicit-instantiation declarations}. Only \emph{after}
the template to be instantiated is determined does the
\lstinline!extern!~\lstinline!template! syntax take effect; see also \intraref{potential-pitfalls-externtemplate}{corresponding-explicit-instantiation-declarations-and-definitions}.

\subsubsection[A more complete illustrative example]{A more complete illustrative example}\label{a-more-complete-illustrative-example}

So far, we have seen the use of \romeogloss{explicit-instantiation declarations} and \romeogloss{explicit-instan\-tiation definitions} applied to only a
(standard) \emph{class} template, \lstinline!std::vector!. The same syntax
shown in the previous code snippet applies also to full specializations of individual
function templates and variable templates. % (see \featureref{\locationb}{variable-templates}).

As a more comprehensive, albeit largely pedagogical example, consider
the overly simplistic \lstinline!my::Vector! class template along with
other related templates defined within a header file
\lstinline!my_vector.h!:

\begin{emcppslisting}[emcppsbatch={e1,e1b}]
// my_vector.h:
#ifndef INCLUDED_MY_VECTOR  // internal include guard
#define INCLUDED_MY_VECTOR

#include <cstddef>  // (ù{\codeincomments{std::size\_t}}ù)
#include <utility>  // (ù{\codeincomments{std::swap}}ù)

namespace my  // namespace for all entities defined within this component
{

template <typename T>
class Vector
{
    static std::size_t s_count; // track number of objects constructed
    T* d_data_p;                // pointer to dynamically allocated memory
    std::size_t d_length;       // current number of elements in the vector
    std::size_t d_capacity;     // number of elements currently allocated

public:
    // ...

    std::size_t length() const { return d_length; }
        // return the number of elements

    // ...
};

// ...          Any partial or full specialization definitions            ...
// ...          of the class template (ù{\codeincomments{Vector}}ù) go here.                     ...

template <typename T>
void swap(Vector<T> &lhs, Vector<T> &rhs) { return std::swap(lhs, rhs); }
    // free function that operates on objects of type (ù{\codeincomments{my::Vector}}ù) via ADL

// ...            Any [full] specialization definitions                   ...
// ...            of free function (ù{\codeincomments{swap}}ù) would go here.                    ...

template <typename T>
const std::size_t vectorSize = sizeof(Vector<T>);  // C++14 variable template
    // This nonmodifiable (ù{\codeincomments{static}}ù) variable holds the size of a (ù{\codeincomments{my::Vector<T>}}ù).

// ...           Any [full] specialization definitions                    ...
// ...           of variable (ù{\codeincomments{vectorSize}}ù) would go here.                    ...

template <typename T>
std::size_t Vector<T>::s_count = 0;
    // definition of (ù{\codeincomments{static}}ù) counter in general template

// ... We may opt to add explicit-instantiation declarations here;
//     see the next code example.

}  // close (ù{\codeincomments{my}}ù) namespace

#endif  // close internal include guard
\end{emcppslisting}

\noindent In the \lstinline!my_vector! component in the code snippet above, we
have defined the following, in the \lstinline!my! namespace:
\begin{enumerate}
\item{a \lstinline!class! template, \lstinline!Vector!, parameterized on element type}
\item{a free-function template, \lstinline!swap!, that operates on objects of corresponding specialized \lstinline!Vector! type}
\item{a \lstinline!const! C++14 variable template,
%(see \featureref{\locationb}{variable-templates}),
\lstinline!vectorSize!, that represents the number of bytes in the \romeogloss{footprint} of an object of the corresponding specialized \lstinline!Vector! type}
\end{enumerate}
Any use of these templates by a client might and typically will trigger
the depositing of equivalent definitions as object code in the client
translation unit's resulting \lstinline!.o! file, irrespective of whether
the definition being used winds up getting inlined.

To eliminate object code for specializations of entities in the
\lstinline!my_vector! component, we must first decide where the unique
definitions will go; see \intraref{annoyances-externtemplate}{no-good-place-to-put-definitions-for-unrelated-classes}. In this
specific case, however, we own the component that requires
specialization, and the specialization is for a ubiquitous built-in type;
hence, the natural place to generate the specialized definitions is in a
\lstinline!.cpp! file corresponding to the component's header:

\begin{emcppslisting}[emcppsbatch=e1,emcppsstandards=c++14]
// my_vector.cpp:
#include <my_vector.h>  // We always include the component's own header first.
    // By including this header file, we have introduced the general template
    // definitions for each of the explicit-instantiation declarations below.

namespace my  // namespace for all entities defined within this component
{

template class Vector<int>;
    // Generate object code for all non(ù{\codeincomments{template}}ù) member-function and (ù{\codeincomments{static}}ù)
    // member-variable definitions of template (ù{\codeincomments{my::Vector}}ù) having (ù{\codeincomments{int}}ù) elements.

template std::size_t Vector<double>::length() const;  // BAD IDEA
    // In addition, we could generate object code for just a particular member
    // function definition of (ù{\codeincomments{my::Vector}}ù) (e.g., (ù{\codeincomments{length}}ù)) for some other
    // parameter type (e.g., (ù{\codeincomments{double}}ù)), which is shown here for pedagogy only.

template void swap(Vector<int>& lhs, Vector<int>& rhs);
    // Generate object code for the full specialization of the (ù{\codeincomments{swap}}ù) free-
    // function template that operates on objects of type (ù{\codeincomments{my::Vector<int>}}ù).

template const std::size_t vectorSize<int>;  // C++14 variable template
    // Generate the object-code-level definition for the specialization of the
    // C++14 variable template instantiated for built-in type (ù{\codeincomments{int}}ù).

//template std::size_t Vector<int>::s_count = 0;
    // Generate the object-code-level definition for the specialization of the
    // (ù{\codeincomments{static}}ù) member variable of (ù{\codeincomments{Vector}}ù) instantiated for built-in type (ù{\codeincomments{int}}ù).

}; // close (ù{\codeincomments{my}}ù) namespace
\end{emcppslisting}

\noindent Each of the constructs introduced by the keyword \lstinline!template!
within the \lstinline!my! namespace in the code snippet above represents a
separate \romeogloss{explicit-instantiation definition}. These
constructs instruct the compiler to generate object-level definitions
for general templates declared in \lstinline!my_vector.h! specialized on
the built-in type \lstinline!int!.

Having installed the necessary \romeogloss{explicit-instantiation
definitions} in the component's \lstinline!my_vector.cpp! file, we
must now go back to its \lstinline!my_vector.h! file and, without altering
any of the previously existing lines of code, \emph{add} the
corresponding \romeogloss{explicit-instantiation declarations} to
suppress redundant local code generation:

\begin{emcppshiddenlisting}[emcppsbatch=e1b]
#undef INCLUDED_MY_VECTOR
\end{emcppshiddenlisting}
\begin{emcppslisting}[emcppsbatch=e1b]
// my_vector.h:
#ifndef INCLUDED_MY_VECTOR  // internal include guard
#define INCLUDED_MY_VECTOR

namespace my  // namespace for all entities defined within this component
{

// ...
// ...  everything that was in the original (ù{\codeincomments{my}}ù) namespace
// ...

extern template class Vector<int>;
    // Suppress object code for this class template specialized for (ù{\codeincomments{int}}ù).

extern template std::size_t Vector<double>::size() const;  // BAD IDEA
    // Suppress object code for this member, only specialized for (ù{\codeincomments{double}}ù).

extern template void swap(Vector<int>& lhs, Vector<int>& rhs);
    // Suppress object code for this free function specialized for (ù{\codeincomments{int}}ù).

extern template std::size_t vectorSize<int>;  // C++14
    // Suppress object code for this variable template specialized for (ù{\codeincomments{int}}ù).

extern template std::size_t Vector<int>::s_count;
    // Suppress object code for this (ù{\codeincomments{static}}ù) member definition w.r.t. (ù{\codeincomments{int}}ù).

}  // close (ù{\codeincomments{my}}ù) namespace

#endif  // close internal include guard
\end{emcppslisting}

\noindent Each of the constructs that begin with \lstinline!extern!~\lstinline!template!
in the example above are \romeogloss{explicit-instantiation
declarations}, which serve only to suppress the generation of any
object code emitted to the \lstinline!.o! file of the current translation
unit in which such specializations are used. These added
\lstinline!extern!~\lstinline!template! declarations must appear in the
\lstinline!my_header.h! source file \emph{after} the declaration of the
corresponding general template and, importantly, before whatever
relevant definitions are ever used.

\subsubsection[The effect on various \lstinline!.o! files]{The effect on various {\SubsubsecCode .o} files}\label{the-effect-on-various-.o-files}

To illustrate the effect of \romeogloss{explicit-instantiation
declarations} and \romeogloss{explicit-instantiation
definitions} on the contents of object and executable files,
we'll use a simple \lstinline!lib_interval! library \romeogloss{component}
consisting of a header file, \lstinline!lib_interval.h!, and an
implementation file, \lstinline!lib_interval.cpp!. The latter, apart from
including its corresponding header, is effectively empty:

\begin{emcppslisting}[emcppsbatch=e2]
// lib_interval.h:
#ifndef INCLUDED_LIB_INTERVAL  // internal include guard
#define INCLUDED_LIB_INTERVAL

namespace lib  // namespace for all entities defined within this component
{

template <typename T>  // elided sketch of a class template
class Interval
{
    T d_low;   // interval's low value
    T d_high;  // interval's high value

public:
    explicit Interval(const T& p) : d_low(p), d_high(p) { }
        // Construct an empty interval.

    Interval(const T& low, const T& high) : d_low(low), d_high(high) { }
        // Construct an interval having the specified boundary values.

    const T& low() const { return d_low; }
        // Return this interval's low value.

    const T& high() const { return d_high; }
        // Return this interval's high value.

    int length() const { return d_high - d_low; }
        // Return this interval's length.

    // ...
};

template <typename T>                   // elided sketch of a function template
bool intersect(const Interval<T>& i1, const Interval<T>& i2)
    // Determine whether the specified intervals intersect ...
{
    bool result = false;  // nonintersecting until proven otherwise
    // ...
    return result;
}

}  // close (ù{\codeincomments{lib}}ù) namespace

#endif  // INCLUDED_LIB_INTERVAL
\end{emcppslisting}
\vspace*{2ex}   %%%% extra space to highlight different snippets
\begin{emcppslisting}[emcppsbatch=e2]
// lib_interval.cpp:
#include <lib_interval.h>
\end{emcppslisting}

\noindent This library component defines, in the namespace \lstinline!lib!, a heavily
elided sketch of an implementation of (1) a class template,
\lstinline!Interval!, and (2) a function template, \lstinline!intersect!, the
only practical purpose of which is to provide some sample template
source code to compile.

Let's also consider a trivial application that uses this library
\romeogloss{component}:

\begin{emcppslisting}[emcppsbatch=e2]
// app.cpp:
#include <lib_interval.h>  // Include the library component's header file.

int main(int argv, const char** argc)
{
    lib::Interval<double> a(0, 5);  // instantiate with (ù{\codeincomments{double}}ù) type parameter
    lib::Interval<double> b(3, 8);  // instantiate with (ù{\codeincomments{double}}ù) type parameter
    lib::Interval<int>    c(4, 9);  // instantiate with (ù{\codeincomments{int}}ù)    type parameter

    if (lib::intersect(a, b))  // instantiate deducing (ù{\codeincomments{double}}ù) type parameter
    {
        return 0;  // return "success" as (0.0, 5.0) does intersect (3.0, 8.0)
    }

    return 1;  // Return "failure" status as function apparently doesn't work.
}
\end{emcppslisting}

\noindent The purpose of this application is merely to exhibit a couple of
instantiations of the library \emph{class} template,
\lstinline!lib::Interval!, for type parameters \lstinline!int! and
\lstinline!double!, and of the library \emph{function} template,
\lstinline!lib::intersect!, for just \lstinline!double!.

Next, we compile the application and library translation units,
\lstinline!app.cpp! and\linebreak[4]
\lstinline!lib_interval.cpp!, and inspect the symbols
in their respective corresponding object files, \lstinline!app.o! and
\lstinline!lib_interval.o!:

\begin{lstlisting}[language=bash]
$ gcc -I. -c app.cpp lib_interval.cpp
$ nm -C app.o lib_interval.o

app.o:
0000000000000000 W lib::Interval<double>::Interval(double const&, double const&)
0000000000000000 W lib::Interval<int>::Interval(int const&, int const&)
0000000000000000 W bool lib::intersect<double>(lib::Interval<double> const&,
                                               lib::Interval<double> const&)
0000000000000000 T main

lib_interval.o:
\end{lstlisting}

\noindent Looking at \lstinline!app.o! in the previous example, the class and
function templates used in the \lstinline!main! function, which is
defined in the \lstinline!app.cpp! file, were instantiated
\emph{implicitly} and the relevant code was added to the resulting
object file, \lstinline!app.o!, with each instantiated function definition
in its own separate \romeogloss{section}. In the \lstinline!Interval!
\emph{class} template, the generated symbols correspond to the two
unique instantiations of the constructor, i.e., for \lstinline!double! and
\lstinline!int!, respectively. The \lstinline!intersect! function
template, however, was implicitly instantiated for only type
\lstinline!double!. Note importantly that all of the implicitly
instantiated functions have the \lstinline!W! symbol type, indicating that
they are \emph{weak} symbols, which are permitted to be present in
multiple object files. By contrast, this file defines the strong symbol
\lstinline!main!, marked here by a \lstinline!T!. Linking this file with any
other file containing such a symbol would cause the linker to report a
multiply-defined-symbol error. On the other hand, the
\lstinline!lib_interval.o! file corresponds to the
\lstinline!lib_interval! library component, whose \lstinline!.cpp! file
served only to include its own \lstinline!.h! file, and is again
effectively empty.

Let's now link the two object files, \lstinline!app.o! and
\lstinline!lib_interval.o!, and inspect the symbols in the resulting
executable, \lstinline!app!{\cprotect\footnote{We have stripped out
extraneous unrelated information that the \lstinline!nm! tool produces;
note that the \lstinline!-C! option invokes the symbol demangler, which
turns encoded names like \lstinline!_ZN3lib8IntervalIdEC1ERKdS3_! into
something more readable like
  \lstinline!lib::Interval<double>::Interval(double!~\lstinline!const&,!~\lstinline!double!~\lstinline!const&)!.}}:

\begin{lstlisting}[language=bash]
$ gcc -o app app.o lib_interval.o
$ nm -C app
000000000040056e T lib::Interval<double>::Interval(double const&, double const&)
00000000004005a2 T lib::Interval<int>::Interval(int const&, int const&)
00000000004005ce T bool lib::intersect<double>(lib::Interval<double> const&,
                                               lib::Interval<double> const&)
00000000004004b7 T main
\end{lstlisting}

\noindent As the textual output above confirms, each of the needed \emph{weak}
template symbols, previously marked with a \lstinline!W!, is bound into the
final program as a \emph{strong} symbol, now --- like \lstinline!main! ---
marked with a \lstinline!T!.{\cprotect\footnote{Whether the symbol is
marked \lstinline!W! or \lstinline!T! in the final executable is
implementation specific and of no consequence here. We present these
  concepts in this particular way to aid cognition.}} In this tiny
illustrative example, only one set of weak symbols appeared in the
combined \lstinline!.o! files.

More generally, if the application comprises multiple object files, each
file will potentially contain their own set of weak symbols,
often leading to duplicate code \romeogloss{sections} for implicitly
instantiated class, function, and variable
templates instantiated on the same parameters. When the linker combines
object files, it will arbitrarily choose at most one of each of these
respective and hopefully identical weak-symbol \romeogloss{sections} to
include in the final executable, now marked as a strong symbol
(\lstinline!T!).

Imagine now that our program includes a large number of \lstinline!.o!
files, many of which make use of our \lstinline!lib_interval! component,
particularly to operate on \lstinline!double! intervals. Suppose, for now,
we decide we would like to suppress the generation of object code for
templates related to just \lstinline!double! types with the intent of later putting them all in one place, i.e., the currently
empty \lstinline!lib_interval.o!. Achieving this objective is precisely what the
\lstinline!extern!~\lstinline!template! syntax is designed to accomplish.

Returning to our \lstinline!lib_interval.h! file, we need not change one
line of code; we need only to \emph{add} two
\romeogloss{explicit-instantiation declarations} --- one for the
template \emph{class}, \lstinline!Interval<double>!, and one for the
template \emph{function},
\lstinline!intersect<double>(const!~\lstinline!double&,!~\lstinline!const!~\lstinline!double&)!
--- to the header file anywhere \emph{after} their respective
corresponding general template declaration and definition:

\begin{emcppshiddenlisting}[emcppsbatch=e3]
// lib_interval.h:

#include <cstddef>  // (ù{\codeincomments{std::size\_t}}ù)

namespace lib {
template <typename T>
class Interval
{};

template <typename T>
bool intersect(const Interval<T>&, const Interval<T>&)
{
    return false;
}

}

\end{emcppshiddenlisting}
\begin{emcppslisting}[emcppsbatch=e3]
// lib_interval.h:  // (no change to existing code)
#ifndef INCLUDED_LIB_INTERVAL  // internal include guard
#define INCLUDED_LIB_INTERVAL

namespace lib  // namespace for all entities defined within this component
{

template <typename T>                    // elided sketch of a class template
class Interval;                          // ...
    // ...                       (same as before)

template <typename T>                    // elided sketch of a function template
bool intersect(const Interval<T>& i1, const Interval<T>& i2);
    // ...                       (same as before)

extern template class Interval<double>;  // explicit-instantiation declaration

extern template                          // explicit-instantiation declaration
bool intersect(const Interval<double>&, const Interval<double>&);

}  // close (ù{\codeincomments{lib}}ù) namespace

#endif  // INCLUDED_LIB_INTERVAL
\end{emcppslisting}

\noindent Let's again compile the two \lstinline!.cpp! files and inspect the
corresponding \lstinline!.o! files:

\begin{lstlisting}[language=bash]
$ gcc -I. -c app.cpp lib_interval.cpp
$ nm -C app.o lib_interval.o

app.o:
                 U lib::Interval<double>::Interval(double const&, double const&)
0000000000000000 W lib::Interval<int>::Interval(int const&, int const&)
                 U bool lib::intersect<double>(lib::Interval<double> const&,
                                               lib::Interval<double> const&)
0000000000000000 T main

lib_interval.o:
\end{lstlisting}

\noindent Notice that this time some of the symbols --- specifically those
relating to the \lstinline!class! and \lstinline!function! templates
instantiated for type \lstinline!double! --- have changed from \lstinline!W!,
indicating a \emph{weak} symbol, to \lstinline!U!, indicating an
\emph{undefined} one. What this means is that, instead of generating a
weak symbol for the explicit specializations for \lstinline!double!, the
compiler left those symbols undefined, as if only the
\emph{declarations} of the member and free-function templates had been
available when compiling \lstinline!app.cpp!, yet inlining of the
instantiated definitions is in no way affected. \romeogloss{Undefined
symbols} are symbols that are expected to be made available to the
linker from other object files. Attempting to link this application
expectedly fails because no object files being linked contain the needed
definitions for those instantiations:

\begin{lstlisting}[language=bash]
$ gcc -o app app.o lib_interval.o

app.o: In function ('main':}
app.cpp:(.text+0x38): undefined reference to
  `lib::Interval<double>::Interval(double const&, double const&)'
app.cpp:(.text+0x69): undefined reference to
  `lib::Interval<double>::Interval(double const&, double const&)'
app.cpp:(.text+0xa1): undefined reference to
  `bool lib::intersect<double>(lib::Interval<double> const&,
                               lib::Interval<double> const&)'

collect2: error: ld returned 1 exit status
\end{lstlisting}

\noindent To provide the missing definitions, we will need to instantiate them
explicitly. Since the type for which the class and function are being
specialized is the ubiquitous built-in type, \lstinline!double!, the ideal
place to sequester those definitions would be within the \lstinline!.o!
file of the \lstinline!lib_interval! library component itself, but see \intraref{annoyances-externtemplate}{no-good-place-to-put-definitions-for-unrelated-classes}. To force the needed
template definitions into the \lstinline!lib_interval.o! file, we will
need to pull out our trusty \romeogloss{explicit-instantiation
definition} syntax, i.e., the \lstinline!template! prefix:

\begin{emcppslisting}[emcppsbatch=e3]
// lib_interval.cpp:
#include <lib_interval.h>

template class lib::Interval<double>;
    // example of an explicit-instantiation definition for a class

template bool lib::intersect(const Interval<double>&, const Interval<double>&);
    // example of an explicit-instantiation definition for a function
\end{emcppslisting}

\noindent We recompile once again and inspect our newly generated object files:

\begin{lstlisting}[language=bash]
$ gcc -I. -c app.cpp lib_interval.cpp
$ nm -C app.o lib_interval.o

app.o:
                 U lib::Interval<double>::Interval(double const&, double const&)
0000000000000000 W lib::Interval<int>::Interval(int const&, int const&)
                 U bool lib::intersect<double>(lib::Interval<double> const&,
                                               lib::Interval<double> const&)
0000000000000000 T main

lib_interval.o:
0000000000000000 W lib::Interval<double>::Interval(double const&)
0000000000000000 W lib::Interval<double>::Interval(double const&, double const&)
0000000000000000 W lib::Interval<double>::low() const
0000000000000000 W lib::Interval<double>::high() const
0000000000000000 W lib::Interval<double>::length() const
0000000000000000 W bool lib::intersect<double>(lib::Interval<double> const&,
                                               lib::Interval<double> const&)
\end{lstlisting}

\noindent The application \lstinline!.o! file, \lstinline!app.o!, naturally remained
unchanged. What's new here is that the functions that were missing from
the \lstinline!app.o! file are now available in the
\lstinline!lib_interval.o! file, again as \emph{weak} (\lstinline!W!), as
opposed to strong (\lstinline!T!), symbols. Notice, however, that explicit
instantiation forces the compiler to generate code for all of the member
functions of the class template for a given specialization. These
symbols might all be linked into the resulting executable unless we take
explicit precautions to exclude those that aren't
needed{\cprotect\footnote{To avoid including the explicitly generated
definitions that are being used to resolve undefined symbols, we have
instructed the linker to remove all unused code \romeogloss{sections} from the
executable. The \lstinline!-Wl! option passes comma-separated options to
the linker. The \lstinline!--gc-sections! option instructs the compiler
to compile and assemble and instructs the linker to omit individual
unused \romeogloss{sections}, where each section contains, for example, its own
  instantiation of a function template.}}:

\begin{lstlisting}[language=bash]
$ gcc -o app app.o lib_interval.o -Wl,--gc-sections
$ nm -C app
00000000004005ca T lib::Interval<double>::Interval(double const&, double const&)
000000000040056e T lib::Interval<int>::Interval(int const&, int const&)
000000000040063d T bool lib::intersect<double>(lib::Interval<double> const&,
                                               lib::Interval<double> const&)
00000000004004b7 T main
\end{lstlisting}

\noindent tl;dr: This \lstinline!extern!~\lstinline!template! feature is provided to
enable software architects to reduce code bloat in individual
\lstinline!.o! files for common instantiations of class,
function, and, as of C++14, variable templates in
large-scale C++ software systems. The practical benefit is in reducing
the physical size of libraries, which \emph{might} lead to improved link
times. \romeogloss{Explicit-instantiation declarations} do \emph{not}
(1) affect the meaning of a program, (2) suppress template
instantiation, (3) impede the compiler's ability to \lstinline!inline!, or
(4) meaningfully improve compile time. To be clear, the \emph{only}
purpose of the \lstinline!extern!~\lstinline!template! syntax is to suppress
object-code generation for the current translation unit, which is then
selectively overridden in the translation unit(s) of choice.

\subsection[Use Cases]{Use Cases}\label{use-cases}

\subsubsection[Reducing template code bloat in object files]{Reducing template code bloat in object files}\label{reducing-template-code-bloat-in-object-files}

The motivation for the \lstinline!extern!~\lstinline!template! syntax is as a
purely \romeogloss{physical} (not \romeogloss{logical}) optimization, i.e., to
reduce the amount of redundant code within individual object files
resulting from common template instantiations in client code. As an
example, consider a fixed-size-array class template,
\lstinline!FixedArray!, that is used widely, i.e., by many clients from
separate translation units, in a large-scale \lstinline!game! project for
both integral and floating-point calculations, primarily with type
parameters \lstinline!int! and \lstinline!double! and array sizes of either
\lstinline!2! or \lstinline!3!:

\begin{emcppslisting}[emcppsbatch=e4]
// game_fixedarray.h:
#ifndef INCLUDED_GAME_FIXEDARRAY  // *internal* include guard
#define INCLUDED_GAME_FIXEDARRAY

#include <cstddef>  // (ù{\codeincomments{std::size\_t}}ù)
namespace game  // namespace for all entities defined within this component
{

template <typename T, std::size_t N>
class FixedArray                                  // widely used class template
{
    // ... (elided private implementation details)
public:
    FixedArray()                                 { /*...*/ }
    FixedArray(const FixedArray<T, N>& other)    { /*...*/ }
    T& operator[](std::size_t index)             { /*...*/ }
    const T& operator[](std::size_t index) const { /*...*/ }
};

template <typename T, std::size_t N>
T dot(const FixedArray<T, N>& a, const FixedArray<T, N>& b) { /*...*/ }

// Explicit-instantiation declarations for full template specializations
// commonly used by the (ù{\codeincomments{game}}ù) project are provided below.

extern template class FixedArray<int, 2>;              // class template
extern template int dot(const FixedArray<int, 2>& a,   // function template
                        const FixedArray<int, 2>& b);  // for (ù{\codeincomments{int}}ù) and (ù{\codeincomments{2}}ù)

extern template class FixedArray<int, 3>;              // class template
extern template int dot(const FixedArray<int, 3>& a,   // function template
                        const FixedArray<int, 3>& b);  // for (ù{\codeincomments{int}}ù) and (ù{\codeincomments{3}}ù)

extern template class FixedArray<double, 2>;           // for (ù{\codeincomments{double}}ù) and (ù{\codeincomments{2}}ù)
extern template double dot(const FixedArray<double, 2>& a,
                           const FixedArray<double, 2>& b);

extern template class FixedArray<double, 3>;           // for (ù{\codeincomments{double}}ù) and (ù{\codeincomments{3}}ù)
extern template double dot(const FixedArray<double, 3>& a,
                           const FixedArray<double, 3>& b);

}  // close (ù{\codeincomments{game}}ù) namespace

#endif  // INCLUDED_GAME_FIXEDARRAY
\end{emcppslisting}

\noindent Specializations commonly used by the \lstinline!game! project are
provided by the \lstinline!game! library. In the component header in the
example above, we have used the \lstinline!extern!~\lstinline!template! syntax
to suppress object-code generation for instantiations of both the class
template \lstinline!FixedArray! and the function template \lstinline!dot! for
element types \lstinline!int! and \lstinline!double!, each for array sizes
\lstinline!2! and \lstinline!3!. To ensure that these specialized definitions
are available in every program that might need them, we use the
\lstinline!template! syntax counterpart to \emph{force} object-code
generation within just the one \lstinline!.o! corresponding to the
\lstinline!game_fixedarray! library component{\cprotect\footnote{Notice
  that we have chosen \emph{not} to nest the explicit specializations
  (or any other definitions) of entities already declared directly
  within the \lstinline!game! namespace, preferring instead to qualify each
  entity explicitly to be consistent with how we render free-function
  definitions (to avoid self-declaration); see \cite{lakos20}, section~2.5, ``Component Source-Code Organization," pp.~333--342, specifically Figure 2-36b, p.~340. See also \intraref{potential-pitfalls-externtemplate}{corresponding-explicit-instantiation-declarations-and-definitions}.}}:

\begin{emcppslisting}[emcppsbatch=e4]
// game_fixedarray.cpp:
#include <game_fixedarray.h>  // included as first substantive line of code

// Explicit-instantiation definitions for full template specializations
// commonly used by the (ù{\codeincomments{game}}ù)] project are provided below.

template class game::FixedArray<int, 2>;              // class template
template int game::dot(const FixedArray<int, 2>& a,   // function template
                       const FixedArray<int, 2>& b);  // for (ù{\codeincomments{int}}ù) and (ù{\codeincomments{2}}ù)

template class game::FixedArray<int, 3>;              // class template
template int game::dot(const FixedArray<int, 3>& a,   // function template
                       const FixedArray<int, 3>& b);  // for (ù{\codeincomments{int}}ù) and (ù{\codeincomments{3}}ù)

template class game::FixedArray<double, 2>;           // for (ù{\codeincomments{double}}ù) and (ù{\codeincomments{2}}ù)
template double game::dot(const FixedArray<double, 2>& a,
                          const FixedArray<double, 2>& b);

template class game::FixedArray<double, 3>;           // for (ù{\codeincomments{double}}ù) and (ù{\codeincomments{3}}ù)
template double game::dot(const FixedArray<double, 3>& a,
                          const FixedArray<double, 3>& b);
\end{emcppslisting}

\noindent Compiling \lstinline!game_fixedarray.cpp! and examining the resulting
object file shows that the code for all explicitly instantiated classes
and free functions was generated and placed into the object file,
\lstinline!game_fixedarray.o!{\cprotect\footnote{Note that only a subset
  of the relevant symbols have been retained.}}:

\begin{lstlisting}[language=bash]
$ gcc -I. -c game_fixedarray.cpp
$ nm -C game_fixedarray.o
0000000000000000 W game::FixedArray<double, 2ul>::FixedArray(
  game::FixedArray<double, 2ul> const&)
0000000000000000 W game::FixedArray<double, 2ul>::FixedArray()
0000000000000000 W game::FixedArray<double, 2ul>::operator[](unsigned long)
0000000000000000 W game::FixedArray<double, 3ul>::FixedArray(
  game::FixedArray<double, 3ul> const&)
0000000000000000 W game::FixedArray<int, 3ul>::FixedArray()
                                    :
0000000000000000 W double game::dot<double, 2ul>(
  game::FixedArray<double, 2ul> const&, game::FixedArray<double, 2ul> const&)
0000000000000000 W double game::dot<double, 3ul>(
  game::FixedArray<double, 3ul> const&, game::FixedArray<double, 3ul> const&)
0000000000000000 W int game::dot<int, 2ul>(
  game::FixedArray<int, 2ul> const&, game::FixedArray<int, 2ul> const&)
                                    :
0000000000000000 W game::FixedArray<int, 2ul>::operator[](unsigned long) const
0000000000000000 W game::FixedArray<int, 3ul>::operator[](unsigned long) const
\end{lstlisting}

\noindent This \lstinline!FixedArray! class template is used in multiple translation
units within the \lstinline!game! project. The first one contains a set
of geometry utilities:

\begin{emcppshiddenlisting}[emcppsbatch=e5]
// game_fixedarray.h:
#pragma once

#include <cstddef>  // (ù{\codeincomments{std::size\_t}}ù)

template <typename T, std::size_t N>
class FixedArray {};

template <typename T, std::size_t N>
T dot(const FixedArray<T, N>& a, const FixedArray<T, N>& b) { /*...*/ }

// game_unit.h:
#include <game_fixedarray.h>

namespace game {

struct Unit {
      FixedArray<double,3> centerOfMass() const;      
};

}
\end{emcppshiddenlisting}
\begin{emcppslisting}[emcppsbatch=e5]
// app_geometryutil.cpp:

#include <game_fixedarray.h>
#include <game_unit.h>

using namespace game;

void translate(game::Unit* object, const FixedArray<double, 2>& dst)
    // Perform precise movement of the object on 2D plane.
{
    FixedArray<double, 2> objectProjection;
    // ...
}

void translate(game::Unit* object, const FixedArray<double, 3>& dst)
    // Perform precise movement of the object in 3D space.
{
    FixedArray<double, 3> delta;
    // ...
}

bool isOrthogonal(const FixedArray<int, 2>& a1, const FixedArray<int, 2>& a2)
    // Return true if 2d arrays are orthogonal.
{
    return 0 == dot(a1, a2);
}

bool isOrthogonal(const FixedArray<int, 3>& a1, const FixedArray<int, 3>& a2)
    // Return true if 3d arrays are orthogonal.
{
    return 0 == dot(a1, a2);
}
\end{emcppslisting}

\noindent The second one deals with physics calculations:

\begin{emcppslisting}[emcppsbatch=e5]
// app_physics.cpp:

#include <game_fixedarray.h>
#include <game_unit.h>

using namespace game;

void collide(game::Unit* objectA, game::Unit* objectB)
    // Calculate the result of object collision in 3D space.
{
    FixedArray<double, 3> centerOfMassA = objectA->centerOfMass();
    FixedArray<double, 3> centerOfMassB = objectB->centerOfMass();
    // ..
}

void accelerate(game::Unit* object, const FixedArray<double, 3>& force)
    // Calculate the position after applying a specified force for the
    // duration of a game tick.
{
    // ...
}
\end{emcppslisting}

\noindent Note that the object files for the application components throughout the
\lstinline!game! project do not contain any of the implicitly
instantiated definitions that we had chosen to uniquely sequester
externally, i.e., within the \lstinline!game_fixedarray.o! file:

\begin{lstlisting}[language=bash]
$ nm -C app_geometryutil.o
000000000000003e T isOrthogonal(game::FixedArray<int, 2ul> const&,
  game::FixedArray<int, 2ul> const&)
0000000000000068 T isOrthogonal(game::FixedArray<int, 3ul> const&,
  game::FixedArray<int, 3ul> const&)
0000000000000000 T translate(game::Unit*, game::FixedArray<double, 2ul> const&)
000000000000001f T translate(game::Unit*, game::FixedArray<double, 3ul> const&)
                 U game::FixedArray<double, 2ul>::FixedArray()
                 U game::FixedArray<double, 3ul>::FixedArray()
                 U int game::dot<int, 2ul>(game::FixedArray<int, 2ul> const&,
  game::FixedArray<int, 2ul> const&)
                 U int game::dot<int, 3ul>(game::FixedArray<int, 3ul> const&,
  game::FixedArray<int, 3ul> const&)

$ nm -C app_physics.o
0000000000000039 T accelerate(game::Unit*,
  game::FixedArray<double, 3ul> const&)
0000000000000000 T collide(game::Unit*, game::Unit*)
                 U game::FixedArray<double, 3ul>::FixedArray()
0000000000000000 W game::Unit::centerOfMass()
\end{lstlisting}

\noindent Whether optimization involving \romeogloss{explicit-instantiation
directives} reduces library sizes on disc has no noticeable effect or
actually makes matters worse will depend on the particulars of the
system at hand. Having this optimization applied to frequently used
templates across a large organization has been known to decrease object
file sizes, storage needs, link times, and overall build times, but see \intraref{potential-pitfalls-externtemplate}{accidentally-making-matters-worse}.

\subsubsection[Insulating template definitions from clients]{Insulating template definitions from clients}\label{insulating-template-definitions-from-clients}

Even before the introduction of \romeogloss{explicit-instantiation
\emph{declarations}}, strategic use of \romeogloss{explicit-instantiation
\emph{definitions}} made it possible to \romeogloss{insulate} the
\emph{definition} of a template from client code, presenting instead
just a limited set of instantiations against which clients may link.
Such insulation enables the definition of the template to change without
forcing clients to recompile. What's more, new explicit instantiations
can be added without affecting existing clients.

As an example, suppose we have a single free-function template,
\lstinline!transform!, that operates on only floating-point values:

\begin{emcppslisting}[emcppsbatch=e6]
// transform.h:
#ifndef INCLUDED_TRANSFORM_H
#define INCLUDED_TRANSFORM_H

template <typename T>  // declaration (only) of free-function template
T transform(const T& value);
    // Return the transform of the specified floating-point (ù{\codeincomments{value}}ù).

#endif
\end{emcppslisting}

\noindent Initially, this function template will support just two built-in types,
\lstinline!float! and \lstinline!double!, but it is anticipated to eventually
support the additional built-in type \lstinline!long!~\lstinline!double! and
perhaps even supplementary user-defined types (e.g., \lstinline!Float128!)
to be made available via separate headers (e.g., \lstinline!float128.h!).
By placing only the declaration of the \lstinline!transform! function
template in its component's header, clients will be able to link against
only two supported explicit specializations provided in the
\lstinline!transform.cpp! file:

\begin{emcppslisting}[emcppsbatch=e6]
// transform.cpp:
#include <transform.h>  // Ensure consistency with client-facing declaration.

template <typename T>   // redeclaration/definition of free-function template
T transform(const T& value)
{
    // insulated implementation of transform function template
}

// explicit-instantiation *definitions*
template float transform(const float&);    // Instantiate for type (ù{\codeincomments{float}}ù).
template double transform(const double&);  // Instantiate for type (ù{\codeincomments{double}}ù).
\end{emcppslisting}

\noindent Without the two \romeogloss{explicit-instantiation definitions} in
the \lstinline!transform.cpp! file above, its corresponding object file,
\lstinline!transform.o!, would be empty.

Note that, as of C++11, we \emph{could} place the corresponding
\romeogloss{explicit-instantiation definitions} in the header file
for, say, documentation purposes:

\begin{emcppslisting}
// transform.h:
#ifndef INCLUDED_TRANSFORM_H
#define INCLUDED_TRANSFORM_H

template <typename T>  // declaration (only) of free-function template
T transform(const T& value);
    // Return the transform of the specified floating-point (ù{\codeincomments{value}}ù).

// explicit-instantiation declarations, available as of C++11
extern template float transform(const float&);    // user documentation only;
extern template double transform(const double&);  // has no effect whatsoever

#endif
\end{emcppslisting}

\noindent But because no definition of the \lstinline!transform! free-function
template is visible in the header, no \emph{implicit} instantiation can
result from client use; hence, the two \romeogloss{explicit-instantiation
declarations} above for \lstinline!float! and \lstinline!double!,
respectively, do nothing.

\subsection[Potential Pitfalls]{Potential Pitfalls}\label{potential-pitfalls-externtemplate}

\subsubsection[Corresponding explicit-instantiation declarations and definitions]{Corresponding explicit-instantiation declarations and definitions}\label{corresponding-explicit-instantiation-declarations-and-definitions}

To realize a reduction in object-code size for individual translation
units and yet still be able to link all valid programs successfully into
a well-formed program, several moving parts have to be brought together
correctly:
\begin{enumerate}
\item{Each general template, \lstinline!C<T>!, whose object code bloat is to be optimized must be declared within some designated component’s header file, \lstinline!c.h!.}
\item{The specific definition of each \lstinline!C<T>! relevant to an explicit specialization being optimized — including general, partial-specialization, and full-specialization definitions — must appear in the header file prior to its corresponding \romeogloss{explicit-instantiation declaration}.}
\item{Each \romeogloss{explicit-instantiation declaration} for each specialization of each separate top-level — i.e., class, function, or variable — template must appear in the component’s \lstinline!.h! file \emph{after} the corresponding general template declaration and the relevant general, partial-specialization, or full-specialization definition, but, in practice, always after \emph{all} such definitions, not just the relevant one.}
\item{Each template specialization having an \romeogloss{explicit-instantiation declaration} in the header file must have a corresponding \romeogloss{explicit-instantiation definition} in the component’s implementation file, \lstinline!c.cpp!.}
\end{enumerate}
Absent items (1) and (2), clients would have no way to safely separate
out the usability and inlineability of the template definitions yet
consolidate the otherwise redundantly generated object-level definitions
within just a single translation unit. Moreover, failing to provide the
relevant definition would mean that any clients using one of these
specializations would either fail to compile or, arguably worse, pick up
the general definitions when a more specialized definition was intended,
likely resulting in an ill-formed program.

Failing item (3), the object code for that particular specialization of
that template will be generated locally in the client's translation unit
as usual, negating any benefits with respect to local object-code size,
irrespective of what is specified in the \lstinline!c.cpp! file.

Finally, unless we provide a matching \romeogloss{explicit-instantiation
definition} in the \lstinline!c.cpp! file for each and every
corresponding \romeogloss{explicit-instantiation declaration} in the
\lstinline!c.h! file as in item (4), our optimization attempts might well
result in a library component that compiles, links, and even passes some
unit tests but, when released to our clients, fails to link.
Additionally, any \romeogloss{explicit-instantiation definition} in
the \lstinline!c.cpp! file that is not accompanied by a corresponding
\romeogloss{explicit-instantiation declaration} in the \lstinline!c.h!
file will inflate the size of the \lstinline!c.o! file with no possibility
of reducing code bloat in client code{\cprotect\footnote{Fortunately,
these extra instantiations do not result in multiply-defined symbols
because they still reside in their own \romeogloss{sections} and are marked
  as \emph{weak} symbols.}}:

\begin{emcppslisting}[emcppsbatch=e7]
// c.h:
#ifndef INCLUDED_C                            // internal include guard
#define INCLUDED_C

template <typename T> void f(T v) {/*...*/};  // general template definition

extern template void f<int>(int v);           // OK, matched in (ù{\codeincomments{c.cpp}}ù)
extern template void f<char>(char c);         // Error, unmatched in (ù{\codeincomments{.cpp}}ù) file

#endif
\end{emcppslisting}
\vspace*{2ex}   %%%% extra space to highlight different snippets
\begin{emcppslisting}[emcppsbatch=e7]
// c.cpp:
#include <c.h>                                 // incorporate own header first

extern template void f<int>(int v);            // OK, matched in (ù{\codeincomments{c.h}}ù)
extern template void f<double>(double v);      // Bug, unmatched in (ù{\codeincomments{c.h}}ù) file
\end{emcppslisting}
\vspace*{2ex}   %%%% extra space to highlight different snippets
\begin{emcppslisting}[emcppsbatch=e7]
// client.cpp:
#include <c.h>

void client()
{
    int    i = 1;
    char   c = 'a';
    double d = 2.0;

    f(i);  // OK, matching explicit-instantiation directives
    f(c);  // Link-Time Error, no matching explicit-instantiation definition
    f(d);  // Bug, size increased due to no matching explicit-instantiation
           // declaration
}
\end{emcppslisting}

\noindent In the example above, \lstinline!f(i)! works as expected, with the linker
finding the definition of \lstinline!f<int>! in \lstinline!c.o!; \lstinline!f(c)!
fails to link, because no definition of \lstinline!f<c>! is guaranteed to
be found anywhere; and \lstinline!f(d)! accidentally works by silently
generating a \emph{redundant} local copy of \lstinline!f<double>! in
\lstinline!client.o! while another, identical definition is generated
explicitly in \lstinline!c.o!. Importantly, note that
\lstinline!extern!~\lstinline!template! has \emph{absolutely no affect} on
overload resolution because the call to \lstinline!f(c)! did \emph{not}
resolve to \lstinline!f<int>!.

\subsubsection[Accidentally making matters worse]{Accidentally making matters worse}\label{accidentally-making-matters-worse}

When making the decision to preinstantiate common specializations of
popular templates within some designated \lstinline!.o! file, one must
consider that not all programs necessarily need every (or even any) such
instantiation. Special consideration should be given to classes that
have many methods but typically use only a few.

The language feature is sufficiently flexible that one can suppress and
preinstantiate just one or a handful of member functions of such a type. Intuition is all well and good, but measurement simply has no
substitute.

If one suspects that \romeogloss{explicit-instantiation directives}
might profitably reduce the size of libraries resulting from object code
that is bloated due to redundant local reinstantiations of popular
templates on common types, measuring before and after and retaining the
change \emph{only} if it offers a significant --- at least measurable
--- improvement avoids complicating the codebase without a verifiable
return on the investment. Finally, remember that one might need to
explicitly tell the linker to strip unused \romeogloss{sections} resulting,
for example, from forced instantiation of common template
specializations, to avoid inadvertently bloating executables, which
could adversely affect load times.

\subsection[Annoyances]{Annoyances}\label{annoyances-externtemplate}

\subsubsection[No good place to put definitions for unrelated classes]{No good place to put definitions for unrelated classes}\label{no-good-place-to-put-definitions-for-unrelated-classes}

When we consider the implications of physical
dependency,\footnote{See \cite{lakos96}}\footnote{See \cite{lakos20}} determining in which component to deposit the
specialized definitions can be problematic. For example, consider a
codebase implementing a core library that provides both a nontemplated
\lstinline!String! class and a \lstinline!Vector! container class template.
These fundamentally unrelated entities would ideally live in separate
physical \romeogloss{components} --- i.e., \lstinline!.h!/\lstinline!.cpp! pairs
--- neither of which depends physically on the other. That is, an
application using just one of these components could ideally be
compiled, linked, tested, and deployed entirely independently of the
other. Now, consider a large codebase that makes heavy use of
\lstinline!Vector<String>!: In what component should the object-code-level
definitions for the \lstinline!Vector<String>! specialization
reside?{\cprotect\footnote{Note that the problem of determining in which
component to instantiate the object-level implementation of a template
for a user-defined type is similar to that of specializing an
  arbitrary user-defined trait for a user-defined type.}} There are two
obvious alternatives:
\begin{enumerate}
\item{\lstinline!vector!: In this case, \lstinline!vector.h! would hold  \lstinline!extern!~\lstinline!template!~\lstinline!class!~\lstinline!Vector<String>;! --- the \romeogloss{explicit-instantiation \emph{declaration}} --- and \lstinline!vector.cpp! would hold\linebreak[4]%%%%%
\lstinline!template!~\lstinline!class!~\lstinline!Vector<String>;! --- the \romeogloss{explicit-instantiation \emph{definition}}. With this approach, we would create a physical dependency of the \lstinline!vector! component on \lstinline!string!. Any client program wanting to use a \lstinline!Vector! would also depend on \lstinline!string! regardless of whether it was needed.}
\item{\lstinline!string!: In this case, \lstinline!string.h! and \lstinline!string.cpp! would instead be modified so as to depend on \lstinline!vector!. Clients wanting to use a \lstinline!string! would also be forced to depend physically on \lstinline!vector! \emph{at} \emph{compile} \emph{time}.}
\end{enumerate}
Another possibility might be to create a third component, call it
\lstinline!stringvector!, that itself depends on both \lstinline!vector!
and \lstinline!string!. By
\romeogloss{escalating}\footnote{\cite{lakos20}, section~3.5.2, ``Escalation," pp.~604--614} the mutual dependency
to a higher level in the physical hierarchy, we avoid forcing any client
to depend on more than what is actually needed. The practical drawback
to this approach is that only those clients that proactively include the
composite \lstinline!stringvector.h! header would realize any benefit;
fortunately, in this case, there is no \romeogloss{one-definition rule (ODR)} issue if they don't.

Finally, complex machinery could be added to both \lstinline!string.h! and
\lstinline!vector.h! to conditionally include \lstinline!stringvector.h!
whenever both of the other headers are included; such heroic efforts
would, nonetheless, involve a \romeogloss{cyclic physical dependency} among
all three of these components. Circular intercomponent collaborations
are best avoided.\footnote{\cite{lakos20}, section~3.4, ``Avoiding Cyclic Link-Time Dependencies," pp.~592--601}

\subsubsection[All members of an explicitly defined template class must be valid]{All members of an explicitly defined template class must be valid}\label{all-members-of-an-explicitly-defined-template-class-must-be-valid}

In general, when using a template class, only those members that are
actually used get implicitly instantiated. This hallmark allows class
templates to provide functionality for parameter types having certain
capabilities (e.g., default constructible) while also providing partial
support for types lacking those same capabilities. When providing an
\romeogloss{explicit-instantiation definition}, however, \emph{all}
members of a template class are instantiated.

Consider a simple class template having a data member that can be either
default-initialized (via the template's default constructor) or
initialized with an instance of the member's type (supplied at
construction):

\begin{emcppslisting}[emcppsbatch=e8]
template <typename T>
class W
{
    T d_t;  // a data member of type (ù{\codeincomments{T}}ù)

public:
    W() : d_t() {}
        // Create an instance of (ù{\codeincomments{W}}ù) with a default-constructed (ù{\codeincomments{T}}ù) member.

    W(const T& t) : d_t(t) {}
        // Create an instance of (ù{\codeincomments{W}}ù) with a copy of the specified (ù{\codeincomments{t}}ù).

    void doStuff() { /* do stuff */ }
};
\end{emcppslisting}

\noindent This class template can be used successfully with a type, such as
\lstinline!U! in the code snippet below, that is not default constructible:

\begin{emcppslisting}[emcppsbatch=e8]
struct U
{
    U(int i) { /* do something with (ù{\codeincomments{i}}ù) */ }
    // ...
};

void useWU()
{
    W<U> wu1(U(17));  // OK, using copy constructor for (ù{\codeincomments{U}}ù)
    wu1.doStuff();
}
\end{emcppslisting}

\noindent As it stands, the code above is well formed even though
\lstinline!W<U>::W()! would fail to compile if instantiated. Consequently,
although providing an \romeogloss{explicit-instantiation declaration}
for \lstinline!W<U>! is valid, a corresponding
\romeogloss{explicit-instantiation definition} for \lstinline!W<U>!
fails to compile, as would an implicit instantiation of
\lstinline!W<U>::W()!:

\begin{emcppslisting}[emcppsbatch=e8]
extern template class W<U>;  // Valid: Suppress implicit instantiation of (ù{\codeincomments{W<U>}}ù).

template class W<U>;         // Error, (ù{\codeincomments{U::U()}}ù) not available for (ù{\codeincomments{W<U>::W()}}ù)

void useWU0()
{
    W<U> wu0{};       // Error, (ù{\codeincomments{U::U()}}ù) not available for (ù{\codeincomments{W<U>::W()}}ù)
}
\end{emcppslisting}

\noindent Unfortunately, the only workaround to achieve a comparable reduction in
code bloat is to provide member-specific \romeogloss{explicit-instantiation
directives} for each valid member of \lstinline!W<U>!, an approach
that would likely carry a significantly greater maintenance burden:

\begin{emcppslisting}[emcppsbatch=e8]
extern template W<U>::W(const U& u);   // suppress individual member
extern template void W<U>::doStuff();  //    "         "        "
// ... Repeat for all other functions in (ù{\codeincomments{W}}ù) except (ù{\codeincomments{W<U>::W()}}ù).

template W<U>::W(const U& u);          // instantiate individual member
template void W<U>::doStuff();         //       "         "        "
// ... Repeat for all other functions in (ù{\codeincomments{W}}ù) except (ù{\codeincomments{W<U>::W()}}ù).
\end{emcppslisting}

\noindent The power and flexibility to make it all work --- albeit annoyingly ---
are there nonetheless.

\subsection[See Also]{See Also}\label{see-also}

\begin{itemize}
\item{\seealsoref{variable-templates}{\locationb}Extension of the template syntax for defining a family of like-named variables or static data members that can be instantiated explicitly}
\end{itemize}

\subsection[Further Reading]{Further Reading}\label{further-reading}

\begin{itemize}
\item{For a different perspective on this feature, see \cite{lakos20},
section 1.3.16, ``\texttt{extern} Templates,'' pp.~183--185.}
\item{For a more complete discussion of how compilers and linkers work with
respect to C++, see \cite{lakos20}, Chapter 1, ``Compilers, Linkers, and
Components,'' pp.~123--268.}
\end{itemize}
 

\newpage
%\section[Forwarding References]{Forwarding {\SecCode \&\&} References}\label{forwardingref}\label{forwarding-references}

% 01 Feb, 2021 JMB - packet 4 compilation fixes

\emcppsFeature{
    short={Forwarding References},
    long={Forwarding {\SecCode \&\&} References},
    toclong={Forwarding \lstinline!\&\&! References},
}{forwardingref}
\setcounter{table}{0}
\setcounter{footnote}{0}
\setcounter{lstlisting}{0}
%\section[Forwarding References]{Forwarding {\SecCode \&\&} References}\label{forwardingref}

A forwarding reference (\texttt{T\&\&}) --- distinguishable from an
rvalue reference (\texttt{\&\&}) (see ``\titleref{Rvalue-References}" on page~\pageref{Rvalue-References}) based only on context
--- is a distinct, special kind of reference that (1) binds
universally to the result of an expression of \emph{any} \textbf{value
category} and (2) preserves aspects of that \textbf{value category} so
that the bound object can be \emph{moved from}, if appropriate.

\subsection[Description]{Description}\label{description-forwardingref}

Sometimes we want the same reference to be able to bind to either an
\emph{lvalue} or an \emph{rvalue} and then later be able to discern,
from the reference itself, whether the result of the original expression
was eligible to be \emph{moved from}. A \emph{forwarding reference}
(e.g., \texttt{forRef} in the example below) used in the interface of a function
\emph{template} (e.g., \texttt{myFunc}) affords precisely this
capability and will prove invaluable for the purpose of conditionally
moving, or else copying, an object from within the function template's
body:

\begin{emcppslisting}
template <typename T>
void myFunc(T&& forRef)
{
    // It is possible to check if (ù{\codeincomments{forRef}}ù) is eligible to be moved from or not
    // from within the body of (ù{\codeincomments{myFunc}}ù).
}
\end{emcppslisting}

\noindent In the definition of the \texttt{myFunc} function template in the example above, the
parameter \texttt{forRef} syntactically appears to be a
non-\texttt{const} reference to an \emph{rvalue} of type \texttt{T}; in
this very precise context, however, the same \texttt{T\&\&} syntax
designates a \textbf{forwarding reference}, with the effect of retaining
the original value category of the object bound to the argument
\texttt{forRef}; see \textit{\titleref{description-forwardingref}: \titleref{identifying-forwarding-references}} on page~\pageref{identifying-forwarding-references}. The \texttt{T\&\&} syntax
represents a \emph{forwarding} reference --- as opposed to an
\emph{rvalue} reference --- whenever an \emph{individual} function template has a type parameter (e.g., \texttt{T}) and
an unqualified function parameter of type that is exactly
\texttt{T\&\&} (e.g., \texttt{const}~\texttt{T\&\&} would be an
  \emph{rvalue} reference, not a forwarding reference).

Consider, for example, a function \texttt{f} that takes a single
argument by reference and then attempts to use it to invoke one of two
overloads of a function \texttt{g}, depending on whether the original
argument was an \emph{lvalue} or \emph{rvalue}:

\begin{emcppslisting}[emcppsbatch=e1]
struct S { /* some UDT that might benefit from being able to be moved */ };

void g(const S&);  // target function - overload for (ù{\codeincomments{const}}ù) (ù{\codeincomments{int}}ù) lvalues
void g(S&&);       // target function - overload for (ù{\codeincomments{int}}ù) rvalues only

template <typename T>
void f(T&& forRef); // forwards to target overload (ù{\codeincomments{g}}ù) based on value category
\end{emcppslisting}

\noindent In
theory, we could have chosen a non-\texttt{const} \emph{lvalue}
reference along with a modifiable \emph{rvalue} reference here for
\emph{pedagogical} symmetry; such an inherently unharmonious overload
set would, however, not typically occur in practice; see ``\titleref{Rvalue-References}" on page~\pageref{Rvalue-References}. In this specific case --- where \texttt{f} is a function template,
\texttt{T} is a template type parameter, and the type of the parameter
itself is exactly \texttt{T\&\&} --- the \texttt{forRef} function
parameter (in the code snippet above) denotes a \emph{forwarding reference}. If \texttt{f}\label{f-invoked-example}
is invoked with an \emph{lvalue}, \texttt{forRef} is an \emph{lvalue}
reference; otherwise \texttt{forRef} is an \emph{rvalue} reference.
Given the dual nature of \texttt{forRef}, one rather verbose way of
determining the original value category of the passed argument would be
to use the \texttt{std::is\_lvalue\_reference} \textbf{type trait} on
\texttt{forRef} itself:

\begin{emcppslisting}[emcppsbatch=e1]
#include <type_traits>  // (ù{\codeincomments{std::is\_lvalue\_reference}}ù)
#include <utility>  // (ù{\codeincomments{std::move}}ù)

template <typename T>
void f(T&& forRef)      // (ù{\codeincomments{forRef}}ù) is a forwarding reference.
{
    if (std::is_lvalue_reference<T>::value)  // using a C++11 type trait
    {
        g(forRef);             // propagates (ù{\codeincomments{forRef}}ù) as an *lvalue*
    }                          // invokes (ù{\codeincomments{g(const S\&)}}ù)
    else
    {
        g(std::move(forRef));  // propagates (ù{\codeincomments{forRef}}ù) as an *rvalue*
    }                          // invokes (ù{\codeincomments{g(S\&\&)}}ù)
}
\end{emcppslisting}

\noindent The \texttt{std::is\_lvalue\_reference<T>::value} predicate above asks
the question, ``Did the object bound to \texttt{fRef} originate
from an lvalue expression?'' and allows the developer to branch on the
answer. A more concise but otherwise equivalent implementation is
generally preferred; see \textit{\titleref{description-forwardingref}: \titleref{the-std::forward-utility}} on page~\pageref{the-std::forward-utility}:

\begin{emcppshiddenlisting}[emcppsbatch=e2]
struct S {};
void g(const S&);
void g(S&&);
\end{emcppshiddenlisting}
\begin{emcppslisting}[emcppsbatch=e2]
#include <utility>  // (ù{\codeincomments{std::forward}}ù)

template <typename T>
void f(T&& forRef)
{
    g(std::forward<T>(forRef));
        // same as (ù{\codeincomments{g(std::move(forRef))}}ù) if and only if (ù{\codeincomments{forRef}}ù) is an *rvalue*
        // reference; otherwise, equivalent to (ù{\codeincomments{g(forRef)}}ù)
}
\end{emcppslisting}

\noindent A \texttt{client} function invoking \texttt{f} will enjoy the same
behavior with either of the two implementation alternatives offered
above:

\begin{emcppslisting}[emcppsbatch=e2]
void client()
{
    S s;
    f(s);    // Instantiates (ù{\codeincomments{f<S\&>}}ù) -- (ù{\codeincomments{fRef}}ù) is an lvalue reference ((ù{\codeincomments{S\&}}ù)).
             // The function (ù{\codeincomments{f<S\&>}}ù) will end up invoking (ù{\codeincomments{g(S\&)}}ù).

    f(S());  // Instantiates (ù{\codeincomments{f<S\&\&>}}ù) -- (ù{\codeincomments{fRef}}ù) is an rvalue reference ((ù{\codeincomments{S\&\&}}ù)).
             // The function (ù{\codeincomments{f<S\&\&>}}ù) will end up invoking (ù{\codeincomments{g(S\&\&)}}ù).
}
\end{emcppslisting}

\noindent Use of \texttt{std::forward} in combination with forwarding references
is typical in the implementation of industrial-strength generic
libraries; see \textit{\titleref{use-cases-forwardingref}} on page~\pageref{use-cases-forwardingref}.

\subsubsection[A brief review of function template argument deduction]{A brief review of function template argument deduction}\label{a-brief-review-of-function-template-argument-deduction}

Invoking a function template without explicitly providing template
arguments at the call site will compel the compiler to attempt, if
possible, to \emph{deduce} those template \emph{type} arguments from the
function arguments:

\begin{emcppslisting}
template <typename T> void f();
template <typename T> void g(T x);
template <typename T> void h(T y, T z);

void example0()
{
    f();        // Error, couldn't infer template argument (ù{\codeincomments{T}}ù)
    g(0);       // OK, (ù{\codeincomments{T}}ù) deduced as (ù{\codeincomments{int}}ù) from literal (ù{\codeincomments{0}}ù) -- (ù{\codeincomments{x}}ù) is an (ù{\codeincomments{int}}ù).
    h(0, 'a');  // Error, deduced conflicting types for (ù{\codeincomments{T}}ù) ((ù{\codeincomments{int}}ù) vs. (ù{\codeincomments{char}}ù))
}
\end{emcppslisting}

\noindent Any \textbf{cv-qualifiers} (\texttt{const}, \texttt{volatile}, or both)
on a \emph{deduced} function parameter will be applied \emph{after} type
deduction is performed:

\begin{emcppslisting}
template <typename T> void cf(const T x);
template <typename T> void vf(volatile T y);
template <typename T> void wf(const volatile T z);

void example1()
{
    cf(0);  // OK, (ù{\codeincomments{T}}ù) deduced as (ù{\codeincomments{int}}ù) -- (ù{\codeincomments{x}}ù) is a (ù{\codeincomments{const int}}ù).
    vf(0);  // OK, (ù{\codeincomments{T}}ù) deduced as (ù{\codeincomments{int}}ù) -- (ù{\codeincomments{y}}ù) is a (ù{\codeincomments{volatile int}}ù).
    wf(0);  // OK, (ù{\codeincomments{T}}ù) deduced as (ù{\codeincomments{int}}ù) -- (ù{\codeincomments{z}}ù) is a (ù{\codeincomments{const}}ù) (ù{\codeincomments{volatile}}ù) (ù{\codeincomments{int}}ù).
}
\end{emcppslisting}

\noindent Similarly, \textbf{ref-qualifiers} other than \texttt{\&\&}
(\texttt{\&} or \texttt{\&\&} along with any cv-qualifier) do not
alter the deduction process, and they too are applied after deduction:

\begin{emcppslisting}
template <typename T> void rf(T& x);
template <typename T> void crf(const T& x);

void example2()
{
    int i;
    rf(i);   // OK, (ù{\codeincomments{T}}ù) is deduced as (ù{\codeincomments{int}}ù) -- (ù{\codeincomments{x}}ù) is an (ù{\codeincomments{int\&}}ù).
    crf(i);  // OK, (ù{\codeincomments{T}}ù) is deduced as (ù{\codeincomments{int}}ù) -- (ù{\codeincomments{x}}ù) is a (ù{\codeincomments{const int\&}}ù).

    rf(0);   // Error, expects an lvalue for 1st argument
    crf(0);  // OK, (ù{\codeincomments{T}}ù) is deduced as (ù{\codeincomments{int}}ù) -- (ù{\codeincomments{x}}ù) is a (ù{\codeincomments{const int\&}}ù).
}
\end{emcppslisting}

\noindent Type deduction works differently for \emph{forwarding} references
where the only qualifier on the template argument is \texttt{\&\&}. For
the sake of exposition, consider a function template declaration,
\texttt{f}, accepting a forwarding reference, \texttt{forRef}:

\begin{emcppslisting}[emcppsbatch=e3]
template <typename T> void f(T&& forRef);
\end{emcppslisting}

 %%%%%%%% PRODUCTION: make sure we are as precise as we can be with this xref on FPPs.
\noindent We have seen in the example on page~\pageref{f-invoked-example} that,
when \texttt{f} is invoked with an
\emph{lvalue} of type \texttt{S}, then \texttt{T} is deduced as
\texttt{S\&} and \texttt{forRef} becomes an \emph{lvalue} reference.
When \texttt{f} is instead invoked with an \emph{xvalue} of type
\texttt{S}, then \texttt{T} is deduced as \texttt{S} and
\texttt{forRef} becomes an \emph{rvalue} reference. The underlying
process that results in this ``duality'' relies on a special rule (known
as \textbf{reference collapsing}; see the next section) introduced as part of
\textbf{type deduction}. When the type \texttt{T} of a
\emph{forwarding} reference is being deduced from an expression
\texttt{E}, \texttt{T} itself will be deduced as an \emph{lvalue}
reference if \texttt{E} is an \emph{lvalue}; otherwise normal
type-deduction rules will apply and \texttt{T} will be deduced as an \emph{rvalue}
reference:

\begin{emcppslisting}[emcppsbatch=e3]
void g()
{
    int i;
    f(i);  // (ù{\codeincomments{i}}ù) is an *lvalue* expression.
           // (ù{\codeincomments{T}}ù) is therefore deduced as (ù{\codeincomments{int\&}}ù) -- special rule!
           // (ù{\codeincomments{T\&\&}}ù) becomes (ù{\codeincomments{int\& \&\&}}ù), which collapses to (ù{\codeincomments{int\&}}ù).

    f(0);  // (ù{\codeincomments{0}}ù) is an *rvalue* expression.
           // (ù{\codeincomments{T}}ù) is therefore deduced as (ù{\codeincomments{int}}ù).
           // (ù{\codeincomments{T\&\&}}ù) becomes (ù{\codeincomments{int\&\&}}ù), which is an *rvalue* reference.
}
\end{emcppslisting}

\noindent For more on general type deduction, see ``\titleref{auto-feature}" on page~\pageref{auto-feature}.

\subsubsection[Reference collapsing]{Reference collapsing}\label{reference-collapsing}

As we saw in the previous section, when a function having a
\emph{forwarding} reference parameter, \texttt{forRef}, is invoked with
a corresponding \emph{lvalue} argument, an interesting phenomenon
occurs: After type deduction, we temporarily get what appears
syntactically to be an \emph{rvalue} reference to an \emph{lvalue}
reference. As references to references are \emph{not} allowed in C++,
the compiler employs \textbf{reference collapsing} to resolve the
\emph{forwarding}-reference parameter, \texttt{forRef}, into a single
reference, thus providing a way to infer, from \texttt{T} itself, the
original \textbf{value category} of the argument passed to \texttt{f}.

The process of \textbf{reference collapsing} takes place automatically
in any situation where a reference to a reference is formed. Table~\ref{forwardingref-table1}
illustrates the simple rules for collapsing ``unstable" references
into ``stable" ones. Notice, in particular, that an \emph{lvalue}
reference always overpowers an \emph{rvalue} reference. The only
situation in which two references collapse into an \emph{rvalue}
reference is when they are both \emph{rvalue} references.
\begin{table}[h!]
\begin{center}
\begin{threeparttable}
\caption{Collapsing ``unstable'' reference pairs into a single
``stable'' one}\label{forwardingref-table1}\vspace{1.5ex}
{\small \begin{tabular}{c|c|c}\thickhline
\rowcolor[gray]{.9}   {\sffamily\bfseries 1st Reference Type} & {\sffamily\bfseries 2nd Reference Type} &
{\sffamily\bfseries Result of Reference Collapsing} \\ \hline
{\tt \&} &{\tt \&} & {\tt \&} \\ \hline
{\tt \&} & {\tt \&\&} &{\tt \&} \\ \hline
{\tt \&\&} & {\tt \&} & {\tt \&} \\ \hline
{\tt \&\&} & {\tt \&\&} & {\tt \&\&} \\ \thickhline
\end{tabular}
}
\end{threeparttable}
\end{center}
\end{table}

Finally, note that it is not possible to write a reference-to-reference
type in C++\linebreak%%%%%
 \mbox{explicitly}:

\begin{emcppslisting}
int    i   = 0;   // OK
int&   ir  = i;   // OK
int& & irr = ir;  // Error, (ù{\codeincomments{irr}}ù) declared as a reference to a reference
\end{emcppslisting}

\noindent It is, however, easy to do so with type aliases and template parameters,
and that is where reference collapsing comes into play:

\begin{emcppslisting}[emcppsbatch=e4]
#include <type_traits>  // (ù{\codeincomments{std::is\_same)}ù)
using i = int&;  // OK
using j = i&;    // OK, (ù{\codeincomments{int\& \&}}ù) becomes (ù{\codeincomments{int\&}}ù).
static_assert(std::is_same<j,int&>::value);
\end{emcppslisting}

\noindent During computations involving \textbf{metafunctions}, or as part of
language rules (such as type deduction), however, references to
references can occur spontaneously:

\begin{emcppslisting}
template <typename T>
struct AddLvalueRef { typedef T& type; };
    // metafunction that transforms to an *lvalue* reference to (ù{\codeincomments{T}}ù)

template <typename T>
void f(T input)
{
    typename AddLvalueRef<T>::type ir1 = input;    // OK, adds (ù{\codeincomments{\&}}ù) to make (ù{\codeincomments{T\&}}ù)
    typename AddLvalueRef<T&>::type ir2 = input;   // OK, collapses to (ù{\codeincomments{T\&}}ù)
    typename AddLvalueRef<T&&>::type ir3 = input;  // OK, collapses to (ù{\codeincomments{T\&}}ù)
}
\end{emcppslisting}

\noindent Notice that we are using the \texttt{typename} keyword in the example above as a
generalized way of indicating, during \textbf{template instantiation},
that a dependent name is a type (as opposed to a value).

\subsubsection[Identifying forwarding references]{Identifying forwarding references}\label{identifying-forwarding-references}

The syntax for a \emph{forwarding} reference (\texttt{\&\&}) is the same
as that for \emph{rvalue} references; the only way to discern one from
the other is by observing the surrounding context. When used in a manner
where \textbf{type deduction} can take place, the \texttt{T\&\&} syntax
does \emph{not} designate an \emph{rvalue} reference; instead, it
represents a \emph{forwarding} reference; for type deduction to be in
effect, an \emph{individual} (possibly member) function \emph{template}
must have a type parameter (e.g., \texttt{T}) and a function parameter of
type that exactly matches that parameter followed by \texttt{\&\&}
(e.g., \texttt{T\&\&}):

\begin{emcppslisting}
struct S0
{
    template <typename T>
    void f(T&& forRef);
        // OK, fully eligible for template-argument type deduction: (ù{\codeincomments{forRef}}ù)
        // is a forwarding reference.
};
\end{emcppslisting}

\noindent Note that if the function parameter is qualified, the syntax reverts to
the usual meaning of \emph{rvalue} reference:

\begin{emcppslisting}
struct S1
{
    template <typename T>
    void f(const T&& crRef);
        // Eligible for type deduction but is not a forwarding reference: due
        // to the (ù{\codeincomments{const}}ù) qualifier, (ù{\codeincomments{crRef}}ù) is an *rvalue* reference.
};
\end{emcppslisting}

\noindent If a member function of a class template is not itself also a template,
then its template type parameter will not be deduced:

\begin{emcppslisting}
template <typename T>
struct S2
{
    void f(T&& rRef);
       // Not eligible for type deduction because (ù{\codeincomments{T}}ù) is fixed and known as part
       // of the instantiation of (ù{\codeincomments{S2}}ù): (ù{\codeincomments{rRef}}ù) is an *rvalue* reference.
};
\end{emcppslisting}

\noindent More generally, note that the \texttt{\&\&} syntax can \emph{never}
imply a \emph{forwarding} reference for a function that is not itself a
template; see \textit{\titleref{annoyances-forwardingref}:
\titleref{forwarding-references-look-just-like-rvalue-references}} on page~\pageref{forwarding-references-look-just-like-rvalue-references}.

\subsubsection[{\tt auto\&\&} --- a forwarding reference in a non-parameter context]{{\SubsubsecCode auto\&\&}~{\subsubemdash}~a forwarding reference in a non-parameter context}\label{auto-a-forwarding-reference-in-a-non-parameter-context}

Outside of template function parameters, \emph{forwarding} references
can also appear in the context of variable definitions using the
\texttt{auto} variable (see ``\titleref{auto-feature}" on page~\pageref{auto-feature}) because they too are subject to type
deduction:

\begin{emcppslisting}
void f()
{
    auto&& i = 0;  // (ù{\codeincomments{i}}ù) is a forwarding reference because the type of (ù{\codeincomments{i}}ù) must
                   // be deduced from the initialization expression (ù{\codeincomments{0}}ù).
}
\end{emcppslisting}

\noindent Just like function parameters, \texttt{auto\&\&} resolves to either an
\emph{lvalue} reference or \emph{rvalue} reference depending on the
\textbf{value category} of the initialization expression:

\begin{emcppslisting}
void g()
{
    int i;
    auto&& lv = i;  // (ù{\codeincomments{lv}}ù) is an (ù{\codeincomments{int\&}}ù).

    auto&& rv = 0;  // (ù{\codeincomments{rv}}ù) is an (ù{\codeincomments{int\&\&}}ù).
}
\end{emcppslisting}

\noindent Similarly to \texttt{const}~\texttt{auto\&}, the \texttt{auto\&\&}
syntax binds to anything. In \texttt{auto\&\&}, however, the
\texttt{const}-ness of the bound object is \emph{preserved} rather than
always enforced:

\begin{emcppslisting}
void h()
{
    int        i = 0;
    const int ci = 0;

    auto&& lv  = i;   // (ù{\codeincomments{lv}}ù) is an (ù{\codeincomments{int\&}}ù).
    auto&& clv = ci;  // (ù{\codeincomments{clv}}ù) is a (ù{\codeincomments{const int\&}}ù).
}
\end{emcppslisting}

\noindent Just as with function parameters, the original \textbf{value category}
of the expression used to initialize a \emph{forwarding} reference
variable can be propagated during subsequent function invocation --
e.g., using \texttt{std::forward} (see \textit{\titleref{description-forwardingref}: \titleref{the-std::forward-utility}} on page~\pageref{the-std::forward-utility}):

\begin{emcppshiddenlisting}[emcppsbatch=e5]
#include <utility>  // (ù{\codeincomments{std::forward}}ù)
template <typename T>
T get();
\end{emcppshiddenlisting}
\begin{emcppslisting}[emcppsbatch=e5]
template <typename T>
void use(T&& t); // Here (ù{\codeincomments{use}}ù) also takes a forwarding reference parameter
                 // to do with as it pleases.

template <typename T>
void l()
{
    auto&& fr = get<T>();
        // (ù{\codeincomments{get<T>()}}ù) might be either an *lvalue* or *rvalue* depending on (ù{\codeincomments{T}}ù).

    use(std::forward<decltype(fr)>(fr));  // (ù{\codeincomments{decltype}}ù) is a C++11 feature.
        // Propagate the original value category of (ù{\codeincomments{get<T>()}}ù) into (ù{\codeincomments{use}}ù).
}
\end{emcppslisting}

\noindent Notice that, because (1) \texttt{std::forward} (see the next section) requires the
type of the object that's going to be forwarded as a user-provided
template argument and (2) it is not possible to name the type of
\texttt{fr}, \texttt{decltype} (see ``\titleref{decltype}" on page~\pageref{decltype}) was used in the example above
to retrieve the type of \texttt{fr}.

\subsubsection[Forwarding references without forwarding]{Forwarding references without forwarding}\label{forwarding-references-without-forwarding}

Sometimes deliberately \emph{not} forwarding (see \textit{\titleref{description-forwardingref}: \titleref{the-std::forward-utility}} on page~\pageref{the-std::forward-utility}) an \texttt{auto\&\&} variable or
a forwarding reference function parameter at all can be useful, instead employing
\emph{forwarding} references solely for their \texttt{const}-preserving
and universal binding semantics. As an example, consider the task of
obtaining iterators over a range of an unknown \textbf{value category}:

\begin{emcppshiddenlisting}[emcppsbatch=e6]
template <typename T>
T getRange();
\end{emcppshiddenlisting}
\begin{emcppslisting}[emcppsbatch=e6]
#include <iterator>  // (ù{\codeincomments{std::begin}}ù), (ù{\codeincomments{std::end}}ù),

template <typename T>
void m()
{
    auto&& r = getRange<T>();
        // (ù{\codeincomments{getRange<T>()}}ù) might be either an lvalue or rvalue depending on (ù{\codeincomments{T}}ù).

    auto b = std::begin(r);
    auto e = std::end(r);

    traverseRange(b, e);
}
\end{emcppslisting}

\noindent Using \texttt{std::forward} in the initialization of both \texttt{b} and
\texttt{e} might result in moving from \texttt{r} twice, which is
potentially unsafe (see ``\titleref{Rvalue-References}" on page~\pageref{Rvalue-References}). Forwarding
\texttt{r} only in the initialization of \texttt{e} might avoid issues
caused by moving an object twice but might result in inconsistent
behavior with \texttt{b}, especially if the implementation of \texttt{r}
makes use of reference qualifiers (see ``\titleref{refqualifiers}" on page~\pageref{refqualifiers}).

\subsubsection[The {\tt std::forward} utility]{The {\SubsubsecCode std::forward} utility}\label{the-std::forward-utility}

The final piece of the forwarding reference puzzle is the
\texttt{std::forward} utility function. Since the expression naming a
forwarding reference \texttt{x} is always an \emph{lvalue} --- due to its
reachability either by name or address in virtual memory --- and since our
intention is to move \texttt{x} in case it was an \emph{rvalue} to begin
with, we need a conditional \emph{move} operation that will move
\texttt{x} only in that case and otherwise let \texttt{x} pass through as
an \emph{lvalue}.

The declaration for \texttt{std::forward<T>} is as follows (in
\texttt{<utility>}):

\begin{emcppshiddenlisting}[emcppsbatch=e7]
#include <type_traits>  // (ù{\codeincomments{std::remove\_reference}}ù)
using namespace std; // to avoid qualifying remove_reference
\end{emcppshiddenlisting}
\begin{emcppslisting}[emcppsbatch=e7]
template <class T> T&& forward(typename remove_reference<T>::type& t) noexcept;
template <class T> T&& forward(typename remove_reference<T>::type&& t) noexcept;
\end{emcppslisting}

\noindent The second overload is ill-formed if invoked when \texttt{T} is an
\emph{lvalue} reference type.

Remember that the type \texttt{T} associated with a forwarding reference
is deduced as a reference type if given an \emph{lvalue} reference and
as a non-reference type otherwise. So for a forwarding reference
\texttt{forRef} of type \texttt{T\&\&}, we have two cases:
\begin{itemize}
\item{An \emph{lvalue} of type \texttt{U} was used for initializing \texttt{forRef}, so \texttt{T} is \texttt{U\&}, thus the first overload of \texttt{forward} will be selected and will be of the form \texttt{U\&}~\texttt{forward(U\&}~\texttt{u)}~\texttt{noexcept}, thus just returning the original \emph{lvalue} reference.}
\item{An \emph{rvalue} of type \texttt{U} was used for initializing \texttt{forRef}, so \texttt{T} is \texttt{U}, so the second overload of \texttt{forward} will be selected and will be of the form \texttt{U\&\&}~\texttt{forward(U\&}~\texttt{u)}~\texttt{noexcept}, essentially equivalent to \texttt{std::move}.}
\end{itemize}
Note that, in the body of a function template accepting a forwarding
reference \texttt{T\&\&} named \texttt{x}, \texttt{std::forward<T>(x)}
could be replaced with \texttt{static\_cast<T\&\&>(x)} to achieve the
same effect. Due to \textbf{reference collapsing} rules, \texttt{T\&\&}
will resolve to \texttt{T\&} whenever the original \textbf{value
category} of \texttt{x} was an \emph{lvalue} and to \texttt{T\&\&}
otherwise, thus achieving the \emph{conditional move} behavior
elucidated in \textit{\titleref{description-forwardingref}} on page~\pageref{description-forwardingref}. Using \texttt{std::forward} over \texttt{static\_cast}
will, however, ensure that the types of \texttt{T} and \texttt{x} match,
preventing accidental unwanted conversions and, separately, perhaps also
more clearly expressing the programmer's intent.

\subsection[Use Cases]{Use Cases}\label{use-cases-forwardingref}

\subsubsection[Perfectly forwarding an expression to a downstream consumer]{Perfectly forwarding an expression to a downstream consumer}\label{perfectly-forwarding-an-expression-to-a-downstream-consumer}

A frequent use of forwarding references and \texttt{std::forward} is to
propagate an object, whose \textbf{value category} is
invocation-dependent, down to one or more service providers that will
behave differently depending on the \textbf{value category} of the
original argument.

As an example, consider an overload set for a function, \texttt{sink},
that accepts a \texttt{std::string} either by \texttt{const}
\emph{lvalue} reference (e.g., with the intention of \textit{copying} from it) or
\emph{rvalue} reference (e.g., with the intention of \textit{moving} from it):

\begin{emcppshiddenlisting}[emcppsbatch=e8]
#include <string>   // (ù{\codeincomments{std::string}}ù)
#include <utility>  // (ù{\codeincomments{std::move}}ù)
std::string target;
\end{emcppshiddenlisting}
\begin{emcppslisting}[emcppsbatch=e8]
void sink(const std::string& s) { target = s; }
void sink(std::string&& s)      { target = std::move(s); }
\end{emcppslisting}

\noindent Now, let's assume that we want to create an intermediary function
template, \texttt{pipe}, that will accept an \texttt{std::string} of any
\textbf{value category} and will dispatch its argument to the
corresponding overload of \texttt{sink}. By accepting a
\emph{forwarding} reference as a function parameter and invoking
\texttt{std::forward} as part of \texttt{pipe}'s body, we can achieve
our original goal without any code duplication:

\begin{emcppslisting}[emcppsbatch=e8]
template <typename T>
void pipe(T&& x)
{
    sink(std::forward<T>(x));
}
\end{emcppslisting}

\noindent Invoking \texttt{pipe} with an \emph{lvalue} will result in \texttt{x}
being an \emph{lvalue} reference and thus\linebreak[4]%%%%%%%%
\mbox{\texttt{sink(const}~\texttt{std::string\&)}'s} being called. Otherwise,
\texttt{x} will be an \emph{rvalue} reference and
\texttt{sink(std::string\&\&)} will be called. This idea of enabling
\emph{move} operations without code duplication (as \texttt{pipe} does)
is commonly referred to as \textit{\titleref{use-cases-forwardingref}: \titleref{perfect-forwarding-for-generic-factory-functions}} on page~\pageref{perfect-forwarding-for-generic-factory-functions}.

\subsubsection[Handling multiple parameters concisely]{Handling multiple parameters concisely}\label{handling-multiple-parameters-concisely}

Suppose you have a \textbf{value-semantic type (VST)} that holds a
collection of attributes where some (not necessarily proper) subset of
them need to be changed together{\cprotect\footnote{This type of
value-semantic type can be classified more specifically as a
\emph{complex}-\emph{constrained} attribute class; see \textbf{lakos2a}, section~4.2.}}:

\begin{emcppslisting}
#include <type_traits>  // (ù{\codeincomments{std::enable\_if}}ù)
#include <utility>      // (ù{\codeincomments{std::forward}}ù)

struct Person { /* UDT that benefits from move semantics */ };

class StudyGroup
{
    Person d_a;
    Person d_b;
    Person d_c;
    Person d_d;
    // ...

public:
    bool isValid(const Person& a, const Person& b, const Person& c, const Person& d);
        // Return (ù{\codeincomments{true}}ù) if these specific people form a valid study group under
        // the guidelines of the study-group commission, and false otherwise.
    // ...

    template <typename PA, typename PB, typename PC, typename PD,
        typename = typename std::enable_if<
            std::is_same<typename std::decay<PA>::type, Person>::value &&
            std::is_same<typename std::decay<PB>::type, Person>::value &&
            std::is_same<typename std::decay<PC>::type, Person>::value &&
            std::is_same<typename std::decay<PD>::type, Person>::value>::type>
    int setPersonsIfValid(PA&& a, PB&& b, PC&& c, PD&& d)
    {
        enum { e_SUCCESS = 0, e_FAIL };

        if (!isValid(a, b, c, d))
        {
            return e_FAIL;  // bad choice; no change
        }

        // Move or copy each person into this object's (ù{\codeincomments{Person}}ù) data members:
        d_a = std::forward<PA>(a);
        d_b = std::forward<PB>(b);
        d_c = std::forward<PC>(c);
        d_d = std::forward<PD>(d);

        return e_SUCCESS;  // Study group was updated successfully.
    }
};
\end{emcppslisting}

\noindent The \texttt{setPersonsIfValid} function is producing the full crossproduct of instantiations for every variation of qualifiers that can be
on a \texttt{Person} object. Any combination of \emph{lvalue} and
\emph{rvalue} \texttt{Person}s can be passed, and a template will be
instantiated that will copy the \emph{lvalues} and move from the
\emph{rvalues}. To make sure the \texttt{Person} objects are
created externally, the function is restricted, using
\texttt{std::enable\_if}, to instantiate only for types that decay to
\texttt{Person} (i.e., types that are cv-qualified or ref-qualified
\texttt{Person}). Because each parameter is a forwarding reference, they
can all implicitly convert to \texttt{const}~\texttt{Person\&} to pass
to \texttt{isValid}, creating no additional temporaries. Finally,
\texttt{std::forward} is then used to do the actual moving or copying as
appropriate to data members.

\subsubsection[Perfect forwarding for generic factory functions]{Perfect forwarding for generic factory functions}\label{perfect-forwarding-for-generic-factory-functions}

Consider the prototypical standard-library generic factory function,
\texttt{std::make\_shared<T>}. On the surface, the requirements for this
function are fairly simple --- allocate a place for a \texttt{T} and
then construct it with the same arguments that were passed to
\texttt{make\_shared}. This, however, gets reasonably complex to
implement efficiently when \texttt{T} can have a wide variety of ways in
which it might be initialized.

For simplicity, we will show how a two-argument \texttt{my::make\_shared}
might be defined, knowing that a full implementation would employ
variadic template arguments for this purpose --- see ``\titleref{variable-templates}" on page~\pageref{variable-templates}. We will also implement a simpler
version of \texttt{make\_shared} that simply creates the element on the
heap with \texttt{new} and constructs a \texttt{std::shared\_ptr} to
manage the lifetime of that element. The declaration of this form of
make shared would be structured like this:

\begin{emcppshiddenlisting}[emcppsbatch=e9]
#include <memory>  // (ù{\codeincomments{std::shared\_ptr}}ù)
\end{emcppshiddenlisting}
\begin{emcppslisting}[emcppsbatch=e9]
namespace my {
template <typename ELEMENT_TYPE, typename ARG1, typename ARG2>
std::shared_ptr<ELEMENT_TYPE> make_shared(ARG1&& arg1, ARG2&& arg2);
}
\end{emcppslisting}

\noindent As you can see, we have two forwarding reference arguments ---
\texttt{arg1} and \texttt{arg2} --- with deduced types \texttt{ARG1} and
\texttt{ARG2}. Now, the body of our function needs to carefully
construct our \texttt{ELEMENT\_TYPE} object on the heap and then create
our output \texttt{shared\_ptr}:

\begin{emcppslisting}[emcppsbatch=e9]
template <typename ELEMENT_TYPE, typename ARG1, typename ARG2>
std::shared_ptr<ELEMENT_TYPE> my::make_shared(ARG1&& arg1, ARG2&& arg2)
{
    ELEMENT_TYPE *element_p = new ELEMENT_TYPE(std::forward<ARG1>(arg1),
                                               std::forward<ARG2>(arg2));
    try
    {
        return std::shared_ptr<ELEMENT_TYPE>(element_p);
    }
    catch (...)
    {
        delete element_p;
        throw;
    }
}
\end{emcppslisting}

\noindent Note that this simplified implementation needs to take care that the
constructor for the return value does not throw, cleaning up the
allocated element if that should happen; normally a \textbf{RAII}
proctor to manage this ownership would be a more robust solution to this
problem.

Importantly, the use of \texttt{std::forward} to construct the element
means that the arguments passed to \texttt{make\_shared} will be used to
find the appropriate matching two-argument constructor of
\texttt{ELEMENT\_TYPE}. When those arguments are \emph{rvalues}, the
constructor found will again search for one that takes an \emph{rvalue}
and the arguments will be moved from. Even more, because this function
wants to forward exactly the \texttt{const}-ness and reference type of
the input arguments, we would have to write 12 distinct overloads for
each argument if we were not using perfect forwarding -- the full cross
product of \texttt{const} (or not), \texttt{volatile} (or not), and
\texttt{\&}, \texttt{\&\&}, (or not). This would mean a full
implementation of just this two-argument variation would require 144
distinct overloads, all almost identical and most never actually
instantiated. The use of forwarding references reduces that to just 1
overload for each number of arguments.

\subsubsection[Wrapping initialization in a generic factory function]{Wrapping initialization in a generic factory function}\label{wrapping-initialization-in-a-generic-factory-function}

Occasionally we might want to initialize an object with an intervening
function call wrapping the actual construction of that object. Suppose
we have a tracking system that we want to use to monitor how many times
certain initializers have been invoked:

\begin{emcppshiddenlisting}[emcppsbatch=e10]
#include <utility>      // (ù{\codeincomments{std::forward}}ù)
struct BigObject {
    BigObject(const char *, const char *);
};
\end{emcppshiddenlisting}
\begin{emcppslisting}[emcppsbatch=e10]
struct TrackingSystem
{
    template <typename T>
    static void trackInitialization(int numArgs);
        // Track the creation of a (ù{\codeincomments{T}}ù) with a constructor taking (ù{\codeincomments{numArgs}}ù)
        // arguments.
};
\end{emcppslisting}

\noindent Now we want to write a general utility function that can be used to
construct an arbitrary object and notify the tracking system of the
construction for us. Here we will use a variadic pack (see
``\titleref{variable-templates}" on page~\pageref{variable-templates}) of forwarding references to handle
calling the constructor for us:

\begin{emcppslisting}[emcppsbatch=e10]
template <class ELEMENT_TYPE, typename... ARGS>
ELEMENT_TYPE trackConstruction(ARGS&&... args)
{
    TrackingSystem::trackInitialization<ELEMENT_TYPE>(sizeof...(args));
    return ELEMENT_TYPE(std::forward<ARGS>(args)...);
}
\end{emcppslisting}

\noindent This lets us add tracking easily to convert any initialization to a
tracked one by inserting a call to this function around the constructor
arguments:

\begin{emcppslisting}[emcppsbatch=e10]
void myFunction()
{
    BigObject untracked("Hello", "World");
    BigObject tracked = trackConstruction<BigObject>("Hello","World");
}
\end{emcppslisting}

\noindent On the surface there does seem to be a difference between how
\texttt{untracked} and \texttt{tracked} are initialized. The first
variable is having its constructor directly invoked, while the second is
being constructed from an object being returned by-value from
\texttt{trackConstruction}. This construction, however, has long been
something that has been optimized away to avoid any additional objects
and construct the object in question just once. In this case, because
the element being returned is initialized by the \texttt{return}
statement of \texttt{trackConstruction}, the optimization is called
\textbf{return value optimization (RVO)}. C++ has always
allowed this optimization by enabling \textbf{copy elision}. In C++17,
this elision can even be guaranteed and is allowed to be done for
objects that have no copy or move constructors. Prior to C++17, this
elision can still be guaranteed (on all compilers that the authors are
aware of) by \textbf{declaring} but not \textbf{defining} the copy
constructor for \texttt{BigObject}. You'll find that this code will
still compile and link with such an object, providing observable proof
that the copy constructor is never actually invoked with this pattern.

\subsubsection[Emplacement]{Emplacement}\label{emplacement}

Prior to C++11, inserting an object into a standard library container
always required the programmer to first create such an object and then
copy it inside the container's storage. As an example, consider
inserting a temporary \texttt{std::string} object in a\linebreak[4]
\mbox{\texttt{std::vector<std::string>}}:

\begin{emcppshiddenlisting}[emcppsbatch=e11]
#include <vector>  // (ù{\codeincomments{std::vector}}ù)
#include <string>  // (ù{\codeincomments{std::string}}ù)
\end{emcppshiddenlisting}
\begin{emcppslisting}[emcppsbatch=e11]
void f(std::vector<std::string>& v)
{
    v.push_back(std::string("hello world"));
        // invokes (ù{\codeincomments{std::string::string(const char*)}}ù) and the copy-constructor
}
\end{emcppslisting}

\noindent In the function above, a temporary \texttt{std::string} object is
created in the stack frame of \texttt{f} and is then copied to the
dynamically allocated buffer managed by \texttt{v}. Additionally, the
buffer might have insufficient capacity and hence might require
reallocation, which would in turn require every element of \texttt{v} to
be somehow copied from the old buffer to the new, larger one.

In C++11, the situation is significantly better thanks to
rvalue references. The temporary will be moved into
\texttt{v}, and any buffer reallocation will \emph{move} the elements
between buffers rather than copy them, assuming that the element's
move-constructor is a noexcept specifier (see ``\titleref{noexcept-specifier}" on page~\pageref{noexcept-specifier}). The amount of work
can, however, be further minimized: What if, instead of first creating
an object externally, we constructed the new \texttt{std::string} object
directly in \texttt{v}'s buffer?

This is where \textbf{emplacement} comes into play. All standard library
containers, including \texttt{std::vector}, now provide an
\textbf{emplacement} API powered by variadic templates (see ``\titleref{variadictemplate}" on page~\pageref{variadictemplate})
and perfect forwarding (see \textit{\titleref{use-cases-forwardingref}: \titleref{perfect-forwarding-for-generic-factory-functions}} on page~\pageref{perfect-forwarding-for-generic-factory-functions}). Rather than accepting a
fully-constructed element, \textbf{emplacement} operations accept an
arbitrary number of arguments, which will in turn be used to construct a
new element directly in the container's storage, thereby avoiding
unnecessary copies or even moves:

\begin{emcppslisting}[emcppsbatch=e11]
void g(std::vector<std::string>& v)
{
    v.emplace_back("hello world");
        // invokes only the (ù{\codeincomments{std::string::string(const char*)}}ù) constructor
}
\end{emcppslisting}

\noindent Calling \texttt{std::vector<std::string>::emplace\_back} with a
\texttt{const}~\texttt{char*} argument results in a new
\texttt{std::string} object being created in-place in the next suitable
spot of the vector's storage. Internally,
\texttt{std::allocator\_traits::construct} is invoked, which typically
employs \textbf{placement \texttt{new}} to construct the object in raw
dynamically allocated memory. As previously mentioned,
\texttt{emplace\_back} makes use of both variadic
templates and forwarding references; it accepts any
number of forwarding references and internally \textit{perfectly forwards} them
to the constructor of \texttt{T} via \texttt{std::forward}:

\begin{emcppshiddenlisting}[emcppsbatch=e12]
#include <new>     // placement (ù{\codeincomments{new}}ù)
namespace std {
template <typename T>
class vector {
    template <typename... Args>
    void emplace_back(Args&&... args);
};
char freeLocationInBuffer[10];
}
\end{emcppshiddenlisting}
\begin{emcppslisting}[emcppsbatch=e12]
template <typename T>
template <typename... Args>
void std::vector<T>::emplace_back(Args&&... args)
{
    // ...
    new (&freeLocationInBuffer) T(std::forward<Args>(args)...);  // pseudocode
    // ...
}
\end{emcppslisting}

\noindent \textbf{Emplacement} operations remove the need for copy or move
operations when inserting elements into containers, potentially
increasing the performance of a program and sometimes --- depending on the
container --- even allowing even noncopyable or nonmovable objects to be stored
in a container.

As previously mentioned,
declaring (but not defining) the \emph{copy} or \emph{move} ctor of a
noncopyable or nonmovable type to be private is often a way to guarantee that
a C++11/14 compiler constructed an object in place.
Containers that might need to move elements around for other operations
(such as \texttt{std::vector} or \texttt{std::deque}) will still need
movable elements, while node-based containers that never move the
elements themselves after initial construction (such as
\texttt{std::list} or \texttt{std::map}) can use emplace along with
noncopyable or nonmovable objects.

\subsubsection[Decomposing complex expressions]{Decomposing complex expressions}\label{decomposing-complex-expressions}

Many modern C++ libraries have adopted a more ``functional'' style of
programming, chaining the output of one function as the arguments of
another function to produce very complex expressions that accomplish a
great deal in relatively concise fashion. Consider the way in which the
C++20 ranges library encapsulates containers and arbitrary pairs of
iterators into objects that can be adapted and manipulated through long
chains of functions. Let's say you have a function that reads a file,
does some spellchecking for every unique word in the file, and gives you a
list of incorrect words and corresponding suggested proper spellings,
and you have a range-like library with common utilities similar to
standard UNIX processing utilities:

\begin{emcppshiddenlisting}[emcppsbatch={e13,e14},emcppsstandards=c++14]
#include <map>     // (ù{\codeincomments{std::map}}ù)
#include <string>  // (ù{\codeincomments{std::string}}ù)
#include <tuple>   // (ù{\codeincomments{std::tuple}}ù)

class SpellingSuggestion {
    bool isCorrect() const;
};
class LineSeq {};
LineSeq openFile(const std::string&);
LineSeq splitRegex(LineSeq seq, const std::string &);
LineSeq filterRegex(LineSeq seq, const std::string &);
LineSeq sort(LineSeq seq);
LineSeq uniq(LineSeq seq);
template <typename F>
LineSeq transform(LineSeq seq, F&& func) { return {}; }
template <typename F>
LineSeq filter(LineSeq seq, F&& pred) { return {}; }
std::map<std::string,SpellingSuggestion> makeMap(LineSeq seq);
SpellingSuggestion checkSpelling(const std::string& word);
\end{emcppshiddenlisting}
\begin{emcppslisting}[emcppsbatch=e13]
SpellingSuggestion checkSpelling(const std::string& word);

std::map<std::string, SpellingSuggestion> checkFileSpelling(
                                                   const std::string& filename)
{
    return makeMap(
        filter(transform(
           uniq(sort(filterRegex(splitRegex(openFile(filename),"\\S+"),"\\w+"))),
        [](const std::string& x)
        {
            return std::tuple<std::string, SpellingSuggestion>(x,
                                                             checkSpelling(x));
        }
   ), [](auto&& x) { return !std::get<1>(x).isCorrect(); }));
}
\end{emcppslisting}

\noindent Upon doing code review for this amazing use of a modern library produced
by the smart, new programmer on your team, you discover that you actually
have a very hard time understanding what is going on. On top of that,
the usual tools you have to poke and prod at the code by adding
\texttt{printf} statements or even breakpoints in your debugger are very
hard to apply to the complex set of nested templates involved.

Each of the functions in this range library -- \texttt{makeMap},
\texttt{transform}, \texttt{uniq}, \texttt{sort}, \texttt{filterRegex},
\texttt{splitRegex}, and \texttt{openFile} -- is a set of complex
templated overloads and deeply subtle metaprogramming that becomes hard
to unravel for a nonexpert C++ programmer. On the other hand, you have
also looked at the code generated for this function and the abstractions
amazingly get compiled away to a very robust implementation.

To better understand, document, and debug what is happening
here, you want to decompose this expression into many, capturing the
implicit temporaries returned by all of these functions and ideally not
changing the actual semantics of what is being done. To do that
properly, you need to capture the type and value category of each
subexpression appropriately, without necessarily being able to easily
decode it manually from the expression. Here is where \texttt{auto\&\&}
forwarding references can be used effectively to decompose and document
this expression while achieving the same:

\begin{emcppslisting}[emcppsbatch=e14]
std::map<std::string, SpellingSuggestion> checkFileSpelling(
                                                   const std::string& filename)
{
    // Create a range over the contents of (ù{\codeincomments{filename}}ù).
    auto&& openedFile = openFile(filename);

    // Split the file by whitespace.
    auto&& potentialWords = splitRegex(
        std::forward<decltype(openedFile)>(openedFile), "\\S+");

    // Filter out only words made from word-characters.
    auto&& words = filterRegex(
        std::forward<decltype(potentialWords)>(potentialWords), "\\w+");

    // Sort all words.
    auto&& sortedWords = sort(std::forward<decltype(words)>(words));

    // Skip adjacent identical words. (This is now a sequence of unique words.)
    auto&& uniqueWords = uniq(std::forward<decltype(sortedWords)>(sortedWords));

    // Get a (ù{\codeincomments{SpellingSuggestion}}ù) for every word.
    auto&& suggestions = transform(
        std::forward<decltype(uniqueWords)>(uniqueWords),
        [](const std::string&x) {
            return std::tuple<std::string,SpellingSuggestion>(
                x,checkSpelling(x));
        });

    // Filter out correctly spelled words, keeping only elements where the
    // second element of the tuple, which is a (ù{\codeincomments{SpellingSuggestion}}ù), is not
    // (ù{\codeincomments{correct}}ù).
    auto&& corrections = filter(
        std::forward<decltype(suggestions)>(suggestions),
        [](auto&& suggestion){ return !std::get<1>(suggestion).isCorrect(); });

    // Return a map made from these 2-element tuples:
    return makeMap(std::forward<decltype(corrections)>(corrections));
}
\end{emcppslisting}

\noindent Now each step of this complex expression is documented, each temporary
has a name, but the net result of the lifetimes of each object is
functionally the same. No new conversions have been introduced, and
every object that was used as an \emph{rvalue} in the original
expression will still be used as an \emph{rvalue} in this much longer
and more descriptive implementation of the same functionality.

\subsection[Potential Pitfalls]{Potential Pitfalls}\label{potential-pitfalls}

\subsubsection[Surprising number of template instantiations with string literals]{Surprising number of template instantiations with string literals}\label{surprising-number-of-template-instantiations-with-string-literals}

When forwarding references are used as a means to avoid code repetition
between exactly two overloads of the same function (one accepting a
\texttt{const}~\texttt{T\&} and the other a \texttt{T\&\&}), it can be
surprising to see more than two template instantiations for that
particular template function, in particular when the function is invoked
using string literals.

Consider, as an example, a \texttt{Dictionary} class containing two
overloads of an \texttt{addWord} member function:
\begin{emcppshiddenlisting}[emcppsbatch={e15,e16}]
#include <string>  // (ù{\codeincomments{std::string}}ù)
\end{emcppshiddenlisting}
\begin{emcppslisting}[emcppsbatch=e15]
class Dictionary
{
    // ...

public:
    void addWord(const std::string& word);  // (0) copy (ù{\codeincomments{word}}ù) in the dictionary
    void addWord(std::string&& word);       // (1) move (ù{\codeincomments{word}}ù) in the dictionary
};

void f()
{
    Dictionary d;

    std::string s = "car";
    d.addWord(s);                    // invokes (0)

    const std::string cs = "toy";
    d.addWord(cs);                   // invokes (0)

    d.addWord("house");              // invokes (1)
    d.addWord("garage");             // invokes (1)
    d.addWord(std::string{"ball"});  // invokes (1)
}
\end{emcppslisting}

\noindent Now, imagine replacing the two overloads of \texttt{addWord} with a
single \textit{perfectly forwarding} template member function, with the intention
of avoiding code repetition between the two overloads:

\begin{emcppslisting}[emcppsbatch=e16]
class Dictionary
{
    // ...

public:
    template <typename T>
    void addWord(T&& word);
};
\end{emcppslisting}

\noindent Perhaps surprisingly, the number of template instantiations skyrockets:

\begin{emcppslisting}[emcppsbatch=e16]
void f()
{
    Dictionary d;

    std::string s = "car";
    d.addWord(s);  // instantiates (ù{\codeincomments{addWord<std::string\&>}}ù)

    const std::string cs = "toy";
    d.addWord(cs);  // instantiates (ù{\codeincomments{addWord<const std::string\&>}}ù)

    d.addWord("house");              // instantiates (ù{\codeincomments{addWord<char const(\&)[6]>}}ù)
    d.addWord("garage");             // instantiates (ù{\codeincomments{addWord<char const(\&)[7]>}}ù)
    d.addWord(std::string{"ball"});  // instantiates (ù{\codeincomments{addWord<std::string\&\&>}}ù)
}
\end{emcppslisting}

\noindent Depending on the variety of argument types supplied to \texttt{addWord},
having many call sites could result in an undesirably large number of
distinct template instantiations, perhaps significantly increasing
object code size, compilation time, or both.

\subsubsection[{\tt std::forward<T>} can enable move operations]{{\SubsubsecCode std::forward<T>} can enable move operations}\label{std::forward<t>-can-enable-move-operations}

Invoking \texttt{std::forward<T>(x)} is equivalent to conditionally
invoking \texttt{std::move} (if \texttt{T} is an
\emph{lvalue} reference). Hence, any subsequent use of \texttt{x} is
subject to the same caveats that would apply to an \emph{lvalue} cast to
an unnamed \emph{rvalue} reference; see ``\titleref{Rvalue-References}" on page~\pageref{Rvalue-References}:

\begin{emcppshiddenlisting}[emcppsbatch=e17]
#include <utility>      // (ù{\codeincomments{std::forward}}ù)
template <typename T>
void g(T&& x);
\end{emcppshiddenlisting}
\begin{emcppslisting}[emcppsbatch=e17]
template <typename T>
void f(T&& x)
{
    g(std::forward<T>(x));  // OK
    g(x);                   // Oops! (ù{\codeincomments{x}}ù) could have already been moved from.
}
\end{emcppslisting}

\noindent Once an object has been passed as an argument using
\texttt{std::forward}, it should typically not be accessed again
without first assigning it a new value because it could now be in a
moved-from state.

\subsubsection[A perfect-forwarding constructor can hijack the copy constructor]{A perfect-forwarding constructor can hijack the copy constructor}\label{a-perfect-forwarding-constructor-can-hijack-the-copy-constructor}

A single-parameter constructor of a class \texttt{S} accepting a
forwarding reference can unexpectedly be a better match during overload
resolution compared to \texttt{S}'s copy constructor:

\begin{emcppslisting}
struct S
{
    S();                           // default constructor
    template <typename T> S(T&&);  // forwarding constructor
    S(const S&);                   // copy constructor
};

void f()
{
    S a;
    const S b;

    S x(a);  // invokes forwarding constructor
    S y(b);  // invokes copy constructor
}
\end{emcppslisting}

\noindent Despite the programmer's intention to copy from \texttt{a} into
\texttt{x}, the forwarding constructor of \texttt{S} was invoked
instead, because \texttt{a} is a non-\texttt{const} \emph{lvalue}
expression, and instantiating the forwarding constructor with
\texttt{T}~\texttt{=}~\texttt{S\&} results in a better match than even
the copy constructor.

This potential pitfall can arise in practice, for example, when writing
a value-semantic wrapper template (e.g., \texttt{Wrapper}) that can be
initialized by \textit{perfectly forwarding} the object to be wrapped into it:

\begin{emcppshiddenlisting}[emcppsbatch=e18]
#include <string>   // (ù{\codeincomments{std::string}}ù)
#include <utility>  // (ù{\codeincomments{std::forward}}ù)
\end{emcppshiddenlisting}
\begin{emcppslisting}[emcppsbatch=e18]
template <typename T>
class Wrapper  // wrapper for an object of arbitrary type 'T'
{
private:
    T d_datum;

public:
    template <typename U>
    Wrapper(U&& datum) : d_datum(std::forward<U>(datum)) { }
        // perfect-forwarding constructor (to optimize runtime performance)

    // ...
};

void f()
{
    std::string s("hello world");
    Wrapper<std::string> w0(s);  // OK, (ù{\codeincomments{s}}ù) is copied into (ù{\codeincomments{d\_datum}}ù).

    Wrapper<std::string> w1(std::string("hello world"));
        // OK, the temporary string is moved into (ù{\codeincomments{d\_datum}}ù).
}
\end{emcppslisting}

\noindent Similarly to the example involving class \texttt{S} in the example above, attempting
to copy-construct a non-\texttt{const} instance of \texttt{Wrapper}
(e.g., \texttt{wr}, above) results in an error:

\begin{emcppslisting}[emcppsbatch=e18]
void g(Wrapper<int>& wr)  // The same would happen if (ù{\codeincomments{wr}}ù) were passed by value.
{
    Wrapper<int> w2(10);  // OK, invokes perfect-forwarding constructor
    Wrapper<int> w3(wr);  // Error, no conversion from (ù{\codeincomments{Wrapper<int>}}ù) to (ù{\codeincomments{int}}ù)
}
\end{emcppslisting}

\noindent The compilation failure above occurs because the perfect-forwarding
constructor template, instantiated with \texttt{Wrapper<int>\&}, is a
better match than the implicitly generated copy constructor, which
accepts a \texttt{const}~\texttt{Wrapper<int>\&}. Constraining the
perfect forwarding constructor via \textbf{SFINAE} (e.g., with
\texttt{std::enable\_if}) to explicitly \emph{not} accept objects whose
type is \texttt{Wrapper} fixes this problem:

\begin{emcppshiddenlisting}[emcppsbatch=e19]
#include <type_traits>  // (ù{\codeincomments{std::enable\_if}}ù), (ù{\codeincomments{std::decay}}ù)
#include <utility>  // (ù{\codeincomments{std::forward}}ù)
\end{emcppshiddenlisting}
\begin{emcppslisting}[emcppsbatch=e19]
template <typename T>
class Wrapper
{
private:
    T d_datum;

public:
    template <typename U,
        typename = typename std::enable_if<
            !std::is_same<typename std::decay<U>::type, Wrapper>::value
        >::type
    >
    Wrapper(U&& datum) : d_datum(std::forward<U>(datum)) { }
        // This constructor participates in overload resolution only if (ù{\codeincomments{U}}ù),
        // after being decayed, is not the same as (ù{\codeincomments{Wrapper}}ù).
};

void h(Wrapper<int>& wr)  // The same would happen if (ù{\codeincomments{wr}}ù) were passed by value.
{
    Wrapper<int> w4(10);  // OK, invokes the perfect-forwarding constructor
    Wrapper<int> w5(wr);  // OK, invokes the copy constructor
}
\end{emcppslisting}

\noindent Notice that the \texttt{std::decay} \textbf{metafunction} was used as
part of the constraint; for more information on the using
\texttt{std::decay}, see \textit{\titleref{annoyances-forwardingref}: \titleref{metafunctions-are-required-in-constraints}} on page~\pageref{metafunctions-are-required-in-constraints}.

\subsection[Annoyances]{Annoyances}\label{annoyances-forwardingref}

\subsubsection[Forwarding references look just like rvalue references]{Forwarding references look just like rvalue references}\label{forwarding-references-look-just-like-rvalue-references}

Despite \emph{forwarding} references and rvalue
references having significantly different semantics, as discussed in \textit{\titleref{description-forwardingref}: \titleref{identifying-forwarding-references}} on page~\pageref{identifying-forwarding-references}, they share the same syntax. For any given type \texttt{T},
whether the \texttt{T\&\&} syntax designates an rvalue
reference or a \emph{forwarding} reference depends entirely on the
surrounding context.{\cprotect\footnote{In C++20, developers might be
subject to additional confusion due to the new terse concept notation
syntax, which allows function templates to be defined without any
explicit appearance of the \texttt{template} keyword. As an example, a
constrained function parameter, like
\texttt{Addable}~\texttt{auto\&\&}~\texttt{a} in the example below, is a forwarding
reference; looking for the presence of the mandatory \texttt{auto}
keyword is helpful in identifying whether a type is a forwarding
reference or \emph{rvalue} reference:

\begin{emcppslisting}[emcppsstandards=c++20,basicstyle={\ttfamily\footnotesize}]
template<typename T>
concept Addable = requires(T a, T b) { a + b; };

void f(Addable auto&& a);  // C++20 terse concept notation

void example()
{
    int i;

    f(i);  // OK, (ù{\codeincomments{decltype(a)}}ù) is (ù{\codeincomments{int\&}}ù) in (ù{\codeincomments{f}}ù).
    f(0);  // OK, (ù{\codeincomments{decltype(a)}}ù) is (ù{\codeincomments{int\&\&}}ù) in (ù{\codeincomments{f}}ù).
}
\end{emcppslisting}
      }}

\begin{emcppslisting}
template <typename T> struct S0 { void f(T&&); };  // rvalue reference
struct S1 { template <typename T> void f(T&&); };  // forwarding reference
\end{emcppslisting}

\noindent Furthermore, even if \texttt{T} is subject to template argument
deduction, the presence of \emph{any} qualifier will suppress the
special \emph{forwarding}-reference deduction rules:

\begin{emcppslisting}
template <typename T> void f(T&&);           // forwarding reference
template <typename T> void g(const T&&);     // (ù{\codeincomments{const}}ù) rvalue reference
template <typename T> void h(volatile T&&);  // (ù{\codeincomments{volatile}}ù) rvalue reference
\end{emcppslisting}

\noindent It is truly remarkable that we still do not have some unique syntax
(e.g., \texttt{\&\&\&}) that we could use, at least optionally, to imply
unequivocally a \emph{forwarding} reference that is independent of its
context.

\subsubsection[Metafunctions are required in constraints]{Metafunctions are required in constraints}\label{metafunctions-are-required-in-constraints}

As we showed in \textit{\titleref{use-cases-forwardingref}} on page~\pageref{use-cases-forwardingref}, being able to
perfectly forward arguments of the same general type and effectively
leave only the value category of the argument up to type deduction is a frequent need. This
is necessary if you do not want to delay construction of the arguments
until they are forwarded, possibly because doing so would produce many
unnecessary temporaries.

The challenge to make this work correctly is significant. The template
must be constrained using \textbf{SFINAE} and the appropriate
\textbf{type traits} to disallow types that aren't some form of
cv-qualified or ref-qualified version of the type that you want to
accept. As an example, let's consider a function intended to \emph{copy}
or \emph{move} a \texttt{Person} object into a data structure:

\begin{emcppshiddenlisting}[emcppsbatch=e20]
#include <type_traits>  // (ù{\codeincomments{std::enable\_if}}ù), (ù{\codeincomments{std::decay}}ù), (ù{\codeincomments{std::is\_same}}ù)
class Person {};
\end{emcppshiddenlisting}
\begin{emcppslisting}[emcppsbatch=e20]
class PersonManager {
// ...
template <typename T, typename = typename std::enable_if<
       std::is_same<typename std::decay<T>::type, Person>::value>::type>
void addPerson(T&& person) {}
    // This function participates in overload resolution only if (ù{\codeincomments{T}}ù) is
    // (possibly cv- or ref-qualified) (ù{\codeincomments{Person}}ù).
// ...
};
\end{emcppslisting}

\noindent This incantation to constrain \texttt{T} has a number of layers to it,
so let's unpack them one at a time.
\begin{itemize}
\item{\texttt{T} is the template argument we are trying to deduce. We’d like to limit it to being a \texttt{Person} that is \texttt{const}, \texttt{volatile}, \texttt{\&}, \texttt{\&\&}, or some (possibly empty) valid combination of those.}
\item{\texttt{std::decay<T>::type} is then the application of the standard metafunction (defined in \texttt{<type\_traits>}) \texttt{std::decay} to \texttt{T}. This metafunction removes all cv-qualifiers and ref-qualifiers from \texttt{T}, and so, for the types to which we want to limit \texttt{T}, this will \emph{always} be \texttt{Person}. Note that \texttt{decay} will also allow some other implicitly convertible transformations, such as converting an array type to the corresponding pointer type. For types we are concerned with --- those that decay to a \texttt{Person} --- this metafunction is equivalent to \texttt{std::remove\_cv<std::remove\_reference<T>::type>::type}, or the equivalent and shorter \texttt{std::remove\_cvref<T>::type>} available in C++20. Due to historical availability and readability, we will continue with our use of \texttt{decay} for this purpose.}
\item{\texttt{std::is\_same<std::decay<T>::type,}~\texttt{Person>::value} is then the application of another metafunction, \texttt{std::is\_same}, to two arguments --- our decay expression and \texttt{Person}, which results in a \texttt{value} that is either \texttt{std::true\_type} or \texttt{std::false\_type} --- special types that can convert, in compile time, expressions to \texttt{true} or \texttt{false}. For the types \texttt{T} that we care about, this expression will be \texttt{true}, and for all other types this expression will be \texttt{false}.}
\item{\texttt{std::enable\_if<X>::type} is yet another metafunction that evaluates to a valid type if and only if \texttt{X} is true. Unlike the \texttt{value} in \texttt{std::is\_same}, this expression is simply not valid if \texttt{X} is false.}
\item{Finally, by using this \texttt{enable\_if} expression as a default-initialized template argument, the expression is going to be instantiated for any deduced \texttt{T} considered during overload resolution for \texttt{addPerson}. This instantiation will fail for any of the types we don’t want to allow (something that is not a cv). Because of this, for any \texttt{T} that isn’t one of the types for which we want to allow \texttt{addPerson} to be invoked, this substitution will fail. Rather than being an error, this just removes \texttt{addPerson} from the overload set being considered, hence the term \textbf{SFINAE}. In this case, that would give us a different error indicating that we attempted to pass a non-\texttt{Person} to \texttt{addPerson}, which is exactly the result we want.}
\end{itemize}
Putting this all together means we get to call \texttt{addPerson} with
\emph{lvalues} and \emph{rvalues} of type \texttt{Person}, and the value
category will be appropriately usable within \texttt{addPerson}
(generally with use of \texttt{std::forward} within that function's
definition).

\subsection[See Also]{See Also}\label{see-also}

\begin{itemize}
\item{\seealsoref{Rvalue-References}{\locationc}Feature that can be confused with forwarding references due to similar syntax.}
\item{`\seealsoref{auto-feature}{\locationc}Feature that can introduce a forwarding reference with the \texttt{auto\&\&} syntax.}
\item{\seealsoref{variadictemplate}{\locationc}Feature commonly used in conjunction with forwarding references to provide highly generic interfaces.}
\end{itemize}

\subsection[Further Reading]{Further Reading}\label{further-reading}

\begin{itemize}
\item{``Item 24: Distinguish universal references from rvalue references," \cite{meyers15}}
\item{\cite{sutter14}}
\item{\cite{niebler13}}
\end{itemize}




\newpage
%\section[Generalized PODs]{Generalized Plain Old Data Types}%\label{}


\emcppsFeature{
    short={Generalized PODs},
    long={Generalized Plain Old Data Types},
}{gpods}
\setcounter{table}{0}
\setcounter{footnote}{0}
\setcounter{lstlisting}{0}
%\section[Generalized PODs]{Generalized Plain Old Data Types}%\label{}

placeholder 

\newpage
%\section[Inheriting Ctors]{Inheriting Base Class Constructors}\label{ctorinheriting}\label{inheriting-constructors}
% 12 Feb 2021, revisions in; ready for Josh's code fixes
% 16 Feb 2021 JMB; code compiles

\emcppsFeature{
    short={Inheriting Ctors},
    long={Inheriting Base-Class Constructors},
}{ctorinheriting}
\label{inheriting-constructors}
\setcounter{table}{0}
\setcounter{footnote}{0}
\setcounter{lstlisting}{0}
%\section[Inheriting Ctors]{Inheriting Base Class Constructors}\label{ctorinheriting}\label{inheriting-constructors}



The term \emph{inheriting constructors} refers to the use of a
\emcppsgloss[using declaration]{\lstinline!using! declaration} to expose nearly all of the
constructors of a base class in the scope of a derived class.

\subsection[Description]{Description}\label{description-inheritingctor}

In a class definition, a \emcppsgloss[using declaration]{\lstinline!using! declaration} naming a
base class's constructor results in the derived class ``inheriting'' all
of the nominated base class's constructors, except for \emph{copy} and
\emph{move} constructors. Just like \lstinline!using! declarations of member functions, the nominated base class's
constructors will be considered when no matching constructor is found in
the derived class. When a base class constructor is selected in this way,
that constructor will be used to construct the base class, and the
remaining bases and data members of the subclass will be initialized as
if by the default constructor (e.g., applying default initializers;
see \featureref{\locationc}{Default-Member-Initializers}).
%``\titleref{Default-Member-Initializers}" on page~\pageref{Default-Member-Initializers}). 

\begin{emcppslisting}[emcppsbatch=e1]
struct B0
{
    B0() = default;      // public, default constructor
    B0(int)         { }  // public, one argument (implicit) value constructor
    B0(int, int)    { }  // public, two argument value constructor

private:
    B0(const char*) { }  // private, one argument (implicit) value constructor
};

struct D0 : B0
{
    using B0::B0;  // (ù{\codeincomments{using}}ù) declaration
    D0(double d);  // suppress implicit default constructor
};

D0 t(1);     // OK, inherited from (ù{\codeincomments{B0::B0(int)}}ù)
D0 u(2, 3);  // OK, inherited from (ù{\codeincomments{B0::B0(int, int)}}ù)
D0 v("hi");  // Error, Base constructor is declared (ù{\codeincomments{private}}ù).
\end{emcppslisting}
    
\noindent The only constructors that are explicitly \emph{not} inheritable by the
derived class are the potentially compiler-generated \emph{copy} and
\emph{move} constructors:

\begin{emcppslisting}[emcppsbatch=e1]
#include <utility>  // (ù{\codeincomments{std::move}}ù)

B0 b1(1);              // OK, base-class object can be created.
B0 b2(2, 3);           // OK, base-class object can be created.
B0 b3(b1);             // OK, base-class object can be copied (from *lvalue*).
B0 b4(std::move(b1));  // OK, base-class object can be moved (from *rvalue*).

D0 w(b1);    // Error, base-class copy constructor is not inherited.
D0 v;        // OK, base-class default constructor is inherited.
D0 x(B0{});  // Error, base-class move constructor is not inherited.

D0 y(B0(4));  // Error, base-class move constructor is not inherited.
D0 z(t);      // OK, uses compiler-generated (ù{\codeincomments{D0::D0(const D0\&)}}ù)
D0 j(D0(5));  // OK, uses compiler-generated (ù{\codeincomments{D0::D0(D\&\&)}}ù)
\end{emcppslisting}

\noindent Note that we use
 braced initialization (see \featureref{\locationc}{bracedinit}) 
% ``\titleref{bracedinit}" on page~\pageref{bracedinit})
  in \lstinline!D0!~\lstinline!x(B0{});!
  to ensure that a variable \lstinline!x! of type \lstinline!D0! is declared.
  \lstinline!D0!~\lstinline!x(B0());! would instead be interpreted as a
  declaration of a function \lstinline!x! returning \lstinline!D0! and
  accepting a pointer to a nullary function returning \lstinline!B0!, which
  is referred to as the \emcppsgloss{most vexing parse}.
    
The constructors inherited by the derived class have the same effect on
whether the compiler implicitly generates special member functions as
explicitly implemented ones would. For example, \lstinline!D0!'s default
constructor would be implicitly \emph{deleted} (see \featureref{\locationa}{deleted-functions}) 
%``\titleref{deleted-functions}" on page~\pageref{deleted-functions}) 
if
\lstinline!B0! doesn't have a default constructor. Note that since the copy
and move constructors are \emph{not} inherited, their presence in the
base class wouldn't suppress implicit generation of copy and move
assignment in the derived class. For instance, \lstinline!D0!'s implicitly
generated assignment operators \emcppsgloss{hide} their
counterparts in \lstinline!B0!:

\begin{emcppslisting}[emcppsbatch=e1]
void f()
{
    B0 b(0), bb(0);  // Create destination and source (ù{\codeincomments{B0}}ù) objects.
    D0 d(0), dd(0);  //   "         "       "    "    (ù{\codeincomments{D0}}ù)    ".

    b = bb;          // OK, assign base from lvalue base.
    b = B0(0);       // OK,   "       "   "  rvalue   "

    d = bb;          // Error, (ù{\codeincomments{B0::operator=}}ù) is hidden by (ù{\codeincomments{D0::operator=}}ù).
    d = B0(0);       // Error,       "         "    "     "        "

    d.B0::operator=(bb);     // OK, explicit slicing is still possible.
    d.B0::operator=(B0(0));  // OK,    "         "      "   "      "

    d = dd;          // OK, assign derived from lvalue derived.
    d = D0(0);       // OK,   "        "     "  rvalue    "
}
\end{emcppslisting}
    
\noindent Note that, when inheriting constructors, private constructors in the
base class are accessed as private constructors of that base class and
are subject to the same access controls; see \intraref{annoyances-inheritingctor}{access-levels-of-inherited-constructors-are-same-as-in-base-class}. 
% \textit{\titleref{annoyances-inheritingctor}: \titleref{access-levels-of-inherited-constructors-are-same-as-in-base-class}} on page~\pageref{access-levels-of-inherited-constructors-are-same-as-in-base-class}.

Inheriting constructors having the same \emcppsgloss{signature} from multiple
base classes leads to ambiguity errors:

\begin{emcppslisting}
struct B1A { B1A(int); }; // Here we have two bases classes, each of which
struct B1B { B1B(int); }; // provides a conversion constructor from an int.

struct D1 : B1A, B1B
{
    using B1A::B1A;
    using B1B::B1B;
};

D1 d1(0);  // Error, Call of overloaded (ù{\codeincomments{D1(int)}}ù) is ambiguous.
\end{emcppslisting}
    
\noindent Each inherited constructor shares the same characteristics as the
corresponding one in the nominated base class's constructor and then
delegates to it. This means the \emcppsgloss[access specifier]{access specifiers}, the
\lstinline!explicit! specifier, the \lstinline!constexpr! specifier, the
default arguments, and the exception specification are also preserved by
constructor inheritance; see \featureref{\locatione}{noexcept-specifier} and \featureref{\locationc}{constexprfunc}.  
%``\titleref{noexcept}" on page~\pageref{noexcept} and ``\titleref{constexprfunc}" on page~\pageref{constexprfunc}. 
For template
constructors, the template parameter list and the default template
arguments are preserved as well:

\begin{emcppslisting}
struct B2
{
    template <typename T = int>
    explicit B2(T) { }
};

struct D2 : B2 { using B2::B2; };
\end{emcppslisting}
    
\noindent The declaration \lstinline!using!~\lstinline!B2::B2! above behaves as if a
constructor template that delegates to its nominated base class's
template was provided in \lstinline!D2!:

\begin{emcppshiddenlisting}[emcppsbatch=e2]
struct B2
{
    template <typename T = int>
    explicit B2(T) { }
};
\end{emcppshiddenlisting}
\begin{emcppslisting}[emcppsbatch=e2]
// pseudocode
struct D2 : B2
{
    template <typename T = int>
    explicit D2(T i) : B2(i) { }
};
\end{emcppslisting}
    
\noindent When deriving from a base class in which inheriting most
(but not all) of its constructors is desirable, suppressing
inheritance of one or more of them is possible by providing constructors in the
derived class having the same signature as the ones that would be
inherited:

\begin{emcppshiddenlisting}[emcppsbatch=e3]
#include <iostream>  // (ù{\codeincomments{std::cout}}ù)
\end{emcppshiddenlisting}
\begin{emcppslisting}[emcppsbatch=e3]
struct B3
{
    B3()         { std::cout << "B3()\n"; }
    B3(int)      { std::cout << "B3(int)\n"; }
    B3(int, int) { std::cout << "B3(int, int)\n"; }
};

struct D3 : B3
{
    using B3::B3;
    D3(int) { std::cout << "D3(int)\n"; }
};

D3 d;        // prints (ù{\codeincomments{"B3()"}}ù)
D3 e(0);     // prints (ù{\codeincomments{"D3(int)"}}ù) --- The derived constructor is invoked.
D3 f(0, 0);  // prints (ù{\codeincomments{"B3(int, int)"}}ù)
\end{emcppslisting}
    
\noindent In other words, we can suppress what would otherwise be an inherited
constructor from a nominated base class by simply declaring a
replacement with the same signature in the derived class. We can then
choose to either implement it ourselves, \lstinline!default! it (see \featureref{\locationa}{defaulted-special-member-functions}), 
%``\titleref{defaulted-special-member-functions}" on page~\pageref{defaulted-special-member-functions}), 
or \lstinline!delete! it (see \featureref{\locationa}{deleted-functions}). 
% ``\titleref{deleted-functions}" on page~\pageref{deleted-functions}).

If we have chosen to inherit the constructors from multiple
base classes, we can disambiguate conflicts by declaring the offending
constructor(s) explicitly in the derived class and then delegating to
the base classes if and as appropriate:

\begin{emcppslisting}
struct B1A { B1A(int); };  // Here we have two base classes, each of which
struct B1B { B1B(int); };  // provides a conversion constructor from an (ù{\codeincomments{int}}ù).

struct D3 : B1A, B1B
{
    using B1A::B1A;  // Inherit the (ù{\codeincomments{int}}ù) constructor from base class (ù{\codeincomments{B1A}}ù).
    using B1B::B1B;  // Inherit the (ù{\codeincomments{int}}ù) constructor from base class (ù{\codeincomments{B1B}}ù).

    D3(int i) : B1A(i), B1B(i) { }  // User-declare (ù{\codeincomments{int}}ù) conversion constructor
};                                  // that delegates to bases.

D3 d3(0);  // OK, calls (ù{\codeincomments{D3(int)}}ù)
\end{emcppslisting}
    
\noindent Lastly, inheriting constructors from a \emcppsgloss{dependent
type} affords a capability over C++03 that is more than just convenience
and avoidance of boilerplate code.
%  \textit{\titleref{annoyances-inheritingctor}: \titleref{access-levels-of-inherited-constructors-are-same-as-in-base-class}} on page~\pageref{access-levels-of-inherited-constructors-are-same-as-in-base-class}.}} 
  In all of the example code in \intrarefsimple{description-inheritingctor} 
%  \textit{\titleref{description-inheritingctor}} 
  thus far, we know how to ``spell" the
base-class constructor; we are simply automating some drudge work. In
the case of a \emph{dependent} base class, however, we do \emph{not}
know how to spell the constructors, so we \emph{must} rely on
\emcppsgloss{inheriting constructors} if that is the forwarding semantic we
seek:

\begin{emcppslisting}
template <typename T>
struct S : T  // The base type, (ù{\codeincomments{T}}ù), is a *dependent type*.
{
    using T::T;  // inheriting constructors generically from a dependent type
};

#include <string>   // (ù{\codeincomments{std::string}}ù)
#include <vector>   // (ù{\codeincomments{std::vector}}ù)

S<std::string>        ss("hello");    // OK, uses constructor from base
S<std::vector<char>> svc("goodbye");  // Error, no suitable constructor in base
\end{emcppslisting}
    
\noindent In this example, we created a class template, \lstinline!S!, that
derives publicly from its template argument, \lstinline!T!. Then, when
creating an object of type \lstinline!S! parameterized by
\lstinline!std::string!, we were able to pass it a string literal via
the inherited \lstinline!std::string! constructor overloaded on a
\lstinline!const!~\lstinline!char*!. Notice, however, that no such constructor
is available in \lstinline!std::vector!; hence, attempting to create the
derived class from a literal string results in a compile-time error. See \intraref{use-cases-ctorinheriting}{incorporating-reusable-functionality-via-a-mix-in-class}. 
%\textit{\titleref{use-cases-ctorinheriting}: \titleref{incorporating-reusable-functionality-via-a-mix-in-class}} on page~\pageref{incorporating-reusable-functionality-via-a-mix-in-class}.

A decidedly more
complex alternative affording a different set of trade-offs would
  involve variadic template constructors (see \featureref{\locationc}{variadictemplate}) 
%  ``\titleref{variadictemplate}" on page~\pageref{variadictemplate}) 
  having
 forwarding references (see \featureref{\locationc}{forwardingref}) 
% ``\titleref{forwardingref}" on page~\pageref{forwardingref}) 
 as parameters. In this
  alternative approach, all of the constructors from the
  \mbox{\lstinline!public!}, \mbox{\lstinline!protected!}, and \mbox{\lstinline!private!} regions of the
  bases class would now appear under the same access specifier --- i.e.,
  the one in which the perfectly forwarding constructor is declared.
  What's more, this approach would not retain other constructor
  characteristics, such as \lstinline!explicit!, \lstinline!noexcept!,
  \lstinline!constexpr!, and so on. The forwarding can, however, be restricted to
  inheriting just the \lstinline!public! constructors (without
  characteristics) by constraining on \lstinline!std::is_constructible!
  using \emcppsgloss{SFINAE}; see \intraref{annoyances-inheritingctor}{access-levels-of-inherited-constructors-are-same-as-in-base-class}.

\subsection[Use Cases]{Use Cases}\label{use-cases-ctorinheriting}
%\subsubsection[Abstract use case]{Abstract use case}
\label{abstract-use-case}

Employing this form of \lstinline!using! declaration to inherit a nominated
base class's constructors --- essentially verbatim --- suggests that one
or more of those constructors is sufficient to initialize the
\emph{entire} derived-class object to a valid useful state. Typically,
such will pertain only when the derived class adds no member data of its
own. While additional derived-class member data could possibly
be initialized if by a \emph{defaulted} default constructor, this state must be \emph{orthogonal} to
any modifiable state initialized in the base class, as such state is
subject to independent change via \emcppsgloss{slicing}, which might in turn
invalidate \emcppsgloss{object invariants}. Derived-class data will be either default-initialized or have its value set using member
initializers (see \featureref{\locationc}{Default-Member-Initializers}).  
%``\titleref{Default-Member-Initializers}" on page~\pageref{Default-Member-Initializers}). 
Hence, most typical use cases will involve wrapping an
existing class by deriving from it (either publicly or privately),
adding only defaulted data members having orthogonal values, and then
adjusting the derived class's behavior via \emcppsgloss{overriding} its virtual
or \emcppsgloss[hide]{hiding} its nonvirtual member functions.

\subsubsection[Avoiding boilerplate code when employing structural inheritance]{Avoiding boilerplate code when employing structural inheritance}\label{avoiding-boilerplate-code-code-when-employing-structural-inheritance}

A key indication for using inheriting constructors is that the
derived class addresses only auxiliary or optional, rather than required
or necessary, functionality to its self-sufficient base class. As an
interesting, albeit mostly pedagogical, example, suppose we
want to provide a proxy for a \lstinline!std::vector! that performs
explicit checking of indices supplied to its index operator:

\begin{emcppslisting}
#include <cassert>                                                              
#include <vector>

template <typename T>
struct CheckedVector : std::vector<T>
{
    using std::vector<T>::vector;      // Inherit (ù{\codeincomments{std::vector}}ù)'s constructors.

    T& operator[](std::size_t index)   // Hide (ù{\codeincomments{std::vector}}ù)'s index operator.
    {
         assert(index < std::vector<T>::size());
         return std::vector<T>::operator[](index);
    }

    const T& operator[](std::size_t index) const  // Hide (ù{\codeincomments{const}}ù) index operator.
    {
         assert(index < std::vector<T>::size());
         return std::vector<T>::operator[](index);
    }
};
\end{emcppslisting}
    
\noindent In the example above, inheriting constructors allowed us to use public
structural inheritance to readily create a distinct new type having
all of the functionality of its base type except for a couple of
functions where we chose to augment the original behavior.

Although this example might be compelling, it suffers from inherent deficiencies making it insufficient for general use in practice: Passing the derived class to a function --- whether by value or reference --- will
strip it of its auxiliary functionality. When we have access to the source, an alternative solution would be to use conditional compilation to add explicit checks in certain build configurations (e.g., using C-style \lstinline!assert! macros).{\cprotect\footnote{A more robust solution along these same lines is anticipated for a future release of the C++ language standard and will be addressed in \cite{lakos23}.}} 

\subsubsection[Avoiding boilerplate code when employing implementation inheritance]{Avoiding boilerplate code when employing implementation inheritance}\label{avoiding-boilerplate-code-when-employing-implementation-inheritance}

Sometimes it can be cost effective to adapt a \emcppsgloss{concrete class}
having virtual functions to a specialized purpose
by using inheritance. Useful design patterns
exist where a \emcppsgloss{partial implementation} class, derived from a
\emcppsgloss{pure abstract interface} (a.k.a. a \emcppsgloss{protocol}), contains data,
constructors, and pure virtual functions.\footnote{See
  \cite{lakos2a}, section~4.7.} Such inheritance, known as
\emcppsgloss{implementation inheritance}, is decidedly distinct from pure \emcppsgloss{interface inheritance}, which is often the preferred design
  pattern in practice.\footnote{See \cite{lakos2b}, section~4.6.} As an example, consider a base class,
\lstinline!NetworkDataStream!, that allows overriding its virtual functions
for processing a stream of data from an expanding variety of arbitrary
sources over the network:

\begin{emcppshiddenlisting}[emcppsbatch=e4]
#include <iostream>   // (ù{\codeincomments{std::ostream}}ù)
class DataPacket;                                                               
class RawDataStreamHandle;                                                      
class TCPConnection;                                                            
class UDPConnection;
std::ostream& operator<<(std::ostream& stream, const DataPacket& data);

struct LogTrace {};
constexpr LogTrace LOG_TRACE; // Quick and dirty syntactic simulation of logger
template <typename T>
const LogTrace& operator<<(const LogTrace& lhs, const T&) { return lhs; }
\end{emcppshiddenlisting}
\begin{emcppslisting}[emcppsbatch=e4]
class NetworkDataStream
{
private:
    // ...                   (member data)

public:
    explicit NetworkDataStream(TCPConnection* tcpConnection);
    explicit NetworkDataStream(UDPConnection* udpConnection);
    explicit NetworkDataStream(RawDataStreamHandle* rawDataStreamHandle);

    virtual ~NetworkDataStream();

    virtual void onPacketReceived(DataPacket& dataPacket) = 0;
        // Derived classes must override this method.
};        
\end{emcppslisting}
    
\noindent The \lstinline!NetworkDataStream! class above provides three constructors, with more
under development, that can be used assuming no per-packet processing is
required. Now, imagine the need for logging information about received
packets (e.g., for auditing purposes). Inheriting constructors make
deriving from \lstinline!NetworkDataStream! and overriding (see \featureref{\locationa}{override}) 
%``\titleref{override}" on page~\pageref{override}) 
\lstinline!onPacketReceived(DataPacket&)! more convenient because we don't need to reimplement each of the constructors, which are
anticipated to increase in number over time:

\begin{emcppslisting}[emcppsbatch=e4]
class LoggedNetworkDataStream : public NetworkDataStream
{
public:
    using NetworkDataStream::NetworkDataStream;

    void onPacketReceived(DataPacket& dataPacket) override
    {
        LOG_TRACE << "Received packet " << dataPacket;    // local log facility
        NetworkDataStream::onPacketReceived(dataPacket);  // Delegate to base.
    }
};
\end{emcppslisting}
    

\subsubsection[Implementing a strong \lstinline!typedef!]{Implementing a strong {\SubsubsecCode typedef}}\label{implementing-a-strong-typedef}

Classic \lstinline!typedef! declarations --- just like C++11
\lstinline!using! declarations (see \featureref{\locationa}{alias-declarations-and-alias-templates}) 
% ``\titleref{alias-declarations-and-alias-templates}" on page~\pageref{alias-declarations-and-alias-templates}) 
--- are just synonyms; they
offer absolutely no additional type safety over using the original type name. A commonly desired capability is to
provide an alias to an existing type \lstinline!T! that is uniquely
interoperable with itself, explicitly convertible from \lstinline!T!, but
not implicitly convertible from \lstinline!T!. This somewhat \emph{more}
type-safe form of alias is sometimes referred to as a \emcppsgloss[strong typedef]{strong \lstinline!typedef!}. A typical implementation of a \emcppsgloss[strong typedef]{strong \lstinline!typedef!} suppresses implicit conversions both from the new type to the type it wraps and vice versa via \lstinline!explicit! converting constructors and \lstinline!explicit! conversion operators.  In this respect, the relationship of \emcppsgloss[strong typedef]{strong \lstinline!typedef!} with the type it wraps is analogous to that of a scoped enumeration (\lstinline!enum!~\lstinline!class!)  to its \emcppsgloss{underlying type}; see \featureref{\locationc}{explicit-enumeration-underlying-type}. 
%A so-called \romeogloss{strong
%  \lstinline!typedef!} is similar to a classic, C-style enumeration in
%  that it is (1) its own type and (2) implicitly convertible to its
%  base type (which for enumerators corresponds to its
%  \romeogloss{underlying type}; see \featureref{\locationc}{explicit-enumeration-underlying-type}). 
%%  ``\titleref{explicit-enumeration-underlying-type}" on page~\pageref{explicit-enumeration-underlying-type}). 
%  Unlike a classic
%  \lstinline!enum!, however, a typical implementation of a \romeogloss{strong
%  \lstinline!typedef!} allows only for explicit conversion from its base
%  type. An analogy to the more strongly typed
%\lstinline!enum!~\lstinline!class! (see \featureref{\locationc}{enumclass}) 
%% ``\titleref{enumclass}" on page~\pageref{enumclass}) 
%would suppress conversion
%  in either direction, e.g., via private inheritance and then
%  \lstinline!explicit! conversion constructors and
% \lstinline!explicit! conversion operators (see \featureref{\locationa}{explicit-conversion-operators}).
% [Underlying Type'11]{.xref}, [`enum class`]{.xref}, and [`explicit` Operators]{.xref}.}}
% ``\titleref{explicit-conversion-operators}" on page~\pageref{explicit-conversion-operators}).}}

As a practical example, suppose we are exposing, to a fairly wide and varied audience, a class,
\lstinline!PatientInfo!, that associates two \lstinline!Date! objects to a
given hospital patient:

\begin{emcppslisting}[emcppsbatch=e5]
class Date
{
    // ...

public:
    Date(int year, int month, int day);

    // ...
};

class PatientInfo
{
private:
    Date d_birthday;
    Date d_appointment;

public:
    PatientInfo(Date birthday, Date appointment);
        // Please pass the birthday as the first date and the appointment as
        // the second one!
};
\end{emcppslisting}
    
\noindent For the sake of argument, imagine that our users are not as diligent as
they should be in reading documentation to know which constructor
argument is which:

\begin{emcppslisting}[emcppsbatch=e5]
PatientInfo client1(Date birthday, Date appointment)
{
    return PatientInfo(birthday, appointment);  // OK
}

int client2(PatientInfo* result, Date birthday, Date appointment)
{
    *result = PatientInfo(appointment, birthday);  // Oops! wrong order
    return 0;
}
\end{emcppslisting}
    
\noindent Now suppose that we continue to get complaints, from folks like
\lstinline!client2! in the example above, that our code doesn't work. What can we
do?

Although this example is presented lightheartedly, misuse by clients is a perennial problem in large-scale
software organizations. Choosing the same type for both arguments
might well be the right choice in some environments but not in others.
We are not advocating use of this technique; we are merely
  acknowledging that it exists.

One way is to force clients to make a conscious and explicit decision in
their own source code as to which \lstinline!Date! is the birthday and
which is the appointment. Employing a \emcppsgloss[strong typedef]{strong \lstinline!typedef!}
can help us to achieve this goal. Inheriting constructors provide a
concise way to define a \emcppsgloss[strong typedef]{strong \lstinline!typedef!}; for
the example above, they can be used to define two new types to
 uniquely represent a birthday and an appointment date:

\begin{emcppslisting}[emcppsbatch=e5]
struct Birthday : Date  // somewhat type-safe alias for a (ù{\codeincomments{Date}}ù)
{
    using Date::Date;  // inherit (ù{\codeincomments{Date}}ù)'s three integer ctor
    explicit Birthday(Date d) : Date(d) { }  // (ù{\codeincomments{explicit}}ù) conversion from (ù{\codeincomments{Date}}ù)
};

struct Appointment : Date  // somewhat type-safe alias for a (ù{\codeincomments{Date}}ù)
{
    using Date::Date;  // inherit (ù{\codeincomments{Date}}ù)'s three integer ctor
    explicit Appointment(Date d) : Date(d) { }  // (ù{\codeincomments{explicit}}ù) conv. from (ù{\codeincomments{Date}}ù)
};
\end{emcppslisting}
    
\noindent The \lstinline!Birthday! and \lstinline!Appointment! types expose the same
interface of \lstinline!Date!, yet, given our inheritance-based design,
\lstinline!Date! is not implicitly convertible to either. Most importantly,
however, these two new types are not implicitly convertible to each
other:

\begin{emcppslisting}[emcppsbatch=e5]
Birthday b0(1994, 10, 4);  // OK, thanks to inheriting constructors
Date d0 = b0;              // OK, thanks to public inheritance
Birthday b1 = d0;          // Error, no implicit conversion from (ù{\codeincomments{Date}}ù)
Appointment a0;            // Error, (ù{\codeincomments{Appointment}}ù) has no default ctor.
Appointment a1 = b0;       // Error, no implicit conversion from (ù{\codeincomments{Birthday}}ù)
Birthday n2(d0);           // OK, thanks to an (ù{\codeincomments{explicit}}ù) constructor in (ù{\codeincomments{Birthday}}ù)
Appointment a2(1999, 9, 17); // OK, thanks to inheriting constructors           
Birthday    b3(a2);          // OK, an (ù{\codeincomments{Appointment}}ù) (unfortunately) is a (ù{\codeincomments{Date.}}ù)
\end{emcppslisting}
    
\noindent We can now reimagine a \lstinline!PatientInfo! class that exploits this
newfound (albeit artificially manufactured) type-safety:

\begin{emcppshiddenlisting}[emcppsbatch={e6,e7}]
class Date
{
public:
    Date(int year, int month, int day);
};
struct Birthday : Date  // somewhat type-safe alias for a (ù{\codeincomments{Date}}ù)
{
    using Date::Date;  // inherit (ù{\codeincomments{Date}}ù)'s three integer ctor
    explicit Birthday(Date d) : Date(d) { }  // (ù{\codeincomments{explicit}}ù) conversion from (ù{\codeincomments{Date}}ù)
};

struct Appointment : Date  // somewhat type-safe alias for a (ù{\codeincomments{Date}}ù)
{
    using Date::Date;  // inherit (ù{\codeincomments{Date}}ù)'s three integer ctor
    explicit Appointment(Date d) : Date(d) { }  // (ù{\codeincomments{explicit}}ù) conv. from (ù{\codeincomments{Date}}ù)
};
\end{emcppshiddenlisting}
\begin{emcppslisting}[emcppsbatch={e6,e7}]
class PatientInfo
{
private:
    Birthday d_birthday;
    Appointment d_appointment;

public:
    PatientInfo(Birthday birthday, Appointment appointment);
        // Please pass the birthday as the first argument and the appointment as
        // the second one! 
};
\end{emcppslisting}
    
\noindent Now our clients have no choice but to make their intentions clear at the
call site. The previous implementation of the client functions no longer compile:

\begin{emcppslisting}[emcppsbatch=e6,emcppserrorlines={1,2,4}]
PatientInfo client1(Date birthday, Date appointment)
{
    return PatientInfo(birthday, appointment);     // Error, doesn't compile.
}

int client2(PatientInfo* result, Date birthday, Date appointment)
{
    *result = PatientInfo(appointment, birthday);  // Error, doesn't compile.
    return 0;
}
\end{emcppslisting}
Because the clients now need to explicitly convert their \lstinline!Date! objects to the appropriate \emcppsgloss[strong typedef]{strong \lstinline!typedef!s}, it is easy to spot and fix the defect in \lstinline!client2!:
\begin{emcppslisting}[emcppsbatch=e7]
PatientInfo client1(Date birthday, Date appointment)
{
    return PatientInfo(Birthday(birthday), Appointment(appointment));  // OK
}

int client2(PatientInfo* result, Date birthday, Date appointment)
{
    Birthday b(birthday);
    Appointment a(appointment);
    *result = PatientInfo(b, a);  // OK
}
\end{emcppslisting}
    
\noindent 
%This example works because the \romeogloss{value constructor} takes three
%arguments and cannot be invoked as part of an implicit conversion
%sequence; see \intraref{potential-pitfalls-ctorinheriting}{beware-of-inheriting-implicit-constructors}. 
%%\textit{\titleref{potential-pitfalls-ctorinheriting}: \titleref{beware-of-inheriting-implicit-constructors}} on page~\pageref{beware-of-inheriting-implicit-constructors}. 
%Note that, in an ideal
%world where thorough unit testing is ubiquitous, such machinations would
%most likely be supererogatory.
In this example, the client functions failed to compile after the introduction of the \emcppsgloss[strong typedef]{strong \lstinline!typedef!s}, which is the intended effect. However, if \lstinline!Date! objects were \emph{implicitly} constructed when client functions created \lstinline!PatientInfo!, the defective code would continue to compile because both \emcppsgloss[strong typedef]{strong \lstinline!typedef!s} can be implicitly constructed from the same arguments; see \intraref{potential-pitfalls-ctorinheriting}{beware-of-inheriting-implicit-constructors}.

Replicating
types that have identical behavior in the name of type safety can run
afoul of interoperability. Distinct types that are otherwise
physically similar are often most appropriate when their respective
behaviors are inherently distinct and unlikely to interact in practice
(e.g., a \mbox{\lstinline!CartesianPoint!} and a \mbox{\lstinline!RationalNumber!},
each implemented as having two integral data members).\footnote{See
  \cite{lakos2a}, section~4.4.}

\subsubsection[Incorporating reusable functionality via a mix-in]{Incorporating reusable functionality via a mix-in}\label{incorporating-reusable-functionality-via-a-mix-in-class}

Some classes are designed to generically enhance the behavior of a class
just by inheriting from it; such classes are sometimes referred to as
\emph{mix-ins}. If we want to adapt a class to support the additional
behavior of the mix-in, with no other change to its behavior, we can use
simple \emcppsgloss{structural inheritance} (e.g., to preserve reference
compatibility through function calls). To preserve the public interface,
however, we will need it to inherit the constructors as well.

Consider, for example, a simple class to track the total number of
objects created:

\begin{emcppslisting}[emcppsbatch=e8]
template <typename T>
struct CounterImpl  // mix-in used to augment implementation of arbitrary type
{
    static int s_constructed;  // count of the number of (ù{\codeincomments{T}}ù) objects constructed

    CounterImpl()                   { ++s_constructed; }
    CounterImpl(const CounterImpl&) { ++s_constructed; }
};

template <typename T>
int CounterImpl<T>::s_constructed;  // required member definition
\end{emcppslisting}
    
\noindent The class template \lstinline!CounterImpl!, in the example above, counts the number of
times an object of type \lstinline!T! was constructed during a run of the
program. We can then write a generic adapter, \lstinline!Counted!, to
facilitate use of \lstinline!CounterImpl! as a \emph{mix-in}:

\begin{emcppslisting}[emcppsbatch=e8]
template <typename T>
struct Counted : T, CounterImpl<T>
{
    using T::T;
};
\end{emcppslisting}
    
\noindent Note that the \lstinline!Counted! adapter class inherits all of the
constructors of the \emph{dependent} class, \lstinline!T!, that it wraps,
without its having to know what those constructors are:

\begin{emcppshiddenlisting}[emcppsbatch=e8]
// myfoo.h:
struct MyFoo {};
\end{emcppshiddenlisting}
\begin{emcppslisting}[emcppsbatch=e8]
#include <string>   // (ù{\codeincomments{std::string}}ù)
#include <vector>   // (ù{\codeincomments{std::vector}}ù)
#include <myfoo.h>  // (ù{\codeincomments{MyFoo}}ù)

Counted<std::string>       cs ("ABC");  // Construct a counted (ù{\codeincomments{string}ù).          
Counted<std::vector<char>> cvc(3, 'a'); // Construct a counted vector of (ù{\codeincomments{char}ù).
Counted<MyFoo>             cmf;         // Construct a counted (ù{\codeincomments{MyFoo}}ù) object.
\end{emcppslisting}
    
\noindent While inheriting constructors are a convenience in nongeneric
programming, they can be an essential tool for generic idioms.

\subsection[Potential Pitfalls]{Potential Pitfalls}\label{potential-pitfalls-ctorinheriting}

\subsubsection[Newly introduced constructors in the base class can silently alter program behavior]{Newly introduced constructors in the base class can silently alter\\[0.5ex] program behavior}\label{newly-introduced-constructors-in-the-base-class-can-silently-alter-program-behavior}

The introduction of a new constructor in a base class might silently
change a program's run-time behavior if that constructor happens to be a
better match during overload resolution of an existing instantiation of
a derived class. Consider a \lstinline!Session! class that initially provides only two constructors:

\begin{emcppshiddenlisting}[emcppsbatch=e9]
class RawSessionHandle;
\end{emcppshiddenlisting}
\begin{emcppslisting}[emcppsbatch=e9]
struct Session
{
    Session();
    explicit Session(RawSessionHandle* rawSessionHandle);
};
\end{emcppslisting}
    
\noindent Now, imagine that a class, \lstinline!AuthenticatedSession!, derived from
\lstinline!Session!, inherits the two constructors of its base class and provides its own constructor that accepts an integral
authentication token:

\begin{emcppslisting}[emcppsbatch=e9]
struct AuthenticatedSession : Session
{
    using Session::Session;
    explicit AuthenticatedSession(long long authToken);
};
\end{emcppslisting}
    
\noindent Finally, consider an instantiation of \lstinline!AuthenticatedSession! in
user-facing code:

\begin{emcppslisting}[emcppsbatch=e9]
AuthenticatedSession authSession(45100);
\end{emcppslisting}
    
\noindent In the example above, \lstinline!authSession! will be initialized by
invoking the constructor accepting a \lstinline!long!~\lstinline!long! (see \featureref{\locationa}{long-long}) 
%``\titleref{long-long}" on page~\pageref{long-long}) 
authentication token. If, however, a new
constructor having the signature \lstinline!Session(int!~\lstinline!fd)! is added to the
base class, it will be invoked instead because it is a better match to the
literal \lstinline!45100! (of type \lstinline!int!) than the constructor
taking a \lstinline!long!~\lstinline!long! supplied explicitly in the derived
class; hence, adding a constructor to a base class might lead to a
potential latent defect that would go unreported at compile
time.

Note that this problem with implicit conversions for function parameters is not unique
to inheriting constructors; any form of \lstinline!using! declaration or
invocation of an overloaded function carries a similar risk. Imposing
stronger typing --- e.g., by using \emcppsgloss[strong typedef]{strong \lstinline!typedef!s} (see \intraref{use-cases-ctorinheriting}{implementing-a-strong-typedef}) 
%\textit{\titleref{use-cases-ctorinheriting}: \titleref{implementing-a-strong-typedef}} on page~\pageref{implementing-a-strong-typedef}) 
--- might sometimes, however, help to prevent such
unfortunate missteps.

\subsubsection[Inheriting \emph{implicit} constructors]{Inheriting {\sfbsubsubsecitalRomeo implicit} constructors}\label{beware-of-inheriting-implicit-constructors}

Inheriting from a class that has implicit constructors can cause
surprises. Consider again using inheriting constructors to
implement a \emcppsgloss[strong typedef]{strong \lstinline!typedef!} from \intraref{use-cases-ctorinheriting}{implementing-a-strong-typedef}. 
%\textit{\titleref{use-cases-ctorinheriting}: \titleref{implementing-a-strong-typedef}} on page~\pageref{implementing-a-strong-typedef}. 
This time, however, let's suppose we are
exposing a class,
\lstinline!PointOfInterest!, that associates the name and address of a
given popular tourist attraction:

\begin{emcppslisting}[emcppsbatch=e10]
#include <string>  // (ù{\codeincomments{std::string}}ù)

class PointOfInterest
{
private:
    std::string d_name;
    std::string d_address;

public:
    PointOfInterest(const std::string& name, const std::string& address);
        // Please pass the (ù{\codeincomments{name}}ù) as the *first* and the (ù{\codeincomments{address}}ù) *second*!
};
\end{emcppslisting}
    
\noindent Again imagine that our users are not always careful about inspecting the
function prototype:

\begin{emcppslisting}[emcppsbatch=e10]
PointOfInterest client1(const std::string& name, const std::string& address)
{
    return PointOfInterest(name, address);  // OK
}

int client2(PointOfInterest*   result,                                          
            const std::string& name,                                            
            const std::string& address)
{
    *result = PointOfInterest(address, name);  // Oops! wrong order
    return 0;
}
\end{emcppslisting}
    
\noindent We might think to again use \emcppsgloss[strong typedef]{strong \lstinline!typedef!s} here as we
did for \lstinline!PatientInfo! in \intraref{use-cases-ctorinheriting}{implementing-a-strong-typedef}: 
%\textit{\titleref{use-cases-ctorinheriting}: \titleref{implementing-a-strong-typedef}} on page~\pageref{implementing-a-strong-typedef}:

\begin{emcppshiddenlisting}[emcppsbatch=e11]
#include <string>   // (ù{\codeincomments{std::string}}ù)
\end{emcppshiddenlisting}
\begin{emcppslisting}[emcppsbatch=e11]
struct Name : std::string  // somewhat type-safe alias for a (ù{\codeincomments{std::string}}ù)
{
    using std::string::string;  // Inherit, as is, all of (ù{\codeincomments{std::string}}ù)'s ctors.
    explicit Name(const std::string& s) : std::string(s) { }  // conversion
};

struct Address : std::string  // somewhat type-safe alias for a (ù{\codeincomments{std::string}}ù)
{
    using std::string::string;  // Inherit, as is, all of (ù{\codeincomments{std::string}}ù)'s ctors.
    explicit Address(const std::string& s) : std::string(s) { }  // conversion
};
\end{emcppslisting}
    
\noindent The \lstinline!Name! and \lstinline!Address! types are not interconvertible; they expose the same interfaces as \lstinline!std::string! but are not
implicitly convertible from it:

\begin{emcppslisting}[emcppsbatch=e11]
Name n0 = "Big Tower";  // OK, thanks to inheriting constructors
std::string s0 = n0;    // OK, thanks to public inheritance
Name n1 = s0;           // Error, no implicit conversion from (ù{\codeincomments{std::string}}ù)
Address a0;             // OK, unfortunately a (ù{\codeincomments{std::string}}ù) has a default ctor.
Address a1 = n0;        // Error, no implicit conversion from (ù{\codeincomments{Name}}ù)
Name n2(s0);            // OK, thanks to an (ù{\codeincomments{explicit}}ù) constructor in (ù{\codeincomments{Name}}ù)
Name b3(a0);            // OK, an (ù{\codeincomments{Address}}ù) (unfortunately) is a (ù{\codeincomments{std::string}}ù).
\end{emcppslisting}
    
\noindent We can rework the \lstinline!PointOfInterest! class to use the
\emcppsgloss[strong typedef]{strong \lstinline!typedef!} idiom:

\begin{emcppslisting}[emcppsbatch=e11]
class PointOfInterest
{
private:
    Name    d_name;
    Address d_address;

public:
    PointOfInterest(const Name& name, const Address& address);
};
\end{emcppslisting}
    
\noindent Now if our clients use the base class itself as a parameter, they will
again need to make their intentions known:

\begin{emcppslisting}[emcppsbatch=e11]
PointOfInterest client1(const std::string& name, const std::string& address)    
{                                                                               
    return PointOfInterest(Name(name), Address(address));
} 

int client2(PointOfInterest*   result,                                          
            const std::string& name,                                            
            const std::string& address)                                         
{
   *result = PointOfInterest(Name(name), Address(address)); // Fix forced.     
    return 0;                                                                   
}
\end{emcppslisting}
    
\noindent But suppose that some clients pass the arguments by
\lstinline!const!~\lstinline!char*! instead of\linebreak[4]
\mbox{\lstinline!const!~\lstinline!std::string&!}:

\begin{emcppslisting}[emcppsbatch=e11]
PointOfInterest client3(const char* name, const char* address)
{
    return PointOfInterest(address, name);  // Bug, compiles but runtime error
}
\end{emcppslisting}
    
\noindent In the case of \lstinline!client3! in the code snippet above, passing the arguments through
\emph{does} compile because the \lstinline!const!~\lstinline!char*!
constructors are inherited; hence, there is no attempt to convert
to a \lstinline!std::string! before matching the \emph{implicit} conversion
constructor. Had the \lstinline!std::string! conversion
constructor been declared to be \lstinline!explicit!, the code would not
have compiled. In short, inheriting constructors from types that perform implicit
conversions can seriously undermine the effectiveness of the \emcppsgloss[strong typedef]{strong
\lstinline!typedef!} idiom.

\subsection[Annoyances]{Annoyances}\label{annoyances-inheritingctor}

\subsubsection[Inherited constructors cannot be selected individually]{Inherited constructors cannot be selected individually}\label{inherited-constructors-cannot-be-selected-individually}

The inheriting-constructors feature does not allow the programmer
to select a subset of constructors to inherit; all of the base
class's eligible constructors are always inherited unless a constructor
with the same signature is provided in the derived class. If the
programmer desires to inherit all constructors of a base class except
for perhaps one or two, the straightforward workaround would be to
declare the undesired constructors in the derived class and then use
deleted functions (see \featureref{\locationa}{deleted-functions}) 
% ``\titleref{deleted-functions}" on page~\pageref{deleted-functions}) 
to explicitly exclude them.

For example, suppose we have a general class, \lstinline!Datum!, that can
be constructed from a variety of types:

\begin{emcppslisting}[emcppsbatch=e12]
struct Datum
{
    Datum(bool);
    Datum(char);
    Datum(short);
    Datum(int);
    Datum(long);
    Datum(long long);
};
\end{emcppslisting}
    
\noindent If we wanted to create a version of \lstinline!Datum!, call it
\lstinline!NumericalDatum!, that inherits all but the one constructor
taking a \lstinline!bool!, our derived class would (1) inherit publicly, 
(2) declare the unwanted constructor, and then (3) mark it with
\lstinline!=!\,\lstinline!delete!:

\begin{emcppslisting}[emcppsbatch=e12]
struct NumericalDatum : Datum
{
    using Datum::Datum;             // Inherit all the constructors...
    NumericalDatum(bool) = delete;  // ...except the one taking a (ù{\codeincomments{bool}}ù).
};
\end{emcppslisting}
    
\noindent Note that the subsequent addition of any non-numerical constructor to
\lstinline!Datum! (e.g., a constructor taking \lstinline!std::string!) would
defeat the purpose of \lstinline!NumericalDatum! unless that inherited constructor were explicitly excluded from \lstinline!NumericalDataum! by use of \lstinline!=!\,\lstinline!delete!.

\subsubsection[Access levels of inherited constructors are the same as in base class]{Access levels of inherited constructors are the same as in base class}\label{access-levels-of-inherited-constructors-are-same-as-in-base-class}

Unlike base-class member functions that can be introduced with a
\lstinline!using! directive with an arbitrary access level into the derived
class (as long as they are accessible by the derived class), the access
level of the \lstinline!using! declaration for inherited constructors is
ignored. The inherited constructor overload is instead
accessible \emph{if} the corresponding base-class constructor would be
accessible:

\begin{emcppslisting}[emcppsbatch=e13]
struct Base
{
private:
    Base(int) { }  // This constructor is declared (ù{\codeincomments{private}}ù) in the base class.
    void pvt0() { }
    void pvt1() { }

public:
    Base() { }     // This constructor is declared (ù{\codeincomments{public}}ù) in the base class.
    void pub0() { }
    void pub1() { }
};
\end{emcppslisting}
    
\noindent Note that, when employing \lstinline!using! to (1) inherit constructors or
(2) elevate base-class definitions in the presence of private
inheritance, public clients of the class might find it necessary 
to look at what are ostensibly private implementation details of the
derived class to make proper use of that type through its
public interface:

\begin{emcppslisting}[emcppsbatch=e13]
struct Derived(ù{\footnotemark}ù) : private Base
{
    using Base::Base;  // OK, inherited (ù{\codeincomments{Base()}}ù) as (ù{\codeincomments{public}}ù) constructor
                       // and (ù{\codeincomments{Base(int)}}ù) as (ù{\codeincomments{private}}ù) constructor

private:
    using Base::pub0;  // OK, (ù{\codeincomments{pub0}}ù) is declared (ù{\codeincomments{private}}ù) in derived class.
    using Base::pvt0;  // Error, (ù{\codeincomments{pvt0}}ù) was declared private in base class.

public:
    using Base::pub1;  // OK, (ù{\codeincomments{pub1}}ù) is declared (ù{\codeincomments{public}}ù) in derived class.
    using Base::pvt1;  // Error, (ù{\codeincomments{pvt1}}ù) was declared (ù{\codeincomments{private}}ù) in base class.
};

void client()
{
     Derived x(0);  // Error, Constructor was declared private in base class.
     Derived d;     // OK, constructor was declared (ù{\codeincomments{public}}ù) in base class.
     d.pub0();      // Error, (ù{\codeincomments{pub0}}ù) was declared (ù{\codeincomments{private}}ù) in derived class.
     d.pub1();      // OK, (ù{\codeincomments{pub1}}ù) was declared (ù{\codeincomments{public}}ù) in derived class.
     d.pvt0();      // Error, (ù{\codeincomments{pvt0}}ù) was declared private in base class.
     d.pvt1();      // Error, (ù{\codeincomments{pvt1}}ù) was declared private in base class.
}
\end{emcppslisting}
{\cprotect\footnotetext{Alisdair Meredith, one of the authors of the
Standards paper that proposed this feature (\cite{meredith08}), suggests that
placing the \lstinline!using! declaration for \emcppsgloss{inheriting
constructors} as the very first member declaration and preceding any
\emcppsgloss[access specifier]{access specifiers} might be the least confusing location. Programmers might still be confused by
the disparate default access levels of \lstinline!class! versus
  \lstinline!struct!.}}
    
\noindent This C++11 feature was itself created because the previously proposed
solution --- which also involved a couple of features new in C++11, namely,
forwarding the arguments to base-class constructors with
forwarding references (see \featureref{\locationc}{forwardingref}) 
% ``\titleref{forwardingref}" on page~\pageref{forwardingref}) 
and variadic
templates (see \featureref{\locationc}{variadictemplate}) 
% ``\titleref{variadictemplate}" on page~\pageref{variadictemplate}) 
--- made somewhat different trade-offs and was
considered too onerous and fragile to be practically useful:

\begin{emcppslisting}
#include <utility>  // (ù{\codeincomments{std::forward}}ù)

struct Base
{
    Base(int) { }
};

struct Derived : private Base
{
protected:
    template <typename... Args>
    Derived(Args&&... args) : Base(std::forward<Args>(args)...)
    {
    }
};
\end{emcppslisting}
    
\noindent In the example above, we have used forwarding references (see \featureref{\locationc}{forwardingref}) 
% ``\titleref{forwardingref}" on page~\pageref{forwardingref})
to properly delegate the implementation of a constructor that is
declared \mbox{\lstinline!protected!} in the derived class to a \lstinline!public!
constructor of a privately inherited base class. Although this approach
fails to preserve many of the characteristics of the inheriting
constructors (e.g., \lstinline!explicit!, \lstinline!constexpr!,
\lstinline!noexcept!, and so on), the functionality described in the code snippet above is
simply not possible using the C++11 inheriting-constructors feature.

\subsubsection[Flawed initial specification led to diverging early implementations]{Flawed initial specification led to diverging early implementations}

The original specification of inheriting constructors in C++11 had a
significant number of problems with general use.\footnote{For the detailed analysis of the issues that were the consequence of the flawed initial C++11 specification of inheriting constructors, see [PRODUCTION: LINK TO BOOK WEBSITE SUPPLEMENTAL MATERIAL.]} As originally specified,
inherited constructors were treated as if they were redeclared in
the derived class. For C++17, a significant rewording of this feature\footnote{\cite{smith15b}}
happened to instead find the base class
constructors and then define how they are used to construct an instance
of the derived class, as we have presented here. With a final fix in
C++20 with the resolution of CWG issue \#2356,\footnote{\cite{smith18}} a complete
working feature was specified. All of these fixes for C++17 were
accepted as defect reports and thus apply retroactively to C++11 and
C++14. For the major compilers, this was either standardizing already
existing practice or quickly adopting the changes.{\cprotect\footnote{For example, GCC versions above 7.0
and Clang versions above 4.0 all have the modern behavior
fully implemented regardless of which standard version is chosen when
  compiling.}}

\subsection[See Also]{See Also}\label{see-also}

\begin{itemize}
\item{%``\titleref{delegating-constructors}" on page~\pageref{delegating-constructors} — 
\seealsoref{delegating-constructors}{\seealsolocationa}are used to call one constructor from another from within the same user-defined type.}
\item{%``\titleref{Defaulted-Special-Member-Functions}" on page~\pageref{Defaulted-Special-Member-Functions} — 
\seealsoref{Defaulted-Special-Member-Functions}{\seealsolocationa}are used to implement functions that might otherwise have been suppressed by inherited constructors.}
\item{%``\titleref{deleted-functions}" on page~\pageref{deleted-functions} — 
\seealsoref{deleted-functions}{\seealsolocationa}can be used to exclude inherited constructors that are unwanted entirely.}
\item{%``\titleref{override}" on page~\pageref{override} — 
\seealsoref{override}{\seealsolocationa}can be used to ensure that a member function intended to override a virtual function actually does so.}
%\item{%``\titleref{Default-Member-Initializers}" on page~\pageref{Default-Member-Initializers} — 
\item{%``\titleref{Default-Member-Initializers}" on page~\pageref{Default-Member-Initializers} — 
\seealsoref{Default-Member-Initializers}{\seealsolocationc}can be used to provide nondefault values for data members in derived classes that make use of inheriting constructors.}
\item{%``\titleref{forwardingref}" on page~\pageref{forwardingref} — 
\seealsoref{forwardingref}{\seealsolocationc}are used as an alternative (workaround) when access levels differ from those for base-class constructors.}
\item{%``\titleref{variadictemplate}" on page~\pageref{variadictemplate} — 
\seealsoref{variadictemplate}{\seealsolocationc}are used as an alternative (workaround) when access levels differ from those for base-class constructors.}
\end{itemize}

\subsection[Further Reading]{Further Reading}\label{further-reading}

\hspace*{\fill}

%%%%%%%%%%% first para moves to be a final annoyance; remainder moves to supplemental material on website. See folder "For Web," file "inheritingctor-appendix"  
%\subsection[Appendix: C++17 Improvements Made Retroactive to C++11/14]{Appendix: C++17 Improvements Made Retroactive to C++11/14}\label{appendix}\label{c++17-improvements-made-retroactive-to-c++11/14}
%
%The original specification of inheriting constructors in C++11 had a
%significant number of problems with general use. As originally specified,
%inherited constructors were treated as if they were redeclared in
%the derived class. For C++17, a significant rewording of this feature\footnote{\cite{smith15b}}
%happened to instead find the base class
%constructors and then define how they are used to construct an instance
%of the derived class, as we have presented here. With a final fix in
%C++20 with the resolution of CWG issue \#2356,\footnote{\cite{smith18}} a complete
%working feature was specified. All of these fixes for C++17 were
%accepted as defect reports and thus apply retroactively to C++11 and
%C++14. For the major compilers, this was either standardizing already
%existing practice or quickly adopting the changes.{\cprotect\footnote{For example, GCC versions above 7.0
%and Clang versions above 4.0 all have the modern behavior
%fully implemented regardless of which standard version is chosen when
%  compiling.}}
%
%The subsections that follow describe the subtle bugs that came with the previous
%specification, both for completeness and to give a better
%understanding of what to expect on very old compilers, though none fully
%implemented the original specification as written.
%
%\subsubsection[Inheriting constructors declared with a C-style ellipsis]{Inheriting constructors declared with a C-style ellipsis}
%
%Forwarding arguments from a constructor declared using a C-style ellipsis cannot forward correctly. Arguments passed through the ellipsis are not available as named arguments but must instead be accessed through the \lstinline!va_arg! family of macros. Without named arguments, no easily supported way is available to call the base-class constructor with the additional arguments:
%\begin{emcppslisting}
%struct Base
%{
%    Base(int x, ...) { }  // constructor taking C-style variadic args
%};
%
%struct Derived : Base
%{
%    using Base::Base;  // Error, Prior to C++17 fixes, standard wording
%                       // does not allow forwarding C-style variadic args.
%};
%\end{emcppslisting}
%    This problem is sidestepped in C++17 because the base-class constructor becomes available just like any other base-class function made available through a \lstinline!using! declaration in the derived class.
%    
%\subsubsection[Inheriting constructors that rely on friendship to declare function parameters]{Inheriting constructors that rely on friendship to declare\\ function parameters}
%
%When a constructor depends on access to a \lstinline!private! member of a class (e.g., a \lstinline!typedef!), an inheriting constructor does not implicitly grant friendship that the base class might have that makes the constructor valid. For example, consider the following class template, which grants friendship to class \lstinline!B!:
%\begin{emcppslisting}
%template <typename T>
%struct S
%{
%private:
%    typedef int X;
%    friend struct B;
%};
%\end{emcppslisting}
%    Then, we can create a class with a constructor that relies on that friendship. In this case, we consider a constructor template using the dependent member \lstinline!X!, assuming that, in the normal case, \lstinline!X! would be publicly accessible:
%\begin{emcppslisting}
%struct B
%{
%    template <typename T>
%    B(T, typename T::X);
%};
%\end{emcppslisting}
%    Now consider class \lstinline!D! derived from \lstinline!B! and inheriting its constructors:
%\begin{emcppslisting}
%struct D : B
%{
%    using B::B;
%};
%\end{emcppslisting}
%    Without friendship, we cannot construct a \lstinline!D! from an \lstinline!S!, but we can construct a \lstinline!B! from an \lstinline!S!, suggesting something is wrong with the inheritance. Note that the SFINAE rules for templates mean that the inheriting constructor is a problem only if we try to construct an \lstinline!S! with the problem type and does not cause a hard error without that use case. The following example illustrates the problematic usage:
%\begin{emcppslisting}
%S<int> s;   // full specialization of (ù{\codeincomments{S}}ù) for type (ù{\codeincomments{int}}ù)
%B b(s, 2);  // OK, thanks to friendship
%D d(s, 2);  // Error, Prior to C++17 fixes, friendship is not inherited.
%\end{emcppslisting}
%    As C++17 redefines the semantics of the inheriting constructor as if the base class’s constructors were merely exposed in the derived one, friendship is evaluated within the scope of the base class.
%    
%\subsubsection[Inheriting constructor templates would be ill formed for a local class]{Inheriting constructor templates would be ill formed for a local class}
%
%A class declared within a function is a \romeogloss{local class}. Local classes have many restrictions, one of which is that they cannot declare member templates. If we inherit constructors from a base class with constructor templates, even \lstinline!private! ones, the implicit declaration of a constructor template to forward arguments to the base-class constructor would be \romeogloss{ill formed}:
%\begin{emcppslisting}
%struct Base
%{
%    template <typename T>
%    Base(T);
%};
%
%void f()
%{
%    class Local : Base
%    {
%        using Base::Base;  // Error, Prior to C++17 fixes, we cannot redeclare
%                           // the constructor template in local class.
%    };
%}
%\end{emcppslisting}
%    C++17 resolves this by directly exposing the base-class constructors, rather than defining new constructors to forward arguments.
%    
%\subsubsection[SFINAE evaluation context with default function arguments]{SFINAE evaluation context with default function arguments}
%
%Constructors that employ \romeogloss{SFINAE} tricks in default function arguments perform \mbox{\romeogloss{SFINAE}} checks in the wrong context and therefore inherit ill-formed constructors. No such issues occur when these \romeogloss{SFINAE} tricks are performed on default template arguments instead. As an example, consider a class template \lstinline!Wrap! that has a template constructor with a \romeogloss{SFINAE} constraint:
%\begin{emcppslisting}
%struct S { };
%
%template <typename T>
%struct Wrap
%{
%    template <typename U>
%    Wrap(U, typename std::enable_if<
%        std::is_constructible<T, U>::value>::type* = nullptr)
%        // This constructor is enabled only if (ù{\codeincomments{T}}ù) is constructible from (ù{\codeincomments{U}}ù).
%    {
%        std::cout << "SFINAE ctor\n";
%    }
%
%    Wrap(S)
%    {
%        std::cout << "S ctor\n";
%    }
%};
%\end{emcppslisting}
%    If we derive from \lstinline!Wrap! and inherit its constructors, we would expect the \romeogloss{SFINAE} constraint to behave exactly as in the base class, i.e., the template constructor overload would be silently discarded if \lstinline!std::is_constructible<T,!~\lstinline!U>::value! evaluates to \lstinline!false!:
%\begin{emcppslisting}
%template <typename T>
%struct Derived : Wrap<T>
%{
%    using Wrap<T>::Wrap;
%};
%\end{emcppslisting}
%    However, prior to C++17’s retroactive fixes, \romeogloss{SFINAE} was triggered only for \lstinline!Wrap!, not for \lstinline!Derived!:
%\begin{emcppslisting}
%void f()
%{
%    S s;
%    Wrap<int> w(s);     // prints (ù{\codeincomments{"S ctor"}}ù)
%    Derived<int> d(s);  // error prior to fixes; prints (ù{\codeincomments{"S ctor"}}ù) afterward
%}
%\end{emcppslisting}
%    
%\subsubsection[Suppression of constructors in the presence of default arguments]{Suppression of constructors in the presence of default arguments}
%
%A constructor having one or more default arguments in the derived class does not suppress any corresponding constructors matching only the nondefaulted arguments in the base class, leading to ambiguities:
%\begin{emcppslisting}
%struct B          // base class
%{
%    B(int, int);  // value constructor with two (required) (ù{\codeincomments{int}}ù) parameters
%};
%
%struct D : B
%{
%    using B::B;
%    D(int, int, int = 0);  // doesn't suppress (ù{\codeincomments{D(int, int)}}ù) from (ù{\codeincomments{B(int, int)}}ù)
%};
%\end{emcppslisting}
%    In the code example above, the original defective behavior was that there would be two overloaded constructors in \lstinline!D!; attempting to construct a \lstinline!D! from two integers became ambiguous. In the corrected behavior, the inheriting \lstinline!D(int,!~\lstinline!int)! from the base-class constructor \lstinline!B(int,!~\lstinline!int)!, whose domain is fully subsumed by the derived class’s explicitly specified constructor \lstinline!D(int,!~\lstinline!int,!~\lstinline!int!~\lstinline!=!~\lstinline!0)!, is suppressed.
%    
%\subsubsection[Suprising behavior with unary constructor templates]{Suprising behavior with unary constructor templates}
%
%Because inherited constructors are redeclarations within the derived class and expect to forward properly to the corresponding base-class constructors, constructor templates may do very surprising things. In particular, a gregarious, templated constructor can appear to cause inheritance of a base-class copy constructor. Consider the following class with a constructor template:
%\begin{emcppslisting}
%struct A
%{
%    A() = default;
%    A(const A&) { std::cout << "copy\n"; }
%
%    template <typename T>
%    A(T) { std::cout << "convert\n"; }
%};
%\end{emcppslisting}
%    This simple class can convert from any type and prints those of its constructors that were called. Now consider we want to make a \romeogloss{strong \lstinline!typedef!} for \lstinline!A!:
%\begin{emcppslisting}
%struct B : A
%{
%    using A::A;  // inherited base class (ù{\codeincomments{A}}ù)'s constructors
%};
%\end{emcppslisting}
%    The problem is that because \lstinline!A! can convert from anything, when \lstinline!B! inherits \lstinline!A!'s constructor template, \lstinline!B! can then use the inherited constructor to construct an instance of \lstinline!B! from \lstinline!A!. Perhaps more surprising, because the definition of the inherited constructor in \lstinline!B! is to initialize the \lstinline!A! subobject with its parameters, the nontemplate inherited constructor will be chosen as the best match, not the templated, converting constructor!\cprotect\footnote{Note that if the template constructor for \lstinline!A! were a \emph{copy} or \emph{move} constructor for \lstinline!A!, then it would be excluded from being an inherited constructor and this odd behavior would be avoided. The by-value parameter of this constructor is also why \lstinline!"copy"! is output twice in this example.
%\begin{emcppslisting}[style=footcode]
%A x;
%B y = x;  // Surprise! This compiles, and it prints (ù{\codeincomments{"copy"}}ù) twice!
%\end{emcppslisting}
%    }
%


\newpage
%\section[{\tt initializer\_list}]{List Initialization: {\SecCode std::initializer\_list<T>}}%\label{}


\emcppsFeature{
    short={\lstinline!initializer_list!},
    tocshort={\TOCCode initializer\_list},
    long={List Initialization: {\SecCode std::initializer\_list<T>}},
    toclong={List Initialization: \lstinline!std::initializer_list<T>!},
    rhshort={\RHCode initializer\_list},
}{initlist}
\setcounter{table}{0}
\setcounter{footnote}{0}
\setcounter{lstlisting}{0}
%\section[{\tt initializer\_list}]{List Initialization: {\SecCode std::initializer\_list<T>}}%\label{}

placeholder

\newpage
%\section[Lambdas]{Unnamed Local Function Objects (Closures)}\label{lambda}
% 20 March 2021, received and cleaned up; updated commands; sent to CE, Steven B for revisions, and Josh for code review. 
% 21 March 2021 JMB - cleaned up code compilation
% 23 March 2021, added noindent commands for each para under code, so this matches other features 
% 25 March 2021, added one code correx from Josh



\emcppsFeature{
    short={Lambdas},
    long={Unnamed Local Function Objects (Closures)},
}{lambda}\label{lambda-expressions}
\setcounter{table}{0}
\setcounter{footnote}{0}
\setcounter{lstlisting}{0}
%\section[Lambdas]{Unnamed Local Function Objects (Closures)}\label{lambda}
%\subsection[Lambda Expressions]{Lambda Expressions}\label{lambda-expressions}

Lambda expressions provide a means of defining function objects at the
point where they are needed, enabling a powerful and convenient way to
specify callbacks or local functions.

\subsection[Description]{Description}\label{description}

Generic, object-oriented, and functional programming paradigms all place
great importance on the ability of a programmer to specify a
\emph{callback} that is passed as an argument to a function. For
example, the Standard-Library algorithm, \lstinline!std::sort!, accepts a
callback argument specifying the sort order:

\begin{emcppslisting}
#include <algorithm>   // (ù{\codeincomments{std::sort}}ù)
#include <functional>  // (ù{\codeincomments{std::greater}}ù)
#include <vector>      // (ù{\codeincomments{std::vector}}ù)

template <typename T>
void sortAscending(std::vector<T>& v)
{
    std::sort(v.begin(), v.end(), std::greater<T>());
}
\end{emcppslisting}
    

\noindent The function object, \lstinline!std::greater<T>()!, is callable with two
arguments of type \lstinline!T! and returns \lstinline!true! if the first is
greater than the second and \lstinline!false! otherwise. The Standard
Library provides a small number of similar functor types, but, for more
complicated cases, the programmer must write a functor themselves. If a
container holds a sequence of \lstinline!Employee! records, for example, we
might want to sort the container by name or by salary:

\begin{emcppslisting}[emcppsbatch=e1]
#include <string>  // (ù{\codeincomments{std::string}}ù)
#include <vector>  // (ù{\codeincomments{std::vector}}ù)

struct Employee
{
    std::string name;
    long        salary;  // in whole dollars
};

void sortByName(std::vector<Employee>& employees);
void sortBySalary(std::vector<Employee>& employees);
\end{emcppslisting}
    

\noindent The implementation of \lstinline!sortByName! can delegate the sorting task
to the standard algorithm, \lstinline!std::sort!. However, because
\lstinline!Employee! does not supply \lstinline!operator<! and to achieve the
correct sorting criteria, we will need to supply \lstinline!std::sort! with
a callback that compares the names of two \lstinline!Employee! objects. We
implement this callback as a pointer to a simple function that we pass
to \lstinline!std::sort!:

\begin{emcppslisting}[emcppsbatch=e1]
#include <algorithm>  // (ù{\codeincomments{std::sort}}ù)

bool nameLt(const Employee& e1, const Employee& e2)
    // returns true if (ù{\codeincomments{e1.name}}ù) is less than (ù{\codeincomments{e2.name}}ù)
{
    return e1.name < e2.name;
}

void sortByName(std::vector<Employee>& employees)
{
    std::sort(employees.begin(), employees.end(), &nameLt);
}
\end{emcppslisting}
    

\noindent The \lstinline!sortBySalary! function can similarly delegate to
\lstinline!std::sort!. For illustrative purposes, we will use a
\emcppsgloss{function object} (a.k.a., \emcppsgloss{functor}) rather than a
function pointer as the callback to compare the salaries of two
\lstinline!Employee! objects. Every \emcppsgloss{functor class} must provide a
\emcppsgloss{call operator} (i.e., \lstinline!operator()!), which, in this case,
compares the salary fields of its arguments:

\begin{emcppslisting}[emcppsbatch=e1]
struct SalaryLt
{
    // Functor whose call operator compares two (ù{\codeincomments{Employee}}ù) objects and returns
    // (ù{\codeincomments{true}}ù) if the first has a lower salary than the second, (ù{\codeincomments{false}}ù) otherwise.

    bool operator()(const Employee& e1, const Employee& e2) const
    {
        return e1.salary < e2.salary;
    }
};

void sortBySalary(std::vector<Employee>& employees)
{
    std::sort(employees.begin(), employees.end(), SalaryLt());
}
\end{emcppslisting}
    

\noindent Although it is a bit more verbose, a call through the \emcppsgloss{function
object} is easier for the compiler to analyze and automatically inline
within \lstinline!std::sort! than is a call through the function pointer.
\emcppsgloss[function object]{Function objects} are also more flexible because they can carry
state, as we'll see shortly. The sorting example illustrates how small
bits of a function's logic must be factored out into special-purpose
auxiliary functions and/or functor classes that are often not re-usable.
It is possible, for example, that the \lstinline!nameLt! function and
\lstinline!SalaryLt! class are not used anywhere else in the program.

When callbacks are tuned to the specific context in which they are used,
they become both more complicated and less re-usable. Let's say, for
example, that we wish to count the number of employees whose salary is
above the average for the collection. Using Standard Library algorithms,
this task seems trivial: (1) sum all of the salaries using
\lstinline!std::accumulate!, (2) calculate the average salary by dividing
this sum by the total number of employees, and (3) count the number of
employees with above-average salaries using \lstinline!std::count_if!.
Unfortunately, both \lstinline!std::accumulate! and \lstinline!std::count_if!
require callbacks to return the salary for an \lstinline!Employee! and to
supply the criterion for counting, respectively. The callback for
\lstinline!std::accumulate! must take two parameters --- the current
running sum and an element from the sequence being summed --- and must
return the new running sum:

\begin{emcppslisting}[emcppsbatch=e1]
struct SalaryAccumulator
{
    long operator()(long currSum, const Employee& e) const
        // returns the sum of (ù{\codeincomments{currSum}}ù) and the (ù{\codeincomments{salary}}ù) field of (ù{\codeincomments{e}}ù)
    {
        return currSum + e.salary;
    }
};
\end{emcppslisting}
    

\noindent The callback for \lstinline!std::count_if! is a \emcppsgloss{predicate} (i.e.,
an expression that yields a Boolean result in response to a yes-or-no
question) that takes a single argument and returns \lstinline!true! if an
element of that value should be counted and \lstinline!false! otherwise. In
this case, we are concerned with \lstinline!Employee! object's having
salaries above the average. Our \emcppsgloss{predicate functor} must,
therefore, carry around that average so that it can compare it to the
salary of the employee that is presented as an argument:

\begin{emcppslisting}[emcppsbatch=e1]
class SalaryIsGreater  // function object constructed with a reference salary
{
    const long d_referenceSalary;

public:
    explicit SalaryIsGreater(long rs) : d_referenceSalary(rs) { }
        // construct with a reference salary, (ù{\codeincomments{rs}}ù)

    bool operator()(const Employee& e) const
        // return (ù{\codeincomments{true}}ù) if the salary for (ù{\codeincomments{Employee}}ù) (ù{\codeincomments{e}}ù) is greater than the
        // reference salary specified on construction, (ù{\codeincomments{false}}ù) otherwise
    {
        return e.salary > d_referenceSalary;
    }
};
\end{emcppslisting}
    

\noindent Note that, unlike our previous \emcppsgloss[functor class]{functor classes},
\lstinline!SalaryIsGreater! has a member variable, i.e., it has
\emph{state}. This member variable must be initialized, necessitating a
constructor. Its \emcppsgloss{call operator} compares its input argument
against this member variable to compute the \emcppsgloss{predicate} value.

With these two \emcppsgloss[functor class]{functor classes} defined, we can finally
implement the simple three-step algorithm for determining the number of
employees with salaries greater than the average:

\begin{emcppslisting}[emcppsbatch=e1]
#include <algorithm>  // (ù{\codeincomments{std::count\_if}}ù)
#include <numeric>    // (ù{\codeincomments{std::accumulate}}ù)

std::size_t numAboveAverageSalaries(const std::vector<Employee>& employees)
{
    const long sum = std::accumulate(employees.begin(), employees.end(), 0L,
                                     SalaryAccumulator());

    const long average = sum / employees.size();
    return std::count_if(employees.begin(), employees.end(),
                         SalaryIsGreater(average));
}
\end{emcppslisting}
    

\noindent The first statement creates an object of the \lstinline!SalaryAccumulator!
class and passes that object to the \lstinline!std::accumulate! algorithm
to produce the sum of all of the salaries. The second statement divides
the sum by the size of the \lstinline!employees! collection to compute the
average salary. The third statement creates an object of the
\lstinline!SalaryIsGreater! class and passes it to the
\lstinline!std::count_if! algorithm to compute the result. Note that the
local variable, \lstinline!average!, is used to initialize the reference
value in the \lstinline!SalaryIsGreater! object.

We now turn our attention to a syntax that allows us to rewrite these
examples much more simply and compactly. Returning to the sorting
example, the rewrite has the name-comparison and salary-comparison
operations expressed in-place, within the call to \lstinline!std::sort!:

\begin{emcppslisting}[emcppsbatch=e1]
void sortByName2(std::vector<Employee>& employees)
{
    std::sort(employees.begin(), employees.end(),
              [](const Employee &e1, const Employee &e2)
              {
                  return e1.name < e2.name;
              });
}

void sortBySalary2(std::vector<Employee>& employees)
{
    std::sort(employees.begin(), employees.end(),
              [](const Employee &e1, const Employee &e2)
              {
                  return e1.salary < e2.salary;
              });
}
\end{emcppslisting}
    

\noindent In each case, the third argument to \lstinline!std::sort! --- beginning
with \lstinline![]! and ending with the nearest closing \lstinline!}! --- is called a
\emcppsgloss[lambda expressions]{lambda expression}. Intuitively, for this case, one can think of
a \emcppsgloss[lambda expressions]{lambda expression} as an \emph{operation} that can be invoked
as a callback by the algorithm. The example shows a function-style
parameter list --- matching that expected by the \lstinline!std::sort!
algorithm --- and a function-like body that computes the needed
predicate. Using \emcppsgloss{lambda expressions}, a developer can express a
desired operation directly at the point of use rather than defining it
elsewhere in the program.

The compactness and simplicity afforded by the use of \emcppsgloss{lambda
expressions} is even more evident when we rewrite the average-salaries
example:

\begin{emcppslisting}[emcppsbatch=e1]
std::size_t numAboveAverageSalaries2(const std::vector<Employee>& employees)
{
    const long sum = std::accumulate(employees.begin(), employees.end(), 0L,
                                     [](long currSum, const Employee& e)
                                     {
                                         return currSum + e.salary;
                                     });

    const long average = sum / employees.size();
    return std::count_if(employees.begin(), employees.end(),
                         [average](const Employee& e)
                         {
                             return e.salary > average;
                         });
}
\end{emcppslisting}
    

\noindent The first \emcppsgloss[lambda expressions]{lambda expression}, above, specifies the operation for
adding another salary to a running sum. The second \emcppsgloss[lambda expressions]{lambda
expression} returns true if the \lstinline!Employee! argument, \lstinline!e!,
has a salary that is larger than \lstinline!average!, which is a local
variable \emph{captured} by the \emcppsgloss[lambda expressions]{lambda expression}. A
\emcppsgloss{lambda capture} is a set of local variables that are usable
within the body of the \emcppsgloss[lambda expressions]{lambda expression}, effectively making
the \emcppsgloss[lambda expressions]{lambda expression} an extension of the immediate
environment. We will look at the syntax and semantics of \emcppsgloss[lambda capture]{lambda
captures} in more detail in {Parts of a Lambda Expression}.

Note that the lambda expressions replaced a significant portion of code
that was previously expressed as separate functions or \emcppsgloss{functor}
classes. The fact that some of that code reduction is in the form of
documentation (comments) increases the appeal of \emcppsgloss{lambda
expressions} to a surprising degree. Creating a named entity such as a
function or class imposes on the developer the responsibility to give
that entity a meaningful name and sufficient documentation for a future
human reader to understand its \emph{abstract} purpose, outside the
context of its use, even for one-off, non-reusable entities. Conversely,
when an entity is defined right at the point of use, it might not need a
name at all, and it is often self-documenting, as in both the sorting
and average-salaries examples above. Both the original creation and
maintenance of the code is simplified.

\subsubsection[Parts of a lambda expression]{Parts of a lambda expression}\label{parts-of-a-lambda-expression}

A \emcppsgloss[lambda expressions]{lambda expression} has a number of parts and subparts, many of
which are optional. For exposition purposes, let's look at a sample
\emcppsgloss[lambda expressions]{lambda expression} that contains all of the parts:

\begin{emcppslisting}[emcppsignore={not code}]
       introducer
     ______^______
    /             \
    [&, cap2, cap3](T1 arg1, T2 arg2) mutable noexcept -> R { /* ... */ }
     \_____ _____/ \__________________ ___________________/ \____ ____/
           V                          V                          V
        capture                   declarator                    body
\end{emcppslisting}
    

\noindent Evaluating a lambda expression creates a temporary \emcppsgloss{closure}
object of an unnamed type called the \emcppsgloss{closure type}. Each part of
a \emcppsgloss[lambda expressions]{lambda expression} is described in detail in the subsections
below.

\subsubsection[Closures]{Closures}\label{closures}

A \emcppsgloss[lambda expressions]{lambda expression} looks a lot like an unnamed function
definition, and it is often convenient to think of it that way, but a
lambda expression is actually more complex than that. First and
foremost, a \emcppsgloss[lambda expressions]{lambda expression}, as the name implies, is an
\emph{expression} rather than a \emph{definition}. The result of
evaluating a \emcppsgloss[lambda expressions]{lambda expression} is a special \emcppsgloss{function
object} called a \emcppsgloss{closure}{\cprotect\footnote{The terms
\emph{lambda} and \emph{closure} are borrowed from \emph{Lambda
Calculus}, a computational system developed by Alonzo Church in the
1930s. Many computer languages have features inspired by Lambda
Calculus, although most (including C++) take some liberties with the
terminology. See \cite{rojas15} and
 \cite{barendregt84}.}}; it is not until the
\emcppsgloss{closure} is \emph{invoked} --- which can happen immediately but
usually occurs later (e.g., as a callback) --- that the actual body of
the \emcppsgloss[lambda expressions]{lambda expression} gets evaluated.

Evaluating a \emcppsgloss[lambda expressions]{lambda expression} creates a temporary
\emcppsgloss{closure object} of an unnamed type called the \emcppsgloss{closure
type}. The \emcppsgloss{closure type} encapsulates captured variables (see\linebreak[4] 
\featureref{\locationd}{lambda-capture-expressions}) and has a \emcppsgloss{call operator} that
executes the body of the \emcppsgloss[lambda expressions]{lambda expression}. Each \emcppsgloss[lambda expressions]{lambda
expression} has a unique \emcppsgloss{closure type}, even if it is identical
to another \emcppsgloss[lambda expressions]{lambda expression} in the program. If the
\emcppsgloss[lambda expressions]{lambda expression} appears within a template, the
\emcppsgloss{closure type} for each instantiation of that template is unique.
Note, however, that, although the \emcppsgloss{closure object} is an unnamed
temporary object, it can be saved in a named variable whose type can be
queried. \emcppsgloss[closure type]{Closure types} are copy constructible and move
constructible, but they have no other constructors and have deleted
assignment operators.{\cprotect\footnote{C++17 provides default
constructors for empty-capture lambdas. Empty-capture lambdas are
  assignable in C++20.}} Interestingly, it is possible to \emph{inherit}
from a \emcppsgloss{closure type}, provided the derived class constructs its
\emcppsgloss{closure type} base class using only the default or move
constructors. This ability to derive from a \emcppsgloss{closure type} is
convenient when implementing certain library features such as
\lstinline!std::bind!, which take advantage of the empty-base optimization:

\begin{emcppslisting}
#include <utility>  // (ù{\codeincomments{std::move}}ù)

template <typename Func>
int callFunc(const Func& f) { return f(); }

void f1()
{
    int   i   = 5;
    auto  c1  = [i]{ return 2 * i; };   // OK, deduced type for (ù{\codeincomments{c1}}ù)
    using C1t = decltype(c1);           // OK, named alias for unnamed type
    C1t   c1b = c1;                     // OK, copy of (ù{\codeincomments{c1}}ù)
    auto  c2  = [i]{ return 2 * i; };   // OK, identical lambda expression
    using C2t = decltype(c2);
    C1t   c2b = c2;                     // Error, different types, (ù{\codeincomments{C1t}}ù) & (ù{\codeincomments{C2t}}ù)
    using C3t = decltype([]{/* ... */});  // Error, lambda expr within (ù{\codeincomments{decltype}}ù)

    class C1Derived : public C1t        // OK, inherit from closure type
    {
        int d_auxValue;

    public:
        C1Derived(C1t c1, int aux) : C1t(std::move(c1)), d_auxValue(aux) { }
        int aux() const { return d_auxValue; }
    };

    int ret = callFunc([i]{ return 2 * i; });  // OK, deduced arg type, (ù{\codeincomments{Func}}ù)

    c1b = c1;  // Error, assignment of closures is not allowed.
}
\end{emcppslisting}
    

\noindent The types of \lstinline!c1! and \lstinline!c2!, above, are different, even
though they are token-for-token identical. As there is no way to
explicitly name a \emcppsgloss{closure type}, we use \lstinline!auto! (in the
case of \lstinline!c1! and \lstinline!c2! in \lstinline!f1!) or template-argument
deduction (in the case of \lstinline!f! in \lstinline!callFunc!) to create
variables directly from the \emcppsgloss[lambda expressions]{lambda expression}, and we use
\lstinline!decltype! to create aliases to the types of existing closure
variables (\lstinline!C1t! and \lstinline!C2t!). Note that using
\lstinline!decltype! directly on a \emcppsgloss[lambda expressions]{lambda expression} is ill
formed, as shown with \lstinline!C3t!, because there would be no way to
construct an object of the resulting unique type. The derived class,
\lstinline!C1Derived!, uses the type alias \lstinline!C1t! to refer to its
base class. Note that its constructor forwards its first argument to the
base-class move constructor.

There is no way to specify a \emcppsgloss{closure type} prior to creating an
actual \emcppsgloss{closure object} of that type. Consequently, there is no
way to declare \lstinline!callFunc! with a parameter of the actual
\emcppsgloss{closure type} that will be passed; hence, it is declared as a
template parameter. As a special case, however, if the \emcppsgloss{lambda
capture} is \emph{empty} (i.e., the \emcppsgloss[lambda expressions]{lambda expression} begins
with \lstinline![]!; see \featureref{\locationd}{lambda-capture-expressions}), then the
\emcppsgloss{closure} is implicitly convertible to an ordinary function
pointer having the same signature as its \emcppsgloss{call operator}:

\begin{emcppslisting}
char callFunc2(char (*f)(const char*)) { return f("x"); }  // not a template

char c = callFunc2([](const char* s) { return s ? s[0] : '\0'; });
    // OK, closure argument is converted to function-pointer parameter

char d = callFunc2([c](const char* s) { /* ... */ });
    // Error, lambda capture is not empty; no conversion to function pointer
\end{emcppslisting}
    

\noindent The \lstinline!callFunc2! function takes a callback in the form of a
pointer to function. Even though it is not a template, it can be called
with a lambda argument having the same parameter types, the same return
type, and an empty \emcppsgloss{lambda capture}; the \emcppsgloss{closure object}
is converted to an ordinary pointer to function. This conversion is
\emph{not} available in the second call to \lstinline!callFunc2! because
the \emcppsgloss{lambda capture} is not empty.

Conversion to function pointer is considered a user-defined conversion
operator and thus cannot be implicitly combined with other conversions
on the same expression. It can, however, be invoked \emph{explicitly},
as needed:

\begin{emcppslisting}
using Fp2 = int(*)(int);  // function-pointer type

struct FuncWrapper
{
    FuncWrapper(Fp2) { /* ... */ }  // implicit conversion from function-pointer
    // ...
};

int f2(FuncWrapper)  { /* ... */ return 0; }
int i2 = f2([](int x) { return x; });  // Error, two user-defined conversions
int i3 = f2(static_cast<Fp2>([](int x) { return x; }));  // OK, explicit cast
int i4 = f2(+[](int x) { return x; });  // OK, forced conversion
\end{emcppslisting}
    

\noindent The first call to \lstinline!f2! fails because it would require two
implicit user-defined conversions: one from the \emcppsgloss{closure type} to
the \lstinline!Fp2! function-pointer type and one from \lstinline!Fp2! to
\lstinline!FuncWrapper!. The second call succeeds because the first
conversion is made explicit with the \lstinline!static_cast!. The third
call is an interesting shortcut that takes advantage of the fact that
unary \lstinline!operator+! is defined as the identity transformation for
pointer types. Thus, the \emcppsgloss{closure}-to-pointer conversion is
invoked for the operand of \lstinline!operator+!, which returns the
unchanged pointer, which, in turn, is converted to \lstinline!FuncWrapper!;
the first and third steps of this sequence use only one user-defined
conversion each. The Standard Library \lstinline!std::function! class
template provides another way to pass a function object of unnamed type,
one that does not require the \emcppsgloss{lambda capture} to be empty; see
\intrarefsimple{use-cases-lambda}.

The compile-time and runtime phases of defining a \emcppsgloss{closure type}
and constructing a \emcppsgloss{closure object} from a single \emcppsgloss[lambda expressions]{lambda
expression} resembles the phases of calling a function template; what
looks like an ordinary function call is actually broken down into a
compile-time instantiation and a runtime call. The \emcppsgloss{closure type}
is deduced when a \emcppsgloss[lambda expressions]{lambda expression} is encountered during
compilation. When the control flow passes through the \emcppsgloss[lambda expressions]{lambda
expression} at run time, the \emcppsgloss{closure object} is
\emph{constructed} from the list of captured local variables. In the
\lstinline!numAboveAverageSalaries! example in the {Description} section,
the \lstinline!SalaryIsGreater! class can be thought of as a
\emcppsgloss{closure type} --- created by hand instead of by the compiler ---
whereas the call to \lstinline!SalaryIsGreater(average)! is analogous to
constructing the \emcppsgloss{closure object} at run time.

Finally, the purpose of a \emcppsgloss{closure} is to be invoked. It can be
invoked immediately by supplying arguments for each of its parameters:

\begin{emcppslisting}
#include <iostream>  // (ù{\codeincomments{std::cot}}ù)
void f3()
{
    [](const char* s) { std::cout << s; }("hello world\n");
        // equivalent to (ù{\codeincomments{std::cout << "hello world$\backslash$n";}}ù)
}
\end{emcppslisting}
    

\noindent The \emcppsgloss{closure object}, in this example, is invoked immediately and
then destroyed, making the above just a complicated way to say
\lstinline!std::cout!~\lstinline!<<!~\lstinline!"hello!~\lstinline!world\n";!.
More commonly, the \emcppsgloss[lambda expressions]{lambda expression} is used as a local
function for convenience and to avoid clutter:

\begin{emcppslisting}
#include <cmath>  // (ù{\codeincomments{std::sqrt}}ù)

double hypotenuse(double a, double b)
{
    auto sqr = [](double x) { return x * x; };
    return std::sqrt(sqr(a) + sqr(b));
}
\end{emcppslisting}
    

\noindent Note that there is no way to overload calls to \emcppsgloss[closure]{closures}:

\begin{emcppslisting}
auto sqr = [](int x) { return x * x; };     // OK, store **closure** in (ù{\codeincomments{sqr}}ù)
auto sqr = [](double x) { return x * x; };  // Error, redefinition of (ù{\codeincomments{sqr}}ù)
\end{emcppslisting}
    

\noindent The most common use of a \emcppsgloss[lambda expressions]{lambda expression}, however, is as a
callback to a function template, e.g., as a functor argument to an
algorithm from the Standard Library:

\begin{emcppslisting}
#include <algorithm>  // (ù{\codeincomments{std::partition}}ù)

template <typename FwdIt>
FwdIt oddEvenPartition(FwdIt first, FwdIt last)
{
    using value_type = decltype(*first);
    return std::partition(first, last, [](value_type v) { return v % 2 != 0; });
}
\end{emcppslisting}
    

\noindent The \lstinline!oddEvenPartition! function template moves odd values to the
start of the sequence and even values to the back. The \emcppsgloss{closure
object} is invoked repeatedly within the \lstinline!std::partition!
algorithm.

\subsubsection[Lambda capture and lambda introducer]{Lambda capture and lambda introducer}\label{lambda-capture-and-lambda-introducer}

The purpose of the \emcppsgloss{lambda capture} is to make certain local
variables from the environment available to be used (or, more precisely,
\emcppsgloss{ODR-used}, which means that they are used in a
potentially-evaluated context) within the \emcppsgloss{lambda body}. Each
local variable can be \emcppsgloss{captured by copy} or \emcppsgloss{captured by
reference}. Orthogonally, each variable can be \emcppsgloss{explicitly
captured} or \emcppsgloss{implicitly captured}. When a \emcppsgloss[lambda expressions]{lambda
expression} appears within a non-static member function, the
\lstinline!this! pointer can be captured as a special case. We'll examine
each of these aspects of \emcppsgloss{lambda capture} in turn. An extension
to \emcppsgloss{lambda capture} in C++14 is discussed in \intrarefsimple{c++14-init-capture}.

Syntactically, the \emcppsgloss{lambda capture} consists of an optional
\emcppsgloss{capture default} followed by a comma-separated list of zero or
more identifiers (or the keyword \lstinline!this!), which are
\emcppsgloss{explicitly captured}. The capture default can be one of
\lstinline!=! or \lstinline!&! for \emcppsgloss[captured by copy]{capture by copy} or
\emcppsgloss[captured by reference]{capture by reference}, respectively. If there is a
\emcppsgloss{capture default}, then \lstinline!this! and any local variables in
scope that are \emcppsgloss{ODR-used} within the \emcppsgloss{lambda body} and not
\emcppsgloss{explicitly captured} will be \emcppsgloss{implicitly captured}.

\begin{emcppslisting}
void f1()
{
    int a = 0, b = 1, c = 2;
    auto c1 = [a, b]{ return a + b; };
        // (ù{\codeincomments{a}}ù) and (ù{\codeincomments{b}}ù) are explicitly captured.
    auto c2 = [&]{ return a + b; };
        // (ù{\codeincomments{a}}ù) and (ù{\codeincomments{b}}ù) are implicitly captured.
    auto c3 = [&, b]{ return a + b; };
        // (ù{\codeincomments{a}}ù) is implicitly captured and (ù{\codeincomments{b}}ù) is explicitly captured.
    auto c4 = [a]{ return a + b; }
        // Error, (ù{\codeincomments{b}}ù) is ODR-used but not captured.
}
\end{emcppslisting}
    

\noindent The Standard defines the \emcppsgloss{lambda introducer} as the
\emcppsgloss{lambda capture} together with its surrounding \lstinline![! and
\lstinline!]!. If the \emcppsgloss{lambda introducer} is an empty pair of
brackets, no variables will be captured.

\begin{emcppslisting}
auto c1 = []{ /* ... */ };  // Empty **lambda capture**
\end{emcppslisting}
    

\noindent The \emcppsgloss{lambda capture} enables access to portions of the local
stack frame. As such, only variables with \emph{automatic storage
duration} --- i.e., non-static local variables --- can be captured, as
we'll see in detail later in this section and in \emcppsgloss{lambda body}.
An \emcppsgloss{explicitly captured} variable whose name is immediately
preceded by an \lstinline!&! symbol in the \emcppsgloss{lambda capture} is
\emcppsgloss{captured by reference}; without the \lstinline!&!, it is
\emcppsgloss{captured by copy}. If the \emcppsgloss{capture default} is
\lstinline!&!, then all \emcppsgloss{implicitly captured} variables are
\emcppsgloss{captured by reference}. Otherwise, if the \emcppsgloss{capture
default} is \lstinline!=!, all \emcppsgloss{implicitly captured} variables are
\emcppsgloss{captured by copy}:

\begin{emcppslisting}
void f2a()
{
    int a = 0, b = 1;
    auto c1 = [&a]{ /* ... */ return a; };  // (ù{\codeincomments{a}}ù) captured by reference
    auto c2 = [a] { /* ... */ return a; };  // (ù{\codeincomments{a}}ù) captured by copy
    auto c3 = [a, &b] { return a + b; };
        // (ù{\codeincomments{a}}ù) is explicitly captured by copy and (ù{\codeincomments{b}}ù) is explicitly
        // captured by reference.
    auto c4 = [=]{ return a + b; };
        // (ù{\codeincomments{a}}ù) and (ù{\codeincomments{b}}ù) are implicitly captured by copy.
    auto c5 = [&]{ return &a; };
        // (ù{\codeincomments{a}}ù) is implicitly captured by reference.
    auto c6 = [&, b]{ return a * b; };
        // (ù{\codeincomments{a}}ù) is implicitly captured by reference and (ù{\codeincomments{b}}ù) is explicitly
        // captured by copy.
    auto c7 = [=, &b]{ return a * b; };
        // (ù{\codeincomments{a}}ù) is implicitly captured by copy and (ù{\codeincomments{b}}ù) is explicitly
        // captured by reference.
    auto c8 = [a]{ return a * b; };
        // Error, (ù{\codeincomments{a}}ù) is explicitly captured by copy, but (ù{\codeincomments{b}}ù) is not captured.
    auto c9 = [this]{ /* ... */ };  // Error, no (ù{\codeincomments{this}}ù) in nonmember function
}

class Class1a
{
public:
    void mf1()
    {
        auto c12 = [this]{ return this; };  // Explicitly capture (ù{\codeincomments{this}}ù).
        auto c13 = [=]   { return this; };  // Implicitly capture (ù{\codeincomments{this}}ù).
    }
};
\end{emcppslisting}
    

\noindent Redundant captures are not allowed; the same name (or \lstinline!this!)
cannot appear twice in the \emcppsgloss{lambda capture}. Moreover, if the
\emcppsgloss{capture default} is \lstinline!&!, then none of the
\emcppsgloss{explicitly captured} variables may be \emcppsgloss{captured by
reference}, and if the \emcppsgloss{capture default} is \lstinline!=!, then any
\emcppsgloss{explicitly captured} entities can be neither \emcppsgloss{explicitly
copied} variables nor \lstinline!this!{\cprotect\footnote{C++20 removed
the prohibition on explicit capture of \lstinline!this! with an
  \lstinline!=! capture default.}}:

\begin{emcppslisting}
class Class1b
{
public:
    void mf1()
    {
        int a = 0;
        auto c1 = [a, &a]{ /* ... */ }; // Error, (ù{\codeincomments{a}}ù) is captured twice.
        auto c2 = [=, a]{ /* ... */ };
            // Error, explicit capture of (ù{\codeincomments{a}}ù) by copy is redundant.
        auto c3 = [&,&a]{ /* ... */ };
            // Error, explicit capture of (ù{\codeincomments{a}}ù) by reference is redundant.
        auto c4 = [=, this]{ return this; };
            // Error, explicit capture of (ù{\codeincomments{this}}ù) with (ù{\codeincomments{=}}ù) capture default
    }
};
\end{emcppslisting}
    

\noindent We'll use the term \emph{primary variable} to refer to the block-scope
local variable outside of the \emcppsgloss[lambda expressions]{lambda expression} and
\emph{captured variable} to refer to the variable of the same name as
viewed from within the \emcppsgloss{lambda body}. For every object that is
\emcppsgloss{captured by copy}, the \emcppsgloss{lambda closure} will contain a
member variable having the same name and type, after stripping any
reference qualifier (except reference-to-function); this member variable
is initialized from the primary variable by direct initialization and is
destroyed when the \emcppsgloss{closure object} is destroyed. Any
\emcppsgloss[ODR-used]{ODR-use} of that name within the \emcppsgloss{lambda body} will refer
to the \emcppsgloss{closure}'s member variable. Thus, for an entity that is
\emcppsgloss{captured by copy}, the primary and captured variables refer to
distinct objects with distinct lifetimes. By default, the \emcppsgloss{call
operator} is \lstinline!const!, providing read-only access to members of
the \emcppsgloss{closure object} (i.e., captured variables that are
\emcppsgloss{captured by copy}); mutable (non\lstinline!const!) \emcppsgloss[call operator]{call
operators} are discussed in \intrarefsimple{lambda-declarator}:

\begin{emcppshiddenlisting}[emcppsbatch=e2]
#include <cassert>  // standard C (ù{\codeincomments{assert}}ù) macro
\end{emcppshiddenlisting}
\begin{emcppslisting}[emcppsbatch=e2]

void f3()
{
    int a = 5;
    auto c1 = [a]        // (ù{\codeincomments{a}}ù) is captured by copy.
    {
        return a;        // return value of copy of (ù{\codeincomments{a}}ù)
    };
    a = 10;              // Modify (ù{\codeincomments{a}}ù) after it was captured by (ù{\codeincomments{c1}}ù).
    assert(5 == c1());   // OK, (ù{\codeincomments{a}}ù) within (ù{\codeincomments{c1}}ù) had value from before the change.

    int& b = a;
    auto c2 = [b]        // (ù{\codeincomments{b}}ù) is (ù{\codeincomments{int}}ù) (not (ù{\codeincomments{int\&}}ù)) **captured by copy**.
    {
        return b;        // return value of copy of (ù{\codeincomments{b}}ù)
    };
    b = 15;              // Modify (ù{\codeincomments{a}}ù) through reference (ù{\codeincomments{b}}ù).
    assert(10 == c2());  // OK, (ù{\codeincomments{b}}ù) within (ù{\codeincomments{c2}}ù) is a copy, not a reference.

    auto c3 = [a]
    {
        ++a;             // Error, (ù{\codeincomments{a}}ù) is (ù{\codeincomments{const}}ù) within the lambda body.
    };
}
\end{emcppslisting}
    

\noindent In the example above, the \emcppsgloss[lambda expressions]{lambda expression} is evaluated to
produce a \emcppsgloss{closure object}, \lstinline!c1!, that captures a
\emph{copy} of \lstinline!a!. Even when the primary \lstinline!a! is
subsequently modified, the captured \lstinline!a! in \lstinline!c1! remains
unchanged. When \lstinline!c1! is invoked, the \emcppsgloss{lambda body} returns
the \emph{copy}, which still has the value 5. The same applies to
\lstinline!c2!, but note that the copy of \lstinline!b! is \emph{not a
reference} even though \lstinline!b! is a reference. Thus, the copy of
\lstinline!b! in \lstinline!c2! is the value of \lstinline!a! that \lstinline!b!
referred to at the time that \lstinline!c2! was created.

When a variable is \emcppsgloss{captured by reference}, the captured variable
is simply an alias to the primary variable; no copies are made. It is,
therefore, possible to modify the primary variable and/or take its
address within the \emcppsgloss{lambda body}:

\begin{emcppshiddenlisting}[emcppsbatch=e3]
#include <cassert>  // standard C (ù{\codeincomments{assert}}ù) macro
\end{emcppshiddenlisting}
\begin{emcppslisting}[emcppsbatch=e3]
void f4()
{
    int a = 5;
    auto c1 = [&a]      // (ù{\codeincomments{a}}ù) is **captured by reference**.
    {
        a = 10;         // Modify (ù{\codeincomments{a}}ù) through the captured variable.
        return &a;      // return address of captured (ù{\codeincomments{a}}ù)
    };
    assert(c1() == &a); // OK, primary and captured (ù{\codeincomments{a}}ù) have the same address.
    assert(10 == a);    // OK, primary (ù{\codeincomments{a}}ù) is now 10.

    int& b = a;
    auto c2 = [&b]      // (ù{\codeincomments{b}}ù) is **captured by reference**
    {
        return &b;      // return address of captured (ù{\codeincomments{b}}ù)
    };
    assert(c2() == &b); // OK, primary and captured (ù{\codeincomments{b}}ù) have the same address.
    assert(c2() == &a); // OK, captured (ù{\codeincomments{b}}ù) is an alias for (ù{\codeincomments{a}}ù).
}
\end{emcppslisting}
    

\noindent In contrast to the \lstinline!f3! example, the \lstinline!c1! \emcppsgloss{closure
object} above does \emph{not} hold a copy of the captured variable,
\lstinline!a!, though the compiler may choose to define a member of type
\lstinline!int&! that refers to \lstinline!a!. Within the \emcppsgloss{lambda
body}, modifying \lstinline!a! modifies the primary variable, and taking
its address returns the address of the primary variable, i.e., the
captured variable is an alias for the primary variable. With respect to
variables that are \emcppsgloss{captured by reference}, the \emcppsgloss{lambda
body} behaves very much as though it were part of the surrounding block.
The lifetime of a variable that is \emcppsgloss{captured by reference} is the
same as that of the primary variable (since they are the same). In
particular, if a copy of the \emcppsgloss{closure object} outlives the
primary variable, then the captured variable becomes a \emph{dangling
reference}; see \intrarefsimple{potential-pitfalls-lambda}.

If \lstinline!this! appears in the \emcppsgloss{lambda capture}, then (1) the
current \lstinline!this! pointer is \emcppsgloss{captured by copy} and (2)
within the \emcppsgloss{lambda body}, member variables accessible through
\lstinline!this! can be used without prefixing them with \lstinline!this->!,
as though the \emcppsgloss{lambda body} were an extension of the surrounding
member function. The \emcppsgloss{lambda body} cannot refer to the
\emcppsgloss{closure} directly; the captured \lstinline!this! does not point to
the \emcppsgloss{closure} but to the \lstinline!*this! object of the function
within which it is defined:

\begin{emcppshiddenlisting}[emcppsbatch=e4]
#include <cassert>  // standard C (ù{\codeincomments{assert}}ù) macro
\end{emcppshiddenlisting}
\begin{emcppslisting}[emcppsbatch=e4]
struct Class1
{
    int d_value;

public:
    // ...
    void mf() const
    {
        auto c1 = []{ return *this; };        // Error, (ù{\codeincomments{this}}ù) is not captured.
        auto c2 = []{ return d_value; };      // Error, (ù{\codeincomments{this}}ù) is not captured.
        auto c3 = [d_value]{ /* ... */; };    // Error, cannot capture member
        auto c4 = [this]{ return this; };     // OK, returns (ù{\codeincomments{this}}ù)
        auto c5 = [this]{ return d_value; };  // OK, returns (ù{\codeincomments{this->d\_value}}ù)
        assert(this == c4());                 // OK, captured (ù{\codeincomments{this}}ù) is (ù{\codeincomments{Class1}}ù).
    }
};
\end{emcppslisting}
    

\noindent Note that \lstinline!c4! returns \lstinline!this!, which is the address of the
\lstinline!Class1! for which \lstinline!mf! was called. This is one way in
which the \emcppsgloss{closure type} is different from a named
\emcppsgloss{functor type} --- there is no way for an object of
\emcppsgloss{closure type} to refer to itself directly. Because the
\emcppsgloss{closure type} is unnamed and because it does not supply its own
\lstinline!this! pointer, it is difficult (but not impossible) to create a
\emph{recursive} \emcppsgloss[lambda expressions]{lambda expression}; see \textbf{Usage examples [there is no section with this name] }.

If \lstinline!this! is captured (implicitly or explicitly), the
\emcppsgloss{lambda body} will behave much like an extension of the member
function in which the \emcppsgloss[lambda expressions]{lambda expression} appears, with direct
access to the class's members:

\begin{emcppshiddenlisting}[emcppsbatch=e5]
#include <algorithm>  // (ù{\codeincomments{std::count\_if}}ù)
#include <vector>     // (ù{\codeincomments{std::vector}}ù)
\end{emcppshiddenlisting}
\begin{emcppslisting}[emcppsbatch=e5]

class Class2
{
    int d_value;

public:
    std::size_t mem(const std::vector<int>& v) const
    {
        auto f = []{ return d_value; };
            // Error, (ù{\codeincomments{this}}ù) not captured; can't see (ù{\codeincomments{d\_value}}ù).
        return std::count_if(v.begin(), v.end(),
                             [this](int element){ return element < d_value; });
            // OK, uses (ù{\codeincomments{this->d\_value}}ù).
    }
};
\end{emcppslisting}
    

\noindent Note that capturing \lstinline!this! does not copy the class object that it
points to; the original \lstinline!this! and the captured \lstinline!this!
will point to the same object:

\begin{emcppshiddenlisting}[emcppsbatch=e6]
#include <cassert>  // standard C (ù{\codeincomments{assert}}ù) macro
\end{emcppshiddenlisting}
\begin{emcppslisting}[emcppsbatch=e6]
class Class3
{
    int d_value;

public:
    void mf()
    {
        auto c1 = [this]{ ++d_value; };  // increment (ù{\codeincomments{this->d\_value}}ù)
        d_value = 1;
        c1();
        assert(2 == d_value);            // change to (ù{\codeincomments{d\_value}}ù) is visible
    }
};
\end{emcppslisting}
    

\noindent Here, we captured \lstinline!this! in \lstinline!c1! but then proceeded to
modify the object pointed to by \lstinline!this! within the \emcppsgloss{lambda
body}.{\cprotect\footnote{In C++17, it is possible to capture
\lstinline!*this!, which results in the entire class object being copied,
  not just the \lstinline!this! pointer.}}

A \emcppsgloss[lambda expressions]{lambda expression} can occur wherever other expressions can
occur, including within other \emcppsgloss{lambda expressions}. The set of
entities that can be captured in a valid \emcppsgloss[lambda expressions]{lambda expression}
depends on the surrounding scope. A \emcppsgloss[lambda expressions]{lambda expression} that does
not occur immediately within \emcppsgloss{block scope} cannot have a
\emcppsgloss{lambda capture}:

\begin{emcppslisting}
namespace ns1
{
    int v = 10;
    int w = [v]{ /* ... */ return 0; }();
        // Error, capture in global/namespace scope

    void f4(int a = [v]{ return v; }());  // Error, capture in default argument
}
\end{emcppslisting}
    

\noindent When a \emcppsgloss[lambda expressions]{lambda expression} occurs in block scope, it can capture
any local variables with \emph{automatic} (i.e., non-static) storage
duration in its \emcppsgloss{reaching scope}. The Standard defines the
\emcppsgloss{reaching scope} of the \emcppsgloss[lambda expressions]{lambda expression} as the set of
enclosing scopes up to and including the innermost enclosing function
and its parameters. Static variables can be used without capturing them;
see \intrarefsimple{lambda-body}:

\begin{emcppslisting}
void f5(const int& a)
{
    int b = 2 * a;
    if (a)
    {
        int c;
        // ...
    }
    else
    {
               int d = 4 * a;
        static int e = 10;

        auto c1 = [a]{ /* ... */ };    // OK, capture argument (ù{\codeincomments{a}}ù) from (ù{\codeincomments{f5}}ù).
        auto c2 = [=]{ return b; };  // OK, implicitly capture local (ù{\codeincomments{b}}ù).
        auto c3 = [&c]{ /* ... */ };   // Error, (ù{\codeincomments{c}}ù) is not in **reaching scope**.
        auto c4 = [&]{ d += 2; };    // OK, implicitly capture local (ù{\codeincomments{d}}ù).
        auto c5 = [e]{ /* ... */ };    // Error, (ù{\codeincomments{e}}ù) has static duration.
    }

    struct LocalClass
    {
        void mf()
        {
            auto c6 = [b]{ /* ... */ };  // Error, (ù{\codeincomments{b}}ù) not in **reaching scope**
        }
    };
}
\end{emcppslisting}
    

\noindent The \emcppsgloss{reaching scope} of the \emcppsgloss{lambda expressions} for
\lstinline!c1! through \lstinline!c5!, above, includes the local variable
\lstinline!d! in the \lstinline!else! block, \lstinline!b! in the surrounding
function block, and \lstinline!a! from \lstinline!f5!'s arguments. The local
variable, \lstinline!c!, is not in their \emcppsgloss{reaching scope} and cannot
be captured. Although \lstinline!e! \emph{is} in their \emcppsgloss{reaching
scope}, it cannot be captured because it does not have automatic storage
duration. Finally, the \emcppsgloss[lambda expressions]{lambda expression} for \lstinline!c6! is
within a member function of a local class. Its \emcppsgloss{reaching scope}
ends with the innermost function, \lstinline!LocalClass::mf!, and does not
include the surrounding block that includes \lstinline!a! and \lstinline!b!.

Only when the innermost enclosing function is a non-static class member
function can \lstinline!this! be captured:

\begin{emcppslisting}
void f5()
{
    auto c1 = [this]{ /* ... */ };  // Error, (ù{\codeincomments{f5}}ù) is not a member function.
}

class Class3
{
    static void mf1()
    {
        auto c2 = [this]{ /* ... */ };  // Error, (ù{\codeincomments{mf1}}ù) is static.
    }

    void mf2()
    {
        auto c3 = [this]{ /* ... */ };  // OK, (ù{\codeincomments{mf2}}ù) is non-static member function

        struct LocalClass
        {
            static void mf3()
            {
                auto c4 = [this]{ /* ... */ };
                    // Error, innermost function, (ù{\codeincomments{mf3}}ù), is static
            }
        };
    }
};
\end{emcppslisting}
    

\noindent When a \emcppsgloss[lambda expressions]{lambda expression} is enclosed within another
\emcppsgloss[lambda expressions]{lambda expression}, then the \emcppsgloss{reaching scope} includes
all of the intervening \emcppsgloss[lambda body]{lambda bodies}. Any variable captured
(implicitly or explicitly) by the inner \emcppsgloss[lambda expressions]{lambda expression} must
be either defined or captured by the enclosing lambda expression:

\begin{emcppslisting}
void f6()
{
    int a, b, c;
    const char* d;
    auto c1 = [&a]                   // capture (ù{\codeincomments{a}}ù) from function block
    {
        int d;                        // local definition of (ù{\codeincomments{d}}ù) hides outer def
        auto c2 = [&a]{ /* ... */ };  // OK, (ù{\codeincomments{a}}ù) is captured in enclosing lambda
        auto c3 = [d]{ /* ... */ };   // OK, capture (ù{\codeincomments{int d}}ù) from enclosing
        auto c4 = [&]{ return d; };   // OK,    "       "     "      "
        auto c5 = [b]{ /* ... */ };   // Error, (ù{\codeincomments{b}}ù) is not captured in enclosing
    };
    auto c6 = [=]
    {
        auto c7 = [&]{ return b; };
            // OK, ODR-use of (ù{\codeincomments{b}}ù) causes implicit capture in (ù{\codeincomments{c7}}ù) and (ù{\codeincomments{c6}}ù).
        auto c8 = [&d]{ return &d; };
            // (ù{\codeincomments{d}}ù) is captured by copy in (ù{\codeincomments{c6}}ù); (ù{\codeincomments{c8}}ù) returns address of copy
    };
}
\end{emcppslisting}
    

\noindent Note that there are two variables named \lstinline!d!: one at function
scope and one within the body of the first \emcppsgloss[lambda expressions]{lambda expression}.
Following normal rules for unqualified name lookup, the inner
\emcppsgloss{lambda expressions} used to initialize \lstinline!c3! and
\lstinline!c4! capture the \emph{inner} \lstinline!d! (of type \lstinline!int!),
not the \emph{outer} \lstinline!d! (of type \lstinline!const!~\lstinline!char*!).
Because it is not captured, primary variable \lstinline!b! is
\emph{visible} but not \emph{usable} --- an important distinction that
we'll discuss in \intrarefsimple{lambda-body} --- within the body of \lstinline!c1! and
cannot, therefore, be captured by \lstinline!c5!.

The \emcppsgloss{lambda body} for \lstinline!c7! \emcppsgloss[ODR-used]{ODR-uses} \lstinline!b!,
thus causing it to be \emcppsgloss{implicitly captured}. This capture by
\lstinline!c7! constitutes an \emcppsgloss[ODR-used]{ODR-use} of \lstinline!b! within the
enclosing \emcppsgloss[lambda expressions]{lambda expression}, \lstinline!c6!, in turn causing
\lstinline!b! to be \emcppsgloss{implicitly captured} by \lstinline!c6!. In this
way, a single \emcppsgloss[ODR-used]{ODR-use} can trigger a \emph{chain} of
\emcppsgloss[implicitly captured]{implicit captures} from an enclosed \emcppsgloss[lambda expressions]{lambda expression}
through its enclosing \emcppsgloss{lambda expressions}. Critically, when a
variable is \emcppsgloss{captured by copy} in one \emcppsgloss[lambda expressions]{lambda expression},
any enclosed \emcppsgloss{lambda expressions} that capture the same name will
capture the \emph{copy}, not the primary variable, as we see in the
\emcppsgloss[lambda expressions]{lambda expression} for \lstinline!c8!.

Note that, when a variable is named in a \emcppsgloss{lambda capture}, it
isn't automatically \emph{captured}. A variable is not captured unless
it is \emcppsgloss{ODR-used} within the \emcppsgloss[lambda expressions]{lambda expression}:

\begin{emcppslisting}
void f7()
{
    int       a = 0;
    int const b = 2;
    auto c1 = [a]{ return 2 * a; };      // OK, (ù{\codeincomments{a}}ù) is explicitly captured.
    auto c2 = [a]{ return 0; };          // Warning, no uses of (ù{\codeincomments{a}}ù)
    auto c3 = [a]{ return sizeof(a); };  // Warning, (ù{\codeincomments{sizeof(a)}}ù) isn't an ODR-use
    auto c4 = [b]{ return b; };          // Warning, value of (ù{\codeincomments{b}}ù) isn't an ODR-use
    auto c5 = [&b]{ return &b; };        // OK, address of (ù{\codeincomments{b}}ù) is an ODR-use
}
\end{emcppslisting}
    

\noindent In the above example, the \emcppsgloss{lambda body} for \lstinline!c1! contains
an \emcppsgloss[ODR-used]{ODR-use} of \lstinline!a! and thus captures \lstinline!a!.
Conversely, \lstinline!c2! does \emph{not} capture \lstinline!a! because the
\emcppsgloss{lambda body} for \lstinline!c2! does not contain an
\emcppsgloss[ODR-used]{ODR-use} of \lstinline!a!; most compilers will issue a warning
diagnostic about the superfluous presence of \lstinline!a! in the
\emcppsgloss{lambda capture} for \lstinline!c2!. The case of \lstinline!c3! is a
bit more subtle. Although \lstinline!a! is \emph{used} within \lstinline!c3!,
it is not \emcppsgloss{ODR-used} --- i.e., it is not used in a
potentially-evaluated context because \lstinline!sizeof! does not evaluate
its argument --- so the warning is the same as for \lstinline!c2!; see
\intrarefsimple{potential-pitfalls-lambda}. The last two cases are a bit more subtle: The use of the
\emph{value} of a \lstinline!const! variable (in \lstinline!c4!) is \emph{not}
an \emcppsgloss[ODR-used]{ODR-use} of that variable, whereas the use its \emph{address}
(in \lstinline!c5!) \emph{is} an \emcppsgloss[ODR-used]{ODR-use}.

Finally, a \emcppsgloss{lambda capture} within a \emcppsgloss{variadic function
template} (see \featureref{\locationc}{variadictemplate}) may contain a
\emcppsgloss[parameter pack]{parameter pack expansion}:

\begin{emcppslisting}
#include <utility>  // (ù{\codeincomments{std::forward}}ù)

template <typename... ArgTypes>
int f8(const char* s, ArgTypes&&... args);

template <typename... ArgTypes>
int f9(ArgTypes&&... args)
{
    const char* s = "Introduction";
    auto c1 = [=]{ return f8(s, args...); };  // OK, (ù{\codeincomments{args...}}ù) captured by copy
    auto c2 = [s,&args...]{ return f8(s, std::forward<ArgTypes>(args)...); };
        // OK, explicit capture of (ù{\codeincomments{args...}}ù) by reference
}
\end{emcppslisting}
    

\noindent In the example above, the variadic arguments to \lstinline!f9! are
\emcppsgloss{implicitly captured} using \emcppsgloss[captured by copy]{capture by copy} in the first
\emcppsgloss[lambda expressions]{lambda expression}. This means that, regardless of the
\emcppsgloss{value category} (\romeovalue{rvalue}, \romeovalue{lvalue}, and so on) of the
original arguments, the captured variables are all \romeovalue{lvalue} members
of the resulting \emph{closure}. Conversely, the second \emcppsgloss[lambda expressions]{lambda
expression} captures the set of arguments using \emcppsgloss[captured by reference]{capture by
reference}, again resulting in captured variables that are
\romeovalue{lvalues}. The \lstinline!ArgTypes! \emcppsgloss[parameter pack]{parameter pack expansion}
designates a list of \emph{types}, not \emph{variables}, and does not,
therefore, need to be captured to be used within the \emcppsgloss[lambda expressions]{lambda
expression}, nor would it be valid to attempt to capture it. Because
\lstinline!ArgTypes! is specified using a forwarding reference
(\lstinline!&&!, see \featureref{\locationc}{forwardingref}), the Standard
Library function, \lstinline!std::forward!, can be used to cast the
captured variables to the \emcppsgloss{value category} of their corresponding
arguments.

\subsubsection[C++14 init capture]{C++14 init capture}\label{c++14-init-capture}

\textbf{TODO VR: this subsection shouldn't be here, we have a specific C++14
feature `lambdacapture' for this stuff.}

Each item in a C++11 \emcppsgloss{lambda capture} is either a \emcppsgloss{capture
default} or a \emcppsgloss{simple capture} consisting of the name of a local
variable, either by itself (for \emcppsgloss[captured by copy]{capture by copy}) or preceded by
an \lstinline!&! (for \emcppsgloss[captured by reference]{capture by reference}). C++14 introduces
another possibility, an \emcppsgloss{init capture}, consisting of a variable
name and an initializer, which creates a new captured variable
initialized to an arbitrary expression. The initializer can be either
preceded by an \lstinline!=! token or can be a braced initialization (see
\featureref{\locationc}{bracedinit}). The newly defined variable does
not necessarily share the name or type of a primary local variable:

\begin{emcppslisting}
void f1(int a, bool bits[])
{
    double g(int);

    int b;
    auto c1 = [i{5}]{ /* ... */ };        // Define (ù{\codeincomments{int i = 5}}ù).
    auto c2 = [b=bits[a]]{ /* ... */ };   // Define (ù{\codeincomments{bool b}}ù) = copy of (ù{\codeincomments{bits[a]}}ù).
    auto c3 = [&r=a, b]{ /* ... */ };     // (ù{\codeincomments{int\& r}}ù) = reference to (ù{\codeincomments{a}}ù). Copy (ù{\codeincomments{b}}ù).
    auto c4 = [&, p=bits+1]{ *p = !a; };  // (ù{\codeincomments{bool* p}}ù) = pointer to (ù{\codeincomments{bits[1]}}ù).
    auto c5 = [x=g(a)]{ return x; };      // (ù{\codeincomments{double x = g(a)}}ù).
}
\end{emcppslisting}
    

\noindent In the \emcppsgloss[lambda expressions]{lambda expression} for \lstinline!c1!, above, the
\emcppsgloss{init capture} defines a new variable, \lstinline!i!, initialized
with the value 5. Within the \emcppsgloss{lambda body}, this \lstinline!i! is
indistinguishable from any other variable captured by the
\emcppsgloss{closure}, but, outside of the \emcppsgloss{lambda body}, it differs
from a \emcppsgloss{simple capture} in that it does not capture a local
variable. The \emcppsgloss{lambda capture} for \lstinline!c2! similarly defines
a new variable, \lstinline!b!, initialized from an expression involving
\lstinline!bits! and \lstinline!a!, but does not capture either of them. Once
captured, \lstinline!bits[a]! can change without affecting the value of
\lstinline!b!. An \emcppsgloss{init capture} can also define a variable of
\emph{reference} type, as shown in the \emcppsgloss{init capture} for
\lstinline!c3!. The \emcppsgloss{lambda capture} for \lstinline!c3! also shows an
\emcppsgloss{init capture} mixed in with \emcppsgloss[simple capture]{simple captures}. Note,
however, that the \emcppsgloss{capture default}, if any, has no effect on the
\emcppsgloss{init capture}, as shown in the \emcppsgloss[lambda expressions]{lambda expression} for
\lstinline!c4!. The \emcppsgloss{init capture} for \lstinline!c5! shows an
arbitrary expression being captured in this the return value of a
function call.

The variable defined in an \emcppsgloss{init capture} is defined as if by the
declaration:

\begin{emcppslisting}[emcppsignore={pseudocode}]
(ù{\codeincomments{auto}}ù) *init-capture* (ù{\codeincomments{;}}ù)
\end{emcppslisting}
    

\noindent The hypothetical variable created by such a definition is unique and
separate from any similarly named variable in the environment. This
uniqueness allows the same name to appear on both the left and right of
the \lstinline!=! symbol:

\begin{emcppslisting}
#include <string>   // (ù{\codeincomments{std::string}}ù)
#include <utility>  // (ù{\codeincomments{std::move}}ù)

void f2(std::string s)
{
   float a = 1.2;
   auto c1 = [&s=s]{ /* ... */ };            // effectively capture by reference
   auto c2 = [a=static_cast<double>(a)]{ /* ... */ };
   auto c3 = [s=std::move(s)]{ /* ... */ };  // effectively capture by move
   std::string s2 = s;                     // BAD IDEA, (ù{\codeincomments{s}}ù) is moved from
}
\end{emcppslisting}
    

\noindent In all of the \emcppsgloss{lambda expressions} above, an \emcppsgloss{init
capture} defines a variable whose name shadows one in the environment.
In fact, the first use, in \lstinline!c1!, defines a captured variable,
\lstinline!s!, that is a reference to the argument variable by the same
name, yielding a behavior within the \emcppsgloss{lambda body} that is
essentially indistinguishable from normal \emcppsgloss[captured by reference]{capture by reference}.
The second \emcppsgloss[lambda expressions]{lambda expression} copies the local variable,
\lstinline!a!, into a captured variable such that, within the
\emcppsgloss{lambda body}, the variable \lstinline!a! is a \lstinline!double!.
Finally, the \emcppsgloss{lambda capture} for \lstinline!c3! creates a captured
variable, \lstinline!s!, initialized from the argument \lstinline!s! using the
\lstinline!std::string! \emph{move constructor} (see \featureref{\locationc}{Rvalue-References}), simulating something that many
considered to be missing in C++11: capture by move. Note, however, that
after constructing the \lstinline!c3! \emcppsgloss{closure}, the local variable,
\lstinline!s!, is in a \emph{moved-from} state and thus has an unspecified
value, even if \lstinline!c3! is never invoked.

In C++14, an \emcppsgloss{init capture} cannot be a \emcppsgloss[parameter pack]{parameter pack
expansion}.{\cprotect\footnote{A syntax for \emcppsgloss[parameter pack]{parameter pack
expansion} of \emcppsgloss[init capture]{init captures} was introduced in C++20:

\begin{emcppslisting}[style=footcode]
template <typename T> class X { /* ... */ };
template <typename T> X<T> f3(T&&);

template <typename... ArgTypes>
void f4(ArgTypes&&... args)
{
    auto c1 = [...x=f3(std::forward<ArgTypes>(args))]{ /* ... */ };
        // Error in C++14.  OK in C++20
}
\end{emcppslisting}
      }}

\subsubsection[Lambda declarator]{Lambda declarator}\label{lambda-declarator}

The \emcppsgloss{lambda declarator} looks a lot like a function declaration
and is effectively the declaration of the \emcppsgloss{closure type}'s
\emcppsgloss{call operator}. The \emcppsgloss{lambda declarator} comprises the
\emcppsgloss{call operator}'s parameter list, mutability, exception
specification, and return type:

\begin{emcppslisting}[emcppsignore={not code}]
[/* ... */](T1 arg1, T2 arg2) mutable   noexcept(/* ... */) -> RetType { /* ... */ }
           \________ _______/ \__ __/   \_______ _______/   \_____ ___/
                    V            V              V                 V
            parameter list   mutability   exception spec     return type
           \___________________________ _____________________________/
                                       V
                               lambda declarator
\end{emcppslisting}
    

\noindent Although the \emcppsgloss{lambda declarator} looks very similar to a function
declaration, we cannot forward declare any part of a \emcppsgloss[lambda expressions]{lambda
expression}; we can only define it.

The entire \emcppsgloss{lambda declarator} is optional. However, if
\emph{any} part is present, the parameter list must be present (even if
it declares no parameters):

\begin{emcppslisting}
auto c1 = [](int x) noexcept { /* ... */ };  // OK, param list and exception spec
auto c2 = []() -> int { return 5; };         // OK, with ret type empty param list
auto c3 = [] -> double { /* ... */ };        // Error, ret type with no param list
\end{emcppslisting}
    

\noindent The parameter list for a \emcppsgloss[lambda expressions]{lambda expression} is the same as a
parameter list for a function declaration, with minor modifications.

\begin{enumerate}
\item{A parameter is not permitted to have the same name as an explicitly-captured variable:

\begin{emcppslisting}
void f1()
{
    int a;
    auto c1 = [a](short* a){ /* ... */ };  // Error, parameter shadows captured (ù{\codeincomments{a}}ù).
    auto c2 = [](short* a){ /* ... */ };   // OK, parameter hides local (ù{\codeincomments{a}}ù).
    auto c3 = [=](short* a){ /* ... */ };  // OK, local (ù{\codeincomments{a}}ù) is not captured.
}
\end{emcppslisting} 

\noindent In the definition of \lstinline!c1!, the \emcppsgloss[lambda expressions]{lambda expression}
\emcppsgloss[explicitly captured]{explicitly captures} \lstinline!a!, then improperly tries to
declare a parameter by the same name. When \lstinline!a! is not captured,
as in the \emcppsgloss[lambda expressions]{lambda expression} for \lstinline!c2!, having a parameter
named \lstinline!a! does not pose a problem; within the \emcppsgloss{lambda
body}, the declaration of \lstinline!a! in the parameter list will prevent
name lookup from finding the declaration in the enclosing scope. The
situation with \lstinline!c3! is essentially the same as for \lstinline!c2!;
because name lookup finds \lstinline!a! in the parameter list rather than
the enclosing scope, it does not attempt to capture it.}

\item{In C++11, none of the parameters may have default arguments. This restriction does not apply to C++14 and after:

\begin{emcppslisting}
auto c4 = [](int x, int y = 0){ /* ... */ };  // Error in C++11.  OK in C++14.
\end{emcppslisting}
}

\item{If the type of any of the parameters contains the keyword \lstinline!auto!, then the \emcppsgloss[lambda expressions]{lambda expression} becomes a \emcppsgloss{generic lambda}; see \featureref{\locationd}{genericlambda}. Everything in this chapter applies to \emcppsgloss{generic lambda} as well as regular \emcppsgloss{lambda expressions}, so it is recommended that you read the rest of this chapter before moving on to the \emcppsgloss[generic lambda]{generic lambdas} feature.}
\end{enumerate}

Note that, unlike the \emcppsgloss{lambda capture}, the parameter list is
usually dictated by the \emph{client} of a \emcppsgloss[lambda expressions]{lambda expression}
rather than by its author. Moreover, the \emcppsgloss{lambda capture} is
evaluated only once when the \emcppsgloss{closure} is created, whereas the
parameter list is bound to actual arguments each time the \emcppsgloss{call
operator} is invoked:

\begin{emcppslisting}[emcppserrorlines={22,24}]
#include <iostream>  // (ù{\codeincomments{std::cout}}ù)
#include <vector>    // (ù{\codeincomments{std::vector}}ù)

template <typename InputIter, typename Func>
void applyToEveryOtherElement(InputIter start, InputIter last, Func f) 
    // For elements in the range (ù{\codeincomments{[start, last)}}ù), invoke (ù{\codeincomments{f}}ù) on the first
    // element, skip the second element, etc., alternating between calling (ù{\codeincomments{f}}ù)
    // and skipping elements.
{
    while (start != last)
    {
        f(*start++);                     // Process one element.
        if (start != last) { start++; }  // Skip one element.
    }
}

void f2(const std::vector<float>& vec)
{
    std::size_t size = vec.size();
    applyToEveryOtherElement(vec.begin(), vec.end(),
                [](float x){ /* ... */ });         // OK, one (ù{\codeincomments{float}}ù) parameter
    applyToEveryOtherElement(vec.begin(), vec.end(),
                [](float x, int y){ /* ... */ });  // Error, too many parameters
    applyToEveryOtherElement(vec.begin(), vec.end(),
                [size](){ /* ... */ });            // Error, too few parameters
    applyToEveryOtherElement(vec.begin(), vec.end(),
                [size](double x){ /* ... */ });    // OK, convertible from (ù{\codeincomments{float}}ù)
}
\end{emcppslisting}
    

\noindent In the above definition of \lstinline!applyToEveryOtherElement!, the
callable argument, \lstinline!f!, is applied to half of the elements of the
input range and has the \emcppsgloss{closure type} resulting from a
\emcppsgloss[lambda expressions]{lambda expression}. In this example, therefore, each
instantiation of \lstinline!applyToEveryOtherElement! calls its \lstinline!f!
parameter multiple times with a single argument of type \lstinline!float!.
In the first call to \lstinline!applyToEveryOtherElement!, the
\emcppsgloss{lambda closure} has a parameter list consisting of a single
parameter of type \lstinline!float! and is thus compatible with the
expected signature of \lstinline!f!. The second and third calls to
\lstinline!applyToEveryOtherElement! supply a \emcppsgloss{lambda closure} with
too many and too few parameters, respectively, resulting in a
compilation error at the point where \lstinline!f! is called. The last call
to \lstinline!applyToEveryOtherElement! supplies a \emcppsgloss{lambda closure}
that takes a single argument of type \lstinline!double!. Since an argument
of type \lstinline!float! is convertible to \lstinline!double!, this
\emcppsgloss[lambda expressions]{lambda expression} is also a valid argument to
\lstinline!applyToEveryOtherElement!. Note that the presence or absence of
a \emcppsgloss{lambda closure} makes no difference to the validity of the
\emcppsgloss[lambda expressions]{lambda expression} from the point of view of the client
function.

The \lstinline!mutable! keyword, if present, indicates that the
\emcppsgloss{call operator} for the \emcppsgloss{closure type} should \emph{not}
be \lstinline!const!. Recall that a normal class member function can modify
the class's members if not declared \lstinline!const!. The \emcppsgloss{call
operator} is just an ordinary member function in this regard:

\begin{emcppslisting}
class Class1
{
    int d_value;

public:
    // ...
    void operator()(int v)       { d_value = v; }  // OK, object is mutable
    void operator()(int v) const { d_value = v; }  // Error, object is (ù{\codeincomments{const}}ù)
};
\end{emcppslisting}
    

\noindent The \lstinline!const! version of \lstinline!operator()! cannot modify the
member variables of the \lstinline!Class1! object, whereas the undecorated
one can. The \emcppsgloss{call operator} for a \emcppsgloss{closure type} has the
inverse default \lstinline!const!ness: the \emcppsgloss{call operator} is
implicitly \lstinline!const! \emph{unless} the \emcppsgloss{lambda declarator}
is decorated with \lstinline!mutable!. In practice, this rule means that
member variables of the \emcppsgloss{closure object}, i.e., variables that
were \emcppsgloss{captured by copy}, are \lstinline!const! by default and cannot
be modified within the \emcppsgloss{lambda body} unless the \lstinline!mutable!
keyword is present:

% Note: this was const int, which made the capture variable const even though c2 was mutable
% We should bring this up as a pitfall/annoyance?
\begin{emcppshiddenlisting}[emcppsbatch=e7]
#include <cassert>  // standard C (ù{\codeincomments{assert}}ù) macro
\end{emcppshiddenlisting}
\begin{emcppslisting}[emcppsbatch=e7]
void f2()
{
    int a = 5;
    auto c1 = [a]()         { return ++a; };  // Error, copy of (ù{\codeincomments{a}}ù) is (ù{\codeincomments{const}}ù)
    auto c2 = [a]() mutable { return ++a; };  // OK, increment *copy* of (ù{\codeincomments{a}}ù)
    assert(6 == c2());                        // OK, captured (ù{\codeincomments{a}}ù) incremented
    assert(5 == a);                           // OK, primary (ù{\codeincomments{a}}ù) not changed
    assert(7 == c2());                        // OK, captured (ù{\codeincomments{a}}ù) incremented
}
\end{emcppslisting}
    

\noindent The two \emcppsgloss{lambda expressions} are identical except for the
\lstinline!mutable! keyword. Both use \emcppsgloss[captured by copy]{capture by copy} to capture
local variable \lstinline!a!, and both try to increment \lstinline!a!, but
only the one decorated with \lstinline!mutable! can perform that
modification. When the \emcppsgloss{call operator} on \emcppsgloss{closure object}
\lstinline!c2! is invoked, it increments the \emph{captured copy} of
\lstinline!a!, leaving the primary \lstinline!a! untouched. If \lstinline!c2()!
is invoked again, it increments its copy of \lstinline!a! a second time.
Using \emcppsgloss[captured by reference]{capture by reference} allows the primary variable to be
changed within the \emcppsgloss{lambda body} regardless of the presence or
absence of the \lstinline!mutable! keyword. Similarly, member variables
accessed through \lstinline!this! are unaffected by the \lstinline!mutable!
keyword, but they \emph{are} affected by the \lstinline!const!ness of the
surrounding member function:

\begin{emcppshiddenlisting}[emcppsbatch=e8]
#include <cassert>  // standard C (ù{\codeincomments{assert}}ù) macro
\end{emcppshiddenlisting}
\begin{emcppslisting}[emcppsbatch=e8]
class Class2
{
    int d_value;
public:
    // ...
    void mf()
    {
        d_value = 1;
        int a = 0;
        auto c1 = [&a,this]{
            a = d_value;   // OK, (ù{\codeincomments{a}}ù) is a reference to a non-(ù{\codeincomments{const}}ù) object.
            d_value *= 2;  // OK, (ù{\codeincomments{this}}ù) points to a non-(ù{\codeincomments{const}}ù) object.
        };
        c1();
        assert(1 == a && 2 == d_value);  // values updated by (ù{\codeincomments{c1}}ù)
        c1();
        assert(2 == a && 4 == d_value);  // values updated by (ù{\codeincomments{c1}}ù) again
    }

    void cmf() const
    {
        int a = 0;
        auto c2 = [=]() mutable {
            ++a;        // OK, increment mutable **captured variable**
            ++d_value;  // Error, (ù{\codeincomments{*this}}ù) is const within (ù{\codeincomments{cmf}}ù).
        };
    }
};
\end{emcppslisting}
    

\noindent The \emcppsgloss[lambda expressions]{lambda expression} for \lstinline!c1! is not decorated with
\lstinline!mutable! yet both \lstinline!a! and \lstinline!d_value! can be
modified; the first because it was \emcppsgloss{captured by reference}, and
the second because it was accessed through \lstinline!this! within a
non\lstinline!const! member function, \lstinline!mf!. Conversely, the
\emcppsgloss[lambda expressions]{lambda expression} for \lstinline!c1! cannot modify
\lstinline!d_value! even though it is declared \lstinline!mutable! because
the captured \lstinline!this! pointer points to a \lstinline!const! object
within the surrounding member function, \lstinline!cmf!.

The \emcppsgloss{lambda declarator} may include an \emcppsgloss{exception
specification}, consisting of a \lstinline!throw! or \lstinline!noexcept!
clause, after the \lstinline!mutable! decoration, if any. The syntax and
meaning of the \emcppsgloss{exception specification} is identical to that of
a normal function; see \featureref{\locatione}{noexcept-specifier}.

The return type of the \emcppsgloss{call operator} can be determined either
by a \emcppsgloss{trailing return type} or by a \emcppsgloss{deduced return type}.
Every example we have seen up to this point has used a \emcppsgloss{deduced
return type} whereby the return type of the \emcppsgloss{closure}'s
\emcppsgloss{call operator} is deduced by the type of object returned by the
return statement(s). If there are no return statements in the
\emcppsgloss{lambda body} or if the return statements have no operands, then
the return type is \lstinline!void!. If there are multiple return
statements, they must agree with respect to the return type. The rules
for a \emcppsgloss{deduced return type} are the same for a \emcppsgloss[lambda expressions]{lambda
expression} as they are in C++14 for an ordinary function; see \featureref{\locationf}{Function-Return-Type-Deduction}:

\begin{emcppslisting}
void f3()
{
    auto c1 = [](int& i){ i = 0; };   // deduced return type is (ù{\codeincomments{void}}ù)
    auto c2 = []{ return "hello"; };  // deduced return type is (ù{\codeincomments{const char*}}ù)
    auto c3 = [](bool c)              // deduced return type is (ù{\codeincomments{int}}ù)
    {
        if (c) { return 5; }
        else   { return 6; }
    };
    auto c4 = [](bool c)
    {
        if (c) { return 5; }          // deduced return type is (ù{\codeincomments{int}}ù)
        else   { return 6.0; }        // Error, (ù{\codeincomments{double}}ù) does not match (ù{\codeincomments{int}}ù).
    };
}
\end{emcppslisting}
    

\noindent All four of the above \emcppsgloss{lambda expressions} have a \emcppsgloss{deduced
return type}. The first one deduces a return type of \lstinline!void!
because the \emcppsgloss{lambda body} has no \lstinline!return! statements. The
next one deduces a return type of \lstinline!const!~\lstinline!char*! because
the string literal, \lstinline!"hello"!, decays to a
\lstinline!const!~\lstinline!char*! in that context. The third one deduces a
return type of \lstinline!int! because all of the \lstinline!return!
statements return values of type \lstinline!int!. The last one fails to
compile because the two branches return values of different
types.{\cprotect\footnote{The original C++11 Standard did not allow a
\emcppsgloss{deduced return type} for a \emcppsgloss{lambda body} containing
anything other than a single \lstinline!return! statement. This
restriction was lifted by a defect report and is no longer part of
C++11. Compiler versions that predate ratification of this defect
report might reject \emcppsgloss{lambda expressions} having multiple
  statements and a \emcppsgloss{deduced return type}.}}

If a \emcppsgloss{deduced return type} is impossible or undesirable (see
\featureref{\locationf}{Function-Return-Type-Deduction} for a description of why this feature
needs to be used with care), a \emcppsgloss{trailing return type} can be
specified (see \featureref{\locationa}{trailing-function-return-types}:

\begin{emcppshiddenlisting}[emcppsbatch=e9]
#include <utility>  // (ù{\codeincomments{std::pair}}ù)
\end{emcppshiddenlisting}
\begin{emcppslisting}[emcppsbatch=e9]
void f4()
{
    auto c1 = [](bool c) -> double {
        if (c) { return 5; }          // OK, (ù{\codeincomments{int}}ù) value converted to (ù{\codeincomments{double}}ù)
        else   { return 6.0; }        // OK, (ù{\codeincomments{double}}ù) return value
    };
    auto c2 = []() -> std::pair<int, int> {
        return { 5, 6 };              // OK, brace-initialize returned (ù{\codeincomments{pair}}ù)
    };
}
\end{emcppslisting}
    

\noindent In the first \emcppsgloss[lambda expressions]{lambda expression} above, we specify a
\emcppsgloss{trailing return type} of \lstinline!double!. The two branches of
the \lstinline!if! statement would return different types (\lstinline!int! and
\lstinline!double!), but, because the return type has been definitively
declared, the compiler converts the return values to the known return
type (\lstinline!double!). The second \emcppsgloss[lambda expressions]{lambda expression} returns a
value by brace initialization, which is insufficient for deducing a
return value. Again, the ambiguity is resolved by declaring the return
value explicitly. Note that, unlike ordinary functions, a \emcppsgloss[lambda expressions]{lambda
expression} cannot have a return type specified before the
\emcppsgloss{lambda introducer} or \emcppsgloss{lambda declarator}:

\begin{emcppslisting}
auto c5 = int [](int x){ return 0; };     // Error, return type misplaced
auto c6 = [] int (int x){ return 0; };    // Error, return type misplaced
auto c7 = [](int x) -> int{ return 0; };  // OK, trailing return type
\end{emcppslisting}
    

\noindent Attributes (see \featureref{\locationa}{attributes}) that appertain to the
\emph{type} of \emcppsgloss{call operator} can be inserted in the
\emcppsgloss{lambda declarator} just before the \emcppsgloss{trailing return
type}. If there is no \emcppsgloss{trailing return type}, the attributes can
be inserted before the open brace of the \emcppsgloss{lambda body}.
Unfortunately, these attributes do not apply to the \emcppsgloss{call
operator} itself, but to its type, ruling out some common attributes:

\begin{emcppslisting}[emcppsignore={invalid construct}]
#include <cstdlib>  // (ù{\codeincomments{std::abort}}ù)
auto c1 = []() noexcept [[noreturn]] {  // Error, (ù{\codeincomments{[[noreturn]]}}ù) on a type
    std::abort();
};
\end{emcppslisting}
    

\subsubsection[Lambda body]{Lambda body}\label{lambda-body}

Combined, the \emcppsgloss{lambda declarator} and the \emcppsgloss{lambda body}
make up the declaration and definition of an \lstinline!inline! class
member function that is the \emcppsgloss{call operator} for the
\emcppsgloss{closure type}. For the purposes of name lookup and the
interpretation of \lstinline!this!, the \emcppsgloss{lambda body} is considered
to be in the context where the \emcppsgloss[lambda expressions]{lambda expression} is evaluated
(independent of the context where the \emcppsgloss{closure}'s \emcppsgloss{call
operator} is invoked).

Critically, the set of entity names that can be used from within the
\emcppsgloss{lambda body} is not limited to captured local variables. Types,
functions, templates, constants, and so on --- just like any other
member function --- do not need to be captured and, in fact,
\emph{cannot} be captured in most cases. To illustrate, let's create a
number of entities in multiple scopes:

\begin{emcppslisting}[emcppsbatch=e10]
#include <iostream>  // (ù{\codeincomments{std::cout}}ù)

namespace ns1
{
    void f1() { std::cout << "ns1::f1" << '\n'; }
    struct Class1 { Class1() { std::cout << "ns1::Class1()" << '\n'; } };
    int g0 = 0;
}

namespace ns2
{
    void f1() { std::cout << "ns2::f1" << '\n'; }

    template <typename T>
    struct Class1 { Class1() { std::cout << "ns2::Class1()" << '\n'; } };

    int const g1 = 1;
    int       g2 = 2;

    class Class2
    {
        int        d_value;  // non-static member variable
        static int s_mem;    // static member variable

        void mf1() { std::cout << "Class2::mf1" << '\n'; }

        struct Nested { Nested() { std::cout << "Nested()" << '\n'; } };

        template <typename T>
        static void print(const T& v) { std::cout << v << '\n'; }

    public:
        explicit Class2(int v) : d_value(v) { }

        void mf2();
        void mf3();
        void mf4();
        void mf5();
    };

    int Class2::s_mem = 0;
}
\end{emcppslisting}
    

\noindent Namespace \lstinline!ns1! contains three global entities: function
\lstinline!f1!, class \lstinline!Class1!, and variable \lstinline!g0!. Namespace
\lstinline!ns2! contains global variables \lstinline!g1! and \lstinline!g2!,
function \lstinline!f1!, and classes \lstinline!Class1! and \lstinline!Class2!.
Within \lstinline!Class2!, we have non-static member variable
\lstinline!d_value!, non-static member function \lstinline!mf1!, static
member function template \lstinline!print!, and public member functions
\lstinline!mf2! through \lstinline!mf5!.

With these declarations, we first demonstrate the use of entities that
are not variables and are accessible within the scope of a
\emcppsgloss{lambda body}:

\begin{emcppslisting}[emcppsbatch=e10]
void ns2::Class2::mf2()
{
    using LocalType = const char*;

    auto c1 = []{
        // Access non-variables in scope
        f1();           // find global (ù{\codeincomments{ns2::f1}}ù) by unqualified name lookup.
        Class1<int> x1; // construct (ù{\codeincomments{ns2::Class1<int>}}ù) object.
        Nested      x2; // construct (ù{\codeincomments{ns2::Class2::Nested}}ù) object.
        print("print"); // call static member function (ù{\codeincomments{ns2::Class2::print}}ù).
        LocalType   x3; // declare object of local type.
        ns1::f1();      // find global (ù{\codeincomments{ns1::f1}}ù) func by qualified name lookup
        ns1::Class1 x4; // find global (ù{\codeincomments{ns1::Class1}}ù) type by qualified lookup
    };
}
\end{emcppslisting}
    

\noindent We can see that, within the \emcppsgloss{lambda body}, non-variables can be
accessed normally, using either unqualified name lookup or, if needed,
qualified name lookup. Unqualified name lookup will find global entities
within the namespace; types and static functions within the class; and
types declared within the enclosing function scope. Qualified name
lookup will find entities in other namespaces.

Variables with static storage duration can also be accessed directly,
without being captured:

\begin{emcppslisting}[emcppsbatch=e10]
void ns2::Class2::mf3()
{
    static int s1 = 3;
    auto c1 = []{
        // Access variables with static storage duration.
        print(g1);      // print global constant (ù{\codeincomments{ns2::g1}}ù)
        print(g2);      // print global variable (ù{\codeincomments{ns2::g2}}ù)
        print(ns1::g0); // print global variable (ù{\codeincomments{ns1::g0}}ù)
        print(s_mem);   // print static member variable (ù{\codeincomments{s\_mem}}ù)
        print(s1);      // print static local variable (ù{\codeincomments{s1}}ù)
    };
}
\end{emcppslisting}
    

\noindent Here we see global constants, global variables, static member variables,
and local static variables being used from the local scope.

Next, we look at uses of variables with \emph{automatic} storage
duration from the \emcppsgloss[lambda expressions]{lambda expression}'s surrounding block scope:

\begin{emcppslisting}[emcppsbatch=e10]
void ns2::Class2::mf4()
{
    int       a = 4;
    int       b = 5;
    int const k = 6;

    auto c1 = [](double b) {
        // Access local variables within the reaching scope.
        print(a);    // Error, (ù{\codeincomments{a}}ù) is ODR-used but not captured.
        print(b);    // OK, (ù{\codeincomments{b}}ù) argument hides (ù{\codeincomments{b}}ù) defined in (ù{\codeincomments{mf1}}ù).
        int kb = k;  // OK, (ù{\codeincomments{const}}ù) variable is not ODR-used.
        print(k);    // Error, (ù{\codeincomments{k}}ù) is ODR-used but not captured.
    };
}
\end{emcppslisting}
    

\noindent The attempt to print \lstinline!a! will fail because \lstinline!a! is a
non-static local variable within the surrounding scope but is not
captured. Printing \lstinline!b! is not a problem because the \lstinline!b!
parameter in the \emcppsgloss{lambda declarator} is local to the
\emcppsgloss{lambda body}; it hides the \lstinline!b! variable in the
surrounding block scope. A \lstinline!const! variable of fundamental type
is not \emcppsgloss{ODR-used} unless its \emph{identity} (i.e., address) is
needed. Hence, \lstinline!k! can be used even though it is not captured.
Conversely, the call to \lstinline!print(k)! \emph{is} an \emcppsgloss[ODR-used]{ODR-use}
of \lstinline!k! because \lstinline!print! takes its argument by
\emph{reference}, which requires taking the \emph{address} of
\lstinline!k!, which, in turn, makes it an \emcppsgloss[ODR-used]{ODR-use}. Since
\lstinline!k! is a local variable with automatic storage duration that was
not captured, \lstinline!print(k)! is ill formed.

Finally, we look at access to (non-static) members of the surrounding
class:

\begin{emcppslisting}[emcppsbatch=e10]
void ns2::Class2::mf5()
{
    auto c1 = []{
        // Access non-static members
        print(d_value); // Error, (ù{\codeincomments{this}}ù) is ODR-used but not captured.
        mf1();          // Error,    "   "        "       "   "     "
    };
    auto c2 = [=] {
        print(d_value); // OK, (ù{\codeincomments{this}}ù) is captured; print (ù{\codeincomments{this->d\_value}}ù).
        mf1();          // OK,   "    "     "    ; call (ù{\codeincomments{this->mf1()}}ù).
    };
}
\end{emcppslisting}
    

\noindent The above shows that member variables and functions can be accessed only
if \lstinline!this! is (implicitly or explicitly) captured by the
\emcppsgloss[lambda expressions]{lambda expression}, as is the case for \lstinline!c2! but not for
\lstinline!c1!.

\subsection[Use Cases]{Use Cases}\label{use-cases-lambda}

\subsubsection[Interface adaptation, partial application, and currying]{Interface adaptation, partial application, and currying}\label{interface-adaptation,-partial-application,-and-currying}

\emcppsgloss[lambda expressions]{Lambda expressions} can be used to adapt the set of arguments
provided by an algorithm to the parameters expected by another facility:

\begin{emcppslisting}
#include <algorithm>  // (ù{\codeincomments{std::count\_if}}ù)
#include <string>     // (ù{\codeincomments{std::string}}ù)
#include <vector>     // (ù{\codeincomments{std::vector}}ù)

extern "C" int f1(const char* s, std::size_t n);

void f2(const std::vector<std::string>& vec)
{
    std::size_t n = std::count_if(vec.begin(), vec.end(),
        [](const std::string& s){ return 0 != f1(s.data(), s.size()); });
    // ...
}
\end{emcppslisting}
    

\noindent Here we have a function, \lstinline!f1!, that takes a C string and length
and computes some predicate, returning 0 for false and nonzero for true.
We want to use this predicate with \lstinline!std::count_if! to count how
many strings in a specified vector match this predicate. The
\emcppsgloss[lambda expressions]{lambda expression} in \lstinline!f2! adapts \lstinline!f1! to the
needs of \lstinline!std::count_if! by converting a \lstinline!std::string!
argument into \lstinline!const!~\lstinline!char*! and \lstinline!std::size_t!
arguments and converting the \lstinline!int! return value to \lstinline!bool!.

A particularly common kind of interface adaptation is \emcppsgloss{partial
application}, whereby we reduce the \emph{parameter count} of a function
by holding one or more of its arguments constant for the duration of the
algorithm:

\begin{emcppslisting}
#include <algorithm>  // (ù{\codeincomments{std::all\_of}}ù)

template <typename InputIter, typename T>
bool all_greater_than(InputIter first, InputIter last, const T& v)
    // returns (ù{\codeincomments{true}}ù) if all the values in the specified range [(ù{\codeincomments{first}}ù), (ù{\codeincomments{last}}ù))
    // are greater than the specified (ù{\codeincomments{v}}ù), and false otherwise
{
    return std::all_of(first, last, [&](const T& i) { return i > v; });
}
\end{emcppslisting}
    

\noindent In the example above, the greater-than operator (\lstinline!>!) takes two
operands, but the \lstinline!std::all_of! algorithm expects a functor
taking a single argument. The \emcppsgloss[lambda expressions]{lambda expression} passes its
single argument as the first operand to \lstinline!operator>! and
\emph{binds} the other operand to the captured \lstinline!v! value, thus
solving the interface mismatch.

Finally, let's touch on \emcppsgloss{currying}, a transformation borrowed
from lambda calculus and functional programming languages. Currying is a
flexible way to get results similar to \emcppsgloss{partial application} by
transforming, e.g., a function taking two parameters into one taking
just the first parameter and returning a function taking just the second
parameter. To apply this technique, we define a \emcppsgloss[lambda expressions]{lambda
expression} whose call operator returns another \emcppsgloss[lambda expressions]{lambda
expression}, i.e., a \emcppsgloss{closure} that returns another
\emcppsgloss{closure}:

\begin{emcppshiddenlisting}[emcppsbatch=e11]
#include <algorithm>  // (ù{\codeincomments{all\_of}}ù)
\end{emcppshiddenlisting}
\begin{emcppslisting}[emcppsbatch=e11]
template <typename InputIter, typename T>
bool all_greater_than2(InputIter first, InputIter last, const T& v)
    // returns (ù{\codeincomments{true}}ù) if all the values in the specified range [(ù{\codeincomments{first}}ù), (ù{\codeincomments{last}}ù))
    // are greater than the specified (ù{\codeincomments{v}}ù), and false otherwise
{
    auto isGreaterThan = [](const T& v){
        return [&v](const T& i){ return i > v; };
    };
    return std::all_of(first, last, isGreaterThan(v));
}
\end{emcppslisting}
    

\noindent The example above is another way to express the previous example. The
call operator for \lstinline!isGreaterThan! takes a single argument,
\lstinline!v!, and returns another single-argument \emcppsgloss{closure object}
that can be used to compare \lstinline!i! to \lstinline!v!. Thus,
\lstinline!isGreaterThan(v)(i)! is equivalent to
\lstinline!i!~\lstinline!>!~\lstinline!v!.

\subsubsection[Emulating local functions]{Emulating local functions}\label{emulating-local-functions}

Local functions in languages other than C++ allow functions to be
defined within other functions. They are useful when the outer function
needs to repeat a set of steps two or more times but where the repeated
steps are meaningless outside of the immediate context and/or require
access to the outer function's local variables. Using a \emcppsgloss[lambda expressions]{lambda
expression} to produce a re-usable \emcppsgloss{closure} provides this
functionality in C++:

\begin{emcppslisting}
class Token { /* ... */ };

bool parseToken(const char*& cursor, Token& result)
    // Parse the token at (ù{\codeincomments{cursor}}ù) up to the next space or end-of-string,
    // setting (ù{\codeincomments{result}}ù) to the resulting token value.  Advance (ù{\codeincomments{cursor}}ù) to the
    // space or to the null terminator and return (ù{\codeincomments{true}}ù) on success.  Reset
    // (ù{\codeincomments{cursor}}ù) to its original value, set (ù{\codeincomments{result}}ù) to an empty token, and
    // return (ù{\codeincomments{false}}ù) on failure.
{
    const char* const initCursor = cursor;
    auto error = [&]
    {
        cursor = initCursor;
        result = Token{};
        return false;
    };
    // ...
    if (*cursor++ != '.')
    {
        return error();
    }
    // ...
}
\end{emcppslisting}
    

\noindent The \lstinline!error! \emcppsgloss{closure object} acts as a local function that
performs all of the necessary error processing and returns
\lstinline!false!. Using this object, every error branch can be reduced to
a single statement, \lstinline!return!~\lstinline!error()!. Without
\emcppsgloss{lambda expressions}, the programmer would likely resort to
defining a custom class to store the parameters, using a \lstinline!goto!,
or, worse, cutting-and-pasting the three statements shown within the
\emcppsgloss{lambda body}.

\subsubsection[Emulate user-defined control constructs]{Emulate user-defined control constructs}\label{emulate-user-defined-control-constructs}

Using a \emcppsgloss[lambda expressions]{lambda expression}, an algorithm can look almost like a
new control construct in the language:

\begin{emcppslisting}
#include <mutex>  // (ù{\codeincomments{std::mutex}}ù)
#include <vector>  // (ù{\codeincomments{std::vector}}ù)

template <typename RandomIter, typename F>
void parallel_foreach(RandomIter first, RandomIter last, const F& op)
    // For each element, (ù{\codeincomments{e}}ù), in [(ù{\codeincomments{first}}ù), (ù{\codeincomments{last}}ù)), create a copy (ù{\codeincomments{opx}}ù) of (ù{\codeincomments{op}}ù),
    // and invoke (ù{\codeincomments{opx(e)}}ù).  Any number of invocations of (ù{\codeincomments{opx(e)}}ù) may occur
    // concurrently, each using a separate copy of (ù{\codeincomments{op}}ù).
{ /*...*/ }

void processData(std::vector<double>& data)
{
    double       beta    = 0.0;
    double const init = 7.45e-4;
    std::mutex m;

    parallel_foreach(data.begin(), data.end(), [&, init](double e) mutable
    {
        if (e < 1.0)
        {
           // ...
        }
        else
        {
           // ...
        }
    });
}
\end{emcppslisting}
    

\noindent The \lstinline!parallel_foreach! algorithm is intended to act like a
\lstinline!for! loop except that all of the elements in the input range may
potentially be processed in parallel. By inserting the ``body'' of this
``parallel for loop'' directly into the call to
\lstinline!parallel_foreach!, the resulting loop looks and feels a lot
like a built-in control construct. Note that the \emcppsgloss{capture
default} is \emcppsgloss[captured by reference]{capture by reference} and will result in all of the
iterations sharing the outer function's call frame, including the mutex
variable, \lstinline!m!, used to prevent race conditions. This
\emcppsgloss{capture default} should be used with care in parallel
computations, which often use \emcppsgloss[captured by copy]{capture by copy} to deliberately
\emph{avoid} sharing. If an asynchronous computation might outlive its
caller, then using \emcppsgloss[captured by copy]{capture by copy} is a must for avoiding
dangling references; see \intrarefsimple{potential-pitfalls-lambda}.

\subsubsection[Variables and control constructs in expressions]{Variables and control constructs in expressions}\label{variables-and-control-constructs-in-expressions}

In situations where a single expression is required --- e.g.,
member-initializers, initializers for \lstinline!const! variables, and so
on --- an \emph{immediately evaluated} \emcppsgloss[lambda expressions]{lambda expression} allows
that expression to include local variables and control constructs such
as loops:

\begin{emcppslisting}
#include <climits>  // (ù{\codeincomments{SHRT\_MAX}}ù)

bool isPrime(long i);
    // Return (ù{\codeincomments{true}}ù) if (ù{\codeincomments{i}}ù) is a prime number.

const short largestShortPrime = []{
    for (short v = SHRT_MAX; ; v -= 2) {
        if (isPrime(v)) return v;
    }
}();
\end{emcppslisting}
    

\noindent The value of \lstinline!largestShortPrime! must be set at initialization
time because it is a \lstinline!const! variable with static duration. The
loop inside of the \emcppsgloss[lambda expressions]{lambda expression} computes the desired
value, using local variable, \lstinline!v!, and a \lstinline!for! loop. Note
that the call operator for the resulting \emcppsgloss{closure object} is
immediately invoked via the \lstinline!()! argument list at the end of the
\emcppsgloss[lambda expressions]{lambda expression}; the \emcppsgloss{closure object} is never stored
in a named variable and goes out of scope as soon as the full expression
is completely evaluated. This computation would formerly have been
possible only by creating a single-use \emph{named} function.

\subsubsection[Use with \lstinline!std::function!]{Use with {\SubsubsecCode std::function}}\label{use-with-std::function}

As convenient as a \emcppsgloss[lambda expressions]{lambda expression} is for passing a functor
to an algorithm \emph{template}, the fact that each \emcppsgloss{closure
type} is unnamed and distinct makes it difficult to use them outside of
a generic context. The C++11 Standard Library class template,
\lstinline!std::function!, bridges this gap by providing a polymorphic
invocable type that can be constructed from any type with a compatible
invocation prototype, including but not limited to \emcppsgloss[closure type]{closure
types}.

A simple interpreter for a postfix input language stores a sequence of
instructions in a \lstinline!std::vector!. Each instruction can be of a
different type, but they all accept the current stack pointer as an
argument and return the new stack pointer as a result. Each instruction
is typically a small operation, ideally suited for being expressed as
\emcppsgloss{lambda expressions}:

\begin{emcppslisting}
#include <cstdlib>     // (ù{\codeincomments{std::strtol}}ù)
#include <functional>  // (ù{\codeincomments{std::function}}ù)
#include <string>      // (ù{\codeincomments{std::string}}ù)
#include <vector>      // (ù{\codeincomments{std::vector}}ù)

using Instruction = std::function<long*(long*& sp)>;

std::vector<Instruction> instructionStream;

std::string nextToken();                 // read the next token
char tokenOp(const std::string& token);  // operator for token

void readInstructions()
{
    std::string token;
    Instruction nextInstr;
    while (!(token = nextToken()).empty())
    {
        switch (tokenOp(token))
        {
            case 'i':
            {
                // Integer literal
                long v = std::strtol(token.c_str(), nullptr, 10);
                nextInstr = [v](long* sp){ *sp++ = v; return sp; };
                break;
            }
            case '+':
            {
                // + operation
                nextInstr = [](long*& sp){
                    long v1 = *--sp;
                    long v2 = *--sp;
                    *sp++ = v1 + v2;
                    return sp;
                };
                break;
            }
            // ... more cases
        }

        instructionStream.push_back(nextInstr);
    }
}
\end{emcppslisting}
    

\noindent The \lstinline!Instruction! type alias is a \lstinline!std::function! that can
hold, through a process called \emcppsgloss{type erasure}, any invocable
object that takes a \lstinline!long*! argument and returns a \lstinline!long*!
result. The \lstinline!readInstructions! function reads successive string
tokens and switches on the operation represented by the token. If the
operation is \lstinline!i!, then the token is an integer literal. The
string token is converted into a \lstinline!long! value, \lstinline!v!, which
is captured in a \emcppsgloss[lambda expressions]{lambda expression}. The resulting
\emcppsgloss{closure} object is stored in the \lstinline!nextInstr! variable;
when called, it will push \lstinline!v! onto the stack. Note that the
\lstinline!nextInstr! variable outlives the primary \lstinline!v! variable,
but, because \lstinline!v! was \emcppsgloss{captured by copy}, the
\emcppsgloss{captured variable}'s lifetime is the same as the \emcppsgloss{closure
object}'s. If the next operation is \lstinline!+!, \lstinline!nextInstr! is
set to the \emcppsgloss{closure object} of an entirely different
\emcppsgloss[lambda expressions]{lambda expression}, one that captures nothing and whose call
operator pops two values from the stack and pushes their sum back onto
the stack.

After the \lstinline!switch! statement, the current value of
\lstinline!nextInstr! is appended to the instruction stream. Note that,
although each \emcppsgloss{closure type} is different, they all can be stored
in an \lstinline!Instruction! object because the prototype for their call
operator matches the prototype specified in the instantiation of
\lstinline!std::function!. The \lstinline!nextInstr! variable can be created
empty, assigned from the value of a \emcppsgloss[lambda expressions]{lambda expression}, and then
later reassigned from the value of a different \emcppsgloss[lambda expressions]{lambda
expression}. This flexibility makes \lstinline!std::function! and
\emcppsgloss{lambda expressions} a potent combination.

One specific use of \lstinline!std::function! worth noting is to return a
\emcppsgloss[lambda expressions]{lambda expression} from a non-template function:

\begin{emcppshiddenlisting}[emcppsbatch=e12]
#include <functional>  // (ù{\codeincomments{std::function}}ù)
\end{emcppshiddenlisting}
\begin{emcppslisting}[emcppsbatch=e12]
std::function <int(int)> add_n(int n)
{
    return [n](int i) { return n + i; };
}

int result = add_n(3)(5);  // (ù{\codeincomments{result}}ù) is 8.
\end{emcppslisting}
    

\noindent The return value of \lstinline!add_n! is a \emcppsgloss{closure object} wrapped
in a \lstinline!std::function! object. Note that \lstinline!add_n! is not a
template and that it is not called in a template or \lstinline!auto!
context. This example illustrates a runtime-polymorphic way to achieve
\emcppsgloss{currying}; see the earlier example in \intrarefsimple{interface-adaptation,-partial-application,-and-currying}.

\subsubsection[Event-driven callbacks]{Event-driven callbacks}\label{event-driven-callbacks}

Event-driven systems tend to have interfaces that are littered with
callbacks:

\begin{emcppslisting}[emcppsbatch=e13]
#include <memory>  // std::unique_ptr

class DialogBox { /* ... */ };

template <typename Button1Func, typename Button2Func>
std::unique_ptr<DialogBox> twoButtonDialog(const char* prompt,
                                           const char* button1text,
                                           Button1Func button1callback,
                                           const char* button2text,
                                           Button2Func button2callback)
{
    // ...
}
\end{emcppslisting}
    

\noindent The \lstinline!twoButtonDialog! factory function takes three strings and
two callbacks and returns a pointer to a dialog box having two buttons.
The dialog-box logic invokes one of the two callbacks, depending on
which of the two buttons is pressed. These callbacks are often quite
small pieces of code that can best be expressed directly in the program
logic using \emcppsgloss{lambda expressions}:

\begin{emcppslisting}[emcppsbatch=e13]
void runModalDialogBox(DialogBox& db);

void launchShuttle(/* ... */)
{
    bool doLaunch = false;

    std::unique_ptr<DialogBox> confirm =
        twoButtonDialog("Are you sure you want to launch the shuttle?",
                        "Yes", [&]{ doLaunch = true; },
                        "No",  []{});

    runModalDialogBox(*confirm);

    if (doLaunch)
    {
        // ... launch the shuttle!
    }
}
\end{emcppslisting}
    

\noindent Here, the user is being prompted as to whether or not to launch a
missile. Since the dialog box is processed entirely within the
\lstinline!launchShuttle! function, it is convenient to express two
callbacks in-place, within the function, using \emcppsgloss{lambda
expressions}. The first \emcppsgloss[lambda expressions]{lambda expression} --- passed as the
callback for when the user clicks ``Yes'' --- captures the
\lstinline!doLaunch! flag by reference and simply sets it to \lstinline!true!.
The second \emcppsgloss[lambda expressions]{lambda expression} --- passed as the callback for
when the user clicks ``No'' --- does nothing, leaving the
\lstinline!doLaunch! flag having its original \lstinline!false! value. The
simplicity of these callbacks come from fact that they are effectively
extensions of the surrounding block and, hence, have access (via the
\emcppsgloss{lambda capture}) to block-scoped variables such as
\lstinline!doLaunch!.

\subsubsection[Recursion]{Recursion}\label{recursion}

A \emcppsgloss[lambda expressions]{lambda expression} cannot refer to itself, so creating one
that is recursive involves using one of a number of different possible
workarounds. If the \emcppsgloss{lambda capture} is empty, recursion can be
accomplished fairly simply by converting the \emcppsgloss[lambda expressions]{lambda expression}
into a plain function pointer stored in a \lstinline!static! variable:

\begin{emcppslisting}
void f1()
{
    static int (*const fact)(int) = [](int i)
    {
        return i < 2 ? 1 : i * fact(i-1);
    };

    int result = fact(4);  // computes 24
}
\end{emcppslisting}
    

\noindent In the above example, \lstinline!fact(n)! returns the factorial of
\lstinline!n!, computed using a recursive algorithm. The variable,
\lstinline!fact!, becomes visible before its initializer is compiled,
allowing it to be called from within the \emcppsgloss[lambda expressions]{lambda expression}. To
enable the conversion to function pointer, the \emcppsgloss{lambda capture}
must be empty; hence, \lstinline!fact! must be static so that it can be
accessed without capturing it.

If a recursive \emcppsgloss[lambda expressions]{lambda expression} is desired with a nonempty
\emcppsgloss{lambda capture}, then the entire recursion can be enclosed in an
outer \emcppsgloss[lambda expressions]{lambda expression}:

\begin{emcppslisting}
void f2(int n)
{
    auto permsN = [n](int m) -> int
    {
        static int (*const imp)(int, int) = [](int x, int m) {
            return m <= x ? m : m * imp(x, m - 1);
        };
        return imp(m - n + 1, m);
    };

    int a = permsN(5);  // permutations of 5 items, (ù{\codeincomments{n}}ù) at a time
    int b = permsN(4);  // permutations of 4 items, (ù{\codeincomments{n}}ù) at a time
}
\end{emcppslisting}
    

\noindent In this example, \lstinline!permsN(m)!, returns the number of permutations
of \lstinline!m! items taken \lstinline!n! at a time, where \lstinline!n! is
captured by the \emcppsgloss{closure object}. The implementation of
\lstinline!permsN! defines a nested \lstinline!imp! function pointer that uses
the same technique as \lstinline!fact!, above, to achieve recursion. Since
\lstinline!imp! must have an empty \emcppsgloss{lambda capture}, everything it
needs is passed in as arguments by the \lstinline!permsN! enclosing
\emcppsgloss[lambda expressions]{lambda expression}. Note that the \lstinline!imp! pointer and the
\emcppsgloss[lambda expressions]{lambda expression} from which it is initialized do not needed to
be scoped inside of the \lstinline!permN! \emcppsgloss[lambda expressions]{lambda expression};
whether such nesting is desirable is a matter of taste.

In C++14, additional approaches to recursion (e.g., the ``Y Combinator''
borrowed from lambda calculus{\cprotect\footnote{\cite{derevenets16}}}) are possible due to
using \emcppsgloss[generic lambda]{generic lambdas}; see \featureref{\locationd}{genericlambda}.

\subsection[Potential Pitfalls]{Potential Pitfalls}\label{potential-pitfalls-lambda}

\subsubsection[Dangling references]{Dangling references}\label{dangling-references}

\emcppsgloss[closure object]{Closure objects} can capture references to local variables and
copies of the \lstinline!this! pointer. If a copy of the \emcppsgloss{closure
object} outlives the stack frame in which it was created, these
references can refer to objects that have been destroyed. The two ways
in which a \emcppsgloss{closure object} can outlive its creation context are
if (1) it is returned from the function or (2) it is stored in a data
structure for later invocation:

\begin{emcppslisting}
#include <functional>  // std::function
#include <vector>      // std::vector

class Class1
{
    int d_mem;

    static std::vector<std::function<double(void*)> > s_workqueue;

    std::function<void(int)> mf1()
    {
        int local;
        return [&](int i) -> void { d_mem = local = i; };  // Bug, dangling refs
    }

    void mf2()
    {
        double local = 1.0;
        s_workqueue.push_back([&,this](void* p) -> double {
                return p ? local : double(d_mem);
            });  // Bug, dangling refs
    }
};
\end{emcppslisting}
    

\noindent The example above uses \lstinline!std::function! to hold \emcppsgloss[closure object]{closure
objects}, as described in \intrarefsimple{use-with-std::function}. In member
function \lstinline!mf1!, the \emcppsgloss{lambda body} modifies both the local
variable and the member variable currently in scope. However, as soon as
the function returns, the local variable goes out of scope and the
\emcppsgloss{closure} contains a dangling reference. Moreover, the object on
which it is invoked can also go out of scope while the \emcppsgloss{closure
object} continues to exist. Modifying either \lstinline!this->d_mem! or
\lstinline!local! through the capture is likely to corrupt the stack,
leading to a crash, potentially much later in the program.

The member function \lstinline!mf2!, rather than \emph{return} a
\emcppsgloss{closure} with dangling references, stores it in a data
structure, i.e., the \lstinline!s_workqueue! static \lstinline!vector!. Once
again, \lstinline!local! and \lstinline!d_mem! become dangling references and
can result in data corruption when the call operator for the stored
\emcppsgloss{closure object} is invoked. It is safest to capture
\lstinline!this! and use \emcppsgloss[captured by reference]{capture by reference} only when the
lifetime of the \emcppsgloss{closure object} is clearly limited to the
current function. \emcppsgloss[implicitly captured]{Implicitly captured} \lstinline!this! is
particularly insidious because, even if the \emcppsgloss{capture default} is
\emcppsgloss[captured by copy]{capture by copy}, member variables are not copied and are often
referenced without the \lstinline!this->! prefix, making them hard to spot
in the source code.

\subsubsection[Overuse]{Overuse}\label{overuse}

The ability to write functions, especially functions with state, at the
point where they are needed and without much of the syntactic overhead
that accompanies normal functions and class methods, can potentially
lead to a style of code that uses ``lambdas everywhere'', losing the
abstraction and well-documented interfaces of separate functions.
\emcppsgloss[lambda expressions]{Lambda expressions} are not intended for large-scale reuse.
Sprinkling \emcppsgloss{lambda expressions} throughout the code can result in
poor software-engineering practices such as cut-and-paste programming
and the absence of cohesive abstractions.

\subsubsection[Mixing captured and non-captured variables]{Mixing captured and non-captured variables}\label{mixing-captured-and-non-captured-variables}

A \emcppsgloss{lambda body} can access both automatic-duration local
variables that were captured from the enclosing block and
static-duration variables that need not and cannot be captured.
Variables \emcppsgloss{captured by copy} are ``frozen'' at the point of
capture and cannot be changed except by the \emcppsgloss{lambda body} (if
mutable), whereas static variables can be changed independent of the
\emcppsgloss[lambda expressions]{lambda expression}. This difference is often useful but can
cause confusion when reasoning about a \emcppsgloss[lambda expressions]{lambda expression}:

\begin{emcppshiddenlisting}[emcppsbatch=e14]
#include <cassert>  // standard C (ù{\codeincomments{assert}}ù) macro
\end{emcppshiddenlisting}
\begin{emcppslisting}[emcppsbatch=e14]

void f1()
{
    static int a;
    int        b;

    a = 5;
    b = 6;

    auto c1 = [b]{ return a + b; };  // OK, (ù{\codeincomments{b}}ù) is **captured by copy**.
    assert(11 == c1());              // OK, (ù{\codeincomments{a == 5}}ù) and (ù{\codeincomments{b == 6}}ù).
    ++b;                             // Increment *primary* (ù{\codeincomments{b}}ù).
    assert(11 == c1());              // OK, captured (ù{\codeincomments{b}}ù) did not change.
    ++a;                             // Increment static-duration (ù{\codeincomments{a}}ù).
    assert(12 == c1());              // Bug, (ù{\codeincomments{a == 6}}ù) and captured (ù{\codeincomments{b == 6}}ù)
}
\end{emcppslisting}
    

\noindent When the \emcppsgloss{closure object} for \lstinline!c1! is created, the
captured \lstinline!b! value is frozen within the \emcppsgloss{lambda body}.
Changing the primary \lstinline!b! has no effect. However, \lstinline!a! is
not captured (nor is it allowed to be). As a result, there is only
\emph{one} \lstinline!a! variable, and modifying that variable outside of
the \emcppsgloss{lambda body} changes the result of invoking the call
operator.

C++14 \emcppsgloss{capture init} can be used to effectively capture a
non-local variable:

\begin{emcppshiddenlisting}[emcppsbatch=e15]
#include <cassert>  // standard C (ù{\codeincomments{assert}}ù) macro
\end{emcppshiddenlisting}
\begin{emcppslisting}[emcppsbatch=e15]
void f2()
{
    static int a;
    int        b;

    a = 5;
    b = 6;

    auto c1 = [a=a,b]{ return a + b; };  // OK, capture both (ù{\codeincomments{a}}ù) and (ù{\codeincomments{b}}ù).
    ++a;                                 // Increment static-duration (ù{\codeincomments{a}}ù).
    assert(11 == c1());                  // OK, captured (ù{\codeincomments{a}}ù) did not change.
}
\end{emcppslisting}
    

\noindent The \emcppsgloss{lambda capture}, \lstinline![a=a,b]!, creates a new capture
variable \lstinline!a! that is initialized from the static variable
\lstinline!a! at the point that the \emcppsgloss[lambda expressions]{lambda expression} is evaluated
to produce a \emcppsgloss{closure object}. The \lstinline!a! variable within the
\emcppsgloss{lambda body} refers to this captured variable, not to the static
one.

\subsubsection[Local variables in unevaluated contexts can yield surprises]{Local variables in unevaluated contexts can yield surprises}\label{local-variables-in-unevaluated-contexts-can-yield-surprises}

To use a local variable, \lstinline!x!, from the surrounding block as part
of an unevaluated operand (e.g., \lstinline!sizeof(x)! or
\lstinline!alignof(x)!), it is generally not necessary to capture
\lstinline!x! because it is not \emcppsgloss{ODR-used} within the \emcppsgloss{lambda
body}. Whether or not \lstinline!x! is captured, most expressions in
unevaluated contexts behave as though \lstinline!x! were \emph{not}
captured and the expression were evaluated directly in the enclosing
block scope. This is itself surprising because, for example, a captured
variable in a non-\lstinline!mutable! \emcppsgloss[lambda expressions]{lambda expression} is
\lstinline!const!, whereas the primary variable might not be:

\begin{emcppslisting}
#include <iostream>  // (ù{\codeincomments{std::cout}}ù)

short s1(int&)       { return 0; }
int   s1(const int&) { return 0; }

void f1()
{
    int x = 0;  // (ù{\codeincomments{x}}ù) is a non-const (ù{\emphincomments{lvalue}}ù).
    [x]{
        // captured (ù{\codeincomments{x}}ù) in non-(ù{\codeincomments{mutable}}ù) lambda is (ù{\emphincomments{lvalue}}ù) of type (ù{\codeincomments{const int}}ù)
        std::cout << sizeof(s1(x)) << '\n';  // prints sizeof(short)
        auto s1x = s1(x);                    // yields an (ù{\codeincomments{int}}ù)
        std::cout << sizeof(s1x) << '\n';    // prints sizeof(int)
    }();
}
\end{emcppslisting}
    

\noindent The first print statement calls \lstinline!s1(x)! in an unevaluated
context, which ignores the captured \lstinline!x! and returns the size of
the result of \lstinline!s1(int&)!. The next statement actually
\emph{evaluates} \lstinline!s1(x)!, passing the \emph{captured} \lstinline!x!
and calling \lstinline!s1(const!~\lstinline!int&)! because the call operator
is not decorated with \lstinline!mutable!.

When using \lstinline!decltype(x)!, the result is the declared type of the
\emph{primary} variable, regardless of whether or not \lstinline!x! was
captured. However, if \lstinline!x! had been \emcppsgloss{captured by copy},
\lstinline!decltype((x))! (with two sets of parentheses) would have yielded
the \romeovalue{lvalue} type of the \emph{captured} variable. There is some dispute
as to what the correct results should be if \lstinline!x! is \emph{not}
captured, with some compilers yielding the type of the primary variable
and others complaining that it was not captured.

\begin{emcppslisting}
void f1()
{
    int x = 0;  // (ù{\codeincomments{x}}ù) is a non-(ù{\codeincomments{const}}ù) }}(ù{\emphincomments{lvalue}}ù).
    auto c1 = [x]{ decltype((x)) y = x; };  // (ù{\codeincomments{y}}ù) has type (ù{\codeincomments{const int\&}}ù).
}
\end{emcppslisting}
    

\noindent Finally, there is an unsettled question as to whether \lstinline!typeid(x)!
is an \emcppsgloss[ODR-used]{ODR-use} of \lstinline!x! and, therefore, requires that
\lstinline!x! be captured. Some compilers will complain about the following
code:

\begin{emcppslisting}
#include <typeinfo>  // (ù{\codeincomments{typeid}}ù)
void f3()
{
    int x = 0;
    auto c1 = []() -> const std::type_info& { return typeid(x); };
        // Error, on some platforms ``(ù{\codeincomments{x}}ù) was not captured''
}
\end{emcppslisting}
    

\noindent One can avoid this pitfall simply by calling \lstinline!typeid! outside of
the lambda, capturing the result if necessary:

\begin{emcppslisting}
#include <typeinfo>  // (ù{\codeincomments{typeid}}ù)
void f3()
{
    int x = 0;
    const std::type_info& xid = typeid(x);
        // OK, (ù{\codeincomments{typeid}}ù) called outside of lambda
    auto c1 = [&]() -> const std::type_info& { return xid; };
        // OK, return captured (ù{\codeincomments{typeinfo}}ù)
}
\end{emcppslisting}
    

\subsection[Annoyances]{Annoyances}\label{annoyances}

\subsubsection[Debugging]{Debugging}\label{debugging}

By definition, lambdas do not have names. Tools such as debuggers and
stack-trace examiners typically display the compiler-generated names of
the closure types instead of names selected by the programmer to clearly
describe the purpose of a function, making it difficult to discern where
a problem occurred.

\subsubsection[Can’t capture \lstinline!*this! by copy]{Can’t capture {\SubsubsecCode *this} by copy}\label{can’t-capture-*this-by-copy}

A \emcppsgloss[lambda expressions]{lambda expression} can freeze the value of a surrounding local
variable by using \emcppsgloss[captured by copy]{capture by copy}, but no such ability is
available directly to copy the object pointed to by \lstinline!this!. In
C++14, this deficiency can be mitigated using \emcppsgloss{capture init}:

\begin{emcppslisting}
class Class1
{
    int d_value;

    void mf1()
    {
        auto c1 = [self=*this]{ return self.d_value; };
    }
};
\end{emcppslisting}
    

\noindent The \emcppsgloss{lambda capture}, \lstinline![self=*this]! creates a new
captured variable, \lstinline!self!, that contains a copy of
\lstinline!*this!. Unfortunately, accessing member variable
\lstinline!d_value! requires explicit use of \lstinline!self.d_value!.

C++11 doesn't have \emcppsgloss{capture init}, so it is necessary to create a
\lstinline!self! variable external to the \emcppsgloss[lambda expressions]{lambda expression} and
capture \emph{that} variable{\cprotect\footnote{As of C++17,
\lstinline!*this! can be captured directly with \lstinline!this! within the
\emcppsgloss{lambda body} pointing to the \emph{copy} rather than the
original:

\begin{emcppslisting}[style=footcode]
class Class2
{
    int d_value;

    void mf1()
    {
        auto c1 = [*this]{ return d_value; };
            // C++17: return (ù{\fncodeincomments{d\_value}}ù) from copy of (ù{\fncodeincomments{*this}}ù)
    }
};
\end{emcppslisting}
      }}:

\begin{emcppslisting}
class Class1
{
    int d_value;

    void mf1()
    {
        Class1& self = *this;
        auto c1 = [self]{ return self.d_value; };
    }
};
\end{emcppslisting}
    

\subsubsection[Confusing mix of immediate and deferred-execution code]{Confusing mix of immediate and deferred-execution code}\label{confusing-mix-of-immediate-and-deferred-execution-code}

The main selling point of \emcppsgloss{lambda expressions} --- i.e., the
ability to define a function object at the point of use --- can
sometimes be a liability. The code within a \emcppsgloss{lambda body} is
typically not executed immediately but is deferred until some other
piece of code, e.g., an algorithm, invokes it as a callback. The code
that is immediately executed and the code whose invocation is deferred
are visually intermixed in a way that could confuse a future maintainer.
For example, let's look at a simplified excerpt from an earlier use
case, \intraref{use-cases-lambda}{use-with-std::function}.

\begin{emcppslisting}
#include <string>      // (ù{\codeincomments{std::string}}ù)
#include <functional>  // (ù{\codeincomments{std::function}}ù)

void readInstructions()
{
    std::string                     token;
    std::function<long*(long*& sp)> nextInstr;

    while ( /* ... */ (!token.empty()))
    {
        switch (token[0])
        {
            // ... more cases
            case '+':
            {
                // + operation
                nextInstr = [](long*& sp){
                    long v1 = *--sp;
                    long v2 = *--sp;
                    *sp++ = v1 + v2;
                    return sp;
                };
                break;
            }
            // ... more cases
        }
       // ...
    }
}
\end{emcppslisting}
    

\noindent A casual reading might lead to the assumption that operations such as
\lstinline!*--sp! are taking place within \lstinline!case!~\lstinline!'+'!, when
the truth is that these operations are encapsulated in a \emcppsgloss[lambda expressions]{lambda
expression} and are not executed until the \emcppsgloss{closure object} is
called (via \lstinline!nextInstr!) in a relatively distant part of the
code.

\subsubsection[Trailing punctuation]{Trailing punctuation}\label{trailing-punctuation}

The body of a \emcppsgloss[lambda expressions]{lambda expression} is a \emph{compound statement}.
When compound statements appear elsewhere in the C++ grammar, e.g., as
the body of a function or loop, they are not followed by punctuation. A
\emcppsgloss[lambda expressions]{lambda expression}, conversely, is invariably followed by some
sort of punctuation, usually a semicolon or parenthesis but sometimes a
comma or binary operator. This difference between a \emcppsgloss{lambda body}
and other compound statements makes this punctuation easy to forget:

\begin{emcppslisting}
auto c1 = []{ /* ... */ };  // <-- Don't forget the semicolon at the end.
\end{emcppslisting}
    

\noindent The extra punctuation can also be unattractive when emulating a control
construct using \emcppsgloss{lambda expressions} as in the
\lstinline!parallel_foreach! example in \intraref{use-cases-lambda}{emulate-user-defined-control-constructs}:

\begin{emcppshiddenlisting}[emcppsbatch=e16]
template <typename RandomIter, typename F>
void parallel_foreach(RandomIter first, RandomIter last, const F& op) {}
#include <vector>  // (ù{\codeincomments{std::vector}}ù)
\end{emcppshiddenlisting}
\begin{emcppslisting}[emcppsbatch=e16]
void f(const std::vector<int>& data)
{
    // ...
    for (int e : data)
    {
        // ...             (ù{\codeincomments{for}}ù) loop body
    }  // <-- no punctuation after the closing brace

    parallel_foreach(data.begin(), data.end(), [&](int e)
    {
        // ...             parallel loop body
    });  // <-- Don't forget the closing parenthesis and semicolon.
    // ...
}
\end{emcppslisting}
    

\noindent In the above code snippet, the programmer would like the
\lstinline!parallel_foreach! algorithm to look as much like the built-in
\lstinline!for! loop as possible. However, the built-in \lstinline!for! loop
doesn't end with a closing parenthesis and a semicolon, whereas the
\lstinline!parallel_foreach! does, so the illusion of a language extension
is incomplete.

\subsection[See Also]{See Also}\label{see-also}

\begin{itemize}
\item{\seealsoref{decltype}{\seealsolocationa}illustrates a form of type inference often used in conjunction with (or in place of) trailing return types.}
\item{\seealsoref{Function-Return-Type-Deduction}{\seealsolocationf}shows a form of type inference that shares syntactical similarities with trailing return types, leading to potential pitfalls when migrating from C++11 to C++14.}
\end{itemize}

\subsection[Further Reading]{Further Reading}\label{further-reading}

TODO



\newpage
%\section[{\tt noexcept} Operator]{The {\SecCode noexcept} Operator}%\label{noexcept}


\emcppsFeature{
    short={\lstinline!noexcept! Operator},
    tocshort={{\TOCCode noexcept} Operator},
    long={Asking if an Expression Cannot {\SecCode throw}},
    toclong={Asking if an Expression Cannot \lstinline!throw!},
    rhshort={{\RHCode noexcept} Operator},
}{noexceptoperator}
\setcounter{table}{0}
\setcounter{footnote}{0}
\setcounter{lstlisting}{0}
%\section[{\tt noexcept} Operator]{The {\SecCode noexcept} Operator}%\label{noexcept}

placeholder text.........
 
 

\newpage
%\section[Opaque {\tt enum}s]{Opaque Enumeration Declarations}\label{enumopaque}\label{opaque-enumeration-declarations}

% 01 Feb, 2021 JMB - packet 4 compilation fixes


\emcppsFeature{
    short={Opaque \lstinline!enum!s},
    tocshort={Opaque {\TOCCode enum}s},
    long={Opaque Enumeration Declarations},
    rhshort={Opaque {\RHCode enum}s},
}{enumopaque}
\label{opaque-enumeration-declarations}
\setcounter{table}{0}
\setcounter{footnote}{0}
\setcounter{lstlisting}{0}
%\section[Opaque {\tt enum}s]{Opaque Enumeration Declarations}\label{enumopaque}\label{opaque-enumeration-declarations}


Enumerated types, such as an \lstinline!enum! or \lstinline!enum! class (see \featureref{\locationc}{enumclass}), whose underlying type (see \featureref{\locationc}{explicit-enumeration-underlying-type}) is well-specified, can be
declared without being defined, i.e., declared without its enumerators.

\subsection[Description]{Description}\label{description}

We identify two distinct forms of \romeogloss{opaque declarations}, i.e.,
declarations that are not also definitions:
\begin{enumerate}
\item{A \romeogloss{forward declaration} has some translation unit in which the full definition and that declaration both appear. This can be a declaration in a header file where the definition is in the same header or in the corresponding implementation file. It can also be a declaration that appears in the same implementation file as the corresponding definition.}
\item{A \romeogloss{local declaration} has no translation unit that includes both that declaration and the corresponding full definition.}
\end{enumerate}

A classic (C++03) C-style \lstinline!enum! cannot have \romeogloss{opaque
declarations}, nor can its definition be repeated within the same
\romeogloss{translation unit} (TU):

\begin{emcppslisting}
enum E0;                 // Error, opaque declaration
enum    { e_A0, e_B0 };  // OK, anonymous classic (only) (ù{\codeincomments{enum}}ù)
enum E1 { e_A1, e_B1 };  // OK, definition
enum E1;                 // OK, redeclaration in same TU
enum E2 { e_A2, e_B2 };  // OK, definition
enum E2 { e_A2, e_B2 };  // Error, redefinition in same TU
\end{emcppslisting}

\noindent The underlying integral types used to represent objects of each of the
(classic) enumerations in the example above is \romeogloss{implementation defined}, making
all of them ineligible for \romeogloss{opaque declaration}. This restriction
on \romeogloss{opaque declarations} exists because the specific values of
the enumerators might affect the underlying type %(see \featureref{\locationc}{explicit-enumeration-underlying-type})
(e.g.,
size, alignment, \lstinline!signed!ness), and therefore the declaration
alone cannot be used to create objects of that type. A declaration that
specifies the underlying type or a full definition can, however, be used
to create objects of that type. Specifying an underlying
type %(see \featureref{\locationc}{explicit-enumeration-underlying-type})
explicitly makes \romeogloss{opaque declaration} possible:

\begin{emcppslisting}[emcppsbatch=e1]
enum    : char { e_A, e_B };     // OK, (anonymous) complete definition
enum E3 : char;                  // OK, forward declaration w/underlying type
enum E3 : char { e_A3, e_B3 };   // OK, compatible complete definition
enum E4 : short { e_A4, e_B4 };  // OK, complete definition
enum E4 : short;                 // OK, compatible opaque redeclaration
enum E4 : int;                   // Error, incompatible opaque redeclaration
enum E5 : int { e_A5, e_B5 };    // OK, complete definition
enum E5 : int { e_A5, e_B5 };    // Error, complete redefinition in same TU
\end{emcppslisting}

\noindent The modern (C++11) \lstinline!enum!~\lstinline!class!, %(see \featureref{\locationc}{enumclass}),
which
provides its enumerators with (1) stronger typing and (2) an enclosing
scope, also comes with a default \romeogloss{underlying type} of
\lstinline!int!, thereby making it eligible to be declared without a
definition (even without explicit qualification):

\begin{emcppslisting}[emcppsbatch=e1]
enum class E6;                 // OK, implicit underlying type ((ù{\codeincomments{int}}ù))
enum class E6 : int;           // OK, explicit matching underlying type
enum class E6 { e_A3, e_B3 };  // OK, compatible complete definition
enum class E7 { e_A4, e_B4 };  // OK, complete definition, (ù{\codeincomments{int}}ù) underlying type
enum class E7;                 // OK, compatible opaque redeclaration
enum class E7 : short;         // Error, incompatible opaque redeclaration
enum class E8 : long;          // OK, opaque declaration, (ù{\codeincomments{long}}ù) underlying type
enum class E8 : long { e_A5 }; // OK, compatible complete definition
enum class E9 { e_A6, e_B7 };  // OK, complete definition
enum class E9 { e_A6, e_B7 };  // Error, complete redefinition in same TU
enum class    { e_A, e_B };    // Error, anonymous (ù{\codeincomments{enum}}ù) classes are not allowed
\end{emcppslisting}

\noindent To summarize, each classical \lstinline!enum! type having an explicitly
specified \romeogloss{{underlying type}} and every modern \lstinline!enum!~\lstinline!class! type can be declared (e.g., for the first
time in a TU) as a \romeogloss{complete type}:

\begin{emcppslisting}[emcppsbatch=e1]
enum E10 : char;  static_assert(sizeof(E10)  == 1);
enum class E11;   static_assert(sizeof(E11) == sizeof(int));

E10 a;            static_assert(sizeof a == 1);
E11 b;            static_assert(sizeof b == sizeof(int));
\end{emcppslisting}

\noindent Typical usage of opaque enumerations often involves placing the
\romeogloss{forward declaration} within a header and sequestering the
complete definition within a corresponding \lstinline!.cpp! (or else a
second header), thereby \romeogloss{insulating} (at least some) clients from
changes to the enumerator list (see \intraref{use-cases-opaqueenum}{using-opaque-enumerations-within-a-header-file}):

\begin{emcppslisting}
// mycomponent.h:
// ...
enum E9 : char;   // forward declaration of (ù{\codeincomments{enum E9}}ù)
enum class E10;   // forward declaration of (ù{\codeincomments{enum class E10}}ù)


// mycomponent.cpp:
#include <mycomponent.h>
// ...
enum E9 : char { e_A9, e_B9, e_C9 };
    // complete definition compatible with forward declaration of (ù{\codeincomments{E9}}ù)

enum class E10 { e_A10, e_B10, e_C10 };
    // complete definition compatible with forward declaration of (ù{\codeincomments{E10}}ù)
\end{emcppslisting}

\noindent Note, however, that clients embedding \emph{local declarations} directly
in their code can be problematic; see \intraref{potential-pitfalls-opaqueenum}{redeclaring-an-externally-defined-enumeration-locally}.

\subsection[Use Cases]{Use Cases}\label{use-cases-opaqueenum}

\subsubsection[Using opaque enumerations within a header file]{Using opaque enumerations within a header file}\label{using-opaque-enumerations-within-a-header-file}

\romeogloss{Physical design} involves two related but distinct notions of
information \emph{hiding}: \romeogloss{encapsulation} and \romeogloss{insulation}.
An implementation detail is \emph{encapsulated} if changing it (in a
semantically compatible way) does not require clients to rework their
code but might require them to recompile it.

An \emph{insulated} implementation detail, on the other hand, can be
altered compatibly \emph{without} forcing clients even to recompile.
The advantages of avoiding such \romeogloss{compile-time coupling} transcend
merely reducing compile time. For larger codebases in which various
layers are managed under different release cycles, making a change to an
\emph{insulated} detail can be done with a \lstinline!.o! patch and a
relink the same day, whereas an \emph{uninsulated} change might
precipitate a waterfall of releases spanning days, weeks, or even
longer.

As a first example of \romeogloss{opaque-enumeration} usage, consider a
non-\emph{value-semantic} \romeogloss{mechanism} class, \lstinline!Proctor!,
implemented as a finite-state machine:

\begin{emcppslisting}[emcppsbatch=e2]
// proctor.h:
// ...
class Proctor
{
    int d_current;  // "opaque" but unconstrained (ù{\codeincomments{int}}ù) type (BAD IDEA)
    // ...

public:
    Proctor();
    // ...
};
\end{emcppslisting}

\noindent Among other private members, \lstinline!Proctor! has a data member,
\lstinline!d_current!, representing the current enumerated state of
the object. We anticipate that the implementation of the underlying
state machine will change regularly over time but that the
\lstinline!public! interface is relatively stable. We will, therefore, want
to ensure that all parts of the implementation that are likely to change
reside outside of the header. Hence, the complete definition of the
enumeration of the states (including the enumerator list itself) is
sequestered within the corresponding \lstinline!.cpp! file:

\begin{emcppslisting}[emcppsbatch=e2]
// proctor.cpp:
#include <proctor.h>

enum State { e_STARTING, e_READY, e_RUNNING, e_DONE };
Proctor::Proctor() : d_current(e_STARTING) { /* ... */ }
// ...
\end{emcppslisting}

\noindent Prior to C++11, enumerations could not be \romeogloss{forward declared}. To
avoid exposing (in a header file) enumerators that were used only
privately (in the \lstinline!.cpp! file), a completely unconstrained
\lstinline!int! would be used as a data member to represent the state. With
the advent of modern C++, we now have better options. First, we might
consider adding an explicit underlying type to the enumeration in the
\lstinline!.cpp! file:

\begin{emcppshiddenlisting}[emcppsbatch=e3]
// proctor.h:
// duplicate for second example cpp file below
class Proctor
{
    int d_current;
public:
    Proctor();
};
\end{emcppshiddenlisting}
\begin{emcppslisting}[emcppsbatch=e3]
// proctor.cpp:
#include <proctor.h>

enum State : int { e_STARTING, e_READY, e_RUNNING, e_DONE };
Proctor::Proctor() : d_current(e_STARTING) { /* ... */ }
// ...
\end{emcppslisting}

\noindent Now that the \romeogloss{component-local} \lstinline!enum! has an explicit
underlying type, we can \romeogloss{forward declare} it in the header file.
The existence of \lstinline!proctor.cpp!, which includes
\lstinline!proctor.h!, makes this declaration a \emph{forward declaration}
and not just an \emph{opaque declaration}. Compilation of
\lstinline!proctor.cpp! guarantees that the declaration and definition are
compatible. Having this \emph{forward declaration} improves (somewhat)
our type safety:

\begin{emcppslisting}[emcppsbatch=e4]
// proctor.h:
// ...
enum State : int;  // opaque declaration of enumeration (new in C++11)

class Proctor
{
    State d_current;  // opaque classical enumerated type (BETTER IDEA)
    // ...

public:
    Proctor();
    // ...
};
\end{emcppslisting}

\noindent But we can do even better. First we will want to nest the enumerated
\lstinline!State! type within the private section of the proctor to avoid
needless namespace pollution. What's more, because the numerical values
of the enumerators are not relevant, we can more closely model our
intent by nesting a more strongly typed \lstinline!enum!~\lstinline!class! instead:

\begin{emcppslisting}[emcppsbatch=e5]
// proctor.h:
// ...
class Proctor
{
    enum class State;  // forward (nested) declaration of type-safe enumeration
    State d_current;   // opaque (modern) enumerated data type (BEST IDEA)
    // ...

public:
    Proctor();
    // ...
};
\end{emcppslisting}

\noindent We would then declare the nested
\lstinline!enum!~\lstinline!class! %(see \featureref{\locationc}{enumclass})
accordingly in the
\lstinline!.cpp! file:

\begin{emcppslisting}[emcppsbatch=e5]
// proctor.cpp:
#include <proctor.h>

enum class Proctor::State { e_STARTING, e_READY, e_RUNNING, e_DONE };
Proctor::Proctor() : d_current(State::e_STARTING) { /* ... */ }
// ...
\end{emcppslisting}

\noindent Finally, notice that in this example we first \emph{forward declared}
the nested \lstinline!enum!~\lstinline!class! type within class scope and
then in a separate statement defined a data member of the opaque
enumerated type. We needed to do this in two statements because
simultaneously \emph{opaquely declaring} either a classic \lstinline!enum!
having an explicit underlying type %(see \featureref{\locationc}{explicit-enumeration-underlying-type})
or an
\lstinline!enum!~\lstinline!class! %(see \featureref{\locationc}{enumclass})
and also defining an object
of that type in a single statement is not possible:

\begin{emcppslisting}
enum E1 : int e1;  // Error, syntax not supported
enum class E2 e2;  // Error,    "    "     "
\end{emcppslisting}

\noindent Fully defining an enumeration and simultaneously defining an object of
the type in one stroke is, however, possible:

\begin{emcppslisting}
enum E3 : int { e_A, e_B } e3;  // OK, full type definition + object definition
enum class E4 { e_A, e_B } e4;  // OK,  "    "       "      "   "        "
\end{emcppslisting}

\noindent Providing such a full definition, however, would have run counter to our
intention to \romeogloss{insulate} the enumerator list of
\lstinline!Proctor::State! from clients \lstinline!#include!ing the header
file defining \lstinline!Proctor!.

\subsubsection[Cookie: Insulating all external clients from the enumerator list]{Cookie: Insulating all external clients from the enumerator list}\label{cookie:-insulating-all-external-clients-from-the-enumerator-list}

A commonly recurring \romeogloss{design pattern}, commonly known as the
``Memento pattern,"\footnote{\cite{gamma95}, Chapter 5, section ``Memento,'' pp.~283--???} manifests when a facility providing a service, often in a
multi-client environment, hands off a packet of opaque information ---
a.k.a. a \romeogloss{cookie} --- to a client to hold and then present back to
the facility to enable resuming operations where processing left off.
Since the information within the cookie will not be used substantively
by the client, any unnecessary compile-time coupling of clients with the
implementation of that cookie serves only to impede fluid
maintainability of the facility issuing the cookie. With respect
to not just \emph{encapsulating} but \emph{insulating} pure
implementation details that are held but not used substantively by
clients, we offer this Memento pattern as a possible use case
for \romeogloss{opaque enumerations}.

Event-driven programming,\footnote{See also \cite{gamma95}, Chapter 5, section ``Observer,''  pp.~293--???} historically implemented using \romeogloss{callback
functions}, introduces a style of programming that is decidedly
different from that to which we might have become accustomed. In this
programming paradigm, a higher-level agent (e.g., \lstinline!main!) would
begin by instantiating an \lstinline!Engine! that will be responsible for
monitoring for events and invoking provided callbacks when appropriate.
Classically, clients might have registered a function pointer and a
corresponding pointer to a client-defined piece of identifying data, but
here we will make use of a C++11 Standard Library type,
\lstinline!std::function!, which can encapsulate arbitrary callable
function objects and their associated state. This callback will be
provided one object to represent the event that just happened and
another object that can be used opaquely to reregister interest in the
same event again.

This opaque cookie and passing around of the client state might seem
like an unnecessary step, but often the event management involved in
software of this sort is wrapping the most often executed code in very
busy systems, and performance of each basic operation is therefore very
important. To maximize performance, every potential branch or lookup in
an internal data structure must be minimized, and allowing clients to
pass back the internal state of the engine when reregistering can
greatly reduce the engine's work to continue a client's processing of
events without tearing down and rebuilding all client state each time an
event happens. More importantly, event managers such as this often
become highly concurrent to take advantage of modern hardware, so
performant manipulation of their own data structures and well-defined
lifetime of the objects they interact with become paramount. This makes
the simple guarantee of, ``If you don't reregister, then the engine will
clean everything up; if you do, then the callback function will continue
its lifetime,'' a very tractable paradigm to follow.

\begin{emcppslisting}[emcppsbatch=e6]
// callbackengine.h:
#include <deque>          // (ù{\codeincomments{std::deque}}ù)
#include <functional>     // (ù{\codeincomments{std::function}}ù)

class EventData;       // information that clients will need to process an event
class CallbackEngine;  // the driver for processing and delivering events

class CallbackData
{
    // This class represents a handle to the shared state associating a
    // callback function object with a (ù{\codeincomments{CallbackEngine}}ù).

public:
    typedef std::function<void(const EventData&, CallbackEngine*,
        CallbackData)> Callback;
        // alias for a function object returning (ù{\codeincomments{void}}ù) and taking, as arguments,
        // the event data to be consumed by the client, the address of the
        // (ù{\codeincomments{CallbackEngine}}ù) object that supplied the event data, and the
        // callback data that can be used to reregister the client, should the
        // client choose to show continued interest in future instances of the
        // same event

    enum class State;  // GOOD IDEA
        // nested forward declaration of insulated enumeration, enabling
        // changes to the enumerator list without forcing clients to recompile

private:
    // ... (a smart pointer to an opaque object containing the state and the
    //      callback to invoke)

public:
    CallbackData(const Callback &cb, State init);

    // ... (constructors, other manipulators and accessors, etc.)

    State getState() const;
        // Return the current (ù{\codeincomments{state}}ù) of this callback.

    void setState(State state) const;
        // Set the current state to the specified (ù{\codeincomments{state}}ù).

    Callback& getCallback() const;
        // Return the (ù{\codeincomments{callback}}ù) function object specified at construction.
};

class CallbackEngine
{
private:
    // ... (other, stable (ù{\codeincomments{private}}ù) data members implementing this object)

    bool d_running;  // active state

    std::deque<CallbackData> d_pendingCallbacks;
        // The collection of clients currently registered for interest, or having
        // callbacks delivered, with this (ù{\codeincomments{CallbackEngine}}ù).
        //
        // Reregistering or skipping reregistering when
        // called back will lead to updating internal data structures based on
        // the current value of this (ù{\codeincomments{State}}ù).

public:
    // ...   (other public member functions, e.g., creators, manipulators)

    void registerInterest(CallbackData::Callback cb);
        // Register (e.g., from (ù{\codeincomments{main}}ù)) a new client with this manager object.

    void reregisterInterest(const CallbackData& callback);
        // Reregister (e.g., from a client) the specified (ù{\codeincomments{callback}}ù) with this
        // manager object, providing the (ù{\codeincomments{state}}ù) contained in the (ù{\codeincomments{CallbackData}}ù)
        // to enable resumption from the same state as processing left off.

    void run();
        // Start this object's event loop.

    // ...  (other public member functions, e.g., manipulators, accessors)
};
\end{emcppslisting}

\noindent A client would, in \lstinline!main!, create an instance of this
\lstinline!CallbackEngine!, define the appropriate functions to be invoked
when events happen, register interest, and then let the engine
\lstinline!run!:

\begin{emcppslisting}[emcppsbatch=e6]
// myapplication.cpp:
// ...
#include <callbackengine.h>

static void myCallback(const EventData&    event,
                       CallbackEngine*     engine,
                       const CallbackData& cookie);
    // Process the specified (ù{\codeincomments{event}}ù), and then potentially reregister the
    // specified (ù{\codeincomments{cookie}}ù) for interest in the same data.

int main()
{
   CallbackEngine e;  // Create a configurable callback engine object.

   //...     (Configure the callback engine, (ù{\codeincomments{e}}ù), as appropriate.)

   e.registerInterest(&myCallback);  // Even a stateless function pointer can
                                     // be used with (ù{\codeincomments{std::function}}ù).

   // ...create and register other clients for interest...

   e.run();    // Cede control to (ù{\codeincomments{e}}ù)'s event loop until complete.

   return 0;
}
\end{emcppslisting}

\noindent The implementation of \lstinline!myCallback!, in the example below, is then
free to reregister interest in the same event, save the cookie elsewhere
to reregister at a later time, or complete its task and let the
\lstinline!CallbackEngine! take care of properly cleaning up all now
unnecessary resources:

\begin{emcppshiddenlisting}[emcppsbatch=e6]
// myapplication.cpp:
struct EventProcessor {
    static int processEvent(const EventData &event);
    static void storeCallback(CallbackEngine *, const CallbackData&);    
};
\end{emcppshiddenlisting}
\begin{emcppslisting}[emcppsbatch=e6,emcppsoutputfile=myapplication.cpp]
void myCallback(const EventData&     event,
                CallbackEngine      *engine,
                const CallbackData&  cookie)
{
    int status = EventProcessor::processEvent(event);

    if (status > 0)  // status is non-zero; continue interest in event
    {
        engine->reregisterInterest(cookie);
    }
    else if (status < 0)  // Negative status indicates (ù{\codeincomments{EventProcessor}}ù) wants
                          // to reregister later.
    {
        EventProcessor::storeCallback(engine,cookie);
                          // Call (ù{\codeincomments{reregisterInterest}}ù) later.
    }

    // Return flow of control to the (ù{\codeincomments{CallbackEngine}}ù) that invoked this
    // callback.  If (ù{\codeincomments{status}}ù) was zero, then this callback should be cleaned
    // up properly with minimal fuss and no leaks.
}
\end{emcppslisting}

\noindent What makes use of the \romeogloss{opaque enumeration} here especially apt is
that the internal data structures maintained by the
\lstinline!CallbackEngine! might be very subtly interrelated, and any
knowledge of a client's relationship to those data structures that can
be maintained through callbacks is going to reduce the amount of lookups
and synchronization that would be needed to correctly reregister a
client without that information. The otherwise wide contract on
\lstinline!reregisterInterest! means that clients have no need themselves
to directly know anything about the actual values of the \lstinline!State!
they might be in. More notably, a component like this is likely to be
very heavily reused across a large codebase, and being able to maintain
it while minimizing the need for clients to recompile can be a huge boon
to deployment times.

To see what is involved, we can consider the business end of the
\lstinline!CallbackEngine! implementation and an outline of what a
single-threaded implementation might involve:

\begin{emcppshiddenlisting}[emcppsbatch=e6]
// callbackengine.cpp:
class EventData
{
};
EventData getNextEvent();
\end{emcppshiddenlisting}
\begin{emcppslisting}[emcppsbatch=e6]
// callbackengine.cpp:
#include <callbackengine.h>

enum class CallbackData::State
{
    // Full (local) definition of the enumerated states for the callback engine.
    e_INITIAL,
    e_LISTENING,
    e_READY,
    e_PROCESSING,
    e_REREGISTERED,
    e_FREED
};

void CallbackEngine::registerInterest(CallbackData::Callback cb)
{
    // Create a (ù{\codeincomments{CallbackData}}ù) instance with a state of (ù{\codeincomments{e\_INITIAL}}ù) and
    // insert it into the set of active clients.
    d_pendingCallbacks.push_back(CallbackData(cb, CallbackData::State::e_INITIAL));
}

void CallbackEngine::run()
{
    // Update all client states to (ù{\codeincomments{e\_LISTENING}}ù) based on the events in which
    // they have interest.

    d_running = true;
    while (d_running)
    {
        // Poll the operating system API waiting for an event to be ready.
        EventData event = getNextEvent();

        // Go through the elements of (ù{\codeincomments{d\_pendingCallbacks}}ù) to deliver this
        // event to each of them.
        std::deque<CallbackData> callbacks = std::move(d_pendingCallbacks);

        // Loop once over the callbacks we are about to notify to update their
        // state so that we know they are now in a different container.
        for (CallbackData& callback : callbacks)
        {
            callback.setState(CallbackData::State::e_READY);
        }

        while (!callbacks.empty())
        {
            CallbackData callback = callbacks.front();
            callbacks.pop_front();

            // Mark the callback as processing and invoke it.
            callback.setState(CallbackData::State::e_PROCESSING);

            callback.getCallback()(event, this, callback);

            // Clean up based on the new (ù{\codeincomments{State}}ù).
            if (callback.getState() == CallbackData::State::e_REREGISTERED)
            {
                // Put the callback on the queue to get events again.
                d_pendingCallbacks.push_back(callback);
            }
            else
            {
                // The callback can be released, freeing resources.
                callback.setState(CallbackData::State::e_FREED);
            }
        }
    }
}

void CallbackEngine::reregisterInterest(const CallbackData& callback)
{
    if (callback.getState() == CallbackData::State::e_PROCESSING)
    {
        // This is being called reentrantly from (ù{\codeincomments{run()}}ù); simply update state.
        callback.setState(CallbackData::State::e_REREGISTERED);
    }
    else if (callback.getState() == CallbackData::State::e_READY)
    {
        // This callback is in the deque of callbacks currently having events
        // delivered to it; do nothing and leave it there.
    }
    else
    {
        // This callback was saved; set it to the proper state and put it in
        // the queue of callbacks.
        if (d_running)
        {
            callback.setState(CallbackData::State::e_LISTENING);
        }
        else
        {
            callback.setState(CallbackData::State::e_INITIAL);
        }

        d_pendingCallbacks.push_back(callback);
    }
}

\end{emcppslisting}

\noindent Note how the definition of \lstinline!CallbackData::State! is visible and
needed only in this implementation file. Also, consider that the set of
states might grow or shrink as this \lstinline!CallbackEngine! is optimized
and extended, and clients can still pass around the object containing
that state in a type-safe manner while remaining insulated from this
definition.

Prior to C++11, we could not have \emph{forward declared} this
enumeration, and so would have had to represent it in a
\emph{type}-\emph{unsafe} way --- e.g., as an \lstinline!int!. Thanks to
the modern \lstinline!enum!~\lstinline!class! (see \featureref{\locationc}{enumclass}), however, we can
conveniently \romeogloss{forward declare} it as a nested type and then,
separately, fully define it inside the \lstinline!.cpp! implementing other
noninline member functions of the \lstinline!CallbackEngine! class. In
this way, we are able to \emph{insulate} changes to the enumerator list
along with any other aspects of the implementation defined outside of
the \lstinline!.h! file without forcing any client applications to
recompile. Finally, the basic design of the hypothetical \lstinline!CallbackEngine!
in the previous code example could have been used for any number of
useful components: a parser or tokenizer, a workflow engine, or even a
more generalized event loop.

\subsubsection[Dual-Access: Insulating some external clients from the enumerator list]{Dual-Access: Insulating some external clients from the enumerator list}\label{dual-access:-insulating-some-external-clients-from-the-enumerator-list}

In previous use cases, the goal has been to \romeogloss{insulate} \emph{all}
external clients from the enumerators of an enumeration that is visible
(but not necessarily programmatically reachable) in the defining
component's header. Consider the situation in which a \romeogloss{component}
(\lstinline!.h!/\lstinline!.cpp! pair) itself defines an enumeration that will
be used by various clients within a single program, some of which will
need access to the enumerators.

When an \lstinline!enum!~\lstinline!class! or a classic \lstinline!enum! having
an explicitly specified underlying type (see \featureref{\locationc}{explicit-enumeration-underlying-type}) is specified in
a header for direct programmatic use, external clients are at liberty to
unilaterally redeclare it \emph{opaquely}, i.e., without its enumerator
list. A compelling motivation for doing so would be for a client who
doesn't make direct use of the enumerators to \romeogloss{insulate} itself
and/or its clients from having to recompile whenever the enumerator list
changes.

Embedding any such \romeogloss{local declaration} in client code, however,
would be highly problematic: If the underlying type of the declaration
(in one translation unit) were somehow to become inconsistent with that
of the definition (in some other translation unit), any program
incorporating both translation units would immediately become silently
\romeogloss{ill-formed, no diagnostic required (IFNDR)}; see \intraref{potential-pitfalls-opaqueenum}{redeclaring-an-externally-defined-enumeration-locally}. Unless a separate ``forwarding'' header file is
provided along with (and ideally included by) the header defining the
full enumeration, any client opting to exploit this
opacity feature of an enumerated type will have no alternative but to
redeclare the enumeration locally; see \intraref{potential-pitfalls-opaqueenum}{inciting-local-enumeration-declarations-an-attractive-nuisance}.

For example, consider an \lstinline!enum!~\lstinline!class!, \lstinline!Event!,
intended for public use by external clients:

\begin{emcppslisting}
// event.h:
// ...
enum class Event : char { /*... changes frequently ...*/ };
// ...             ^^^^
\end{emcppslisting}

\noindent Now imagine some client header file, \lstinline!badclient.h!, that makes
use of the \lstinline!Event! enumeration and chooses to avoid compile-time
coupling itself to the enumerator list by
embedding, for whatever reason, a \romeogloss{local declaration} of \lstinline!Event! instead:

\begin{emcppslisting}
// badclient.h:
// ...
enum class Event : char;  // BAD IDEA: local external declaration
// ...
struct BadObject
{
    Event d_currentEvent;  // object of locally declare enumeration
    // ...
};
// ...
\end{emcppslisting}

\noindent Imagine now that the number of events that can fit in a \lstinline!char! is
exceeded and we decide to change the definition to have an underlying
type of \lstinline!short!:

\begin{emcppslisting}
// event.h:
// ...
enum class Event : short { /*... changes frequently ...*/ };
// ...             ^^^^^
\end{emcppslisting}

\noindent Client code, such as in \lstinline!badclient.h!, that fails to include the
\lstinline!event.h! header will have no automatic way of knowing that it
needs to change, and recompiling the code for all cases where
\lstinline!event.h! isn't also included in the translation unit will not
fix the problem. Unless every such client is somehow updated manually, a
newly linked program comprising them will be \romeogloss{IFNDR} with the likely
consequence of either a crash or (worse) when the program runs and
misbehaves. When providing a programmatically accessible definition of
an enumerated type in a header where the \romeogloss{underlying type} is
specified either explicitly or implicitly, we can give external
clients a \emph{safe} alternative to local declaration by also providing
an auxiliary header containing \emph{just} the corresponding opaque
declaration:

\begin{emcppslisting}[emcppsbatch=e7]
// event.fwd.h:
// ...
enum class Event : char;
// ...
\end{emcppslisting}

\noindent Here we have chosen to treat the forwarding header file as part of the same event component as the principal header but with an injected descriptive suffix field, \lstinline!.fwd!; this approach, as opposed to, say, \lstinline!file_fwd.h!, \lstinline!filefwd.h!, or \lstinline!file.hh!, was chosen so as not to (1) encroach on a disciplined, component-naming scheme involving reserved use of underscores\footnote{\cite{lakos20}, section~2.4, pp.~297--333} or (2) confuse tools and scripts that expect header-file names to end with a \lstinline!.h! suffix.

In general, having a forwarding header always included in its
corresponding full header facilitates situations such as default
template arguments where the declaration can appear at most once in any
given translation unit; the only drawback being that the
comparatively small forwarding header file must now also be opened and
parsed if the full header file is included in a given translation
unit. To ensure consistency, we thus \lstinline!#include! this forwarding
header in the original header defining the full enumeration:

\begin{emcppslisting}[emcppsbatch=e7]
// event.h:
// ...                  // Ensure opaque declaration (included here) is
#include <event.fwd.h>  // consistent with complete definition (below).
// ...
enum class Event : char { /*... changes frequently ...*/ };
// ...
\end{emcppslisting}

\noindent In this way, every translation unit that includes the definition will
serve to ensure that the forward declaration and definition match;
hence, clients can incorporate safely only the presumably more stable
forwarding header:

\begin{emcppslisting}[emcppsbatch=e7]
// goodclient.h:
// ...
#include <event.fwd.h>  // GOOD IDEA: consistent opaque declaration
// ...
class Client
{
    Event d_currentEvent;
    // ...
};
\end{emcppslisting}

\noindent Note that we have consistently
employed angle brackets exclusively for all include directives used
throughout this book to maximize flexibility of deployment presuming a
  regimen for unique naming.\footnote{See \cite{lakos20}, section~1.5.1,
  pp.~201--203.}

  To illustrate real-world practical use of the opaque-enumerations
feature, consider the various components{\cprotect\footnote{See
  \cite{lakos20}, sections~1.6 and 1.11, pp.~209--216 and
  pp.~256--259, respectively.}} that might depend
on{\cprotect\footnote{\cite{lakos20}, section~1.9?, pp.~237--243  JOHN: Please consult the book and correct. Section 1.9 is pp 243--251. Section 1.8 is pp. 237--243.}} an \lstinline!Event! enumeration such as that above:
\begin{itemize}
\item{\textbf{\romeogloss{\lstinline!message!}} — The component provides a \emph{value-semantic} \lstinline!Message! class consisting of just raw data,\cprotect\footnote{We sometimes refer to data that is meaningful only in the context of a higher-level entity as \romeogloss{dumb data}; see \cite{lakos20}, section~3.5.5, pp.~629--633.} including an \lstinline!Event! field representing the type of event. This component never makes direct use of any enumeration values and thus needs to include only \lstinline!event.fwd.h! and the corresponding opaque \emph{forward} declaration of the \lstinline!Event! enumeration.}
\item{\textbf{\romeogloss{\lstinline!sender!} and \romeogloss{\lstinline!receiver!}} — These are a pair of components that, respectively, create and consume \lstinline!Message! objects. To populate a \lstinline!Message! object, a \lstinline!Sender! will need to provide a valid value for the \lstinline!Event! member. Similarly, to process a \lstinline!Message!, a \lstinline!Receiver! will need to understand the potential individual enumerated values for the \lstinline!Event! field. Both of these components will include the primary \lstinline!event.h! header and thus have the complete definition of \lstinline!Event! available to them.}
\item{\textbf{\romeogloss{\lstinline!messenger!}} — The final component, a general engine capable of being handed \lstinline!Message! objects by a \lstinline!Sender! and then delivering those objects in an appropriate fashion to a \lstinline!Receiver!, needs a complete and usable definition of \lstinline!Message! objects — possibly copying them or storing them in containers before delivery — but has no need for understanding the possible values of the \lstinline!Event! member within those \lstinline!Message! objects. This component can participate fully and correctly in the larger system while being completely \emph{insulated} from the enumeration values of the \lstinline!Event! enumeration.}
\end{itemize}
tl;dr: By factoring out the \lstinline!Event! enumeration into its own
separate component and providing two distinct but compatible headers,
one containing the opaque declaration and the other (which includes the
first) providing the complete definition, we enable having different
components choose not to compile-time couple themselves with the
enumerator list without forcing them to (\emph{unsafely}) redeclare the
enumeration locally.

\subsection[Potential Pitfalls]{Potential Pitfalls}\label{potential-pitfalls-opaqueenum}

\subsubsection[Redeclaring an externally defined enumeration locally]{Redeclaring an externally defined enumeration locally}\label{redeclaring-an-externally-defined-enumeration-locally}

An opaque enumeration declaration enables the use of that
enumeration without granting visibility to its enumerators, reducing
physical coupling between components. Unlike a \romeogloss{forward class
declaration}, an opaque enumeration declaration produces a complete
type, sufficient for substantive use (e.g., via the linker):

\begin{emcppshiddenlisting}[emcppsbatch=e8]
// client.cpp:
#include <cstdint>  // (ù{\codeincomments{std::uint8\_t}}ù)
\end{emcppshiddenlisting}
\begin{emcppslisting}[emcppsbatch=e8]
// client.cpp:
enum Event : std::uint8_t;
Event e;  // OK, (ù{\codeincomments{Event}}ù) is a complete type.
\end{emcppslisting}

\noindent The \emph{underlying type} specified in an opaque \lstinline!enum!
declaration must exactly match the full definition; otherwise a program
incorporating both is automatically \romeogloss{IFNDR}. Updating an
\lstinline!enum!'s underlying type to accommodate additional values can
lead to latent defects when these changes are not propagated to all
local declarations:

\begin{emcppshiddenlisting}[emcppsbatch=e9]
// library.h:
#include <cstdint>  // (ù{\codeincomments{std::uint16\_t}}ù)
\end{emcppshiddenlisting}
\begin{emcppslisting}[emcppsbatch=e9]
// library.h:
enum Event : std::uint16_t { /* now more than 256 events */ };
\end{emcppslisting}

\noindent Consistency of a local opaque \lstinline!enum! declaration's underlying
type with that of its complete definition in a separate translation unit
cannot be enforced by the compiler, potentially leading to a program
that is \romeogloss{IFNDR}. In the \lstinline!client.cpp! example shown above, if the opaque
declaration in \lstinline!client.cpp! is not somehow updated to reflect the
changes in \lstinline!event.h!, the program will compile, link, and run,
but its behavior has silently become undefined. The only robust solution
to this problem is for \lstinline!library.h! to provide two separate header
files; see \intrarefsimple{inciting-local-enumeration-declarations-an-attractive-nuisance}.

The problem with local declarations is by no means limited to opaque
enumerations. Embedding a local declaration for any object whose use
might be associated with its definition in a separate translation unit
via just the linker invites instability:

\noindent\begin{minipage}{.45\textwidth}
\begin{emcppslisting}
// main.cpp:
extern int x;  // BAD IDEA!
// ...
\end{emcppslisting}
\end{minipage}\hfill
\begin{minipage}{.45\textwidth}
\begin{emcppslisting}
// library.cpp:
int x;
// ...
\end{emcppslisting}
\end{minipage}

\noindent The definition of object \lstinline!x! (in the code snippets above) resides
in the \lstinline!.cpp! file of the library component while a supposed
declaration of \lstinline!x! is embedded in the file defining
\lstinline!main!. Should the type of just the definition of \lstinline!x!
change, both translation units will continue to compile but, when
linked, the resulting program will be \romeogloss{IFNDR}:

\noindent\begin{minipage}{.65\textwidth}
\begin{emcppslisting}
// main.cpp:
extern int x;  // ILL-FORMED PROGRAM
// ...
\end{emcppslisting}
\end{minipage}\hfill
\begin{minipage}{.35\textwidth}
\begin{emcppslisting}
// library.cpp:
double x;
// ...
\end{emcppslisting}
\end{minipage}

\noindent To ensure consistency across translation units, the time-honored
tradition is to place, in a header file managed by the supplier, a
declaration of each external-linkage entity intended for use outside of
the translation unit in which it is defined; that header is then
included by both the supplier and each consumer:

\noindent\begin{minipage}{.3\textwidth}
\begin{emcppslisting}[emcppsbatch=e10]
// main.cpp:
#include <library.h>
// ...
\end{emcppslisting}
\end{minipage}\hfill
\begin{minipage}{.3\textwidth}
\begin{emcppslisting}[emcppsbatch=e10]
// library.h:
// ...
extern int x;
// ...
\end{emcppslisting}
\end{minipage}
\begin{minipage}{.3\textwidth}
\begin{emcppslisting}[emcppsbatch=e10]
// library.cpp:
#include <library.h>
int x;
// ...
\end{emcppslisting}
\end{minipage}

\noindent In this way, any change to the definition of \lstinline!x! in
\lstinline!library.cpp! --- the supplier --- will trigger a compilation error
when \lstinline!library.cpp! is recompiled, thereby forcing a corresponding
change to the declaration in \lstinline!library.h!. When that happens,
typical build tools will take note of the change in the header file's
timestamp relative to that of the \lstinline!.o! file corresponding to
\lstinline!main.cpp! --- the consumer --- and indicate that it too needs to be
recompiled. Problem solved.

The maintainability pitfall associated with opaque enumerations,
however, is qualitatively more severe than for other external-linkage
types, such as a global \lstinline!int!: (1) the full definition for the
enumeration type itself needs to reside in a header for \emph{any}
external client to make use of its individual enumerators and (2)
typical components consist of just a \lstinline!.h!/\lstinline!.cpp! pair,
i.e., exactly one \lstinline!.h! file and usually just one \lstinline!.cpp!
file.{\cprotect\footnote{\cite{lakos20}, sections~2.2.11--2.2.13,
  pp.~280--281}}

Exposing, within a library header file, an opaquely declarable
enumeration that is programmatically accessible by external clients
without providing some maintainable way for those clients to keep their
elided declarations in sync with the full definition introduces what we
are calling an \emph{attractive nuisance}: the client is forced to
choose between (a) introducing the added risk and maintenance burden of
having to manually maintain consistency between the underlying types for
all its separate opaque uses and the one full definition or (b)
forgo use of this opaque-enumeration feature entirely, forcing
gratuitous compile-time coupling with
the unused and perhaps unstable enumerators. At even moderate scale, excessive
compile-time coupling can adversely affect projects in ways that are far
more insidious than just increased compile times during development ---
e.g., any emergency changes that might need to occur and be deployed
quickly to production without forcing all clients to recompile and then
be retested and then, eventually, be rereleased.\footnote{For a complete
real-world example of how compile-time coupling can delay a ``hot fix''
by weeks, not just hours, see \cite{lakos20}, section~3.10.5,
pp.~783--789.}

\subsubsection[Inciting local enumeration declarations: an attractive nuisance]{Inciting local enumeration declarations: an attractive nuisance}\label{inciting-local-enumeration-declarations-an-attractive-nuisance}

Whenever we, as library component authors, provide the complete
definition of an \lstinline!enum!~\lstinline!class! or a classic enumeration
with an explicitly specified underlying type and fail to provide a
corresponding header having just the opaque declaration, we confront our
clients with the unenviable conundrum of whether to needlessly
compile-time couple themselves and/or their clients with the details of
the enumerator list or to make the dubious choice to unilaterally
redeclare that enumeration locally.

The problems associated with local declarations of data whose types are
maintained in separate translation units is not limited to enumerations;
see \intrarefsimple{redeclaring-an-externally-defined-enumeration-locally}. The maintainability pitfall associated with
\romeogloss{opaque enumerations}, however, is qualitatively more severe than
for other external-linkage types, such as a global \lstinline!int!, in that
the ability to elide the enumerators amounts to an \emph{attractive
nuisance} wherein a client --- wanting to do so and having access to
only a single header containing the unelided definition (i.e.,
comprising the enumeration name, underlying integral type, and
enumerator list) --- might be persuaded into providing an elided copy of
the \lstinline!enum!'s definition (i.e., one omitting just the enumerators)
locally!

Ensuring that library components that define
enumerations (e.g.,
\lstinline!enum!~\lstinline!class!~\lstinline!Event!) whose
enumerators can be elided also consistently provide a second
forwarding header file containing the opaque declaration of each such
enumeration (i.e., enumeration name and underlying integral type only)
would be one generally applicable way to sidestep this often
surprisingly insidious maintenance burden; see \intrarefsimple{dual-access:-insulating-some-external-clients-from-the-enumerator-list}. Note that the attractive
nuisance potentially exists even when the primary intent of the component is not to make the enumeration generally
available.\footnote{\cite{wight}}

\subsection[Annoyances]{Annoyances}\label{annoyances}

\subsubsection[Opaque enumerations are not completely type safe]{Opaque enumerations are not completely type safe}\label{opaque-enumerations-are-not-completely-type-safe}

Making an enumeration opaque does not stop it from being used to create
an object that is initialized opaquely to a zero value and then
subsequently used (e.g., in a function call):

\begin{emcppslisting}
enum Bozo : int;  // forward declaration of enumeration (ù{\codeincomments{Bozo}}ù)
void f(Bozo);     // forward declaration of function (ù{\codeincomments{f}}ù)

void g()
{
    Bozo clown{};
    f(clown);      // OK, who knows if zero is a valid value?!
}
\end{emcppslisting}

\noindent Though creating a zero-valued enumeration variable by default is not
new, allowing one to be created without even knowing what enumerated
values are valid is arguably dubious.

\subsection[See Also]{See Also}\label{see-also}

\begin{itemize}
\item{\seealsoref{explicit-enumeration-underlying-type}{\locationc}Discusses the underlying integral representation for enumeration variables and their values.}
\item{\seealsoref{enumclass}{\locationc}Introduces an implicitly scoped, more strongly typed enumeration.}
\end{itemize}

\subsection[Further Reading]{Further Reading}\label{further-reading}

\begin{itemize}
\item{For more on internal versus external linkage, see \cite{lakos20}, section~1.3.1, pp.~154--159.}
\item{For more on the use of header files to ensure consistency across translation units, see \cite{lakos20}, section~1.4, pp.~190--201, especially Figure~1-35, p.~197.}
\item{For more on the use of \lstinline!#include! directives and \lstinline!#include! guards, see \cite{lakos20}, section~1.5, pp.~201--209.}
\item{For a complete delineation of inherent properties that belong to every well-conceived \lstinline!.h!/\lstinline!.cpp! pair, see \cite{lakos20}, sections~1.6 and 1.11, pp.~219--216 and 256--259, respectively.}
\item{For an introduction to physical dependency, see \cite{lakos20}, section~1.8, pp.~237--243.}
\item{For a suggestion on how to achieve unique naming of files, see \cite{lakos20}, section~2.4, pp.~297--333.}
\item{For a thorough treatment of \romeogloss{architectural insulation}, see \cite{lakos20}, sections~3.10--3.11, pp.~773--835.}
\end{itemize}



\newpage
%\section[Range {\tt for}]{Range-Based {\SecCode for} Loops}%\label{}
% 26 March 2021, cleaned up laTeX conversion
% 27 March 2021, NEED info for citations; out to CE, Josh, external review
% 27 March 2021 JMB, code fixed
% 28 March 2021, out to SB for tech review


\emcppsFeature{
    short={Range \lstinline!for!},
    tocshort={Range {\TOCCode for}},
    long={Range-Based {\SecCode for} Loops},
    toclong={Range-Based \lstinline!for! Loops},
    rhshort={Range {\RHCode for}},
}{rangefor}\label{range-based-for-loops}
\setcounter{table}{0}
\setcounter{footnote}{0}
\setcounter{lstlisting}{0}
%\section[Range {\tt for}]{Range-Based {\SecCode for} Loops}%\label{}
%\subsection[Range-Based \tt{for} Loops]{Range-Based {\SubsecCode for} Loops}\label{range-based-for-loops}

Range-Based \lstinline!for! loops provide a simplified, more compact,
syntax for iterating through a range of elements.

\subsection[Description]{Description}\label{description-rangefor}

Iterating over the elements of a collection is a fundamental operation
usually performed with a \lstinline!for! loop:

\begin{emcppslisting}[emcppsbatch=e1]
#include <vector>  // (ù{\codeincomments{std::vector}}ù)
#include <string>  // (ù{\codeincomments{std::string}}ù)

void f1(const std::vector<std::string>& vec)
{
    for (std::vector<std::string>::const_iterator i = vec.begin();
         i != vec.end(); ++i)
    {
        // ...
    }
}
\end{emcppslisting}
    

\noindent The code above iterates over the strings in a \lstinline!std::vector!. It
is significantly more verbose than similar code in other languages
because it uses a very general-purpose construct, the \lstinline!for! loop,
to perform the specialized but common task of traversing a collection.
In C++11, the definition of \lstinline!i! can be simplified somewhat by
using \lstinline!auto!:

\begin{emcppslisting}[emcppsbatch=e1]
void f2(const std::vector<std::string>& vec)
{
    for (auto i = vec.begin(); i != vec.end(); ++i)
    {
        const std::string& s = *i;
        // ...
    }
}
\end{emcppslisting}
    

\noindent Although \lstinline!auto! does have a number of potential pitfalls, this
use of \lstinline!auto! to deduce the return type of \lstinline!vec.begin()!
is one of the safer idiomatic uses; see \featureref{\locationc}{auto}. While this
version of the loop is simpler to write, it still uses the fully
general, three-part \lstinline!for! construct. Moreover, it evaluates
\lstinline!vec.end()! each time through the loop.

The C++11 \emcppsgloss[range based for loop]{ranged-based \lstinline!for! loop} (sometimes colloquially
referred to as the ``foreach'' loop) is a more concise loop notation
tuned for traversing the elements of a container or other sequential
range. A \emcppsgloss[range based for loop]{ranged-based \lstinline!for! loop} works with \emph{ranges}
and \emph{elements} rather than \emph{iterators} or \emph{indexes}:

%   (LABEL: ``iterating over \lstinline!std::vector! with a range-based \lstinline!for! loop'')
\begin{emcppslisting}[emcppsbatch=e1,label={iterating-over-std::vector-with-a-range-based-for-loop}]
void f3(const std::vector<std::string>& vec)
{
    for (const std::string& s : vec)
    {
        // ...
    }
}
\end{emcppslisting}
    

\noindent The loop in the above example can be read as ``for each element
\lstinline!s! in \lstinline!vec! \ldots''. There is no need to specify the
name or type of the iterator, the loop-termination condition, or the
increment clause; the syntax is focused purely on yielding (in order)
each element of the collection for processing within the body of the
loop. {\cprotect\footnote{In C++20, the syntax has been enhanced
slightly to permit an optional leading variable declaration clause,
e.g.,
  \lstinline!for!~\lstinline!(std::lock_guard!~\lstinline!g(myMutex);!~\lstinline!const!~\lstinline!std::string&!~\lstinline!s!~\lstinline!:!~\lstinline!vec)!~\lstinline!{!~\lstinline!/*!~\lstinline!...!~\lstinline!*/!~\lstinline!}!.}}

\subsubsection[Specification]{Specification}\label{specification}

The syntax for a \emcppsgloss[range based for loop]{ranged-based \lstinline!for! loop} declares a loop
variable and specifies a range of elements to be traversed:

%%%% LORI: please make the code below between asterisks in italics

\begin{emcppslisting}[emcppsignore={grammar, not code, but leave as c++}]
for ( *(ù{\it{for-range-declaration}}ù)* : *(ù{\it{range-expression}}ù)* ) *(ù{\it{statement}}ù)*
\end{emcppslisting}
    

\noindent The compiler treats this high-level construct as though it were
transformed into a lower-level \lstinline!for! loop with the following
pseudocode:

%%%%%% LORI: please make \emph{for-range-declaration}, \emph{range-expression}
%and \emph{statement} in the code snippet below in italics

\begin{emcppslisting}[emcppsignore={grammar, not code, but leave as c++}]
{
    auto&& __range = (ù{\it{range-expression}}ù);
    for (auto __begin = begin-expr, __end = end-expr;
        __begin != __end;
        ++__begin)
    {
        (ù{\it{for-range-declaration}}ù) = *__begin;
        (ù{\it{statement}}ù)
    }
}
\end{emcppslisting}
    

\noindent The variables \lstinline!__range!, \lstinline!__begin!, and
\lstinline!__end!, above, are for \emph{exposition only}, i.e., the
compiler does not necessarily generate variables with those names and
user code is not permitted to access those variables directly.

The \lstinline!__range! variable is defined as a \emcppsgloss{forwarding
reference} (see {\featureref{\locationc}{forwardingref}); it
will bind to any type of \emcppsgloss{range expression}, regardless of its
\emcppsgloss{value category} (\romeovalue{lvalue} or \romeovalue{rvalue}). If the
\emcppsgloss{range expression} yields a temporary object, its lifetime is
extended, if necessary, until \lstinline!__range! goes out of scope.
While this \emcppsgloss{lifetime extension} of temporary objects works in
most cases, it is insufficient when \lstinline!__range! doesn't bind
directly to the temporary created by the \emcppsgloss{range
expression}, potentially resulting in subtle bugs; see {Potential
Pitfalls}{Lifetime of temporaries in the range expression}.

The \emph{begin-expr} and \emph{end-expr} expressions used to initialize
the \lstinline!__begin! and \lstinline!__end! variables, respectively,
define a half-open range of elements starting with \lstinline!*__begin!
and including all of the elements in the \lstinline!__range! up to but
not including \lstinline!*__end!. The precise meaning of
\emph{begin-expr} and \emph{end-expr} were clarified in C++14 but were
essentially the same in C++11{\cprotect\footnote{The rules for
interpreting \emph{begin-expr} and \emph{end-expr} were slightly
unclear in C++11. A defect report, CWG 1442, clarified the wording
  retroactively. C++14 clarified the wording further.}}:

\begin{itemize}
\item{If \lstinline!__range! refers to an array, then \emph{begin-expr} is the address of the first element of the array and \emph{end-expr} is the address of one past the last element of the array.}
\item{If \lstinline!__range! refers to a class object and \lstinline!begin! and/or \lstinline!end! are members of that class, then \emph{begin-expr} is \lstinline!__range.begin()! and \emph{end-expr} is \lstinline!__range.end()!. Note that if \lstinline!begin! or \lstinline!end! are found in the class, then both of these expressions must be valid or else the program is ill formed.}
\item{Otherwise, \emph{begin-expr} is \lstinline!begin(__range)! and \emph{end-expr} is \lstinline!end(__range)!, where \lstinline!begin! and \lstinline!end! are found using \emcppsgloss{argument-dependent lookup (ADL)}. Note that \lstinline!begin! and \lstinline!end! are looked up only in the namespaces associated with the expressions; names that are local to the context of the \emcppsgloss[range based for loop]{ranged-based \lstinline!for! loop} are not considered; see Annoyances.}
\end{itemize}

Thus, a container such as \lstinline!vector!, with conventional
\lstinline!begin! and \lstinline!end! member functions, provides everything
necessary for a \emcppsgloss[range based for loop]{ranged-based \lstinline!for! loop}, as we saw in the
\lstinline!f3! example on \pageref{iterating-over-std::vector-with-a-range-based-for-loop}. Note that \emph{end-expr} ---
\lstinline!__range.end()! in the case of the \lstinline!vector! --- is
evaluated only once, unlike the idiomatic low-level \lstinline!for! loop,
where it is evaluated prior to every iteration.

In C++11 and C++14, \lstinline!__begin! and \lstinline!__end! are required
to have the same type.{\cprotect\footnote{The C++17 Standard changes the
defining code transformation of the \emcppsgloss[range based for loop]{range-based \lstinline!for!
loop} so as to allow \lstinline!__begin! and \lstinline!__end! to have
different types as long as they are comparable using
  \lstinline!__begin!~\lstinline|!=|~\lstinline!__end!.}} Although the
\lstinline!__begin! and \lstinline!__end! variables look and act like
iterators, they need not conform to all of the iterator requirements in
the Standard. Specifically, the type of \lstinline!__begin! and
\lstinline!__end! must support prefix \lstinline!operator++! but not
necessarily postfix \lstinline!operator++!, and it must support
\lstinline|operator!=| but not necessarily \lstinline|operator==|.

The \emph{for-range-declaration} declares the loop variable. Any
declaration that can be initialized with \lstinline!*__begin! will work.
For instance, if \lstinline!*__begin! returns a reference to a modifiable
object of, e.g., \lstinline!int! type, then \lstinline!int!~\lstinline!j!,
\lstinline!int&!~\lstinline!j!, \lstinline!const!~\lstinline!int&!~\lstinline!j!, and
\lstinline!long!~\lstinline!j! would all be valid
\emph{for-range-declaration}s declaring a loop variable \lstinline!j!.
Alternatively, the type of the loop variable can be deduced using
\lstinline!auto! --- i.e., \lstinline!auto!~\lstinline!j!,
\lstinline!auto&!~\lstinline!j!, \lstinline!const!~\lstinline!auto&!~\lstinline!j!,
or \lstinline!auto&&!~\lstinline!j! (see {\featureref{\locationc}{auto}).

The sequence being traversed can be modified through the loop variable
only if \lstinline!__begin! returns a reference to a modifiable type and
the loop variable is similarly declared as a reference to a modifiable
type (e.g., \lstinline!int&!, \lstinline!auto&!, or \lstinline!auto&&!). Note
that the \emcppsgloss[for range declaration]{for-range declaration} must define a \emph{new}
variable; unlike a traditional \lstinline!for! loop, it cannot name an
existing variable already in scope:

\begin{emcppslisting}
#include <vector>  // (ù{\codeincomments{std::vector}}ù)

void f1(std::vector<int>& vec)
{
    const std::vector<int>& cvec = vec;

    for (auto&& i : cvec)
    {
        i = 0;  // Error, (ù{\codeincomments{i}}ù) is a reference to (ù{\codeincomments{const int}}ù).
    }

    for (int j : vec)
    {
        j = 0;  // Bug, (ù{\codeincomments{j}}ù) is a loop-local variable; (ù{\codeincomments{vec}}ù) is not modified.
    }

    for (int& k : vec)
    {
        k = 0;  // OK, set element of (ù{\codeincomments{vec}}ù) to 0.
    }

    int m;
    for (     m : vec) { /* ... */ }  // Error, (ù{\codeincomments{m}}ù) does not define a variable.
    for (int& m : vec) { /* ... */ }  // OK, loop (ù{\codeincomments{m}}ù) hides function-scope (ù{\codeincomments{m}}ù).
}
\end{emcppslisting}
    

\noindent Since \lstinline!cvec! is \lstinline!const!, the element type returned by
\lstinline!*begin(cvec)! is \lstinline!const!~\lstinline!int&!. Thus, \lstinline!i!
is deduced as \lstinline!const!~\lstinline!int&!, making invalid any attempt
to modify an element through \lstinline!i!. The second loop is valid C++11
code but has a subtle bug: \lstinline!j! is not a reference --- it contains
a \emph{copy} of the current element in the \lstinline!vector! --- so
modifying \lstinline!j! has no effect on the vector. The third loop
correctly sets all of the elements of \lstinline!vec! to zero; the loop
variable \lstinline!k! is a reference to the current element, so setting it
to zero modifies the original vector. The first \lstinline!m! loop attempts
to re-use local variable \lstinline!m! as a loop variable; while this would
be legal in a traditional \lstinline!for! loop, it is ill formed in a
\emcppsgloss[range based for loop]{ranged-based \lstinline!for! loop}. Finally, the last loop re-uses
the name \lstinline!m! from the surrounding scope, hiding the old name for
the duration of the loop, just as it would in a traditional \lstinline!for!
loop.

The \emph{statement} that makes up the loop body can contain anything
that is valid within a traditional \lstinline!for! loop body. In
particular, a \lstinline!break! statement will exit the loop immediately
and a \lstinline!continue! statement will skip to the next iteration.

Applying this transformation to the \lstinline!f3! example (see \pageref{iterating-over-std::vector-with-a-range-based-for-loop}) from the previous section, we can see how the
\emcppsgloss[range based for loop]{ranged-based \lstinline!for! loop} hooks into the iterator idiom to
traverse a \lstinline!vector! of \lstinline!string! elements:

\begin{emcppslisting}
#include <string>  // (ù{\codeincomments{std::string}}ù)
#include <vector>  // (ù{\codeincomments{std::vector}}ù)

void f3(const std::vector<std::string>& vec)
{
    // for (const std::string& s : vec) { /* ... */ }
    {
        auto&& __range = vec;  // reference to the (ù{\codeincomments{std::vector}}ù)
        for (auto __begin = begin(__range), __end = end(__range);
            __begin != __end;
            ++__begin)
        {
            const std::string& s = *__begin;  // Get current (ù{\codeincomments{string}}ù) element.
            {
                // ...
            }
        }
    }
}
\end{emcppslisting}
    

\noindent In this expansion, \lstinline!__range! has type
\lstinline!const!~\lstinline!std::vector<std::string>&! while
\lstinline!__begin! and \lstinline!__end! have type
\lstinline!std::vector<std::string>::const_iterator!.

\subsubsection[Traversing arrays and initializer lists]{Traversing arrays and initializer lists}\label{traversing-arrays-and-initializer-lists}

The \lstinline!<iterator>! standard header defines array overloads for
\lstinline!std::begin! and \lstinline!std::end! such that, when applied to a
C-style array, \lstinline!arr!, having a known number of elements,
\lstinline!__bound!, \lstinline!std::begin(arr)! returns the address of the
first element of \lstinline!arr! and \lstinline!std::end(arr)! returns the
address of one past that of the last element of \lstinline!arr!, i.e.,
\lstinline!arr!~\lstinline!+!~\lstinline!__bound!. This functionality is built
into the initialization of \lstinline!__begin! and \lstinline!__end! as a
special case, in the expansion of a \emcppsgloss[range based for loop]{range-based \lstinline!for!
loop}, so that it is possible to traverse the elements of an array
without needing to \lstinline!#include!~\lstinline!<iterator>!:

\begin{emcppslisting}
void f1()
{
    double data[] = {1.9, 2.8, 4.7, 7.6, 11.5, 16.4, 22.3, 29.2, 37.1, 46.0};
    for (double& d : data)
    {
        d *= 3.0;  // triple every element in the array
    }
}
\end{emcppslisting}
    

\noindent In the above example, the reference \lstinline!d! is bound, in turn, to
each element of the array. The size of the array is not encoded anywhere
in the loop syntax, either as a literal or as a symbolic value,
simplifying the specification of the loop and preventing errors. Note
that only arrays whose size is known at the point where the loop occurs
can be traversed this way:

\begin{emcppslisting}
extern double data[];  // array of unknown size

void f2()
{
    for (double& d : data)  // Error, (ù{\codeincomments{data}}ù) is an incomplete type.
    {
        // ...
    }
}

double data[10] = { /* ... */ };  // too late to make the above compile
\end{emcppslisting}
    

\noindent The above example would compile if \lstinline!data! were declared having a
size, e.g.,\linebreak[4] \lstinline!extern!~\lstinline!double!~\lstinline!data[10]!, as that
would be a complete type and provide sufficient information to traverse
the array. The second definition of \lstinline!data! in the example
\emph{is} complete but is not visible at the point that the loop is
compiled.

An \emcppsgloss{initializer list} is typically used to initialize an array or
container using \emcppsgloss{braced initialization}; see \featureref{\locationc}{bracedinit}. The \lstinline!initializer_list! template does,
however, provide its own \lstinline!begin! and \lstinline!end! member
functions and is, therefore, directly usable as the
\emph{range-expression} in a \emcppsgloss[range based for loop]{ranged-based \lstinline!for! loop}:

\begin{emcppslisting}
#include <initializer_list>  // (ù{\codeincomments{std::initializer\_list}}ù)

void f3()
{
    for (double v : {1.9, 2.8, 4.7, 7.6, 11.5, 16.4, 22.3, 29.2, 37.1, 46.0})
    {
        // ...
    }
}
\end{emcppslisting}
    

\noindent The example above shows how a series of \lstinline!double! values can be
embedded right within the loop header.

\subsection[Use Cases]{Use Cases}\label{use-cases}

\subsubsection[Iterating over all elements of a container]{Iterating over all elements of a container}\label{iterating-over-all-elements-of-a-container}

The motivating use case for this feature is to loop over the elements in
a container:

\begin{emcppslisting}
#include <list>  // (ù{\codeincomments{std::list}}ù)

void process(int* p);

void f1()
{
    std::list<int> aList{ 1, 2, 4, 7, 11, 16, 22, 29, 37, 46 };

    for (int& i : aList)
    {
        process(&i);
    }
}
\end{emcppslisting}
    

\noindent This idiom takes advantage of all STL-compliant container types
providing \lstinline!begin! and \lstinline!end! operations, which may be used
to delimit a range encompassing the entire container. Thus, the loop
above iterates from \lstinline!aList.begin()! to \lstinline!aList.end()!,
calling \lstinline!process! on each element encountered.

When iterating over a \lstinline!std::map<Key,!~\lstinline!Value>! or
\lstinline!std::unordered_map<Key,!~\lstinline!Value>!, each element has type
\lstinline!std::pair<const!~\lstinline!Key,!~\lstinline!Value>!. To save typing
and to avoid errors related to the first member of the pair being
\lstinline!const! we use the \lstinline!value_type! alias to refer to each
element's type:

\begin{emcppslisting}
#include <iostream>  // (ù{\codeincomments{std::cout}}ù)
#include <map>       // (ù{\codeincomments{std::map}}ù)
#include <string>    // (ù{\codeincomments{std::string}}ù)

using MapType = std::map<std::string, int>;

MapType studentScores
{
    {"Emily", 89},
    {"Joel",  85},
    {"Bud",   86},
};

void printScores()
{
    for (MapType::value_type& studentScore : studentScores)
    {
        const std::string& student = studentScore.first;
        int&               score   = studentScore.second;
        std::cout << student << "\t scored " << score << '\n';
    }
}
\end{emcppslisting}
    

\noindent This example prints each key/value pair in the \lstinline!map!. We create
two aliases, \lstinline!student! for \lstinline!studentScore.first! and
\lstinline!score! for \lstinline!studentScore.second!, to better express the
intent of the code.{\cprotect\footnote{In C++17, \emcppsgloss[structured binding]{structured
bindings} allow two variables to be initialized from a single
\lstinline!pair!, each variable being initialized by the respective
\lstinline!first! and \lstinline!second! members of the \lstinline!pair!. A
\emcppsgloss[range based for loop]{ranged-based \lstinline!for! loop} using a \emcppsgloss{structured
binding} for the loop variables yields a very clean and expressive way
to traverse containers like \lstinline!map! and \lstinline!unordered_map!,
e.g., using
  \lstinline!for!~\lstinline!(auto&!~\lstinline![student,!~\lstinline!score]!~\lstinline!:!~\lstinline!studentScores)!.}}

\subsubsection[Subranges]{Subranges}\label{subranges}

Using a classic \lstinline!for! loop to traverse a container, \lstinline!c!,
allows a subrange of \lstinline!c! to be specified beginning at some point
after \lstinline!c.begin()! --- e.g., \lstinline!++c.begin()! --- and/or ending at
some point before \lstinline!c.end()! --- e.g.,
\lstinline!std::next(c.end(),!~\lstinline!-3)!. To specify a subrange for a
\emcppsgloss[range based for loop]{ranged-based \lstinline!for! loop}, we create a simple adapter to
hold two iterators (or iterator-like objects) that define the desired
subrange:

\begin{emcppslisting}[emcppsbatch=e2]
template <typename Iter>
class Subrange
{
    Iter d_begin, d_end;

public:
    using iterator = Iter;

    Subrange(Iter b, Iter e) : d_begin(b), d_end(e) { }

    iterator begin() const { return d_begin; }
    iterator end()   const { return d_end;   }
};

template <typename Iter>
Subrange<Iter> makeSubrange(Iter beg, Iter end) { return {beg, end}; }
\end{emcppslisting}
    

\noindent The \lstinline!Subrange! class above is a primitive start to a potentially
rich library of range-based utilities.{\cprotect\footnote{The C++20
Standard introduces a new Ranges Library that provides powerful
  features for defining, combining, filtering, and manipulating ranges.}}
It holds two externally-supplied iterators that it can supply to a
\emcppsgloss[range based for loop]{ranged-based \lstinline!for! loop} via its \lstinline!begin! and
\lstinline!end! accessor members. The \lstinline!makeSubrange! factory uses
function template argument deduction to return a \lstinline!Subrange! of
the correct type.{\cprotect\footnote{C++17 introduced \emcppsgloss{class
template argument deduction}, which significantly reduces the need for
  factory templates like \lstinline!makeSubrange!.}}

Let's use \lstinline!Subrange! to traverse a vector in reverse, omitting
its first element:

\begin{emcppslisting}[emcppsbatch=e2]
#include <vector>    // (ù{\codeincomments{std::vector}}ù)
#include <iostream>  // (ù{\codeincomments{std::cout}}ù), (ù{\codeincomments{std::endl}}ù)

template <typename Range>
void printRange(const Range& r)
{
    for (const auto& elem : r)
    {
        std::cout << elem << ' ';
    }

    std::cout << std::endl;
}

std::vector<int> vec{16, 3, 1, 8, 99};

void f1()
{
    printRange(makeSubrange(vec.rbegin(), vec.rend() - 1));
        // print "99 8 1 3"
}
\end{emcppslisting}
    

\noindent The \lstinline!printRange! function template will print out the elements of
any range, provided the element type supports printing to a
\lstinline!std::ostream!. In \lstinline!f1!, we use reverse iterators to
create a \lstinline!Subrange! starting from the last element of
\lstinline!vec! and iterating backwards. By subtracting \lstinline!1! from
\lstinline!vec.rend()!, we exclude the last element of the sequence, which
is the first element of \lstinline!vec!.

In fact, the iterators need not refer to a container at all. For
example, we can use \lstinline!std::istream_iterator! to iterate over
``elements'' in an input stream:

\begin{emcppslisting}[emcppsbatch=e2]
#include <istream>   // (ù{\codeincomments{std::istream}}ù)
#include <iterator>  // (ù{\codeincomments{std::istream\_iterator}}ù)
#include <sstream>   // (ù{\codeincomments{std::istringstream}}ù)

void f2()
{
    std::istringstream inStream("1 2 4 7 11 16 22 29 37 bad 46");
    printRange(makeSubrange(std::istream_iterator<int>(inStream),
                            std::istream_iterator<int>()));
}
\end{emcppslisting}
    

\noindent In \lstinline!f2!, the range being printed uses the
\lstinline!istream_iterator<T>! adapter template. Each time through the
loop, the adapter reads another \lstinline!T! item from its input stream.
At end-of-file or if a read error occurs, the iterator becomes equal to
the sentinel iterator, \lstinline!istream_iterator<T>()!. Note that the
\emcppsgloss[range based for loop]{ranged-based \lstinline!for! loop} feature and the \lstinline!Subrange!
class template do not require that the size of the subrange be known in
advance.

\subsubsection[Range generators]{Range generators}\label{range-generators}

Iterating over a range does not necessarily entail traversing existing
data elements. A range expression could yield a type that
\emph{generates} elements as it goes. A useful example is the
\lstinline!ValueGenerator!, an iterator-like class that produces a sequence
of sequential values{\cprotect\footnote{The \lstinline!iota_view! and
\lstinline!iota! entities from the Ranges Library in the C++20 Standard
provide a more sophisticated version of the \lstinline!ValueGenerator!
  and \lstinline!valueRange! facility described here.}}:

\begin{emcppslisting}[emcppsbatch=e2]
template <typename T>
class ValueGenerator
{
    T d_value;

  public:
    explicit ValueGenerator(const T& v) : d_value(v) { }

    T operator*() const { return d_value; }
    ValueGenerator& operator++() { ++d_value; return *this; }

    friend bool operator!=(const ValueGenerator& a, const ValueGenerator& b)
    {
        return a.d_value != b.d_value;
    }
};

template <typename T>
Subrange<ValueGenerator<T>> valueRange(const T& b, const T& e)
{
    return { ValueGenerator<T>(b), ValueGenerator<T>(e) };
}
\end{emcppslisting}
    

\noindent Instead of referring to an element within a container,
\lstinline!ValueGenerator! is an iterator-like type that \emph{generates}
the value returned by \lstinline!operator*!. \lstinline!ValueGenerator! can be
instantiated for any type that can be incremented, e.g., integral types,
pointers, or iterators. The \lstinline!valueRange! function template is a
simple factory to create a range comprising two \lstinline!ValueGenerator!
objects, using the \lstinline!Subrange! class template defined in the use
case described in {Subranges}. Thus, to print the numbers from 1 to 10,
simply use a \emcppsgloss[range based for loop]{ranged-based \lstinline!for! loop}, employing a call to
\lstinline!valueRange! as the \emcppsgloss{range expression}:

\begin{emcppslisting}[emcppsbatch=e2]
void f3()
{
    // prints "1 2 3 4 5 6 7 8 9 10 "
    for (unsigned i : valueRange(1, 11))
    {
        std::cout << i << ' ';
    }
    std::cout << std::endl;
}
\end{emcppslisting}
    

\noindent Note that the second argument to \lstinline!valueRange! is one \emph{past}
the last item we want to iterate on, i.e., \lstinline!11! instead of
\lstinline!10!. With something like \lstinline!ValueGenerator! as part of a
reusable utility library, this formulation expresses the intent of the
loop more cleanly and concisely than the classic \lstinline!for! loop.

The ability to generate numbers means that a range need not be finite.
For example, we might want to generate a sequence of random numbers of
indefinite length:

\begin{emcppslisting}[emcppsbatch=e2]
#include <random>  // (ù{\codeincomments{std::default\_random\_engine}}ù), (ù{\codeincomments{std::uniform\_int\_distribution}}ù)

template <typename T = int>
class RandomIntSequence
{
    std::default_random_engine       d_generator;
    std::uniform_int_distribution<T> d_uniformDist;

public:
    class iterator
    {
        RandomIntSequence* d_sequence;

        explicit iterator(RandomIntSequence* s) : d_sequence(s) { }
        friend class RandomIntSequence;

    public:
        iterator& operator++() { return *this; }
        T operator*() const { return d_sequence->next(); }

        friend bool operator!=(iterator, iterator) { return true; }
    };

    RandomIntSequence(T min, T max, unsigned seed = 0)
        : d_generator(seed ? seed : std::random_device()())
        , d_uniformDist(min, max) { }

    T next() { return d_uniformDist(d_generator); }

    iterator begin() { return iterator(this); }
    iterator end()   { return iterator(nullptr); }
};

template <typename T>
RandomIntSequence<T> randomIntSequence(T min, T max, unsigned seed = 0)
{
    return {min, max, seed};
}
\end{emcppslisting}
    

\noindent The \lstinline!RandomIntSequence! class template uses the C++11
random-number library to generate high-quality pseudorandom
numbers.{\cprotect\footnote{An introduction to the C++11 random-number
library can be found in Stephan T. Lavavej's excellent talk,
  \lstinline!rand()! Considered Harmful.}} Each call to its \lstinline!next!
member function produces a new random number of integral type,
\lstinline!T!, within the inclusive range specified to the
\lstinline!RandomIntSequence! constructor. The nested \lstinline!iterator!
type holds a pointer to a \lstinline!RandomIntSequence! and simply calls
\lstinline!next! each time it is dereferenced (i.e., via a call to
\lstinline!operator*!).

Of particular interest is \lstinline|operator!=|, which returns
\lstinline!true! when comparing any two
\lstinline!RandomIntSequence<T>::iterator! objects. Thus, any
\emcppsgloss[range based for loop]{ranged-based \lstinline!for! loop} that iterates over a
\lstinline!RandomIntSequence! is an infinite loop unless it terminates by
some other means:

\begin{emcppslisting}[emcppsbatch=e2]
void f4()
{
    for (int rand : randomIntSequence(1, 10))
    {
        std::cout << rand << ' ';
        if (rand == 10) { break; }
    }

    std::cout << std::endl;
}
\end{emcppslisting}
    

\noindent This example prints a list of random numbers in the range 1 through 10,
inclusive. The loop terminates after printing \lstinline!10! for the first
(and only) time.

\subsubsection[Iterating over simple values]{Iterating over simple values}\label{iterating-over-simple-values}

The ability to iterate over a \lstinline!std::initializer_list! can be
very useful for processing a list of simple values or simple objects
without first storing them in a container. Such a use case arises
frequently when testing:

\begin{emcppslisting}
#include <limits>            // (ù{\codeincomments{std::numeric\_limits}}ù)
#include <initializer_list>  // (ù{\codeincomments{std::initializer\_list}}ù)

#define TEST_ASSERT(expr)  // ... assert that (ù{\codeincomments{expr}}ù) is true.

bool isEven(int i)
{
    return i % 2 == 0;
}

void testIsEven()
{
    // ...

    const int minInt = std::numeric_limits<int>::min();
    const int maxInt = std::numeric_limits<int>::max();

    for (int testValue : {minInt, -256, -2, 0, 2, 4, maxInt - 1})
    {
        TEST_ASSERT(isEven(testValue));
        TEST_ASSERT(!isEven(testValue + 1));
    }
}
\end{emcppslisting}
    

\noindent The \lstinline!testIsEven! function iterates over a sample of numbers
within the domain of \lstinline!isEven!, including boundary conditions,
testing that each number is correctly reported as being even and that
adding 1 to the number produces a result that is correctly reported as
not being even.

Initializer lists are not limited to primitive types, so the test data
set can contain more complex values:

\begin{emcppslisting}
#include <initializer_list>  // (ù{\codeincomments{std::initializer\_list}}ù)

#define TEST_ASSERT_EQ(expr1, expr2)  // ... assert that (ù{\codeincomments{expr1 == expr2}}ù).

int half(int i)
{
    return i / 2;
}

struct TestCase
{
    int value;
    int expected;
};

void testHalf()
{
    for (TestCase test : std::initializer_list<TestCase>{
        {-2, -1}, {-1, 0}, {0, 0}, {1, 0}, {2, 1}
    })
    {
        TEST_ASSERT_EQ(test.expected, half(test.value));
    }
}
\end{emcppslisting}
    

\noindent In this case, the \emcppsgloss[range based for loop]{ranged-based \lstinline!for! loop} iterates over a
\lstinline!std::initializer_list! holding \lstinline!TestCase! structures.
This paring of input(s) with expected output(s) of a component under
test is very common in unit tests.

\subsection[Potential Pitfalls]{Potential Pitfalls}\label{potential-pitfalls}

\subsubsection[Lifetime of temporary objects in the range expression]{Lifetime of temporary objects in the range expression}\label{lifetime-of-temporary-objects-in-the-range-expression}

As described in \intrarefsimple{description-rangefor} section, if the \emcppsgloss{range
expression} evaluates to a temporary object, that object remains valid,
as a result of \emcppsgloss{lifetime extension}, for the duration of the
\emcppsgloss[range based for loop]{ranged-based \lstinline!for! loop}. Unfortunately, there are some
subtle ways in which \emcppsgloss{lifetime extension} is not always
sufficient.

The basic notion of \emcppsgloss{lifetime extension} is that, when bound to a
reference, the lifetime of a \romeovalue{prvalue} --- i.e., an object
created by a literal, constructed in place, or returned (by value) from
a function --- is \emph{extended} to match the lifetime of the reference
to which it is bound:

\begin{emcppslisting}[emcppsbatch=e3]
#include <string>  // (ù{\codeincomments{std::string}}ù)

std::string strFromInt(int);

void f1()
{
    const std::string& s1 = std::string('a', 2);
    std::string&&      s2 = strFromInt(9);
    auto&&             i  = 5;

    // (ù{\codeincomments{s1}}ù), (ù{\codeincomments{s2}}ù), and (ù{\codeincomments{i}}ù) are "live" here.

    // ...

}  // (ù{\codeincomments{s1}}ù), (ù{\codeincomments{s2}}ù), and (ù{\codeincomments{i}}ù) are destroyed at end of enclosing block.
\end{emcppslisting}
    

\noindent The first string is constructed in place. The resulting temporary string
would normally be destroyed as soon as the expression was complete but,
because it is bound to a reference, its lifetime is extended; its
destructor is not called and its memory footprint is not reused until
\lstinline!s1! goes out of scope, i.e., at the end of the enclosing block.
The \lstinline!strFromInt! function returns by value; the result of calling
it in the second statement produces a temporary variable whose lifetime
is similarly extended until \lstinline!s2! goes out of scope. Finally, the
\emcppsgloss{forwarding reference}, \lstinline!i!, ensures that space in the
current frame is allocated to hold the temporary copy of the (deduced)
\lstinline!int! value, \lstinline!5!; such space cannot be reused until
\lstinline!i! goes out of scope at the end of the enclosing block. (See
{\featureref{\locationc}{forwardingref}).

When the \emcppsgloss{range expression} for a \emcppsgloss[range based for loop]{range-based
\lstinline!for! loop} is a \romeovalue{prvalue}, \emcppsgloss{lifetime extension} is
vital to keeping the range object live for the duration of the loop:

\begin{emcppslisting}[emcppsbatch=e3]
void f2(int i)
{
    for (char c : strFromInt(i))
    {
        // ...
    }
}
\end{emcppslisting}
    

\noindent The return value from \lstinline!strFromInt! is stored in a temporary
variable of type \lstinline!std::string!. The temporary string is destroyed
when the loop completes, not when the expression evaluation completes.
If the string were to go out of scope immediately, it would not be
possible to iterate over its characters. This code would have
\emcppsgloss{undefined behavior} were it not for the \emcppsgloss{lifetime
extension} harnessed by the \emcppsgloss[range based for loop]{ranged-based \lstinline!for! loop}.

The limitation of \emcppsgloss{lifetime extension} is that it applies only if
the reference is bound \emph{directly} to the temporary variable itself
or to a subobject (e.g., a member variable) of the temporary variable,
in which case the lifetime of the entire temporary variable is extended.
Note that initializing a reference from a reference or a pointer to
either the temporary or one of its subobjects does not count as binding
\emph{directly} to the temporary variable and does not trigger
\emcppsgloss{lifetime extension}. The danger of an object getting destroyed
prematurely is generally seen when the \emcppsgloss{full expression} returns
a reference, pointer, or iterator into a temporary object:

\begin{emcppslisting}[emcppsbatch=e4]
#include <vector>   // (ù{\codeincomments{std::vector}}ù)
#include <string>   // (ù{\codeincomments{std::string}}ù)
#include <utility>  // (ù{\codeincomments{std::pair}}ù)
#include <tuple>    // (ù{\codeincomments{std::tuple}}ù)

struct Point
{
    double x, y;
    Point(double ax, double ay) : x(ax), y(ay) { }
};

struct SRef
{
    const std::string& str;
    SRef(const std::string& s) : str(s) { }
};

std::vector<int> getValues();  // Return a (ù{\codeincomments{vector}}ù) by value.

void f3()
{
    const Point& p1 = Point(1.2, 3.4);    // OK, extend (ù{\codeincomments{Point}}ù) lifetime.
    double&&     d1 = Point(1.2, 3.4).x;  // OK, extend (ù{\codeincomments{Point}}ù) lifetime.
    double&      d2 = Point(1.2, 3.4).y;  // Error, non-(ù{\codeincomments{const}}ù) lvalue ref, (ù{\codeincomments{d2}}ù)

    using ICTuple = std::tuple<int, char>;
    const int&  i1 = getValues()[0];               // Bug, dangling reference
    const int&  i2 = std::get<0>(ICTuple{0,'a'});  // Bug,     "        "
    auto&&      i3 = getValues().begin();          // Bug,     "    iterator
    const auto& s1 = std::string("abc").c_str();   // Bug,     "    pointer
    const auto& i4 = std::string("abc").length();  // OK, (ù{\codeincomments{std::size\_t}}ù) extended

    SRef&&      sr = SRef("hello");  // Bug, (ù{\codeincomments{string}}ù) lifetime is not extended.
    std::string s2 = sr.str;         // Bug, (ù{\codeincomments{string}}ù) has been destroyed.
}
\end{emcppslisting}
    

\noindent The first invocation of the \lstinline!Point! constructor creates a
temporary object that is bound to reference \lstinline!p1!. The lifetime of
this temporary object is extended to match the lifetime of the
reference. Similarly, the lifetime of the second \lstinline!Point! object
is extended because a subobject, \lstinline!x!, is bound to reference
\lstinline!d1!. Note that it is not permitted to bind a temporary to a
non\lstinline!const! \romeovalue{lvalue} reference, as is being attempted with
\lstinline!d2!, above.

The next four definitions do not result in useful \emcppsgloss{lifetime
extension} at all.
\begin{enumerate}
\item{In the case of \lstinline!i1!, \lstinline!getValues()! returns a \romeovalue{prvalue} of type \lstinline!std::vector<int>!, resulting in the creation of a temporary variable. That temporary variable, however, is \emph{not} the value being bound to the \lstinline!i1! reference; rather, the reference is bound to the result of the array-access operator (\lstinline!operator[]!), which returns a reference into the temporary \lstinline!vector! returned by \lstinline!getValues()!. While we might consider an element of a \lstinline!vector! logically to be a subobject of the vector, \lstinline!i1! is not bound directly to that subobject but rather to the reference returned by \lstinline!operator[]!. The vector goes out of scope immediately at the end of the statement, leaving \lstinline!i1! to refer to an element of a deleted object.}
\item{The identical situation occurs with \lstinline!i2! when accessing the member of a temporary \lstinline!std::tuple!, this time via the nonmember function \lstinline!std::get<0>!.}
\item{Rather than a reference, \lstinline!i3! is deduced to be an \emph{iterator} as the result of the expression. The iterator’s lifetime is extended, but the lifetime of the object to which it refers is not.}
\item{Similarly, for \lstinline!s!, the expression \lstinline!std::string("abc").c_str()! yields a pointer into a temporary C-style string. Once again, the temporary \lstinline!std::string! variable is not the object that is bound to the reference \lstinline!s1!, so it gets destroyed at the end of the statement, invalidating the pointer.}
\end{enumerate}

Conversely, \lstinline!i4! binds directly to the temporary object returned
by \lstinline!length!, extending its life even though the string itself
gets destroyed as before. Unlike \lstinline!i3! and \lstinline!s1!, however,
\lstinline!i4! is not an iterator or pointer and so does not retain an
implicit reference to the defunct string object.

The last two definitions, for \lstinline!sr! and \lstinline!s2!, show how
subtle the rules for \emcppsgloss{lifetime extension} can be. The
\lstinline!"hello"! literal is converted into a temporary variable of type
\lstinline!std::string! and passed to the constructor of \lstinline!SRef!,
which \emph{also} creates a temporary object. It is only the
\lstinline!SRef! object that is bound to the \lstinline!sr! reference, so it
is only the \lstinline!SRef! object whose lifetime is extended. The
\lstinline!std::string("hello")! temporary variable gets destroyed when the
constructor finishes executing, leaving the object referenced by
\lstinline!sr! with a member, \lstinline!str!, that refers to a destroyed
object.

There are good reasons why \emcppsgloss{lifetime extension} applies only to
the temporary object being bound to a reference. A lot of code depends
on temporary objects going out of scope immediately, i.e., to release a
lock, memory, or some other resource. For \emcppsgloss[range based for loop]{range-based \lstinline!for! loops}, however, a compelling argument has been made that
the correct behavior would be to extend the lifetime of \emph{all} of
the temporaries constructed while evaluating the \emcppsgloss{range
expression}.{\cprotect\footnote{At the time of writing the P2012 paper
seeks to solve the issue when a \emcppsgloss{range expression} is a
reference into a temporary. See \emph{Fix the range-based for loop},
by Nicolai Josuttis, Victor Zverovich, Filipe Mulonde, and Arthur
O'Dwyer, which references an original paper
\emph{Embracing Modern C++ Safely} by Rostislav Khlebnikov and John
  Lakos.}} Unless and until this behavior is changed in a future
Standard, beware of using a \emcppsgloss{range expression} that returns a
reference to a temporary variable:

\begin{emcppslisting}[emcppsbatch=e4]
#include <iostream>  // (ù{\codeincomments{std::istream}}ù), (ù{\codeincomments{std::cout}}ù)

class RecordList
{
    std::vector<std::string> d_names;
    // ...

public:
    explicit RecordList(std::istream& is);
        // Create a (ù{\codeincomments{RecordList}}ù) with data read from (ù{\codeincomments{is}}ù).

    // ...

    const std::vector<std::string>& names() const { return d_names; }
};

void printNames(std::istream& is)
{
    // Bug, (ù{\codeincomments{RecordList}}ù)'s lifetime is not extended.
    for (const std::string& name : RecordList(is).names())
    {
        std::cout << name << '\n';
    }
}
\end{emcppslisting}
    

\noindent The \lstinline!RecordList! constructed in the \emcppsgloss{range expression} is
not bound to the implied \lstinline!__range! reference within the
\emcppsgloss[range based for loop]{ranged-based \lstinline!for! loop}, so its lifetime will end before
the loop actually begins. Thus, the
\lstinline!const!~\lstinline!std::vector<std::string>&! returned by its
\lstinline!names! method becomes a dangling reference, leading to
\emcppsgloss{undefined behavior} (such as a segmentation fault).

To avoid this pitfall, create a named object for each temporary that you
need to preserve, unless that temporary is the \emcppsgloss{full expression}
for the \emcppsgloss{range expression}:

\begin{emcppslisting}[emcppsbatch=e4]
void printNames2(std::istream& is)
{
    {
        RecordList records(is);  // named variable
        for (const std::string& name : records.names())
        {
            std::cout << name << '\n';
        }

        // safe for (ù{\codeincomments{records}}ù) to go out of scope now
    }

    // ...
}
\end{emcppslisting}
    

\noindent This minor rewrite of \lstinline!printNames! creates an extra block scope
in which we declare \lstinline!records! as a named variable. The inner
scope ensures that \lstinline!records! gets destroyed immediately after the
loop terminates.

\subsubsection[Inadvertent copying of elements]{Inadvertent copying of elements}\label{inadvertent-copying-of-elements}

When iterating through a container with a classic \lstinline!for! loop,
elements are typically referred to through an iterator:

\begin{emcppshiddenlisting}[emcppsbatch=e5]
#include <vector>  // (ù{\codeincomments{std::vector}}ù)
#include <string>  // (ù{\codeincomments{std::string}}ù)
\end{emcppshiddenlisting}
\begin{emcppslisting}[emcppsbatch=e5]
void process(std::string&);

void f1(std::vector<std::string>& vec)
{
    for (std::vector<std::string>::iterator i = vec.begin();
         i != vec.end(); ++i)
    {
        process(*i);  // refer to element via iterator
    }
}
\end{emcppslisting}
    

\noindent The \emcppsgloss[range based for loop]{ranged-based \lstinline!for! loop} gives the element a name and
a type. If the type is not a reference, then each iteration of the loop
will \emph{copy} the current element. In many cases, this copy is
inadvertent:

\begin{emcppslisting}[emcppsbatch=e5]
void f2(std::vector<std::string>& vec)
{
    for (std::string s : vec)
    {
        process(s);  // call (ù{\codeincomments{process}}ù) on *copy* of string element, potential
                     // bug.
    }
}
\end{emcppslisting}
    

\noindent The example above illustrates two issues: (1) there is an unnecessary
expense in copying each string, and (2) \lstinline!process! may modify or
take the address of its argument, in which case it will modify or take
the address of the copy, rather than the original element; the strings
in \lstinline!vec! will remain unchanged.

This error appears to be especially common when using \lstinline!auto! to
deduce the loop variable's type:

\begin{emcppslisting}[emcppsbatch=e5]
void f3(std::vector<std::string>& vec)
{
    for (auto s : vec)
    {
        process(s);  // call (ù{\codeincomments{process}}ù) on *copy* of deduced string element,
                     // potential bug.
    }
}
\end{emcppslisting}
    

\noindent Copying an element is not always erroneous, but it may be wise to
habitually declare the loop variable as a reference, making deliberate
exceptions when needed:

\begin{emcppslisting}[emcppsbatch=e5]
void f4(std::vector<std::string>& vec)
{
    for (std::string& s : vec)
    {
        process(s);  // OK, call (ù{\codeincomments{process}}ù) on *reference* to string element
    }
}
\end{emcppslisting}
    

\noindent If we wish to avoid copying elements but also wish to avoid modifying
them, then a \lstinline!const! reference will provide a good balance. Note,
however, that if the type being iterated over is not the same as the
type of the reference, a conversion might quietly produce the
(undesired) copy anyway:

\begin{emcppslisting}[emcppsbatch=e5]
void f5(std::vector<char*>& vec)
{
    for (const std::string& s : vec)
    {
        // (ù{\codeincomments{s}}ù) is a reference to a *copy* of an element of (ù{\codeincomments{vec}}ù).
    }
}
\end{emcppslisting}
    

\noindent In this example, the elements of \lstinline!vec! have type \lstinline!char*!.
The use of \lstinline!const!~\lstinline!std::string&! to declare the loop
variable \lstinline!s! correctly prevents modification of any elements of
\lstinline!vec!, but there is still a copy being made because each member
access is converted to an object of type \lstinline!std::string!.

For generic code that modifies a container, \lstinline!auto&&! is the
most general way to declare the loop variable. For generic code that
does not modify the container, \lstinline!const!~\lstinline!auto&! is safer:

\begin{emcppslisting}[emcppsbatch=e5]
template <typename Rng>
void f6(Rng& r)
{
    for (auto&& e : r)
    {
        // ...
    }

    for (const auto& cr : r)
    {
        // ...
    }
}
\end{emcppslisting}
    

\noindent Because \lstinline!e! is a \emcppsgloss{forwarding reference} and \lstinline!cr! is
a \lstinline!const! reference, they will both correctly bind to the return
type of \lstinline!*begin(Rng)!, even if that type is a \romeovalue{prvalue}.
Note that the use of \lstinline!auto! can obfuscate code by hiding the
underlying types of objects. Obfuscated code is prone to bugs so these
uses of \lstinline!auto! are recommended chiefly for \emph{generic} code or
where other trade-offs favor using this short-cut; see \featureref{\locationc}{auto}.

\subsubsection[Simple and reference-proxy behaviors can be be different]{Simple and reference-proxy behaviors can be be different}\label{simple-and-reference-proxy-behaviors-can-be-be-different}

Some containers have iterators that return proxies rather than
references to their elements. Depending on how the loop variable is
declared, the unwary programmer may get surprising results when the
container's iterator type returns reference proxies.

An example of such a container is \lstinline!std::vector<bool>!. The type,\linebreak[4]
\lstinline!std::vector<bool>::iterator::reference! is a reference-proxy
class that emulates a reference to a single bit within the
\lstinline!vector!. The proxy class provides an
\lstinline!operator!~\lstinline!bool()! that returns the bit when the proxy is
converted to \lstinline!bool! and an \lstinline!operator=(bool)! that modifies
the bit when assigned a Boolean value.

Let's consider a set of loops, each of which iterates over a
\lstinline!vector! and attempts to set each element of the vector to
\lstinline!true!. We'll embed the loops in a function template so that we
can compare the behavior of instantiating with a normal container
(\lstinline!std::vector<int>!) and with one whose iterator uses a reference
proxy (\lstinline!std::vector<bool>!):

\begin{emcppslisting}[emcppsbatch=e6,emcppserrorlines={7,8,10}]
#include <vector>  // (ù{\codeincomments{std::vector}}ù)

template <typename T>
void f1(std::vector<T>& vec)
{
    for (T      v : vec) { v = true; }  // (1)
    for (T&     v : vec) { v = true; }  // (2)
    for (T&&    v : vec) { v = true; }  // (3)
    for (auto   v : vec) { v = true; }  // (4)
    for (auto&  v : vec) { v = true; }  // (5)
    for (auto&& v : vec) { v = true; }  // (6)
}

void f2()
{
    using IntVec  = std::vector<int>;   // has normal iterator
    using BoolVec = std::vector<bool>;  // has iterator with reference proxy

    IntVec  iv{ /* ... */ };
    BoolVec bv{ /* ... */ };

    f1(iv);
    f1(bv);
}
\end{emcppslisting}
    

\noindent For each of the loops in \lstinline!f1!, the difference in behavior between
the \lstinline!IntVec! and \lstinline!BoolVec! instantiations hinges on what
happens when \lstinline!v! is initialized from \lstinline!*__begin! within
the loop transformation. For the \lstinline!IntVec! iterator,
\lstinline!*__begin! returns a \emph{reference} to the element within the
container. For the \lstinline!BoolVec! iterator, it returns an
\emph{object} of the reference-proxy type.
\begin{itemize}
\item{\emph{Loop (1)} produces identical behavior from both instantiations. The loop makes a local copy of each element, and then modifies the copy. The only difference is that the \lstinline!BoolVec! version performs a conversion to \lstinline!bool! to initialize \lstinline!v!, whereas the \lstinline!IntVec! version initializes \lstinline!v! directly from the element reference. For both the \lstinline!IntVec! or \lstinline!BoolVec! version, the fact that the original vector is unchanged is a potential bug (see Inadvertent copying of elements, above).}
\item{\emph{Loop (2)} modifies the container elements in the \lstinline!IntVec! instantiation but fails to compile for the \lstinline!BoolVec! instantiation. The compilation error comes from trying to initialize the non\lstinline!const! \emcppsgloss[lvalue reference]{\romeovalue{lvalue} reference}, \lstinline!v!, from an \romeovalue{rvalue} of the proxy type. The \lstinline!bool! conversion operator does not help since the result would still be an \romeovalue{rvalue}.}
\item{\emph{Loop (3)} fails to compile for the \lstinline!IntVec! iterator because the \emcppsgloss[rvalue reference]{\romeovalue{rvalue} reference}, \lstinline!v!, cannot be initialized from \lstinline!*__begin!, which is an \emcppsgloss[lvalue reference]{\romeovalue{lvalue} reference}. Surprisingly, the \lstinline!BoolVec! instantiation \emph{does} compile, but the loop does not modify the container. Here, \lstinline!operator!~\lstinline!bool! is invoked on the proxy object returned by \lstinline!*__begin!. The resulting temporary object is bound to \lstinline!v!, and its lifetime is extended for the duration of the iteration. Because \lstinline!v! is bound to a temporary variable, modifying \lstinline!v! modifies only the temporary, not the original, element, resulting in a likely bug as in the case of loop (1).}
\item{\emph{Loop (4)} compiles for both the \lstinline!BoolVec! and \lstinline!IntVec! cases but produces different results for each. For \lstinline!IntVec! iterators, \lstinline!auto! deduces the type of \lstinline!v! as \lstinline!int!, so assigning to \lstinline!v! modifies a local copy of the element, as in loop (1). For \lstinline!BoolVec! iterators, the deduced type of \lstinline!v! is the proxy type rather than \lstinline!bool!. Assigning to the proxy \emph{does} change the element of the container.}
\item{\emph{Loop (5)}, like loop (2), works as expected for \lstinline!IntVec! instantiation but fails to compile for the \lstinline!BoolVec! instantiation. As before, the problem is that the \lstinline!BoolVec! iterator yields an \romeovalue{rvalue} that cannot be used to initialize an \emcppsgloss[lvalue reference]{\romeovalue{lvalue} reference}.}
\item{\emph{Loop (6)} produced identical behavior from both instantiations, modifying each of the \lstinline!vector! elements. The type of \lstinline!v! is deduced to be \lstinline!int&! for \lstinline!IntVec! instantiation and the proxy type for the \lstinline!BoolVec! instantiation. Assigning through either the real reference or the reference proxy modifies the element in the container.}
\end{itemize}

Let's now look at the the situation with \lstinline!const!-qualified loop
variables:

\begin{emcppslisting}[emcppsbatch=e6,emcppserrorlines={6,9}]
template <typename T>
void f3(std::vector<T>& vec)
{
    for (const T      v : vec) { /* ... */ }  // (7)
    for (const T&     v : vec) { /* ... */ }  // (8)
    for (const T&&    v : vec) { /* ... */ }  // (9)
    for (const auto   v : vec) { /* ... */ }  // (10)
    for (const auto&  v : vec) { /* ... */ }  // (11)
    for (const auto&& v : vec) { /* ... */ }  // (12)
}

void f4()
{
    using IntVec  = std::vector<int>;   // has normal iterator
    using BoolVec = std::vector<bool>;  // has iterator with reference proxy

    IntVec  iv{ /* ... */ };
    BoolVec bv{ /* ... */ };

    f3(iv);
    f3(bv);
}
\end{emcppslisting}
    
\begin{itemize}
\item{\emph{Loop (7)} works identically for both instantiations, converting the proxy reference to \lstinline!bool! in the \lstinline!BoolVec! case.}
\item{\emph{Loop (8)} works identically for both instantiations. For the \lstinline!IntVec! case, the result of \lstinline!*__begin! is bound directly to \lstinline!v!. For the \lstinline!BoolVec! case, \lstinline!*__begin! produces proxy reference that is converted to a \lstinline!bool! temporary that is then bound to \lstinline!v!. \emcppsgloss[lifetime extension]{Lifetime extension} keeps the bool value alive.}
\item{\emph{Loop (9)} fails to compile for \lstinline!IntVec! but succeeds for \lstinline!BoolVec! exactly as for loop (3) except that the temporary \lstinline!bool! bound to \lstinline!v! is \lstinline!const! and thus does not risk giving programmers the false belief that they are modifying the container.}
\item{\emph{Loop (10)} has the same behavior for both the \lstinline!IntVec! and \lstinline!BoolVec! instantiations. That mechanism is the same behavior as for loop (4) except that, because \lstinline!v! is \lstinline!const!, \emph{neither} instantiation can modify the container.}
\item{\emph{Loop (11)} also works for both instantiations. For the \lstinline!IntVec! case, the result of \lstinline!*__begin! is bound directly to \lstinline!v!. For the \lstinline!BoolVec! case, \lstinline!v! is deduced to be a \lstinline!const! reference to the proxy type; \lstinline!*__begin! produces a temporary variable of the proxy type, which is then bound to \lstinline!v!. \emcppsgloss[lifetime extension]{Lifetime extension} keeps the proxy alive. In most contexts, a \lstinline!const! proxy reference is an effective stand-in for a \lstinline!const!~\lstinline!bool&!.}
\item{\emph{Loop (12)} fails to compile for \lstinline!IntVec! but succeeds for \lstinline!BoolVec!. The error with \lstinline!IntVec! occurs because \lstinline!const!~\lstinline!auto&&! is always a \lstinline!const! \emcppsgloss[rvalue reference]{\romeovalue{rvalue} reference} (not a \emcppsgloss{forwarding reference}) and cannot be bound to the \emcppsgloss[lvalue reference]{\romeovalue{lvalue} reference}, \lstinline!*__begin!. For \lstinline!BoolVec!, the mechanism is identical to loop (11) except that loop (11) binds the temporary object to an \emcppsgloss[lvalue reference]{\romeovalue{lvalue} reference} whereas loop (12) uses an \emcppsgloss[rvalue reference]{\romeovalue{rvalue} reference}. When the references are \lstinline!const!, however, there is little practical differences between them.}
\end{itemize}

Note that loops 4, 6, 10, 11, and 12 in the \lstinline!BoolVec!
instantiations bind a reference to a temporary proxy reference object,
so taking the address of \lstinline!v! in these situations is likely not to
produce useful results. Additionally, loops 3, 8, and 9 bind \lstinline!v!
to a temporary \lstinline!bool!. Users must be mindful of the lifetime of
these temporary objects (a single-loop iteration) and not allow the
address of \lstinline!v! to escape the loop.

\noindent Proxy objects emulating references to non-class elements within a
container are surprisingly effective, but their limitations are exposed
when they are bound to references. In generic code, as a rule of thumb,
\lstinline!const!~\lstinline!auto&! is the safest way to declare a read-only
loop variable if a reference proxy might be in use, while
\lstinline!auto&&! will give the most consistent results for a loop that
modifies its container. Similar issues, unrelated to \emcppsgloss[range based for loop]{range-based
\lstinline!for! loops}, occur when passing a proxy reference to a function
taking a reference argument.

\subsection[Annoyances]{Annoyances}\label{annoyances}

\subsubsection[No access to the state of the iteration]{No access to the state of the iteration}\label{no-access-to-the-state-of-the-iteration}

When traversing a range with a classic \lstinline!for! loop, the loop
variable is typically an iterator or array index. Within the loop, we
can modify that variable to repeat or skip iterations. Similarly, the
loop-termination condition is usually accessible so that it is possible
to, for example, insert or remove elements and then recompute the
condition:

\begin{emcppslisting}[emcppsbatch=e7]
#include <string>  // (ù{\codeincomments{std::string}}ù)
#include <cctype>  // (ù{\codeincomments{std::isupper}}ù)

void spaceBeforeCaps(std::string& s)
{
    // Insert a space before each capital letter in (ù{\codeincomments{s}}ù).
    using IdxType = std::string::size_type;
    for (IdxType i = 0; i < s.size(); ++i)
    {
        if (std::isupper(s[i]))
        {
            s.insert(i, 1, ' ');  // Insert one space at (ù{\codeincomments{i}}ù).
            ++i;                  // Skip the space.
        }
    }
}
\end{emcppslisting}
    

\noindent The code above depends on (1) having access to the iteration index, (2)
being able to change the iteration index, and (3) recomputing the size
of the collection each time through the loop. No similar function could
be written using a \emcppsgloss[range based for loop]{ranged-based \lstinline!for! loop} since the
\lstinline!__range!, \lstinline!__begin!, and \lstinline!__end! variables
are for exposition only and are not accessible from within the code:

\begin{emcppslisting}[emcppsbatch=e7]
void spaceBeforeCaps2(std::string& s)
{
    // Insert a space before each capital letter in (ù{\codeincomments{s}}ù).
    for (char c : s)
    {
        if (std::isupper(c))
        {
            __begin = s.insert(__begin, ' ');  // Error, no (ù{\codeincomments{\_\_begin}}ù) variable
            ++__begin;                         // Error, no (ù{\codeincomments{\_\_begin}}ù) variable
            __end = s.end();                   // Error, no (ù{\codeincomments{\_\_end}}ù) variable
        }
    }
}
\end{emcppslisting}
    

\noindent A classic \lstinline!for! loop can traverse more than one container at a
time (e.g., to add corresponding elements from two containers and store
them into a third). It accomplishes this feat by either incrementing
multiple iterators on each iteration or keeping a single index that is
used to access multiple, random-access iterators concurrently. Trying to
accomplish something similar with a \emcppsgloss[range based for loop]{range-based \lstinline!for! loop} usually involves using a hybrid approach:

\begin{emcppslisting}
#include <vector>   // (ù{\codeincomments{std::vector}}ù)
#include <cassert>  // standard C (ù{\codeincomments{assert}}ù) macro

void addVectors(std::vector<int>&       result,
                const std::vector<int>& a,
                const std::vector<int>& b)
    // For each element (ù{\codeincomments{ea}}ù) of (ù{\codeincomments{a}}ù) and corresponding element (ù{\codeincomments{eb}}ù) of (ù{\codeincomments{b}}ù), set
    // the corresponding element of (ù{\codeincomments{result}}ù) to (ù{\codeincomments{ea + eb}}ù).  The behavior is
    // undefined unless (ù{\codeincomments{a}}ù) and (ù{\codeincomments{b}}ù) have the same length.
{
    assert(a.size() == b.size());
    result.resize(a.size());

    std::vector<int>::const_iterator ia = a.begin();
    std::vector<int>::const_iterator ib = b.begin();
    for (int& sum : result)
    {
        sum = *ia++ + *ib++;
    }
}
\end{emcppslisting}
    

\noindent Although \lstinline!result! is traversed using the \emcppsgloss[range based for loop]{range-based
\lstinline!for! loop}, \lstinline!a! and \lstinline!b! are effectively traversed
manually using iterators. It is arguable as to whether the code is any
clearer or simpler to write than it would have been using a classic
\lstinline!for! loop.

This situation can be improved through the use of a ``zip'' iterator ---
a type that holds multiple iterators and increments them in lock-step.
Using a ``zip'' iterator, all three containers can be traversed using a
single \emcppsgloss[range based for loop]{ranged-based \lstinline!for! loop}:

\begin{emcppslisting}
#include <cassert>  // standard C (ù{\codeincomments{assert}}ù) macro
#include <tuple>    // (ù{\codeincomments{std::tuple}}ù)
#include <utility>  // (ù{\codeincomments{std::declval}}ù)
#include <vector>   // (ù{\codeincomments{std::vector}}ù)

template <typename... Iter>
class ZipIterator
{
    std::tuple<Iter...> d_iters;

    // ...

public:
    using reference = std::tuple<decltype(*std::declval<Iter>())...>;

    ZipIterator(const Iter&... i);

    reference operator*() const;
    ZipIterator& operator++();
    friend bool operator!=(const ZipIterator& a, const ZipIterator& b);
};

template <typename... Range>
class ZipRange
{
    using ZipIter =
        ZipIterator<decltype(begin(std::declval<Range>()))...>;

    // ...

public:
    ZipRange(const Range&... ranges);

    ZipIter begin() const;
    ZipIter end() const;
};

template <typename... Range>
ZipRange<Range...> makeZipRange(Range&&... r);

void addVectors2(std::vector<int>&       result,
                 const std::vector<int>& a,
                 const std::vector<int>& b)
{
    assert(a.size() == b.size());
    result.resize(a.size());

    for (std::tuple<int, int, int&> elems : makeZipRange(a, b, result))
    {
        std::get<2>(elems) = std::get<0>(elems) + std::get<1>(elems);
    }
}
\end{emcppslisting}
    
\noindent Each iteration, instead of yielding a single element, yields a
\lstinline!std::tuple! of elements resulting from the traversal of multiple
ranges simultaneously. To be used, the elements must be unpacked from
the \lstinline!std::tuple! using \lstinline!std::get!. Zip iterators become
much more attractive in C++17 with the advent of \emcppsgloss[structured binding]{structured
bindings}, which allow multiple loop variables to be declared at once,
without the need to directly unpack the \lstinline!std::tuple!s. The above
implementation and usage of \lstinline!ZipRange! is just a rough sketch;
the full design and implementation of zip iterators and zip ranges is
beyond the scope of this section.

\subsubsection[Adapters are required for many tasks]{Adapters are required for many tasks}\label{adapters-are-required-for-many-tasks}

In the usage examples above, we have seen a number of adapters (e.g., to
traverse subranges, to traverse a container in reverse, to generate
sequential values), and zip iterators to iterate over multiple ranges at
once.

None of these adapters would be required for a classic \lstinline!for!
loop. On the one hand, one-off situations are expressed more simply with
a classic \lstinline!for! loop. On the other hand, the adapters that we
create to make \emcppsgloss[range based for loop]{range-based \lstinline!for! loops} usable in more
situations can lead to the development of a reusable \emph{library} of
adapters. Using the \lstinline!ValueGenerator! class from {Range
generators}, for example, produces simpler and more expressive code than
using a classic \lstinline!for! loop would.{\cprotect\footnote{The
Standard's Ranges Library, introduced in C++20, provides a
  sophisticated algebra for working with and adapting ranges.}}

\subsubsection[No support for sentinel iterator types]{No support for sentinel iterator types}\label{no-support-for-sentinel-iterator-types}

For a given \emcppsgloss{range expression}, \lstinline!__range!,
\lstinline!begin(__range)!, and \lstinline!end(__range)! must return the
same type to be usable with a \emcppsgloss[range based for loop]{ranged-based \lstinline!for! loop}.
This limitation is problematic for ranges of indeterminate length, where
the condition for ending a loop is not determined by comparing two
iterators. For example, in the \lstinline!RandomIntSequence! example (see
{Range generators}), the end iterator for the infinite random sequence
holds a null pointer and is never used, not even within
\lstinline|operator!=|. It would be more efficient and convenient if the
end iterator were a special, empty \emph{sentinel} type. Comparing any
iterator to the sentinel would determine whether the loop should
terminate:

\begin{emcppslisting}
template <typename T = int>
class RandomIntSequence2
{
    // ...

public:
    class sentinelIterator { };

    class iterator {
        // ...
        friend bool operator!=(iterator, sentinelIterator) { return true; }
    };

    iterator         begin() { /* ... */ }
    sentinelIterator end() const { return {}; }
};
\end{emcppslisting}
    

\noindent The above code shows an example of \lstinline!begin! and \lstinline!end!
returning different types, where \lstinline!end! returns an empty sentinel
type. Unfortunately, using this formulation of
\lstinline!RandomIntSequence2!\linebreak[4]\newpage 
with a C++11 \emcppsgloss[range based for loop]{range-based
\lstinline!for! loop} will result in a compilation error complaining that
\lstinline!begin! and \lstinline!end! return inconsistent
types.{\cprotect\footnote{This limitation on the use of sentinel
iterators was lifted in C++17. Sentinel iterators are supported
directly by the C++20 Ranges Library. In C++17, the specification was
modified to:

\begin{emcppslisting}[style=footcode,emcppsignore={grammar, not code, but leave as c++}]
{
    auto&& __range = range-expression;
    auto __begin   = begin-expr;
    auto __end     = end-expr;
    for (; __begin != __end; ++__begin)
    {
        for-range-declaration = *__begin;
        statement
    }
}
\end{emcppslisting}
      }}

Another type that could benefit from a sentinel iterator is
\lstinline!std::istream_iterator!, since the state of the end iterator is
never used. It is unlikely that this interface will change, however, as
\lstinline!std::istream_iterator! has been with us since the first C++
Standard.

\subsubsection[Only ADL lookup]{Only ADL lookup}\label{only-adl-lookup}

The free functions \lstinline!begin! and \lstinline!end! are found using
\emcppsgloss{argument-dependent lookup (ADL)} only. File-scope
(\lstinline!static!) functions are not considered. If we wish to add
\lstinline!begin! and \lstinline!end! functions for a range-like type that we
do not own, we need to put those functions into the same namespace as
the range-like type, inviting a potential name collision with other
compilation units attempting to do the same thing:

\begin{emcppslisting}[emcppsbatch=e8]
// third_party_library.h:

namespace third_party
{

    class IteratorLike { /* ... */ };

    class RangeLike
    {
        // ... does not provide (ù{\codeincomments{begin}}ù) and (ù{\codeincomments{end}}ù) members
    };

    // ... does not provide (ù{\codeincomments{begin}}ù) and (ù{\codeincomments{end}}ù) free functions
}
\end{emcppslisting}
%\vspace*{2ex}
\newpage%%%%%%%%%%%%%
\begin{emcppslisting}[emcppsbatch=e8]
// myclient.cpp:

#include <third_party_library.h>

static third_party::IteratorLike begin(third_party::RangeLike&);
static third_party::IteratorLike end(third_party::RangeLike&);

void f()
{
    third_party::RangeLike rl;
    for (auto&& e : rl)  // Error, (ù{\codeincomments{begin}}ù) not found by ADL
    {
        // ...
    }
}
\end{emcppslisting}
    

\noindent The code above attempts to work around the absence of
\lstinline!begin(RangeLike&)! and \lstinline!end(RangeLike&)! in the
third-party library by defining them locally within
\lstinline!myclient.cpp!. This attempt fails because static functions are
not found via ADL. A better workaround that does work is to create a
range adapter for the range-like class:

\begin{emcppshiddenlisting}[emcppsbatch=e8]
#include <third_party_library.h>
\end{emcppshiddenlisting}
\begin{emcppslisting}[emcppsbatch=e8]
class RangeLikeAdapter
{
    // ...

public:
    RangeLikeAdapter(third_party::RangeLike&);
    third_party::IteratorLike begin() { /* ... */ }
    third_party::IteratorLike end()   { /* ... */ }
};
\end{emcppslisting}
    

\noindent The adapter wraps the range-like type and provides the missing features.
Beware, however, that if the wrapper stores a pointer or reference to a
temporary \lstinline!RangeLike! object, you don't run into the pitfall
where the lifetime of temporary objects is not always extended; see
{Lifetime of temporaries in the range expression}.

\subsection[See Also]{See Also}\label{see-also}

\begin{itemize}
\item{\seealsoref{auto}{\seealsolocationc}explains \lstinline!auto!, often used in a \emcppsgloss[range based for loop]{ranged-based \lstinline!for! loop} to determine the type of the loop variable, and many of the pitfalls of \lstinline!auto! apply when using it for that purpose.}
\item{\seealsoref{refqualifiers}{\seealsolocatione}show how to overload member functions to work differently on \romeovalue{rvalues} and \romeovalue{lvalues}.}
\end{itemize}

\subsection[Further Reading]{Further Reading}\label{further-reading}

No further reading.




\newpage
%\section[{\itshape rvalue} References]{Rvalue References: {\SecCode \&\&}}\label{Rvalue-References}


placeholder text..........




\newpage
%\section[Underlying Type '11]{Explicit Enumeration Underlying Type}\label{explicit-enumeration-underlying-type}

% 01 Feb, 2021 JMB - packet 4 compilation fixes
% 10 Feb 2021, some cleanup before sending for revisions
% 4 March 2021, revisions entered and proofed
% 5 March 2021, glossary entries updated
% 5 March 2021, sent to copyediting
% 13 March 2021, copyedits in and proofed


\emcppsFeature{
    short={Underlying Type '11},
    long={Explicit Enumeration Underlying Type},
}{explicit-enumeration-underlying-type}
\setcounter{table}{0}
\setcounter{footnote}{0}
\setcounter{lstlisting}{0}
%\section[Underlying Type '11]{Explicit Enumeration Underlying Type}\label{explicit-enumeration-underlying-type}


The underlying type of an enumeration is the fundamental
\emcppsgloss{integral type} used to represent its enumerated values, which
can be specified explicitly in C++11.

\subsection[Description]{Description}\label{description}

Every enumeration employs an integral type, known as its
\emcppsgloss[underlying type (UT)]{underlying type}, to represent its compile-time-enumerated
values. By default, the \emcppsgloss[underlying type (UT)]{underlying type} of a C++03
\lstinline!enum! is chosen by the
implementation to be large enough to represent all of the values in the
enumeration and is allowed to exceed the size of an \lstinline!int!
\emph{only} if there are enumerators having values that cannot be
represented as an \lstinline!int! or \lstinline!unsigned!~\lstinline!int!:

\begin{emcppslisting}
enum RGB { e_RED, e_GREEN, e_BLUE };                  // OK, fits any (ù{\codeincomments{char}}ù)

enum Port { e_LEFT = -81, e_RIGHT = -82 };            // OK, fits (ù{\codeincomments{signed char}}ù)

enum Mask { e_LOW = 32767, e_HIGH = 65535 };          // OK, fits (ù{\codeincomments{unsigned short}}ù)

enum Big { e_31 = 1U<<31 };                           // OK, fits (ù{\codeincomments{unsigned int}}ù)

enum Err { K = 1024, M = K*K, G = M*K, T = G*K };     // Error, (ù{\codeincomments{G*K}}ù) overflows (ù{\codeincomments{int}}ù)...

enum OK { K = 1<<10, M = 1<<20, G = 1<<30, T = 1LL<<40 }; // OK
\end{emcppslisting}

\noindent The default underlying type chosen for an \lstinline!enum! is always
sufficiently large to represent \emph{all} enumerator values defined for
that \lstinline!enum!. If the value doesn't fit in an \lstinline!int!, it will
be selected deterministically as the first type able to represent all
values from the sequence: \lstinline!unsigned!~\lstinline!int!, \lstinline!long!,
\lstinline!unsigned!~\lstinline!long!, \lstinline!long!~\lstinline!long!,
\lstinline!unsigned!~\lstinline!long!~\lstinline!long!; see \featureref{\locationa}{long-long}.

While specifying an enumeration’s underlying type was impossible before C++11, the compiler could be forced to choose at least a 32-bit or 64-bit signed integral type by adding an enumerator having a sufficiently large negative value --- e.g., \lstinline!-1 << 31! for a 32-bit and \lstinline!-1LL << 63! for a 64-bit signed integer (assuming such is available on the target platform).

The above applies only to C++03 \lstinline!enum!s; the default underlying type of an \lstinline!enum! \lstinline!class! is ubiquitously \lstinline!int!, and it is not implementation defined; see \featureref{\locationc}{enumclass}.

Note that \lstinline!char! and \lstinline!wchar_t!, like enumerations, are their own distinct types (as opposed to \lstinline!typedef!-like aliases such as \lstinline!std::uint8_t!) and have their own implementation-defined underlying integral types. With \lstinline!char!, for example, the underlying type will always be either \lstinline!signed! \lstinline!char! or \lstinline!unsigned! \lstinline!char! (both of which are also distinct C++ types). The same is true in C++11 for \lstinline!char16_t! and \lstinline!char32_t!.{\protect\footnote{C++20 adds \lstinline!char8_t!, which is a distinct type that has \lstinline!unsigned! \lstinline!char! as its underlying type.}}

\subsubsection[Specifying underlying type explicitly]{Specifying underlying type explicitly}\label{specifying-underlying-type-explicitly}

As of C++11, we have the ability to specify the \emcppsgloss{integral type}
that is used to represent an \lstinline!enum!. This is achieved by
providing the type explicitly in the \lstinline!enum!'s declaration
following the enumeration's (optional) name and preceded by a colon:

\begin{emcppslisting}
enum Port : unsigned char
{
    // Each enumerator of (ù{\codeincomments{Port}}ù) is represented as an (ù{\codeincomments{unsigned char}}ù) type.

    e_INPUT        =  37,  // OK, would have fit in a (ù{\codeincomments{signed char}}ù) too
    e_OUTPUT       = 142,  // OK, would not have fit in a (ù{\codeincomments{signed char}}ù)
    e_CONTROL      = 255,  // OK, barely fits in an 8-bit unsigned integer
    e_BACK_CHANNEL = 256,  // Error, doesn't fit in an 8-bit unsigned integer
};
\end{emcppslisting}

\noindent If any of the values specified in the definition of the \lstinline!enum! is
outside the boundaries of what the provided \emcppsgloss[underlying type (UT)]{underlying type} is
able to represent, the compiler will emit an error, but see \intraref{potential-pitfalls-underlyingenum}{subtleties-of-integral-promotion}. 
%\textit{\titleref{potential-pitfalls-underlyingenum}: \titleref{subtleties-of-integral-promotion}} on page~\pageref{subtleties-of-integral-promotion}.

\subsection[Use Cases]{Use Cases}\label{use-cases}

\subsubsection[Ensuring a compact representation where enumerator values are salient]{Ensuring a compact representation where enumerator values are salient}\label{ensuring-a-compact-representation-where-enumerator-values-are-salient}

When an enumeration needs to have an efficient representation, e.g.,
when it is used as a data member of a widely replicated type,
restricting the width of the underlying type to something smaller than
would occur by default on the target platform might be justified.

As a concrete example, suppose that we want to enumerate the months of
the year, in anticipation of placing that enumeration
inside a date class having an internal representation that maintains the
year as a two-byte signed integer, the month as an enumeration, and the
day as an 8-bit signed integer:

\begin{emcppshiddenlisting}[emcppsbatch=e1]
enum Month { January=1 };
\end{emcppshiddenlisting}
\begin{emcppslisting}[emcppsbatch=e1]
#include <cstdint>  // (ù{\codeincomments{std::int8\_t}}ù), (ù{\codeincomments{std::int16\_t}}ù)

class Date
{
    std::int16_t d_year;
    Month        d_month;
    std::int8_t  d_day;

public:
    Date();
    Date(int year, Month month, int day);

    // ...

    int year() const    { return d_year; }
    Month month() const { return d_month; }
    int day() const     { return d_day; }
};
\end{emcppslisting}

\noindent Suppose that, within an application, a \lstinline!Date! is typically constructed using
values obtained through a GUI, where the month is always selected
from a drop-down menu. The month is
supplied to the constructor as an \lstinline!enum! to avoid recurring
defects where the individual fields of the date are supplied in
month/day/year format. New functionality will be written to expect the
month to be enumerated. Still, the date class will be used in contexts
where the numerical value of the month is significant, such as in calls
to legacy functions that accept the month as an integer. Moreover,
iterating over a range of months is common and requires that the
enumerators convert automatically to their integral \emcppsgloss[underlying type (UT)]{underlying
type}, thus contraindicating use of the more strongly typed
\lstinline!enum! \lstinline!class!:

\begin{emcppslisting}
enum Month // defaulted underlying type (BAD IDEA)
{
    e_JAN = 1, e_FEB, e_MAR, e_APR, e_MAY, e_JUN,
    e_JUL , e_AUG, e_SEP, e_OCT, e_NOV, e_DEC
};
static_assert(sizeof(Month) <= 4 && alignof(Month) <= 4, "");
\end{emcppslisting}

%\begin{emcppslisting}
%enum Month  // defaulted underlying type (BAD IDEA)
%{
%    enum Month {
%    e_JAN = 1, e_FEB, e_MAR, e_APR, e_MAY, e_JUN,
%    e_JUL    , e_AUG, e_SEP, e_OCT, e_NOV, e_DEC
%    };
%};
%
%static_assert(sizeof(Month) <= 4 && alignof(Month) <= 4, "");
%\end{emcppslisting}

%\begin{emcppslisting}
%enum Month  // defaulted underlying type (BAD IDEA)
%{
%    e_JAN = 1, e_FEB, e_MAR,    // winter
%    e_APR    , e_MAY, e_JUN,    // spring
%    e_JUL    , e_AUG, e_SEP,    // summer
%    e_OCT    , e_NOV, e_DEC     // autumn
%};
%
%static_assert(sizeof(Month) == 4 && alignof(Month) == 4, "");
%\end{emcppslisting}

\noindent As it turns out, date values are used widely throughout this codebase,
and the proposed \lstinline!Date! type is expected to be used in large
aggregates. The underlying type of the \lstinline!enum! in the code snippet
above is implementation-defined and could be as small as a \lstinline!char!
or as large as an \lstinline!int! despite all the values fitting in a
\lstinline!char!. Hence, if this enumeration were used as a data member in
the \lstinline!Date! class, \lstinline!sizeof(Date)! would likely balloon to
12 bytes on some relevant platforms due to \emcppsgloss{natural alignment}!
(See \featureref{\locationc}{alignas}.)
%``\titleref{alignas}" on page~\pageref{alignas}.)

While reordering the data members of \lstinline!Date! such that \lstinline!d_year! and
\lstinline!d_day! were adjacent would ensure that \lstinline!sizeof(Date)!
would not exceed 8 bytes, a better approach is to explicitly specify the
enumeration's underlying type to ensure \lstinline!sizeof(Date)! is exactly
the 4 bytes needed to accurately represent the value of a
\lstinline!Date! object. Given that the values in this enumeration fit in
an 8-bit signed integer, we can specify its \emcppsgloss[underlying type (UT)]{underlying type} to
be, e.g., \lstinline!std::int8_t! or \lstinline!signed!~\lstinline!char!, on
every platform:

\begin{emcppslisting}
#include <cstdint>  // (ù{\codeincomments{std::int8\_t}}ù)

enum Month : std::int8_t  // user-provided underlying type (GOOD IDEA)
{
    e_JAN = 1, e_FEB, e_MAR,   
    e_APR    , e_MAY, e_JUN,   
    e_JUL    , e_AUG, e_SEP,   
    e_OCT    , e_NOV, e_DEC     
};

static_assert(sizeof(Month) == 1 && alignof(Month) == 1, "");
\end{emcppslisting}

\noindent With this revised definition of \lstinline!Month!, the size of a
\lstinline!Date! class is 4 bytes, which is especially valuable for large
aggregates:

\begin{emcppslisting}[emcppsbatch=e1]
Date timeSeries[1000 * 1000];  // (ù{\codeincomments{sizeof(timeSeries)}}ù) is now 4Mb (not 12Mb)
\end{emcppslisting}


\subsection[Potential Pitfalls]{Potential Pitfalls}\label{potential-pitfalls-underlyingenum}

\subsubsection[External use of opaque enumerators]{External use of opaque enumerators}\label{external-use-of-opaque-enumerators-enumunderlying}

Providing an explicit underlying type to an \lstinline!enum! enables
clients to declare it as a complete type without
its enumerators. Unless the opaque form of its definition is exported in
a header file separate from its full definition, external clients
wishing to exploit the opaque version will be forced to locally declare
it with its \emcppsgloss[underlying type (UT)]{underlying type} but without its enumerator list. If
the underlying type of the full definition were to change, any program
incorporating \emph{its own} original and now inconsistent elided
definition and the \emph{new} full one would become silently \emcppsgloss{ill formed, no diagnostic required (IFNDR)}. (See \featureref{\locationc}{enumopaque}.)
%``\titleref{enumopaque}" on page~\pageref{enumopaque}.)

\subsubsection[Subtleties of integral promotion]{Subtleties of integral promotion}\label{subtleties-of-integral-promotion}

When used in an arithmetic context, one might naturally assume that the
type of a classic \lstinline!enum! will first convert to its
\emcppsgloss[underlying type (UT)]{underlying type}, which is not always the case. When used in a
context that does not explicitly operate on the \lstinline!enum! type
itself, such as a parameter to a function that takes that enum type,
\emcppsgloss{integral promotion} comes into play. For unscoped enumerations
without an explicitly specified underlying type and for character types
such as \lstinline!wchar_t!, \lstinline!char16_t!, and \lstinline!char32_t!,
integral promotion will directly convert the value to the first type in
the list \lstinline!int!, \lstinline!unsigned!~\lstinline!int!, \lstinline!long!,
\lstinline!unsigned!~\lstinline!long!, \lstinline!long!~\lstinline!long!, and
\lstinline!unsigned!~\lstinline!long!~\lstinline!long! that is sufficiently large
to represent all of the values of the underlying type. Enumerations
having a fixed underlying type will, as a first step, behave as if they
had decayed to their underlying type.

In most arithmetic expressions, this difference is irrelevant.
Subtleties arise, however, when one relies on overload resolution for
identifying the underlying type:

\begin{emcppslisting}[emcppsbatch=e2]
void f(signed char x);
void f(short x);
void f(int x);
void f(long x);
void f(long long x);

enum E1         { q, r, s, t, u };
enum E2 : short { v, w, x, y, z };

void test()
{
   f(E1::q);  // always calls (ù{\codeincomments{f(int)}}ù) on all platforms
   f(E2::v);  // always calls (ù{\codeincomments{f(short)}}ù) on all platforms
}
\end{emcppslisting}

\noindent The overload resolution for \lstinline!f! considers the type to which each
\emph{individual} enumerator can be directly integrally promoted. This
conversion for \lstinline!E1! can be only to \lstinline!int!. For \lstinline!E2!,
the conversion will consider \lstinline!int! \emph{and} \lstinline!short!, and
\lstinline!short!, being an exact match, will be selected. Note that even
though both enumerations are small enough to fit into a
\lstinline!signed!~\lstinline!char!, that overload of \lstinline!f! will never be
selected.

One might want to get to the implementation-defined underlying type,
though, and the Standard does provide a trait to do that:
\lstinline!std::underlying_type! in C++11 and the corresponding
\lstinline!std::underlying_type_t! alias in C++14. This trait can safely
be used in a cast without risking loss of value (see \featureref{\locationc}{auto}):
%  ``\titleref{auto}" on page~\pageref{auto}):

\begin{emcppslisting}[emcppsbatch=e2]
#include <type_traits>  // (ù{\codeincomments{std::underlying\_type}}ù)

template <typename E>
typename std::underlying_type<E>::type toUnderlying(E value)
{
    return static_cast<typename std::underlying_type<E>::type>(value);
}

void h()
{
    auto e1 = toUnderlying(E1::q); // might be anywhere from (ù{\codeincomments{signed char}}ù) to (ù{\codeincomments{int}}ù)
    auto e2 = toUnderlying(E2::v); // always deduced as (ù{\codeincomments{short}}ù)
}
\end{emcppslisting}

%\noindent As of C++20, however, the use of a classic enumerator in a context in
%which it is compared to or otherwise used in a binary operation with
%either an enumerator of another type or a nonintegral type (i.e., a
%floating-point type, such as \lstinline!float!, \lstinline!double!, or
%\lstinline!long!~\lstinline!double!) is deprecated, with the the possibility
%of being removed in C++23. Platforms might decide to warn against such
%uses retroactively:
%
%\begin{emcppslisting}
%enum { k_GRAMS_PER_OZ = 28 };  // not the best idea
%
%double gramsFromOunces(double ounces)
%{
%    return ounces * k_GRAMS_PER_OZ;  // deprecated in C++20; might warn
%}
%\end{emcppslisting}
%
%\noindent Casting to the \emcppsgloss{underlying type} is \emph{not} necessarily the
%same as direct integral promotion. In this context, we might want to
%change our \lstinline!enum! to a
%\lstinline!constexpr!~\lstinline!int! in the long
%term (see \featureref{\locationc}{constexprvar}): 
%%  ``\titleref{constexprvar}" on page~\pageref{constexprvar}):
%
%\begin{emcppslisting}
%constexpr int k_GRAMS_PER_OZ = 28;  // future proof
%\end{emcppslisting}

Casting to the underlying type is not necessarily the same as direct integral promotion. If
an enumerator is intended to be used in arithmetic operations,\footnote{As of C++20, using an expression of an unscoped enumerated type in a \emph{binary} operation along with 
an expression of either some \emph{other} enumerated type or any nonintegral type (e.g., 
\lstinline!double!) is deprecated, with the possibility of being removed in C++23. Platforms 
may decide to warn against such uses retroactively.} a \lstinline!constexpr! variable might 
be a better alternative (see \featureref{\locationc}{constexprvar}):

\begin{emcppslisting}
// enum { k_GRAMS_PER_OZ = 28 };    // not the best idea
constexpr int k_GRAMS_PER_OZ = 28;  // better idea

double gramsFromOunces(double ounces)
{
    return ounces * k_GRAMS_PER_OZ;
}
\end{emcppslisting}


\subsection[See Also]{See Also}\label{see-also}

\begin{itemize}
\item{\seealsoref{constexprvar}{\seealsolocationc}introduces an alternative way of declaring compile-time constants.}
\item{\seealsoref{enumclass}{\seealsolocationc}introduces a scoped, more strongly typed enumeration for which the default underlying type is \lstinline!int! and can also be specified explicitly.}
\item{\seealsoref{enumopaque}{\seealsolocationc}offers a means to insulate individual enumerators from clients.}
\end{itemize}

\subsection[Further Reading]{Further Reading}\label{further-reading}

\begin{itemize}
\item{``Item 10: Prefer scoped \texttt{enum}s to unscoped \texttt{enum}s," \cite{meyers15}}
\item{\cite{grimm17}}
\end{itemize}



\newpage
%\section[User-Defined Literals]{User-Defined Literal Operators}%\label{}
% 8 March 2021, ready for Josh's code check
% 9 March 2021 JMB - code compiles, corrected u8 bugs
% 17 March 2021, copyedits in and proofed


\emcppsFeature{
    short={User-Defined Literals},
    long={User-Defined Literal Operators},
}{userdeflit}
\setcounter{table}{0}
\setcounter{footnote}{0}
\setcounter{lstlisting}{0}
%\section[User-Defined Literals]{User-Defined Literal Operators}%\label{}


C++11 allows developers to define a new suffix for a \emph{numeric},
\emph{character}, or \emph{string} literal, enabling a convenient
lexical representation of the \emph{value} of a \emcppsgloss[user defined type (UDT)]{user-defined
type (UDT)} or even a novel notation for the value of a built-in type.

\subsection[Description]{Description}\label{description-userdeflit}

A \emcppsgloss{literal} is a single token in a program that represents, in C and classic C++, a
value of an integer, floating-point, character, string, Boolean, or
pointer type.

Examples of familiar literal tokens are integer literals \lstinline!19! and
\lstinline!0x13!, each representing an \lstinline!int! having a value of 19;
floating-point literals \lstinline!0.19! and \lstinline!1.9e-1!, each
representing a \lstinline!double! having a value of 0.19; character
literals \lstinline!'a'! and \lstinline!\141!, 
each representing a
\lstinline!char! having the (ASCII) value for the letter ``a''; string
literal, \lstinline!"hello"!, representing a null-terminated array
containing the six characters \lstinline!'h'!, \lstinline!'e'!, \lstinline!'l'!,
\lstinline!'l'!, \lstinline!'o'!, and \lstinline!'\0'!; 
and Boolean
keyword literals \lstinline!true! and \lstinline!false!, representing the
corresponding Boolean values. C++11 added the keyword literal,
\lstinline!nullptr! (see \featureref{\locationa}{null-pointer-literal-(nullptr)}), representing the
null pointer value.

Both integer and floating-point literals have always had suffixes to
identify other numeric C++ types. For example, \lstinline!123L! and
\lstinline!123ULL! are literals of type \lstinline!signed!~\lstinline!long! and\linebreak[4]%%%%%   
\lstinline!unsigned!~\lstinline!long!~\lstinline!long!, respectively, both having
a decimal value 123, whereas \lstinline!123.f! is a literal of type
\lstinline!float! having precisely the decimal value 123. We can easily
distinguish programmatically between these different types of literals
using, e.g., overload resolution:

\begin{emcppslisting}
void f(const int&);            // (1) overload for type (ù{\codeincomments{int}}ù)
void f(const long&);           // (2)    "      "   "   (ù{\codeincomments{long}}ù)
void f(const double&);         // (3)    "      "   "   (ù{\codeincomments{double}}ù)
void f(const float&);          // (4)    "      "   "   (ù{\codeincomments{float}}ù)
void f(const unsigned int&);   // (5)    "      "   "   (ù{\codeincomments{unsigned int}}ù)
void f(const unsigned long&);  // (6)    "      "   "   (ù{\codeincomments{unsigned long}}ù)

void test0()
{
    f(123);    // OK, calls (1)
    f(123L);   // OK, calls (2)
    f(123.);   // OK, calls (3)
    f(123.f);  // OK, calls (4)
    f(123U);   // OK, calls (5)
    f(123Lu);  // OK, calls (6)
    f(123.L);  // Error, call to (ù{\codeincomments{f(long double)}}ù) is ambiguous
    f(123f);   // Error, invalid hex digit (ù{\codeincomments{f}}ù) in decimal constant
}
\end{emcppslisting}
    
\noindent Notice that applying an \lstinline!L! or \lstinline!l! suffix to a
floating-point literal (of default type \lstinline!double!) identifies it
as being of type \lstinline!long!~\lstinline!double!, which is a
\emcppsgloss{standard conversion} away from both \lstinline!float! and
\lstinline!double!, making the call ambiguous. Applying \lstinline!F! or
\lstinline!f! to an integer literal is, by default, not permitted unless a
user-defined literal (UDL) of a compatible type can be found.

Classic C++ allowed values of only built-in types to be represented as
compile-time literals. To express a hard-coded value of a UDT, a
developer would need to use a \emcppsgloss{value constructor} or
\emcppsgloss{factory function} and, unlike literals of built-in types, these
\emph{runtime} workarounds could never be used in a \emcppsgloss{constant
expression}. For example, we might want to create a user-defined type,
\lstinline!Name!, that can construct itself from a null-terminated string:

\begin{emcppslisting}[emcppsbatch=e1]
class Name  // user-defined type constructible from a literal string
{
    // ...

public:
    Name(const char*);  // value constructor taking a null-terminated string
    //...
};
\end{emcppslisting}
    
\noindent We can then initialize a variable of type \lstinline!Name! from a string
literal using the value constructor:

\begin{emcppslisting}[emcppsbatch=e1]
Name nameField("Maria");  // (ù{\codeincomments{Name}}ù) object having value (ù{\codeincomments{"Maria"}}ù)
\end{emcppslisting}
    
\noindent Alternatively, we could create one or more \emcppsgloss[factory function]{factory functions}
that return a constructed object appropriately configured with the
desired value. Multiple factory functions having different names can be
created without necessarily adding new constructors to the definition of
the returned type though, in some cases, a factory function might be a
\lstinline!friend! of the type it configures:

\begin{emcppshiddenlisting}[emcppsbatch=e2]
class Temperature;
bool operator==(const Temperature&,const Temperature&);
\end{emcppshiddenlisting}
\begin{emcppslisting}[emcppsbatch=e2]
#include <cassert>  // standard C (ù{\codeincomments{assert}}ù) macro

class Temperature { /*...*/ };

Temperature fahrenheit(double degrees);  // configured from degrees Fahrenheit
Temperature celsius(double degrees);     // configured from degrees Celsius

void test1()
{
    Temperature t1 = fahrenheit(32);  // water freezes at this temperature
    Temperature t2 = celsius(0.0);    //   "      "     "  "        "
    assert(t1 == t2);                 // expect same type and same value
}
\end{emcppslisting}
    
\noindent  Note that, as of C++11, the two functional constructs
described above can be declared \lstinline!constexpr! and thus be eligible
to be evaluated as part of a \emcppsgloss{constant expression}; see
\featureref{\locationc}{constexprfunc}. 

Although usable, the aforementioned C++03 workarounds for representing
literal values of a UDT lack the compactness and expressiveness of those
for built-in types. A fundamental design goal of C++ has always been to
minimize such differences. To that end, C++11 extends the notion of
\emcppsgloss[type suffix]{type suffix} to include user-definable identifiers with
a leading underscore (e.g., \lstinline!_name!, \lstinline!_time!,
\lstinline!_temp!):

\begin{emcppslisting}[emcppsbatch=e2]
Temperature operator""_F(long double degrees) { /*...*/ }  // define suffix (ù{\codeincomments{\_F}}ù)
Temperature operator""_C(long double degrees) { /*...*/ }  // define suffix (ù{\codeincomments{\_C}}ù)

void test2()  // same as (ù{\codeincomments{test1}}ù) above, but this time with user-defined literals
{
    Temperature t1 = 32.0_F; // water freezes at 32 degrees Fahrenheit
    Temperature t2 = 0.0_C;  //   "      "     "  0 degrees Celsius
    assert(t1 == t2);        // expect same type and same value
}
\end{emcppslisting}
    
\noindent The example above demonstrates the basic idea of a \emcppsgloss[user defined
literal (UDL)]{UDL} implemented as a new kind of operator:
\lstinline!operator""! followed by a suffix name. A \emcppsgloss[user defined
literal (UDL)]{UDL} is a
literal token having a \emcppsgloss{UDL suffix} that names a \emcppsgloss{UDL
operator}. A \emcppsgloss[user defined
literal (UDL)]{UDL} \emph{uses} a suffix, whereas a \emcppsgloss{UDL
operator} \emph{defines} a suffix. A \emcppsgloss{UDL suffix} must be a valid
identifier that begins with an underscore (\lstinline!_!). Note that the
leading underscore is \emph{not} required for \emcppsgloss[UDL suffix]{UDL suffixes}
defined by the C++ Standard Library. A \emcppsgloss[user defined
literal (UDL)]{UDL} is formed by
appending a \emcppsgloss{UDL suffix} to a native literal of one of four
\emcppsgloss[type category]{type categories}: \emcppsgloss[integer literal]{integer literals} (e.g.,
\lstinline!2020_year!), \emcppsgloss[floating point literal]{floating-point literals} (e.g.,
\lstinline!98.6_F!), \emcppsgloss[character literal]{character literals} (e.g.,
\lstinline!'x'_ebcdic!), and \emcppsgloss[string literal]{string literals} (e.g.,
\lstinline!"1!~\lstinline!Pennsylvania!~\lstinline!Ave"_validated!). Regardless
of the type category, a \emcppsgloss[user defined literal (UDL)]{UDL} can evaluate to any built-in or
user-defined type. What's more, the same \emcppsgloss{UDL suffix} can apply
to more than one of the four UDL type categories enumerated above,
potentially yielding a different type for each category.

Each UDL operator is effectively a \emcppsgloss{factory operator} that takes
a highly constrained parameter list (see \intrarefsimple{user-defined-literal-(udl)-operators}) 
and returns an
appropriately configured object. The defined \emcppsgloss{UDL suffix}, like
postfix operators \lstinline!++! and \lstinline!--!, names a function to be
called, but is parsed as part of the preceding literal (with no
intervening whitespace) and cannot be applied to an arbitrary run-time
expression. In fact, the compiler does not always produce an \emph{a
priori} interpretation of the literal before invoking the \emcppsgloss{UDL
operator}.

Let's now see how we might create a \emcppsgloss[user defined literal (UDL)]{UDL} for our original
user-defined type, \lstinline!Name!. In this example, the \emcppsgloss[user defined literal (UDL)]{UDL}
``function'' operates on a string literal, rather than a floating-point
literal:

\begin{emcppshiddenlisting}[emcppsbatch=e3]
#include <cstddef>  // (ù{\codeincomments{std::size\_t}}ù)
struct Name
{
public:
    Name(const char*);
};
\end{emcppshiddenlisting}
\begin{emcppslisting}[emcppsbatch=e3]
Name operator""_Nm(const char* n, std::size_t /* length */) { return Name(n); }
    // user-defined literal (UDL) operator for UDL suffix (ù{\codeincomments{\_Nm}}ù)

Name nameField = "Maria"_Nm;  // (ù{\codeincomments{Name}}ù) object having value "Maria"
\end{emcppslisting}
    
\noindent The UDL definition for string literals above, \lstinline!operator""_Nm! in
this case, takes two arguments: a \lstinline!const!~\lstinline!char*!
representing a null-terminated character string and a
\lstinline!std::size_t! representing the length of that string excluding
the null terminator. In our example, we ignore the second argument in
the body of the UDL operator, but it must nonetheless be present in the
parameter list.

\emcppsgloss[UDL operator]{UDL operators}, like \emcppsgloss[factory function]{factory functions}, can return a
value of any type, including built-in types. Unit-conversion functions,
for example, often return built-in types normalized to specific units:

\begin{emcppslisting}[emcppsbatch={e4,e5}]
#include <ctime>  // (ù{\codeincomments{std::time\_t}}ù)

constexpr std::time_t minutes(int m) { return m * 60; }    // minutes to seconds
constexpr std::time_t hours(int h)   { return h * 3600; }  // hours to seconds
\end{emcppslisting}
    
\noindent Each of the unit-conversion functions above returns an
\lstinline!std::time_t! (a standard type alias for a built-in integral
type) representing a duration in seconds. We can combine such uniform
quantities initialized from values in disparate units as needed:

\begin{emcppslisting}[emcppsbatch=e4]
std::time_t duration = hours(3) + minutes(15);   // 3.25 hours as seconds
\end{emcppslisting}
    
\noindent Replacing the unit-conversion functions with UDLs allows us to express
the desired value with a more natural-looking syntax. We now define two
new UDL operators, creating suffixes \lstinline!_min! for minutes and
\lstinline!_hr! for hours, respectively:

\begin{emcppslisting}[emcppsbatch=e5]
std::time_t operator""_min(unsigned long long m)
{
    return static_cast<std::time_t>(m * 60);  // minutes-to-seconds conversion
}

std::time_t operator""_hr(unsigned long long h)
{
    return static_cast<std::time_t>(h * 3600);  // hours-to-seconds conversion
}

std::time_t duration = 3_hr + 15_min;  // 3.25 hours as seconds
\end{emcppslisting}
    
\noindent We are not done yet. Unlike built-in literals, \emcppsgloss[user defined literal (UDL)]{UDLs} that use the suffixes defined above cannot express arbitrary literal values
that can be treated as compile-time constants usable in
\emcppsgloss[constant expression]{constant expressions} such as sizing an array or within a
\lstinline!static_assert!:

\begin{emcppslisting}
int a1[5_hr];                    // Error, (ù{\codeincomments{5\_hr}}ù) is not a compile-time constant.
static_assert(1_min == 60, "");  // Error, (ù{\codeincomments{1\_min}}ù) "  "  "    "      "     "
\end{emcppslisting}
    
\noindent Typical definitions of UDLs will, therefore, also involve another C++11
feature, \lstinline!constexpr! functions (see \featureref{\locationc}{constexprfunc}). By simply adding
\lstinline!constexpr! to the declaration of our UDL operators, we enable
them to be evaluated at compile-time and, hence, usable in
\emcppsgloss[constant expression]{constant expressions}:

\begin{emcppshiddenlisting}[emcppsbatch=e6]
#include <ctime>  // (ù{\codeincomments{std::time\_t}}ù)
\end{emcppshiddenlisting}
\begin{emcppslisting}[emcppsbatch=e6]
constexpr std::time_t operator""_Min(unsigned long long m)
{
    return static_cast<std::time_t>(m * 60);  // minutes-to-seconds conversion
}

constexpr std::time_t operator""_Hr(unsigned long long h)
{
    return static_cast<std::time_t>(h * 3600);  // hours-to-seconds conversion
}

int a2[5_Hr];                    // OK, (ù{\codeincomments{5\_Hr}}ù) is a compile-time constant.
static_assert(1_Min == 60, "");  // OK, (ù{\codeincomments{1\_Min}}ù) " "    "     "      "
\end{emcppslisting}
    
\noindent In short, a UDL operator is a new kind of \emcppsgloss{free operator} that
(1) may be defined with certain specific signatures limited to a subset
of the built-in types and (2) is invoked automatically when used as a
suffix of a built-in literal. There are, however, multiple ways to
define UDL operators --- \emph{precomputed-argument}, \emph{raw}, and
\emph{template} --- each more expressive and complex than the previous.
The allowed variations on the definitions of \emcppsgloss[UDL operator]{UDL operators}
are elucidated below.

\subsubsection[Restrictions on UDLs]{Restrictions on UDLs}\label{restrictions-on-user-defined-literals-(udls)}

A UDL suffix can be defined for only integer, floating-point, character,
and string literals. Left deliberately unsupported are the two
\emph{Boolean} literals, \lstinline!true! and \lstinline!false!, and the
\emph{pointer} literal, \lstinline!nullptr! (see \featureref{\locationa}{null-pointer-literal-(nullptr)}). These omissions serve to sidestep
lexical ambiguities --- e.g., the Boolean literal, \lstinline!true!,
combined with the \emcppsgloss{UDL suffix} \lstinline!_fact! would be
indistinguishable from the identifier, \lstinline!true_fact!.

We'll refer to the sequence of characters that make up a literal
excluding any suffix as a \emcppsgloss{naked literal} --- e.g., for literal
\lstinline!"abc"_udl!, the \emcppsgloss{naked literal} is \lstinline!"abc"!. UDL
suffixes can be appended to otherwise valid lexical literal tokens only.
That is, appending a UDL suffix to a token that
wouldn't be considered a valid lexical literal without the suffix is not permitted.
Though creating a suffix, \lstinline!_ipv4!, to
represent an IPv4 Internet address consisting of four octets separated
by periods might be tempting, \lstinline!192.168.0.1! would not be a valid lexical token in a
program and, hence, neither would \lstinline!192.168.0.1_ipv4!.
Interestingly, creating a \emcppsgloss[user defined literal (UDL)]{UDL} that
would cause an overflow if interpreted without the suffix \emph{is} permitted. Thus, for
example, the \emcppsgloss[user defined literal (UDL)]{UDL} \lstinline!0x123456789abcdef012345678_verylong!,
which comprises 24 hex digits and the \emcppsgloss{UDL suffix}
\lstinline!_verylong!, would be valid even on an architecture whose native
integers cannot exceed 64 bits (16 hex digits); see \featureref{\locationa}{long-long}.

Note that there are no \emph{negative} numeric literals in C++. The
negative numeric value \lstinline!-123! is represented as two separate
tokens: the negation operator and a positive literal, \lstinline!123!.
Similarly, an expression like \lstinline!-3_t1! will attempt to apply the
negation operator to the object of, say, \lstinline!Type1! produced by the
\lstinline!3_t1! \emcppsgloss[user defined literal (UDL)]{UDL}, which in turn comes from passing an
\lstinline!unsigned!~\lstinline!long!~\lstinline!long! having the positive value
3 to the \emcppsgloss{UDL operator} \lstinline!operator""_t1!. Such an
expression will be \emcppsgloss{ill formed} unless there is a negation
operator that operates on \emph{rvalues} of type \lstinline!Type1!. None of
the three forms of \emcppsgloss[UDL operator]{UDL operators} (see \intrarefsimple{prepared-argument-udl-operators}, \intrarefsimple{raw-udl-operators}, and \intrarefsimple{templated-udl-operators}) needs, or
is able, to handle a \emcppsgloss{naked literal} representing a negative
number.

When two or more string literals appear without intervening tokens, they
are concatenated and treated as a single string literal. If at least one of
those strings has a UDL suffix, then the suffix is applied to the concatenation
of the \emcppsgloss{naked literal} strings. If more than one of the strings has a
suffix, then all such suffixes must be the same \emcppsgloss{UDL suffix};
concatenating string literals having different suffixes is not
permitted, but strings having no suffix may be concatenated with strings
having a \emcppsgloss{UDL suffix}:

\begin{emcppslisting}
#include <cstddef>  // (ù{\codeincomments{std::size\_t}}ù)

struct XStr { /*...*/ };
XStr operator""_X(const char* n, std::size_t length);
XStr operator""_Y(const char* n, std::size_t length);

char a[] = "hello world";            // single native string literal
char b[] = "hello"      " world";    // native equivalent to "hello world"
XStr c   = "hello world"_X;          // user-defined string literal
XStr d   = "hello"_X    " world";    // UDL equivalent to "hello world"_X
XStr e   = "hello"      " world"_X;  //  "      "      "     "     "
XStr f   = "hello"_X    " world"_X;  //  "      "      "     "     "
XStr g   = "hel"_X "lo" " world"_X;  //  "      "      "     "     "
XStr h   = "hello"_X    " world"_Y;  // Error, mixing UDL suffixes (ù{\codeincomments{\_X}}ù) and (ù{\codeincomments{\_Y}}ù)
\end{emcppslisting}
    
\noindent Finally, combining a \emcppsgloss{UDL suffix} with a
second built-in or user-defined suffix on a single token is not possible. Writing
\lstinline!45L_Min!, for example, in an attempt to combine the \lstinline!L!
suffix (for \lstinline!long!) with the \lstinline!_Min! suffix (for the
user-defined minutes suffix described earlier) will
simply yield the undefined and invalid suffix, \lstinline!L_Min!.

\subsubsection[UDL operators]{UDL operators}\label{user-defined-literal-(udl)-operators}

A \emcppsgloss{UDL suffix} (e.g., \lstinline!_udl!) is created by defining a
\emcppsgloss{UDL operator} (e.g., \lstinline!operator""_udl!) that follows a
strict set of rules described in this section. In the declaration and
definition of a \emcppsgloss{UDL operator}, the name of the \emcppsgloss{UDL
suffix} \emph{may} be separated from the quotes by whitespace; in fact,
some older compilers might even \emph{require} such whitespace due to a
defect (since corrected and retroactively applied)  in the original C++11 specification. Thus, for all but the oldest C++11 compilers,
\lstinline!operator""_udl!, \lstinline!operator""!~\lstinline!_udl!, and
\lstinline!operator!~\lstinline!""!~\lstinline!_udl! are all valid spellings of
the same \emcppsgloss{UDL operator} name. Note that whitespace is \emph{not}
permitted between a literal and its suffix in the \emph{use} of a
\emcppsgloss[user defined literal (UDL)]{UDL}. For example, \lstinline!1.2!~\lstinline!_udl! is ill formed;
\lstinline!1.2_udl! must appear as a single token with no whitespace.

A \emcppsgloss[user defined literal (UDL)]{UDL} generally consists of two parts: (1) a valid lexical
literal token and (2) a user-defined suffix. The signature of each
\emcppsgloss{UDL operator} must conform to one of three patterns,
distinguished by the way the compiler supplies the \emcppsgloss{naked
literal} to the \emcppsgloss{UDL operator}:

\begin{enumerate}
\item{\emcppsgloss[prepared-argument UDL operator]{Prepared-argument UDL operator} — The \emcppsgloss{naked literal} is evaluated at compile time and passed into the operator as a value:
\begin{emcppshiddenlisting}[emcppsbatch=e7]
struct Type1 {};
struct Type2 {};
struct Type3 {};
\end{emcppshiddenlisting}
\begin{emcppslisting}[emcppsbatch=e7]
Type1 operator"" _t1(unsigned long long n);
Type1 t1 = 780_t1;  // calls (ù{\codeincomments{operator""\_t1(780ULL)}}ù)
\end{emcppslisting}
    }
\item{\emcppsgloss[raw UDL operator]{Raw UDL operator} — The characters that make up the \emcppsgloss{naked literal} are passed to the operator as a \emph{raw}, unevaluated string (for numeric literals only):
\begin{emcppslisting}[emcppsbatch=e7]
Type2 operator"" _t2(const char *token);
Type2 t2 = 780_t2;  // calls (ù{\codeincomments{operator""\_t2("780")}}ù)
\end{emcppslisting}
    }
\item{\emcppsgloss[templated UDL operator]{Templated UDL operator} — The UDL operator is a template whose parameter list is a variadic sequence of \lstinline!char! values (see \featureref{\locationc}{variadictemplate}) that make up the \emcppsgloss{naked literal} (for numeric literals only):
\begin{emcppslisting}[emcppsbatch=e7]
template <char...> Type3 operator"" _t3();
Type3 t3 = 780_t3;  // calls (ù{\codeincomments{operator""\_t3<'7', '8', '0'>()}}ù)
\end{emcppslisting}
    }
\end{enumerate}

Each of these three forms of \emcppsgloss[UDL operator]{UDL
operators} is expounded in more detail in its own separate section; see \intrarefsimple{prepared-argument-udl-operators}, \intrarefsimple{raw-udl-operators}, and \intrarefsimple{templated-udl-operators}.

When a \emcppsgloss[user defined literal (UDL)]{UDL} is encountered, the compiler prioritizes a
\emcppsgloss{prepared-argument UDL operator} over the other two. Given a
\emcppsgloss[user defined literal (UDL)]{UDL} having suffix \lstinline!_udl!, the compiler will look for
any \lstinline!operator""_udl! in the local scope (\emcppsgloss{unqualified
name lookup}). If, among the operators found, there is a
\emcppsgloss{prepared-argument UDL operator} that exactly matches the type of
the \emcppsgloss{naked literal}, then that \emcppsgloss{UDL operator} is called.
Otherwise, for numeric literals only, the \emcppsgloss[raw UDL operator]{raw} or
\emcppsgloss{templated UDL operator} --- only one of which is permitted to exist
for a given suffix --- is invoked. This set of lookup rules is deliberately
short and rigid. Importantly, this lookup sequence differs from other
operator invocations in that it does \emph{not} involve \emcppsgloss{overload
resolution} or argument conversions, nor does it employ
\emcppsgloss[argument-dependent lookup (ADL)]{argument-dependent lookup (ADL)} to find operators in
other namespaces.

Although \emcppsgloss[argument-dependent lookup (ADL)]{ADL} is never an issue for \emcppsgloss[user defined literal (UDL)]{UDLs}, common
practice is to gather related \emcppsgloss[UDL operator]{UDL operators} into a namespace
(whose name often contains the word ``literals''). This namespace is
then typically nested within the namespace containing the definitions of
the user-defined types that the \emcppsgloss[UDL operator]{UDL operators} return. These
literals-only nested namespaces enable a user to import, via a single
\lstinline!using! directive, just the literals into their scope, thereby
substantially decreasing the likelihood of name collisions:

\begin{emcppshiddenlisting}[emcppsbatch=e8]
#include <cstddef>  // (ù{\codeincomments{std::size\_t}}ù)
struct Type1 {};
\end{emcppshiddenlisting}
\begin{emcppslisting}[emcppsbatch=e8]
namespace ns1  // namespace containing types returned by UDL operators
{
    struct Type1 { };
    bool check(const Type1&);

    namespace literals  // nested namespace for UDL operators returning (ù{\codeincomments{ns1}}ù) types
    {
        Type1 operator"" _t1(const char*, std::size_t);
    }

    using namespace literals;  // Make literals available in namespace (ù{\codeincomments{ns1}}ù).
}

void test1()  // file scope: finds UDL operator via (ù{\codeincomments{using}}ù) directive
{
    using namespace ns1::literals;  // OK, imports only the *inner* UDL operators
    check("xyzzy"_t1);              // OK, finds (ù{\codeincomments{ns1::check}}ù) via ADL
}
\end{emcppslisting}
    
\noindent To use the \lstinline!_t1! \emcppsgloss{UDL suffix} above, \lstinline!test1! must
somehow be able to find the declaration of its corresponding \emcppsgloss{UDL
operator} locally, which is accomplished by placing the operator in a
nested namespace and importing the entire namespace via a \emcppsgloss[using
directive]{\lstinline!using! directive}. We could have avoided the nested namespace and, instead, required each needed operator to be imported individually:

\begin{emcppshiddenlisting}[emcppsbatch=e9]
#include <cstddef>  // (ù{\codeincomments{std::size\_t}}ù)
struct Type2 {};
\end{emcppshiddenlisting}
\begin{emcppslisting}[emcppsbatch=e9]
namespace ns2  // namespace defining types returned by non-nested UDL operators
{
    struct Type2 { };
    bool check(const Type2&);

    Type2 operator"" _t2(const char*, std::size_t);  // BAD IDEA: not nested
}

void test2()  // file scope: finds UDL operator via (ù{\codeincomments{using}}ù) declaration
{
    using ns2::operator"" _t2;  // OK, imports just the needed UDL operator
    check("xyzzy"_t2);          // OK, finds (ù{\codeincomments{ns2::check}}ù) via ADL
}
\end{emcppslisting}
    
\noindent When multiple \emcppsgloss[UDL operator]{UDL operators} are provided for a collection of
types, however, the idiom of placing just the \emcppsgloss[UDL operator]{UDL operators} in
a nested namespace (typically incorporating the name ``literals'')
obviates most of the commonly cited ill effects (e.g., accidental
unwanted name collisions) attributed to more general use of
\lstinline!using! directives. In the interest of brevity, we will freely
omit the nested-literal namespaces in expository-only examples.

Finally, despite its use in the Standard for this specific purpose,
there is never a need for a namespace comprising only UDLs to be
declared \lstinline!inline! and doing so is contraindicated; see \featureref{\locatione}{inline-namespaces}.

\subsubsection[Prepared-argument UDL operators]{Prepared-argument UDL operators}\label{prepared-argument-udl-operators}

The \emph{prepared-argument} pattern for \emcppsgloss[UDL operator]{UDL operators} is
supported for all four of the \emcppsgloss[UDL type category]{UDL type categories}: integer,
floating-point, character, and string. If this form of \emcppsgloss{UDL
operator} is selected (see \intrarefsimple{user-defined-literal-(udl)-operators}), the compiler
first determines which of the four \emcppsgloss[UDL type category]{UDL type categories} applies
to the \emcppsgloss{naked literal}, evaluates it without regard to the
\emcppsgloss{UDL suffix}, and then passes the prepared (i.e., precomputed)
value to the \emcppsgloss{UDL operator}, which further processes its argument
and returns the value of the UDL. Note that the \emcppsgloss{UDL type
category} of the \emcppsgloss{UDL operator} refers to only the \emcppsgloss{naked
literal}; its return value can be any arbitrary type, with or without an
obvious relationship to its type category:

\begin{emcppshiddenlisting}[emcppsbatch=e10]
#include <cstddef>  // (ù{\codeincomments{std::size\_t}}ù)
\end{emcppshiddenlisting}
\begin{emcppslisting}[emcppsbatch=e10]
struct Smile { /* ... */ };  // arbitrary user-defined type

Smile operator"" _fx(long double);  // floating-point literal returning (ù{\codeincomments{Smile}}ù)
float operator"" _ix(unsigned long long);  // integer literal returning (ù{\codeincomments{float}}ù)
int operator"" _sx(const char*, std::size_t);  // string literal returning (ù{\codeincomments{int}}ù)
\end{emcppslisting}
    
\noindent The \emcppsgloss{UDL type category} for a \emcppsgloss{prepared-argument UDL
operator} is determined by the operator's signature; integer and
floating-point \emcppsgloss[UDL operator]{UDL operators} each take a single argument of,
respectively, the largest unsigned integer or floating-point type
defined by the language. Multiple \emcppsgloss[prepared-argument UDL
operator]{prepared-argument UDL
operators} may be declared in the same scope for the same \emcppsgloss{UDL
suffix} (e.g., \lstinline!_aa!) and each may return a distinct arbitrary
type (e.g., \lstinline!short!, \lstinline!Smile!):

\begin{emcppslisting}[emcppsbatch=e10]
short operator"" _aa(unsigned long long n);  // integer literal operator
Smile operator"" _aa(long double n);         // floating-point literal operator
bool  operator"" _bb(long double n);         // floating-point literal operator
\end{emcppslisting}
    
\noindent The \emcppsgloss{naked literal} is evaluated as an
\lstinline!unsigned!~\lstinline!long!~\lstinline!long! for integer literals (see \featureref{\locationa}{long-long}) or a
\lstinline!long!~\lstinline!double! for floating-point literals; the prepared
value is then passed via the \lstinline!n! parameter to the \emcppsgloss{UDL
operator}. The compiler does the work of parsing and evaluating the
sequence of digits, radix prefixes (e.g., \lstinline!0! for octal and
\lstinline!0x! for hexadecimal), decimal points, and exponents. C++14
provides additional features related to built-in
literals; see \featureref{\locationb}{binary-literals} and \featureref{\locationb}{digitseparator}:

\begin{emcppslisting}[emcppsbatch=e10]
short v1 = 123_aa;  // OK, invokes (ù{\codeincomments{operator"" \_aa(unsigned long long)}}ù)
Smile v2 = 1.3_aa;  // OK, invokes (ù{\codeincomments{operator"" \_aa(long double)}}ù)
\end{emcppslisting}
    
\noindent There is no \emcppsgloss{overload resolution},
\emcppsgloss{integer-to-floating-point conversion}, or
\emcppsgloss{floating-point-to-integer conversion}, nor are \emcppsgloss[UDL operator]{UDL
operators} having numerically lower precision permitted:

\begin{emcppslisting}[emcppsbatch=e10]
Smile operator"" _cc(double n);  // Error, invalid parameter list
\end{emcppslisting}
    
\noindent If lookup does not find a \emcppsgloss{UDL operator} that matches the type
category of the \emcppsgloss{naked literal} exactly, the match fails:

\begin{emcppslisting}[emcppsbatch=e10]
bool v3 = 123_bb;  // Error, unable to find integer literal (ù{\codeincomments{\_bb}}ù)
bool v4 = 1.3_bb;  // OK, invokes (ù{\codeincomments{operator"" \_bb(long double)}}ù)
\end{emcppslisting}
    
\noindent Because the \emcppsgloss{naked literal} is fully evaluated by the compiler,
overflow and precision loss could become issues just as they are for
native literals. These limitations vary by platform, but typical
platforms are limited to 64-bit
\lstinline!unsigned!~\lstinline!long!~\lstinline!long! or 64-bit
\lstinline!long!~\lstinline!double! types in IEEE
754 format:

\begin{emcppslisting}[emcppsbatch=e10]
Smile v5 = 1.2e310_aa;              // Bug, argument evaluates to infinity
Smile v6 = 2.5e-310_aa;             // Bug, argument evaluates to denormalized
short v7 = 0x1234568790abcdef0_aa;  // Error, doesn't fit in any integer type
\end{emcppslisting}
    
\noindent Note that the oversized integer initializer for \lstinline!v7! in the code snippet above
results in an error on some compilers but only a warning on others.

A rarely used feature of C++03 is the \emcppsgloss{encoding prefix},
\lstinline!L!, on a character or string literal. The literals \lstinline!'x'!
and \lstinline!"hello"! have (built-in) types \lstinline!char! and
\lstinline!char[6]!, respectively, whereas \lstinline!L'x'! and
\lstinline!L"hello"! have the respective types \lstinline!wchar_t! and
\lstinline!wchar_t[6]!. C++11 added three more string 
\emcppsgloss[encoding prefix]{encoding prefixes}: \lstinline!u! to indicate \emcppsgloss{UTF-16}
with a type of \lstinline!const char16_t*!, \lstinline!U! to indicate \emcppsgloss{UTF-32} 
with a type \lstinline!const char32_t*!, and \lstinline!u8! to indicate \emcppsgloss{UTF-8}
with a type of \lstinline!const char*!.  The \lstinline!u! and \lstinline!U! prefixes can also
be used on character literals; see  \featureref{\locationa}{unicode-string-and-character-literals}.{\cprotect\footnote{C++17 allows the \lstinline!u8! prefix on
  character literals as well as on string literals, while C++20 changes the type for
  \lstinline!u8! prefixed string and character literals from \lstinline!char! to \lstinline!char8_t!.}}

The four C++11 \emcppsgloss[encoding prefix]{encoding prefixes} for character literals can
each be supported by a distinct \emcppsgloss{UDL operator} signature (e.g.,
\lstinline!_dd! below), each of which might return a distinct, arbitrary
type:

\begin{emcppslisting}[emcppsbatch=e10]
int         operator"" _dd(char     ch);  //  'x'
double      operator"" _dd(char16_t ch);  // u'x'
const char* operator"" _dd(char32_t ch);  // U'x'
Smile       operator"" _dd(wchar_t  ch);  // L'x'
\end{emcppslisting}
    
\noindent Any or all of the above forms can co-exist. A character \emcppsgloss{naked
literal} (e.g., \lstinline!'Q'!) is translated into the appropriate
character type and value in the \emcppsgloss{execution character set} (i.e.,
the set of characters used at run time on the target operating system)
and passed via the \lstinline!ch! parameter to the body of the \emcppsgloss{UDL
operator}. As ever, there are no narrowing or widening conversions, so
the program is ill formed if the precise signature of the needed
\emcppsgloss{UDL operator} is not found:

\begin{emcppslisting}[emcppsbatch=e10]
int    operator"" _ee(char);      //  'x'
double operator"" _ee(char32_t);  // U'x'

int         c1 =  'Q'_ee;  // OK, matches (ù{\codeincomments{char}}ù) parameter type
double      c2 = U'Q'_ee;  // OK, matches (ù{\codeincomments{char32\_t}}ù) parameter type
const char* c3 = u'Q'_ee;  // Error, no match for (ù{\codeincomments{char16\_t}}ù) parameter type
Smile       c4 = L'Q'_ee;  // Error, no match for (ù{\codeincomments{wchar\_t}}ù) parameter type
\end{emcppslisting}
    
\noindent Similarly, there are four valid \emcppsgloss{UDL operator} signatures for
string literals in C++11, each of which again might
return a different type{\cprotect\footnote{C++20 allows a fifth signature for
string literals with the \lstinline!u8! prefix that takes a \lstinline!const char8_t*!.}}:

\begin{emcppslisting}[emcppsbatch=e10]
bool  operator"" _dd(const char*     str, std::size_t len);  //   "str"
int   operator"" _dd(const char16_t* str, std::size_t len);  //  u"str"
float operator"" _dd(const char32_t* str, std::size_t len);  //  U"str"
Smile operator"" _dd(const wchar_t*  str, std::size_t len);  //  L"str"
\end{emcppslisting}
    
\noindent The string \emcppsgloss{naked literal} evaluates to a null-terminated
character array. The address of the first element of that array is
passed to the \emcppsgloss{UDL operator} via the \lstinline!str! parameter, and
its length --- excluding the null terminator --- is passed via \lstinline!len!.

To recap, multiple \emcppsgloss[prepared-argument UDL operator]{prepared-argument UDL operators} can co-exist
for a single \emcppsgloss{UDL suffix}, each having a distinct type category
and each potentially returning a different C++ type. To show another
example, the suffix \lstinline!_s! on a floating-point literal could
return an \lstinline!double! to mean seconds, whereas the same suffix on a
string literal might return an \lstinline!std::string!. Moreover, string and character 
\emcppsgloss[UDL operator]{UDL operators} that differ by character type (\lstinline!char!,
\lstinline!char16_t!, and so on) often have different return types.
Similar string \emcppsgloss[UDL operator]{UDL operators} typically --- but not necessarily
--- return similar types, such as \lstinline!std::string! and
\lstinline!std::u16string!, that differ only in their underlying character
type:

\begin{emcppslisting}
#include <string>   // (ù{\codeincomments{std::string}}ù)
#include <utility>  // (ù{\codeincomments{std::pair}}ù)

double operator"" _s(unsigned long long);  // integer UDL operator
double operator"" _s(long double);         // floating-point UDL operator

std::string    operator"" _s(const char*,     std::size_t);  // string UDL
std::u16string operator"" _s(const char16_t*, std::size_t);  //   "     "

double         d = 12_s;       // yields (ù{\codeincomments{double}}ù) having value (ù{\codeincomments{12.0}}ù)
std::u16string w = u"Hola"_s;  // yields (ù{\codeincomments{std::u16string}}ù) having value (ù{\codeincomments{Hola}}ù)
std::string    s = "Hello"_s;  // yields (ù{\codeincomments{std::string}}ù) having value (ù{\codeincomments{Hello}}ù)
\end{emcppslisting}
    
\noindent Note again that these operators are invariably declared
\lstinline!constexpr! (see \featureref{\locationc}{constexprfunc})
because they can \emph{always} be evaluated at compile time because
their arguments are necessarily expressed \emph{literally} in the source
code.

\subsubsection[Raw UDL operators]{Raw UDL operators}\label{raw-udl-operators}

The \emph{raw} pattern for \emcppsgloss[UDL operator]{UDL operators} is supported for only
the integer and floating-point \emcppsgloss[UDL type category]{UDL type categories}. If this
form of \emcppsgloss{UDL operator} is selected (see \intrarefsimple{user-defined-literal-(udl)-operators}), the compiler packages up the \emcppsgloss{naked literal} as an
unprocessed string --- i.e., a sequence of raw characters transferred
from the source --- and passes it to the \emcppsgloss{UDL operator} as a
null-terminated character string. All \emcppsgloss[raw UDL operator]{raw UDL operators} (e.g.,
for suffix \lstinline!_rl! below) have the same signature:

\begin{emcppslisting}
struct Type { /*...*/ };

Type operator"" _rl(const char*);

Type t1 = 425_rl;  // invokes (ù{\codeincomments{operator ""\_rl("425")}}ù)
\end{emcppslisting}
    
\noindent This signature can be distinguished from a \emcppsgloss{prepared-argument UDL
operator} for string literals by the \emph{absence} of an
\lstinline!std::size_t! parameter representing its \emph{length}.

The raw string argument will be verified by the compiler to be a
well-formed integer or floating-point literal token but is otherwise
untouched by the compiler. For any given \emcppsgloss{UDL suffix}, at most
one matching \emcppsgloss{raw UDL operator} may be in scope at a time; hence,
the \emph{return} type cannot vary based on, e.g., whether the
\emcppsgloss{naked literal} contains a decimal point. Such a capability
\emph{is}, however, available; see \intrarefsimple{templated-udl-operators}.

In particular, it might be the case that not all valid tokens accepted
by the compiler will satisfy the \emcppsgloss{narrow contract} offered by a
specific \emcppsgloss{UDL operator}. For example, we can define a \emcppsgloss{UDL
suffix}, \lstinline!_3!, to express base-3 integers using a \emcppsgloss{raw
UDL operator}:

\begin{emcppslisting}[emcppsbatch=e11]
int operator"" _3(const char* digits)
{
    int result = 0;

    while (*digits)
    {
        result *= 3;
        result += *digits - '0';
        ++digits;
    }

    return result;
}
\end{emcppslisting}
    
\noindent We can now test this function at run time using the standard
\lstinline!assert! macro:

\begin{emcppslisting}[emcppsbatch=e11]
#include <cassert>  // standard C (ù{\codeincomments{assert}}ù) macro

void test()
{
    assert( 0 ==   0_3);
    assert( 1 ==   1_3);
    assert( 2 ==   2_3);
    assert( 3 ==  10_3);
    assert( 4 ==  11_3);
    // ...
    assert( 8 ==  22_3);
    assert( 9 == 100_3);
    assert(10 == 101_3);
}
\end{emcppslisting}
    
\noindent Note that we could have declared our \lstinline!_3! \emcppsgloss{raw
UDL operator} \lstinline!constexpr! and replaced all of the (runtime)
\lstinline!assert! statements with (compile-time) \lstinline!static_assert! declarations. Also observe that, as written, \lstinline!constexpr!~\lstinline!operator""!~\lstinline!_3! requires C++14's relaxed \lstinline!constexpr! function restrictions (see \featureref{\locationc}{constexprfunc}), but it could be reimplemented to work in C++11.

Let's now consider valid lexical integer literals representing values
outside of what would be considered valid of base-3 integers:

\begin{emcppslisting}[emcppsbatch=e11]
int i1 = 22_3;                    // (1) OK, returns (ù{\codeincomments{(int) 8}}ù)
int i2 = 23_3;                    // (2) Bug, returns (ù{\codeincomments{(int) 9}}ù)
int i3 = 21.1_3;                  // (3) Bug, returns (ù{\codeincomments{(int) 58}}ù)
int i4 = 22211100022211100022_3;  // (4) Bug, too big for 32-bit (ù{\codeincomments{int}}ù)
\end{emcppslisting}
    
\noindent In the example code above, (1) is a valid base-3 integer; (2) is a valid
integer literal but contains the digit 3, which is \emph{not} a valid
base-3 digit; (3) is a valid floating-point literal, but the \emcppsgloss{UDL
operator} returns values of only type \lstinline!int!; and (4) is in
principle a valid integer literal but represents a value that is too
large to fit in a 32-bit \lstinline!int!. Because cases (2), (3), and (4)
are valid lexical literals, it is up to the implementation of the
\emcppsgloss{UDL operator} to reject invalid values.

Let's now consider a more robust implementation of a base-3 integer
\emcppsgloss[user defined literal (UDL)]{UDL}, \lstinline!_3b!, that throws an exception when the literal
fails to represent a valid base-3 integer:

\begin{emcppslisting}
#include <stdexcept>  // (ù{\codeincomments{std::out\_of\_range}}ù), (ù{\codeincomments{std::overflow\_error}}ù)
#include <limits>     // (ù{\codeincomments{std::numeric\_limits}}ù)

int operator"" _3b(const char *digits)
{
    int ret = 0;

    for (char c = *digits; c; c = *++digits)
    {
        if ('\'' == c)  // Ignore the C++14 digit separator.
        {
            continue;
        }

        if (c < '0' || '2' < c)  // Reject non-base-3 characters.
        {
            throw std::out_of_range("Invalid base-3 digit");
        }

        if (ret >= (std::numeric_limits<int>::max() - (c - '0')) / 3)
        {
            // Reject if (ù{\codeincomments{3 * ret + (c - '0')}}ù) would overflow.
            throw std::overflow_error("Integer too large");
        }

        ret = 3 * ret + (c - '0');  // Consume (ù{\codeincomments{c}}ù).
    }

    return ret;
}
\end{emcppslisting}
    
\noindent In this implementation of a \emcppsgloss{raw UDL operator} for a suffix
\lstinline!_3b!, the first \lstinline!if! statement looks for the C++14-only digit separator and ignores it. The second \lstinline!if! statement
looks for characters not in the valid range for base-3 and, for cases (2) and (3), throws
an \lstinline!out_of_range! exception. The third
\lstinline!if! statement determines whether the computation is about to
overflow and, for case (4), throws an \lstinline!overflow_error! exception.
The absence of compiler interpretation makes \emcppsgloss[raw UDL operator]{raw UDL operators}
potentially difficult to write but also powerful. Interpreting
existing literal characters in a new way and accepting literals that
would otherwise overflow or lose precision confers additional expressiveness on raw UDLs; e.g., the base-3 \emcppsgloss{raw UDL operator} shown above could
not be expressed using a \emcppsgloss{prepared-argument UDL operator}. The downside of this additional expressiveness is the need to implement, debug, and maintain custom token parsing, including error checking.

If the intent is to consume integer literals, a \emcppsgloss{raw UDL
operator} needs to  either process or reject not only the
decimal digits \lstinline!'0'! to \lstinline!'9'!, but also the hex digits
\lstinline!'a'!--\lstinline!'f'! and \lstinline!'A'!--\lstinline!'F'! as well as
radix prefixes \lstinline!0!, \lstinline!0x!, and \lstinline!0X!. In C++14, the operator 
must also handle the \lstinline!0b! (binary) radix prefix and the digit
separator \lstinline!'\''!.  

For floating-point literals, the \emcppsgloss{raw UDL operator} needs to
handle the decimal digits, decimal point, exponent prefix (\lstinline!'e'!
or \lstinline!'E'!), and optional exponent sign (\lstinline!'+'! or
\lstinline!'-'!). It is possible to handle both integer and floating-point
literals in a single \emcppsgloss{raw UDL operator}, provided the return type
is the same for both. In all cases, it is usually wise to reject
\emph{any} unexpected character, including characters that are not
currently legal within numeric literals, in case the set of legal
characters is enlarged in the future, e.g., by adding a new radix.

A \emcppsgloss{raw UDL operator} can be declared \lstinline!constexpr!, but be
aware that the rules for what is allowed in a \lstinline!constexpr!
function differ significantly between C++11 and C++14, as described in \featureref{\locationc}{constexprfunc}. In particular, the
loop-based implementation of \lstinline!operator""_3! (above) can be
declared \lstinline!constexpr! in C++14 but not in C++11, though it is
possible to define a \lstinline!constexpr! \emcppsgloss{UDL operator} in C++11
that has the same behavior by using a recursive implementation. If the
literal is evaluated as part of a constant expression and if the literal
contains errors that would result in exceptions being thrown (i.e., an
invalid character or overflow), the compiler will reject the invalid
literal at compile time. If, however, the literal is not part of a
constant expression, an exception will still be thrown at run time:

\begin{emcppslisting}[emcppsbatch=e11]
constexpr int i4 = 25_3;  // Error, "throw" not allowed in constant expression
          int i5 = 25_3;  // Bug, exception thrown at run time
\end{emcppslisting}
    
\noindent To ensure that every invalid literal is detected at compile time, use
templated UDL operators, as described in the next section.
%\intrarefnopage{templated-udl-operators}.

\subsubsection[Templated UDL operators]{Templated UDL operators}\label{templated-udl-operators}

A \emcppsgloss{templated UDL operator}, known as a \emph{literal operator
template} in the Standard, is a variadic template (see \featureref{\locationc}{variadictemplate}) having a template parameter list
consisting of a pack of an arbitrary number of \lstinline!char! parameters
and an empty runtime parameter list:

\begin{emcppslisting}[emcppsbatch=e12]
struct Type { /*...*/ };

template <char...> Type operator"" _udl();
\end{emcppslisting}
    
\noindent The \emcppsgloss{templated UDL operator} pattern supports only the integer
and floating-point type categories.{\cprotect\footnote{C++20 added
support for user-defined string literal templates, albeit with a
  different syntax.}} If this form of \emcppsgloss{UDL operator} is selected
(see \intrarefsimple{user-defined-literal-(udl)-operators}), the compiler breaks up the
\emcppsgloss{naked literal} into a sequence of raw characters and passes each
one as a separate template argument to the instantiation of the
\emcppsgloss{UDL operator}:

\begin{emcppslisting}[emcppsbatch=e12]
Type t1 = 42.5_udl;  // calls (ù{\codeincomments{operator""\_udl<'4', '2', '.', '5'>()}}ù)
\end{emcppslisting}
    
\noindent As in the case of \emcppsgloss[raw UDL operator]{raw UDL operators}, the raw sequence of
characters will be verified by the compiler to be a well-formed integer
or floating-point literal token, but the \emcppsgloss{UDL operator} must
deduce meaning from those characters. Unlike \emcppsgloss[raw UDL operator]{raw UDL operators},
a \emcppsgloss{templated UDL operator} can return different types based on
the content of the \emcppsgloss{naked literal}. For example, electrical
resistance might be expressed as \lstinline!Resistance<float>! or
\lstinline!Resistance<int>!, where the intention is to express resistance
less than 10 ohms using the \lstinline!float! specialization and larger
resistance using the \lstinline!int! specialization:

\begin{emcppslisting}[emcppsbatch=e12]
template <class T> class Resistance;
template <> class Resistance<int>   { /* ... (resistance >= 10 ohms) */ };
template <> class Resistance<float> { /* ... (resistance <  10 ohms) */ };
\end{emcppslisting}
    
\noindent The \lstinline!_ohms! \emcppsgloss{UDL operator} (in the code snippet below) determines the correct
return type and value based on the characters making up the naked
numeric literal. The compile-time template logic needed to make this
selection requires template metaprogramming, which in turn requires
partial template specialization. Function templates, including
\emcppsgloss[templated UDL operator]{templated UDL operators}, cannot have partial specializations,
so the selection logic is delegated to a helper class template,
\lstinline!MakeResistance!:

\begin{emcppslisting}[emcppsbatch=e12]
template <char C, char... Cs> struct MakeResistance;

template <char... Cs>
constexpr typename MakeResistance<Cs...>::ReturnType
operator"" _ohms() { return MakeResistance<Cs...>::factory(); }
\end{emcppslisting}
    
\noindent \lstinline!MakeResistance<c...>::ReturnType! is the template metafunction used to compute the return type. In C++14, the return type can
be deduced directly from the \lstinline!return! statement; see \featureref{\locationf}{Function-Return-Type-Deduction}.
The implementation of \lstinline!MakeResistance! uses partial
specializations to detect specific numeric literal patterns:

\begin{emcppslisting}[emcppsbatch=e12]
template <char d0, char... d>
struct MakeResistance  // primary template (for value >= 10)
{    
    using ReturnType = Resistance<int>;
    static constexpr ReturnType factory();
};

template <char d0>
struct MakeResistance<d0>  // specialize for single digit ('0' - '9')
{    
    using ReturnType = Resistance<float>;
    static constexpr ReturnType factory(); // resistance 0 to 9 ohms
};

template <char d0, char... d>
struct MakeResistance<d0, '.', d...>  // specialize for decimal point after 
                                      // first digit (e.g., (ù{\codeincomments{1.23}}ù))
{    
    using ReturnType = Resistance<float>;
    static constexpr ReturnType factory(); // resistance 0.0 to < 10.0
};
\end{emcppslisting}
    
\noindent The primary \lstinline!MakeResistance! template will return a
\lstinline!Resistance<int>! for any sequence of characters that does not
match one of the partial specializations. The first partial
specialization matches a single-digit integer literal (i.e., an integer
value from 0 through 9). The second partial specialization matches a
sequence of two or more characters where the second character is a
decimal point. Thus, any character sequence representing a value less
than 10 will return \lstinline!Resistance<float>!:

\begin{emcppslisting}[emcppsbatch=e12]
Resistance<int>   r1 = 200_ohms;
Resistance<float> r2 = 5_ohms;
Resistance<float> r3 = 2.5_ohms;
\end{emcppslisting}
    
\noindent Note that this simplified example does not recognize floating-point
literals such as \lstinline!12.5! or \lstinline!.04! that don't have their
decimal point as the second character. The description of fixed-point
literals in \intraref{use-cases-userdeflit}{user-defined-numeric-types} provides a more complete exposition of this sort
of return-type selection, including details of the recursive template
metaprogramming used to determine the return type and its value.

As template arguments, the characters that make up the \emcppsgloss{naked
literal} are constant expressions and can be used with
\lstinline!static_assert! to force error detection at compile time. Unlike
\emcppsgloss[raw UDL operator]{raw UDL operators}, there is no risk of throwing an exception at
run time, even when initializing a value in a non-\lstinline!constexpr!
context:

\begin{emcppslisting}[emcppsbatch=e12]
constexpr auto r4 = 12.5_ohms;  // Error, (ù{\codeincomments{constexpr}}ù) context
          auto r5 = 12.5_ohms;  // Bug, non-(ù{\codeincomments{constexpr}}ù) context
\end{emcppslisting}
    
\noindent Being able to select a context-specific return type and to force
compile-time error checking makes \emcppsgloss[templated UDL operator]{templated UDL operators} the
most expressive pattern for defining \emcppsgloss[UDL operator]{UDL operators}. These
capabilities come at the cost, however, of having to develop them using
the less-than-readable template sublanguage in C++.

\subsubsection[UDLs in the C++14 Standard Library]{UDLs in the C++14 Standard Library}\label{udls-in-the-c++14-standard-library}

This book is primarily about modern C++ \emph{language} features, but a
short description of \emcppsgloss[UDL suffix]{UDL suffixes} in the Standard
\emph{Library} provides context for better understanding and
appreciating the \emcppsgloss[user defined literal (UDL)]{UDL} language feature. These new suffixes
(starting with C++14) make it easier to write software using standard
strings, units of time, and complex numbers. Note that, because these
are \emph{standard} \emcppsgloss[UDL suffix]{UDL suffixes}, their names do not have a
leading underscore.

A native string literal, without a suffix, describes an C-style array of
characters, which decays to a pointer-to-character when passed as a
function argument. The C++ Standard Library has had, from the start,
string classes (\lstinline!std::basic_string!, \lstinline!std::string!, and
\lstinline!std::wstring!) that improve on C-style character arrays by
providing proper \emcppsgloss{copy semantics}, equality comparison, variable sizing,
and so on. With the advent of \emcppsgloss[user defined literal (UDL)]{UDLs}, we can finally create
literals of these library string types by utilizing the standard
\lstinline!s! suffix. The \emcppsgloss[UDL operator]{UDL operators} for string literals are in
header file \lstinline!<string>! in namespace\linebreak%%%%%
\lstinline!std::literals::string_literals!{\cprotect\footnote{C++20 adds \lstinline!char8_t!
and changes the type of \lstinline!s3! in the coding example from \lstinline!std::string! to
\lstinline!std::u8string!.}}:

\begin{emcppslisting}[emcppsstandards=c++14]
#include <string>  // (ù{\codeincomments{std::basic\_string}}ù), related types, and UDL operators

using namespace std::literals::string_literals;  // (ù{\codeincomments{std::basic\_string}}ù) UDLs
const char*    s1 =   "hello";    // Value decays to (ù{\codeincomments{(const char *) "hello"}}ù).
std::string    s2 =   "hello"s;   // value (ù{\codeincomments{std::string("hello")}}ù)
std::string    s3 = u8"hello"s;   // value (ù{\codeincomments{std::string(u8"hello")}}ù)
std::u16string s4 =  u"hello"s;   // value (ù{\codeincomments{std::u16string(u"hello")}}ù)
std::u32string s5 =  U"hello"s;   // value (ù{\codeincomments{std::u32string(U"hello")}}ù)
std::wstring   s6 =  L"hello"s;   // value (ù{\codeincomments{std::wstring(L"hello")}}ù)
\end{emcppslisting}
    
\noindent Complex numbers can also be expressed using a more natural style,
mimicking the notation used in mathematics. Within namespace
\lstinline!std::literals::complex_literals!, the suffixes \lstinline!i!,
\lstinline!il!, and \lstinline!if! are used to name \lstinline!double!,
\lstinline!long!~\lstinline!double!, and \lstinline!float! imaginary numbers,
respectively. Note that all three suffixes work for both integer and
floating-point literals:

\begin{emcppslisting}[emcppsstandards=c++14]
#include <complex>  // (ù{\codeincomments{std::complex}}ù) and UDL operators

using namespace std::literals::complex_literals; // (ù{\codeincomments{std::complex}ù) UDLs
std::complex<double>      c1 = 2.4 + 3i;    // value (ù{\codeincomments{{2.4, 3.0}}}ù)
std::complex<long double> c2 = 1.2 + 5.1l;  // value (ù{\codeincomments{{1.2L, 5.1L}}}ù)
std::complex<float>       c3 = 0.1f + 2.if;  // value (ù{\codeincomments{{0.1F, 2.0F}}}ù)
\end{emcppslisting}
    
\noindent The time utilities in the standard header, \lstinline!<chrono>!, contain an
elaborate and flexible system of units of duration. Each unit is a
specialization of the class template\linebreak%%%%%
 \lstinline!std::chrono::duration!,
which is instantiated with a representation (either integral or
floating-point) and a ratio relative to seconds. Thus,
\lstinline!duration<long,!~\lstinline!ratio<3600,!~\lstinline!1>>! can represent
an integral number of hours.

The \lstinline!<chrono>! header also defines literal suffixes (in namespace\linebreak%%%%% 
\lstinline!std::literals::chrono_literals!) having familiar names for time
units such as \lstinline!s! for seconds, \lstinline!min! for minutes, and so
on. Integer literals will yield a \lstinline!duration! having an integral
internal representation, and floating-point literals will yield one
having a floating-point internal representation:

\begin{emcppslisting}[emcppsbatch=e13,emcppsstandards=c++14]
#include <chrono>  // (ù{\codeincomments{std::literals::chrono\_literals}}ù)

using namespace std::literals::chrono_literals;  // (ù{\codeincomments{std::chrono::duration}}ù) UDLs
auto d1 = 2h;      // 2 hours   (integral internal representation)
auto d2 = 1.3h;    // 1.3 hours (floating-point internal representation)
auto d3 = 10min;   // 10 minutes (integral)
auto d4 = 30s;     // 30 seconds (integral)
auto d5 = 250ms;   // 250 milliseconds (integral)
auto d6 = 90us;    // 90 microseconds (integral)
auto d7 = 104.ns;  // 104.0 nanoseconds (floating-point)
\end{emcppslisting}
    
\noindent In the example above, the \lstinline!auto! keyword (see \featureref{\locationc}{auto}) is used to allow the compiler to deduce
the correct type from the literal expression. Although simple integer
duration types have convenient aliases such as
\lstinline!std::chrono::hours!, some types do not have a standard name for
the corresponding \lstinline!duration! specialization; e.g., \lstinline!1.2hr!
returns a value of type
\lstinline!std::chrono::duration<T,!~\lstinline!std::ratio<3600>>! where \lstinline!T! is a signed integer of at least 23 bits and 
where the actual integer representation is implementation-defined.
Naming a duration is even more complex when adding \lstinline!duration!s
together; the resulting \lstinline!duration! type is selected by the
library to minimize loss of precision:

\begin{emcppslisting}[emcppsbatch=e13]
auto d8 = 2h + 35min + 20s;  // integral, 9320 seconds (2:35:20 in seconds)
auto d9 = 2.4s + 100ms;       // floating-point, 2500.0 milliseconds
\end{emcppslisting}
    
\noindent We are certain to see more \emcppsgloss[UDL suffix]{UDL suffixes} defined in future
Standards.

\subsection[Use Cases]{Use Cases}\label{use-cases-userdeflit}

\subsubsection[Wrapper classes]{Wrapper classes}\label{wrapper-classes}

Wrappers can be used to add or remove capabilities for their (often
built-in) underlying type. They assign \emph{meaning} to a type and are
thus useful in preventing programmer confusion or ambiguity in overload
resolution. For example, an inventory-control system might track items
by both part number and model number. Both numbers could be simple
integers, but they have very different meanings. To prevent programming
errors, we create wrapper classes, \lstinline!PartNumber! and
\lstinline!ModelNumber!, each holding an \lstinline!int! value:

\begin{emcppslisting}[emcppsbatch=e14]
class PartNumber
{
    int d_value;

public:
    constexpr PartNumber(int v) : d_value(v) { }
    // ...
};

class ModelNumber
{
    int d_value;

public:
    constexpr ModelNumber(int v) : d_value(v) { }
    // ...
};
\end{emcppslisting}
    
\noindent Neither \lstinline!PartNumber! nor \lstinline!ModelNumber! defines integer
operations such as addition or multiplication, so any attempt to modify
one (other than by assignment) or add two such values would result in a
compile-time error. Moreover, having wrapper classes allows us to
overload on the different types, preventing overload resolution
ambiguities. Without \emcppsgloss[user defined literal (UDL)]{UDLs}, however, we must represent
\lstinline!PartNumber! or \lstinline!ModelNumber! literals by explicitly
casting \lstinline!int! literals to the correct type:

\begin{emcppslisting}[emcppsbatch=e14]
// operations on model and part numbers:
int inventory(ModelNumber n) { int count = 0; /*...*/ return count; }
int inventory(PartNumber n)  { int count = 0; /*...*/ return count; }
void registerPart(const char* shortName, ModelNumber mn, PartNumber pn) { }

PartNumber pn1 = PartNumber(77) + 90;  // Error, no (ù{\codeincomments{operator+(PartNumber, int)}}ù)

int c1 = inventory(77);               // Error, ambiguous overload
int c2 = inventory(ModelNumber(77));  // OK, call (ù{\codeincomments{inventory(ModelNumber)}}ù)
int c3 = inventory(PartNumber(77));   // OK, call (ù{\codeincomments{inventory(PartNumber)}}ù)

void registerParts1()
{
    registerPart("Bolt", PartNumber(77), ModelNumber(77));  // Error, reversed
    registerPart("Bolt", ModelNumber(77), PartNumber(77));  // OK, correct args
}
\end{emcppslisting}
    
\noindent The code above allows the compiler to detect errors that
would be easy to make had part and model numbers been represented as raw
\lstinline!int! values. Attempting to add to a part number is rejected,
as is the attempt to call \lstinline!inventory! without specifying whether
part-number inventory or model-number inventory is desired. Finally, the
\lstinline!registerPart! function cannot be called with its arguments
accidentally reversed.

We can now create \emcppsgloss[UDL suffix]{UDL suffixes}, \lstinline!_part! and
\lstinline!_model!, to simplify our use of hard-coded part and model
numbers, making the code more readable:

\begin{emcppslisting}[emcppsbatch=e14]
namespace inventory_literals
{
    constexpr ModelNumber operator"" _model(unsigned long long v) { return v; }
    constexpr PartNumber  operator"" _part (unsigned long long v) { return v; }
}

using namespace inventory_literals;  // Make literals available.
int c5 = inventory(77);              // Error, ambiguous overload
int c6 = inventory(77_model);        // OK, call (ù{\codeincomments{inventory(ModelNumber)}}ù).
int c7 = inventory(77_part);         // OK, call (ù{\codeincomments{inventory(PartNumber)}}ù).

void registerParts2()
{
    registerPart("Bolt", 77_part, 77_model); // Error, reversed model & part
    registerPart("Bolt", 77_model, 77_part); // OK, arguments in correct order
}
\end{emcppslisting}
    
\noindent A wrapper class can also be useful for tracking certain compile-time
attributes of a general-purpose type such as \lstinline!std::string!. For
example, a system that reads input from a user must sanitize each input
string before passing it to, for example, a database. Raw input and
sanitized input are both strings, but an unsanitized string must never
be confused for a sanitized one. Thus, we create a wrapper class,
\lstinline!SanitizedString!, that can be constructed only by a member
factory function, to which we give the easily-searchable name
\lstinline!fromRawString!:

\begin{emcppslisting}[emcppsbatch=e15]
#include <string>  // (ù{\codeincomments{std::string}}ù)

class SanitizedString
{
    std::string d_value;

    explicit SanitizedString(const std::string& value) : d_value(value) { }

public:
    static SanitizedString fromRawString(const std::string& rawStr)
    {
        return SanitizedString(rawStr);
    }

    // ...

    friend SanitizedString operator+(const SanitizedString& s1,
                                     const SanitizedString& s2)
    {
        return SanitizedString(s1.d_value + s2.d_value);
    }

    // ...
};
\end{emcppslisting}
    
\noindent Calling \lstinline!SanitizedString::fromRawString! is deliberately
cumbersome in an attempt to make developers think carefully before using
it. This call might, however, be \emph{too} cumbersome in situations where the safety
of a literal string is not in question:

\begin{emcppslisting}[emcppsbatch=e15]
std::string getInput();  // Read (unsanitized) string from input.
bool isSafeString(const std::string& s);  // Determine whether (ù{\codeincomments{s}}ù) is safe.

void process(const SanitizedString& instructions);
    // Run the specified (ù{\codeincomments{instructions}}ù).

void processInstructions1()
    // Read instructions from input and process them.
{
    // Read instructions from input.
    std::string instructions = getInput();

    if (isSafeString(instructions))
    {
       // String is considered safe; sanitize it.
       SanitizedString sanInstr = SanitizedString::fromRawString(instructions);
       // ...

       // Prepend a "begin" instruction, then process the instructions.

       process("Instructions = begin\n" + sanInstr);
           // Error, no (ù{\codeincomments{operator+(const char*, SanitizedString)}}ù)

       process(SanitizedString::fromRawString("Instructions = begin\n") +
               sanInstr);
           // OK, but cumbersome
    }
    else
    {
        // ...                  (error handling)
    }
}
\end{emcppslisting}
    
\noindent The first call to \lstinline!process! does not compile because, by design,
we cannot concatenate a raw string and a sanitized string. The second
call works but is unnecessarily cumbersome. Literal strings are
\emph{always} assumed to be safe (if there is proper code review)
because they cannot originate from outside the program; there should be
no need to call \lstinline!SanitizedString::fromRawString!. Again, a
\emcppsgloss[user defined literal (UDL)]{UDL} can make the code more compact and readable:

\begin{emcppslisting}[emcppsbatch=e15]
namespace sanitized_string_literals
{
    SanitizedString operator""_san(const char *str, std::size_t len)
        // Create a sanitized string.    
    {
        return SanitizedString::fromRawString(std::string(str, len));
    }
}

void processInstructions2()
    // Read instructions from input and process them.
{
    using namespace sanitized_string_literals;

    // Read instructions from input.
    std::string instructions = getInput();

    if (isSafeString(instructions))
    {
       // The (ù{\codeincomments{instructions}}ù) string is considered safe; sanitize it.
       SanitizedString sanInstr = SanitizedString::fromRawString(instructions);
       // ...

       // Prepend a "begin" instruction, then process the instructions.

       process("Instructions = begin\n"_san + sanInstr);
           // OK, concatenate two sanitized strings.
    }
    // ...
}
\end{emcppslisting}
    
\noindent This usage shows a case where the \emcppsgloss[user defined literal (UDL)]{UDL} is more than just
convenient; because a \emcppsgloss[user defined literal (UDL)]{UDL} applies \emph{only} to literals, it
is largely immune to accidental misuse.

\subsubsection[User-defined numeric types]{User-defined numeric types}\label{user-defined-numeric-types}

There is sometimes a need to represent an integer of indefinite
magnitude, i.e., where computations are immune from overflow (within the
bounds of available memory). A \lstinline!BigNum! class, along with
associated arithmetic operators, can represent such indefinite-magnitude
integers:

\begin{emcppslisting}[emcppsbatch=e16]
namespace bignum
{
class BigNum
{
    // ...
};

BigNum operator+(const BigNum&);
BigNum operator-(const BigNum&);
BigNum operator+(const BigNum&, const BigNum&);
BigNum operator-(const BigNum&, const BigNum&);
BigNum operator*(const BigNum&, const BigNum&);
BigNum operator/(const BigNum&, const BigNum&);
BigNum abs(const BigNum&);
// ...
\end{emcppslisting}
    
\noindent A \lstinline!BigNum! literal must be able to represent a value larger than
would fit in the largest built-in integral type, so we define the suffix
using a \emcppsgloss{raw UDL operator}:

\begin{emcppslisting}[emcppsbatch=e16]
namespace literals
{
BigNum operator"" _bignum(const char *digits)  // raw literal
{
    BigNum value;
    // ...        (Compute (ù{\codeincomments{BigNum}}ù) from (ù{\codeincomments{digits}}ù).)
    return value;
}
}  // close namespace (ù{\codeincomments{literals}}ù) 

using namespace literals;
}  // close namespace (ù{\codeincomments{bignum}}ù) 

using namespace bignum::literals;  // Make (ù{\codeincomments{\_bignum}}ù) literal available.
bignum::BigNum bnval  = 587135094024263344739630005208021231626182814_bignum;
bignum::BigNum bigone = 1_bignum;  // small value, but still has type (ù{\codeincomments{BigNum}}ù)
\end{emcppslisting}
    
\noindent The \lstinline!BigNum! class is appropriate for large integers, but numbers having 
fractional parts have a different problem: The IEEE standard
\lstinline!double! floating-point type cannot exactly represent certain
values, e.g., \lstinline!24692134.03!. When \lstinline!double!s are used to
represent values, long summations may eventually produce an error in the
hundredths place; i.e., the result will be off by \lstinline!.01! or more.
In financial calculations, where values represent money, errors of just
one or two pennies might be unacceptable. For this problem, we turn to
decimal fixed-point (rather than binary floating-point) arithmetic.

A decimal fixed-point representation for a number is one where the
number of decimal places of precision is chosen by the programmer and
fixed at compile time. Within the specified size and precision, every
decimal value can be represented exactly --- e.g., a fixed-point number
with two decimal digits of precision can represent the value
\lstinline!24692134.03! exactly but cannot represent the value
\lstinline!24692134.035!. We'll define our \lstinline!FixedPoint! class as a
template, where the \lstinline!Precision! parameter specifies the number of
decimal places{\cprotect\footnote{A more complete and powerful
fixed-point class template was proposed for standardization in \mbox{\cite{mcfarlane19}}.}}:

\begin{emcppslisting}[emcppsbatch=e17]
#include <limits>  // (ù{\codeincomments{std::numeric\_limits}}ù)
#include <string>  // (ù{\codeincomments{std::string}}ù), (ù{\codeincomments{std::to\_string}}ù)

namespace fixedpoint
{

template <unsigned Precision>
class FixedPoint
{
    long long d_data;  // integral data = value * pow(10, Precision)

public:
    constexpr FixedPoint() : d_data(0) { } // zero value
    constexpr FixedPoint(long long);       // Convert from (ù{\codeincomments{long long}}ù).
    constexpr FixedPoint(double);          // Convert from (ù{\codeincomments{double}}ù).

    constexpr FixedPoint(long long data, std::true_type /*isRaw*/)
        // Create a (ù{\codeincomments{FixedPoint}}ù) object with the specified (ù{\codeincomments{data}}ù).  Note that
        // this is a "raw" constructor and no precision adjustment is made to
        // the (ù{\codeincomments{data}}ù).
        : d_data(data) { }

    friend std::ostream& operator<<(std::ostream& stream, const FixedPoint& v)
        // Format the specified fixed-point number (ù{\codeincomments{v}}ù), write it to the
        // specified (ù{\codeincomments{stream}}ù), and return (ù{\codeincomments{stream}}ù).
    {
        std::string str = std::to_string(v.d_data);
        // Insert leading '0's, if needed.
        if (str.length() < Precision)
            str.insert(0, (Precision - str.length()), '0');
        str.insert(str.length() - Precision, 1, '.');
        return stream << str;
    }
};
\end{emcppslisting}
    
\noindent Our data representation is a \lstinline!long!~\lstinline!long! (see \featureref{\locationa}{long-long}), making the
largest value that can be represented
\lstinline!std::numeric_limits<long!~\lstinline!long>::max()!~\lstinline!/!~\lstinline!pow(10,!~\lstinline!Precision)!, assuming an integer \lstinline!pow! function. The output function,
\lstinline!operator<<!, converts \lstinline!d_data! to a string and
then inserts the decimal point into the correct location. The special
``raw'' constructor exists so that our \emcppsgloss{UDL operator} can easily
construct a value without losing precision, as we'll see later; the
unused second parameter is a dummy to distinguish it from other
constructors.

We want to define a \emcppsgloss{templated UDL operator} to return a
\lstinline!FixedPoint! type such that, for example, \lstinline!12.34! would
return a value of type \lstinline!FixedPoint<2>! (for the two decimal
places), whereas \lstinline!12.340! would return a value of type
\lstinline!FixedPoint<3>!. We must first define a variadic helper template (see \featureref{\locationc}{variadictemplate}) to
compute both the type and the raw value of the fixed point number, given a
sequence of digits:

\begin{emcppslisting}[emcppsbatch=e17]
namespace literals  // fixed-point literals defined in this namespace
{

template <long long rawVal, int precision, char... c>
struct MakeFixedPoint;
\end{emcppslisting}
    
\noindent This helper template will be recursively instantiated; at each level of
recursion \lstinline!rawVal! is the value computed so far,
\lstinline!precision! is the number of decimal places seen so far, and
\lstinline!c...! is the list of literal characters to be consumed. A
special value of \lstinline!-1! for \lstinline!precision! indicates that the decimal
point has not yet been consumed.

The base case of our recursive template occurs when the parameter pack,
\lstinline!c...!, is empty --- i.e., there are no more characters to
consume. In this case, the computed type is simply
\lstinline!FixedPoint<precision>!, and the value of the \emcppsgloss{UDL
operator} is computed from \lstinline!rawVal!. We define the base case as a
partial specialization of \lstinline!MakeFixedPoint! where there is no
character parameter pack:

\begin{emcppslisting}[emcppsbatch=e17]
template <long long rawVal, int precision>
struct MakeFixedPoint<rawVal, precision> 
{
    // base case when there are no more characters
    using type = FixedPoint<(precision < 0) ? 0 : precision>;
    static constexpr type makeValue() { return { rawVal, std::true_type{} }; }
        // Return the computed fixed-point number.
};
\end{emcppslisting}
    
\noindent The other case occurs when there are one or more characters yet to be
consumed. The helper must perform error checking for bad input
characters and overflow before consuming the character and instantiating
itself recursively:

\begin{emcppslisting}[emcppsbatch=e17]
template <long long rawVal, int precision, char c0, char... c>
struct MakeFixedPoint<rawVal, precision, c0, c...>
{
private:
    static constexpr long long maxData = std::numeric_limits<long long>::max();
    static constexpr bool      c0isdig = ('0' <= c0 && c0 <= '9');

    // Check for out-of-range characters and overflow.
    static_assert(c0isdig || '\'' == c0 || '.' == c0,
                  "Invalid fixed-point digit");
    static_assert(!c0isdig || (maxData - (c0 - '0')) / 10 >= rawVal,
                  "Fixed-point overflow");

    // (ù{\codeincomments{precision}}ù) is
    // (1) < 0 if a decimal point was not seen,
    // (2) 0   if a decimal point was seen but no digits after the decimal point,
    // (3) > 0 otherwise, incremented once for each digit after the decimal point.
    static constexpr int nextPrecision = ('.' == c0     ?  0 :
                                          precision < 0 ? -1 :
                                          precision + 1);

    // Instantiate this template recursively to consume remaining characters.
    using RecurseType = MakeFixedPoint<(c0isdig ? 10 * rawVal + c0 - '0' :
                                        rawVal), nextPrecision, c...>;

public:
    using                 type = typename RecurseType::type;
    static constexpr type makeValue() { return RecurseType::makeValue(); }
        // Return the computed fixed-point number.
};
\end{emcppslisting}
    
\noindent This specialization consumes one character, \lstinline!c0!, from the
parameter pack. The first\linebreak[4]%%%%%%  
\lstinline!static_assert! checks that
\lstinline!c0! is either a digit, digit separator
(\lstinline!'\''!), 
or decimal point (\lstinline!'.'!). The second
\lstinline!static_assert! checks that the computation is not in danger of
overflowing the maximum value of a \lstinline!long!~\lstinline!long!. The
constant, \lstinline!nextPrecision!, which will be passed to the recursive
instantiation of this template, keeps track of how many digits have been
consumed after the decimal point (or \lstinline!-1! if the decimal point has not yet
been consumed). The \lstinline!RecurseType! alias is the recursive
instantiation of this template with the updated raw value (after
consuming \lstinline!c0!), the updated precision, and the input character
sequence after having dropped \lstinline!c0!. Thus, each recursion gets a
potentially larger \lstinline!rawVal!, a potentially larger
\lstinline!precision!, and a shorter list of unconsumed characters. The
definitions of \lstinline!type! and \lstinline!makeValue! simply defer to the
definitions in the recursive instantiation.

Finally, we define the \lstinline!_fixed! \emcppsgloss{templated UDL operator},
instantiating \lstinline!MakeFixedPoint! with an initial \lstinline!rawVal! of
\lstinline!0!, an initial \lstinline!precision! of \lstinline!-1!, and with a \lstinline!c...!
parameter pack consisting of all of the characters in the \emcppsgloss{naked
literal}:

\begin{emcppslisting}[emcppsbatch=e17]
template <char... c>
constexpr typename MakeFixedPoint<0, -1, c...>::type operator"" _fixed()
{
    return MakeFixedPoint<0, -1, c...>::makeValue();
}

}  // Close namespace (ù{\codeincomments{literals}}ù). 
}  // Close namespace (ù{\codeincomments{fixedpoint}}ù). 
\end{emcppslisting}
    
\noindent Now the \lstinline!_fixed! suffix can be used for fixed-point
\emcppsgloss[user defined literal (UDL)]{UDLs} where the precision of the returned type is automatically
deduced based on the number of decimal places in the literal. Note that
the literal can be used in a \lstinline!constexpr! context:

\begin{emcppshiddenlisting}[emcppsbatch=e17]
#include <iostream>  // (ù{\codeincomments{std::cout}}ù)
template <unsigned Precision>
std::ostream& operator<<(std::ostream&,const fixedpoint::FixedPoint<Precision>&);
\end{emcppshiddenlisting}
\begin{emcppslisting}[emcppsbatch=e17]
int fixedTest()
{
    using namespace fixedpoint::literals;

    constexpr auto fx1 = 123.45_fixed;   // return type (ù{\codeincomments{FixedPoint<2>}}ù)
    constexpr auto fx2 = 123.450_fixed;  // return type (ù{\codeincomments{FixedPoint<3>}}ù)
    std::cout << fx1 << '\n';            // prints "123.45"
    std::cout << fx2 << '\n';            // prints "123.450"
}
\end{emcppslisting}
    
\noindent An effort is underway to define a standard \emcppsgloss[decimal floating point]{decimal
floating-point} type, which, like our decimal fixed-point type, retains
the benefit of precisely representing decimal fractions but where the
precision is variable at run time. If implemented as a library type, a
\emcppsgloss{UDL suffix} would allow such type to have a natural
representation in code.\footnote{\cite{kuhl12}}

\subsubsection[User-defined string types]{User-defined string types}\label{user-defined-string-types}

A universally unique identifier (UUID) is a 128-bit number that
identifies specific pieces of data in computer systems. It can be
readily expressed as an array of two 64-bit integers:

\begin{emcppshiddenlisting}[emcppsbatch=e18]
#include <cstddef>  // (ù{\codeincomments{std::size\_t}}ù)
\end{emcppshiddenlisting}
\begin{emcppslisting}[emcppsbatch=e18]
#include <cstdint>  // std::uint64_t
class UUIDv4
{
    std::unit64_t d_value[2];
    // ...
};
\end{emcppslisting}
    
\noindent A version-4 UUID has a canonical human-readable format consisting of
five groups of hex digits separated by hyphens, e.g.,
\lstinline!ed66b67a-f593-4def-9a9b-e69d1d6295ef!. Although storing this
representation as a string would be easy, converting it using the
packed, 128-bit integer format of the above \lstinline!UUIDv4! class is
more efficient and convenient. Moreover, UUIDs that are hard-coded into
software are often generated by external tools --- i.e., to identify the
exact product build --- and should therefore be compile-time constants.
Using \emcppsgloss[user defined literal (UDL)]{UDLs}, we can readily express a UUID literal using the
human-readable string format, converting it to a compile-time constant
in the packed format:

\begin{emcppslisting}[emcppsbatch=e18]
namespace uuid_literals
{
    constexpr UUIDv4 operator""_uuid(const char* s, std::size_t len)
    {
        return { /* ... (decode UUID expressed in canonical format) */ };
    }
}

using namespace uuid_literals;
constexpr UUIDv4 buildId = "eeec1114-8078-49c5-93ca-fea6fbd6a280"_uuid;
\end{emcppslisting}
    

\subsubsection[Unit conversions and dimensional units]{Unit conversions and dimensional units}\label{unit-conversions-and-dimensional-units}

\emcppsgloss[user defined literal (UDL)]{UDLs} can be convenient for specifying a unit name on a numeric
literal, providing a concise way both to convert the number to a
normalized unit and to annotate a value's unit within the code. For
example, the standard trigonometric functions all operate on
\lstinline!double! values where angles are expressed in radians. However,
many people are more comfortable with degrees than radians, especially
when expressing the value directly as a handwritten number:

\begin{emcppslisting}[emcppsbatch=e19]
#include <cmath>  // (ù{\codeincomments{std::sin}}ù), (ù{\codeincomments{std::cos}}ù)

constexpr double pi = 3.14159265358979311599796346854;

double s1 = std::sin(30.0);    // Bug, intended (ù{\codeincomments{sin(30 deg)}}ù) but got (ù{\codeincomments{sin(30 rad)}}ù)
double s2 = std::sin(pi / 6);  // OK, returns (ù{\codeincomments{sin(30 deg)}}ù)
\end{emcppslisting}
    
\noindent The normalized unit in this case is a radian, expressed as a
\lstinline!double!, but radians are generally fractions of $\pi$ and are thus
inconvenient to write. \emcppsgloss[user defined literal (UDL)]{UDLs} can provide convenient
normalization from degrees or gradians to radians:

\begin{emcppslisting}[emcppsbatch=e19]
namespace trig_literals {

constexpr double operator"" _rad(long double r)  { return r; }
constexpr double operator"" _deg(long double d)  { return pi * d / 180.0; }
constexpr double operator"" _grad(long double d) { return pi * d / 200.0; }

}

using namespace trig_literals;
double s3 = std::sin(30.0_deg);     // OK, returns (ù{\codeincomments{sin(30 deg)}}ù)
double s4 = std::sin(4.7124_rad);   // OK, returns approx (ù{\codeincomments{-1.0}}ù)
double s5 = std::cos(50.0_grad);    // OK, returns (ù{\codeincomments{sin(50 grad) == sin(45 deg)}}ù)
\end{emcppslisting}
    
\noindent Unfortunately, the applicability of the above approach to unit normalization
is very limited. First, the conversion is one way --- e.g., the
expression \lstinline!std::cout!~\lstinline!<<!~\lstinline!30.0_deg! will print
out \lstinline!0.524!, not \lstinline!30.0!, necessitating a call to a
radians-to-degrees conversion function when a human-readable value is
desired. Second, a \lstinline!double! does not encode any information about
the units that it holds, so \lstinline!double!~\lstinline!inputAngle! doesn't
tell the reader (or program) whether the angle is expected to be input
in degrees or radians.

A more robust way to use \emcppsgloss[user defined literal (UDL)]{UDLs} to express units is to
define them as part of a comprehensive library of unit classes.
Dimensional quantities (length, temperature, currency, and so on) often
benefit from being represented by \emcppsgloss[dimensional unit type]{dimensional unit types} that
prevent confusion as to both the units and dimension of numeric values.
For example, creating a function to compute kinetic energy from speed
and mass seems simple:

\begin{emcppslisting}[emcppsbatch=e20]
double kineticE(double speed, double mass)
    // Return kinetic energy in joules given (ù{\codeincomments{speed}}ù) in mps and (ù{\codeincomments{mass}}ù) in kg.
{
    return speed * speed * mass;
}
\end{emcppslisting}
    
\noindent Yet this simple function can be called incorrectly in numerous ways,
with no compiler diagnostics to help prevent the errors:

\begin{emcppslisting}[emcppsbatch=e20]
double d1   = 15;             // distance in meters
double t1   = 4;              // time in seconds
double s1   = d1 / t1;        // speed in m/s (meters/second)
double m1   = 2045;           // mass in g
double m1Kg = 2045.0 / 1000;  // mass in kg

double x1 = kineticE(d1, m1Kg);  // Bug, distance instead of speed
double x2 = kineticE(m1Kg, s1);  // Bug, arguments reversed
double x3 = kineticE(s1, m1);    // Bug, mass should be in kg, not g
double x4 = kineticE(s1, m1Kg);  // OK, meters per sec and kg units
\end{emcppslisting}
    
\noindent One way to detect some of these errors at compile time is to use a
wrapper for each dimension:

\begin{emcppslisting}[emcppsbatch=e20]
struct Time     { constexpr Time(double sec);        /* ... */ };
struct Distance { constexpr Distance(double meters); /* ... */ };
struct Speed    { constexpr Speed(double mps);       /* ... */ };
struct Mass     { constexpr Mass(double kg);         /* ... */ };
struct Energy   { constexpr Energy(double joules);   /* ... */ };

// Compute speed from distance and time:
Speed operator/(Distance, Time);

Distance d2(15.0);             // distance in meters
Time     t2(4.0);              // time in seconds
Speed    s2(d2 / t2);          // speed in m/s (meters/second)
Mass     m2(2045.0);           // Bug, trying to get g, got kg instead
Mass     m2Kg(2045.0 / 1000);  // OK, mass in kg

Energy kineticE(Speed s, Mass m);
Energy x5 = kineticE(d2, m2Kg); // Error, 1st argument has an incompatible type.
Energy x6 = kineticE(m2Kg, s2); // Error, reversed arguments, incompatible types
Energy x7 = kineticE(s2, m2);   // Bug, mass should be in kg, not g.
Energy x8 = kineticE(s2, m2Kg); // OK, m/s and Kg units
\end{emcppslisting}
    
\noindent Note that the compiler correctly diagnoses an error in the
initializations of \lstinline!x5! and \lstinline!x6! but still fails to
diagnose the unit error in the initialization of \lstinline!x7!.
User-defined literals can amplify the benefits of dimensional unit
classes by adding unit suffixes to numeric literals, eliminating
implicit unit assumptions:

\begin{emcppslisting}[emcppsbatch=e20]
namespace si_literals
{
    constexpr Distance operator"" _m  (long double meters);
    constexpr Distance operator"" _cm (long double centimeters);
    constexpr Time     operator"" _s  (long double seconds);
    constexpr Speed    operator"" _mps(long double mps);
    constexpr Mass     operator"" _g  (long double grams);
    constexpr Mass     operator"" _kg (long double kg);
    constexpr Energy   operator"" _j  (long double joules);
}

using namespace si_literals;
auto d3   = 15.0_m;    // distance in meters
auto t3   = 4.0_s;     // time in seconds
auto s3   = d3 / t3;   // speed in m/s (meters/second)
auto m3   = 2045.0_g;  // mass expressed as g but stored as Kg
auto m3Kg = 2.045_kg;  // mass expressed as kg

Energy x9  = kineticE(s3, m3);    // OK, (ù{\codeincomments{m3}}ù) has been normalized to Kg units.
Energy x10 = kineticE(s3, m3Kg);  // OK, m/s and Kg units
\end{emcppslisting}
    
\noindent Note that there are two \emcppsgloss[user defined literal (UDL)]{UDLs} that yield \lstinline!Distance! and
two \emcppsgloss[user defined literal (UDL)]{UDLs} that yield \lstinline!Mass!. Typically, the internal
representation of each of these dimensional types has a normalized
representation; e.g., \lstinline!Distance! might be represented in meters
internally, so \lstinline!25_cm! would be represented by a \lstinline!double!
data member with value \lstinline!0.25!. It is also possible, however, for
the unit to be stored alongside the value, thus avoiding rounding errors
in certain cases. Better yet, the unit can be encoded as a template
parameter at compile time:

\begin{emcppslisting}
#include <ratio>  // (ù{\codeincomments{std::ratio}}ù)
template <class Ratio> class MassUnit;
using Grams     = MassUnit<std::ratio<1, 1000>>;
using Kilograms = MassUnit<std::ratio<1>>;

namespace unit_literals
{
    constexpr Grams     operator"" _g  (long double grams);
    constexpr Kilograms operator"" _kg (long double kg);
}
\end{emcppslisting}
    
\noindent We now get \emph{different} types for \lstinline!100_g! and
\lstinline!0.1_kg!, but we can define \lstinline!MassUnit! in such a way that
they interoperate. The time-interval units that we saw previously in
{\emcppsgloss[user defined literal (UDL)]{UDLs} in the C++14 Standard Library} provide a taste of what is
possible with this approach.{\cprotect\footnote{Mateusz Pusz explores
the topic of a comprehensive physical units library in \cite{pusz20b}.}}

\subsubsection[Test drivers]{Test drivers}\label{test-drivers}

Because code is easier to maintain when ``magic'' values are expressed
as named constants rather than literal values, a typical program does
not contain many literals (see \intraref{potential-pitfalls-defmemberinit}{overuse-defmemberinit}). The
exception to this general rule is in unit tests where many different values are
successively passed to a subsystem to test its behavior. For example, we define a \lstinline!Date! class
that supplies a subtraction operator returning the number of days
between two \lstinline!Date!s:

\begin{emcppslisting}[emcppsbatch=e21]
// date.h: (component header file)

class Date
{
    // ...
public:
    constexpr Date(int year, int month, int day);
    // ...
};

int operator-(const Date& lhs, const Date& rhs);
    // Return the number of days from (ù{\codeincomments{rhs}}ù) to (ù{\codeincomments{lhs}}ù).
\end{emcppslisting}
    
\noindent To test the subtraction operator, we need to feed it combinations of
dates and compare the result with the expected result. We do this by
creating an array where each row holds a pair of dates and the expected
result from subtracting them. Due to the large number of hard-coded
values in the array, having a literal representation for the
\lstinline!Date! class would be convenient, even if the author of
\lstinline!Date! did not see fit to provide one:

\begin{emcppslisting}[emcppsbatch=e21]
// date.t.cpp: (component test driver)

#include <date.h>   // (ù{\codeincomments{Date}}ù)
#include <cstdlib>  // (ù{\codeincomments{std::size\_t}}ù)
#include <cassert>  // standard C (ù{\codeincomments{assert}}ù) macro

namespace test_literals
{
    constexpr Date operator"" _date(const char*, std::size_t);
        // UDL to convert date in "yyyy-mm-dd" format to a (ù{\codeincomments{Date}}ù) object
}

void testSubtraction()
{
    using namespace test_literals;  // Import (ù{\codeincomments{\_date}}ù) UDL suffix.

    struct TestRow
    {
        Date lhs;  // left operand
        Date rhs;  // right operand
        int  exp;  // expected result
    };

    const TestRow testData[] =
    {
        { "2021-01-01"_date, "2021-01-01"_date,  0 },
        { "2021-01-01"_date, "2020-12-31"_date,  1 },
        { "2021-01-01"_date, "2021-01-02"_date, -1 },
        // ...
    };

    const std::size_t testDataSize = sizeof(testData) / sizeof(TestRow);

    for (std::size_t i = 0; i < testDataSize; ++i)
    {
        assert(testData[i].lhs - testData[i].rhs == testData[i].exp);
    }
}
\end{emcppslisting}
    

\subsection[Potential Pitfalls]{Potential Pitfalls}\label{potential-pitfalls-defmemberinit}

\subsubsection[Unexpected characters can yield bad values]{Unexpected characters can yield bad values}\label{unexpected-characters-can-yield-bad-values}

\emcppsgloss[raw UDL operator]{Raw UDL operators} and \emcppsgloss[templated UDL operator]{templated UDL operators} must
parse and handle every character from the \emph{union} of the set of
legal characters in integer and floating-point literals, even if the
\emcppsgloss{UDL operator} is expecting only one of the two numeric type
categories. Failure to generate an error for an invalid character is
likely to produce an incorrect value, rather than a program crash or
compilation error:

\begin{emcppslisting}
short operator"" _short(const char *digits)
{
    short result = 0;
    for (; *digits; ++digits)
    {
        result = result * 10 + *digits - '0';
    }

    return result;
}

short s1 = 123_short;   // OK, value 123
short s2 = 123._short;  // Bug, (ù{\codeincomments{'.'}}ù) treated as digit value (ù{\codeincomments{-2}}ù)
\end{emcppslisting}
    
\noindent Testing only for the \emph{expected} characters and rejecting any others
is better than checking for \emph{invalid} characters and accepting the
rest:

\begin{emcppslisting}
#include <stdexcept>  // (ù{\codeincomments{std::out\_of\_range}}ù)

short operator"" _shrt2(const char *digits)
{
    short result = 0;
    for (; *digits; ++digits)
    {
        if (*digits == '.')
        {
            throw std::out_of_range("Bad digit");  // BAD IDEA
        }

        if (!std::isdigit(*digits))
        {
            throw std::out_of_range("Bad digit");  // BETTER
        }

        result = result * 10 + *digits - '0';
    }

    return result;
}

short s3 = 123_shrt2;    // OK, value 123
short s4 = 123._shrt2;   // Error (detected), throws (ù{\codeincomments{out\_of\_range("Bad digit")}}ù)
short s5 = 0x123_shrt2;  // Error (detected), throws (ù{\codeincomments{out\_of\_range("Bad digit")}}ù)
\end{emcppslisting}
    
\noindent The first \lstinline!if! will catch an unexpected decimal point but not an
unexpected characters such as \lstinline!'e'!, \lstinline!'x'!, \lstinline!'\'!. The second \lstinline!if! will catch all unexpected characters. If a
new radix or other currently illegal character is introduced in a future
Standard, the second \lstinline!if! will avoid processing it incorrectly.
Note that, for example, after \emcppsgloss[user defined literal (UDL)]{UDLs} were added in C++11, both
the \lstinline!0b! radix and the digit separator (\lstinline!'!) were
introduced in C++14, potentially breaking any C++11-compliant
\emcppsgloss{UDL operator} that didn't properly handle those characters.


%%%%%%%% PRODUCTION: May want to soften/neutralize the language on FPPs. 
%\emph{TODO (VR): this annoyance seems quite subjective.}
\subsubsection[Overuse]{Overuse}\label{overuse-defmemberinit}

Although UDLs offer conciseness, they aren't always necessarily the most effective way to create
a literal value in a program. Sometimes a regular constructor or
function call might be almost as concise, simpler, and more flexible. For
example, \lstinline!deg(90)! is as readable as \lstinline!90_deg! and can be
applied to runtime values as well as to literals. If the constructor for
a type naturally takes two or more arguments, a string \emcppsgloss{UDL
operator} could theoretically parse a comma-separated list of arguments
--- e.g., \lstinline!"(2.0,!~\lstinline!6.0)"_point! to represent a 2-D
coordinate --- but is the literal really easier to read than
\lstinline!Point(2.0,!~\lstinline!6.0)!?

Finally, even for simple cases, consider how \emph{often} a literal is
likely to be used. The use of ``magic numbers'' in code is widely
discouraged. Numeric literals other than, say, \lstinline!-1!, \lstinline!0!,
\lstinline!1!, \lstinline!2!, and \lstinline!10! or string literals other than
\lstinline!""! are typically used only to initialize named constants. The
speed of sound, for example,  would probably be written as a constant, e.g., 
\lstinline!speedOfSound!, rather than a literal, \lstinline!343_mps!. Using literals to supply the special values used to initialize named
constants provides little benefit by way of overall program readability:

\begin{emcppshiddenlisting}[emcppsbatch=e22]
struct Speed    { constexpr Speed(long double mps) {} };
constexpr Speed operator"" _mps(unsigned long long speed) { return Speed(speed); }
constexpr Speed operator*(long d, const Speed&rhs) { return Speed(d); }
\end{emcppshiddenlisting}
\begin{emcppslisting}[emcppsbatch=e22]
constexpr Speed operator"" _mps(unsigned long long speed);  // meters per second

constexpr Speed speedOfSound1 = 343_mps;   // OK, very clear
constexpr Speed speedOfSound2(343);        // OK, almost as clear

constexpr Speed mach2 = 2 * speedOfSound1; // Literal is irrelevant.
constexpr Speed mach3 = 3 * speedOfSound2; // Literal is irrelevant.
constexpr Speed mach4 = 4 * 343_mps;       // Bad style: "magic" number
\end{emcppslisting}
    
\noindent Often, the most common literal is the one that expresses the notion of
an \emph{empty} or \emph{zero} value. A clearer, more descriptive alternative might be to create named constants, such as  \lstinline!constexpr!~\lstinline!Thing!~\lstinline!k_EMPTY_THING!, instead of
defining a \emcppsgloss{UDL operator} just to be able to write
\lstinline!""_thing! or \lstinline!0_thing!.


%%%%%%%%%% REMOVE for this edition but keep for perhaps using for later editions. 
%\subsubsection[Obfuscation]{Obfuscation}\label{obfuscation}
%
%\emcppsgloss[user defined literal (UDL)]{UDLs} are mostly \emcppsgloss{syntactic sugar}, although they can
%sometimes help with correctness (e.g., for expressions that should never
%be invoked on user input). As in the case of operator overloading,
%overuse of \emcppsgloss[user defined literal (UDL)]{UDLs} to create ``cute'' syntax can result in
%hard-to-understand code. For example, although
%\lstinline!192.168.0.1_ipv4! is not a valid token (as we mentioned in \ref{}
%{Restrictions on \emcppsgloss[user defined literal (UDL)]{UDLs}), the same token \emph{would be} valid
%(in C++14 and later) if the dots are replaced by digit separators,
%yielding \lstinline!192'168'0'1_ipv4!. Whether such a potentially
%appealing notation is of benefit to users is a matter of debate. If a
%\emcppsgloss[user defined literal (UDL)]{UDL} returns a new interpretation of an existing type, the
%meaning might be muddled rather than clarified.
%
%\emph{TODO (VR): this annoyance seems quite subjective.}
%%%%%%%%%%%%%%%%%%%%%%%%%%%%%%%%%

\subsubsection[Preprocessor surprises]{Preprocessor surprises}\label{preprocessor-surprises}

A string literal with a suffix, e.g., \lstinline!"hello"_wrld!, is a
single token in C++11 but was two tokens, \lstinline!"hello"! and
\lstinline!_wrld!, in previous versions of the language. This change could
manifest in a subtle difference in meaning, usually resulting in a
compilation error, if \lstinline!_wrld! is a macro:

\begin{emcppslisting}[emcppserrorlines=2]
#define _wrld " world"
const char* s = "hello"_wrld;  // "hello world" in C++03, UDL in C++11
\end{emcppslisting}

\subsubsection[Verbose usage]{Verbose usage}

One of the main selling points of user-defined literals is enabling developers to write concise and expressive code. However, since literal operators are commonly bundled together in a separate namespace, the required \emcppsgloss[using directive]{\lstinline!using! directive} on the caller side can become a burden if only one or few uses of a literal operator appear in a scope:

\begin{emcppshiddenlisting}[emcppsbatch={e23,e24}]
#include <cstddef>  // (ù{\codeincomments{std::size\_t}}ù)
struct Date {
constexpr Date(const char *) {};
};
namespace test_literals {
constexpr Date operator"" _date(const char *c, std::size_t s) { return Date(c); }
}
bool operator >(const Date&,const Date&);
\end{emcppshiddenlisting}
\begin{emcppslisting}[emcppsbatch=e23]
bool isAfterLunarLanding(Date d)
{
    using namespace test_literals;  // Import (ù{\codeincomments{\_date}}ù) UDL suffix.
    return d > "1969-07-20"_date;
}
\end{emcppslisting}

\noindent The authors of the \lstinline!_date! UDL should provide its same functionality as a function or constructor to ensure that users do not always require a possibly over-verbose \emcppsgloss[using directive]{\lstinline!using! directive} to achieve their goals:

\begin{emcppslisting}[emcppsbatch=e24]
bool isAfterLunarLanding(Date d)
{
    return d > Date("1969-07-20");
}
\end{emcppslisting}

    

\subsection[Annoyances]{Annoyances}\label{annoyances}

\subsubsection[No conversion from floating-point to integer \lstinline!UDL!]{No conversion from floating-point to integer {\SubsubsecCode UDL}}\label{no-conversion-from-floating-point-to-integer-udl}

Defining a prepared-argument floating-point \emcppsgloss{UDL operator} does
not make the corresponding suffix available to numeric literals that
look like integers and vice versa:

\begin{emcppslisting}
double operator"" _mpg(long double v);

double v1 = 12_mpg;   // Error, no integer UDL operator for (ù{\codeincomments{\_mpg}}ù)
double v2 = 12._mpg;  // OK, floating-point UDL operator for (ù{\codeincomments{\_mpg}}ù) found
\end{emcppslisting}
    
\noindent If the intention is to define a prepared-argument UDL where numbers with and without a decimal point are accepted, then both forms of the \emcppsgloss{UDL operator} must be defined.

\subsubsection[Potential suffix-name collisions]{Potential suffix-name collisions}\label{potential-suffix-name-collisions}

Using a \emcppsgloss{UDL suffix} requires bringing the corresponding
\emcppsgloss{UDL operator} into the current scope, e.g., by means of a
\lstinline!using! directive. If the scope is large enough and if more than
one imported namespace contains a \emcppsgloss{UDL operator} with the same
name, a name collision can result:

\begin{emcppshiddenlisting}[emcppsbatch=e25]
namespace trig_literals {
bool operator"" _deg(long double d);
}
namespace temperature_literals {
bool operator"" _deg(long double d);
}
\end{emcppshiddenlisting}
\begin{emcppslisting}[emcppsbatch=e25]
using namespace trig_literals;        // (ù{\codeincomments{\_deg}}ù), (ù{\codeincomments{\_rad}}ù), and (ù{\codeincomments{\_grad}}ù) suffixes
using namespace temperature_literals; // colliding (ù{\codeincomments{\_deg}}ù) suffix

auto d = 12.0_deg;  // Error, ambiguous use of suffix, (ù{\codeincomments{\_deg}}ù)
\end{emcppslisting}
    
%\noindent Furthermore, there is no straightforward way to disambiguate this case
%because qualified name lookup does not work for \emcppsgloss[UDL suffix]{UDL suffixes}
%(e.g., \lstinline!12.0trig_literals::_deg! does not compile).

\noindent While it is possible to disambiguate the colliding suffixes via qualified name lookup, the ensuing verbosity might defeat the purpose of using a UDL in the first place:

\begin{emcppslisting}[emcppsbatch=e25]
auto a = trig_literals::operator""_deg(12.0);
auto b = temperature_literals::operator""_deg(12.0);
\end{emcppslisting}
   

\subsubsection[Confusing raw with string UDL operators]{Confusing raw with string UDL operators}\label{easy-to-confuse-raw-with-string-udl-operators}

A \emcppsgloss{UDL operator} that takes a single
\lstinline!const!~\lstinline!char*! argument is a \emcppsgloss{raw UDL operator}
for numeric literals but can be easily confused for a
\emcppsgloss{prepared-argument UDL operator} for string literals:

\begin{emcppslisting}
int operator"" _udl(const char *);

int s = "hello"_udl;  // Error, no match for (ù{\codeincomments{operator""(const char*, size\_t)}}ù)
\end{emcppslisting}
    
\noindent Fortunately, such a problem will typically result in a compile-time
error.

\subsubsection[No templated UDL operators for string literals]{No templated UDL operators for string literals}\label{no-templated-udl-operators-for-string-literals}

\emcppsgloss[templated UDL operator]{Templated UDL operators} are called only for numeric literals;
string literals are limited to the prepared-argument pattern. It is thus
not possible to choose, at compile time, different return types based on
the contents of a string literal.\footnote{As of C++20, new syntax is added that removes the limitation of not being able to affect the return-type base on the contents of a literal string.}

\subsubsection[No way to parse a leading \lstinline!-! or \lstinline!+!]{No way to parse a leading {\SubsubsecCode -} or {\SubsubsecCode +}}\label{no-way-to-parse-a-leading---or-+}

As described in \intraref{description-userdeflit}{restrictions-on-user-defined-literals-(udls)}, a
\lstinline!-!~or~\lstinline!+! before a numeric literal is a separate negation
operator and not part of the literal. There are occasions, however,
where it would be convenient to know whether the literal value is being
negated. For example, if temperatures are being stored as
\lstinline!double! values in Kelvin and if the \emcppsgloss{UDL suffix}
\lstinline!_C! converts a floating-point literal from Celsius to Kelvin by
calling a function, \lstinline!cToK(double)!, then the expression
\lstinline!-10.0_C! produces the nonsensical value \lstinline!-283.15!
(\lstinline!-cToK(10.0)!) rather than the intuitive value of \lstinline!+263.15!
(\lstinline!cToK(-10.0)!). Alas, parsing the
\lstinline!-!~sign as part of the literal is simply not possible.

\subsubsection[Parsing numbers is hard]{Parsing numbers is hard}\label{parsing-numbers-is-hard}

Many of the benefits of \emcppsgloss[raw UDL operator]{raw UDL operators} and \emcppsgloss[templated
UDL operator]{templated
UDL operators} require parsing integer and/or floating-point values
manually, in code, often using recursion. Getting this right is tedious
at best. The Standard Library does not provide much support, especially
for \lstinline!constexpr! parsing.



\subsection[See Also]{See Also}\label{see-also}

\begin{itemize}
\item{\seealsoref{decltype}{\seealsolocationa}introduces a keyword often helpful for deducing the return type of a \emcppsgloss{templated UDL operator}.}
\item{\seealsoref{null-pointer-literal-(nullptr)}{\seealsolocationa}describes a keyword that unambiguously denotes the null pointer literal.}
\item{\seealsoref{auto}{\seealsolocationc} shows how type inference can be used to declare a variable to hold the value of a \emcppsgloss[user defined literal (UDL)]{UDL} when the type of the \emcppsgloss[user defined literal (UDL)]{UDL} varies based on its contents.}
\item{\seealsoref{constexprfunc}{\seealsolocationc}describes how most \emcppsgloss[user defined literal (UDL)]{UDLs} can be used as part of a constant expression.}
\item{\seealsoref{variadictemplate}{\seealsolocationc}shows how templates can take an infinite number of parameters, which is required for implementing \emcppsgloss[templated UDL operator]{templated UDL operators}.}
\item{\seealsoref{inline-namespaces}{\seealsolocatione}describes a feature not recommended for \emcppsgloss[UDL operator]{UDL operators}, yet the C++14 Standard Library puts \emcppsgloss[UDL operator]{UDL operators} into \lstinline!inline! namespaces.}
\end{itemize}

\subsection[Further Reading]{Further Reading}\label{further-reading}

\begin{itemize}
\item{\cite{mcintosh08}}
\item{For a discussion of user-defined literals with recommendations for idiomatic use, see \cite{dewhurst19}.}
\end{itemize}




\newpage
%\section[Variadic Templates]{Variable-Argument-Count Templates}\label{variadictemplate}




\emcppsFeature{
    short={Variadic Templates},
    long={Variable-Argument-Count Templates},
}{variadictemplate}
\setcounter{table}{0}
\setcounter{footnote}{0}
\setcounter{lstlisting}{0}
%\section[Variadic Templates]{Variable-Argument-Count Templates}\label{variadictemplate}

placeholder text.........
 
 


%%%%%%%%%%%%%%% C++14
\stepcounter{cppxx}
\addtocontents{toc}{\protect\renewcommand*\protect\cftsecindent{1.25em}}
\cftaddnumtitleline{toc}{section}{\thechapter.\thecppxx}{C++14}{\thepage}
% to add to the TOC at the section level
\renewcommand{\cppxx}{C++14}


\newpage
\addtocontents{toc}{\protect\renewcommand*\protect\cftsecindent{4.5em}}
%\section[{\tt constexpr} Functions '14]{Relaxed Restrictions on {\SecCode constexpr} Functions}\label{relaxed-constexpr-restrictions}%\sectionmark{Relaxed {\RHCode constexpr} Restrictions}
% 24 March 2021, no longer holding for revisions. Sending to Josh for code review. 
% 24 March 2021 JMB - compile check done
% 25 March 2021, commands updated
% 25 March 2021, sent to copyediting


\emcppsFeature{
    short={\lstinline!constexpr! Functions '14},
    tocshort={{\TOCCode constexpr} Functions '14},
    long={Relaxed Restrictions on {\SecCode constexpr} Functions},
    toclong={Relaxed Restrictions on \lstinline!constexpr! Functions},
    rhshort={{\RHCode constexpr} Functions '14},
}{relaxed-constexpr-restrictions}
\setcounter{table}{0}
\setcounter{footnote}{0}
\setcounter{lstlisting}{0}
%\section[{\tt constexpr} Functions '14]{Relaxed Restrictions on {\SecCode constexpr} Functions}
%\section[{\tt constexpr} Functions '14]{Relaxed Restrictions on {\SecCode constexpr} Functions}\label{relaxed-constexpr-restrictions}%\sectionmark{Relaxed {\RHCode constexpr} Restrictions}


C++14 lifts restrictions regarding use of many language features in the
body of a \lstinline!constexpr! function (see \featureref{\locationc}{constexprfunc}).

\subsection[Description]{Description}\label{description}

The cautious introduction (in C++11) of \lstinline!constexpr!
functions --- i.e., functions eligible for compile-time evaluation
--- was accompanied by a set of strict rules that, despite making life
easier for compiler implementers, severely narrowed the breadth of valid
use cases for the feature. In C++11, \lstinline!constexpr!
function bodies were restricted to essentially a single
\lstinline!return! statement and were not permitted to have any modifiable
local state (variables) or \emcppsgloss{imperative} language constructs
(e.g., assignment), thereby greatly reducing their \mbox{usefulness}:

\begin{emcppslisting}
constexpr int fact11(int x)
{
    static_assert(x >= 0, "");
        // Error, (ù{\codeincomments{x}}ù) is not a constant expression.

    static_assert(sizeof(x) >= 4, "");  // OK in C++11/14

    return x < 2 ? 1 : x * fact11(x - 1);  // OK in C++11/14
}
\end{emcppslisting}
    
\noindent Notice that recursive calls were supported, often leading to convoluted
implementations of algorithms (compared to an \emcppsgloss{imperative}
counterpart); see \intraref{use-cases-relaxedconstexpr}{non-recursive-constexpr-algorithms}. 

The C++11 \lstinline!static_assert! feature (see \featureref{\locationa}{compile-time-assertions-(static_assert)}) was always
permitted in a C++11 \lstinline!constexpr! function body.
However, because the input variable \lstinline!x! in \lstinline!fact11!
(in the code snippet above) is inherently not a compile-time constant expression, it can
never appear as part of a \lstinline!static_assert! predicate. Note that a
\lstinline!constexpr! function returning \lstinline!void! was also not permitted:

\begin{emcppslisting}[emcppsstandards={c++14}]
constexpr void no_op() { }  // Error in C++11; OK in C++14
\end{emcppslisting}
    
\noindent Experience gained from the release and subsequent real-world use of
C++11 emboldened the standard committee to lift most of these (now
seemingly arbitrary) restrictions for C++14, allowing use of (nearly)
\emph{all} language constructs in the body of a \lstinline!constexpr!
function. In C++14, familiar non-expression-based control-flow
constructs, such as \lstinline!if! statements and \lstinline!while! loops, are
also available, as are modifiable local variables and assignment
operations:

\begin{emcppshiddenlisting}[emcppsbatch=e1]
// --- Replace
    int temp = x - 1;  // Error in C++11; OK in C++14
    int temp = x - 1;
// --- End
\end{emcppshiddenlisting}
\begin{emcppslisting}[emcppsstandards={c++14},emcppsbatch=e1]
constexpr int fact14(int x)
{
    if (x <= 2)        // Error in C++11; OK in C++14
    {
        return 1;
    }

    int temp = x - 1;  // Error in C++11; OK in C++14
    return x * fact14(temp);
}
\end{emcppslisting}
    
\noindent Some useful features remain disallowed in C++14; most notably, any form
of dynamic allocation is not permitted, thereby preventing the use of
common standard container types, such as \lstinline!std::string! and
\lstinline!std::vector!{\cprotect\footnote{In C++20, even more
restrictions were lifted, allowing, for example, some limited forms of
  dynamic allocation, \lstinline!try! blocks, and uninitialized variables.}}:
\begin{enumerate}
\item{\lstinline!asm! declarations}
\item{\lstinline!goto! statements}
\item{Statements with labels other than \lstinline!case! and \lstinline!default!}
\item{\lstinline!try! blocks}
\item{Definitions of variables
\begin{enumerate}
\item{of other than a \emcppsgloss{literal type} (i.e., fully processable at compile time)}
\item{decorated with either \lstinline!static! or \lstinline!thread_local!}
\item{left uninitialized}
\end{enumerate}
}
\end{enumerate}
The restrictions on what can appear in the body of a \lstinline!constexpr!
that remain in C++14 are reiterated here in codified
form\footnote{Note that the degree to which these remaining forbidden features are reported varies substantially from one popular compiler to the next.}:

\begin{emcppslisting}[emcppsstandards={c++14}]
template <typename T>
constexpr void f()
{
try {                  // Error, (ù{\codeincomments{try}}ù) outside body isn't allowed (until C++20).
    std::ifstream is;  // Error, objects of *non-literal* types aren't allowed.
    int x;             // Error, uninitialized vars. disallowed (until C++20)
    static int y = 0;  // Error, (ù{\codeincomments{static}}ù) variables are disallowed.
    thread_local T t;  // Error, (ù{\codeincomments{thread\_local}}ù) variables are disallowed.
    try{}catch(...){}  // Error, (ù{\codeincomments{try}}ù)/(ù{\codeincomments{catch}}ù) disallowed (until C++20)
    if (x) goto here;  // Error, (ù{\codeincomments{goto}}ù) statements are disallowed.
    []{};              // Error, lambda expressions are disallowed (until C++17).
here: ;                // Error, labels (except (ù{\codeincomments{case}}ù)/(ù{\codeincomments{default}}ù)) aren't allowed.
    asm("mov %r0");    // Error, (ù{\codeincomments{asm}}ù) directives are disallowed.
} catch(...) { }       // Error, (ù{\codeincomments{try}}ù) outside body disallowed (until C++20)
}
\end{emcppslisting}
    

\subsection[Use Cases]{Use Cases}\label{use-cases-relaxedconstexpr}

\subsubsection[Nonrecursive \lstinline!constexpr! algorithms]{Nonrecursive {\SubsubsecCode constexpr} algorithms}\label{non-recursive-constexpr-algorithms}

The C++11 restrictions on the use of \lstinline!constexpr! functions often
forced programmers to implement algorithms (that would otherwise be
implemented iteratively) in a recursive manner. Consider, as a familiar
example, a naive{\cprotect\footnote{For a more efficient (yet less
intuitive) C++11 algorithm, see \intraref{appendix-relaxedconstexpr}{appendix:-optimized-c++11-example-algorithms}, \intrarefsimple{recursive-fibonacci}.\textbf{AUs: The Appendix may need to be slightly restructured. We have no current method for an xref to this depth of subsectioning. }}}
C++11-compliant \lstinline!constexpr! implementation of a function,
\lstinline!fib11!, returning the \emph{n}th Fibonacci number\footnote{We used \lstinline!long!~\lstinline!long! (instead of \lstinline!long!)
here to ensure a unique C++ type having at least 8 bytes on all
conforming platforms for simplicity of exposition (avoiding an internal
copy). We deliberately chose \emph{not} to make the value returned
\lstinline!unsigned! because the extra bit does not justify changing the
\emcppsgloss{algebra} (from \lstinline!signed! to \lstinline!unsigned!). For more
discussion on these specific topics, see \featureref{\locationa}{long-long}.}:
  
\begin{emcppslisting}
constexpr long long fib11(long long x)
{
    return
        x == 0 ? 0
               : (x == 1 || x == 2) ? 1
                                    : fib11(x - 1) + fib11(x - 2);
}
\end{emcppslisting}
    
\noindent The implementation of the \lstinline!fib11! function (above) has various
undesirable properties.
\begin{enumerate}
\item{\emph{Reading difficulty} — Because it must be implemented using a single \lstinline!return! statement, branching requires a chain of \emph{ternary operators}, leading to a single long expression that might impede human comprehension.}\pagebreak%%%%%%%%
\item{\emph{Inefficiency and lack of scaling} — The explosion of recursive calls is taxing on compilers: (1) the time to compile is markedly slower for the \emph{recursive} (C++11) algorithm than it would be for its \emph{iterative} (C++14) counterpart, even for modest inputs,{\cprotect\footnote{As an example, Clang 10.0.0, running on an x86-64 machine, required more than 80 times longer to evaluate \lstinline!fib(27)! implemented using the \emph{recursive} (C++11) algorithm than to evaluate the same functionality implemented using the \emph{iterative} (C++14) algorithm.}} and (2) the compiler might simply refuse to complete the compile-time calculation if it exceeds some internal (platform-dependent) \emph{threshold} number of operations.{\cprotect\footnote{The same Clang 10.0.0 compiler discussed in the previous footnote failed to compile \lstinline!fib11(28)!:

\begin{lstlisting}[language=bash,style=footcodeplain]
error: static_assert expression is not an integral constant expression
    static_assert(fib11(28) == 317811, "");
                  ^~~~~~~~~~~~~~~~~~~

note: constexpr evaluation hit maximum step limit; possible infinite loop?
\end{lstlisting}

\noindent GCC 10.x fails at \lstinline!fib(36)!, with a similar diagnostic:

\begin{lstlisting}[language=bash,style=footcodeplain]
error: 'constexpr' evaluation operation count exceeds limit of 33554432
       (use '-fconstexpr-ops-limit=' to increase the limit)
\end{lstlisting}

\noindent Clang 10.x fails to compile any attempt at constant evaluating \lstinline!fib(28)!, with the following diagnostic \mbox{message}:

\begin{lstlisting}[language=bash,style=footcodeplain]
note: constexpr evaluation hit maximum step limit; possible infinite loop?
\end{lstlisting} 
}} % end footnote, end cprotect
} % end item
\item{\emph{Redundancy} — Even if the recursive implementation were suitable for small input values during compile-time evaluation, it would be unlikely to be suitable for any run-time evaluation, thereby requiring programmers to provide and maintain \emph{two} separate versions of the same algorithm: a compile-time \emph{recursive} one and a runtime \emph{iterative} one.}
\end{enumerate}

In contrast, an \emph{imperative} implementation of a \lstinline!constexpr!
function implementing a function returning the $n$th Fibonacci
number in C++14, \lstinline!fib14!, does not suffer from any of the three
issues discussed above:

\begin{emcppslisting}[emcppsstandards={c++14}]
constexpr long long fib14(long long x)
{
    if (x == 0) { return 0; }

    long long a = 0;
    long long b = 1;

    for (long long i = 2; i <= x; ++i)
    {
        long long temp = a + b;
        a = b;
        b = temp;
    }

    return b;
}
\end{emcppslisting}
    
\noindent As one would expect, the compile time required to evaluate the iterative
implementation (above) is manageable{\cprotect\footnote{Both GCC 10.x
and Clang 10.x evaluated \lstinline!fib14(46)! 1836311903 correctly in
under 20ms on a machine running Windows 10 x64 and equipped with a
  Intel Core i7-9700k CPU.}}; of course, far more
computationally efficient (e.g., closed form{\cprotect\footnote{E.g.,
see
  http://mathonline.wikidot.com/a-closed-form-of-the-fibonacci-sequence.}})
solutions to this classic exercise are available.

\subsubsection[Optimized metaprogramming algorithms]{Optimized metaprogramming algorithms}\label{optimized-metaprogramming-algorithms}

C++14's relaxed \lstinline!constexpr! restrictions enable the use of
modifiable local variables and \emcppsgloss{imperative} language constructs
for metaprogramming tasks that were historically often implemented by
using (Byzantine) recursive template instantiation (notorious for their
voracious consumption of compilation time).

Consider, as the simplest of examples, the task of counting the number
of occurrences of a given type inside a \emcppsgloss{type list} represented
here as an empty variadic template (see \featureref{\locationc}{variadictemplate}) that can be
instantiated using a variable-length sequence of arbitrary C++ types{\cprotect\footnote{Variadic templates are a C++11
  feature having many valuable and practical uses. In this case, the
  variadic feature enables us to easily describe a template that takes
  an arbitrary number of C++ type arguments by specifying an ellipsis
  (\lstinline!...!) immediately following \lstinline!typename!. Emulating such
  functionality in C++98/03 would have required significantly more
  effort: A typical workaround for this use case would have been to
  create a template having some fixed maximum number of arguments (e.g.,
  20), each defaulted to some unused (incomplete) type (e.g.,
  \lstinline!Nil!):

  \begin{emcppslisting}[style=footcode]
  struct Nil;  // arbitrary unused (incomplete) type

  template <typename = Nil, typename = Nil, typename = Nil, typename = Nil>
  struct TypeList { };
      // emulates the variadic (ù{\fncodeincomments{TypeList}}ù) template (ù{\fncodeincomments{struct}}ù) for up to four
      // type arguments
  \end{emcppslisting}

\vspace*{-1ex}      
\noindent Another theoretically appealing approach is to implement a Lisp-like
  recursive data structure; the compile-time overhead for such
  implementations, however, often makes them impractical.}}:

\begin{emcppslisting}[emcppsbatch=e2]
template <typename...> struct TypeList { };
    // empty variadic template instantiable with arbitrary C++ type sequence
\end{emcppslisting}
    
\noindent Explicit instantiations of this variadic template could be used to
create objects:

\begin{emcppslisting}[emcppsbatch=e2]
TypeList<>                  emptyList;
TypeList<int>               listOfOneInt;
TypeList<int, long, double> listOfThreeIntLongDouble;
\end{emcppslisting}
    
\noindent A naive C++11-compliant implementation of a \emcppsgloss{metafunction}
\lstinline!Count!, used to ascertain the (order-agnostic) number of times a
given C++ type was used when creating an instance of the
\lstinline!TypeList! template (above), would usually make recursive use of
(baroque) \emcppsgloss{partial class template
specialization}{\cprotect\footnote{The use of class-template
specialization (let alone partial specialization) might be unfamiliar
to those not accustomed to writing low-level template metaprograms, but the point of this use case is to obviate such
unfamiliar use. As a brief refresher, a general class template is what
the client typically sees at the user interface. A specialization is
typically an implementation detail consistent with the
\emcppsgloss{contract} specified in the general template but somehow more
restrictive. A partial specialization (possible for \emph{class} but
not \emph{function} templates) is itself a template but with one or
more of the general template parameters resolved. An \emcppsgloss{explicit}
or \emcppsgloss{full specialization} of a template is one in which
\emph{all} of the template parameters have been resolved and, hence, is
not itself a template. Note that a \emcppsgloss{full specialization} is a
stronger candidate for a match than a partial specialization, which is
a stronger match candidate than a simple template specialization,
which, in turn, is a better match than the general template (which, in
  this example, happens to be an \emcppsgloss{incomplete type}).}} to satisfy the single-return-statement requirements{\cprotect\footnote{Notice that this \lstinline!Count!
\emcppsgloss{metafunction} also makes use (in its implementation) of
variadic class templates to parse a \emcppsgloss{type list} of unbounded
depth. Had this been a C++03 implementation, we would have been forced
to create an approximation (to the simple class-template
specialization containing the \emcppsgloss{parameter pack}
\lstinline!Tail...!) consisting of a bounded number (e.g., 20) of simple
(class) template specializations, each one taking an increasing number
of template arguments:

\begin{emcppshiddenlisting}[emcppsbatch=e3]
#include <type_traits>  // (ù{\codeincomments{std::integral\_constant}}ù), (ù{\codeincomments{std::is\_same}}ù)
template <typename...> struct TypeList { };
template <typename X, typename list> struct Count {};
\end{emcppshiddenlisting}
\begin{emcppslisting}[emcppsbatch=e3,style=footcode]
template <typename X, typename Y>
struct Count<X, TypeList<Y>>
    : std::integral_constant<int, std::is_same<X, Y>::value> { };
    // (class) template specialization for one argument

template <typename X, typename Y, typename Z>
struct Count<X, TypeList<Y, Z>>
    : std::integral_constant<int,
        std::is_same<X, Y>::value + std::is_same<X, Z>::value> { };
    // (class) template specialization for two arguments

template <typename X, typename Y, typename Z, typename A>
struct Count<X, TypeList<Y, Z, A>>
    : std::integral_constant<int,
        std::is_same<X, Y>::value + Count<X, TypeList<Z, A>>::value> { };
    // recursive (class) template specialization for three arguments

// ...
\end{emcppslisting}
      }}:

\begin{emcppshiddenlisting}[emcppsbatch=e4]
template <typename...> struct TypeList { };
\end{emcppshiddenlisting}
\begin{emcppslisting}[emcppsbatch=e4,label={relaxedconstexpr-countcode}]
#include <type_traits>  // (ù{\codeincomments{std::integral\_constant}}ù), (ù{\codeincomments{std::is\_same}}ù)

template <typename X, typename List> struct Count;
    // general template used to characterize the interface for the (ù{\codeincomments{Count}}ù)
    // metafunction
    // Note that this general template is an incomplete type.

template <typename X>
struct Count<X, TypeList<>> : std::integral_constant<int, 0> { };
    // partial (class) template specialization of the general (ù{\codeincomments{Count}}ù) template
    // (derived from the integral-constant type representing a compile-time
    // (ù{\codeincomments{0}}ù)), used to represent the base case for the recursion --- i.e., when
    // the supplied (ù{\codeincomments{TypeList}}ù) is empty 
    // The payload (i.e., the enumerated (ù{\codeincomments{value}}ù) member of the base class) 
    // representing the number of elements in the list is (ù{\codeincomments{0}}ù).

template <typename X, typename Head, typename... Tail>
struct Count<X, TypeList<Head, Tail...>>
    : std::integral_constant<int,
        std::is_same<X, Head>::value + Count<X, TypeList<Tail...>>::value> { };
    // simple (class) template specialization of the general (ù{\codeincomments{count}}ù) template
    // for when the supplied list is not empty 
    // In this case, the second parameter will be partitioned as the first 
    // type in the sequence and the (possibly empty) remainder of the 
    // (ù{\codeincomments{TypeList}}ù). The compile-time value of the base class will be either the 
    // same as or one greater than the value accumulated in the (ù{\codeincomments{TypeList}}ù) so 
    // far, depending on whether the first element is the same as the one 
    // supplied as the first type to (ù{\codeincomments{Count}}ù).

static_assert(Count<int, TypeList<int, char, int, bool>>::value == 2, "");
\end{emcppslisting}
    
\noindent Notice that we made use of a C++11 \emcppsgloss{parameter pack},
\lstinline!Tail...! (see \featureref{\locationc}{variadictemplate}), in the
implementation of the simple template specialization to package up and
pass along any remaining types.

As should be obvious by now, the C++11 restriction encourages both
somewhat rarified metaprogramming-related knowledge and a
\emph{recursive} implementation that can be compile-time intensive in
practice.{\cprotect\footnote{For a more efficient C++11 version of
\lstinline!Count!, see \intraref{appendix-relaxedconstexpr}{appendix:-optimized-c++11-example-algorithms}, \intrarefsimple{constexpr-typelist-count-algorithm}.\textbf{AUs: another nested appendix xref}}} By exploiting C++14's relaxed
\lstinline!constexpr! rules, a simpler and typically more compile-time
friendly \emph{imperative} solution can be realized:

\begin{emcppshiddenlisting}[emcppsbatch=e5]
#include <type_traits>  // (ù{\codeincomments{std::integral\_constant}}ù), (ù{\codeincomments{std::is\_same}}ù)
\end{emcppshiddenlisting}
\begin{emcppslisting}[emcppsbatch=e5,emcppsstandards={c++14}]
template <typename X, typename... Ts>
constexpr int count()
{
    bool matches[sizeof...(Ts)] = { std::is_same<X, Ts>::value... };
        // Create a corresponding array of bits where (ù{\codeincomments{1}}ù) indicates sameness.

    int result = 0;
    for (bool m : matches)  // (C++11) range-based (ù{\codeincomments{for}}ù) loop
    {
        result += m;        // Add up (ù{\codeincomments{1}}ù) bits in the array.
    }

    return result;  // Return the accumulated number of matches.
}
\end{emcppslisting}
    
\noindent The implementation above --- though more efficient and comprehensible
--- will require some initial learning for those unfamiliar with modern
C++ variadics. The general idea here \pagebreak%%%%% 
 is to use \emcppsgloss{pack expansion}
in a nonrecursive manner{\cprotect\footnote{\emcppsgloss[pack expansion]{Pack expansion} is
a language construct that expands a \emcppsgloss{variadic pack} during
compilation, generating code for each element of the pack. This
construct, along with a \emcppsgloss{parameter pack} itself, is a
  fundamental building block of variadic templates,
  introduced in C++11. As a minimal example, consider the variadic
  function template, \lstinline!e!:

  \begin{emcppslisting}[style=footcode]
  template <int... Is> void e() { f(Is...); }
  \end{emcppslisting}
      
\noindent \lstinline!e! is a function template that can be instantiated
  with an arbitrary number of compile-time-constant integers. The
  \lstinline!int...!~\lstinline!Is! syntax declares a \emcppsgloss{variadic pack}
  of compile-time-constant integers. The \lstinline!Is...! syntax (used to
  invoke \lstinline!f!) is a basic form of pack expansion that will resolve
  to all the integers contained in the \lstinline!Is! pack, separated by
  commas. For instance, invoking
  \lstinline!e<0,!~\lstinline!1,!~\lstinline!2,!~\lstinline!3>()! results in the
  subsequent invocation of
  \lstinline!f(0,!~\lstinline!1,!~\lstinline!2,!~\lstinline!3)!. Note that --- as
  seen in the \lstinline!count! example (which starts on page~\pageref{relaxedconstexpr-countcode}) --- any arbitrary
  expression containing a variadic pack can be expanded:

  \begin{emcppslisting}[style=footcode]
  template <int... Is> void g() { h((Is > 0)...); }
  \end{emcppslisting}
      
\noindent The \lstinline!(Is!~\lstinline!>!~\lstinline!0)...! expansion (above) will
  resolve to \lstinline!N! comma-separated Boolean values, where \lstinline!N!
  is the number of elements contained in the \lstinline!Is!
  \emcppsgloss{variadic pack}. As an example of this expansion, invoking
  \lstinline!g<5,!~\lstinline!-3,!~\lstinline!9>()! results in the subsequent invocation of \lstinline!h(true,!~\lstinline!false,!~\lstinline!true)!.}} to
initialize the \lstinline!matches! array with a sequence of zeros and ones
(representing, respectively, mismatch and matches between \lstinline!X! and
a type in the \lstinline!Ts...! pack) and then iterate over the array to
accumulate the number of ones as the final \lstinline!result!. This
\lstinline!constexpr!-based solution is both easier to understand and
typically faster to compile.{\cprotect\footnote{For a type list
containing 1024 types, the imperative (C++14) solution compiles about
twice as fast on GCC 10.x and roughly 2.6 times faster on Clang
  10.x.}}

\subsection[Potential Pitfalls]{Potential Pitfalls}\label{potential-pitfalls}

None so far

\subsection[Annoyances]{Annoyances}\label{annoyances}

None so far

\subsection[See Also]{See Also}\label{see-also}

\begin{itemize}
\item{\seealsoref{constexprfunc}{\seealsolocationc}Conditionally safe C++11 feature that first introduced compile-time evaluations of functions.}
\item{\seealsoref{constexprvar}{\seealsolocationc}Conditionally safe C++11 feature that first introduced variables usable as constant expressions.}
\item{\seealsoref{variadictemplate}{\seealsolocationc}Conditionally safe C++11 feature allowing templates to accept an arbitrary number of parameters.}
\end{itemize}

\subsection[Further Reading]{Further Reading}\label{further-reading}

None so far

\subsection[Appendix]{Appendix}\label{appendix-relaxedconstexpr}

\subsubsection[Optimized C++11 example algorithms]{Optimized C++11 example algorithms}\label{appendix:-optimized-c++11-example-algorithms}

\paragraph[Recursive Fibonacci]{Recursive Fibonacci}\label{recursive-fibonacci}

Even with the restrictions imposed by C++11, we can write a
more efficient recursive algorithm to calculate the \emph{n\/}th
Fibonacci number:

\begin{emcppslisting}
#include <utility>  // (ù{\codeincomments{std::pair}}ù)

constexpr std::pair<long long, long long> fib11NextFibs(
    const std::pair<long long, long long> prev,  // last two calculations
    int count)                                   // remaining steps
{
    return (count == 0) ? prev : fib11NextFibs(
        std::pair<long long, long long>(prev.second,
                                        prev.first + prev.second),
        count - 1);
}

constexpr long long fib11Optimized(long long n)
{
    return fib11NextFibs(
        std::pair<long long, long long>(0, 1), // first two numbers
        n                                      // number of steps
    ).second;
}
\end{emcppslisting}
    

\paragraph[\lstinline!constexpr! type list \lstinline!Count! algorithm]{{\ParaCode constexpr} type list {\ParaCode Count} algorithm}\label{constexpr-typelist-count-algorithm}

As with the \lstinline!fib11Optimized! example, providing a more efficient version of the \lstinline!Count! algorithm in
C++11 is also possible, by accumulating the final result through recursive
\lstinline!constexpr! function invocations:

\begin{emcppslisting}
#include <type_traits>  // (ù{\codeincomments{std::is\_same}}ù)

template <typename>
constexpr int count11Optimized() { return 0; }
    // Base case: always return (ù{\codeincomments{0}}ù).

template <typename X, typename Head, typename... Tail>
constexpr int count11Optimized()
    // Recursive case: compare the desired type ((ù{\codeincomments{X}}ù)) and the first type in
    // the list ((ù{\codeincomments{Head}}ù)) for equality, turn the result of the comparison
    // into either (ù{\codeincomments{1}}ù) (equal) or (ù{\codeincomments{0}}ù) (not equal), and recurse with the rest
    // of the type list ((ù{\codeincomments{Tail...}}ù)).
{
    return (std::is_same<X, Head>::value ? 1 : 0)
        + count11Optimized<X, Tail...>();
}
\end{emcppslisting}
    
\noindent This algorithm can be optimized even further in C++11 by using
a technique similar to the one shown for the iterative C++14
implementation. By leveraging a \lstinline!std::array! as compile-time
storage for bits where \lstinline!1! indicates equality between types, we can compute the final result with a fixed number of template
instantiations:

\begin{emcppslisting}
#include <array>        // (ù{\codeincomments{std::array}}ù)
#include <type_traits>  // (ù{\codeincomments{std::is\_same}}ù)

template <int N>
constexpr int count11VeryOptimizedImpl(
    const std::array<bool, N>& bits,  // storage for "type sameness" bits
    int i)                            // current array index
{
    return i < N
        ? bits[i] + count11VeryOptimizedImpl<N>(bits, i + 1)
            // Recursively read every element from the (ù{\codeincomments{bits}}ù) array and
            // accumulate into a final result.
        : 0;
}

template <typename X, typename... Ts>
constexpr int count11VeryOptimized()
{
    return count11VeryOptimizedImpl<sizeof...(Ts)>(
        std::array<bool, sizeof...(Ts)>{ std::is_same<X, Ts>::value... },
            // Leverage pack expansion to avoid recursive instantiations.
        0);
}
\end{emcppslisting}
    
\noindent Note that, despite being recursive, \lstinline!count11VeryOptimizedImpl!
will be instantiated only once with \lstinline!N! equal to the number of
elements in the \lstinline!Ts...! pack.





\newpage
%\section[{\itshape Generic} Lambdas]{Lambdas Having a Templated Call Operator}\label{genericlambda}



\emcppsFeature{
    short={{\itshape Generic} Lambdas},
    long={Lambdas Having a Templated Call Operator},
}{genericlambda}
\setcounter{table}{0}
\setcounter{footnote}{0}
\setcounter{lstlisting}{0}
%\section[{\itshape Generic} Lambdas]{Lambdas Having a Templated Call Operator}\label{genericlambda}



placeholder text.........
 
 


\newpage
%\section[Lambda Captures]{Lambda-Capture Expressions}\label{lambda-capture-expressions}
% 24 March 2021, no longer holding for revisions. Sending to Josh for code review. 
% 24 March 2021 JMB - compile check done
% 25 March 2021, commands updated



\emcppsFeature{
    short={Lambda Captures},
    long={Lambda-Capture Expressions},
}{lambda-capture-expressions}
\setcounter{table}{0}
\setcounter{footnote}{0}
\setcounter{lstlisting}{0}
%\section[Lambda Captures]{Lambda-Capture Expressions}\label{lambda-capture-expressions}

Lambda-capture expressions enable \emcppsgloss{synthetization} (spontaneous implicit creation) of
arbitrary data members within \emcppsgloss[closure]{closures} generated by
lambda expressions (see \featureref{\locationc}{lambda}).

\subsection[Description]{Description}\label{description}

In C++11, lambda expressions can capture variables in the
surrounding scope either \emph{by value} or \emph{by reference}{\cprotect\footnote{We use the familiar (C++11) feature
  \lstinline!auto! (see \featureref{\locationc}{auto-feature}) to deduce a closure's type since there is no
  way to name such a type explicitly.}}:

\begin{emcppslisting}
void test()
{
    int i = 0;
    auto f0 = [i]{ };   // Create a copy of (ù{\codeincomments{i}}ù) in the closure named (ù{\codeincomments{f0}}ù).
    auto f1 = [&i]{ };  // Store a reference to (ù{\codeincomments{i}}ù) in the closure named (ù{\codeincomments{f1}}ù).
}
\end{emcppslisting}
    
\noindent Although one could specify \emph{which} and \emph{how} existing
variables were captured, the programmer had no control over the creation
of new variables within a \emcppsgloss{closure}. C++14 extends the
\emcppsgloss[lambda introducer]{lambda-introducer} syntax to support implicit creation of
arbitrary data members inside a \emcppsgloss{closure} via either \emcppsgloss{copy
initialization} or \emcppsgloss{list initialization}:

\begin{emcppslisting}[emcppsstandards={c++14}]
auto f2 = [i = 10]{ /* body of closure */ };
    // Synthesize an (ù{\codeincomments{int}}ù) data member, (ù{\codeincomments{i}}ù), initialized with (ù{\codeincomments{10}}ù) in the closure.

auto f3 = [c{'a'}]{ /* body of closure */ };
    // Synthesize a (ù{\codeincomments{char}}ù) data member, (ù{\codeincomments{c}}ù), initialized with (ù{\codeincomments{'a'}}ù) in the closure.
\end{emcppslisting}
    
\noindent Note that the identifiers \lstinline!i! and \lstinline!c! above do not refer
to any existing variable; they are specified by the programmer creating
the closure. For example, the \emcppsgloss{closure} type assigned (i.e.,
bound) to \lstinline!f2! (above) is similar in functionality to an
\emcppsgloss{invocable} \lstinline!struct! containing an \lstinline!int! data
member:

\begin{emcppslisting}[emcppsstandards={c++14}]
// pseudocode
struct f2LikeInvocableStruct
{
    int i = 10;  // The type (ù{\codeincomments{int}}ù) is deduced from the initialization expression.
    auto operator()() const { /* closure body */ }  // The (ù{\codeincomments{struct}}ù) is invocable.
};
\end{emcppslisting}
    
\noindent The type of the data member is deduced from the initialization
expression provided as part of the capture in the same vein as
\lstinline!auto! (see \featureref{\locationc}{auto-feature}) type deduction; hence, it's not possible to
synthesize an uninitialized \emcppsgloss{closure} data member:

\begin{emcppslisting}
auto f4 = [u]{ };    // Error, (ù{\codeincomments{u}}ù) initializer is missing for lambda capture.
auto f5 = [v{}]{ };  // Error, (ù{\codeincomments{v}}ù)'s type cannot be deduced.
\end{emcppslisting}
    
\noindent It is possible, however, to use variables outside the scope of the
lambda as part of a lambda-capture expression (even capturing them \emph{by
reference} by prepending the \lstinline!&! token to the name of the
synthesized data member):

\begin{emcppslisting}[emcppsstandards={c++14}]
int i = 0;  // zero-initialized (ù{\codeincomments{int}}ù) variable defined in the enclosing scope

auto f6 = [j   = i]{ };  // OK, the local (ù{\codeincomments{j}}ù) data member is a copy of (ù{\codeincomments{i}}ù).
auto f7 = [&ir = i]{ };  // OK, the local (ù{\codeincomments{ir}}ù) data member is an alias to (ù{\codeincomments{i}}ù).
\end{emcppslisting}
    
\noindent Though capturing \emph{by reference} is possible, enforcing \lstinline!const! on a lambda-capture expression is not:

\begin{emcppslisting}[emcppsstandards={c++14}]
auto f8 = [const i = 10]{ };                    // Error, invalid syntax
auto f9 = [const auto i = 10]{ };               // Error, invalid syntax
auto fA = [i = static_cast<const int>(10)]{ };  // OK, (ù{\codeincomments{const}}ù) is ignored.
\end{emcppslisting}
    
\noindent The initialization expression is evaluated during the \emph{creation} of
the closure, not its \emph{invocation}:

\begin{emcppslisting}[emcppsstandards={c++14}]
#include <cassert>  // standard C (ù{\codeincomments{assert}}ù) macro

void g()
{
    int i = 0;

    auto fB = [k = ++i]{ };  // (ù{\codeincomments{++i}}ù) is evaluated at creation only.
    assert(i == 1);  // OK

    fB();  // Invoke (ù{\codeincomments{fB}}ù) (no change to (ù{\codeincomments{i}}ù)).
    assert(i == 1);  // OK
}
\end{emcppslisting}
    
\noindent Finally, using the same identifier as an existing
variable is possible for a synthesized capture, resulting in the original variable
being \emcppsgloss{shadowed} (essentially hidden) in the lambda expression's
body but not in its \emcppsgloss{declared interface}. In the example below,
we use the (C++11) compile-time operator
\lstinline!decltype! (see \featureref{\locationa}{decltype}) to infer the C++ type from the
initializer in the capture to create a parameter of that same type as
that part of its \emcppsgloss{declared interface}{\cprotect\footnote{Note
that, in the shadowing example defining \lstinline!fC!, GCC version 10.x
incorrectly evaluates \lstinline!decltype(i)! inside the body of the
lambda expression as \lstinline!const!~\lstinline!char!, rather than
  \lstinline!char!; see \intraref{potential-pitfalls-lambdacapture}{forwarding-an-existing-variable-into-a-closure-always-results-in-an-object-(never-a-reference)}.}}{\cprotect\footnote{Here we are using the (C++14) variable
  template (see \featureref{\locationb}{variable-templates}) version of the standard \lstinline!is_same! metafunction where \lstinline!std::is_same<A,!~\lstinline!B>::value! is replaced with
  \lstinline!std::is_same_v<A,!~\lstinline!B>!.}}:

\begin{emcppslisting}[emcppserrorlines={10},emcppsstandards={c++14}]
#include <type_traits>  // (ù{\codeincomments{std::is\_same}}ù)

int i = 0;

auto fC = [i = 'a'](decltype(i) arg)
{
    static_assert(std::is_same<decltype(arg), int>::value, "");
         // (ù{\codeincomments{i}}ù) in the interface (same as (ù{\codeincomments{arg}}ù)) refers to the (ù{\codeincomments{int}}ù) parameter.

    static_assert(std::is_same<decltype(i), char>::value, "");
        // (ù{\codeincomments{i}}ù) in the body refers to the (ù{\codeincomments{char}}ù) data member deduced at capture.
};
\end{emcppslisting}
    
\noindent Notice that we have again used \lstinline!decltype!, in conjunction with
the standard \lstinline!is_same! metafunction (which is \lstinline!true! if
and only if its two arguments are the same C++ type). This time, we're using \lstinline!decltype! to
demonstrate that the type (\lstinline!int!), extracted from the local
variable \lstinline!i! within the declared-interface portion of
\lstinline!fC!, is distinct from the type (\lstinline!char!) extracted from
the \lstinline!i! within \lstinline!fC!'s body. In other words, the effect
of initializing a variable in the capture portion of the lambda is to
hide the name of an existing variable that would otherwise be accessible
in the lambda's body.{\cprotect\footnote{Also note that, since the
deduced \lstinline!char! member variable, \lstinline!i!, is not materially
used (\emcppsgloss{ODR-used}) in the body of the lambda expression assigned
(bound) to \lstinline!fC!, some compilers, e.g., Clang, may warn:

\begin{lstlisting}[language=bash,style=footcodeplain]
warning: lambda capture 'i' is not required to be captured for this use
\end{lstlisting}
      }}

\pagebreak%%%%%%%
\subsection[Use Cases]{Use Cases}\label{use-cases-lambdacapture}

\subsubsection[Moving (as opposed to copying) objects into a closure]{Moving (as opposed to copying) objects into a closure}\label{moving-(as-opposed-to-copying)-objects-into-a-closure}

Lambda-capture expressions can be used to \emph{move} (see
\featureref{\locationc}{Rvalue-References}) an existing variable into a
closure{\cprotect\footnote{Though possible, it is surprisingly difficult
in C++11 to \emph{move} from an existing variable into a closure.
Programmers are either forced to pay the price of an unnecessary copy or to employ esoteric and fragile techniques, such as writing a wrapper
that hijacks the behavior of its copy constructor to do a \emph{move}
instead:

\begin{emcppslisting}[style=footcode]
#include <utility>  // (ù{\codeincomments{std::move}}ù)
#include <memory>   // (ù{\codeincomments{std::unique\_ptr}}ù)

template <typename T>
struct MoveOnCopy  // wrapper template used to hijack copy ctor to do move
{
    T d_obj;

    MoveOnCopy(T&& object) : d_obj{std::move(object)} { }
    MoveOnCopy(MoveOnCopy& rhs) : d_obj{std::move(rhs.d_obj)} { }
};

void f()
{
    std::unique_ptr<int> handle{new int(100)};  // move-only
        // Create an example of a handle type with a large body.

    MoveOnCopy<decltype(handle)> wrapper(std::move(handle));
        // Create an instance of a wrapper that moves on copy.

    auto &&lambda = [wrapper](){ /* use (ù{\codeincomments{wrapper.d\_obj}}ù) */ };
        // Create a "copy" from a wrapper that is captured by value.
}
\end{emcppslisting}
    
\noindent In the example above, we make use of the bespoke (``hacked'')
\lstinline!MoveOnCopy! class template to wrap a movable object;
when the lambda-capture expression tries to \emph{copy} the wrapper (\emph{by value}),
the wrapper in turn \emph{moves} the wrapped
  \lstinline!handle! into the body of the closure.}} (as opposed to
capturing it \emph{by copy} or \emph{by reference}). As an example of
\emph{needing} to move from an existing object into a closure, consider
the problem of accessing the data managed by
\emcppsgloss[std::uniqueptr]{\lstinline!std::unique_ptr!} (movable but not copyable) from a
separate thread --- for example, by enqueuing a task in a \emcppsgloss{thread
pool}:

\begin{emcppshiddenlisting}[emcppsbatch=e1,emcppsstandards={c++14}]
#include <memory>   // (ù{\codeincomments{std::unique\_ptr}}ù)
struct ThreadPool {
  using Handle = int;
  template <typename T>
  Handle enqueueTask(T&& t) { return 0; }
};
ThreadPool& getThreadPool();
struct Dataset {};
int processDataset(const std::unique_ptr<Dataset> &);
\end{emcppshiddenlisting}
\begin{emcppslisting}[emcppsbatch=e1]
ThreadPool::Handle processDatasetAsync(std::unique_ptr<Dataset> dataset)
{
    return getThreadPool().enqueueTask([data = std::move(dataset)]
    {
        return processDataset(data);
    });
}
\end{emcppslisting}
    
\noindent As illustrated above, the \lstinline!dataset! smart pointer is moved into
the closure passed to\linebreak[4] \mbox{\lstinline!enqueueTask!} by leveraging lambda-capture
expressions --- the \emcppsgloss[std::uniqueptr]{\lstinline!std::unique_ptr!} is \emph{moved}
to a different thread because a copy would have not been possible.

\subsubsection[Providing mutable state for a closure]{Providing mutable state for a closure}\label{providing-mutable-state-for-a-closure}

Lambda-capture expressions can be useful in conjunction with
\lstinline!mutable! lambda expressions to provide an initial state that
will change across invocations of the closure. Consider, for instance,
the task of logging how many TCP packets have been received on a socket
(e.g., for debugging or monitoring purposes){\cprotect\footnote{In
this example, we are making use of the (C++11) \lstinline!mutable! feature
  of lambdas to enable the counter to be modified on each invocation.}}:

\begin{emcppshiddenlisting}[emcppsbatch=e2]
#include <iostream>  // (ù{\codeincomments{std::cout}}ù)
struct TcpSocket {
    TcpSocket(int s);
    template <typename T>
    void onPacketReceived(T&& t) {}
};

\end{emcppshiddenlisting}
\begin{emcppslisting}[emcppsbatch=e2,emcppsstandards={c++14}]
void listen()
{
    TcpSocket tcpSocket(27015);  // some well-known port number
    tcpSocket.onPacketReceived([counter = 0]() mutable
    {
        std::cout << "Received " << ++counter << " packet(s)\n";
        // ...
    });
}
\end{emcppslisting}
    
\noindent Use of \lstinline!counter!~\lstinline!=!~\lstinline!0! as part of the
\emcppsgloss{lambda introducer} tersely produces a \emcppsgloss{function object}
that has an internal counter (initialized with zero), which is
incremented on every received packet. Compared to, say, capturing a
\lstinline!counter! variable \emph{by reference} in the closure, the solution
above limits the scope of \lstinline!counter! to the body of the lambda
expression and ties its lifetime to the closure itself, thereby
preventing any risk of dangling references.

\subsubsection[Capturing a modifiable copy of an existing \lstinline!const! variable]{Capturing a modifiable copy of an existing {\SubsubsecCode const} variable}\label{capturing-a-modifiable-copy-of-an-existing-const-variable}

Capturing a variable \emph{by value} in C++11 does allow the
programmer to control its \lstinline!const! qualification; the generated
closure data member will have the same \lstinline!const! qualification as
the captured variable, irrespective of whether the lambda is decorated
with \lstinline!mutable!:

\begin{emcppshiddenlisting}[emcppsbatch=e3,emcppsstandards={c++14}]
#include <type_traits>   // (ù{\codeincomments{std::is\_same}}ù)
\end{emcppshiddenlisting}
\begin{emcppslisting}[emcppsbatch=e3]
void f()
{
    int i = 0;
    const int ci = 0;

    auto lc = [i, ci]             // This lambda is not decorated with (ù{\codeincomments{mutable}}ù).
    {
        static_assert(std::is_same<decltype(i), int>::value, "");
        static_assert(std::is_same<decltype(ci), const int>::value, "");
    };

    auto lm = [i, ci]() mutable       // Decorating with (ù{\codeincomments{mutable}}ù) has no effect.
    {
        static_assert(std::is_same<decltype(i), int>::value, "");
        static_assert(std::is_same<decltype(ci), const int>::value, "");
    };
}
\end{emcppslisting}
    
\noindent In some cases, however, a lambda capturing a \lstinline!const! variable
\emph{by value} might need to modify that value when invoked. As an
example, consider the task of comparing the output of two Sudoku-solving
algorithms, executed in parallel:

\begin{emcppshiddenlisting}[emcppsbatch=e4,emcppsstandards={c++14}]
struct Puzzle {};
Puzzle generateRandomSudokuPuzzle();
struct ThreadPool {
  using Handle = int;
  template <typename T>
  Handle enqueueTask(T&& t) { return 0; }
};
ThreadPool& getThreadPool();
struct NaiveAlgorithm {};
struct FastAlgorithm {};
template <typename U, typename V>
void waitForCompletion(U&&,V&&);
\end{emcppshiddenlisting}
\begin{emcppslisting}[emcppsbatch=e4]
template <typename Algorithm> void solve(Puzzle&);
    // This (ù{\codeincomments{solve}}ù) function template mutates a Sudoku grid in place to solution.

void performAlgorithmComparison()
{
    const Puzzle puzzle = generateRandomSudokuPuzzle();
        // (ù{\codeincomments{const}}ù)-correct: (ù{\codeincomments{puzzle}}ù) is not going to be mutated after being
        // randomly generated.

    auto task0 = getThreadPool().enqueueTask([puzzle]() mutable
    {
        solve<NaiveAlgorithm>(puzzle);  // Error, (ù{\codeincomments{puzzle}}ù) is (ù{\codeincomments{const}}ù)-qualified.
        return puzzle;
    });

    auto task1 = getThreadPool().enqueueTask([puzzle]() mutable
    {
        solve<FastAlgorithm>(puzzle);  // Error, (ù{\codeincomments{puzzle}}ù) is (ù{\codeincomments{const}}ù)-qualified.
        return puzzle;
    });

    waitForCompletion(task0, task1);
    // ...
}
\end{emcppslisting}
    
\noindent The code above will fail to compile as capturing \lstinline!puzzle! will
result in a \lstinline!const!-qualified closure data member, despite the
presence of \lstinline!mutable!. A convenient workaround is to use a
(C++14) lambda-capture expression in which a local modifiable copy is deduced:

\begin{emcppslisting}[emcppsbatch=e4]
void performAlgorithmComparison2()
{
    // ...

    const Puzzle puzzle = generateRandomSudokuPuzzle();

    auto task0 = getThreadPool().enqueueTask([p = puzzle]() mutable
    {
        solve<NaiveAlgorithm>(p);  // OK, (ù{\codeincomments{p}}ù) is now modifiable.
        return p;
    });

    // ...
}
\end{emcppslisting}
    
\noindent Note that use of \lstinline!p!~\lstinline!=!~\lstinline!puzzle! (above) is
roughly equivalent to the creation of a new variable using
\lstinline!auto! (i.e.,
\lstinline!auto!~\lstinline!p!~\lstinline!=!~\lstinline!puzzle;!), which guarantees
that the type of \lstinline!p! will be deduced as a non-\lstinline!const!
\lstinline!Puzzle!. Capturing an existing \lstinline!const! variable as a mutable copy is
possible, but doing the opposite is not easy; see \intraref{annoyances-lambdacapture}{there’s-no-easy-way-to-synthesize-a-const-data-member}.

\subsection[Potential Pitfalls]{Potential Pitfalls}\label{potential-pitfalls-lambdacapture}

\subsubsection[Forwarding an existing variable into a closure always results in an object (never a reference)]{Forwarding an existing variable into a closure always results in\\ an object (never a reference)}\label{forwarding-an-existing-variable-into-a-closure-always-results-in-an-object-(never-a-reference)}

Lambda-capture expressions allow existing variables to be
\emcppsgloss{perfectly forwarded} (see \featureref{\locationc}{forwardingref})
into a closure:

\begin{emcppslisting}[emcppsstandards={c++14}]
#include <utility>  // (ù{\codeincomments{std::forward}}ù)

template <typename T>
void f(T&& x)  // (ù{\codeincomments{x}}ù) is of type forwarding reference to (ù{\codeincomments{T}}ù).
{
    auto lambda = [y = std::forward<T>(x)]  
        // Perfectly forward (ù{\codeincomments{x}}ù) into the closure.
    {
        // ... (use (ù{\codeincomments{y}}ù) directly in this lambda body)
    };
}
\end{emcppslisting}
    
\noindent Because \lstinline!std::forward<T>! can evaluate to a reference (depending
on the nature of \lstinline!T!), programmers might incorrectly assume that
a capture such as \lstinline!y!~\lstinline!=!~\lstinline!std::forward<T>(x)!
(above) is somehow either a capture \emph{by value} or a capture
\emph{by reference}, depending on the original \emcppsgloss{value category}
of \lstinline!x!.

Remembering that lambda-capture expressions work similarly to
\lstinline!auto! type deduction for variables, however,
reveals that such captures will \emph{always} result in an object, \emph{never} a reference:

\begin{emcppslisting}[emcppsignore={pseudocode}]
// pseudocode ((ù{\codeincomments{auto}}ù) is not allowed in a lambda introducer.)
auto lambda = [auto y = std::forward<T>(x)] { };
    // The capture expression above is semantically similar to an (ù{\codeincomments{auto}}ù)
    // (deduced-type) variable.
\end{emcppslisting}
    
\noindent If \lstinline!x! was originally an \romeovalue{lvalue}, then \lstinline!y! will be
equivalent to a \emph{by-copy} capture of \lstinline!x!. Otherwise,
\lstinline!y! will be equivalent to a \emph{by-move} capture of
\lstinline!x!.{\cprotect\footnote{Note that both \emph{by-copy} and
\emph{by-move} capture communicate \emcppsgloss{value} for
  \emcppsgloss[value semantic type (VST)]{value-semantic types}.}}

If the desired semantics are to capture \lstinline!x! \emph{by move} if it
originated from \emcppsgloss{rvalue} and \emph{by reference} otherwise, then
the use of an extra layer of indirection (using, e.g.,
\lstinline!std::tuple!) is required:

\begin{emcppslisting}[emcppsstandards={c++14}]
#include <tuple>  // (ù{\codeincomments{std::tuple}}ù)

template <typename T>
void f(T&& x)
{
    auto lambda = [y = std::tuple<T>(std::forward<T>(x))]
    {
        // ... (Use (ù{\codeincomments{std::get<0>(y)}}ù) instead of (ù{\codeincomments{y}}ù) in this lambda body.)
    };
}
\end{emcppslisting}
    
\noindent In the revised code example above, \lstinline!T! will be an \emcppsgloss[lvalue reference]{\romeovalue{lvalue} reference} if \lstinline!x! was originally an \romeovalue{lvalue}, resulting in
the \emcppsgloss{synthetization} of a \lstinline!std::tuple! containing an
\emcppsgloss[lvalue reference]{\romeovalue{lvalue} reference}, which --- in turn --- has semantics
equivalent to \lstinline!x!'s being captured \emph{by reference}.
Otherwise, \lstinline!T! will not be a reference type, and \lstinline!x! will
be \emph{moved} into the closure.

\subsection[Annoyances]{Annoyances}\label{annoyances-lambdacapture}

\subsubsection[There’s no easy way to synthesize a \lstinline!const! data member]{There’s no easy way to synthesize a {\SubsubsecCode const} data member}\label{there’s-no-easy-way-to-synthesize-a-const-data-member}

Consider the (hypothetical) case where the programmer desires to capture
a copy of a non-\lstinline!const! integer \lstinline!k! as a \lstinline!const!
closure data member:

\begin{emcppslisting}[emcppsstandards={c++14}]
void test1()
{
    int k;
    [k = static_cast<const int>(k)]() mutable  // (ù{\codeincomments{const}}ù) is ignored
    {
        ++k;  // "OK" -- i.e., compiles anyway even though we don't want it to
    };
}
\end{emcppslisting}

\begin{emcppslisting}[emcppsstandards={c++14}]
void test2()
{
    int k;
    [const k = k]() mutable  // Error, invalid syntax
    {
        ++k;  // no easy way to force this variable to be (ù{\codeincomments{const}}ù)
    };
}
\end{emcppslisting}
    
\noindent The language simply does not provide a convenient mechanism for
synthesizing, from a modifiable variable, a \lstinline!const! data member.
If such a \lstinline!const! data member somehow proves to be necessary, we
can either create a \lstinline!ConstWrapper! \lstinline!struct! (that adds
\lstinline!const! to the captured object) or write a full-fledged
\emcppsgloss{function object} in lieu of the leaner \emcppsgloss[lambda expressions]{lambda
expression}. Alternatively, a \lstinline!const! copy of the object can be
captured with traditional (C++11) lambda-capture expressions:

\begin{emcppslisting}
int test3()
{
    int k;
    const int kc = k;

    auto l = [kc]() mutable
    {
        ++kc;  // Error, increment of read-only variable (ù{\codeincomments{kc}}ù)
    };
}
\end{emcppslisting}
    

\subsubsection[\lstinline!std::function! supports only copyable callable objects]{{\SubsubsecCode std::function} supports only copyable callable objects}\label{std::function-supports-only-copyable-callable-objects}

Any lambda expression capturing a move-only object produces a closure
type that is itself movable but \emph{not} copyable:

\begin{emcppshiddenlisting}[emcppsbatch=e4,emcppsstandards={c++14}]
#include <memory>   // (ù{\codeincomments{std::unique\_ptr}}ù)
#include <utility>  // (ù{\codeincomments{std::move}}ù)
\end{emcppshiddenlisting}
\begin{emcppslisting}[emcppsbatch=e4]
void f()
{
    std::unique_ptr<int> moo(new int);    // some move-only object
    auto la = [moo = std::move(moo)]{ };  // lambda that does move capture

    static_assert(false == std::is_copy_constructible<decltype(la)>::value, "");
    static_assert( true == std::is_move_constructible<decltype(la)>::value, "");
}
\end{emcppslisting}
    
\noindent Lambdas are sometimes used to initialize instances of
\lstinline!std::function!, which requires the stored \emcppsgloss{callable
object} to be copyable:

\begin{emcppslisting}
std::function<void()> f = la;  // Error, (ù{\codeincomments{la}}ù) must be copyable.
\end{emcppslisting}
    
\noindent Such a limitation --- which is more likely to be encountered when using
lambda-capture expressions --- can make \lstinline!std::function!
unsuitable for use cases where move-only closures might conceivably be
reasonable. Possible workarounds include (1) using a different
type-erased, \emcppsgloss{callable object} wrapper type that supports
move-only callable objects,{\cprotect\footnote{The
\lstinline!any_invocable! library type, proposed for C++23, is an
example of a type-erased wrapper for move-only callable objects; see
  \cite{calabrese20}.}} (2) taking a performance hit by wrapping the
desired \emcppsgloss{callable object} into a copyable wrapper (such as
\lstinline!std::shared_ptr!), or (3) designing software such that
noncopyable objects, once constructed, never need to
move.\footnote{For an in-depth discussion of how large systems can
benefit from a design that embraces local arena memory allocators and,
thus, minimizes the use of moves across natural memory boundaries
identified throughout the system, see \cite{lakos22}.}

\subsection[See Also]{See Also}\label{see-also}

\begin{itemize}

\item{\seealsoref{auto-feature}{\seealsolocationc}offers a model with the same type deduction rules.}
\item{\seealsoref{bracedinit}{\seealsolocationc}illustrates one possible way of initializing the captures.}
\item{\seealsoref{forwardingref}{\seealsolocationc}describes a feature that contributes to a source of misunderstanding of this feature.}
\item{\seealsoref{lambda}{\seealsolocationc}provides the needed background for understanding the feature in general.}
\item{\seealsoref{Rvalue-References}{\seealsolocationc}gives a full description of an important feature used in conjunction with movable types.}
\end{itemize}

\subsection[Further Reading]{Further Reading}\label{further-reading}

None so far








%%%%%%%%%%%%%%%% MATERIAL FOR DESIGN TESTING
%\newpage
%\addtocontents{toc}{\protect\renewcommand*\protect\cftsecindent{4.5em}}
%\section[Short Title]{Longer Wordier Title for Test}\label{justfortest}
%
%Section text; just to test where/how the C++ change (from C++11 to C++14) in the RH must be placed in relation to sectioning commands such that the correct one displays
%
%\begin{table}[h!] 
%\begin{center}
%\begin{threeparttable}
%\caption{Testing Table Numbering}\label{test-table3}\vspace{1.5ex} 
%{\small \begin{tabular}{c|c|c}\thickhline 
%\rowcolor[gray]{.9}   {\sffamily\bfseries Compiler} & {\sffamily\bfseries Compiler-Specific} &
%{\sffamily\bfseries Standard-Conforming} \\ \hline 
%GCC &\texttt{\_\_attribute\_\_((pure))} & \texttt{[[gnu::pure]]} \\ \hline
%Clang & \texttt{\_\_attribute\_\_((no\_sanitize))} &\texttt{[[clang::no\_sanitize]]} \\ \hline 
%MSVC & \texttt{declspec(deprecated)} & \texttt{[[deprecated]]} \\ \thickhline
%\end{tabular}
%}
%\end{threeparttable} 
%\end{center}
%\end{table}
%
%Donec mattis, ex a ornare auctor, tortor felis pulvinar odio, eu mollis tortor orci id ligula. Vestibulum dignissim magna vitae lectus cursus, ac sollicitudin mi pulvinar. Praesent vulputate lorem lectus, eu pulvinar mi finibus eu. Donec sit amet massa massa. Sed eget molestie quam.\footnote{testing a footnote} Praesent iaculis diam ut nunc condimentum auctor. Morbi porttitor varius ante vel venenatis. Integer ut porttitor elit, id lacinia turpis. Integer id iaculis ante, placerat ornare tortor. Sed et consectetur est, sit amet rhoncus turpis.
%
%\begin{table}[h!] 
%\begin{center}
%\begin{threeparttable}
%\caption{Testing Table Numbering}\label{test-table4}\vspace{1.5ex} 
%{\small \begin{tabular}{c|c|c}\thickhline 
%\rowcolor[gray]{.9}   {\sffamily\bfseries Compiler} & {\sffamily\bfseries Compiler-Specific} &
%{\sffamily\bfseries Standard-Conforming} \\ \hline 
%GCC &\texttt{\_\_attribute\_\_((pure))} & \texttt{[[gnu::pure]]} \\ \hline
%Clang & \texttt{\_\_attribute\_\_((no\_sanitize))} &\texttt{[[clang::no\_sanitize]]} \\ \hline 
%MSVC & \texttt{declspec(deprecated)} & \texttt{[[deprecated]]} \\ \thickhline
%\end{tabular}
%}
%\end{threeparttable} 
%\end{center}
%\end{table}
%
%This feature, when used in conjunction with \emph{explicit instantiation
%definitions}, can significantly improve compilation times for a set of
%translation units that often instantiate common templates:
%
%\noindent\begin{minipage}{.45\textwidth}
%\begin{lstlisting}[language=C++, caption=code 1,frame=tb]{Name}
%void code()
%{
%
%}
%\end{lstlisting}
%\end{minipage}\hfill
%\begin{minipage}{.45\textwidth}
%\begin{lstlisting}[language=C++, caption=code 2,frame=tb]{Name}
%void code()
%{
%
%}
%\end{lstlisting}
%\end{minipage}
%
%Nullam nibh tortor, finibus ut lectus eu, convallis vehicula libero. Cras maximus ligula nisl, a eleifend mauris venenatis nec. Suspendisse potenti. Maecenas viverra laoreet mauris ut laoreet. Donec non felis risus. Ut faucibus, justo id luctus lobortis, nisi lacus pharetra est, id pretium arcu ligula ac magna. Morbi a nisl sit amet lacus vehicula tempor nec sit amet nibh. 
%
%\noindent\begin{minipage}{.45\textwidth}
%\begin{lstlisting}[language=C++, caption=code 1 with a long and wrapping title,frame=tb]{Name}
%void code()
%{
%
%}
%\end{lstlisting}
%\end{minipage}\hfill
%\begin{minipage}{.45\textwidth}
%\begin{lstlisting}[language=C++, caption=code 2 with a long and wrapping title,frame=tb]{Name}
%void code()
%{
%
%}
%\end{lstlisting}
%\end{minipage}
%
%Pellentesque id lorem ac sem ullamcorper semper. Donec imperdiet sapien in nisi faucibus, et tempus lacus hendrerit. Nulla laoreet risus eu tortor ultrices, sed bibendum neque sodales. Ut faucibus ipsum id convallis convallis. Vivamus vehicula rutrum metus in semper. Morbi fringilla ex vel vulputate vehicula. Maecenas ut porttitor massa.
%




  % Chapter 2
\cleardoublepage

\addtocontents{toc}{\protect\newpage}
\chapter[Unsafe Features]{Unsafe Features}\label{ch-unsafe}
\emcppschapterstart{}

PRODUCTION: Material copied from Ch 0; AUs will edit/replace on FPPs. 

AUTHORS: Remember to add material about whether these features are easy to teach. 

When an expert programmer uses any C++ feature appropriately, the feature typically does no direct harm. Yet other developers --- seeing the feature’s use in the code base but failing to appreciate the highly specialized or nuanced reasoning justifying it --- might attempt to use it in what they perceive to be a similar way, yet with profoundly less desirable results. Similarly, maintainers may change the use of a fragile feature altering its semantics in subtle but damaging ways.

Features that are classified as unsafe are those that might have valid, and even very important, use cases, yet our experience indicates that routine or widespread use thereof would be counterproductive in a typical large-scale software-development enterprise.

For example, we deem the final contextual keyword an unsafe feature because the situations in which it would be misused overwhelmingly outnumber those vanishingly few isolated cases in which it is appropriate, let alone valuable. Furthermore, its widespread use would inhibit fine-grained (e.g., hierarchical) reuse, which is critically important to the success of a large organization.

%%%%%%%%%%%%%%% C++11
\stepcounter{cppxx}
\addtocontents{toc}{\protect\renewcommand*\protect\cftsecindent{1.25em}}
\cftaddnumtitleline{toc}{section}{\thechapter.\thecppxx}{C++11}{\thepage}
% to add to the TOC at the section level
\renewcommand{\cppxx}{C++11}

\newpage
\addtocontents{toc}{\protect\renewcommand*\protect\cftsecindent{4.5em}}
% \section[{\tt carries\_dependency}]{The {\SecCode [[carries\_dependency]]} Attribute}\label{carriesdependency}
% 19 Feb 2021, ready for Josh's code fixes
% 20 Feb 2021 JMB - code cleaned up, compiles.
% 2 March 2021, copyedits entered and proofed. 
% 18 March 2021, author revisions made after copyediting; entered and proofed.
% 26 March 2021, in FPPs 


\emcppsFeature{
    short={\lstinline!carries_dependency!},
    tocshort={\TOCCode carries\_dependency},
    long={The {\SecCode [[carries\_dependency]]} Attribute},
    toclong={The \lstinline![[carries\_dependency]]! Attribute},
    rhshort={\RHCode carries\_dependency},
}{carriesdependency}
\label{the-carries_dependency-attribute}
\setcounter{table}{0}
\setcounter{footnote}{0}
\setcounter{lstlisting}{0}
 %\section[{\tt carries\_dependency}]{The {\SecCode [[carries\_dependency]]} Attribute}\label{carriesdependency}
%\subsection[The \lstinline![[carries_dependency]]! Attribute]{The {\SecCode [[carries\_dependency]]} Attribute}\label{the-[[carries_dependency]]-attribute}

The\lstinline![[carries_dependency]]! attribute provides a means to
manually identify function parameters and \lstinline!return! values as
components of \emcppsgloss[data dependency chain]{data dependency chains} to enable (primarily 
across translation units) use of the lighter-weight
\emcppsgloss[release consume]{release-consume} \emcppsgloss{synchronization paradigm} as an
optimization over the more conservative \emcppsgloss[release acquire]{release-acquire}
paradigm.{\cprotect\footnote{The authors would like to thank Michael
Wong, Paul McKenney, and Maged Michael for reviewing and contributing to this feature section.}}

\subsection[Description]{Description}\label{description}

C++11 ushered in support for multithreading by introducing a rigorously
specified memory model. The Standard Library provides support for
managing threads, including their execution, synchronization, and
intercommunication. As a part of the new memory model, the Standard
defines various \emcppsgloss[synchronization operation]{synchronization operations} that are classified as either \emph{sequentially
consistent}, \emph{release}, \emph{acquire}, 
\emph{release-and-acquire}, or \emph{consume} operations. These
operations play a key role in making changes in data in one thread
visible in another.

The modern C++ memory model describes two \emcppsgloss[synchronization
paradigm]{synchronization paradigms} that are used
to coordinate data flow among concurrent threads of execution. The current suite of supported
\emcppsgloss[synchronization paradigm]{synchronization paradigms} comprises \emcppsgloss[release acquire]{release-acquire} and \emcppsgloss[release consume]{release-consume}, although in practice
\emcppsgloss[release consume]{release-consume} is implemented in terms of \emcppsgloss[release acquire]{release-acquire} in all known implementations. In
particular, the \emcppsgloss[release consume]{release-consume} paradigm requires that the
compiler be given fine-grained understanding of the \emcppsgloss[intra thread dependency]{intra-thread dependencies} among the reads and writes within a program
and relates those to atomic \emph{release stores} and \emph{consume
loads} that happen concurrently across multiple threads of execution.
Dependency chains in the \emcppsgloss[release consume]{release-consume} synchronization
paradigm specify which evaluations following the \emph{consume load} are
\emcppsgloss{ordered after} a corresponding \emph{release store}.

\subsubsection[The release-acquire paradigm]{The release-acquire paradigm}\label{the-release-acquire-paradigm}

A \emph{release} operation writes a value to a memory location, and an
\emph{acquire} operation reads a value from a memory location. Although many have
referred to the release-acquire paradigm as \emph{acquire-release}, the proper, standard,
  time-ordered nomenclature is \emph{release-acquire}. In a \emcppsgloss[release acquire]{release-acquire} pair, the acquire operation reads the value written by
the release operation, which means that all of the reads and writes to
\emph{any} memory location \emph{before the release operation} happen
before \emph{all} of the reads and writes \emph{after the acquire operation}. Note that this paradigm does \emph{not} use 
dependency chains or the \lstinline![[carries_dependency]]! attribute. See \intraref{use-cases-carriesdependency}{producer-consumer-programming-pattern} 
%{\intraref{{Use Cases: A producer-consumer implementation}}} below 
for a
complete example that implements this paradigm.

\subsubsection[Data dependency]{Data dependency}\label{data-dependency}

In the current revisions of C++, \emcppsgloss{data dependency} is defined as
existing whenever the output of one evaluation is used as the input of
another. 
 When one evaluation has a data dependency on another
evaluation, the second evaluation is said to \emcppsgloss{carry dependency}
to the other. The Standard Library function
\lstinline!std::kill_dependency! is also related and can be used to
  \emph{break} a data dependency chain. Naturally the compiler must
ensure that any evaluation that depends on another must not be started
until the first evaluation is complete. A \emcppsgloss{data dependency chain}
is formed when multiple evaluations carry dependency transitively; the
output of one evaluation is used as the input of the next evaluation in
the chain.

\subsubsection[The release-consume paradigm]{The release-consume paradigm}\label{the-release-consume-paradigm}

Some systems use the read-copy-update (RCU) synchronization mechanism.
This approach preserves the order of \emph{loads} and \emph{stores} that
form in a \emcppsgloss{data dependency chain}, which is a sequence of
\emph{loads} and \emph{stores} in which the input to one operation is an
output of another. A compiler can use guaranteed order of loads
and stores provided by the RCU synchronization mechanism for performance purposes by omitting certain
\emcppsgloss[memory fence instruction]{memory-fence instructions} that would otherwise be required to
enforce the correct ordering. In such cases, however, ordering is
guaranteed only between those operations making up the relevant
\emcppsgloss{data dependency chain}. The C++ definition
of data dependency is intended to mimic the data dependency on RCU
systems. Note, however, that C++ currently defines data dependency in
terms of evaluations, while RCU data dependency is defined in terms of
  loads and stores.

This optimization was intended to be available in C++ through use of a
\emph{release-consume} pair, which, as its name
suggests, consists of a \emph{release-store} operation and a
\emph{consume-load} operation. A \emph{consume} operation is much like
an \emph{acquire} operation, except that it guarantees only the ordering
of those evaluations in a \emcppsgloss{data dependency chain}, starting with
the consume-load operation.

Note, however, that currently no known implementation is able to take
advantage of the current C++ \emph{consume} semantics; hence, all
current compilers promote \emph{consume} loads to \emph{acquire} loads,
effectively making the \lstinline![[carries_dependency]]! attribute
redundant. Revisions to render this feature implementable and therefore
usable are currently under consideration by the C++ Standards Committee.
Prototypes for various approaches have been produced. When a usable
feature with real implementations is delivered, it quite possibly will
not work exactly as described in the examples here; see \intrarefsimple{use-cases-carriesdependency}. 
% {\intraref{{Use Cases}}}, below.

\subsubsection[Using the {\protect\lstinline![[carries_dependency]]!} attribute]{Using the {\SubsubsecCode [[carries\_dependency]]} attribute}\label{using-the-[[carries_dependency]]-attribute}

\emcppsgloss[data dependency chain]{Data dependency chains} can and do propagate into and out of
called functions. If one of these interoperating functions is in a
separate translation unit, the compiler will have no way of seeing the
dependency chain. In such cases, the user can apply the
\lstinline![[carries_dependency]]! attribute to imbue the necessary
information for the compiler to track the propagation of dependency
chains into and out of functions across translation units, thus possibly
avoiding unnecessary memory-fence instructions; see \intrarefsimple{use-cases-carriesdependency}. Note that the Standard Library function
\lstinline!std::kill_dependency! is also related and can be used to
  \emph{break} a data dependency chain.
% {\intrarefsub{{Using the \lstinline![[carries_dependency]]! attribute}}}, below.

The \lstinline![[carries_dependency]]! attribute can be applied to a
function declaration as a whole by placing it in front of the function declaration,
in which case the attribute applies to the \lstinline!return! value:

\begin{emcppslisting}
[[carries_dependency]] int* f();  // attribute applied to entire function (ù{\codeincomments{f}}ù)
\end{emcppslisting}
    
\noindent In the example above, the \lstinline![[carries_dependency]]! attribute was
applied to the declaration of function \lstinline!f! to indicate that the
\lstinline!return! value carries a dependency out of the function. The
compiler may now be able to avoid emitting a memory-fence instruction
for the return value of \lstinline!f!.

The \lstinline![[carries_dependency]]! attribute can also be applied
to one or more of the function's parameter declarations by placing it
immediately after the parameter name:

\begin{emcppslisting}
void g(int* input [[carries_dependency]]); // attribute applied to (ù{\codeincomments{input}}ù)
\end{emcppslisting}
    
\noindent In the declaration of function \lstinline!g! in the example above, the
\lstinline![[carries_dependency]]! attribute is applied to the
\lstinline!input! parameter to indicate that a dependency is carried
through that parameter into the function, which may obviate the
compiler's having to emit an unnecessary memory-fence instruction for
the \lstinline!input! parameter; see \featureref{\locationa}{attributes}.
%{\intraref{{Attribute Syntax}}}.

In both cases, if a function or a parameter declaration specifies the
\lstinline![[carries_dependency]]! attribute, the first declaration of
that function shall specify that \lstinline![[carries_dependency]]!
attribute. Similarly, if the first declaration of a function or one of
its parameters specifies the \lstinline![[carries_dependency]]! attribute
in one translation unit and the first declaration of the same function
in another translation unit doesn't, the program is \emcppsgloss{ill formed, no diagnostic required (IFNDR)}.

%The user is responsible for ensuring that an existing dependency chain
%is available if needed for synchronization purposes. The
%\lstinline![[carries_dependency]]! attribute will not create a \mbox{dependency}.

It is important to note that while the \lstinline![[carries_dependency]]! attribute informs the compiler about the presence of a dependency chain, it does not itself create one.  The dependency chain must be present in the implementation to have any effect on synchronization.

\subsection[Use Cases]{Use Cases}\label{use-cases-carriesdependency}

\subsubsection[Producer-consumer programming pattern]{Producer-consumer programming pattern}\label{producer-consumer-programming-pattern}

The popular producer-consumer programming pattern uses
\emph{release-acquire} pairs to synchronize between threads:

%\begin{emcppslisting}
%#include <cassert>  // standard C (ù{\codeincomments{assert}}ù) macro
%#include <atomic>   // (ù{\codeincomments{std::atomic}}ù), (ù{\codeincomments{std::memory\_order\_release}}ù)
%                    // (ù{\codeincomments{std::memory\_order\_acquire}}ù)
%
%struct S
%{
%    int i;
%    char c;
%    double d;
%};
%
%S data;
%std::atomic<int> guard;
%
%void producerThread()
%{
%    data.i = 42;
%    data.c = 'c';
%    data.d = 5.0;
%    guard.store(1, std::memory_order_release);
%}
%
%void consumerThread()
%{
%    if (guard.load(std::memory_order_acquire) == 1)
%    {
%        // By dint of the release-acquire guarantee, we know that all the
%        // data changes are visible if the guard change is visible.
%        assert(data.i == 42);
%        assert(data.c == 'c');
%        assert(data.d == 5.0);
%    }
%}
%\end{emcppslisting}
\newpage%%%%%%%
\begin{emcppslisting}
// my_shareddata.h                                                              
void initSharedData();                                                          
    // Initialize the shared data of (ù{\codeincomments{my\_shareddata.o}}ù) to a well-known          
    // aggregation of values.                                                   
                                                                                
void accessSharedData();                                                        
    // Confirm that the shared data of (ù{\codeincomments{my\_shareddata.o}}ù) have been initialized  
    // and have their expected values.
\end{emcppslisting}
\vspace*{2ex}
\begin{emcppslisting}
// my_shareddata.cpp
#include <my_shareddata.h>

#include <atomic>  // (ù{\codeincomments{std::atomic}}ù), (ù{\codeincomments{std::memory\_order\_release}}ù), and
                   // (ù{\codeincomments{std::memory\_order\_acquire}}ù)
#include <cassert> // standard C (ù{\codeincomments{assert}}ù) macro

struct S
{
    int    i;
    char   c;
    double d;
};

static S                data;     // (ù{\codeincomments{static}}ù) for insulation
static std::atomic<int> guard(0); // (ù{\codeincomments{static}}ù) for insulation

void initSharedData()
{
    data.i = 42;
    data.c = 'c';
    data.d = 5.0;

    guard.store(1, std::memory_order_release);
}

void accessSharedData()
{
    while(0 == guard.load(std::memory_order_acquire))
        /* empty */ ;

    assert(42  == data.i);
    assert('c' == data.c);
    assert(5.0 == data.d);
}
\end{emcppslisting}
\newpage%%%%%%%
\begin{emcppslisting}
// my_app.cpp
#include <my_shareddata.h>
#include <thread>  // (ù{\codeincomments{std::thread}}ù)

int main()
{
    std::thread t2(accessSharedData);
    std::thread t1(  initSharedData);

    t1.join();
    t2.join();
}
\end{emcppslisting}
    
\noindent When this \emph{release-acquire} \emph{synchronization paradigm} is
used, the compiler must maintain the ordering of the statements to avoid
breaking the \emph{release-acquire} guarantee; the compiler will also
need to insert memory-fence instructions to prevent the hardware from
breaking this guarantee.

If we wanted to modify the example above to use \emph{release-consume}
semantics, we would somehow need to make the \lstinline!assert! statements
a part of the dependency chain on the \lstinline!load! from the
\lstinline!guard! object. We can accomplish this because reading data
through a pointer establishes a dependency chain between the reading of
that pointer value and the reading of the referenced data.  Since \emph{release-consume} allows the developer to specify that data of concern, using that policy instead of \emph{release-acquire} policy (in the code example above) allows the compiler to be more selective in its use of memory fences::

%\begin{emcppslisting}[emcppsbatch=e1]
%#include <cassert>  // standard C (ù{\codeincomments{assert}}ù) macro
%#include <atomic>   // (ù{\codeincomments{std::atomic}}ù), (ù{\codeincomments{std::memory\_order\_release}}ù)
%                    // (ù{\codeincomments{std::memory\_order\_consume}}ù)
%
%struct S
%{
%    int i;
%    char c;
%    double d;
%};
%
%S data;
%std::atomic<S*> guard(nullptr);
%
%void producerThread()
%{
%    data.i = 42;
%    data.c = 'c';
%    data.d = 5.0;
%    guard.store(&data, std::memory_order_release);
%}
%
%void consumerThread()
%{
%    S* setData = guard.load(std::memory_order_consume);
%    if (setData)
%    {
%        assert(setData->i == 42);
%        assert(setData->c == 'c');
%        assert(setData->d == 5.0);
%    }
%}
%\end{emcppslisting}
\enlargethispage*{2ex}
\begin{emcppslisting}
// my_shareddata.cpp (use (ù{\codeincomments{\_*consume}}ù), not (ù{\codeincomments{*\_acquire}}ù))                                                          
#include <my_shareddata.h>                                                      
                                                                                
#include <atomic>  // (ù{\codeincomments{std::atomic}}ù), (ù{\codeincomments{std::memory\_order\_release}}ù), and               
                   // (ù{\codeincomments{std::memory\_order\_consume}}ù) (not (ù{\codeincomments{*\_acquire}}ù))                 
#include <cassert> // standard C (ù{\codeincomments{assert}}ù) macro                                   
                                                                                
struct S                                                                        
{                                                                               
    /* definition not changed */                                                
};                                                                              
                                                                                
static S               data;           // (ù{\codeincomments{static}}ù) for insulation (as before)   
static std::atomic<S*> guard(nullptr); // guards just one (ù{\codeincomments{struct S}}ù).           
                                                                                
void initSharedData()                                                           
{                                                                               
    data.i = 42;   // as before                                                 
    data.c = 'c';  // as before                                                 
    data.d = 5.0;  // as before                                                 
                                                                                
    guard.store(&data, std::memory_order_release);  // Set (ù{\codeincomments{\&data}}ù), not 1.      
}                                                                               
                                                                                
void accessSharedData()                                                         
{                                                                               
    S *sharedDataPtr = nullptr;                                                 
                                                                                
    // Load using (ù{\codeincomments{*\_consume}}ù), not (ù{\codeincomments{*\_acquire}}ù).                                 
    while (nullptr == (sharedDataPtr = guard.load(std::memory_order_consume)))  
        /* empty */ ;                                                           
                                                                                
    assert(&data == sharedDataPtr);                                             
                                                                                
    assert(42  == sharedDataPtr->i);                                            
    assert('c' == sharedDataPtr->c);                                            
    assert(5.0 == sharedDataPtr->d);                                            
}
\end{emcppslisting}
    
\noindent Finally, if we want to start to refactor the work of
the \lstinline!my_sharedata! component into multiple 
 functions across different
translation units, we would want to carefully apply the 
\lstinline![[carries_dependency]]! attribute to the newly refactored
functions, so calling into these functions might conceivably be better
optimized:

%\begin{emcppslisting}[emcppsbatch=e1]
%[[carries_dependency]] S* loadData()
%{
%    return guard.load(std::memory_order_consume);
%}
%
%void checkData(S* s [[carries_dependency]])
%{
%    assert(s->i == 42);
%    assert(s->c == 'c');
%    assert(s->d == 5.0);
%}
%
%void betterThreadB()
%{
%    S* setData = loadData();
%    if (setData)
%    {
%        checkData(setData);
%    }
%}
%\end{emcppslisting}
\begin{emcppslisting}
// (ù{\codeincomments{my\_shareddataimpl.h}}ù)                                                          
                                                                                
struct S                                                                        
{                                                                               
    int    i;                                                                   
    char   c;                                                                   
    double d;                                                                   
};                                                                              
                                                                                
[[carries_dependency]] S* getSharedDataPtr();                                   
    // Return the address of the shared data in this translation unit.          
                                                                                
void releaseSharedData(S *sharedDataPtr [[carries_dependency]]);                
    // Release the shared data in this translation unit.  The behavior is       
    // undefined unless (ù{\codeincomments{getSharedDataPtr() == sharedDataPtr}}ù).                  
                                                                                
[[carries_dependency]] S* accessInitializedSharedData();                        
    // Return the address of the initialized shared data in this translation    
    // unit.                                                                    
                                                                                
void checkSharedDataValue(S* s [[carries_dependency]],                          
                          int    i,                                             
                          char   c,                                             
                          double d);                                            
    // Confirm that data at the specified (ù{\codeincomments{s}}ù) has the specified (ù{\codeincomments{i}}ù), (ù{\codeincomments{c}}ù), and   
    // (ù{\codeincomments{d}}ù) as constituent values.
\end{emcppslisting}
\newpage%%%%%%
\begin{emcppslisting}
// (ù{\codeincomments{my\_shareddataimpl.cpp}}ù)                                                        
                                                                                
#include <my_shareddataimpl.h>                                                  
                                                                                
#include <cassert>                                                              
#include <atomic>                                                               
                                                                                
static S               data;           // (ù{\codeincomments{static}}ù) for insulation               
static std::atomic<S*> guard(nullptr); // guards one (ù{\codeincomments{struct S}}ù).                
                                                                                
[[carries_dependency]] S* getSharedDataPtr()                                    
{                                                                               
    return &data;                                                               
}                                                                               
                                                                                
void releaseSharedData(S *sharedDataPtr [[carries_dependency]])                 
{                                                                               
    assert(&data == sharedDataPtr);                                             
                                                                                
    guard.store(sharedDataPtr, std::memory_order_release);                      
}                                                                               
                                                                                
[[carries_dependency]] S* accessInitializedSharedData()                         
{                                                                               
    S *sharedDataPtr = nullptr;                                                 
                                                                                
    while (nullptr == (sharedDataPtr = guard.load(std::memory_order_consume)))  
        /* empty */ ;                                                           
                                                                                
    assert(&data == sharedDataPtr);                                             
                                                                                
    return sharedDataPtr;                                                       
}  

void checkSharedDataValue(S*     s [[carries_dependency]],                      
                          int    i,                                             
                          char   c,                                             
                          double d)                                             
{                                                                               
    assert(i == s->i);                                                          
    assert(c == s->c);                                                          
    assert(d == s->d);                                                          
}    
\end{emcppslisting}
\newpage%%%%%
\begin{emcppslisting}
// (ù{\codeincomments{my\_shareddata.cpp}}ù) (re-factored to use (ù{\codeincomments{*impl}}ù))                              
#include <my_shareddata.h>                                                      
#include <my_shareddataimpl.h>                                                  
                                                                                
void initSharedData()                                                           
{                                                                               
    S *sharedDataPtr = getSharedDataPtr();                                      
                                                                                
    sharedDataPtr->i = 42;                                                      
    sharedDataPtr->c = 'c';                                                     
    sharedDataPtr->d = 5.0;                                                     
                                                                                
    releaseSharedData(sharedDataPtr);                                           
}                                                                               
                                                                                
void accessSharedData()                                                         
{                                                                               
    S *sharedDataPtr = accessInitializedSharedData();                           
    checkSharedDataValue(sharedDataPtr, 42, 'c', 5.0);                          
}       
\end{emcppslisting}

    
%\noindent Again, as of this writing, all known compilers implement \emph{consume}
%loads as \emph{acquire} loads and thus fail to provide the desired
%optimization.

\subsection[Potential Pitfalls]{Potential Pitfalls}\label{potential-pitfalls}

\subsubsection[No practical use on current platforms]{No practical use on current platforms}\label{no-practical-use-on-current-platforms}

All known compilers promote \emph{consume} loads to \emph{acquire}
loads, thus failing to omit superfluous memory-fence instructions.
Developers writing code with the expectation that it will be run under
the more efficient \emcppsgloss[release consume]{release-consume} \emcppsgloss{synchronization paradigm} will find that their code will continue to work --- as
expected --- under the more conservative \emcppsgloss[release acquire]{release-acquire}
guarantees until such time as a theoretical, not-yet-existent compiler
that properly supports the \emcppsgloss[release consume]{release-consume}
\emcppsgloss{synchronization paradigm} becomes widely available. In the
meantime, applications that require the potential performance benefits
of \emph{consume} semantics typically make careful use of platform-specific functionality instead.{\cprotect\footnote{Since C++17, the
use of \lstinline!memory_order_consume! has been explicitly
  discouraged after the acceptance of \cite{boehm16}. The specific
  note in the standard now says, ``Prefer
  \lstinline!memory_order_acquire!, which provides stronger guarantees
  than \lstinline!memory_order_consume!. Implementations have found it infeasible to provide performance better
  than that of \lstinline!memory_order_acquire!. Specification revisions
  are under consideration'' (\cite{iso17}, section~32.4 ``Order and Consistency," paragraph~1.3, Note~1, p.~1346). }}

\subsection[Annoyances]{Annoyances}\label{annoyances}

\hspace*{\fill}\\


%\subsubsection[Ill formed when inconsistently applied]{Ill formed when inconsistently applied}\label{ill-formed-when-inconsistently-applied}
%
%Like many aspects of a declaration, such as the
%\lstinline![[noreturn]]! attribute (see\linebreak%%%%%
%\featureref{\locationa}{the-noreturn-attribute}),
%\lstinline!alignas! specifier (see \featureref{\locationa}{alignas}), language linkage, and so on,
%the \lstinline![[carries_dependency]]! attribute must be applied to a
%function's declaration consistently across all translation units.
%Failing to apply it on the first declaration in a translation unit and
%then later to a (re)declaration is ill formed. Such uniqueness issues
%are readily dispatched when (1) each function's owner supplies a
%corresponding header having the canonical declarations for that function
%and (2) every client includes that corresponding header rather than
%attempting to redefine the function locally.

\newpage%%%%%%
\subsection[See Also]{See Also}\label{see-also}

\begin{itemize}
\item{%\intraref{Attribute Syntax}
\seealsoref{attributes}{\seealsolocationa}provides an in-depth discussion of how attributes pertain to C++ language entities.}
\item{%\intraref{\texttt{[[noreturn]]} Attribute}
\seealsoref{the-noreturn-attribute}{\seealsolocationa}offers an example of another \emph{attribute} that \emph{is} implemented ubiquitously.}
\end{itemize}

\subsection[Further Reading]{Further Reading}\label{further-reading}

\begin{itemize}
\item{\cite{marton17}}
\item{\cite{marton18}}
\end{itemize}




\newpage
%\section[{\tt final}]{Preventing Overriding and Derivation}\label{final}



\emcppsFeature{
    short={\lstinline!final!},
    tocshort={\TOCCode final},
    long={Preventing Overriding and Derivation},
    rhshort={\RHCode final},
}{final}
\setcounter{table}{0}
\setcounter{footnote}{0}
\setcounter{lstlisting}{0}
%\section[{\tt final}]{Preventing Overriding and Derivation}\label{final}


placeholder


\newpage
%\section[{\tt friend} '11]{Extended {\SecCode friend} Declarations}\label{extended-friend-declarations}
% 12 Feb 2021, revisions in; ready for Josh's code fixes
% 16 Feb 2021 JMB; code compiles
% 24 Feb 2021, copyedits in and proofed

\emcppsFeature{
    short={\lstinline!friend! '11},
    tocshort={{\TOCCode friend} '11},
    long={Extended {\SecCode friend} Declarations},
    toclong={Extended \lstinline!friend! Declarations},
    rhshort={{\RHCode friend} '11},
}{extended-friend-declarations}
\setcounter{table}{0}
\setcounter{footnote}{0}
\setcounter{lstlisting}{0}
%\section[{\tt friend} '11]{Extended {\SecCode friend} Declarations}\label{extended-friend-declarations}


Extended \lstinline!friend! declarations enable a class's author to
designate a type alias, a template parameter, or any other previously
declared type as a \lstinline!friend! of that class.

\subsection[Description]{Description}\label{description-extendedfriend}

A \lstinline!friend! declaration located within a given
\emcppsgloss[user defined type (UDT)]{user-defined type (UDT)} grants a designated type (or
\emph{free} function) access to private and protected members of that
class. Because the extended \lstinline!friend! syntax does not affect
\emph{function} friendships, this feature section addresses extended
friendship only between \emph{types}.

Prior to C++11, the Standard required an \emph{elaborated type
specifier} to be provided after the \lstinline!friend! keyword to designate
some other \emph{class} as being a \lstinline!friend! of a given type. An
elaborated type specifier for a class is a syntactical element having the form\linebreak[4]
\mbox{\lstinline!<class|struct|union>!~\lstinline!<identifier>!}. Elaborated
type specifiers can be used to refer to a previously declared entity or
to declare a new one, with the restriction that such an entity is
one of \lstinline!class!, \lstinline!struct!, or \lstinline!union!:

\begin{emcppslisting}
// C++03

struct S;
class C;
enum E { };

struct X0
{
    friend S;         // Error, not legal C++98/03
    friend struct S;  // OK, refers to (ù{\codeincomments{S}}ù) above
    friend class S;   // OK, refers to (ù{\codeincomments{S}}ù) above (might warn)
    friend class C;   // OK, refers to (ù{\codeincomments{C}}ù) above
    friend class C0;  // OK, declares (ù{\codeincomments{C0}}ù) in (ù{\codeincomments{X0}}ù)'s namespace
    friend union U0;  // OK, declares (ù{\codeincomments{U0}}ù) in (ù{\codeincomments{X0}}ù)'s namespace
    friend enum E;    // Error, (ù{\codeincomments{enum}}ù) cannot be a friend.
    friend enum E2;   // Error, (ù{\codeincomments{enum}}ù) cannot be forward-declared.
};
\end{emcppslisting}
    
\noindent This restriction prevents other potentially useful entities, e.g., type
aliases and template parameters, from being designated as friends:

\begin{emcppslisting}
// C++03

struct S;
typedef S SAlias;

struct X1
{
    friend struct SAlias;  // Error, using typedef-name after (ù{\codeincomments{struct}}ù)
};

template <typename T>
struct X2
{
    friend class T;        // Error, using template type parameter after (ù{\codeincomments{class}}ù)
};
\end{emcppslisting}
    
\noindent Furthermore, even though an entity belonging to a namespace other than
the class containing a \lstinline!friend! declaration might be visible,
explicit qualification is required to avoid unintentionally declaring a new type:

\begin{emcppslisting}
// C++03

struct S;  // This (ù{\codeincomments{S}}ù) resides in the global namespace.

namespace ns
{
    class X3
    {
        friend struct S;
            // OK, but declares a new (ù{\codeincomments{ns::S}}ù) instead of referring to (ù{\codeincomments{::S}}ù)
    };
}
\end{emcppslisting}
    
\noindent C++11 relaxes the aforementioned \emph{elaborated type specifier}
requirement and extends the classic \lstinline!friend! syntax by instead
allowing either a \emph{simple type specifier}, which is any unqualified
type or type alias, or a \emph{typename specifier}, e.g., the name of a
template \emph{type} parameter or dependent type thereof:

\begin{emcppslisting}
struct S;
typedef S SAlias;

namespace ns
{
    template <typename T>
    struct X4
    {
        friend T;           // OK
        friend S;           // OK, refers to (ù{\codeincomments{::S}}ù)
        friend SAlias;      // OK, refers to (ù{\codeincomments{::S}}ù)
        friend decltype(0); // OK, equivalent to (ù{\codeincomments{friend int;}}ù)
        friend C;           // Error, (ù{\codeincomments{C}}ù) does not name a type.
    };
}
\end{emcppslisting}
    
\noindent Notice that now it is again possible to declare as a \lstinline!friend! a
type that is expected to have already been declared, e.g., \lstinline!S!,
without having to worry that a typo in the spelling of the type would
silently introduce a new type declaration, e.g., \lstinline!C!, in the
enclosing scope.

Finally, consider the hypothetical case in which a class template,
\lstinline!C!, befriends a \emph{dependent} (e.g., nested) type, \lstinline!N!,
of its type parameter, \lstinline!T!:

\begin{emcppslisting}[emcppsbatch=e6]
template <typename T>
class C
{
    friend typename T::N;       // (ù{\codeincomments{N}}ù) is a *dependent* *type* of parameter (ù{\codeincomments{T}}ù).
    enum { e_SECRET = 10022 };  // This information is (ù{\codeincomments{private}}ù) to class (ù{\codeincomments{C}}ù).
};

struct S
{
    struct N
    {
        static constexpr int f()  // (ù{\codeincomments{f}}ù) is eligible for compile-time computation.
        {
            return C<S>::e_SECRET;  // Type (ù{\codeincomments{S::N}}ù) is a (ù{\codeincomments{friend}}ù) of (ù{\codeincomments{C<S>}}ù).
        }
    };
};

static_assert(S::N::f() == 10022, "");  // (ù{\codeincomments{N}}ù) has (ù{\codeincomments{private}}ù) access to (ù{\codeincomments{C<S>}}ù).
\end{emcppslisting}
    
\noindent In the example above, the nested type \lstinline!S::N! --- but not
\lstinline!S! itself --- has private access to \lstinline!C<S>::e_SECRET!.{\cprotect\footnote{Note that the need for \lstinline!typename! in the \lstinline!friend!
declaration in the example above to introduce the dependent type \lstinline!N! is relaxed
in C++20. For information on other contexts in which
\lstinline!typename! will eventually no longer be required, see~\cite{meredith20}.}}

\subsection[Use Cases]{Use Cases}\label{use-cases}

\subsubsection[Safely declaring a previously declared type to be a friend]{Safely declaring a previously declared type to be a friend}\label{safely-declaring-a-previously-declared-type-to-be-a-friend}

In C++98/03, to befriend a type that was already declared required
\emph{redeclaring} it. If the type were misspelled in the friend
declaration, a new type would be declared:

\begin{emcppslisting}
class Container { /* ... */ };

class ContainerIterator
{
    friend class Contianer;  // Compiles but wrong: (ù{\codeincomments{ia}}ù) should have been (ù{\codeincomments{ai}}ù).
    // ...
};
\end{emcppslisting}
    
\noindent The code above will compile and appear to be correct until
\lstinline!ContainerIterator! attempts to access a \lstinline!private! or
\lstinline!protected! member of \lstinline!Container!. At that point, the
compiler will surprisingly produce an error. As of C++11, we have the
option of preventing this mistake by using extended \lstinline!friend!
declarations:

\begin{emcppslisting}
class Container { /* ... */ };

class ContainerIterator
{
    friend Contianer;  // Error, (ù{\codeincomments{Contianer}}ù) not found
    // ...
};
\end{emcppslisting}
    

\subsubsection[Befriending a type alias used as a customization point]{Befriending a type alias used as a customization point}\label{befriending-a-type-alias-used-as-a-customization-point}

In C++03, the only option for friendship was to specify a particular
\lstinline!class! or \lstinline!struct! when granting private access. Let's
begin by considering a scenario in which we have an
\emcppsgloss{in-process} \emcppsgloss[value semantic type (VST)]{value-semantic type (VST)}
that serves as a \emph{handle} to a platform-specific object, such as a
\lstinline!Window! in a graphical application. (When used to qualify a VST, the
term \emcppsgloss{in-process}, also called \emph{in-core}, refers to a type
that has typical value-type--like operations but does not refer to a
value that is meaningful outside of the current process.\footnote{See
  \cite{lakos2a}, section~4.2.}) Large parts of a codebase
may seek to interact with \lstinline!Window! objects without needing or
obtaining access to the internal representation.

A very small part of the codebase that handles platform-specific window
management, however, needs privileged access to the internal
representation of \lstinline!Window!. One way to achieve this goal is to
make the platform-specific \lstinline!WindowManager! a \lstinline!friend! of
the \lstinline!Window! class; however, see \intraref{potential-pitfalls-extendedfriend}{long-distance-friendship}. 
%\textit{\titleref{potential-pitfalls-extendedfriend}: \titleref{long-distance-friendship}} on page~\pageref{long-distance-friendship}.

\begin{emcppslisting}
class WindowManager;  // forward declaration enabling extended (ù{\codeincomments{friend}}ù) syntax

class Window
{
private:
    friend class WindowManager;  // could instead use (ù{\codeincomments{friend WindowManager;}}ù)
    int d_nativeHandle;          // in-process (only) value of this object

public:
    // ... all the typical (e.g., special) functions we expect of a value type
};
\end{emcppslisting}
    
\noindent In the example above, class \lstinline!Window! befriends class
\lstinline!WindowManager!, granting it private access. Provided that the
implementation of \lstinline!WindowManager! resides in the same physical
\emcppsgloss[components]{component} as that of class \lstinline!Window!, no
\emcppsgloss{long-distance friendship} results. The consequence of such a
monolithic design would be that every client that makes use of the
otherwise lightweight \lstinline!Window! class would necessarily depend
physically on the presumably heavier-weight \lstinline!WindowManager! class.

Now consider that the \lstinline!WindowManager! implementations on
different platforms might begin to diverge significantly. To keep the
respective implementations maintainable, one might choose to factor them
into distinct C++ types, perhaps even defined in separate files, and to
use a \emph{type alias} determined using platform-detection preprocessor
macros to configure that alias:

%%%%%%%%%%%%%%%%%%%%% all separate. intended? 
%%% yes, and VR requests some separation

\begin{emcppslisting}[emcppsbatch=e1]
// windowmanager_win32.h:

#ifdef WIN32
class Win32WindowManager { /* ... */ };
#endif
\end{emcppslisting}

\vspace*{2ex}    

\begin{emcppshiddenlisting}[emcppsbatch=e1]
// windowmanager_unix.h:
#define UNIX
\end{emcppshiddenlisting}
\begin{emcppslisting}[emcppsbatch=e1]
// windowmanager_unix.h:

#ifdef UNIX
class UnixWindowManager { /* ... */ };
#endif
\end{emcppslisting}
    
\vspace*{2ex}

\begin{emcppslisting}[emcppsbatch=e1]
// windowmanager.h:

#ifdef WIN32
#include <windowmanager_win32.h>
typedef Win32WindowManager WindowManager;
#else
#include <windowmanager_unix.h>
typedef UnixWindowManager WindowManager;
#endif
\end{emcppslisting}
    
\vspace*{2ex}

\begin{emcppslisting}[emcppsbatch=e1]
// window.h:
#include <windowmanager.h>

class Window
{
private:
    friend WindowManager;  // C++11 extended friend declaration
    int d_nativeHandle;

public:
    // ...
};
\end{emcppslisting}
   
%%%%%%%%%%%%%%%%%%%%%%%      

In this example, class \lstinline!Window! no longer befriends a specific
class named \lstinline!WindowManager!; instead, it befriends the
\lstinline!WindowManager! type alias, which in turn has been set to the
correct platform-specific window manager implementation. Such extended
use of \lstinline!friend! syntax was not available in C++03.

Note that this use case involves \emcppsgloss{long-distance friendship}
inducing an implicit cyclic dependency between the \emcppsgloss[components]{component}
implementing \lstinline!Window! and those implementing
\lstinline!WindowManager!; see \intraref{potential-pitfalls-extendedfriend}{long-distance-friendship}. 
%\textit{\titleref{potential-pitfalls-extendedfriend}: \titleref{long-distance-friendship}} on page~\pageref{long-distance-friendship}. 
Such designs, though
undesirable, can result from an emergent need to add new platforms while
keeping tightly related code sequestered within smaller, more manageable
physical units. An alternative design would be to obviate the
\emcppsgloss{long-distance friendship} by widening the API for the
\lstinline!Window! class, the natural consequence of which would be to
invite public client abuse vis-a-vis \emcppsgloss{Hyrum's law}.

\subsubsection[Using the \lstinline!PassKey! idiom to enforce initialization]{Using the {\SubsubsecCode PassKey} idiom to enforce initialization}\label{using-the-passkey-idiom-to-enforce-initialization}

Prior to C++11,
efforts to grant private access to a class defined in a separate
physical unit required declaring the higher-level type itself to be a
\lstinline!friend!, resulting in this highly undesirable form of
friendship; see \intraref{potential-pitfalls-extendedfriend}{long-distance-friendship}.  
%\textit{\titleref{potential-pitfalls-extendedfriend}: \titleref{long-distance-friendship}} on page~\pageref{long-distance-friendship}. 
The ability in C++11 to
declare a template \emph{type} parameter or any other type specifier to
be a friend affords new opportunities to enforce selective
private access (e.g., to one or more individual functions) without
explicitly declaring another type to be a \lstinline!friend!; see also \intrarefsimple{granting-a-specific-type-access-to-a-single-private-function}. 
%\textit{\titleref{description-extendedfriend}: \titleref{granting-a-specific-type-access-to-a-single-private-function}} on page~\pageref{granting-a-specific-type-access-to-a-single-private-function}. 
In this use case, however, our use of extended
\lstinline!friend! syntax to befriend a template parameter is unlikely to run
afoul of sound physical design.

Let's say we have a commercial library, and we want it to verify a
software-license key in the form of a C-style string, prior to allowing use of 
other parts of the API:


\begin{emcppshiddenlisting}[emcppsbatch=e2]
class LibPassKey;
\end{emcppshiddenlisting}
\begin{emcppslisting}[emcppsbatch=e2]
// simplified pseudocode
LibPassKey initializeLibrary(const char* licenseKey);
int utilityFunction1(LibPassKey object /*, ... (other parameters) */);
int utilityFunction2(LibPassKey object /*, ... (other parameters) */);
\end{emcppslisting}
    
\noindent Knowing full well that this is not a \emph{secure} approach and that
innumerable deliberate, malicious ways exist to get around the C++ type
system, we nonetheless want to create a plausible regime where no
\emph{well-formed} code can \emph{accidentally} gain access
to library functionality other than by legitimately initializing the
system using a valid license key. We could easily cause a
function to \lstinline!throw!, \lstinline!abort!, and so on, at run time when
the function is called prior to the client's license key being
authenticated. However, part of our goal, as a friendly library vendor, is to
ensure that clients do not \emph{inadvertently} call other library
functions prior to initialization. To that end, we propose the following protocol:
\begin{enumerate}
\item{use an instantiation of the \lstinline!PassKey! class template\cprotect\footnote{\cite{mayrand15}} that only our API \emph{utility}\linebreak[4] \emph{\lstinline!struct!}\cprotect\footnote{\cite{lakos20}, section~2.4.9, pp.~312--321, specifically Figure~2-23, p.~316} can create}
\item{return a constructed object of this type only upon successful validation of the license key}
\item{require that clients present this (constructed) passkey \emph{object} every time they invoke any other function in the API}
\end{enumerate}
Here's an example that encompasses all three aforementioned points:

\begin{emcppslisting}
template <typename T>
class PassKey  // reusable standard utility type
{
    PassKey() { }  // private default constructor (no aggregate initialization)
    friend T;      // Only (ù{\codeincomments{T}}ù) is allowed to create this object.
};

struct BestExpensiveLibraryUtil
{
    class LicenseError { /*...*/ };  // thrown if license string is invalid

    using LibPassKey = PassKey<BestExpensiveLibraryUtil>;
        // This is the type of the (ù{\codeincomments{PassKey}}ù) that will be returned when this
        // utility is initialized successfully, but only this utility is able 
        // to construct an object of this type. Without a valid license string,
        // the client will have no way to create such an object and thus no way
        // to call functions within this library.

    static LibPassKey initializeLibrary(const char* licenseKey)
        // This function must be called with a valid (ù{\codeincomments{licenseKey}}ù) string prior
        // to using this library; if the supplied license is valid, a
        // (ù{\codeincomments{LibPassKey}}ù) *object* will be returned for mandatory use in *all*
        // subsequent calls to useful functions of this library. This function
        // throws (ù{\codeincomments{LicenseError}}ù) if the supplied (ù{\codeincomments{licenseKey}}ù) string is invalid.
    {
        if (isValid(licenseKey))
        {
            // Initialize library properly.

            return LibPassKey();
                // Return a default-constructed (ù{\codeincomments{LibPassKey}}ù). Note that only
                // this utility is able to construct such a key.
        }

        throw LicenseError();  // supplied license string was invalid
    }

    static int doUsefulStuff(LibPassKey key /*,...*/);
        // The function requires a (ù{\codeincomments{LibPassKey}}ù) object, which can be constructed
        // only by invoking the (ù{\codeincomments{static initializeLibrary}}ù) function, to be
        // supplied as its first argument. ...

private:
    static bool isValid(const char* key);
        // externally defined function that returns true if (ù{\codeincomments{key}}ù) is valid
};
\end{emcppslisting}
    
\noindent Other than going outside the language with invalid constructs or circumventing the
type system with esoteric tricks, this approach, among other things,
prevents invoking the\linebreak[4]%%%%%%%%  
\lstinline!doUsefulStuff! function without a proper
license. What's more, the C++ type system \emph{at compile time} forces
a prospective client to have initialized the library before any attempt
is made to use any of its other functionality.

\subsubsection[Granting a specific type access to a single \lstinline!private! function]{Granting a specific type access to a single {\SubsubsecCode private} function}\label{granting-a-specific-type-access-to-a-single-private-function}

When designing in purely logical terms, wanting to grant some other
logical entity special access to a type that no other entity enjoys is a
common situation. Doing so does not necessarily become problematic until
that friendship spans physical boundaries; see \intraref{potential-pitfalls-extendedfriend}{long-distance-friendship}. 
%\textit{\titleref{potential-pitfalls-extendedfriend}: \titleref{long-distance-friendship}} on page~\pageref{long-distance-friendship}.

As a simple approximation to a real-world use
case,\footnote{For an example of a real-world database
implementation that requires managed objects to befriend that database
manager, see \cite{csb15}, section~2.1.} suppose we have a lightweight object-database class, \lstinline!Odb!, that is designed to operate
collaboratively with objects, such as \lstinline!MyWidget!, that are
themselves designed to work collaboratively with \lstinline!Odb!. Every compliant UDT
suitable for management by \lstinline!Odb! will need to maintain an integer
object ID that is read/write accessible by an \lstinline!Odb! object. Under
no circumstances is any other object permitted to access, let alone
modify, that ID independently of the \lstinline!Odb! API.

Prior to C++11, the design of such a feature might require every
participating class to define a data member named \lstinline!d_objectId!
and to declare the \lstinline!Odb! class a \lstinline!friend! (using old-style
\lstinline!friend! syntax):

\begin{emcppslisting}
class MyWidget  // grants just (ù{\codeincomments{Odb}}ù) access to *all* of its private data
{
    int d_objectId;    // required by our collaborative-design strategy
    friend class Odb;  //    "     "   "        "        "       "
    // ...

public:
    // ...
};

class Odb
{
    // ...

public:
    template <typename T>
    void processObject(T& object)
        // This function template is generally callable by clients.
    {
        int& objId = object.d_objectId;
        // ... (process as needed)
    }

    // ...
};
\end{emcppslisting}
    
\noindent In this example, the \lstinline!Odb! class implements the public member
function template,\linebreak[4] \lstinline!processObject!, which then extracts the
\lstinline!objectId! field for access. The collateral damage is that we
have exposed all of our private details to \lstinline!Odb!, which is at
best a gratuitous widening of our sphere of encapsulation.

Using the \lstinline!PassKey! pattern allows us to be more selective with
what we share:

\begin{emcppslisting}[emcppsbatch=e3]
template <typename T>
class Passkey
    // Implement this eminently reusable (ù{\codeincomments{Passkey}}ù) class template again here.
{
    Passkey() { }  // prevent aggregate initialization
    friend T;      // Only the (ù{\codeincomments{T}}ù) in (ù{\codeincomments{PassKey<T>}}ù) can create a (ù{\codeincomments{PassKey}}ù) object.
    Passkey(const Passkey&) = delete;             // no copy/move construction
    Passkey& operator=(const Passkey&) = delete;  // no copy/move assignment
};
\end{emcppslisting}
    
\noindent We are now able to adjust the design of our systems such that only the
minimum private functionality is exposed to \lstinline!Odb!:

\begin{emcppslisting}[emcppsbatch=e3]
class Odb;      // Objects of this class have special access to other objects.

class MyWidget  // grants just (ù{\codeincomments{Odb}}ù) access to only its (ù{\codeincomments{objectId}}ù) member function
{
    int d_objectId;  // must have an (ù{\codeincomments{int}}ù) data member of any name we choose
    // ...

public:
    int& objectId(const Passkey<Odb>&) { return d_objectId; }
        // Return a non-(ù{\codeincomments{const}}ù) reference to the mandated (ù{\codeincomments{int}}ù) data member.
        // (ù{\codeincomments{objectId}}ù) is callable only within the scope of (ù{\codeincomments{Odb}}ù).

    // ...
};

class Odb
{
    // ...

public:
    template <typename T>
    void processObject(T& object)
        // This function template is generally callable by clients.
    {
        int& objId = object.objectId(Passkey<Odb>());
        // ...
    }

    // ...
};
\end{emcppslisting}
    
\noindent Instead of granting \lstinline!Odb! private access to \emph{all}
encapsulated implementation details of\linebreak[4] \mbox{\lstinline!MyWidget!}, this example
uses the \lstinline!PassKey! idiom to enable just \lstinline!Odb! to call the
(syntactically \lstinline!public!) \lstinline!objectId! member function of
\lstinline!MyWidget! with no private access whatsoever. As a further
demonstration of the efficacy of this approach, consider that we are
able to create and invoke the \lstinline!processObject! method of an
\lstinline!Odb! object from a function, \lstinline!f!, but we are blocked from
calling the \lstinline!objectId! method of a \lstinline!MyWidget! object
directly:

\begin{emcppslisting}[emcppsbatch=e3]
void f()
{
    Odb mgr;          // object receiving fine-grained privileged access
    MyWidget widget;  // object granting selective private access to just (ù{\codeincomments{Odb}}ù)
    mgr.processObject(widget);

    int& objId = widget.objectId(PassKey<Odb>());  // cannot call out of (ù{\codeincomments{Odb}}ù)
       // Error, (ù{\codeincomments{Passkey<T>::Passkey()}}ù) [with(ù{\codeincomments{T}}ù) = (ù{\codeincomments{Odb}}ù)] is private within
       // this context.
}
\end{emcppslisting}
    
\noindent Notice that use of the extended \lstinline!friend! syntax to befriend a
template parameter and thereby enable the \lstinline!PassKey! idiom here
improved the granularity with which we effectively grant privileged
access to an individually named type but didn't fundamentally alter the
testability issues that result when private access to specific C++
types is allowed to extend across physical boundaries; again, see \intraref{potential-pitfalls-extendedfriend}{long-distance-friendship}. 
%\textit{\titleref{potential-pitfalls-extendedfriend}: \titleref{long-distance-friendship}} on page~\pageref{long-distance-friendship}.

\subsubsection[Curiously recurring template pattern]{Curiously recurring template pattern}\label{curiously-recurring-template-pattern}

Befriending a template parameter via extended \lstinline!friend!
declarations can be helpful when implementing the \emcppsgloss[curiously recurring template pattern (CRTP)]{curiously
recurring template pattern (CRTP)}. For use-case examples and more
information on the pattern itself, see \intrarefsimple{appendix:-crtp-use-cases}. 
%\textit{\titleref{appendix:-crtp-use-cases}} on page~\pageref{appendix:-crtp-use-cases}.

\subsection[Potential Pitfalls]{Potential Pitfalls}\label{potential-pitfalls-extendedfriend}

\subsubsection[Long-distance friendship]{Long-distance friendship}\label{long-distance-friendship}

Since before C++ was standardized, granting private access via a
\lstinline!friend! declaration across physical boundaries, known as
\emcppsgloss{long-distance friendship}, was
observed\footnote{\cite{lakos96}, section~3.6.1, pp.~141--144}\footnote{\cite{lakos20}, section~2.6, pp.~342--370, specifically p.~367 and p.~362} to potentially lead
to designs that are qualitatively more difficult to understand, test,
and maintain. When a user-defined type, \lstinline!X!, befriends some other
specific type, \lstinline!Y!, in a separate, higher-level translation unit,
testing \lstinline!X! thoroughly without also testing \lstinline!Y! is no
longer possible. The effect is a test-induced cyclic dependency between
\lstinline!X! and \lstinline!Y!. Now imagine that \lstinline!Y! depends on a
sequence of other types, \lstinline!C1!, \lstinline!C2!, \ldots,
\lstinline!CN-2!, each defined in its own physical \emcppsgloss[components]{component},
\lstinline!CI!, where \lstinline!CN-2! depends on \lstinline!X!. The result is a
physical design cycle of size \emph{N}. As \emph{N} increases, the
ability to manage complexity quickly becomes intractable. Accordingly,
the two design imperatives that were most instrumental in shaping the
C++20 \emcppsgloss{modules} feature were (1) to have no cyclic module
dependencies and (2) to avoid intermodule friendships.

\subsection[See Also]{See Also}\label{see-also}

\begin{itemize}
\item{\seealsoref{alias-declarations-and-alias-templates}{\seealsolocationa}are among the beneficiaries of extended friend declarations.}
\end{itemize}

\subsection[Further Reading]{Further Reading}\label{further-reading}

\begin{itemize}
\item{For yet more potential uses of the extended friend pattern in
metaprogramming contexts, such as using CRTP, see
\cite{alexandrescu01}.}
\item{\cite{lakos96}, section~3.6, pp.~136--146, is dedicated to the
classic use (and misuse) of friendship.}
\item{\cite{miller05}}
\item{\cite{{lakos20}} provides extensive advice on \emph{sound}
\emcppsgloss{physical design}, which generally precludes \emcppsgloss{long-distance
friendship}.}
\end{itemize}

\subsection[Appendix: Curiously Recurring Template Pattern Use Cases]{Appendix: Curiously Recurring Template Pattern Use Cases}\label{appendix:-crtp-use-cases}

\subsubsection[Refactoring using the curiously recurring template pattern]{Refactoring using the curiously recurring template pattern}\label{refactoring-using-the-curiously-recurring-template-pattern}

Avoiding code duplication across disparate classes can sometimes be
achieved using a strange template pattern first recognized in the
mid-90s, which has since become known as the \emcppsgloss{curiously recurring
template pattern (CRTP)}. The pattern is \emph{curious} because it
involves the surprising step of declaring as a base class, such as
\lstinline!B!, a template that \emph{expects} the derived class, such as
\lstinline!C!, as a template argument, such as \lstinline!T!:

\begin{emcppslisting}
template <typename T>
class B
{
    // ...
};

class C : public B<C>
{
   // ...
};
\end{emcppslisting}
    
\noindent As a trivial illustration of how the CRTP can be used as a refactoring tool,
suppose that we have several classes for which we would like to track,
say, just the number of active instances:

\begin{emcppslisting}
class A
{
    static int s_count;  // declaration
    // ...

public:
    static int count() { return s_count; }

    A()          { ++s_count; }
    A(const A&)  { ++s_count; }
    A(const A&&) { ++s_count; }
    ~A()         { --s_count; }

    A& operator=(A&)  = default;  // see special members
    A& operator=(A&&) = default;  //  "     "       "
    // ...
};

int A::s_count;  // definition (in (ù{\codeincomments{.cpp}}ù) file)

class B { /* similar to A (above) */ };
// ...

void test()
{          // (ù{\codeincomments{A::s\_count}}ù) = (ù{\codeincomments{0}}ù), (ù{\codeincomments{B::s\_count}}ù) = (ù{\codeincomments{0}}ù)
    A a1;  // (ù{\codeincomments{A::s\_count}}ù) = (ù{\codeincomments{1}}ù), (ù{\codeincomments{B::s\_count}}ù) = (ù{\codeincomments{0}}ù)
    B b1;  // (ù{\codeincomments{A::s\_count}}ù) = (ù{\codeincomments{1}}ù), (ù{\codeincomments{B::s\_count}}ù) = (ù{\codeincomments{1}}ù)
    A a2;  // (ù{\codeincomments{A::s\_count}}ù) = (ù{\codeincomments{2}}ù), (ù{\codeincomments{B::s\_count}}ù) = (ù{\codeincomments{1}}ù)
}          // (ù{\codeincomments{A::s\_count}}ù) = (ù{\codeincomments{0}}ù), (ù{\codeincomments{B::s\_count}}ù) = (ù{\codeincomments{0}}ù)
\end{emcppslisting}
    
\noindent In this example, we have multiple classes, each repeating the same
common machinery. Let's now explore how we might refactor this example
using the CRTP:

\begin{emcppslisting}[emcppsbatch=e4]
template <typename T>
class InstanceCounter
{
protected:
    static int s_count;  // declaration

public:
    static int count() { return s_count; }
};

template <typename T>
int InstanceCounter<T>::s_count;  // definition (in same file as declaration)

struct A : InstanceCounter<A>
{
    A()          { ++s_count; }
    A(const A&)  { ++s_count; }
    A(const A&&) { ++s_count; }
    ~A()         { --s_count; }

    A& operator=(const A&)  = default;
    A& operator=(A&&)       = default;
    // ...
};
\end{emcppslisting}
    
\noindent Notice that we have factored out a common counting mechanism into an
\lstinline!InstanceCounter! class template and then derived our
representative class \lstinline!A! from \lstinline!InstanceCounter<A>!, and we
would do similarly for classes \lstinline!B!, \lstinline!C!, and so on. This
approach works because the compiler does not need to see the derived
type until the point at which the template is instantiated, which will
be \emph{after} it has seen the derived type.

Prior to C++11, however, there was plenty of room for user error.
Consider, for example, forgetting to change the base-type parameter when
copying and pasting a new type:

\begin{emcppslisting}[emcppsbatch=e4]
struct B : InstanceCounter<A>  // Oops! We forgot to change (ù{\codeincomments{A}}ù) to (ù{\codeincomments{B}}ù) in
                               // (ù{\codeincomments{InstanceCounter}}ù): The wrong count will be
                               // updated!
{
    B() { ++s_count; }
};
\end{emcppslisting}
    
\noindent Another problem is that a client deriving from our class can mess with
our protected \lstinline!s_count!:

\begin{emcppslisting}[emcppsbatch=e4]
struct AA : A
{
    AA() { s_count = -1; }  // Oops! *Hyrum's Law* is at work again!
};
\end{emcppslisting}
    
\noindent We could inherit from the \lstinline!InstanceCounter! class privately, but
then \lstinline!InstanceCounter! would have no way to add to the derived
class's public interface, for example, the public \lstinline!count! static
member function.

As it turns out, however, both of these missteps can be erased simply by
making the internal mechanism of the \lstinline!InstanceCounter! template
private and then having \mbox{\lstinline!InstanceCounter!} befriend its template
parameter, \lstinline!T!:

\begin{emcppslisting}
template <typename T>
class InstanceCounter
{
    static int s_count;  // Make this static data member private.
    friend T;            // Allow access only from the derived (ù{\codeincomments{T}}ù).

public:
    static int count() { return s_count; }
};

template <typename T>
int InstanceCounter<T>::s_count;
\end{emcppslisting}
    
\noindent Now if some other class does try to derive from this type, it cannot
access this type's counting mechanism. If we want to suppress even that
possibility, we can declare and default (see \featureref{\locationa}{defaulted-special-member-functions}) 
% ``\titleref{defaulted-special-member-functions}" on page~\pageref{defaulted-special-member-functions}) 
the \lstinline!InstanceCounter!
class constructors to be private as well.

\subsubsection[Synthesizing equality using the curiously recurring template pattern]{Synthesizing equality using the curiously recurring template pattern}\label{synthesizing-equality-using-crtp}

As a second example of code factoring using the CRTP, suppose that we want
to create a factored way of synthesizing \lstinline!operator==! for types
that implement just an \lstinline!operator<!.{\cprotect\footnote{This
example is based on a similar one found on stackoverflow.com:\linebreak[3]
  https://\linebreak[3]stackoverflow.\linebreak[3]com/\linebreak[3]questions/\linebreak[3]4173254/what-is-the-curiously-recurring-template-pattern-crtp}}
In this example, the CRTP base-class template, \lstinline!E!, will
synthesize the homogeneous \lstinline!operator==! for its parameter type,
\lstinline!D!, by returning \lstinline!false! if either argument is \emph{less
than} the other:

\begin{emcppslisting}[emcppsbatch=e5]
template <typename D>
class E { }; // CRTP base class used to synthesize (ù{\codeincomments{operator==}}ù) for (ù{\codeincomments{D}}ù)

template <typename D>
bool operator==(const E<D>& lhs, const E<D>& rhs)
{
    const D& d1 = static_cast<const D&>(lhs);  // derived type better be (ù{\codeincomments{D}}ù)
    const D& d2 = static_cast<const D&>(rhs);  //    "     "     "    "  "
    return !(d1 < d2) && !(d2 < d1);           // assuming (ù{\codeincomments{D}}ù) has an (ù{\codeincomments{operator<}}ù)
}
\end{emcppslisting}
    
\noindent A client that implements an \lstinline!operator<! can now reuse this CRTP
base case to synthesize an \lstinline!operator==!:

\begin{emcppshiddenlisting}[emcppsbatch={e5,e6}]
#include <cassert>  // standard C (ù{\codeincomments{assert}}ù) macro
\end{emcppshiddenlisting}

\begin{emcppslisting}[emcppsbatch=e5]
struct S : E<S>
{
    int d_size;
};

bool operator<(const S& lhs, const S& rhs)
{
    return lhs.d_size < rhs.d_size;
}

void test1()
{
    S s1; s1.d_size = 10;
    S s2; s2.d_size = 10;

    assert(s1 == s2);  // compiles and passes
}
\end{emcppslisting}
    
\noindent As this code snippet suggests, the base-class template, \lstinline!E!, is
able to use the template parameter, \lstinline!D! (representing the derived
class, \lstinline!S!), to synthesize the homogeneous free
\lstinline!operator==! function for \lstinline!S!.

Prior to C++11, no means existed to guard against accidents, such as
inheriting from the wrong base and then perhaps even forgetting to
define the \lstinline!operator<!:

\begin{emcppslisting}[emcppsbatch=e5]
struct P : E<S>  // Oops! should have been (ù{\codeincomments{E(P)}}ù) -- a serious latent defect
{
    int d_x;
    int d_y;
};

void test2()
{
    P p1; p1.d_x = 10; p1.d_y = 15;
    P p2; p2.d_x = 10; p2.d_y = 20;

    assert( !(p1 == p2) );  // Oops! This fails because of (ù{\codeincomments{E(S)}}ù) above.
}
\end{emcppslisting}
    
\noindent Again, thanks to C++11's extended \lstinline!friend! syntax, we can defend
against these defects at compile time simply by making the CRTP base
class's default constructor \emph{private} and befriending its template
parameter:

\begin{emcppslisting}
template <typename D>
class E
{
     E() = default;
     friend D;
};
\end{emcppslisting}
    
\noindent Note that the goal here is not security but simply guarding against
accidental typos, copy-paste errors, and other occasional  human errors. By making this change, we will soon realize that there
is no \lstinline!operator<! defined for \lstinline!P!.

\subsubsection[Compile-time polymorphism using the curiously recurring template pattern]{Compile-time polymorphism using the curiously recurring template pattern}\label{compile-time-polymorphism-using-crtp}

Object-oriented programming provides certain flexibility that at times
might be supererogatory. Here we will exploit the familiar domain of
abstract/concrete shapes to demonstrate a mapping between runtime
polymorphism using virtual functions and compile-time polymorphism using
the CRTP. We begin with a simple abstract \lstinline!Shape! class that
implements a single, pure, virtual \lstinline!draw! function:

\begin{emcppslisting}[emcppsbatch=e6]
class Shape
{
public:
    virtual void draw() const = 0;  // abstract (ù{\codeincomments{draw}}ù) function (interface)
};
\end{emcppslisting}
    
\noindent From this abstract \lstinline!Shape! class, we now derive two concrete
shape types, \lstinline!Circle! and \mbox{\lstinline!Rectangle!}, each implementing
the \emph{abstract} \lstinline!draw! function:

\begin{emcppslisting}[emcppsbatch=e6]
#include <iostream>  // (ù{\codeincomments{std::cout}}ù)

class Circle : public Shape
{
    int d_radius;

public:
    Circle(int radius) : d_radius(radius) { }

    void draw() const  // concrete implementation of abstract (ù{\codeincomments{draw}}ù) function
    {
        std::cout << "Circle(radius = " << d_radius << ")\n";
    }
};

class Rectangle : public Shape
{
    int d_length;
    int d_width;

public:
    Rectangle(int length, int width) : d_length(length), d_width(width) { }

    void draw() const  // concrete implementation of abstract (ù{\codeincomments{draw}}ù) function
    {
        std::cout << "Rectangle(length = " << d_length << ", "
                                "width = " << d_width  << ")\n";
    }
};
\end{emcppslisting}
    
\noindent Notice that a \lstinline!Circle! is constructed with a single integer
argument, i.e., \lstinline!radius!, and a \lstinline!Rectangle! is constructed
with two integers, i.e., \lstinline!length! and \lstinline!width!.

We now implement a function that takes an arbitrary shape, via a
\lstinline!const! \romeovalue{lvalue} reference to its abstract base class, and
prints it:

\begin{emcppslisting}[emcppsbatch=e6]
void print(const Shape& shape)
{
    shape.draw();
}

void testShape()
{
    print(Circle(1));        // OK, prints: (ù{\codeincomments{Circle(radius = 1)}}ù)
    print(Rectangle(2, 3));  // OK, prints: (ù{\codeincomments{Rectangle(length = 2, width = 3)}}ù)
    print(Shape());          // Error, (ù{\codeincomments{Shape}}ù) is an abstract class.
}
\end{emcppslisting}
    
\noindent Now suppose that we didn't need all the runtime flexibility offered by
this system and wanted to map just what we have in the previous code
snippet onto templates that avoid the spatial and runtime overhead of
virtual-function tables and dynamic dispatch. Such transformation again
involves creating a CRTP base class, this time in lieu of our abstract
interface:

\begin{emcppslisting}[emcppsbatch=e7]
template <typename T>
struct Shape
{
    void draw() const
    {
        static_cast<const T*>(this)->draw();  // assumes (ù{\codeincomments{T}}ù) derives from (ù{\codeincomments{Shape}}ù)
    }
};
\end{emcppslisting}
    
\noindent Notice that we are using a \lstinline!static_cast! to the address of an
object of the \lstinline!const! template parameter type, \lstinline!T!,
assuming that the template argument is of the same type as some derived
class of this object's type. We now define our types as before, the only
difference being the form of the base type:

\begin{emcppshiddenlisting}[emcppsbatch=e7]
#include <iostream>  // (ù{\codeincomments{std::cout}}ù)

class Circle : public Shape<Circle>
{
    int d_radius;

public:
    Circle(int radius) : d_radius(radius) { }

    void draw() const  // concrete implementation of abstract (ù{\codeincomments{draw}}ù) function
    {
        std::cout << "Circle(radius = " << d_radius << ")\n";
    }
};

class Rectangle : public Shape<Rectangle>
{
    int d_length;
    int d_width;

public:
    Rectangle(int length, int width) : d_length(length), d_width(width) { }

    void draw() const  // concrete implementation of abstract (ù{\codeincomments{draw}}ù) function
    {
        std::cout << "Rectangle(length = " << d_length << ", "
                                "width = " << d_width  << ")\n";
    }
};
\end{emcppshiddenlisting}
\begin{emcppslisting}[emcppsbatch=e7,emcppsignore={expanded as hidden listing}]
class Circle : public Shape<Circle>
{
    // same as above
};

class Rectangle : public Shape<Rectangle>
{
    // same as above
};
\end{emcppslisting}
    
\noindent We now define our \lstinline!print! function, this time as a function
template taking a \lstinline!Shape! of arbitrary type \lstinline!T!:

\begin{emcppslisting}[emcppsbatch=e7]
template <typename T>
void print(const Shape<T>& shape)
{
    shape.draw();
}
\end{emcppslisting}
    
\noindent The result of compiling and running \lstinline!testShape! above is the
same, including that \lstinline!Shape()! doesn't compile.

However, opportunities for undetected failure remain. Suppose we decide
to add a third shape, \lstinline!Triangle!, constructed with three sides:

\begin{emcppslisting}[emcppsbatch=e7]
class Triangle : public Shape<Rectangle>  // Oops!
{
    int d_side1;
    int d_side2;
    int d_side3;

public:
    Triangle(int side1, int side2, int side3)
        : d_side1(side1), d_side2(side2), d_side3(side3) { }

    void draw() const
    {
        std::cout << "Triangle(side1 = " << d_side1 << ", "
                              "side2 = " << d_side2 << ", "
                              "side3 = " << d_side3 << ")\n";
    }
};
\end{emcppslisting}
    
\noindent Unfortunately, we forgot to change the base-class type parameter when we
copy-pasted from \lstinline!Rectangle!.

Let's now create a new test that exercises all three and see what
happens on our platform:

\begin{emcppslisting}[emcppsbatch=e7]
void test2()
{
    print(Circle(1));          // prints: (ù{\codeincomments{Circle(radius = 1)}}ù)
    print(Rectangle(2, 3));    // prints: (ù{\codeincomments{Rectangle(length = 2, width = 3)}}ù)
    print(Triangle(4, 5, 6));  // prints: (ù{\codeincomments{Rectangle(length = 4, width = 5)}}ù) ?!
    Shape<int> bug;            // Compiles?!
}
\end{emcppslisting}
    
\noindent As should by now be clear, a defect in our \lstinline!Triangle!
implementation results in \emph{hard} \emcppsgloss{undefined behavior} that
could have been prevented at compile time by using the extended
\lstinline!friend! syntax. Had we defined the CRTP base-class template's
default constructor to be \emph{private} and made its type parameter a
\lstinline!friend!, we could have prevented the copy-paste error with
\lstinline!Triangle! and suppressed the ability to create a \lstinline!Shape!
object without deriving from it (e.g., see \lstinline!bug! in the previous
code snippet):

\begin{emcppslisting}[emcppsbatch=e8]
template <typename T>
class Shape
{
    Shape() = default;  // Default the (ù{\codeincomments{default}}ù) constructor to be (ù{\codeincomments{private}}ù).
    friend T;           // Ensure only a type derived from (ù{\codeincomments{T}}ù) has access.
};
\end{emcppslisting}
    
\noindent Generally, whenever we are using the CRTP, making just the
default constructor of the base-class template \lstinline!private! and
having it befriend its type parameter is typically a trivial local
change, is helpful in avoiding various forms of accidental misuse and
is unlikely to induce long-distance friendships where none previously
existed: Applying extended \lstinline!friend! syntax to an existing CRTP
is typically \emph{safe}.

\subsubsection[Compile-time visitor using the curiously recurring template pattern]{Compile-time visitor using the curiously recurring template pattern}\label{compile-time-visitor-using-crtp}

As more real-world applications of compile-time polymorphism using the CRTP,
consider implementing traversal and visitation of complex data
structures. In particular, we want to facilitate employing
\emph{default-action} functions, which allow for simpler code from the
point of view of the programmer who needs the results of the traversal.
We illustrate our compile-time visitation approach using binary trees as
our data structure.

We begin with the traditional node structure of a binary tree, where
each node has a left and right subtree plus a label:

\begin{emcppslisting}[emcppsbatch=e9]
struct Node
{
    Node* d_left;
    Node* d_right;
    char  d_label;  // (ù{\codeincomments{label}}ù) will be used in the pre-order example.

    Node() : d_left(0), d_right(0), d_label(0) { }
};
\end{emcppslisting}
    
\noindent Now we wish to have code that traverses the tree in one of the three
traditional ways: \emph{pre-order}, \emph{in-order}, \emph{post-order}.
Such traversal code is often intertwined with the actions to be taken.
In our implementation, however, we will write a CRTP-like base-class
template, \lstinline!Traverser!, that implements empty stub functions for
each of the three traversal types, relying on the CRTP-derived
type to supply the desired functionality:

\begin{emcppslisting}[emcppsbatch=e9]
template <typename T>
class Traverser
{
private:
    Traverser() = default;  // Make the default constructor (ù{\codeincomments{private}}ù).
    friend T;               // Grant access only to the derived class.

public:
    void visitPreOrder(Node*)  { }  // stub-functions & placeholders
    void visitInOrder(Node*)   { }  // (Each of these three functions
    void visitPostOrder(Node*) { }  // defaults to an inline "no-op.")

    void traverse(Node* n)  // factored subfunctionality
    {
        T *t = static_cast<T*>(this);  // Cast (ù{\codeincomments{this}}ù) to the derived type.

        if (n) { t->visitPreOrder(n);     }  // optionally defined in derived
        if (n) { t->traverse(n->d_left);  }  //     "         "    "     "
        if (n) { t->visitInOrder(n);      }  //     "         "    "     "
        if (n) { t->traverse(n->d_right); }  //     "         "    "     "
        if (n) { t->visitPostOrder(n);    }  //     "         "    "     "
    }
};
\end{emcppslisting}
    
\noindent The factored traversal mechanism is implemented in the
\lstinline!Traverser! base-class template. A proper subset of the four
customization points, that is, the four member functions invoked from
the \emph{public} \lstinline!traverse! function of the \lstinline!Traverser!
base class, is implemented as appropriate in the derived class,
identified by \lstinline!T!. Each of these customization functions is
invoked in order. Notice that the \lstinline!traverse! function is safe to call on a
\lstinline!nullptr! as each individual customization-function invocation
will be independently bypassed if its supplied \lstinline!Node! pointer is
null. If a customization function is defined in the derived class, that
version of it is invoked; otherwise, the corresponding empty
\lstinline!inline! base-class version of that function is invoked instead.
This approach allows for any of the three traversal orders to be
implemented simply by supplying an appropriately configured derived type
where clients are obliged to implement only the portions they need. Even
the traversal itself can be modified, as we will soon see, where we
create the very data structure we're traversing.

Let's now look at how derived-class authors might use this pattern.
First, we'll write a traversal class that fully populates a tree to a
specified depth:

\begin{emcppslisting}[emcppsbatch=e9]
struct FillToDepth : Traverser<FillToDepth>
{
    using Base = Traverser<FillToDepth>;  // similar to a local (ù{\codeincomments{typedef}}ù)

    int d_depth;            //  final "height" of the tree
    int d_currentDepth;     //  current distance from the root

    FillToDepth(int depth) : d_depth(depth), d_currentDepth(0) { }

    void traverse(Node*& n)
    {
        if (d_currentDepth++ < d_depth && !n)  // descend; if not balanced...
        {
            n = new Node;     // Add (ù{\codeincomments{node}}ù) since it's not already there.
        }

        Base::traverse(n);    // Recurse by invoking the *base* version.

        --d_currentDepth;     // Ascend.
    }
};
\end{emcppslisting}
    
\noindent The derived class's version of the \lstinline!traverse! member function
acts as if it overrides the \lstinline!traverse! function in the base-class
template and then, as part of its re-implementation, defers to the base-class version to perform the actual traversal.

Importantly, note that we have re-implemented \lstinline!traverse! in the
derived class with a function by the same name but having a
\emph{different} \emcppsgloss{signature} that has more capability (i.e., it's able
to modify its immediate argument) than the one in the base-class
template. In practice, this signature modification is something we would
do rarely, but part of the flexibility of this design pattern, as with
templates in general, is that we can take advantage of \emcppsgloss{duck
typing} to achieve useful functionality in somewhat unusual ways. For
this pattern, the designers of the base-class template and the designers
of the derived classes are, at least initially, likely to be the same
people, and they will arrange for these sorts of signature variants to
work correctly if they need such functionality. Or they may decide that
overridden methods should follow a proper contract and signature that
they determine is appropriate, and they may declare improper overrides
to be undefined behavior. In this example, we aim for illustrative
flexibility over rigor.

\begin{emcppslisting}[emcppsbatch=e9]
void traverse(Node* n);   // as declared in the (ù{\codeincomments{Traverser}}ù) base-class template
void traverse(Node*& n);  // as declared in the (ù{\codeincomments{FillToDepth}}ù) derived class
\end{emcppslisting}
    
\noindent Unlike virtual functions, the signatures of corresponding functions in
the base and derived classes need not match exactly \emph{provided} the
derived-class function can be called in the same way as the
corresponding one in the base class. In this case, the compiler has all
the information it needs to make the call properly:

\begin{emcppslisting}[emcppsignore={illustrative only}]
static_cast<FillToDepth *>(this)->traverse(n);  // what the compiler sees
\end{emcppslisting}
    
\noindent Suppose that we now want to create a type that labels a \emph{small}
tree, balanced or not, according to its pre-order traversal:

\begin{emcppslisting}[emcppsbatch=e9]
struct PreOrderLabel : Traverser<PreOrderLabel>
{
    char d_label;

    PreOrderLabel() : d_label('a') { }

    void visitPreOrder(Node* n)  // This choice controls traversal order.
    {
        n->d_label = d_label++;   
            // Each successive label is sequential alphabetically.
    }
};
\end{emcppslisting}
    
\noindent The simple pre-order traversal class, \lstinline!PreOrderLabel!, labels the
nodes such that it visits each parent \emph{before} it visits either of
its two children.

Alternatively, we might want to create a read-only derived class,
\lstinline!InOrderPrint!, that simply prints out the sequence of labels
resulting from an \emph{in-order} traversal of the, e.g., previously
pre-ordered, labels:

\begin{emcppslisting}[emcppsbatch=e9]
#include <cstdio>  // (ù{\codeincomments{std::putchar}}ù)

struct InOrderPrint : Traverser<InOrderPrint>
{
    ~InOrderPrint()
    {
        std::putchar('\n');  // Print single newline at the end of the string.
    }

    void visitInOrder(const Node* n) const
    {
        std::putchar(n->d_label);  // Print the label character exactly as is.
    }
};
\end{emcppslisting}
    
\noindent The simple \lstinline!InOrderPrint!-derived class, shown in the example above, prints out the
labels of a tree \emph{in order}: left subtree, then node, then right
subtree. Notice that since we are only examining the tree here --- not
modifying it --- we can declare the overriding method to take a
\lstinline!const!~\lstinline!Node*! rather than a \lstinline!Node*! and make the
method itself \lstinline!const!. Once again, compatibility of signatures,
not identity, is the key.

Finally, we might want to clean up the tree. We do so in
\emph{post-order} since we do not want to delete a node before we have
cleaned up its children!

\begin{emcppslisting}[emcppsbatch=e9]
struct CleanUp : Traverser<CleanUp>
{
    void visitPostOrder(Node*& n)
    {
        delete n;  // always necessary
        n = 0;     // might be omitted in a "raw" version of the type
    }
};
\end{emcppslisting}
    
\noindent Putting it all together, we can create a \lstinline!main! program that
creates a balanced tree to a depth of four and then labels it in
\emph{pre-order}, prints those labels in \emph{in-order}, and destroys
it in \emph{post-order}:

\begin{emcppslisting}[emcppsbatch=e9]
int main()
{
    Node* n = 0;                  // tree handle

    FillToDepth(4).traverse(n);   // (1) Create balanced tree.
    PreOrderLabel().traverse(n);  // (2) Label tree in pre-order.
    InOrderPrint().traverse(n);   // (3) Print labels in order.
    CleanUp().traverse(n);        // (4) Destroy tree in post-order.
    return 0;
}
\end{emcppslisting}
    
\noindent Running this program results in a binary tree of height 4, as
illustrated in the code snippet below, and has reliably consistent output:

\begin{lstlisting}[style=plain]
dcebgfhakjlinmo
\end{lstlisting}
    
\begin{lstlisting}[style=plain]
Level 0:                                a
                              .    '         '    .
Level 1:                b '                           ' i
                   .  '    '  .                   .  '    '  .
Level 2:        c                f              j              m
              /   \            /   \          /   \          /   \
Level 3:    d       e        g       h      k       l      n       o
\end{lstlisting}
    
\noindent This use of the CRTP for traversal truly shines when the data structure to
be traversed is especially complex, such as an abstract-syntax-tree
(AST) representation of a computer program, where tree nodes have many
different types, with each type having custom ways of representing the
subtrees it contains. For example, a translation unit is a sequence of
declarations; a declaration can be a type, a variable, or a function;
functions have return types, parameters, and a compound statement; the
statement has substatements, expressions, and so on. We would not want
to rewrite the traversal code for each new application. Given a reusable
CRTP-based traverser for our AST, we don't have to.

For example, consider writing a type that visits each integer literal
node in a given AST:

\begin{emcppshiddenlisting}[emcppsbatch=e10]
template <typename T>
class AstTraverser {};
class IntegerLiteral {};
\end{emcppshiddenlisting}
\begin{emcppslisting}[emcppsbatch=e10]
struct IntegerLiteralHandler : AstTraverser<IntegerLiteralHandler>
{
    void visit(IntegerLiteral* iLit)
    {
        // ... (do something with this integer literal)
    }
};
\end{emcppslisting}
    
\noindent The AST traverser, which would implement a separate empty \lstinline!visit!
overload for each syntactic node type in the grammar, would invoke our
derived \lstinline!visit! member function with every integer literal in the
program, regardless of where it appeared. This CRTP-based traverser
would also call many other \lstinline!visit! methods, but each of those would perform no action at all by default and would likely be elided at even modest compiler-optimization levels. Be aware, however, that although we ourselves are
not rewriting the traversal code each time, the compiler is still doing
it because every CRTP instantiation produces a new copy of the traversal
code. If the traversal code is large and complex, the consequence might
be increased program size, that is, \emcppsgloss{code bloat}.

Finally, the CRTP can be used in a variety of situations for many
purposes,\footnote{\cite{fluentcpp17}} which explains both its \emph{curiously recurring} nature and nomenclature. Those uses invariably
benefit from (1) declaring the base-class template's default constructor
\emph{private} and (2) having that template befriend its type parameter,
which is possible only by means of the extended \lstinline!friend! syntax.
Thus, the CRTP base-class template can ensure, at compile time, that its
type argument is actually derived from the base class as required by the
pattern.




\newpage
%\section[{\tt inline} {\tt namespace}]{Transparently Nested Namespaces}\label{inline-namespaces}%\sectionmark{{\RHCode inline} namespaces}
% 14 Feb 2021, revisions in; ready for Josh's code fixes
% 16 Feb 2021 JMB; code compiles

\emcppsFeature{
    short={\lstinline!inline!~\lstinline!namespace!},
    tocshort={{\TOCCode inline}~{\TOCCode namespace}},
    long={Transparently Nested Namespaces},
    rhshort={{\RHCode inline}~{\RHCode namespace}},
}{inline-namespaces}
\setcounter{table}{0}
\setcounter{footnote}{0}
\setcounter{lstlisting}{0}
%\section[{\tt inline} {\tt namespace}]{Transparently Nested Namespaces}\label{inline-namespaces}%


An \lstinline!inline! namespace is a nested namespace whose member entities
closely behave as if they were declared directly within the enclosing
namespace.

\subsection[Description]{Description}\label{description-inlinenamespace}

To a first approximation, an \emcppsgloss[inline namespace]{\lstinline!inline!~\lstinline!namespace!}
(e.g., \lstinline!v2! in the code snippet below) acts a lot like a
conventional nested namespace (e.g., \lstinline!v1!) followed by a
\lstinline!using! directive for that namespace in its enclosing namespace{\cprotect\footnote{C++17 allows developers to concisely declare nested
namespaces with shorthand notation:
\begin{emcppslisting}[style=footcode]
namespace a::b { /* ... */ }
// is the same as
namespace a { namespace b { /* ... */ } }
\end{emcppslisting}
C++20 expands on the above syntax by allowing the insertion of the
\lstinline!inline! keyword in front of any of the namespaces, except the
first one:
\begin{emcppslisting}[style=footcode]
namespace a::inline b::inline c { /* ... */ }
// is the same as
namespace a { inline namespace b { inline namespace c { /* ... */ } } }

inline namespace a::b { }  // Error, cannot start with (ù{\codeincomments{inline}}ù) for compound namespace names
namespace inline a::b { }  // Error, (ù{\codeincomments{inline}}ù) at front of sequence explicitly disallowed
\end{emcppslisting}
      }}:

\begin{emcppslisting}
// example.cpp:
namespace n
{
    namespace v1  // conventional nested namespace followed by (ù{\codeincomments{using}}ù) directive
    {
        struct T { };     // nested type declaration (identified as (ù{\codeincomments{::n::v1::T}}ù))
        int d;            // (ù{\codeincomments{::n::v1::d}}ù) at, e.g., (ù{\codeincomments{0x01a64e90}}ù)
    }

    using namespace v1;   // import names (ù{\codeincomments{T}}ù) and (ù{\codeincomments{d}}ù) into (ù{\codeincomments{namespace n}}ù)
}

namespace n
{
    inline namespace v2   // similar to being followed by (ù{\codeincomments{using namespace v2}}ù)
    {
        struct T { };     // nested type declaration (identified as (ù{\codeincomments{::n::v2::T}}ù))
        int d;            // (ù{\codeincomments{::n::v2::d}}ù) at, e.g., (ù{\codeincomments{0x01a64e94}}ù)
    }

    // using namespace v2;  // redundant when used with an (ù{\codeincomments{inline namespace}}ù)
}
\end{emcppslisting}
    
 
%%%%% because the long list of code terms in the FN below cause a persistent end-of-line error, those long lists are being set as a command:
%\newcommand{inlinelonglistone}{\lstinline!namespace!~\lstinline!n!~\lstinline!{!~\lstinline!inline!~\lstinline!namespace~\lstinline!v!~\lstinline!{!~\lstinline!int!~\lstinline!d;!~\lstinline!}!~\lstinline!}!}
%\newcommand{inlinelonglisttwo}{\lstinline!namespace!~\lstinline!n!~\lstinline!{!~\lstinline!namespace}!~\lstinline!v!~\lstinline!{!~\lstinline!int!~\lstinline!d;!~\lstinline!}!~\lstinline!using!~\lstinline!namespace!~\lstinline!v;!~\lstinline!}! }
\noindent Four subtle details distinguish these approaches:
\begin{enumerate}
\item{Name collisions with existing names behave differently due to differing name-lookup rules.}
\item{\emcppsgloss[argument-dependent lookup (ADL)]{Argument-dependent lookup (ADL)} gives special treatment to \lstinline!inline! namespaces.}
\item{Template specializations can refer to the primary template in an \lstinline!inline! namespace even if written in the enclosing namespace.}
\item{Reopening namespaces might reopen an \lstinline!inline! namespace.}
\end{enumerate}

One important aspect that all forms of namespaces share, however, is
that (1) nested symbolic names (e.g., \lstinline!n::v1::T!) at the
\emcppsgloss{API} level, (2) \emcppsgloss{mangled names} (e.g.,
\lstinline!_ZN1n2v11dE!, \lstinline!_ZN1n2v21dE!), and (3) assigned
relocatable addresses (e.g., \lstinline!0x01a64e90!, \lstinline!0x01a64e94!)
at the \emcppsgloss{ABI} level remain unaffected by the use of either
\lstinline!inline! or \lstinline!using! or both. To be precise, source files containing, alternately, \lstinline!namespace!~\lstinline!n!~\lstinline!{!~\lstinline!inline!~\lstinline!namespace!~\lstinline!v!~\lstinline!{!~\lstinline!int!~\lstinline!d;!~\lstinline!}!~\lstinline!}! 
and 
\lstinline!namespace!~\lstinline!n!~\lstinline!{!~\lstinline!namespace!~\lstinline!v!~\lstinline!{!~\lstinline!int!~\lstinline!d;!~\lstinline!}!~\lstinline!using!~\lstinline!namespace!~\lstinline!v;!~\lstinline!}!, will produce identical assembly.{\cprotect\footnote{These mangled names can be seen with GCC by running \lstinline!g++!~\lstinline!-S!~\lstinline!<file>.cpp! and viewing the contents of the generated \lstinline!<file>.s!. Note that Compiler Explorer is another valuable tool for learning about what comes out the other end of a C++ compiler: see https://godbolt.org/.}} Note that a \lstinline!using! directive immediately following an \lstinline!inline!
namespace is superfluous; name lookup will always consider names in
\lstinline!inline! namespaces before those imported by a \lstinline!using!
directive. Such a directive can, however, be used to import the contents
of an \lstinline!inline! namespace to some other namespace, albeit only in
the conventional, \emcppsgloss[using directive]{\lstinline!using! directive} sense; see \intraref{annoyances-inlinenamespace}{only-one-namespace-can-contain-any-given-inline-namespace}. 

More generally, each namespace has what is called its
\emph{\lstinline!inline! namespace set}, which is the transitive closure of
all \lstinline!inline! namespaces within the namespace. All names in the
\lstinline!inline! namespace set are roughly intended to behave as if they
are defined in the enclosing namespace. Conversely, each \lstinline!inline!
namespace has an \emph{enclosing namespace set} that comprises all
enclosing namespaces up to and including the first non-\lstinline!inline!
namespace.

\subsubsection[Loss of access to duplicate names in enclosing namespace]{Loss of access to duplicate names in enclosing namespace}\label{loss-of-access-to-duplicate-names-in-enclosing-namespace}

When both a type and a variable are declared with the same name in the
same scope, the variable name hides the type name --- such behavior can
be demonstrated by using the form of \lstinline!sizeof! that accepts a
nonparenthesized \emph{expression} (recall that the form of \lstinline!sizeof! that accepts a \emph{type} as its argument requires parentheses):

\begin{emcppslisting}
struct A { double d; };  static_assert(sizeof(  A) == 8, "");  // type
                      // static_assert(sizeof   A  == 8, "");  // Error

int A;                   static_assert(sizeof(  A) == 4, "");  // data
                         static_assert(sizeof   A  == 4, "");  // OK
\end{emcppslisting}
    
\noindent Unless both type and variable entities are declared within the same
scope, no preference is given to variable names; the name of an entity
in an inner scope hides a like-named entity in an enclosing scope:

\begin{emcppslisting}
void f()
{
    double B;                 static_assert(sizeof(B) == 8, "");  // variable
    {                         static_assert(sizeof(B) == 8, "");  // variable
        struct B { int d; };  static_assert(sizeof(B) == 4, "");  // type
    }                         static_assert(sizeof(B) == 8, "");  // variable
}
\end{emcppslisting}
    
\noindent When an entity is declared in an enclosing \lstinline!namespace! and
another entity having the same name hides it in a \emph{lexically}
nested scope, then (apart from \lstinline!inline! namespaces) access to a
hidden element can generally be recovered by using scope resolution:

\begin{emcppslisting}
struct C { double d; };  static_assert(sizeof(  C) == 8, "");

void g()
{                        static_assert(sizeof(  C) == 8, "");  // type
    int C;               static_assert(sizeof(  C) == 4, "");  // variable
                         static_assert(sizeof(::C) == 8, "");  // type
}                        static_assert(sizeof(  C) == 8, "");  // type
\end{emcppslisting}
    
\noindent A conventional nested namespace behaves as one might expect:

\begin{emcppslisting}
namespace outer
{
    struct D { double d; }; static_assert(sizeof(       D) == 8, ""); // type

    namespace inner
    {                       static_assert(sizeof(       D) == 8, ""); // type
        int D;              static_assert(sizeof(       D) == 4, ""); // var
    }                       static_assert(sizeof(       D) == 8, ""); // type
                            static_assert(sizeof(inner::D) == 4, ""); // var
                            static_assert(sizeof(outer::D) == 8, ""); // type
    using namespace inner;//static_assert(sizeof(       D) == 0, ""); // Error
                            static_assert(sizeof(inner::D) == 4, ""); // var
                            static_assert(sizeof(outer::D) == 8, ""); // type
}                           static_assert(sizeof(outer::D) == 8, ""); // type
\end{emcppslisting}
    
\noindent In the example above, the inner variable name, \lstinline!D!, hides the
outer type with the same name, starting from the point of \lstinline!D!'s
declaration in \lstinline!inner! until \lstinline!inner! is closed, after
which the unqualified name \lstinline!D! reverts to the type in the
\lstinline!outer! namespace. Then, right after the subsequent
\lstinline!using!~\lstinline!namespace!~\lstinline!inner;! directive, the meaning
of the unqualified name \lstinline!D! in \lstinline!outer! becomes ambiguous,
shown here with a \lstinline!static_assert! that is commented out; any
attempt to refer to an unqualified \lstinline!D! from here to the end of
the scope of \lstinline!outer! will fail to compile. The type entity
declared as \lstinline!D! in the \lstinline!outer! namespace can, however,
still be accessed --- from inside or outside of the \lstinline!outer!
namespace, as shown in the example --- via its qualified name,
\lstinline!outer::D!.

If an \lstinline!inline! namespace were used instead of a nested namespace
followed by a \mbox{\lstinline!using!} directive, however, the ability to recover
by name the hidden entity in the enclosing namespace is lost.
Unqualified name lookup considers the inline namespace set and the used
namespace set simultaneously. Qualified name lookup first considers the
\lstinline!inline! namespace set and \emph{then} goes on to look into used
namespaces. This means we can still refer to \lstinline!outer::D! in the
example above, but doing so would still be ambiguous if \lstinline!inner!
were an inline namespace. This subtle difference in behavior is a
byproduct of the highly specific use case that motivated this feature
and for which it was explicitly designed; see \intraref{use-cases-inlinenamespace}{link-safe-abi-versioning}. 
%\textit{\titleref{use-cases-inlinenamespace}: \titleref{link-safe-abi-versioning}} on page~\pageref{link-safe-abi-versioning}.

\subsubsection[Argument-dependent–lookup interoperability across \lstinline!inline! \lstinline!namespace! boundaries]{Argument-dependent–lookup interoperability across\\[0.5ex] {\SubsubsecCode inline} {\SubsubsecCode namespace} boundaries}\label{argument-dependent–lookup-interoperability-across-inline-namespace-boundaries}

Another important aspect of \lstinline!inline! namespaces is that they
allow \emcppsgloss[argument-dependent lookup (ADL)]{ADL} to work seamlessly across \lstinline!inline! namespace
boundaries. Whenever unqualified function names are being resolved, a
list of \emph{associated namespaces} is built for each argument of the
function. This list of associated namespaces comprises the namespace of the
argument, its enclosing namespace set, plus the \lstinline!inline!
namespace set.

Consider the case of a type, \lstinline!U!, defined in an \lstinline!outer!
namespace, and a function, \lstinline!f(U)!, declared in an \lstinline!inner!
namespace nested within \lstinline!outer!. A second type, \lstinline!V!, is
defined in the \lstinline!inner! namespace, and a function, \lstinline!g!, is
declared, after the close of \lstinline!inner!, in the \lstinline!outer!
namespace:

\begin{emcppslisting}
namespace outer
{
    struct U { };

    // inline               // Uncommenting this line fixes the problem.
    namespace inner
    {
         void f(U) { }
         struct V { };
    }

    using namespace inner;  // If we inline (ù{\codeincomments{inner}}ù), we don't need this line.

    void g(V) { }
}

void client()
{
    f(outer::U());         // Error, (ù{\codeincomments{f}}ù) is not declared in this scope.
    g(outer::inner::V());  // Error, (ù{\codeincomments{g}}ù) is not declared in this scope.
}
\end{emcppslisting}
    
\noindent In the example above, a \lstinline!client! invoking \lstinline!f! with an
object of type \lstinline!outer::U! fails to compile because
\lstinline!f(outer::U)! is declared in the nested \lstinline!inner! namespace,
which is not the same as declaring it in \lstinline!outer!. Because
\emcppsgloss[argument-dependent lookup (ADL)]{ADL} does not look into namespaces added with the \lstinline!using!
directive, \emcppsgloss[argument-dependent lookup (ADL)]{ADL} does not find the needed
\lstinline!outer::inner::f! function. Similarly, the type \lstinline!V!,
defined in namespace \lstinline!outer::inner!, is not declared in the same
namespace as the function \lstinline!g! that operates on it. Hence, when
\lstinline!g! is invoked from within \lstinline!client! on an object of type
\lstinline!outer::inner::V!, \emcppsgloss[argument-dependent lookup (ADL)]{ADL} again does not find the needed
function \lstinline!outer::g(outer::V)!.

Simply making the \lstinline!inner! namespace \lstinline!inline! solves both
of these \emcppsgloss[argument-dependent lookup (ADL)]{ADL}-related problems. All transitively nested
\lstinline!inline! namespaces --- up to and including the most proximate
non-\lstinline!inline! enclosing namespace --- are treated as one with
respect to \emcppsgloss[argument-dependent lookup (ADL)]{ADL}.

\subsubsection[The ability to specialize templates declared in a nested \lstinline!inline! namespace]{The ability to specialize templates declared in a nested {\SubsubsecCode inline} namespace}\label{the-ability-to-specialize-templates-declared-in-a-nested-inline-namespace}

The third property that distinguishes \lstinline!inline! namespaces from
conventional ones, even when followed by a \lstinline!using! directive, is
the ability to specialize a class template defined within an
\lstinline!inline! namespace from within an enclosing one; this ability
holds transitively up to and including the most proximate
non-\lstinline!inline! namespace:

\begin{emcppslisting}[emcppserrorlines={24}]
namespace out                      // proximate non-(ù{\codeincomments{inline}}ù) outer namespace
{
    inline namespace in1           // first-level nested (ù{\codeincomments{inline}}ù) namespace
    {
        inline namespace in2       // second-level nested (ù{\codeincomments{inline}}ù) namespace
        {
            template <typename T>  // primary class template general definition
            struct S { };

            template <>            // class template *full* specialization
            struct S<char> { };
        }

        template <>                // class template *full* specialization
        struct S<short> { };
    }

    template <>                    // class template *full* specialization
    struct S<int> { };
}

using namespace out;               // conventional using directive

template <>
struct S<int> { };                 // Error, cannot specialize from this scope
\end{emcppslisting}
    
\noindent Note that the conventional nested namespace \lstinline!out! followed by a
\lstinline!using! directive in the enclosing namespace does not admit
specialization from that outermost namespace, whereas all of the
\lstinline!inline! namespaces do. Function templates behave similarly
except that --- unlike class templates, whose definitions must reside
entirely within the namespace in which they are declared --- a function
template can be \emph{declared} within a nested namespace and then be
\emph{defined} from anywhere via a \emcppsgloss{qualified name}:

\begin{emcppslisting}
namespace out                      // proximate non-(ù{\codeincomments{inline}}ù) outer namespace
{
    inline namespace in1           // first-level nested (ù{\codeincomments{inline}}ù) namespace
    {
        template <typename T>      // function template declaration
        void f();

        template <>                // function template (full) specialization
        void f<short>() { }
    }

    template <>                    // function template (full) specialization
    void f<int>() { }
}

template <typename T>              // function template general definition
void out::in1::f() { }
\end{emcppslisting}
    
\noindent An important takeaway from the examples above is that every template
entity --- be it class or function --- \emph{must} be declared
in \emph{exactly} one place within the collection of namespaces that
comprise the \lstinline!inline! namespace set. In particular, declaring a
class template in a nested \lstinline!inline! namespace and then
subsequently defining it in a containing namespace is not possible
because, unlike a function definition, a type definition cannot be
placed into a namespace via name qualification alone:

\begin{emcppslisting}[emcppserrorlines={17}]
namespace outer
{
    inline namespace inner
    {
        template <typename T>      // class template declaration
        struct Z;                  // (if defined, must be within same namespace)

        template <>                // class template full specialization
        struct Z<float> { };
    }

    template <typename T>          // inconsistent declaration (and definition)
    struct Z { };                  // (ù{\codeincomments{Z}}ù) is now ambiguous in namespace (ù{\codeincomments{outer}}ù).

    const int i = sizeof(Z<int>);  // Error, Reference to (ù{\codeincomments{Z}}ù) is ambiguous.

    template <>                    // attempted class template full specialization
    struct Z<double> { };          // Error, (ù{\codeincomments{outer::Z}}ù) or (ù{\codeincomments{outer::inner::Z}}ù)?
}
\end{emcppslisting}
    

\subsubsection[Reopening namespaces can reopen nested \lstinline!inline! ones]{Reopening namespaces can reopen nested {\SubsubsecCode inline} ones}\label{reopening-namespaces-can-reopen-nested-inline-ones}

Another subtlety specific to \lstinline!inline! namespaces is related to
reopening namespaces. Consider a namespace \lstinline!outer! that declares
a nested namespace \lstinline!outer::m! and an \lstinline!inline!
namespace \lstinline!inner! that, in turn, declares a nested namespace
\lstinline!outer:inner::m!. In this case, subsequent attempts to reopen
namespace \lstinline!m! cause an ambiguity error:

\begin{emcppshiddenlisting}[emcppsbatch=e1]
#include <type_traits>  // (ù{\codeincomments{std::is\_same}}ù)
\end{emcppshiddenlisting}
\begin{emcppslisting}[emcppsbatch=e1,emcppserrorlines={17,18,19}]
namespace outer
{
    namespace m { }      // opens and closes (ù{\codeincomments{::outer::m}}ù)

    inline namespace inner
    {
        namespace n { }  // opens and closes (ù{\codeincomments{::outer::inner::n}}ù)
        namespace m { }  // opens and closes (ù{\codeincomments{::outer::inner::m}}ù)
    }

    namespace n          // OK, reopens (ù{\codeincomments{::outer::inner::n}}ù)
    {
        struct S { };    // defines (ù{\codeincomments{::outer::inner::n::S}}ù)
    }

    namespace m          // Error, (ù{\codeincomments{namespace m}}ù) is ambiguous.
    {
        struct T { };    // with clang defines (ù{\codeincomments{::outer::m::T}}ù)
    }
}

static_assert(std::is_same<outer::n::S, outer::inner::n::S>::value, "");
\end{emcppslisting}
    
\noindent In the code snippet above, no issue occurs with reopening
\lstinline!outer::inner::n! and no issue would have occurred with reopening
\lstinline!outer::m! but for the \lstinline!inner! namespaces having been
declared \lstinline!inline!. When a new namespace declaration is
encountered, a lookup determines if a matching namespace having that
name appears anywhere in the \emph{\lstinline!inline! namespace set} of the
current namespace. If the namespace is ambiguous, as is the case with
\lstinline!m! in the example above, one can get the surprising error shown.{\cprotect\footnote{Note that reopening already declared
namespaces, such as \lstinline!m! and \lstinline!n! in the \lstinline!inner!
and \lstinline!outer! example, is handled incorrectly on several
popular platforms. Clang, for example, will perform a name lookup when
encountering a new namespace declaration and give preference to the
outermost namespace found, causing the last declaration of \lstinline!m!
to reopen \lstinline!::outer::m! instead of being ambiguous. GCC, prior
to version~8.1, will not perform name lookup and will place \emph{any}
nested namespace declarations directly within their enclosing
namespace. This compiler defect causes the last declaration of
\lstinline!m! to reopen \lstinline!::outer::m! instead of
\lstinline!::outer::inner::m! and the last declaration of \lstinline!n! to
open a new namespace, \lstinline!::outer::n!, instead of reopening
  \lstinline!::outer::inner::n!.}} If a matching namespace is found
unambiguously inside an \lstinline!inline! namespace, \lstinline!n! in this
case, then it is that nested namespace that is reopened --- here,
\lstinline!::outer::inner::n!. The inner namespace is reopened even though
the last declaration of \lstinline!n! is not lexically scoped within
\lstinline!inner!. Notice that the definition of \lstinline!S! is perhaps
surprisingly defining \lstinline!::outer::inner::n::S!, not
\lstinline!::outer::n::S!. For more on what is \emph{not} supported by this
feature, see \intraref{annoyances-inlinenamespace}{inability-to-redeclare-across-namespaces-impedes-code-factoring}. 
%\textit{\titleref{annoyances-inlinenamespace}: \titleref{inability-to-redeclare-across-namespaces-impedes-code-factoring}} on page~\pageref{inability-to-redeclare-across-namespaces-impedes-code-factoring}.

\subsection[Use Cases]{Use Cases}\label{use-cases-inlinenamespace}

\subsubsection[Facilitating API migration]{Facilitating API migration}\label{facilitating-api-migration}

Getting a large codebase to \emph{promptly} upgrade to a new version of
a library in any sort of timely fashion can be challenging. As a
simplistic illustration, imagine that we have just developed a new
library, \lstinline!parselib!, comprising a class template, \lstinline!Parser!,
and a function template, \lstinline!analyze!, that takes a \lstinline!Parser!
object as its only argument:

\begin{emcppslisting}[emcppsbatch=e2]
namespace parselib
{
    template <typename T>
    class Parser
    {
        // ...

    public:
        Parser();
        int parse(T* result, const char* input);
            // Load (ù{\codeincomments{result}}ù) from null-terminated (ù{\codeincomments{input}}ù); return (ù{\codeincomments{0}}ù) (on
            // success) or nonzero (with no effect on (ù{\codeincomments{result}}ù)).
    };

    template <typename T>
    double analyze(const Parser<T>& parser);
}
\end{emcppslisting}
    
\noindent To use our library, clients will need to specialize our \lstinline!Parser!
class directly within the \lstinline!parselib! namespace:

\begin{emcppslisting}[emcppsbatch=e2]
struct MyClass { /*...*/ };  // end-user-defined type

namespace parselib  // necessary to specialize (ù{\codeincomments{Parser}}ù)
{
    template <>            // Create *full* specialization of class
    class Parser<MyClass>  // (ù{\codeincomments{Parser}}ù) for user-type (ù{\codeincomments{MyClass}}ù).
    {
        // ...

    public:
        Parser();
        int parse(MyClass* result, const char* input);
            // The *contract* for a specialization typically remains the same.
    };

    double analyze(const Parser<MyClass>& parser);
};
\end{emcppslisting}
    
\noindent Typical \lstinline!client! code will also look for the \lstinline!Parser!
class directly within the \lstinline!parselib! namespace:

\begin{emcppslisting}[emcppsbatch=e2]
void client()
{
    MyClass result;
    parselib::Parser<MyClass> parser;

    int status = parser.parse(&result, "...( MyClass value )...");
    if (status != 0)
    {
        return;
    }

    double value = analyze(parser);
    // ...
}
\end{emcppslisting}
    
\noindent Note that invoking \lstinline!analyze! on objects of some instantiated type
of the \lstinline!Parser! class template will rely on \emcppsgloss[argument-dependent lookup (ADL)]{ADL} to find
the corresponding overload.

We anticipate that our library's API will evolve over time, so we want to
enhance the design of \lstinline!parselib! accordingly. One of our goals is
to somehow encourage clients to move essentially all at once, yet also to
accommodate both the early adopters and the inevitable stragglers that
make up a typical adoption curve. Our approach will be to create, within
our outer \lstinline!parselib! namespace, a nested \lstinline!inline!
namespace, \lstinline!v1!, which will hold the current implementation of
our library software:

\begin{emcppslisting}[emcppsbatch=e3]
namespace parselib
{
    inline namespace v1             // Note our use of (ù{\codeincomments{inline}}ù) namespace here.
    {
        template <typename T>
        class Parser
        {
            // ...

        public:
            Parser();
            int parse(T* result, const char* input);
                // Load (ù{\codeincomments{result}}ù) from null-terminated (ù{\codeincomments{input}}ù); return 0 (on
                // success) or nonzero (with no effect on (ù{\codeincomments{result}}ù)).
        };

        template <typename T>
        double analyze(const Parser<T>& parser);
    }
}
\end{emcppslisting}
    
\noindent As suggested by the name \lstinline!v1!, this namespace serves primarily as
a mechanism to support library evolution through \emcppsgloss{API} and
\emcppsgloss{ABI} versioning (see \intraref{use-cases-inlinenamespace}{link-safe-abi-versioning} 
%\textit{\titleref{use-cases-inlinenamespace}: \titleref{link-safe-abi-versioning}} on page~\pageref{link-safe-abi-versioning} 
and \intraref{use-cases-inlinenamespace}{build-modes-and-abi-link-safety} ). 
%\textit{\titleref{use-cases-inlinenamespace}: \titleref{build-modes-and-abi-link-safety}} on page~\pageref{build-modes-and-abi-link-safety}). 
The need to specialize \lstinline!class!~\lstinline!Parser! and,
independently, the reliance on ADL to find the free function template
\lstinline!analyze! require the use of \lstinline!inline! namespaces, as
opposed to a conventional namespace followed by a \lstinline!using!
directive.

Note that, whenever a subsystem starts out directly in a first-level
namespace and is subsequently moved to a second-level nested namespace
for the purpose of versioning, declaring the inner namespace
\lstinline!inline! is the most reliable way to avoid inadvertently
destabilizing existing clients; see also \intraref{potential-pitfalls-inlinenamespace}{enabling-selective-using-directives-for-short-named-entities}.  
%\textit{\titleref{potential-pitfalls-inlinenamespace}: \titleref{enabling-selective-using-directives-for-short-named-entities}} on page~\pageref{enabling-selective-using-directives-for-short-named-entities}.

Now suppose we decide to enhance \lstinline!parselib! in a
non--backwards-compatible manner, such that the signature of
\lstinline!parse! takes a second argument \lstinline!size! of type
\lstinline!std::size_t! to allow parsing of non--null-terminated strings
and to reduce the risk of buffer overruns. Instead of unilaterally
removing all support for the previous version in the new release, we can
create a second namespace, \lstinline!v2!, containing the new
implementation and then, at some point, make \lstinline!v2! the
\lstinline!inline! namespace instead of \lstinline!v1!:

\begin{emcppslisting}
#include <cstddef>  // (ù{\codeincomments{std::size\_t}}ù)

namespace parselib
{
    namespace v1  // Notice that (ù{\codeincomments{v1}}ù) is now just a nested namespace.
    {
        template <typename T>
        class Parser
        {
            // ...

        public:
            Parser();
            int parse(T* result, const char* input);
                // Load (ù{\codeincomments{result}}ù) from null-terminated (ù{\codeincomments{input}}ù); return 0 (on
                // success) or nonzero (with no effect on (ù{\codeincomments{result}}ù)).
        };

        template <typename T>
        double analyze(const Parser<T>& parser);
    }

    inline namespace v2    // Notice that use of (ù{\codeincomments{inline}}ù) keyword has moved here.
    {
        template <typename T>
        class Parser
        {
            // ...

        public:  // note incompatible change to (ù{\codeincomments{Parser}}ù)'s essential API
            Parser();
            int parse(T* result, const char* input, std::size_t size);
                // Load (ù{\codeincomments{result}}ù) from (ù{\codeincomments{input}}ù) of specified (ù{\codeincomments{size}}ù); return 0
                // on success) or nonzero (with no effect on (ù{\codeincomments{result}}ù)).
        };

        template <typename T>
        double analyze(const Parser<T>& parser);
    }
}
\end{emcppslisting}
    
\noindent When we release this new version with \lstinline!v2! made \lstinline!inline!,
all existing clients that rely on the version supported directly in
\lstinline!parselib! will, by design, break when they recompile. At
that point, each client will have two options. The first one is
to upgrade the code immediately by passing in the size of the input
string (e.g., \lstinline!23!) along with the address of its first
character:

\begin{emcppslisting}[emcppserrorlines={4}]
void client()
{
    // ...
    int status = parser.parse(&result, "...( MyClass value )...", 23);
    // ...                                                      ^^^^ Look here!
}
\end{emcppslisting}
    
\noindent The second option is to change all references to
\lstinline!parselib! to refer to the original version in \lstinline!v1!
explicitly:

\begin{emcppshiddenlisting}[emcppsbatch=e3]
struct MyClass { /*...*/ };  // end-user-defined type
\end{emcppshiddenlisting}
\begin{emcppslisting}[emcppsbatch=e3]
namespace parselib
{
    namespace v1  // specializations moved to nested namespace
    {
        template <>
        class Parser<MyClass>
        {
            // ...

        public:
            Parser();
            int parse(MyClass* result, const char* input);
        };

        double analyze(const Parser<MyClass>& parser);
    }
};

void client1()
{
    MyClass result;
    parselib::v1::Parser<MyClass> parser;  // reference nested namespace (ù{\codeincomments{v1}}ù)

    int status = parser.parse(&result, "...( MyClass value )...");
    if (status != 0)
    {
        return;
    }

    double value = analyze(parser);
    // ...
}
\end{emcppslisting}
    
\noindent Providing the updated version in a new \lstinline!inline! namespace
\lstinline!v2! provides a more flexible migration path --- especially for a
large population of independent client programs --- compared to manual
targeted changes in client code.

Although new users would pick up the latest version automatically either
way, existing users of \lstinline!parselib! will have the option of
converting immediately by making a few small syntactic changes or
opting to remain with the original version for a while longer by making
all references to the library namespace refer explicitly to the desired
version. If the library is released before the \lstinline!inline! keyword
is moved, early adopters will have the option of opting in by referring
to \lstinline!v2! explicitly until it becomes the default. Those who have
no need for enhancements can achieve stability by referring to a
particular version in perpetuity or until it is physically removed from
the library source.

Although this same functionality can sometimes be realized without using \lstinline!inline! namespaces (i.e., by adding a
\lstinline!using!~\lstinline!namespace! directive at the end of the
\lstinline!parselib! namespace), any benefit of ADL and the ability to
specialize templates from within the enclosing \lstinline!parselib!
namespace itself would be lost. Note that, because
specialization doesn't kick in until overload resolution is completed,
specializing overloaded functions is dubious at best; see \intraref{potential-pitfalls-inlinenamespace}{specializing-templates-in-std-can-be-problematic}.
%\textit{\titleref{potential-pitfalls-inlinenamespace}: \titleref{specializing-templates-in-std-can-be-problematic}} on page~\pageref{specializing-templates-in-std-can-be-problematic}.}}

Providing separate namespaces for each successive version has an
additional advantage in an entirely separate dimension: avoiding inadvertent, difficult-to-diagnose, latent linkage defects. Though not
demonstrated by this specific example,
cases do arise where simply changing which of the version namespaces is
declared \lstinline!inline! might lead to an \emcppsgloss[ill formed, no diagnostic required (IFNDR)]{ill formed, no-diagnostic
required (IFNDR)} program. This might happen when one or more of
its translation units that use the library are not recompiled before the
program is relinked to the new static or dynamic library containing the
updated version of the library software; see \intraref{use-cases-inlinenamespace}{link-safe-abi-versioning}. 
%\textit{\titleref{use-cases-inlinenamespace}: \titleref{link-safe-abi-versioning}} on page~\pageref{link-safe-abi-versioning}.

For distinct
nested namespaces to guard effectively against accidental link-time
errors, the symbols involved have to (1) reside in object code (e.g.,
a \emcppsgloss{header-only library} would fail this requirement) and (2)
have the same \emcppsgloss{name mangling} (i.e., linker symbol) in both
versions. In this particular instance, however, the signature of the
\lstinline!parse! member function of \lstinline!parser! did change, and its
mangled name will consequently change as well; hence the same
  \lstinline!undefined!~\lstinline!symbol! link error would result either way.

\subsubsection[Link-safe ABI versioning]{Link-safe ABI versioning}\label{link-safe-abi-versioning}

\lstinline!inline! namespaces are not intended as a mechanism for
source-code versioning; instead, they prevent programs from being
\emcppsgloss[ill formed]{ill formed} due to linking some version of a library with client
code compiled using some other, typically older version of the same
library. Below, we present two examples: a simple pedagogical example
to illustrate the principle followed by a more real-world example.
Suppose we have a library component \lstinline!my_thing! that implements
an example type, \lstinline!Thing!, which wraps an \lstinline!int! and
initializes it with some value in its default constructor defined
out-of-line in the \lstinline!cpp! file:

\begin{emcppslisting}
struct Thing  // version 1 of class (ù{\codeincomments{Thing}}ù)
{
    int i;    // integer data member (size is 4)
    Thing();  // original non-(ù{\codeincomments{inline}}ù) constructor (defined in (ù{\codeincomments{.cpp}}ù) file)
};
\end{emcppslisting}
    
\noindent Compiling a source file with this version of the header included might produce an object file that can be incompatible yet linkable with an object file resulting from compiling a different source file with a different version of this header included:

\begin{emcppslisting}
struct Thing   // version 2 of class (ù{\codeincomments{Thing}}ù)
{
    double d;  // double-precision floating-point data member (size is 8)
    Thing();   // updated non-(ù{\codeincomments{inline}}ù) constructor (defined in (ù{\codeincomments{.cpp}}ù) file)
};
\end{emcppslisting}
    
\noindent To make the problem that we are illustrating concrete, let's represent
the client as a \lstinline!main! program that does nothing but create a
\lstinline!Thing! and print the value of its only data member, \lstinline!i!.

\begin{emcppslisting}[emcppsbatch=e4]
// main.cpp:
#include <my_thing.h>  // (ù{\codeincomments{my::Thing}}ù) (version 1)
#include <iostream>    // (ù{\codeincomments{std::cout}}ù)

int main()
{
     my::Thing t;
     std::cout << t.i << '\n';
}
\end{emcppslisting}
    
\noindent If we compile this program, a reference to a locally undefined linker
symbol, such as
\lstinline!_ZN2my7impl_v15ThingC1Ev!,{\cprotect\footnote{On a Unix
machine, typing \lstinline!nm!~\lstinline!main.o! reveals the symbols used
in the specified object file. A symbol prefaced with a capital
\lstinline!U! represents an undefined symbol that must be resolved by the
linker. Note that the linker symbol shown here incorporates an
intervening \lstinline!inline! namespace, \lstinline!impl_v1!, as will be
  explained shortly.}} which represents the \lstinline!my::Thing::Thing!
constructor, will be generated in the \lstinline!main.o! file:

\begin{lstlisting}[language=bash,style=plain]
$ g++ -c main.cpp
\end{lstlisting}
    
\noindent Without explicit intervention, the spelling of this linker symbol would
be unaffected by any subsequent changes made to the implementation of
\lstinline!my::Thing!, such as its data members or implementation of its
default constructor, even after recompiling. The same, of course,
applies to its definition in a separate translation unit.

We now turn to the translation unit implementing type
\lstinline!my::Thing!. The \lstinline!my_thing! \emcppsgloss{component} consists
of a \lstinline!.h!/\lstinline!.cpp! pair: \lstinline!my_thing.h! and
\lstinline!my_thing.cpp!. The header file \lstinline!my_thing.h! provides
the physical interface, such as the definition of the principal type,
\lstinline!Thing!, its member and associated free function declarations,
plus definitions for inline functions and function templates, if any:

\begin{emcppslisting}[emcppsbatch=e4]
// my_thing.h:
#ifndef INCLUDED_MY_THING
#define INCLUDED_MY_THING

namespace my                  // outer namespace (used directly by clients)
{
    inline namespace impl_v1  // inner namespace (for implementer use only)
    {
        struct Thing
        {
            int i;    // original data member, size = 4
            Thing();  // default constructor (defined in (ù{\codeincomments{my\_thing.cpp}}ù))
        };
    };
}

#endif
\end{emcppslisting}
    
\noindent The implementation file \lstinline!my_thing.cpp! contains all of the
non-\lstinline!inline! function bodies that will be translated separately
into the \lstinline!my_thing.o! file:

\begin{emcppslisting}[emcppsbatch=e4]
// my_thing.cpp:
#include <my_thing.h>

namespace my                   // outer namespace (used directly by clients)
{
    inline namespace impl_v1   // inner namespace (for implementer use only)
    {
        Thing::Thing() : i(0)  // load a 4-byte value into (ù{\codeincomments{Thing}}ù)'s data member
        {
        }
    }
}
\end{emcppslisting}
    
\noindent Observing common good practice, we include the header file of the component
as the first substantive line of code to ensure that --- irrespective of
anything else --- the header always compiles in isolation, thereby
avoiding insidious include-order dependencies.{\cprotect\footnote{See
  \cite{lakos20}, section~1.6.1, ``Component Property 1," pp.~210--212.}} When we compile the source file \lstinline!my_thing.cpp!,
we produce an object file \lstinline!my_thing.o! containing the definition
of the very same linker symbol, such as
\lstinline!_ZN2my7impl_v15ThingC1Ev!, for the default constructor of
\lstinline!my::Thing! needed by the client:

\begin{lstlisting}[language=bash,style=plain]
$ g++ -c my_thing.cpp
\end{lstlisting}
    
\noindent We can then link \lstinline!main.o! and \lstinline!my_thing.o! into an
executable and run it:

\begin{lstlisting}[language=bash,style=plain]
$ g++ -o prog main.o my_thing.o
$ ./prog

0
\end{lstlisting}
    
\noindent Now, suppose we were to change the definition of \lstinline!my::Thing! to hold a \lstinline!double! instead of an \lstinline!int!, recompile
\lstinline!my_thing.cpp!, and then relink with the original
\lstinline!main.o! without recompiling \lstinline!main.cpp! first. None of the
relevant linker symbols would change, and the code would recompile and link
just fine, but the resulting binary \lstinline!prog! would be
\emcppsgloss[ill formed, no diagnostic required (IFNDR)]{IFNDR}: the client would be trying to print a 4-byte,
\lstinline!int! data member, \lstinline!i!, in \lstinline!main.o! that was loaded
by the library component as an 8-byte, \lstinline!double! into \lstinline!d!
in \lstinline!my_thing.o!. We can resolve this problem by changing --- or,
if we didn't think of it in advance, by adding --- a new \lstinline!inline!
namespace and making that change there:

\begin{emcppshiddenlisting}[emcppsbatch=e5]
// my_thing.h:
#ifndef INCLUDED_MY_THING
#define INCLUDED_MY_THING

namespace my                  // outer namespace (used directly by clients)
{
    inline namespace impl_v2  // inner namespace (for implementer use only)
    {
        struct Thing
        {
            double d; // new data member, size = 8
            Thing();  // default constructor (defined in (ù{\codeincomments{my\_thing.cpp}}ù))
        };
    };
}

#endif
\end{emcppshiddenlisting}
\begin{emcppslisting}[emcppsbatch=e5]
// my_thing.cpp:
#include <my_thing.h>

namespace my                     // outer namespace (used directly by clients)
{
    inline namespace impl_v2     // inner namespace (for implementer use only)
    {
        Thing::Thing() : d(0.0)  // load 8-byte value into (ù{\codeincomments{Thing}}ù)'s data member
        {
        }
    }
}
\end{emcppslisting}
    
\noindent Now clients that attempt to link against the new library will not find
the linker symbol, such as \lstinline!_Z...impl_v1...v!, and the link
stage will fail. Once clients recompile, however, the undefined linker
symbol will match the one available in the new \lstinline!my_thing.o!,
such as \lstinline!_Z...impl_v2...v!, the link stage will succeed, and
the program will again work as expected. What's more, we have the option
of keeping the original implementation. In that case, existing clients
that have not as yet recompiled will continue to link against the old
version until it is eventually removed after some suitable deprecation
period.

As a more realistic second example of using \lstinline!inline! namespaces
to guard against linking incompatible versions, suppose we have two
versions of a \lstinline!Key! class in a security library in the enclosing
namespace, \lstinline!auth! --- the original version in a regular nested
namespace \lstinline!v1!, and the new current version in an \lstinline!inline!
nested namespace \lstinline!v2!:

\begin{emcppslisting}
#include <cstdint>  // (ù{\codeincomments{std::uint32\_t}}ù), (ù{\codeincomments{std::unit64\_t}}ù)

namespace auth      // outer namespace (used directly by clients)
{
    namespace v1    // inner namespace (optionally used by clients)
    {
        class Key
        {
        private:
            std::uint32_t d_key;
                // (ù{\codeincomments{sizeof(Key)}}ù) is 4 bytes

        public:
            std::uint32_t key() const;  // stable interface function

            // ...
        };
    }

    inline namespace v2    // inner namespace (default current version)
    {
        class Key
        {
        private:
            std::uint64_t d_securityHash;
            std::uint32_t d_key;
                // (ù{\codeincomments{sizeof(Key)}}ù) is 16 bytes

        public:
            std::uint32_t key() const;  // stable interface function

            // ...
        };
    }
}
\end{emcppslisting}
    
\noindent Attempting to link together older binary artifacts built against version
1 with binary artifacts built against version 2 will result in a
link-time error rather than allowing an \emcppsgloss[ill formed]{ill formed} program to
be created. Note, however, that this approach works only if functionality essential
to typical use is defined out of line in a \lstinline!.cpp! file. For example, it would add absolutely no value for libraries that are
shipped entirely as header files, since the versioning offered here
occurs strictly at the binary level (i.e., between object files)
during the link stage.

\subsubsection[Build modes and ABI link safety]{Build modes and ABI link safety}\label{build-modes-and-abi-link-safety}

In certain scenarios, a class might have two different memory layouts
depending on compilation flags. For instance, consider a low-level
\lstinline!ManualBuffer! class template in which an additional data member
is added for debugging purposes:

\begin{emcppslisting}
template <typename T>
struct ManualBuffer
{
private:
    alignas(T) char d_data[sizeof(T)];  // aligned and big enough to hold a (ù{\codeincomments{T}}ù)

#ifndef NDEBUG
    bool d_engaged;  // tracks whether buffer is full (debug builds only)
#endif

public:
    void construct(const T& obj);
        // Emplace (ù{\codeincomments{obj}}ù). (Engage the buffer.) The behavior is undefined unless
        // the buffer was not previously engaged.

    void destroy();
        // Destroy the current (ù{\codeincomments{obj}}ù). (Disengage the buffer.) The behavior is
        // undefined unless the buffer was previously engaged.

    // ...
};
\end{emcppslisting}
    
\noindent Note that we have
  employed the C++11 \lstinline!alignas! attribute (see \featureref{\locationc}{alignas})   
 % ``\titleref{alignas}" on page~\pageref{alignas}) 
  here
  because it is exactly what's needed for this usage example.

The \lstinline!d_engaged! flag in the example above serves as a way to detect misuse of
the \mbox{\lstinline!ManualBuffer!} class but only in debug builds. The extra
space and run time required to maintain this Boolean flag is undesirable
in a release build because \mbox{\lstinline!ManualBuffer!} is intended to be an
efficient, lightweight abstraction over the direct use of
\emcppsgloss[placement new]{placement \lstinline!new!} and explicit destruction.

The linker symbol names generated for the methods of
\lstinline!ManualBuffer! are the same irrespective of the chosen build
mode. If the same program links together two object files where
\lstinline!ManualBuffer! is used --- one built in debug mode and one built
in release mode --- the \emcppsgloss[one-definition rule (ODR)]{one-definition rule (ODR)} will be violated,
and the program will again be \emcppsgloss[ill formed, no diagnostic required (IFNDR)]{IFNDR}.

Prior
to \lstinline!inline! namespaces, it was possible to control the
ABI-level name of linked symbols by creating separate template
instantiations on a per-build-mode basis:

\begin{emcppslisting}[]
#ifndef NDEBUG
enum { is_debug_build = 1 };
#else
enum { is_debug_build = 0 };
#endif

template <typename T, bool Debug = is_debug_build>
struct ManualBuffer { /* ... */ };
\end{emcppslisting}
    
\noindent While the code above changes the interface of \lstinline!ManualBuffer! to
accept an additional template parameter, it also allows debug and
release versions of the same class to coexist in the same program,
  which might prove useful, e.g., for testing.
  
Another way of avoiding incompatibilities at link time is to
introduce two \lstinline!inline! namespaces, the entire purpose of which is
to change the ABI-level names of the linker symbols associated with
\lstinline!ManualBuffer! depending on the build mode:

\begin{emcppslisting}
#ifndef NDEBUG            // perhaps a BAD IDEA
inline namespace release
#else
inline namespace debug
#endif
{
    template <typename T>
    struct ManualBuffer
    {
        // ... (same as above)
    };
}
\end{emcppslisting}
    
\noindent The approach demonstrated in this example tries to ensure that a linker
error will occur if any attempt is made to link objects built with a
build mode different from that of \mbox{\lstinline!manualbuffer.o!}. Tying it to
the \lstinline!NDEBUG! flag, however, might have unintended consequences;
we might introduce unwanted restrictions in what we call
\emcppsgloss{mixed-mode builds}. Most modern platforms support the notion of
linking a collection of object files irrespective of their
optimization levels. The same is certainly true for whether or not
C-style \lstinline!assert! is enabled. In other words, we may want to have
a mixed-mode build where we link object files that differ in their optimization and assertion options, as long as they are binary compatible --- i.e., in this case, they all must be uniform with respect to the implementation of \lstinline!ManualBuffer!.  Hence, a more general, albeit more complicated and
manual, approach would be to tie the non-interoperable behavior
associated with this ``safe'' or ``defensive'' build mode to a different
switch entirely. Another consideration would be to avoid ever inlining a
namespace into the global namespace since no method is available to
recover a symbol when there is a collision:

\begin{emcppslisting}
namespace buflib  // GOOD IDEA: enclosing namespace for nested (ù{\codeincomments{inline}}ù) namespace
{
#ifdef SAFE_MODE  // GOOD IDEA: separate control of non-interoperable versions
    inline namespace safe_build_mode
#else
    inline namespace normal_build_mode
#endif
    {
        template <typename T>
        struct ManualBuffer
        {
        private:
            alignas(T) char d_data[sizeof(T)];  // aligned/sized to hold a (ù{\codeincomments{T}}ù)

#ifdef SAFE_MODE
            bool d_engaged;  // tracks whether buffer is full (safe mode only)
#endif

        public:
            void construct(const T& obj);  // sets (ù{\codeincomments{d\_engaged}}ù) (safe mode only)
            void destroy();                // sets (ù{\codeincomments{d\_engaged}}ù) (safe mode only)
            // ...
        };
    }
}
\end{emcppslisting}
    
\noindent And, of course, the appropriate conditional compilation within the
function bodies would need to be in the corresponding \lstinline!.cpp!
file.

Finally, if we have two implementations of a particular entity that are
sufficiently distinct, we might choose to represent them in their
entirety, controlled by their own bespoke conditional-compilation
switches, as illustrated here using the \lstinline!my::VersionedThing! type
(see \intraref{use-cases-inlinenamespace}{link-safe-abi-versioning}): 
% \textit{\titleref{use-cases-inlinenamespace}: \titleref{link-safe-abi-versioning}} on page~\pageref{link-safe-abi-versioning}:

\begin{emcppshiddenlisting}[emcppsbatch=e7]
// my_versionedthing.h:
#define MY_THING_VERSION_1
\end{emcppshiddenlisting}
\begin{emcppslisting}[emcppsbatch=e7]
// my_versionedthing.h:
#ifndef INCLUDED_MY_VERSIONEDTHING
#define INCLUDED_MY_VERSIONEDTHING

namespace my
{
#ifdef MY_THING_VERSION_1  // bespoke switch for this component version
    inline
#endif
    namespace v1
    {
        struct VersionedThing
        {
            int d_i;
            VersionedThing();
        };
    }

#ifdef MY_THING_VERSION_2  // bespoke switch for this component version
    inline
#endif
    namespace v2
    {
        struct VersionedThing
        {
            double d_i;
            VersionedThing();
        };
    }
}
#endif
\end{emcppslisting}
    
\noindent However, see \intraref{potential-pitfalls-inlinenamespace}{inline-namespace-based-versioning-doesn’t-scale}. 
%\textit{\titleref{potential-pitfalls-inlinenamespace}: \titleref{inline-namespace-based-versioning-doesn’t-scale}} on page~\pageref{inline-namespace-based-versioning-doesn’t-scale}.

\subsubsection[Enabling selective \lstinline!using! directives for short-named entities]{Enabling selective {\SubsubsecCode using} directives for short-named entities}\label{enabling-selective-using-directives-for-short-named-entities}

Introducing a large number of small names into client code that doesn't
follow rigorous nomenclature can be problematic. Hoisting these names
into one or more nested namespaces so that they are easier to identify
as a unit and can be used more selectively by clients, such as through
explicit qualification or using directives, can sometimes be an
effective way of organizing shared codebases. For example,
\lstinline!std::literals! and its nested namespaces, such as
\lstinline!chrono_literals!, were introduced as \lstinline!inline! namespaces
in C++14. As it turns out, clients of these nested namespaces have no
need to specialize any templates defined in these namespaces nor do they
define types that must be found through \emcppsgloss[argument-dependent lookup (ADL)]{ADL}, but one can at
least imagine special circumstances in which such tiny-named entities
are either templates that require specialization or operator-like
functions, such as \lstinline!swap!, defined for local types within those
nested namespaces. In those cases, \lstinline!inline! namespaces would be
required to preserve the desired ``as if'' properties.

Even without either of these two needs, another property of an
\lstinline!inline! namespace differentiates it from a non-\lstinline!inline!
one followed by a \lstinline!using! directive. Recall from\intraref{description-inlinenamespace}{loss-of-access-to-duplicate-names-in-enclosing-namespace}  
%\textit{\titleref{description-inlinenamespace}: \titleref{loss-of-access-to-duplicate-names-in-enclosing-namespace}} on page~\pageref{loss-of-access-to-duplicate-names-in-enclosing-namespace} 
that a name in an outer namespace will
hide a duplicate name imported via a \lstinline!using! directive, whereas
any access to that duplicate name within the enclosing namespace would
be ambiguous when that symbol is installed by way of an \lstinline!inline!
namespace. To see why this more forceful clobbering behavior might be
preferred over hiding, suppose we have a communal namespace \lstinline!abc!
that is shared across multiple disparate headers. The first header,
\lstinline!abc_header1.h!, represents a collection of logically related
small functions declared directly in \lstinline!abc!:

\begin{emcppslisting}[emcppsbatch=e6]
// abc_header1.h:
namespace abc
{
    int i();
    int am();
    int smart();
}
\end{emcppslisting}
    
\noindent A second header, \lstinline!abc_header2.h!, creates a suite of many
functions having tiny function names. In a perhaps misguided effort to
avoid clobbering other symbols within the \lstinline!abc! namespace having
the same name, all of these tiny functions are sequestered within a
\lstinline!nested! namespace:

\begin{emcppslisting}[emcppsbatch=e6]
// abc_header2.h:
namespace abc
{
    namespace nested  // Should this instead have been an (ù{\codeincomments{inline}}ù) namespace?
    {
        int a();  // lots of functions with tiny names
        int b();
        int c();
        // ...
        int h();
        int i();  // might collide with another name declared in (ù{\codeincomments{abc}}ù)
        // ...
        int z();
    }

    using namespace nested;  // becomes superfluous if (ù{\codeincomments{nested}}ù) is made (ù{\codeincomments{inline}}ù)
}
\end{emcppslisting}
    
\noindent Now suppose that a client application includes
both of these headers to accomplish some task:

\begin{emcppslisting}[emcppsbatch=e6]
// client.cpp:
#include <abc_header1.h>
#include <abc_header2.h>

int function()
{
    if (abc::smart() < 0) { return -1; }  // uses (ù{\codeincomments{smart()}}ù) from (ù{\codeincomments{abc\_header1.h}}ù)
    return abc::z() + abc::i() + abc::a() + abc::h() + abc::c();  // Oops!
        // Bug, silently uses the (ù{\codeincomments{abc::i()}}ù) defined in (ù{\codeincomments{abc\_header1.h}}ù)
}
\end{emcppslisting}
    
\noindent In trying to cede control to the client as to whether the declared or
imported \lstinline!abc::i()! function is to be used, we have, in effect,
invited the defect illustrated in the above example whereby the client
was expecting the \lstinline!abc::i()! from \lstinline!abc_header2.h! and yet
picked up the one from \lstinline!abc_header1.h! by default. Had the
\lstinline!nested! namespace in \lstinline!abc_header2.h! been declared
\lstinline!inline!, the qualified name \lstinline!abc::i()! would have
automatically been rendered \emph{ambiguous} in namespace \lstinline!abc!,
the translation would have failed \emph{safely}, and the defect would
have been exposed at compile time. The downside, however, is that no
method would be available to recover nominal access to the
\lstinline!abc::i()! defined in \lstinline!abc_header1.h! once
\lstinline!abc_header2.h! is included, even though the two functions
(e.g., including their \emcppsgloss{mangled names} at the ABI level) remain
distinct.

\subsection[Potential Pitfalls]{Potential Pitfalls}\label{potential-pitfalls-inlinenamespace}

\subsubsection[\lstinline!nline!-namespace-based versioning doesn’t scale]{{\SubsubsecCode inline}-namespace-based versioning doesn’t scale}\label{inline-namespace-based-versioning-doesn’t-scale}

The problem with using \lstinline!inline! namespaces for ABI link safety is
that the protection they offer is only partial; in a few major places,
critical problems can linger until run time instead of being caught at
compile time.

Controlling which namespace is \lstinline!inline! using macros, such as was
done in the\linebreak[4] \lstinline!my::VersionedThing! example in \intraref{use-cases-inlinenamespace}{link-safe-abi-versioning},  
%\textit{\titleref{use-cases-inlinenamespace}: \titleref{link-safe-abi-versioning}} on page~\pageref{link-safe-abi-versioning}, 
will result in code that
directly uses the unversioned name, \lstinline!my::VersionedThing! being
bound directly to the versioned name \lstinline!my::v1::VersionedThing! or
\lstinline!my::v2::VersionedThing!, along with the class layout of that
particular entity. Sometimes details of using the \lstinline!inline!
namespace member are not resolved by the linker, such as the object
layout when we use types from that namespace as member variables in
other objects:

\begin{emcppslisting}[emcppsbatch=e7]
// my_thingaggregate.h:

// ...
#include <my_versionedthing.h>
// ...

namespace my
{
    struct ThingAggregate
    {
        // ...
        VersionedThing d_thing;
        // ...
    };
}
\end{emcppslisting}
    
\noindent This new \lstinline!ThingAggregate! type does not have the versioned
\lstinline!inline! namespace as part of its mangled name; it does, however,
have a completely different layout if built with
\lstinline!MY_THING_VERSION_1! defined
versus \lstinline!MY_THING_VERSION_2! defined. Linking a program with mixed
versions of these flags will result in runtime failures that are
decidedly difficult to diagnose.

This same sort of problem will arise for functions taking arguments of
such types; calling a function from code that is wrong about the layout
of a particular type will result in stack corruption and other undefined
and unpredictable behavior. This macro-induced problem will also arise in cases where an old object
file is linked against new code that changes which namespace is
\lstinline!inline!d but still provides the definitions for the old version
namespace. The old object file for the client can still link, but new
object files using the headers for the old objects might attempt to
manipulate those objects using the new namespace.

The only viable workaround for this approach is to propagate the
\lstinline!inline! namespace hierarchy through the entire software stack.
Every object or function that uses\linebreak[4] \lstinline!my::VersionedThing! needs to
also be in a namespace that differs based on the same control macro. In
the case of \lstinline!ThingAggregate!, one could just use the same
\lstinline!my::v1! and \lstinline!my::v2! namespaces, but higher-level
libraries would need their own \lstinline!my!-specific nested namespaces.
Even worse, for higher-level libraries, every lower-level library having
a versioning scheme of this nature would need to be considered,
resulting in having to provide the full cross-product of nested
namespaces to get link-time protection against mixed-mode builds.

This need for layers above a library to be aware of and to integrate
into their own structure the same namespaces the library has removes all
or most of the benefits of using \lstinline!inline! namespaces for
versioning. For an authentic real-world case study of heroic industrial use --- and
eventual disuse --- of \lstinline!inline!-namespaces for versioning, see \intrarefsimple{appendix:-case-study-of-using-inline-namespaces-for-versioning}. 
% \textit{\titleref{appendix:-case-study-of-using-inline-namespaces-for-versioning}} on page~\pageref{appendix:-case-study-of-using-inline-namespaces-for-versioning}.

\subsubsection[Relying on \lstinline!inline! namespaces to solve library evolution]{Relying on {\SubsubsecCode inline} namespaces to solve library evolution}\label{specializing-templates-in-std-can-be-problematic}

%Inline namespaces might be perceived as a complete solution for the owner of a library to evolve its API. Such, however, is not the case. As an example, consider the case of the C++ Standard Library, which does not use \lstinline!inline! namespaces for versioning. Instead, to allow for its necessary evolution, the Standard
%Library --- the \lstinline!std! namespace, in particular --- carries
%certain special restrictions on what one can do with it that are
%enforced by deeming certain constructs \romeogloss{ill formed} or
%engendering \romeogloss{undefined behavior}.

Inline namespaces might be misperceived as a complete solution for the owner of a library to evolve its API. As an especially relevant example, consider the C++ Standard Library, which itself does not use inline namespaces for versioning. Instead, to allow for its anticipated essential evolution, the Standard Library imposes certain special restrictions on what is permitted to occur within its own \lstinline!std! namespace by dint of deeming certain problematic uses as either \emcppsgloss{ill formed} or otherwise engendering \emcppsgloss{undefined behavior}.

Since C++11, several restrictions related to the Standard Library were
put in place:
\begin{itemize}
\item{Users may not add any new declarations within namespace \lstinline!std!. This means that users cannot add new \emph{functions}, \emph{overloads}, \emph{types}, or \emph{templates} to \lstinline!std!. This restriction gives the Standard Library freedom to add new \emph{names} in future versions of the Standard.}
\item{Users may not specialize member functions, member function templates, or member class templates. Specializing any of those entities might significantly inhibit a Standard Library vendor’s ability to maintain its otherwise encapsulated implementation details.}
\item{Users may add specializations of top-level Standard Library templates only if the declaration depends on the name of a nonstandard user-defined type and only if that user-defined type meets all requirements of the original template. Specialization of function templates is allowed but generally discouraged because this practice doesn’t scale since function templates cannot be partially specialized. Specializing of standard class templates when the specialization names a nonstandard user-defined type, such as \lstinline!std::vector<MyType*>!, is allowed but also problematic when not explicitly supported. While certain specific types, such as \lstinline!std::hash!, are designed for user specialization, steering clear of the practice for any other type helps to avoid surprises.}
\end{itemize}

Several other good practices facilitate smooth evolution for the
Standard Library{\cprotect\footnote{These restrictions are normative in
C++20, having finally formalized what were long identified as best
practices. Though these restrictions might not be codified in the
Standard for pre-C++20 software, they have been recognized best
practices for as long as the Standard Library has existed and
adherence to them will materially improve the ability of software to
migrate to future language standards irrespective of what version of
  the language standard is being targeted.}}:
\begin{itemize}
\item{Avoid specializing variable templates, even if dependent on user-defined types, except for those variable templates where specialization is explicitly allowed.\cprotect\footnote{C++20 limits the specialization of variable templates to only those instances where specialization is explicitly allowed and does so only for the mathematical constants in \lstinline!<numbers>!.}}
\item{Other than a few very specific exceptions, avoiding the forming of pointers to Standard Library functions — either explicitly or implicitly — allows the library to add overloads, either as part of the Standard or as an implementation detail for a particular Standard Library, without breaking user code.\cprotect\footnote{C++20 identifies these functions as \lstinline!addressable! and gives that property to only \lstinline!iostream! manipulators since those are the only functions in the Standard Library for which taking their address is part of normal usage.}}
\item{Overloads of Standard Library functions that depend on user-defined types are permitted, but, as with specializing Standard Library templates, users must still meet the requirements of the Standard Library function. Some functions, such as \lstinline!std::swap!, are designed to be customization points via overloading, but leaving functions not specifically designed for this purpose to vendor implementations only helps to avoid surprises.}
\end{itemize}

Finally, upon reading about this \lstinline!inline! namespace feature, one
might think that all names in namespace \lstinline!std! could be made
available at a global scope simply by inserting an\linebreak[4]%%%%%%
\lstinline!inline!~\lstinline!namespace!~\lstinline!std!~\lstinline!{}! before
including any standard headers. This practice is, however, explicitly
called out as ill-formed within the C++11 Standard. Although not uniformly diagnosed as an error by
all compilers, attempting this forbidden practice is apt to lead to
surprising problems even if not diagnosed as an error immediately.

\subsubsection[Inconsistent use of \lstinline!inline! keyword is ill formed, no diagnostic required]{Inconsistent use of {\SubsubsecCode inline} keyword is ill formed, no diagnostic required}\label{inconsistent-use-of-inline-keyword-is-ifndr}

It is an \emcppsgloss[one-definition rule (ODR)]{ODR} violation, \emcppsgloss[ill formed, no diagnostic required (IFNDR)]{IFNDR}, for a nested namespace
to be \lstinline!inline! in one translation unit and non-\lstinline!inline! in
another. And yet, the motivating use case of this feature relies on the
linker to actively complain whenever different, incompatible versions
--- nested within different, possibly \lstinline!inline!-inconsistent,
namespaces of an ABI --- are used within a single executable. Because
declaring a nested namespace \lstinline!inline! does not, by design, affect
linker-level symbols, developers must take appropriate care, such as
effective use of header files, to defend against such preventable
inconsistencies.

\subsection[Annoyances]{Annoyances}\label{annoyances-inlinenamespace}

\subsubsection[Inability to redeclare across namespaces impedes code factoring]{Inability to redeclare across namespaces impedes code factoring}\label{inability-to-redeclare-across-namespaces-impedes-code-factoring}

An essential feature of an \lstinline!inline! namespace is the ability to
declare a template within a nested \lstinline!inline! namespace and then
specialize it within its enclosing namespace. For example, we can
declare
\begin{itemize}
\item{a type template, \lstinline!S0!}
\item{a couple of function templates, \lstinline!f0! and \lstinline!g0!}
\item{and a member function template \lstinline!h0!, which is similar to \lstinline!f0!}
\end{itemize}
in an \lstinline!inline! namespace, \lstinline!inner!, and specialize each of
them, such as for \lstinline!int!, in the enclosing namespace,
\lstinline!outer!:

\begin{emcppslisting}
namespace outer                                          // enclosing namespace
{
    inline namespace inner                               // nested namespace
    {
        template<typename T> struct S0;                  // declarations of
        template<typename T> void f0();                  // various class
        template<typename T> void g0(T v);               // and function
        struct A0 { template <typename T> void h0(); };  // templates
    }

    template<> struct S0<int> { };                       // specializations
    template<> void f0<int>() { }                        // of the various
    void g0(int) { }  /* overload not specialization */  // class and function
    template<> void A0::h0<int>() { }                    // declarations above
}                                                        // in (ù{\codeincomments{outer}}ù) namespace
\end{emcppslisting}
    
\noindent Note that, in the case of \lstinline!g0! in this example, the
``specialization'' \lstinline!void!~\lstinline!g0(int)! is a non-template
\emph{overload} of the function template \lstinline!g0! rather than a
specialization of it. We \emph{cannot}, however,
portably{\cprotect\footnote{GCC provides the \lstinline!-fpermissive! flag,
which allows the example containing specializations within the inner
namespace to compile with warnings. Note again that \lstinline!g1(int)!,
being an \emph{overload} and not a \emph{specialization}, wasn't an
  error and, therefore, isn't a warning either.}} declare these
templates within the \lstinline!outer! namespace and then specialize them
within the \lstinline!inner! one, even though the \lstinline!inner! namespace
is \lstinline!inline!:

\begin{emcppslisting}
namespace outer                                     // enclosing namespace
{
    template<typename T> struct S1;                 // class template
    template<typename T> void f1();                 // function template
    template<typename T> void g1(T v);              // function template

    struct A1 { template <typename T> void h1(); }; // member function template

    inline namespace inner                          // nested namespace
    {                                               // BAD IDEA
        template<> struct S1<int> { };              // Error, (ù{\codeincomments{S1}}ù) not a template
        template<> void f1<int>() { }               // Error, (ù{\codeincomments{f1}}ù) not a template
        void g1(int) { }                            // OK, overloaded function
        template<> void A1::h1<int>() { }           // Error, (ù{\codeincomments{h1}}ù) not a template
    }
}
\end{emcppslisting}
    
\noindent Attempting to declare a template in the \lstinline!outer! namespace and
then define, effectively redeclaring, it in an \lstinline!inline! inner one
causes the name to be inaccessible within the \lstinline!outer! namespace:

\begin{emcppslisting}[emcppserrorlines=13]
namespace outer                                          // enclosing namespace
{                                                        // BAD IDEA
    template<typename T> struct S2;                      // declarations of
    template<typename T> void f2();                      // various class and
    template<typename T> void g2(T v);                   // function templates
    
    inline namespace inner                               // nested namespace
    {
        template<typename T> struct S2 { };              // definitions of
        template<typename T> void f2() { }               // unrelated class and
        template<typename T> void g2(T v) { }            // function templates
    }

    template<> struct S2<int> { };     // Error, (ù{\codeincomments{S2}}ù) is ambiguous in (ù{\codeincomments{outer}}ù).
    template<> void f2<int>() { }      // Error, (ù{\codeincomments{f2}}ù) is ambiguous in (ù{\codeincomments{outer}}ù).
    void g2(int) { }                   // OK, (ù{\codeincomments{g2}}ù) is an overload definition.
}
\end{emcppslisting}
    
\noindent Finally, declaring a template in the nested \lstinline!inline! namespace
\lstinline!inner! in the example above and then subsequently defining it in
the enclosing \lstinline!outer! namespace has the same effect of making
declared symbols ambiguous in the \lstinline!outer! namespace:

\begin{emcppslisting}[emcppsbatch=e8]
namespace outer                                          // enclosing namespace
{                                                        // BAD IDEA
    inline namespace inner                               // nested namespace
    {
        template<typename T> struct S3;                  // declarations of
        template<typename T> void f3();                  // various class
        template<typename T> void g3(T v);               // and function
        struct A3 { template <typename T> void h3(); };  // templates
    }

    template<typename T> struct S3 { };                  // definitions of
    template<typename T> void f3() { }                   // unrelated class
    template<typename T> void g3(T v) { }                // and function
    template<typename T> void A3::h3() { };              // templates

    template<> struct S3<int> { };     // Error, (ù{\codeincomments{S3}}ù) is ambiguous in (ù{\codeincomments{outer}}ù).
    template<> void f3<int>() { }      // Error, (ù{\codeincomments{f3}}ù) is ambiguous in (ù{\codeincomments{outer}}ù).
    void g3(int) { }                   // OK, (ù{\codeincomments{g3}}ù) is an *overload* definition.
    template<> void A3::h3<int>() { }  // Error, (ù{\codeincomments{h2}}ù) is ambiguous in (ù{\codeincomments{outer}}ù).
}
\end{emcppslisting}
    
\noindent Note that, although the definition for a member function template must
be located directly within the namespace in which it is declared, a
class or function template, once declared, may instead be defined in a different
scope by using an appropriate name qualification:

\begin{emcppshiddenlisting}[emcppsbatch=e9]
namespace outer
{
    inline namespace inner
    {
        template<typename T> struct S3;
        template<typename T> void f3();
        template<typename T> void g3(T v);
        struct A3 { template <typename T> void h3(); };
    }
}
\end{emcppshiddenlisting}
\begin{emcppslisting}[emcppsbatch=e9]
template <typename T> struct outer::S3 { };        // OK, enclosing namespace
template <typename T> void outer::inner::f3() { }  // OK, nested namespace
template <typename T> void outer::g3(T v) { }      // OK, enclosing namespace
template <typename T> void outer::A3::h3<T>() { }  // Error, ill-formed

namespace outer
{
    inline namespace inner
    {
        template <typename T> void A3::h3() { }    // OK, within same namespace
    }
}
\end{emcppslisting}
    
\noindent Also note that, as ever, the corresponding definition of the declared
template must have been seen before it can be used in a context
requiring a complete type. The importance of ensuring that all
specializations of a template have been seen before it is used
substantively (i.e., \emcppsgloss[one-definition rule (ODR)]{ODR}-used) cannot be overstated, giving
rise to the only limerick, which is actually part of the normative text,
in the C++ Language Standard{\cprotect\footnote{See \cite{iso11},
  section~14.7.3.7, pp.~375--375, specifically p.~376.}}:
\begin{quote}
When writing a specialization,\\
be careful about its location;\\
or to make it compile\\
will be such a trial\\
as to kindle its self-immolation.
\end{quote}

\subsubsection[Only one namespace can contain any given \lstinline!inline! namespace]{Only one namespace can contain any given {\SubsubsecCode inline} namespace}\label{only-one-namespace-can-contain-any-given-inline-namespace}

Unlike conventional \lstinline!using! directives, which can be used to
generate arbitrary many-to-many relationships between different
namespaces, \lstinline!inline! namespaces can be used only to contribute
names to the sequence of enclosing namespaces up to the first
non-\lstinline!inline! one. In cases in which the names from a namespace
are desired in multiple other namespaces, the classical \lstinline!using!
directive must be used, with the subtle differences between the two
modes properly addressed.

As an example, the C++14 Standard Library provides a hierarchy of nested
\lstinline!inline! namespaces for literals of different sorts within
namespace\linebreak[4] \lstinline!std!: \lstinline!std::literals::complex_literals!,
\lstinline!std::literals::chrono_literals!,\linebreak[4]
\lstinline!std::literals::string_literals!, and
\lstinline!std::literals::string_view_literals!.\linebreak[4] 
These namespaces can
be imported to a local scope in one shot via a
\lstinline!using!~\lstinline!std::literals! or instead, more selectively, by
\lstinline!using! the nested namespaces directly. This separation of the
types used with user-defined literals, which are all in namespace
\lstinline!std!, from the user-defined literals that can be used to create
those types led to some frustration; those who had a
\lstinline!using!~\lstinline!namespace!~\lstinline!std;! could reasonably have
expected to get the user-defined literals associated with their
\lstinline!std! types. However, the types in the nested namespace
\lstinline!std::chrono! did \emph{not} meet this
expectation.{\cprotect\footnote{\cite{hinnant17}}}

Eventually \emph{both} solutions for incorporating literal namespaces,
\lstinline!inline! from\linebreak[4] \lstinline!std::literals! and non-\lstinline!inline! from
\lstinline!std::chrono!, were pressed into service when, in C++17, a
\mbox{\lstinline!using!~\lstinline!namespace!~\lstinline!literals::chrono_literals;!}
was added to the\linebreak[4]
 \lstinline!std::chrono! namespace. The Standard does not, however, benefit in any objective way from any of
these namespaces being \lstinline!inline! since the artifacts in the
\lstinline!literals! namespace neither depend on ADL nor are templates in
need of user-defined specializations; hence, having all
non-\lstinline!inline! namespaces with appropriate \lstinline!using!
declarations would have been functionally indistinguishable from the
bifurcated approach taken.

\subsection[See Also]{See Also}\label{see-also}

\begin{itemize}
\item{%``\titleref{alignas}" on page~\pageref{alignas} — 
\seealsoref{alignas}{\seealsolocationc}provides properly aligned storage for an object of arbitrary type \lstinline!T! in the example in  \intraref{use-cases-inlinenamespace}{build-modes-and-abi-link-safety}.}
\end{itemize}

\subsection[Further Reading]{Further Reading}\label{further-reading}

\begin{itemize}
\item{\cite{sutter14b} uses inline namespaces as part of a proposal for a portable ABI across compilers.}
\item{\cite{lopez-gomez20} uses inline namespaces as part of a solution to avoid ODR violation in an interpreter.}
\end{itemize}

\subsection[Appendix: Case study of using \lstinline!inline! namespaces for versioning]{Appendix: Case study of using {\SubsecCode inline} namespaces for versioning}\label{appendix:-case-study-of-using-inline-namespaces-for-versioning}

\noindent\textbf{By Niall Douglas}\\[.5ex]

\noindent Let me tell you what I (don't) use them for. It is not a conventional
opinion.

At a previous well-regarded company, they were shipping no less than
forty-three copies of Boost in their application. Boost was not on the
approved libraries list, but the great thing about header-only libraries
is that they don't obviously appear in final binaries, unless you look
for them. So each individual team was including bits of Boost quietly
and without telling their legal department. Why? Because it saved time.
(This was C++98, and \lstinline!boost::shared_ptr! and
\lstinline!boost::function! are both extremely attractive facilities.)

Here's the really interesting part: Most of these copies of Boost were
not the same version. They were varying over a five-year release period.
And, unfortunately, Boost makes no API or ABI guarantees. So,
theoretically, you could get two different incompatible versions of
Boost appearing in the same program binary, and BOOM! there goes memory
corruption.

I advocated to Boost that a simple solution would be for Boost to wrap
up their implementation into an internal inline namespace. That inline
namespace ought to mean something:
\begin{itemize}
\item{\lstinline!lib::v1! is the \emph{stable}, version-1 ABI, which is guaranteed to be compatible with all past and future \lstinline!lib::v1! ABIs, forever, as determined by the ABI-compliance-check tool that runs on \emcppsgloss{CI}. The same goes for \lstinline!v2!, \lstinline!v3!, and so on.}
\item{\lstinline!lib::v2_a7fe42d! is the \emph{unstable}, version-2 ABI, which may be incompatible with any other \lstinline!lib::*!~\lstinline!ABI!; hence, the seven hex chars after the underscore are the git short \emcppsgloss{SHA}, permuted by every commit to the git repository but, in practice, per CMake configure, because nobody wants to rebuild everything per commit. This ensures that no symbols from any revision of \lstinline!lib! will \emph{ever} silently collide or otherwise interfere with any other revision of \lstinline!lib!, when combined into a single binary by a dumb linker.}
\end{itemize}

I have been steadily making progress on getting Boost to avoid putting
anything in the global namespace, so a straightforward find-and-replace
can let you ``fix'' on a particular version of Boost.

That's all the same as the pitch for \lstinline!inline! namespaces. You'll
see the same technique used in \lstinline!libstdc++! and many other major modern C++
codebases.

But I'll tell you now, I don't use \lstinline!inline! namespaces anymore.
Now what I do is use a macro defined to a uniquely named namespace. My
build system uses the git SHA to synthesize namespace macros for my
namespace name, beginning the namespace and ending the namespace.
Finally, in the documentation, I teach people to always use a namespace
alias to a macro to denote the namespace:

\begin{emcppshiddenlisting}[emcppsbatch=e10]
#define OUTCOME_V2_NAMESPACE ov2ns
namespace OUTCOME_V2_NAMESPACE {
}
\end{emcppshiddenlisting}
\begin{emcppslisting}[emcppsbatch=e10]
namespace output = OUTCOME_V2_NAMESPACE;
\end{emcppslisting} 

\noindent That macro expands to something like \lstinline!::outcome_v2_ee9abc2!;
that is, I don't use \lstinline!inline! namespaces anymore.

Why?

Well, for \emph{existing} libraries that don't want to break backward
source compatibility, I think \lstinline!inline! namespaces serve a need.
For \emph{new} libraries, I think a macro-defined namespace is clearer.
\begin{itemize}
\item{It causes users to publicly commit to ``I know what you’re doing here, what it means, and what its consequences are.''}
\item{It declares to \emph{other} users that something unusual (i.e., go read the documentation) is happening here, instead of silent magic behind the scenes.}
\item{It prevents accidents that interfere with ADL and other customization points, which induce surprise, such as accidentally injecting a customization point into \lstinline!lib!, not into \lstinline!lib::v2!.}
\item{Using macros to denote namespace lets us reuse the preprocessor machinery to generate C++ modules using the exact same codebase; C++ modules are used if the compiler supports them, else we fall back to inclusion.}
\end{itemize}

Finally, and here's the real rub, because we now have namespace aliases,
if I were tempted to use an \lstinline!inline! namespace, nowadays I probably would
instead use a uniquely named namespace instead, and, in the \lstinline!include! file,
I'd alias a user-friendly name to that uniquely named namespace. I think
that approach is less likely to induce surprise in the typical
developer's likely use cases than \lstinline!inline! namespaces, such as
injecting customization points into the wrong namespace.

So now I hope you've got a good handle on \lstinline!inline! namespaces: I
was once keen on them, but after some years of experience, I've gone off
them in favor of better-in-my-opinion alternatives. Unfortunately, if your type \lstinline!x::S! has members of type
\lstinline!a::T! and macros decide if that is \lstinline!a::v1::T! or
\lstinline!a::v2::T!, then no linker protects the higher-level types from
ODR bugs, unless you also version \lstinline!x!.




\newpage
%\section[{\tt noexcept} Specifier]{The {\SecCode noexcept} Function Specification}\label{noexcept-specifier}


\emcppsFeature{
    short={\lstinline!noexcept! Specifier},
    tocshort={{\TOCCode noexcept} Specifier},
    long={The {\SecCode noexcept} Function Specification},
    toclong={The \lstinline!noexcept! Function Specification},
    rhshort={{\RHCode noexcept} Specifier},
}{noexcept-specifier}
\setcounter{table}{0}
\setcounter{footnote}{0}
\setcounter{lstlisting}{0}
%\section[{\tt noexcept} Specifier]{The {\SecCode noexcept} Function Specification}\label{noexcept-specifier}

placeholder

%%%%%%% NOTE: Labels will need to be changed since this is now two different sections, not one. 

\newpage
%\section[Ref-Qualifiers]{Reference Qualified Member Functions}\label{refqualifiers}



\emcppsFeature{
    short={Ref-Qualifiers},
    long={Reference-Qualified Member Functions},
}{refqualifiers}
\setcounter{table}{0}
\setcounter{footnote}{0}
\setcounter{lstlisting}{0}
%\section[Ref-Qualifiers]{Reference Qualified Member Functions}\label{refqualifiers}


placeholder


\newpage
%\section[{\tt union} '11]{Unions Having Non-Trivial Members\sectionmark{{\RHCode union}~'11}}\label{unrestricted-unions}\sectionmark{{\RHCode union}~'11}
%%%%% copyedits in and proofed
%% 27 Jan, code checked and replaced as needed by Josh



\emcppsFeature{
    short={\lstinline!union!~'11},
    tocshort={{\TOCCode union}~'11},
    rhshort={{\RHCode union}~'11},
    long={Unions Having Non-Trivial Members},
}{unrestricted-unions}
%\section[{\tt union} '11]{Unions Having Non-Trivial Members\sectionmark{{\RHCode union}~'11}}\label{unrestricted-unions}\sectionmark{{\RHCode union}~'11}


Any nonreference type is permitted to be a member of a \lstinline!union!.

\subsection[Description]{Description}\label{unrestrictedunion-description}

Prior to C++11, only \romeogloss{trivial types} --- e.g.,
\romeogloss{fundamental types}, such as \lstinline!int! and \lstinline!double!,
enumerated or pointer types, or a C-style array or \lstinline!struct!
(a.k.a. a \romeogloss{POD}) --- were allowed to be members of a
\lstinline!union!. This limitation prevented any user-defined type having
a \romeogloss{non-trivial special member function} from being a member of a
\lstinline!union!:

\begin{emcppshiddenlisting}[emcppsbatch=e1]
#include <string>  // (ù{\codeincomments{std::string}}ù)
\end{emcppshiddenlisting}
\begin{emcppslisting}[emcppsbatch=e1]
union U0
{
    int         d_i;  // OK
    std::string d_s;  // compile-time error in C++03 (OK as of C++11)
};
\end{emcppslisting}

\noindent C++11 relaxes such restrictions on \lstinline!union! members, such as
\lstinline!d_s! above, allowing any type other than a \romeogloss{reference
type} to be a member of a \lstinline!union!.

A \lstinline!union! type is permitted to have user-defined special member
functions but --- by design --- does not initialize any of its members
automatically. Any member of a \lstinline!union! having a
\romeogloss{non-trivial constructor}, such as \lstinline!struct!~\lstinline!Nt!
below, must be constructed manually (e.g., via \romeogloss{placement
\lstinline!new!}) before it can be used:

\begin{emcppslisting}[emcppsbatch=e2]
struct Nt  // used as part of a (ù{\codeincomments{union}}ù) (below)
{
    Nt();   // non-trivial default constructor
    ~Nt();  // non-trivial destructor

    // Copy construction and assignment are implicitly defaulted.
    // Move construction and assignment are implicitly deleted.
};
\end{emcppslisting}

\noindent As an added safety measure, any non-trivial \romeogloss{special member
function} defined --- either implicitly or explicitly --- for any
\romeogloss{member} of a \lstinline!union! results in the compiler implicitly
deleting (see %Section~\ref{deleted-functions},
``\titleref{deleted-functions}" on page~\pageref{deleted-functions}) the corresponding \romeogloss{special
member function} of the \lstinline!union! itself:

\begin{emcppslisting}[emcppsbatch=e2]
union U1
{
    int d_i;   // fundamental type having all trivial special member functions
    Nt  d_nt;  // user-defined type having non-trivial special member functions

    // Implicitly deleted special member functions of (ù{\codeincomments{U1}}ù):
    /*
        U1()                     = delete; // due to explicit (ù{\codeincomments{Nt::Nt()}}ù)
        U1(const U1&)            = delete; // due to implicit (ù{\codeincomments{Nt::Nt(const Nt\&)}}ù)
        ~U1()                    = delete; // due to explicit (ù{\codeincomments{Nt::}}ù)~(ù{\codeincomments{Nt()}}ù)
        U1& operator=(const U1&) = delete; // due to implicit
                                           // (ù{\codeincomments{Nt::operator=(const Nt\&)}}ù)
    */
};
\end{emcppslisting}

%\noindent This same sort of precautionary deletion also occurs for any class
%containing such a union as a data member (see {\it\titleref{unrestrictedunion-use-cases}: \titleref{implementing-a-sum-type-as-a-discriminating-(or-tagged)-union}} on page~\pageref{implementing-a-sum-type-as-a-discriminating-(or-tagged)-union}).

A special member function of a \lstinline!union! that is implicitly deleted
can be restored via explicit declaration, thereby forcing a programmer
to consider how non-trivial members should be managed. For example,
we can start providing a \emph{value constructor} and corresponding
\emph{destructor}:

%begin{emcppslisting}
%struct U2
%{
%    union
%    {
%        int  d_i;   // fundamental type (trivial)
%        Nt   d_nt;  // non-trivial user-defined type
%    };
%
%    bool d_useInt;  // discriminator
%
%    U2(bool useInt) : d_useInt(useInt)       // value constructor
%    {
%        if (d_useInt) { new (&d_i) int(); }  // value initialized (to (ù{\codeincomments{0}}ù))
%        else          { new (&d_nt) Nt(); }  // default constructed in place
%    }
%
%    ~U2()  // destructor
%    {
%        if (!d_useInt) { d_nt.~Nt(); }
%    }
%};
%\end{emcppslisting}
\begin{emcppslisting}[emcppsbatch=e2]
#include <new>  // placement (ù{\codeincomments{new}}ù)

struct U2
{
    union
    {
        int  d_i;   // fundamental type (trivial)
        Nt   d_nt;  // non-trivial user-defined type
    };

    bool d_useInt;  // discriminator

    U2(bool useInt) : d_useInt(useInt)
    {
        if (d_useInt) { new (&d_i) int(); }  // value initialized (to (ù{\codeincomments{0}}ù))
        else          { new (&d_nt) Nt(); }  // default constructed in place
    }

    ~U2()  // destructor
    {
        if (!d_useInt) { d_nt.~Nt(); }
    }
};
\end{emcppslisting}


\noindent Notice that we have employed \romeogloss{placement \lstinline!new!} syntax to
control the lifetime of both member objects. Although assignment would
be permitted for the trivial \lstinline!int! type, it would be
\romeogloss{undefined behavior} for the non-trivial \lstinline!Nt! type:

%begin{emcppslisting}
%U2(bool useInt) : d_useInt(useInt)  // value constructor
%{
%    if (d_useInt) { d_i = int(); }  // value initialized (to (ù{\codeincomments{0}}ù))
%    else          { d_nt = Nt(); }  // undefined behavior
%}
%\end{emcppslisting}
\begin{emcppshiddenlisting}[emcppsbatch=e3]
// duplicate the requirements for U2 so that we can show an alternate version
// of the constructor
struct Nt  // used as part of a (ù{\codeincomments{union}}ù) (below)
{
    Nt();   // non-trivial default constructor
    ~Nt();  // non-trivial destructor
};
struct U2
{
    union
    {
        int  d_i;   // fundamental type (trivial)
        Nt   d_nt;  // non-trivial user-defined type
    };

    bool d_useInt;  // discriminator
\end{emcppshiddenlisting}
\begin{emcppslisting}[emcppsbatch=e3]
    U2(bool useInt) : d_useInt(useInt)
    {
        if (d_useInt) { d_i = int(); }  // value initialized (to (ù{\codeincomments{0}}ù))
        else          { d_nt = Nt(); }  // BAD IDEA: undefined behavior (no
                                        // lhs object)
    }
\end{emcppslisting}
\begin{emcppshiddenlisting}[emcppsbatch=e3]
};
\end{emcppshiddenlisting}


\noindent Now if we were to try to copy-construct or assign one object of type
\lstinline!U2! to another, the operation would fail because we have not
yet specifically addressed those \romeogloss{special member functions}:

\begin{emcppslisting}[emcppsbatch=e2]
void f()
{
    U2 a(false), b(true);  // OK (construct both instances of (ù{\codeincomments{U2}}ù))
    U2 c(a);               // Error, no (ù{\codeincomments{U2(const U2\&)}}ù)
    a = b;                 // Error, no (ù{\codeincomments{U2\& operator=(const U2\&)}}ù)
}
\end{emcppslisting}

\noindent We can restore these implicitly deleted special member functions too,
simply by adding appropriate copy-constructor and assignment-operator
definitions for \lstinline!U2! explicitly:

%begin{emcppslisting}
%union U2
%{
%    // ... (everything in (ù{\codeincomments{U2}}ù) above)
%
%    U2(const U2& original) : d_useInt(original.d_useInt)
%    {
%        if (d_useInt) { new (&d_i) int(original.d_i);  }
%        else          { new (&d_nt) Nt(original.d_nt); }
%    }
%
%    U2& operator=(const U2& rhs)
%    {
%        if (this == &rhs) // Prevent self-assignment.
%        {
%            return *this;
%        }
%
%        // Resolve all possible combinations of active types between the
%        // left-hand side and right-hand side of the assignment:
%
%        if (d_useInt)
%        {
%            if (rhs.d_useInt) { d_i = rhs.d_i; }
%            else              { new (&d_nt) Nt(rhs.d_nt); }
%        }
%        else
%        {
%            if (rhs.d_useInt) { d_nt.~Nt(); new (&d_i) int(rhs.d_i); }
%            else              { d_nt = rhs.d_nt; }
%        }
%
%        return *this;
%    }
%};
%\end{emcppslisting}
\begin{emcppshiddenlisting}[emcppsbatch=e4]
#include <new>  // placement (ù{\codeincomments{new}}ù)
// duplicate the requirements for U2 so that we can show an alternate version
// of the constructor
struct Nt  // used as part of a (ù{\codeincomments{union}}ù) (below)
{
    Nt();   // non-trivial default constructor
    ~Nt();  // non-trivial destructor
};
// --- Replace
// ... (everything in (ù{\codeincomments{U2}}ù) above)
    union
    {
        int  d_i;   // fundamental type (trivial)
        Nt   d_nt;  // non-trivial user-defined type
    };

    bool d_useInt;  // discriminator    
// --- End
\end{emcppshiddenlisting}
\begin{emcppslisting}[emcppsbatch=e4]
class U2
{
    // ... (everything in (ù{\codeincomments{U2}}ù) above)

    U2(const U2& original) : d_useInt(original.d_useInt)
    {
        if (d_useInt) { new (&d_i) int(original.d_i);  }
        else          { new (&d_nt) Nt(original.d_nt); }
    }

    U2& operator=(const U2& rhs)
    {
        if (this == &rhs) // Prevent self-assignment.
        {
            return *this;
        }

        // Resolve all possible combinations of active types between the
        // left-hand side and right-hand side of the assignment:

        if (d_useInt)
        {
            if (rhs.d_useInt) { d_i = rhs.d_i; }
            else              { new (&d_nt) Nt(rhs.d_nt); }  // (ù{\codeincomments{int}}ù) DTOR trivial
        }
        else
        {
            if (rhs.d_useInt) { d_nt.~Nt(); new (&d_i) int(rhs.d_i); }
            else              { d_nt = rhs.d_nt; }
        } d_useInt = rhs.d_useInt;

        // Resolve all possible combinations of active types between the
        // left-hand side and right-hand side of the assignment.  Use the
        // corresponding assignment operator when they match; otherwise,
        // if the old member is (ù{\codeincomments{d\_nt}}ù), run its non-trivial destructor, and
        // then copy-construct the new member in place:

        return *this;
    }
};
\end{emcppslisting}
Note that in the code example above, we ignore exceptions for exposition simplicity. Note also that attempting to
restore a \lstinline!union!'s implicitly deleted special member
functions by using the \lstinline!=!~\lstinline!default! syntax (see
%``\titleref{Defaulted-Special-Member-Functions}" on page~\pageref{Defaulted-Special-Member-Functions}
\featureref{\locationa}{Defaulted-Special-Member-Functions}) will still result in their being deleted because
the compiler cannot know which member of the union is active.

\subsection[Use Cases]{Use Cases}\label{unrestrictedunion-use-cases}

\subsubsection[Implementing a \romeogloss{sum type} as a discriminated {\tt union}]{Implementing a sum type as a discriminated {\SubsubsecCode union}}\label{implementing-a-sum-type-as-a-discriminating-(or-tagged)-union}

A \romeogloss{sum type} is an algebraic data type that provides a choice
among a fixed set of specific types. A C++11 unrestricted union can serve as a convenient and efficient way to define storage for a sum type (also called a \emph{tagged} or \emph{discriminated} union) because the alignment and size calculations are performed automatically by the compiler.

As an example, consider writing a parsing function \lstinline!parseInteger!
that, given a\linebreak[4] \lstinline!std::string! \lstinline!input!, will return, as a
\romeogloss{sum type} \lstinline!ParseResult! (see below), containing either an
\lstinline!int! result (on success) or an informative error message
on failure:

\begin{emcppshiddenlisting}[emcppsbatch=e5]
#include <sstream>  // (ù{\codeincomments{std::ostringstream}}ù)
#include <string>   // (ù{\codeincomments{std::string}}ù)
struct ParseResult {
    explicit ParseResult(int values);
    explicit ParseResult(const std::string& error);
};

// --- Replace
    if (/* Failure case (1). */)
    if (true)
// --- End

// --- Replace
    if (/* Failure case (2). */)
    if (true)
// --- End

\end{emcppshiddenlisting}
\begin{emcppslisting}[emcppsbatch=e5]
ParseResult parseInteger(const std::string& input)  // Return a sum type.
{
    int result;     // accumulate (ù{\codeincomments{result}}ù) as we go
    std::size_t i;  // current character index

    // ...

    if (/* Failure case (1). */)
    {
        std::ostringstream oss;
        oss << "Found non-numerical character '" << input[i]
            << "' at index '" << i << "'.";

        return ParseResult(oss.str());
    }

    if (/* Failure case (2). */)
    {
        std::ostringstream oss;
        oss << "Accumulating '" << input[i]
            << "' at index '" << i
            << "' into the current running total '" << result
            << "' would result in integer overflow.";

        return ParseResult(oss.str());
    }

    // ...

    return ParseResult(result);  // Success!
}
\end{emcppslisting}

\noindent The implementation above relies on \lstinline!ParseResult! being able to
hold a value of type either \lstinline!int! or \lstinline!std::string!. By
encapsulating a C++ \lstinline!union! and a \emph{discriminator} as part
of the \lstinline!ParseResult! \romeogloss{sum type}, we can achieve the
desired semantics:

%begin{emcppslisting}
%class ParseResult
%{
%    union  // storage for either the result or the error
%    {
%        int         d_value;  // trivial result type
%        std::string d_error;  // non-trivial error type
%    };
%
%    bool d_isError;  // discriminator
%
%public:
%    explicit ParseResult(int value);                 // value constructor (1)
%    explicit ParseResult(const std::string& error);  // value constructor (2)
%
%    ParseResult(const ParseResult& rhs);             // copy constructor
%    ParseResult& operator=(const ParseResult& rhs);  // copy assignment
%
%    ~ParseResult();                                  // destructor
%};
%\end{emcppslisting}
\begin{emcppshiddenlisting}[emcppsbatch=e6]
#include <string>   // (ù{\codeincomments{std::string}}ù)
\end{emcppshiddenlisting}
\begin{emcppslisting}[emcppsbatch=e6]
class ParseResult
{
    union  // storage for either the result or the error
    {
        int         d_value; // result type (trivial)
        std::string d_error; // error  type (non-trivial)
    };

    bool d_isError;  // discriminator

public:
    explicit ParseResult(int value);                 // value constructor (1)
    explicit ParseResult(const std::string& error);  // value constructor (2)

    ParseResult(const ParseResult& rhs);             // copy constructor
    ParseResult& operator=(const ParseResult& rhs);  // copy assignment

    ~ParseResult();                                  // destructor
};
\end{emcppslisting}


\noindent If a \romeogloss{sum type} comprised more than two types, the discriminator would be an appropriately-sized integral or enumerated type instead of a Boolean.

As discussed in
%{\it\titleref{unrestrictedunion-description}} on page~\pageref{unrestrictedunion-description}
\intrarefsimple{unrestrictedunion-description}, having a non-trivial
type within a \lstinline!union! forces the programmer to provide each
desired special member function and define it manually; note
that the use of placement \lstinline!new! is not required for either of the
two \emph{value constructors} (above) because the initializer syntax
(below) is sufficient to begin the lifetime of even a non-trivial
object:

%begin{emcppslisting}
%ParseResult::ParseResult(double value) : d_value(value), d_isError(false)
%{
%}
%
%ParseResult::ParseResult(const std::string& error)
%    : d_error(error), d_isError(true)
%    // Note that placement (ù{\codeincomments{new}}ù) was not necessary here because a new
%    // (ù{\codeincomments{std::string}}ù) object will be created as part of the initialization of
%    // (ù{\codeincomments{d\_error}}ù).
%{
%}
%\end{emcppslisting}
\begin{emcppslisting}[emcppsbatch=e6]
ParseResult::ParseResult(int value) : d_value(value), d_isError(false)
{
}

ParseResult::ParseResult(const std::string& error)
    : d_error(error), d_isError(true)
    // Note that placement (ù{\codeincomments{new}}ù) was not necessary here because a new
    // (ù{\codeincomments{std::string}}ù) object will be created as part of the initialization of
    // (ù{\codeincomments{d\_error}}ù).
{
}
\end{emcppslisting}


\noindent Placement \lstinline!new! and explicit destructor calls are still, however,
required for destruction and both copy operations{\cprotect\footnote{For
more information on initiating the lifetime of an object, see \cite{iso14}, section 3.8, ``Object Lifetime," pp. 66--69.}}:

%begin{emcppslisting}
%ParseResult::~ParseResult()
%{
%    if(d_isError)
%    {
%        d_error.std::string::~string();
%            // An explicit destructor call is required for (ù{\codeincomments{d\_error}}ù) because its
%            // destructor is non-trivial.
%    }
%}
%
%ParseResult::ParseResult(const ParseResult& rhs) : d_isError(rhs.d_isError)
%{
%    if (d_isError)
%    {
%        new (&d_error) std::string(rhs.d_error);
%            // Placement (ù{\codeincomments{new}}ù) is necessary here to begin the lifetime of a
%            // (ù{\codeincomments{std::string}}ù) object at the address of (ù{\codeincomments{d\_error}}ù).
%    }
%    else
%    {
%        d_value = rhs.d_value;
%            // Placement (ù{\codeincomments{new}}ù) is not necessary here as (ù{\codeincomments{int}}ù) is a trivial type.
%    }
%}
%
%ParseResult& ParseResult::operator=(const ParseResult& rhs)
%{
%    // Destroy (ù{\codeincomments{lhs}}ù)'s error string if existent:
%    if (d_isError) { d_error.std::string::~string(); }
%
%    // Copy (ù{\codeincomments{rhs}}ù)'s object:
%    if (rhs.d_isError) { new (&d_error) std::string(rhs.d_error); }
%    else               { d_value = rhs.d_value; }
%
%    d_isError = rhs.d_isError;
%    return *this;
%}
%\end{emcppslisting}
\begin{emcppshiddenlisting}[emcppsbatch=e6]
// The calls to d_error.std::string::~string() below do not compile on clang
// before clang 11.  Invoking it as ~basic_string instead works, but that
// is a confusing visible change.  Adding this using fixes the problem invisibly.
using std::string;
\end{emcppshiddenlisting}
\begin{emcppslisting}[emcppsbatch=e6]
ParseResult::~ParseResult()
{
    if (d_isError)
    {
        d_error.std::string::~string();
            // An explicit destructor call is required for (ù{\codeincomments{d\_error}}ù) because its
            // destructor is non-trivial.
    }
}

ParseResult::ParseResult(const ParseResult& rhs) : d_isError(rhs.d_isError)
{
    if (d_isError)
    {
        new (&d_error) std::string(rhs.d_error);
            // Placement (ù{\codeincomments{new}}ù) is necessary here to begin the lifetime of a
            // (ù{\codeincomments{std::string}}ù) object at the address of (ù{\codeincomments{d\_error}}ù).
    }
    else
    {
        d_value = rhs.d_value;
            // Placement (ù{\codeincomments{new}}ù) is not necessary here as (ù{\codeincomments{int}}ù) is a trivial type.
    }
}

ParseResult& ParseResult::operator=(const ParseResult& rhs)
{
    if (this == &rhs) // Prevent self-assignment.
    {
        return *this;
    }
    // Destroy (ù{\codeincomments{lhs}}ù)'s error string if existent:
    if (d_isError) { d_error.std::string::~string(); }

    // Copy (ù{\codeincomments{rhs}}ù)'s object:
    if (rhs.d_isError) { new (&d_error) std::string(rhs.d_error); }
    else               { d_value = rhs.d_value; }

    d_isError = rhs.d_isError;
    return *this;
}
\end{emcppslisting}
In practice, \lstinline!ParseResult! would typically use a more general \romeogloss{sum type}{\cprotect\footnote{\lstinline!std::variant!, introduced
in C++17, is the standard construct used to represent a \romeogloss{sum
type} as a \emph{discriminated union}. Prior to C++17,
\lstinline!boost::variant! was the most widely used \emph{tagged} union
  implementation of a \romeogloss{sum type}.}} abstraction to support arbitrary value types and provide proper exception safety.

%\noindent In practice, \texttt{ParseResult} would typically be defined as a
%template and renamed to allow any arbitrary result type \texttt{T} to be
%returned or else implemented in terms of a more general \textbf{sum
%type} abstraction.{\cprotect\footnote{\texttt{std::variant}, introduced
%in C++17, is the standard construct used to represent a \textbf{sum
%type} as a \emph{discriminating union}. Prior to C++17,
%\texttt{boost::variant} was the most widely used \emph{tagged} union
%  implementation of a \textbf{sum type}.}}

\subsection[Potential Pitfalls]{Potential Pitfalls}\label{potential-pitfalls}

\subsubsection[Inadvertent misuse can lead to latent \romeogloss{undefined behavior} at runtime]{Inadvertent misuse can lead to latent \romeogloss{undefined behavior} at runtime}\label{inadvertent-misuse-can-lead-to-latent-undefined-behavior-at-runtime}

When implementing a type that makes use of an unrestricted union,
forgetting to initialize a non-trivial object (using either a
\emph{member initialization list} or \romeogloss{placement \lstinline!new!}) or
accessing a different object than the one that was actually initialized
can result in tacit \romeogloss{undefined behavior}. Although forgetting to
destroy an object does not necessarily result in \romeogloss{undefined
behavior}, failing to do so for any object that manages a resource such
as dynamic memory will result in a \emph{resource leak} and/or lead to
unintended behavior. Note that destroying an
object having a trivial destructor is never necessary; there are, however, rare cases where
we may choose not to destroy an object having a non-trivial
one.
%{\cprotect\footnote{A specific example of where one might
%deliberately choose \emph{not} to destroy an object occurs when a
%collection of related objects are allocated from the same local memory
%resource and then deallocated unilaterally by releasing the memory
%back to the resource. No issue arises if the only resource that is ``leaked''
%by not invoking each individual destructor is the memory allocated
%from that memory resource, and that memory can be
%reused without resulting in \textbf{undefined behavior} if it
%is not subsequently referenced in the context of the deallocated
%  objects.}}

\subsection[Annoyances]{Annoyances}\label{annoyances}

\hspace*{\fill}

\subsection[See Also]{See Also}\label{see-also}

\begin{itemize}
\item{%Section~\ref{deleted-functions}, ``\titleref{deleted-functions}" on page~\pageref{deleted-functions} —
\seealsoref{deleted-functions}{\locationa}Any special member function of a \lstinline!union! that corresponds to a non-trivial one in any of its member elements will be implicitly deleted.} % "or when" changed to "and when" per JD Garcia and Slava
\end{itemize}

\subsection[Further Reading]{Further Reading}\label{further-reading}

\begin{itemize}
\item{\cite{goldthwaite07}}
\item{\cite{ouellet16}}
\end{itemize}





 %%%%%%%%%%%%%%% C++14 starts
\stepcounter{cppxx}
\addtocontents{toc}{\protect\renewcommand*\protect\cftsecindent{1.25em}}
\cftaddnumtitleline{toc}{section}{\thechapter.\thecppxx}{C++14}{\thepage}
% to add to the TOC at the section level
\renewcommand{\cppxx}{C++14}
%%%%%%%%%%%%%%

\newpage
\addtocontents{toc}{\protect\renewcommand*\protect\cftsecindent{4.5em}}
\renewcommand{\cppxx}{C++14}
%\section[{\tt decltypeauto}]{Deducing Types Using {\SecCode decltype} Sematics}\label{decltypeauto}



\emcppsFeature{
    short={\lstinline!decltypeauto!},
    tocshort={\TOCCode decltypeauto},
    long={Deducing Types Using {\SecCode decltype} Semantics},
    toclong={Deducing Types Using \lstinline!decltype! Semantics},
    rhshort={\RHCode decltypeauto},
}{decltypeauto}
\setcounter{table}{0}
\setcounter{footnote}{0}
\setcounter{lstlisting}{0}
%\section[{\tt decltypeauto}]{Deducing Types Using {\SecCode decltype} Semantics}\label{decltypeauto}


placeholder




\newpage
% \section[Deduced Return Type]{Function ({\SecCode auto}) {\SecCode return}-Type Deduction}\label{Function-Return-Type-Deduction}



\emcppsFeature{
    short={Deduced Return Type},
    long={Function ({\SecCode auto}) {\SecCode return}-Type Deduction},
    toclong={Function (\lstinline!auto!) \lstinline!return!-Type Deduction},
}{Function-Return-Type-Deduction}
\setcounter{table}{0}
\setcounter{footnote}{0}
\setcounter{lstlisting}{0}
%\section[Deduced Return Type]{Function ({\SecCode auto}) {\SecCode return}-Type Deduction}\label{Function-Return-Type-Deduction}



placeholder text........
 
 



  % Chapter 3
\cleardoublepage
%
%%%%%% RHs revert for Ch 4 + BM
\fancyhead[LO]{}
%\renewcommand{\sectionmark}[1]{\markright{{\sfnine\thesection\ #1}}}
%\renewcommand{\chaptermark}[1]{\markboth{{\sfb\thechapter\quad {\sfb #1}}}{}}
%%%%%%%%%%


\chapter[Parting Thoughts]{Parting Thoughts}\label{ch-parting}

\addtocontents{toc}{\protect\renewcommand*\protect\cftsecindent{1.25em}}

\section{Testing Section}

\section{Testing Another Section}

  % Chapter 4
\cleardoublepage
%
%
\cleardoublepage
%%\bibliographystyle{ieeetr-hacked}
\bibliographystyle{lakos}
\bibliography{emcs}
\nocite{*}
%\cleardoublepage
%
\chapter*{Bibliography}

\noindent\textbf{alexandrescu01}\\
\hspace*{2em}\parbox[t]{\textwidth}{Andrei Alexandrescu, \textit{Modern C++ Design: Generic Programming and Design Patterns Applied} (Reading, MA: Addison-Wesley, 2001)}\\

\noindent\textbf{ballman}\\
\hspace*{2em}\parbox[t]{\textwidth}{Aaron Ballman. ``Rule 03. Integers (INT)," \textit{ SEI CERT C++ Coding Standard} (Pittsburgh, PA: Carnegie Mellon University Software Engineering Institute)\\
https://wiki.sei.cmu.edu/confluence/pages/viewpage.action?pageId=88046333}\\

\noindent\textbf{calabrese20}\\
\hspace*{2em}\parbox[t]{\textwidth}{Matt Calabrese and Ryan McDougall. ``{\ttfamily any\_invocable}."  Technical Report P0288R6. (Geneva, Switzerland: International Standards Organization, 2020)\\
http://www.open-std.org/jtc1/sc22/wg21/docs/papers/2020/p0288r6.html}\\

\noindent\textbf{cpp11}\\
\hspace*{2em}\parbox[t]{\textwidth}{\textit{Information Technology — Programming
Languages — C++} (Geneva, Switzerland: International Standards Organization, 2011)\\
http://www.open-std.org/jtc1/sc22/wg21/docs/papers/2011/n3242.pdf}\\

\noindent\textbf{cppreferencea}\\
\hspace*{2em}\parbox[t]{\textwidth}{``std::all\_of, std::any\_of, std::none\_of," \textit{C++ Algorithm Library}, cppreference.com\\
https://en.cppreference.com/w/cpp/algorithm/all\_any\_none\_of}\\

\noindent\textbf{cppreferenceb}\\
\hspace*{2em}\parbox[t]{\textwidth}{``std::unordered\_map," \textit{C++ Containers Library}, cppreference.com\\
https://en.cppreference.com/w/cpp/container/unordered\_map}\\

\noindent\textbf{csb15}\\
\hspace*{2em}\parbox[t]{\textwidth}{\textit{C++ Object Persistence with ODB}, Revision 2.4, February 2015 (Code Synthesis Tools CC, 2015)\\
https://www.codesynthesis.com/products/odb/doc/manual.xhtml}\\

\noindent\textbf{cwg20}\\
\hspace*{2em}\parbox[t]{\textwidth}{``"Allow Duplicate Attributes," Technical Report, P2156r0, (Geneva, Switzerland: C++ Standards Committee Working Group, 2020)\\
https://wg21.link/p2156r0}\\

\newpage%%%%%%
\noindent\textbf{freesoftwarefdn20}\\
\hspace*{2em}\parbox[t]{\textwidth}{\textit{Using the GNU Compiler Collection (GCC)} (Boston, MA: Free Software Foundation, Inc., 2020)\\
https://gcc.gnu.org/onlinedocs/gcc/}\\
%\textbf{RETAIN IN CASE WE NEED THIS; USED IN ATTRIBUTES FEATURE}\qquad  section~6.33.1, ``Common Function Attributes," \\
%https://gcc.gnu.org/onlinedocs/gcc/Common-Function-Attributes.html\#Common-Function-Attribut}\\

\noindent\textbf{gustedt13}\\
\hspace*{2em}\parbox[t]{\textwidth}{Jens Gustedt. ``right angle brackets: shifting semantics,"  \textit{Jens Gustedt's Blog}, blog (WordPress, December 18, 2013)\\
https://gustedt.wordpress.com/2013/12/18/right-angle-brackets-shifting-semantics/}\\

\noindent\textbf{hinnant02}\\
\hspace*{2em}\parbox[t]{\textwidth}{Howard Hinnant. \textit{A Proposal to Add Move Semantics Support to the C++ Language} (Geneva, Switzerland: International Standards Organization, 2002)\\
http://www.open-std.org/jtc1/sc22/wg21/docs/papers/2002/n1377.htm}\\

\noindent\textbf{hinnant14}\\
\hspace*{2em}\parbox[t]{\textwidth}{Howard Hinnant. ``Everything You Ever Wanted To Know About Move Semantics (and Then Some)," \textit{Conference of the ACCU}, Bristol, England, April 8--12, 2014\\
https://accu.org/content/conf2014/Howard\_Hinnant\_Accu\_2014.pdf}\\

\noindent\textbf{hinnant16}\\
\hspace*{2em}\parbox[t]{\textwidth}{Howard Hinnant. ``Everything You Ever Wanted to Know About Move Semantics" \textit{Bloomberg Engineering Distinguished Speaker Series}, New York, NY, 2016\\
https://www.youtube.com/watch?v=vLinb2fgkHk\&t=28s }\\

\noindent\textbf{hinnant17}\\
\hspace*{2em}\parbox[t]{\textwidth}{Howard Hinnant. ``Issue 2278: User-Defined Literals for Standard Library Types," 10 Sep 2020, Technical Report CWG2278 (Geneva, Switzerland: International Standards Organization, 2017)\\
https://cplusplus.github.io/LWG/issue2278}\\

\noindent\textbf{hruska20}\\
\hspace*{2em}\parbox[t]{\textwidth}{Joel Hruska. ``How L1 and L2 CPU Caches Work, and Why They’re an Essential Part of Modern Chips," \textit{Extreme Tech}. April 14, 2020\\
https://www.extremetech.com/extreme/188776-how-l1-and-l2-cpu-caches-work-and-why-theyre-an-essential-part-of-modern-chips}\\

\noindent\textbf{ieee19}\\
\hspace*{2em}\parbox[t]{\textwidth}{\textit{IEEE Standard for Floating-Point Arithmetic} (New York, NY: Institute of Electrical and Electronics Engineers, Inc., 2019)}\\

\noindent\textbf{intel16}\\
\hspace*{2em}\parbox[t]{\textwidth}{Intel Corporation. \textit{Intel 64 and IA-32 Architectures Optimization Reference Manual}, Publication Number 248966-033 (Santa Clara, CA: Intel Corporation, 2016)\\
https://www.intel.com/content/dam/www/public/us/en/documents/manuals/64-ia-32-architectures-optimization-manual.pdf}\\

\newpage%%%%
\noindent\textbf{inteliig}\\
\hspace*{2em}\parbox[t]{\textwidth}{Intel. \textit{The Intel Intrinsics Guide} (Santa Clara, CA: Intel Corporation)\\
https://software.intel.com/sites/landingpage/IntrinsicsGuide/\#}\\

\noindent\textbf{iso11}\\
\hspace*{2em}\parbox[t]{\textwidth}{\textit{ISO/IEC 14882:2011 Programming Language C++} (Geneva, Switzerland: International Standards Organization, 2011)\\ 
NEED URL}\\

\noindent\textbf{iso14}\\
\hspace*{2em}\parbox[t]{\textwidth}{\textit{ISO/IEC 14882:2014 Programming Language C++} (Geneva, Switzerland: International Standards Organization, 2014)\\ 
http://www.open-std.org/JTC1/SC22/WG21/docs/papers/2013/n3797.pdf}\\

\noindent\textbf{iso18a}\\
\hspace*{2em}\parbox[t]{\textwidth}{\textit{C++ Standard Core Language Active Issues, Revision 100} (Geneva, Switzerland: International Standards Organization, 2018)\\
http://www.open-std.org/jtc1/sc22/wg21/docs/cwg\_active.html}\\

\noindent\textbf{iso18b}\\
\hspace*{2em}\parbox[t]{\textwidth}{\textit{ISO/IEC 9899:2018 Information Technology ---
Programming Languages --- C} (Geneva, Switzerland: International Standards Organization, 2018)\\ 
https://www.iso.org/standard/74528.html}\\

\noindent\textbf{iso20a}\\
\hspace*{2em}\parbox[t]{\textwidth}{\textit{ISO/IEC P2156R1:2020 Allow Duplicate Attributes} (Geneva, Switzerland: International Standards Organization, 2020)\\
http://www.open-std.org/jtc1/sc22/wg21/docs/papers/2020/p2156r1.pdf}\\

\noindent\textbf{iso20b}\\
\hspace*{2em}\parbox[t]{\textwidth}{\textit{ISO/IEC 14882:2020 Programming Languages --- C++} (Geneva, Switzerland: International Standards Organization, 2020)\\ 
https://www.iso.org/standard/79358.html}\\

\noindent\textbf{kahan97}\\
\hspace*{2em}\parbox[t]{\textwidth}{W. Kahan, \textit{Lecture Notes on the Status of IEEE Standard 754 for Binary Floating-Point Arithmetic} (Berkeley, CA: Electrical Engineering and Computer Science Department, University of California, 1997)\\
https://people.eecs.berkeley.edu/~wkahan/ieee754status/IEEE754.PDF }\\

\noindent\textbf{lakos16}\\
\hspace*{2em}\parbox[t]{\textwidth}{John Lakos, Jeffrey Mendelsohn, Alisdair Meredith, and Nathan Myers, ``On Quantifying Memory-Allocation Strategies (Revision 2),"  P0089R1, February 12, 2016\\
http://www.open-std.org/jtc1/sc22/wg21/docs/papers/2016/p0089r1.pdf}\\

\newpage%%%%%%
\noindent\textbf{lakos17}\\
\hspace*{2em}\parbox[t]{\textwidth}{John Lakos, ``Local (`Arena') Memory Allocators --- Part II," \textit{Meeting C++}, 2017 \\
https://www.youtube.com/watch?v=fN7nVzbRiEk}\\

\noindent\textbf{lakos19}\\
\hspace*{2em}\parbox[t]{\textwidth}{John Lakos, ``Value Proposition: Allocator-Aware (AA) Software," \textit{C++Now}, 2019\\
https://www.youtube.com/watch?v=dDR93TfacHc}\\

\noindent\textbf{lakos20}\\
\hspace*{2em}\parbox[t]{\textwidth}{John Lakos, \textit{Large-Scale C++ --- Volume I: Process and Architecture} (Reading, MA: Addison-Wesley, 2019)}\\

\noindent\textbf{lakos22}\\
\hspace*{2em}\parbox[t]{\textwidth}{John Lakos and Joshua Berne, \textit{C++ Allocators for the Working Programmer}, manuscript} \\

\noindent\textbf{lakos23}\\
\hspace*{2em}\parbox[t]{\textwidth}{John Lakos and Rostislava Klebnikhov. \textit{Design by Contract for Large-Scale Software}, manuscript} \\

\noindent\textbf{lakos2a}\\
\hspace*{2em}\parbox[t]{\textwidth}{John Lakos. \textit{Large-Scale C++ --- Volume II: Design and Implementation}, manuscript} \\

\noindent\textbf{lakos2b}\\
\hspace*{2em}\parbox[t]{\textwidth}{John Lakos. \textit{Large-Scale C++ --- Volume III: Verification and Testing} manuscript} \\

\noindent\textbf{lakos96}\\
\hspace*{2em}\parbox[t]{\textwidth}{John Lakos, \textit{Large-Scale C++ Software Design} (Reading, MA: Addison-Wesley, 1996)} \\

\noindent\textbf{mayrand15}\\
\hspace*{2em}\parbox[t]{\textwidth}{Fran\c{c}ois Mayrand. ``Passkey Idiom and Better Friendship in C++," \textit{Desktop Application Development}, blog, 21 May 2015 (Spiria Digital Inc., 2015)\\
https://www.spiria.com/en/blog/desktop-software/passkey-idiom-and-better-friendship-c/}\\

\noindent\textbf{meredith08}\\
\hspace*{2em}\parbox[t]{\textwidth}{A. Meredith, M. Wong, and J. Maurer. ``Inheriting Constructors (Revision 4),” N2512, 02 April 2008  (Geneva, Switzerland: International Standards Organization, 2008)\\
http://www.open-std.org/jtc1/sc22/wg21/docs/papers/2008/n2512.html}\\

\noindent\textbf{meredith20}\\
\hspace*{2em}\parbox[t]{\textwidth}{A. Meredith. ``Down with \sout{typename} in the Library!” P2150R0, 14 April 2020 (Geneva, Switzerland: International Standards Organization, 2020)\\
http://open-std.org/JTC1/SC22/WG21/docs/papers/2020/p2150r0.html}\\

\noindent\textbf{meyers97}\\
\hspace*{2em}\parbox[t]{\textwidth}{Scott Meyers. \textit{Effective C++}, second ed. (Boston, MA: Addison Wesley, 1997)} \\

\noindent\textbf{microsoft}\\
\hspace*{2em}\parbox[t]{\textwidth}{Microsoft. \textit{Guidelines Support Library}\\
https://github.com/Microsoft/GSL}\\

\noindent\textbf{microsoftC26481}\\
\hspace*{2em}\parbox[t]{\textwidth}{Microsoft. ``C26481 NO\_POINTER\_ARITHMETIC," \textit{Guidelines Support Library}, 2020
https://docs.microsoft.com/en-us/cpp/code-quality/c26481?view=vs-2019}\\

\noindent\textbf{miller13}\\
\hspace*{2em}\parbox[t]{\textwidth}{Mike Miller. ``Issue 1655: Line Endings in Raw String Literals,” 26 April 2013, included within \textit{C++ Standard Core Language Active Issues, Revision 100} (Geneva, Switzerland: International Standards Organization, 2018)\\
http://www.open-std.org/jtc1/sc22/wg21/docs/cwg\_active.html\#1655}

\noindent\textbf{miller17}\\
\hspace*{2em}\parbox[t]{\textwidth}{Mike Miller. ``Issue 2354: Extended Alignment and Object Representation,” 05 September 2017, included within \textit{Core Language Working Group Tentatively Ready Issues for the February, 2019 (Kona) Meeting} (Geneva, Switzerland: International Standards Organization, 2019)\\
http://www.open-std.org/jtc1/sc22/wg21/docs/papers/2019/p1359r0.html\#2354}\\

\noindent\textbf{meyers05}\\
\hspace*{2em}\parbox[t]{\textwidth}{Scott Meyers. \textit{Effective C++} (Boston, MA: Addison-Wesley, 2005)}\\

\noindent\textbf{pusz20}\\
\hspace*{2em}\parbox[t]{\textwidth}{Mateusz Pusz.\textit{Enable variable template template Parameters}, Technical Report P2008R0 (Geneva, Switzerland: C++ Standards Committee Working Group, 2005)}\\

\noindent\textbf{solihin15}\\
\hspace*{2em}\parbox[t]{\textwidth}{Yan Solihin. \textit{Fundamentals of Parallel Multi-Core Architecture} (Boca Raton, FL: CRC Press, 2015)}\\

\noindent\textbf{seacord13}\\
\hspace*{2em}\parbox[t]{\textwidth}{Robert C. Seacord. \textit{Secure Coding in C and C++}, 2nd ed. (Boston, MA: Addison Wesley, 2013)}\\

\noindent\textbf{sharpe13}\\
\hspace*{2em}\parbox[t]{\textwidth}{Chris Sharpe. ``Contextually converted to bool," \textit{Chris's C++ Thoughts}, blog (Blogspot, July 28, 2013)\\
http://chris-sharpe.blogspot.com/2013/07/contextually-converted-to-bool.html}\\

\noindent\textbf{smith15a}\\
\hspace*{2em}\parbox[t]{\textwidth}{Richard Smith. \textit{Attributes for namespaces and enumerators},  Technical Report N4196 (Geneva, Switzerland: C++ Standards Committee Working Group, 2015)}\\

\noindent\textbf{smith15b}\\
\hspace*{2em}\parbox[t]{\textwidth}{Richard Smith. \textit{Rewording Inheriting Constructors (Core Issue 1941 et al.)},  19 Oct 2020, Technical Report P0136R1 (Geneva, Switzerland: C++ Standards Committee Working Group, 2015)}\\

\noindent\textbf{smith18}\\
\hspace*{2em}\parbox[t]{\textwidth}{Richard Smith. \textit{Base Class Copy and Move Constructors Should Not Be Inherited}, 26 Feb 2018, Technical Report CWG2356 (Geneva, Switzerland: C++ Standards Committee Working Group, 2018)}\\


\noindent\textbf{stevens93}\\
\hspace*{2em}\parbox[t]{\textwidth}{W. Richard Stevens. \textit{Advanced Programming in the UNIX Environment} (Boston, MA: Addison-Wesley, 1993)} \\

\noindent\textbf{stroustrup13}\\
\hspace*{2em}\parbox[t]{\textwidth}{Bjarne Stroustrup. \textit{The C++ Programming Language}, 4th ed. (Boston, MA: Addison-Wesley, 2013)} \\

\noindent\textbf{stroustrup20}\\
\hspace*{2em}\parbox[t]{\textwidth}{Bjarne Stroustrup and Herb Sutter, Eds. \textit{C++ Core Guidelines} (\textbf{CITY}, WA: Standard C++ Foundation), 2020\\
https://isocpp.github.io/CppCoreGuidelines/CppCoreGuidelines} \\

\noindent\textbf{vandevoorde05}\\
\hspace*{2em}\parbox[t]{\textwidth}{Daveed Vandevoorde. \textit{Right Angle Brackets}, Technical Report N1757, Revision 2 (Geneva, Switzerland: C++ Standards Committee Working Group, 2005)\\
http://www.open-std.org/jtc1/sc22/wg21/docs/papers/2005/n1757.html}\\

\noindent\textbf{wight}\\
\hspace*{2em}\parbox[t]{\textwidth}{Hyrum Wight. \textit{Hyrum's Law}.\\
https://www.hyrumslaw.com/}\\





\cleardoublepage
\cleardoublepage % This command ensures that each chapter starts on a right-hand page 

\chapter*{Glossary}\label{Glossary}

% A new command, gloss, handles the glossary entry. \gloss{word to be defined}{definition}

\gloss{ABI}{}

\gloss{access modifier}{}

\gloss{access specifiers}{}

\gloss{acquire/release memory barrier}{}

\gloss{aggregate}{}

\gloss{aggregates}{}

\gloss{aggregate initialization}{}

\gloss{algebra}{}

\gloss{alias template}{}

\gloss{alignment}{}

\gloss{alignment requirements}{}

\gloss{API}{}

\gloss{argument-dependent lookup (ADL)}{}

\gloss{arithmetic type}{}

\gloss{array type}{}

\gloss{atomic}{}

\gloss{automatic variables}{}

\gloss{barriers}{}

\gloss{basic source character set}{}

\gloss{benchmark test}{}

\gloss{B0}{}

\gloss{boilerplate code}{}

\gloss{brace elision}{}

\gloss{bytes}{}

\gloss{C-style functions}{}

\gloss{cache hit}{}

\gloss{cache line}{}

\gloss{cache miss}{}

\gloss{callback functions}{}

\gloss{callable object}{}

\gloss{cast}{}

\gloss{CI}{}

\gloss{Class member access expression}{TODO (include any expression that is used to refer to a class member, such as \texttt{object.member}, \texttt{object->member}, \texttt{object.*member}.)}

\gloss{closure}{}

\gloss{closures}{}

\gloss{code bloat}{}

\gloss{compile-time constant}{}

\gloss{compile-time coupling}{}

\gloss{Compile-time dispatch}{TODO}

\gloss{complete type}{}

\gloss{component}{}

\gloss{component-local}{}

\gloss{Concepts}{}

\gloss{concrete class}{}

\gloss{conditionally supported}{}

\gloss{Constant expression}{An expression that can be evaluated at compile time.  Mention \texttt{constexpr} and state that \texttt{const} variables that are initialized from a compile-time constants are themselves required to be compile-time constants.  New info to me in June 2020, worthwhile to have. }

\gloss{constant initialization}{}

\gloss{contextual convertibility to {\tt bool}}{}

\gloss{Contextual keyword}{A \textit{contextual keyword} is a special identifier that acts like a \textit{keyword} when used in particular contexts. \texttt{override} is an example as it can be used as a regular identifier outside of member-function declarators.}

\gloss{continuous refactoring}{}

\gloss{contract}{}

\gloss{conventional string literals}{}

\gloss{converting constructors}{}

\gloss{conversion operators}{}

\gloss{copy elision}{}

\gloss{copy initialization}

\gloss{Copy semantics}{TODO}

\gloss{critical section}{}

\gloss{\textbf{curiously recurring template pattern (CRTP)}}{}

\gloss{cv-qualifiers}{}

\gloss{cyclically dependent}{}

\gloss{cyclic physical dependency}{}

\gloss{data member}{}

\gloss{data races}{}

\gloss{declare}{}

\gloss{Declaration}{TODO}

\gloss{declared interface}{}

\gloss{Declared type}{TODO The type of the *entity* named by the given expression.}

\gloss{declaring}{}

\gloss{defensive}{}

\gloss{define}{}

\gloss{defining}{}

\gloss{definition}{}

\gloss{delegating constructor}{}

\gloss{deleted function}{}

\gloss{dependent type}{}

\gloss{design pattern}{}

\gloss{diffusion}{}

\gloss{direct initialization}{}

\gloss{direct mapped}{}

\gloss{duck typing}{}

\gloss{dumb data}{}

\gloss{embedded systems}{}

\gloss{emplacement}{}

\gloss{encapsulation}{}

\gloss{Entity}{TODO}

\gloss{escalating}{}

\gloss{excess-\textit{n}}{}

\gloss{Expression SFINAE}{TODO}

\gloss{explicit}{}

\gloss{{\tt explicit} specifier}{}

\gloss{extended alignment}{}

\gloss{explicit-instantiation \emph{declaration}}{}

\gloss{explicit-instantiation \emph{definition}}{}

\gloss{explicit-instantiation \emph{directive}}{}

\gloss{false sharing}{}

\gloss{fences}{}

\gloss{flow of control}{}

\gloss{footprint}{}

\gloss{forward declaration}{}

\gloss{forward declared}{}

\gloss{forwarding reference}{}

\gloss{fragmentation}{}

\gloss{free function}{}

\gloss{full specialization}{}

\gloss{fully associative}{}

\gloss{fully constructed}{fully constructed}

\gloss{function object}{}

\gloss{function scope}{}

\gloss{fundamental alignment}{}

\gloss{fundamental integral types}{}

\gloss{Fundamental type}{TODO}

\gloss{garbage value}{}

\gloss{general purpose machines}{General-purpose machines are what Big Tech, the financial industry, and many other large companies use exclusively.}

\gloss{golden output}{}

\gloss{header-only library}{}

\gloss{hide or hiding}{Function-name \textbf{hiding} occurs when a member function in a derived class has the same name as one in the base class, but it is not overriding it due to a difference in the function signature or because the member function in the base class is not \texttt{virtual}. The hidden member function will \textbf{not} participate in dynamic dispatch; the member function of the base class will be invoked instead when invoked via a pointer or reference to the base class . The same code would have invoked the derived class's implementation had the member function of the base class had been \textbf{overridden} rather than \textbf{hidden}.}

\gloss{higher-order function}{}

\gloss{Hyrum's law}{}

\gloss{Id-expression}{TODO are most commonly \textbf{Identifiers}; other forms include overloaded operator names (in function notation), names of user-defined-conversion or literal operators, and destructor names, and ??template names followed by their argument lists??.}

\gloss{ill formed}{TODO (\texttt{[temp.res]p8})}

\gloss{ill formed, no diagnostic required (IFNDR)}{}

\gloss{imperative}{}

\gloss{implementation-defined}{}

\gloss{implementation inheritance}{}

\gloss{incomplete type}{}

\gloss{inheriting constructors}{}

\gloss{{\tt inline}~{\tt namespace}}{}

\gloss{in-process}{}

\gloss{instantiation time}{}

\gloss{insulate}{}

\gloss{insulating}

\gloss{Insulation}{TODO}

\gloss{Integer literal}{TODO}

\gloss{integral constant expression}{}

\gloss{integral promotion}{}

\gloss{integral type}{}

\gloss{interface inheritance}{}

\gloss{invocable}{}

\gloss{join}{}

\gloss{lambda expression}{}

\gloss{lambda expressions}{}

\gloss{lambda-introducer (adj)}{}

\gloss{leaks}{}

\gloss{linkage}{}

\gloss{list initialization}{}

\gloss{literal type}{}

\gloss{local class}{}

\gloss{local declaration}{}

\gloss{local scope}{}

\gloss{locality of reference}{}

\gloss{logical}{}

\gloss{long-distance friendship}{}

\gloss{lvalue}{}

\gloss{lvalue reference}{}

\gloss{mangled names}{}

\gloss{mantissa}{}

\gloss{maximal alignment}{}

\gloss{maximally aligned}{}

\gloss{mechanism}{}

\gloss{member initialization lists}{}

\gloss{member initializer list}{}

\gloss{memory leak}{}

\gloss{metafunction}{}

\gloss{mixed-mode builds}{}

\gloss{mix-in}{}

\gloss{modules}{}

\gloss{monotonic allocator}{}

\gloss{most vexing parse}{}

\gloss{move operations}{}

\gloss{Move semantics}{TODO}

\gloss{multithreading contexts}{}

\gloss{Meyers Singleton}{}

\gloss{name mangling}{}

\gloss{natural alignment}{}

\gloss{naturally aligned}{}

\gloss{{\tt new} handler}{}

\gloss{nibbles}{}

\gloss{nonprimitive functionality}{}

\gloss{non-trivial constructor}{}

\gloss{non-trivial special member function}{TODO}

\gloss{null address}{}

\gloss{object invariants}{}

\gloss{one-definition rule (ODR)}{}

\gloss{ODR-used}{}

\gloss{one definition rule}{}

\gloss{opaque declarations}{}

\gloss{opaque enumeration}{}

\gloss{Overriding}{TODO}

\gloss{parameter pack}{}

\gloss{partial class template specialization}{}

\gloss{partial implementation}{}

\gloss{partially constructed}{}

\gloss{perfectly forwarded}{}

\gloss{physical}{}

\gloss{physical design}{}

\gloss{Placement}{TODO}

\gloss{placement {\tt new}}{TODO}

\gloss{POD type}{TODO}

\gloss{pointer to member}{}

\gloss{polymorphic memory resource}{}

\gloss{precondition}{}

\gloss{predicate function}{}

\gloss{protocol}{}

\gloss{\romeovalue{prvalue}}{}

\gloss{qualified name}{}

\gloss{RAII}{``Resource Acquisition is Initialization"}

\gloss{Range}{TODO}

\gloss{raw string literals}{}

\gloss{Redundant check}{TODO}

\gloss{reference collapsing}{}

\gloss{Reference type}{TODO}

\gloss{ref-qualifiers}{}

\gloss{regular type}{}

\gloss{return value optimization (RVO)}{}

\gloss{\romeovalue{rvalue}}{}

\gloss{safe-bool-idiom}{}

\gloss{section}{}

\gloss{set associative}{}

\gloss{SFINAE}{TODO}

\gloss{SHA}{}

\gloss{shadowed}{}

\gloss{side effects}{}

\gloss{Signature}{TODO}

\gloss{signed integer overflow}{}

\gloss{slicing}{}

\gloss{special member function}{TODO}

\gloss{special member functions}{TODO}

\gloss{standard conversion}{}

\gloss{standard-layout types}{}

\gloss{static data space}{}

\gloss{\lstinline!static! duration}{}

\gloss{{\ttfamily std::unique\_ptr}}{}

\gloss{String literal}{TODO}

\gloss{strong \lstinline!typedef!}{}

\gloss{strong-\lstinline!typedef! idiom}{}

\gloss{Structural inheritance}{TODO}

\gloss{Sum type}{Abstract data type allowing the representation of one of multiple possible alternative types. Each alternative has its own type (and state), and only one alternative can be ``active" at any given point in time. Sum types automatically keep track of which choice is ``active," and properly implement value-sematic special member functions (even for non-trivial types). They can be implemented efficiently as a C++ \texttt{class} using a C++ \texttt{union} and a separate (integral) discriminator. This sort of implementation is commonly referred to as a discriminating (or ``tagged") union.}

\gloss{synthetization}{}

\gloss{template-head}{}

\gloss{Template instantiation time}{TODO}

\gloss{Template instantiation}{TODO}

\gloss{template template parameter}{}

\gloss{test driver}{}

\gloss{thrashing}{}

\gloss{thread pool}{}

\gloss{TLB}{}

\gloss{translation unit (TU)}{}

\gloss{trivial}{}

\gloss{triviality}{}

\gloss{trivial copy constructor}{}

\gloss{trivial operation}{}

\gloss{trivially-constructible}{}

\gloss{trivially copyable}{}

\gloss{Trivial type}{TODO}

\gloss{type alias}{}

\gloss{type deduction}{}

\gloss{typedef}{}

\gloss{type expression}{}

\gloss{type inference}{}

\gloss{type list}{}

\gloss{Type trait}{TODO}

\gloss{UDT}{}

\gloss{Undefined behavior}{TODO}

\gloss{undefined symbols}{}

\gloss{underlying type}{}

\gloss{user-defined type (UDT)}{}

\gloss{user-provided}{}

\gloss{user-provided special member function}{}

\gloss{{\tt using}-declaration}{}

\gloss{{\tt using} directive}{}

\gloss{value}{}

\gloss{value category}{}

\gloss{value constructor}{}

\gloss{value-semantic}{}

\gloss{value-semantic type (VST)}{}

\gloss{variable}{}

\gloss{variadic pack}{}

\gloss{vocabulary type}{}

\gloss{vocabulary types}{}

\gloss{well formed}{}

\gloss{working set}{}

\gloss{zero initialized}{}

%
%%\nocite{*}
%
\cleardoublepage

\newcommand\see[2]{\emph{\seename} #1}  
\newcommand\also[2]{\emph{See also} #1} 
\newcommand\seename{See} 
\newcommand\seealso[2]{\emph{See also} #1}

\begin{theindex}
\footnotesize
{\sffamily\bfseries{Symbols}}\nopagebreak

  \item \$, 31
  \item \%, 12
  \item \&, 31

  \indexspace
{\sffamily\bfseries\sffamily\bfseries{A}}\nopagebreak

  \item a very long entry to test the column width, 31
  \item anise, 31
  \item apple
    \subitem braebrun, 31
    \subitem cameo, 31
    \subitem fuji, 31
    \subitem gala, 31
    \subitem granny smith, 31
    \subitem red delicious, 31
  \item apricots, 31
  \item avocado, 31

  \indexspace
{\sffamily\bfseries{B}}\nopagebreak

  \item banana, 31
  \item basil, 31
  \item bibendum, 12
    \subitem dapibus, 12
  \item blackberry, 31

  \indexspace
{\sffamily\bfseries{C}}\nopagebreak

  \item cabbage, 31
  \item celery, 31
  \item chervil, 31
  \item chives, 31
  \item cilantro, 31
  \item \protect{\tt codeword}, 12
  \item corn, 31
  \item cucumber, 31

  \indexspace
{\sffamily\bfseries{D}}\nopagebreak

  \item dates, 31
  \item dill, 31

  \indexspace
{\sffamily\bfseries{E}}\nopagebreak

  \item eggplant, 31
  \item endive, 31

  \indexspace
{\sffamily\bfseries{F}}\nopagebreak

  \item fennel, 31
  \item fig, 31
  \item \protect{\tt function}, 31

  \indexspace
{\sffamily\bfseries{G}}\nopagebreak

  \item garlic, 31
  \item grapes, 31

  \indexspace
{\sffamily\bfseries{H}}\nopagebreak

  \item horseradish, 31
  \item how about another long entry, 31
  \item huckleberry, 31

  \indexspace
{\sffamily\bfseries{J}}\nopagebreak

  \item jicama, 31

  \indexspace
{\sffamily\bfseries{K}}\nopagebreak

  \item kale, 31
  \item kiwi, 31

  \indexspace
{\sffamily\bfseries{L}}\nopagebreak

  \item leeks, 31
  \item lemon, 31
  \item lettuce
    \subitem boston bibb, 31
    \subitem iceberg, 31
    \subitem mesclun, 31
    \subitem red leaf, 31
  \item lime, 31
  \item lorem, \see{lobortis}{12}

  \indexspace
{\sffamily\bfseries{M}}\nopagebreak

  \item majoram, 31
  \item mango, 31
  \item maybe another long entry for this test, 31
  \item melon
    \subitem canary, 31
    \subitem cantaloupe, 31
    \subitem honeydew, 31
    \subitem watermelon, 31
  \item mushrooms
    \subitem button, 31
    \subitem porcini, 31
    \subitem portabella, 31
    \subitem shitake, 31

  \indexspace
{\sffamily\bfseries{N}}\nopagebreak

  \item nectarine, 31
  \item nutmeg, 31

  \indexspace
{\sffamily\bfseries{O}}\nopagebreak

  \item okra, 31\pagebreak % you'll need to manually break pages or columns at times to get a nice layout
  \item onion
    \subitem red, 31
    \subitem vidalia, 31
    \subitem yellow, 31
  \item orange, 31

  \indexspace
{\sffamily\bfseries{P}}\nopagebreak

  \item papaya, 31
  \item parsley, 31
  \item peaches, 31
  \item peppers
    \subitem ancho, 31
    \subitem bell, 31
    \subitem habe\~{n}eros, 31
    \subitem jalape\~{n}os, 31
    \subitem pabla\~{n}os, 31
  \item perhaps yet another long entry for this test, 31
  \item plantains, 31
  \item plums, 31
  \item potatoes
    \subitem red-skinned, 31
    \subitem russet, 31
    \subitem yukon gold, 31
  \item pumpkin, 31

  \indexspace
{\sffamily\bfseries{Q}}\nopagebreak

  \item quince, 31

  \indexspace
{\sffamily\bfseries{R}}\nopagebreak

  \item radicchio, 31
  \item radish, 31
  \item raspberry, 31
  \item rosemary, 31
  \item rutabaga, 31\newpage % to end this column short so that it and the last column are balanced

  \indexspace
{\sffamily\bfseries{S}}\nopagebreak

  \item shallots, 31
  \item spinach, 31
  \item squash, 31
  \item still another long entry for column width testing, 31

  \indexspace
{\sffamily\bfseries{T}}\nopagebreak

  \item thyme, 31
  \item tomatillo, 31
  \item tomatoes
    \subitem cherry, 31
    \subitem grape, 31
    \subitem heirloom, 31
    \subitem hybrid, 31
    \subitem roma, 31
  \item \protect{\tt typeof}, 31

  \indexspace
{\sffamily\bfseries{U}}\nopagebreak

  \item ugly fruit, 31

  \indexspace
{\sffamily\bfseries{V}}\nopagebreak

  \item verbena, 31
  \item viverra, \seealso{neque}{12}

  \indexspace
{\sffamily\bfseries{X}}\nopagebreak

  \item xacuti masala, 31

  \indexspace
{\sffamily\bfseries{Y}}\nopagebreak

  \item yams, 31
  \item yet another very long entry for column width test, 31

  \indexspace
{\sffamily\bfseries{Z}}\nopagebreak

  \item zucchini, 31

\end{theindex}


% 
% %
% 
% 
\end{document}
